\documentclass[11pt,a4paper]{article}
\usepackage[T1]{fontenc}
\usepackage[utf8]{inputenc}
\usepackage{lmodern,microtype}
\usepackage[margin=26mm,headheight=15pt]{geometry}
\usepackage{amsmath,amssymb,amsthm,mathtools}
\usepackage{booktabs,longtable,array,enumitem}
\usepackage[dvipsnames]{xcolor}
\usepackage{xurl}
\usepackage[colorlinks=true,linkcolor=MidnightBlue,citecolor=MidnightBlue,urlcolor=MidnightBlue]{hyperref}
\usepackage{bookmark,fancyhdr}
\pagestyle{fancy}\fancyhf{}
\fancyhead[L]{\small Observation and Recovery: Collected Sequels}
\fancyhead[R]{\small KEIO}\fancyfoot[C]{\thepage}
\renewcommand{\headrulewidth}{0.3pt}
\setlength{\emergencystretch}{3em}
\setlength{\parskip}{3pt}
\setlist{itemsep=3pt,topsep=5pt}
\setcounter{tocdepth}{2}
\makeatletter
\renewcommand*\l@section{\@dottedtocline{1}{0em}{3.2em}}
\renewcommand*\l@subsection{\@dottedtocline{2}{1.5em}{4.8em}}
\makeatother
\renewcommand{\thesection}{\arabic{part}.\arabic{section}}
\renewcommand{\theHsection}{\arabic{part}.\arabic{section}}
\numberwithin{equation}{section}
\newtheorem{theorem}{Theorem}[section]
\newtheorem{proposition}[theorem]{Proposition}
\newtheorem{lemma}[theorem]{Lemma}
\newtheorem{corollary}[theorem]{Corollary}
\theoremstyle{definition}
\newtheorem{definition}[theorem]{Definition}
\newtheorem{example}[theorem]{Example}
\theoremstyle{remark}
\newtheorem{remark}[theorem]{Remark}
\DeclareMathOperator{\supp}{supp}
\DeclareMathOperator{\dist}{dist}
\DeclareMathOperator{\diag}{diag}
\DeclareMathOperator{\rank}{rank}
\DeclareMathOperator{\Var}{Var}
\DeclareMathOperator{\Cov}{Cov}
\DeclareMathOperator{\tr}{tr}
\newcommand{\R}{\mathbb R}
\newcommand{\C}{\mathbb C}
\newcommand{\N}{\mathbb N}
\newcommand{\Z}{\mathbb Z}
\newcommand{\T}{\mathbb T}
\newcommand{\E}{\mathbb E}
\hypersetup{pdftitle={Observation and Recovery: Geometry, Mixtures, and Spectral Inversion},pdfauthor={KEIO},pdfsubject={Unified edition of the ten sequels following the 223-page treatise; the original treatise is excluded}}
\begin{document}

\hypersetup{pageanchor=false}
\begin{titlepage}
\centering
\vspace*{25mm}
{\LARGE\bfseries Observation and Recovery\par}
\vspace{6mm}
{\Large Geometry, Mixtures, and Spectral Inversion\par}
\vspace{12mm}
{\large A Unified Edition of the Ten Subsequent Papers\par}
\vspace{5mm}
{\normalsize Following the 223-page treatise\par}
\vspace{20mm}
{\Large KEIO\par}
\vspace{6mm}
{\large 16 October 2026\par}
\vfill
\begin{minipage}{.82\textwidth}
\small
This volume collects the ten subsequent research manuscripts in one
continuous document. The earlier 223-page integrated treatise is not
reproduced. The original arguments, hypotheses, proofs, research ledgers,
and limitations are retained, with unified navigation and references.
This is an editorial consolidation of AI-assisted research manuscripts,
not an independent certification of their mathematical claims.
\end{minipage}
\vspace*{15mm}
\end{titlepage}

\pagenumbering{roman}
\hypersetup{pageanchor=true}
\section*{Abstract}
\addcontentsline{toc}{section}{Abstract and editorial guide}
This volume assembles ten successive investigations of observation and
recovery in arithmetic, geometry, probability, and statistics. Its central
objects are adaptive continuous measurements, repeated-observation mixture
experiments, spectral and zero-set inverse maps, and factorization
boundaries for probability distributions. The sequence develops recovery
widths and restricted information classes; population estimation under
replication, sampling budgets, and contamination; Gaussian and stable
mixture identification; prime and exponential-sum zero phenomena;
continuous spectral inversion; and Bessel--gamma positivity and extraction
problems. Complete source arguments are retained in chronological parts,
so that assumptions, intermediate advances, counterexamples, and remaining
obstructions can be read together. A common table of contents, globally
distinct theorem and equation numbers, cross-links, and a consolidated
bibliography make the sequence accessible as one document.

\paragraph{Scope of this edition.}
The source interval begins with \emph{Exact Recovery, Optimal Replication,
and Bessel Classification} and ends with \emph{Observation Budgets and
Hidden Parameters}. It contains all ten later manuscripts available in
this research sequence as of 16 October 2026. The 223-page
\emph{Observation and Recovery in Arithmetic and Probability: Prime Flows,
Rare Events, and Infinite Mixtures} remains an external predecessor.
References to that work are retained, but none of its chapters is imported.

\paragraph{How to read the chronology.}
The parts preserve the source order and the original mathematical text.
Each begins with its source abstract and keeps its stated hypotheses and
claim accounting. Statements about what is new, previously known, or still
open are dated statements within that part; later parts may extend the
range, weaken an assumption, or settle a question left there. Counts of
targets are local research ledgers, not additive counts of independent
discoveries or historical conjectures. Repeated statements are retained
when they support a later proof, comparison, or extension.

\paragraph{Numbering and references.}
Number \(i.j.k\) denotes item \(k\) in section \(j\) of Part \(i\).
Machine-linked references have been renumbered throughout. Unprefixed
section or theorem numbers inside bibliographic annotations refer to the
original cited publication or manuscript. Companion-paper bibliography
entries link to the corresponding part when the source is included here.
The source abstracts are historical abstracts, not a fresh claim of
priority by this editorial edition. Symbols are defined within their
respective parts; a symbol reused in another model has its local meaning.

\paragraph{Proof and verification status.}
This compilation preserves the status of the source manuscripts. It adds
no new proof of a named conjecture and does not turn a conditional result
into an unconditional one. Its production checks concern source coverage,
reference integrity, compilation, and presentation. Mathematical assessment
should use the full hypotheses and proofs, rather than the source titles
or target counts alone.

\paragraph{A specific overlap in the budget results.}
The general regular-component budget theorem in
Part~\ref{part:01}, Theorem~\ref{p01:info:thm:general-budget},
already contains the \(B^{-1/3}\) order for Bernoulli components.
Part~\ref{part:10}, Theorem~\ref{p10:ml16:budget}, revisits this
rate with a direct Bernoulli proof and an explicit finite linear-program
implementation. These are not two independent minimax-rate discoveries.
The source text is retained so the two arguments can be compared.

\paragraph{Progression and model distinctions.}
The restricted likelihood ranges in Parts III--IV precede the complete
curve stated in Part VI. The prime-zeta endpoint discussion in Part I
precedes the counting results in Parts VI--VII. Part VII treats the
constant-one coefficient conjecture separately from the still-distinct
optimal-exponent question. Part V already gives logarithmic constructions
for convex polynomial and analytic measurements; Part X adds fixed
two-sided Hessian bounds in the smooth class. Its remaining analytic
questions should not be read as withdrawing Part V's upper bound.
The positive finite-query noise gap of Part IX is quantified in Part X;
matching decay rates remain a separate question.

The zero-set counterexamples with additive frequency relations coexist
with the finite, internally rationally independent inverse theorem of
Part X. Similarly, separate terminal laws under a two-sided scale-moment
condition, separate terminal laws without moments, and joint replicates
are different experiments. Their identification thresholds should not be
combined without preserving those assumptions. Five or six sufficient
replicates are not claimed to be minimal for every fully mixed model.

\paragraph{Meaning of a single-file source.}
The accompanying TEX file contains the complete typesetting source and
requires no external chapter files, images, or bibliography database.
Mathematical dependencies on the excluded 223-page predecessor are still
dependencies; compilability does not assert that every foundational
result has been reproved here.


\clearpage
\section*{The ten constituent papers}
\addcontentsline{toc}{section}{The ten constituent papers}
{\small\setlength{\tabcolsep}{4pt}\renewcommand{\arraystretch}{1.25}
\begin{longtable}{@{}p{.06\textwidth}p{.67\textwidth}p{.15\textwidth}r@{}}
\toprule Part & Source title & Date & Page\\\midrule\endfirsthead
\toprule Part & Source title & Date & Page\\\midrule\endhead
I & \hyperref[part:01]{Exact Recovery, Optimal Replication, and Bessel Classification} & 14 Oct. 2026 & \pageref{part:01}\\
II & \hyperref[part:02]{Observation Geometry and Exact Recovery} & 14 Oct. 2026 & \pageref{part:02}\\
III & \hyperref[part:03]{Binomial Mixture Likelihood and Exact Observation Thresholds} & 14 Oct. 2026 & \pageref{part:03}\\
IV & \hyperref[part:04]{Transport, Spectral Positivity, and Sharp Observation Boundaries} & 14 Oct. 2026 & \pageref{part:04}\\
V & \hyperref[part:05]{Adaptive Analytic Information and Sharp Recovery Boundaries} & 14 Oct. 2026 & \pageref{part:05}\\
VI & \hyperref[part:06]{Boundary Laws for Observation and Recovery} & 15 Oct. 2026 & \pageref{part:06}\\
VII & \hyperref[part:07]{Exact Recovery Laws and a Coefficient Conjecture} & 16 Oct. 2026 & \pageref{part:07}\\
VIII & \hyperref[part:08]{Geometry and Rigidity of Inverse Observation} & 16 Oct. 2026 & \pageref{part:08}\\
IX & \hyperref[part:09]{Sharp Boundaries for Adaptive Recovery and Mixture Identification} & 16 Oct. 2026 & \pageref{part:09}\\
X & \hyperref[part:10]{Observation Budgets and Hidden Parameters} & 16 Oct. 2026 & \pageref{part:10}\\
\bottomrule\end{longtable}
}

\paragraph{Suggested thematic routes.}
For measurement complexity, follow the information sections in Parts
IV--X. For likelihood, corruption, and budget allocation, follow the
statistical sections in Parts I--VII and IX--X. For joint and separate
mixture observations, follow the Gaussian and stable sections throughout
the volume. For zero dynamics and inverse spectra, follow the arithmetic
and spectral sections; for distribution factorization, follow the Bessel,
Thorin, and gamma sections. The detailed contents below provides the
precise section locations.

\clearpage
\tableofcontents
\clearpage
\pagenumbering{arabic}


% ===== BEGIN SOURCE 01: KEIO_Exact_Recovery_and_Fifteen_Questions.tex =====
\clearpage
\part{Exact Recovery, Optimal Replication, and Bessel Classification}
\label{part:01}
\setcounter{section}{0}
\begin{center}\large\itshape Extensions of Observation and Recovery, with an Audit of Fifteen Published Questions\end{center}
\begin{center}\small Source manuscript: 14 October 2026.\end{center}
\begingroup
\renewcommand{\arraystretch}{1}
\begin{quote}\small\textbf{Source abstract.}\quad
We develop several extensions of an integrated study of arithmetic and
probabilistic observation maps. For repeated observations with an arbitrary
mixing distribution on an interval, regular component estimation and a
local quadratic relative-entropy bound imply the sharp total-budget
Wasserstein rate $B^{-1/3}$, attained by a fixed unequal allocation and
unimprovable by adaptive policies. In a factorial Legendre experiment this
strictly improves every equal-group allocation. For each fixed stable
shape, we obtain necessary and sufficient conditions on the observed
partial-sum indices; the converse constructs one signed mixing density
annihilating all selected observations simultaneously. We classify exact
real Mellin zeros and integer moment-sign patterns, identify the precise
quasianalytic boundary of sparse-moment recovery, and use quadrature zeros
to separate individual identification from uniform statistical learning.
For the published McKay classification problem, complex Bessel zeros
exclude hyperbolic complete monotonicity for four families, while the
second family admits an exact self-decomposability threshold and is never
a generalized gamma convolution. Further results concern correlated blind
calibration and uniformly translated prime-flow zero counts. An explicit
ledger treats fifteen questions from published papers or public preprints.
It separates inherited proofs, classical consequences, partial answers,
and unresolved targets. In particular, no Riemann-zeta conjecture is
claimed solved, and no numerical count of new independent conjecture
resolutions is asserted.
\end{quote}



\section{Scope, principal results, and the problem ledger}
\label{p01:intro:scope}

This paper is a sequel to \emph{Observation and Recovery in Arithmetic
and Probability: Prime Flows, Rare Events, and Infinite Mixtures}
\cite{p01:KEIOIntegrated}, a 223-page integrated working manuscript. The
question motivating this sequel is concrete: which externally stated
problems can its methods actually settle, and which internal sufficient
conditions can be sharpened to exact boundaries? We give complete proofs
of the principal extensions rather than infer additional conclusions
from the length or ambition of the earlier manuscript.

Here ``new extension'' means new relative to that manuscript. It does
not establish historical priority. The literature audit uses specific
primary-source formulations and the follow-ups located as of
14 October 2026; it is not an exhaustive certification that no other
resolution exists. The mathematical proofs and the priority assessment
are separate questions. The manuscript has received internal independent
proof checks, but not external peer review or formal verification.

\subsection{The most substantial extensions}

\begin{enumerate}
\item \textbf{A complete allocation rate.}
Theorem~\ref{p01:info:thm:general-budget} proves
$\mathfrak R_B^{\rm ad}\asymp\mathfrak R_B^{\rm fix}\asymp B^{-1/3}$
under explicit component-estimation and local-KL assumptions. The
factorial example gives a strict improvement over the lower bound
$(B\log(eB))^{-1/4}$ for every equal-group design
(Corollary~\ref{p01:info:cor:factorial-budget}). The multilevel idea itself
has established predecessors.

\item \textbf{Exact sparse stable sampling.}
For fixed stable index $\alpha\ne1$, arbitrary mixing laws are recovered
from the separate partial-sum laws at indices $\mathcal N$ exactly when
$\sum_{n\in\mathcal N}n^{-|1-\alpha|}=\infty$. At index one with
nonzero skewness, the exact condition is
$\sum_{n\in\mathcal N}(1+(\log n)^2)^{-1}=\infty$.
Theorem~\ref{p01:mix:thm:sharp-stable} includes genuine probability-law
counterexamples whenever the appropriate sum converges.

\item \textbf{A concrete published classification.}
Corollary~\ref{p01:rare:cor:four-mckay} answers the HCM portion of the McKay
question negatively throughout all four stated families. For
$\psi_{\mu,a,b}$, Theorem~\ref{p01:rare:thm:psi-classification} proves
\[
 \mathrm{SD}\quad\Longleftrightarrow\quad
 (2\mu+3)\sqrt{b^2-a^2}\ge b,
 \qquad \mathrm{GGC}\text{ never holds when }a>0.
\]
This is a complete answer to those subquestions, not to every clause
of the source's multipart problem.

\item \textbf{Exact recovery boundaries beyond a single example.}
Sections~\ref{p01:mom:sec:main}, \ref{p01:qa:section}, and
\ref{p01:mix:sec:quadrature-constraints} give exact moment-sign conditions,
a Denjoy--Carleman criterion, and density constraints that restore
identification. The last result also proves that individual identification
alone does not imply uniform learnability.
\end{enumerate}

\subsection{Fifteen externally stated targets}

The entries below are \emph{targets}, not fifteen claims of new solutions.
Three entries are distinct numbered prime-zeta conjectures; the multipart
McKay question is kept as one entry. The stable, regularity, calibration,
and allocation extensions are not counted as extra published conjectures.
The two inherited resolutions are proved again in
Section~\ref{p01:mix:sec:source-audit}.

\begingroup
\small
\setlength{\tabcolsep}{4pt}
\renewcommand{\arraystretch}{1.18}
\begin{longtable}{@{}p{0.045\linewidth}>{\raggedright\arraybackslash}p{0.285\linewidth}>{\raggedright\arraybackslash}p{0.61\linewidth}@{}}
\toprule
No. & Published target & Outcome and exact boundary \\
\midrule
\endfirsthead
\toprule
No. & Published target & Outcome and exact boundary \\
\midrule
\endhead
\bottomrule
\endfoot
1 & Lin: arbitrary normal directing laws \cite[\S2.6]{p01:mix:Lin2015}
& Positive answer without variance moments. Inherited proof rechecked;
Proposition~\ref{p01:mix:prop:normal-audit}. No new priority claim. \\
2 & Le et al.: converse to the differential barrier \cite[\S6]{p01:mix:Le2026}
& Negative for the two stated finite operator classes. Inherited entire
kernel counterexample rechecked; Proposition~\ref{p01:mix:prop:pde-audit}. \\
3 & Le et al.: constraints restoring estimation \cite[\S6]{p01:mix:Le2026}
& Partial answer for finite polynomial defects: prescribed quadrature
zeros give identification; a common smoothness bound gives uniform
learning. The general research program remains open. \\
4 & B\"urgin--Epperlein--Wirth, Conjecture 17 \cite{p01:mom:BEW2022}
& Correct, but already a consequence of Luxemburg--Korevaar (1971).
Alternative proof and stronger formulations; not a new historical closure. \\
5 & Belovas et al., prime-zeta Conjecture 1 \cite{p01:pf:BCS}
& Sharp right boundary already follows from Sepulcre--Vidal (2022).
Not a new closure. \\
6 & Belovas et al., prime-zeta Conjecture 2 \cite{p01:pf:BCS}
& Uniform normalized ordinates follow from classical almost-periodic
zero density in the specified strip. A uniform twist/translate
strengthening is proved here. \\
7 & Belovas et al., prime-zeta Conjecture 3 \cite{p01:pf:BCS}
& Positive linear counting follows in $1<\gamma<S$. Classical
consequence; the endpoint $\gamma=1$ is not settled. \\
8 & Ford--Zaharescu / FSZ: secondary interval term \cite{p01:zet:FZ,p01:zet:FSZ}
& Unresolved here. An $o(N(T))$ error does not give the required $o(T)$
error for sharp interval tests. \\
9 & \c{C}i\c{c}ek--Gonek: average-gap normalization \cite{p01:zet:CG}
& Unresolved here. The Fourier test requires a growing frequency
$x=T^k$, outside the fixed-frequency estimates. \\
10 & Odlyzko / Arias de Reyna: Gram-phase distribution \cite{p01:zet:AR}
& Unresolved here. Nonlinear phase and growing local frequency leave
an unproved exponential-sum estimate. \\
11 & Baricz et al., \S4.2: noncentral chi-square HCM \cite{p01:rare:Baricz2026}
& HCM classification already answered by Salazar (2026). An elementary
GGC obstruction is also proved here, without a priority claim. \\
12 & Baricz et al., \S4.3: McKay classification \cite{p01:rare:Baricz2026}
& HCM subquestion closed negatively for all four families; exact SD/GGC
classification of $\psi$. Other SD/GGC cases and the extra Theorem 12
question remain unresolved. \\
13 & Baricz et al., \S4.4: SD of three further laws \cite{p01:rare:Baricz2026}
& Unresolved here. One case reduces to the explicit Bessel heat-trace
inequality \eqref{p01:rare:eq:third-unsolved-inequality}. \\
14 & Lau--Scarlett: interaction in one-bit mean estimation \cite{p01:info:LauScarlett2026}
& Unresolved here. Repeated observations of one latent individual
are unavailable in that observation model. \\
15 & Vinayak et al.: larger-range binomial MLE bound \cite[Remark 3.2]{p01:info:VKVK2019}
& Unresolved here. The unequal-allocation theorem is not an equal-group
MLE theorem; verified 2025 progress is cited explicitly. \\
\end{longtable}
\endgroup

There are six targets unresolved here (8--10 and 13--15), two partly
answered targets (3 and 12), two inherited proof resolutions (1 and 2),
and five targets treated as classical consequences or previously answered
questions (4--7 and 11). This bookkeeping is intentionally conservative:
the exact new McKay subcase results and internal sharp classifications
have mathematical value without being relabeled as fifteen independent
open-problem resolutions.

\subsection{Notation and observation models}

The Wasserstein distance $W_1$ is taken with the usual Euclidean cost;
on a real interval it equals the $L^1$ distance between distribution
functions. Total variation uses the probability convention
$d_{\rm TV}(P,Q)=\sup_A|P(A)-Q(A)|$, while the variation norm of a
signed measure is its total mass after absolute variation. Constants in
$\lesssim$ and $\asymp$ may depend on a fixed model or interval, but not
on the sample size or budget. ``Separate partial-sum laws'' do not mean
jointly observed sums. ``Two replicates'' means their joint distribution
under a shared latent parameter. The distinction is essential.


% ===== Consolidated body: rare_events =====
\section{Bessel laws: published classification questions and exact answers}
\label{p01:rare:sec:bessel}

The three problems in Sections 4.2--4.4 of Baricz, Prabhu K, Singh,
and Vijesh V \cite{p01:rare:Baricz2026} provide external, explicitly stated
targets.  Their published theorem numbers are 1--4, 12, 14, 18, and 22;
the corresponding arXiv numbers are 2.1--2.4, 3.12, 3.14, 3.18, and
3.22.  We use the published numbering.  The first problem concerns
hyperbolically completely monotone densities for noncentral chi-square
laws, with generalized gamma convolutions discussed immediately before
the question.  The second is a multipart classification of McKay laws.
The third asks for self-decomposability of three further Bessel laws.
They are three source problems, not a separate conjecture for every
parameter or every distribution appearing in them.

The first problem has a recent HCM solution in Salazar
\cite{p01:rare:Salazar2026}; we make no priority claim for that answer.
Below we prove the negative HCM answer for all four McKay families,
and an exact self-decomposability and GGC classification for the second
family.  These conclusions answer explicit parts of the second source
problem.  They do not settle its remaining SD/GGC cases or its additional
question about the density in Theorem 12.  We also give the elementary
GGC obstruction for noncentral chi-square laws and a precise reduction
of one case of the third problem.  The latter remains unresolved here.
The source and follow-up checks were made on 14 October 2026; the
classification proofs are independent of any claim of historical priority.

\subsection{A zero-free obstruction for HCM functions}

A positive function $f$ is \emph{hyperbolically completely monotone}
(HCM) when $f(uv)f(u/v)$ is completely monotone as a function of
$v+v^{-1}$ for every $u>0$.  A \emph{generalized gamma convolution}
(GGC) is a law whose Laplace exponent has a Thorin representation
\begin{equation}\label{p01:rare:eq:thorin}
 h(s)=d s+\int_{(0,\infty)}\log(1+s/t)\,U(dt),\qquad d\geq0,
\end{equation}
with a nonnegative measure $U$ and the integrability needed for a finite
exponent.  In particular,
\begin{equation}\label{p01:rare:eq:stieltjes}
 h'(s)=d+\int_{(0,\infty)}\frac{U(dt)}{s+t}
\end{equation}
is a Stieltjes function and has nonpositive imaginary part in the upper
half-plane.  The classical implication from an HCM density to a GGC
law is stronger than infinite divisibility and will not be reversed.

\begin{lemma}[Zero-free continuation of HCM functions]
\label{p01:rare:lem:hcm-zero-free}
A strictly positive HCM function has a holomorphic, nowhere vanishing
extension to $D=\mathbb C\setminus(-\infty,0]$.  Consequently, if a
positive real function already has a holomorphic continuation to $D$
with a zero in $D$, it is not HCM.
\end{lemma}
\begin{proof}
The classical representation recalled in
\cite[eq.~(7)]{p01:rare:BondessonSimon2016} is
\[
 f(x)=C x^{\beta-1}h_1(x)h_2(1/x),\qquad
 - (\log h_i)'\text{ is Stieltjes},
\]
where $C>0$ and $\beta\in\mathbb R$ in the nonzero case.  Each
Stieltjes function $S_i$ extends holomorphically to $D$: the usual
integral representation converges uniformly on its compact subsets.
Since $D$ is simply connected, it has an analytic primitive.  Thus
\[
 \widetilde h_i(z)=h_i(1)
 \exp\left\{-\int_1^z S_i(w)\,dw\right\}
\]
is holomorphic and nowhere zero on $D$, and agrees with $h_i$ on the
positive axis.  Both $z\mapsto1/z$ and the principal power
$z^{\beta-1}$ are holomorphic on $D$, and the latter never vanishes.
Their product with these two functions gives the claimed extension.
Uniqueness of holomorphic continuation proves the last assertion.
\end{proof}

For $\mu>-1$ define the normalized entire Bessel factor
\begin{equation}\label{p01:rare:eq:normalized-bessel}
 F_{\mu,a}(z)=\sum_{k=0}^{\infty}
 \frac{(a^2z^2/4)^k}{k!(\mu+1)_k}
 =\Gamma(\mu+1)(2/(az))^\mu I_\mu(az).
\end{equation}
The series is the definition at the origin and removes any ambiguity
from fractional powers.  If $j_{\mu,n}$ is a positive zero of $J_\mu$,
then $F_{\mu,a}(i j_{\mu,n}/a)=0$.  Existence of these positive zeros
and the Bessel connection formulas are classical
\cite[\S\S10.21,10.27,10.42]{p01:rare:DLMF}.

\begin{theorem}[Finite products of modified Bessel factors]
\label{p01:rare:thm:bessel-products}
Let $r\geq1$, $a_j>0$, $\mu_j>-1$, $m_j\in\mathbb N$, and
$\beta\in\mathbb R$.  Suppose
\[
 \beta+\sum_{j=1}^r m_j\mu_j>-1,
 \qquad b>\sum_{j=1}^r m_j a_j.
\]
The normalized probability density proportional to
\begin{equation}\label{p01:rare:eq:bessel-product-density}
 x^\beta e^{-b x}\prod_{j=1}^r I_{\mu_j}(a_jx)^{m_j},\qquad x>0,
\end{equation}
is not HCM.
\end{theorem}
\begin{proof}
The local Bessel series and its exponential asymptotic at infinity
show integrability under the displayed assumptions.  Every factor is
positive for $x>0$, so the normalizing constant is positive and finite.
On $D$ the density extends as
\[
 C' z^{\beta+\sum_jm_j\mu_j}e^{-b z}
       \prod_{j=1}^r F_{\mu_j,a_j}(z)^{m_j},
\]
using the principal power.  All factors are holomorphic there.  At
$z=i j_{\mu_1,1}/a_1$ the first entire factor vanishes; every other
factor is finite, and the power and exponential are nonzero.  Hence the
extension has a zero in $D$.  Lemma~\ref{p01:rare:lem:hcm-zero-free} applies.
\end{proof}

\begin{corollary}[The HCM part of the McKay question]
\label{p01:rare:cor:four-mckay}
Each of the following densities, with its positive normalizing constant
suppressed, fails to be HCM throughout the indicated range:
\[
\begin{array}{lll}
 \varphi_{\mu,a,b}(x)&\propto x^\mu e^{-b x}I_\mu(ax),
   &\mu>-\tfrac12,\quad b>a>0,\\[2pt]
 \psi_{\mu,a,b}(x)&\propto x^{\mu+1}e^{-b x}I_\mu(ax),
   &\mu>-1,\quad b>a>0,\\[2pt]
 \varphi_{\mu,\nu,a,b}(x)&\propto x^{\nu-1}e^{-b x}I_\mu(ax),
   &\mu>-1,\quad \mu+\nu>0,\quad b>a>0,\\[2pt]
 \xi_{\mu,a,b}(x)&\propto x^{2\mu}e^{-b x}I_\mu(ax)^2,
   &\mu>-\tfrac14,\quad b>2a>0.
\end{array}
\]
In particular, this answers the HCM question for all parameters covered
by Theorems 1--4 of \cite{p01:rare:Baricz2026}.  The last conclusion holds
beyond the infinite-divisibility range proved in that source.
\end{corollary}
\begin{proof}
Apply Theorem~\ref{p01:rare:thm:bessel-products} with $(\beta,m)$ equal to
$(\mu,1)$, $(\mu+1,1)$, $(\nu-1,1)$, and $(2\mu,2)$, respectively.
\end{proof}

The obstruction concerns the density, not its Laplace transform.
The first McKay law is nevertheless a GGC, as already proved in
\cite[Theorem~1]{p01:rare:Baricz2026}.  Thus its non-HCM classification is
consistent with its positive GGC classification.  Salazar's recent
result treats a gamma density multiplied by ${}_0F_1$ with argument
linear in $x$; the factors here have argument proportional to $x^2$.
The above proof does not use a purported identification of these two
families.

\subsection{An exact SD/GGC classification for the second McKay family}

\begin{theorem}[Sharp classification of $\psi_{\mu,a,b}$]
\label{p01:rare:thm:psi-classification}
Let $\mu>-1$ and $b>a>0$, and let $X$ have density
\begin{equation}\label{p01:rare:eq:psi-density}
 \psi_{\mu,a,b}(x)=
 \frac{\sqrt\pi(b^2-a^2)^{\mu+3/2}}
 {2b(2a)^\mu\Gamma(\mu+3/2)}
 x^{\mu+1}e^{-b x}I_\mu(ax),\qquad x>0.
\end{equation}
This law is infinitely divisible, is never a GGC, and is
self-decomposable if and only if
\begin{equation}\label{p01:rare:eq:sd-threshold}
 (2\mu+3)\sqrt{b^2-a^2}\geq b.
\end{equation}
Its density is never HCM.  The infinite-divisibility assertion is
already present in \cite[Theorem~2]{p01:rare:Baricz2026}; the point here is
the exact additional classification.
\end{theorem}
\begin{proof}
Put $\alpha=\mu+3/2>1/2$.  Termwise integration of the nonnegative
series for $I_\mu$ and the gamma duplication formula give, for $c>a$,
\[
 \int_0^\infty x^{\mu+1}e^{-c x}I_\mu(ax)\,dx
 =\frac{2(2a)^\mu\Gamma(\alpha)}{\sqrt\pi}
     c(c^2-a^2)^{-\alpha}.
\]
Consequently the Laplace transform and exponent satisfy
\begin{align}
 L(s)&=(1+s/b)
  \left(\frac{b^2-a^2}{(s+b)^2-a^2}\right)^\alpha,
       \label{p01:rare:eq:psi-laplace}\\
 h'(s)&=\frac{\alpha}{s+b-a}
       +\frac{\alpha}{s+b+a}-\frac1{s+b}
       =\int_0^\infty e^{-s t}k(t)\,dt,
       \label{p01:rare:eq:psi-logderivative}\\
 k(t)&=e^{-bt}\bigl(2\alpha\cosh(at)-1\bigr).
       \label{p01:rare:eq:psi-k}
\end{align}
Here $k>0$ and $k(t)/t$ is a L\'evy density.  Integration in $s$ gives
$h(s)=\int_0^\infty(1-e^{-s t})k(t)\,dt/t$, proving infinite
divisibility directly.

If the law were GGC, its logarithmic derivative would have the
Stieltjes property \eqref{p01:rare:eq:stieltjes}.  At $z=-b+i\eta$,
\[
 \operatorname{Im}h'(z)=\frac1\eta
             -\frac{2\alpha\eta}{a^2+\eta^2}>0
 \quad\text{when}\quad
 0<\eta^2<\frac{a^2}{2\alpha-1}.
\]
Analytic continuation of \eqref{p01:rare:eq:psi-logderivative} is unique
in the upper half-plane, so this contradicts the Stieltjes property.

For this smooth, exponentially decreasing $k$, self-decomposability is
equivalent to $k$ being nonincreasing.  We recall the argument to make
the necessity as well as sufficiency explicit.  If $k$ is
nonincreasing and $0<c<1$, then
\[
 h(s)-h(cs)=\int_0^\infty(1-e^{-s t})
              \frac{k(t)-k(t/c)}t\,dt
\]
is a Bernstein function vanishing at zero.  Hence $L(s)/L(cs)$ is a
Laplace transform, which is exactly self-decomposability.  Conversely,
if the law is self-decomposable, each $R_c(s)=L(s)/L(cs)$ is a Laplace
transform.  Thus $(1-R_c)/(1-c)$ is Bernstein, and its limit as
$c\uparrow1$ is $s h'(s)$.  Closure of Bernstein functions under finite
pointwise limits shows that $s h'(s)$ is Bernstein.  Integration by
parts gives
\[
 s h'(s)=\int_0^\infty(1-e^{-s t})(-k'(t))\,dt.
\]
The right side is a finite signed L\'evy representation, and its limit
at infinity is $k(0)$.  The positive Bernstein representation therefore
has zero drift and a finite measure.  Uniqueness of Laplace transforms
of finite measures forces that measure to be $-k'(t)\,dt$, so $k'\leq0$.

Finally,
\[
 k'(t)=e^{-bt}\{b-2\alpha[b\cosh(at)-a\sinh(at)]\}.
\]
The minimum of the bracketed hyperbolic expression on $t\geq0$ is
$\sqrt{b^2-a^2}$, attained at
$t=a^{-1}\operatorname{arctanh}(a/b)$.  Therefore $k'\leq0$ is
equivalent to $2\alpha\sqrt{b^2-a^2}\geq b$, which is
\eqref{p01:rare:eq:sd-threshold}.  The non-HCM assertion follows from
Corollary~\ref{p01:rare:cor:four-mckay}.
\end{proof}

The exclusion of GGC holds even in the nonempty SD region.  For
example, $\mu=0$, $b=2$, and $a=1$ satisfy
\eqref{p01:rare:eq:sd-threshold}.  This explicitly prevents the invalid
inference that failure of a Stieltjes test proves failure of
self-decomposability.  In the limiting case $a\downarrow0$, the law
converges to a gamma law of shape $2\mu+2$ and rate $b$; the strict
assumption $a>0$ in the negative GGC and HCM assertions is essential.

\subsection{Noncentral chi-square laws and the unresolved third problem}

\begin{proposition}[The noncentral chi-square GGC obstruction]
\label{p01:rare:prop:noncentral-ggc}
A noncentral chi-square law with degrees of freedom $\kappa>0$ and
noncentrality $\lambda\geq0$ is GGC if and only if $\lambda=0$.
Consequently its density is HCM if and only if $\lambda=0$.
The latter classification is already given by \cite{p01:rare:Salazar2026};
the following is a short independent GGC argument, without a priority
claim.
\end{proposition}
\begin{proof}
Its Laplace transform is
$L(s)=(1+2s)^{-\kappa/2}\exp\{-\lambda s/(1+2s)\}$, so
\[
 h'(s)=\frac{\kappa}{1+2s}+\frac{\lambda}{(1+2s)^2}.
\]
At $z=-1/2-\varepsilon+i\varepsilon$ one has
\[
 \operatorname{Im}h'(z)
    =-\frac{\kappa}{4\varepsilon}
       +\frac{\lambda}{8\varepsilon^2}>0
       \quad\text{if }0<\varepsilon<\frac{\lambda}{2\kappa}.
\]
For $\lambda>0$ this violates the Stieltjes condition necessary for
GGC.  An HCM density would imply GGC and is therefore impossible.
For $\lambda=0$ the law is gamma, whose density is HCM and whose
Laplace exponent has a one-atom Thorin representation.
\end{proof}

The third problem in \cite[\S4.4]{p01:rare:Baricz2026} concerns
self-decomposability of the laws in its Theorems 14, 18, and 22.
One can make a concrete reduction for Theorem 18.  With
$\mu,\nu>1/2$ and $a,b>0$, its Laplace transform is
\begin{equation}\label{p01:rare:eq:third-laplace}
 L(s)=e^{-(a+b)\sqrt s}
 {}_0F_1(;\mu+1;a^2s/4)
 {}_0F_1(;\nu+1;b^2s/4).
\end{equation}
The Bessel product formula gives
\begin{equation}\label{p01:rare:eq:third-logderivative}
 h'(s)=\frac{a+b}{2\sqrt s}
    -\sum_{n\geq1}\frac1{s+j_{\mu,n}^2/a^2}
    -\sum_{n\geq1}\frac1{s+j_{\nu,n}^2/b^2}.
\end{equation}
Hence the L\'evy spectral function is
\begin{equation}\label{p01:rare:eq:third-k}
 k(t)=\frac{a+b}{2\sqrt{\pi t}}
    -\sum_{n\geq1}e^{-j_{\mu,n}^2t/a^2}
    -\sum_{n\geq1}e^{-j_{\nu,n}^2t/b^2}.
\end{equation}
The series and all their derivatives converge locally uniformly for
$t>0$, by $j_{\eta,n}\sim\pi n$.  Their Laplace transforms may be
integrated termwise for $s>0$, since
$\sum_n(s+j_{\eta,n}^2/c^2)^{-1}<\infty$.  Thus
\eqref{p01:rare:eq:third-k} is the inverse Laplace transform of
\eqref{p01:rare:eq:third-logderivative}.  The already proved infinite
divisibility in the source ensures that it is nonnegative, by
uniqueness of the Bernstein measure.  The standard L\'evy criterion
then says that self-decomposability is equivalent to
\begin{equation}\label{p01:rare:eq:third-unsolved-inequality}
 \sum_{n\geq1}\frac{j_{\mu,n}^2}{a^2}e^{-j_{\mu,n}^2t/a^2}
 +\sum_{n\geq1}\frac{j_{\nu,n}^2}{b^2}e^{-j_{\nu,n}^2t/b^2}
 \leq\frac{a+b}{4\sqrt\pi\,t^{3/2}}
 \qquad(t>0).
\end{equation}
This identifies the missing heat-trace inequality exactly.  Neither
the negative residues in \eqref{p01:rare:eq:third-logderivative} nor
positivity of $k$ proves it.  No proof of
\eqref{p01:rare:eq:third-unsolved-inequality}, and no classification for
all three families in the third source problem, is asserted here.

\section{Blind calibration with arbitrary local correlations}
\label{p01:rare:sec:blind-correlated}

The scalar calibration argument in the integrated rare-event theory
uses signs of mixed second derivatives that can fail for correlated
local vectors.  There is another exact route: use first derivatives
at several points instead of insisting on a second-order calculation
at one point.  This is an extension of that manuscript, not a claim
to resolve an externally documented longstanding conjecture.

Let $Y=(Y_1,\ldots,Y_m)$ be a bounded nonnegative random vector such
that
\begin{equation}\label{p01:rare:eq:gated-assumptions}
 p_0=\mathbb P(Y=0)>0,
 \qquad \mathbb P(Y_j>0)>0\quad(1\leq j\leq m).
\end{equation}
Its coordinates may have arbitrary dependence.  For $p_j>1$ put
\begin{equation}\label{p01:rare:eq:blind-H}
 h(v)=-\log\mathbb E e^{-v\cdot Y},\qquad
 H(s)=\int_0^\infty h(s_1x^{-p_1},\ldots,s_mx^{-p_m})\,dx.
\end{equation}
For an original counting exponent $\rho>0$ and channel powers $F_j$,
the substitution $x=y^\rho$ gives precisely $p_j=F_j/\rho$.
Thus these ratios, rather than an unnormalized $\rho$, are the
identifiable parameters in the next theorem.

\begin{lemma}[Analyticity and visible axes]
\label{p01:rare:lem:analytic-axes}
Under \eqref{p01:rare:eq:gated-assumptions}, $H$ is finite and continuous
on $[0,\infty)^m$, real analytic on $(0,\infty)^m$, and satisfies
\begin{align}
 H(c^{p_1}s_1,\ldots,c^{p_m}s_m)&=c H(s),\qquad c>0,
       \label{p01:rare:eq:homogeneity}\\
 H(te_j)&=C_jt^{1/p_j},\qquad t\geq0,
       \quad 0<C_j<\infty.\label{p01:rare:eq:axis-law}
\end{align}
\end{lemma}
\begin{proof}
The estimates
$0\leq h(v)\leq-\log p_0$ and
$h(v)\leq\sum_j v_j\mathbb E Y_j$ control $x\leq1$ and $x\geq1$,
respectively.  They give convergence and, uniformly for bounded $s$,
an integrable majorant; continuity follows by dominated convergence.
The substitution $x=cy$ proves \eqref{p01:rare:eq:homogeneity}, and the
substitution $x=t^{1/p_j}y$ on the $j$th axis proves
\eqref{p01:rare:eq:axis-law}.  Nontriviality of $Y_j$ gives $C_j>0$.

For analyticity fix a positive $s^0$, put
$S=\sum_j s_j^0x^{-p_j}Y_j$, and allow complex perturbations satisfying
$|s_j-s_j^0|\leq\delta s_j^0$.  If
$\Delta=\sum_j(s_j-s_j^0)x^{-p_j}Y_j$, then $|\Delta|\leq\delta S$.
Writing $L_0=\mathbb E e^{-S}\geq p_0$, we have
\[
 |\mathbb E e^{-(S+\Delta)}-L_0|
 \leq\delta\mathbb E[S e^{-(1-\delta)S}]
 \leq\frac{\delta}{e(1-\delta)}.
\]
Choose $\delta$ so that the last bound is below $p_0/2$.  A single
local analytic logarithm is then defined by expansion about $L_0$,
uniformly in $x$.  For $x\leq1$ its absolute value is bounded by a
constant depending on $p_0$ and $\delta$.  For large $x$, boundedness
of $Y$ gives the sharper bound $O(\sum_jx^{-p_j})$, uniformly in
this complex polydisc.  Holomorphic parameter integration proves
analyticity of $H$.
\end{proof}

\begin{theorem}[Finite-point blind calibration without covariance signs]
\label{p01:rare:thm:correlated-calibration}
For every nonempty open set $U\subset(0,\infty)^m$, the functions
\[
 \phi_j(s)=s_j\partial_jH(s),\qquad 1\leq j\leq m,
\]
are linearly independent on $U$.  Consequently, for an open dense
set of $m$-tuples $(s^{(1)},\ldots,s^{(m)})\in U^m$, whose complement
has Lebesgue measure zero, the matrix and vector
\begin{equation}\label{p01:rare:eq:calibration-matrix}
 A_{\ell j}=s_j^{(\ell)}\partial_jH(s^{(\ell)}),
 \qquad b_\ell=H(s^{(\ell)})
\end{equation}
satisfy
\begin{equation}\label{p01:rare:eq:calibration-solution}
 \det A\ne0,\qquad (p_1,\ldots,p_m)^T=A^{-1}b.
\end{equation}
In particular, the unknown power ratios are determined by the exact
first derivatives and values of $H$ at $m$ generic points of any
open observation patch.
\end{theorem}
\begin{proof}
Differentiating \eqref{p01:rare:eq:homogeneity} at $c=1$ gives
\begin{equation}\label{p01:rare:eq:euler}
 \sum_jp_j\phi_j(s)=H(s).
\end{equation}
Suppose $\sum_jv_j\phi_j=0$ on $U$ for constants $v_j$.  Analyticity
extends this identity to the connected positive orthant.  For every
$s>0$ the derivative of
$t\mapsto H(e^{t v_1}s_1,\ldots,e^{t v_m}s_m)$ is zero for all
$t\in\mathbb R$.  Hence this function is constant.  Passing to the
$j$th axis using continuity and \eqref{p01:rare:eq:axis-law} yields
\[
 C_j(e^{t v_j}s_j)^{1/p_j}=C_js_j^{1/p_j}
 \qquad(t\in\mathbb R).
\]
Since $C_j>0$, one has $v_j=0$ for every $j$.

The vectors $(\phi_1(s),\ldots,\phi_m(s))$, $s\in U$, therefore
span $\mathbb R^m$.  One may choose $m$ of them as a basis, proving
that the determinant in \eqref{p01:rare:eq:calibration-matrix} is not
identically zero as an analytic function on $((0,\infty)^m)^m$.
Its zero set has empty interior and Lebesgue measure zero, which proves
the generic assertion on $U^m$.  Applying \eqref{p01:rare:eq:euler} at
these points gives $Ap=b$ and the formula follows.
\end{proof}

No pairwise distinction among the exponents is needed in this theorem.
Its visibility condition is essential: if $Y_j=0$ almost surely, the
response is independent of the $j$th channel and cannot determine its
power.  The result provides exact identifiability, not a lower bound
for $|\det A|$ and not stability under measurement noise.

\begin{corollary}[Calibration from a smooth active rate patch]
\label{p01:rare:cor:rate-calibration}
Suppose $-\nabla^2H$ is positive definite on an open positive patch,
and let $e=\nabla H(s)$ there.  For the lower-tail dual
$I(e)=\sup_{u\geq0}\{H(u)-u\cdot e\}$, an open active patch determines
all $p_j$.  Generically $m$ first jets of $I$ suffice: use
\begin{equation}\label{p01:rare:eq:rate-matrix}
 A_{\ell j}=-e_j^{(\ell)}\partial_jI(e^{(\ell)}),\qquad
 b_\ell=I(e^{(\ell)})-e^{(\ell)}\cdot\nabla I(e^{(\ell)}).
\end{equation}
A sufficient condition for the required strict Hessian inequality on
the whole positive orthant is that the powers $p_j$ be pairwise
distinct, in addition to \eqref{p01:rare:eq:gated-assumptions}.
\end{corollary}
\begin{proof}
The inverse function theorem and strict concavity give
$s=-\nabla I(e)$, $\nabla H(s)=e$, and
$H(s)=I(e)-e\cdot\nabla I(e)$.  Substitution in
\eqref{p01:rare:eq:calibration-matrix} proves the assertion.

For the sufficient condition, differentiation under the integral gives
\[
 -v^T\nabla^2H(s)v
 =\int_0^\infty
 \operatorname{Var}_{s,x}\left(\sum_jv_jx^{-p_j}Y_j\right)dx,
\]
where the tilt has density proportional to
$\exp(-\sum_js_jx^{-p_j}Y_j)$.  The derivatives are integrable by
the analytic bounds in Lemma~\ref{p01:rare:lem:analytic-axes}.
If the integral is zero, the variance vanishes for almost every $x$.
The tilted law is equivalent to the original law and retains a positive
atom at $Y=0$, so the almost sure constant is zero.  Fubini's theorem
then shows, for almost every $Y$, that
$\sum_jv_jY_jx^{-p_j}=0$ for almost every $x>0$.  This continuous
function vanishes everywhere, and linear independence of distinct
powers gives $v_jY_j=0$ almost surely for every $j$.  The coordinate
visibility assumptions imply $v=0$.
\end{proof}

\begin{corollary}[Some finite jet at one point determines the powers]
\label{p01:rare:cor:one-point-finite-jet}
At every fixed $s^0>0$ there is a finite integer $r=r(H,s^0)$ such
that the derivatives of $H$ at $s^0$ through order $r$ determine all
power ratios.  No uniform bound on $r$ is asserted.
\end{corollary}
\begin{proof}
If the vectors
$\bigl(\partial^\alpha\phi_1(s^0),\ldots,
\partial^\alpha\phi_m(s^0)\bigr)$ did not span $\mathbb R^m$,
a nonzero vector would annihilate all of them.  The Taylor series and
analytic continuation would give a nontrivial constant linear relation
among the $\phi_j$, contrary to
Theorem~\ref{p01:rare:thm:correlated-calibration}.  Select finitely many
multi-indices giving a basis and differentiate \eqref{p01:rare:eq:euler}
at those multi-indices.  This gives a nonsingular linear system for
$p$, involving only finitely many derivatives of $H$.
\end{proof}

\paragraph{Connection with exact recovery of the local law.}
In the repeated-coordinate design of the integrated manuscript, use
$Y=(X_1,\ldots,X_d,X_1)$ and assign the last coordinate a power
different from the first.  For a law with a positive atom at zero and
remaining support in a compact box $[a,b]^d$, $a>0$, the preceding
calibration identifies all the power ratios.  The known-power
repeated-coordinate Mellin inversion can then be applied to the exact
response.  Thus arbitrary local correlations are no obstruction to
blind calibration when an open response patch is available.  This
combination still requires a function-valued observation to recover
the entire microscopic law; the finite calibration data alone do not
identify that law.  The unavoidable common rescaling of all original
powers and the counting exponent is also unchanged.


% ===== Consolidated body: mixtures =====
\section{Exchangeable mixtures and differential obstructions: a source audit}
\label{p01:mix:sec:source-audit}

This branch concerns three externally posed targets. Two are precise
yes-or-no questions; the third is a broad question about restrictions
that restore statistical estimation. They must not be counted as four
questions by splitting ordinary and shifted differential operators.

\begin{enumerate}
\item \textbf{Normal directing distributions.}
Jeff Lin, \emph{A Negative Answer to a Problem of Aldous on Determination
of Exchangeable Sequences}, arXiv:1507.00678v1, Section~2.6, printed
page~17 (PDF page index~16)\cite{p01:mix:Lin2015}, explicitly leaves open whether the class
of all normal distributions is good. Proposition~2.6.4 gives a positive
answer when the random variance has an exponential moment.
The exact source inspected is
\url{https://arxiv.org/pdf/1507.00678v1}.
The journal article is \emph{Electronic Journal of Probability} 21
(2016), DOI \texttt{10.1214/16-EJP4403}; its publisher PDF could not be
inspected in this audit. The complete no-moment solution is already
in the integrated draft, and is rederived below. It is not a new
result of the present attack.

\item \textbf{The converse to the differential barrier.}
Le, Bariletto, Rinaldo, and Ho, \emph{Partial Differential Equation
Barriers to Identifiability in Infinite Mixture Models}, author
manuscript dated 8 August 2026, remark following Proposition~6 and the
first question of Section~6, printed page~13 (PDF page index~12)\cite{p01:mix:Le2026}, ask
whether absence of the stated differential structure implies
identifiability. The inspected source is
\url{https://nhatptnk8912.github.io/Identifiability_Infinite_Mixture_final.pdf};
the current arXiv record is
\url{https://arxiv.org/abs/2608.08597}.
The integrated draft already contains a complete negative answer for
both finite operator classes in Definitions~1 and~4. Its proof is
rechecked below. Arbitrary nonlocal operators are outside the claim.

\item \textbf{Restrictions restoring consistent estimation.}
The second question in Section~6 of the same source asks which model
or truth constraints rule out invisible perturbations and permit
consistent estimation. This is a research program, not one precise
universal conjecture. The quadrature-node restriction proved below
gives a new special-case answer for kernels with a finite polynomial
defect. It does not close the full question for all kernel families.
\end{enumerate}

Targeted web and SciSpace searches located the two source papers and
nearby finite-mixture literature but no subsequent resolution of these
exact formulations. This bounded search does not establish priority or
the present status of every version. Lin's general Aldous question,
its unrestricted continuous-time analogue, and the unrestricted random
moment question already receive negative answers in his paper; they
are not counted as new open targets here.

\subsection{Independent rederivation of the two precise draft closures}

\begin{proposition}[The normal case without exponential moments]
\label{p01:mix:prop:normal-audit}
Let $(A,V)$ have any Borel probability law on
$\mathbb R\times[0,\infty)$ and, conditionally on it, let
$X_j=A+\sqrt V Z_j$ with independent standard normal $Z_j$.
If $\mathcal N\subset\mathbb N$ satisfies
$\sum_{n\in\mathcal N}1/n=\infty$, the separate laws of
$S_n=\sum_{j=1}^nX_j$, $n\in\mathcal N$, determine the joint law
of $(A,V)$.
\end{proposition}
\begin{proof}
For the signed difference $\rho$ of two candidate laws and fixed
$t\ne0$, define
\[
 H_t(z)=\int e^{ita-zv}\,d\rho(a,v),\qquad \operatorname{Re}z>0.
\]
It is bounded by $\|\rho\|_{\rm TV}$ and holomorphic: on
$\operatorname{Re}z\ge\epsilon>0$, every derivative of the integrand
is bounded by a constant times $v^ke^{-\epsilon v}$. Equality of the
laws of $S_n/n$ gives $H_t(t^2/(2n))=0$.
A nonzero bounded holomorphic function in the right half-plane has
zeros satisfying
$\sum_j\operatorname{Re}z_j/(1+|z_j|^2)<\infty$.
Indeed, the map $w=(z-1)/(z+1)$ and Jensen's formula give the disk
condition $\sum_j(1-|w_j|)<\infty$, and
$1-|w|^2=4\operatorname{Re}z/|z+1|^2$ makes the two conditions
equivalent. Here the sum diverges, so $H_t$ is zero. Continuity in
$t$ includes $t=0$.
Consequently the finite measure $e^{-v}\rho(da,dv)$ has Fourier
transform $H_t(1-iu)=0$ for every $t,u$. Fourier uniqueness gives
this measure, and hence $\rho$, equal to zero. No extension through
the boundary point $z=0$ and no variance moment were used.
\end{proof}

\begin{proposition}[A counterexample to the finite differential converse]
\label{p01:mix:prop:pde-audit}
For $x,\theta\in[0,1]$, let $L_n(t)=P_n(2t-1)$ be shifted Legendre
polynomials and set
\[
 k(x\mid\theta)=1+\sum_{n=2}^{\infty}\frac{L_n(x)L_n(\theta)}{n!}.
\]
This is a normalized, strictly positive, jointly entire kernel. Its
signed mixing nullspace is exactly
$\{cL_1(\theta)d\theta:c\in\mathbb R\}$, yet no nonzero finite
point-derivative functional annihilates all its parameter sections.
In particular it has neither a finite-order parameter differential
annihilator nor a finite-shift, finite-order differential annihilator
as defined in the cited source.
\end{proposition}
\begin{proof}
On the unit square $|L_n|\le1$, so $k\ge1-(e-2)=3-e>0$.
Legendre orthogonality gives normalization. On a complex compact
set, $|L_n|$ grows at most exponentially in $n$, so the factorial
series and every fixed derivative converge locally uniformly.
For a finite signed measure $\nu$, coefficient extraction gives its
mass and every Legendre moment except the first. Thus $A\nu=0$
implies
\[
 \int p\,d\nu=3\left(\int L_1\,d\nu\right)\int_0^1pL_1
 \quad\hbox{for every polynomial }p.
\]
Polynomial density identifies the asserted nullspace. Conversely
orthogonality proves that every displayed measure is invisible.
The distinct positive mixing densities $1\pm cL_1$, $0<c<1$,
therefore produce the same uniform observation density.

Suppose
$T(p)=\sum_{i=1}^J\sum_{j=0}^{r_i}b_{ij}p^{(j)}(\theta_i)$,
with distinct nodes, annihilates $k(x\mid\cdot)$ for almost every
$x$. Coefficient extraction, including parameter differentiation,
gives $T(1)=0$ and $T(L_n)=0$ for $n\ge2$. Hence
$T(p)=3T(L_1)\int pL_1$ on polynomials. With
$Q(t)=\prod_i(t-\theta_i)^{r_i+1}$,
\[
 0=T(Q^2L_1)=3T(L_1)\int_0^1Q^2L_1^2.
\]
The integral is strictly positive, so $T(L_1)=0$. Hermite
interpolation then forces every $b_{ij}=0$. Complex coefficients
are handled by real and imaginary parts. At each fixed base
parameter, a finite differential or finite shifted differential
relation is such a functional, after combining coincident nodes.
Its coefficients must therefore vanish. The nullspace above remains
an explicit nonlocal integral obstruction; it is not being denied.
\end{proof}

No fatal mathematical gap was found in these two arguments. Their
historical-priority status remains separate from their validity.

\section{A new sharpening: exact sparse sampling for every fixed stable shape}
\label{p01:mix:sec:stable-sharp}

The integrated draft proves the sufficient conditions in the following
theorem and proves the converse only for the Gaussian case. The
necessity for every other fixed stable shape is the new extension here.
It addresses a natural sharpening of the normal-case target, not an
additional independently posed published conjecture.

For $0<\alpha\le2$, $-1\le\beta\le1$, put
\[
 \psi_{\alpha,\beta}(t)=
 \begin{cases}
 |t|^\alpha\{1-i\beta\tan(\pi\alpha/2)\operatorname{sgn}t\},
       &\alpha\ne1,\\
 |t|\{1+i(2\beta/\pi)\operatorname{sgn}t\log|t|\},
       &\alpha=1,
 \end{cases}
 \qquad\psi_{\alpha,\beta}(0)=0.
\]
Write $K_{a,b}$ for the stable law with characteristic function
$\exp\{ita-b\psi_{\alpha,\beta}(t)\}$, where
$(a,b)\in\mathbb R\times[0,\infty)$. Conditional on one draw
$(A,B)$ from the mixing law, draw the entire iid sequence from
$K_{A,B}$. The index and skewness are fixed throughout a model.

\begin{theorem}[Necessary and sufficient sparse stable sampling]
\label{p01:mix:thm:sharp-stable}
Let $\mathcal N\subset\mathbb N$. If $\alpha\ne1$, the separate
partial-sum laws $(\mathcal L(S_n))_{n\in\mathcal N}$ identify
every Borel mixing probability on $\mathbb R\times[0,\infty)$
if and only if
\begin{equation}\label{p01:mix:eq:stable-power-criterion}
 \sum_{n\in\mathcal N}n^{-|1-\alpha|}=\infty.
\end{equation}
If $\alpha=1$ and $\beta\ne0$, the corresponding necessary and
sufficient condition is
\begin{equation}\label{p01:mix:eq:stable-log-criterion}
 \sum_{n\in\mathcal N}\frac1{1+(\log n)^2}=\infty.
\end{equation}
For $(\alpha,\beta)=(1,0)$, no choice of $\mathcal N$ gives
injectivity, even when all positive integers are observed.

Whenever the appropriate sum converges in either of the first two
cases, nonidentifiability has witnesses that are absolutely continuous,
have all mixed polynomial moments finite and equal, and have the same
$B$ marginal. For $\alpha\ne1$ they also have the same $A$ marginal.
Any prescribed lower bound $B\ge b_0>0$ can be imposed. No moment
assumption is required in the positive directions.
\end{theorem}

\subsection{A regularized Blaschke-product construction}

\begin{lemma}[Half-line Schwartz realizations of the required zeros]
\label{p01:mix:lem:blaschke-realization}
Each of the following zero sets admits a nonzero real Schwartz
function $f$ on $\mathbb R$, supported in $[0,\infty)$, whose
Laplace transform $F$ vanishes on that set and satisfies $F(0)=0$:
\begin{enumerate}
\item $\{c\lambda_n,\overline c\lambda_n:n\in\mathcal N\}$,
where $\operatorname{Re}c=1$, and either
$0<\lambda_n\le1$, $\sum_n\lambda_n<\infty$, or
$\lambda_n\ge1$, $\sum_n\lambda_n^{-1}<\infty$;
\item $\{1+iy_n,1-iy_n:n\in\mathcal N\}$, where
$\sum_n(1+y_n^2)^{-1}<\infty$.
\end{enumerate}
Finite or empty sets are allowed.
\end{lemma}
\begin{proof}
For a zero $\zeta$ in the right half-plane let
$b_\zeta(z)=(z-\zeta)/(z+\overline\zeta)$. Its modulus is at most
one there and exactly one on the imaginary axis. Pair conjugate zeros
and form
\[
 P(z)=\prod_n b_{\zeta_n}(z)b_{\overline\zeta_n}(z)
 =\prod_n\frac{z^2-2(\operatorname{Re}\zeta_n)z+|\zeta_n|^2}
                  {z^2+2(\operatorname{Re}\zeta_n)z+|\zeta_n|^2}.
\]
For the three cases of the statement, the difference of each factor
from one on a fixed compact subset of $\operatorname{Re}z>0$ is,
respectively, $O(\lambda_n)$, $O(\lambda_n^{-1})$, or
$O((1+y_n^2)^{-1})$. Thus the product converges locally uniformly,
is not identically zero, has the required zeros, obeys $|P|\le1$,
and has the conjugation symmetry $P(\overline z)=\overline{P(z)}$.

We justify the boundary regularity needed for a joint signed-density
construction. Put $Q=P'/P$ away from the zeros. For one factor,
\[
 \frac{b_\zeta'(z)}{b_\zeta(z)}
 =\frac1{z-\zeta}-\frac1{z+\overline\zeta}
 =\frac{2\operatorname{Re}\zeta}
        {(z-\zeta)(z+\overline\zeta)}.
\]
If $\zeta=(1+i\gamma)\lambda$, then
\[
 |i\xi-\zeta|=|i\xi+\overline\zeta|
 =\{\lambda^2+(\xi-\gamma\lambda)^2\}^{1/2}
 \ge C_\gamma|\xi|,
 \qquad |i\xi-\zeta|\ge\lambda.
\]
Differentiating the rational expression therefore gives, for every
$r\ge0$, an upper bound $C_r\lambda|\xi|^{-r-2}$ in the small
$\lambda$ case, and $C_r\lambda^{-r-1}$ in the large $\lambda$
case. Summation is valid locally uniformly on the imaginary axis
away from zero in the first case, and on the whole axis in the
second. It follows that every fixed boundary derivative of $Q$
has at most polynomial growth at zero and infinity.

For $\zeta=1+iy$, the same derivative is at most
$C_r/(1+(\xi-y)^2)$, because both denominator moduli are at least
one. The elementary inequality
\[
 \frac1{1+(\xi-y)^2}
 \le\frac{2(1+\xi^2)}{1+y^2}
\]
therefore bounds each derivative of the sum by
$C_r(1+\xi^2)\sum_n(1+y_n^2)^{-1}$, including the conjugate terms.
The recursion
$P^{(j+1)}/P=(P^{(j)}/P)'+Q(P^{(j)}/P)$ shows that boundary
derivatives of $P$ have at most polynomial growth in the same sense;
we use $|P(i\xi)|=1$ wherever the boundary product is defined.

With principal square roots set
\[
 F(z)=P(z)\exp\{-\sqrt z-1/\sqrt z\},
 \qquad\operatorname{Re}z>0.
\]
The boundary function $H(\xi)=F(i\xi)$, extended by $H(0)=0$, is
Schwartz and flat at zero, since
\[
 \operatorname{Re}\sqrt{i\xi}=\sqrt{|\xi|/2},\qquad
 \operatorname{Re}(i\xi)^{-1/2}=1/\sqrt{2|\xi|}
\]
dominate all the preceding polynomial derivative bounds. Define
\[
 f(b)=\frac1{2\pi}\int_{\mathbb R}e^{i\xi b}H(\xi)\,d\xi.
\]
Fourier inversion makes $f$ real, nonzero and Schwartz.
For $b<0$, shift the vertical contour to $R+i\mathbb R$.
The horizontal sides vanish because, for $x\ge0$,
$\operatorname{Re}\sqrt{x+i\xi}\ge\sqrt{|\xi|/2}$.
Starting at $\epsilon+i\mathbb R$ and then taking
$\epsilon\downarrow0$ justifies the boundary contour. Hence
\[
 |f(b)|\le\frac{e^{bR}}{2\pi}
              \int_{\mathbb R}e^{-\sqrt{|\xi|/2}}\,d\xi.
\]
Letting $R\to\infty$ proves $f(b)=0$ for $b<0$.
Fubini and Cauchy's formula on right half-disks give
\[
 \int_0^\infty e^{-sb}f(b)\,db
 =\frac1{2\pi}\int_{\mathbb R}\frac{H(\xi)}{s-i\xi}\,d\xi
 =F(s),\qquad s>0.
\]
The large semicircle vanishes by the square-root exponential bound;
the origin can be removed by a small semicircle, since $F$ tends to
zero there. Finally $F(s)\to0$ as $s\downarrow0$, so dominated
convergence gives $\int f=0$. The formula also holds for complex
$s$ with positive real part by holomorphy.
\end{proof}

\subsection{Proof of the exact stable criteria}

\begin{proof}[Proof of Theorem~\ref{p01:mix:thm:sharp-stable}]
For the difference $\rho$ of two mixing laws, retain
$H_t(z)=\int e^{ita-zb}\,d\rho(a,b)$.
The law of $S_n/n$ supplies the value at
$z_n(t)=n\psi_{\alpha,\beta}(t/n)$. If $\alpha\ne1$, write
$c_t=1-i\beta\tan(\pi\alpha/2)\operatorname{sgn}t$. Then
\[
 z_n(t)=|t|^\alpha c_t n^{1-\alpha},\qquad
 \frac{\operatorname{Re}z_n(t)}{1+|z_n(t)|^2}
 \asymp_t n^{-|1-\alpha|}.
\]
These points are distinct and lie in the open right half-plane.
The half-plane zero condition in
Proposition~\ref{p01:mix:prop:normal-audit}, followed by Fourier--Laplace
uniqueness, proves sufficiency of
\eqref{p01:mix:eq:stable-power-criterion}.
For $\alpha=1$, put $\kappa=2\beta/\pi$. When $\kappa\ne0$,
\[
 z_n(t)=|t|+i\kappa t\log|t|-i\kappa t\log n,
 \qquad
 \frac{\operatorname{Re}z_n(t)}{1+|z_n(t)|^2}
 \asymp_t \frac1{1+(\log n)^2}.
\]
The same proof gives sufficiency of
\eqref{p01:mix:eq:stable-log-criterion}.

For necessity when $\alpha\ne1$, set
$\lambda_n=n^{1-\alpha}$ and
$c=1-i\beta\tan(\pi\alpha/2)$. Convergence of the appropriate
sum is exactly one of the first two hypotheses of
Lemma~\ref{p01:mix:lem:blaschke-realization}. Obtain its real function
$f$ and Laplace transform $F$ with zeros $c\lambda_n$ and
$\overline c\lambda_n$. Choose a nonzero real even
$h\in C_c^\infty(\{t:1<|t|<2\})$ and put
\[
 G(t,b)=h(t)|t|^\alpha f(b|t|^\alpha),\qquad
 g(a,b)=\frac1{2\pi}\int_{\mathbb R}e^{-ita}G(t,b)\,dt.
\]
All derivatives of $G$ are rapidly decreasing in $b$, uniformly on
its compact $t$-support. Fourier integration by parts makes
$g\in\mathcal S(\mathbb R^2)$. It is real, nonzero, supported in
$b\ge0$, and has total integral zero because $G$ vanishes near $t=0$.
Its Fourier--Laplace transform is
\begin{equation}\label{p01:mix:eq:homogeneous-transform}
 \int e^{ita-sb}g(a,b)\,da\,db
       =h(t)F(s/|t|^\alpha),\qquad t\ne0,
       \quad\operatorname{Re}s>0.
\end{equation}
Substituting $s=z_n(t)$ makes this zero for every selected $n$.

For $\alpha=1$, $\kappa\ne0$, use part~(2) of the lemma with
$y_n=\kappa\log n$, and define instead
\[
 G(t,b)=h(t)|t| f(b|t|)\exp\{i\kappa t b\log|t|\},\qquad
 g(a,b)=\frac1{2\pi}\int e^{-ita}G(t,b)\,dt.
\]
Here the displayed formula is used for $t\ne0$, and $G$ is set
equal to zero near $t=0$, as permitted by the support of $h$.
The modulation introduces only polynomial factors in $b$ upon
differentiation, so the same Schwartz and support conclusions hold.
Moreover $G(-t,b)=\overline{G(t,b)}$, which makes $g$ real.
Now
\begin{equation}\label{p01:mix:eq:logarithmic-transform}
 \int e^{ita-sb}g(a,b)\,da\,db
 =h(t)F\left(s/|t|-i\kappa\operatorname{sgn}(t)\log|t|\right),
 \quad t\ne0,\quad\operatorname{Re}s>0.
\end{equation}
At $s=z_n(t)$ the argument is
$1-i\kappa\operatorname{sgn}(t)\log n$, one of the prescribed
zeros. This frequency-dependent modulation is needed: simply
reusing the homogeneous scaling would not cancel the logarithmic
location term.

In both constructions write $g=g_+-g_-$ and
$m=\int g_+=\int g_->0$. The normalized Jordan parts
$m^{-1}g_+\,da\,db$ and $m^{-1}g_-\,da\,db$ are distinct
absolutely continuous mixing probabilities with identical selected
partial-sum laws. Schwartz decay gives every polynomial moment.
For each fixed $b$,
\[
 \int a^j g(a,b)\,da=i^{-j}\partial_t^jG(0,b)=0
 \qquad(j\ge0).
\]
Thus the $B$ marginals and all mixed polynomial moments agree.
In the homogeneous case, $\int G(t,b)db=h(t)\int f=0$, so the
$A$ marginals also agree by Fourier uniqueness. That extra marginal
claim is not asserted for the modulated index-one construction.
Replacing $B$ by $B+b_0$ multiplies the signed transform by
$e^{-b_0s}$ and preserves its zeros, the equalities of moments,
and the equal coordinate marginals that have been established.

For the remaining symmetric Cauchy case choose $b,c>0$.
The mixing law of $(A,B)=(cZ,b)$, with $Z$ standard symmetric
Cauchy, and the point mass at $(0,b+c)$ give the same law of
$S_n$: both are centered Cauchy with scale $n(b+c)$.
The laws of the exchangeable sequences differ, because the
within-sequence difference of the first two coordinates has scales
$2b$ and $2(b+c)$, respectively. This completes the proof.
\end{proof}

\begin{remark}[Full sampling and examples]
\label{p01:mix:rem:stable-examples}
For $\alpha\ne1$, observing every positive integer satisfies the
criterion because $0<|1-\alpha|\le1$. For $\alpha=1$, $\beta\ne0$,
the logarithmic criterion also diverges on all positive integers.
At index $\alpha=3/2$, observations at $n=k^2$ suffice, because
$\sum_k(k^2)^{-1/2}$ diverges, whereas observations at $n=k^3$
do not. For a skew index-one family, a sampling sequence with
$\log n_k\asymp k^a$ identifies all mixing laws exactly when
$a\le1/2$. These are consequences of the new necessity proof,
not additional externally posed problems.
\end{remark}

\subsection{A separate observation: two joint replicates always suffice}

\begin{proposition}[Exact replication threshold for fixed stable shapes]
\label{p01:mix:prop:stable-two-replicates}
For every fixed $\alpha,\beta$ above, two conditionally independent
observations sharing $(A,B)$ identify its arbitrary Borel mixing law.
One observation does not. Thus the exact shared-latent replication
threshold is two, including the symmetric Cauchy family.
\end{proposition}
\begin{proof}
Equality of two-observation laws gives
\[
 H_u\bigl(\psi(t)+\psi(u-t)\bigr)=0\qquad(t,u\in\mathbb R)
\]
for the signed mixing transform $H_u$. Fix $u\ne0$ and let
$t>\max\{u,0\}+1$. The real part of the argument is
$t^\alpha+(t-u)^\alpha$, a strictly increasing positive function
of $t$. Thus the zeros include distinct points accumulating at an
interior point of the right half-plane. The identity theorem gives
$H_u=0$. Continuity includes $u=0$, and Fourier--Laplace uniqueness
finishes the argument.
For one observation, fix $b,c>0$, take $A$ with law $K_{0,c}$ and
$B=b$, or instead take $(A,B)=(0,b+c)$. The convolution-semigroup
identity makes the observed law $K_{0,b+c}$ in both cases, although
the mixing laws are distinct. This proves both assertions.
\end{proof}

This proposition is an elementary transform consequence and its
historical novelty has not been established. It also explains why
the Cauchy counterexample for separate partial-sum marginals does not
contradict identification from a joint pair of observations.

\section{A new constraint mechanism: prescribed density zeros repair finite defects}
\label{p01:mix:sec:quadrature-constraints}

This section gives a concrete special-case answer to the second
question in Section~6 of Le et al. The restriction is on the unknown
mixing density, not an extra observable moment. Its statistical
effect depends on whether a common smoothness bound is imposed.

Let $S\subset\mathbb N$ be finite and nonempty, put $M=\max S$,
and let
\[
 k_S(x\mid\theta)
 =1+\sum_{\substack{n\ge1\\n\notin S}}
       2^{-n^2-2}L_n(x)L_n(\theta),\qquad x,\theta\in[0,1].
\]
This is a positive jointly entire kernel. Write $A$ for its
one-observation mixing operator and
$\mathcal H_S=\operatorname{span}\{L_s:s\in S\}$.
Coefficient extraction and polynomial density, exactly as in
Proposition~\ref{p01:mix:prop:pde-audit}, give
\begin{equation}\label{p01:mix:eq:polynomial-fiber}
 A\mu=A\nu\quad\Longleftrightarrow\quad
 \nu-\mu=h(\theta)d\theta\quad\hbox{for some }h\in\mathcal H_S.
\end{equation}
Every $h\in\mathcal H_S$ has degree at most $M$ and integral zero.

Put $J=\lceil(M+1)/2\rceil$ and let $z_1,\ldots,z_J$ be the roots
of $L_J$ in $(0,1)$. These nodes admit positive quadrature weights
$w_i$ such that
\begin{equation}\label{p01:mix:eq:positive-quadrature}
 \int_0^1p(t)dt=\sum_{i=1}^Jw_ip(z_i)
       \qquad(\deg p\le2J-1),\qquad w_i>0.
\end{equation}
For completeness, orthogonality proves that $L_J$ has $J$ distinct
interior roots: multiplying by the polynomial formed from its
sign-changing roots would otherwise contradict orthogonality.
For any $p$ of degree at most $2J-1$, division by $L_J$ gives
$p=qL_J+r$ with $\deg q,\deg r<J$. Integrate and interpolate $r$
at the nodes to obtain the formula with
$w_i=\int\ell_i$, where $\ell_i$ is the $i$th Lagrange polynomial.
Applying the formula to $\ell_i^2$ gives
$w_i=\int\ell_i^2>0$.

\begin{theorem}[Quadrature zeros identify a density among all competitors]
\label{p01:mix:thm:quadrature-zero-identification}
If $g\in C^1([0,1])$ is a nonnegative probability density satisfying
$g(z_i)=0$ for $1\le i\le J$, then $\mu=g(t)dt$ is identifiable
among all Borel mixing probabilities from its one-observation law
under $k_S$. In particular the polynomial density
\[
 g_*(t)=(2J+1)L_J(t)^2
\]
is an individually identifiable absolutely continuous mixing law
inside this globally nonidentifiable model.
\end{theorem}
\begin{proof}
If $A\nu=A\mu$, \eqref{p01:mix:eq:polynomial-fiber} says
$\nu=(g+h)dt$ for $h\in\mathcal H_S$. Since $g+h$ is continuous,
nonnegativity almost everywhere implies nonnegativity everywhere.
Thus $h(z_i)\ge0$. But \eqref{p01:mix:eq:positive-quadrature} gives
\[
 0=\int_0^1h=\sum_{i=1}^Jw_ih(z_i),
\]
so all these values are zero. Both $g$ and $g+h$ have differentiable
local minima at each interior node, hence $g'(z_i)=0$ and
$h'(z_i)=0$. The polynomial $h$ has at least $2J$ zeros counted
with multiplicity and degree at most $M<2J$, so $h=0$.
The example is normalized by Legendre orthogonality.
\end{proof}

\begin{theorem}[The same density zeros do not suffice for uniform learning]
\label{p01:mix:thm:zero-constraint-learning}
Let $\mathcal D$ be the class of entire nonnegative probability
densities on $[0,1]$ vanishing at all $z_i$. Every member is
identified among all Borel competitors by one observation, but
there is a constant $c_S>0$ such that
\[
 \inf_{\widehat\mu}\sup_{g\in\mathcal D}
      \mathbb E_{(A(gdt))^{\otimes N}}
          W_1(\widehat\mu,gdt)\ge c_S\qquad(N\ge1).
\]
In contrast, for every fixed $L<\infty$ for which the following class
is nonempty, uniform $W_1$ consistency holds on
\[
 \mathcal D_L=\left\{g\in C^2([0,1]):g\ge0,\ \int_0^1g=1,
           \ g(z_i)=0\ (1\le i\le J),\ \|g\|_{C^2}\le L\right\}.
\]
The last assertion requires no analyticity of $g$.
\end{theorem}
\begin{proof}
For $a\in[0,1]\setminus\{z_i\}$, normalize
\[
 g_{a,t}(x)=C_{a,t}L_J(x)^2e^{-(x-a)^2/t},\qquad t>0.
\]
It is an entire probability density with the prescribed zeros and
converges weakly to $\delta_a$ as $t\downarrow0$. The same holds
at endpoints, using one-sided concentration. Nodes themselves are
approximated by points outside the finite node set. Finite convex
combinations stay in $\mathcal D$. Since finite atomic probabilities
are weakly dense in all probabilities on the interval, the weak
closure of $\{gdt:g\in\mathcal D\}$ is all of
$\mathcal P([0,1])$.

Choose $s\in S$ and the distinct mixing laws
$\mu_\pm=(1\pm\tfrac12L_s)dt$, which have equal observed laws.
Approximate both weakly by members $\mu_{\pm,j}$ of the class.
The kernel is continuous on a compact rectangle, so weak convergence
of mixing laws gives uniform convergence of their observed densities
and hence total-variation convergence. For each fixed $N$,
\[
 \operatorname{TV}\big((A\mu_{+,j})^{\otimes N},
                       (A\mu_{-,j})^{\otimes N}\big)\longrightarrow0.
\]
The two-point bound follows by integrating the triangle inequality
for $W_1$ against the overlap of the two data laws. Letting
$j\to\infty$ yields the constant
$c_S=\tfrac12W_1(\mu_+,\mu_-)>0$. For $S=\{1\}$ it is $1/12$.

For $\mathcal D_L$, the uniform second-derivative bound makes its
density functions precompact in $C^1$ by Arzel\`a--Ascoli.
Every weak limit consequently has a $C^1$ nonnegative density with
the same prescribed zeros and integral one. The previous theorem
therefore makes $A$ injective on the weak closure $K$ of this class.
That closure is compact; continuity of $A$ makes its inverse on
$A(K)$ uniformly continuous for weak metrics.
To see uniform estimation directly, take finite $1/N$-nets of $K$
in $W_1$, estimate the observed law by its empirical measure, and
choose a net element minimizing bounded-Lipschitz distance to that
empirical law. The bounded-Lipschitz unit ball on $[0,1]$ has
finite uniform nets, so the expected empirical distance tends to
zero uniformly over all observation laws. Uniform continuity of
$A$ and of its inverse, together with the bounded diameter of the
parameter interval, gives the required uniform expected $W_1$
convergence. Fixed tie-breaking makes the finite-net estimator
measurable.
\end{proof}

The quadrature-zero argument and its consistency distinction are new
extensions of the integrated draft. They do not solve the broad
constraint question for arbitrary standard families, and no historical
novelty claim is made. A fixed smoothness bound alone is insufficient
in a nonidentifiable family; the prescribed zeros remove the finite
polynomial ambiguity, while the common bound prevents that ambiguity
from reappearing in the closure.


% ===== Consolidated body: information =====
\section{Optimal replication budgets through multilevel distribution recovery}
\label{p01:info:sec:budget}

An equal number of repeated observations for every individual need not be
an efficient use of a total observation budget. We prove a sharp separation
for the factorial Legendre experiment: an unequal, predetermined allocation
attains a Wasserstein risk of order $B^{-1/3}$, and no adaptive allocation
can improve this order. The previously established equal-group lower
bound is of order $(B\log(eB))^{-1/4}$. Thus the change of allocation
improves a power of the budget, not merely a logarithm.

The upper bound uses the classical multilevel Monte Carlo telescoping
principle. Distribution-function versions of that principle were developed
by Giles, Nagapetyan and Ritter~\cite{p01:info:GNR2015}; nested conditional
expectations are treated in~\cite{p01:info:Giles2018}. The proof below uses
integrated distribution-function error, needs no density or smoothness of
the mixing law, and is paired with an adaptive statistical lower bound.
Lee and Valiant~\cite{p01:info:LeeValiant2021} already use unequal coin-sampling
depths and telescoping corrections in a population problem; their target
is the fraction of coins above a known gap, rather than the complete
mixing law in $W_1$. See also Lee's thesis~\cite{p01:info:Lee2021},
Section~4.3, pp.~65--67. The generic multilevel technique and the idea
of combining unequal replication depths are not claimed as new.

\subsection{The experiment and a general allocation theorem}

Let $(P_\theta)_{\theta\in[0,1]}$ be a known family of probability laws.
Starting a new group draws a fresh latent parameter independently from an
unknown $G\in\mathcal P([0,1])$. Every observation within that group has
law $P_\theta$, independently conditional on its latent parameter. A
policy may start new groups, revisit previous groups, randomize, and
adapt every choice to the entire observed history. Each observation costs
one unit. The total number of observations is at most the deterministic
integer budget $B$. The latent parameters themselves are never observed.
Write $\mathfrak R_B^{\rm ad}$ for the minimax expected $W_1$ risk over
all such policies and estimators, and $\mathfrak R_B^{\rm fix}$ when all
group sizes are chosen in advance, but need not be equal.

\begin{theorem}[Sharp adaptive budget under regular component estimation]
\label{p01:info:thm:general-budget}
Suppose there are an integer $b_0\geq1$ and estimators $T_r$ valued in
$[0,1]$, based on $b_0r$ observations from a single group, such that
\begin{equation}\label{p01:info:eq:component-risk}
 \sup_{\theta\in[0,1]}\mathbb E_\theta|T_r-\theta|
 \leq A r^{-1/2},\qquad r\geq1.
\end{equation}
The estimators can be randomized, with their randomization independent of
the unknown parameter. Then
\begin{equation}\label{p01:info:eq:general-upper}
 \mathfrak R_B^{\rm ad}\leq\mathfrak R_B^{\rm fix}
 \leq C B^{-1/3},
\end{equation}
where $C$ depends only on $A,b_0$.

Suppose additionally that on some nondegenerate closed interval
$J\subset(0,1)$ there is a finite $D$ with
\begin{equation}\label{p01:info:eq:local-kl}
 \operatorname{KL}(P_\theta\Vert P_{\theta'})
 \leq D|\theta-\theta'|^2\qquad(\theta,\theta'\in J).
\end{equation}
Then constants $c,C>0$ satisfy
\begin{equation}\label{p01:info:eq:adaptive-rate}
 c B^{-1/3}\leq\mathfrak R_B^{\rm ad}
 \leq\mathfrak R_B^{\rm fix}\leq C B^{-1/3}
\end{equation}
for all sufficiently large $B$. Thus adaptive group selection cannot
improve the minimax order attained by a fixed unequal allocation.
\end{theorem}

\begin{proof}[Proof of the upper bound]
Put $r_\ell=2^\ell$. For a fresh group at level $\ell\geq1$, compute
$T_{r_\ell}$ from its $b_0r_\ell$ observations and compute
$T_{r_{\ell-1}}$ from the first $b_0r_{\ell-1}$ observations. Independent
auxiliary randomization may be used for the two statistics if needed.
They share the same latent parameter, and both have the required
marginal laws. Consequently
\begin{equation}\label{p01:info:eq:level-coupling}
 \mathbb E_G|T_{r_\ell}-T_{r_{\ell-1}}|
 \leq A(1+\sqrt2)r_\ell^{-1/2}.
\end{equation}
At level zero use $T_1$. Use $N_\ell$ independent groups at each level,
and use disjoint groups across levels. Define the random signed
distribution function
\begin{align}
 H_L(t)={}&\frac1{N_0}\sum_{i=1}^{N_0}
       \mathbf1_{\{T_{1,i}\leq t\}}\notag\\
 &+\sum_{\ell=1}^L\frac1{N_\ell}\sum_{i=1}^{N_\ell}
   \left(\mathbf1_{\{T_{r_\ell,i}\leq t\}}
       -\mathbf1_{\{T_{r_{\ell-1},i}\leq t\}}\right).
 \label{p01:info:eq:ml-cdf}
\end{align}
Within a parenthesis the two estimates come from the same group.
Let $F_{G,L}$ be the distribution function of $T_{r_L}$ when its latent
parameter has law $G$. Telescoping gives $\mathbb EH_L=F_{G,L}$.
No independence between the two estimates in a parenthesis is needed.

For any $u,v\in[0,1]$,
\[
 \int_0^1|\mathbf1_{\{u\leq t\}}-\mathbf1_{\{v\leq t\}}|^2\,dt
 =|u-v|.
\]
The independence of the groups and levels, followed by
\eqref{p01:info:eq:level-coupling}, therefore yields
\[
 \int_0^1\operatorname{Var}(H_L(t))\,dt
 \leq C_A\sum_{\ell=0}^L\frac{r_\ell^{-1/2}}{N_\ell},
 \qquad C_A=\max\{1,A(1+\sqrt2)\}.
\]
Cauchy--Schwarz first in probability and then in $t$ gives
\begin{equation}\label{p01:info:eq:ml-stochastic}
 \mathbb E\|H_L-F_{G,L}\|_{L^1[0,1]}
 \leq\left(C_A\sum_{\ell=0}^L
                   \frac{r_\ell^{-1/2}}{N_\ell}\right)^{1/2}.
\end{equation}
Coupling $T_{r_L}$ with its own latent parameter and using the
one-dimensional distribution-function formula for $W_1$ also gives
\begin{equation}\label{p01:info:eq:ml-bias}
 \|F_{G,L}-F_G\|_1\leq A r_L^{-1/2}.
\end{equation}

Here is a deterministic allocation with explicit integer rounding. Put
$b=\lfloor B/b_0\rfloor$ and, for $b\geq64$, choose $L\geq0$ so that
\[
 R:=2^L\leq b^{2/3}/16<2^{L+1},\qquad
 S_L=\sum_{\ell=0}^L r_\ell^{1/4},\qquad
 N_\ell=\left\lceil\frac{b}{2S_L}r_\ell^{-3/4}\right\rceil.
\]
The allocation costs at most $B$, since
\[
 \sum_{\ell=0}^L r_\ell N_\ell
 \leq b/2+\sum_{\ell=0}^Lr_\ell
 <b/2+2R\leq b.
\]
Furthermore,
\[
 \sum_{\ell=0}^L\frac{r_\ell^{-1/2}}{N_\ell}
 \leq\frac{2S_L^2}{b},\qquad
 S_L\leq\frac{R^{1/4}}{1-2^{-1/4}},\qquad
 b^{2/3}/32<R\leq b^{2/3}/16.
\]
Combining these bounds with \eqref{p01:info:eq:ml-stochastic} and
\eqref{p01:info:eq:ml-bias} shows
\begin{equation}\label{p01:info:eq:raw-cdf-rate}
 \sup_G\mathbb E\|H_L-F_G\|_1\leq C b^{-1/3}.
\end{equation}

Finally convert $H_L$, which need not be monotone, to a probability law.
Enumerate a countable $W_1$-dense family $(Q_j)$ of rational atomic
probabilities on $[0,1]$. Select the first index satisfying
\[
 \|F_{Q_j}-H_L\|_1
 \leq\inf_k\|F_{Q_k}-H_L\|_1+B^{-1},
\]
and output $\widehat G=Q_j$. The objective functions and their countable
infimum are measurable, so this defines a measurable estimator.
Density of the family implies that the infimum is at most
$\|F_G-H_L\|_1$. Hence
\[
 W_1(\widehat G,G)\leq2\|H_L-F_G\|_1+B^{-1}.
\]
Equation~\eqref{p01:info:eq:raw-cdf-rate} proves the claim. For the finitely
many smaller budgets, increase the constant and use any fixed law.
\end{proof}

\begin{proof}[Proof of the adaptive lower bound]
Write $J=[a,a+s]$ with $s>0$. For an integer $K\geq1$, set
\[
 a_j=a+\frac{s(j-1/2)}K,\qquad h=\frac{s}{4K},\qquad
 G_v=\frac1K\sum_{j=1}^K\delta_{a_j+hv_j},
 \qquad v\in\{-1,1\}^K.
\]
The intervals containing different atoms are disjoint and ordered.
Matching sorted atoms therefore gives the exact formula
\begin{equation}\label{p01:info:eq:assouad-transport}
 W_1(G_v,G_w)=\frac{2h}{K}\,d_H(v,w).
\end{equation}

Consider a more informative oracle: it reveals the locations $a_j$ and
the common weights $1/K$, and at each step allows the statistician to
choose any index $j$ and receive an independent observation with law
$P_{a_j+hv_j}$. The statistician still has at most $B$ observations.
Every original grouped policy can be simulated in this oracle experiment:
when the policy starts a group, draw an independent uniform index in
$\{1,\ldots,K\}$, and query that same index whenever the group is revisited.
The simulator may conceal the index from the original policy. Thus a
lower bound for this oracle experiment is also a lower bound for the
original one.

For any adaptive randomized oracle policy, let $\mathbb P_v$ denote its
full transcript law, including its random seed, and let $N_j$ be the
number of queries to index $j$. If $v^{(j)}$ flips only the $j$th sign,
the chain rule for relative entropy and \eqref{p01:info:eq:local-kl} imply
\begin{equation}\label{p01:info:eq:adaptive-kl}
 \operatorname{KL}(\mathbb P_v\Vert\mathbb P_{v^{(j)}})
 \leq4Dh^2\mathbb E_vN_j.
\end{equation}
Indeed the conditional decision rule is identical in the two models at
every fixed past transcript; only observations at index $j$ contribute.
Stopping before $B$ causes no difficulty, since subsequent dummy
observations can be identical under all models. Since $\sum_jN_j\leq B$,
averaging \eqref{p01:info:eq:adaptive-kl} over $v,j$ gives
\[
 \frac1{K2^K}\sum_v\sum_j
 \operatorname{KL}(\mathbb P_v\Vert\mathbb P_{v^{(j)}})
 \leq\frac{4Dh^2B}{K}.
\]
Pinsker's inequality and Jensen's inequality consequently imply
\begin{equation}\label{p01:info:eq:adaptive-tv}
 \frac1{K2^K}\sum_v\sum_j
 d_{\rm TV}(\mathbb P_v,\mathbb P_{v^{(j)}})
 \leq\left(\frac{2Dh^2B}{K}\right)^{1/2}.
\end{equation}
For any sign estimator $\widehat v$, the binary testing inequality for
each pair $v,v^{(j)}$, averaged over all such pairs, gives
\[
 \frac1{2^K}\sum_v\mathbb E_vd_H(\widehat v,v)
 \geq\frac K2\left(1-
 \frac1{K2^K}\sum_v\sum_j
          d_{\rm TV}(\mathbb P_v,\mathbb P_{v^{(j)}})\right).
\]
Choose $K=\max\{1,\lceil(2Ds^2B)^{1/3}\rceil\}$.
The right side of \eqref{p01:info:eq:adaptive-tv} is at most $1/4$,
so the average Hamming risk is at least $3K/8$.

Given any probability-valued estimator $\widehat G$, choose a nearest
$G_{\widehat v}$ among this finite family, using a fixed tie rule.
The triangle inequality and nearest-point property imply
\[
 W_1(\widehat G,G_v)
 \geq\frac12W_1(G_{\widehat v},G_v)
 =\frac hK d_H(\widehat v,v).
\]
Averaging shows that its worst-case risk is at least $3h/8$,
which is bounded below by $c B^{-1/3}$ for all sufficiently large $B$.
The argument includes randomized estimators and every adaptive policy.
\end{proof}

\subsection{Application to the factorial Legendre kernel}

Let $L_j$ denote the shifted Legendre polynomials and let
\[
 k_*(x\mid\theta)=1+\sum_{j\geq2}\frac{L_j(x)L_j(\theta)}{j!},
 \qquad x,\theta\in[0,1].
\]
This is the strictly positive entire kernel already considered in the
replication analysis. Its one-observation mixture map has a nontrivial
null direction, whereas pairs identify arbitrary mixing laws.

\begin{lemma}[Equal-group obstruction for the factorial kernel]
\label{p01:info:lem:equal-group}
Let $R_{N,q}$ be the minimax expected $W_1$ risk from $N$ independent
groups of $q$ observations each in the factorial experiment. Uniformly
over integers $N,q\geq1$,
\[
 R_{N,q}\geq c\max\left\{N^{-1/2},
 \min\left\{\frac1{\log(eN)},
              \frac1{\sqrt{q\log(eN)}}\right\}\right\}.
\]
Consequently $\inf_{Nq\leq B}R_{N,q}\geq
c\{B\log(eB)\}^{-1/4}$ for sufficiently large $B$.
\end{lemma}
\begin{proof}
An oracle revealing the $N$ latent parameters is more informative.
The laws $G_\pm=(1/2\pm a)\delta_1+(1/2\mp a)\delta_0$ with
$a=1/(16\sqrt N)$ have transport distance $2a$ and $N$-sample
relative entropy bounded by an absolute constant below $1/2$.
Pinsker's inequality and the two-point testing bound give
$R_{N,q}\geq cN^{-1/2}$.

For the other term put $\theta_0=1/2$,
$\mathcal H=L^2(P_{\theta_0};\mathbb C)$, and
$f_z(x)=k_*(x\mid\theta_0+z)/k_*(x\mid\theta_0)$.
The factorial series converges locally uniformly for complex parameters,
and the denominator is at least $3-e$. Thus for constants $R\in(0,1/4]$
and $C_0\geq1$, $f_z$ is $\mathcal H$-holomorphic and
\[
 \|f_z-1\|_{\mathcal H}\leq\sqrt{C_0}|z|,
 \qquad \int f_z\,dP_{\theta_0}=1 \quad (|z|\leq R).
\]
The last identity cancels the cross term in the squared Hilbert norm.
Writing $F_z=f_z^{\otimes q}=\sum_{j\geq0}T_jz^j$ in
$L^2(P_{\theta_0}^{\otimes q};\mathbb C)$ gives, for $0<h\leq R/2$,
\[
 \|F_z\|\leq e^{C_0q|z|^2/2},\qquad
 \|T_j\|\leq e^{2C_0qh^2}(2h)^{-j},
\]
by Cauchy's coefficient estimate on $|z|=2h$.

For an integer $m\geq1$, set $n=m+1$ and
$\theta_j=\theta_0-h+2hj/n$. Define
\[
 G=2^{1-n}\sum_{j\ {
m even}}\binom nj\delta_{\theta_j},
 \qquad
 H=2^{1-n}\sum_{j\ {
m odd}}\binom nj\delta_{\theta_j},
 \qquad 0\leq j\leq n.
\]
Finite differences show that their moments agree through degree $m$.
On the $j$th grid interval the difference of their cumulative
distribution functions is $2^{1-n}(-1)^j\binom{n-1}{j}$.
Its absolute values sum to one, so $W_1(G,H)=2h/(m+1)$.
Writing $Q_q(G)=\int P_\theta^{\otimes q}\,dG(\theta)$ and using
$|\int(\theta-\theta_0)^j\,d(G-H)|\leq2h^j$ yields
\[
 \left\|\frac{dQ_q(G)-dQ_q(H)}{dP_{\theta_0}^{\otimes q}}\right\|_2
 \leq2e^{2C_0qh^2}\sum_{j>m}2^{-j}
 =2e^{2C_0qh^2}2^{-m}.
\]
Hence the one-group total variation distance is at most
$e^{2C_0qh^2}2^{-m}$. Take
\[
 m=\left\lceil\frac4{3\log2}\log(4N)\right\rceil,
 \qquad
 h=\min\left\{R/2,\sqrt{\frac{m\log2}{8C_0q}}\right\}.
\]
The one-group distance is at most $1/(4N)$, so the $N$-group distance
is at most $1/4$. The two-point bound implies
\[
 R_{N,q}\geq c\frac h{m+1}
 \geq c'\min\{\log(eN)^{-1},(q\log(eN))^{-1/2}\}.
\]
Finally set $L=\log(eN)$. If $q\geq L$, the maximum of the two
lower bounds is at least
$(NqL)^{-1/4}\geq\{B\log(eB)\}^{-1/4}$, up to a constant.
If $q<L$, the bound $c/L\geq c/\log(eB)$ is larger than the asserted
bound for all sufficiently large $B$.
\end{proof}

\begin{corollary}[Unequal replication is strictly better than equal groups]
\label{p01:info:cor:factorial-budget}
For the factorial Legendre experiment,
\[
 \mathfrak R_B^{\rm ad}\asymp\mathfrak R_B^{\rm fix}\asymp B^{-1/3}.
\]
If $\mathfrak R_B^{\rm eq}$ denotes the corresponding infimum restricted
to equal groups, then
\[
 \mathfrak R_B^{\rm eq}\geq c(B\log(eB))^{-1/4},\qquad
 \frac{\mathfrak R_B^{\rm fix}}{\mathfrak R_B^{\rm eq}}
 \leq C B^{-1/12}(\log(eB))^{1/4}\longrightarrow0.
\]
Thus no equal-group allocation is order optimal under the unrestricted
total-budget criterion. Adaptive selection is unnecessary for the optimal
order, but unequal group sizes are essential for attaining it.
\end{corollary}

\begin{proof}
Put $A_j(x)=j!(2j+1)L_j(x)$. The bounded pair statistic
\[
 Y(x,y)=\frac12+\frac{35}{18}
 \left(A_2(x)A_3(y)-\frac{10}{21}A_5(x)-\frac4{15}A_3(x)\right)
\]
satisfies $\mathbb E_\theta Y(X_1,X_2)=\theta$ and $|Y|<2100$.
For completeness, $\mathbb E_\theta A_j(X)=L_j(\theta)$ and
\[
 L_2L_3=\frac{10}{21}L_5+\frac4{15}L_3+\frac9{35}L_1,
 \qquad L_1(\theta)=2\theta-1,
\]
prove unbiasedness. The bounds $|L_j|\leq1$ give
$|A_2A_3|\leq420$, $|A_5|\leq1320$, and $|A_3|\leq42$,
which prove the stated numerical bound. Average $Y$ over $r$ disjoint
pairs and clip to $[0,1]$. The resulting $T_r$ satisfies
\eqref{p01:info:eq:component-risk} with $b_0=2$ and $A=2100$.

The kernel is bounded below by $c_0=3-e>0$. Uniform convergence of its
derivative on the compact square gives
$M=\sup_{x,\theta}|\partial_\theta k_*(x\mid\theta)|<\infty$.
Therefore
\[
 \operatorname{KL}(P_\theta\Vert P_{\theta'})
 \leq\chi^2(P_\theta\Vert P_{\theta'})
 \leq\frac{M^2}{c_0}|\theta-\theta'|^2.
\]
Theorem~\ref{p01:info:thm:general-budget} applies. Taking the ratio of its
upper bound with Lemma~\ref{p01:info:lem:equal-group} proves the strict
separation.
\end{proof}

\begin{remark}[What this closes]
The complete total-budget minimax order is now determined for arbitrary
unequal and adaptive allocations in this kernel model. This closes the
allocation problem left open by the equal-group analysis. It does not
determine the sharp growing-replication curve $R_{N,q}$ for each fixed pair
$(N,q)$, and it does not close the logarithmic gap within the smaller
equal-group class itself. The result is a proved resolution of that
internally posed extension, not a claim to have resolved a named published
conjecture or established historical priority.
\end{remark}

\subsection{An exact Bernoulli reduction and the remaining equal-group gap}

\begin{proposition}[Bounded unbiased statistics transfer population bounds]
\label{p01:info:prop:bernoulli-reduction}
Suppose a statistic $Y$ based on $b_0$ observations has
$|Y|\leq M$ and $\mathbb E_\theta Y=\theta$ for $\theta\in[0,1]$.
With an independent uniform $U$, define
\[
 Z=\mathbf1_{\{U\leq(Y+M)/(2M)\}}.
\]
Then conditionally on $\theta$, $Z$ is Bernoulli with parameter
$p(\theta)=(\theta+M)/(2M)$. Disjoint blocks produce conditionally
independent Bernoulli variables. Consequently every distribution-learning
upper bound for $N$ groups of $t$ Bernoulli observations transfers to
$N$ groups of $b_0t$ observations from the original experiment, with the
$W_1$ risk multiplied by at most $2M$.
\end{proposition}
\begin{proof}
Conditional expectation gives the Bernoulli parameter, and independent
blocks and auxiliary uniforms give conditional independence. If a
Bernoulli-population estimator produces a law $\widehat H$ on $[0,1]$,
push it forward by the clipped inverse map
$u\mapsto\min\{1,\max\{0,2Mu-M\}\}$. This map is $2M$-Lipschitz
and is the true inverse on $p([0,1])$, so contraction of transport under
Lipschitz maps gives the stated risk bound. Randomization causes no
statistical difficulty; it is part of the estimator, not additional data.
\end{proof}

Vinayak, Kong, Valiant and Kakade~\cite{p01:info:VKVK2019} prove a sharp
binomial-mixture rate $(t\log N)^{-1/2}$ in the range
$t\in[\Omega(\log N),O(N^{2/9-\varepsilon})]$; see their
Section~1.1 and Theorem~3.2. Their Remark~3.2 explicitly conjectures
extending the MLE bound to $t=O(N^{2/3-\varepsilon})$. That theorem
cannot be extrapolated to
$t\asymp N/\log N$, which is the range relevant to the optimal
equal-group total budget. The exact reduction above is valid, but an
appropriate all-regime upper bound or a direct proof would still be
needed to close that logarithmic gap. Musco, Musco, Rosenblatt and
Singh~\cite{p01:info:MMRS2025}, Theorem~20, improve the proved MLE range
to $t=O(N^{1/4-\varepsilon})$. Their discussion of conjectured
improvements (published Appendix~C.2; Section~6.2 of the cited arXiv
version) still requires sharper coefficient and likelihood bounds.
Neither proved range reaches the budget endpoint. No closure of that
gap is claimed here.

\subsection{Published questions checked against the new result}

Lau and Scarlett's 2026 COLT open problem~\cite{p01:info:LauScarlett2026}
asks whether arbitrary nonadaptive one-bit quantizers can attain the
optimal adaptive mean-estimation rate over nonparametric finite-moment
classes, or whether one adaptive transition is necessary. This is an
explicit published open question; the present replication theorem does
not resolve it. The observation is different: their quantizer receives
one fresh unbounded sample, while the present experiment can repeatedly
observe the same latent individual. The bounded Bernoulli reduction also
does not remove that distinction or supply a bounded unbiased statistic
for their unrestricted class.

The second explicit published question is Vinayak et al.'s
Remark~3.2 conjecture just described. It concerns the particular MLE
under equal group sizes, and neither our reduction nor our different
allocation proves it. Their Remark~3.3 and supplementary Section~A.6
also outline local moment matching for a broader range, but the supplement
calls this a high-level construction and does not state a proved
all-regime bound sufficient for the present equal-budget endpoint.
The 2025 paper just cited provides a verified later improvement and
still discusses the stronger assertions as conjectural. No subsequent
resolution was verified in the targeted search. In particular, the
coefficient assertion $|C(t,m,j)|\leq\exp(m^2/t)$, where $C(t,m,j)$
are the degree-$t$ Bernstein coefficients of the shifted degree-$m$
Chebyshev polynomial, is a precise related approximation question;
we have not proved it and do not count it as an additional resolution.
Finally,
Lee and Kennedy's October 2026 preprint~\cite{p01:info:LeeKennedy2026} already
treats heterogeneous trial profiles for pointwise estimation of smooth
binomial mixing densities. Its loss and smoothness class differ from the
arbitrary-law $W_1$ budget problem proved here; trial heterogeneity by
itself cannot be described as an untouched research direction.


% ===== Consolidated body: moments =====
\section{Exact Mellin zeros, moment comparisons, and the tail-order conjecture}
\label{p01:mom:sec:main}

\paragraph{Problem and priority ledger.}
The explicitly published unresolved assertion addressed in this section is
Conjecture~17 of B\"urgin, Epperlein, and Wirth~\cite{p01:mom:BEW2022}:
on a prescribed positive compact interval, reciprocal-summable integer
exponents should admit a nonzero continuous signed density annihilating
those moments. The conjecture is located on printed page~15 of the
cited arXiv version. Its assertion is correct, but it already follows by
a logarithmic change of variables from Luxemburg--Korevaar's 1971
construction~\cite[Theorem~5.2]{p01:mom:LK1971}, restated with its product
formula in Wellershoff~\cite[Theorem~21, printed p.~18]{p01:mom:Wellershoff2021}.
Accordingly, this is an explicit deduction and an alternative proof,
not a claim of a historically new resolution. The exact real zero and
sign classifications below are further deductions from that classical
construction. Their precise formulations were not located in the targeted
search, which does not establish priority. They are not counted as three
additional published open problems.

For a finite signed measure $\nu$ supported on a compact subinterval of
$(0,\infty)$, write
\[
 \mathcal M_\nu(z)=\int x^z\,d\nu(x),\qquad z\in\mathbb C,
\]
where $x^z=e^{z\log x}$. This is an entire function. For an integrable
function $h$, write $\mathcal M_h=\mathcal M_{h(x)dx}$.
A function is in $G^\sigma(\mathbb R)$ if it is smooth and satisfies
$\|h^{(k)}\|_\infty\le CB^k(k!)^\sigma$ for some $C,B>0$ and every
integer $k\ge0$.

\subsection{An alternative proof of the published conjecture}
\label{p01:mom:sec:smoothing}

\begin{proposition}[Localized smoothing of real M\"untz annihilators]
\label{p01:mom:prop:smoothing}
Let $0<a<b$, let $I$ be a nonempty open interval with
$\overline I\subset(a,b)$, and let
$0=\lambda_0<\lambda_1<\lambda_2<\cdots$ tend to infinity.
For every $\sigma>1$ the following are equivalent:
\begin{enumerate}
\item $\sum_{j\ge1}\lambda_j^{-1}<\infty$;
\item there exists a nonzero real $h\in G^\sigma(\mathbb R)$ with
compact support in $I$ and $\mathcal M_h(\lambda_j)=0$ for every $j$.
\end{enumerate}
The same assertion holds for a finite sequence, with the sum interpreted
as finite. The annihilator space in the second assertion is infinite
dimensional.
\end{proposition}

\begin{proof}
Choose a nondegenerate interval $[c,d]\subset I$. The M\"untz theorem
on positive intervals~\cite[Theorem~3.2]{p01:mom:BE1995} gives a nonzero
real signed measure $\nu$ supported on $[c,d]$ that annihilates all the
specified powers when the reciprocal sum converges. Choose a nonnegative,
nonzero $\varphi\in G^\sigma(\mathbb R)$ supported in a sufficiently
small neighborhood of $1$, and set
\begin{equation}\label{p01:mom:eq:smoothing}
 h(x)=\int t^{-1}\varphi(x/t)\,d\nu(t).
\end{equation}
The neighborhood can be chosen so that $\operatorname{supp}h\Subset I$.
Moreover
\[
 \|h^{(k)}\|_\infty\le\|\nu\|_{\rm TV}c^{-k-1}
                         \|\varphi^{(k)}\|_\infty,
\]
which proves the Gevrey bound. Fubini's theorem gives the entire identity
\begin{equation}\label{p01:mom:eq:factorization}
 \mathcal M_h(z)=\mathcal M_\nu(z)\mathcal M_\varphi(z).
\end{equation}
For every real $s$, the second factor is strictly positive. A nonzero
compactly supported signed measure has a nonzero integer moment by
polynomial density, so $h\ne0$. For necessity, the divergent M\"untz
condition makes the specified powers dense in $C[a,b]$. An annihilator
must therefore vanish. Disjoint subintervals of $I$ give infinitely many
linearly independent annihilators.

For completeness, the required Gevrey bump can be obtained by multiplying
two translates of $e^{-t^{-1/(\sigma-1)}}\mathbf1_{t>0}$. Cauchy's
estimate on disks of radius proportional to $t>0$, followed by maximizing
$t^{-k}e^{-ct^{-1/(\sigma-1)}}$, gives the stated Gevrey order; extension
by zero is smooth.
\end{proof}

Taking integer $\lambda_j=n_j$ proves the precise assertion of the
published conjecture, with arbitrary interior localization and every
Gevrey order above one. There is no obstruction arising from total mass:
adding the exponent zero costs only a finite extra constraint.

\subsection{The complete positive-real zero set and its multiplicities}
\label{p01:mom:sec:real-zeros}

We use the following classical external input. In the notation adapted
to the present proof, Luxemburg--Korevaar's theorem states that if
$d_j\ne0$ and $\sum_j|d_j|^{-1}<\infty$, then for any $u<v$ and any
integer $m\ge0$ there is a nonzero $g\in C_c^\infty([u,v])$ with
\begin{equation}\label{p01:mom:eq:LK-product}
 \widehat g(z)=C_* z^m e^{-i\tau z}
       \prod_j(1-z^2/d_j^2)\prod_k\cos(\epsilon_k z),
 \qquad
 \widehat g(z)=\int g(t)e^{-itz}\,dt,
\end{equation}
where $C_*\ne0$, $\epsilon_k>0$, $\sum_k\epsilon_k<\infty$, and
$\tau$ is real. Repetitions of $d_j$ encode multiplicity. The theorem and
factorization used here were checked in the cited 2021 author-hosted
restatement; the original 1971 publisher PDF was not retrievable in this
verification pass. The existence theorem itself is an imported classical
result, not proved anew here.

\begin{theorem}[Exact real Mellin zeros with multiplicity]
\label{p01:mom:thm:real-zeros}
Let $\Lambda\subset(0,\infty)$ have no finite accumulation point, let
$m_\lambda\in\mathbb N$ for $\lambda\in\Lambda$, and fix an integer
$m_0\ge1$. The following conditions are equivalent:
\begin{enumerate}
\item $\displaystyle\sum_{\lambda\in\Lambda}m_\lambda/\lambda<\infty$;
\item for every $0<a<b$, every nonempty open $I$ with
$\overline I\subset(a,b)$, and every $\sigma>1$, there is a nonzero
real $h\in G^\sigma(\mathbb R)$ compactly supported in $I$ such that
$\mathcal M_h$ has a zero of order $m_0$ at zero, and its positive-real
zeros are exactly $\Lambda$, with multiplicities $m_\lambda$.
\end{enumerate}
The construction may additionally arrange that its negative-real zeros
are exactly $-\Lambda$, with the same multiplicities.
\end{theorem}

\begin{proof}
For necessity, if $\operatorname{supp}h\subset[a,b]$, then
\[
 F(z)=b^{-z}\mathcal M_h(z)=\int_a^b(x/b)^z h(x)\,dx
\]
is bounded and holomorphic on $\operatorname{Re}z>0$ and is not
identically zero. The disk map $w=(z-1)/(z+1)$ and the Blaschke zero
condition imply
\[
 \sum_{\lambda\in\Lambda}
       m_\lambda\frac{\lambda}{1+\lambda^2}<\infty.
\]
Since the zero set has no finite accumulation and is bounded away from
zero, this is equivalent to the required reciprocal sum. One can obtain
the disk condition directly from Jensen's formula after removing a zero
at the center: boundedness bounds the sums
$\sum_{|w_j|<r}\log(r/|w_j|)$ uniformly as $r\uparrow1$, and hence
$\sum_j(1-|w_j|)<\infty$.

For sufficiency, first suppose $\Lambda$ is infinite. Choose
$[c,d]\Subset I$, put $u=\log c$, $v=\log d$, and apply
\eqref{p01:mom:eq:LK-product} to the sequence consisting of $i\lambda$
repeated $m_\lambda$ times, and to $m=m_0$. Multiplying $g$ by a
nonzero constant permits the normalization $C_*=i^{-m_0}$. All other
factors have the conjugation symmetry
$\widehat g(-\overline z)=\overline{\widehat g(z)}$, so uniqueness of
the Fourier transform shows that $g$ is real. Set
$h_0(x)=x^{-1}g(\log x)$ on $(0,\infty)$ and zero elsewhere. It is
real and smooth, with compact support in $[c,d]$, and
\begin{equation}\label{p01:mom:eq:positive-real-product}
 \mathcal M_{h_0}(s)
 =s^{m_0}e^{\tau s}
   \prod_{\lambda\in\Lambda}(1-s^2/\lambda^2)^{m_\lambda}
   \prod_k\cosh(\epsilon_k s),\qquad s\in\mathbb R.
\end{equation}
The last product is finite and strictly positive for every real $s$:
its logarithm converges because $\sum_k\epsilon_k^2<\infty$.
The middle product converges locally uniformly, and away from its
specified zeros it is nonzero. Indeed, all sufficiently late factors
are near one and $\sum_\lambda m_\lambda/\lambda^2<\infty$.
Thus the real zero set and all its multiplicities are exactly those
asserted.

If $\Lambda$ is finite, start instead with a nonzero nonnegative smooth
bump $g_0$ supported in $(u,v)$. The constant-coefficient differential
operator corresponding to the polynomial
$s^{m_0}\prod_{\lambda\in\Lambda}(1-s^2/\lambda^2)^{m_\lambda}$
produces $g$ whose bilateral Laplace transform is that polynomial
times the strictly positive real transform of $g_0$. Specifically,
replace $s$ by $-d/dt$ in this polynomial. This also treats
$\Lambda=\varnothing$.

Finally apply \eqref{p01:mom:eq:smoothing} to the signed measure
$h_0(x)\,dx$, choosing the bump support near $1$ sufficiently small.
The result belongs to $G^\sigma$, is still supported in $I$, and its
Mellin transform is multiplied by an entire function strictly positive
on the real axis. Hence no real zero is created or removed, and the
multiplicities are unchanged. This proves the assertion.
\end{proof}

\begin{corollary}[Exact real moment coincidences with shape constraints]
\label{p01:mom:cor:real-logconcave}
Under the summability condition in Theorem~\ref{p01:mom:thm:real-zeros},
there are distinct probability densities $g_+,g_-\in G^\sigma(\mathbb R)$
with support $[a,b]$, positive and strictly log-concave on $(a,b)$,
such that their real-order moments agree at a positive exponent $s$
if and only if $s\in\Lambda$. The multiplicity of that agreement as a
zero of the moment-difference function is $m_s$. Their difference is
supported in any prescribed interior interval.
\end{corollary}

\begin{proof}
Take the mass-zero $h$ of Theorem~\ref{p01:mom:thm:real-zeros}. With
$p=(\sigma-1)^{-1}$, normalize the probability density
\[
 g_0(x)=C\exp\bigl(-(x-a)^{-p}-(b-x)^{-p}\bigr)
 \mathbf1_{a<x<b}.
\]
It lies in $G^\sigma$ and has strictly negative logarithmic second
derivative on $(a,b)$. On the compact support of $h$, both positivity
and the inequality $g_0g_0''-(g_0')^2<0$ have a strict margin.
Consequently $g_\pm=g_0\pm\varepsilon h$ preserve them for sufficiently
small $\varepsilon>0$. Outside that compact set the densities equal
$g_0$. Their masses are one, and their moment difference is
$2\varepsilon\mathcal M_h$.
\end{proof}

\subsection{A complete classification of integer moment sign patterns}
\label{p01:mom:sec:sign-patterns}

This records more than equality or inequality in the eventual tail
order. It specifies every comparison of the two integer moment sequences.

\begin{theorem}[Exact signs of moment differences]
\label{p01:mom:thm:sign-patterns}
Let $(\eta_n)_{n\ge1}$ be a sequence in $\{-1,0,1\}$ that is not
identically zero. Let $Z=\{n\ge1:\eta_n=0\}$, and enumerate its
complement as $n_1<n_2<\cdots$ when the complement is infinite. The
following are equivalent:
\begin{enumerate}
\item There are two distinct probability measures supported in a
compact subinterval of $(0,\infty)$ such that
\[
 \operatorname{sgn}\left(\int x^n\,d\mu_+(x)
                         -\int x^n\,d\mu_-(x)\right)=\eta_n
 \quad(n\ge1).
\]
\item The set $Z$ satisfies $\sum_{n\in Z}n^{-1}<\infty$ (so its
complement is infinite), and
\begin{equation}\label{p01:mom:eq:sign-cost}
 \sum_{j:\,\eta_{n_j}\ne\eta_{n_{j+1}}}\frac1{n_{j+1}}<\infty.
\end{equation}
\item For every $0<a<b$ and $\sigma>1$, such a pair can be chosen
with strictly log-concave $G^\sigma$ probability densities supported
on $[a,b]$, positive on $(a,b)$, whose difference is supported in any
prescribed nonempty interior interval.
\end{enumerate}
\end{theorem}

\begin{proof}
Let $F(z)=\int x^z\,d(\mu_+-\mu_-)(x)$. It is nonzero, real on
real arguments, and its positive zeros obey the Blaschke summability
condition proved above. In particular $\sum_{n\in Z}1/n<\infty$.
If consecutive nonzero samples have opposite signs, the intermediate
value theorem gives a zero $r_j\in(n_j,n_{j+1})$. These intervals are
disjoint, and $1/n_{j+1}\le1/r_j$. Summing gives
\eqref{p01:mom:eq:sign-cost}. This proves necessity, including the fact that
a finite nonzero complement cannot occur.

Conversely, put a root of multiplicity two at every integer in $Z$.
For each sign transition put one simple root at
$r_j=n_{j+1}-1/2$. This point lies strictly between $n_j$ and
$n_{j+1}$, is not an integer, and satisfies
$1/r_j\le2/n_{j+1}$. The combined root multiset therefore satisfies
Theorem~\ref{p01:mom:thm:real-zeros}. Use that theorem with a simple root
at zero. Its transform has exactly those positive roots. Even roots
at $Z$ do not change the sign; the selected simple roots change it
exactly once between consecutive nonzero samples when required.
Multiplying $h$ by $-1$ if necessary matches the first prescribed
nonzero sign. Thus every sign is correct. The perturbation in
Corollary~\ref{p01:mom:cor:real-logconcave} produces the required probability
densities. The last assertion plainly implies the first.
\end{proof}

\begin{example}[Sparse and dense switching behave differently]
\label{p01:mom:ex:switching}
If no zero moments are prescribed, signs can switch precisely at
indices $2,4,8,16,\ldots$, or at $2^2,3^2,4^2,\ldots$, because the
reciprocal sums converge. They cannot switch at every integer, or
at all prime indices, because the corresponding sums diverge.
Thus smooth strict log-concavity permits infinitely many alternating
moment blocks, but it does not permit arbitrary switching frequency.
This is an exact restriction, not merely an upper bound on the number
of changes. It refines the smooth alternating examples in the published
tail-order literature, rather than treating those already established
examples as unresolved problems.
\end{example}

\subsection{Which ultrafilters preserve all moment information?}
\label{p01:mom:sec:ultrafilters}

For $S\subset\mathbb N_0$ set
$w(S)=\sum_{n\in S\setminus\{0\}}1/n$, and let
\[
 \mathcal F_{\rm M}=\{A\subset\mathbb N_0:w(\mathbb N_0\setminus A)<\infty\}.
\]
For a free ultrafilter $\mathcal U$, two real sequences are identified
if their equality set belongs to $\mathcal U$. Let $\Phi_\mathcal U$
map a measure to the equivalence class of its integer moment sequence.
B\"urgin--Epperlein--Wirth's Theorem~6 gives the sufficiency of
$\mathcal F_{\rm M}\subset\mathcal U$. The next result records the
converse, including a strong regularity restriction on the test class.

\begin{theorem}[Exact ultrafilter criterion]
\label{p01:mom:thm:ultrafilter}
Fix $0<a<b$ and $\sigma>1$, and let $\mathcal L_{a,b,\sigma}$ be the
probability measures with $G^\sigma(\mathbb R)$ densities supported on
$[a,b]$, positive and strictly log-concave on $(a,b)$. For a free
ultrafilter $\mathcal U$, the following are equivalent:
\begin{enumerate}
\item $\mathcal F_{\rm M}\subset\mathcal U$;
\item $\Phi_\mathcal U$ is injective on all finite positive measures
with compact support in $(0,\infty)$;
\item $\Phi_\mathcal U$ is injective on $\mathcal L_{a,b,\sigma}$.
\end{enumerate}
\end{theorem}

\begin{proof}
Assume the first condition. If two measures have equality set $E$ in
$\mathcal U$ and $w(E)<\infty$, then $\mathbb N_0\setminus E$ belongs
to $\mathcal F_{\rm M}\subset\mathcal U$, a contradiction. Thus
$w(E)=\infty$, and M\"untz uniqueness gives equality of the measures.
The second condition plainly implies the third.

If the first condition fails, choose $A\in\mathcal F_{\rm M}$ with
$A\notin\mathcal U$. Then $S=\mathbb N_0\setminus A$ belongs to
$\mathcal U$ and satisfies $w(S)<\infty$. Since the ultrafilter is
free, $S$ is infinite. Corollary~\ref{p01:mom:cor:real-logconcave}, with
$\Lambda=S\setminus\{0\}$, gives distinct measures in
$\mathcal L_{a,b,\sigma}$ whose moments agree at least on $S$.
Their equality set belongs to $\mathcal U$, so the third condition
fails. This proves the equivalence.
\end{proof}

\paragraph{What is and is not closed.}
The mathematical assertion called Conjecture~17 in the 2022 paper is
proved, with a clear classical antecedent dating to 1971. Exact real
Mellin-zero sets, multiplicities, and the full sign-pattern criterion
supply useful stronger formulations. The ultrafilter theorem completes
a published one-way criterion, but the cited source did not pose that
converse as a separate open conjecture. No claim is made here that
three or four distinct published open problems have been resolved.


% ===== Consolidated body: quasianalytic =====
\section{A complete regularity criterion for sparse-moment identification}
\label{p01:qa:section}

The Gevrey boundary in the integrated manuscript is one slice of a more
general dichotomy. The ingredients below are classical Denjoy--Carleman
and M\"untz theorems; the conclusion is their application to the same
inverse observation problem. It is an extension of an internal research
question, not the resolution of a separately identified famous conjecture.

Let $M=(M_k)_{k\ge0}$ be a positive log-convex sequence with $M_0=1$,
$M_k\ge k!$, and $M_k^{1/k}\to\infty$. Write $\mathcal C_M[a,b]$ for
the smooth real functions on the closed interval such that
\[
 \|f^{(k)}\|_\infty\le A B^kM_k\qquad(k\ge0)
\]
for some $A,B>0$. We use the classical Denjoy--Carleman criterion:
this class has uniqueness from the full jet at one point precisely when
\begin{equation}\label{p01:qa:DC}
 \sum_{k\ge1}\frac{M_{k-1}}{M_k}=\infty.
\end{equation}
When this sum converges, nonzero nonnegative compactly supported
$\mathcal C_M$ functions exist in every open interval. These standard
facts are stated, with proofs and the relevant endpoint interpretation,
in \cite[Theorem 3.6 and Corollary 3.12]{p01:Rainer2022}.

\begin{theorem}[Quasianalyticity is the exact sparse-moment boundary]
\label{p01:qa:criterion}
Fix $0<a<b$ and an unbounded set $T\subset\mathbb N_0$ containing $0$
with $\sum_{n\in T\setminus\{0\}}1/n<\infty$.
The following conditions are equivalent:
\begin{enumerate}
\item The sum in \eqref{p01:qa:DC} diverges.
\item The $T$-moments identify probability densities within
$\mathcal C_M[a,b]$.
\end{enumerate}
If the sum in \eqref{p01:qa:DC} converges, the signed kernel of the
$T$-moment map contains infinitely many linearly independent
$\mathcal C_M$ functions supported in any prescribed nonempty
interior interval. In particular nonidentification persists among
strictly positive, strictly log-concave densities on $[a,b]$.
\end{theorem}

\begin{proof}
Suppose first that \eqref{p01:qa:DC} holds. If $f,g$ have equal $T$-moments,
put $h=f-g$. If some derivative at $b$ is nonzero, let $m$ be its least
order. Taylor's formula and $t=s(1-x/b)$ give
\begin{equation}\label{p01:qa:watson}
 \int_a^b x^s h(x)\,dx
 \sim (-1)^m h^{(m)}(b)\frac{b^{s+m+1}}{s^{m+1}}
 \qquad(s\to\infty).
\end{equation}
Indeed $|h(x)|\le C(b-x)^m$ throughout $[a,b]$. After multiplication
by $s^{m+1}b^{-s-m-1}$, the transformed integrand is dominated by
$Ct^m e^{-t}$ and converges to
$(-1)^m h^{(m)}(b)t^m e^{-t}/m!$. Dominated convergence proves the
displayed asymptotic. It precludes moment zeros along an unbounded
set. Hence every derivative of $h$ at $b$ is zero. Extend $h$ by zero
just to the right of $b$; it remains smooth and obeys the same derivative
bounds. Quasianalytic uniqueness on this larger interval gives $h=0$.

Conversely suppose the sum in \eqref{p01:qa:DC} converges. Choose a closed
interval $[c,d]$ strictly inside the requested support interval $I$.
M\"untz non-density on $[c,d]$, Hahn--Banach, and Riesz representation
give a nonzero finite real signed measure $\nu$ annihilating $x^n$ for
$n\in T$. Choose a nonnegative nonzero $\mathcal C_M$ bump $\varphi$
supported sufficiently near $1$ that
$[c,d]\operatorname{supp}\varphi\subset I$. Set
\[
 h(x)=\int t^{-1}\varphi(x/t)\,d\nu(t).
\]
Its derivatives satisfy
\[
 \|h^{(k)}\|_\infty
 \le \|\nu\|_{\rm var}c^{-k-1}\|\varphi^{(k)}\|_\infty
 \le A'(B/c)^k M_k.
\]
The moment identity
\[
 \int x^s h(x)\,dx
 =\left(\int t^s\,d\nu(t)\right)
  \left(\int u^s\varphi(u)\,du\right)
\]
annihilates $T$ and also proves $h\ne0$: a nonzero compactly supported
measure has a nonzero integer moment, and the second factor is positive.
Using disjoint subintervals of $I$ gives linear independence.

For the density conclusion take the normalized analytic baseline
$g_0(x)=C\exp(-x^2)$ on $[a,b]$. It belongs to $\mathcal C_M$ because
$M_k\ge k!$. For small $\varepsilon>0$, $g_\pm=g_0\pm\varepsilon h$
are positive and integrate to one. On the compact support of $h$,
$g_0g_0''-(g_0')^2=-2g_0^2$ is bounded strictly below zero, so a
sufficiently small $C^2$ perturbation preserves strict log-concavity;
outside the support there is no change. These densities are smooth
up to the endpoints; no globally smooth extension by zero is claimed.
\end{proof}

If the reciprocal sum for $T$ diverges, classical M\"untz completeness
identifies all finite measures, regardless of regularity. Thus there
are two independent ways to obtain uniqueness: completeness of the
observations, or quasianalyticity of the unknown density.

\begin{corollary}[Fixed derivative bounds and uniform learning]
\label{p01:qa:learning}
Let $S\subset\mathbb N_0$ be unbounded and contain $0$. For the
interval kernel
\[
 k_S(x\mid\theta)=1+\sum_{s\in S\setminus\{0\}}
  2^{-s^2-s-2}L_s(x)\theta^s,
 \quad x\in[0,1],\quad\theta\in[1,2],
\]
where $L_s$ is the shifted Legendre polynomial normalized by
$L_s(1)=1$, consider the probability densities satisfying fixed bounds
$\|g^{(k)}\|_\infty\le A B^kM_k$ for every $k$.
If \eqref{p01:qa:DC} holds, there are estimators from independent one-view
observations whose worst-case expected $W_1$ error tends to zero.
If \eqref{p01:qa:DC} fails and $\sum_{s\in S\setminus\{0\}}1/s<\infty$,
there exist fixed finite $A,B$ for which the worst-case risk is bounded
away from zero for every sample size. No rate for the first assertion
is claimed.
\end{corollary}

\begin{proof}
Legendre orthogonality gives
$\mathbb E_\theta L_s(X)=2^{-s^2-s-2}\theta^s/(2s+1)$ for $s\in S$;
absolute uniform convergence of the series gives the converse.
Thus equality of one-view laws is exactly equality of the $S$-moments.
If the reciprocal sum over $S\setminus\{0\}$ diverges, classical
M\"untz completeness gives injectivity. Otherwise
Theorem~\ref{p01:qa:criterion} gives injectivity in the quasianalytic case.

The set $K$ of densities with the fixed bounds is compact in $C[1,2]$:
successive Arzel\`a--Ascoli subsequences give convergence of every
derivative, and the derivative bounds and probability constraints
are closed under these limits. The observation map $A:K\to
\mathcal P([0,1])$ is continuous. Choose a countable dense family
$(v_j)$ in the unit ball of $C[0,1]$, and metrize weak convergence by
\[
 d(P,Q)=\sum_{j\ge1}2^{-j}|Pv_j-Qv_j|.
\]
For the empirical law $P_N$, boundedness gives
$\sup_P\mathbb E d(P_N,P)\le N^{-1/2}$. A measurable approximate
minimum-distance estimator over a countable dense subset of $K$
satisfies $d(A\widehat g_N,Ag)\le2d(P_N,Ag)+N^{-1}$.
Compactness and injectivity make $A^{-1}$ uniformly continuous on
$A(K)$. Its inverse modulus, Markov's inequality, and boundedness of
$W_1$ on $[1,2]$ yield uniform expected consistency.

In the nonquasianalytic case, the two distinct densities supplied by
Theorem~\ref{p01:qa:criterion} satisfy common finite derivative constants
$A,B$ and induce identical observation laws. The triangle inequality
forces minimax risk at least half their positive $W_1$ distance for
all $N$.
\end{proof}

For example, $M_k=(k!)^\sigma$ gives the familiar threshold $\sigma=1$.
The criterion also handles scales arbitrarily close to analyticity:
if $M_k=k!\prod_{j=1}^k\log(e+j)^\beta$, then
$M_{k-1}/M_k=1/[k\log(e+k)^\beta]$. Identification within the
corresponding class therefore holds for $\beta\le1$, whereas localized
ambiguities exist for $\beta>1$ whenever the observed moments are
M\"untz-sparse. Here $\beta\ge0$ ensures the stated log-convex and
factorial lower-bound hypotheses.


% ===== Consolidated body: prime_flow =====
\section{Prime-zeta questions: a prior-art audit and a uniform counting theorem}
\label{p01:pf:sec:prime-questions}

Belovas, \v{C}epaityt\.{e}, and Sabaliauskas
\cite[Section~3]{p01:pf:BCS} formulate three conjectures in the absolutely
convergent half-plane. The first asks whether their right boundary
$S=1.7795446535\ldots$ is sharp. The second asks whether the ordinates,
normalized by the observation height, become uniform when zeros are
restricted to $\Re s>\gamma$, with $1<\gamma<S$. The third asks for two-sided linear
bounds for their number. We treat these as three separate targets,
not as three newly solved open problems: the first is already implied
by Sepulcre--Vidal \cite[Theorem~4.3]{p01:pf:SV}, and the classical
Jessen--Tornehave theory \cite[Sections~25--26 and Theorems~5--7]{p01:pf:JT}
already supplies the density mechanism behind the other two.
The following direct proof records the additional uniform quantifiers
and removes an ambiguity about counting distinct or multiple zeros.

For positive coefficients $a_p$ satisfying $0<a_-\le a_p\le a_+<\infty$,
put
\[
 F_\omega(s)=\sum_p a_p\omega_pp^{-s},\qquad
 \Omega=\prod_p\{z\in\mathbb C:|z|=1\}.
\]
Let $S>1$ be the unique solution of
\[
 a_2 2^{-S}=\sum_{p\ge3}a_pp^{-S}.
\]
For $1<\gamma<S$, denote by $N_\omega(\gamma;A,T)$ the number of
\emph{distinct} zeros in
$\Re s>\gamma$, $A<\Im s\le A+T$, and by
$N_\omega^{\rm mult}(\gamma;A,T)$ the count with multiplicity.

\begin{theorem}[Uniform linear density and negligible multiple zeros]
\label{p01:pf:thm:uniform-density}
For every fixed $1<\gamma<S$ there is $c_\gamma>0$, depending on the
coefficients but not on their phases, such that
\begin{equation}\label{p01:pf:eq:uniform-density}
 \sup_{\omega\in\Omega,\,A\in\mathbb R}
 \left|\frac{N_\omega(\gamma;A,T)}T-c_\gamma\right|
 \longrightarrow0\qquad(T\longrightarrow\infty).
\end{equation}
The same limit holds for $N_\omega^{\rm mult}$. In particular the
multiplicity excess is $o(T)$ uniformly in the twist and initial height.
For each fixed $x\in(0,1)$,
\begin{equation}\label{p01:pf:eq:uniform-ordinates}
 \frac{N_\omega(\gamma;A,xT)}{N_\omega(\gamma;A,T)}
 \longrightarrow x
\end{equation}
uniformly in $\omega,A$.
\end{theorem}

\begin{proof}
We give the compact-hull argument, including the exceptional-set check.

\emph{Uniform local zero bounds.}
The family $F_\omega$ is compact for locally uniform convergence on
$\Re s>1$: coordinatewise convergence of phases and the summable
majorant give convergence on each compact set. No hull member is
identically zero, since
$2^sF_\omega(s)\to a_2\omega_2\ne0$ along the positive real axis.
Consequently every fixed compact rectangle in this half-plane has a
uniform bound on the number of zeros, counted with multiplicity.
Indeed, otherwise a convergent sequence of hull members would have
unboundedly many zeros in that rectangle; choose finitely many disks
covering it whose boundaries avoid the zeros of the nonzero limit,
and apply Rouch\'e on those boundaries. This gives a contradiction.
The same bound holds after any vertical translation, which merely
changes the phases. The triangle inequality shows that every hull
member is zero-free in $\Re s>S$.

\emph{Uniform bounded densities.}
Fix a compact real interval $I\subset(1,\infty)$ and give $\Omega$
product Haar probability. The random variable $F_\omega(\sigma)$
has a planar density bounded uniformly for $\sigma\in I$.
To see this, retain any five primes. The characteristic function of
their sum is
$\prod_{j=1}^5 J_0(a_{p_j}p_j^{-\sigma}|u|)$, where
$J_0(v)=(2\pi)^{-1}\int_0^{2\pi}e^{iv\cos\theta}\,d\theta$.
The elementary estimate $|J_0(v)|\le C(1+|v|)^{-1/2}$ makes this
product integrable on $\mathbb R^2$, uniformly in $\sigma\in I$.
For completeness, the estimate follows by splitting the integral
into intervals of length $|v|^{-1/2}$ around the stationary points
of $\cos\theta$ and integrating by parts on their complement.
Fourier inversion gives a bounded density; convolution with the
independent remaining sum does not increase its supremum.

The pair $(F_\omega(\sigma),F_\omega'(\sigma))$ likewise has a
bounded density on $\mathbb R^4$, uniformly on $I$. Use ten fixed
primes and put $\lambda_j=\log p_j$. Its characteristic function,
with dual variables $u,v\in\mathbb C$, is bounded by
\[
 \prod_{j=1}^{10}
 \left|J_0\!\left(a_{p_j}p_j^{-\sigma}|u-\lambda_jv|\right)\right|.
\]
There is $b>0$ such that at least nine of the ten quantities
$|u-\lambda_jv|$ are at least $b(|u|^2+|v|^2)^{1/2}$.
Otherwise two of them would both be smaller than this bound;
invertibility of the two distinct linear forms, and a sufficiently
small common $b$, would contradict the norm of $(u,v)$.
Thus the product is $O((1+\|(u,v)\|)^{-9/2})$, which is
integrable in dimension four. Fourier inversion and convolution
again give the claimed bounded density. Haar invariance makes both
density bounds uniform also in the ordinate.

\emph{Boundary and multiple zeros are Haar-null events.}
All derivatives of the hull are uniformly bounded on a fixed compact
rectangle and a neighborhood of it. On any fixed line segment, if a
zero exists, a mesh of spacing $h$ contains a point where
$|F_\omega|\le C h$. The planar density bound and the union bound
give probability $O(h^{-1}h^2)=O(h)$; letting $h\downarrow0$ proves
that the probability of a zero on that segment is zero.

Similarly, if a multiple zero exists in a compact rectangle, a
square mesh of size $h$ contains a point where simultaneously
$|F_\omega|\le C h^2$ and $|F_\omega'|\le C h$.
The four-dimensional density bound gives probability $O(h^6)$ at
one mesh point. There are $O(h^{-2})$ points, so the probability
is $O(h^4)$ and tends to zero. We used the second-derivative bound
and $F'=0$ at the multiple zero in obtaining the stronger $h^2$
bound for $F$. Thus almost every hull member has only simple zeros
in every fixed compact rectangle.

\emph{Counting as an integrable observable.}
Choose $R>S$ and let $n(\omega)$ count distinct zeros in
$\gamma<\Re s<R$, $0<\Im s\le1$; define $n_{\rm mult}$ analogously.
Both are uniformly bounded by the local zero bound. The possible
discontinuities of $n$ are contained in the phases having a boundary
zero or a multiple zero in the closed rectangle; this is a closed
Haar-null subset by the preceding argument. Rouch\'e shows local
constancy off it. The discontinuities of $n_{\rm mult}$ only require
the boundary exceptional set. Consequently both observables admit
continuous upper and lower approximations whose Haar integrals differ
arbitrarily little. This last assertion follows by taking a small
open neighborhood of their closed discontinuity set, covering its
complement by finitely many neighborhoods of constancy, and using a
partition of unity.

The flow $\Phi_t\omega=(\omega_pp^{-it})_p$ has uniform Haar
averages for every continuous function on $\Omega$. For finite
characters this follows directly from
$\sum k_p\log p\ne0$ for every nonzero finite integer vector, by
unique factorization. Finite character polynomials are uniformly
dense in $C(\Omega)$, so the conclusion extends to continuous
functions. Applying the continuous upper and lower approximations
now gives, uniformly in the initial phase,
\[
 \frac1T\int_0^T n(\Phi_t\omega)\,dt\longrightarrow\int n\,d\mu,
 \qquad
 \frac1T\int_0^T n_{\rm mult}(\Phi_t\omega)\,dt
 \longrightarrow\int n_{\rm mult}\,d\mu.
\]
The two integrals are equal because Haar-almost every hull member
has only simple zeros in the rectangle.
Each zero with ordinate between $1$ and $T$ contributes exactly
one to the first sliding-window integral. Only the two end windows
contribute a discrepancy, bounded uniformly by the local zero bound.
The same assertion holds with multiplicity. Thus both counting
functions divided by $T$ converge uniformly to the common value
$c_\gamma=\int n\,d\mu$.

\emph{Positivity.}
The frozen phases $\omega_2=1$, $\omega_p=-1$ for $p\ge3$ give a
simple zero at $S$, since the normalized germ has derivative
$\sum_{p\ge3}(a_p/a_2)(2/p)^S\log(p/2)>0$.
A vertical phase translation places this zero at ordinate $1/2$.
It is inside the rectangle because $\gamma<S<R$. Rouch\'e then
gives a nonempty open neighborhood of phases for which the rectangle
contains a zero. Every nonempty open subset of the product torus has
positive Haar measure, so $c_\gamma>0$.
Initial height is absorbed into the phase; this proves all
uniformities in \eqref{p01:pf:eq:uniform-density}. Applying the same
limit to $T$ and $xT$ proves \eqref{p01:pf:eq:uniform-ordinates}.
\end{proof}

\begin{corollary}[Status of the three published prime-zeta targets]
\label{p01:pf:cor:three-targets}
For the ordinary prime zeta function, the assertions of
\cite[Conjectures~1--3]{p01:pf:BCS} hold, with Conjectures~2 and~3
restricted to the nonempty strips $1<\gamma<S$. The source first
specifies $0<\gamma<S$ and then imposes $\gamma>1$, so this is
precisely its stated admissible range. For $\gamma\ge S$ the ordinary
prime zeta function has no zeros in $\Re s>\gamma$ and a positive
linear lower bound would be false.
Conjecture~1 follows already from \cite[Theorem~4.3]{p01:pf:SV};
Conjectures~2 and~3 follow from Theorem~\ref{p01:pf:thm:uniform-density}.
Their density mechanism is classical, so these implications are
literature reconciliation rather than a claimed new major theorem.
\end{corollary}
\begin{proof}
For the first assertion, the frozen simple zero and Kronecker
approximation produce zeros of the ordinary prime zeta function
arbitrarily near $S$, while the triangle inequality excludes the
half-plane to its right. At $\Re s=S$, equality would force
$-e^{-it\log(p/2)}=1$ for every odd prime. The coordinates $p=3,5$
already contradict unique factorization, so the boundary itself is
not attained. The first conjecture follows.
The other two assertions are respectively
\eqref{p01:pf:eq:uniform-ordinates} and
$N(\gamma;0,T)=c_\gamma T+o(T)$ with $c_\gamma>0$.
\end{proof}

\paragraph{The unresolved endpoint is not absorbed by this argument.}
The additional statement $M_1(T)=O(T)$ in the discussion of the third
conjecture is not proved here. The constants above depend on a compact
rectangle strictly inside $\Re s>1$; taking $\gamma\downarrow1$ after
$T\to\infty$ does not provide a bound at $\gamma=1$. In particular,
a bound for limiting densities alone would not justify interchanging
these limits. This is an explicit failed endpoint attempt, with the
missing uniformity identified.

\section{An internal extension: logarithmic weights at the kernel-budget boundary}
\label{p01:pf:sec:log-budget}

This section addresses an internal quantitative question, not a
separately verified published conjecture. The power-weight calculation
in the integrated manuscript has a transition at exponent $3$.
Logarithmic weights reveal a second transition at logarithmic exponent
$-2$, including an iterated logarithm exactly at that boundary.

For positive summable weights, define
\[
 \Psi_w(B)=\inf_{\substack{D_p\in\mathbb Z_{\ge0},\ D\text{ finite}\\
                         \sum_pD_p\log p\le B}}
 \sum_pw_p\left(1-\cos\frac{\pi}{D_p+2}\right).
\]
The rectangular positive-kernel extremality theorem in the integrated
manuscript identifies this with the least possible mean energy among
nonnegative normalized kernels with that degree budget.

\begin{theorem}[Complete logarithmic transition at exponent three]
\label{p01:pf:thm:log-budget}
If $w_p\asymp p^{-3}(\log p)^\kappa$, then, as $B\to\infty$,
\begin{equation}\label{p01:pf:eq:log-budget-orders}
 \Psi_w(B)\asymp B^{-2}
 \begin{cases}
 (\log B)^{\kappa+2},&\kappa>-2,\\
 (\log\log B)^3,&\kappa=-2,\\
 1,&\kappa<-2.
 \end{cases}
\end{equation}
For $\kappa<-2$, writing
$L=\sum_pw_p^{1/3}(\log p)^{2/3}<\infty$, the stronger limit holds:
\[
 \lim_{B\to\infty}B^2\Psi_w(B)=\frac{\pi^2}{2}L^3.
\]
\end{theorem}
\begin{proof}
Put $c(d)=1-\cos(\pi/(d+2))$ and
$L(x)=\sum_{p\le x}w_p^{1/3}(\log p)^{2/3}$.
The elementary bounds
$2/(d+2)^2\le c(d)\le\pi^2/(2(d+2)^2)$ and H\"older give
\begin{equation}\label{p01:pf:eq:holder-budget}
 \sum_pw_pc(D_p)
 \ge \frac{2L(x)^3}{(B+2\sum_{p\le x}\log p)^2}.
\end{equation}
Indeed, write each summand of $L(x)$ as
$(w_p/(D_p+2)^2)^{1/3}((D_p+2)\log p)^{2/3}$ before
applying H\"older. Prime partial summation yields
\[
 L(x)\asymp
 \begin{cases}
 (\log x)^{(\kappa+2)/3},&\kappa>-2,\\
 \log\log x,&\kappa=-2,
 \end{cases}
\]
and $L(\infty)<\infty$ exactly when $\kappa<-2$.
Taking $x=B$ in \eqref{p01:pf:eq:holder-budget} proves the first two
lower estimates; a single prime proves the remaining $B^{-2}$
lower estimate.

Suppose first $\kappa>-2$, and set
\[
 A=\frac{B}{C(\log B)^{(\kappa+2)/3}},\qquad
 D_p=\left\lfloor A p^{-1}(\log p)^{(\kappa-1)/3}\right\rfloor.
\]
The active range is $p\le x$ up to a fixed finite initial adjustment,
where $x(\log x)^{-(\kappa-1)/3}=A$ and
$x\asymp B/\log B$. The budget satisfies
\[
 \sum_pD_p\log p
 \le A\sum_{p\le x}\frac{(\log p)^{(\kappa+2)/3}}p
 \ll A(\log x)^{(\kappa+2)/3}\le B
\]
for sufficiently large fixed $C$. As $D_p+1$ is at least the
quantity rounded down, the active energy is bounded by
\[
 C A^{-2}\sum_{p\le x}\frac{(\log p)^{(\kappa+2)/3}}p
 \ll B^{-2}(\log B)^{\kappa+2}.
\]
The inactive tail is
$O(x^{-2}(\log x)^{\kappa-1})
 =O(B^{-2}(\log B)^{\kappa+1})$, which is smaller.
The finitely many initially adjusted coordinates obey the same
bound after increasing the constant.

For $\kappa=-2$, take
\[
 A=\frac{B}{C\log\log B},\qquad
 D_p=\left\lfloor\frac{A}{p\log p}\right\rfloor.
\]
Now $x\log x=A$ and
$x\asymp B/(\log B\log\log B)$.
The budget is $O(A\log\log x)\le B$, while the active energy is
$O(A^{-2}\log\log x)=O(B^{-2}(\log\log B)^3)$.
The tail is
$O(x^{-2}(\log x)^{-3})
 =O(B^{-2}(\log\log B)^2/\log B)$.
This proves the matching upper estimate at the boundary.

Finally suppose $L<\infty$ and put
$t_p=w_p^{1/3}(\log p)^{-1/3}$,
$D_p=\lfloor Bt_p/L\rfloor$. These degrees have finite support
and budget at most $B$. For each fixed prime,
\[
 B^2w_pc(D_p)\longrightarrow
 \frac{\pi^2}{2}L^2w_p^{1/3}(\log p)^{2/3}.
\]
The same expression without $\pi^2/2$ and with an absolute constant
is a summable majorant, since $D_p+1\ge Bt_p/L$.
Dominated convergence proves the upper limiting constant.
For the lower limit, fix finitely many primes. Along any sequence
with bounded scaled energy, their degrees tend to infinity, so
$c(D_p)=(\pi^2/2+o(1))/(D_p+2)^2$.
The H\"older argument then gives the lower constant
$(\pi^2/2)L(x)^3$ on that fixed set. Let the set increase to all
primes. This proves the asserted exact limit.
\end{proof}

\begin{corollary}[Critical-weight recurrence bounds]
\label{p01:pf:cor:log-recurrence}
For exact prime frequencies, or relative prime frequencies with a
fixed prime omitted, let $\mathcal R_w(\varepsilon)$ be the uniform
energy covering time and $\ell=\log(1/\varepsilon)$.
Under the preceding weight assumption,
\[
 \log\mathcal R_w(\varepsilon)\ll\varepsilon^{-1/2}
 \begin{cases}
 \ell^{(\kappa+2)/2},&\kappa>-2,\\
 (\log\ell)^{3/2},&\kappa=-2,\\
 1,&\kappa<-2.
 \end{cases}
\]
Together with the small-ball lower bound
$\log\mathcal R_w(\varepsilon)\gg
\varepsilon^{-1/2}\ell^{(\kappa-3)/2}$, this yields
\[
 \lim_{\varepsilon\downarrow0}
 \frac{\log\log\mathcal R_w(\varepsilon)}
      {\log(1/\varepsilon)}=\frac12
 \qquad\text{for every real }\kappa.
\]
\end{corollary}
\begin{proof}
Use product extremizers at a budget for which
$\Psi_w(B)\le\varepsilon/4$. Their inactive tail has mass at most
$\varepsilon/4$, their Fourier norm has logarithm $O(B)$, and the
linear certificate $(\varepsilon/2-\mathcal E_Q)K$ has arithmetic
height $O(B)$. The elementary rational-separation argument in the
integrated manuscript therefore gives a covering time $\exp(CB)$.
Invert the three orders in \eqref{p01:pf:eq:log-budget-orders}.
The lower bound is the logarithmic-weight small-ball avoidance
estimate. All displayed logarithmic factors are subpowers, proving
the final limit.
\end{proof}

This sharp optimization is a statement about the specified positive
Fourier certificate. It does not close the logarithmic discrepancy
between the actual arithmetic covering time and the entropy lower
bound. No additional conjecture is declared solved by that optimization.


% ===== Consolidated body: zeta =====
\section{Zeta-phase targets: a critical audit and three unresolved transfers}
\label{p01:zet:section}

This section attacks three published distribution problems using the
prime-phase framework. None of the three is claimed to be solved.
The useful output is an independent audit of the elementary and
arithmetic mechanisms, together with an exact obstruction to deducing
equidistribution from phase domination alone. The global constant
\[
 c=\frac32-\frac1{\sqrt2}\cot\frac1{\sqrt2}
   =0.672500703679\ldots
\]
is the previously published Montgomery--Taylor value in the
unconditional theorem of Alp\"oge--Furman, not a new global bound.
Their Theorem A and Sections 2--5 are the relevant primary comparison
\cite{p01:zet:AF}. The checks below do not independently validate a full
formalization or all estimates in the earlier research drafts.

\subsection{The published targets and their exact scales}

\paragraph{Target Z1: interval statistics at the secondary scale.}
Ford--Zaharescu, Conjecture A, and Ford--Soundararajan--Zaharescu,
Conjecture 1, ask for
\begin{equation}\label{p01:zet:eq:FSZ}
 \#\{0<\gamma\le T:\{\alpha\gamma\}\in I\}
 = |I|N(T)+T\int_I g_\alpha(u)\,du+o(T),
\end{equation}
uniformly in circular intervals $I$, for fixed $\alpha>0$.
The zeros are counted with multiplicity. For example, at
$\alpha=\log p/(2\pi)$,
\[
 g_\alpha(u)=-\frac{\log p}{\pi}
       \sum_{k\ge1}p^{-k/2}\cos(2\pi ku).
\]
The primary locators are \cite[Conjecture A, pp.~5--6]{p01:zet:FZ} and
\cite[Conjecture 1, p.~2; Theorems 1--3]{p01:zet:FSZ}.
The latter also distinguishes the sharp discrepancy conjecture from
ordinary uniform distribution. Our shrinking-box assertions concern
relative errors against $|I|N(T)$; even for a fixed interval,
$o(N(T))$ is insufficient for the $o(T)$ remainder in
\eqref{p01:zet:eq:FSZ}.

\paragraph{Target Z2: distribution at the average-gap normalization.}
\c{C}i\c{c}ek--Gonek explicitly state, immediately before Conjecture 1,
that uniform distribution of $\gamma\log T/(2\pi)$ for
$0<\gamma\le T$ is not known. Their Theorem 1.3 proves it under RH
and Gonek's exponential-sum conjecture
\begin{equation}\label{p01:zet:eq:Gonek}
 \sum_{0<\gamma\le T}x^{i\gamma}
 \ll_\varepsilon T x^{-1/2+\varepsilon}
                    +T^{1/2}x^\varepsilon
 \qquad(x,T\ge2).
\end{equation}
See \cite[Section 1, Conjecture 1 and Theorem 1.3]{p01:zet:CG}.
The $k$th Weyl test of the normalized sequence is
$\sum_{0<\gamma\le T}(T^k)^{i\gamma}$, for fixed integer $k\ne0$.
Thus a fixed-prime phase theorem is not the required assertion.
Neither \eqref{p01:zet:eq:Gonek} nor its normalized-zero consequence is
proved by the arguments audited here.

\paragraph{Target Z3: distribution modulo Gram points.}
Arias de Reyna records Odlyzko's conjecture that
\[
 x_n=\frac{\vartheta(\gamma_n)}\pi
\]
is uniformly distributed modulo one; here $\vartheta$ is the
Riemann--Siegel theta function. See \cite[Section 1 and the questions
at the end of Section 1.1]{p01:zet:AR}. This nonlinear normalization is
distinct from that in Target Z2. Its local frequency is
$2k\vartheta'(T)=k\log(T/(2\pi))+O(T^{-2})$.
The paper proves a restricted range of complex-zero exponential
tests in Theorem 3.1, which must not be confused with all integer
Weyl tests for the real ordinates. No completion of the latter
tests follows here.

The literature search performed for this audit found no primary
source resolving Z1--Z3. This is a search outcome, not a proof of
historical priority or an exhaustive current-status certificate.

\subsection{What survives the audit of the localization mechanism}

The following ingredients can be checked without presuming that a
new zeta-distribution theorem has already been established.

\begin{lemma}[Rank and inertia]\label{p01:zet:lem:rank}
Suppose $G=P+Q$ is Hermitian, $P\succeq0$,
$\operatorname{rank}P\le s$, $\operatorname{tr}P\le s$,
$n_+(Q)\le b$, and $s+2b\le M$. Then
\[
 s\ge4\operatorname{tr}G-\operatorname{tr}G^2-2M.
\]
\end{lemma}
\begin{proof}
Let $\lambda_1\ge\cdots\ge\lambda_r>0$ be the positive
eigenvalues of $G$. Since $G\preceq P+Q_+$, min--max gives
$\lambda_{b+j}(G)\le\lambda_j(P)$. Thus the positive tail after
the first $b$ entries has at most $s$ entries and sum at most $s$.
For $q(x)=4x-x^2$, use $q(x)\le0$ on $x\le0$,
$q(x)\le4$ everywhere, and $q(x)\le2x+1$ everywhere. Consequently
$\operatorname{tr}q(G)\le4b+3s\le2M+s$.
\end{proof}

\begin{lemma}[Nontrivial rational resonances are sparse]
\label{p01:zet:lem:resonances}
Let $S$ be a finite set of primes and let $q\ne1$ be a positive
rational supported on $S$. If $n=p^j$, $m=r^h$, with $j,h\ge1$,
and $qn^\epsilon m^\eta=1$, where $\epsilon,\eta\in\{-1,1\}$,
then both $p$ and $r$ belong to $S$. Consequently
\[
 \sum_{qn^\epsilon m^\eta=1}
   \frac{\Lambda(n)\Lambda(m)}{\sqrt{nm}}
 \le\left(\sum_{p\in S}\frac{\log p}{\sqrt p-1}\right)^2.
\]
\end{lemma}
\begin{proof}
If $p\notin S$, its valuation must be cancelled by $r$, so $r=p$.
The remaining exponent must vanish, whence $q=1$, a contradiction.
The same argument applies to $r$. Dropping the equation then
factorizes the positive sum; for each prime,
$\sum_{j\ge1}(\log p)p^{-j/2}=(\log p)/(\sqrt p-1)$.
\end{proof}

For a real even $f$ supported strictly inside $(-1/2,1/2)$, put
$\phi(u)=f(u/L)$ and $\alpha_k=T+2\pi k/L$. Direct Fourier
inversion and Poisson summation give, for fixed $\lambda$ and
sufficiently large $L$,
\begin{equation}\label{p01:zet:eq:weighted-Poisson}
 \sum_{k\in\mathbb Z}e^{i\lambda\alpha_k}
   \widehat\phi(t-\alpha_k)\widehat\phi(t'-\alpha_k)
 =Le^{i\lambda(t+t')/2}
   \int\phi(u-\lambda/2)\phi(u+\lambda/2)e^{-i(t-t')u}\,du.
\end{equation}
There are no aliases because the sum of the two frequency supports,
translated by $\lambda$, contains no nonzero multiple of $L$.
At $\lambda=0,t=t'$, the normalized full-grid squared norm is one.
This confirms the sign, sampling normalization, and location of the
rational twist used in the weighted compression.

Writing $q=b/c$ in lowest terms, the difference $bn-cm$ is a
nonzero integer in a nonresonant opposite-sign term. With
$b,c\le H$ and $n,m\le X$, the relevant frequencies have separation
$\gg(HX)^{-1}$. In an equal-sign term, closeness of $qnm$ to one
instead implies $bnm\asymp c\le H$, which gives an adequate lower
bound as well. Thus the stated rational separation has no missing
prime-pair hypothesis. This verification applies to rational twists;
it does not extend automatically to an arbitrary real frequency.

The exact zero-side decomposition also passes the elementary check:
an off-line reflected pair contributes
\[
 uu^{\mathsf T}+\bar u\bar u^{\mathsf T}
      =2vv^{\mathsf T}-2ww^{\mathsf T},\qquad u=v+iw,
\]
and has at most one positive direction. A small operator norm does
not imply a small positive index. The localization argument must
therefore retain the deep zeros and charge them by their count,
while subtracting the shallow exterior contribution exactly.
Selberg's density bound permits a depth
$d_0=B\log\log T/\log T$ at a cost
$O(N(T,2T)(\log T)^{-B/4})$; this estimate is stated in
\cite[Lemma 2]{p01:zet:FZ}. For a fixed phase buffer, repeated integration
by parts gives arbitrary inverse powers of the physical separation
times $\log T$, enough for the corresponding trace and
Hilbert--Schmidt tail errors.

\begin{proposition}[The exact remaining implication]
\label{p01:zet:prop:conditional-localization}
Let $B$ be a fixed Haar-continuity phase region and let $K$ be a
smooth inner cutoff with $0\le K\le1$. Suppose the exact retained
Weil matrix $H_K$ satisfies
\[
 \operatorname{tr}H_K=(\mu(K)+o(1))N,\qquad
 \operatorname{tr}H_K^2\le(R\mu(K)+o(1))N,
\]
and its simple-zero count $s$ admits the rank--inertia decomposition
of Lemma~\ref{p01:zet:lem:rank} with budget
$M=(\mu(B)+o(1))N$. Then
\[
 \liminf\frac{s}{N}\ge(4-R)\mu(K)-2\mu(B).
\]
Taking inner cutoffs and phase buffers to their limits gives
$\liminf s/N\ge(2-R)\mu(B)$.
\end{proposition}
\begin{proof}
Substitute the stated traces and budget in
Lemma~\ref{p01:zet:lem:rank}. All cutoffs and buffers are fixed before
the height limit. Inner approximation then gives the last assertion.
\end{proof}

No contradiction or sign error was found in these central joints.
That is a narrower conclusion than a new independent certification
of the complete uniform arithmetic proposition and complex Gaussian
box sandwich. In particular, subsequent closures should state
explicitly whether they need the fixed-region theorem, the stronger
shrinking-region theorem, or only the unconditional linear algebra.
The results in the next subsection require none of the zeta estimates.

\subsection{A sharp Fourier obstruction to an equidistribution upgrade}

\begin{theorem}[Sharp character bound under measure domination]
\label{p01:zet:thm:sharp-character}
Let $\mu$ be Haar probability measure on a finite or countable
product of circles, and let $\chi$ be a nontrivial continuous
character. Suppose
\[
 c\mu\le\nu\le\mu,\qquad 0<c<1,
 \qquad m=\nu(\Omega),\qquad
 t=\frac{m-c}{1-c}\in[0,1].
\]
Then
\begin{equation}\label{p01:zet:eq:sharp-character}
 \left|\int\chi\,d\nu\right|
     \le\frac{1-c}{\pi}\sin(\pi t).
\end{equation}
For every prescribed $m\in[c,1]$ equality is attainable. In
particular, domination by fixed positive constants does not imply
vanishing Fourier coefficients or equidistribution after
normalization.
\end{theorem}
\begin{proof}
Write $d\nu=(c+(1-c)g)d\mu$, where $0\le g\le1$ and
$\int g\,d\mu=t$. Haar invariance gives $\int\chi\,d\mu=0$.
The pushforward of $\mu$ by $\chi$ is uniform on the circle: a
nontrivial character of a product of circles is onto. Rotate by a
unit complex number so that the integral to be bounded is real and
nonnegative. Conditional expectation on the character coordinate
reduces the problem to maximizing
\[
 \frac1{2\pi}\int_{-\pi}^{\pi}g_0(\theta)\cos\theta\,d\theta,
 \quad0\le g_0\le1,\quad\frac1{2\pi}\int g_0=t.
\]
Let $A=[-\pi t,\pi t]$ and $a=\cos(\pi t)$. Pointwise,
$(g_0-1_A)(\cos\theta-a)\le0$. The integral of $g_0-1_A$
vanishes, so the displayed objective is at most
$\sin(\pi t)/\pi$. Equality holds for $g_0=1_A$, including the
endpoint cases by continuity. Pulling this density back through
$\chi$ proves attainability.
\end{proof}

For the normalized probability $\nu/m$, the sharp upper bound is
\[
 \frac{(1-c)\sin(\pi t)}{\pi(c+(1-c)t)}.
\]
At the Montgomery--Taylor value of $c$, its maximum over $t$ is
approximately $0.1256444247$, strictly positive. This number is an
information bound for the abstract dominated class; it is not a
claim that zeta phases attain it. The argument is an elementary
rearrangement principle, and no historical novelty is asserted for
that principle.

\begin{corollary}[Lower domination leaves room for fixed angular bias]
\label{p01:zet:cor:not-uniform}
Even if every weak limit of a family of simple-zero phase
submeasures satisfies $c\mu\le\nu\le\mu$, this property alone
cannot prove uniform distribution of the full selected family.
\end{corollary}
\begin{proof}
Take $0<t<1$ and the attaining density from
Theorem~\ref{p01:zet:thm:sharp-character}. Its nonzero character
coefficient remains nonzero after normalization. Hence the desired
conclusion is false for the class specified solely by the measure
inequalities. This counterexample is an obstruction to inference,
not a counterexample involving actual zeta zeros.
\end{proof}

\subsection{Why the three attacks stop at different boundaries}

For Z1, the boundary smoothing and the depth scale are incompatible
with simply upgrading the current proof to its secondary term. A
cutoff boundary layer of width $\eta$ has a Haar error of order
$\eta$. Its Haar-predicted count error is of order $N\eta$;
making this $o(T)$ in the sandwich requires
$\eta=o(1/\log T)$. The existing
complex-box argument removes deep zeros at depth
$d_0\asymp\log\log T/\log T$ and uses an analytic smoothing width
$s$ with $d_0/s=o(1)$ and $s$ smaller than the boundary buffer.
These conditions force $\eta\gg d_0$, the opposite range. This
is an obstruction to this particular smoothing proof, not a lower
bound on the true discrepancy or a disproof of Z1.

For Z2, the rational-character proposition requires
$\log H=o(\log T)$ if its bandwidth is to remain asymptotically
$\log T$. The Weyl test $x=T^k$ lies outside that range. Even when
$T^k$ happens to be rational, its rational height is of order $T^k$,
not $T^{o(1)}$. Treating the associated phase as another fixed
prime coordinate would therefore be an invalid substitution.

For Z3, linearizing the theta phase over a short height interval
does not remove the same obstruction: its slope still has order
$\log T$. The derivative varies slowly, but slow variation is not
small frequency. A short-interval argument would also need a new
uniform zero-count and twisted-moment input at that frequency.

Finally, the abstract fourth-moment certificate remains valid
without giving a zeta closure. The polynomial identities
\[
 1+\frac{4x}{\theta}-q_\theta(x)
   =\frac{(x^2-2x+\theta)^2}{\theta^2},\qquad
 \frac8\theta-q_\theta(x)
   =\frac{(x-2)^2(x^2+2\theta)}{\theta^2},
\]
with
\[
 q_\theta(x)=
 \frac{8\theta x-(4+2\theta)x^2+4x^3-x^4}{\theta^2},
\]
yield $s\ge(L_+)^2/D$ in the rank--inertia model, where
$L=4\operatorname{tr}G-\operatorname{tr}G^2-2N$ and
$D=\operatorname{tr}(G(2I-G))^2$. The input moment vector
$(1,4/3,2,13/4)$ would give $16/21$, but the identities do not
establish that vector for a zeta compression. None of Z1--Z3 is
closed by postulating it.

\paragraph{Outcome of this branch.}
There are three documented published targets and zero claimed
closures. The rigorously established addition is the sharp
domination-to-Fourier bound, which identifies exactly what one
cannot infer from the phase-measure comparison. The fixed-rational
resonance argument and elementary localization reduction survive
the present audit; their strongest analytic extensions still need
a complete independent verification before being used to advertise
a solution to an external open problem.


\clearpage
\begingroup
\small

\endgroup
\endgroup
% ===== END SOURCE 01 =====


% ===== BEGIN SOURCE 02: KEIO_Observation_Geometry_and_Twenty_Targets.tex =====
\clearpage
\part{Observation Geometry and Exact Recovery}
\label{part:02}
\setcounter{section}{0}
\begin{center}\large\itshape Replication Rank, Thorin Measures, and Zero Patterns\end{center}
\begin{center}\small Source manuscript: 14 October 2026.\end{center}
\begingroup
\renewcommand{\arraystretch}{1}
\begin{quote}\small\textbf{Source abstract.}\quad
We extend two earlier manuscripts through several exact recovery and
classification problems. For arbitrary Gaussian mixing laws with a linear
covariance constraint, the sharp replication threshold is the smallest
rank at which a projected Gram map has interior, plus one for an unknown
mean. For covariance graphs this is the generic completion rank plus one.
A multilevel Fourier construction gives the sharp adaptive total-budget
Wasserstein rate $B^{-1/(d+2)}$ for regular $d$-dimensional components,
including dimension two without a logarithm; support-dependent
uniform components give $B^{-1/(d+1)}$ for $d\ge2$.
Symmetry and forced atoms of Thorin measures yield complete GGC
classifications for two generalized McKay families. A recently published
Bessel-zero order theorem implies self-decomposability in all three
families of a stated external problem.
Additional results include exact early-time heat sampling criteria,
Cantor families of shape-constrained phase ambiguities, an exact random
recurrence entropy constant, and inverse recovery of all prime phases
from the right-edge critical-point divisor of an Euler product.
For finite exponential sums we prove equidistribution of zero ordinates
modulo every length outside a countable resonance set, answering a
published existence question. We explicitly separate these conclusions
from remaining conjectures and from established inputs. Twenty targets
are assessed: fifteen have complete answers in their specified scope,
one has a special-case answer, and four remain open here.
These counts refer to research targets, not to twenty independent
historically open conjectures.
\end{quote}
\thispagestyle{empty}


\section{The research question and the scope of the answers}
\label{p02:intro:scope}

The starting point is the 223-page integrated manuscript
\emph{Observation and Recovery in Arithmetic and Probability}
\cite{p01:KEIOIntegrated}, followed by the 42-page manuscript
\emph{Exact Recovery, Optimal Replication, and Bessel Classification}
\cite{comp:part:01}. The present paper is a separate, self-contained
sequel. Its purpose is to take the tools that survived those drafts,
formulate further questions with explicit quantifiers, and determine
which implications can actually be proved.

The common object is an observation map from a hidden parameter or
law to a transform, a zero pattern, or a finite collection of experiments.
Three mechanisms recur. Algebraic equations satisfied by attainable
transform arguments produce compactly supported invisible perturbations.
Analytic uniqueness turns a sufficiently rich observation set into
identification. Positive representing measures turn classification into
support and mass constraints. These mechanisms then transfer to graph
completion, statistical allocation, heat observations, and zero dynamics.
No claim that all these applications are instances of one universal
theorem is needed.

\subsection{What changes from the two predecessors}

The previous scalar replication result does not decide arbitrary random
covariance models. Theorems~\ref{p02:frontmix:thm:fixed-threshold} and
\ref{p02:graph:threshold} replace a scalar count by the interior rank of
a projected Gram map. This yields an exact combinatorial interpretation
for graph-constrained covariances, and counterexamples that survive
smoothness, strong log-concavity, and finitely many moment constraints.
Theorem~\ref{p02:frontmix:thm:terminal-threshold} further distinguishes a
fixed experiment from changing terminal lengths.

The previous budget exponent $1/3$ becomes $1/(d+2)$ for general
regular $d$-dimensional components, with matching lower bounds against
every adaptive policy. The proof requires a transport-controlling
Fourier norm in place of the one-dimensional distribution function.
The different exponent $1/(d+1)$ for uniform endpoints exposes the role
of component information, rather than dimension alone.

For the two remaining McKay GGC classifications, we give both positive
Thorin measures and necessary mass inequalities. This supplies the
missing directions of the classification, not just further sufficient
examples. For the separate SD question we use an external order theorem
posted on 7 October 2026 \cite{p02:bes:BariczPrabhu2026}; that recent theorem
is an input, and its authors receive explicit credit.

The recurrence theorem identifies a leading constant for rapidly
decreasing weights, a stronger conclusion in a different spectral class
from the earlier order bounds. The zero-divisor theorem is an inverse
statement: even discarding the values of the Euler logarithmic derivative
retains every prime phase if its near-edge zeros are kept. Finally, the
modular-ordinate theorem answers an external existence question for finite
exponential sums and supplies equidistribution for generic lengths.

\subsection{Twenty targets, with no conjecture-count inflation}

In the following register, \textbf{E} means a question or a clause located
in an external primary source; \textbf{N} means a question formulated
within this project. A complete entry means that the precise statement
in its row is proved below, including the negative direction when claimed.
It does not certify historical novelty. In particular E1--E3 are clauses
of one multipart source problem, and E4 contains three distributional
cases of another source problem.

\begingroup
\fontsize{9.5}{11.2}\selectfont
\setlength{\tabcolsep}{4pt}
\renewcommand{\arraystretch}{1.08}
\begin{longtable}{@{}p{0.65cm}>{\raggedright\arraybackslash}p{6.3cm}>{\raggedright\arraybackslash}p{7.45cm}@{}}
\toprule
ID & Precisely scoped target & Result and proof location\\
\midrule
\endfirsthead
\toprule
ID & Precisely scoped target & Result and proof location\\
\midrule
\endhead
E1 & GGC classification of the third McKay density over its full density domain.
& Complete: $\nu\le\mu+1$ is necessary and sufficient.
Theorem~\ref{p02:bes:thm:third-family}.\\
E2 & GGC classification of the squared-Bessel fourth McKay density.
& Complete: $-1/4<\mu\le1/2$ is exactly the GGC range.
Theorem~\ref{p02:bes:thm:fourth-family}.\\
E3 & HCM status of the additional density in source Theorem 12.
& Complete: no $\mu>0$ is HCM.
Proposition~\ref{p02:bes:prop:theorem12}.\\
E4 & SD of the three families asked about in source \S4.4.
& Complete on all stated source ranges, using the credited Bessel-zero
order theorem. Theorem~\ref{p02:bes:thm:third-source-problem}.\\
E5 & Existence of lengths giving dense zero ordinates modulo the length
for rationally independent exponential sums.
& Complete, with a stronger result: equidistribution for every
nonresonant length. Theorem~\ref{p02:zp:thm:podesta}.\\
N6 & Fixed joint replication needed for arbitrary Gaussian mean/covariance mixing.
& Complete: $d+1$ observations, or $d$ for a known mean.
Theorems~\ref{p02:frontmix:thm:fixed-threshold} and \ref{p02:frontmix:thm:wishart}.\\
N7 & Exact replication count for a linear covariance space or covariance graph.
& Complete: $r(V)+1$, respectively $\operatorname{gcr}(G)+1$, with unknown
mean. Theorem~\ref{p02:graph:threshold} and Corollary~\ref{p02:graph:gcr}.\\
N8 & Minimal vector-output count and terminal-length criterion with bounded covariance.
& Complete: $d$ outputs; an infinite set of terminal lengths is necessary
and sufficient. Theorems~\ref{p02:frontmix:thm:terminal-threshold} and
\ref{p02:frontmix:thm:rank-barrier}.\\
N9 & Optimal total observation budget for regular multivariate mixing laws.
& Complete under stated component assumptions:
$B^{-1/(d+2)}$ for every $d$.
Theorems~\ref{p02:bud:thm:upper}--\ref{p02:bud:thm:lower}.\\
N10 & Optimal grouped budget for unknown uniform-rectangle endpoint populations.
& Complete for $d\ge2$: $B^{-1/(d+1)}$.
Corollary~\ref{p02:bud:cor:uniform-endpoints}.\\
N11 & Whether strong log-concavity, Gevrey regularity, and finitely many
moments eliminate continuous Fourier-intensity ambiguity.
& Complete negative answer: a compact Cantor family remains, even after
translation and reflection are removed.
Theorem~\ref{p02:mt:phase-cantor}.\\
N12 & Exact identification from early-time heat observations with derivative multiplicities.
& Complete: $\sum_j r_jt_j=\infty$; localized shape-constrained
counterexamples otherwise. Theorem~\ref{p02:mt:heat-criterion}.\\
N13 & Exact leading logarithmic covering constant for rapidly decreasing weights
and independent random frequencies.
& Complete under the stated width condition: covering entropy divided by
Haar small-ball entropy tends to one.
Theorem~\ref{p02:pt:thm:exact-entropy}.\\
N14 & Whether one exact radial law determines a summable weighted product of circles.
& Complete: every positive weight and its multiplicity are recoverable.
Theorem~\ref{p02:pt:thm:radial-rigidity}.\\
N15 & Whether the near-edge Euler critical-point divisor determines the entire prime twist.
& Complete: coordinatewise inverse recovery and topological conjugacy.
Theorem~\ref{p02:zp:thm:inverse}.\\
E16 & Limiting zero density for arbitrary atomic stationary Gaussian analytic spectra.
& Special case: an explicit conditional density for three frequencies
whose two nonzero values have irrational ratio. The general program remains open here.
Theorem~\ref{p02:zp:thm:three-atoms}.\\
N17 & Matching entropy-order covering bounds for the exact prime flow with power exponent
$\sigma>3$.
& Open here. One anchor prime's exponential degree cost can be removed,
but the exponent gap remains.
Proposition~\ref{p02:pt:prop:anchor-height}.\\
N18 & SD classification on the extended density domains of the third and fourth McKay laws.
& Open outside the source's affirmative parameter ranges.
GGC classification does not decide every non-GGC SD case.\\
E19 & Whether bounded maximum likelihood threshold forces bounded generic completion rank.
& Open here. The replication interpretation does not supply the missing
graph rank bound; see \cite[p.~4]{p02:graph:BS2017}.\\
E20 & The original equal-group binomial population-MLE range question.
& Open here. Optimal unequal allocation does not settle the prescribed
equal-group estimator; see \cite[Remark 3.2]{p01:info:VKVK2019}.\\
\bottomrule
\end{longtable}
\endgroup

The fifteen completed rows consist of five external-question clauses
across three source groups and ten newly formulated targets.
They are not fifteen independent previously recognized conjectures.
Among the three external source groups, the SD question in
\cite[\S4.4]{p01:rare:Baricz2026} and the modular-density existence question
in \cite[Section 1]{p02:zp:Podesta2026} receive affirmative answers throughout
their stated scope. The broader McKay program retains the SD extension
in N18. E16 is only a special-case result, and rows 17--20 are not counted
as completed.

The source audit was targeted, with a cutoff of 14 October 2026.
Some arguments may already be implicit in older general results.
In particular, Wishart deconvolution, hypergeometric ratio measures,
multilevel Monte Carlo, almost-periodic spectra, and convolution phase
ambiguities are established theory. Their use here is distinguished
from the additional statements and proofs. A claim of an Annals-level
breakthrough, historical priority, or the resolution of a famous
longstanding conjecture would require evidence beyond this manuscript.

\subsection{Reading map and notation}

Sections~\ref{p02:frontmix:sec:main} and \ref{p02:graph:section} develop exact
observability and its graph interpretation.
Section~\ref{p02:bud:sec:main} studies finite-budget statistical recovery.
Sections~\ref{p02:bes:sec:mckay} and \ref{p02:bes:sec:sd} concern positive
distributional representations.
Section~\ref{p02:mt:section} treats exact ambiguity and heat observations.
Sections~\ref{p02:pt:sec:main} and \ref{p02:pt:sec:radial-inverse} connect
recurrence and radial geometry.
Section~\ref{p02:zp:sec:scope} studies zero patterns and the external
exponential-sum question.

All observation-identification assertions concern exact probability laws,
not one finite sample. Statistical assertions explicitly use the total
number $B$ of observations and expected Wasserstein loss.
Zero divisors always retain multiplicity. Constants may depend on fixed
model parameters but not on the asymptotic budget or tolerance.
Detailed assumptions appear in each theorem.


% ===== Consolidated body: mixtures =====
\section{Multivariate Gaussian environments: replication, rank, and changing time scales}
\label{p02:frontmix:sec:main}

The fixed univariate stable-shape results do not describe an arbitrary
random covariance matrix. In dimension $d$, the relevant observation
geometry is the rank of a matrix of Fourier frequencies. This section
derives three classifications: the exact fixed-group threshold $d+1$
for unknown mean and covariance; the threshold $d$ when changing
terminal-time designs are available; and the threshold $d$ for centered
covariance mixtures, equivalently integer-degree Wishart mixtures.
These are internally posed extensions. No source inspected in the
priority audit explicitly posed these exact classifications as open
conjectures, and they are not counted as additional published open
problems. The scalar normal question is not counted again.

Let $\mathbb S_d$ be the real symmetric matrices, with positive
semidefinite and positive definite cones $\mathbb S_d^+$ and
$\mathbb S_d^{++}$. Draw $(A,\Sigma)$ from an arbitrary Borel
probability $\mu$ on $\mathbb R^d\times\mathbb S_d^+$ and, conditional
on this draw, let $X_j$ be iid $N_d(A,\Sigma)$. Degenerate covariance
matrices are permitted in the positive results. Write
$S_n=\sum_{j=1}^nX_j$. Equality of observation laws always means
equality of their complete distributions, rather than finitely many
moments or one finite dataset.

\subsection{The Fourier--Laplace transform and its attainable arguments}

For a finite signed mixing measure $\rho$, define
\begin{equation}\label{p02:frontmix:eq:transform}
 H_u(Z)=\int \exp\{iu^Ta-\operatorname{tr}(Z\Sigma)\}
                  \,d\rho(a,\Sigma),
 \qquad u\in\mathbb R^d,\quad \operatorname{Re}Z\in\mathbb S_d^{++}.
\end{equation}
For each fixed $u$, this is bounded and holomorphic on the indicated
tube domain. Indeed, on a compact subset of the domain,
$\operatorname{tr}((\operatorname{Re}Z)\Sigma)\ge c\operatorname{tr}\Sigma$;
every derivative is bounded by a polynomial in
$\operatorname{tr}\Sigma$ times $e^{-c\operatorname{tr}\Sigma}$.
No moment of $A$ or $\Sigma$ is required.

We use the following elementary uniqueness fact. If $H_u(T)=0$ for
every real $T\in\mathbb S_d^{++}$ and every $u$, then $\rho=0$.
Holomorphic uniqueness first gives $H_u(I-iV)=0$ for all
$V\in\mathbb S_d$. These are all Fourier coefficients of the finite
signed measure $e^{-\operatorname{tr}\Sigma}\rho$ on the Euclidean
space $\mathbb R^d\times\mathbb S_d$. Fourier uniqueness and the
strict positivity of the tilt prove the assertion. Holomorphic
uniqueness on a real open set can be justified by iterating the
one-variable identity theorem in a small coordinate box and then
continuing on the connected tube domain.

\begin{lemma}[Weighted Gram realization]
\label{p02:frontmix:lem:gram}
Let $w_1,\ldots,w_q>0$ and $\tau=\sum_jw_j$. The possible pairs
\[
 u=\sum_{j=1}^q w_jt_j,\qquad
 T=\frac12\sum_{j=1}^q w_jt_jt_j^T,
 \qquad t_j\in\mathbb R^d,
\]
are exactly those for which
\begin{equation}\label{p02:frontmix:eq:gram-cone}
 C=T-\frac{uu^T}{2\tau}\succeq0,
 \qquad\operatorname{rank}C\le q-1.
\end{equation}
\end{lemma}
\begin{proof}
Put $v=(\sqrt{w_1},\ldots,\sqrt{w_q})^T$ and let the $j$th
column of $Y$ be $\sqrt{w_j}t_j$. Then $Yv=u$. The orthogonal
decomposition
\[
 Y=\frac{uv^T}{\tau}+Z,\qquad Zv=0,
 \qquad YY^T=\frac{uu^T}{\tau}+ZZ^T
\]
proves necessity. Conversely, factor $2C=LL^T$ with $L$ having
$r=\operatorname{rank}C$ columns. Choose a $q\times r$ matrix $V$
with orthonormal columns in $v^\perp$ and set $Z=LV^T$.
The resulting $Y$ gives the required frequencies. The case $r=0$
uses the zero matrix $Z$.
\end{proof}

For a fixed group $(X_1,\ldots,X_q)$ the joint characteristic
function is $H_u(T)$ with $w_j=1$. For independent Gaussian blocks
of lengths $w_j$ the same formula holds with the displayed weights:
conditionally a block has mean $w_jA$ and covariance $w_j\Sigma$.

\subsection{Determinantal differential operators give local counterexamples}

Coordinates on $\mathbb S_d$ are $\sigma_{ij}$, $i\le j$, with
Lebesgue measure in these independent coordinates. Define the
symmetric matrix of commuting differential operators
\[
 (\nabla_\Sigma)_{ii}=\partial_{\sigma_{ii}},\qquad
 (\nabla_\Sigma)_{ij}=\tfrac12\partial_{\sigma_{ij}}
       \quad(i<j),\qquad
 (\nabla_\Sigma)_{ji}=(\nabla_\Sigma)_{ij}.
\]
Integration by parts against the transform in
\eqref{p02:frontmix:eq:transform} replaces $\nabla_\Sigma$ by $T$,
and replaces $\nabla_a\nabla_a^T$ by $-uu^T$. Consequently
\begin{equation}\label{p02:frontmix:eq:operator-symbols}
 \mathcal D_\tau=\det(2\tau\nabla_\Sigma+
                       \nabla_a\nabla_a^T),\qquad
 \mathcal E=\det(\nabla_\Sigma)
\end{equation}
have transform multipliers
$\det(2\tau T-uu^T)$ and $\det T$, respectively.
Every determinant here is a polynomial in commuting constant
coefficient operators; it is not a determinant of noncommuting
operators. In particular, these are nonzero operators. For example,
at $u=0$ the first multiplier is $(2\tau)^d\det T$.

\begin{lemma}[Local smooth perturbations with protected marginals]
\label{p02:frontmix:lem:local-perturbations}
Let $\mathcal P$ be any nonzero real constant coefficient
differential operator on $\mathbb R^d\times\mathbb S_d$.
For every nonempty open set
$U\subset\mathbb R^d\times\mathbb S_d^{++}$ and every integer
$K\ge0$, there are distinct compactly supported $C^\infty$
probability densities $p_+,p_-$ inside $U$ whose difference is
\[
 p_+-p_-=2\epsilon\,
 \partial_{a_1}^{K+1}\partial_{\sigma_{11}}^{K+1}
                       \mathcal P f
\]
for a real $f\in C_c^\infty(U)$ and $\epsilon>0$.
They have the same full $A$ marginal, the same full $\Sigma$
marginal, and the same joint polynomial moments of total degree at
most $K$. They are strictly positive on a neighborhood of the
support of their nonzero difference.
\end{lemma}
\begin{proof}
Choose a nonzero real bump $f$ in a sufficiently small ball with
closure in $U$ and denote the displayed differentiated function by
$g$. A nonzero constant coefficient operator cannot kill a nonzero
compactly supported smooth function: its Fourier transform would
give $P(i\xi)\widehat f(\xi)=0$, and the complement of the zero
set of the nonzero polynomial is dense. Continuity would imply
$\widehat f=0$. Hence $g\ne0$.
Choose a smooth nonnegative probability density $p$ supported in
a slightly larger ball in $U$, with $p$ bounded below by a positive
constant on $\operatorname{supp}g$. For sufficiently small
$\epsilon>0$, put $p_\pm=p\pm\epsilon g$. Integration of a
compactly supported derivative shows $\int g=0$, so these are
probabilities. Integrating first in $a_1$ or in $\sigma_{11}$
shows that both coordinate marginals of $g$ vanish. Integration
by parts $K+1$ times in $a_1$ kills every polynomial of total
degree at most $K$. The same argument shows that $\epsilon$ can
be made arbitrarily small while the two densities remain distinct.
\end{proof}

Thus the negative results below hold in every small interior
parameter neighborhood, with covariance uniformly bounded and
bounded away from singularity. They do not arise from heavy tails,
zero covariance, or a failure to know the separate parameter
marginals. Compact support precludes matching all polynomial
moments for distinct measures; only an arbitrarily prescribed
finite degree is asserted.

\subsection{Exact fixed-group replication threshold}

\begin{theorem}[Unknown mean and full covariance require exactly $d+1$ views]
\label{p02:frontmix:thm:fixed-threshold}
The map
\[
 \mu\longmapsto\int N_d(a,\Sigma)^{\otimes q}\,d\mu(a,\Sigma)
\]
is injective on all Borel probabilities on
$\mathbb R^d\times\mathbb S_d^+$ if and only if $q\ge d+1$.
For every $q\le d$, every neighborhood $U$ and every finite moment
degree $K$ as in Lemma~\ref{p02:frontmix:lem:local-perturbations}, the
nonidentifiability witnesses can have all the properties in that
lemma.
\end{theorem}
\begin{proof}
For $q\ge d+1$, Lemma~\ref{p02:frontmix:lem:gram} gives, for every
fixed $u$, all $T$ in the real open set
$T-uu^T/(2q)\succ0$. Equality of observed laws therefore makes
$H_u$ vanish on this open set. Holomorphic uniqueness gives its
vanishing on the whole tube domain, and Fourier--Laplace
uniqueness gives equality of the mixing laws.

For $q\le d$, every attainable pair satisfies
$\det(2qT-uu^T)=0$, because the centered Gram matrix has rank at
most $q-1<d$. Apply Lemma~\ref{p02:frontmix:lem:local-perturbations}
with $\mathcal P=\mathcal D_q$. The transform of the difference
of its two densities vanishes at every attainable pair. Thus all
joint characteristic functions of the $q$ observations agree,
although the mixing densities are distinct.
\end{proof}

This gives a finite certificate for equality of entire Gaussian
exchangeable laws: equality of their $(d+1)$-coordinate marginals
forces equality of their directing parameter laws and hence of all
their finite-dimensional distributions. No smaller fixed marginal
has that property, even under the local smooth restrictions above.
The number $d+1$ counts observations sharing one latent parameter;
it does not count the number of independently sampled groups used
in statistical estimation.

\subsection{Changing the terminal length lowers the dimension threshold}

For $n\ge d$ define the $d$-vector observation group
\begin{equation}\label{p02:frontmix:eq:terminal-design}
 \mathcal O_n=(X_1,\ldots,X_{d-1},S_n).
\end{equation}
When $d=1$, this means only $S_n$. The separate laws of these
groups for different $n$ are given; no joint coupling between
experiments with different terminal lengths is assumed observed.

\begin{theorem}[A rank-completing sequence of partial-sum experiments]
\label{p02:frontmix:thm:terminal-threshold}
Let $\mathcal N\subset\{d,d+1,\ldots\}$.
\begin{enumerate}
\item If $\sum_{n\in\mathcal N}n^{-1}=\infty$, the laws
$\{\mathcal L(\mathcal O_n):n\in\mathcal N\}$ identify every
Borel mixing law on $\mathbb R^d\times\mathbb S_d^+$, without
moments.
\item On the class whose covariance support is bounded, the same
design is injective if and only if $\mathcal N$ is infinite. No
tail restriction on the mean is needed for this assertion.
\item For every finite $\mathcal N$, there are local smooth
nonidentifiability witnesses with identical parameter marginals
and moments through any fixed finite degree, as in
Lemma~\ref{p02:frontmix:lem:local-perturbations}.
\end{enumerate}
\end{theorem}
\begin{proof}
An invertible linear transformation turns $\mathcal O_n$ into the
independent conditional blocks
\[
 X_1,\ldots,X_{d-1},\quad S_n-S_{d-1},
\]
with positive lengths $1,\ldots,1,n-d+1$ and total length $n$.
Fix $u\ne0$ and a positive semidefinite matrix $C$ of rank $d-1$
such that $u$ is not orthogonal to $\ker C$. Lemma~\ref{p02:frontmix:lem:gram}
shows that equality of the observed laws gives
\[
 H_u\left(C+\frac{uu^T}{2n}\right)=0
             \qquad(n\in\mathcal N).
\]
The function
\[
 F(z)=H_u(C+zuu^T),\qquad\operatorname{Re}z>0,
\]
is bounded and holomorphic, because $C+(\operatorname{Re}z)uu^T$
is positive definite. A nonzero bounded holomorphic function in a
half-plane has zeros satisfying the Blaschke sum
$\sum\operatorname{Re}z_j/(1+|z_j|^2)<\infty$. Applying this
condition to $z_n=1/(2n)$ proves $F=0$ under the condition in (1).

Every $T\in\mathbb S_d^{++}$ is reached in this way. Indeed, put
\[
 \lambda=(u^TT^{-1}u)^{-1},\qquad C=T-\lambda uu^T.
\]
Conjugation by $T^{-1/2}$ shows that $C\succeq0$ and has rank
$d-1$. Its kernel is spanned by $T^{-1}u$, whose inner product
with $u$ is positive. Since $T=C+\lambda uu^T$, we get
$H_u(T)=0$. Continuity in $u$ includes $u=0$, and transform
uniqueness proves (1).

If both candidate covariance supports are bounded, then $F$
extends to an entire function of $z$: the quantity $u^T\Sigma u$
is uniformly bounded, while $e^{iu^Ta}$ has modulus one. An
infinite $\mathcal N$ gives zeros accumulating at zero, so the
ordinary identity theorem suffices. The same argument proves
the positive half of (2).

For a finite nonempty $\mathcal N$, use in
Lemma~\ref{p02:frontmix:lem:local-perturbations} the nonzero operator
\[
 \mathcal P=\prod_{n\in\mathcal N}\mathcal D_n.
\]
At the Fourier--Laplace arguments of experiment $n$, its factor
$\det(2nT-uu^T)$ vanishes by the weighted Gram relation. The same
single pair of compactly supported densities is therefore
invisible to every experiment in the finite list. For the empty
list take $\mathcal P$ to be the identity. This proves (3) and
the negative half of (2).
\end{proof}

The harmonic condition in (1) is only asserted to be sufficient
for unrestricted covariance laws. The exact sparse condition in
that unrestricted multivariate setting is not settled here. The
``infinite if and only if'' statement is restricted to bounded
covariance support.

\begin{theorem}[Fewer than $d$ vector outputs cannot recover an arbitrary covariance]
\label{p02:frontmix:thm:rank-barrier}
Suppose $d\ge2$ and $q<d$. Allow every experiment consisting of
$q$ vector outputs
\[
 Y_\ell=\sum_{j=1}^m w_{\ell j}X_j,
       \qquad 1\le\ell\le q,
\]
with arbitrary finite $m$ and arbitrary known real scalar weights
$w_{\ell j}$. Even the collection of all these experiment laws
does not identify the Gaussian mixing law. A single pair of
local smooth mixing densities, with identical separate parameter
marginals and any prescribed finite degree of equal joint moments,
is invisible to every one of these experiments simultaneously.
In particular, all separate partial-sum laws fail to identify the
mixing law in every dimension $d\ge2$.
\end{theorem}
\begin{proof}
For output frequencies $t_1,\ldots,t_q$, put
$v_j=\sum_\ell w_{\ell j}t_\ell$. The transform is evaluated at
\[
 u=\sum_jv_j,\qquad T=\tfrac12\sum_jv_jv_j^T.
\]
Every $v_j$ lies in the span of $t_1,\ldots,t_q$, so
$\operatorname{rank}T\le q<d$ and $\det T=0$. Apply the local
perturbation lemma with $\mathcal P=\mathcal E$.
Its transform vanishes at every such argument, independently of
the number or values of the weights. For separate partial sums,
take $q=1$ and all weights equal to one.
\end{proof}

The weights in this theorem are scalar and each output is a
$d$-dimensional vector. Matrix-valued weights can alter the rank
geometry and are not covered. Combining the last two theorems
gives the sharp output count $d$ for this class of changing-length
experiments, compared with $d+1$ for one fixed iid group.

\subsection{Centered mixtures and singular Wishart observations}

\begin{theorem}[Exact centered covariance and integer Wishart thresholds]
\label{p02:frontmix:thm:wishart}
If the mean is known and fixed, $q$ shared-latent Gaussian
replicates identify an arbitrary covariance mixing law on
$\mathbb S_d^+$ if and only if $q\ge d$.
For covariance laws supported in $\mathbb S_d^{++}$, the same
threshold holds if only
\[
 W_q=\sum_{j=1}^q(X_j-A)(X_j-A)^T
\]
is observed, that is, for mixtures of Wishart laws with known
integer degrees of freedom $q$. When $q<d$, compactly supported
smooth covariance mixing densities already give counterexamples.
\end{theorem}
\begin{proof}
Subtract the known mean. For the full group the characteristic
transform of a signed covariance mixing measure is
$L(T)=\int e^{-\operatorname{tr}(T\Sigma)}d\rho(\Sigma)$ with
$T=\tfrac12\sum_jt_jt_j^T$. If $q\ge d$, all positive definite
$T$ are attainable, and Laplace uniqueness proves identification.
If $q<d$, take a nonzero real bump $f$ compactly supported in any
small positive definite neighborhood. The nonzero function
$g=\det(\nabla_\Sigma)f$ has total integral zero and transform
$\det(T)L_f(T)$, which vanishes on every rank-$q$ probe.
Adding and subtracting a sufficiently small multiple of $g$ to
a smooth probability bump positive on its support produces the
claimed covariance densities. This already makes their $W_q$
laws equal.

It remains to prove sufficiency from $W_q$ for $q\ge d$ and
positive definite covariance. Its conditional density, with
respect to Lebesgue measure on the positive definite cone, is
\[
 c_{d,q}|W|^{(q-d-1)/2}|\Sigma|^{-q/2}
          e^{-\operatorname{tr}(\Sigma^{-1}W)/2}.
\]
Equality of two mixture laws therefore gives a zero difference
of these densities almost everywhere. The integral after
dividing by the positive factor in $W$ is continuous on the
positive definite cone: on a compact set with $W\succeq cI$,
the quantity
$|\Sigma|^{-q/2}e^{-c\operatorname{tr}(\Sigma^{-1})/2}$ is
uniformly bounded. Thus the equality holds everywhere.

Push the signed covariance difference forward under
$B=\Sigma^{-1}$, and fix $W_0\succ0$. The signed measure
\[
 d\eta(B)=|B|^{q/2}e^{-\operatorname{tr}(BW_0)/2}\,d\rho_*(B)
\]
is finite, by the same eigenvalue bound. Its Laplace transform
at $V\succ0$ is the vanished mixture difference at $W_0+2V$.
Laplace uniqueness gives $\eta=0$, and strict positivity of the
weight implies $\rho_*=0$ and $\rho=0$.
\end{proof}

The sufficiency side belongs naturally to classical Wishart
deconvolution, rather than constituting a newly solved conjecture.
Haff, Kim, Koo, and Richards develop nonparametric estimation of
Wishart mixing densities on the positive definite cone using
Helgason--Fourier analysis\cite{p02:frontmix:Haff2011}. The rank
argument above explicitly exhibits the obstruction in the
singular integer-degree range.

\subsection{A transfer to random anisotropic diffusions}

\begin{corollary}[Changing observation times versus one fixed time grid]
\label{p02:frontmix:cor:diffusion}
Let $B_t=tA+\Sigma^{1/2}W_t$, with standard $d$-dimensional
Brownian motion independent of an arbitrary random pair
$(A,\Sigma)$. The collection of joint laws
\[
 \big\{\mathcal L(B_h,B_{2h},\ldots,B_{dh}):h>0\big\}
\]
identifies the environment law of $(A,\Sigma)$. No collection
of observations at fewer than $d$ time points per experiment,
even allowing all possible positive times, can have this
property. On one fixed time grid, exactly $d+1$ time points
are necessary and sufficient.
\end{corollary}
\begin{proof}
Successive increments on a time grid are independent conditional
on $(A,\Sigma)$, with lengths equal to the time increments.
For a fixed total length $\tau$ and $q$ increments, the weighted
Gram lemma applies verbatim. Thus one fixed grid has an open
set of transform arguments exactly when $q\ge d+1$; the proof
of Theorem~\ref{p02:frontmix:thm:fixed-threshold} applies, using
$\mathcal D_\tau$ for the negative direction.

For the varying equal-step grid with $q=d$, fix $u\ne0$ and
$T\succ0$. Set
\[
 \tau=\tfrac12u^TT^{-1}u,\qquad h=\tau/d.
\]
Then $T-uu^T/(2\tau)$ is positive semidefinite of rank $d-1$,
so $H_u(T)$ is directly observed at this time grid. Continuity
includes $u=0$ and transform uniqueness gives the positive
claim. For $q<d$, every grid produces a covariance frequency
matrix of rank at most $q$, so the single covariance-determinant
perturbation used in Theorem~\ref{p02:frontmix:thm:rank-barrier}
is invisible on all the grids simultaneously.
\end{proof}

This separates recovery of a random matrix environment from
recovery of its one-time predictive distributions. In particular,
for $d\ge2$ all one-time laws of such random-drift diffusions can
coincide while their environment laws and full process laws differ.

\subsection{Why finite Gaussian mixtures behave differently}

\begin{proposition}[The classical finite-mixture contrast]
\label{p02:frontmix:prop:finite}
If both candidate mixing laws have finite support on distinct
Gaussian parameters, one observation already identifies them.
This is identification within the finite-support class, not among
all Borel competitors.
\end{proposition}
\begin{proof}
For the finitely many distinct parameter pairs choose a vector
$v$ for which the projected pairs
$(v^Ta_j,v^T\Sigma_jv)$ are distinct. Such a $v$ exists: for
each distinct pair of parameters, at least one of the linear
mean difference or quadratic covariance difference is a nonzero
polynomial, and finitely many of their common zero sets cannot
cover $\mathbb R^d$.

It remains to prove linear independence of the finitely many
functions $e^{im_jt-v_jt^2/2}$ with distinct $(m_j,v_j)$ and
$v_j\ge0$. In a zero linear combination, let $v_*$ be the
smallest variance and multiply by $e^{v_*t^2/2}$. The terms of
larger variance tend to zero as $t\to\infty$. The remaining
trigonometric polynomial $P(t)=\sum_{v_j=v_*}c_je^{im_jt}$
therefore tends to zero. Its mean square on $[0,R]$ tends, as
$R\to\infty$, to $\sum_{v_j=v_*}|c_j|^2$, because the means
in this group are distinct. Since $P$ tends to zero and is
bounded, this mean square tends to zero. All coefficients in
the smallest-variance group vanish. Iterating proves the claim.
\end{proof}

The proposition is classical; it is associated with the multivariate
Gaussian case of Yakowitz and Spragins\cite{p02:frontmix:Yakowitz1968}
and is rederived only to prevent a finite-versus-arbitrary-mixture
confusion. A single Gaussian can already have a nontrivial
continuous location-mixture representation if arbitrary competitors
are allowed: mix $N_d(a,\Sigma/2)$ over
$a\sim N_d(0,\Sigma/2)$ to obtain $N_d(0,\Sigma)$.

\subsection{Source status and limits of the priority check}

Vandermeulen and Scott's grouped-sample paper\cite{p02:frontmix:VS2015},
Theorem~4 on printed page~5 of the inspected 16-page preprint,
establishes a sharp universal $2m-1$ group-size bound for mixtures
with $m$ arbitrary component measures. That finite-component
statement neither asserts nor contradicts the present thresholds
for arbitrary mixing laws restricted to a Gaussian parameter
family. Haff et al.\cite{p02:frontmix:Haff2011}, Sections~2 and~3
(reprint pages~2--6), already analyze invertible Wishart
convolution on the positive definite cone. These are related
published results, not open problems claimed solved here.

Searches also located Barndorff-Nielsen's classical paper on
mixtures of exponential families\cite{p02:frontmix:BN1965}, but its
full text could not be retrieved in this audit. We therefore do
not claim that the fixed-group threshold is historically new or
absent as a corollary of that literature. Searches of the exact
rank and varying-terminal-length formulations did not locate
an exact prior statement. This bounded negative search is not
a proof of priority. The statements and counterexamples above
are proved here independently; historical novelty remains a
separate question.


% ===== Consolidated body: graph_replication =====
\section{Covariance replication thresholds are generic completion ranks}
\label{p02:graph:section}

The finite differential obstructions in the integrated manuscript lead
to a different question: can the exact number of observations needed
for an arbitrary Gaussian mixing law be read from an algebraic invariant
of its covariance constraints? The answer is affirmative. For a graph
constraint, the invariant is the classical generic completion rank,
not the maximum likelihood threshold. This is a transfer between two
observation problems; it does not solve the open combinatorial problem
of computing these ranks for every graph.

Let $V$ be a linear subspace of the real symmetric $d\times d$ matrices
containing a positive definite matrix. Equip it with the Frobenius inner
product and let $\pi_V$ be orthogonal projection. Put
$C=V\cap\mathbb S_+^d$ and $C^\circ=V\cap\mathbb S_{++}^d$.
For $r\ge0$ define the polynomial Gram map
\[
 \Gamma_r:\mathbb R^{d\times r}\longrightarrow V,
 \qquad \Gamma_r(Y)=\pi_V(YY^T),\qquad \Gamma_0=0,
\]
and let $r(V)$ be the least $r$ for which its image has nonempty
interior in $V$. Equivalently, its Jacobian is surjective at some point.
In particular $1\le r(V)\le d$.

Conditionally on $(A,\Sigma)$, the observations are iid
$N_d(A,\Sigma)$, and the mixing law is an arbitrary Borel probability
on $\mathbb R^d\times C^\circ$. A $q$-view law is the joint law of
$q$ such observations sharing $(A,\Sigma)$.

\begin{theorem}[Exact threshold for a linear covariance model]
\label{p02:graph:threshold}
The $q$-view law identifies every such mixing probability if and only if
\[
 q\ge r(V)+1.
\]
If the mean is fixed and known, the exact threshold is instead
$q\ge r(V)$.
Below either threshold, distinct compactly supported smooth mixing
densities with identical observations exist in every open ball whose
closure lies inside the parameter domain. They can be chosen strongly
log-concave on their common convex support and can agree on all
polynomial moments up to any prescribed finite degree.
The positive assertions require no moment or support assumptions.
\end{theorem}

\begin{lemma}[The Gram image and polynomial annihilators]
\label{p02:graph:gram-lemma}
If $r<r(V)$, a nonzero real homogeneous polynomial $P$ of positive
degree on $V$ vanishes on $\Gamma_r(\mathbb R^{d\times r})$.
If $r\ge r(V)$, that image contains a nonempty open subset of the
interior of the dual cone
$C^*=\{Z\in V:\langle Z,\Sigma\rangle\ge0\text{ for all }\Sigma\in C\}$.
\end{lemma}
\begin{proof}
The image is semialgebraic by the real projection theorem. A polynomial
map whose image has empty interior has image of dimension less than
$\dim V$; its Zariski closure is then proper, so a nonzero polynomial
vanishes on it. Equivalently one can use the characteristic-zero
Jacobian criterion for algebraic dependence of the coordinates of
$\Gamma_r$. Since the image is a cone, every homogeneous part of an
annihilating polynomial also annihilates it. Its constant part is zero,
giving the first assertion. These elementary real-algebraic facts are
used only for polynomial maps, where dimension and Jacobian rank agree.

Every Gram image lies in $C^*$, because
$\langle\pi_V(YY^T),\Sigma\rangle=\operatorname{tr}(Y^T\Sigma Y)\ge0$.
An open subset of this image is therefore in $\operatorname{int}C^*$.
Surjectivity at one point gives such an open subset by the submersion
theorem. The reverse implication follows either from Sard's theorem
or the polynomial dimension statement. Finally, $\Gamma_d$ has
surjective differential at every invertible $Y$, proving $r(V)\le d$.
\end{proof}

\begin{proof}[Proof of the identification direction]
Choose orthonormal coordinates on $V$. For the signed difference
$\rho$ of two candidate mixing laws and fixed $u\in\mathbb R^d$, put
\[
 H_u(Z)=\int e^{iu\cdot a-\langle Z,\Sigma\rangle}
             \,d\rho(a,\Sigma),\qquad
 \operatorname{Re}Z\in\operatorname{int}C^*.
\]
This is bounded and holomorphic on that tube. Indeed, on a compact
subset of its real base there is $c>0$ with
$\langle\operatorname{Re}Z,\Sigma\rangle\ge c\|\Sigma\|$ for
$\Sigma\in C$. All covariance derivatives are dominated by
polynomials times $e^{-c\|\Sigma\|}$. The mean factor has modulus one.

The joint characteristic function at $t_1,\ldots,t_q$ gives $H_u(Z)$
with
\[
 u=\sum_i t_i,\qquad
 Z=\frac12\pi_V\left(\sum_i t_it_i^T\right)
   =\frac1{2q}\pi_V(uu^T)+\frac12\Gamma_{q-1}(Y).
\]
Here $Y$ ranges over all $d\times(q-1)$ real matrices: write the
centered columns $t_i-u/q$ in an orthonormal basis of the sum-zero
subspace of $\mathbb R^q$.
For each fixed $u$, if $q-1\ge r(V)$, these probes contain a real
open subset of $\operatorname{int}C^*$. Translation by
$\pi_V(uu^T)/(2q)\in C^*$ preserves that interior.
Equality of the observations makes $H_u$ zero on that open subset.
The identity theorem, first in a complex polydisc and then on the
connected tube, makes $H_u$ identically zero.

Fix $Z_0\in\operatorname{int}C^*$. The Fourier transform of the finite
signed measure $e^{-\langle Z_0,\Sigma\rangle}\rho$ is
$H_u(Z_0-iW)$, zero for all $u\in\mathbb R^d$ and $W\in V$.
Fourier uniqueness proves $\rho=0$. With known mean, no sum constraint
is present and the covariance probes are $\tfrac12\Gamma_q(Y)$;
the same argument gives the threshold $r(V)$.
\end{proof}

\begin{proof}[Proof of the sharp counterexamples]
Suppose $q-1<r(V)$ and take $P$ from
Lemma~\ref{p02:graph:gram-lemma}. The real polynomial
\[
 Q(u,Z)=P\left(2Z-\frac1q\pi_V(uu^T)\right)
\]
is nonzero, vanishes on all $q$-view probes, is even in $u$, and
satisfies $Q(0,0)=0$. In orthonormal covariance coordinates $\theta$,
define the constant-coefficient operator
\[
 D=Q(i\partial_a,\partial_\theta).
\]
All its coefficients are real: every monomial has even total degree
in $u$. It is nonzero. For a nonzero real bump $\phi$ supported in
a smaller ball inside the requested parameter ball, set $h=D\phi$.
Integration by parts gives
\[
 \int e^{iu\cdot a-\langle Z,\Sigma\rangle}h(a,\Sigma)
                \,da\,d\Sigma
 =Q(u,Z)\int e^{iu\cdot a-\langle Z,\Sigma\rangle}
                 \phi(a,\Sigma)\,da\,d\Sigma.
\]
Thus $h$ has mass zero and annihilates the entire observation map.
Moreover $h\ne0$. A nonzero constant-coefficient operator cannot
annihilate a nonzero compactly supported smooth function: on taking
the ordinary Fourier transform, its nonzero polynomial symbol is
nonzero on a real open set, forcing an entire Fourier transform to
vanish there and hence everywhere if the result were zero.

To make all moments through degree $K$ vanish, replace $D$ by
$D\partial_{a_1}^{2M}$ with $2M>K$; its symbol remains nonzero and
its observation zeros persist. Integration by parts annihilates each
polynomial of degree at most $K$.

Choose a slightly larger ball of radius $R$ still inside the requested
parameter ball, with Euclidean coordinates $x$ about its center, and use the
normalized baseline
\[
 g_0(x)=c\exp\left(-\frac1{1-\|x\|^2/R^2}\right)
                  \mathbf1_{\{\|x\|<R\}}.
\]
It is globally smooth, positive in the ball, and its log-Hessian is
bounded above by $-2R^{-2}I$. For sufficiently small $\varepsilon>0$,
$g_0\pm\varepsilon h$ are distinct normalized positive densities
with the same support and observations. On the compact support of
$h$, smallness in $C^2$ preserves a strictly negative log-Hessian;
outside it the baseline is unchanged. The moment constraints persist.

For known mean use $P(2\partial_\theta)$ and, for moment matching,
an additional even high derivative in one covariance coordinate.
This proves the stated negative conclusions in both models.
\end{proof}

\subsection{An exact graph-theoretic interpretation}

For a finite simple graph $G$ on $d$ vertices, let
\[
 V_G=\{\Sigma\in\mathbb S^d:
       \Sigma_{ij}=0\text{ when }i\ne j\text{ and }ij\notin E(G)\}.
\]
This is a covariance-zero model; it is not the precision-zero
graphical model used to define the usual maximum likelihood threshold.
The same graph nevertheless indexes the same matrix-entry projection.
Generic completion rank is classically defined as the least rank at
which this projection of the symmetric determinantal variety is
dominant; see \cite[Definition 1.1]{p02:graph:BS2017}.

\begin{corollary}[Replication equals generic completion rank, plus one]
\label{p02:graph:gcr}
For arbitrary Gaussian mixtures with covariance in $V_G\cap\mathbb S_{++}^d$,
the exact joint replication thresholds are
\[
 q_{\rm unknown\ mean}(G)=\operatorname{gcr}(G)+1,
 \qquad q_{\rm known\ mean}(G)=\operatorname{gcr}(G).
\]
In particular, with unknown mean the thresholds are $2$ for an edgeless
graph, $3$ for a forest containing an edge, $4$ for a cycle, and
$\omega(G)+1$ for a chordal graph. For the complete graph they are
$d+1$ and $d$, respectively.
\end{corollary}
\begin{proof}
Over the complex numbers every symmetric matrix of rank at most $r$
can be written $YY^T$. Dominance of the projection is therefore
equivalent to a nonzero maximal Jacobian minor of $\Gamma_r$.
Its coefficients are real, so a nonzero minor is also nonzero at some
real $Y$. Thus the classical generic completion rank equals $r(V_G)$.
Apply Theorem~\ref{p02:graph:threshold}.

Here are elementary checks of the displayed graph classes. A clique
of size $k$ forces $r(V_G)\ge k$, because its specified principal
block must have rank at most $r$ and cannot contain an open set if
$r<k$. If $G$ has an elimination ordering with at most $k$ later
neighbors per vertex, take $r=k+1$ and choose row vectors $y_i$ such
that every set of at most $r$ of them is linearly independent.
A linear dependence among the differentials of the specified Gram
entries is a stress with vertex equations
\[
 2a_i y_i+\sum_{j:ij\in E(G)}b_{ij}y_j=0.
\]
At the first vertex in the ordering, independence of its vector and
its remaining neighbors forces every incident coefficient to zero.
Remove it and continue. All stress coefficients vanish, so the
Jacobian is surjective. Consequently
$r(V_G)\le\operatorname{degeneracy}(G)+1$.
This proves the forest claim and, by a perfect elimination ordering,
the chordal claim. An edgeless graph has independent positive diagonal
entries already at rank one. For a cycle, degeneracy gives $r\le3$;
at rank two the symmetric rank-two variety on its $d\ge3$ vertices
has dimension $2d-1$, whereas $d$ diagonals and $d$ edges give $2d$
specified entries. Hence $r\ge3$. The complete graph follows from
the clique bound.
\end{proof}

The graph invariant is classical; the asserted transfer concerns the
identification of an entire probability law on mean and covariance
space. Generic matrix completion, finite Gaussian-mixture
identifiability, and arbitrary mixing-law identification have different
quantifiers. In particular the counterexamples above are smooth
continuous mixtures, not finite collections of Gaussian components.

\subsection{A second transfer: likelihood existence does not measure recovery}

Blekherman--Sinn prove that $\operatorname{mlt}(K_{5,5})=4$ whereas
$\operatorname{gcr}(K_{5,5})=5$ and, more generally,
$\operatorname{gcr}(K_{m,m})=m$ for $m\ge3$ while
$\operatorname{mlt}(K_{m,m})$ has order $\sqrt m$
\cite[Corollary 2.11 and Theorems 2.1, 2.7]{p02:graph:BS2017}.
The preceding corollary therefore gives an exact threshold of six
joint replicates in the unknown-mean $K_{5,5}$ covariance mixture,
and a threshold $m+1$ for $K_{m,m}$.
The gap from the likelihood-existence threshold is unbounded.
This is a transfer of their established graph computations, not a
new computation of those graph invariants.

The conjectures in \cite[Section 2]{p02:graph:BDG2024} concerning equality
of MLT and GCR for sparse random graphs remain combinatorial questions.
Theorem~\ref{p02:graph:threshold} gives a new observation-theoretic
interpretation of GCR but supplies no proof that it equals MLT on
those graphs. Similarly, the general question in
\cite[p.~4]{p02:graph:BS2017} asking whether bounded MLT forces bounded
GCR is not settled here. These are explicit external targets checked
against the transfer; the missing step is a bound on the Gram-map
rank, not an additional transform uniqueness argument.

\subsection{Quantitative consequence of a fixed replication obstruction}

\begin{corollary}[A sample-size-independent lower bound below the threshold]
\label{p02:graph:risk-floor}
Fix any of the subthreshold models in Theorem~\ref{p02:graph:threshold},
and any finite number of prescribed moment constraints. There is a
compact class of smooth, strongly log-concave mixing densities inside
the parameter domain and a constant $c>0$ such that the minimax
expected $W_1$ error is at least $c$, regardless of the number of
independent groups, provided every group has at most $q$ observations.
The conclusion also allows randomized adaptive choice of group sizes
up to this same cap.
\end{corollary}
\begin{proof}
Take the two densities supplied by the theorem. Their $q$-view laws
agree, hence so do all their lower-view marginals. Each adaptive
group experiment can be simulated from an independent full $q$-tuple
for each fresh individual by revealing coordinates only when requested.
The transcript laws therefore agree for the two models. The triangle
inequality gives a worst-case risk at least half the positive $W_1$
distance between their mixing laws. The two-point class itself is compact.
\end{proof}

This lower bound concerns a fixed cap on within-individual replication.
It is compatible with the polynomial total-budget rates obtained later
by allowing the replication depths to grow.


% ===== Consolidated body: budgets =====
\section{Dimension and component regularity in optimal replication budgets}
\label{p02:bud:sec:main}

The scalar allocation theorem extends to an arbitrary latent dimension,
but its distribution-function proof does not. A smoothed Fourier
construction provides the missing replacement. It gives the exact
total-budget rate $B^{-1/(d+2)}$ for regular $d$-dimensional components,
including $d=2$ without a logarithmic loss. A directed relative-entropy
argument also covers support-dependent components: mixtures of
$d$-dimensional uniform rectangles have rate $B^{-1/(d+1)}$ for $d\geq2$.
Both statements allow every adaptive randomized replication policy in
their lower bounds and use predetermined unequal group sizes in their
upper bounds.

Multilevel telescoping, negative Sobolev bounds for transport, and
population learning are established methods; see
\cite{p01:info:GNR2015,p01:info:Giles2018,p02:bud:WeedBerthet2019,p01:info:LeeValiant2021}.
The results below are complete deductions for the stated grouped
experiments. Historical priority for this general combination is not
asserted. The loss is Wasserstein distance between arbitrary mixing
laws, including singular laws. No density, separation, or finite-support
assumption on the mixing law is used.

\subsection{A general rate theorem}

Let $\Theta=[0,1]^d$ with its Euclidean metric, and let
$(P_\theta)_{\theta\in\Theta}$ be a known measurable family of
probability laws on a standard Borel observation space. A new individual
draws a latent $\theta$ independently from an unknown
$G\in\mathcal P(\Theta)$. Repeated observations of that individual
are independent with law $P_\theta$, conditional on its latent parameter.
A policy can start individuals, revisit them, and randomize using the
entire observed history. At most $B$ observations are allowed, with
$B$ deterministic. Write $\mathcal R_B^{\rm ad}$ for the minimax
expected $W_1$ risk over all policies and probability-valued estimators,
and $\mathcal R_B^{\rm fix}$ for the restriction to predetermined,
possibly unequal group sizes.

We impose the following component-estimation assumption: for an integer
$b_0\geq1$ and constants $A,\alpha>0$, there are measurable estimators
$T_r\in\Theta$, using $b_0r$ observations from one individual, such that
\begin{equation}\label{p02:bud:eq:rms}
 \sup_{\theta\in\Theta}
 \mathbb E_\theta\|T_r-\theta\|^2\leq A^2r^{-2\alpha},
 \qquad r\in\mathbb N.
\end{equation}
Independent auxiliary randomization is allowed.

\begin{theorem}[Dimension-dependent multilevel upper bound]
\label{p02:bud:thm:upper}
Under \eqref{p02:bud:eq:rms}, for all sufficiently large $B$,
\begin{equation}\label{p02:bud:eq:upper-general}
 \mathcal R_B^{\rm ad}\leq\mathcal R_B^{\rm fix}\leq C
 \begin{cases}
 B^{-\alpha/(1+d\alpha)},&d\geq2,\\
 B^{-\alpha/(1+\alpha)},&d=1,\ 0<\alpha<1,\\
 B^{-1/2}\log(eB),&d=1,\ \alpha=1,\\
 B^{-1/2},&d=1,\ \alpha>1.
 \end{cases}
\end{equation}
The constants depend only on $d,\alpha,A,b_0$. In particular,
root-mean-square component estimation, $\alpha=1/2$, gives
$CB^{-1/(d+2)}$ in every dimension.
\end{theorem}

For the lower bound, a one-sided condition suffices. Suppose some cube
$J=a+[0,s]^d\Subset\Theta$ and constants $D,\kappa>0$ satisfy
\begin{equation}\label{p02:bud:eq:directed-kl}
 \operatorname{KL}(P_\theta\Vert P_{\theta+t e_1})
 \leq Dt^\kappa
 \quad\text{whenever }t>0\text{ and }\theta,\theta+te_1\in J.
\end{equation}
The reverse relative entropy can be infinite.

\begin{theorem}[Adaptive lower bound from directed local information]
\label{p02:bud:thm:lower}
Under \eqref{p02:bud:eq:directed-kl},
\begin{equation}\label{p02:bud:eq:lower-general}
 \mathcal R_B^{\rm ad}\geq
 c\max\{B^{-1/2},B^{-1/(d+\kappa)}\}.
\end{equation}
Consequently, if \eqref{p02:bud:eq:rms} holds with
$\alpha=1/\kappa$, the following rates are exact for both adaptive
and fixed unequal allocations:
\begin{equation}\label{p02:bud:eq:phase}
 \mathcal R_B^{\rm ad}\asymp\mathcal R_B^{\rm fix}\asymp
 \begin{cases}
 B^{-1/(d+\kappa)},&d+\kappa>2,\\
 B^{-1/2},&d+\kappa<2.
 \end{cases}
\end{equation}
At $d+\kappa=2$, necessarily $d=1,\kappa=1$, the proved bounds are
$cB^{-1/2}$ and $CB^{-1/2}\log(eB)$. No sharp claim is made at that
critical point.
\end{theorem}

The reciprocal relation $\alpha=1/\kappa$ in the last assertion is an
assumption tying component estimation to local information. Without it,
the separately stated upper and lower bounds still hold, but need not
match.

\subsection{A Fourier transport bound and the critical dimension}

Embed $[0,1]^d$ in the flat torus
$\mathbb T_d=\mathbb R^d/(4\mathbb Z)^d$. The torus distance agrees
with Euclidean distance between any two points of this cube. For a
mass-zero finite signed measure $\rho$ on the torus, define
\[
 \|\rho\|_{\rm KR}=\sup_{\operatorname{Lip}(f)\leq1}
          \left|\int f\,d\rho\right|.
\]
Constants in $f$ do not affect this expression. For probability measures
it is the torus $W_1$ distance. In particular it is the original
Euclidean $W_1$ when both laws are supported on $\Theta$.
For $k\in\mathbb Z^d$, put
\[
 \widehat\rho(k)=\int e^{-\pi i k\cdot x/2}\,d\rho(x),
 \qquad
 \|\rho\|_{-1}^2=
      \sum_{k\ne0}\frac{|\widehat\rho(k)|^2}{\|k\|^2}.
\]

\begin{lemma}[Smoothed point masses in a transport-controlling Hilbert norm]
\label{p02:bud:lem:fourier}
There is $C_d<\infty$ such that every smooth mass-zero signed density
satisfies
\begin{equation}\label{p02:bud:eq:kr-hminus}
 \|\rho\|_{\rm KR}\leq C_d\|\rho\|_{-1}.
\end{equation}
Let $K_h$ be the law on $\mathbb T_d$ of $hZ$ modulo $4\mathbb Z^d$,
where $Z\sim N(0,I_d)$ and $0<h\leq1$. Then
\[
 W_1(K_h*\mu,\mu)\leq\sqrt d\,h
 \quad\text{for every probability law }\mu.
\]
For any coupled $X,Y$ with $\mathbb E\|X-Y\|^2\leq a^2$, where
the norm uses representatives in $\Theta$,
\begin{equation}\label{p02:bud:eq:spectral-variance}
 \mathbb E\|K_h*(\delta_X-\delta_Y)\|_{-1}^2
 \leq C_d\sum_{k\ne0}
       e^{-c h^2\|k\|^2}\|k\|^{-2}
       \min\{1,a^2\|k\|^2\}.
\end{equation}
\end{lemma}
\begin{proof}
For a smooth $1$-Lipschitz periodic function $f$, Parseval's identity
and its gradient bound imply
$\sum_{k\ne0}\|k\|^2|\widehat f(k)|^2\leq C_d$, with Fourier
coefficients of functions normalized by torus volume. Fourier duality
and Cauchy--Schwarz prove \eqref{p02:bud:eq:kr-hminus}. A general Lipschitz
function is uniformly approximated by its convolution with smooth
probability kernels, without increasing its Lipschitz constant, so the
same conclusion holds for the supremum defining $\|\cdot\|_{\rm KR}$.

Couple $X\sim\mu$ to $X+hZ$ modulo the torus. Their torus distance
is at most $h\|Z\|$, which proves the smoothing bound. The Fourier
coefficient of $K_h$ is $\exp(-\pi^2h^2\|k\|^2/8)$.
Also
\[
 |e^{-\pi i k\cdot X/2}-e^{-\pi i k\cdot Y/2}|^2
 \leq\min\{4,\tfrac{\pi^2}{4}\|k\|^2\|X-Y\|^2\}.
\]
Taking expectations and using Tonelli's theorem proves
\eqref{p02:bud:eq:spectral-variance}.
\end{proof}

\begin{lemma}[The allocation sum]
\label{p02:bud:lem:allocation-sum}
Set $r_\ell=2^\ell$, $R=2^L$ and $h=R^{-\alpha}$. Consider a
same-individual coupling of $T_{r_\ell}$ and $T_{r_{\ell-1}}$,
using a prefix for the coarser statistic. The smoothed signed differences
have second moments in $\|\cdot\|_{-1}$ bounded by $Cv_\ell$, where
one may choose
\begin{equation}\label{p02:bud:eq:vlevel}
 v_\ell=
 \begin{cases}
 r_\ell^{-\alpha},&d=1,\\
 1+\log(R/r_\ell),&d=2,\\
 R^{\alpha(d-2)},&d\geq3.
 \end{cases}
\end{equation}
At level zero the centered random measure
$K_h*(\delta_{T_1}-\mathcal L(T_1))$ obeys the same bound.
Writing $S_L=\sum_{\ell=0}^L\sqrt{r_\ell v_\ell}$ gives
\begin{equation}\label{p02:bud:eq:Slevel}
 S_L\leq C
 \begin{cases}
 R^{(1-\alpha)/2},&d=1,\ \alpha<1,\\
 1+\log R,&d=1,\ \alpha=1,\\
 1,&d=1,\ \alpha>1,\\
 R^{[1+\alpha(d-2)]/2},&d\geq2.
 \end{cases}
\end{equation}
All bounds are uniform in $G$.
\end{lemma}
\begin{proof}
The coupling shares its latent parameter, and \eqref{p02:bud:eq:rms} gives
\[
 \mathbb E_G\|T_{r_\ell}-T_{r_{\ell-1}}\|^2
 \leq2A^2(1+2^{2\alpha})r_\ell^{-2\alpha}.
\]
Independence between these two statistics is neither required nor used.
The number of lattice points with $n\leq\|k\|<n+1$ is at most
$C_dn^{d-1}$ for $n\geq1$. Lemma~\ref{p02:bud:lem:fourier} reduces the
second-moment estimate, after changing constants, to
\begin{equation}\label{p02:bud:eq:shell-sum}
 \sum_{n\geq1}n^{d-3}e^{-c h^2n^2}
       \min\{1,n^2r_\ell^{-2\alpha}\}.
\end{equation}
In dimension one, split at $n=r_\ell^\alpha$: the initial sum is
at most $Cr_\ell^{-\alpha}$ and the tail of $\sum n^{-2}$ has
the same bound. In dimension two, split at $r_\ell^\alpha$ and
$h^{-1}=R^\alpha$. The first and last pieces are bounded, and the
middle piece is at most
$C\log(R^\alpha/r_\ell^\alpha)$.
For $d\geq3$ drop the minimum and bound the remaining Gaussian sum
by $Ch^{2-d}=CR^{\alpha(d-2)}$. For level zero, use
$\mathbb E|\widehat{\delta_{T_1}}(k)
-\mathbb E\widehat{\delta_{T_1}}(k)|^2\leq1$ and the same sums;
in dimension one their bound is constant.

The one-dimensional assertions about $S_L$ are geometric sums.
For $d=2$, write $j=L-\ell$ to obtain
\[
 S_L=\sqrt R\sum_{j=0}^L2^{-j/2}\sqrt{1+j\log2}
 \leq C\sqrt R.
\]
For $d\geq3$, factor out $R^{\alpha(d-2)/2}$ and sum
$\sum_{\ell=0}^L2^{\ell/2}\leq C\sqrt R$.
This proves \eqref{p02:bud:eq:Slevel}. In particular the logarithm in the
two-dimensional variance does not survive the allocation sum.
\end{proof}

\subsection{Construction, integer costs, and probability-valued output}

\begin{proof}[Proof of Theorem~\ref{p02:bud:thm:upper}]
Use $N_\ell$ fresh independent individuals at each level, disjoint
across levels. Each level-$\ell$ individual supplies $b_0r_\ell$
observations. For $\ell\geq1$, the two statistics are the fine
statistic and its coarse prefix version. Define the smooth signed measure
of total mass one
\begin{align}
 H_L={}&\frac1{N_0}\sum_{i=1}^{N_0}K_h*\delta_{T_{1,i}}\notag\\
 &+\sum_{\ell=1}^L\frac1{N_\ell}\sum_{i=1}^{N_\ell}
       K_h*(\delta_{T_{r_\ell,i}}
                  -\delta_{T_{r_{\ell-1},i}}).
 \label{p02:bud:eq:multilevel-measure}
\end{align}
Let $G_R$ be the law of $T_R$ when its latent parameter has law $G$.
Correct marginal laws of every prefix give $\mathbb EH_L=K_h*G_R$.
For a mean-zero random vector in a Hilbert space, second moments of
independent sums add. Lemma~\ref{p02:bud:lem:allocation-sum} therefore gives
\[
 \mathbb E\|H_L-\mathbb EH_L\|_{-1}^2
 \leq C\sum_{\ell=0}^L\frac{v_\ell}{N_\ell}.
\]
This identity can alternatively be checked for every finite Fourier
truncation and passed to the limit by monotone convergence, avoiding any
integrability ambiguity. Coupling $T_R$ with its latent parameter and
applying Lemma~\ref{p02:bud:lem:fourier} gives
\begin{equation}\label{p02:bud:eq:raw-risk}
 \mathbb E\|H_L-G\|_{\rm KR}
 \leq C\left(\sum_{\ell=0}^L\frac{v_\ell}{N_\ell}\right)^{1/2}
       +AR^{-\alpha}+\sqrt d\,h.
\end{equation}

Put $b=\lfloor B/b_0\rfloor$. For $R=2^L\leq b/4$, set
\begin{equation}\label{p02:bud:eq:integer-allocation}
 N_\ell=\left\lceil\frac b{2S_L}
                       \sqrt{\frac{v_\ell}{r_\ell}}\right\rceil.
\end{equation}
The deterministic number of blocks is bounded by
\[
 \sum_{\ell=0}^Lr_\ell N_\ell
 \leq b/2+\sum_{\ell=0}^Lr_\ell<b/2+2R\leq b,
\]
and
$\sum_\ell v_\ell/N_\ell\leq2S_L^2/b$.
Thus \eqref{p02:bud:eq:raw-risk} is at most
$C(S_L/\sqrt b+R^{-\alpha})$.
Choose $R$ to be the largest dyadic integer at most the corresponding
quantity below:
\[
 \frac1{16}
 \begin{cases}
 b^{1/(1+d\alpha)},&d\geq2,\\
 b^{1/(1+\alpha)},&d=1,\ \alpha<1,\\
 b^{1/2},&d=1,\ \alpha=1,\\
 b^{1/(2\alpha)},&d=1,\ \alpha>1.
 \end{cases}
\]
For all sufficiently large $b$ these quantities exceed one and satisfy
$R\leq b/4$. Equation~\eqref{p02:bud:eq:Slevel} proves the rates in
\eqref{p02:bud:eq:upper-general} for the signed output.

To obtain a probability-valued estimator, enumerate a countable
$W_1$-dense collection $(Q_j)$ of rational atomic probabilities on
$\Theta$. Choose the first $j$ satisfying
\[
 \|Q_j-H_L\|_{\rm KR}
 \leq\inf_k\|Q_k-H_L\|_{\rm KR}+B^{-1}.
\]
This is measurable. Indeed, after fixing $f(0)=0$, the Lipschitz unit
ball on the compact torus has a countable dense subset in the uniform
norm; its supremum represents each displayed norm. The objective and
its countable infimum are therefore measurable. Density implies the
infimum is no larger than $\|G-H_L\|_{\rm KR}$. For
$\widehat G=Q_j$, the triangle inequality gives
\[
 W_1(\widehat G,G)=\|\widehat G-G\|_{\rm KR}
 \leq2\|H_L-G\|_{\rm KR}+B^{-1}.
\]
This proves the theorem. No assertion of computational optimality is
needed for this statistical minimax construction.
\end{proof}

\subsection{An adaptive lower bound with only directed finite entropy}

\begin{proof}[Proof of Theorem~\ref{p02:bud:thm:lower}]
Partition $J=a+[0,s]^d$ into $K=m^d$ congruent cubes and write
$a_1,\ldots,a_K$ for their centers. Set $h=s/(8m)$ and
\[
 G_v=K^{-1}\sum_{j=1}^K\delta_{a_j+h v_j e_1},
 \qquad v\in\{-1,1\}^K.
\]
All these points are in $J$. Distances between points from different
cells are at least $6h$. Within each cell the two possible points are
$2h$ apart. Consequently
\begin{equation}\label{p02:bud:eq:cube-transport}
 W_1(G_v,G_w)=\frac{2h}{K}d_H(v,w).
\end{equation}
The upper bound matches atoms within cells. For the lower bound assign
the values $h,-h$ to the two points in every differing cell, with $h$
assigned to the $G_v$ point; assign zero in unchanged cells. These
assignments are $1$-Lipschitz on the finite set, because cross-cell
distances are at least $6h$. Their McShane extension is a valid
Kantorovich test function and proves equality.

Give the statistician a stronger oracle: the cells, centers, and equal
weights are known, and a query to any chosen cell $j$ returns a fresh
observation from $P_{a_j+hv_je_1}$. Signs remain unknown. An ordinary
grouped policy can be simulated exactly: every time it starts an
individual, independently draw a uniform cell index; query that cell on
every revisit, and conceal the index from the original policy. Thus any
lower bound for this oracle is a lower bound for the original experiment.

Fix an adaptive randomized oracle policy, and include its parameter-free
random seed in the transcript. Write $\mathbb P_v$ for its transcript
law and $N_j$ for its number of queries to cell $j$. Only use oriented
edges with $v_j=-1$ and let $v^{(j)}$ flip that sign to $+1$.
The chain rule for relative entropy and \eqref{p02:bud:eq:directed-kl} give
\[
 \operatorname{KL}(\mathbb P_v\Vert\mathbb P_{v^{(j)}})
 \leq D(2h)^\kappa\mathbb E_vN_j.
\]
At every fixed past transcript the policy's conditional choice is the
same in both models; only queries to $j$ contribute entropy. An early
stop can be padded with common dummy observations.
Since $\sum_jN_j\leq B$,
\[
 \frac1{K2^{K-1}}\sum_{j=1}^K\sum_{v:v_j=-1}
 \operatorname{KL}(\mathbb P_v\Vert\mathbb P_{v^{(j)}})
 \leq\frac{2D(2h)^\kappa B}{K}.
\]
Pinsker and Jensen therefore bound the average total variation over
these edges by
\[
 \left(\frac{D(2h)^\kappa B}{K}\right)^{1/2}
 =\left(\frac{D(s/4)^\kappa B}{m^{d+\kappa}}\right)^{1/2}.
\]
Choose
$m=\max\{1,\lceil(16D(s/4)^\kappa B)^{1/(d+\kappa)}\rceil\}$.
The average is at most $1/4$. Binary testing on each edge then gives,
for every sign estimator,
\[
 2^{-K}\sum_v\mathbb E_vd_H(\widehat v,v)\geq3K/8.
\]
For an arbitrary probability-valued estimate $\widehat G$, choose a
nearest law $G_{\widehat v}$ in the finite family. The triangle
inequality and \eqref{p02:bud:eq:cube-transport} imply
\[
 W_1(\widehat G,G_v)\geq\tfrac12W_1(G_{\widehat v},G_v)
       =\frac hK d_H(\widehat v,v).
\]
Its maximum risk is at least $3h/8\geq cB^{-1/(d+\kappa)}$.
The argument uses finite entropy only in the displayed orientation.

For the $B^{-1/2}$ term, an oracle that reveals $B$ independent latent
parameters can simulate any allowed grouped experiment: use the next
latent draw whenever a previously unobserved individual is first
queried, and generate its conditional observations as required.
At most $B$ individuals can be queried. For distinct fixed
$u,w\in\Theta$, consider
$G_\pm=(1/2\pm\varepsilon)\delta_u
 +(1/2\mp\varepsilon)\delta_w$ with
$\varepsilon=1/(16\sqrt B)$. Their $B$-sample relative entropy is
bounded by an absolute constant below $1/2$, and their $W_1$ separation
is $2\varepsilon\|u-w\|$. The two-point testing inequality proves
the remaining term. Combining with Theorem~\ref{p02:bud:thm:upper} yields
the stated exact rates and critical bounds.
\end{proof}

\subsection{Regular observation channels: three concrete applications}

\begin{corollary}[Gaussian random effects and joint Bernoulli traits]
\label{p02:bud:cor:regular-channels}
For each of the following experiments, with arbitrary latent mixing law
on $[0,1]^d$,
\[
 \mathcal R_B^{\rm ad}\asymp\mathcal R_B^{\rm fix}
       \asymp B^{-1/(d+2)}:
\]
\begin{enumerate}
\item $P_\theta=N_d(\theta,\Sigma)$ with a known positive-definite
covariance matrix $\Sigma$;
\item $P_\theta=\bigotimes_{j=1}^d\operatorname{Bernoulli}(\theta_j)$.
\end{enumerate}
An observation is one $d$-vector. If individual coordinate measurements
are charged separately, the fixed factor $d$ changes constants, not
the exponent. No independence of the latent coordinates under $G$ is
assumed.
\end{corollary}
\begin{proof}
For the Gaussian experiment, project the sample mean onto $[0,1]^d$.
Euclidean projection cannot increase its distance to $\theta$, and
its mean squared error is at most $\operatorname{tr}(\Sigma)/r$.
The Gaussian relative entropy is
$\frac12(\theta-\theta')^T\Sigma^{-1}(\theta-\theta')$.
Thus the two main theorems apply with $\alpha=1/2,\kappa=2$.

For the Bernoulli experiment the vector of coordinate sample means has
mean squared error at most $d/(4r)$. On $J=[1/4,3/4]^d$,
relative entropy is bounded above by chi-squared divergence coordinate
by coordinate, giving
\[
 \operatorname{KL}(P_\theta\Vert P_{\theta'})
 \leq\sum_{j=1}^d
       \frac{(\theta_j-\theta'_j)^2}
            {\theta'_j(1-\theta'_j)}
 \leq\frac{16}3\|\theta-\theta'\|^2.
\]
The same theorems apply.
\end{proof}

\begin{corollary}[The tensor factorial kernel]
\label{p02:bud:cor:factorial}
Let
\[
 k_*(x\mid t)=1+\sum_{j\geq2}\frac{L_j(x)L_j(t)}{j!},
 \qquad
 P_\theta(dx)=\prod_{a=1}^d k_*(x_a\mid\theta_a)\,dx,
\]
where $L_j(x)=P_j(2x-1)$ are shifted Legendre polynomials.
Then arbitrary joint mixing laws on $[0,1]^d$ have exact adaptive
and fixed unequal budget risk $B^{-1/(d+2)}$.
\end{corollary}
\begin{proof}
For $A_j(x)=j!(2j+1)L_j(x)$, define the scalar pair statistic
\[
 Y(x,y)=\frac12+\frac{35}{18}
   \left(A_2(x)A_3(y)-\frac{10}{21}A_5(x)-\frac4{15}A_3(x)\right).
\]
Legendre orthogonality gives $\mathbb E_tA_j(X)=L_j(t)$, and
\[
 L_2L_3=\frac{10}{21}L_5+\frac4{15}L_3+\frac9{35}L_1
\]
implies $\mathbb E_tY=t$. Since $|L_j|\leq1$ on $[0,1]$,
$|Y|<2100$. Apply this statistic coordinatewise to $r$ disjoint
pairs and project its average onto the cube. Its mean squared error
is at most $d\,2100^2/r$, so \eqref{p02:bud:eq:rms} holds with $b_0=2$.

The scalar density is at least $3-e>0$, and its parameter derivative
is bounded on the compact square by uniform convergence of the
factorial series. If this derivative bound is $M$, then
\[
 \operatorname{KL}(P_t^*\Vert P_{t'}^*)
 \leq\chi^2(P_t^*\Vert P_{t'}^*)
 \leq M^2(3-e)^{-1}|t-t'|^2.
\]
Relative entropy adds over coordinates. Hence $\kappa=2$, and the
claimed rate follows.
\end{proof}

\subsection{Unknown Gaussian covariance: statistical dimension versus observability}

Let $V$ be a $p$-dimensional linear subspace of the symmetric
$m\times m$ matrices, with fixed basis $S_1,\ldots,S_p$.
Consider an injective affine parametrization of a full-dimensional
compact box of mean and covariance parameters,
\[
 a=a(u),\qquad
 \Sigma=\Sigma_0+\sum_{j=1}^p v_jS_j,
 \qquad (u,v)\in[0,1]^{m+p},
\]
where the mean map is nonsingular and every covariance satisfies
$\lambda_0I\preceq\Sigma\preceq\lambda_1I$ for fixed
$0<\lambda_0\leq\lambda_1<\infty$. Rescaling box coordinates is
absorbed into the basis and mean map.

\begin{corollary}[Budget law for structured Gaussian mixing]
\label{p02:bud:cor:structured-gaussian}
For $P_{u,v}=N_m(a(u),\Sigma(v))$ and arbitrary mixing laws on the
parameter box, put $k=m+p$. Then
\begin{equation}\label{p02:bud:eq:gaussian-budget}
 \mathcal R_B^{\rm ad}\asymp\mathcal R_B^{\rm fix}
       \asymp B^{-1/(k+2)}.
\end{equation}
The loss may use either Euclidean box coordinates or the parameter
metric $[\|a-a'\|^2+\|\Sigma-\Sigma'\|_F^2]^{1/2}$, since these
fixed metrics are equivalent.

For a covariance graph $H$ with $m$ vertices and free diagonal and
edge entries, $p=m+|E(H)|$ and the exponent is
$1/(2m+|E(H)|+2)$. This depends on the number of covariance parameters,
whereas the exact finite-replication observability threshold depends on
the rank geometry of the projected Gram map in
Theorem~\ref{p02:graph:threshold}. Here the graph constrains covariance
entries, not precision-matrix entries. The increasing replication depths
used in \eqref{p02:bud:eq:gaussian-budget} also explain why this result is
consistent with the risk floor under a fixed cap below the observability
threshold.
\end{corollary}
\begin{proof}
Use $n=2r$ observations from one individual. The sample mean
$\overline X$ and unbiased sample covariance $\widehat\Sigma$ satisfy
\[
 \mathbb E\|\overline X-a\|^2=\frac{\operatorname{tr}\Sigma}{n},
 \qquad
 \mathbb E\|\widehat\Sigma-\Sigma\|_F^2
 =\frac{(\operatorname{tr}\Sigma)^2+
              \operatorname{tr}(\Sigma^2)}{n-1}.
\]
For the second identity, an orthogonal change of the sample indices
expresses $(n-1)\widehat\Sigma$ as the sum of $n-1$ independent
matrices $ZZ^T$ with $Z\sim N_m(0,\Sigma)$.
The Gaussian fourth-moment identity gives
$\operatorname{Var}(Z_iZ_j)=\Sigma_{ii}\Sigma_{jj}+\Sigma_{ij}^2$;
summing over $i,j$ proves the formula. Thus both errors are uniformly
$O(r^{-1})$.

Project $\widehat\Sigma-\Sigma_0$ orthogonally onto $V$, recover its
basis coefficients, apply the inverse affine map to the sample mean,
and project the resulting parameter vector onto $[0,1]^k$.
All intervening linear maps have fixed finite norms, so this gives
\eqref{p02:bud:eq:rms} with $d=k,\alpha=1/2,b_0=2$.

The Gaussian entropy formula is
\[
 \operatorname{KL}(N(a,\Sigma)\Vert N(a',\Sigma'))
 =\frac12\left\{
 \operatorname{tr}(\Sigma'^{-1}\Sigma)-m
 -\log\det(\Sigma'^{-1}\Sigma)
 +(a-a')^T\Sigma'^{-1}(a-a')\right\}.
\]
The eigenvalues of $\Sigma'^{-1/2}\Sigma\Sigma'^{-1/2}$ lie in a
fixed compact subinterval of $(0,\infty)$. On that interval,
$x-1-\log x\leq C(x-1)^2$. The entropy is therefore at most
$C(\|a-a'\|^2+\|\Sigma-\Sigma'\|_F^2)$, hence at most a
constant times squared parameter distance. Apply the two main theorems
with $\kappa=2$. The graph dimension count is immediate.
\end{proof}

\subsection{A nonregular application: populations of unknown endpoints}

\begin{corollary}[Exact budget for uniform-rectangle mixtures]
\label{p02:bud:cor:uniform-endpoints}
For an endpoint vector $\theta\in[1,2]^d$, let the component law be
\[
 P_\theta=\bigotimes_{j=1}^d\operatorname{Uniform}[0,\theta_j].
\]
For every $d\geq2$ and arbitrary $G\in\mathcal P([1,2]^d)$,
\begin{equation}\label{p02:bud:eq:uniform-rate}
 \mathcal R_B^{\rm ad}\asymp\mathcal R_B^{\rm fix}
       \asymp B^{-1/(d+1)}.
\end{equation}
In dimension one the present argument gives only
$cB^{-1/2}\leq\mathcal R_B^{\rm ad}
\leq\mathcal R_B^{\rm fix}\leq CB^{-1/2}\log(eB)$.
\end{corollary}
\begin{proof}
Let $M_{r,j}$ be the maximum of the $r$ observations of coordinate $j$.
Its normalized distribution satisfies
$\mathbb P_\theta(M_{r,j}/\theta_j\leq t)=t^r$ for $0\leq t\leq1$.
Direct integration gives
\[
 \mathbb E_\theta\left(\frac{r+1}{r}M_{r,j}\right)=\theta_j,
 \qquad
 \operatorname{Var}_\theta\left(\frac{r+1}{r}M_{r,j}\right)
 =\frac{\theta_j^2}{r(r+2)}.
\]
Clip these coordinate estimates to $[1,2]$ and subtract one to work on
the unit cube. Their mean squared error is at most $4d/r^2$.
Thus $\alpha=1$, with $b_0=1$.

If only the first endpoint increases by $t>0$, then
\[
 \operatorname{KL}(P_\theta\Vert P_{\theta+te_1})
 =\log\left(1+\frac t{\theta_1}\right)\leq t.
\]
This is a directed bound with $\kappa=1$. In the opposite direction
the relative entropy is infinite, because the larger rectangle has
mass outside the smaller rectangle. The orientation in
Theorem~\ref{p02:bud:thm:lower} is therefore necessary for this application.
Since $d+1>2$ precisely when $d\geq2$, the theorem proves the exact
rate in those dimensions and the stated bounds when $d=1$.
\end{proof}

The mixture density for a single observation is
\[
 p_G(x)=\int\prod_{j=1}^d
       \frac{\mathbf1_{\{0\leq x_j\leq\theta_j\}}}{\theta_j}
       \,dG(\theta),
\]
the classical scale-mixture-of-uniforms model. Pavlides and
Wellner~\cite[Section~5, p.~88]{p02:bud:PW2012} posed density-MLE rate
conjectures for that model. Kim, Kur and Guntuboyina
\cite[Corollary~4.4]{p02:bud:KKG2025} later proved the proposed cube-root Hellinger rate
up to logarithms under an additional lower-bound condition. Our
experiment permits repeated measurements of the same endpoint vector,
and the target is the latent $G$ in $W_1$, rather than $p_G$ in
Hellinger distance. Corollary~\ref{p02:bud:cor:uniform-endpoints} therefore
solves a specified replication extension, not those original MLE
conjectures.

\subsection{What the dimensional result determines}

For regular components the allocation mechanism has no logarithmic
exception at $d=2$: its level variances grow logarithmically, but their
cost-weighted square roots sum geometrically. This is a statement about
the grouped experiment, and does not settle critical-dimensional
questions for direct empirical measures or an equal-group restriction.
For Gaussian covariance mixing, the latent dimension $k=m+\dim V$
controls the statistical budget, while algebraic rank controls exact
observability from a fixed number of replicates. For uniform endpoints,
the change from a quadratic to a linear local information exponent
improves the budget exponent from $1/(d+2)$ to $1/(d+1)$.

The one-dimensional endpoint logarithm is left open. Likewise, the
equal-group binomial MLE range conjecture of Vinayak et al.
\cite[Remark~3.2]{p01:info:VKVK2019}, and its later refinements in
\cite[Theorem~20 and Section~6.2 of the author version]{p01:info:MMRS2025}, are not resolved by
switching to a different allocation class. The completed results here
are the multidimensional arbitrary-allocation minimax theorems and
their fully specified applications.


% ===== Consolidated body: bessel =====
\section{Thorin symmetry and the remaining McKay classifications}
\label{p02:bes:sec:mckay}

Two externally posed problems guide this section: the multipart McKay
classification in \cite[\S4.3]{p01:rare:Baricz2026}, and the
self-decomposability problem in \cite[\S4.4]{p01:rare:Baricz2026}.
The published theorem numbers 3, 4, 12, 14, 18, and 22 correspond to
arXiv numbers 2.3, 2.4, 3.12, 3.14, 3.18, and 3.22.
The previous manuscript treated the HCM obstruction for the four McKay
densities and classified the second family.  Those results are not
counted as new here.  We now classify the GGC property of the third
and fourth families over their entire admissible parameter ranges,
answer the additional HCM question for Theorem 12, and prove SD for
all three families singled out in the third source problem.

The McKay argument uses the classical positive hypergeometric ratio
measure, followed by a new application of its inverse first moment.
The necessity argument extracts the symmetry, total mass, and forced
atoms of a possible Thorin measure.  The SD argument transfers the
order monotonicity theorem of Baricz and Prabhu K
\cite{p02:bes:BariczPrabhu2026} to a heat-trace bound.  That theorem was
posted on 7 October 2026, and is explicitly credited; its proof and
discussion contain no SD classification of the three families below.
These are two source problems with several cases, not a separate
conjecture for every case.

\subsection{A positive Thorin measure from a hypergeometric ratio}

Write $F(A,B;C;z)={}_2F_1(A,B;C;z)$.  We use the following classical
representation, stated explicitly as equation (5-1) in the published
version of \cite{p01:rare:Baricz2026} and credited there to K\"ustner.
For $0<A,B\leq C$ there is a probability measure $q$ on $[0,1]$ with
\begin{equation}\label{p02:bes:eq:kustner}
 \frac{F(A+1,B+1;C+1;z)}{F(A,B;C;z)}
 =\int_{[0,1]}\frac{q(dt)}{1-tz},
 \qquad z\in\mathbb C\setminus[1,\infty).
\end{equation}
This representation is a classical input, not a new assertion of the
present work.

\begin{lemma}[The inverse first moment of the ratio measure]
\label{p02:bes:lem:inverse-moment}
Under these hypotheses, with $m=\min(A,B)$,
\begin{equation}\label{p02:bes:eq:inverse-moment}
 q(\{0\})=0,\qquad
 \int_{(0,1]}\frac{q(dt)}t=\frac{C}{\max(A,B)}.
\end{equation}
\end{lemma}
\begin{proof}
By symmetry suppose $A\leq B$.  If $A=C$, then $A=B=C$ and the
identity $F(A,A;A;z)=(1-z)^{-A}$ proves the claim directly.
Otherwise $A<C$, and Euler's integral gives
\[
 F(A,B;C;-R)=D\int_0^1
 t^{A-1}(1-t)^{C-A-1}(1+Rt)^{-B}\,dt,
 \qquad D=\frac{\Gamma(C)}{\Gamma(A)\Gamma(C-A)}.
\]
For $B>A$, substitution $u=Rt$ shows
\[
 F(A,B;C;-R)\sim D R^{-A}\mathrm B(A,B-A).
\]
For completeness, dominated convergence applies on $u\leq R/2$;
the contribution from $u>R/2$, after multiplication by $R^A$, is
$O(R^{A-B})$, by integrability of $(1-t)^{C-A-1}$ at one.  Applying
the same argument to the differentiated integral gives
\begin{equation}\label{p02:bes:eq:negative-axis-asymptotic}
 \frac{R F'(A,B;C;-R)}{F(A,B;C;-R)}\longrightarrow A.
\end{equation}
Indeed, the ratio of its two limiting integrals is
$B\mathrm B(A+1,B-A)/\mathrm B(A,B-A)=A$.
When $B=A$, the denominator after the substitution is asymptotic to
$\log R$.  In the numerator the extra factor is $A u/(1+u)$;
replacing it by $A$ changes the integral by a bounded amount, since
the integrand then gains a factor $(1+u)^{-1}$.  The contribution
near $u=R$ is bounded in the same way as above.  Thus
\eqref{p02:bes:eq:negative-axis-asymptotic} remains valid.

Using $F'=(AB/C)F(A+1,B+1;C+1;\cdot)$ in
\eqref{p02:bes:eq:kustner}, monotone convergence now gives
\[
 \int\frac{q(dt)}t
 =\lim_{R\to\infty}\int\frac{R}{1+Rt}\,q(dt)
 =\frac{C}{AB}A=\frac CB.
\]
Its finiteness also rules out an atom at zero.
\end{proof}

\begin{theorem}[A symmetric hypergeometric Thorin measure]
\label{p02:bes:thm:hypergeometric-thorin}
Let $b>d>0$, $0<A,B\leq C$, and $M\geq2m$, where
$m=\min(A,B)$.  Normalize
\begin{equation}\label{p02:bes:eq:hypergeometric-laplace}
 L(s)=\left(\frac b{s+b}\right)^M
 \frac{F(A,B;C;d^2/(s+b)^2)}{F(A,B;C;d^2/b^2)},
 \qquad s\geq0.
\end{equation}
Then $L$ is the Laplace transform of a GGC.  Its Thorin measure is
\begin{equation}\label{p02:bes:eq:explicit-thorin}
 U=(M-2m)\delta_b+(T_-)_{\#}w+(T_+)_{\#}w,
 \qquad
 w(dt)=\frac{AB}{C}\frac{q(dt)}t,
 \quad T_\pm(t)=b\pm d\sqrt t.
\end{equation}
It has total mass $M$ and is supported in $[b-d,b+d]$.
Here $(T)_\#w$ denotes the pushforward measure.
\end{theorem}
\begin{proof}
Put $u=s+b$ and $h=-\log L$.  Differentiation and
\eqref{p02:bes:eq:kustner} give
\[
 h'(s)=\frac M u+
 \frac{2ABd^2}{Cu}\int\frac{q(dt)}{u^2-d^2t}.
\]
For $t>0$ the elementary identity
\[
 \frac{2d^2}{u(u^2-d^2t)}
 =\frac1t\left\{\frac1{u-d\sqrt t}
                      +\frac1{u+d\sqrt t}-\frac2u\right\}
\]
and Lemma~\ref{p02:bes:lem:inverse-moment} yield
\[
 h'(s)=\frac{M-2m}{s+b}
 +\int\left\{\frac1{s+b-d\sqrt t}
                   +\frac1{s+b+d\sqrt t}\right\}w(dt).
\]
All measures and coefficients are nonnegative, $w((0,1])=m$, and
all rates are bounded below by $b-d>0$.  Integrating from zero to $s$
therefore gives the Thorin representation
$h(s)=\int\log(1+s/r)\,U(dr)$ with \eqref{p02:bes:eq:explicit-thorin}.
This is the defining GGC representation and proves all assertions.
\end{proof}

\subsection{A necessary symmetry and atom budget}

\begin{lemma}[Symmetry forced by an odd logarithmic derivative]
\label{p02:bes:lem:thorin-symmetry}
Let a probability density $f$ have exponential decay rate $r_0=b-d>0$,
in the sense that its Laplace integral converges throughout
$\operatorname{Re}s>-r_0$ and is analytic there. Suppose its logarithmic derivative
$g(s)=-L'(s)/L(s)$ has, for sufficiently large $|s+b|$, a convergent
Laurent expansion
\begin{equation}\label{p02:bes:eq:odd-expansion}
 g(s)=\frac M{s+b}
       +\sum_{n\geq1}\frac{c_n}{(s+b)^{2n+1}},
 \qquad M>0,
\end{equation}
with real coefficients.  If the law is GGC, its Thorin measure $U$
has total mass $M$, is symmetric about $b$, and is supported in
$[b-d,b+d]$.  If, in addition,
\begin{equation}\label{p02:bes:eq:tail-power}
 f(x)\sim C x^{q-1}e^{-(b-d)x},\qquad C>0,\quad q>0,
\end{equation}
then
\begin{equation}\label{p02:bes:eq:endpoint-atoms}
 U(\{b-d\})=U(\{b+d\})=q.
\end{equation}
Finally, whenever $\lim_{\eta\downarrow0}i\eta g(-b+i\eta)$ exists,
that limit is $U(\{b\})$.
\end{lemma}
\begin{proof}
For a GGC one has
$g(s)=a+\int(s+r)^{-1}\,U(dr)$ with $a\geq0$.
The expansion at positive infinity gives $a=0$, and monotone
convergence in $s/(s+r)$ gives $U((0,\infty))=M$.
Since $L$ is positive on the real interval $s>-r_0$, $g$ is locally
analytic and real across that interval.  Stieltjes inversion implies
$U((0,r_0))=0$.

The Laurent series is a single-valued analytic continuation to the
exterior of a sufficiently large circle centered at $-b$, and is real
on its real part.  In the upper half-plane it agrees with the
Stieltjes representation by the identity theorem.  Stieltjes inversion
on the sufficiently negative real axis therefore implies that $U$
has bounded support.  Consequently all centered moments exist, and
expansion of $\int(s+b+(r-b))^{-1}\,U(dr)$ shows from
\eqref{p02:bes:eq:odd-expansion} that every odd centered moment vanishes.
The reflection of $U$ about $b$ has the same moments as $U$; uniqueness
of moments for compactly supported measures makes the two measures
equal.  Reflection of the lower support bound $r\geq b-d$ gives the
upper bound $r\leq b+d$.

Assumption \eqref{p02:bes:eq:tail-power}, applied also to $xf(x)$, gives
$L(-r_0+\varepsilon)\sim C\Gamma(q)\varepsilon^{-q}$ and
$-L'(-r_0+\varepsilon)\sim
C\Gamma(q+1)\varepsilon^{-q-1}$.  These elementary Abelian limits
follow by scaling $y=\varepsilon x$ and bounding the tail above and
below by $(1\pm\delta)Cx^{q-1}e^{-r_0x}$.
Thus $\varepsilon g(-r_0+\varepsilon)\to q$.
In the Stieltjes representation the same limit, by dominated
convergence and the lower support bound, is $U(\{r_0\})$.
Symmetry gives the other endpoint atom.
For the center, $|i\eta/(r-b+i\eta)|\leq1$, and dominated
convergence directly gives
$i\eta g(-b+i\eta)\to U(\{b\})$.
\end{proof}

\begin{theorem}[Complete GGC classification of the generalized McKay law]
\label{p02:bes:thm:third-family}
Let $\mu>-1$, $\mu+\nu>0$, and $b>a>0$.  The law with density
\begin{equation}\label{p02:bes:eq:third-density}
 f_{\mu,\nu,a,b}(x)=C_{\mu,\nu,a,b}
 x^{\nu-1}e^{-bx}I_\mu(ax),\qquad x>0,
\end{equation}
where $C_{\mu,\nu,a,b}>0$ normalizes the integral, is GGC if and only if
\begin{equation}\label{p02:bes:eq:third-ggc-boundary}
 \nu\leq\mu+1.
\end{equation}
It is therefore self-decomposable throughout this range, which is
exactly the range covered by \cite[Theorem~3]{p01:rare:Baricz2026}.
Its Thorin measure there has total mass $\mu+\nu$, is symmetric
about $b$, has no central atom, and has endpoint atoms of mass
$(\nu-1/2)_+$ each at $b\pm a$.
\end{theorem}
\begin{proof}
The density is integrable by its Bessel asymptotics at zero and
infinity.  Termwise integration of its nonnegative Bessel series gives
\eqref{p02:bes:eq:hypergeometric-laplace} with
\[
 M=\mu+\nu,\quad A=M/2,\quad B=(M+1)/2,
 \quad C=\mu+1,\quad d=a.
\]
The condition $\nu\leq\mu+1$ is precisely $B\leq C$;
then $0<A<B\leq C$ and $M=2A$.
Theorem~\ref{p02:bes:thm:hypergeometric-thorin} proves sufficiency and
the stated total mass and absence of a central atom.

For necessity the hypergeometric function is analytic and nonzero
near zero, regardless of \eqref{p02:bes:eq:third-ggc-boundary}.  Its
logarithmic derivative consequently has an odd Laurent expansion
\eqref{p02:bes:eq:odd-expansion} with leading coefficient $M$.
At infinity the density satisfies
\[
 f_{\mu,\nu,a,b}(x)\sim C' x^{\nu-3/2}e^{-(b-a)x}.
\]
If $\nu>1/2$, Lemma~\ref{p02:bes:lem:thorin-symmetry} forces two
atoms of mass $\nu-1/2$.  Their total cannot exceed $M$, so
$2\nu-1\leq\mu+\nu$, or $\nu\leq\mu+1$.
If $\nu\leq1/2$, the admissibility condition $\mu+\nu>0$ implies
$\mu+1>1-\nu\geq\nu$, so the same necessary inequality is automatic.
This proves the equivalence.

The endpoint mass is $\nu-1/2$ when positive, as just proved.
For $\nu<1/2$ the transform stays finite at its exponential boundary,
so $\varepsilon g(-b+a+\varepsilon)\to0$; for $\nu=1/2$ its
growth is logarithmic and the same limit is zero.  This proves the
remaining endpoint claims.  Every GGC is SD, either by the classical
inclusion or by its decreasing L\'evy spectral function
$k(t)=\int e^{-rt}\,U(dr)$.
\end{proof}

\begin{theorem}[Complete GGC classification of the squared-Bessel McKay law]
\label{p02:bes:thm:fourth-family}
Let $\mu>-1/4$ and $b>2a>0$.  The density
\begin{equation}\label{p02:bes:eq:fourth-density}
 f_{\mu,a,b}(x)=C_{\mu,a,b}x^{2\mu}e^{-bx}I_\mu(ax)^2
\end{equation}
defines a GGC if and only if
\begin{equation}\label{p02:bes:eq:fourth-boundary}
 -\tfrac14<\mu\leq\tfrac12.
\end{equation}
Hence every law covered by \cite[Theorem~4]{p01:rare:Baricz2026} is SD.
Its Thorin measure in this range has total mass $4\mu+1$, central
atom $2\max(\mu,0)$ at $b$, and endpoint atoms $2\max(\mu,0)$
each at $b\pm2a$.
\end{theorem}
\begin{proof}
The source's Laplace calculation, or multiplication and integration
of the Bessel series, gives \eqref{p02:bes:eq:hypergeometric-laplace} with
\[
 M=4\mu+1,\quad A=\mu+\tfrac12,
 \quad B=2\mu+\tfrac12,
 \quad C=\mu+1,
 \quad d=2a.
\]
For $-1/4<\mu\leq1/2$ one has $0<A,B\leq C$.
If $\mu\leq0$, $M=2B=2\min(A,B)$; if $\mu\geq0$,
$M-2\min(A,B)=2\mu$.  Theorem~\ref{p02:bes:thm:hypergeometric-thorin}
therefore proves sufficiency and the central mass formula.

Suppose now $\mu>1/2$ and the law were GGC.
Its odd Laurent expansion has total Thorin mass $M=4\mu+1$.
Its tail is $C'x^{2\mu-1}e^{-(b-2a)x}$, so
Lemma~\ref{p02:bes:lem:thorin-symmetry} forces endpoint atoms of mass
$2\mu$ each.  There is also a forced central atom.
Indeed, put $u=i\eta$ and $z=d^2/u^2=-d^2/\eta^2$.  The
logarithmic derivative satisfies
\[
 u g(u-b)=M+2z\frac{F'(A,B;C;z)}{F(A,B;C;z)}.
\]
Here $0<A<C$ and $B>A$, so the Euler-integral proof of
\eqref{p02:bes:eq:negative-axis-asymptotic}, which does not require
$B\leq C$, gives $zF'/F\to-A$.  The displayed expression tends to
$M-2A=2\mu$.  The analytic continuation along the imaginary $u$-axis
is legitimate: its hypergeometric argument is negative, and Euler's
integral is strictly positive there; continuation from large $|u|$
agrees with the Stieltjes transform by uniqueness.  The last clause
of Lemma~\ref{p02:bes:lem:thorin-symmetry} identifies the limit with
$U(\{b\})$.

The two endpoint atoms and this central atom are at three distinct
points and have total mass $6\mu$.  Positivity of $U$ would imply
$6\mu\leq4\mu+1$, contrary to $\mu>1/2$.
This proves necessity.  The endpoint masses in the positive range
follow from the tail in the same way as in
Theorem~\ref{p02:bes:thm:third-family}; they vanish at and below $\mu=0$.
The SD conclusion again follows from the GGC representation.
\end{proof}

At $\mu=1/2$ the fourth family has the especially transparent Thorin
measure $\delta_{b-2a}+\delta_b+\delta_{b+2a}$: it is a sum of three
independent exponential variables with those rates.  The necessity
proof identifies why extending the GGC range beyond $1/2$ is
impossible: the forced masses at the center and both endpoints would
exceed the entire available Thorin mass.  A failure of GGC outside
the stated ranges is not a proof of failure of SD; the latter remains
a separate question outside the parameters in the source theorems.

\subsection{The additional HCM density in the source problem}

\begin{proposition}[Theorem 12 has no HCM parameter]
\label{p02:bes:prop:theorem12}
For every $\mu>0$, the density in
\cite[Theorem~12]{p01:rare:Baricz2026}, namely
\begin{equation}\label{p02:bes:eq:theorem12-density}
 \varsigma_\mu(x)=\frac\mu x e^{-1/(2x)}I_\mu(1/(2x)),
 \qquad x>0,
\end{equation}
is not HCM.
\end{proposition}
\begin{proof}
The classical HCM representation
$f(x)=C x^{\beta-1}h_1(x)h_2(1/x)$ with
$-(\log h_i)'$ Stieltjes implies that a positive HCM function has a
zero-free holomorphic continuation to
$D=\mathbb C\setminus(-\infty,0]$: integrate the two Stieltjes
functions on this simply connected domain and exponentiate.
See \cite[eq.~(7)]{p01:rare:BondessonSimon2016}.
On the other hand, setting
$\mathcal F_\mu(z)={}_0F_1(;\mu+1;z^2/4)$, the density in
\eqref{p02:bes:eq:theorem12-density} extends to
\[
 C' x^{-\mu-1}e^{-1/(2x)}\mathcal F_\mu(1/(2x))
\]
on $D$, with the principal power.  The Bessel zeros give
\[
 \mathcal F_\mu(\pm i j_{\mu,n})=0,
 \qquad x=\pm i/(2j_{\mu,n})\in D.
\]
The continuation vanishes at these points.  Its other factors are holomorphic
and nonzero there, contradicting the HCM continuation property.
\end{proof}

\section{A heat-trace transfer closes the three Bessel SD cases}
\label{p02:bes:sec:sd}

Let $j_{\nu,n}$ be the positive zeros of $J_\nu$ and put
$Q_\nu(t)=\sum_{n\geq1}j_{\nu,n}^2e^{-t j_{\nu,n}^2}$.
The following recent theorem is an external input, not a new solution
of the Ismail--Miller conjecture in this manuscript:
\begin{equation}\label{p02:bes:eq:known-order-theorem}
 Q_\nu(t)\leq Q_\mu(t)
 \quad\text{for }t>0,\quad\nu\geq\mu\geq0.
\end{equation}
It is Theorem 1.1 of Baricz and Prabhu K
\cite{p02:bes:BariczPrabhu2026}.  The primary proof derives a positive
integral representation for an order derivative of $I_{\nu+1}/I_\nu$,
identifies a completely monotone function by Bessel products, and
uses a probability-tail representation to obtain
\eqref{p02:bes:eq:known-order-theorem}.  We apply it to a different
question, the SD problem recorded in \cite[\S4.4]{p01:rare:Baricz2026}.

\begin{lemma}[A sharp continuum bound for the sine heat trace]
\label{p02:bes:lem:theta-bound}
For every $t>0$,
\begin{equation}\label{p02:bes:eq:theta-bound}
 Q_{1/2}(t)=\pi^2\sum_{n\geq1}n^2e^{-\pi^2n^2t}
 <\frac1{4\sqrt\pi\,t^{3/2}}.
\end{equation}
The ratio of the two sides tends to one as $t\downarrow0$.
\end{lemma}
\begin{proof}
The Jacobi transformation, obtained by Poisson summation for a
Gaussian, reads
\[
 \sum_{n\geq1}e^{-\pi^2 n^2t}
 =\frac1{2\sqrt{\pi t}}-\frac12
   +\frac1{\sqrt{\pi t}}\sum_{n\geq1}e^{-n^2/t}.
\]
Termwise differentiation, justified locally uniformly for $t>0$, gives
\[
 Q_{1/2}(t)=\frac1{4\sqrt\pi\,t^{3/2}}
 -\frac1{2\sqrt\pi\,t^{5/2}}
       \sum_{n\geq1}(2n^2-t)e^{-n^2/t}.
\]
For $0<t\leq1$ every term in the last sum is positive.  This proves
the strict inequality there, as well as the asymptotic at zero.
For $t\geq1$, the ratio of successive summands of
$\sum n^2e^{-\pi^2n^2t}$ is at most $4e^{-3\pi^2t}$, and therefore
\[
 4\sqrt\pi\,t^{3/2}Q_{1/2}(t)
 \leq\frac{4\pi^{5/2}t^{3/2}e^{-\pi^2t}}
                 {1-4e^{-3\pi^2t}}<1.
\]
For an explicit elementary bound, $t^{3/2}e^{-\pi^2t}$ decreases
on $[1,\infty)$, the denominator is greater than $1/2$, and the
last expression is less than $8\pi^{5/2}e^{-\pi^2}<256e^{-9}<1$.
\end{proof}

\begin{theorem}[A self-decomposable normalized $I$-Bessel factor]
\label{p02:bes:thm:I-factor-sd}
For $\nu\geq1/2$ and $a>0$, the function
\begin{equation}\label{p02:bes:eq:I-factor}
 P_{\nu,a}(s)=e^{-a\sqrt s}
 {}_0F_1(;\nu+1;a^2s/4),\qquad s\geq0,
\end{equation}
is the Laplace transform of a self-decomposable probability law.
Its L\'evy spectral function is
\begin{equation}\label{p02:bes:eq:I-k}
 k_{\nu,a}(t)=\frac{a}{2\sqrt{\pi t}}
     -\sum_{n\geq1}e^{-j_{\nu,n}^2t/a^2},
\end{equation}
and $k_{\nu,a}'(t)<0$ for every $t>0$.
\end{theorem}
\begin{proof}
The Bessel product formula gives, for $h=-\log P_{\nu,a}$,
\[
 h'(s)=\frac a{2\sqrt s}
    -\sum_{n\geq1}\frac1{s+j_{\nu,n}^2/a^2}
    =\int_0^\infty e^{-st}k_{\nu,a}(t)\,dt.
\]
The separate Laplace transforms are finite for $s>0$, so termwise
integration is justified.  Differentiating \eqref{p02:bes:eq:I-k} and
using \eqref{p02:bes:eq:known-order-theorem} and
Lemma~\ref{p02:bes:lem:theta-bound} gives
\[
 k_{\nu,a}'(t)
 =-\frac a{4\sqrt\pi\,t^{3/2}}
      +\frac1{a^2}Q_\nu(t/a^2)<0.
\]
The spectral function tends to zero as $t\to\infty$, so its strict
decrease implies $k_{\nu,a}(t)>0$.  Moreover
$k_{\nu,a}(t)\leq a/(2\sqrt{\pi t})$.
Thus $k(t)/t$ integrates $\min(1,t)$, and integrating the displayed
formula for $h'$ yields
\[
 h(s)=\int_0^\infty(1-e^{-st})\frac{k_{\nu,a}(t)}t\,dt,
\]
since $P_{\nu,a}(0)=1$.  This proves infinite divisibility without
assuming it in advance.  For $0<c<1$, monotonicity of $k$ gives
\[
 h(s)-h(cs)=\int_0^\infty(1-e^{-st})
              \frac{k(t)-k(t/c)}t\,dt,
\]
again a Bernstein exponent.  Hence $P_{\nu,a}(s)/P_{\nu,a}(cs)$ is
a Laplace transform, which proves self-decomposability.
\end{proof}

We next record the two $K$-Bessel factors used to transfer this result.
\begin{lemma}[Two classical GGC $K$-factors]
\label{p02:bes:lem:K-factors}
The function
\begin{equation}\label{p02:bes:eq:K-factor}
 G_{\nu,b}(s)=\frac{2^{1-\nu}b^\nu}{\Gamma(\nu)}
      s^{\nu/2}K_\nu(b\sqrt s)
\end{equation}
is GGC for $\nu>0$, $b>0$.  The reciprocal factor
\begin{equation}\label{p02:bes:eq:R-factor}
 R_{\nu,b}(s)=2^{\nu-1}\Gamma(\nu)b^{-\nu}
       s^{-\nu/2}\frac{e^{-b\sqrt s}}{K_\nu(b\sqrt s)}
\end{equation}
is GGC for $\nu\geq1/2$, $b>0$, with the convention $R_{1/2,b}=1$.
\end{lemma}
\begin{proof}
The first is the transform of an inverse gamma density, proportional
to $x^{-\nu-1}e^{-b^2/(4x)}$.  The product of its values
at $uv$ and $u/v$ is a positive constant times
$\exp\{-(b^2/(4u))(v+v^{-1})\}$, so that density is HCM and hence GGC.

For the second, write $D_\nu(t)=J_\nu(t)^2+Y_\nu(t)^2$.
The classical Stieltjes representation of $K_{\nu-1}/K_\nu$, also
used in the proof of Theorem 22 of \cite{p01:rare:Baricz2026}, yields
\begin{equation}\label{p02:bes:eq:R-stieltjes}
 -\frac{d}{ds}\log R_{\nu,b}(s)
 =\frac1\pi\int_0^\infty
 \frac{1-2/(\pi tD_\nu(t))}{s+t^2/b^2}\,dt.
\end{equation}
Here the first contribution is $b/(2\sqrt s)$ and the subtracted
one is $bK_{\nu-1}(b\sqrt s)/(2\sqrt sK_\nu(b\sqrt s))$;
the recurrence for $K_\nu'$ cancels the apparent $\nu/(2s)$ term.
Nicholson's integral \cite[eq.~10.9.30]{p01:rare:DLMF}
\[
 D_\nu(t)=\frac8{\pi^2}\int_0^\infty
 K_0(2t\sinh u)\cosh(2\nu u)\,du
\]
shows that $D_\nu(t)\geq D_{1/2}(t)=2/(\pi t)$ for $\nu\geq1/2$.
Thus the numerator in \eqref{p02:bes:eq:R-stieltjes} is nonnegative.
The standard small- and large-argument Bessel asymptotics ensure
convergence; alternatively it follows directly from the convergent
Stieltjes representation before subtraction.  At $s=0$ the
normalization in \eqref{p02:bes:eq:R-factor} has limit one.  Integrating
the positive Stieltjes derivative from zero proves the GGC
representation.  For $\nu=1/2$ the numerator vanishes identically.
\end{proof}

\begin{theorem}[The three SD cases in the third source problem]
\label{p02:bes:thm:third-source-problem}
All three laws asked about in \cite[\S4.4]{p01:rare:Baricz2026} are
self-decomposable throughout the parameter ranges of the source
theorems.  More precisely:
\begin{enumerate}
 \item The law in Theorem 14 is SD for $a,b,\nu>0$ and $\mu>1/2$.
 \item The law in Theorem 18 is SD for $a,b>0$ and $\mu,\nu>1/2$.
 \item The law in Theorem 22 is SD for $a,b>0$ and $\mu,\nu>1/2$.
\end{enumerate}
The same assertions extend to the non-strict endpoints of any of
these order restrictions where the displayed factors are defined.
\end{theorem}
\begin{proof}
With the normalizations \eqref{p02:bes:eq:I-factor}--\eqref{p02:bes:eq:R-factor},
the three Laplace transforms in the source factor exactly as
\begin{equation}\label{p02:bes:eq:three-factorizations}
 L_{14}(s)=P_{\mu,a}(s)G_{\nu,b}(s),\qquad
 L_{18}(s)=P_{\mu,a}(s)P_{\nu,b}(s),\qquad
 L_{22}(s)=P_{\mu,a}(s)R_{\nu,b}(s).
\end{equation}
For example, the first product is
\[
 \frac{2^{\mu-\nu+1}\Gamma(\mu+1)b^\nu}{a^\mu\Gamma(\nu)}
 e^{-a\sqrt s}s^{(\nu-\mu)/2}
 I_\mu(a\sqrt s)K_\nu(b\sqrt s),
\]
which is precisely the source's normalized transform.
The other two products follow by the same cancellation of powers
and constants.  Theorem~\ref{p02:bes:thm:I-factor-sd} treats every $P$
factor, and Lemma~\ref{p02:bes:lem:K-factors} treats every $G$ and $R$
factor.  A GGC is SD, and a product of SD Laplace transforms is SD:
for each $c\in(0,1)$ its ratio at $s$ and $cs$ is a product of
Laplace transforms.  This proves all three statements, including
their defined endpoint extensions.
\end{proof}

\paragraph{What is closed and what is not.}
For the parameter sets actually stated in Theorems 1--4 and 12 of
the source, the preceding manuscript and this section together give
the full HCM/GGC/SD answers requested in its second problem.
In addition, Theorems~\ref{p02:bes:thm:third-family} and
\ref{p02:bes:thm:fourth-family} classify GGC on the larger full domains
where those densities exist.  SD outside the source's ranges for
the third and fourth McKay families has not been classified here.
Theorem~\ref{p02:bes:thm:third-source-problem} answers the third source
problem on all its stated ranges, using the explicitly credited
recent order theorem.  The classical ratio representation, HCM
zero-free principle, and order theorem themselves are not claimed as
new discoveries.  Historical priority for the consequences proved
here still requires specialist literature review; their full proofs
and precise scope do not depend on a priority claim.


% ===== Consolidated body: moment_transfer =====
\section{Moment cancellation transferred to diffraction and heat observation}
\label{p02:mt:section}

This section treats two different observation maps.  Its first result
constructs a compact Cantor family of strongly log-concave objects with
the same complete diffraction data and any prescribed finite number of
common moments.  Its second result gives an exact criterion for observing
a one-sided heat source through time jets at a single sensor.  The
underlying convolution ambiguity, marginal log-concavity, and entire-function
construction are classical; the simultaneous constraints and the precise
transfer statements are proved below.  Historical priority for these
formulations is not asserted.

\subsection{A Cantor family invisible to complete Fourier intensity}
\label{p02:mt:phase-section}

For a probability density $f$ on $\mathbb R$, put
$\widehat f(\xi)=\int e^{-i\xi x}f(x)\,dx$.  We call $f$
$\kappa$-strongly log-concave if $e^{\kappa x^2/2}f(x)$ is
log-concave, with zero values allowed outside its interval of support.
For smooth positive densities this means $(\log f)''\le-\kappa$.
Write $G^\sigma$ for the global Gevrey class defined by
$\|f^{(r)}\|_\infty\le AB^r(r!)^\sigma$ for some $A,B>0$.

\begin{theorem}[A compact ambiguity family under simultaneous shape priors]
\label{p02:mt:phase-cantor}
Fix $L>0$, $\kappa>0$, $\sigma>1$, and an integer $K\ge1$.
There is a family
\[
 \{f_\omega:\omega\in\{-1,1\}^{\mathbb N},\ \omega_1=1\}
 \subset G^\sigma(\mathbb R)
\]
of probability densities with all of the following properties:
\begin{enumerate}
\item Every density has support exactly $[-L,L]$, is positive on
its interior, has mean zero, and is $\kappa$-strongly log-concave.
All members obey common Gevrey constants $A,B$.
\item The Fourier magnitude $|\widehat f_\omega(\xi)|$ is independent
of $\omega$ for every real $\xi$.  The moments
$\int x^r f_\omega(x)\,dx$, $0\le r\le K$, are also independent
of $\omega$.
\item Distinct members are not translates or reflected translates
of one another.  The displayed parametrization is a homeomorphism
from the Cantor space onto its image in $C^\infty(\mathbb R)$.
\end{enumerate}
In particular, the ambiguity is uncountable even after the usual
translation, reflection, and constant-phase symmetries are removed.
\end{theorem}

\begin{proof}
Choose $a_j>0$ with $\sum_j a_j=L$, and $c_j>0$ with
$\sum_jc_j^{-1}\le\kappa^{-1}$; for example,
$a_j=L2^{-j}$ and $c_j=\kappa2^j$.  Choose increasing odd integers
$m_j>K$.  We first construct a density $q_j$ on $[-a_j,a_j]$
which is $c_j$-strongly log-concave and whose first nonzero odd
cumulant has order $m_j$.

Put $p=(\sigma-1)^{-1}$ and normalize the even density
\[
 b_j(x)=C_j\exp\!\left(-c_jx^2
          -(a_j-x)^{-p}-(a_j+x)^{-p}\right)
          \mathbf1_{|x|<a_j}.
\]
It belongs to $G^\sigma$, is positive on its interior, and satisfies
$(\log b_j)''\le-2c_j$.  Let $\psi_j$ be an even nonnegative
nonzero Gevrey bump supported in $(-a_j/2,a_j/2)$, and set
$h_j=\psi_j^{(m_j)}$.  Integration by parts gives
\begin{equation}\label{p02:mt:cancel}
 \int x^r h_j(x)\,dx=0\quad(0\le r<m_j),\qquad
 \int x^{m_j}h_j(x)\,dx=(-1)^{m_j}m_j!\int\psi_j\ne0.
\end{equation}
The function $h_j$ is odd.  For a sufficiently small nonzero
$\varepsilon_j$, the density $q_j=b_j+\varepsilon_jh_j$ remains
positive and $c_j$-strongly log-concave: on the compact perturbation
support this follows by continuity in $C^2$, and elsewhere nothing
changes.  It has mass one and mean zero.  Its odd moments below
$m_j$ vanish, whereas its $m_j$th moment does not.  The
moment--cumulant identities therefore show that its odd cumulants
below $m_j$ vanish and that its $m_j$th cumulant is nonzero.
Write this last cumulant as $\gamma_j\ne0$.

Let $X_j$ be independent with density $q_j$.  For each sign sequence
$\omega$ the series $S_\omega=\sum_j\omega_jX_j$ converges absolutely
since $|X_j|\le a_j$.  Denote its density by $f_\omega$.  This density
exists because it is the convolution of $q_1(\omega_1\,\cdot)$ with
the law of the remaining sum.  Hence
\[
 \|f_\omega^{(r)}\|_\infty\le\|q_1^{(r)}\|_\infty,
\]
which gives common global Gevrey constants.  The support of the
remaining sum is its full interval
$[-\sum_{j\ge2}a_j,\sum_{j\ge2}a_j]$: finite sums have interval
support, and the omitted tails are uniformly bounded by
$\sum_{j>n}a_j\to0$.  Convolving with $q_1$ therefore gives support
$[-L,L]$ and strict positivity inside it.

We recall the precise curvature rule for convolution.  If $q,r$ are
$c,d$-strongly log-concave, write
$q(x)=e^{-cx^2/2}Q(x)$ and $r(y)=e^{-dy^2/2}R(y)$ with $Q,R$
log-concave.  Completing the square gives
\[
 e^{\frac{cd}{c+d}z^2/2}(q*r)(z)
 =\int e^{-\frac{c+d}{2}(x-\frac{d}{c+d}z)^2}
          Q(x)R(z-x)\,dx.
\]
The integrand is jointly log-concave in $(z,x)$.  The marginal theorem
of Pr\'ekopa~\cite[Theorem~6]{p02:mt:Prekopa1973} shows that $q*r$ is
$(c^{-1}+d^{-1})^{-1}$-strongly log-concave.  Consequently every
finite partial convolution here is $\kappa$-strongly log-concave.
Those convolutions converge uniformly to $f_\omega$ by uniform
continuity of $q_1$ and the deterministic tail bound.  Passing the
log-concavity inequality to the limit proves the same curvature
bound for $f_\omega$.

Reflection conjugates each factor's Fourier transform on the real
axis, so
\[
 |\widehat f_\omega(\xi)|
    =\lim_{n\to\infty}\prod_{j=1}^n|\widehat q_j(\xi)|
\]
is independent of $\omega$.  Equation~\eqref{p02:mt:cancel} says that
$q_j(x)$ and $q_j(-x)$ have the same moments through $K$.
Expanding moments of each finite sum and then passing to the bounded
limit gives the common moments asserted in the theorem.

For distinct sign sequences let $j$ be their first different index.
Cumulants add for independent bounded random variables.  This also
holds for the infinite sum: bounded partial
sums converge uniformly, and their moment-generating functions and
logarithms converge holomorphically on a common neighborhood of zero.
At the odd order $m_j$, all later factors have zero cumulant, while all earlier
factors agree between the two sequences.  The difference between the
two cumulants is therefore $2\omega_j\gamma_j\ne0$.  This proves
injectivity.  All means are zero, so a translating parameter would
have to vanish.  Reflection changes every sign; injectivity on the
full sign space shows that no two sequences with first sign $+1$
represent reflections of one another.  Positive mass one excludes
nontrivial constant phases.

Finally, if two sequences with fixed first sign agree through index
$n$, couple all $X_j$ identically.  Their remaining sums differ by
at most $2\sum_{j>n}a_j$.  For every $r\ge0$ the mean-value theorem
applied to $q_1^{(r)}$ yields
\[
 \|f_\omega^{(r)}-f_\eta^{(r)}\|_\infty
 \le 2\|q_1^{(r+1)}\|_\infty\sum_{j>n}a_j.
\]
The parametrization is continuous into $C^\infty$; a continuous
injection from a compact space into a Hausdorff space is a
homeomorphism onto its image.
\end{proof}

\paragraph{Published target and exact scope.}
The 2025 \emph{Phasebook} survey poses questions about inferring
object complexity and certifying reconstructions from diffraction
data~\cite[Questions~3.1 and Section~3.4]{p02:mt:Phasebook2025}.
Theorem~\ref{p02:mt:phase-cantor} supplies a rigorous obstruction in the
one-dimensional continuous setting even after substantial smoothness,
curvature, support, and moment information is supplied.  It does not
resolve the survey's broad algorithmic or higher-dimensional questions.
Nonnegativity alone was already shown insufficient in the discrete
setting by Beinert~\cite{p02:mt:Beinert2017}; reflection of convolution
factors is also classical.  The theorem's proved content is the
simultaneous Cantor-family construction under the stated constraints,
whose historical priority has not been established.

The obstruction applies to uncertainty quantification as well.  If a
data distribution depends on $f$ only through its Fourier intensity,
choose two members a positive distance $\delta$ apart in $L^1$ (or
in $W_1$).  Under their common data law every estimator has worst-case
expected error at least $\delta/2$.  Every confidence set covering
each object with probability at least $1-\alpha$ has diameter at
least $\delta$ with probability at least $1-2\alpha$.  These statements
follow respectively from the triangle inequality and the union bound;
they remain true with arbitrarily many or noiseless intensity data.

\subsection{An exact criterion for early-time heat jets at one sensor}
\label{p02:mt:heat-section}

Let $0<a<b$ be known.  For a finite signed measure $\mu$ supported
on $[a,b]$, its heat evolution observed at the origin is
\[
 U_\mu(t)=\frac1{\sqrt{4\pi t}}
             \int_a^b e^{-x^2/(4t)}\,d\mu(x),\qquad t>0.
\]
The one-sided support assumption removes the elementary ambiguity
between sources at $x$ and $-x$.  Let $t_j\downarrow0$ be distinct
and let $r_j\ge1$ be integers.  The observations are the time jets
$U_\mu^{(\ell)}(t_j)$ for $0\le\ell<r_j$.

\begin{theorem}[Weighted sampling criterion and exact hidden coincidences]
\label{p02:mt:heat-criterion}
The following statements are equivalent:
\begin{enumerate}
\item $\sum_j r_jt_j=\infty$.
\item The displayed observations identify every finite signed measure
supported on $[a,b]$.
\item They identify all compactly supported, smooth, strongly
log-concave probability densities with support $[a,b]$.
\end{enumerate}
If the sum is finite, then for every $\sigma>1$, every $\kappa>0$,
and every nonempty open interval $I\Subset(a,b)$, there are distinct
$\kappa$-strongly log-concave probability densities
$f_+,f_-\in G^\sigma(\mathbb R)$, positive on $(a,b)$, whose
difference is supported in $I$, such that
\[
 U_{f_+}(t)=U_{f_-}(t)\quad\Longleftrightarrow\quad
       t\in\{t_j:j\ge1\},\qquad t>0.
\]
At $t_j$ this difference has a zero of exactly order $r_j$.
\end{theorem}

\begin{proof}
Write $s=(4t)^{-1}$ and
$H_\mu(s)=\int e^{-sx^2}\,d\mu(x)$.  Multiplication by the nonzero
factor $(4\pi t)^{-1/2}$ and the analytic change of variable
$t\mapsto(4t)^{-1}$ preserve zero multiplicities.  A difference of
two observation sequences vanishes precisely when $H_\mu$ has zeros
of orders at least $r_j$ at $s_j=(4t_j)^{-1}$.

The function $e^{a^2z}H_\mu(z)$ is bounded and holomorphic on
$\operatorname{Re}z>0$.  If it is nonzero, its zeros obey the
half-plane Blaschke condition
\[
 \sum_j r_j\frac{s_j}{1+s_j^2}<\infty.
\]
Since $s_j\to\infty$, this is equivalent to
$\sum_jr_jt_j<\infty$.  Divergence therefore forces
$H_\mu\equiv0$.  Differentiating at zero gives all polynomial
moments of the pushforward under $x\mapsto x^2$; polynomial density
and injectivity of this map on $[a,b]$ give $\mu=0$.

For the converse, we use the exact-real-zero consequence of the
Luxemburg--Korevaar construction, as verified in
Wellershoff's explicit restatement~\cite[Theorem~21]{p01:mom:Wellershoff2021}.
For clarity the precise consequence needed here is the following:
if $s_j>0$ tend to infinity and $\sum_jr_j/s_j<\infty$, every
positive interval $J$ contains the support of a nonzero real
$q\in G^\sigma$ such that
\[
 \int q(y)\,dy=0,\qquad
 Z_{(0,\infty)}\left(s\mapsto\int y^s q(y)\,dy\right)
       =\{s_j\},
\]
with the prescribed multiplicities $r_j$.  To see the exactness,
the classical product construction, on a logarithmic interval,
has bilateral Laplace transform
\[
 s e^{\tau s}\prod_j(1-s^2/s_j^2)^{r_j}
                      \prod_k\cosh(\epsilon_ks),
 \qquad \epsilon_k>0,\quad\sum_k\epsilon_k<\infty.
\]
The last product is finite and strictly positive for real $s$;
the first product has exactly the claimed real zeros.  The logarithmic
pushforward gives a smooth $q$.  Multiplicative convolution with a
nonnegative Gevrey bump supported sufficiently near $1$ makes it
Gevrey while preserving its localization and every real zero, since
that convolution multiplies its Mellin transform by a strictly
positive real factor.  This is the classical imported existence
input, not a newly proved entire-function theorem.

Choose $J\Subset\{e^{-x^2}:x\in I\}$ and take such a $q$.
Set
\[
 h(x)=2xe^{-x^2}q(e^{-x^2})\quad(x>0),
 \qquad h(x)=0\quad(x\le0).
\]
It is real, compactly supported in $I$, and Gevrey of order $\sigma$;
Gevrey classes are preserved by multiplication and analytic changes
of variable on compact sets, as follows directly from the
Fa\`a di Bruno formula.  Substitution gives
\begin{equation}\label{p02:mt:heat-mellin}
 \int_a^b e^{-sx^2}h(x)\,dx
       =\int y^s q(y)\,dy.
\end{equation}
In particular $\int h=0$, and the heat transform has exactly the
desired positive-time zeros and multiplicities.

Finally normalize
\[
 f_0(x)=C\exp\!\left(-\kappa x^2
        -(x-a)^{-p}-(b-x)^{-p}\right)\mathbf1_{a<x<b},
 \qquad p=(\sigma-1)^{-1}.
\]
Its logarithmic second derivative is at most $-2\kappa$.
For sufficiently small $\varepsilon>0$, the densities
$f_\pm=f_0\pm\varepsilon h$ are positive on $(a,b)$,
$\kappa$-strongly log-concave, and of mass one, by the same compact
$C^2$ perturbation argument used above.  Equation~\eqref{p02:mt:heat-mellin}
proves every required coincidence and excludes all others.
\end{proof}

For values alone, $t_j=j^{-\alpha}$ identifies all one-sided sources
exactly when $\alpha\le1$.  Geometric sampling $t_j=2^{-j}$ does
not identify even the strongly log-concave Gevrey class.  If time
jets are available, $r_j=\lceil2^j/j\rceil$ restores identification
at these same geometric times.  These are consequences of one
weighted criterion, not separate problem resolutions.

\begin{proposition}[Quasianalytic source priors override sparse times]
\label{p02:mt:heat-qa}
Let $M_0=1$ and $M$ be positive and log-convex, with $M_k\ge k!$,
and suppose $\sum_{k\ge1}M_{k-1}/M_k=\infty$.  Among densities
smooth on $[a,b]$ satisfying $\|f^{(k)}\|_\infty\le AB^kM_k$
for some finite $A,B$, the values $U_f(t_j)$ identify $f$ for every
sequence $t_j\downarrow0$, independently of $\sum_jt_j$.
No smooth extension by zero at the endpoints is assumed.
\end{proposition}

\begin{proof}
Let $h$ be the difference of two such densities.  If its first nonzero
derivative at $a$ has order $m$, elementary endpoint Laplace asymptotics
give
\[
 \int_a^b e^{-sx^2}h(x)\,dx
 \sim \frac{h^{(m)}(a)e^{-sa^2}}{(2a)^{m+1}s^{m+1}}
 \qquad(s\to\infty).
\]
Indeed, after $u=s(x-a)$ the normalized integrand is bounded by
$Cu^me^{-2au}$ and converges to
$h^{(m)}(a)u^me^{-2au}/m!$.  This rules out zeros along an
unbounded sequence.  Thus all derivatives of $h$ at $a$ vanish.
Its zero extension just to the left of $a$ obeys the same derivative
bounds.  Denjoy--Carleman uniqueness on the enlarged
interval~\cite[Theorem~3.6 and Corollary~3.12]{p01:Rainer2022} forces
$h=0$.
\end{proof}

\paragraph{Published target and exact scope.}
Discrete recovery of heat sources and its connection with exponential
systems are established subjects; Choulli's paper gives the relevant
primary literature and analyzes point measurements through Dirichlet
series~\cite[Sections~1.1--1.3]{p02:mt:Choulli2016}.  Its bounded-domain
stability results are not claimed to be replaced or improved here.
Our target is the explicitly formulated whole-line, one-sided-source,
early-time jet problem, with an exact sampling criterion and localized
shape-constrained counterexamples.  No exact published open question
with this complete formulation was found in the targeted search.
The result is therefore reported as a proved transfer theorem, not as
a newly closed named open problem.  Neither theorem in this section
adds another resolution of the old moment-tail conjecture.


% ===== Consolidated body: prime_transfer =====
\section{From arithmetic certificates to exact recurrence entropy}
\label{p02:pt:sec:main}

This section separates three questions suggested by the 223-page integrated
manuscript and its 42-page sequel.  First, can the cost of a fixed prime
coordinate be removed from the exponential arithmetic estimate?  Second,
can a random-frequency theorem identify the leading constant of the covering
time, rather than only its order?  Third, how much of a weighted phase space
can be recovered from its one-dimensional radial distribution?  We give
precise answers below.  None is presented as the solution of a separately
verified longstanding conjecture.  In particular, the exact-prime exponent
gap for power weights with exponent greater than three remains open.

The finite-dimensional connection between ergodization and small divisors
is classical; see Bounemoura~\cite{p02:pt:Bounemoura}.  Small deviations for
rapidly decreasing weighted sums are also an established subject; for
geometric weights the constant used below is already a special case of
Aurzada's Theorem~4~\cite{p02:pt:Aurzada}, and broader fast-weight results appear
in Rozovsky~\cite{p02:pt:Rozovsky}.  Our small-ball calculation is included for
self-containment.  The additional conclusion is an almost-sure equality,
to leading constant, between that entropy and the logarithm of a uniform
covering time.  The proof uses a sum of inverse small divisors, rather than
the product of a worst small divisor and a Fourier norm.

Let $\Omega=\mathbb T^{\mathbb N}$, where angles have period $2\pi$, and
let $\mu$ be product Haar probability.  For positive weights with
$W=\sum_jw_j<\infty$, put
\[
 E_w(\theta)=\sum_jw_j(1-\cos\theta_j),\qquad
 \Phi_t\theta=(\theta_j-t\lambda_j)_j.
\]
The associated product metric is
$d_w(\theta,\varphi)^2=2E_w(\theta-\varphi)$.
Write $R_{w,\lambda}(\varepsilon)$ for the infimum of the times $T$ such
that every starting point enters $\{E_w<\varepsilon\}$ during $[0,T]$.
Translation invariance makes the same bound uniform in the target and in
the initial time of an observation interval.

\subsection{A summed-divisor certificate on compact abelian groups}

\begin{lemma}[Averaging the actual inverse divisors]
\label{p02:pt:lem:sum-certificate}
Let $G$ be a compact abelian group with Haar probability, let
$t\mapsto a_t$ be a continuous one-parameter subgroup, and write
$\chi(a_t)=e^{-it\omega_\chi}$ for its characters.  Let $F$ be a real
trigonometric polynomial with positive mean $m=\widehat F(1)$, and suppose
$\omega_\chi\ne0$ at every nontrivial character in its Fourier support.
Set
\[
 S(F)=\sum_{\chi\ne1}
             \frac{|\widehat F(\chi)|}{|\omega_\chi|}.
\]
Every orbit segment of length $T>2S(F)/m$ meets $\{F>0\}$.
\end{lemma}
\begin{proof}
Termwise integration of the finite Fourier expansion gives, for every
$x\in G$ and every initial time $A$,
\[
 \left|\frac1T\int_A^{A+T}F(xa_t)\,dt-m\right|
 \le\frac2T\sum_{\chi\ne1}
                \frac{|\widehat F(\chi)|}{|\omega_\chi|}.
\]
The average is positive under the stated inequality, so at least one
point of the segment has $F>0$.
\end{proof}

The lemma applies to nonproduct Fourier supports without alteration.
It is an elementary finite Fourier identity, not a new general
ergodization theorem.  The following probabilistic use of the identity
is what improves the leading constant.

\begin{lemma}[Fractional moments avoid the worst-divisor loss]
\label{p02:pt:lem:fractional-divisors}
Let $\lambda_1,\ldots,\lambda_n$ be independent, with densities bounded
by $L\ge1$.  Let $\mathcal K\subset\mathbb Z^n\setminus\{0\}$ be a
finite deterministic set of cardinality at most $N$, and let
$|c_k|\le B$.  For $0<q<1$,
\begin{equation}\label{p02:pt:eq:fractional-divisors}
 \mathbb E\left(\sum_{k\in\mathcal K}
                       \frac{|c_k|}{|k\cdot\lambda|}\right)^q
 \le C_q B^q N L^q.
\end{equation}
The same conclusion holds for $k\cdot\lambda-\lambda_0\sum_jk_j$,
if $\lambda_0$ is any additional independent random variable.
\end{lemma}
\begin{proof}
For each nonzero $k$, choose $j$ with $k_j\ne0$ and condition on every
coordinate except $\lambda_j$.  The conditional density of the linear
form is at most $L/|k_j|\le L$.  Consequently
\[
 \mathbb P(|k\cdot\lambda|<u)\le\min(1,2Lu).
\]
Integration of this bound, splitting at $u=(2L)^{-1}$, gives
$\mathbb E|k\cdot\lambda|^{-q}\le C_qL^q$.
For nonnegative numbers, $(\sum a_k)^q\le\sum a_k^q$; applying this
inequality proves \eqref{p02:pt:eq:fractional-divisors}.  In the relative
model condition also on $\lambda_0$.  No independence between distinct
linear forms is asserted or needed.
\end{proof}

\subsection{Rapidly decreasing spectra: small balls with the exact constant}

Fix $a,b>0$ and assume
\begin{equation}\label{p02:pt:eq:fast-weights}
 \log w_j=-a j^b+o(j^b)\qquad(j\to\infty).
\end{equation}
The weights need not be monotone.  Set
\begin{equation}\label{p02:pt:eq:entropy-scale}
 \ell=\log(1/\varepsilon),\qquad
 H(\varepsilon)=\frac{b}{2(b+1)}a^{-1/b}\ell^{1+1/b}.
\end{equation}

\begin{lemma}[Elementary spectral sums]
\label{p02:pt:lem:spectral-sums}
Under \eqref{p02:pt:eq:fast-weights},
\begin{align}
 \log\sum_{j>n}w_j&=-a n^b+o(n^b),\label{p02:pt:eq:tail-log}\\
 \frac12\sum_{j\ge1}(L+\log w_j)_+
 &\sim\frac{b}{2(b+1)}a^{-1/b}L^{1+1/b}
       \qquad(L\to\infty).\label{p02:pt:eq:spectral-sum}
\end{align}
If $n(\varepsilon)$ is the least integer with
$\sum_{j>n}w_j\le\varepsilon/8$, then
$n(\varepsilon)\sim(\ell/a)^{1/b}$.
\end{lemma}
\begin{proof}
For any $0<\delta<a$, all sufficiently late weights lie between
$e^{-(a+\delta)j^b}$ and $e^{-(a-\delta)j^b}$.
The lower bound for the tail follows from its first term.  Integral
comparison bounds the upper tail by a polynomial factor times
$e^{-(a-\delta)n^b}$, with a harmless adjustment of that polynomial
when $b>1$.  Divide logarithms by $n^b$ and let $\delta\downarrow0$.

For the second assertion, the same comparison reduces the calculation
to $w_j=e^{-c j^b}$.  Summing through
$n=(L/c)^{1/b}+O(1)$ and using
$\sum_{j\le n}j^b=n^{b+1}/(b+1)+o(n^{b+1})$ gives
\[
 \frac12\sum_{j\le n}(L-cj^b)
 =\frac{b}{2(b+1)}c^{-1/b}L^{1+1/b}(1+o(1)).
\]
Finitely many initial weights contribute only $O(L)$, which is smaller
than the displayed scale.  Letting $\delta\downarrow0$ proves the
formula for the original weights.  The last assertion follows from
\eqref{p02:pt:eq:tail-log}, including the minimality of $n(\varepsilon)$.
\end{proof}

\begin{proposition}[The Haar small-ball constant]
\label{p02:pt:prop:fast-small-ball}
For weights satisfying \eqref{p02:pt:eq:fast-weights},
\begin{equation}\label{p02:pt:eq:fast-small-ball}
 -\log\mu\{E_w\le\varepsilon\}\sim H(\varepsilon).
\end{equation}
\end{proposition}
\begin{proof}
If $\Theta$ is uniform on $[-\pi,\pi]$, then
$\mathbb P(1-\cos\Theta\le u)\asymp\sqrt u$ for $0<u\le1$.
Also
\[
 \mathbb E e^{-t(1-\cos\Theta)}=e^{-t}I_0(t),\qquad
 -\log(e^{-t}I_0(t))\ge\tfrac12\log t-C\quad(t\ge1).
\]
The last bound follows directly by comparing $1-\cos\theta$ with
$\theta^2$ on $[-\pi,\pi]$ and integrating a Gaussian.
Exponential Markov with $t=1/\varepsilon$ therefore gives
\[
 -\log\mu\{E_w\le\varepsilon\}
 \ge\frac12\sum_j(\ell+\log w_j)_+
       -C\#\{j:w_j\ge\varepsilon\}-1.
\]
The count is $O(\ell^{1/b})=o(H)$, so
Lemma~\ref{p02:pt:lem:spectral-sums} proves the lower asymptotic bound.

For the other direction take $n=n(\varepsilon)$.  The inactive tail
contributes at most $2\sum_{j>n}w_j\le\varepsilon/4$, deterministically.
For each $j\le n$ require
\[
 |\Theta_j|\le \min\{1,\sqrt{\varepsilon/(n w_j)}\}.
\]
The active energy is at most $\varepsilon/2$.  By independence, the
negative logarithm of the probability of this event is at most
\[
 Cn+\frac12\sum_{j\le n}
                     (\ell+\log n+\log w_j)_+.
\]
Extending the sum to all $j$ and applying
\eqref{p02:pt:eq:spectral-sum}, with $\log n=o(\ell)$, gives
$(1+o(1))H(\varepsilon)$.  This proves
\eqref{p02:pt:eq:fast-small-ball}.
\end{proof}

\subsection{An exact leading constant for random-frequency covering}

Let $\lambda_j=c_j+\xi_j$, where the $\xi_j$ are independent uniforms
on $[-r_j,r_j]$, with $0<r_j\le1$.  Assume that the centers have at
most polynomial growth and that
\begin{equation}\label{p02:pt:eq:width-assumption}
 L_n:=\max\{1,(2r_j)^{-1}:1\le j\le n\},\qquad
 \log L_n=o(n^{b+1}).
\end{equation}
The conclusion also holds for the relative frequencies
$\lambda_j-\lambda_0$, where $\lambda_0$ is an independent bounded
random variable.  All assertions are for a fixed deterministic weight
sequence satisfying \eqref{p02:pt:eq:fast-weights}.

\begin{theorem}[Exact recurrence entropy]
\label{p02:pt:thm:exact-entropy}
With probability one,
\begin{equation}\label{p02:pt:eq:exact-entropy}
 \log R_{w,\lambda}(\varepsilon)
 \sim-\log\mu\{E_w\le\varepsilon\}
 \sim\frac{b}{2(b+1)}a^{-1/b}
           (\log(1/\varepsilon))^{1+1/b}.
\end{equation}
The covering estimate is uniform in starting points, targets, and initial
times.  The same asymptotic holds for the supremum of the lengths of
intervals avoiding $\{E_w<\varepsilon\}$ along any one orbit.
\end{theorem}
\begin{proof}
We first construct a certificate with the correct number of modes.
For an integer $D\ge0$ put
\[
 u_k=\sqrt{\frac2{D+2}}\sin\frac{(k+1)\pi}{D+2}
 \quad(0\le k\le D),\qquad
 K_D(\theta)=\left|\sum_{k=0}^D u_ke^{ik\theta}\right|^2.
\]
The elementary sine eigenvector identity gives
\begin{equation}\label{p02:pt:eq:sine-kernel}
 \int K_D=1,\quad K_D\ge0,\quad
 \int(1-\cos\theta)K_D
       =1-\cos\frac\pi{D+2}\le\frac{\pi^2}{2(D+2)^2}.
\end{equation}
Indeed $\sum u_k^2=1$ and
$\sum_{k=0}^{D-1}u_ku_{k+1}=\cos(\pi/(D+2))$; expanding the square
proves the identities.  Its Fourier support is $[-D,D]$, each
coefficient has modulus at most one, and
$\|K_D\|_A\le(\sum|u_k|)^2\le D+1$.

Take $n=n(\varepsilon)$ and, for $1\le j\le n$, set
\[
 D_j=\left\lceil\sqrt{64\pi^2n w_j/\varepsilon}\right\rceil,
 \quad K=\prod_{j\le n}K_{D_j},\quad
 F=(\varepsilon/2-E_{w,\le n})K.
\]
Then $\widehat F(0)\ge\varepsilon/2-\varepsilon/128>
\varepsilon/3$, and $F>0$ implies $E_w<\varepsilon$.
The Fourier support is contained in $|k_j|\le D_j+1$, so its
cardinality $N_\varepsilon$ satisfies
\begin{align}
 \log N_\varepsilon
 &\le\sum_{j\le n}\log(2D_j+3)\notag\\
 &\le Cn+\frac12\sum_{j\le n}
       (\ell+\log n+\log(64\pi^2)+\log w_j)_+
 \le(1+o(1))H(\varepsilon).\label{p02:pt:eq:mode-entropy}
\end{align}
Moreover $|\widehat F(k)|\le\varepsilon/2+2W$, by convolution with
the Fourier coefficients of the energy and $|\widehat K|\le1$.
In particular the individual coefficient bound stays bounded as
$\varepsilon\downarrow0$.

Use the sequence $\varepsilon_m=e^{-m}$.  Since
$n(\varepsilon_m)\sim(m/a)^{1/b}$, assumption
\eqref{p02:pt:eq:width-assumption} yields
$\log L_{n(\varepsilon_m)}=o(H(\varepsilon_m))$.
For any fixed $\eta>0$, choose $q<1$ so close to one that
$q(1+\eta)>1$.  Lemma~\ref{p02:pt:lem:fractional-divisors} and Markov's
inequality show that
\[
 \mathbb P\{S(F)>e^{(1+\eta)H(\varepsilon_m)}\}
 \le\exp\{[1-q(1+\eta)+o(1)]H(\varepsilon_m)\}.
\]
These probabilities are summable because
$H(\varepsilon_m)\asymp m^{1+1/b}$.  Borel--Cantelli and
Lemma~\ref{p02:pt:lem:sum-certificate} give almost surely
\[
 \log R_{w,\lambda}(\varepsilon_m)
 \le(1+\eta)H(\varepsilon_m)+m+O(1).
\]
Here $m=o(H(\varepsilon_m))$.  Intersecting over positive rational
$\eta$ gives the upper constant one.  Monotonicity of the covering
time and $H(e^{-(m+1)})/H(e^{-m})\to1$ extend it to every small
$\varepsilon$.  The relative-frequency version follows from the
last assertion of Lemma~\ref{p02:pt:lem:fractional-divisors}.

For the lower bound, polynomial growth of the frequencies and
\eqref{p02:pt:eq:fast-weights} ensure
$V=\sum_jw_j|\lambda_j|<\infty$.  Hence $E_w(\Phi_t\theta)$ is
uniformly $V$-Lipschitz in $t$.  If every starting point has energy
below $\varepsilon$ at some time in $[0,T]$, a time mesh of spacing
$\varepsilon/V$ covers $\Omega$ by at most $2+TV/\varepsilon$
translates of $\{E_w\le2\varepsilon\}$.  Haar measure implies
\[
 1\le(2+TV/\varepsilon)\mu\{E_w\le2\varepsilon\}.
\]
By Proposition~\ref{p02:pt:prop:fast-small-ball}, taking logarithms gives
$\log T\ge(1-o(1))H(\varepsilon)$, since
$\log(1/\varepsilon)=o(H)$ and $H(2\varepsilon)\sim H(\varepsilon)$.
This proves \eqref{p02:pt:eq:exact-entropy}.

Almost surely all finite integer relations among the chosen frequencies
are nonzero, by conditioning on one coordinate and then taking a
countable union.  Every orbit is therefore dense by the finite-dimensional
Kronecker theorem.  To obtain long avoidance intervals on every orbit,
apply the preceding mesh argument at tolerance $2\varepsilon$ and
$T=\exp((1-\eta)H(\varepsilon))$.  For all small $\varepsilon$ its
union has measure less than one, so there is a segment whose energy
stays above $2\varepsilon$.  Compactness in time and continuity give
an open neighborhood of starting points whose segments stay above
$\varepsilon$.  Every dense orbit meets that neighborhood.  The
uniform upper covering bound supplies the opposite bound, and then
$\eta\downarrow0$ proves the final assertion.
\end{proof}

\begin{example}[Arbitrarily small perturbations of prime logarithms]
\label{p02:pt:ex:prime-perturbations}
Enumerate the primes as $p_j$, put $c_j=\log p_j$, and let
$r_j=\rho p_j^{-A}$ with any fixed $A>0$ and $0<\rho\le1$.
The width hypothesis holds because $\log L_n=O(\log n)$.
Thus \eqref{p02:pt:eq:exact-entropy} holds for weights satisfying
\eqref{p02:pt:eq:fast-weights}, for both full and relative prime
frequencies.  For $w_j=e^{-a j}$ the leading constant is $1/(4a)$;
for $w_j=e^{-a j^2}$ it is $1/(3\sqrt a)$.
This does not imply the same statement at the unperturbed vector:
a full-measure set need not contain that prescribed point.
\end{example}

\begin{proposition}[Perturbation widths can change the leading constant]
\label{p02:pt:prop:thin-width-obstruction}
Let $w_j=e^{-a j^b}$ and suppose $|\lambda_j|\le e^{-c j^{b+1}}$
for all $j$, with $a,b,c>0$.  Then
\[
 \liminf_{\varepsilon\downarrow0}
 \frac{\log R_{w,\lambda}(\varepsilon)}
      {(\log(1/\varepsilon))^{1+1/b}}
 \ge c\,a^{-(b+1)/b}.
\]
Consequently, if $c>ab/[2(b+1)]$, the ratio-one conclusion in
\eqref{p02:pt:eq:exact-entropy} fails for every such frequency vector.
This includes independent uniform perturbations of the zero vector
with these radii, which are nonresonant with probability one.
\end{proposition}
\begin{proof}
Choose the largest integer $j$ for which $w_j\ge8\varepsilon$.
Then $j\sim(\ell/a)^{1/b}$.  Start its phase at $\pi$.
During $0\le t\le\pi/(2|\lambda_j|)$ that coordinate stays in the
opposite closed semicircle, so its energy is at least $w_j$ and
the target cannot be reached.  If $\lambda_j=0$ this gives an
infinite obstruction.  Otherwise
\[
 \log R_{w,\lambda}(\varepsilon)
 \ge c j^{b+1}+\log(\pi/2).
\]
Divide by $\ell^{1+1/b}$ and compare the result with
\eqref{p02:pt:eq:entropy-scale}.  Nonresonance in the uniform model
again follows by conditioning on one coordinate for each of the
countably many nonzero integer characters.
\end{proof}

Thus control of concentration of the frequency law is substantive.
The proposition demonstrates a sharp relevant growth scale for that
control; it does not determine the full transition at the boundary
$\log L_n\asymp n^{b+1}$, and it concerns zero-centered rather than
prime-centered perturbations.

\paragraph{Precise improvement over the two previous manuscripts.}
The integrated manuscript proved matching upper and lower orders for
randomly perturbed prime flows with polynomial-logarithmic weights,
and explicitly did not claim a limiting leading constant.  The sequel
optimized a deterministic rectangular budget at its logarithmic
transition.  Theorem~\ref{p02:pt:thm:exact-entropy} treats a different,
rapidly decreasing spectral class and proves a ratio-one limit for
the actual uniform covering time.  The fractional-moment argument
is essential to this constant: bounding all divisors by their worst
value and then multiplying by the Fourier norm can pay the mode
entropy twice.  We do not claim a ratio-one limit for polynomial
weights, where the $O(n)$ terms in the mode count are of leading size.

\subsection{One exact anchor and the remaining arithmetic obstruction}

The original actual-support obstruction used a pure high-frequency
prime mode to force a large arithmetic height.  Such a pure mode has
a large, rather than small, temporal frequency.  The following
elementary refinement makes this distinction quantitative.

\begin{proposition}[An anchor prime has no exponential degree cost]
\label{p02:pt:prop:anchor-height}
Fix a prime $q$.  For a finite prime-indexed integer vector $k$ put
$v(k)=\sum_p k_p\log p$.  If
$\sum_{p\ne q}|k_p|\log p\le B$, then every $k\ne0$ satisfies
\begin{equation}\label{p02:pt:eq:anchor-separation}
 |v(k)|\ge e^{-1-2B},
\end{equation}
independently of the size of $k_q$.
Consequently a real finite polynomial $F$ with positive mean $m$,
whose Fourier modes all obey that tail-height bound, gives a uniform
prime-flow hitting time
\begin{equation}\label{p02:pt:eq:anchor-certificate}
 T>2e^{1+2B}\|F\|_A/m
\end{equation}
for $\{F>0\}$.
\end{proposition}
\begin{proof}
If $|v(k)|>1$ there is nothing to prove.  Otherwise
$|k_q|\log q\le1+B$, and therefore
$\sum_p|k_p|\log p\le1+2B$.
Write $v(k)=\log(A/C)$ with coprime positive integers $A,C$.
Unique prime factorization gives $A\ne C$, and
$\max(A,C)\le e^{1+2B}$.  Since
$|\log(A/C)|\ge|A-C|/\max(A,C)$, the separation follows.
Apply Lemma~\ref{p02:pt:lem:sum-certificate} and
$S(F)\le e^{1+2B}\|F\|_A$ for the second assertion.
\end{proof}

For a product of the kernels in \eqref{p02:pt:eq:sine-kernel}, the norm
cost of the anchor is at most $D_q+1$.  Thus resolving that coordinate
to energy $O(\varepsilon)$ costs only $O(\log(1/\varepsilon))$ in
the logarithm of \eqref{p02:pt:eq:anchor-certificate}, even though its
degree is of order $\varepsilon^{-1/2}$.  This is a strict removal of
one artificial height cost; it is not a change in the known
$\varepsilon^{-1/2}$ power for a full power-weight spectrum.

\begin{proposition}[Quantitative return-section reduction]
\label{p02:pt:prop:return-section}
Index one distinguished coordinate by $0$, suppose $\lambda_0>0$,
and assume $V_2^2=\sum_{j\ge1}w_j\lambda_j^2<\infty$.
Let $P=2\pi/\lambda_0$, and let $Q(\varepsilon)$ be the least integer
$N$ such that the discrete tail translation with angles
$(2\pi\lambda_j/\lambda_0)_{j\ge1}$ hits its tail energy ball from
every starting point at some iterate $0\le n\le N$.  Then
\begin{align}
 R_{w,\lambda}(\varepsilon)&\le P\,[Q(\varepsilon)+1],
       \label{p02:pt:eq:return-upper}\\
 Q(C\varepsilon)&\le\lceil R_{w,\lambda}(\varepsilon)/P\rceil+1,
 \qquad
 C=\left(1+\frac{\pi V_2}{2\lambda_0\sqrt{w_0}}\right)^2.
       \label{p02:pt:eq:return-lower}
\end{align}
The inequalities are interpreted by approximation if an infimum is
not attained, and are trivial if the relevant time is infinite.
\end{proposition}
\begin{proof}
From any starting point the distinguished coordinate first reaches
zero at a time $t_0\in[0,P)$.  Its subsequent exact returns are
$t_0+nP$, and the tail at those times follows the stated discrete
translation.  This proves \eqref{p02:pt:eq:return-upper}.

For the other direction start the distinguished coordinate at zero,
with an arbitrary tail point, and take a continuous-time hit at $t$.
Choose the nearest nonnegative integer $n$ to $t/P$.  The elementary
circle estimate $\theta^2\le(\pi^2/2)(1-\cos\theta)$ gives
\[
 |t-nP|\le\frac\pi{\sqrt2\lambda_0}
                    \sqrt{\varepsilon/w_0}.
\]
In the tail metric $d^2=2E$, translation over a time $u$ moves a
point by at most $|u|V_2$.  The triangle inequality therefore bounds
the energy at $nP$ by
\[
 \frac12\left(\sqrt{2\varepsilon}
       +\frac{\pi V_2}{\sqrt2\lambda_0}
                              \sqrt{\varepsilon/w_0}\right)^2
 =C\varepsilon.
\]
The index is no larger than the stated integer bound.  Taking the
supremum over tail starting points proves the result.
\end{proof}

\paragraph{Why this does not close the exact-prime gap.}
Two free anchor degrees cannot both be removed from
\eqref{p02:pt:eq:anchor-separation}.  Indeed $\log2/\log3$ is irrational
by unique factorization, and the pigeonhole proof of Dirichlet
approximation produces nonzero integers $r,s$ with
$|r\log2-s\log3|\to0$.  If both primes were omitted from the height,
the proposed bound would be a positive constant for these forms,
which is false.  The return section likewise replaces one continuous
coordinate by a simultaneous discrete approximation problem; it
does not make the remaining coordinates independent.  Removing
finitely many primes leaves the power-weight tail in the same
asymptotic class.  Therefore neither proposition establishes
\[
 \log R_{w,\log p}(\varepsilon)
 \asymp\varepsilon^{-1/(\sigma-1)}
          (\log(1/\varepsilon))^{-\sigma/(\sigma-1)}
 \quad\text{for }w_p\asymp p^{-\sigma},\ \sigma>3.
\]
That matching entropy estimate remains an arithmetic question.

\section{Recovering the weighted phase geometry from one radial law}
\label{p02:pt:sec:radial-inverse}

The leading small-ball entropy in the preceding theorem records two
tail parameters.  The full radial distribution contains strictly more
information: it determines every weight, including finite
multiplicities.  This is an inverse statement about the product
geometry, not about its temporal frequencies.

\begin{theorem}[Exact radial rigidity of weighted circle products]
\label{p02:pt:thm:radial-rigidity}
Let $(w_j)$ and $(v_j)$ be finite or countably infinite positive
summable sequences.  If
\[
 \sum_jw_j(1-\cos\Theta_j)
 \quad\hbox{and}\quad
 \sum_jv_j(1-\cos\Theta_j)
\]
have the same distribution under independent uniform angles, then
their multisets of weights agree, including multiplicities.
Thus their weighted product-circle metric spaces are isometric by
a coordinate permutation.  In fact, equality of all even cumulants
already forces equality of the weight multisets.
\end{theorem}
\begin{proof}
Put $X_w=\sum_jw_j\cos\Theta_j$.  This is bounded, and its moment
generating function is
\[
 M_w(t)=\prod_j I_0(tw_j).
\]
The classical Bessel product~\cite[eq.~10.21.15]{p01:rare:DLMF} gives
\[
 I_0(z)=\prod_{r\ge1}(1+z^2/j_{0,r}^2),\qquad
 \log I_0(z)=\sum_{m\ge1}\frac{(-1)^{m+1}}m
                  \left(\sum_{r\ge1}j_{0,r}^{-2m}\right)z^{2m}
\]
near zero.  All coefficients in the latter series are nonzero.
Local uniform convergence, justified by $\sum_jw_j^2<\infty$,
therefore gives
\begin{equation}\label{p02:pt:eq:cumulant-spectrum}
 \kappa_{2m}(X_w)
 =(2m)!\frac{(-1)^{m+1}}m
        \left(\sum_{r\ge1}j_{0,r}^{-2m}\right)
        \sum_jw_j^{2m}.
\end{equation}
To see that the interchange is legitimate, take
$|t|<j_{0,1}/(2\sup_jw_j)$; the logarithmic products converge
absolutely there and are dominated by a constant times
$\sum_jw_j^2$.  Centering and changing sign do not change even
cumulants of order at least two.  Hence the radial law, or its even
cumulants alone, determines every number
$P_m=\sum_jw_j^{2m}$, $m\ge1$.

Now define the finite positive measure
\[
 \nu_w=\sum_jw_j^2\,\delta_{w_j^2}.
\]
Its moments of orders $0,1,\ldots$ are $P_1,P_2,\ldots$.
The measures associated with two candidate sequences are supported
in one common compact interval.  Polynomial approximation of
continuous functions shows that equality of these moments implies
equality of the measures.  At a positive atom $x$, its multiplicity
as a squared weight is exactly $\nu_w\{x\}/x$.  Zero has no atom
because the weights are positive.  This recovers the full multiset.
A coordinate permutation preserving weights preserves the metric,
which proves the geometric assertion.
\end{proof}

\begin{corollary}[A constructive descending recovery rule]
\label{p02:pt:cor:descending-recovery}
From the power traces in \eqref{p02:pt:eq:cumulant-spectrum}, the largest
weight $u$ and its multiplicity $r$ are
\[
 u=\lim_{m\to\infty}P_m^{1/(2m)},\qquad
 r=\lim_{m\to\infty}P_m/u^{2m}.
\]
Subtracting $r u^{2m}$ from each $P_m$ and repeating recovers every
distinct positive weight and its multiplicity.
\end{corollary}
\begin{proof}
Summability implies that the largest positive weight is attained
and has finite multiplicity.  The inequalities
$u^{2m}\le P_m\le u^{2m-2}P_1$ prove the first limit.  For the second,
dominated convergence applies to
$\sum_j(w_j/u)^{2m}$, dominated for $m\ge1$ by the summable series
$\sum_j(w_j/u)^2$.  After deleting the largest weights the same
argument applies.  There are only finitely many weights above any
positive threshold, so descending recovery reaches every weight.
\end{proof}

\begin{proposition}[Leading entropy is not a complete invariant]
\label{p02:pt:prop:entropy-not-rigid}
Changing finitely many weights in a sequence satisfying
\eqref{p02:pt:eq:fast-weights} does not change the leading small-ball
constant \eqref{p02:pt:eq:fast-small-ball}.  Such a change can nevertheless
change the radial law and the metric space.
\end{proposition}
\begin{proof}
The changed sequence still satisfies \eqref{p02:pt:eq:fast-weights} with
the same $a,b$, so Proposition~\ref{p02:pt:prop:fast-small-ball} applies
with the same constant.  Increase, for example, one largest weight
strictly.  The multiset changes, and so does $P_1$ and the variance
in \eqref{p02:pt:eq:cumulant-spectrum}.  Thus the radial law changes.
The diameter squared of the weighted product metric is $4\sum_jw_j$,
which also increases, so the two metric spaces are not isometric.
\end{proof}

This inverse theorem extends the earlier manuscript's use of small-ball
entropy as a coarse scale.  It needs the full law, or infinitely many
exact cumulants.  It provides no finite-data stability estimate and
does not identify the frequencies: Haar radial laws are unchanged
when the temporal frequency vector changes.


% ===== Consolidated body: zero_process =====
\section{What a zero pattern remembers: phases, spectra, and modular ordinates}
\label{p02:zp:sec:scope}

The objects in this section are zero divisors in a region of absolute
convergence, with vertical translation as the dynamics. They are not the
nontrivial zeros of the Riemann zeta function. In particular, the Euler
product below is zero-free; its critical points are the zeros under study.
Three questions organize the argument.
\begin{enumerate}
\item \emph{An inverse question formulated here.} Does the critical-point
pattern near the right edge retain the infinitely many prime phases, or
does passing from a function to its zeros destroy some of this information?
Theorem~\ref{p02:zp:thm:inverse} proves complete recovery and identifies the
entire dynamical system, not just a factor of it.
\item \emph{A published modular-density question.} Podesta
\cite[Section~1, paragraphs following the average vertical zero density]{p02:zp:Podesta2026} asks whether, for an exponential sum with rationally
independent positive exponents, there exist lengths $L$ for which the
imaginary parts of its zeros are dense modulo $L$. The paper states that
a solution will be presented in a future publication. We prove the
stronger equidistribution statement for every $L$ outside an explicit
countable resonance set. This answers the existence question as printed;
it does not classify the resonant lengths.
\item \emph{A special case of an external research program.} Feldheim's
research statement \cite[p.~1, paragraph beginning with the unresolved
atomic case]{p02:zp:FeldheimStatement} asks for an understanding of the random
limiting zero density when a stationary Gaussian analytic function has
atoms in its spectral measure. Her paper
\cite[Theorem~1 and Remark~2.3]{p02:zp:FeldheimDensity} proves randomness of
that limit and describes the two-atom example. We compute the conditional
density explicitly for three frequencies with irrational ratio, and
express the random limit and its quadratic variance coefficient through
three independent Rayleigh variables. This is a special-case computation,
not a solution for arbitrary atomic spectra.
\end{enumerate}
The general relation between almost periodicity and pure point spectra
is classical; see \cite{p02:zp:LenzSpindelerStrungaru}. The results below are
proved directly to isolate the additional inverse statement and the
precise scope of the applications. No historical priority claim is made
for the elementary torus arguments or the three-term intensity formula.
A targeted search through 14 October 2026 located the cited modular
question but no subsequent primary-source resolution of that exact
formulation; this is a bounded search, not a priority certificate.

\subsection{Compact analytic families and their zero divisors}

Let $\mathbb T=\{z\in\mathbb C:|z|=1\}$ and
$\Omega=\prod_p\mathbb T$, with product topology and Haar probability $m$.
Put
\[
 R_t\omega=(\omega_p p^{-it})_p,\qquad
 Q_\omega(s)=\sum_p\frac{(\log p)\omega_pp^{-s}}{1-\omega_pp^{-s}}
            =-\frac{Z_\omega'(s)}{Z_\omega(s)},\qquad \Re s>1.
\]
The product $Z_\omega=\prod_p(1-\omega_pp^{-s})^{-1}$ never vanishes
there. All sums and their derivatives converge locally uniformly,
uniformly in $\omega$. Thus $\omega\mapsto Q_\omega$ is continuous for
locally uniform convergence, and
$Q_{R_t\omega}(s)=Q_\omega(s+it)$. No $Q_\omega$ is identically zero:
as $\sigma\to+\infty$, its first Dirichlet term is
$(\log2)\omega_2 2^{-\sigma}$ and the remaining terms are uniformly
smaller in relative size.

For a vertical open strip $D$ compactly contained horizontally in
$\{\Re s>1\}$, denote the zero divisor by
\[
 \mathcal Z_D(\omega)=\sum_{\rho\in D:\ Q_\omega(\rho)=0}
       \operatorname{mult}_{Q_\omega}(\rho)\,\delta_\rho.
\]
We use the vague topology of locally finite measures on $D$, and let
$T_t$ translate a divisor down by $it$. A point carries its real
coordinate as a mark. Forgetting that mark is a separate operation.

\begin{lemma}[Continuity and uniform local zero bounds]
\label{p02:zp:lem:divisor-continuity}
The map $\mathcal Z_D$ is continuous and satisfies
$\mathcal Z_D(R_t\omega)=T_t\mathcal Z_D(\omega)$.
For every compact real interval $I\subset(1,\infty)$ there is $C_I<\infty$
such that, uniformly in $\omega$ and $u\in\mathbb R$, the number of zeros
in $I+i[u,u+1]$, counted with multiplicity, is at most $C_I$.
The same statements hold for any compact continuous family of
nonzero holomorphic functions with the stated translation covariance.
\end{lemma}
\begin{proof}
Suppose $f_j\to f\not\equiv0$ locally uniformly. Around the finitely
many zeros relevant to a compactly supported continuous test function,
choose disjoint small disks with zero-free boundaries. Rouche's theorem
preserves the number of zeros, with multiplicity, in each disk. On the
remaining compact set the limit has a positive lower modulus, so no
additional zeros occur. Uniform continuity of the test function then
gives convergence of its integrals against the divisors. This argument
does not require simple zeros. It proves continuity of $\mathcal Z_D$;
covariance follows from the analytic identity.

For a fixed compact rectangle, continuity and compactness give a uniform
bound on zero counts: use a nonnegative compactly supported test function
equal to one on the rectangle. A slightly larger open strip may be used
when the rectangle touches the boundary of the original $D$. Translation
covariance gives the same bound at every height.
\end{proof}

\subsection{The edge is a complete phase code}

Let $E>1$ be the unique solution
\begin{equation}\label{p02:zp:eq:edge}
 \frac{\log2}{2^E+1}=\sum_{p\ge3}\frac{\log p}{p^E-1},
 \qquad z_2^*=-1,\quad z_p^*=1\ (p\ge3),
 \qquad w_p=(\log p)p^{-E}.
\end{equation}
For the untwisted product $E=2.813014\ldots$ is the classical right edge
for zeros of $\zeta'$; its value and qualitative phase freezing appear
in \cite[Theorems~5.1 and~6.1]{p02:zp:AriasLune}. The following short argument
also establishes the uniform input needed for the inverse problem.

\begin{lemma}[Uniform edge recovery estimate]
\label{p02:zp:lem:edge-code}
There exist $\epsilon_0,C>0$ such that every zero
$\rho=\sigma+i\gamma$ of $Q_\omega$ with $E-\epsilon_0<\sigma\le E$
satisfies
\begin{equation}\label{p02:zp:eq:recovery-error}
 \sum_p w_p\left|\omega_p-z_p^*p^{i\gamma}\right|^2
       \le C(E-\sigma).
\end{equation}
There are no zeros with $\sigma>E$. For every $\omega$, there exist
zeros with real parts tending to $E$ and ordinates tending to $+\infty$.
\end{lemma}
\begin{proof}
For real $\sigma>1$, write $r_p=p^{-\sigma}$ and
\[
 c_p=\frac{(\log p)r_p^2}{1-r_p^2},\quad
 R_p=\frac{(\log p)r_p}{1-r_p^2},\quad
 \xi_p=\frac{z_p-r_p}{1-r_pz_p}.
\]
For $|z_p|=1$ we have $|\xi_p|=1$ and
$(\log p)r_pz_p/(1-r_pz_p)=c_p+R_p\xi_p$.
Set $C_0=\sum_pc_p$, $b_0=R_2$, $R_0=\sum_{p\ge3}R_p$, and
\[
 \Delta(\sigma)=C_0+R_0-b_0
       =\sum_{p\ge3}\frac{\log p}{p^\sigma-1}
          -\frac{\log2}{2^\sigma+1}.
\]
The ratio of the first sum to the final positive term is strictly
decreasing from infinity to zero: the logarithmic derivative of its
$p$th summand is
$(\log2)2^\sigma/(2^\sigma+1)-(\log p)p^\sigma/(p^\sigma-1)<0$.
This proves uniqueness of $E$ and negativity of $\Delta$ to its right.
At a zero use $z_p=\omega_pp^{-i\gamma}$. The vector equation
\[
 C_0+b_0\xi_2=-\sum_{p\ge3}R_p\xi_p
\]
rules out $\Delta<0$. Close to $E$ it implies
\[
 C_0b_0|\xi_2+1|^2
 \le\Delta(\sigma)(R_0+b_0-C_0),\qquad
 \sum_{p\ge3}R_p|\xi_p-1|^2
       =2\Delta(\sigma)+b_0|\xi_2+1|^2.
\]
Here $C_0,b_0,R_0$ and $b_0-C_0$ stay positive and bounded away from
zero. Since $\Delta(\sigma)=O(E-\sigma)$, and the circle maps satisfy
\[
 |\xi-1|=\frac{(1+r)|z-1|}{|1-rz|},\qquad
 |\xi+1|=\frac{(1-r)|z+1|}{|1-rz|},
\]
uniform comparison and $R_p\ge w_p$ give
\eqref{p02:zp:eq:recovery-error}.

For existence, the frozen function is $Q_{z^*}(s)=\Delta(s)$.
Writing $h=(\log2)/(2^E+1)$ and
$a_p=(\log p)/(p^E-1)$ gives
\[
 \Delta'(E)=h\frac{(\log2)2^E}{2^E+1}
    -\sum_{p\ge3}a_p\frac{(\log p)p^E}{p^E-1}
    <h(\log2-\log3)<0.
\]
Thus $E$ is a simple frozen zero. Integer independence of the prime
logarithms makes every forward $R_t$ orbit dense in $\Omega$.
Choose $t_j\to\infty$ with $R_{t_j}\omega\to z^*$. Local uniform
convergence and Rouche's theorem produce zeros
$\rho_j=E+o(1)+it_j$ of $Q_\omega$. Their real parts cannot exceed $E$.
At equality the displayed defects force $R_{\Im\rho_j}\omega=z^*$;
there is at most one such ordinate, since two would give a nonzero
common period for all prime logarithms. Hence the sequence can be
chosen with $\Re\rho_j<E$.
\end{proof}

\begin{theorem}[Inverse phase recovery and exact zero-hull conjugacy]
\label{p02:zp:thm:inverse}
Fix $0<\epsilon<E-1$ and $\delta>0$, and put
$D=(E-\epsilon,E+\delta)+i\mathbb R$.
The map $\omega\mapsto\mathcal Z_D(\omega)$ is a homeomorphism of
$\Omega$ onto its image. It conjugates $R_t$ with vertical translation
of the zero pattern. Its image is the translation hull of the divisor
of the zeros of $\zeta'$ in $D$.

More explicitly, any sequence of zeros
$\rho_j=\sigma_j+i\gamma_j$ with $\sigma_j\to E$ recovers the phase by
\begin{equation}\label{p02:zp:eq:phase-limit}
 \omega_p=z_p^*\lim_{j\to\infty}p^{i\gamma_j}\qquad(p\text{ prime}).
\end{equation}
A single such zero recovers it with the weighted error
\eqref{p02:zp:eq:recovery-error}. If two twists have sequences of zeros with
real parts tending to $E$ and differences of ordinates tending to zero,
then the twists are identical.
\end{theorem}
\begin{proof}
Since each $w_p>0$, \eqref{p02:zp:eq:recovery-error} implies
\eqref{p02:zp:eq:phase-limit}, independently of the chosen sequence.
If two divisors in $D$ coincide, apply this formula to a sequence given
by Lemma~\ref{p02:zp:lem:edge-code}; their phases coincide coordinatewise.
The divisor map is therefore continuous and injective. Its domain is
compact and the vague measure space is Hausdorff, so it is a
homeomorphism onto its image. Covariance proves conjugacy. The orbit of
$\omega=(1)_p$ is dense, so its divisor has precisely this image as
its translation hull. Finally $p^{i(\gamma_j-\gamma'_j)}\to1$ for each
fixed prime, which proves the last assertion by the same recovery formula.
\end{proof}

\begin{remark}[Prime sums and the unavoidable scalar ambiguity]
\label{p02:zp:rem:prime-sum}
For $F_\omega(s)=\sum_p a_p\omega_pp^{-s}$ with positive bounded
amplitudes, assume that there exists $S>1$ satisfying the edge equation
$a_2 2^{-S}=\sum_{p\ge3}a_pp^{-S}$. The same argument recovers the
relative phases $\eta_p=-\omega_p/\omega_2$, $p\ge3$:
\[
 \eta_p=\lim_j(p/2)^{i\gamma_j}.
\]
Indeed, normalization gives
$1-\sum_{p\ge3}b_p\eta_p(p/2)^{-(s-S)}$ with
$\sum b_p=1$. At a zero $\sigma+i\gamma$ near $S$, taking real parts
and using a first-order expansion in $S-\sigma$ bounds
$\sum b_p|\eta_p(p/2)^{-i\gamma}-1|^2$ by $C(S-\sigma)$.
The frozen zero is simple, and relative prime logarithms are integer
independent, so the existence and compactness proof applies verbatim.
A common scalar multiplying all $\omega_p$ leaves the divisor unchanged
and cannot be recovered. The resulting hull is the relative phase torus,
with frequencies generated by $\log(p/2)$, not the full Euler phase torus.
\end{remark}

\subsection{Exact spectra, correlations, and what hyperuniformity still needs}

Write $k=(k_p)_p\in\mathbb Z^{(\mathcal P)}$ for a finitely supported
integer vector, and put
\[
 \chi_k(\omega)=\prod_p\omega_p^{k_p},\qquad
 \lambda_k=\sum_pk_p\log p.
\]
The characters form an orthonormal basis of $L^2(\Omega,m)$ and
$\chi_k(R_t\omega)=e^{-i\lambda_kt}\chi_k(\omega)$.
Distinct $k$ give distinct $\lambda_k$, by unique factorization.

\begin{theorem}[The spectrum of the marked edge-zero process]
\label{p02:zp:thm:spectrum}
The Haar random marked divisor of Theorem~\ref{p02:zp:thm:inverse} is
stationary and ergodic, with simple pure point dynamical spectrum
\[
                 \{\log q:q\in\mathbb Q_{>0}\}.
\]
It is not mixing. For any continuous observable $A$ of that divisor,
write
\[
 a_k=\int A(\mathcal Z_D(\omega))\overline{\chi_k(\omega)}\,dm.
\]
Its centered covariance is exactly
\begin{equation}\label{p02:zp:eq:covariance}
 C_A(t)=\sum_{k\ne0}|a_k|^2e^{-i\lambda_kt}.
\end{equation}
This series converges absolutely and uniformly. The same correlations
are deterministic vertical averages for every initial twist:
\begin{align*}
 C_A(t)=\lim_{T\to\infty}\frac1T\int_0^T
 &(A(\mathcal Z_D(R_{u+t}\omega))-a_0)\\
 &\qquad\overline{(A(\mathcal Z_D(R_u\omega))-a_0)}\,du.
\end{align*}
In particular, randomizing a long vertical shift of the actual
$\zeta'$ edge divisor converges to this Haar law.
\end{theorem}
\begin{proof}
The inverse theorem transfers the character basis to the divisor
space, giving exactly the stated spectrum. No nonconstant character
has zero frequency, proving ergodicity. A character has correlation
$e^{-i\lambda_kt}$, which does not tend to zero, so mixing fails.
Parseval gives \eqref{p02:zp:eq:covariance} and
$\sum_{k\ne0}|a_k|^2<\infty$, proving uniform absolute convergence.

For a continuous function on $\Omega$, finite character sums approximate
it uniformly by Stone--Weierstrass. The average of every nontrivial
character along $R_u$ tends to zero uniformly in its starting point.
Thus all continuous observables have uniform Haar time averages.
Apply this first to the product in the displayed correlation, and
then to an arbitrary continuous test on the divisor hull.
\end{proof}

\begin{proposition}[Bragg weights and an exact variance test]
\label{p02:zp:prop:variance}
Fix a compact real interval $I\subset(1,\infty)$ and project the zeros
in $I+i\mathbb R$ to their ordinates, retaining multiplicity. Let this
random measure be $\nu_\omega$, and write
$N_T=\nu_\omega([0,T])$, $\rho=\mathbb E N_1$.
There are constants $c_k\in\mathbb C$, with $c_0=\rho$, such that
\begin{equation}\label{p02:zp:eq:bragg-form}
 \operatorname{Var}\!\left(\int\psi\,d\nu_\omega\right)
   =\sum_{k\ne0}|c_k|^2|\widehat\psi(\lambda_k)|^2,
 \qquad \widehat\psi(\lambda)=\int e^{-i\lambda t}\psi(t)\,dt,
\end{equation}
for compactly supported bounded measurable $\psi$. Consequently
\begin{equation}\label{p02:zp:eq:number-variance}
 \operatorname{Var}N_T
 =\sum_{k\ne0}|c_k|^2
             \frac{4\sin^2(\lambda_kT/2)}{\lambda_k^2}.
\end{equation}
Define $B(r)=\sum_{0<|\lambda_k|\le r}|c_k|^2$. Then
\begin{align}
 \operatorname{Var}N_T=o(T)
   &\quad\Longleftrightarrow\quad B(r)=o(r)\quad(r\downarrow0),
       \label{p02:zp:eq:hyperuniform-test}\\
 \sup_{T>0}\operatorname{Var}N_T<\infty
   &\quad\Longleftrightarrow\quad
        \sum_{k\ne0}\frac{|c_k|^2}{\lambda_k^2}<\infty.
       \label{p02:zp:eq:bounded-variance-test}
\end{align}
The same statements hold for a bounded nonnegative real-coordinate
weight of compact support in place of $1_I$.
\end{proposition}
\begin{proof}
Uniform local zero bounds imply that $\nu_\omega$ is uniformly
translation bounded. Define the complex locally finite measures
\[
 \nu_k(B)=\int_\Omega\nu_\omega(B)\overline{\chi_k(\omega)}\,dm(\omega).
\]
Covariance and Haar invariance give
$\nu_k(B+u)=e^{-i\lambda_ku}\nu_k(B)$.
Multiplication by $e^{i\lambda_kt}$ makes this a translation-invariant
measure, so $d\nu_k(t)=c_ke^{-i\lambda_kt}dt$. Thus the $k$th Fourier
coefficient of $\int\psi\,d\nu_\omega$ is
$c_k\widehat\psi(\lambda_k)$. Parseval proves
\eqref{p02:zp:eq:bragg-form}, including indicator functions, and then
\eqref{p02:zp:eq:number-variance}.

The measure $\sum_{k\ne0}|c_k|^2\delta_{\lambda_k}$ is finite near
zero, because the Fourier transform of a sufficiently narrow positive
test function is bounded away from zero there. Also
\[
 \sum_{|\lambda_k|\ge1}\frac{|c_k|^2}{\lambda_k^2}<\infty.
\]
Indeed, integrate \eqref{p02:zp:eq:number-variance} over $0<T<1$; its left
side is uniformly bounded, whereas the integral of
$4\sin^2(\lambda T/2)$ is bounded below by a positive constant for
$|\lambda|\ge1$.

For $|\lambda|T\le1$ the sine factor bounds the variance below by
$cT^2B(1/T)$. This proves necessity in
\eqref{p02:zp:eq:hyperuniform-test}. Conversely suppose $B(r)=o(r)$.
The part $|\lambda|\le1/T$, divided by $T$, is at most $TB(1/T)$.
For $1/T<|\lambda|<\delta$, use dyadic shells and
$B(r)\le\eta r$ for $r\le2\delta$ to bound the same quotient by
$C\eta\sum_{j\ge0}2^{-j}$. The part $|\lambda|\ge\delta$ tends
to zero by the established inverse-square summability away from zero.
Let $\eta\downarrow0$.

Inverse-square summability directly implies bounded variance.
Conversely, if the variance is bounded by $M$, average over $0<T<R$.
For every nonzero fixed $\lambda$ the average of
$4\sin^2(\lambda T/2)$ tends to $2$. Fatou's lemma gives
$2\sum_{k\ne0}|c_k|^2/\lambda_k^2\le M$.
\end{proof}

The measure in \eqref{p02:zp:eq:bragg-form}, including the atom
$\rho^2\delta_0$, is the pure point diffraction measure in the
Fourier convention used there. Individual $c_k$ may vanish. Complete
dynamical spectrum does not assert a nonzero Bragg peak at every
$\log q$: nonlinear observables can recover frequencies absent from
a particular linear count. Nor does pure point spectrum by itself
prove \eqref{p02:zp:eq:hyperuniform-test}. Estimating these low-frequency
weights for the Euler critical-point process remains unresolved here.

Every nonempty projected process of positive intensity is nevertheless
nonmixing. Otherwise its bounded local-count observables would have
decaying correlations; their pure point covariance expansions force
them all to be constant. That would make the random measure deterministic
and invariant under every real translation, hence a multiple of
Lebesgue measure, contrary to its nonzero locally finite atomic nature.
Thus its law cannot equal a positive-intensity Poisson process or the
sine-kernel determinantal process. For example their compact-test
covariances tend to zero, by the Poisson covariance formula and the
decay of the squared sine kernel, respectively. This is a fixed-strip,
fixed-scale statement. It imposes no obstruction to a different limit
under a changing strip or rescaling, and none to the GUE conjecture for
critical-line zeta zeros.

\subsection{A resolution of the generic-period modular-density question}

\begin{theorem}[Equidistribution of exponential-sum zeros modulo a length]
\label{p02:zp:thm:podesta}
Let
\[
 f(s)=A_0+\sum_{j=1}^N A_je^{\omega_js},\qquad
 A_j\ne0,\quad 0<\omega_1<\cdots<\omega_N,
\]
where $\omega_1,\ldots,\omega_N$ are linearly independent over
$\mathbb Q$. Count zeros with multiplicity. If
\begin{equation}\label{p02:zp:eq:nonresonance}
 \frac{2\pi}{L}\notin\operatorname{span}_{\mathbb Q}
                       \{\omega_1,\ldots,\omega_N\},
\end{equation}
then, for every interval $J\subset[0,L)$,
\begin{equation}\label{p02:zp:eq:modular-equidistribution}
 \frac{\#\{\rho:f(\rho)=0,\ 0<\Im\rho\le T,
                              \Im\rho\bmod L\in J\}}
      {\#\{\rho:f(\rho)=0,\ 0<\Im\rho\le T\}}
       \longrightarrow\frac{|J|}{L}.
\end{equation}
The denominator is $(\omega_N/(2\pi))T+o(T)$.
Both conclusions are uniform when only the arguments of the $A_j$
vary and their nonzero moduli are fixed. In particular, the ordinates
are dense modulo $L$, and the assertion holds for all but countably
many $L>0$.
\end{theorem}
\begin{proof}
Multiplication by a nonzero constant does not change zeros, so write
the phase family as
$f_\theta(s)=a_0+\sum_{j=1}^Na_j\theta_je^{\omega_js}$ with $a_j>0$.
The constant term dominates uniformly for sufficiently negative real
part, and the final term dominates uniformly for sufficiently positive
real part. Hence every member has all zeros in one common bounded real
strip. None is identically zero. Lemma~\ref{p02:zp:lem:divisor-continuity}
applies, with flow $\theta_j\mapsto\theta_je^{i\omega_jt}$, and gives
a uniform bound on the number of zeros in every interval of unit height.

Choose $g\in C_c(\mathbb R)$ equal to one throughout the common real
zero strip, and $\eta\in C_c^\infty(\mathbb R)$. The local observable
\[
 A_\eta(\theta)=\sum_{f_\theta(\rho)=0}
                     g(\Re\rho)\eta(\Im\rho)
\]
is continuous on $\mathbb T^N$. For any
$\alpha\notin\Gamma:=\sum_j\mathbb Z\omega_j$, character approximation
gives, uniformly in $\theta$,
\begin{equation}\label{p02:zp:eq:twisted-average}
 \frac1T\int_0^T e^{i\alpha t}A_\eta(R_t\theta)\,dt\longrightarrow0.
\end{equation}
Indeed the integral of each approximating character has frequency
$\alpha+\sum_jk_j\omega_j\ne0$.
Let $H$ contain the support of $\eta$. By a change of variable for
each zero, and the uniform zero bound in the two endpoint regions,
\begin{equation}\label{p02:zp:eq:unsmoothing}
 \int_0^T e^{i\alpha t}A_\eta(R_t\theta)\,dt
 =\widehat\eta(\alpha)
       \sum_{\substack{f_\theta(\rho)=0\\0<\Im\rho\le T}}
                       e^{i\alpha\Im\rho}+O_\eta(1).
\end{equation}
The error is uniform in $T$ and the initial phases. Choose $\eta$ with
$\widehat\eta(\alpha)\ne0$, which is possible by taking a positive
function supported in a sufficiently short interval. Equations
\eqref{p02:zp:eq:twisted-average}--\eqref{p02:zp:eq:unsmoothing} prove that
the twisted zero sum is $o(T)$.

For completeness we compute the untwisted mean. Define
\[
 J(\sigma)=\int_{\mathbb T^N}\log|f_\theta(\sigma)|\,d\theta.
\]
Circle Jensen integration and the positive constant term imply
$\log a_0\le J(\sigma)\le\log\sum_ja_je^{\omega_j\sigma}$,
where $\omega_0=0$. On compact real intervals these bounds also
give integrability of the absolute value of the logarithm, since
its positive part is uniformly bounded. Poincare--Lelong and Fubini
show that the mean zero measure is
$(2\pi)^{-1}J''(\sigma)\,d\sigma\,dt$, with $J''$ understood as a
measure. Where the constant term dominates, the convergent logarithm
series has no constant phase character, so $J=\log a_0$. Where the
final term dominates, the same argument gives
$J=\log a_N+\omega_N\sigma$. Therefore the total mean number per unit
height is $\rho=\omega_N/(2\pi)$.
Taking $\alpha=0$ and $\int\eta=1$ in
\eqref{p02:zp:eq:unsmoothing}, unique ergodicity gives
$N(T)=\rho T+o(T)$ uniformly in the phases.

Condition~\eqref{p02:zp:eq:nonresonance} says exactly that
$2\pi m/L\notin\Gamma$ for every nonzero integer $m$.
The established twisted sums, divided by $N(T)$, are all the Weyl
tests for uniform distribution modulo $L$. Trigonometric majorants
and minorants of interval indicators prove
\eqref{p02:zp:eq:modular-equidistribution}, including uniformity in the
initial phases. The excluded rational span is countable.
\end{proof}

There is also a short proof of the existence-and-density conclusion
alone. A frozen phase choice has a simple real zero, obtained from
$a_0=\sum_{j\ge1}a_je^{\omega_j\sigma}$. Under
\eqref{p02:zp:eq:nonresonance}, the joint flow of the $N$ phases and
$t\bmod L$ is dense. Approximate the frozen phase and any desired
residue simultaneously, then use Rouche's theorem near that real
zero. This observation is sufficient for the existence question in
\cite{p02:zp:Podesta2026}; the longer argument supplies uniform distribution
and the exact density.

The nonresonance condition is sufficient, not asserted necessary.
For instance $1+e^{\omega s}$ has zeros on an arithmetic progression,
so appropriate resonant lengths fail even density. In the Euler
edge process there are explicit resonant obstructions: if
$L_p=2\pi/\log p$, \eqref{p02:zp:eq:recovery-error} confines ordinates of
zeros with $E-\epsilon<\Re\rho\le E$ modulo $L_p$ to an arc of
length $O_p(\sqrt\epsilon)$, centered according to
$\omega_p/z_p^*$. Thus the exceptional periods are mathematically
meaningful. Recovering the full classification of resonant $L$ is
outside Theorem~\ref{p02:zp:thm:podesta}.

\subsection{Three atomic Gaussian frequencies: an explicit conditional law}

The phase hull describes randomness generated by phases. Gaussian
coefficients add invariant random amplitudes, and this distinction
changes the macroscopic variance. We make the separation explicit.
Let $0<a<b$ with $a/b\notin\mathbb Q$ and
\[
       F(s)=\xi_0+\xi_1e^{-as}+\xi_2e^{-bs},
\]
where the $\xi_j$ are independent standard circular complex Gaussians.
Write $R_j=|\xi_j|$; their densities are $2r e^{-r^2}\,dr$ on
$(0,\infty)$, and their phases are independent uniform variables.
Conditioning on $R=(R_0,R_1,R_2)$ and removing a common phase leaves
a two-torus Kronecker system.

For fixed positive $r=(r_0,r_1,r_2)$ set
\[
 A=r_0,\qquad B=r_1e^{-a\sigma},\qquad C=r_2e^{-b\sigma},
 \qquad
 D_r(\sigma)=a^2B^2+b^2C^2+ab(A^2-B^2-C^2).
\]
When $A,B,C$ are the sides of a strict triangle, let
\[
 \mathcal A_r(\sigma)=\frac14
 \sqrt{(A+B+C)(-A+B+C)(A-B+C)(A+B-C)}
\]
be its area, and define
\begin{equation}\label{p02:zp:eq:triangle-intensity}
 I_r(\sigma)=\frac{D_r(\sigma)}{4\pi^2\mathcal A_r(\sigma)}
       \quad\hbox{in the strict triangle region},
 \qquad I_r(\sigma)=0\quad\hbox{outside its closure}.
\end{equation}
Values at its finitely many boundary points may be assigned arbitrarily.

\begin{theorem}[Atomic conditional density and random limiting measure]
\label{p02:zp:thm:three-atoms}
For every fixed positive $r$, the mean zero measure conditional on
$R=r$ has planar density $I_r(\sigma)$, independent of height. It has
no mass on the boundary lines of the strict triangle region.
For every finite real interval $J$ and every choice of the coefficient
phases,
\begin{equation}\label{p02:zp:eq:atomic-limit}
 \frac1T\#\{\rho:F(\rho)=0,\ \Re\rho\in J,\ 0<\Im\rho\le T\}
       \longrightarrow M_r(J):=\int_J I_r(\sigma)\,d\sigma.
\end{equation}
Consequently the Gaussian limiting real-coordinate measure has the
explicit law $I_R(\sigma)\,d\sigma$. For a bounded interval of positive length
$J$, writing $N_T(J)$ for the numerator above,
\begin{equation}\label{p02:zp:eq:atomic-variance}
 \lim_{T\to\infty}\frac{\operatorname{Var}N_T(J)}{T^2}
   =\operatorname{Var}_R M_R(J)>0.
\end{equation}
The expectation and variance on the right are ordinary integrals
against $\prod_{j=0}^2 2r_je^{-r_j^2}dr_j$.
\end{theorem}
\begin{proof}
At a fixed point $s=\sigma+it$, absorb $t$ into the two uniform phases.
The value of the function is
$A+B e^{i\theta}+C e^{i\phi}$.
In the strict triangle region its zeros as a function of
$(\theta,\phi)$ are exactly two noncollinear configurations. At either
configuration the absolute real Jacobian is
\[
 |\operatorname{Im}(B e^{i\theta}\overline{C e^{i\phi}})|
       =2\mathcal A_r(\sigma).
\]
Also
\[
 |F'(s)|^2=|-aB e^{i\theta}-bC e^{i\phi}|^2=D_r(\sigma)>0,
\]
because $2\operatorname{Re}(B e^{i\theta}\overline{C e^{i\phi}})
=A^2-B^2-C^2$.
The change-of-variables formula for these two local phase branches
therefore gives the zero intensity
\[
 \frac1{(2\pi)^2}\sum_{\text{two branches}}
       \frac{|F'(s)|^2}{2\mathcal A_r(\sigma)}
       =I_r(\sigma).
\]
This is the two-dimensional Kac--Rice calculation, and can equivalently
be obtained by applying the area formula on compact subsets where the
two displayed Jacobians are nonzero. If a side is strictly larger than
the sum of the others, cancellation is impossible.

The degenerate-triangle equations are real exponential polynomials
with three distinct exponents. Each has finitely many real roots
(at most two), by Rolle's theorem after division by an extreme
exponential. At each such real coordinate there is only one phase
configuration giving cancellation, modulo $2\pi$.
For bounded height, the set of initial phases which reaches this
configuration is a smooth curve in the two-torus and has Haar measure
zero. A countable union over height intervals still has measure zero.
Thus the mean zero measure has no mass on these boundary lines, and
\eqref{p02:zp:eq:triangle-intensity} accounts for all of it.

For fixed amplitudes the divisor map is continuous on the compact
two-torus and the flow is uniquely ergodic. Smooth real-coordinate
test functions consequently have deterministic limiting densities
given by the mean zero measure, using the smoothing and bounded-endpoint
argument from Theorem~\ref{p02:zp:thm:podesta}. Since this mean measure is
absolutely continuous, continuous majorants and minorants at the
endpoints of $J$ extend the limit to $1_J$. This proves
\eqref{p02:zp:eq:atomic-limit} for every initial phase, not just almost
every one.

To pass to Gaussian variances, a uniform integrability argument is
needed. Normalize any nonzero coefficient vector to have Euclidean
norm one. The normalized vectors form a compact sphere; because the
three exponentials are linearly independent, none gives an identically
zero function. Lemma~\ref{p02:zp:lem:divisor-continuity} gives a zero bound
on $J+i[0,1]$ uniform even when some coefficients vanish. Vertical
translation preserves this sphere. Therefore $N_T(J)\le C_J(T+1)$
for every coefficient vector, and dominated convergence proves
\eqref{p02:zp:eq:atomic-variance} with a nonnegative constant.
It is strictly positive: there is an open set of amplitudes on which
the constant term dominates throughout $J$, so $M_R(J)=0$; another
open set gives a strict triangle at an interior point of $J$, and then
$M_R(J)>0$. Both open sets have positive Rayleigh probability.
\end{proof}

Feldheim already proves that an atomic spectral component gives a
positive quadratic variance coefficient in this Gaussian setting
\cite[Theorem~1]{p02:zp:FeldheimVariance}. Equation~\eqref{p02:zp:eq:atomic-variance}
is not claimed as a new growth law; its contribution here is the
explicit conditional coefficient supplied by
\eqref{p02:zp:eq:triangle-intensity}. Conditional on the amplitudes, the
limiting density is deterministic and the phase dynamics is pure point.
Unconditionally, the random invariant amplitudes produce the quadratic
term. One must therefore not transfer that quadratic term to the
fixed-amplitude Euler phase hull.

These arguments give a complete answer to the new inverse question,
an affirmative answer to the printed generic-period existence question
with a stronger uniform-distribution theorem, and an explicit
three-frequency instance of the atomic-density program. They leave
the actual Euler Bragg weights, its fixed-strip hyperuniformity, the
resonant-period classification, and general atomic Gaussian spectra
undetermined.


\clearpage
\section{Connections, limits, and the next mathematical bottlenecks}
\label{p02:end:connections}

\subsection{One graph gives two genuinely different invariants}

Theorems~\ref{p02:graph:threshold} and \ref{p02:bud:thm:lower} can be applied
to the same covariance model, but they answer different questions.
For an $m$-vertex covariance graph $G$, a compact parameter box
of full dimension has
\[
 k=2m+|E(G)|,\qquad
 q_*=\operatorname{gcr}(G)+1,\qquad
 \mathcal R_B\asymp B^{-1/(k+2)}.
\]
Here $q_*$ is the exact number of shared-latent observations needed
to identify an arbitrary mean/covariance law from one fixed group law.
The dimension $k$ governs the rate at which an arbitrary mixing
law can be learned under a total budget with growing replication depths.
This distinction is operational: infinitely many groups below $q_*$
leave a positive risk floor, whereas unequal depths growing with $B$
give consistency. A graph's likelihood-existence threshold is a third
quantity. Existing computations for complete bipartite graphs show
that it cannot replace $q_*$.

Thus the transfer produces an exact answer without computing every
graph invariant. A substantive further advance would give new bounds
on $\operatorname{gcr}(G)$, or prove a relation to the maximum likelihood
threshold on a class where that relation is presently unknown.
Transform uniqueness alone cannot supply those combinatorial bounds.

\subsection{Shape constraints do not determine the information content}

Strong log-concavity and arbitrarily high Gevrey regularity coexist
with the nonidentifiability constructions in
Theorems~\ref{p02:graph:threshold}, \ref{p02:mt:phase-cantor}, and
\ref{p02:mt:heat-criterion}. The three failures have different causes:
a polynomial relation among transform probes; reflection of independent
convolution factors; and a summable analytic zero set.
These are explicit exact counterexamples, not examples of numerical
instability.

On the other hand a quasianalytic restriction changes the heat problem
completely: Proposition~\ref{p02:mt:heat-qa} identifies sources from any
sequence of positive times tending to zero. Finite moments and very
strong smoothness do not replace quasianalyticity. This identifies
which qualitative prior changes the uniqueness argument, instead of
merely improving constants in an estimate.

\subsection{A probability class can be excluded by a mass budget}

For the McKay classifications the decisive inverse object is a Thorin
measure. If a candidate GGC law has an odd logarithmic derivative
about a center, the representing measure must be symmetric about that
center. Endpoint singularities and a central asymptotic then force
specific atoms. Their sum cannot exceed the total mass.
For the squared-Bessel family beyond the claimed range, this becomes
the explicit contradiction
\[
                  6\mu\le4\mu+1\qquad(\mu>1/2).
\]
For the third family it becomes $\nu\le\mu+1$.
The positive half of each classification constructs a measure satisfying
the same constraints through a classical hypergeometric ratio.
This paired construction and obstruction is stronger than verifying
a finite list of derivative signs.

The unresolved SD extension requires a different argument. A law can
be self-decomposable without being GGC; a failed positive Thorin
representation cannot by itself exclude self-decomposability.
One must study the monotonicity of its Levy kernel, or find another
valid SD criterion, on the remaining domain.

\subsection{Full observations, coarse invariants, and the arithmetic gap}

The exact radial law determines every metric weight, while its leading
small-ball entropy retains only a coarse asymptotic scale.
The near-edge critical-point divisor determines every Euler phase,
while its projected count spectrum can have missing Bragg peaks.
These inverse statements specify what is lost by each further
compression of the observations.

The random recurrence theorem does not settle its deterministic
prime-frequency analogue. Its proof uses independent anti-concentration
to control a sum of inverse divisors. At one prescribed vector of
prime logarithms that probabilistic step is unavailable.
A theorem holding on a set of full measure cannot be substituted for
an arithmetic estimate at that vector. Likewise, the exact criterion
for hyperuniformity in \eqref{p02:zp:eq:hyperuniform-test} is not a proof
that the Euler zero process satisfies it: the low-frequency Bragg
weights still have to be estimated.

\subsection{Four concrete continuation questions}

The following problems isolate missing steps without claiming that
they have been solved.
\begin{enumerate}
\item \textbf{The scalar endpoint boundary.}
For arbitrary mixing of $\operatorname{Uniform}[0,\theta]$,
$\theta\in[1,2]$, is the total-budget $W_1$ rate $B^{-1/2}$, or is
a logarithmic factor necessary? The present proved interval is
\[
 cB^{-1/2}\le\mathcal R_B^{\rm ad}
 \le\mathcal R_B^{\rm fix}\le CB^{-1/2}\log(eB).
\]
The lower bound for this particular experiment must improve, or the
critical multilevel summation must be bypassed.
\item \textbf{Sparse terminal lengths without covariance bounds.}
For $d$-output terminal experiments with unbounded covariance laws,
is harmonic divergence also necessary, or does multivariate coupling
change the exact sparse criterion? The bounded-covariance case is
settled, and harmonic divergence is sufficient without boundedness.
\item \textbf{Resonant modular lengths.}
Theorem~\ref{p02:zp:thm:podesta} settles all nonresonant lengths.
Determine which resonant lengths still have density or equidistribution.
The arithmetic-progression example and Euler edge phase confinement
show why a universal assertion over resonant lengths would be false.
\item \textbf{Stability of phase recovery from perturbed zeros.}
Theorem~\ref{p02:zp:thm:inverse} gives an exact inverse and a weighted
single-zero error bound. A finite-height statistical recovery theorem
would also have to control how close to the edge one can get, how many
zeros are available, and the sensitivity to uncertainty in ordinates.
The deterministic arithmetic return problem is part of that difficulty.
\end{enumerate}

\subsection{Validation and interpretation of this draft}

Each main branch has been read by a second reasoning agent in addition
to its author; the assembled statements were also reviewed for matching
assumptions, transform signs, endpoints, and scope. These are internal
checks, not independent human peer review or a formal proof assistant
verification. Numerical checks, where used, served only to detect
normalization mistakes; no numerical scan replaces an analytic proof.

The source audit establishes connections to printed questions, not
historical priority. That requires specialist review, including older
exponential-family completeness theory for the Gaussian branch.



\clearpage
\begingroup
\small

\endgroup
\endgroup
% ===== END SOURCE 02 =====


% ===== BEGIN SOURCE 03: KEIO_Binomial_Likelihood_and_Exact_Observation_Thresholds.tex =====
\clearpage
\part{Binomial Mixture Likelihood and Exact Observation Thresholds}
\label{part:03}
\setcounter{section}{0}
\begin{center}\large\itshape A Conjectured Range, Sparse Gaussian Sampling, and Resonant Zero Patterns\end{center}
\begin{center}\small Source manuscript: 14 October 2026.\end{center}
\begingroup
\renewcommand{\arraystretch}{1}
\begin{quote}\small\textbf{Source abstract.}\quad
For an arbitrary mixing law on $[0,1]$, we prove that every
nonparametric maximum likelihood estimator from $N$ independent
binomial counts with $t$ trials has expected Wasserstein error
$O_\varepsilon((t\log N)^{-1/2})$ throughout
$\log N\le t\le N^{2/3-\varepsilon}$. This is the range conjectured
by Vinayak, Kong, Valiant and Kakade in 2019. Two estimates enter:
a contour integral removes a polynomial prefactor from the
Chebyshev--Bernstein change of basis, and an aggregate ellipse bound
gives finite likelihood brackets of logarithmic cardinality
$O(\sqrt t\log^{3/2}(t/\delta))$. We derive consequences for
replication caps, and develop separate exact results for sparse
Gaussian sampling, self-decomposability on Bessel parameter families,
and resonant periods of three-term exponential sums. Further results
identify a deterministic small-divisor criterion for recurrence and
prove limitations of finite continuous phase calibration and noisy
analytic heat observations. The research register distinguishes
published conjectures, new questions, partial classifications, and
unresolved arithmetic obstacles. Earlier results are not counted as
new resolutions. Historical priority of individual auxiliary lemmas
is not asserted.
\end{quote}
\thispagestyle{empty}


\section{What is proved, and what the research programme did not close}
\label{p03:intro:scope}

This manuscript continues the 223-page integrated working paper
\cite{p01:KEIOIntegrated}, its 42-page sequel \cite{comp:part:01}, and
the 56-page sequel \cite{comp:part:02}. The new work starts at the
unresolved boundaries of those drafts. None of their completed
results is included in the present count of research targets.
The most directly externally testable conclusion is
Theorem~\ref{p03:bnew:thm:mle}: it has the same unrestricted mixing-law
MLE, support interval, loss, and trial-count regime as the original
2019 conjecture. It does not replace the estimator or change the
sampling design.

The improvement is a change in the range of a theorem, not a
constant-factor gain. The originally proved exponent for the upper
trial count was $2/9$; the later argument of Musco et al. reaches
$1/4$; the theorem here reaches $2/3$, with an arbitrarily small
fixed loss in each exponent. The target Wasserstein error stays
$(t\log N)^{-1/2}$. The still larger near-linear range discussed in
the 2026 author version of Musco et al. remains open here.

The other sections answer different questions about observation
maps. They should not be represented as consequences of one
universal theorem. The Gaussian section uses holomorphic functions
on a positive matrix tube to construct a single joint invisible
perturbation. The Bessel section uses polynomial root locations to
make Levy kernels explicit. The zero section uses a lifted torus
and triangle geometry to retain information lost by a nonresonant
equidistribution argument. Topology and analytic approximation give
negative results for two inverse problems with strong shape
constraints.

\subsection{Interpretation and validation}

All the mathematical statements below are statements of this draft,
not a claim of community certification. The main proofs have been
checked by a second reasoning agent and by the assembling reviewer;
these checks are internal, not external peer review or formal proof
verification. We give full arguments so that the results can be
independently assessed. The literature comparison uses specific
primary-source formulations available on 14 October 2026 and does
not certify an exhaustive historical search.

The goal of an advance comparable to major work in analysis is a
research direction, not a conclusion justified by ambition or by
page count. In particular, no Riemann-hypothesis consequence,
unrestricted McKay self-decomposability classification, or
deterministic prime-flow recurrence asymptotic is asserted.

\subsection{Research register}

The register at the end records the precise scope of each attempted
route and its outcome. A theorem for a new internal question is not
counted as a solution of a previously published conjecture; two
lemmas in one proof are not counted as two conjecture resolutions.


% Consolidated section: budget_new
\section{A wider validity range for the binomial-mixture maximum likelihood estimator}
\label{p03:bnew:sec:mle}

Consider $N$ independent counts $X_i$ obtained by drawing
$\theta_i\sim G$ and then $X_i\mid\theta_i\sim\operatorname{Binomial}(t,\theta_i)$,
where $G$ is an arbitrary probability measure on $[0,1]$. Put
\[
 B_j^t(x)=\binom tj x^j(1-x)^{t-j},\qquad
 p_G(j)=\int B_j^t(x)\,dG(x),\quad 0\leq j\leq t.
\]
Let $\widehat G$ be any maximum likelihood solution over all probability
measures on $[0,1]$. Such a solution exists: the attainable probability
vectors form the compact convex hull of the continuous binomial curve,
and the likelihood is upper semicontinuous. Nonuniqueness of the mixing
measure is allowed throughout.

Vinayak, Kong, Valiant and Kakade
\cite[Theorem~3.2 and Remark~3.2]{p01:info:VKVK2019} proved the medium-sample
rate below for $t\leq N^{2/9-\varepsilon}$ and conjectured the range
$t\leq N^{2/3-\varepsilon}$. Musco et al.
\cite[Theorem~20 and Section~6.2, author version]{p01:info:MMRS2025}
increased the proved exponent to $1/4$ and discussed further conditional
improvements. The following argument reaches the originally conjectured
range. It uses a bound with unspecified absolute constants for
Chebyshev--Bernstein coefficients; it does not prove the sharper
constant-one coefficient inequality discussed in those papers.

\begin{theorem}[The conjectured medium-sample MLE range]
\label{p03:bnew:thm:mle}
For every $0<\varepsilon<2/3$ there are constants $C_\varepsilon,N_\varepsilon$
such that, whenever $N\geq N_\varepsilon$ and
\begin{equation}\label{p03:bnew:eq:range}
 \log N\leq t\leq N^{2/3-\varepsilon},
\end{equation}
every mixing-measure MLE satisfies, uniformly over $G$,
\[
 \mathbb P_G\left\{W_1(\widehat G,G)
             >\frac{C_\varepsilon}{\sqrt{t\log N}}\right\}
 \leq N^{-2},\qquad
 \mathbb E_G W_1(\widehat G,G)
 \leq\frac{C_\varepsilon}{\sqrt{t\log N}}.
\]
The same conclusion, with constants also depending on $a>0$, holds if
the lower limit in \eqref{p03:bnew:eq:range} is $a\log N$.
\end{theorem}

The matching lower bound in \cite[Theorem~3.3]{p01:info:VKVK2019}
therefore makes the MLE minimax optimal throughout this range.

The proof separates two estimates that control different directions of
the observation map: an aggregate analytic coefficient bound makes its
image statistically small, whereas a contour bound controls polynomial
inversion without a polynomial prefactor in $t$.

\subsection{A contour estimate without a degree prefactor}

Write $\widetilde T_m(x)=T_m(2x-1)$ and define the coefficients
$C(t,m,j)$ by
\[
 \widetilde T_m(x)=\sum_{j=0}^t C(t,m,j)B_j^t(x),\qquad m\leq t.
\]

\begin{lemma}[A uniform coefficient bound]
\label{p03:bnew:lem:bernstein}
There is an absolute $C$ such that, for $0\leq m\leq t/2$,
\begin{equation}\label{p03:bnew:eq:bernstein}
 \max_{0\leq j\leq t}|C(t,m,j)|\leq C\exp(2m^2/t).
\end{equation}
The bound includes $j=0,t$; no restriction away from the parameter
endpoints is imposed.
\end{lemma}
\begin{proof}
The assertion is immediate for $m=0$ and for $j=0,t$, where the
coefficients have absolute value one. For $0<j<t$, set
\[
 r=\frac j{t-j},\qquad q=\frac r{(1+r)^2},\qquad z=re^{i\phi}.
\]
The polynomial
\[
 F(z)=(1+z)^t\widetilde T_m\left(\frac z{1+z}\right)
      =\sum_{j=0}^t\binom tj C(t,m,j)z^j
\]
has the algebraic expression
\begin{equation}\label{p03:bnew:eq:F-polynomial}
 F(z)=\frac{(1+z)^{t-m}}2
       \left\{(\sqrt z+i)^{2m}+(\sqrt z-i)^{2m}\right\}.
\end{equation}
Although a square root is used in this identity, the sum is a
polynomial in $z$ and is independent of its choice. The identity follows
from $T_m(w)=[(w+\sqrt{w^2-1})^m+(w-\sqrt{w^2-1})^m]/2$.

For $u=2\sqrt q\,|\sin(\phi/2)|\in[0,1]$, one has
\[
 |1+z|=(1+r)\sqrt{1-u^2},\qquad
 \max\{|\sqrt z+i|^2,|\sqrt z-i|^2\}=(1+r)(1+u).
\]
Consequently, using $m\leq t/2$,
\begin{align*}
 \frac{|F(re^{i\phi})|}{(1+r)^t}
 &\leq(1-u^2)^{(t-m)/2}(1+u)^m\\
 &\leq\exp(-tu^2/4+mu)
 \leq\exp(2m^2/t)\exp(-tu^2/8).
\end{align*}
The last step is the elementary maximum
$\sup_{u\in\mathbb R}\{mu-tu^2/8\}=2m^2/t$.
Because $\sin(|\phi|/2)\geq|\phi|/\pi$ for $|\phi|\leq\pi$,
\[
 \frac1{2\pi}\int_{-\pi}^\pi
       \exp\{-tq\sin^2(\phi/2)/2\}\,d\phi
 \leq\frac{C}{\sqrt{1+tq}}.
\]
Cauchy's coefficient formula now yields
\begin{equation}\label{p03:bnew:eq:cauchy-coeff}
 |C(t,m,j)|\leq
 \frac{(1+r)^t}{\binom tj r^j}\,
 \frac{C\exp(2m^2/t)}{\sqrt{1+tq}}.
\end{equation}
For completeness the binomial probability at its mean obeys
\[
 \binom tj(j/t)^j(1-j/t)^{t-j}
       \geq c\sqrt{\frac{t}{j(t-j)}}
       \geq\frac{c}{\sqrt{1+tq}}.
\]
The first inequality follows directly by applying upper and lower
Stirling bounds $c\sqrt n(n/e)^n\leq n!\leq C\sqrt n(n/e)^n$
to $t!,j!,(t-j)!$; here $j,t-j\geq1$.
This probability equals $\binom tjr^j/(1+r)^t$ and $tq=j(t-j)/t$.
It cancels the square-root factor in \eqref{p03:bnew:eq:cauchy-coeff},
proving \eqref{p03:bnew:eq:bernstein}.
\end{proof}

The weighted square-summability of the coefficients used next is
explicitly developed in \cite[Lemma~13]{p01:info:MMRS2025}; we give the
short periodic-function argument and the approximation step to make
the dependence on constants self-contained.

\begin{lemma}[Lipschitz approximation with bounded coefficient sum]
\label{p03:bnew:lem:jackson}
For every $1$-Lipschitz $f:[0,1]\to\mathbb R$ with $f(0)=0$ and
integer $k\geq2$, there is
$Q_k=\sum_{m=0}^k a_m\widetilde T_m$ such that
\[
 \|Q_k-f\|_\infty\leq C/k,\qquad
 \sum_{m=0}^k|a_m|\leq C.
\]
\end{lemma}
\begin{proof}
The even periodic function $g(\phi)=f((1+\cos\phi)/2)$ has
$\|g\|_\infty\leq1$ and almost-everywhere derivative bounded by $1/2$.
If $g=\sum_{m\geq0}c_m\cos(m\phi)$ in $L^2$, Parseval applied
to its weak derivative gives $\sum_{m\geq1}m^2|c_m|^2\leq C$.
Hence $|c_0|+\sum_{m\geq1}|c_m|\leq C$ by Cauchy--Schwarz.

Put $s=\lfloor k/2\rfloor+1$. Normalize the nonnegative even
trigonometric polynomial
\[
 J_s(u)=c_s\left(\frac{\sin(su/2)}{\sin(u/2)}\right)^4
\]
to integral one on $[-\pi,\pi]$. It has degree $2(s-1)\leq k$,
and its first absolute moment is at most $C/s$. To check this, the
unnormalized integral is comparable to $s^3$: on $|u|\leq1/s$ the
quotient is comparable to $s$, while everywhere it is at most
$C\min(s,|u|^{-1})$. The same bound gives an unnormalized first
moment at most $Cs^2$.
Convolution $J_s*g$ therefore differs from $g$ uniformly by $C/k$.
It is even of degree at most $k$; each cosine coefficient is the
corresponding $c_m$ multiplied by a Fourier coefficient of a probability
kernel, of absolute value at most one. Replacing $\cos(m\phi)$ by
$\widetilde T_m(x)$ proves the assertion.
\end{proof}

For probability vectors put
$h(p,q)^2=\sum_j(\sqrt{p(j)}-\sqrt{q(j)})^2$.

\begin{corollary}[An inverse modulus]
\label{p03:bnew:cor:inverse}
For all mixing laws $G,H$ and $2\leq k\leq t/2$,
\begin{equation}\label{p03:bnew:eq:inverse}
 W_1(G,H)\leq\frac Ck+C\exp(2k^2/t)h(p_G,p_H).
\end{equation}
\end{corollary}
\begin{proof}
Represent the polynomial from Lemma~\ref{p03:bnew:lem:jackson} in the
degree-$t$ Bernstein basis. Lemma~\ref{p03:bnew:lem:bernstein} bounds
every resulting coefficient by $C\exp(2k^2/t)$. Therefore
\[
 \left|\int f\,d(G-H)\right|
 \leq C/k+C\exp(2k^2/t)\|p_G-p_H\|_1.
\]
Use $\|p-q\|_1\leq2h(p,q)$ and Kantorovich--Rubinstein duality.
\end{proof}

\subsection{An aggregate ellipse bound and finite likelihood brackets}

Expand the whole probability vector in shifted Chebyshev polynomials:
\[
 B_j^t(x)=\sum_{m=0}^t a_{jm}\widetilde T_m(x).
\]

\begin{lemma}[Aggregate Chebyshev decay]
\label{p03:bnew:lem:ellipse}
For $m\geq1$,
\begin{equation}\label{p03:bnew:eq:aggregate}
 \sum_{j=0}^t|a_{jm}|\leq2e^{-m^2/(2t)};
 \qquad \sum_{j=0}^t|a_{j0}|=1.
\end{equation}
\end{lemma}
\begin{proof}
On the shifted Bernstein ellipse
$z=(2+w+w^{-1})/4$, $|w|=e^u$, its focal-distance identity is
$|z|+|1-z|=\cosh u$. Thus
\[
 \sum_{j=0}^t|B_j^t(z)|=(|z|+|1-z|)^t
              =(\cosh u)^t\leq e^{tu^2/2}.
\]
The Laurent coefficient of $w^m$ in $B_j^t((2+w+w^{-1})/4)$ is
$a_{jm}/2$ for $m\geq1$. Cauchy's formula, followed by summation
over $j$ inside the integral, gives
$\sum_j|a_{jm}|\leq2\exp(tu^2/2-mu)$.
Choose $u=m/t$. For $m=0$, the real cosine integral expresses every
$a_{j0}$ as a nonnegative average of $B_j^t$; summing gives one.
\end{proof}

\begin{lemma}[Small upper brackets for all binomial mixtures]
\label{p03:bnew:lem:brackets}
For $0<\delta<1/2$, there is a finite collection $\mathcal U_\delta$
of nonnegative vectors such that every $p_G$ is coordinatewise bounded
above by some $u\in\mathcal U_\delta$ with
\[
 1\leq\sum_j u(j)\leq1+\delta,
\]
and
\begin{equation}\label{p03:bnew:eq:entropy}
 \log|\mathcal U_\delta|
 \leq C\sqrt t\,[\log(Ct/\delta)]^{3/2}.
\end{equation}
\end{lemma}
\begin{proof}
Set
\[
 K=\min\left\{t,
   \left\lceil\sqrt{2t\log(16t/\delta)}\right\rceil\right\}.
\]
Every Chebyshev moment
$\mu_m(G)=\int\widetilde T_m\,dG$ lies in $[-1,1]$, and $\mu_0=1$.
Round $\mu_1,\ldots,\mu_K$ onto a grid of mesh
$\xi=\delta/(8K)$, giving values $\nu_m$ within $\xi$ of the
respective moments. For each possible grid vector define
\[
 v(j)=a_{j0}+\sum_{m=1}^Ka_{jm}\nu_m,
 \qquad
 e(j)=\xi\sum_{m=1}^K|a_{jm}|+\sum_{m=K+1}^t|a_{jm}|.
\]
Then $|v(j)-p_G(j)|\leq e(j)$ for the law whose moments were rounded.
Lemma~\ref{p03:bnew:lem:ellipse} gives
\[
 \sum_j e(j)\leq2K\xi+2t e^{-(K+1)^2/(2t)}\leq\delta/2,
\]
where the tail is zero when $K=t$.
Take $u(j)=\max\{v(j)+e(j),0\}$. For each covered law, $u\geq p_G$
and $u-p_G\leq2e$ coordinatewise, so its total mass lies in
$[1,1+\delta]$. Discard grid vectors whose upper vectors fail this
mass condition. This does not discard the covering vector of any law.
The number of grid vectors is at most $(CK/\delta)^K$. The definition
of $K$ implies \eqref{p03:bnew:eq:entropy} (also when $K=t$).
\end{proof}

\begin{theorem}[Predictive Hellinger accuracy of the full MLE]
\label{p03:bnew:thm:hellinger}
There is an absolute $C$ such that, for $0<\eta<1/2$,
\begin{equation}\label{p03:bnew:eq:hellinger}
 \mathbb P_G\left\{
 h(p_{\widehat G},p_G)^2>
 C\frac{\sqrt t[\log(CNt)]^{3/2}+\log(1/\eta)}{N}
 \right\}\leq\eta.
\end{equation}
This result holds for every $N,t\geq1$, with a vacuous bound allowed
when its right-hand threshold exceeds the maximal squared distance two.
\end{theorem}
\begin{proof}
Here is a direct finite-bracket likelihood argument. Fix a prospective
error $0<\rho\leq\sqrt2$ and set $\delta=\rho^2/100$.
For an upper bracket $u\geq p$, let $S=\sum u\in[1,1+\delta]$
and $q=u/S$. Then
$\|q-p\|_1\leq2\delta$, hence $h(q,p)\leq\sqrt{2\delta}$.
If $h(p,p_G)\geq\rho$, it follows that $h(q,p_G)\geq3\rho/4$.
The affinity identity and $1-x\leq e^{-x}$ imply
\[
 \sum_j\sqrt{q(j)p_G(j)}
       =1-\tfrac12h(q,p_G)^2\leq e^{-9\rho^2/32}.
\]
The MLE has likelihood at least that of the truth. If a bad MLE lies
under $u$, then $q\geq p_{\widehat G}/(1+\delta)$ and therefore
\[
 \prod_{i=1}^N\frac{q(X_i)}{p_G(X_i)}\geq(1+\delta)^{-N}.
\]
Zero true probabilities cause no difficulty: almost surely no such
coordinate is observed, and the displayed expectation calculation is
over the support of $p_G$. Markov's inequality applied to the square
root likelihood ratio bounds the probability for this fixed bracket by
\[
 (1+\delta)^{N/2}
 \left(\sum_j\sqrt{q(j)p_G(j)}\right)^N
 \leq e^{-N\rho^2/4}.
\]
A union bound over Lemma~\ref{p03:bnew:lem:brackets} proves
\[
 \mathbb P_G\{h(p_{\widehat G},p_G)\geq\rho\}
 \leq\exp\{C\sqrt t[\log(Ct/\rho^2)]^{3/2}-N\rho^2/4\}.
\]
For $\rho^2$ a sufficiently large absolute multiple of the threshold
in \eqref{p03:bnew:eq:hellinger}, one has $\rho^2\geq1/N$ and
$\log(Ct/\rho^2)\leq\log(CNt)$, proving the claim. If this choice
exceeds two, the theorem is automatic.
\end{proof}

\begin{proof}[Proof of Theorem~\ref{p03:bnew:thm:mle}]
Use Theorem~\ref{p03:bnew:thm:hellinger} with $\eta=N^{-2}$. Throughout
\eqref{p03:bnew:eq:range}, it gives, except on an event of probability
$N^{-2}$,
\[
 h(p_{\widehat G},p_G)
 \leq C\frac{t^{1/4}(\log N)^{3/4}}{\sqrt N}.
\]
Choose $k=\lfloor\sqrt{c\varepsilon t\log N}\rfloor$ for a fixed
sufficiently small absolute $c>0$, for example $c=1/100$.
For all large $N$, $2\leq k\leq t/2$. The inverse modulus gives
\[
 W_1(\widehat G,G)
 \leq\frac{C_\varepsilon}{\sqrt{t\log N}}
 +C N^{2c\varepsilon}
       \frac{t^{1/4}(\log N)^{3/4}}{\sqrt N}.
\]
Multiplying the second term by $\sqrt{t\log N}$ bounds it by
\[
 C N^{2c\varepsilon-1/2}t^{3/4}(\log N)^{5/4}
 \leq C N^{-(3/4-2c)\varepsilon}(\log N)^{5/4},
\]
which is bounded as $N\to\infty$. This proves the probability bound.
Since the diameter of $[0,1]$ is one, the exceptional event contributes
at most $N^{-2}$ to expected $W_1$, proving the expectation bound.
For a lower limit $a\log N$, decrease $c$ depending on $a$ to ensure
$k\leq t/2$; the rest of the argument is unchanged.
\end{proof}

\subsection{Scope of the resolution}

Theorem~\ref{p03:bnew:thm:mle} concerns the original unrestricted mixing-law
MLE, not a moment estimator or a changed replication design. The
parameter support includes both endpoints and arbitrary singular laws.
It proves the range stated in Remark~3.2 of \cite{p01:info:VKVK2019},
provided the preceding self-contained argument withstands external
review. The larger near-linear range discussed conditionally in
\cite[Section~6.2, author version]{p01:info:MMRS2025} is not claimed.
Nor do we prove the dimension-free $N^{-1/2}$ fingerprint-discrepancy
inequality conjectured there: the predictive Hellinger bound here still
contains the factor $t^{1/4}$, apart from logarithms.
The generic contour, Jackson and likelihood-testing methods are
classical; historical novelty of the individual coefficient or entropy
lemmas is not asserted.

\section{Adaptive replication caps: an exact low-cap phase}
\label{p03:bnew:sec:caps}

In the binomial experiment, now count each individual Bernoulli toss
as one observation. A policy can start and revisit individuals
adaptively, has at most $B$ total tosses, and may toss each individual
at most $Q$ times. Write $R_{B,Q}$ for its minimax expected $W_1$ risk,
allowing arbitrary randomized policies and arbitrary probability-valued
estimators. The following concerns the full adaptive experiment, even
though its upper bound uses equal groups.

\begin{theorem}[Exact capped-budget rate in a polynomial range]
\label{p03:bnew:thm:caps}
For each $0<\gamma<2/5$, uniformly over integers
$1\leq Q\leq B^{2/5-\gamma}$ and all sufficiently large $B$,
\begin{equation}\label{p03:bnew:eq:cap-rate}
 R_{B,Q}\asymp_\gamma
 \max\left\{\frac1{Q+1},\frac1{\sqrt{Q\log(eB)}}\right\}.
\end{equation}
The lower bound, with an absolute constant, holds for every
$1\leq Q\leq B$ without the polynomial-range restriction.
\end{theorem}

\begin{lemma}[Explicit moment-matching laws]
\label{p03:bnew:lem:finite-difference}
For every integer $K\geq0$ and interval of length $a$, there are two
probability laws on that interval whose moments through order $K$
agree and whose $W_1$ distance is $a/(K+1)$.
\end{lemma}
\begin{proof}
On $[0,1]$, split the signed measure
\[
 2^{-K}\sum_{j=0}^{K+1}(-1)^j\binom{K+1}{j}
                         \delta_{j/(K+1)}
\]
into its positive and negative parts. Each part has total mass one.
The $(K+1)$st finite difference annihilates polynomials of degree at
most $K$. Between consecutive support points the absolute cumulative
signed mass is $2^{-K}\binom Kj$, by the identity
$\sum_{i=0}^j(-1)^i\binom{K+1}{i}=(-1)^j\binom Kj$.
Integration of the absolute cumulative mass gives $1/(K+1)$.
An affine change of variables proves the general assertion.
\end{proof}

\begin{proof}[Adaptive lower bound in Theorem~\ref{p03:bnew:thm:caps}]
Degree-$Q$ moment-matching laws have exactly the same law of a full
$Q$-toss vector from an individual. They are therefore indistinguishable
under every capped adaptive policy: prepare independent full vectors
when each individual is first queried, and reveal the requested prefixes.
Lemma~\ref{p03:bnew:lem:finite-difference} and the triangle inequality
give $R_{B,Q}\geq1/[2(Q+1)]$.

For the second bound suppose first that
$Q\geq K:=\lceil4\log(eB)\rceil$ and put
\[
 h^2=\frac{K}{64eQ}.
\]
Take the two laws from Lemma~\ref{p03:bnew:lem:finite-difference} on
$[1/2-h,1/2+h]$, with moments matched through order $K$.
Under the fair-coin law for the full $Q$-toss vector, write its signs
as independent Rademacher variables $Z_1,\ldots,Z_Q$.
For a bias $1/2+a$, its likelihood ratio is
\[
 \prod_{i=1}^Q(1+2aZ_i)
 =\sum_{S\subseteq\{1,\ldots,Q\}}(2a)^{|S|}
                                      \prod_{i\in S}Z_i.
\]
Walsh orthogonality and moment matching show that the squared $L^2$
distance between the two mixture likelihood ratios is at most
\begin{align*}
 4\sum_{\ell>K}\binom Q\ell(4h^2)^\ell
 &\leq4\sum_{\ell>K}
           \left(\frac{4eQh^2}{\ell}\right)^\ell\\
 &\leq4\sum_{\ell>K}16^{-\ell}
 \leq C16^{-K}.
\end{align*}
Thus the total variation distance between their full-vector laws is at
most $C4^{-K}$. An oracle revealing $B$ independent full $Q$-toss
vectors dominates every allowed policy: at most $B$ individuals can
ever be queried, and each one's observed history can be read from its
own oracle vector. The total variation between the two oracle laws is
at most $CB4^{-K}\leq1/4$ for all sufficiently large $B$.
Their latent-law separation is $2h/(K+1)$, so the two-point testing
inequality gives $R_{B,Q}\geq c/\sqrt{Q\log(eB)}$.
If $Q<K$, the first bound already implies this second bound up to an
absolute constant. This proves the lower bound for all caps.
\end{proof}

\begin{proof}[Upper bound in Theorem~\ref{p03:bnew:thm:caps}]
Take $N=\lfloor B/Q\rfloor$ independent individuals, each observed
$t=Q$ times, and use their MLE. For $Q\geq c_0\log B$ with any
fixed $c_0>0$, the assumptions of Theorem~\ref{p03:bnew:thm:mle} hold
for all large $B$, after changing its fixed $\varepsilon$ and lower
constant. In fact $\log N\asymp_\gamma\log B$ and
$Q\leq B^{2/5-\gamma}$ implies
$Q\leq N^{2/3-\gamma}$ for all sufficiently large $B$, since
\[
 (2/5-\gamma)(5/3-\gamma)<2/3-\gamma.
\]
It follows that the expected risk is at most
$C_\gamma/\sqrt{Q\log B}$.

For $4\leq Q\leq c_0\log B$, use the inverse modulus with
$k=\lfloor Q/2\rfloor$. The Hellinger theorem with exceptional
probability $N^{-2}$ gives
\[
 \mathbb E W_1(\widehat G,G)
 \leq C/Q+C e^{Q/2}
          Q^{1/4}[\log(CNQ)]^{3/4}/\sqrt N+N^{-2}.
\]
Choose, for example, $c_0=1/10$. Since $N\asymp B/Q$, the second
term is $o(1/Q)$ uniformly in this range. The caps $Q\leq3$ are
handled by the diameter bound one. The two ranges prove
\eqref{p03:bnew:eq:cap-rate}.
\end{proof}

The fully uncapped regular-channel rate proved earlier is $B^{-1/3}$.
The new lower bound shows that a cap satisfying
$Q\gtrsim B^{2/3}/\log(eB)$ is necessary to attain that rate within
a fixed constant. The earlier multilevel allocation attains it with
$Q\asymp B^{2/3}$. This leaves a logarithmic difference between a
necessary and a sufficient depth. The intermediate region beyond the
exact range of Theorem~\ref{p03:bnew:thm:caps} is not declared solved.

\section{The scalar endpoint critical case: a precise limitation of the current method}
\label{p03:bnew:sec:uniform-critical}

For $P_\theta=\operatorname{Uniform}[0,\theta]$, $\theta\in[1,2]$,
the arbitrary-adaptive total-budget minimax bounds remain
\[
 cB^{-1/2}\leq R_B\leq CB^{-1/2}\log(eB).
\]
We do not close this gap here. The following calculation establishes a
specific obstruction to removing its logarithm merely by choosing a
better allocation for the same independent-level CDF estimator.

\begin{proposition}[A critical multilevel variance barrier]
\label{p03:bnew:prop:uniform-barrier}
Fix a component endpoint $\theta\in[1,2]$. At level $\ell\geq1$,
let $r_\ell=2^{\ell-1}$, couple the maxima $M_{r_\ell}$ and
$M_{2r_\ell}$ by using the same initial observations, and put
\[
 D_\ell(x)=\mathbf1_{\{M_{2r_\ell}\leq x\}}
                 -\mathbf1_{\{M_{r_\ell}\leq x\}},\qquad x\in[0,2].
\]
Use $n_\ell$ independent copies at level $\ell$, with independence
across levels. For any independent base-level estimator, the sum of
that base estimator and the level averages through $L$ has integrated
variance at least
\begin{equation}\label{p03:bnew:eq:barrier}
 \frac{cL^2}{B}
 \quad\text{if}\quad\sum_{\ell=1}^L2r_\ell n_\ell\leq B.
\end{equation}
No choice of these allocation integers removes the squared logarithm
in integrated variance when $L$ is of order $\log B$.
\end{proposition}
\begin{proof}
For $0\leq x\leq\theta$ write $z=x/\theta$. Monotonicity of the
maxima gives
\[
 \mathbb E D_\ell(x)=z^{2r}-z^r,
 \qquad \mathbb E D_\ell(x)^2=z^r-z^{2r},\quad r=r_\ell.
\]
Direct integration therefore gives
\begin{align*}
 \int_0^2\operatorname{Var}(D_\ell(x))\,dx
 =\theta\left\{
 \frac1{r+1}-\frac2{2r+1}+\frac2{3r+1}-\frac1{4r+1}\right\}
 \geq\frac c r.
\end{align*}
For an elementary bound without cancellation, integrate only over
$z\in[1-1/(4r),1-1/(8r)]$. On this interval $z^r$ stays in a
fixed compact subinterval of $(0,1)$, so
$\operatorname{Var}(D)=a(1-a)$ with $a=z^r-z^{2r}$ is bounded
below by an absolute positive constant. The interval length is
$1/(8r)$, proving the inequality.
Independence gives integrated variance at least
$c\sum_{\ell=1}^L(r_\ell n_\ell)^{-1}$.
Cauchy--Schwarz bounds this below by
$cL^2/\sum_\ell r_\ell n_\ell\geq cL^2/B$.
\end{proof}

This is an obstruction for a specified estimator architecture and for
integrated squared CDF loss. It is not an unrestricted minimax lower
bound in $W_1$: another estimator can exploit information not represented
by these independent increments, and squared $L^2$ loss cannot simply
be substituted for expected $L^1$ loss. In particular the proposition
does not justify declaring the critical endpoint logarithm necessary.
The familiar critical multilevel cost--variance phenomenon and its
randomized counterparts have substantial prior literature
\cite{p01:info:Giles2018,p03:bnew:RheeGlynn2015}; our contribution here is the
explicit calculation locating that obstruction in this endpoint model.


% Consolidated section: mixtures_new
\section{An exact sparse terminal-design criterion without covariance moments}
\label{p03:mixnew:sec:terminal}

The previous manuscript \cite[terminal-design theorem]{comp:part:02}
left one direction open for multivariate
Gaussian mixtures with unbounded covariance: harmonic divergence was
sufficient for a sequence of $d$-output terminal designs, but necessity
was proved only for finite design sets. We now prove necessity for
every convergent harmonic subseries. The new ingredient is a
Fourier--Laplace construction on the cone of positive semidefinite
matrices. It produces one joint signed density, not a separate formal
Laplace transform for each mean frequency.

Let $(A,\Sigma)$ have an arbitrary Borel law on
$\mathbb R^d\times\mathbb S_d^+$, and conditionally let $X_j$ be
iid $N_d(A,\Sigma)$. For $n\ge d$ put
\[
 \mathcal O_n=(X_1,\ldots,X_{d-1},S_n),\qquad
 S_n=X_1+\cdots+X_n.
\]
Different terminal lengths supply separate group laws. No joint law
between distinct terminal-length experiments is assumed given.

\begin{theorem}[Exact harmonic criterion for unrestricted covariance]
\label{p03:mixnew:thm:exact-terminal}
For every $d\ge1$ and every
$\mathcal N\subset\{d,d+1,\ldots\}$, the collection
$\{\mathcal L(\mathcal O_n):n\in\mathcal N\}$ identifies every
Borel mixing law if and only if
\[
 \sum_{n\in\mathcal N}\frac1n=\infty.
\]
Whenever the sum converges, there are distinct absolutely continuous
mixing laws with all polynomial moments finite and equal, with the
same full $A$ marginal and the same full $\Sigma$ marginal, that
produce the same observations for every selected $n$. Any prescribed
lower covariance bound $\Sigma\succeq b_0I$, $b_0>0$, can be
imposed. No bounded upper covariance support is asserted.
\end{theorem}

We give the analytic realization in detail. On $\mathbb S_d$ use
orthonormal coordinates for the Frobenius inner product
$\langle V,S\rangle=\operatorname{tr}(VS)$, and write
$p=d(d+1)/2$. All matrix Fourier integrals use the corresponding
Lebesgue measure. For a complex symmetric matrix $T=X+iV$ with
$X\succ0$, all square roots are principal analytic matrix roots.

\begin{lemma}[A matrix regularizer and a cone-supported inverse]
\label{p03:mixnew:lem:matrix-realization}
Let $F$ be a nonzero bounded holomorphic function on the right
half-plane, satisfying $F(\overline z)=\overline{F(z)}$.
Suppose $F$ and all its holomorphic derivatives extend continuously
and boundedly to the closed right half-plane. Suppose moreover that
its boundary values $F(iy)$ are a Schwartz function of
$y\in\mathbb R$, flat at zero.
Choose a nonzero real even
$h\in C_c^\infty(\mathbb R^d\setminus\{0\})$. Then
\begin{equation}\label{p03:mixnew:eq:matrix-transform}
 \mathcal H(u,T)=h(u)F(u^TT^{-1}u)
     \exp\{-\operatorname{tr}(T^{1/2})
                  -\operatorname{tr}(T^{-1/2})\}
\end{equation}
is the Fourier--Laplace transform
\[
 \mathcal H(u,T)=\int e^{iu^Ta-\operatorname{tr}(T\Sigma)}
                         g(a,\Sigma)\,da\,d\Sigma
\]
of a nonzero real Schwartz function $g$ on
$\mathbb R^d\times\mathbb S_d$, supported where $\Sigma\succeq0$.
Both coordinate marginals of $g$ vanish. Every joint polynomial
moment of $g$ is zero.
\end{lemma}
\begin{proof}
The inverse of an accretive matrix is accretive:
\[
 \operatorname{Re}(T^{-1})=T^{-*}XT^{-1}\succ0.
\]
Thus $\operatorname{Re}(u^TT^{-1}u)>0$ for every real $u\ne0$,
and the composition in \eqref{p03:mixnew:eq:matrix-transform} is
well defined and holomorphic in $T$.

We first prove a uniform decay estimate on the tube. The principal
square root of an accretive matrix is accretive. One direct
justification uses
\[
 T^{1/2}=\frac1\pi\int_0^\infty
                 t^{-1/2}T(tI+T)^{-1}\,dt.
\]
Indeed $\|t(tI+T)^{-1}\|<1$, so
$\operatorname{Re}\{I-t(tI+T)^{-1}\}\succ0$; the convergent
integral has positive real part. The same reasoning applies to
$T^{-1/2}$. Since $T$ and its square root are symmetric, write
$T^{1/2}=A+iB$ with real symmetric matrices $A,B$ and $A\succ0$.
Squaring gives
\[
 A^2-B^2=X\succ0,\qquad AB+BA=V.
\]
Consequently $\|B\|\le\|A\|$ in operator norm and
\[
 \|V\|\le2\|A\|\|B\|
          \le2(\operatorname{tr}A)^2.
\]
With $\Psi(T)$ denoting the exponential factor in
\eqref{p03:mixnew:eq:matrix-transform}, this proves the uniform bound
\begin{equation}\label{p03:mixnew:eq:tube-decay}
 |\Psi(X+iV)|\le
       \exp\{-\sqrt{\|V\|/2}\}\qquad(X\succ0).
\end{equation}
In particular its right side is integrable in all matrix-frequency
coordinates, independently of $X$.

Next consider the boundary $T=iV$ with $V$ invertible and real
symmetric. If its eigenvalues are $\lambda_1,\ldots,\lambda_d$,
\begin{equation}\label{p03:mixnew:eq:boundary-decay}
 |\Psi(iV)|=\exp\left\{-\frac1{\sqrt2}
       \sum_{j=1}^d\bigl(|\lambda_j|^{1/2}
                         +|\lambda_j|^{-1/2}\bigr)\right\}.
\end{equation}
All derivatives of the logarithm of $\Psi(iV)$ grow at most
polynomially in $1+\|V\|$ and $\|V^{-1}\|$. Here is a
functional-calculus justification that also handles repeated
eigenvalues. Put $\delta=\min_j|\lambda_j|$. Take contours
surrounding the spectrum of the normal matrix $iV$, at distance
comparable to $\delta$, grouped into unions of disks if necessary.
They avoid zero and the negative real axis. The resolvent bound
$\|(zI-iV)^{-1}\|\le\operatorname{dist}(z,i\operatorname{spec}V)^{-1}$,
the Cauchy integral formula for matrix functions, and its
differentiated formulas bound each derivative of $T^{\pm1/2}$
by a polynomial in $1+\|V\|$ and $\delta^{-1}$.

On this boundary $u^T(iV)^{-1}u=-iu^TV^{-1}u$ is purely
imaginary. The bounded boundary derivatives of $F$ and the rational
derivatives of $V^{-1}$ give the same polynomial derivative bounds
for its composition, uniformly for $u$ in the compact support of
$h$. At a zero of $u^TV^{-1}u$, use the flat extension of $F$
at zero. Formula \eqref{p03:mixnew:eq:boundary-decay} dominates all
these polynomial factors. It follows that the boundary function
\[
 H(u,V)=\mathcal H(u,iV)
\]
extends by zero across $\det V=0$ to a Schwartz function of
$(u,V)$, compactly supported in $u$ and flat on the singular set.
The same follows directly for each derivative from the displayed
bounds, since distance to the singular set is comparable to
$\delta$.

Let
\[
 g_u(\Sigma)=(2\pi)^{-p}\int e^{i\operatorname{tr}(V\Sigma)}
                                      H(u,V)\,dV.
\]
We prove its cone support, rather than merely invoking a tube
Paley--Wiener theorem. For $X\succ0$, define $g_{u,X}$ by the
same inverse Fourier integral with $\mathcal H(u,X+iV)$ in
place of $H$. Contour shifting inside the convex positive tube gives
\begin{equation}\label{p03:mixnew:eq:shift}
 g_{u,X}(\Sigma)=e^{-\operatorname{tr}(X\Sigma)}g_u(\Sigma).
\end{equation}
For clarity, this identity can be proved without a multidimensional
contour convention. Differentiate the inverse Fourier integral
along a straight path of positive definite real parts. The
Cauchy--Riemann equations and integration by parts in $V$ give
the differential equation
$\partial_Dg_{u,X}=-\operatorname{tr}(D\Sigma)g_{u,X}$.
On a compact path inside the cone, Cauchy estimates and
\eqref{p03:mixnew:eq:tube-decay} justify differentiation and make
all boundary terms vanish. Start at $X=\epsilon I$ and let
$\epsilon\downarrow0$; the same uniform bound gives dominated
convergence to $g_u$ for almost every invertible boundary $V$.
The exceptional singular matrices, and the possible zero set of
$u^TV^{-1}u$, have measure zero. At the latter set the prescribed
flat limit agrees as well.

Boundedness of $F$, compact support of $h$, and
\eqref{p03:mixnew:eq:tube-decay} bound $|g_{u,X}(\Sigma)|$ by a
constant independent of $X\succ0$. If $\Sigma$ is not positive
semidefinite, choose $X_0\succ0$ with
$\operatorname{tr}(X_0\Sigma)<0$. Taking $X=RX_0$ in
\eqref{p03:mixnew:eq:shift} and letting $R\to\infty$ proves
$g_u(\Sigma)=0$. Thus its support is in the required cone.

Finally set
\[
 g(a,\Sigma)=(2\pi)^{-d}\int e^{-iu^Ta}g_u(\Sigma)\,du.
\]
It is Schwartz by the boundary estimates. The symmetry
$H(-u,-V)=\overline{H(u,V)}$ makes it real. It is nonzero:
the scalar argument takes an interval of positive real values
when $T$ varies through positive scalar matrices, and $F$ is not
identically zero. Fourier inversion of \eqref{p03:mixnew:eq:shift}
proves the stated Fourier--Laplace identity, since $g_u$ is
Schwartz and supported on the cone.

The function $g_u$ vanishes near $u=0$. Fourier differentiation
therefore gives
$\int a^\gamma g(a,\Sigma)\,da=0$ for every multi-index
$\gamma$ and every $\Sigma$. This proves the vanishing covariance
marginal and every joint polynomial moment. For the other marginal,
let $T=\epsilon I$. The exponential factor is
$\exp\{-d(\sqrt\epsilon+\epsilon^{-1/2})\}$, while $F$ is
bounded. Thus $\mathcal H(u,\epsilon I)\to0$ for every $u$.
Dominated convergence against the integrable $g$, followed by
Fourier uniqueness in $a$, proves the vanishing mean marginal.
\end{proof}

\begin{lemma}[Scalar input with prescribed sparse zeros]
\label{p03:mixnew:lem:scalar}
If $\sum_{n\in\mathcal N}n^{-1}<\infty$, the function
\[
 F(z)=\left\{\prod_{n\in\mathcal N}\frac{2n-z}{2n+z}\right\}
                    \exp\{-\sqrt z-z^{-1/2}\},
 \qquad\operatorname{Re}z>0,
\]
satisfies all hypotheses of Lemma~\ref{p03:mixnew:lem:matrix-realization}
and vanishes at every $2n$, $n\in\mathcal N$. Empty and finite
sets are allowed.
\end{lemma}
\begin{proof}
The factors differ from one by $O_K(1/n)$ on each compact set
of the right half-plane. The product is therefore nonzero
holomorphic, bounded by one, and invariant under conjugation.
These assertions hold continuously at every finite point of the
closed half-plane. To control derivatives uniformly there, write
$b_n(z)=(2n-z)/(2n+z)$. For $\operatorname{Re}z\ge0$,
$|b_n(z)|\le1$, and direct differentiation gives
$|b_n^{(r)}(z)|\le C_r n^{-r}$ for $r\ge1$.
The product rule for finite products, together with
$\sum_n n^{-1}<\infty$, bounds every derivative uniformly in
the number of factors and in the entire closed half-plane.
The locally uniform differentiated products then give the same
bounds for the infinite product, including its continuous
boundary derivatives. There is no division by a factor that
might vanish in this argument.

For $\operatorname{Re}z\ge0$ and $z\ne0$, the principal roots satisfy
\[
 \operatorname{Re}\sqrt z\ge |z|^{1/2}/\sqrt2,
 \qquad
 \operatorname{Re}z^{-1/2}\ge |z|^{-1/2}/\sqrt2.
\]
Every derivative of the root exponential is that exponential
times a finite sum of powers of $z^{1/2}$ and $z^{-1/2}$.
It therefore decays faster than every power at zero and at
infinity, uniformly on the closed half-plane. The product rule
proves all the required bounded continuous derivative extensions,
as well as boundary Schwartz decay and flatness at zero.
\end{proof}

\begin{proof}[Proof of Theorem~\ref{p03:mixnew:thm:exact-terminal}]
For completeness, recall the sufficient direction. For the signed
difference of two candidate laws, put
\[
 H_u(T)=\int e^{iu^Ta-\operatorname{tr}(T\Sigma)}\,d\rho.
\]
The experiment is equivalent to independent conditional blocks of
lengths $1,\ldots,1,n-d+1$. Their Fourier arguments are exactly
\[
 T=\frac{uu^T}{2n}+C,\qquad C\succeq0,
                          \operatorname{rank}C\le d-1.
\]
This follows by decomposing the weighted frequency matrix into
its projection on the vector of square roots of block lengths
and its orthogonal complement. Fix $u\ne0$ and a rank-$d-1$
matrix $C\succeq0$ with $u$ not orthogonal to its kernel.
Then $z\mapsto H_u(C+zuu^T)$ is bounded holomorphic on
$\operatorname{Re}z>0$ and vanishes at $1/(2n)$.
Harmonic divergence violates the classical half-plane Blaschke
condition for a nonzero bounded function \cite{p03:mixnew:Garnett2007},
so this function is zero.
Every $T\succ0$ has the decomposition
\[
 T=C+\lambda uu^T,
 \quad\lambda=(u^TT^{-1}u)^{-1},
 \quad C=T-\lambda uu^T\succeq0,\quad\operatorname{rank}C=d-1.
\]
This gives $H_u(T)=0$. Continuity supplies $u=0$; analytic
continuation in $T$ followed by Fourier uniqueness after the
tilt $e^{-\operatorname{tr}\Sigma}$ gives $\rho=0$.

For necessity use the scalar input in Lemma~\ref{p03:mixnew:lem:scalar}
and the joint signed density in
Lemma~\ref{p03:mixnew:lem:matrix-realization}. At a positive definite
argument of experiment $n$, the determinant identity gives
\[
 0=\det(2nT-uu^T)
   =(2n)^d\det T\left(1-\frac{u^TT^{-1}u}{2n}\right).
\]
Thus $F(u^TT^{-1}u)=F(2n)=0$, and the signed observation
transform vanishes. Singular arguments follow by continuity:
full-rank $d$-column frequency matrices are dense, and the signed
density has finite total variation. At $u=0$ it vanishes by the
support of $h$.

Write $g=g_+-g_-$ and $m=\int g_+=\int g_->0$.
Then $m^{-1}g_+$ and $m^{-1}g_-$ are distinct probability
densities with the same laws of every selected experiment.
Schwartz decay, vanishing marginals, and vanishing polynomial
moments give the additional claims. Translating both covariance
variables by $b_0I$ multiplies the signed transform by
$e^{-b_0\operatorname{tr}T}$, preserving its zeros and all the
asserted equalities.
\end{proof}

For bounded covariance support, any infinite terminal set was
already sufficient because the scalar function in the positive
proof becomes entire. For an infinite convergent harmonic subseries,
the witnesses consequently require unbounded covariance support.
The theorem closes the exact
unrestricted sparse-design classification left open in the
previous manuscript; it is not presented as a newly located
published conjecture.

\subsection{Two graph-theoretic routes that remain open}
\label{p03:mixnew:sec:graph-barrier}

The analytic completion above does not settle the separate relation
between generic completion rank and maximum likelihood threshold.
Blekherman and Sinn ask whether some function $f$ satisfies
\[
 \operatorname{mlt}(G)\le k
 \quad\Longrightarrow\quad
 \operatorname{gcr}(G)\le f(k)
 \qquad\text{for every finite graph }G
\]
\cite[Introduction, final question]{p02:graph:BS2017}.
Their complete bipartite examples have generic completion rank of
order $m$ and likelihood threshold of order $\sqrt m$ for $K_{m,m}$.
They rule out a linear comparison, but do not supply bounded-threshold,
unbounded-rank examples. The signed polynomial annihilators underlying
the earlier observation-rank criterion do not provide such examples
either: the likelihood threshold asks for positive semidefinite
certificates, whereas a nonzero annihilator can be indefinite. Thus
an algebraic rank calculation alone does not decide this question.

A second published target is equality of the two parameters with high
probability for $G(n,c/n)$ at every fixed $c>0$
\cite[Conjecture~2.4]{p02:graph:BDG2024}.
Peled's August 2026 revision establishes the generic-rigidity phase
transition and sharp asymptotic rank information. Its concluding
discussion still identifies a globally rigid large core as a missing
step toward the likelihood-threshold conjecture
\cite[Section~7, item~3, p.~30]{p03:mixnew:Peled2026}.
Our transform construction adds neither that global rigidity nor
the necessary positivity. Both graph targets therefore remain open
in this work. This is an explicit boundary of the method, rather than
a claimed consequence of the harmonic classification.

The literature check for the new terminal-design theorem was bounded:
the primary sources above and grouped-sample identifiability sources
were inspected, and no exact earlier statement for this sparse
terminal experiment was located. This does not certify historical
priority. In particular, neither the scalar Blaschke condition nor
the general Fourier--Laplace principle is claimed as new.


% Consolidated section: bessel_new
\section{Self-decomposability beyond the McKay Thorin boundaries}
\label{p03:bn:section}

Two different questions in the published source motivate this section.
Baricz, Prabhu, Singh, and Vijesh~\cite{p01:rare:Baricz2026}, Section~2.3,
ask for the largest infinite-divisibility domain of the generalized
McKay distribution.  Their Section~4.3, pp.~282--283, poses a multipart
problem asking, among other things, for the self-decomposability domains
of the third and fourth McKay families.  The two families are cases of
that one multipart problem, not two separately posed conjectures.

The full GGC classifications established in the preceding working
manuscript~\cite{comp:part:02} were
\[
 x^{\nu-1}e^{-bx}I_\mu(ax):\quad \nu\leq\mu+1,
 \qquad
 x^{2\mu}e^{-bx}I_\mu(ax)^2:\quad \mu\leq\tfrac12.
\]
Those classifications are not new results of this section.  We now
enter their complementary parameter regions.  We prove infinite
divisibility and determine the exact self-decomposability boundary
on infinitely many parameter rays.  For the fourth family at
$\mu=3/2$, the boundary is specified by a unique critical point, and an
explicit rational choice of parameters gives an infinitely divisible
law that is not self-decomposable.  The remaining noninteger parameter
regions are explicitly left open.

All densities below are on $(0,\infty)$ and are normalized by a positive
constant.  Write ID for infinite divisibility, SD for
self-decomposability, and GGC for generalized gamma convolution.  The
positive scale parameter $a$ is always strictly positive.

\subsection{An exact criterion for exponential tilts}

\begin{lemma}[The tilt boundary from a positive spectral kernel]
\label{p03:bn:tilt}
Let $d>0$, and let $K$ be a positive real $C^1$ function on
$[0,\infty)$ such that $K$ and $K'$ are $O(e^x)$.
Suppose a probability law with Laplace transform $L_b$ satisfies
\begin{equation}\label{p03:bn:log-derivative}
 -\frac{d}{ds}\log L_b(s)
   =\int_0^\infty e^{-st}k_b(t)\,dt,
 \qquad k_b(t)=e^{-bt}K(dt),\qquad b>d.
\end{equation}
Then the law is ID, with L\'evy density $k_b(t)/t$.  It is SD if and
only if
\begin{equation}\label{p03:bn:beta-criterion}
 \frac bd\geq\sup_{x\geq0}\frac{K'(x)}{K(x)}.
\end{equation}
\end{lemma}

\begin{proof}
The assumptions imply
$\int_0^\infty(1\wedge t)k_b(t)\,dt/t<\infty$.
Integrating \eqref{p03:bn:log-derivative} from $0$ to $s$, and using
$L_b(0)=1$, gives
\[
 h(s):=-\log L_b(s)
   =\int_0^\infty(1-e^{-st})\frac{k_b(t)}t\,dt.
\]
This is the L\'evy--Khintchine representation of a positive ID law.

For completeness, in this smooth, exponentially decaying setting the
SD criterion is exactly $k_b'\leq0$.  If $k_b$ is nonincreasing, then
for each $0<c<1$,
\[
 h(s)-h(cs)
   =\int_0^\infty(1-e^{-st})
          \frac{k_b(t)-k_b(t/c)}t\,dt
\]
is a L\'evy exponent.  Therefore $L_b(s)/L_b(cs)$ is a Laplace
transform, which is the defining SD property.

Conversely, suppose all the ratios
$R_c(s)=L_b(s)/L_b(cs)$ are Laplace transforms.  Each function
$-R_c'(s)/(1-c)$ is completely monotone.  Differentiating the explicit
smooth expression for $R_c$ shows that these functions and all their
derivatives converge locally on $(0,\infty)$, as $c\uparrow1$, to
the corresponding derivatives of $(s h'(s))'$.  Hence this last
function is completely monotone.  Integration by parts gives
\[
 (s h'(s))'=-\int_0^\infty e^{-st}t k_b'(t)\,dt.
\]
The signed measure on the right has finite total variation.
Bernstein's representation theorem and uniqueness of Laplace
transforms therefore imply $-t k_b'(t)\geq0$.  Continuity yields
$k_b'\leq0$ everywhere.  Finally,
$k_b'(t)=e^{-bt}\{dK'(dt)-bK(dt)\}$, proving
\eqref{p03:bn:beta-criterion}.
\end{proof}

\begin{lemma}[Simultaneous returns of finitely many phases]
\label{p03:bn:returns}
For real numbers $\theta_1,\ldots,\theta_m$, there are arbitrarily
large positive $x$ for which all $e^{i\theta_jx}$ are arbitrarily
close to $1$.
\end{lemma}

\begin{proof}
Partition the $m$-torus into $N^m$ equal boxes and apply the
pigeonhole principle to the first $N^m+1$ integer multiples of
$(\theta_1,\ldots,\theta_m)$ modulo $2\pi$.  There is a positive
integer $q_N\leq N^m$ for which all phases are within $2\pi/N$ of
$1$.  If these integers have an unbounded subsequence, the assertion
follows.  Otherwise a fixed positive integer occurs along a sequence
with $N\to\infty$, is an exact common period, and its multiples
give the assertion.
\end{proof}

\subsection{All integer power biases of the third family}

\begin{theorem}[ID and an exact SD threshold on every integer-bias ray]
\label{p03:bn:third}
Let $\mu>-1/2$, $n\geq2$ be an integer, and
$\nu=\mu+1+n$.  Put $\alpha=\mu+1/2>0$.  For every $b>a>0$, the
law with density
\begin{equation}\label{p03:bn:third-density}
 f_{\mu,n,a,b}(x)=C x^{\mu+n}e^{-bx}I_\mu(ax)
\end{equation}
is ID.  Define polynomials by
\begin{equation}\label{p03:bn:P-recurrence}
 P_0(z)=1,\qquad
 P_{j+1}(z)=2(\alpha+j)zP_j(z)-(1+z^2)P_j'(z).
\end{equation}
The $n$ roots $\omega_1,\ldots,\omega_n$ of $P_n$ are real and
simple.  Set
\begin{equation}\label{p03:bn:third-K}
 K_{\alpha,n}(x)=2(\alpha+n)\cosh x
                       -\sum_{j=1}^n e^{i\omega_jx}.
\end{equation}
This is a positive real even function.  The constant
\begin{equation}\label{p03:bn:third-beta}
 \beta_{\alpha,n}
   :=\max_{x\geq0}\frac{K_{\alpha,n}'(x)}{K_{\alpha,n}(x)}
\end{equation}
exists, is attained at a finite positive point, and satisfies
\begin{equation}\label{p03:bn:third-bound}
 1<\beta_{\alpha,n}\leq
 \frac{2(\alpha+n)+n\sqrt{(n-1)/(2\alpha+1)}}{2\alpha+n}.
\end{equation}
The exact classification on this ray is
\begin{equation}\label{p03:bn:third-SD}
 f_{\mu,n,a,b}\text{ is SD}
       \quad\Longleftrightarrow\quad b/a\geq\beta_{\alpha,n}.
\end{equation}
In particular, each ray contains both ID laws that are not SD and
SD laws outside the GGC domain.
\end{theorem}

\begin{proof}
The elementary Bessel Laplace integral gives, with a positive constant
$C_0$ independent of $u$,
\[
 J_0(u):=\int_0^\infty x^\mu e^{-ux}I_\mu(ax)\,dx
       =C_0(u^2-a^2)^{-\alpha},\qquad u>a.
\]
For example, this follows directly by integrating the positive power
series of $I_\mu$ and using the gamma duplication formula.  All
differentiations under the integral are justified by exponential
decay.  Consequently
\begin{equation}\label{p03:bn:third-J}
 J_n(u):=(-1)^nJ_0^{(n)}(u)
       =C_0\frac{Q_n(u)}{(u^2-a^2)^{\alpha+n}},
\end{equation}
where
\[
 Q_0(u)=1,\qquad
 Q_{j+1}(u)=2(\alpha+j)uQ_j(u)-(u^2-a^2)Q_j'(u).
\]
The leading coefficient of $Q_j$ is $(2\alpha)_j>0$.
The substitution $Q_j(iaz)=i^j a^jP_j(z)$ gives
\eqref{p03:bn:P-recurrence}.

We verify the root assertion rather than assume a special-function
zero theorem.  The polynomial $P_1(z)=2\alpha z$ has a simple real
root and positive leading coefficient.  Suppose $P_j$ has ordered
simple real roots $r_1<\cdots<r_j$ and positive leading coefficient.
At each root,
\[
 P_{j+1}(r_\ell)=-(1+r_\ell^2)P_j'(r_\ell).
\]
These signs alternate.  There is a root in each of the $j-1$
intervening intervals.  The signs at the two outer roots, together
with the positive leading coefficient $(2\alpha)_{j+1}$, give one
root on each outer interval.  These $j+1$ distinct roots exhaust
the degree.  Induction proves simplicity, reality, and interlacing.
Parity in \eqref{p03:bn:P-recurrence} also gives invariance of the roots
under $\omega\mapsto-\omega$.

The Laplace transform of \eqref{p03:bn:third-density} is
$L_b(s)=J_n(s+b)/J_n(b)$.  Its negative logarithmic derivative,
with $u=s+b$, is
\begin{equation}\label{p03:bn:third-rational}
 -\frac{d}{ds}\log L_b(s)
 =\frac{\alpha+n}{u-a}+\frac{\alpha+n}{u+a}
               -\sum_{j=1}^n\frac1{u-ia\omega_j}.
\end{equation}
Taking elementary Laplace transforms yields the kernel
$k_b(t)=e^{-bt}K_{\alpha,n}(at)$ in
\eqref{p03:bn:third-K}.  Root symmetry makes it real and even, and
\begin{equation}\label{p03:bn:third-positive}
 K_{\alpha,n}(x)\geq2(\alpha+n)\cosh x-n
                         \geq2\alpha+n>0.
\end{equation}
Lemma~\ref{p03:bn:tilt} proves ID and the asserted SD equivalence, once
the statements about the maximum are verified.

The logarithmic derivative $K'/K$ is continuous, equals $0$ at
$0$, and tends to $1$ at infinity.  Moreover,
\begin{equation}\label{p03:bn:third-overshoot}
 K'(x)-K(x)=-2(\alpha+n)e^{-x}
                   +\sum_{j=1}^n(1-i\omega_j)e^{i\omega_jx}.
\end{equation}
At simultaneous returns of the phases, the sum on the right is
arbitrarily close to $n$, since $\sum_j\omega_j=0$.
Lemma~\ref{p03:bn:returns} therefore gives arbitrarily large $x$ for
which $K'(x)>K(x)$.  It follows that the supremum is strictly
larger than $1$ and is attained in a compact subinterval of
$(0,\infty)$.

For the quantitative upper bound, put $S=\sum_j|\omega_j|$ and
$A=2(\alpha+n)$.  Then
\[
 \frac{K'(x)}{K(x)}
   \leq\frac{A\sinh x+S}{A\cosh x-n}
   \leq\frac{A+S}{A-n}.
\]
Comparison of the two leading coefficients in the recurrence for
$Q_n$ gives
\[
 Q_n(u)=(2\alpha)_n
 \left(u^n+\frac{n(n-1)a^2}{2(2\alpha+1)}u^{n-2}
                       +\text{lower powers}\right).
\]
This coefficient identity follows inductively: if the second-to-first
coefficient ratio is $r_j$, then
$r_{j+1}=\{(2\alpha+j+2)r_j+j\}/(2\alpha+j)$, starting at
$r_1=0$.  Vieta's formulas now give
$\sum_j\omega_j^2=n(n-1)/(2\alpha+1)$.
Cauchy--Schwarz yields the bound for $S$ in
\eqref{p03:bn:third-bound}.

Finally, $\nu=\mu+1+n>\mu+1$ puts the entire ray strictly beyond
the previously established GGC boundary.  The intervals
$1<b/a<\beta_{\alpha,n}$ and
$b/a\geq\beta_{\alpha,n}$ supply the two asserted subclasses.
\end{proof}

The first ray covered here, $n=2$, already has the explicit kernel
\begin{equation}\label{p03:bn:n2}
 K_{\alpha,2}(x)=2(\alpha+2)\cosh x
                  -2\cos\!\left(\frac{x}{\sqrt{2\alpha+1}}\right).
\end{equation}
The case $n=1$ is the second McKay family, whose SD classification was
obtained in the earlier research round; it is deliberately excluded
from the new-result count.  Theorem~\ref{p03:bn:third} also supplies new
ID parameter rays for the separate maximal-domain question in
Section~2.3 of the source.  It does not determine that maximal domain.

The maximization in \eqref{p03:bn:third-beta} has a certified finite
search region.  If $K'(x_0)/K(x_0)=1+\delta$ with $\delta>0$,
then \eqref{p03:bn:third-overshoot} gives
\[
 \frac{K'(x)}{K(x)}-1
 \leq\frac{n+S}{2(\alpha+n)\cosh x-n}.
\]
Once the right side falls below $\delta$, no subsequent point can
maximize the ratio.  This is an exact tail estimate, not an appeal
to numerical truncation.

\subsection{All half-integer orders of the fourth family}

\begin{theorem}[An ID extension and exact SD thresholds for squared Bessel laws]
\label{p03:bn:fourth}
Let $n\geq1$ be an integer, $\mu=n+1/2$, $d=2a$, and $b>d$.
The law with density
\begin{equation}\label{p03:bn:fourth-density}
 g_{n,a,b}(x)=C x^{2n+1}e^{-bx}I_{n+1/2}(ax)^2
\end{equation}
is ID.  Define
\begin{equation}\label{p03:bn:fourth-Q}
 q_n(z)=\sum_{k=0}^n(-1)^k a_k z^{2(n-k)},\qquad
 a_k=\binom{n}{k}\frac{(1/2)_k}{(n+3/2)_k},
\end{equation}
and let $\lambda_1,\ldots,\lambda_{2n}$ denote all its roots,
counted with multiplicity.  Put
\begin{align}
 \rho_n^2&=\max_{1\leq k\leq n}
       \frac{(n-k+1)(k-1/2)}{k(n+k+1/2)}<1,
       \label{p03:bn:rho}\\
 K_n(x)&=(2n+1)(1+2\cosh x)
                      -\sum_{j=1}^{2n}e^{\lambda_jx}.
       \label{p03:bn:fourth-K}
\end{align}
All $|\lambda_j|\leq\rho_n$, and $K_n$ is a positive real even
function.  Its exact SD threshold is
\begin{equation}\label{p03:bn:fourth-beta}
 \beta_n=\max_{x\geq0}\frac{K_n'(x)}{K_n(x)},\qquad
 g_{n,a,b}\text{ is SD}\quad\Longleftrightarrow\quad
                         \frac{b}{2a}\geq\beta_n.
\end{equation}
The maximum is attained at a finite positive point and
\begin{equation}\label{p03:bn:fourth-bound}
 1<\beta_n\leq\frac{4n+2+2n\rho_n}{2n+2}<3.
\end{equation}
Thus every such order has a nonempty ID but non-SD interval, and
$b\geq6a$ is a uniform sufficient condition for SD at all these
orders.  All the laws in this theorem lie outside the GGC domain.
\end{theorem}

\begin{proof}
The Laplace integral in the source, equation~(2-4), gives
\[
 J(u)=C_1u^{-(4n+3)}
  {}_2F_1(n+1,2n+3/2;n+3/2;d^2/u^2),\qquad u>d,
\]
where $C_1>0$ does not depend on $u$.
Euler's hypergeometric transformation~\cite{p01:rare:DLMF} gives
\[
 {}_2F_1(n+1,2n+3/2;n+3/2;z)
  =(1-z)^{-(2n+1)}{}_2F_1(-n,1/2;n+3/2;z).
\]
Consequently, with $Q_n(u)=d^{2n}q_n(u/d)$,
\begin{equation}\label{p03:bn:fourth-J}
 J(u)=C_1\frac{Q_n(u)}{u^{2n+1}(u^2-d^2)^{2n+1}}.
\end{equation}

Here is a direct root bound.  A polynomial
$p(z)=\sum_{j=0}^m c_jz^j$ with
$0<c_0\leq\cdots\leq c_m$ has no root with $|z|>1$.
Indeed,
\[
 (1-z)p(z)=c_0+\sum_{j=1}^m(c_j-c_{j-1})z^j-c_mz^{m+1},
\]
and at such a root the absolute value of the last term strictly
exceeds the sum of the preceding absolute-value bounds.  Rescaling
shows that a polynomial with arbitrary positive coefficients has all
roots in $|z|\leq\max_j c_j/c_{j+1}$.
Apply this fact to
$p(v)=\sum_{k=0}^n a_kv^{n-k}$, whose roots are the negatives of
the squared roots of $q_n$.  Since
\[
 \frac{a_k}{a_{k-1}}
  =\frac{(n-k+1)(k-1/2)}{k(n+k+1/2)}<1,
\]
the claimed bound $|\lambda_j|\leq\rho_n<1$ follows.

Set $B=2n+1$.  Logarithmic differentiation of
$L_b(s)=J(s+b)/J(b)$ gives
\begin{equation}\label{p03:bn:fourth-rational}
 -\frac{d}{ds}\log L_b(s)
   =\frac B u+\frac B{u-d}+\frac B{u+d}
                     -\sum_{j=1}^{2n}\frac1{u-d\lambda_j},
 \qquad u=s+b.
\end{equation}
The roots are invariant under conjugation and negation, so
\eqref{p03:bn:fourth-K} is real and even.  Pairing opposite roots and
using $|\cosh z|\leq\cosh|z|$ gives
\begin{equation}\label{p03:bn:fourth-positive}
 K_n(x)\geq B+2B\cosh x-2n\cosh(\rho_nx)
       \geq B+(2n+2)\cosh x>0.
\end{equation}
Equation~\eqref{p03:bn:fourth-rational} is precisely the Laplace
transform of $e^{-bt}K_n(dt)$.  Lemma~\ref{p03:bn:tilt} proves ID
and the threshold equivalence.

To prove the strict lower bound for $\beta_n$, first observe that
some root has positive real part.  Indeed, Vieta's formulas give
\[
 \sum_{j=1}^{2n}\lambda_j=0,\qquad
 \sum_{j=1}^{2n}\lambda_j^2=2a_1>0.
\]
If every root had real part zero, the last sum would be nonpositive;
negation symmetry excludes the possibility that all real parts are
nonpositive with one strictly negative.  Thus
$r:=\max_j\operatorname{Re}\lambda_j$ satisfies $0<r<1$.
Let $\mathcal J$ be the multiset of roots with real part $r$ and
$m=|\mathcal J|$.  The real trigonometric polynomial
\[
 T(x)=\sum_{\lambda\in\mathcal J}
                 (1-\lambda)e^{i\operatorname{Im}\lambda\,x}
\]
satisfies $T(0)=m(1-r)>0$.  Simultaneous returns from
Lemma~\ref{p03:bn:returns} give arbitrarily large $x$ with
$T(x)>T(0)/2$.  At these points,
\[
 K_n'(x)-K_n(x)
   =-B(1+2e^{-x})+e^{rx}T(x)+O(e^{r'x})>0,
\]
where $r'<r$ bounds the real parts of the remaining roots; an absent
remainder is simply omitted.  As $K_n'/K_n$ equals $0$ at the
origin and tends to $1$ at infinity, its maximum is greater than
$1$ and is attained at a finite positive point.

Opposite-root pairing also gives
\[
 \left|\sum_j\lambda_je^{\lambda_jx}\right|
       \leq2n\rho_n\sinh(\rho_nx)
       \leq2n\rho_n\sinh x.
\]
Together with \eqref{p03:bn:fourth-positive}, this proves
\[
 \frac{K_n'(x)}{K_n(x)}
 \leq\frac{(2B+2n\rho_n)\sinh x}{B+(2B-2n)\cosh x}
 \leq\frac{2B+2n\rho_n}{2B-2n}<3.
\]
The final GGC statement uses the earlier boundary
$\mu\leq1/2$, since $\mu=n+1/2>1/2$ here.
\end{proof}

The exact thresholds again admit finite tail certification.  From
\eqref{p03:bn:fourth-positive}, $K_n(x)\geq e^x$, and hence
\begin{equation}\label{p03:bn:fourth-tail}
 \left|\frac{K_n'(x)}{K_n(x)}-1\right|
 \leq\{3(2n+1)+2n(1+\rho_n)\}e^{-(1-\rho_n)x}.
\end{equation}
A single point with ratio $1+\delta>1$ therefore supplies an
explicit finite interval containing every global maximizer.

\subsection{A uniquely determined boundary at order three halves}

\begin{theorem}[The order-$3/2$ SD transition and an exact counterexample]
\label{p03:bn:three-halves}
Let $a>0$, $b>2a$, and let $X$ have density
\[
 f(x)=C x^3e^{-bx}I_{3/2}(ax)^2.
\]
Put $c=1/\sqrt5$ and
\begin{align}
 K(x)&=3+6\cosh x-2\cosh(cx),\label{p03:bn:explicit-K}\\
 D(x)&=36+4c^2+18\cosh x-6c^2\cosh(cx)\notag\\
     &\hspace{8mm}-6(1-c)^2\cosh((1+c)x)
                  -6(1+c)^2\cosh((1-c)x).
                  \label{p03:bn:explicit-D}
\end{align}
There is exactly one positive zero $\tau$ of $D$.  With
\begin{equation}\label{p03:bn:explicit-beta}
 \beta=\frac{6\sinh\tau-2c\sinh(c\tau)}
                 {3+6\cosh\tau-2\cosh(c\tau)},
\end{equation}
the law of $X$ is ID for every $b>2a$, and is SD exactly when
$b/(2a)\geq\beta$.  In particular:
\begin{enumerate}
\item $b=1001a/500$ gives an ID law that is not SD;
\item $b=4a$ gives an SD law that is not GGC.
\end{enumerate}
The boundary is larger than $1001/1000$.
\end{theorem}

\begin{proof}
For $n=1$, \eqref{p03:bn:fourth-Q} is $q_1(z)=z^2-1/5$,
which gives \eqref{p03:bn:explicit-K}.  Direct differentiation and
the addition identities for hyperbolic functions show that
$K''K-(K')^2=D$ as displayed in \eqref{p03:bn:explicit-D}.  Thus
the sign of $D$ is the sign of $(K'/K)'$.

Write
$D(x)=d_0+\sum_{m\geq1}d_mx^{2m}/(2m)!$.
Then $d_0=196/5$, and
\begin{equation}\label{p03:bn:D-coefficients}
 d_m=6\left[3-c^{2m+2}-(1-c^2)^2
       \{(1+c)^{2m-2}+(1-c)^{2m-2}\}\right],\qquad m\geq1.
\end{equation}
For $m=1,2,3,4$ the bracketed quantities are respectively
\[
 \frac{42}{25},\qquad \frac{182}{125},\qquad
 \frac{82}{625},\qquad -\frac{9058}{3125}.
\]
They strictly decrease with $m$.  To see this, set
$S_j=(1+c)^j+(1-c)^j$.  Expansion by the binomial theorem gives
$S_{2m}-S_{2m-2}\geq2c^2$ for every $m\geq1$.  Therefore the
increase in the subtracted quantity in
\eqref{p03:bn:D-coefficients} is at least
\[
 (1-c^2)^2(2c^2)-(1-c^2)c^{2m+2}
 \geq\frac{16}{25}\frac25-\frac45\frac1{25}
 =\frac{28}{125}>0.
\]
It follows that $d_1,d_2,d_3>0$ and $d_m<0$ for all $m\geq4$.

On $(0,\infty)$, the function
\[
 \frac{D(x)}{x^6}=\frac{d_0}{x^6}
       +\frac{d_1}{2!x^4}+\frac{d_2}{4!x^2}
       +\frac{d_3}{6!}
       +\sum_{m\geq4}\frac{d_mx^{2m-6}}{(2m)!}
\]
is strictly decreasing.  Termwise differentiation is justified on
each compact subinterval by convergence of the entire power series,
and every nonconstant summand has negative derivative.  The limits
are $+\infty$ at zero and $-\infty$ at infinity, the latter also
following from the negative dominant term
$-6(1-c)^2\cosh((1+c)x)$ in \eqref{p03:bn:explicit-D}.
Thus $D$ has exactly one positive zero $\tau$.  The ratio $K'/K$
strictly increases before $\tau$ and strictly decreases afterwards.
Theorem~\ref{p03:bn:fourth} now proves the exact threshold formula.

We prove the non-SD example without decimal approximations.
Let $r=1001/1000$ and take $x=\log81=4\log3$.
The inequalities $4/9<c<9/20$ imply
\[
 e^{cx}=3^{4c}>3^{16/9}>7,
\]
where the last inequality is the integer comparison
$3^{16}=43046721>40353607=7^9$.  Hence
$\cosh(cx)>7/2$ and $\sinh(cx)<\cosh(cx)$.
Using $e^x=81$, we obtain
\begin{align*}
 K'(x)-rK(x)
 &=6\sinh x-6r\cosh x
      +2r\cosh(cx)-2c\sinh(cx)-3r\\
 &>-\frac6{81}
      -\frac6{1000}\frac{81+1/81}{2}
      +7\left(r-\frac9{20}\right)-3r
   =\frac{604}{1125}>0.
\end{align*}
Therefore $\beta>r$, proving the first example.
For the second, \eqref{p03:bn:fourth-bound} gives
\[
 \beta\leq\frac{6+2/\sqrt5}{4}<2.
\]
Thus $b=4a$, for which $b/(2a)=2$, is SD.  The preceding GGC
classification excludes GGC at $\mu=3/2$.
\end{proof}

Numerical root finding, provided only as an illustration of the
uniquely specified constant, gives
\[
 \tau\approx4.988970295384644,\qquad
 \beta\approx1.005210740181144.
\]
No assertion above depends on these decimal values.  In particular,
the existence of the non-SD interval and the rational counterexample
were proved by exact inequalities.

\subsection{What remains unresolved}

These results resolve the SD question on the rays
$\nu-\mu-1\in\{2,3,\ldots\}$ with $\mu>-1/2$, and on the
orders $\mu\in\{3/2,5/2,\ldots\}$ of the squared-Bessel family.
They do not classify the full real parameter domains in the source.
For the third family, the regions with noninteger
$\nu-\mu-1>0$, and the admissible negative orders
$-1<\mu\leq-1/2$, are not settled here.  For the fourth family,
non-half-integer orders $\mu>1/2$ remain outside the argument.

The specific obstruction is mathematical: the proofs above turn the
Laplace integrals into a power of a quadratic times a finite
polynomial.  Interlacing or coefficient ratios then control every
zero and produce a positive, finite exponential expression for the
L\'evy spectral kernel.  At noninteger parameters the terminating
polynomials become nonterminating hypergeometric functions.  Their
logarithmic derivatives can have additional branch contributions,
and the present root bounds provide no sign control for the full
inverse Laplace transform.  Exponential tilting alone does not
remove this gap.  No full solution of either source question, and no
priority claim for the entire literature, is inferred from the
partial classifications proved here.


% Consolidated section: zeros_new
\section{All modular periods for a three-term exponential sum}
\label{p03:zn:sec:modular}

This section advances the modular-ordinate problem rather than the
previous edge-phase reconstruction argument.  Podesta
\cite[Section~1, paragraphs after the mean vertical zero density]{p02:zp:Podesta2026}
asks whether the imaginary parts of zeros of a rationally independent
exponential sum can be dense modulo a suitable length, and explicitly
defers the solution.  The nonresonant periods follow from a compact-torus
argument.  Here we classify \emph{every} period for a three-term sum,
including the resonant periods.  The proof uses triangle geometry and a
three-dimensional orbit closure.  These ingredients are classical; the
claimed contribution is the explicit complete classification below.
A targeted primary-source search on 14 October 2026 located the printed
question but no subsequent resolution of this exact classification.
This limited search is not a claim of historical priority.

Fix
\begin{equation}\label{p03:zn:eq:model}
 f(s)=-1+B e^{as+i\phi}+C e^{bs+i\psi},
 \qquad B,C>0,\quad 0<a<b,\quad a/b\notin\mathbb Q .
\end{equation}
For $L>0$, put $q=2\pi/L$ and let
\[
 \mathcal R_L=\overline{\{q\,\operatorname{Im}\rho\bmod 2\pi:
                   f(\rho)=0,\ \operatorname{Im}\rho>0\}}
       \subset\mathbb R/(2\pi\mathbb Z).
\]
Thus $\mathcal R_L$ is the circle precisely when the positive ordinates
are dense modulo $L$.  Neither multiplicity nor finitely many low zeros
affects this question.

\begin{theorem}[Complete trinomial classification]\label{p03:zn:thm:classification}
Let $r=a/b$ and $\kappa=BC^{-r}$.  There is a unique $x_0>0$ satisfying
\begin{equation}\label{p03:zn:eq:minimum-equation}
 x_0^{\,2-2r}-x_0^{-2r}=\kappa^2(2r-1).
\end{equation}
Define
\begin{equation}\label{p03:zn:eq:theta}
 u(x)=\frac{x+x^{-1}-\kappa^2x^{2r-1}}2,\qquad
 \Theta=\arccos\!\bigl(\max\{-1,u(x_0)\}\bigr)\in(0,\pi].
\end{equation}
Then the following statements hold.
\begin{enumerate}
\item If $q\notin\operatorname{span}_{\mathbb Q}\{a,b\}$, then
      $\mathcal R_L=\mathbb R/(2\pi\mathbb Z)$.
\item Otherwise write, uniquely,
      \[
       q=(ma+nb)/d,\qquad m,n\in\mathbb Z,\quad d\ge1,\quad
                        \gcd(m,n,d)=1.
      \]
      If $m\ne0$, then $\mathcal R_L=\mathbb R/(2\pi\mathbb Z)$.
\item If $m=0$, necessarily $n>0$ and $\gcd(n,d)=1$, and
      \begin{equation}\label{p03:zn:eq:exact-arcs}
       \mathcal R_L
       =\{z\bmod2\pi:\ dz+n\psi\bmod2\pi
                         \in[-n\Theta,n\Theta]\bmod2\pi\}.
      \end{equation}
      In particular, this set is the circle if and only if $n\Theta\ge\pi$.
\end{enumerate}
Equivalently, the complete set of nondense periods is
\begin{equation}\label{p03:zn:eq:exceptional}
 \left\{\frac{2\pi d}{nb}:\ n,d\ge1,\ \gcd(n,d)=1,\ n\Theta<\pi\right\}.
\end{equation}
This criterion does not depend on $\phi$ or $\psi$, although the location
of the missing arcs depends on $\psi$.  The assertion concerns density,
not equidistribution at resonant periods.
\end{theorem}

\begin{lemma}[Triangle interval and endpoint phases]\label{p03:zn:lem:triangle}
Put $x=Ce^{b\sigma}$ and $U=\kappa x^r$.  The numbers $1,U,x$ are the
side lengths of a possibly degenerate triangle exactly for
\[
 K=\{x>0:-1\le u(x)\le1\}.
 \]
The set $K$ is either one compact interval or two compact intervals,
and its nondegenerate part is dense in it.  On $K$ define the continuous
angles
\begin{equation}\label{p03:zn:eq:angles}
 \theta(x)=\arccos\frac{1+U^2-x^2}{2U},\qquad
 \eta(x)=-\arccos\frac{1+x^2-U^2}{2x}=-\arccos u(x).
\end{equation}
All solutions of $1=Ue^{i\theta'}+xe^{i\eta'}$ are, modulo $2\pi$,
$(\theta',\eta')=(\theta,\eta)$ or $(-\theta,-\eta)$.
At the leftmost endpoint of $K$ these angles are $(0,0)$, and at
the rightmost endpoint they are $(\pi,0)$.  If $K$ has two components,
the two endpoints facing its gap both have angles $(0,-\pi)$.
Moreover,
\[
                    \max_{x\in K}|\eta(x)|=\Theta .
 \]
\end{lemma}
\begin{proof}
The triangle inequalities are $|1-x|\le U\le1+x$, which are equivalent
to $-1\le u(x)\le1$.  Differentiation gives
\[
 2x^{2-2r}u'(x)
   =x^{2-2r}-x^{-2r}-\kappa^2(2r-1).
 \]
The function $x^{2-2r}-x^{-2r}$ is strictly increasing from
$-\infty$ to $+\infty$.  Consequently $u$ decreases up to the unique
$x_0$ in \eqref{p03:zn:eq:minimum-equation} and increases after it.
Also $u(x)\to+\infty$ at both ends of $(0,\infty)$ and
$u(1)=1-\kappa^2/2<1$.  This proves the claimed structure of $K$,
including the case $\min u=-1$, when the strict triangle set has two
components meeting at a degenerate limiting point.  If $\min u<-1$,
the interval where $u<-1$ is the only gap.

The cosine rule gives \eqref{p03:zn:eq:angles}; the signs are opposite
because the sum is positive real.  In the strict triangle case these
are exactly the two conjugate configurations.  Continuity gives the
same assertion at degeneracy, with the configurations possibly equal.
The two outer endpoints solve $U=|1-x|$.  The left one is below $1$,
so $U+x=1$ and its angles are $(0,0)$.  The right one is above $1$,
so $x=1+U$ and its angles are $(\pi,0)$.
The endpoints of a gap solve $U=1+x$, yielding $(0,-\pi)$.
Finally, the smallest value of $u$ attained on $K$ is
$\max\{-1,u(x_0)\}$, proving the formula for the maximum angle.
\end{proof}

\begin{lemma}[Orbit realization of every compatible triangle]
\label{p03:zn:lem:realization}
For $x$ in the strict triangle set, take either angle pair
$(\theta',\eta')=(\theta(x),\eta(x))$ or its negative.
If $q\notin\operatorname{span}_{\mathbb Q}\{a,b\}$, every
$z\in\mathbb R/(2\pi\mathbb Z)$ is a limit of residues of positive
ordinates of zeros with real parts tending to $b^{-1}\log(x/C)$.
In the resonant case $q=(ma+nb)/d$, the same conclusion holds exactly
for the compatible residues
\begin{equation}\label{p03:zn:eq:compatibility}
                 dz=m(\theta'-\phi)+n(\eta'-\psi)\pmod{2\pi}.
\end{equation}
The closure of all zero residues is obtained by allowing all $x\in K$,
both angle pairs, and all compatible residues.
\end{lemma}
\begin{proof}
Consider the orbit
\[
       t\longmapsto (at,bt,qt)\pmod{2\pi}
\]
in the three-dimensional torus.  In the nonresonant case its coordinates
are rationally independent, so the orbit is dense in the entire torus.
In the resonant case its integer relation lattice is exactly
$\mathbb Z(m,n,-d)$.  Indeed, substituting
$q=(ma+nb)/d$ into any relation and using the rational independence of
$a,b$ proves this claim; primitivity of $(m,n,d)$ makes the multiplier
an integer.  The orbit closure is therefore the subtorus
\[
              H=\{(v,w,z):mv+nw-dz=0\pmod{2\pi}\}.
\]
This is the elementary character form of Kronecker's theorem.  Both
claims also follow by averaging characters, then approximating
continuous functions by trigonometric polynomials.  In particular,
each point of the orbit closure can be approached along times
$t_j\to+\infty$, since discarding a bounded initial segment leaves the
same orbit closure.

For a compatible target, choose these $t_j$ so that
$at_j\to\theta'-\phi$, $bt_j\to\eta'-\psi$, and $qt_j\to z$,
all modulo $2\pi$.  Locally uniformly in $s$,
\[
 f(s+it_j)\longrightarrow
 f_*(s)=-1+B e^{as+i\theta'}+C e^{bs+i\eta'}.
\]
At $\sigma=b^{-1}\log(x/C)$ the function $f_*$ vanishes.
This zero is simple: if both $Ue^{i\theta'}+xe^{i\eta'}=1$ and
$aUe^{i\theta'}+bxe^{i\eta'}=0$, solving the two linear equations
forces both summands to be real, contrary to a strict triangle.
Rouche's theorem on a sufficiently small disk gives zeros
$s_j\to\sigma$ of $f(s+it_j)$.  Thus $\rho_j=s_j+it_j$ are zeros of
$f$, their ordinates tend to $+\infty$, and
$q\operatorname{Im}\rho_j\to z$ modulo $2\pi$.

Conversely, every zero gives a triangle, and its angles satisfy
\eqref{p03:zn:eq:compatibility} in the resonant case.  The possible real
parts form the compact set corresponding to $K$.  Taking convergent
subsequences therefore proves containment of the residue closure in
the stated compact set.  Strict triangle points are dense in $K$,
so the preceding realization argument supplies all its endpoint
residues as well.
\end{proof}

\begin{proof}[Proof of Theorem~\ref{p03:zn:thm:classification}]
Lemma~\ref{p03:zn:lem:realization} proves the nonresonant assertion.
For a resonance define the real-valued function
$g(x)=m\theta(x)+n\eta(x)$ on $K$.  Its image is an interval containing
$0$ and $m\pi$.  To see connectedness when $K$ has two components,
observe that the image of each component is an interval, and both
images contain the value $-n\pi$ at the two inner endpoints.
Consequently, with $M=\max_{x\in K}|g(x)|$,
\[
              g(K)\cup[-g(K)]=[-M,M].
\]
The orbit realization lemma now gives the exact general formula
\begin{equation}\label{p03:zn:eq:all-resonances}
 \mathcal R_L=\{z\bmod2\pi:\ dz+m\phi+n\psi\bmod2\pi
                                     \in[-M,M]\bmod2\pi\}.
\end{equation}
Multiplication by the positive integer $d$ is surjective on the circle,
so the set in \eqref{p03:zn:eq:all-resonances} is the whole circle if and
only if $M\ge\pi$.  If $m\ne0$, the endpoint value $m\pi$ proves this
inequality.  If $m=0$, positivity of $q$ gives $n>0$, and
Lemma~\ref{p03:zn:lem:triangle} gives $M=n\Theta$.  This proves
\eqref{p03:zn:eq:exact-arcs} and the complete list \eqref{p03:zn:eq:exceptional}.
\end{proof}

\begin{corollary}[A finite numerator obstruction]\label{p03:zn:cor:finite}
Only finitely many positive numerators $n$ can appear in the exceptional
list \eqref{p03:zn:eq:exceptional}, namely $n<\pi/\Theta$.  If $\Theta=\pi$,
every period $L>0$ has dense ordinates.  If $\Theta<\pi$, the periods
$2\pi d/b$, $d=1,2,\ldots$, all fail density.  At such a period the
closure consists of $d$ closed arcs whose total relative circle length
is $\Theta/\pi$.
\end{corollary}
\begin{proof}
The first three statements are direct substitutions in
\eqref{p03:zn:eq:exceptional}.  In the last case $n=1$ and the arc
$[-\Theta,\Theta]$ has length $2\Theta<2\pi$; its inverse image under
multiplication by $d$ has $d$ components, each of length $2\Theta/d$.
\end{proof}

\begin{corollary}[Closed coefficient thresholds]\label{p03:zn:cor:threshold}
For $n\ge1$, set $\tau_n=\pi/n$, $c_n=\cos\tau_n$ and
\[
 x_n=
 \frac{(1-2r)c_n+
       \sqrt{(1-2r)^2c_n^2+4r(1-r)}}{2(1-r)},\qquad
 K_n(r)=x_n^{-r}\sqrt{1+x_n^2-2c_nx_n}.
\]
Then the resonant periods $L=2\pi d/(nb)$ in lowest terms have dense
ordinates if and only if $\kappa\ge K_n(r)$.  In particular,
every positive period has dense ordinates if and only if
\begin{equation}\label{p03:zn:eq:all-periods-threshold}
 BC^{-a/b}\ge
 \frac{1}{r^r(1-r)^{1-r}},\qquad r=a/b.
\end{equation}
\end{corollary}
\begin{proof}
The inequality $\Theta\ge\tau_n$ is equivalent to the existence of
$x>0$ such that $u(x)\le c_n$, hence to
\[
 \kappa^2\ge\min_{x>0}x^{-2r}(1+x^2-2c_nx).
\]
The function minimized tends to infinity at both endpoints.  Its
derivative has the sign of
$(1-r)x^2-(1-2r)c_nx-r$, which has exactly one positive root,
the displayed $x_n$.  This proves the formula.
When $n=1$, $c_1=-1$, $x_1=r/(1-r)$, and
$K_1(r)=r^{-r}(1-r)^{-(1-r)}$.  Since the length $2\pi/b$ is dense
exactly when $\Theta=\pi$, the last assertion follows.
\end{proof}

\begin{remark}[What the classification does and does not close]
The theorem answers the printed existence question in the three-term
case and additionally decides every resonant length.  The nonresonant
existence argument was already present in the preceding work and is
not counted as a new result.  The all-resonance formula and the finite
numerator obstruction are the increment.  General sums with four or
more terms have higher-dimensional polygon configuration spaces;
the present proof does not classify their resonant lengths.
The proof also gives no equidistribution assertion at a resonant length.
\end{remark}

\section{Two attempted extensions: a precise obstruction and an unresolved estimate}
\label{p03:zn:sec:limits}

\subsection{Unknown frequencies cannot be recovered in an unrestricted Dirichlet class}

One must specify the model before asking whether rightmost zeros recover
unknown frequencies.  The following elementary obstruction survives
both a fixed constant coefficient and a fixed first frequency and
coefficient.  It is recorded as a counterexample, not as a new general
inverse theorem and not as an Euler-product counterexample.

\begin{proposition}[Normalized right-divisor nonidentifiability]
\label{p03:zn:prop:nonidentifiability}
Fix $c<R$, $a>0$, and
$F(s)=1-e^{aR}e^{-as}$.  For every $\lambda>a$ with $\lambda/a$
irrational, choose a nonzero $\epsilon$ such that
$|\epsilon|e^{-\lambda c}<1$, and put
\[
 G_{\lambda,\epsilon}(s)
  =F(s)(1+\epsilon e^{-\lambda s})
  =1-e^{aR}e^{-as}+\epsilon e^{-\lambda s}
            -\epsilon e^{aR}e^{-(a+\lambda)s}.
\]
Then all these finite Dirichlet sums have constant coefficient $1$,
the same smallest positive frequency $a$, and the same coefficient
$-e^{aR}$ at that frequency.  Their complete zero divisors in
$\{\Re s>c\}$ agree, including multiplicities, and their rightmost real
zero coordinate is the common value $R$.  Nevertheless their remaining
frequencies and coefficients vary.
\end{proposition}
\begin{proof}
The factor $1+\epsilon e^{-\lambda s}$ has no zero in $\Re s>c$ by
the displayed strict modulus bound.  It thus preserves the entire
zero divisor there.  Its own zeros all have real part
$\lambda^{-1}\log|\epsilon|<c$, whereas the zeros of $F$ are
$R+2\pi i k/a$, $k\in\mathbb Z$, and are simple.
The frequency and coefficient statements follow from the explicit
expansion and $\lambda>a$.
\end{proof}

This example shows that even the entire right divisor, stronger than an
edge-approaching subset, does not identify a normalized general Dirichlet
sum.  It does not settle the narrower problem for logarithmic derivatives
of Euler products with prescribed local factor form, nor for sums whose
full list of nonconstant frequencies is rationally independent.  The
product construction introduces an additive frequency relation.  Those
restrictions must be retained in any proposed inverse theorem.

\subsection{The actual Euler hyperuniformity problem remains open here}
\label{p03:zn:subsec:hyperuniformity}

For independent Haar phases $\omega_p$, consider the actual Euler hull
\[
 Q_\omega(s)=\sum_p\frac{(\log p)\omega_pp^{-s}}
                            {1-\omega_pp^{-s}},\qquad\Re s>1,
\]
and count zeros whose real parts lie in a fixed compact interval
$I\subset(1,\infty)$.  Write $N_T$ for the count with ordinates in
$[0,T]$.  The previously established pure point expansion gives
\begin{equation}\label{p03:zn:eq:variance}
 \operatorname{Var}N_T
  =\sum_{k\ne0}|c_k|^2
       \frac{4\sin^2(\lambda_kT/2)}{\lambda_k^2},
 \qquad \lambda_k=\sum_p k_p\log p .
\end{equation}
The $c_k$ are the Bragg amplitudes of this particular zero-count
observable.  The exact missing estimate is
\begin{equation}\label{p03:zn:eq:missing}
 \sum_{0<|\lambda_k|\le r}|c_k|^2=o(r)\qquad(r\downarrow0),
\end{equation}
which is equivalent to $\operatorname{Var}N_T=o(T)$.  No estimate
establishing \eqref{p03:zn:eq:missing} for the actual Euler amplitudes is
proved in this work.

Three attempted shortcuts fail for specific reasons.  First, pure point
spectrum alone controls the \emph{location} of the spectral mass, not
its accumulation at zero.  Second, an argument-principle contour can
place its right boundary in a zero-free region, but a left boundary
through the zero-bearing strip has a nontrivial winding count; one
cannot replace it by a globally bounded continuous logarithm.
Third, a finite prime cutoff permits Diophantine lower bounds for
nonzero frequencies, but those constants depend on the cutoff and
the Fourier degrees.  Neither their passage to infinitely many primes
nor the necessary regularity of the sharp zero-count observable is
established.

For clarity, the first failure can be demonstrated on the \emph{same}
prime torus, although with a random density measure instead of a zero
point process.  This is only a test of the proposed spectral argument.

\begin{proposition}[Pure point spectrum is insufficient even on the prime torus]
\label{p03:zn:prop:spectral-obstruction}
There exists a positive continuous function $A$ on
$\Omega=\prod_p\mathbb T$ such that the stationary random measure
$A(R_t\omega)\,dt$, where
$R_t\omega=(\omega_pp^{-it})_p$, is a factor of the prime Kronecker
flow and has pure point spectrum, but its mass $M_T$ in $[0,T]$
satisfies
\[
             \limsup_{T\to\infty}\operatorname{Var}(M_T)/T=\infty.
\]
\end{proposition}
\begin{proof}
Choose distinct rationals $r_j>1$ with
$\lambda_j=\log r_j<2^{-4j}$.  Unique factorization gives distinct
characters $\chi_{k_j}$ with frequency $\lambda_j$.  Set
\[
 A(\omega)=2+\sum_{j\ge1}2^{-j}\operatorname{Re}\chi_{k_j}(\omega).
\]
The series converges uniformly, so $A$ is continuous and $1\le A\le3$.
Orthogonality of the distinct positive and negative frequency characters
gives
\[
 \operatorname{Var}M_T
   =\sum_{j\ge1}\frac{2^{-2j}}2
                    \frac{4\sin^2(\lambda_jT/2)}{\lambda_j^2}.
\]
At $T_j=1/\lambda_j$, its $j$th summand implies
\[
 \frac{\operatorname{Var}M_{T_j}}{T_j}
       \ge 2\sin^2(1/2)\frac{2^{-2j}}{\lambda_j}
       >2\sin^2(1/2)\,2^{2j}.
\]
Stationarity follows from Haar invariance.  Every square-integrable
observable of a factor of this Kronecker flow lies in a pure point
unitary representation, establishing the spectral claim.
\end{proof}

This counterexample concerns a continuous random density, not Euler
zeros.  It prevents an invalid inference from the previously proved
pure point theorem but does not decide either direction of
\eqref{p03:zn:eq:missing}.  Work on Gaussian Weyl--Heisenberg zeros
\cite{p03:zn:FeldheimWH} exhibits both hyperuniform and non-hyperuniform
behaviour in other models; its hypotheses do not apply to the Euler
series.  The actual Euler count-variance question is a new formulation
in this project, not a published conjecture claimed to be solved.


% Consolidated section: arithmetic_new
\section{A deterministic resonance criterion and a computable prime-like flow}
\label{p03:an:sec:arithmetic}

The random-frequency theorem in the preceding manuscript~\cite{comp:part:02}
identifies uniform covering time with small-ball entropy for rapidly
decreasing weights. Its almost-sure conclusion
does not identify any specified frequency sequence. We now give a necessary
and sufficient arithmetic condition for that conclusion and a deterministic,
computable sequence satisfying it. The sequence approximates every prime
logarithm with any prescribed fixed polynomial accuracy. It is not the
sequence of prime logarithms itself.

The underlying finite-dimensional transference principle is classical.
Bounemoura~\cite{p02:pt:Bounemoura} treats ergodization by geometry of numbers,
and the Euclidean covering-radius inequality used below is due to
Banaszczyk~\cite{p03:an:Banaszczyk}; see the formulation and improvement in
Aggarwal--Stephens-Davidowitz~\cite{p03:an:ASD}. The conclusions specific to this
section are the weighted infinite-dimensional formulation, its exact-entropy
criterion, and an effective deterministic realization. Their historical
priority requires further investigation.

\subsection{A weighted dual profile controls actual covering}

Let $w_j>0$, $\sum_jw_j<\infty$, and $\lambda_j\in\mathbb R$. On the
product torus with angles of period $2\pi$, put
\[
 E_w(\theta)=\sum_jw_j(1-\cos\theta_j),\qquad
 \Phi_t\theta=\theta-t\lambda.
\]
Let $R_{w,\lambda}(\varepsilon)$ be the infimum of the times $T$ for which
every starting point meets $\{E_w<\varepsilon\}$ during $[0,T]$; it is
infinite if no such time exists. For finitely supported integer vectors set
\begin{equation}\label{p03:an:eq:profile}
 Q_w(k)^2=\sum_j\frac{k_j^2}{w_j},\qquad
 \Psi_{w,\lambda}(u)=
 \max_{0<Q_w(k)\le u^{-1/2}}|k\cdot\lambda|^{-1}.
\end{equation}
Exact resonances contribute infinity and the empty maximum is zero. The
character set is finite: a nonzero $k_j$ requires $w_j\ge u$. Let
$n(\varepsilon)$ be the least positive integer with
$\sum_{j>n}w_j\le\varepsilon/8$.

\begin{theorem}[Weighted resonance transference]\label{p03:an:thm:transfer}
There is an absolute constant $K\ge1$ such that
\begin{equation}\label{p03:an:eq:transfer}
 \frac\pi2\Psi_{w,\lambda}(\varepsilon)
 \le R_{w,\lambda}(\varepsilon/8),\qquad
 R_{w,\lambda}(\varepsilon)
 \le2Kn(\varepsilon)\Psi_{w,\lambda}
       \left(\frac{\varepsilon}{K^2n(\varepsilon)^2}\right),
\end{equation}
provided the character set in the upper bound is nonempty.
The bound concerns actual covering, without any product-kernel restriction.
\end{theorem}
\begin{proof}
For principal angles, $1-\cos\theta\ge2\theta^2/\pi^2$, whence
\[
 |k\cdot\theta|\le(\pi/\sqrt2)Q_w(k)\sqrt{E_w(\theta)}.
\]
For $Q_w(k)\le\varepsilon^{-1/2}$ the character is within angular distance
$\pi/4$ of $1$ throughout $\{E_w<\varepsilon/8\}$. Starting the character
at $-1$ keeps its angular distance from $1$ at least $\pi/2$ for
$0\le t\le\pi/(2|k\cdot\lambda|)$. This proves the lower bound, including
the infinite bound for an exact resonance.

Classical Euclidean transference supplies an absolute $C_0$ with
\begin{equation}\label{p03:an:eq:classical}
 \mu(\Lambda)\lambda_1(\Lambda^*)\le C_0n
\end{equation}
for every full-rank lattice in $\mathbb R^n$. Here $\mu$ is covering radius
and $\lambda_1$ the shortest nonzero Euclidean vector. This weaker constant
form follows from the cited Banaszczyk theorem; its sharp constant is not
needed. Fix $K\ge\max(1,2\pi C_0)$, write $n=n(\varepsilon)$ and
$P=\Psi_{w,\lambda}(\varepsilon/(K^2n^2))$, and put $T=KnP$.
We may assume $P<\infty$. Set
$D=\operatorname{diag}(\sqrt{\varepsilon/w_1},\ldots,
\sqrt{\varepsilon/w_n})$. The full-dimensional ellipsoid
\[
 \mathcal E=\{tT\lambda_{\le n}+Du:t^2+\|u\|_2^2\le1\}
\]
has support function
$h_{\mathcal E}(k)^2=T^2(k\cdot\lambda)^2+\varepsilon Q_w(k)^2$.
For every nonzero integer $k$, either $Q_w(k)\ge Kn/\sqrt\varepsilon$
or $|k\cdot\lambda|\ge1/P$. In either case
$h_{\mathcal E}(k)\ge Kn$.

Write $\mathcal E=A B_2^n$ with $A$ symmetric positive definite. The dual
lattice of $A^{-1}(2\pi\mathbb Z^n)$ is $A\mathbb Z^n/(2\pi)$ and has
shortest vector at least $Kn/(2\pi)$. Equation~\eqref{p03:an:eq:classical}
therefore gives covering radius at most $2\pi C_0/K\le1$. It follows that
$\mathcal E+2\pi\mathbb Z^n=\mathbb R^n$. Every head target is approximated
at a signed time in $[-T,T]$ with error $Du$, $\|u\|_2\le1$, whose energy
is at most $\frac12\sum_{j\le n}w_j(Du)_j^2\le\varepsilon/2$.
The arbitrary tail has energy at most $2\sum_{j>n}w_j\le\varepsilon/4$.
Thus total energy is strictly less than $\varepsilon$. Translating the
signed segment gives the same covering conclusion for $[0,2T]$.
\end{proof}

\subsection{The exact deterministic entropy criterion}

Assume
\begin{equation}\label{p03:an:eq:fast}
 \log w_j=-a j^b+o(j^b),\quad a,b>0,\qquad
 |\lambda_j|\le C(1+j)^d,
\end{equation}
and write
\[
 \ell=\log(1/\varepsilon),\qquad
 H(\varepsilon)=\frac{b}{2(b+1)}a^{-1/b}\ell^{1+1/b}.
\]
Nonmonotone weights are allowed. The elementary spectral estimates are
\begin{equation}\label{p03:an:eq:spectral}
 n(\varepsilon)\sim(\ell/a)^{1/b},\qquad
 \frac12\sum_j(L+\log w_j)_+
 \sim\frac{b}{2(b+1)}a^{-1/b}L^{1+1/b}.
\end{equation}
To verify them, sandwich late weights between $e^{-(a+\delta)j^b}$ and
$e^{-(a-\delta)j^b}$. Tail comparison gives
$\log\sum_{j>n}w_j=-an^b+o(n^b)$; summing $L-cj^b$ through $(L/c)^{1/b}$
and then sending $\delta$ to zero gives the second estimate.

For context, the established small-ball scale is
\begin{equation}\label{p03:an:eq:smallball}
 -\log\mu\{E_w<\varepsilon\}\sim H(\varepsilon).
\end{equation}
It is not a new small-deviation theorem. One direct verification uses
$\mathbb E e^{-v(1-\cos\Theta)}\le C(1+v)^{-1/2}$ for a uniform angle:
Markov's inequality at parameter $1/\varepsilon$, applied to indices with
$w_j\ge\varepsilon$, gives the lower bound for the negative log probability
with error $O(\ell^{1/b})$. For the other bound, constrain the first
$n=n(\varepsilon)$ angles by
$|\theta_j|\le\min(1,\sqrt{\varepsilon/(2nw_j)})$ and leave the tail free.
The resulting box lies in the ball. Its negative log probability is
$H(\varepsilon)(1+o(1))$ by~\eqref{p03:an:eq:spectral}, because
$n\log n=o(H(\varepsilon))$.

\begin{lemma}[Universal dual entropy obstruction]\label{p03:an:lem:lower}
Under~\eqref{p03:an:eq:fast},
\[\liminf_{\varepsilon\downarrow0}\log\Psi_{w,\lambda}(\varepsilon)
 /H(\varepsilon)\ge1,\]
with infinity allowed.
\end{lemma}
\begin{proof}
Set $n=n(\varepsilon)$ and
$D_j=\lfloor\sqrt{w_j/(n\varepsilon)}\rfloor$, $j\le n$.
There are $N=\prod_{j\le n}(D_j+1)$ integer vectors with
$0\le k_j\le D_j$. Equation~\eqref{p03:an:eq:spectral} shifted by
$\log n=o(\ell)$ gives $\log N\sim H(\varepsilon)$. Their linear forms
lie in an interval of length at most
\[
 M=\sum_{j\le n}D_j|\lambda_j|
 \le(n\varepsilon)^{-1/2}\sum_j\sqrt{w_j}|\lambda_j|.
\]
The last sum is finite. Coincident forms give an exact resonance. Otherwise
some two consecutive values differ by at most $M/(N-1)$. Their difference
vector has $Q_w(k)^2\le\sum_{j\le n}D_j^2/w_j\le1/\varepsilon$.
Hence $\log\Psi\ge\log(N-1)-\log M\ge(1-o(1))H$, since $\ell=o(H)$.
\end{proof}

\begin{theorem}[Necessary and sufficient arithmetic condition]
\label{p03:an:thm:criterion}
Under~\eqref{p03:an:eq:fast}, the following conditions are equivalent:
\begin{align}
 \log R_{w,\lambda}(\varepsilon)&\sim H(\varepsilon),\label{p03:an:eq:c1}\\
 \log\Psi_{w,\lambda}(\varepsilon)&\sim H(\varepsilon),\label{p03:an:eq:c2}\\
 \limsup_{\varepsilon\downarrow0}
 \frac{\log\Psi_{w,\lambda}(\varepsilon)}{H(\varepsilon)}&\le1.
 \label{p03:an:eq:c3}
\end{align}
\end{theorem}
\begin{proof}
Lemma~\ref{p03:an:lem:lower} makes the last two equivalent. The lower bound
in~\eqref{p03:an:eq:transfer} shows that the first implies the last, since
$H(\varepsilon/8)\sim H(\varepsilon)$. Conversely,
$\log(K^2n(\varepsilon)^2/\varepsilon)=\ell+O(\log\ell)$, so the rescaling
in the upper bound of~\eqref{p03:an:eq:transfer} leaves $H$ unchanged to leading
order. That bound proves $\limsup\log R/H\le1$. The lower bound and
Lemma~\ref{p03:an:lem:lower}, with $8\varepsilon$ in place of $\varepsilon$,
prove the matching lower estimate.
\end{proof}

\subsection{A computable sequence near the prime logarithms}

Here computable means that a finite algorithm, given $j,q$, outputs a
rational number within $2^{-q}$ of $\lambda_j$. No feasible running-time
bound is asserted.

\begin{theorem}[Effective realization without random choices]
\label{p03:an:thm:computable}
Let $0<\rho\le1$ be rational and $A\ge1$ an integer. There is a computable
sequence with $|\lambda_j-\log p_j|\le\rho p_j^{-A}$ for all $j$, having
no nonzero finitely supported integer resonance. With $w_j=e^{-j}$,
\begin{equation}\label{p03:an:eq:computable-cover}
 \log R_{w,\lambda}(\varepsilon)
 =\tfrac14\ell^2+O(\ell\log\ell),\qquad \ell=\log(1/\varepsilon).
\end{equation}
The implicit constant may depend on $\rho,A$.
\end{theorem}
\begin{proof}
On $X=[-1,1]^{\mathbb N}$ with product uniform probability put
$r_j=\rho p_j^{-A}$ and $\lambda_j(x)=\log p_j+r_jx_j$. For $m\ge1$ define
\[
 \mathcal K_m=\{k\ne0:k_j=0\ (j>m),\ |k_j|\le D_{mj}\ (j\le m)\},
 \quad D_{mj}=\lceil u_{mj}\rceil,
\]
where $u_{mj}$ is a specified rational upper approximation to
$e^{(m-j)/2}$ with error at most $1/4$. Such approximations are obtained
by summing the Taylor series and using its geometric tail bound; fix
the first truncation which certifies the stated error. In particular,
$e^{(m-j)/2}\le D_{mj}\le e^{(m-j)/2}+5/4$.
Set $N_m=|\mathcal K_m|$, $L_m=p_m^A/(2\rho)$ and
\[
 \delta_m=\frac{2^{-m-4}}{L_m(N_m+1)},\qquad
 B_m=\bigcup_{k\in\mathcal K_m}
              \{x:|k\cdot\lambda(x)|<\delta_m\}.
\]
These are computable data, with only ceilings of rational numbers used.
Conditioning on any coordinate with $k_j\ne0$
bounds the density of the form by $L_m$. Consequently
\begin{equation}\label{p03:an:eq:effective-tail}
 \mathbb P(B_m)\le2L_mN_m\delta_m\le2^{-m-3},\qquad
 \mathbb P\left(\bigcup_{m>M}B_m\right)\le2^{-M-3}.
\end{equation}
The closed set $G=X\setminus\bigcup_{m\ge1}B_m$ has probability at least
$7/8$.

To obtain a computable point, positive measure alone is insufficient; we
use the effective tail bound. For every closed rational cylinder $C$, the
number $\mathbb P(G\cap C)$ is computable. Truncate the union at $m=M$;
its omitted probability is bounded by~\eqref{p03:an:eq:effective-tail}.
The remaining finite union of linear slabs intersected with a rational box
has computable volume. Indeed approximate the computable slab offsets by
rationals and compute inner and outer rational-polytope volumes. Their
difference tends effectively to zero: every hyperplane boundary strip of
thickness $\eta$ in the linear form has probability at most $2L_M\eta$,
and there are only finitely many boundaries. Rational box boundaries have
measure zero. This supplies certified upper and lower approximations to
$\mathbb P(G\cap C)$ to any requested absolute error.

Start with $C_0=X$. At step $s$, partition $C_{s-1}$ into finitely many
closed dyadic cylinders whose first $s$ coordinates have lengths at most
$2^{-s}$. Boundaries have probability zero, so some child has positive
$G$-intersection measure. Dovetail the computable volume approximations
until a strictly positive lower bound is certified for one child, using a
fixed enumeration to resolve ties, and call it $C_s$. This finite search
terminates. The nested cylinders specify a unique point $x^*$ coordinate
by coordinate with an effective error bound. The compact sets $G\cap C_s$
are nonempty; compactness gives $x^*\in G$. Thus $\lambda_j(x^*)$ is
computable and has the stated approximation property.

For $w_j=e^{-j}$, every integer vector with $Q_w(k)\le e^{m/2}$ belongs
to $\mathcal K_m$, so $\Psi_{w,\lambda}(e^{-m})\le\delta_m^{-1}$.
Furthermore,
\[
 \log(N_m+1)=\sum_{j=1}^m\log(2D_{mj}+1)
            =m(m-1)/4+O(m).
\]
Bertrand's postulate gives $p_m\le2^m$, hence $\log L_m=O(m)$ and
$\log\delta_m^{-1}=m^2/4+O(m)$. Every fixed nonzero finite character
belongs to some $\mathcal K_m$, proving nonresonance. Monotonicity gives
\begin{equation}\label{p03:an:eq:computable-profile}
 \log\Psi_{w,\lambda}(\varepsilon)\le\ell^2/4+O(\ell).
\end{equation}
For these exact weights the pigeonhole proof of Lemma~\ref{p03:an:lem:lower}
also gives $\log\Psi\ge\ell^2/4-O(\ell\log\ell)$: here $n=\ell+O(1)$,
and the spectral sum with $L=\ell-\log n$ differs from $\ell^2/4$ by
$O(\ell\log\ell)$. Theorem~\ref{p03:an:thm:transfer} changes $\ell$ by only
$O(\log\ell)$ and proves~\eqref{p03:an:eq:computable-cover}.
\end{proof}

The proof specifies a deterministic terminating rule at every requested
precision. Its potentially enormous computational cost is a separate issue.
Closeness to prime logarithms does not justify a limit passage to the exact
prime flow: the construction controls one fixed sequence, and its constants
and time scales need not be uniform as the perturbations tend to zero.

\subsection{The unresolved arithmetic obstacle at exact primes}

For the actual prime flow with $w_p=p^{-\sigma}$, $\sigma>3$, put
$S_\sigma=\sum_p p^{-\sigma}(\log p)^2<\infty$. Unique factorization
associates to each nonzero finite integer vector distinct positive integers
$N,D$ with $k\cdot\lambda=\log(N/D)$. Cauchy--Schwarz gives
\[
 \log\max(N,D)\le\sum_p|k_p|\log p\le S_\sigma^{1/2}Q_w(k).
\]
The mean-value theorem for $\log$ implies
\begin{equation}\label{p03:an:eq:prime-barrier}
 |k\cdot\lambda|\ge e^{-S_\sigma^{1/2}Q_w(k)},\qquad
 \log\Psi_{w,\log p}(\varepsilon)\le S_\sigma^{1/2}\varepsilon^{-1/2}.
\end{equation}
This is the elementary height barrier, not an improvement. The preceding
manuscripts give entropy exponent $1/(\sigma-1)$ for the lower covering
bound but upper exponent $1/2$. The new transference theorem does not close
that gap: the head dimension now grows as a power of $1/\varepsilon$, so
rescaling by $n(\varepsilon)^2$ is no longer negligible. Substitution of
\eqref{p03:an:eq:prime-barrier} into~\eqref{p03:an:eq:transfer} is actually weaker
than the earlier optimized positive-kernel upper bound.

The exact-prime attempt therefore remains open. The proved new conclusions
within this project are the deterministic fast-weight criterion and its
computable prime-like realization. Neither proves a new zero-free region
for the Riemann zeta function or a theorem about its critical-line zeros.


% Consolidated section: analytic_new
\section{Two limits of supplementary information in analytic inverse problems}
\label{p03:an:section}

The results here address two precisely stated inverse problems.  Neither
is advertised as a resolution of a previously named open problem.  The
first combines convolution ambiguity with a topological obstruction to
arbitrary finite continuous side information.  The second separates
analytic uniqueness from statistical recovery when the analytic class
has no common derivative envelope.  The principal imported ingredients
are classical: the Borsuk--Ulam theorem and localized smooth
annihilators for a summable M\"untz sequence.

\subsection{Complete diffraction and arbitrary finite continuous side data}

For $f\in L^1(\mathbb R)$ put
$\widehat f(\xi)=\int e^{-i\xi x}f(x)\,dx$ and
$\mathcal I(f)=|\widehat f|^2$, with the whole real line observed.
A density is $\kappa$-strongly log-concave when
$e^{\kappa x^2/2}f(x)$ is log-concave.  Gevrey order $\sigma$ means
$\|f^{(r)}\|_\infty\le AB^r(r!)^\sigma$ for some $A,B$.

\begin{theorem}[A uniform obstruction to finite continuous calibration]
\label{p03:an:finite-side}
Fix $N\ge1$, $L,\kappa>0$, $\sigma>1$, and $K\ge1$.
There exist a compact family $\mathcal F\subset C^\infty(\mathbb R)$
and $d_N>0$, chosen independently of the observation map below, with
the following properties.  Every member of $\mathcal F$ is a mean-zero
probability density supported exactly on $[-L,L]$, positive inside,
$\kappa$-strongly log-concave, and Gevrey of order $\sigma$ with common
constants.  All members have the same moments through order $K$.
For every continuous map $A:\mathcal F\to\mathbb R^N$ there exist
$f,g\in\mathcal F$ such that
\[
 \mathcal I(f)=\mathcal I(g),\qquad A(f)=A(g),\qquad
 \|f-g\|_1\ge d_N.
\]
These two densities are not translates or reflected translates of
one another.  Continuity of $A$ is required for the entire observation
map, with $\mathcal F$ carrying its $C^\infty$ topology.
\end{theorem}

\begin{proof}
Set $d=N+1$, and use $d+1$ convolution factors, indexed by $0,\ldots,d$.
Let $a=L/(d+1)$ and $c=(d+1)\kappa$.  Choose strictly increasing odd
integers $K<m_0<m_1<\cdots<m_d$.  Put $p=(\sigma-1)^{-1}$ and let
\[
 b(x)=C\exp\bigl(-cx^2-(a-x)^{-p}-(a+x)^{-p}\bigr)
                 \mathbf1_{|x|<a}
\]
have integral one.  This is an even Gevrey density with
$(\log b)''\le-2c$ in its support.  For each $j$ choose an even,
nonnegative, nonzero Gevrey bump $\psi_j$ supported in $(-a/2,a/2)$,
and put $h_j=\psi_j^{(m_j)}$.  Integration by parts gives
\[
 \int x^r h_j(x)\,dx=0\quad(r<m_j),\qquad
 \int x^{m_j}h_j(x)\,dx=-m_j!\int\psi_j\ne0.
\]
All $h_j$ are odd.  Choose positive $\varepsilon_j$ sufficiently small
that $b+u\varepsilon_j h_j$ is positive and $c$-strongly log-concave
for every $|u|\le1$.  This follows from a uniform $C^2$ perturbation
on the compact support of $h_j$, with the strict curvature reserve
provided by $b$.  Mass and mean remain one and zero.

Fix $q_0=b+\varepsilon_0h_0$.  For
$t=(t_1,\ldots,t_d)\in S^{d-1}=S^N$ define
\[
 q_{j,t_j}=b+t_j\varepsilon_jh_j\quad(1\le j\le d),\qquad
 f_t=q_0*q_{1,t_1}*\cdots*q_{d,t_d}.
\]
Convolution of densities of curvatures $c_0,\ldots,c_d$ has curvature
at least $(\sum_jc_j^{-1})^{-1}$: completing the square and applying
the log-concavity theorem for marginals proves this rule.  Thus $f_t$
has curvature at least $c/(d+1)=\kappa$.  Its support and interior
positivity follow by convolution of interval-supported positive
densities.  Differentiation of the fixed factor gives
$\|f_t^{(r)}\|_\infty\le\|q_0^{(r)}\|_\infty$, which proves common
Gevrey bounds.  The map $t\mapsto f_t$ is continuous in $C^\infty$
(indeed polynomial in the finitely many coordinates), so its image
$\mathcal F$ is compact.  The moment cancellations give common moments
through $K$.

Since $q_{j,-t_j}(x)=q_{j,t_j}(-x)$, reflection conjugates each varying
factor's Fourier transform.  Consequently
$\mathcal I(f_t)=\mathcal I(f_{-t})$ for every $t$.
It remains to check that this antipodal pair is uniformly nontrivial.
The first nonzero odd cumulant of $q_0$ has order $m_0$ and is nonzero.
For $j\ge1$, all odd cumulants of $q_{j,t_j}$ below $m_j$ vanish and
its $m_j$th cumulant is $t_j\beta_j$ for a fixed $\beta_j\ne0$.
This follows from the moment--cumulant identities and the displayed
moment cancellations.  Cumulants add under convolution.  If $j$ is
the first index with $t_j\ne0$, the $m_j$th cumulants of $f_t$ and
$f_{-t}$ differ by $2t_j\beta_j\ne0$.  Their common nonzero $m_0$th
odd cumulant rules out reflection, while their common mean rules out
translation.  In particular,
\[
 d_N:=\min_{t\in S^N}\|f_t-f_{-t}\|_1>0
\]
by compactness and continuity.  Finally the continuous map
$t\mapsto A(f_t)$ from $S^N$ to $\mathbb R^N$ has an antipodal
coincidence by the Borsuk--Ulam theorem.  That pair proves the claim.
\end{proof}

\begin{corollary}[Finite sensor values and a decision-theoretic floor]
\label{p03:an:side-risk}
The theorem applies to any fixed finite list of complex Fourier values,
heat values at arbitrary sensor locations and positive times, moments,
or continuous nonlinear summaries, counting each real component as
one observation.  Any reconstruction from $(\mathcal I(f),A(f))$ has
worst-case $L^1$ error at least $d_N/2$.  The same lower bound holds
for worst-case expected error in any statistical experiment whose law
depends only on these data.
\end{corollary}
\begin{proof}
The listed evaluations are continuous on $\mathcal F$.  The two
densities in the theorem have identical data, hence identical data
laws.  The triangle inequality, integrated under that common law,
gives the risk bound.
\end{proof}

\paragraph{Scope and prior-art boundary.}
Convolution-factor reflection and Borsuk--Ulam information lower bounds
are classical.  The additional construction realizes a uniformly
separated antipodal obstruction under simultaneous support, curvature,
smoothness, and moment constraints while the entire Fourier intensity
is supplied for free.  Priority for this combination is unverified.
It strengthens the earlier fixed-moment ambiguity construction by
allowing arbitrary finite continuous side data, but supplies no rate
for $d_N$ as $N$ grows.  It is not a claim about arbitrary adaptive
or randomized sensing.  In particular, Krieg--Novak--Ullrich
\cite[Theorem~1 and Section~3]{p03:an:KNU} show that individually continuous
adaptive measurements can have a discontinuous joint transcript and
escape naive dimension bounds.  Fixed sensors, and feedback designs
whose complete transcript is continuous, are covered here; arbitrary
switching is not.  Nor does the result cover a second complete
diffraction pattern, which is an infinite-dimensional observation.

\subsection{Analytic uniqueness without any uniform statistical modulus}

Fix $0<a<b$ and distinct $t_j\downarrow0$ with $\sum_jt_j<\infty$.
Write
\[
 (\mathcal Hf)_j=(4\pi t_j)^{-1/2}
                    \int_a^b e^{-x^2/(4t_j)}f(x)\,dx.
\]
Analyticity below is a property on a neighborhood of the closed source
interval.  The density is set to zero outside $[a,b]$ and is not
asserted to be globally analytic across its support endpoints.

\begin{theorem}[A statistical discontinuity inside the analytic class]
\label{p03:an:analytic-barrier}
Fix $\kappa>0$ and an integer $R\ge2$.  There exist $m,M,d>0$ such
that the class $\mathcal A$ of probability densities on $[a,b]$
which extend to entire functions and satisfy
\[
 f\ge m,\qquad \|f\|_{C^R([a,b])}\le M,\qquad
                  (\log f)''\le-\kappa
\]
has the following properties:
\begin{enumerate}
\item $\mathcal H$ is injective on $\mathcal A$.
\item There are $f_n^+,f_n^-\in\mathcal A$ with
\[
 \|f_n^+-f_n^-\|_1\ge d,
 \qquad
 \|\mathcal Hf_n^+-\mathcal Hf_n^-\|_{\ell^2}
 +\|\mathcal Hf_n^+-\mathcal Hf_n^-\|_{\ell^\infty}\longrightarrow0.
\]
\item With adversarial coordinatewise noise bounded by any $\delta>0$,
the minimax $L^1$ reconstruction error is at least $d/2$, even if
the whole infinite sequence is observed.  With independent Gaussian
noise of any standard deviation $\delta>0$ at every sampled time,
the minimax expected $L^1$ error is at least $d/4$.
\end{enumerate}
The constants do not depend on the noise level or the number of
observed coordinates.  No common complex-neighborhood norm or
quasianalytic derivative envelope is imposed.
\end{theorem}

\begin{proof}
For injectivity let $h=f-g$ be nonzero and analytic near $[a,b]$.
There is a least $r\ge0$ with $h^{(r)}(a)\ne0$.  Substituting
$x=a+u/s$ and using Taylor expansion with exponential domination gives
\[
 \int_a^b e^{-sx^2}h(x)\,dx
 \sim \frac{h^{(r)}(a)e^{-sa^2}}{(2a)^{r+1}s^{r+1}}
                  \qquad(s\to\infty).
\]
For completeness, on a fixed endpoint neighborhood the leading term
is $h^{(r)}(a)u^r/(r!s^r)$ and the remaining exponential tends to
$e^{-2au}$; the integral of $u^re^{-2au}$ is $r!/(2a)^{r+1}$.
The part away from $a$ is exponentially smaller.  The transform
therefore has no sufficiently large positive zeros.  Since
$s_j=(4t_j)^{-1}\to\infty$, equality of all observations is impossible.

The smooth annihilator construction for summable M\"untz sequences
provides a nonzero real $h\in C_c^\infty((a,b))$ with
\[
 \int h=0,\qquad \int_a^b e^{-x^2/(4t_j)}h(x)\,dx=0
                 \quad(j\ge1).
\]
One may obtain this directly from the localized Mellin annihilator
with zeros $s_j$, by the substitution $y=e^{-x^2}$; its existence
uses $\sum_j1/s_j<\infty$.  This is the classical existence input
of Luxemburg--Korevaar, in the explicit form restated by
Wellershoff~\cite[Theorem~21]{p01:mom:Wellershoff2021}.

Let $f_0=C e^{-\kappa x^2}$ on $[a,b]$, normalized to have mass one.
Its logarithmic curvature is $-2\kappa$.  Choose $\varepsilon>0$
small enough that $f_0\pm\varepsilon h$ are positive and have
logarithmic curvature at most $-3\kappa/2$.  There are polynomials
$p_n$ with $p_n\to h$ in $C^R([a,b])$ and $\int p_n=0$.
Indeed, approximate $h^{(R)}$ uniformly by polynomials, integrate
$R$ times with the prescribed endpoint derivatives, and subtract
the vanishing constant mass correction.  Put
$f_n^\pm=f_0\pm\varepsilon p_n$.
These functions are entire, have mass one, and for all sufficiently
large $n$ satisfy the asserted bounds with fixed $m,M$ and curvature
at most $-\kappa$.  Their $L^1$ separation tends to
$2\varepsilon\|h\|_1>0$; reduce $d>0$ and discard a finite prefix.

Let
$k_j=(4\pi t_j)^{-1/2}e^{-a^2/(4t_j)}$.
Then $k\in\ell^2\cap\ell^\infty$: for small $t$,
$t^{-1}e^{-a^2/(2t)}\le C_at$, and $\sum_jt_j<\infty$.
Consequently
\[
 \|\mathcal H(f_n^+-f_n^-)\|_{\ell^q}
 \le 2\varepsilon\|k\|_{\ell^q}\|p_n-h\|_1\longrightarrow0,
                       \qquad q=2,\infty.
\]
The midpoint of two data sequences whose sup distance is at most
$2\delta$ is admissible under both hypotheses.  The triangle
inequality proves the adversarial-noise lower bound.

For Gaussian noise choose $n$ so large that the $\ell^2$ distance
of the two mean sequences is at most $\delta$.  Their product Gaussian
laws have relative entropy at most $1/2$: the finite-dimensional
formula is $\|\Delta\|_2^2/(2\delta^2)$, and the same formula for
the infinite product follows by convergence of the Gaussian
likelihood ratios when $\Delta\in\ell^2$.  Pinsker's inequality gives
total variation at most $1/2$.  Integration against the common part
of the two laws gives worst-case expected metric loss at least
$d(1-\mathrm{TV})/2\ge d/4$.
\end{proof}

\paragraph{What this does and does not settle.}
The theorem gives a precise obstruction to seeking a noise rate from
analytic or quasianalytic membership alone.  It persists with any
fixed finite number of uniform derivative bounds.  A fixed full
quasianalytic envelope, by contrast, can yield a compact class and a
uniformly continuous inverse; its sharp modulus is a different problem
and is not computed here.  The mechanism is density of analytic
functions in a smooth nonidentifiable class, so this is a rigorous
barrier, not a new solution of a general inverse-heat stability problem.

\clearpage
\section{Seventeen research routes: exact outcomes and remaining barriers}
\label{p03:end:register}

This register counts investigations, not independent historically open
conjectures. ``Complete'' means complete for the explicitly stated
question in that row. ``Partial'' records a proved subdomain of a
larger question. A supporting lemma, a numerical illustration, or a
corollary is not an additional investigation. The seven complete rows,
three partial rows, and seven unresolved rows below are therefore not
a claim to have solved seven or seventeen famous conjectures.

\begingroup
\fontsize{9}{10.8}\selectfont
\setlength{\tabcolsep}{4pt}
\renewcommand{\arraystretch}{1.12}
\begin{longtable}{@{}p{0.6cm}>{\raggedright\arraybackslash}p{5.2cm}>{\raggedright\arraybackslash}p{8.5cm}@{}}
\toprule
ID & Exact target & Outcome and scope\\
\midrule\endfirsthead
\toprule ID & Exact target & Outcome and scope\\\midrule\endhead
1 & Published 2019 binomial-mixture MLE range conjecture &
\textbf{Complete proof in this draft.} Theorem~\ref{p03:bnew:thm:mle}
reaches $t\le N^{2/3-\varepsilon}$ for the same unrestricted MLE
and Wasserstein loss. This is the clearest direct external conjecture
resolution claimed by the manuscript.\\
2 & Near-linear MLE range discussed in the 2026 author version &
\textbf{Unresolved.} The obtained predictive Hellinger rate still
has a $t^{1/4}$ factor. Neither the exact constant-one coefficient
inequality nor the proposed dimension-free fingerprint bound is proved.\\
3 & Adaptive Bernoulli recovery under a per-individual cap $Q$ &
\textbf{Partial full-phase answer.} Theorem~\ref{p03:bnew:thm:caps}
gives matching risk throughout $Q\le B^{2/5-\gamma}$ and lower
bounds for every cap. The remaining intermediate-cap range is open.\\
4 & Scalar uniform-endpoint critical logarithm &
\textbf{Unresolved.} Proposition~\ref{p03:bnew:prop:uniform-barrier}
proves only a limitation of the specified independent-level CDF
estimator, not a Wasserstein minimax lower bound.\\
5 & Sparse Gaussian terminal designs with unrestricted covariance &
\textbf{Complete.} Theorem~\ref{p03:mixnew:thm:exact-terminal}
establishes harmonic divergence as necessary and sufficient, with
joint invisible laws in every dimension when it fails. This closes
a question left in the previous draft.\\
6 & Does bounded maximum-likelihood threshold force bounded generic
completion rank? &
\textbf{Unresolved.} The Gaussian identification transfer computes
neither invariant in the required generality. The original graph
question is identified in the Gaussian section.\\
7 & Sparse random-graph likelihood/completion-rank conjectures &
\textbf{Unresolved.} The latest rigidity result located in the
source audit still leaves the global-rigidity step needed by the
likelihood conjecture. No graph-conjecture closure is inferred.\\
8 & Full ID/SD domain of the third McKay family &
\textbf{Partial.} Theorem~\ref{p03:bn:third} proves ID and an exact
SD threshold on every integer-bias ray $n\ge2$, $\mu>-1/2$.
Noninteger parameters and further negative orders remain open.\\
9 & Full SD domain of the fourth McKay family &
\textbf{Partial.} Theorem~\ref{p03:bn:fourth} treats all half-integer
orders above $1/2$. Theorem~\ref{p03:bn:three-halves} uniquely specifies
the $3/2$ threshold and proves an exact rational non-SD example.
These are parts of the same published multipart problem as row 8.\\
10 & Every modular period for irrational-frequency trinomials &
\textbf{Complete for three terms.}
Theorem~\ref{p03:zn:thm:classification} determines density and missing
arcs for every resonance. General four-or-more-term classification
and resonant equidistribution are not included.\\
11 & Unknown Euler frequencies from a right zero divisor &
\textbf{Unresolved in the Euler class.}
Proposition~\ref{p03:zn:prop:nonidentifiability} disproves the broader
normalized Dirichlet-sum formulation. The counterexample does not
preserve the required Euler-factor structure.\\
12 & Actual Euler critical-zero hyperuniformity &
\textbf{Unresolved.} The required low-frequency Bragg-weight
estimate is displayed in \eqref{p03:zn:eq:missing}. The accompanying
counterexample invalidates a spectral shortcut, not the conjectured
Euler conclusion itself.\\
13 & A deterministic exact-entropy criterion for fast weights &
\textbf{Complete.} Theorem~\ref{p03:an:thm:criterion} gives a
necessary and sufficient dual-resonance condition for actual
uniform covering. The finite-dimensional transference input is
classical and explicitly attributed.\\
14 & A computable deterministic entropy-optimal frequency sequence &
\textbf{Complete.} Theorem~\ref{p03:an:thm:computable} constructs
one with prescribed polynomial closeness to every prime logarithm.
It is not the sequence of prime logarithms; feasible runtime is
not claimed.\\
15 & The exact-prime recurrence exponent gap for $\sigma>3$ &
\textbf{Unresolved.} The height bound in
\eqref{p03:an:eq:prime-barrier} and the dimension loss in transference
do not improve the previous optimized upper bound.\\
16 & Complete Fourier intensity plus finite continuous side data &
\textbf{Complete negative answer in the stated class.}
Theorem~\ref{p03:an:finite-side} constructs uniformly separated
strongly log-concave Gevrey ambiguities for every jointly continuous
finite-dimensional side map. Arbitrary switching is not covered.\\
17 & A uniform noise modulus from analytic heat-source membership &
\textbf{Complete negative answer without a common full derivative
envelope.} Theorem~\ref{p03:an:analytic-barrier} gives a fixed minimax
floor despite uniqueness and common finite-order derivative bounds.
A sharp modulus under a fixed full envelope remains a different,
unresolved question.\\
\bottomrule
\end{longtable}
\endgroup

\subsection{Which changes are substantial, and which are supporting tools?}

The binomial theorem is the strongest link to a previously printed
conjecture. It combines a forward complexity estimate with an inverse
coefficient estimate. Controlling only one side would not give the
stated range: the generic simplex complexity retains a larger trial
factor, while a polynomial prefactor in the inverse map consumes the
gain from the smaller observation class. Their simultaneous removal
to the stated orders is the relevant proof mechanism.

The sparse Gaussian theorem has a different source of difficulty.
For each fixed mean frequency, constructing a scalar vanishing
function is not enough: such functions might fail to arise from any
single joint law. The matrix-tube realization solves this consistency
problem and gives a Schwartz signed density supported on the
covariance cone. Its vanishing marginals show that the hidden
information can lie entirely in dependence between mean and
covariance, even after every joint polynomial moment is also supplied.

The three-term zero theorem classifies all exceptional periods
instead of merely proving existence of good periods. Its geometry
is elementary once the appropriate lifted torus is used. This makes
it a concise complete classification; it does not establish a new
general method for the critical strip of the Riemann zeta function.

The Bessel results are genuine extensions beyond the previous GGC
boundary, but they are partial classifications of the full published
questions. The exact threshold expressed by a positive finite
spectral kernel is meaningful because all its zeros and its
positivity have been proved. The unique order-$3/2$ maximum and the
rational counterexample give an independently checkable instance.
Counting every integer order as a separate solved conjecture would
misstate the contribution.

The deterministic recurrence and inverse-problem results are new
formulations within this programme. They specify when a transfer is
valid and when extra information remains insufficient. Their
classical ingredients receive credit. They should be judged by the
precise conclusions and their usefulness, not by a nominal count
of solved problems.

\subsection{The next bottlenecks are not removed by aggregation}

For the binomial MLE, reaching a near-linear number of trials needs
more than the global Hellinger estimate proved here. In the present
argument the comparison to the target loss has the factor
$t^{3/4}/\sqrt N$, up to logarithms and a small adjustable power of
$N$. This locates the $2/3$ boundary of this particular proof.
Functional estimates sharper than global Hellinger control, or a
different inverse argument, would be needed to go beyond it.

For the McKay laws, nonterminating hypergeometric logarithmic
derivatives replace the finite root sums used here. Their additional
contributions must be controlled before a full real-parameter SD
classification is justified. For exact prime frequencies, an
almost-sure theorem or a computable nearby sequence supplies no
estimate at the prescribed arithmetic point. For Euler variance,
pure point spectrum supplies no bound on the mass accumulating at
arbitrarily small nonzero frequencies.

These are concrete missing estimates. Combining the completed
sections does not by itself prove them, and no Annals-level or
Fields-level assessment is inferred from combining the papers.
The immediate review priority is the complete binomial MLE argument,
followed by the matrix-tube construction and the resonant-period
classification.


\clearpage
\begingroup\small

\endgroup
\endgroup
% ===== END SOURCE 03 =====


% ===== BEGIN SOURCE 04: KEIO_Transport_Spectral_Positivity_and_Observation.tex =====
\clearpage
\part{Transport, Spectral Positivity, and Sharp Observation Boundaries}
\label{part:04}
\setcounter{section}{0}
\begin{center}\large\itshape Near-linear Mixture Likelihood, Full McKay Domains, and Critical Sampling Budgets\end{center}
\begin{center}\small Source manuscript: 14 October 2026.\end{center}
\begingroup
\renewcommand{\arraystretch}{1}
\begin{quote}\small\textbf{Source abstract.}\quad
We develop six continuations of a program on observation and recovery.
For unrestricted binomial mixtures we prove that every maximum likelihood
estimator attains the minimax Wasserstein order $(t\log N)^{-1/2}$ whenever
$a\log N\le t\le N^{1-\varepsilon}$, for fixed $a,\varepsilon>0$.
The proof transfers kernel transportation inequalities through the
likelihood score equations and then uses derivative-stable polynomial
inversion. Under a total observation budget we determine the full
predetermined Bernoulli replication-cap curve and prove that the rate
$(\log B)/\sqrt B$ is unavoidable for uniform-endpoint mixtures even
under arbitrary adaptive sampling. For two modified-Bessel McKay families
we establish infinite divisibility throughout their admissible parameter
domains; the first answers a maximal-domain question posed in the
published literature. Signed hypergeometric spectral measures have
positive inverse Laplace kernels despite their changing signs.
Further results locate a logarithmic tail threshold for sparse Gaussian
identification, characterize random-clock L\'evy snapshot identification,
classify all modular zero periods of independent exponential sums, and
recover a hidden finite spectrum from its density-transition curve.
Finally, a fixed analytic envelope gives exponential exact-data recovery
but no sampling-sequence-uniform noise rate. Full arguments and an
eighteen-target register distinguish established statements, partial
advances, and unresolved transfers. No Riemann-hypothesis or general
prime-statistics conclusion is claimed.
\end{quote}


\section{What changes from the previous program}
\label{p04:nxo:introduction}

The 223-page manuscript \emph{Observation and Recovery in Arithmetic
and Probability} and its three sequels
\cite{p01:KEIOIntegrated,comp:part:01,comp:part:02,comp:part:03}
developed inverse questions in probability, special functions, and
exponential-sum dynamics. The present paper contains new continuations,
rather than another concatenation of those earlier texts. Definitions
and the arguments for the new claims are given here; the earlier papers
are used to identify which boundaries were previously unresolved.
Their results are not assumed to be externally certified.

The most substantial advances come from changing which intermediate
quantity is controlled. In the likelihood problem, total variation of
the count fingerprint is unnecessarily strong. The transport geometry
of that fingerprint retains exactly the information needed by a
Lipschitz inverse test. In the special-function problem, positivity of
each spectral contribution is unnecessarily restrictive. A signed
representation can still have a positive inverse Laplace kernel, once
its total mass and pole locations are controlled. In adaptive sampling,
a bound on the information supplied by an entire individual can obscure
the decisive issue: the conditional information gained by one additional
observation. We obtain an exact posterior identity for that increment.
These are three distinct mechanisms; the paper does not assert a single
universal theorem subsuming all six branches.

\subsection{The principal quantitative advances}

Let $N$ be the number of independent binomial groups and $t$ their
common size. For their latent mixing law $G$, Theorem~\ref{p04:nxm:thm:main}
proves, for every unrestricted MLE $\widehat G$,
\[
 \sup_G\mathbb E_G W_1(\widehat G,G)
 \le \frac{C_{\varepsilon,a}}{\sqrt{t\log N}},
 \qquad a\log N\le t\le N^{1-\varepsilon}.
\]
This improves the exponent $2/3$ in our immediate predecessor to every
exponent below one. It reaches the near-linear range discussed
conditionally by Musco et al.~\cite[Section~6.2]{p01:info:MMRS2025}, with a
matching lower bound from Vinayak et al.~\cite{p01:info:VKVK2019}.
It does not prove the separate constant-one coefficient conjecture.
The all-$N,t$ interpolation estimate is an upper bound, and its entire
transition is not asserted to be minimax sharp.

For a total raw-observation budget $B$ and per-individual cap $Q$, our
predetermined Bernoulli design has the exact risk order
\[
 \max\left\{B^{-1/3},\ Q^{-1},\
       [Q\log(e+B/Q^{3/2})]^{-1/2}\right\}.
\]
The proof optimizes a frequency-dependent allocation; it is not a claim
about the equal-group MLE. For adaptive designs the same analysis gives
an exact order through every $Q\le B^{2/3-\gamma}$, extending the
earlier $2/5$ exponent, but leaves a critical window open.
For $\operatorname{Uniform}[0,\theta]$ individuals with
$\theta\in[1,2]$, Theorem~\ref{p04:nxb:thm:uniform} closes the logarithmic
gap completely:
\[
 U_B^{\mathrm{ad}}\asymp\frac{\log B}{\sqrt B}.
\]
The lower bound holds for every adaptive policy, even when the latent
log-parameter has a positive uniformly Lipschitz density.

Theorems~\ref{p04:nxs:full-id} and~\ref{p04:nxs:fourth-full} establish ID for
the normalized densities
\begin{align*}
 Cx^{\nu-1}e^{-bx}I_\mu(ax),&\quad
     \mu>-1,\quad \mu+\nu>0,\quad b>a>0,\\
 Cx^{2\mu}e^{-bx}I_\mu(ax)^2,&\quad
     \mu>-\tfrac14,\quad b>2a>0.
\end{align*}
The first gives the whole domain requested in
\cite[Section~2.3]{p01:rare:Baricz2026}. The second extends that paper's
Theorem~4 to every admissible real order. Explicit L\'evy kernels also
give exact spectral SD tests. They do not give an elementary expression
for every threshold, nor settle every strictness question.

\subsection{Identification, zero geometry, and stability}

The Gaussian threshold uses
$\omega_a(v)=\sqrt v/\log^a(e+v)$. An exponential moment in
$\omega_a(\operatorname{tr}\Sigma)$ with $a\le1$ makes every infinite
set of terminal lengths identifying. For $a>1$, a scalar construction
with dyadic terminal lengths defeats identification while preserving
both marginals and every joint polynomial moment. The sharpness is
scalar; additional anchors obstruct a direct counterexample in a fixed
matrix dimension greater than one. A parallel Fourier--Laplace argument
gives the exact sparse-time criterion for dependent random locations
and L\'evy clocks.

For finite rationally independent exponential sums, every period is
classified, including its missing arcs. Only the largest frequency
can create an exception. Removing the largest frequency in a normally
convergent infinite sum removes all exceptions. An amplitude-dependent
transition then determines all magnitudes and relative frequencies of
an unknown finite spectrum. Its observation is an entire exact
transition curve, and several phases remain invisible. It is not a
reconstruction of primes from a single rightmost zero.

The analytic inverse results separate exact information from noisy
information on one fixed class. Any $n$ distinct positive heat
parameters give error $\exp(-\Theta(n))$, uniformly over their
locations. Yet no noise-convergence modulus works uniformly over all
sampling sequences, even under that same analytic envelope. For every
fixed sequence the risk still tends to zero. Thus the negative claim
is not a positive asymptotic error floor on a compact analytic class.

\subsection{How to read the claims}

Section~\ref{p04:nxo:register} records eighteen targets: twelve completed
statements at their specified scope, three partial advances, and three
unresolved transfers. These are research targets, not eighteen named
conjectures. Among the completed statements, the third-McKay domain
problem is a directly located published open question, and the MLE
result reaches a published conditional target. Other rows are full
extensions or internally formulated inverse and minimax problems.
Consequences of the same theorem are not counted as additional famous
problem resolutions.

Primary-source comparisons were checked through 14 October 2026,
including the May 2026 version of \cite{p01:info:MMRS2025}, the published
PJM version of \cite{p01:rare:Baricz2026}, and the recent distinct smooth-density
problem in \cite{p01:info:LeeKennedy2026}. This is a bounded literature
comparison, not a certification of priority. The proofs have undergone
cross-checks within this AI-assisted drafting process. External expert
review, and therefore any claim of final acceptance or exceptional
historical importance, remain outstanding.


% ---- BEGIN mle ----
\section{Near-linear optimality of the unrestricted binomial-mixture MLE}
\label{p04:nxm:sec:main}

The earlier manuscript established optimality through the replication range
$t\le N^{2/3-\varepsilon}$ by controlling the entire predictive distribution
in Hellinger distance. Here we change the statistical geometry: the relevant
fingerprint error is transportation distance on the normalized count grid,
not total variation. Likelihood self-consistency transfers a transportation
inequality for each binomial kernel to the fitted-versus-empirical discrepancy.
Together with derivative-stable polynomial inversion, this gives the
near-linear range discussed in \cite[Section~6.2]{p01:info:MMRS2025}.

We use natural logarithms. For a probability measure $G$ on $[0,1]$, write
\[
 B_j^t(x)=\binom tjx^j(1-x)^{t-j},\qquad
 p_G(j)=\int B_j^t(x)\,dG(x),\quad 0\le j\le t.
\]
The observations $X_1,\ldots,X_N$ are independent with common law $p_G$.
Their empirical probability vector is $q_j=N^{-1}\#\{i:X_i=j\}$.
Let $\widehat G$ be any maximizer of
$\sum_jq_j\log p_H(j)$ over all probability measures $H$ on $[0,1]$,
and write $\widehat p=p_{\widehat G}$.
Such maximizers exist by compactness of the binomial moment body. Every
maximizer assigns positive probability to each observed count.
Write $W_{1,t}$ for Wasserstein distance on $\{0,\ldots,t\}$ with cost
$|j-l|/t$, and $D$ for relative entropy, with the usual extended-value
conventions.

\begin{theorem}[Near-linear MLE optimality and a uniform interpolation bound]
\label{p04:nxm:thm:main}
There are absolute constants $C,A$ such that for all integers $N,t\ge1$,
all mixing laws $G$, and every unrestricted mixing-measure MLE,
\begin{equation}\label{p04:nxm:eq:allrange}
 \mathbb E_G W_1(\widehat G,G)
 \le C\left\{\frac1t+
       \frac1{\sqrt{t\log(1+N/t)}}\right\}.
\end{equation}
More specifically, for $2\le k\le t/2$,
\begin{equation}\label{p04:nxm:eq:masterexpectation}
 \mathbb E_G W_1(\widehat G,G)
 \le \frac Ck+\frac{C\exp(Ak^2/t)}{\sqrt N}.
\end{equation}
For every fixed $\varepsilon\in(0,1)$ and $a>0$, there are constants
$C_{\varepsilon,a},N_{\varepsilon,a}$ such that, whenever
\[
 N\ge N_{\varepsilon,a},\qquad
 a\log N\le t\le N^{1-\varepsilon},
\]
\begin{equation}\label{p04:nxm:eq:nearlinear}
 \mathbb P_G\left\{W_1(\widehat G,G)>
       \frac{C_{\varepsilon,a}}{\sqrt{t\log N}}\right\}\le N^{-2},
 \qquad
 \mathbb E_G W_1(\widehat G,G)
       \le\frac{C_{\varepsilon,a}}{\sqrt{t\log N}}.
\end{equation}
The probability statement holds simultaneously for all MLE solutions.
\end{theorem}

The theorem makes no smoothness, separation, finite-support, or
interiority assumption on $G$. Its estimator is the original unrestricted
MLE. The bound in \eqref{p04:nxm:eq:nearlinear} has the matching minimax order
from \cite[Theorem~3.3]{p01:info:VKVK2019}. We do not claim a matching lower
bound for the entire interpolation formula \eqref{p04:nxm:eq:allrange}; in
particular, its transition when $t$ is close to $N$ may be suboptimal.
The exact constant-one Chebyshev--Bernstein inequality appearing in
\cite[Section~6.2]{p01:info:MMRS2025} is not needed or proved here.
The recent smooth-density analysis of Lee and Kennedy
\cite{p01:info:LeeKennedy2026} concerns pointwise density loss under H\"older
regularity and a regularized series estimator. The unrestricted-law,
Wasserstein-loss MLE question addressed here is different.

\subsection{An entropy--transport principle for mixture likelihood}

The next statement isolates the structural step. The likelihood directional
condition is classical mixture geometry \cite{p04:nxm:Lindsay1983}; the
exponential-moment/transport duality used below is classical
\cite{p04:nxm:BobkovGoetze1999}. We include their elementary finite-space
arguments. The transfer does not assume that the mixture itself has the
same concentration as its individual kernels.

\begin{lemma}[Likelihood transfers kernel transportation inequalities]
\label{p04:nxm:lem:generaltransport}
Let $(\mathcal X,d)$ be a finite metric space and $K_\theta$ a family of
probability vectors. Suppose that, for every $\theta$ and every probability
vector $r$,
\begin{equation}\label{p04:nxm:eq:kernelT1}
 W_{1,d}(r,K_\theta)^2\le 2c D(r\Vert K_\theta).
\end{equation}
For a probability vector $q$, let $\widehat H$ maximize
$\sum_xq(x)\log p_H(x)$ over mixtures
$p_H=\int K_\theta\,dH(\theta)$, and suppose this maximum is finite.
Then, with $\widehat p=p_{\widehat H}$,
\begin{equation}\label{p04:nxm:eq:generaltransport}
 W_{1,d}(q,\widehat p)^2\le2cD(q\Vert\widehat p).
\end{equation}
Consequently $W_{1,d}(q,\widehat p)^2\le2cD(q\Vert p_H)$ for every
comparison mixture $p_H$.
\end{lemma}
\begin{proof}
The first directional derivative of likelihood in the direction
$\delta_\theta-\widehat H$ gives
\[
 S(\theta):=\sum_{x:q(x)>0}\frac{q(x)K_\theta(x)}{\widehat p(x)}\le1.
\]
All denominators on the displayed support are positive. Since
$\int S\,d\widehat H=\sum_xq(x)=1$, one has $S=1$ almost surely under
$\widehat H$. Thus
\[
 R_\theta(x)=\frac{q(x)K_\theta(x)}{\widehat p(x)}
 \quad(q(x)>0),\qquad R_\theta(x)=0\quad(q(x)=0)
\]
is a probability vector for almost every $\theta$, and
$\int R_\theta\,d\widehat H(\theta)=q$. The entropy chain rule reduces here
to the finite-sum identity
\begin{align*}
 \int D(R_\theta\Vert K_\theta)\,d\widehat H(\theta)
 &=\sum_x\int R_\theta(x)
       \log\frac{q(x)}{\widehat p(x)}\,d\widehat H(\theta)\\
 &=D(q\Vert\widehat p).
\end{align*}
Terms with $K_\theta(x)=R_\theta(x)=0$ are zero.
By the dual formula for Wasserstein distance, followed by
\eqref{p04:nxm:eq:kernelT1} and Cauchy--Schwarz,
\begin{align*}
 W_{1,d}(q,\widehat p)
 &\le\int W_{1,d}(R_\theta,K_\theta)\,d\widehat H(\theta)\\
 &\le\sqrt{2c\int D(R_\theta\Vert K_\theta)\,d\widehat H(\theta)}.
\end{align*}
Finally, maximizing likelihood is equivalent to minimizing
$D(q\Vert p_H)$, because $\sum_xq(x)\log q(x)$ is independent of $H$.
\end{proof}

\begin{lemma}[Binomial transportation]
\label{p04:nxm:lem:binomialT1}
For every probability vector $r$ on $\{0,\ldots,t\}$ and
$\theta\in[0,1]$,
\begin{equation}\label{p04:nxm:eq:binomialT1}
 W_{1,t}(r,\operatorname{Binomial}(t,\theta))^2
 \le \frac1{2t}D(r\Vert\operatorname{Binomial}(t,\theta)).
\end{equation}
Hence every binomial-mixture MLE obeys the deterministic inequality
\begin{equation}\label{p04:nxm:eq:deterministicprediction}
 W_{1,t}(q,\widehat p)
 \le\sqrt{\frac{D(q\Vert\widehat p)}{2t}}
 \le\sqrt{\frac{D(q\Vert p_G)}{2t}}.
\end{equation}
\end{lemma}
\begin{proof}
Let $Z_1,\ldots,Z_t$ be independent Bernoulli variables of parameter
$\theta$. For any function $u$ that is $1$-Lipschitz on the normalized
count grid, changing one $Z_i$ changes $u(\sum_iZ_i)$ by at most $1/t$.
The Doob martingale which reveals these variables one at a time has
conditional increment range at most $1/t$. Hoeffding's elementary
bounded-variable inequality, applied conditionally and iterated, gives
\[
 \log\mathbb E\exp\{\lambda(u-\mathbb Eu)\}
 \le\lambda^2/(8t),\qquad\lambda\in\mathbb R.
\]
For clarity, the bounded-variable inequality says that a centered variable
with range length $b$ has logarithmic moment generating function at most
$\lambda^2b^2/8$; it follows by convexity from the two-point distribution
on the range endpoints and two integrations of its variance bound
$b^2/4$.
For $r$ absolutely continuous with respect to the binomial law $K$,
nonnegativity of relative entropy with respect to the tilted law
$e^{\lambda u}K/\mathbb E_K e^{\lambda u}$ gives
\[
 \mathbb E_r u-\mathbb E_Ku
 \le D(r\Vert K)/\lambda+\lambda/(8t),\qquad\lambda>0.
\]
Minimization over $\lambda$ and the Wasserstein dual formula prove
\eqref{p04:nxm:eq:binomialT1}. If entropy is infinite the assertion is
automatic; the degenerate kernels at $\theta=0,1$ cause no exception.
Lemma~\ref{p04:nxm:lem:generaltransport}, with $c=1/(4t)$, proves
\eqref{p04:nxm:eq:deterministicprediction}.
\end{proof}

\begin{lemma}[Dimension-free predictive transportation error]
\label{p04:nxm:lem:predictive}
Uniformly in $G,N,t$,
\begin{equation}\label{p04:nxm:eq:predictiveexpectation}
 \mathbb E_G W_{1,t}(q,\widehat p)\le\frac1{\sqrt{2N}},\qquad
 \mathbb E_G W_{1,t}(q,p_G)\le\frac1{2\sqrt N}.
\end{equation}
Moreover, for $0<\eta<1/2$, except on an event of probability $\eta$,
simultaneously for all MLEs,
\begin{equation}\label{p04:nxm:eq:predictiveprobability}
 W_{1,t}(q,\widehat p)+W_{1,t}(q,p_G)
 \le C\sqrt{\frac{\log(N+1)+\log(2/\eta)}{N}}.
\end{equation}
\end{lemma}
\begin{proof}
If $p_G$ has support size $s\le t+1$, then $\log x\le x-1$ implies
\[
 D(q\Vert p_G)\le\sum_{j:p_G(j)>0}\frac{(q_j-p_G(j))^2}{p_G(j)},
 \qquad
 \mathbb E_G D(q\Vert p_G)\le\frac{s-1}{N}\le\frac tN.
\]
Apply \eqref{p04:nxm:eq:deterministicprediction} and Jensen's inequality.
For the second expectation bound, identify the counts with $j/t\in[0,1]$.
Their empirical and population distribution functions satisfy
\[
 \mathbb E|F_N(x)-F(x)|\le\sqrt{F(x)(1-F(x))/N}\le1/(2\sqrt N).
\]
Integrating the one-dimensional formula
$W_1(q,p_G)=\int_0^1|F_N-F|$ gives the result.

For the probability assertion, a multinomial type $v$ has probability at
most $\exp\{-ND(v\Vert p_G)\}$. Indeed, its multinomial coefficient is
at most $\exp\{NH(v)\}$, as follows by taking its term in the multinomial
expansion of $(\sum_jv_j)^N=1$. There are at most $(N+1)^s$ types.
Consequently, with probability at least $1-\eta/2$,
\[
 D(q\Vert p_G)\le\frac{(t+1)\log(N+1)+\log(2/\eta)}N.
\]
In addition, changing one observation changes $W_{1,t}(q,p_G)$ by at
most $1/N$. The same conditional bounded-variable argument used above,
followed by Chernoff's inequality, bounds it by its expectation plus
$\sqrt{\log(2/\eta)/(2N)}$ except on an event of probability $\eta/2$.
Combine these estimates with \eqref{p04:nxm:eq:deterministicprediction} and
$(t+1)/t\le2$.
\end{proof}

\subsection{Polynomial inversion that preserves a derivative bound}

The following coefficient estimate applies to every bounded polynomial,
not just to individual Chebyshev polynomials. Its contour mechanism is
the one used in the preceding manuscript, with the elementary
Bernstein--Walsh principle replacing a sum over Chebyshev modes.

\begin{lemma}[Bernstein coefficients of an arbitrary bounded polynomial]
\label{p04:nxm:lem:coefficient}
There is an absolute constant $C$ with the following property. If $P$ is
a real polynomial of degree at most $m\le n/2$ and
\[
 P(x)=\sum_{j=0}^n b_jB_j^n(x),
\]
then
\begin{equation}\label{p04:nxm:eq:coefficient}
 \max_j|b_j|\le C e^{2m^2/n}\|P\|_{L^\infty[0,1]}.
\end{equation}
\end{lemma}
\begin{proof}
Write $M=\|P\|_\infty$. For a complex number $w$ with $|w|\ge1$, the
maximum principle applied to the polynomial
$v^mP((2+v+v^{-1})/4)$ on $|v|\le1$ gives
\[
 \left|P\left(\frac{2+w+w^{-1}}4\right)\right|\le M|w|^m.
\]
For $0<j<n$, set $r=j/(n-j)$, $q_*=r/(1+r)^2$, $z=re^{i\phi}$,
and $u=2\sqrt{q_*}|\sin(\phi/2)|$. The two inverse Joukowski choices
corresponding to $z/(1+z)$ have reciprocal moduli, with the larger equal to
\[
 \frac{\max\{|\sqrt z+i|^2,|\sqrt z-i|^2\}}{|1+z|}
 =\frac{1+u}{\sqrt{1-u^2}}.
\]
Therefore the polynomial $F(z)=(1+z)^nP(z/(1+z))$ satisfies
\begin{align*}
 \frac{|F(re^{i\phi})|}{M(1+r)^n}
 &\le(1-u^2)^{(n-m)/2}(1+u)^m\\
 &\le e^{-nu^2/4+mu}\le e^{2m^2/n}e^{-nu^2/8}.
\end{align*}
Values where a displayed rational expression is undefined follow by
continuity of $F$. Since
$\sin(|\phi|/2)\ge|\phi|/\pi$ for $|\phi|\le\pi$, the average of the
last Gaussian factor over $[-\pi,\pi]$ is at most
$C/\sqrt{1+nq_*}$. Cauchy's coefficient formula and
$F(z)=\sum_j\binom njb_jz^j$ now yield
\[
 |b_j|\le Me^{2m^2/n}\frac{(1+r)^n}{\binom njr^j}
                      \frac C{\sqrt{1+nq_*}}.
\]
Stirling's upper and lower factorial bounds give
\[
 \binom nj(j/n)^j(1-j/n)^{n-j}\ge\frac c{\sqrt{1+nq_*}},
\]
which cancels the displayed square-root factor. At the two endpoints,
$b_0=P(0)$ and $b_n=P(1)$; the constant case is also immediate.
\end{proof}

\begin{lemma}[A derivative-stable algebraic Jackson approximation]
\label{p04:nxm:lem:stableJackson}
For every $1$-Lipschitz function $f:[0,1]\to\mathbb R$ and integer
$k\ge2$, there is a real polynomial $P_k$ of degree at most $k$ such that
\begin{equation}\label{p04:nxm:eq:stableJackson}
 \|P_k-f\|_\infty\le C/k,\qquad \|P_k'\|_\infty\le2.
\end{equation}
\end{lemma}
\begin{proof}
Subtract $f(0)$, so that $|f|\le1$. Reflect $f$ from $[0,1]$ to $[1,2]$
by $F(x)=f(2-x)$ and extend with period two. This is a globally
$1$-Lipschitz function with $\|F\|_\infty\le1$.
Convolve it with the nonnegative normalized, period-two Jackson kernel
\[
 J_s(u)=c_s\left(\frac{\sin(s\pi u/2)}{\sin(\pi u/2)}\right)^4,
 \qquad -1\le u\le1.
\]
The quotient is interpreted continuously at zero. The resulting real
trigonometric polynomial $Q_s$ has frequencies $|l|\le2(s-1)$, satisfies
$\|Q_s'\|_\infty\le1$ by positivity and translation invariance of
convolution, and has error $\|Q_s-F\|_\infty\le C/s$.
Here the kernel has first absolute moment $O(1/s)$: its unnormalized
integral is comparable to $s^3$, while its unnormalized first absolute
moment is at most $Cs^2$, using
$|\sin(s\pi u/2)/\sin(\pi u/2)|\le C\min(s,|u|^{-1})$.

Write $Q_s(x)=\sum_{|l|\le2s}a_l e^{i\pi l(x-1/2)}$. Each $|a_l|\le1$,
because it is a Fourier coefficient of $F$ multiplied by a Fourier
coefficient of a probability kernel. Replace each exponential by its
Taylor polynomial at $x=1/2$ of degree $M=100s$, and call the resulting
algebraic polynomial $P$. Conjugate symmetry makes $P$ real. On $[0,1]$,
$|\pi l(x-1/2)|\le\pi s$, so Taylor's remainder gives
\begin{align*}
 \|P-Q_s\|_\infty
 &\le(4s+1)e^{\pi s}\frac{(\pi s)^{M+1}}{(M+1)!},\\
 \|P'-Q_s'\|_\infty
 &\le(4s+1)2\pi s e^{\pi s}\frac{(\pi s)^M}{M!}.
\end{align*}
The bound $M!\ge(M/e)^M$ shows that both expressions are at most
$e^{-100s}$ for $s\ge1$. Thus $\|P'\|_\infty\le2$ and
$\|P-f\|_\infty\le C/s$.
For $k\ge100$, take $s=\lfloor k/100\rfloor$, which has $100s\le k$
and $s\ge k/200$. For $2\le k<100$, use the constant polynomial $f(0)$
and enlarge the absolute approximation constant. Finally add back $f(0)$.
\end{proof}

\begin{corollary}[Lipschitz Bernstein inversion]
\label{p04:nxm:cor:inverseLip}
For a $1$-Lipschitz $f$ and $2\le k\le t/2$, its approximating polynomial
can be written $P_k=\sum_{j=0}^t b_jB_j^t$ with
\begin{equation}\label{p04:nxm:eq:discreteLip}
 \max_{0\le j<t} t|b_{j+1}-b_j|\le C e^{4k^2/t}.
\end{equation}
Consequently, for any mixing laws $G,H$,
\begin{equation}\label{p04:nxm:eq:transportinverse}
 W_1(G,H)\le C/k+C e^{4k^2/t} W_{1,t}(p_G,p_H).
\end{equation}
\end{corollary}
\begin{proof}
Differentiate the Bernstein expansion:
\[
 P_k'(x)=t\sum_{j=0}^{t-1}(b_{j+1}-b_j)B_j^{t-1}(x).
\]
Its degree is at most $k-1\le(t-1)/2$ and its supremum norm is at most two.
Lemma~\ref{p04:nxm:lem:coefficient} applies, with
$2(k-1)^2/(t-1)\le4k^2/t$, proving \eqref{p04:nxm:eq:discreteLip}.
The coefficient vector $j\mapsto b_j$ is therefore Lipschitz on the
normalized count grid with the displayed constant. The approximation
error contributes at most $2\|P_k-f\|_\infty$, and the remaining integral
is $\sum_jb_j(p_G(j)-p_H(j))$. Apply Wasserstein duality on both spaces.
\end{proof}

\subsection{Proof of the rate statements}

\begin{proof}[Proof of Theorem~\ref{p04:nxm:thm:main}]
Insert $H=\widehat G$ into \eqref{p04:nxm:eq:transportinverse} and use
\[
 W_{1,t}(p_G,\widehat p)\le W_{1,t}(p_G,q)+W_{1,t}(q,\widehat p).
\]
Lemma~\ref{p04:nxm:lem:predictive} proves \eqref{p04:nxm:eq:masterexpectation}
with $A=4$. It also proves, with probability at least $1-\eta$,
\begin{equation}\label{p04:nxm:eq:masterprobability}
 W_1(\widehat G,G)\le\frac Ck+
 C e^{4k^2/t}\sqrt{\frac{\log(N+1)+\log(2/\eta)}N},
\end{equation}
simultaneously for all MLEs.

For the near-linear range, set $\eta=N^{-2}$ and
$k=\lfloor\sqrt{c\varepsilon t\log N}\rfloor$, where $c>0$ is sufficiently
small, depending only on the stated lower-limit constant $a$ when necessary.
For large $N$, this is in $[2,t/2]$. The second term of
\eqref{p04:nxm:eq:masterprobability}, multiplied by $\sqrt{t\log N}$, is at most
\[
 C N^{4c\varepsilon-1/2}\sqrt t\log N
 \le C N^{-(1/2-4c)\varepsilon}\log N.
\]
Choose $c<1/8$. This expression is bounded for large $N$, proving the
probability claim. The expectation claim follows directly from
\eqref{p04:nxm:eq:masterexpectation}, or from integration since the loss is
at most one and the exceptional probability is $N^{-2}$.

For completeness, the uniform interpolation bound can be obtained
without asymptotic restrictions. Put $x=N/t$ and $L=\log(1+x)$, and fix
$0<c<1/32$. For all parameter choices for which
$z=\min\{t/2,\sqrt{ctL}\}\ge4$, choose $k=\lfloor z\rfloor$.
Then $k\ge z/2$, so
\[
 1/k\le C\{1/t+1/\sqrt{tL}\}.
\]
If $z=\sqrt{ctL}$, the exponential term in
\eqref{p04:nxm:eq:masterexpectation} satisfies
\[
 \frac{e^{4k^2/t}}{\sqrt N}\sqrt{tL}
 \le (1+x)^{4c}\sqrt{\frac{\log(1+x)}x}\le C.
\]
The last expression is bounded on $(0,\infty)$ because $4c<1/2$.
If $z=t/2$, then $L\ge t/(4c)$ and
\[
 \frac{e^{4k^2/t}}{\sqrt N}
 \le\frac{e^t}{\sqrt{t(e^{t/(4c)}-1)}}\le C/t.
\]
If $z<4$, at least one of $1/t$ and $1/\sqrt{tL}$ is bounded below by an
absolute positive constant, and the trivial loss bound one suffices.
This proves \eqref{p04:nxm:eq:allrange} for every $N,t\ge1$.
\end{proof}

\subsection{Why a uniform bounded-fingerprint shortcut is false}
\label{p04:nxm:sec:obstruction}

Section~6.2 of \cite{p01:info:MMRS2025} discusses a dimension-free bound for the
fingerprint term tested against Bernstein coefficients. Our result proves
the needed uniform control in a Lipschitz test norm after constructing a
suitable approximation. It must not be interpreted as a uniform
$N^{-1/2}$ total-variation bound over all bounded coefficient vectors.
The latter strengthening is false, even if the fitted mixture is chosen by
an oracle rather than maximum likelihood.

\begin{proposition}[An obstruction for all bounded fingerprint tests]
\label{p04:nxm:prop:TVobstruction}
Let the true mixing law be uniform on $[0,1]$, so that $p_G(j)=1/(t+1)$.
There are absolute constants $c,C>0$ such that, whenever $N\ge t+1$ and
$t\ge8$,
\begin{equation}\label{p04:nxm:eq:TVlower}
 \mathbb E_G\inf_H\|q-p_H\|_1
 \ge c\sqrt{t/N}-C/\sqrt t.
\end{equation}
In particular, along $t=\lfloor N^{3/4}\rfloor$ the left side is at least
$cN^{-1/8}$ for all sufficiently large $N$. Consequently there is no
uniform $O(N^{-1/2})$ expectation bound for
$\sup_{\|b\|_\infty\le1}|\sum_jb_j(q_j-p_H(j))|$, for any fitting rule
whose output is a binomial mixture.
\end{proposition}
\begin{proof}
Set $L=\lceil t/4\rceil$, $U=\lfloor3t/4\rfloor$, and
$V(v)=\sum_{j=L}^{U-1}|v_{j+1}-v_j|$.
Each binomial probability vector is unimodal. Its variation on this
central interval is at most twice its largest value there, which is at
most $C/\sqrt t$. Indeed, for each central $j$, $B_j^t(\theta)$ is largest
at $\theta=j/t$, and Stirling's factorial bounds give
$B_j^t(j/t)\le C/\sqrt t$. Integrating over $H$ therefore gives
\begin{equation}\label{p04:nxm:eq:centralvariation}
 V(p_H)\le C/\sqrt t
\end{equation}
for every mixture.

For adjacent bins under the uniform count law, let
$Z=N(q_{j+1}-q_j)=\sum_{i=1}^N Y_i$, where the independent variables
$Y_i$ take values $1,-1$ each with probability $1/(t+1)$ and zero otherwise.
They are centered, and $v=\mathbb EZ^2=2N/(t+1)\ge2$.
Their fourth moment satisfies
\[
 \mathbb EZ^4=v+3N(N-1)\left(\frac2{t+1}\right)^2
 \le v+3v^2\le\frac72v^2.
\]
For a nonnegative random variable $Y$, Cauchy--Schwarz gives
$\mathbb P\{Y\ge\tfrac12\mathbb EY\}\ge
(\mathbb EY)^2/(4\mathbb EY^2)$: split its first moment over that event
and its complement. Apply this with $Y=Z^2$ to obtain
$\mathbb E|Z|\ge c\sqrt v$. There are at least $ct$ adjacent central
pairs, and hence $\mathbb EV(q)\ge c\sqrt{t/N}$.
Finally, the triangle inequality yields
\[
 V(q)\le V(p_H)+2\|q-p_H\|_1.
\]
Use \eqref{p04:nxm:eq:centralvariation}, take the infimum and expectation,
and absorb absolute constants. The final assertion follows from the dual
identity between the $\ell_1$ norm and bounded coefficient vectors.
\end{proof}

This proposition concerns a supremum over all bounded tests, including
tests selected after seeing the data. It does not disprove a conjecture
with one fixed deterministic coefficient vector, or a conjecture restricted
to a carefully chosen low-degree approximation class. The distinction is
essential: the transportation norm used in our proof exploits cancellation
between nearby count bins, whereas unrestricted bounded tests can select
the signs of the local sampling fluctuations.

% ---- BEGIN budget ----
\section{Exact replication caps and the critical endpoint experiment}
\label{p04:nxb:sec:budget}

This section advances two total-budget questions that were left open in
our preceding observation papers.  For Bernoulli populations we determine
the complete cap-dependent minimax curve for predetermined, possibly
unequal, replication counts, and extend the exact adaptive curve to every
polynomially subcritical cap.  For scalar uniform-endpoint populations we
prove that the logarithm at the critical component-estimation exponent is
necessary even under fully adaptive sampling.  These are classification
problems for the experiments specified below; we do not describe them as
previously named external conjectures.

Let an individual draw an independent latent parameter from an unknown
probability law $G$. A raw observation costs one unit, and repeated
observations of that individual are conditionally independent. A policy
may start individuals, revisit them, and randomize using all previous
observations. Its total cost is at most a deterministic integer $B$.
A cap $Q$ means that no individual is observed more than $Q$ times.
The superscripts $\mathrm{ad}$ and $\mathrm{fix}$ respectively allow all
such policies and restrict the replication counts to a possibly randomized
schedule chosen independently of the data. Estimators are arbitrary
probability-valued measurable procedures; the loss is expected $W_1$.

\subsection{The Bernoulli cap curve}

For $P_\theta=\operatorname{Bernoulli}(\theta)$, $0\leq\theta\leq1$,
write $R^{\mathrm{ad}}_{B,Q}$ and $R^{\mathrm{fix}}_{B,Q}$ for these risks.
Set
\begin{equation}\label{p04:nxb:eq:cap-function}
 \Psi(B,Q)=\max\left\{ B^{-1/3},\frac1Q,
 \frac1{\sqrt{Q\log(e+B/Q^{3/2})}}\right\}.
\end{equation}

\begin{theorem}[The complete predetermined cap curve]
\label{p04:nxb:thm:fixed-caps}
There are absolute constants $0<c<C<\infty$ such that, for all integers
$B\geq2$ and $1\leq Q\leq B$,
\[
 c\Psi(B,Q)\leq R^{\mathrm{fix}}_{B,Q}\leq C\Psi(B,Q).
\]
The upper bound uses a predetermined collection of binomial observations
at geometrically varying replication depths, followed by a weighted
Chebyshev-moment projection. It does not assert a property of the MLE.
In particular the optimal uncapped order $B^{-1/3}$ is attained within a
fixed constant precisely when $Q\gtrsim B^{2/3}$, with constants allowed
to depend on that fixed risk constant.
\end{theorem}

\begin{theorem}[Adaptive caps below the critical power]
\label{p04:nxb:thm:adaptive-caps}
For every fixed $0<\gamma<2/3$, uniformly for
$1\leq Q\leq B^{2/3-\gamma}$ as $B\to\infty$,
\[
 R^{\mathrm{ad}}_{B,Q}\asymp_\gamma
 \max\left\{\frac1Q,\frac1{\sqrt{Q\log(eB)}}\right\}.
\]
For arbitrary $1\leq Q\leq B$ the following valid sandwich remains:
\[
 c\max\left\{B^{-1/3},Q^{-1},[Q\log(eB)]^{-1/2}\right\}
 \leq R^{\mathrm{ad}}_{B,Q}\leq C\Psi(B,Q).
\]
In particular the fully adaptive cap curve in a shrinking neighborhood
of $B^{2/3}$ is not claimed to be settled.
\end{theorem}

The preceding cap theorem only reached $Q\leq B^{2/5-\gamma}$, by
substituting an equal-group MLE bound. The new procedure estimates
different frequency bands at different costs. The weighted moment
inequality below is an established Chebyshev-moment principle, explicitly
developed by Musco, Musco, Rosenblatt and Singh
\cite{p01:info:MMRS2025}. Trial-profile-dependent spectral estimation also
appears in Lee and Kennedy \cite{p01:info:LeeKennedy2026}; their target is pointwise
estimation of a H\"older mixing density, whereas ours is arbitrary-law
$W_1$ recovery under an optimized observation budget and a cap.
Neither statement is presented as a special case of the other.

\begin{lemma}[Moment recovery and binomial statistics]
\label{p04:nxb:lem:moments}
Write $\widetilde T_m(x)=T_m(2x-1)$ and
$\mu_m(G)=\int\widetilde T_m\,dG$. For $K\geq2$,
\begin{equation}\label{p04:nxb:eq:weighted-modulus}
 W_1(G,H)\leq\frac CK+
 C\left\{\sum_{m=1}^K\frac{|\mu_m(G)-\mu_m(H)|^2}{m^2}\right\}^{1/2}.
\end{equation}
For $m\leq r/2$ there is a statistic $A_{r,m}$ of
$X\sim\operatorname{Binomial}(r,\theta)$ satisfying
\[
 \mathbb E_\theta A_{r,m}(X)=\widetilde T_m(\theta),\qquad
 \max_{0\leq j\leq r}|A_{r,m}(j)|\leq C\exp(2m^2/r).
\]
\end{lemma}
\begin{proof}
For a $1$-Lipschitz $f$ the periodic function
$g(u)=f((1+\cos u)/2)$ has bounded weak derivative. Its cosine
coefficients satisfy $\sum_{m\geq1}m^2|c_m|^2\leq C$ by Parseval.
Convolution with the normalized nonnegative Jackson kernel
$[\sin(su/2)/\sin(u/2)]^4$, $s=\lfloor K/2\rfloor+1$, has degree
at most $K$, changes $g$ uniformly by at most $C/K$, and multiplies
its coefficients by factors of absolute value at most one. The
first-moment estimate for this kernel follows by integrating
$\min(s^4,|u|^{-4})$ and $|u|\min(s^4,|u|^{-4})$.
Cauchy--Schwarz applied to the resulting polynomial proves
\eqref{p04:nxb:eq:weighted-modulus} by Kantorovich duality.

Take $A_{r,m}(j)$ to be the degree-$r$ Bernstein coefficients of
$\widetilde T_m$. We include the coefficient estimate to make the
cost calculation self-contained. For $0<j<r$ put
$z=te^{iu}$, $t=j/(r-j)$, $q=t/(1+t)^2$. The generating polynomial
\[
 (1+z)^r\widetilde T_m\!\left(\frac z{1+z}\right)
 =\frac{(1+z)^{r-m}}2
       \{(\sqrt z+i)^{2m}+(\sqrt z-i)^{2m}\}
\]
has coefficient $\binom rj A_{r,m}(j)$ at $z^j$.
With $v=2\sqrt q|\sin(u/2)|$, its modulus divided by $(1+t)^r$
is at most
\[
 (1-v^2)^{(r-m)/2}(1+v)^m
 \leq e^{2m^2/r}e^{-rv^2/8}.
\]
The Cauchy integral supplies a factor $C(1+rq)^{-1/2}$.
The binomial probability $\binom rj t^j/(1+t)^r$ is at least
$c/\sqrt{1+rq}$ by Stirling's inequalities. Its reciprocal cancels
the preceding factor.
The endpoint coefficients have absolute value one. The Bernstein
identity proves unbiasedness.
\end{proof}

\begin{lemma}[A spectral allocation bound]
\label{p04:nxb:lem:allocation}
There are absolute constants such that, whenever $K\geq2$ is a power
of two, $K\leq Q/2$, and $K^2\leq cB$, a predetermined capped
procedure satisfies
\begin{equation}\label{p04:nxb:eq:spectral-upper}
 \sup_G\mathbb E_G W_1(\widehat G,G)
 \leq C\left\{K^{-1}+\sqrt{K/B}\,e^{2K^2/Q}\right\}.
\end{equation}
\end{lemma}
\begin{proof}
Partition $1\leq m\leq K$ into dyadic bands with upper endpoints
$k=1,2,4,\ldots,K$. At band $k$ use
$r_k=\min\{Q,\max(2,k^2)\}$ observations per fresh individual.
Then $r_k\geq2m$ throughout that band. Estimate every moment in
it by the average of $A_{r_k,m}$ over its $n_k$ individuals.
The expected squared weighted moment error in the band is at most
$a_k/n_k$, where
\[
 a_k=Ck^{-1}e^{4k^2/r_k}.
\]
Let $S=\sum_k\sqrt{r_k a_k}$ and take
$n_k=\lceil B\sqrt{a_k/r_k}/(2S)\rceil$.
The total cost is at most $B/2+\sum_k r_k\leq B$ after reducing the
absolute constant in $K^2\leq cB$. Moreover
\[
 \sum_k\frac{a_k}{n_k}\leq\frac{2S^2}{B},\qquad
 S\leq C\sum_k\sqrt{k}\,e^{2+2K^2/Q}
 \leq C\sqrt K e^{2K^2/Q}.
\]
Project the estimated moment vector onto the compact moment set of
probability laws on $[0,1]$, in the norm
$(\sum_{m=1}^K|z_m|^2/m^2)^{1/2}$. Its distance to the true moment
vector is at most twice the estimation error. A measurable approximate
projection over a countable dense set of finitely supported rational
laws suffices and adds an arbitrarily small error. Apply
\eqref{p04:nxb:eq:weighted-modulus} and Jensen's inequality.
\end{proof}

\begin{proof}[Upper bound in Theorem~\ref{p04:nxb:thm:fixed-caps}]
Put $L=\log(e+B/Q^{3/2})$ and
$H=\min\{B^{1/3},Q,\sqrt{QL}\}$. Choose the largest power of two
$K\leq c_0H$, where $c_0>0$ is a sufficiently small absolute constant.
If there is no $K\geq2$, the diameter bound gives the assertion by
increasing $C$. Otherwise Lemma~\ref{p04:nxb:lem:allocation} applies,
with finitely many small budgets absorbed into the constant.
If $Q\geq B^{2/3}$, then $K^2/Q\leq c_0^2$ and $K^3\leq c_0^3B$.
If $Q<B^{2/3}$, write $x=B/Q^{3/2}>1$. Then
\[
 \frac{K^{3/2}}{\sqrt B}e^{2K^2/Q}
 \leq c_0^{3/2}\frac{\log(e+x)^{3/4}}{\sqrt x}
                       (e+x)^{2c_0^2}\leq C,
\]
provided, for example, $2c_0^2<1/4$.
Thus both terms in \eqref{p04:nxb:eq:spectral-upper} are at most $C/H$,
and $H^{-1}=\Psi(B,Q)$.
\end{proof}

The projection can be implemented by a finite convex quadratic program.
Indeed $|\widetilde T_m'(x)|\leq2m^2$ on $[0,1]$, so rounding a
probability law to a grid of mesh $K^{-3}$ changes its weighted moment
vector by at most $CK^{-3}(\sum_{m\leq K}m^2)^{1/2}\leq C/K$.
Minimizing the squared weighted discrepancy over the simplex on this
grid, to additive objective error $O(K^{-2})$, therefore preserves the
rate. The grid has $O(K^3)$ points; its entries and the Bernstein
statistics are rational and have polynomial encoding length. Thus the
construction admits a polynomial-time implementation in the usual
rational-input model; no efficient evaluation of a continuum argmin is
being assumed.

\begin{lemma}[Localized moment matching with an information charge]
\label{p04:nxb:lem:localized}
There are absolute constants $c_0,C_0>0$ and $0<\rho<1$ with the
following property. Let $K\geq1$, $Q\geq K$, and
$h^2=c_0K/Q$, with $h\leq1/32$. For any center
$u\in[1/4,3/4]$ there are laws $\nu_+,\nu_-$ supported on
$[u-h,u+h]$ such that
\[
 W_1(\nu_+,\nu_-)=\frac{h}{K+1},\qquad
 \operatorname{KL}(P_{r,+}\Vert P_{r,-})
 \leq C_0 r h^2\rho^K\quad(1\leq r\leq Q),
\]
where $P_{r,\pm}$ is the full $r$-toss mixture law.
\end{lemma}
\begin{proof}
Split the signed measure
$2^{-K}\sum_{j=0}^{K+1}(-1)^j\binom{K+1}{j}\delta_{j/(K+1)}$
into two probability laws, and transport them to $[u-h,u+h]$.
They match moments through order $K$ and have $W_1$ distance
$2h/(K+1)$: integrate their cumulative signed mass, using
$\sum_{i=0}^j(-1)^i\binom{K+1}{i}=(-1)^j\binom Kj$.
Mix each law with weight $1/2$ with the common atom $\delta_u$.
This gives the displayed separation and bounds each observation-law
likelihood ratio below by $1/2$ relative to $r$ independent
$\operatorname{Bernoulli}(u)$ observations.

Set $v=u(1-u)\geq3/16$ and $Z_i=(X_i-u)/\sqrt v$ under that
reference law. Its product likelihood for bias $u+a$ is
$\prod_i(1+aZ_i/\sqrt v)$. Orthogonality of the products
$\prod_{i\in S}Z_i$ and moment matching imply
\[
 \operatorname{KL}(P_{r,+}\Vert P_{r,-})
 \leq\chi^2(P_{r,+},P_{r,-})
 \leq C\sum_{j>K}\binom rj(h^2/v)^j
 \leq Crh^2\left(\frac{C Qh^2}{K}\right)^K.
\]
The last inequality follows from $\binom rj\leq(er/j)^j$;
when $r\leq K$ the sum is zero. Choose $c_0$ small enough.
\end{proof}

\begin{proof}[Lower bound in Theorem~\ref{p04:nxb:thm:fixed-caps}]
Moment-matching laws through degree $Q$ have identical full $Q$-toss
laws and separation $1/(Q+1)$, proving the $Q^{-1}$ term even for
adaptive policies.

We recall a short fully adaptive argument for the $B^{-1/3}$ term.
Place $M\asymp h^{-1}$ separated cells in $[1/4,3/4]$, each with
two possible point biases separated by $h$, and give every cell mass
$1/M$. Strengthen the experiment by allowing the statistician to
select a known cell for each toss. This oracle simulates an original
policy by drawing a uniform cell label whenever an individual starts.
Two parameter vectors differing at cell $j$ have transcript relative
entropy at most $Ch^2\mathbb E N_j$, where $N_j$ is its number of
oracle queries. Sum this inequality over hypercube edges and use
$\sum_jN_j\leq B$; Pinsker and Cauchy--Schwarz bound the mean edge
variation by $C\sqrt{Bh^2/M}$. For $h=cB^{-1/3}$ this is below
$1/2$. The separated-cell $W_1$ identity and nearest-hypercube decoding
then give risk at least $ch$. This proof allows arbitrary adaptation.

For the remaining term put $L=\log(e+B/Q^{3/2})$ and choose
$K=\lceil A L\rceil$ with a sufficiently large fixed $A$.
If $Q<C_1L$, the desired term is already at most
$\sqrt{C_1}/Q$. Otherwise take $h^2=c_0K/Q$ as in
Lemma~\ref{p04:nxb:lem:localized}, increasing $C_1$ so $h\leq1/32$,
and pack $M\asymp h^{-1}$ disjoint intervals of radius $h$ inside
$[1/4,3/4]$. At each interval choose either law from that lemma and
assign it mass $1/M$. Reveal the interval label of every individual
for free. For any predetermined group sizes $r_i\leq Q$ with
$\sum_i r_i\leq B$, neighboring hypercube laws in this stronger
experiment have relative entropy at most
\[
 \frac1M\sum_i C_0r_i h^2\rho^K
 \leq C B h^3\rho^K
 \leq C\frac B{Q^{3/2}}K^{3/2}\rho^K\leq\frac18.
\]
The last inequality holds by increasing $A$, since $K\asymp L$.
It is uniform over randomized schedules after including their
independent seed in the transcript.
For two hypercube vertices the cumulative mass difference returns to
zero at each cell boundary, and hence their $W_1$ distance is exactly
$h/[M(K+1)]$ times their Hamming distance. Pinsker, binary testing,
and nearest-hypercube decoding give minimax risk at least
$ch/(K+1)\asymp[QL]^{-1/2}$. This completes the proof.
\end{proof}

\begin{proof}[Theorem~\ref{p04:nxb:thm:adaptive-caps}]
Besides the preceding adaptive lower bounds, take
$K=\lceil4\log(eB)\rceil$ and $h^2=K/(64eQ)$ when $Q\geq K$.
Two moment-matching laws through degree $K$ supported on
$[1/2-h,1/2+h]$ have $W_1$ separation $2h/(K+1)$.
The Walsh expansion used above bounds the squared $L^2$ difference
of their full $Q$-toss likelihood ratios, relative to the fair coin,
by $C16^{-K}$. Their total variation is at most $C4^{-K}$.
An oracle of $B$ independent full $Q$-toss vectors simulates every
adaptive capped policy by revealing requested prefixes. Its total
variation is at most $CB4^{-K}<1/4$, proving the
$[Q\log(eB)]^{-1/2}$ lower bound. If $Q<K$, the moment floor
already gives it up to a constant. For $Q\leq B^{2/3-\gamma}$,
\[
 \log(e+B/Q^{3/2})\asymp_\gamma\log(eB),
\]
and $B^{-1/3}$ is bounded by the other rate terms for all sufficiently
large $B$. The predetermined upper bound therefore matches.
\end{proof}

The narrow remaining adaptive cap region is consequential. At
$Q=B^{2/3}/\log B$, predetermined designs have risk of order
$B^{-1/3}\sqrt{\log B/\log\log B}$, whereas the present adaptive
lower bound is only $cB^{-1/3}$. We neither claim equality of these
risks nor infer an adaptive advantage from the gap.

\subsection{The scalar uniform-endpoint critical rate}

Now let $P_\theta=\operatorname{Uniform}[0,\theta]$ with
$1\leq\theta\leq2$, without a replication cap, and write $U_B$ for
the fully adaptive minimax risk. The previously known multilevel upper
bound and one-scale lower bound left a logarithmic gap.

\begin{theorem}[The critical logarithm is unavoidable]
\label{p04:nxb:thm:uniform}
For absolute constants $0<c<C<\infty$ and all sufficiently large $B$,
\[
 c\frac{\log B}{\sqrt B}\leq U_B
 \leq C\frac{\log B}{\sqrt B}.
\]
The lower bound allows arbitrary adaptive randomized allocation of the
$B$ observations. It already holds on a class for which the density of
$\log\theta/\log2$ is bounded above and below by positive constants
and has a common Lipschitz constant. A predetermined multilevel design
attains the upper bound.
\end{theorem}

The main new step is an information bound for the entire adaptive
experiment, not merely an integrated-variance calculation for one
multilevel architecture. Bayesian information inequalities and
multilevel distribution recovery are established methods
\cite{p04:nxb:Clarkson2020,p01:info:GNR2015}. We supply the required absolute-loss
inequality and the adaptive information calculation in full.

\begin{lemma}[An absolute-loss Bayesian information inequality]
\label{p04:nxb:lem:absolute-fisher}
Let $(\beta,Y)$ have a joint density $p(\beta,y)$ on
$\mathbb R^D\times\mathcal Y$ with integrable weak derivatives in
$\beta$ and positive definite finite Fisher matrix
\[
 J=\int (\nabla_\beta\log p)(\nabla_\beta\log p)^T p.
\]
For every $v\in\mathbb R^D$ and estimator $T(Y)$,
\[
 \mathbb E|T(Y)-v^T\beta|\geq\frac14\sqrt{v^TJ^{-1}v}.
\]
\end{lemma}
\begin{proof}
For $p_h(\beta,y)=p(\beta+h,y)$, integration along translations and
Cauchy--Schwarz give
\[\operatorname{TV}(p,p_h)\leq\tfrac12\sqrt{h^TJh}.\]
The triangle inequality, integrated against $\min(p,p_h)$, gives
\[
 2\mathbb E|T(Y)-v^T\beta|
 \geq |v^Th|\{1-\operatorname{TV}(p,p_h)\}.
\]
Take $h=J^{-1}v/\sqrt{v^TJ^{-1}v}$; the case $v=0$ is immediate.
\end{proof}

\begin{lemma}[A Fourier trace inequality]
\label{p04:nxb:lem:rotation}
For $k=1,\ldots,K$ let
\[
 v_{k,c}(x)=\frac{\sqrt2\sin(2\pi kx)}{2\pi k},\qquad
 v_{k,s}(x)=-\frac{\sqrt2\cos(2\pi kx)}{2\pi k},
\]
and concatenate them into $v(x)\in\mathbb R^{2K}$. For every positive
definite $J$,
\[
 \int_0^1\sqrt{v(x)^TJ^{-1}v(x)}\,dx
 \geq\frac c{\sqrt{\operatorname{tr}J}}
                 \sum_{k=1}^K\frac1k.
\]
\end{lemma}
\begin{proof}
The map $J\mapsto\sqrt{v^TJ^{-1}v}$ is convex, since it is the
supremum over $z\ne0$ of
$|v^Tz|/(z^TJz)^{1/2}$, each of which is convex in $J$.
The integral is invariant under conjugation by the block rotations
corresponding to $x\mapsto x+t$. Average these conjugates over
$0\leq t\leq1$ and apply Jensen. The average matrix has blocks
$j_k I_2$, where $2\sum_kj_k=\operatorname{tr}J$.
Its integrand is the constant
$(\sum_k1/(2\pi^2k^2j_k))^{1/2}$.
Cauchy--Schwarz gives
$\sum_k1/(k^2j_k)\geq(\sum_k1/k)^2/\sum_kj_k$.
\end{proof}

\begin{lemma}[A constant information charge for an adaptive endpoint query]
\label{p04:nxb:lem:adaptive-trace}
Let $a=\log2$, $U=\log\theta/a\in[0,1]$, and let its density be
\[
 w_\beta(u)=1+\sum_{k=1}^K\{\beta_{k,c}\sqrt2\cos(2\pi ku)
                         +\beta_{k,s}\sqrt2\sin(2\pi ku)\}.
\]
Suppose $1/2\leq w_\beta\leq3/2$ and
$\operatorname{Lip}(w_\beta)\leq L_0$ throughout the parameter set.
For every adaptive policy of at most $B$ raw observations, the trace
of its Fisher information matrix in $\beta$ is at most
$C(1+L_0^2)B$, uniformly in $K$ and $\beta$.
\end{lemma}
\begin{proof}
After $r$ observations of an individual, its posterior under the
baseline density $w_0=1$ is a truncated exponential
\[
 q(u)=\frac{\lambda e^{-\lambda(u-b)}}{1-e^{-\lambda d}}
                  \mathbf1_{[b,1]}(u),\qquad
 \lambda=ar,\quad d=1-b,
\]
where $b$ is the maximum of zero and the scaled logarithms of that
individual's observations. The uniform case is the limit $\lambda=0$.
This remains its baseline posterior after conditioning on the entire
adaptive history: the action probabilities contribute no parameter
factor once the observed history is fixed.
A fresh individual has $q=\mathbf1_{[0,1]}$.

A new scaled log-observation is $Y=U-E$ with $E$ exponential of rate
$a$. Let $q_Y$ be the updated baseline posterior. Direct integration
shows the exact conditional identity
\begin{equation}\label{p04:nxb:eq:posterior-charge}
 \mathbb E_0\|q_Y-q\|_2^2=\frac a2.
\end{equation}
Here is the calculation, after translating $b$ to zero. Put
$Z=1-e^{-\lambda d}$. The current squared norm is
$\lambda(1+e^{-\lambda d})/(2Z)$. The event $Y\leq0$ has probability
$\lambda(1-e^{-(\lambda+a)d})/[(\lambda+a)Z]$, and its posterior
norm is $(\lambda+a)\coth((\lambda+a)d/2)/2$.
For $0<y<d$ the predictive density is
\[
 \frac{a\lambda}{(\lambda+a)Z}e^{-\lambda y}
                       (1-e^{-(\lambda+a)(d-y)}).
\]
Multiply it by $(\lambda+a)\coth((\lambda+a)(d-y)/2)/2$
and integrate. The sum with the $Y\leq0$ contribution is
$a/2+\|q\|_2^2$. The martingale identity
$\mathbb E_0 q_Y=q$ proves \eqref{p04:nxb:eq:posterior-charge}.
The limit $\lambda=0$ follows by continuity.

Uniformly over these posteriors,
\begin{equation}\label{p04:nxb:eq:norm-var}
 \|q\|_2^2\operatorname{Var}_q(U)\leq C.
\end{equation}
If $\lambda d\leq1$, the density is at most $C/d$ and the variance
at most $d^2$. If $\lambda d>1$, the squared norm is at most
$C\lambda$ and the variance at most $C/\lambda^2$; since
$\lambda>1/d\geq1$, their product is again bounded.

Set $z=\int w_\beta q$ and $z_Y=\int w_\beta q_Y$; both lie in
$[1/2,3/2]$. The predictive density under $\beta$, relative to its
baseline counterpart, is $z_Y/z$. The score increment for this
observation is the vector of Fourier coefficients of
$q_Y/z_Y-q/z$, in the $2K$ functions occurring in $w_\beta$.
Bessel's inequality therefore bounds its squared length by
\[
 \|q_Y/z_Y-q/z\|_2^2
 \leq C\|q_Y-q\|_2^2+C\|q\|_2^2|z_Y-z|^2.
\]
Kantorovich duality and a product coupling give
\[
 \mathbb E_0|z_Y-z|^2
 \leq L_0^2\mathbb E_0W_1(q_Y,q)^2
 \leq2L_0^2\operatorname{Var}_q(U).
\]
For the last step, couple conditionally independent draws from
$q_Y$ and $q$, and then average over $Y$; the first draw has marginal
$q$ by the posterior martingale property. Equations
\eqref{p04:nxb:eq:posterior-charge}--\eqref{p04:nxb:eq:norm-var}, and the
bounded predictive likelihood ratio, prove the conditional trace
bound $C(1+L_0^2)$ under the actual parameter $\beta$.
Score increments are martingale differences. Conditional independence
of action randomization from the parameter, and the deterministic
budget bound, allow their squared lengths to be summed through all
$B$ adaptive queries. This proves the assertion.
\end{proof}

\begin{proof}[Lower bound in Theorem~\ref{p04:nxb:thm:uniform}]
Use the $2K$-parameter family in Lemma~\ref{p04:nxb:lem:adaptive-trace}.
Independently put on every coefficient the density
\[
 \pi_b(t)=b^{-1}\cos^2(\pi t/(2b))\mathbf1_{[-b,b]}(t),
 \qquad b=c_1K^{-2}.
\]
Choose the absolute $c_1$ small enough. Then every density in the prior
support lies between $3/4$ and $5/4$, and
$\operatorname{Lip}(w_\beta)\leq Cb\sum_{k\leq K}k\leq C$.
Each prior coordinate has Fisher information $\pi^2/b^2$, so the
prior information trace is at most $CK^5$.
The Bayesian joint information matrix, with the policy's independent
random seed included in its data, consequently satisfies
\[
 \operatorname{tr}J\leq C(B+K^5).
\]
The likelihood and prior score cross term vanishes. Compact support
causes no boundary term because $\pi_b$ and its first derivative
vanish at the endpoints, and its Fisher information is finite.

For a law on $[0,1]$, write $F^\circ(x)=F(x)-\int_0^1F(t)dt$.
For two laws,
\[
 \|\widehat F^\circ-F^\circ\|_1
 \leq2\|\widehat F-F\|_1=2W_1(\widehat G,G).
\]
The centered CDF of $w_\beta$ is
$x-1/2+v(x)^T\beta$, with the vector from
Lemma~\ref{p04:nxb:lem:rotation}; centering removes the otherwise
non-rotating primitive constants. Apply
Lemma~\ref{p04:nxb:lem:absolute-fisher} at every $x$ to the centered
estimated CDF minus $x-1/2$, then integrate and use
Lemma~\ref{p04:nxb:lem:rotation}. Every policy and estimator have Bayes,
and hence worst-case, $W_1$ risk in the $U$ parameter at least
\[
 c\frac{\sum_{k=1}^K k^{-1}}{\sqrt{B+K^5}}.
\]
Taking $K=\lfloor B^{1/5}\rfloor$ gives $c\log B/\sqrt B$.
Finally $|e^{au}-e^{av}|\geq a|u-v|$ transfers the same lower
order to $\theta\in[1,2]$.
\end{proof}

\begin{proof}[Upper bound in Theorem~\ref{p04:nxb:thm:uniform}]
Let $T_r=\max(1,M_r)$, where $M_r$ is an individual's sample
maximum after $r$ observations. Uniformly in $\theta\in[1,2]$,
$\mathbb E_\theta|T_r-\theta|\leq2/(r+1)$.
Take dyadic depths $r_\ell=2^\ell$ up to
$R=2^L\asymp\sqrt B$. Estimate the CDF of $T_1$ at the base level;
at level $\ell\geq1$ average independent copies of
\[
 D_\ell(x)=\mathbf1_{\{T_{r_\ell}\leq x\}}
              -\mathbf1_{\{T_{r_{\ell-1}}\leq x\}},
 \qquad 1\leq x\leq2,
\]
using prefixes of the same individual's observations within a copy,
and fresh individuals across copies and levels. Its integrated variance
is at most
$\mathbb E|T_{r_\ell}-T_{r_{\ell-1}}|\leq C/r_\ell$.
Use $n_\ell=\lceil cB/[(L+1)r_\ell]\rceil$ copies, with a small
fixed $c>0$. The total cost is at most $B$ for all large $B$, because
the rounding overhead is $O(R)=O(\sqrt B)$.
The sum of the level means has integrated variance at most
$C(L+1)^2/B$ and mean equal to the CDF of $T_R$.
Jensen and Cauchy--Schwarz give expected integrated absolute error
$C(L+1)/\sqrt B$; its bias to the latent CDF is at most $2/(R+1)$.
Project onto probability CDFs in $L^1([1,2])$, or use a measurable
approximate projection over a countable dense set, losing at most a
factor two and an arbitrarily small error. The asserted upper bound
follows.
\end{proof}

The result closes the scalar endpoint critical case for this observation
model. It does not claim a universal critical-rate theorem for every
component family with a $1/r$ estimator: the posterior trace identity
and its stable adaptive extension use the uniform-endpoint likelihood.
The capped Bernoulli proof is independent of the near-linear MLE result
in Section~\ref{p04:nxm:sec:main}. A change of estimator and observation
design is essential to the cap theorem.

% ---- BEGIN bessel ----
\section{The full ID domains of two McKay families}
\label{p04:nxs:section}

The parameter-domain question in Section~2.3 of
Baricz, Prabhu K, Singh, and Vijesh~\cite{p01:rare:Baricz2026} asks how far
infinite divisibility extends for
\[
 f_{\mu,\nu,a,b}(x)=C x^{\nu-1}e^{-bx}I_\mu(ax),\qquad x>0.
\]
The answer proved below is the entire domain where the source defines
this probability density.  This replaces the integer-bias restriction
in the preceding manuscript.  The proof uses the classical zero-count
of Runckel, Jacobi Sturm oscillation, and a signed Cauchy representation.
The new step is that the sign-changing spectral representation has a
positive inverse Laplace transform throughout the admissible domain.
A second application reaches the entire parameter domain specified in
the same source for the squared-Bessel McKay density as well.
Throughout this section, ``admissible domain'' means the domain stated
in that source; exceptional negative-integer Bessel orders outside its
convention are not included in a maximality claim.

Write $F(a,b;c;z)={}_2F_1(a,b;c;z)$.  ID means infinitely divisible
and SD means self-decomposable, in each case for probability laws on
$[0,\infty)$.

\subsection{Negative zeros and a signed spectral measure}

\begin{lemma}[A special case of the classical zero count]
\label{p04:nxs:zeros}
Suppose
\[
 \alpha>-\tfrac12,\qquad \lambda>0,
 \qquad \alpha+\lambda>0,
 \qquad c=\alpha+\tfrac12.
\]
If $2\lambda$ is not an integer, the function
\[
 V(z)=F(-\lambda,\tfrac12-\lambda;c;z)
\]
has exactly $N=\lceil\lambda-\tfrac12\rceil$ zeros in the cut plane
$\mathbb C\setminus[1,\infty)$.  All are simple and negative real.
There are no zeros on either bank of $(1,\infty)$.
\end{lemma}

\begin{proof}
Runckel's zero-count theorem, in the form of
\cite[Theorem~3]{p04:nxs:DyachenkoKarp}, gives
\[
 N=\lfloor\lambda\rfloor+
       \frac{1+\operatorname{sgn}\{\Gamma(-\lambda)
                       \Gamma(\tfrac12-\lambda)\}}2.
\]
Indeed, the other two gamma arguments in that theorem are
$c+\lambda>0$ and $\alpha+\lambda>0$; its smallest argument is
$-\lambda$, and its largest is positive.  The duplication formula
reduces the sign in the display to that of $\Gamma(-2\lambda)$.
The resulting integer is $\lceil\lambda-\tfrac12\rceil$.
The nonvanishing on the cut follows from the nondegenerate case of
Runckel's theorem, as recorded in
\cite[Lemma~3]{p04:nxs:DyachenkoKarp}.

It remains to locate every counted zero.  Pfaff's transformation~\cite[Eq.~15.8.1]{p01:rare:DLMF} gives,
for $z<0$ and $y=z/(z-1)\in(0,1)$,
\[
 V(z)=(1-z)^\lambda v_\lambda(y),\qquad
 v_\lambda(y)=F(-\lambda,\alpha+\lambda;c;y).
\]
Put $\mu=\alpha-1/2>-1$.  The latter function is the solution regular
at zero of the Jacobi Sturm equation
\begin{equation}\label{p04:nxs:sturm}
 -(p v')'=E_\lambda w v,\qquad
 p(y)=y^{\mu+1}(1-y)^{1/2},\quad
 w(y)=y^\mu(1-y)^{-1/2},\quad
 E_\lambda=\lambda(\lambda+\alpha),
\end{equation}
normalized by $v(0)=1$.  The left boundary condition is
$\lim_{y\downarrow0}p(y)v'(y)=0$.  At the right endpoint take the
Dirichlet condition $v(1)=0$.  These are the usual separated Jacobi
boundary conditions; the positive quadratic form is
$\int_0^1p|v'|^2$ in $L^2(w\,dy)$.
Gauss's evaluation gives
\begin{equation}\label{p04:nxs:gauss}
 v_\lambda(1)=
 \frac{\Gamma(c)\Gamma(1/2)}
      {\Gamma(c+\lambda)\Gamma(1/2-\lambda)}.
\end{equation}
The complete eigenfunctions can also be written explicitly:
\[
 v_{n+1/2}(y)=(1-y)^{1/2}
 F(-n,n+\mu+3/2;\mu+1;y),\qquad n\geq0.
\]
They are positive scalar multiples of
$(1-y)^{1/2}P_n^{(\mu,1/2)}(1-2y)$.
Multiplication by $(1-y)^{1/2}$ is a unitary map from the Jacobi
space with weight $y^\mu(1-y)^{1/2}$ onto $L^2(w\,dy)$.
Jacobi completeness therefore supplies the entire spectrum, with
eigenvalues $E_{n+1/2}$, $n=0,1,2,\ldots$.
The zero-flux condition at zero excludes the other local solution
$y^{-\mu}$, or the logarithmic solution when $\mu=0$.
There are no nonpositive eigenvalues:
the quadratic form is nonnegative, and its only zero-energy solution
satisfying the left condition is constant, which cannot satisfy the
right condition.  Also $E_\lambda$ is strictly increasing on
$\lambda>\max(0,-\alpha)$, the parameter range under consideration.

Sturm oscillation therefore says that $v_\lambda$ has as many zeros
in $(0,1)$ as there are eigenvalues $E_{n+1/2}<E_\lambda$.
This is precisely $N$.  The statement applies to the possibly singular
left endpoint by the stated zero-flux boundary condition; equivalently
one can exhaust by regular separated problems and take the regular
solution limit.  All interior zeros are simple by uniqueness for the
second-order equation.  Their images under $z=y/(y-1)$ are $N$
simple negative zeros of $V$, exhausting Runckel's total count.
\end{proof}

\begin{lemma}[The signed Cauchy measure]
\label{p04:nxs:signed}
Under the noninteger hypotheses of Lemma~\ref{p04:nxs:zeros}, put
\[
 q=2\lambda,\qquad A=\alpha+q,
 \qquad
 R(w)=\frac{\lambda}{w}
       -\frac{V'(1/w)}{w^2V(1/w)}.
\]
Write the negative zeros of $V(1/w)$ as $-\tau_1,\ldots,-\tau_N$,
where each $\tau_j>0$.  There is a finite real measure $\rho$ on
$(0,1)$ such that
\begin{equation}\label{p04:nxs:cauchy}
 R(w)=\sum_{j=1}^{N}\frac1{w+\tau_j}
                  +\int_0^1\frac{\rho(dt)}{w-t}.
\end{equation}
It has constant sign, total mass $\lambda-N$, and density
\begin{equation}\label{p04:nxs:rho-density}
 \frac{d\rho}{dt}=
 -\frac{\Gamma(c)^2}
 {\Gamma(-\lambda)\Gamma(1/2-\lambda)
  \Gamma(c+\lambda)\Gamma(\alpha+\lambda)}
 \frac{t^{-q-1/2}(1-t)^{A-1}}
      {|V(1/t+i0)|^2},\qquad 0<t<1.
\end{equation}
In particular $\rho$ has no endpoint atoms.
\end{lemma}

\begin{proof}
Although $w^\lambda V(1/w)$ itself need not be single-valued around
the cut, its logarithmic derivative $R$ is a well-defined real
meromorphic function on $\mathbb C\setminus[0,1]$.
Its only poles there are the simple zeros identified in
Lemma~\ref{p04:nxs:zeros}, each with residue one, and
$R(w)=\lambda/w+O(w^{-2})$ at infinity.

For clarity we derive the jump.  Write $V(x+i0)=U(x)+iW(x)$ for
$x>1$.  The hypergeometric equation implies
\[
 (UW'-U'W)(x)=D\,x^{-c}(x-1)^{A-1}.
\]
The connection formula at one gives
\[
 D=\frac{\pi\Gamma(c)^2}
 {\Gamma(-\lambda)\Gamma(1/2-\lambda)
  \Gamma(c+\lambda)\Gamma(\alpha+\lambda)}.
\]
For example, for noninteger $A$ its two leading connection coefficients
are $P=\Gamma(c)\Gamma(A)/[\Gamma(c+\lambda)
\Gamma(\alpha+\lambda)]$ and
$Q=\Gamma(c)\Gamma(-A)/[\Gamma(-\lambda)
\Gamma(1/2-\lambda)]$.  On the upper bank the second term has phase
$e^{-i\pi A}$, and $D=-PQ A\sin(\pi A)$ gives the displayed value by
reflection.  Integer $A>0$ follows by continuity.
Thus $\operatorname{Im}(V'/V)(x+i0)=
D x^{-c}(x-1)^{A-1}/|V(x+i0)|^2$.
The map $w\mapsto1/w$ swaps the banks, so
$-\operatorname{Im}R(t+i0)/\pi$ is exactly
\eqref{p04:nxs:rho-density}.  Its sign is constant.

The connection formula at infinity, whose exponent difference is
$1/2$, gives a nonzero constant term followed by a square-root term
in $w^\lambda V(1/w)$ near zero.  The constant is nonzero because
$2\lambda$ is not an integer and
$c+\lambda,\alpha+\lambda>0$.  Hence
$R(w)=O(|w|^{-1/2})$ there.
At one, Gauss's evaluation gives a finite nonzero value since
$A>0$.  The derivative has at worst
$O(1+|1-w|^{A-1}(1+|\log(1-w)|))$ growth.
These estimates make both small-circle integrals vanish and the jump
absolutely integrable.  Cauchy's formula on a keyhole contour around
$[0,1]$, after subtracting the residue-one poles, proves
\eqref{p04:nxs:cauchy}.  Comparison of the coefficients of $1/w$ at
infinity gives $\rho((0,1))=\lambda-N$.
The endpoint estimates also exclude endpoint atoms.
\end{proof}

\subsection{The maximal ID domain}

\begin{theorem}[Every admissible generalized McKay law is ID]
\label{p04:nxs:full-id}
For every
\begin{equation}\label{p04:nxs:domain}
 \mu>-1,\qquad \mu+\nu>0,\qquad b>a>0,
\end{equation}
the normalized density
\[
 f_{\mu,\nu,a,b}(x)=C x^{\nu-1}e^{-bx}I_\mu(ax),\qquad x>0,
\]
is infinitely divisible.  Thus the maximal parameter domain asked for
in \cite[Section~2.3]{p01:rare:Baricz2026}, within the density family and
normalization convention of that source, is the whole domain
\eqref{p04:nxs:domain}.

More explicitly, if $q=\nu-\mu-1>0$, let
$\alpha=\mu+1/2$, $M=\mu+\nu=2\alpha+q$, and $A=\alpha+q$.
For noninteger $q$, use the zeros and signed measure in
Lemma~\ref{p04:nxs:signed} and define
\begin{equation}\label{p04:nxs:kernel}
 K_{\alpha,q}(x)=2A\cosh x
       -2\sum_{j=1}^{N}\cos(\sqrt{\tau_j}x)
       -2\int_0^1\cosh(\sqrt t\,x)\,\rho(dt).
\end{equation}
Then the L\'evy density is $e^{-bt}K_{\alpha,q}(at)/t$, and
\begin{equation}\label{p04:nxs:positive}
 K_{\alpha,q}(x)\geq M\cosh x>0\qquad(x\geq0).
\end{equation}
\end{theorem}

\begin{proof}
For $q\leq0$, the source's Theorem~3 proves ID on precisely this
part of the admissible domain.  Consider $q>0$.
Scaling permits $a=1$ in the transform calculation.  The positive
Bessel series, integrated termwise, gives
\[
 J(u):=\int_0^\infty x^{\nu-1}e^{-ux}I_\mu(x)\,dx
   =C_0u^{-M}F(M/2,(M+1)/2;\mu+1;u^{-2}),\qquad u>1,
\]
where $C_0>0$.  Euler's transformation yields
\begin{equation}\label{p04:nxs:J}
 J(u)=C_0\frac{u^q V(u^{-2})}{(u^2-1)^A}.
\end{equation}
Notice that $\alpha>-1/2$ and $\alpha+q/2=M/2>0$, exactly the
hypotheses needed in the two lemmas.
For noninteger $q$, logarithmic differentiation gives
\[
 -\frac{J'(u)}{J(u)}=
 \frac{A}{u-1}+\frac{A}{u+1}
 -2\sum_{j=1}^{N}\frac{u}{u^2+\tau_j}
 -2\int_0^1\frac{u}{u^2-t}\,\rho(dt).
\]
Each term is the Laplace transform of its corresponding term in
\eqref{p04:nxs:kernel}.  The signed integral is legitimate because
$\rho$ has finite total variation and $u>1$.

The essential positivity argument uses only the total mass and sign
of $\rho$.  If $d=\lambda-N\geq0$, then
\[
 K_{\alpha,q}(x)\geq2A\cosh x-2N-2d\cosh x
       =M\cosh x+2N(\cosh x-1)\geq M\cosh x.
\]
If $d\leq0$, then
\[
 K_{\alpha,q}(x)\geq2A\cosh x-2N+2(N-\lambda)
       =2A\cosh x-q\geq M\cosh x.
\]
Thus the signed spectral measure cannot produce a negative L\'evy
kernel.  Since $K=O(e^x)$ and $K(0)=M$, the measure
$e^{-bt}K(at)\,dt/t$ satisfies the L\'evy integrability condition for
$b>a$.  Integrating the logarithmic derivative and using $L(0)=1$
gives
\[
 -\log L(s)=\int_0^\infty(1-e^{-st})
             e^{-bt}\frac{K_{\alpha,q}(at)}t\,dt.
\]
This proves ID.

For integer $q\geq1$, one may take a sequence of admissible
noninteger parameters tending to $q$: the normalized densities
converge in $L^1$ by the usual Bessel bounds at zero and infinity,
and the ID class is closed under weak convergence.
There is also a direct formula.  If $q=2m$, $V$ is a degree-$m$
polynomial with simple negative zeros, $N=m$, and $\rho=0$.
If $q=2m+1$, $V$ is a degree-$m$ polynomial with simple negative
zeros and the logarithmic derivative of $u^qV(u^{-2})$ has one
additional $1/u$ term.  The negative-zero assertion follows from
the same Jacobi Sturm argument, including the right-endpoint
eigenvalue case.  Hence the integer kernel is
\[
 K_{\alpha,q}(x)=2A\cosh x
       -2\sum_{j=1}^{\lfloor q/2\rfloor}\cos(\sqrt{\tau_j}x)
       -\mathbf1_{\{q\text{ odd}\}},
\]
and it too satisfies $K\geq M\cosh x$.
\end{proof}

\begin{corollary}[A spectral SD boundary on the full excess-power domain]
\label{p04:nxs:sd-general}
For every admissible $(\alpha,q)$ with $q>0$ define the explicit
kernel of Theorem~\ref{p04:nxs:full-id}, using its integer convention
when necessary, and put
\[
 \beta_{\alpha,q}=\sup_{x\geq0}
                \frac{K_{\alpha,q}'(x)}{K_{\alpha,q}(x)}.
\]
Then $1\leq\beta_{\alpha,q}<\infty$, and the generalized McKay law is
SD if and only if $b/a\geq\beta_{\alpha,q}$.
Moreover $\beta_{\alpha,q}>1$ whenever $q$ is a positive integer or
$q\in(2m,2m+1)$ for an integer $m\geq0$; in these cases the supremum
is attained at a finite positive point.
\end{corollary}

\begin{proof}
The kernel is positive and smooth.  Its cosine terms are bounded,
and $\rho$ has no atom at one, so dominated convergence gives
\[
 e^{-x}K_{\alpha,q}(x)\longrightarrow A,
 \qquad e^{-x}K_{\alpha,q}'(x)\longrightarrow A.
\]
Thus $K'/K\to1$, is bounded, and its supremum is at least one.
The SD criterion for a positive ID law with smooth spectral function
$k(t)$ is $k'(t)\leq0$.  Applied to $k(t)=e^{-bt}K(at)$ this is
exactly the stated inequality.

For completeness this criterion follows directly from the defining
SD ratios.  A nonincreasing $k$ makes
$-\log L(s)+\log L(cs)$ a Bernstein exponent for each $0<c<1$,
with L\'evy density $[k(t)-k(t/c)]/t$.
Conversely, if all $L(s)/L(cs)$ are Laplace transforms, their
negative derivatives divided by $1-c$ converge, with derivatives
locally, to $(s[-\log L(s)]')'$.  This is completely monotone, while
integration by parts identifies it as the Laplace transform of
$-t k'(t)$.  Uniqueness of the Laplace transform forces $k'\leq0$.

For noninteger $q\in(2m,2m+1)$ the measure $\rho$ is strictly
positive on $(0,1)$.  In the difference $K'-K$, its contribution is
\[
 2\int_0^1[\cosh(\sqrt t\,x)
                  -\sqrt t\sinh(\sqrt t\,x)]\,\rho(dt).
\]
Any compact subinterval of $(0,1)$ has positive mass, so this
quantity grows exponentially.  The cosine contribution is bounded,
and the remaining contribution is $-2A e^{-x}$.
It follows that $K'-K>0$ for all sufficiently large $x$.
For positive integer $q$, simultaneous returns of the finitely many
cosine phases make the corresponding bounded contribution to
$K'-K$ arbitrarily close to $q>0$ at arbitrarily large times.
This gives $\beta>1$ in both cases.  Since $K'/K\to1$, the larger
maximum is attained on a compact interval.
\end{proof}

The supremum criterion is a spectral classification, not a claim that
its maximizing point has a closed elementary expression.  For
$q\in(2m+1,2m+2)$ the preceding theorem proves ID and the same exact
SD test, but does not assert that $\beta_{\alpha,q}>1$ for every
$\alpha$.  That strictness question is kept separate.

\subsection{The full ID domain of the squared-Bessel family}

\begin{lemma}[A zero-free disk for the transformed squared-Bessel factor]
\label{p04:nxs:disk}
For every $r>0$, the function
\[
 F_r(z)=F(-r,1/2;r+3/2;z)
\]
has no zero in the closed unit disk.  If $r$ is not an integer,
its total number of zeros in $\mathbb C\setminus[1,\infty)$ is
\[
 N_r=2\left\lfloor\frac{r+1}{2}\right\rfloor,
\]
and it has no zero on either bank of $(1,\infty)$.
\end{lemma}

\begin{proof}
The classical quadratic transformation~\cite[Eq.~15.8.16]{p01:rare:DLMF} gives
\begin{equation}\label{p04:nxs:quadratic}
 F_r(z)=(1-z)^{-1/2}
 F\left(1/4,r+3/4;r+3/2;
                 \frac{-4z}{(1-z)^2}\right).
\end{equation}
The map $z\mapsto-4z/(1-z)^2$ maps the open unit disk onto
$\mathbb C\setminus[1,\infty)$.  This follows either from the
usual slit-plane map or by solving its quadratic inverse; on
$z=e^{i\theta}\ne1$ its value is $1/\sin^2(\theta/2)$.
Euler's beta integral for the last hypergeometric function is an
average of $(1-ZT)^{-1/4}$, with
$T\sim\operatorname{Beta}(r+3/4,3/4)$.
For nonreal $Z$, these integrands lie in a common sector of width
at most $\pi/4$, so their average is nonzero.  For real $Z<1$
they are positive.  This proves nonvanishing in the open disk.

For finite $Z>1$, the upper and lower boundary integrals also
exist: the interior singularity has exponent $-1/4$ and is
integrable.  All their integrands have arguments either zero or
$\pm\pi/4$, so their real parts are positive.  At $Z=1$ the beta
weight gives an integrable endpoint singularity and the integral
is positive.  Formula~\eqref{p04:nxs:quadratic} therefore excludes
zeros on the unit circle except possibly at $z=1$.  At that point,
Gauss's formula gives the positive value
\[
 F_r(1)=\frac{\Gamma(r+3/2)\Gamma(2r+1)}
                  {\Gamma(2r+3/2)\Gamma(r+1)}>0.
\]
Finally Runckel's theorem gives, for noninteger $r$,
\[
 N_r=\lfloor r\rfloor+
                  \frac{1+\operatorname{sgn}\Gamma(-r)}2
       =2\left\lfloor\frac{r+1}{2}\right\rfloor.
\]
Its nondegenerate-cut assertion excludes all zeros on the banks.
All off-cut zeros are simple: a simultaneous zero of $F_r$ and
$F_r'$, away from the regular singular points, would contradict
uniqueness for the hypergeometric differential equation.
The values at zero and one are nonzero.
\end{proof}

\begin{theorem}[Every admissible squared-Bessel McKay law is ID]
\label{p04:nxs:fourth-full}
For every $\mu>-1/4$ and $b>2a>0$, the normalized density
\begin{equation}\label{p04:nxs:fourth-density}
 g_{\mu,a,b}(x)=C x^{2\mu}e^{-bx}I_\mu(ax)^2,
 \qquad x>0,
\end{equation}
is infinitely divisible.

For an explicit formula when $r=\mu-1/2>0$ is not an integer,
put $B=2r+1$, $c=r+3/2$, and let $w_1,\ldots,w_N$ be the
reciprocals of all zeros of $F_r$, with $N=N_r$ as in
Lemma~\ref{p04:nxs:disk}.  In particular $|w_j|<1$.
Define the real signed measure
\begin{equation}\label{p04:nxs:fourth-rho}
 \frac{d\sigma_r}{dt}=
 -\frac{\Gamma(c)^2}
 {\Gamma(-r)\Gamma(1/2)\Gamma(2r+3/2)\Gamma(r+1)}
 \frac{t^{-r-1/2}(1-t)^{2r}}
      {|F_r(1/t+i0)|^2},\qquad 0<t<1.
\end{equation}
It has constant sign and total mass $r-N$.  The L\'evy spectral
function of \eqref{p04:nxs:fourth-density} is
$e^{-bt}K_r(2at)$, where
\begin{equation}\label{p04:nxs:fourth-kernel}
 K_r(x)=B(1+2\cosh x)
       -2\sum_{j=1}^{N}\cosh(\sqrt{w_j}\,x)
       -2\int_0^1\cosh(\sqrt t\,x)\,\sigma_r(dt).
\end{equation}
The choice of either square root gives the same summand.
The kernel is real, positive, and even.
\end{theorem}

\begin{proof}
For $-1/4<\mu\leq1/2$, ID follows from the source's Theorem~4
(and the GGC representation in the earlier manuscript).
Let $r=\mu-1/2>0$.  The Bessel-square Laplace integral, with
scale $2a$ set to one, and Euler's transformation give
\begin{equation}\label{p04:nxs:fourth-J}
 J(u)=C_1u^{-4r-3}F(r+1,2r+3/2;r+3/2;u^{-2})
   =C_1\frac{u^{2r}F_r(u^{-2})}
                  {u^{2r+1}(u^2-1)^{2r+1}}.
\end{equation}
For noninteger $r$, apply the same logarithmic-derivative keyhole
argument as in Lemma~\ref{p04:nxs:signed} to
$H_r(w)=w^rF_r(1/w)$.  It gives
\[
 \frac{H_r'(w)}{H_r(w)}=
       \sum_{j=1}^{N}\frac1{w-w_j}
              +\int_0^1\frac{\sigma_r(dt)}{w-t}.
\]
The Wronskian computation with parameters $(-r,1/2;c)$ gives
exactly \eqref{p04:nxs:fourth-rho}.  Near zero the leading term of
$H_r$ is a nonzero constant; its next nonanalytic exponent is
$r+1/2>0$, with a possible logarithm when this is an integer.
Thus there is no pole of its logarithmic derivative at zero, and
the jump is integrable there.  At one $H_r(1)=F_r(1)>0$ and the
singular exponent is $2r+1>1$.  There is no endpoint atom.
Comparison at infinity gives the mass $r-N$.
Taking Laplace transforms in \eqref{p04:nxs:fourth-J} now gives
\eqref{p04:nxs:fourth-kernel}.

Complex conjugation pairs the $w_j$, making the kernel real.
Let $\varrho=\max_j\sqrt{|w_j|}<1$ if there are any zeros;
otherwise omit all zero terms.  The bound
$|\cosh z|\leq\cosh|z|$ gives
\[
 \left|\sum_{j=1}^{N}\cosh(\sqrt{w_j}\,x)\right|
       \leq N\cosh(\varrho x)\leq N\cosh x.
\]
If $r-N\geq0$, use the positive measure's mass to obtain
\begin{equation}\label{p04:nxs:fourth-positive1}
 K_r(x)\geq B+2(B-r)\cosh x
               =B+(2r+2)\cosh x>0.
\end{equation}
If $r-N\leq0$, instead use $\cosh(\sqrt t\,x)\geq1$ for its
positive contribution, giving
\begin{equation}\label{p04:nxs:fourth-positive2}
 K_r(x)\geq B+2(B-N)\cosh x+2(N-r)
             =2N+1+2(B-N)\cosh x>0.
\end{equation}
Here $N<r+1$ for noninteger $r$, so $B-N>r>0$.
The kernel has growth $O(e^x)$ and $K_r(0)=4r+3$.
The positive L\'evy representation, with $b>2a$, follows exactly
as in Theorem~\ref{p04:nxs:full-id}.

For integer $r=n\geq1$, $F_r$ is a polynomial of degree $n$.
Lemma~\ref{p04:nxs:disk} still puts all reciprocal roots strictly
inside the unit disk.  Use these $n$ roots and set $\sigma_r=0$
in \eqref{p04:nxs:fourth-kernel}.  Its logarithmic derivative is a
finite partial-fraction sum, and
$K_r(x)\geq B+2(B-n)\cosh x>0$ proves ID directly.
\end{proof}

\begin{corollary}[The squared-Bessel SD boundary above order one half]
\label{p04:nxs:fourth-sd}
For every $r>0$, use the explicit kernel in
Theorem~\ref{p04:nxs:fourth-full}, with its polynomial convention at
integers, and set
\[
 \beta_r=\sup_{x\geq0}\frac{K_r'(x)}{K_r(x)}.
\]
Then $1\leq\beta_r<\infty$, and
\[
 g_{r+1/2,a,b}\text{ is SD}\quad\Longleftrightarrow\quad
                          b/(2a)\geq\beta_r.
\]
For $0<r<1$ one has the quantitative strict bound
\[
 1<\beta_r\leq\frac{2r+1}{r+1}<\frac32,
\]
and the maximum is attained at a finite positive point.
More generally $\beta_r>1$ for $r\in(2m,2m+1)$, $m\geq0$.
\end{corollary}

\begin{proof}
All zeros have square roots of modulus less than one, and the
finite measure $\sigma_r$ has no atom at one.  Therefore
$e^{-x}K_r(x)\to B$ and $e^{-x}K_r'(x)\to B$.
The positivity bounds in the theorem show that $K'/K$ is continuous
and bounded; its limit is one.  The SD criterion proved in
Corollary~\ref{p04:nxs:sd-general} gives the equivalence.

If $0<r<1$, there are no zero terms and $\sigma_r$ is a positive
measure of mass $r$.  Thus $K_r'\leq2B\sinh x$, and
\eqref{p04:nxs:fourth-positive1} gives the stated upper bound.
In general, on $r\in(2m,2m+1)$ the density of $\sigma_r$ is
strictly positive throughout $(0,1)$.  The positive integral
contribution to $K_r'-K_r$ contains exponential growth of every
rate less than one.  Choose such a rate strictly larger than
$\max_j\sqrt{|w_j|}$.  It dominates all zero contributions and
the negative term $-B(1+2e^{-x})$.  Hence $K_r'/K_r>1$ at some
finite point and the supremum is attained.
\end{proof}

\subsection{A quantitative tail for the spectral threshold}

\begin{proposition}[Polynomially certified tail control]
\label{p04:nxs:tail}
For fixed admissible parameters with noninteger $q>0$ in the third
family or noninteger $r>0$ in the fourth family, there exist finite
positive constants $C,c_0$ such that, for $x\geq1$,
\[
 \left|\frac{K'(x)}{K(x)}-1\right|
 \leq C\{e^{-c_0x}+x^{-\gamma}\},
 \qquad
 \gamma=\begin{cases}
          \alpha+q,&\text{third family},\\
          2r+1,&\text{fourth family}.
        \end{cases}
\]
The constants can be bounded from the displayed cut densities and
the finite zero sets.  In particular, a certified point with
$K'(x)/K(x)>1+\delta$ reduces the search for the exact maximum to
a finite interval.
\end{proposition}

\begin{proof}
In both cases the explicit density $h(t)$ of the cut measure obeys
$|h(t)|\leq C_1(1-t)^{\gamma-1}$ on $[1/2,1)$:
the denominator has a positive continuous modulus, including its
nonzero limit at one.  On $(0,1/2)$ its absolute mass is finite.
The kernel is bounded below by $c_1\cosh x$ from
\eqref{p04:nxs:positive}, \eqref{p04:nxs:fourth-positive1}, or
\eqref{p04:nxs:fourth-positive2}.  After division by this bound, the
absolute cut contribution to $K'-K$ is at most a constant times
\begin{align*}
 \int_0^1e^{-(1-\sqrt t)x}|h(t)|\,dt
 &\leq C_2 e^{-(1-1/\sqrt2)x}
     +C_1\int_0^1 e^{-(1-t)x/2}(1-t)^{\gamma-1}\,dt\\
 &\leq C_2 e^{-(1-1/\sqrt2)x}
       +C_1\Gamma(\gamma)(2/x)^\gamma.
\end{align*}
The finite zero terms are bounded in the third family, and grow
at a rate strictly smaller than one in the fourth family.
The remaining elementary terms decay exponentially after the same
division.  These estimates prove the claim.
\end{proof}

\paragraph{Scope and priority.}
The full third-family ID statement closes the maximal-domain question
in the stated density domain, subject to independent verification of
this proof.  The fourth-family theorem also reaches its entire
admissible real-order domain.  Runckel's zero count, Jacobi oscillation,
and the hypergeometric transformation identities are classical inputs
and are not new results.  The SD answers are explicit spectral
threshold classifications; they do not provide elementary closed
forms for every threshold, and universal strictness above one in the
negative-cut intervals has not been established.  These distinctions
are necessary when counting externally posed problems rather than
individual consequences.

% ---- BEGIN gaussian ----
\section{A sharp logarithmic tail boundary for sparse Gaussian observation}
\label{p04:nxg:sec:tails}

The unrestricted terminal-design theorem in the preceding manuscript
requires divergence of a harmonic subseries. Bounded covariance support, in
contrast, makes every infinite set of terminal lengths sufficient.
The results below locate a sharp boundary between these two regimes.
The boundary is finer than the square-root exponential moment: it has
a logarithmic correction. The uniqueness argument uses classical
Stieltjes moment determinacy. The converse requires more than a pair
of moment-equivalent measures: it constructs one joint mixing density
whose whole observation transforms vanish at the selected designs.

Write
\[
 \omega_a(v)=\frac{\sqrt v}{\{\log(e+v)\}^{a}},\qquad v\ge0,
 \quad a\ge0.
\]
Let $(A,\Sigma)$ be a random element of
$\mathbb R^d\times\mathbb S_d^+$ and, conditionally on it, let
$X_1,X_2,\ldots$ be iid $N_d(A,\Sigma)$. As before the separate
experiment of terminal length $n\ge d$ observes
\[
 \mathcal O_n=(X_1,\ldots,X_{d-1},S_n),\qquad
 S_n=\sum_{j=1}^nX_j.
\]
Separate experiments supply their separate population laws; their
joint cross-design distribution is not assumed available.

\begin{theorem}[Logarithmic tail threshold]
\label{p04:nxg:thm:log-threshold}
\begin{enumerate}
\item For every $d\ge1$, every infinite set
$\mathcal N\subset\{d,d+1,\ldots\}$, and every $0\le a\le1$,
the laws $\{\mathcal L(\mathcal O_n):n\in\mathcal N\}$ identify
the mixing law within the class
\begin{equation}\label{p04:nxg:eq:tail-class}
 \mathbb E\exp\{c\omega_a(\operatorname{tr}\Sigma)\}<\infty
 \quad\text{for some }c>0.
\end{equation}
No moment or support assumption on $A$ is needed.
\item For $d=1$ and $\mathcal N=\{2^j:j\ge0\}$, the assertion
fails for every $a>1$. There are distinct absolutely continuous
mixing laws with identical $A$ marginals, identical variance
marginals, and identical finite joint polynomial moments, producing
the same law of $S_n$ for every $n\in\mathcal N$. Both laws satisfy
\begin{equation}\label{p04:nxg:eq:tail-all-c}
 \mathbb E\exp\{c\omega_a(\Sigma)\}<\infty
 \qquad\text{for every }c>0.
\end{equation}
Any fixed lower bound $\Sigma\ge v_0>0$ may be imposed.
\end{enumerate}
\end{theorem}

Thus $a=1$ is the exact boundary for the assertion that this tail
condition makes \emph{every} infinite terminal set identifying.
Sharpness for the full matrix-covariance experiment in each fixed
dimension $d>1$ is not asserted. A scalar counterexample cannot simply
be placed in one matrix coordinate: the additional anchors would then
provide extra scalar observations.

\subsection{Boundary zeros, moments, and the positive result}

We use the classical Stieltjes--Carleman criterion: a finite positive
measure on $[0,\infty)$ with moments $m_k$ is determined by its
moments if $\sum_{k\ge1}m_k^{-1/(2k)}=\infty$.
Hardy's square-root exponential condition and its relation to this
criterion are classical \cite{p04:nxg:StoyanovLin,p04:nxg:Lin}.

\begin{lemma}[A logarithmic moment bound]
\label{p04:nxg:lem:moment}
Suppose $V\ge0$ and $\mathbb E e^{c\omega_a(V)}<\infty$ for
some $c>0$ and $a\ge0$. Then a finite constant $C$ satisfies
\[
 (\mathbb E V^k)^{1/(2k)}
       \le Ck\{\log(e+k)\}^{a},\qquad k\ge1.
\]
Consequently the Stieltjes--Carleman condition holds when $a\le1$.
The same conclusion holds for any finite positive measure dominated
by a finite sum of laws satisfying the displayed exponential bound,
possibly with different positive constants $c$.
\end{lemma}
\begin{proof}
Put $Y=\sqrt V$. Since $\log(e+Y^2)$ and $\log(e+Y)$ are
comparable, $\mathbb E\exp\{c_1Y/\log^a(e+Y)\}=M<\infty$
for some $c_1>0$. Set $b_k=k\log^a(e+k)$. For $y=b_kw$ with
$w\ge1$,
\[
 \log(e+b_kw)\le C_a\log(e+k)(1+\log w),
 \qquad
 \frac{y}{\log^a(e+y)}
       \ge c_a\frac{kw}{(1+\log w)^a}.
\]
It follows that
\[
 \sup_{y\ge0}y^{2k}
     \exp\{-c_1y/\log^a(e+y)\}\le (C_1b_k)^{2k},
\]
because $w^2\exp\{-c_1c_aw/(1+\log w)^a\}$ is bounded on
$[1,\infty)$; for $w<1$ the claim is immediate. Hence
$\mathbb E Y^{2k}\le M(C_1b_k)^{2k}$. Absorb
$\max(1,\sqrt M)$ into $C$. The series
$\sum[k\log^a(e+k)]^{-1}$ diverges precisely for $a\le1$.
A finite sum only changes the constants.
\end{proof}

\begin{lemma}[Accumulating positive zeros under a moment envelope]
\label{p04:nxg:lem:zeros}
Let $\nu$ be a finite complex measure on $[0,\infty)$, dominated
in total variation by a finite positive measure $\mu$ with all moments
finite and satisfying the Stieltjes--Carleman condition. If
\[
 \int e^{-t_jv}\,d\nu(v)=0
 \quad\text{for distinct }t_j>0\text{ tending to }0,
\]
then $\nu=0$.
\end{lemma}
\begin{proof}
The Laplace transform $f(t)$ is $C^\infty$ on $[0,\infty)$,
with $f^{(k)}(0)=(-1)^k\int v^k\,d\nu(v)$, by domination.
Continuity and the zeros first give $f(0)=0$. Inductively, if the
derivatives of orders below $k$ vanish at zero, Taylor's formula at
$t_j$ gives $0=f^{(k)}(0)t_j^k/k!+o(t_j^k)$, so the $k$th
derivative also vanishes. This argument applies directly to a complex
function. The nonnegative measures $\mu\pm\operatorname{Re}\nu$
have the same moments as $\mu$, so moment determinacy gives
$\operatorname{Re}\nu=0$. The same argument gives
$\operatorname{Im}\nu=0$.
\end{proof}

\begin{proof}[Proof of the positive assertion of
Theorem~\ref{p04:nxg:thm:log-threshold}]
Let $\rho$ be the signed difference of two candidate mixing laws.
The blocks
\[
 X_1,\ldots,X_{d-1},S_n-X_1-\cdots-X_{d-1}
\]
are conditionally independent with respective lengths
$m_1=\cdots=m_{d-1}=1$, $m_d=n-d+1$. Their Fourier arguments
for the mixing transform are exactly
\begin{equation}\label{p04:nxg:eq:arguments}
 u\in\mathbb R^d,\qquad
 T=\frac{uu^T}{2n}+C,
 \quad C\succeq0,\quad\operatorname{rank}C\le d-1.
\end{equation}
Indeed, let the frequency columns be $\sqrt{m_i}w_i$. Project
their matrix onto the unit vector
\[
 n^{-1/2}(\sqrt{m_1},\ldots,\sqrt{m_d}).
\]
The parallel part
gives $uu^T/(2n)$ and the orthogonal part gives $C$. Conversely
every such $C$ is realized by a matrix on that orthogonal complement.

Fix $u\ne0$ and a rank-$d-1$ positive semidefinite matrix $C$ with
$u$ not orthogonal to its kernel. Define a complex measure $\nu$ by
pushing forward
\[
 e^{iu^TA-\operatorname{tr}(C\Sigma)}\,d\rho(A,\Sigma)
 \quad\text{under }(A,\Sigma)\longmapsto v=u^T\Sigma u.
\]
It is dominated by the pushforward of the sum of the two candidate
laws. Since $v\le\|u\|^2\operatorname{tr}\Sigma$, the latter
measure satisfies the moment bound of Lemma~\ref{p04:nxg:lem:moment}.
Equality of observations and \eqref{p04:nxg:eq:arguments} give zeros
at $t=1/(2n)$ for every selected $n$. Therefore
Lemma~\ref{p04:nxg:lem:zeros} gives
\[
 \int e^{iu^TA-\operatorname{tr}((C+zuu^T)\Sigma)}\,d\rho=0,
 \qquad z>0.
\]
For any $T\succ0$ the choice
\[
 z=(u^TT^{-1}u)^{-1},\qquad C=T-zuu^T
\]
has precisely the required properties. Thus the full transform
vanishes for all real $u\ne0$ and $T\succ0$. Continuity gives
$u=0$. Holomorphy in the positive matrix tube, followed by Fourier
uniqueness after multiplication by $e^{-\operatorname{tr}\Sigma}$,
gives $\rho=0$.
\end{proof}

The proof also gives a more general assertion: a fixed covariance
moment envelope $\mathbb E(\operatorname{tr}\Sigma)^k\le M_k$
with $\sum M_k^{-1/(2k)}=\infty$ suffices. Specifying a common
envelope avoids an unwarranted closure assertion about arbitrary
unions of moment-determinate classes.

\subsection{A sector construction at the logarithmic endpoint}

The following scalar construction is also used for the transfer to
non-Gaussian noise. Throughout, inverse Laplace transforms are ordinary
integrable functions, not distributions supported at the boundary.

\begin{lemma}[Lacunary zeros with nearly square-root decay]
\label{p04:nxg:lem:sector}
Let $b_j=b_0 2^{-j}$, $j\ge0$, with $b_0>0$. For every $q>1$
there exists a nonzero real Schwartz function $f$ on $\mathbb R$,
supported in $[0,\infty)$, such that
\begin{align}
 &F(z):=\int_0^\infty e^{-zv}f(v)\,dv,
       &&\operatorname{Re}z>0,\label{p04:nxg:eq:laplace-f}\\
 &F(b_j)=0\quad(j\ge0),\qquad
       \int_0^\infty v^kf(v)\,dv=0\quad(k\ge0),
       \label{p04:nxg:eq:zero-moments}\\
 &|f^{(m)}(v)|\le C_m
       \exp\{-c_m\omega_q(v)\},
       &&v\ge1,\quad m=0,1,2,\ldots .\label{p04:nxg:eq:tail-f}
\end{align}
\end{lemma}
\begin{proof}
Put $p=2(q-1)>0$, choose $K>1$ with $\log K\ge\max(2p,2)$,
and use principal branches. In the cut plane
$\mathbb C\setminus(-\infty,0]$ set
\begin{equation}\label{p04:nxg:eq:F-log}
 B(z)=\prod_{j\ge0}\frac{z-b_j}{z+b_j},\qquad
 Q(w)=\frac{w}{\{\log(K+w)\}^{p}},\qquad
 F(z)=B(z)\exp\{-z^{1/2}-Q(1/z)\}.
\end{equation}
The product converges locally uniformly away from zero and its
negative poles. The other functions in \eqref{p04:nxg:eq:F-log} are
holomorphic in the cut plane: $K+w$ never lies on the logarithm's
cut there, and $\log(K+w)$ can lie on the nonpositive real axis
only when $w$ is negative real.

We record a uniform bound that is special to geometric zeros. For
every fixed $\theta_0<\pi$,
\begin{equation}\label{p04:nxg:eq:B-sector}
 |B(z)|\le C_{\theta_0},\qquad
 0<|z|<\infty,\quad |\arg z|\le\theta_0.
\end{equation}
To prove it, split the factors into $b_j\le|z|/2$,
$b_j\ge2|z|$, and the remaining at most three indices. For the
first group the positive part of the logarithm of the factor's
modulus is at most $C_{\theta_0}b_j/|z|$; for the second it is
at most $C_{\theta_0}|z|/b_j$. Both geometric sums are bounded.
The middle factors are bounded because $z$ stays a fixed angular
distance from their poles. This proves \eqref{p04:nxg:eq:B-sector}.
On the right half-plane the sharper bound $|B|\le1$ holds.

For $\operatorname{Re}w\ge0$ one has $\operatorname{Re}Q(w)\ge0$.
For example, if $\arg w=\vartheta\in[0,\pi/2]$, write
$\log(K+w)=a+i\varphi$. Then $a\ge\log K$ and
$0\le\varphi\le\vartheta$, so
\[
 0\le p\arg(a+i\varphi)\le(p/\log K)\vartheta
       \le\vartheta.
\]
Hence $0\le\arg Q(w)\le\vartheta$; the lower half-plane follows
by conjugation. In particular $F$ is bounded on the right half-plane
and is conjugation invariant.

There is stronger decay near zero, including in a narrow region to
the left of the imaginary axis. Put $r=e^{-L}$ and
\[
 \delta(r)=\frac{p\pi}{8L}.
\]
For all sufficiently large $L$, uniformly for
$|\arg z|\le\pi/2+\delta(r)$ and $|z|=r$,
\begin{equation}\label{p04:nxg:eq:nearzero}
 \operatorname{Re}Q(1/z)
       \ge \frac{c_p}{rL^{p+1}}.
\end{equation}
Here is a direct check of the endpoint angular estimate. Write
$1/z=e^{L+i\vartheta}$, $|\vartheta|\le\pi/2+\delta(r)$.
Uniformly over this range,
\[
 \log(K+e^{L+i\vartheta})
       =L+i\vartheta+O(e^{-L}),\qquad
 \arg Q(e^{L+i\vartheta})
       =\vartheta-p\arctan(\vartheta/L)+O(e^{-L}/L).
\]
The expression before the error is increasing in $\vartheta$ for
large $L$; its endpoint is
$\pi/2-3p\pi/(8L)+O(L^{-2})$. The modulus of $Q$ is
comparable to $e^L/L^p$. Its real part therefore has the lower
bound \eqref{p04:nxg:eq:nearzero}.

On the imaginary axis all derivatives of $B$ near zero have at
most polynomial growth in $1/r$: this follows either by Cauchy's
inequality in disks of radius proportional to $r$ inside a fixed
sector and \eqref{p04:nxg:eq:B-sector}, or by differentiating the
product. Derivatives of $Q(1/z)$ have the same property after
allowing powers of $\log(1/r)$. Estimate
\eqref{p04:nxg:eq:nearzero} dominates all of them. At infinity,
$\exp(-\sqrt z)$ and its derivatives provide rapid decay, while
$Q(1/z)$ and $B$ remain bounded with polynomially bounded
derivatives. Thus $F(iy)$ extends across $y=0$ to a Schwartz
function, flat at zero. Define
\[
 f(v)=\frac1{2\pi}\int_{\mathbb R}e^{iyv}F(iy)\,dy.
\]
It is real, Schwartz, and nonzero. It is supported in $[0,\infty)$:
for $x>0$, Fourier inversion on the line $x+i\mathbb R$ gives
$f_x(v)=e^{-xv}f(v)$ by contour shifting, and
\[
 |F(x+iy)|\le e^{-\operatorname{Re}\sqrt{x+iy}}
       \le e^{-\sqrt{|y|/2}}
\]
gives a bound on $f_x(v)$ independent of $x$. If $v<0$, letting
$x\to\infty$ forces $f(v)=0$. The same line-shift identity gives
\eqref{p04:nxg:eq:laplace-f}. Flatness at $z=0$ and dominated
differentiation give \eqref{p04:nxg:eq:zero-moments}.

It remains to prove the quantitative tail; Schwartz decay alone
would not suffice. Deform the inverse contour near zero onto
\[
 z_\pm(r)=r\exp\{\pm i(\pi/2+\delta(r))\},
             \qquad0<r\le r_0,
\]
where $r_0$ is sufficiently small. The line element is bounded by
a constant times $dr$. Continue both curves to infinity along a
fixed angle in $(\pi/2,\pi)$, joined in a compact region entirely
within the left half-plane. No negative pole or branch cut is
crossed. The connecting compact pieces contribute $O(e^{-c v})$.
The large-radius pieces are controlled by $e^{-cvr-c\sqrt r}$.
Small arcs around zero disappear by \eqref{p04:nxg:eq:nearzero};
large closing arcs disappear by $e^{-\sqrt z}$ and
$\operatorname{Re}z\le0$. These bounds justify the deformation
for the inverse of $z^mF(z)$ as well, for every fixed $m\ge0$.

For $0<r\le r_0$, writing $L=\log(1/r)$, the integrand is bounded
by a constant times
\begin{equation}\label{p04:nxg:eq:bent-exponent}
 r^m\exp\left\{-c\left(\frac{vr}{L}
                  +\frac1{rL^{p+1}}\right)\right\}.
\end{equation}
For all sufficiently large $v$ and all $L\ge L_0$,
\begin{equation}\label{p04:nxg:eq:minimum}
 \frac{ve^{-L}}{L}+\frac{e^L}{L^{p+1}}
       \ge c_p\frac{\sqrt v}{(\log v)^{1+p/2}}.
\end{equation}
Indeed, on $\tfrac14\log v\le L\le\log v$ the arithmetic--
geometric mean inequality gives the bound. Below that interval the
first term is at least a constant times $v^{3/4}/\log v$.
Above it the second term, eventually increasing in $L$, is at least
$v/(\log v)^{p+1}$. Reserve half the exponent in
\eqref{p04:nxg:eq:bent-exponent} and use the other half with
\eqref{p04:nxg:eq:minimum}; the remaining $r$ integral is bounded
uniformly. The same estimates apply to every $m$, proving
\eqref{p04:nxg:eq:tail-f} because $q=1+p/2$.
\end{proof}

\begin{lemma}[Joint realization with equal marginals]
\label{p04:nxg:lem:joint}
Let $f$ and $F$ be as in Lemma~\ref{p04:nxg:lem:sector}, and let
$\psi:\mathbb R^D\setminus\{0\}\to(0,\infty)$ be even and
$C^\infty$. Choose a nonzero real even
$h\in C_c^\infty(\mathbb R^D\setminus\{0\})$. Then
\begin{equation}\label{p04:nxg:eq:joint-density}
 g(a,v)=(2\pi)^{-D}\int_{\mathbb R^D}
       e^{-iu^Ta}h(u)\psi(u)f(\psi(u)v)\,du
\end{equation}
is a nonzero real Schwartz function supported in $v\ge0$, with
Fourier--Laplace transform
\begin{equation}\label{p04:nxg:eq:joint-transform}
 \int e^{iu^Ta-zv}g(a,v)\,da\,dv
       =h(u)F(z/\psi(u)).
\end{equation}
The right side is taken to be zero wherever $h$ vanishes. Both
coordinate marginals and all joint polynomial moments of $g$ vanish.
Moreover, for every $M\ge0$ there are $C_M,c_M>0$ such that
\begin{equation}\label{p04:nxg:eq:joint-tail}
 |g(a,v)|\le C_M(1+\|a\|)^{-M}
          e^{-c_M\omega_q(v)},\qquad v\ge1.
\end{equation}
\end{lemma}
\begin{proof}
On the compact support of $h$, $\psi$ is bounded above and bounded
away from zero, and all of its derivatives are bounded. Every fixed
number of $u$ derivatives of
$h(u)\psi(u)f(\psi(u)v)$ is therefore bounded by a polynomial in
$v$ times $e^{-c\omega_q(v)}$. Scaling the argument by a fixed
positive constant preserves this estimate. Polynomial factors can
be absorbed by decreasing $c$. Repeated integration by parts in
$u$ gives \eqref{p04:nxg:eq:joint-tail}; applying the same argument to
all $a,v$ derivatives proves the Schwartz assertion. Reality follows
from evenness, and nonzeroness from Fourier uniqueness. Fubini and
the change of variable $w=\psi(u)v$ give
\eqref{p04:nxg:eq:joint-transform}.

The transform in $a$ vanishes in a neighborhood of $u=0$, so its
derivatives there give $\int a^\alpha g(a,v)\,da=0$ for every
multi-index $\alpha$. This proves the vanishing $v$ marginal and
all joint polynomial moments. The $a$ marginal vanishes because
$F(0)=\int f=0$; equivalently integrate the Fourier transform in
$v$. All these differentiations and integrations are justified by
\eqref{p04:nxg:eq:joint-tail} and the Schwartz bounds.
\end{proof}

\begin{proof}[Proof of the negative assertion of
Theorem~\ref{p04:nxg:thm:log-threshold}]
Fix $a>1$ and choose $1<q<a$. Apply
Lemma~\ref{p04:nxg:lem:sector} with $b_j=2^{-j-1}$ and then
Lemma~\ref{p04:nxg:lem:joint} with $D=1$ and $\psi(u)=u^2$.
Since $g$ is real, nonzero, and of total integral zero, its positive
and negative parts have the same strictly positive mass $m$.
They define probability densities $g_+/m$ and $g_-/m$ for
$(A,V)$. The characteristic function of $S_n/n$, conditionally
on $(A,V)$, is $e^{iuA-u^2V/(2n)}$. The difference of the two
unconditional characteristic functions is
\[
 m^{-1}h(u)F(1/(2n))=0\qquad(n=2^j).
\]
The equal marginals and moments follow from
Lemma~\ref{p04:nxg:lem:joint}. Integrating \eqref{p04:nxg:eq:joint-tail}
in $a$ shows that both variance marginals have density bounded by
$Ce^{-c_0\omega_q(v)}$ for $v\ge1$. Since
$\omega_a(v)/\omega_q(v)\to0$, this implies
\eqref{p04:nxg:eq:tail-all-c}. Replacing $V$ by $V+v_0$ multiplies
the signed transform by $e^{-v_0z}$ and preserves all observed
zeros and equalities. The tail condition is unchanged by a fixed
translation.
\end{proof}

As a simpler consequence, every tail condition
$\mathbb E e^{c\Sigma^\beta}<\infty$ with $\beta<1/2$ fails to
ensure identification from an arbitrary infinite terminal set.
The preceding witnesses satisfy this condition for every $c>0$.
At $\beta=1/2$, identification holds; it even holds after weakening
the exponent to $\sqrt\Sigma/\log(e+\Sigma)$.

\section{Transfer to random clocks of symmetric L\'evy noise}
\label{p04:nxg:sec:levy}

The joint realization is not specific to Gaussian kernels. It gives
an exact sparse-time observation criterion for a class of random-clock
models. Let $(L_t)_{t\ge0}$ be a known symmetric L\'evy process in
$\mathbb R^D$ with
\[
 \mathbb E e^{iu^TL_t}=e^{-t\psi(u)},
\]
where $\psi$ is continuous, even, strictly positive away from zero,
and $C^\infty$ away from zero. These assumptions include Brownian
motion and the isotropic symmetric stable symbols
$\psi(u)=\|u\|^\alpha$, $0<\alpha\le2$.
The latent pair $(A,V)$ has an arbitrary Borel law on
$\mathbb R^D\times[0,\infty)$, independent of the driving process.
Only the separate marginal laws of
\[
 Y_t=A+L_{tV}
\]
are observed. Dependence between the latent position $A$ and clock
rate $V$ is unrestricted. The given data are not joint trajectory laws.

\begin{theorem}[An exact sparse-time criterion and its tail refinement]
\label{p04:nxg:thm:levy}
Let $t_j>0$ be distinct and decrease to zero.
\begin{enumerate}
\item The collection $\{\mathcal L(Y_{t_j}):j\ge1\}$ identifies
every Borel law of $(A,V)$ if and only if
\[
 \sum_jt_j=\infty.
\]
When the sum converges, distinct absolutely continuous mixing laws
with identical full $A$ and $V$ marginals and identical finite joint
polynomial moments produce the same observations.
\item For every sequence $t_j\downarrow0$, identification holds
within each class $\mathbb E e^{c\omega_a(V)}<\infty$ for some
$c>0$, with $0\le a\le1$.
\item For the geometric times $t_j=2^{-j}$ and every $a>1$, the
counterexamples in part~1 can satisfy
$\mathbb E e^{c\omega_a(V)}<\infty$ for every $c>0$.
This statement holds in every spatial dimension $D$ and for every
symbol $\psi$ just specified.
\end{enumerate}
\end{theorem}

\begin{proof}
For the difference $\rho$ of candidate laws and fixed $u\ne0$, set
\[
 H_u(z)=\int e^{iu^Ta-zv}\,d\rho(a,v),\qquad\operatorname{Re}z>0.
\]
It is bounded and holomorphic, with zeros $z=t_j\psi(u)$ if the
observations agree. If $\sum_jt_j=\infty$, these zeros violate
the right-half-plane Blaschke condition for a nonzero bounded
holomorphic function \cite{p03:mixnew:Garnett2007}. Hence $H_u$ vanishes.
Continuity gives $u=0$; evaluating at $z=1+i\tau$ and using
Fourier uniqueness for $e^{-v}\rho$ gives $\rho=0$.
Under the tail condition, apply Lemmas~\ref{p04:nxg:lem:moment}
and~\ref{p04:nxg:lem:zeros} to the pushforward
$\int e^{iu^Ta}\,d\rho(a,v)$ instead; this proves part~2.

For completeness, necessity in part~1 needs a scalar realization
for every summable sequence, not only geometric sequences. Suppose
$\sum_jt_j<\infty$ and set
\begin{equation}\label{p04:nxg:eq:F-general}
 F_0(z)=\left\{\prod_j\frac{z-t_j}{z+t_j}\right\}
               \exp\{-\sqrt z-z^{-1/2}\}.
\end{equation}
The product is nonzero, locally uniformly convergent in the right
half-plane, bounded by one, and conjugation invariant. Its boundary
derivatives near $iy=0$ have at most polynomial growth in
$1/|y|$. Indeed, for each factor $b_j(z)=(z-t_j)/(z+t_j)$,
\[
 |b_j(iy)|=1,\qquad
 |b_j^{(k)}(iy)|\le 2k!t_j|y|^{-k-1},\quad k\ge1.
\]
The product rule and $\sum_jt_j<\infty$ bound each derivative
by a finite sum of powers of $1/|y|$. The same estimates, or local
uniform convergence away from zero, justify differentiated limits.
Thus the exponential in \eqref{p04:nxg:eq:F-general} makes its
boundary values Schwartz and flat at zero. Uniform integrable
right-tube bounds follow from
$|F_0(x+iy)|\le e^{-\sqrt{|y|/2}}$. The inverse Fourier and
line-shift argument used in Lemma~\ref{p04:nxg:lem:sector} produces
a nonzero real Schwartz inverse Laplace transform $f_0$, supported
on $[0,\infty)$, with every moment zero. No exponential tail is
claimed for this arbitrary summable design.

Use \eqref{p04:nxg:eq:joint-density} with $f_0$ and the present
symbol $\psi$. The proof of Lemma~\ref{p04:nxg:lem:joint}, without
the quantitative tail estimate, gives a real nonzero joint Schwartz
density with transform $h(u)F_0(z/\psi(u))$ and vanishing
marginals and moments. Its positive and negative parts, normalized
to probability laws, have the same observed characteristic functions,
because at $z=t_j\psi(u)$ this transform is zero. This proves
necessity and the supplementary assertions in part~1.

For part~3 choose $1<q<a$, use the sector construction with zeros
$b_j=2^{-j}$, and repeat the same joint realization. Estimate
\eqref{p04:nxg:eq:joint-tail} gives the claimed exponential integrability
exactly as in the proof of Theorem~\ref{p04:nxg:thm:log-threshold}.
\end{proof}

For Brownian noise this identifies a precise limitation of recovering
a population's joint starting-position and diffusivity distribution
from separate spatial snapshots. All coordinate marginals and all
polynomial mixed moments can agree while every selected sparse
snapshot agrees as well. The logarithmic tail boundary gives a
condition under which this ambiguity disappears. For isotropic stable
noise the same conclusion applies to a fractional diffusion clock.
These are exact-data identifiability statements; no finite-sample
estimation rate or stable numerical inversion is claimed.

\subsection{Priority and unresolved transfers}

The classical moment-determinacy criteria, the half-plane Blaschke
condition, and Fourier--Laplace uniqueness are credited as inputs.
The new claims here are their exact observation consequences together
with the joint-density converse, especially its critical logarithmic
tail estimate. The bounded literature check included the primary
moment-problem sources and grouped-mixture identification
\cite{p02:frontmix:VS2015}. It did not locate the sparse random-clock
criterion above, but does not certify historical priority. These
results settle questions formulated within the present program;
they are not described as resolutions of a located, previously
published conjecture.

The graph-theoretic comparison between maximum likelihood threshold
and generic completion rank remains outside these arguments. The
question of a graph-independent bound on generic completion rank in
terms of maximum likelihood threshold was posed by Blekherman and
Sinn \cite{p02:graph:BS2017}. Neither a moment-determinacy argument
nor an annihilating transform supplies the positive semidefinite
stress certificates needed there. We do not count this graph problem
as solved. For the matrix experiment, a converse with the precise
logarithmic tail bound in each fixed dimension $d>1$, and sharp
statistical stability at the tail boundary, also remain open here.

% ---- BEGIN zeros ----
\section{Every modular period of a rationally independent exponential sum}
\label{p04:nxz:sec:classification}

The three-term restriction in the preceding manuscript can be removed.
The key observation is that it is unnecessary to describe the entire
polygon configuration space: configurations in which the summands form
only two aligned groups already force every nonexceptional character to
cover the circle.  Only the largest frequency can produce a missing arc.
This answers the existence question in Podesta
\cite[Introduction]{p02:zp:Podesta2026} for arbitrary finite sums satisfying its
rational-independence hypothesis, and gives the stronger classification
of every period.  The existence part was already obtained by the
nonresonant torus argument in our preceding work; the new assertion is
the complete classification for any number of terms.  Polygon criteria
and Kronecker--Rouche arguments for real projections are classical; see
Sepulcre \cite[Theorems 5--6]{p04:nxz:Sepulcre}.  Historical priority of the
modular classification beyond these checked sources is not asserted.

Let $m\ge2$, and consider
\begin{equation}\label{p04:nxz:eq:model}
 f(s)=-1+\sum_{j=1}^m B_j e^{\lambda_j s+i\phi_j},\qquad
 B_j>0,\quad 0<\lambda_1<\cdots<\lambda_m,
\end{equation}
where $\lambda_1,\ldots,\lambda_m$ are linearly independent over
$\mathbb Q$.  Write $\mathbb T=\mathbb R/(2\pi\mathbb Z)$ and
\[
 \mathcal R_L(f)=\overline{\{(2\pi/L)\operatorname{Im}\rho\bmod2\pi:
                    f(\rho)=0,\ \operatorname{Im}\rho>0\}}\subset\mathbb T.
\]
Set $r_j=\lambda_j/\lambda_m$, $\kappa_j=B_jB_m^{-r_j}$ for $j<m$,
and define
\begin{equation}\label{p04:nxz:eq:theta}
 S(x)=\sum_{j<m}\kappa_jx^{r_j},\qquad
 u(x)=\frac{1+x^2-S(x)^2}{2x},\qquad
 \Theta=\arccos\!\left(\max\{-1,\min_{x>0}u(x)\}\right).
\end{equation}
The minimum exists, and $0<\Theta\le\pi$.

\begin{theorem}[Only the largest frequency can obstruct modular density]
\label{p04:nxz:thm:all-periods}
For \eqref{p04:nxz:eq:model}, every $L>0$ has
$\mathcal R_L(f)=\mathbb T$, except possibly
\[
                    L=\frac{2\pi d}{n\lambda_m},
       \qquad n,d\in\mathbb N,\quad\gcd(n,d)=1.
\]
At such a period the exact residue closure is
\begin{equation}\label{p04:nxz:eq:arcs}
 \mathcal R_L(f)=
 \{z\in\mathbb T:dz+n\phi_m\in[-n\Theta,n\Theta]\pmod{2\pi}\}.
\end{equation}
Consequently the complete exceptional set is
\begin{equation}\label{p04:nxz:eq:exceptions}
 \left\{\frac{2\pi d}{n\lambda_m}:
          n,d\in\mathbb N,\ \gcd(n,d)=1,\ n\Theta<\pi\right\}.
\end{equation}
In particular, only finitely many integers $n$ can occur.  The
classification and every residue closure are independent of
$\phi_1,\ldots,\phi_{m-1}$.  If $n\Theta<\pi$, the closure in
\eqref{p04:nxz:eq:arcs} consists of $d$ closed arcs, whose total relative
circle length is $n\Theta/\pi$.
\end{theorem}

\begin{lemma}[Realization of compatible polygon configurations]
\label{p04:nxz:lem:realization}
Put
\[
 \mathcal Z=\{(\sigma,\theta)\in\mathbb R\times\mathbb T^m:
              \textstyle\sum_j B_je^{\lambda_j\sigma+i\theta_j}=1\}.
\]
This is a nonempty compact set.  For $q=2\pi/L$, if
$q\notin\operatorname{span}_{\mathbb Q}\{\lambda_j\}$, then
$\mathcal R_L(f)=\mathbb T$.  Otherwise write
\[
 q=(k\cdot\lambda)/d,\qquad k\in\mathbb Z^m,\quad d\ge1,
          \quad\gcd(k_1,\ldots,k_m,d)=1.
\]
Then
\begin{equation}\label{p04:nxz:eq:characters}
 \mathcal R_L(f)=
 \{z\in\mathbb T:dz+k\cdot\phi\in
                 \{k\cdot\theta:(\sigma,\theta)\in\mathcal Z\}\}.
\end{equation}
Every residue on the right is approached by ordinates tending to
$+\infty$.  Multiple limiting zeros are allowed.
\end{lemma}
\begin{proof}
For sufficiently negative $\sigma$, the sum of all lengths is less
than $1$; for sufficiently positive $\sigma$, the last length exceeds
$1$ plus all the others.  Neither region contains a configuration.
This proves compactness.  The sum of the lengths is strictly increasing
from zero to infinity, so there is a unique $\sigma_0$ where its value
is $1$; the all-zero angle vector is then a configuration.

The orbit $t\mapsto(t\lambda,qt)$ has closure $\mathbb T^{m+1}$ in
the nonresonant case.  In the resonant case its character relations
are exactly the multiples of $(k,-d)$, so its closure is the subtorus
$\{(v,z):k\cdot v-dz=0\}$.  This is the character form of
Kronecker's theorem.  Removing any bounded initial part of the orbit
does not change its closure.  Given a configuration and a compatible
residue, choose $t_\nu\to+\infty$ such that
$t_\nu\lambda\to\theta-\phi$ and $qt_\nu\to z$ on the tori.
Locally uniformly,
\[
 f(s+it_\nu)\longrightarrow
             -1+\sum_jB_je^{\lambda_js+i\theta_j}=:f_*(s).
\]
The limit has a zero at $\sigma$ and is not identically zero, since its
constant term is $-1$ and all its exponents are positive.  Rouche's
theorem, on successively smaller circles around this isolated zero,
therefore supplies zeros $s_\nu\to\sigma$ along a subsequence.  The
zeros $s_\nu+it_\nu$ have the required residues.  Conversely, every
actual zero defines a configuration satisfying the character relation.
Compactness proves the converse containment after taking closures.
\end{proof}

\begin{lemma}[Two aligned groups exhaust every nonextreme character]
\label{p04:nxz:lem:groups}
If $k\in\mathbb Z^m\setminus\{0\}$ is not supported solely at the
index $m$, then
\[
                 \{k\cdot\theta:(\sigma,\theta)\in\mathcal Z\}
                         =\mathbb T.
\]
\end{lemma}
\begin{proof}
For a nonempty proper subset $A\subset\{1,\ldots,m\}$ set
\[
 U(\sigma)=\sum_{j\in A}B_je^{\lambda_j\sigma},\quad
 V(\sigma)=\sum_{j\notin A}B_je^{\lambda_j\sigma},\quad
 K_A=\sum_{j\in A}k_j,\quad K_{A^c}=\sum_{j\notin A}k_j.
\]
Whenever $|U-V|\le1\le U+V$, the aligned-group configurations are
obtained from
\[
 \alpha=\arccos\frac{1+U^2-V^2}{2U},\qquad
 \beta=-\arccos\frac{1+V^2-U^2}{2V},
\]
by setting $\theta_j=\alpha$ on $A$ and $\theta_j=\beta$ on $A^c$,
or by conjugating all angles.  Their character is
$K_A\alpha+K_{A^c}\beta$.

Suppose first that $k$ has at least two nonzero entries.  There is an
index $j$ for which $k_j\ne0$ and $k_j\ne\sum_i k_i$.  Indeed, if
the total is zero every nonzero entry works; if the total is nonzero
and every nonzero entry equals it, there can only be one nonzero entry.
Take $A=\{j\}$.  Both $K_A$ and $K_{A^c}$ are nonzero.  Starting at
the unique $\sigma_0$ with $U+V=1$, let $\sigma_1>\sigma_0$ be the
first point where $|U-V|=1$.  It exists because the group containing
the largest exponent eventually dominates.  On $[\sigma_0,\sigma_1]$
the angles are continuous, start at $(0,0)$, and end at either
$(\pi,0)$ or $(0,-\pi)$.  The character therefore joins zero to a
nonzero integral multiple of $\pi$.  Its image and the conjugate image
together cover the circle.

It remains to consider $k$ supported at one $j<m$.  Again take
$A=\{j\}$; now $V-U\to+\infty$.  Let $\sigma_1>\sigma_0$ be the
first point where $V-U=1$.  If there is a point in
$[\sigma_0,\sigma_1)$ where $U-V=1$, let $\sigma_2$ be the last such
point; otherwise put $\sigma_2=\sigma_0$.  The triangle inequalities
hold on $[\sigma_2,\sigma_1]$, and $\alpha$ runs continuously from
$0$ to $\pi$.  The character $k_j\alpha$, together with its negative,
again covers the circle.  This proves the lemma.
\end{proof}

\begin{lemma}[The exact range of the largest angle]
\label{p04:nxz:lem:largest}
The projection of $\mathcal Z$ onto $\theta_m$ is exactly
$[-\Theta,\Theta]\pmod{2\pi}$, with \eqref{p04:nxz:eq:theta}.
\end{lemma}
\begin{proof}
Put $x=B_me^{\lambda_m\sigma}$.  The other lengths sum to $S(x)$.
A necessary condition for a prescribed angle $\eta=\theta_m$ is
\begin{equation}\label{p04:nxz:eq:triangle-bound}
                 |1-xe^{i\eta}|\le S(x)
                 \quad\text{for some }x>0.
\end{equation}
It is also sufficient after allowing the real part to vary.  The
continuous difference $|1-xe^{i\eta}|-S(x)$ tends to $1$ as $x\to0$
and to $+\infty$ as $x\to\infty$, since $0<r_j<1$.  If it is
nonpositive anywhere, it vanishes somewhere.  At such a point all
the smaller summands can be aligned with the nonzero vector
$1-xe^{i\eta}$, giving the required configuration.  Thus no lower
polygon-radius condition has been discarded.

Squaring \eqref{p04:nxz:eq:triangle-bound} gives
$\cos\eta\ge u(x)$.  The function $u$ tends to $+\infty$ at both
ends and $u(1)=1-S(1)^2/2<1$.  The stated arc and $0<\Theta\le\pi$
follow.
\end{proof}

\begin{proof}[Proof of Theorem~\ref{p04:nxz:thm:all-periods}]
Apply Lemma~\ref{p04:nxz:lem:realization}.  A resonant character not
supported at $m$ covers the circle by Lemma~\ref{p04:nxz:lem:groups}.
A character supported at $m$ has $k_m=n>0$, since $q>0$, and
$\gcd(n,d)=1$.  Lemma~\ref{p04:nxz:lem:largest} then gives
\eqref{p04:nxz:eq:arcs}.  The interval $[-n\Theta,n\Theta]$ covers the
circle exactly when $n\Theta\ge\pi$.  If it is shorter, its inverse
image under multiplication by $d$ has $d$ components of equal length
$2n\Theta/d$.  All assertions follow.
\end{proof}

\begin{corollary}[Each resonant decision is a strictly concave scalar problem]
\label{p04:nxz:cor:scalar}
Let $n\ge1$, $c_n=\cos(\pi/n)$ and
\[
 R_n(x)=\sqrt{1+x^2-2c_nx},\qquad D_n(x)=S(x)-R_n(x).
\]
The function $D_n$ is strictly concave on $(0,\infty)$ and has a unique
maximizer $x_n$, characterized by
\begin{equation}\label{p04:nxz:eq:scalar}
 \sum_{j<m}r_j\kappa_jx_n^{r_j-1}
       =\frac{x_n-c_n}{\sqrt{1+x_n^2-2c_nx_n}}.
\end{equation}
Every period $2\pi d/(n\lambda_m)$ in lowest terms is dense if and
only if $D_n(x_n)\ge0$.  In particular, all periods are dense if and
only if
\begin{equation}\label{p04:nxz:eq:all-dense}
                \max_{x>0}\{S(x)-x\}\ge1.
\end{equation}
The unique maximizer in \eqref{p04:nxz:eq:all-dense} solves $S'(x)=1$.
\end{corollary}
\begin{proof}
The criterion $n\Theta\ge\pi$ is equivalent to
$S(x)\ge R_n(x)$ for some $x$.  Because $S''<0$ and $R_n$ is convex,
$D_n''<0$; for $n=1$, $R_1(x)=1+x$ and the same conclusion holds.
Moreover $D_n'(x)\to+\infty$ as $x\to0$ and
$D_n'(x)\to-1$ as $x\to\infty$.  Hence its unique critical point is
the unique maximum, giving \eqref{p04:nxz:eq:scalar}.  The case $n=1$
proves \eqref{p04:nxz:eq:all-dense} and the last assertion.
\end{proof}

\subsection{An infinite spectrum without a largest frequency has no exceptional period}
\label{p04:nxz:subsec:infinite}

\begin{theorem}[Universal modular density for normally convergent infinite sums]
\label{p04:nxz:thm:infinite}
Let $\{\lambda_j:j\ge1\}\subset(0,\infty)$ be linearly independent
over $\mathbb Q$ and have no largest member.  Suppose $B_j>0$ and
\begin{equation}\label{p04:nxz:eq:normal}
                    \sum_{j\ge1}B_j e^{R\lambda_j}<\infty
                         \qquad\text{for every }R>0.
\end{equation}
For arbitrary real phases $\phi_j$, the entire function
\[
                   f(s)=-1+\sum_{j\ge1}B_je^{\lambda_js+i\phi_j}
\]
has positive zero ordinates dense modulo every $L>0$.  In particular,
this holds for the smoothed prime sum
\begin{equation}\label{p04:nxz:eq:prime-example}
                       -1+\sum_{p\ \mathrm{prime}}e^{-p}p^s.
\end{equation}
\end{theorem}
\begin{proof}
Condition \eqref{p04:nxz:eq:normal} gives locally uniform convergence,
uniformly over all choices of coefficient phases.  The positive length
sum $M(\sigma)=\sum_jB_je^{\lambda_j\sigma}$ is continuous and
strictly increasing, tends to zero at $-\infty$ by dominated
convergence, and tends to infinity at $+\infty$.  Hence there is again
a unique $\sigma_0$ with $M(\sigma_0)=1$.

The continuous-time orbit on the countable phase torus, augmented by
$qt$ for $q=2\pi/L$, has the same character description as in
Lemma~\ref{p04:nxz:lem:realization}: its continuous characters involve
only finitely many coordinates.  It is either the entire product
torus, or the subtorus with the single primitive relation
$k\cdot v-dz=0$, where $k$ has finite support and
$q=k\cdot\lambda/d$.  This follows by applying the finite-dimensional
Kronecker theorem to each finite collection of coordinates and then
using the product topology.  Every point can be approached at times
tending to $+\infty$.  Coordinatewise phase convergence implies local
uniform convergence of the entire functions, by the uniform tail
bound in \eqref{p04:nxz:eq:normal}.  The same Rouche argument therefore
realizes each compatible polygon configuration.  The limiting function
is not identically zero, since it tends to $-1$ on the negative real
axis.

The nonresonant case follows at once from the configuration at
$\sigma_0$.  In the resonant case the proof of
Lemma~\ref{p04:nxz:lem:groups} still applies.  When $k$ has at least two
nonzero entries choose its singleton group exactly as there.  When
it has only one, choose that singleton group.  Its length $U$ is one
exponential, whereas the complementary group $V$ contains a term
with a strictly larger frequency, so $V-U\to+\infty$.  The two
continuous triangle-angle paths used in that lemma are therefore
available in both cases, and every nonzero finite character covers
the circle.  This proves density for every period.  Finally,
$\{\log p\}$ is rationally independent by unique factorization, has
no largest element, and $\sum_p e^{-p}p^R<\infty$ for every $R$.
Thus \eqref{p04:nxz:eq:prime-example} satisfies the hypotheses.
\end{proof}

This theorem removes the finite-sum exceptional direction by removing
the largest frequency itself.  Its normal-convergence assumption is
essential to the argument; it does not cover an arbitrary Dirichlet
series converging only on a half-plane.  In particular the smoothed
prime example is not the Riemann zeta function or an Euler product.

\section{Recovering an unknown spectrum from density transitions}
\label{p04:nxz:sec:inverse}

The modular classification permits a different inverse problem from
phase recovery at an Euler edge.  Here the relative frequencies are
unknown, while an external amplitude can be varied.  The observation
is only the transition between dense and nondense modular ordinates.
The following problem is formulated here; it is not presented as a
second previously published conjecture.

Fix $b>0$ and distinct unknown numbers
$0<r_1<\cdots<r_{m-1}<1$ such that $br_1,\ldots,br_{m-1},b$ are
rationally independent.  For $A,C>0$ consider
\begin{equation}\label{p04:nxz:eq:family}
 f_{A,C}(s)=-1+A\sum_{j<m}B_j e^{br_js+i\phi_j}
                         +C e^{bs+i\phi_m},\qquad B_j>0.
\end{equation}
The normalizations of $A,C$ and the constant coefficient are part of
the model.  Observe, for each $C$, the smallest common amplitude $A$
at which every period has dense ordinates:
\[
 T(C)=\inf\{A>0:\mathcal R_L(f_{A,C})=\mathbb T
                              \text{ for every }L>0\}.
\]
Let $S_0(x)=\sum_{j<m}B_jx^{r_j}$.

\begin{theorem}[Exact spectral recovery from a binary transition curve]
\label{p04:nxz:thm:tomography}
The transition is attained, and
\begin{equation}\label{p04:nxz:eq:transition}
 T(C)=\min_{x>0}\frac{1+Cx}{S_0(x)},\qquad
 \mathcal R_L(f_{A,C})=\mathbb T\ \text{for every }L>0
                 \ \Longleftrightarrow\ A\ge T(C).
\end{equation}
Its entire curve determines the positive power sum by the exact inverse
\begin{equation}\label{p04:nxz:eq:inverse}
                         S_0(x)=\inf_{C>0}\frac{1+Cx}{T(C)}.
\end{equation}
It therefore determines the number of smaller frequencies, all ratios
$r_j$, and all magnitudes $B_j$, up to relabeling.  The curve is smooth,
strictly increasing and strictly concave.  If $x(C)$ is the unique
minimizer in \eqref{p04:nxz:eq:transition}, then
\begin{equation}\label{p04:nxz:eq:parametric}
 x(C)=\frac{T'(C)}{T(C)-CT'(C)},\qquad
 S_0(x(C))=\frac1{T(C)-CT'(C)}.
\end{equation}
As $C$ ranges over $(0,\infty)$, $x(C)$ ranges bijectively over
$(0,\infty)$ in the opposite direction.
\end{theorem}
\begin{proof}
The all-period criterion says that $AS_0(x)\ge1+Cx$ for some $x$.
The quotient in \eqref{p04:nxz:eq:transition} tends to infinity at both
ends, so its infimum is attained.  This proves the transition formula.
For every $C,x>0$ it immediately gives
\[
                         S_0(x)\le\frac{1+Cx}{T(C)}.
\]
To obtain equality at a prescribed $x$, put
\[
 H(x)=S_0(x)-xS_0'(x)=\sum_{j<m}(1-r_j)B_jx^{r_j}>0,
              \qquad C_x=\frac{S_0'(x)}{H(x)}.
\]
Strict concavity gives, for every $y>0$,
\[
 S_0(y)\le S_0(x)+S_0'(x)(y-x)
                  =H(x)(1+C_xy),
\]
with equality only at $y=x$.  It follows that
$T(C_x)=1/H(x)$ and that $x$ is its unique minimizing point.  This
proves \eqref{p04:nxz:eq:inverse}.  Direct differentiation yields
\[
             C_x'=\frac{S_0''(x)S_0(x)}{H(x)^2}<0.
\]
The smallest and largest power terms show respectively that
$C_x\sim r_1/((1-r_1)x)$ as $x\to0$ and
$C_x\sim r_{m-1}/((1-r_{m-1})x)$ as $x\to\infty$.
Thus $x\mapsto C_x$ is a smooth decreasing bijection.  The envelope
identity gives
\[
             T'(C)=\frac{x(C)}{S_0(x(C))}>0,
 \qquad T(C)-CT'(C)=\frac1{S_0(x(C))}>0.
\]
Also the derivative of $x/S_0(x)$ with respect to $x$ is
$H(x)/S_0(x)^2>0$, while $x'(C)<0$, proving $T''(C)<0$ and
\eqref{p04:nxz:eq:parametric}.

Finally, a finite positive power sum has a unique representation with
distinct exponents.  Explicitly,
\[
 r_1=\lim_{x\downarrow0}\frac{\log S_0(x)}{\log x},\qquad
 B_1=\lim_{x\downarrow0}x^{-r_1}S_0(x).
\]
Subtract $B_1x^{r_1}$ and repeat until the remainder is zero.  This
establishes uniqueness of the term count, exponents and coefficients;
it makes no finite-data numerical claim about detecting the last term.
\end{proof}

The inverse is a supporting-line reconstruction of a strictly concave
function.  This convex-duality step is classical; the specific content
here is that the observable $T(C)$ is exactly the modular-zero density
transition, so it determines an unknown finite spectrum.

\begin{proposition}[Relative stability of the recovered power-sum profile]
\label{p04:nxz:prop:stability}
Suppose $0<\epsilon<1$ and a positive observed curve satisfies
$(1-\epsilon)T(C)\le\widetilde T(C)\le(1+\epsilon)T(C)$ for all
$C>0$.  Define
$\widetilde S(x)=\inf_{C>0}(1+Cx)/\widetilde T(C)$.  Then
\[
 \frac{S_0(x)}{1+\epsilon}\le\widetilde S(x)
                    \le\frac{S_0(x)}{1-\epsilon}\qquad(x>0).
\]
If the curve is observed only on $[C_-,C_+]$, the same bounds hold for
the infimum restricted to that interval at every $x$ with
$C_x\in[C_-,C_+]$.
\end{proposition}
\begin{proof}
Apply the pointwise multiplicative inequalities to
$(1+Cx)/\widetilde T(C)$ and take infima.  On a restricted interval
the true infimum is still $S_0(x)$ whenever its tangent value $C_x$
belongs to the observed interval.
\end{proof}

\begin{proposition}[An exact phase blind spot]
\label{p04:nxz:prop:blind}
Even knowing all closures $\mathcal R_L(f_{A,C})$, for every
$A,C,L>0$, cannot identify any of the phases
$\phi_1,\ldots,\phi_{m-1}$ when $B_j,r_j,b,\phi_m$ are fixed.
\end{proposition}
\begin{proof}
Theorem~\ref{p04:nxz:thm:all-periods} expresses every such closure in
terms of the magnitudes, frequencies and $\phi_m$ alone.  Arbitrarily
changing the remaining phases leaves all these data unchanged.
\end{proof}

This is a statement about closures, not equality of the zero divisors,
the empirical residue distributions, or zero correlations.  The stable
reconstruction above concerns $S_0$; extracting closely spaced
exponents or very small coefficients from noisy restricted data is a
separate inverse problem.  Neither theorem applies directly to an
infinite Euler product: it has no largest frequency, and its prime-power
frequencies are not rationally independent.  The actual Euler
hyperuniformity estimate and the exact-prime recurrence exponent remain
unproved in this branch.  The two completed targets here are the
finite-sum all-period classification and the newly formulated
transition-curve inverse problem; the scalar test and phase blind spot
are consequences rather than additional conjecture closures.

% ---- BEGIN inverse ----
\section{A common analytic envelope: exact information versus noisy information}
\label{p04:nxi:section}

This section addresses two problems left separate by the earlier analytic
instability construction.  Does a \emph{fixed} complex-neighborhood bound
yield a quantitative finite-sample recovery rate?  Does that same bound
yield a noise rate uniform over all sequences of sampling times?  The
answers are respectively yes and no.  Exact-data recovery has uniform
exponential order, while noise convergence can be arbitrarily slow along
one fixed sampling sequence, without changing the analytic class.
These are internally formulated inverse problems.  No previously named
open conjecture is claimed to be resolved here.

Fix $I=[a,b]$ with $0<a<b$, put $\ell=b-a$, and fix $\rho>1$.
Let $E_\rho$ be the open Bernstein ellipse with foci $a,b$ and parameter
$\rho$: its boundary is the image of $|w|=\rho$ under
\[
 w\longmapsto\frac{a+b}{2}+\frac{\ell}{4}(w+w^{-1}).
\]
Write $\|f\|_{E_\rho}=\sup_{z\in E_\rho}|f(z)|$.
For positive constants $m,M,\kappa$ let $\mathcal A_\rho$ be the class
of functions holomorphic in $E_\rho$, real on $I$, satisfying
\begin{equation}\label{p04:nxi:class}
 \int_I f=1,\qquad f\ge m,\qquad
 (\log f)''\le-\kappa\quad\hbox{on }I,\qquad
 \|f\|_{E_\rho}\le M.
\end{equation}
The functions are probability densities on $I$; no analyticity across
the endpoints is claimed for their zero extensions.  The common bound
in \eqref{p04:nxi:class}, unlike finitely many derivative bounds, controls
\emph{every} derivative uniformly on $I$.
Use inverse-time coordinates
\begin{equation}\label{p04:nxi:heat}
 H_s f=\sqrt{s/\pi}\int_I e^{-s x^2}f(x)\,dx,
 \qquad s>0,
\end{equation}
so that $H_s f$ is the heat observation at location zero and time
$t=1/(4s)$.

\subsection{Exponential exact-data recovery, uniformly over the times}

For distinct $s_1,\ldots,s_n>0$, let $e_n(s_1,\ldots,s_n)$ be the
smallest worst-case $L^1(I)$ error over all reconstructions from the
$n$ exact real numbers $H_{s_j}f$.  No computational restriction is
imposed on a reconstruction.

\begin{theorem}[Uniform exponential order]\label{p04:nxi:exact}
For every $a,b,\rho,\kappa$ as above, there exist $m,M,c>0$ and
$0<\vartheta<1$, depending only on these four parameters, such that
\begin{equation}\label{p04:nxi:exact-bounds}
 \frac{c\rho^{-(n+1)}}{(n+2)^2}
 \ \le\ e_n(s_1,\ldots,s_n)
 \ \le\ 2M\ell\,\vartheta^{n+1}
\end{equation}
for every $n\ge1$ and every collection of distinct positive sampling
parameters.  In particular $e_n=\exp(-\Theta(n))$, with constants
independent of the sampling parameters.

The lower bound also holds for every deterministic procedure using at
most $n$ adaptively selected real linear measurements on the unknown
density, even if those measurements are not heat observations and their
selection depends discontinuously on previous exact answers.
\end{theorem}

\begin{proof}
We give the upper and lower arguments separately.

\emph{Sign changes from the data.}
Suppose that $f,g\in\mathcal A_\rho$ have the same $n$ observations,
and write $h=f-g$.  Then $h$ is orthogonal in the ordinary integral
pairing to
\[
 1,e^{-s_1x^2},\ldots,e^{-s_nx^2}.
\]
This is an extended Chebyshev system on $I$.  To see this directly,
make the strictly increasing change of variable $y=x^2$ and order the
distinct exponents $0,-s_1,\ldots,-s_n$.  Any nonzero combination of
$r+1$ exponentials with distinct real exponents has at most $r$ zeros,
counted with multiplicity: divide by one exponential, differentiate,
and apply Rolle's theorem inductively.

If $h$ is nonzero and has at most $n$ sign changes, say $r$ changes at
$x_1<\cdots<x_r$, choose a nonzero combination $v$ of any $r+1$ of the
listed functions vanishing at those $r$ points.  The zero bound shows
that these zeros are simple and are its only zeros on $I$.  After
changing its overall sign, $h v\ge0$ everywhere and $h v>0$ on an open
subinterval.  This contradicts $\int_I hv=0$.  Consequently a nonzero
$h$ has at least $n+1$ distinct real zeros in $I$.

\emph{Zeros in a common analytic neighborhood.}
Let $\phi:E_\rho\to\mathbb D$ be a conformal map and set
\[
 \vartheta=\max_{x,y\in I}
 \left|\frac{\phi(x)-\phi(y)}{1-\overline{\phi(y)}\phi(x)}\right|.
\]
Compact containment gives $0<\vartheta<1$.
The bounded holomorphic function $h/(2M)$, transferred to the disk,
has the $n+1$ zeros just found.  Dividing by their finite Blaschke
product leaves a function bounded by one.  Equivalently, repeated
Schwarz factorization gives
\[
 |h(x)|\le 2M\prod_{k=1}^{n+1}
 \left|\frac{\phi(x)-\phi(x_k)}
 {1-\overline{\phi(x_k)}\phi(x)}\right|
 \le2M\vartheta^{n+1},\qquad x\in I.
\]
Integration bounds the diameter of every exact-data fiber.  Selecting
any member of a nonempty fiber proves the asserted upper bound.

\emph{An analytic perturbation ball.}
Normalize $f_0(x)=C\exp(-\kappa x^2)$ to have integral one on $I$.
Then $(\log f_0)''=-2\kappa$.  Cauchy estimates on disks of a fixed
radius around $I$ show that there is $\epsilon>0$ such that every real
holomorphic $h$ satisfying
\[
 \|h\|_{E_\rho}\le\epsilon,\qquad\int_I h=0
\]
has $f_0\pm h\ge m$ and $(\log(f_0\pm h))''\le-\kappa$, where
$m>0$ is fixed.  Indeed the $C^2(I)$ norm of $h$ is at most a fixed
multiple of $\epsilon$, and the strict curvature and positivity
reserves of $f_0$ persist under sufficiently small $C^2$ perturbations.
Choose $M>\|f_0\|_{E_\rho}+\epsilon$.  Thus every such perturbation
belongs to the same class \eqref{p04:nxi:class}.

\emph{A polynomial null direction.}
For fixed measurements, choose a nonzero real polynomial $p$ of degree
at most $d=n+1$ annihilated by all $n$ measurements and by integration.
There are at most $n+1$ homogeneous linear constraints on a space of
dimension $n+2$.  Normalize $\|p\|_{I,\infty}=1$.
The polynomial growth inequality
\begin{equation}\label{p04:nxi:polynomial-growth}
 \|p\|_{E_\rho}\le\rho^d
\end{equation}
follows from the maximum principle applied outside the unit disk to
$w^{-d}p((a+b)/2+\ell(w+w^{-1})/4)$.
Also
\begin{equation}\label{p04:nxi:polynomial-mass}
 \|p\|_{L^1(I)}\ge \frac{\ell}{(d+1)^2}.
\end{equation}
For a short proof, expand in the orthonormal shifted Legendre
polynomials.  Their reproducing kernel has diagonal at most
$(d+1)^2/\ell$, since $|P_k(x)|\le1$ on $[-1,1]$.
Hence $1\le(d+1)\|p\|_2/\sqrt\ell$; combining this with
$\|p\|_2^2\le\|p\|_\infty\|p\|_1=\|p\|_1$ proves the claim.

Set $h=\epsilon\rho^{-d}p$.  The pair $f_0\pm h$ is admissible,
has identical data, and has $L^1$ distance at least
$2\epsilon\ell\rho^{-(n+1)}/(n+2)^2$.  The triangle inequality
proves the lower bound with $c=\epsilon\ell$.

For adaptive linear measurements, run the procedure at $f_0$ first
and record the at most $n$ linear functionals selected along that
transcript.  Choose $p$ annihilating those functionals and integration.
An induction over the query index shows that $f_0+h$, $f_0-h$, and
$f_0$ generate exactly the same transcript and therefore the same
measurement choices.  This proof makes no continuity assumption on the
selection rule.
\end{proof}

\paragraph{Interpretation and boundary.}
The exponent constants in \eqref{p04:nxi:exact-bounds} are not matched;
``exponential order'' means upper and lower bounds $\exp(-c_1n)$ and
$\exp(-c_2n)$, not a sharp exponential constant.  The upper bound is
uniform even when observations are taken at extraordinarily small or
nearly identical positive times.  It uses exact real numbers.  It is
not a numerically stable reconstruction algorithm.  The lower bound
addresses linear adaptive information; it does not settle the problem
of arbitrary adaptive analytic functionals posed by
Krieg--Novak--Ullrich~\cite[Section~3]{p03:an:KNU}.

\subsection{A quantitative noise obstruction inside the same analytic class}

For an increasing sequence $\mathbf s=(s_j)_{j\ge1}$ tending to
infinity, define the deterministic minimax risk
\[
 R_{\mathbf s}^{\infty}(\delta)
 =\inf_{\widehat f}\sup_{f\in\mathcal A_\rho}
 \sup_{\|e\|_{\ell^\infty}\le\delta}
 \|\widehat f((H_{s_j}f+e_j)_{j\ge1})-f\|_1.
\]
Define $R_{\mathbf s}^{G}(\delta)$ similarly using expected loss and
independent $N(0,\delta^2)$ errors in all the coordinates.  The estimator
receives the entire infinite noisy sequence in both definitions.

\begin{theorem}[Tail-controlled minimax lower bounds]\label{p04:nxi:noisy}
Use the fixed class from Theorem~\ref{p04:nxi:exact}.  Suppose
$s_1\ge1$ and $s_{j+1}\ge s_j+1$.  There is a constant $C_a>0$
such that, for
\begin{equation}\label{p04:nxi:delta-n}
 \delta_n=C_a\sqrt{s_{n+1}+1}\,e^{-a^2s_{n+1}},
\end{equation}
and all sufficiently large $n$,
\begin{equation}\label{p04:nxi:noisy-lower}
 \min\{R_{\mathbf s}^{\infty}(\delta_n),
             R_{\mathbf s}^{G}(\delta_n)\}
 \ge \frac{c\rho^{-(n+1)}}{2(n+2)^2}.
\end{equation}
The constant $C_a$ can be chosen independently of the sequence and
$n$; $c$ is the constant of Theorem~\ref{p04:nxi:exact}.
\end{theorem}

\begin{proof}
Use the polynomial construction in the preceding proof to obtain
$f_n^\pm=f_0\pm h_n$ with identical first $n$ observations and
separation
\[
 D_n:=\|f_n^+-f_n^-\|_1
 \ge\frac{2c\rho^{-(n+1)}}{(n+2)^2}.
\]
Both are probability densities, so $D_n\le2$.  For $j>n$,
\[
 |H_{s_j}(f_n^+-f_n^-)|
 \le2\sqrt{s_j/\pi}\,e^{-a^2s_j}.
\]
Writing $S=s_{n+1}$, unit spacing and eventual monotonicity of
$u e^{-2a^2u}$ give
\[
 \sum_{j>n}s_j e^{-2a^2s_j}
 \le\sum_{k=0}^{\infty}(S+k)e^{-2a^2(S+k)}
 \le C'_a(S+1)e^{-2a^2S}.
\]
Increase $C_a$ in \eqref{p04:nxi:delta-n} so that the entire mean-vector
difference has $\ell^2$ norm at most $\delta_n$, and hence supremum
norm at most $\delta_n$.
For adversarial noise the midpoint of the two mean sequences is
compatible with both laws, giving risk at least $D_n/2$.
For Gaussian noise the two product Gaussian measures have relative
entropy at most $1/2$; their mean difference is in $\ell^2$, so the
usual finite-dimensional identity passes to the product by the
convergent Gaussian likelihood ratio.  Pinsker's inequality gives
total variation at most $1/2$.  Integration against the common part
of the measures gives risk at least $D_n/4$, which is the displayed
lower bound.
\end{proof}

\begin{corollary}[Geometric sampling rules out a H\"older noise modulus]
\label{p04:nxi:geometric}
If $s_j=q^j$ with $q>1$, then for all sufficiently small $\delta$,
\[
 \min\{R_{\mathbf s}^{\infty}(\delta),R_{\mathbf s}^{G}(\delta)\}
 \ge
 c_q\frac{[\log(1/\delta)]^{-\log\rho/\log q}}
 {[\log\log(1/\delta)]^2}.
\]
Consequently no positive power of $\delta$ bounds either minimax risk
on this class.  The displayed logarithmic exponent is a lower-bound
exponent; matching upper bounds are not asserted.
\end{corollary}
\begin{proof}
A finite initial segment may be discarded so that the spacing condition
holds.  The proof of Theorem~\ref{p04:nxi:noisy} itself is unchanged by
that initial segment, which remains among the annihilated coordinates.
For $L=\log(1/\delta)$ choose the smallest $n$ such that
$q^{n+1}\ge 2L/a^2$.  Then $\delta_n\le\delta$ for large $L$ and
$n+1\le\log_q L+C_{a,q}$.  Apply the two-point argument at noise
level $\delta$; its mean difference is at most $\delta$ in $\ell^2$.
Substitution in \eqref{p04:nxi:noisy-lower} proves the result.
\end{proof}

\begin{theorem}[Arbitrarily slow consistency with a fixed analytic envelope]
\label{p04:nxi:arbitrary}
Fix the class $\mathcal A_\rho$ of Theorem~\ref{p04:nxi:exact} once and
for all.  For every function $\omega:(0,1)\to(0,\infty)$ with
$\omega(\delta)\to0$ as $\delta\downarrow0$, there exists one sequence
$s_j\uparrow\infty$, with $s_{j+1}\ge2s_j$, such that
\[
 \limsup_{\delta\downarrow0}
 \frac{R_{\mathbf s}^{\infty}(\delta)}{\omega(\delta)}
 =\limsup_{\delta\downarrow0}
 \frac{R_{\mathbf s}^{G}(\delta)}{\omega(\delta)}=\infty.
\]
Nevertheless, for this same sequence,
\[
 R_{\mathbf s}^{\infty}(\delta)\longrightarrow0,
 \qquad R_{\mathbf s}^{G}(\delta)\longrightarrow0.
\]
In particular $t_j=1/(4s_j)$ is summable, the noiseless observation
map is injective, and there is no noise modulus uniform over all such
sampling sequences, even under a common full analytic envelope and
strong log-concavity.
\end{theorem}

\begin{proof}
Recursively choose $s_{n+1}\ge2s_n$ so large that
\[
 \omega\left(C_a\sqrt{s_{n+1}+1}e^{-a^2s_{n+1}}\right)
 \le\frac{c\rho^{-(n+1)}}{2n(n+2)^2}.
\]
This is possible because the argument tends to zero.  Theorem
\ref{p04:nxi:noisy} gives a risk ratio at least $n$ along the corresponding
sequence $\delta_n\downarrow0$, for both observation models.

It remains to prove the consistency assertions rather than infer them
from membership in an analytic class.  The class $\mathcal A_\rho$ is
compact in $C^2(I)$ by normal-family compactness and Cauchy estimates.
The constraints in \eqref{p04:nxi:class} are closed under that convergence;
a limiting holomorphic function on $E_\rho$ still has bound $M$.
In particular the class is compact in $L^1(I)$.
The complete map $f\mapsto(H_{s_j}f)_{j\ge1}$ is injective, because
Theorem~\ref{p04:nxi:exact} bounds the diameter of every first-$n$ fiber
by a quantity tending to zero.

For any $\eta>0$, compactness and this injectivity give a finite $J$
and $r>0$ such that
\begin{equation}\label{p04:nxi:finite-separation}
 \max_{1\le j\le J}|H_{s_j}f-H_{s_j}g|\le r
 \quad\Longrightarrow\quad\|f-g\|_1<\eta.
\end{equation}
Otherwise pairs separated by at least $\eta$ and converging in their
first $J$ data coordinates for every growing $J$ would have convergent
subsequences with different limits and identical full data, a
contradiction.
For deterministic noise, choose any class member fitting the first
$J$ observations within $\delta$; if $2\delta\le r$, its error is
less than $\eta$.
For Gaussian noise, minimize the maximum residual over the first $J$
coordinates on the compact class (a measurable approximate minimizer
also suffices).  On the event $\max_{j\le J}|Z_j|\le r/(3\delta)$,
its data discrepancy from the truth is at most $r$ if the optimization
error is at most $r/3$.  The event's complement has probability tending
to zero as $\delta\downarrow0$.  The loss is at most two, since both
outputs and truths are densities, so the expected risk tends to at
most $\eta$.  Let $\eta\downarrow0$.
\end{proof}

\paragraph{What has changed from the earlier barrier.}
Previously a positive noise-risk floor was obtained by allowing the
full analytic envelope to vary.  Here the entire class, the complex
neighborhood, its norm bound, the density lower bound, and its
logarithmic curvature are fixed.  Every fixed sampling sequence is
consistent, yet a sequence can force convergence slower than any
prescribed candidate rate.  At the same time \emph{every} sequence
has exponential exact-data recovery in the number of distinct samples.
The distinction is between exact real-number information and accuracy
of that information, rather than a failure of analyticity or
identifiability.

\paragraph{Unresolved directions and prior-art boundary.}
The arguments combine classical Chebyshev-system oscillation,
Schwarz--Blaschke factorization, polynomial growth, and two-point
statistical lower bounds.  Priority for this particular combined
statement has not been established.  The sharp noise rate even for
$s_j=q^j$ remains open in this work, as does the optimal exponential
constant in Theorem~\ref{p04:nxi:exact}.  A second whole exponentially
weighted diffraction curve is not counted as a new discovery:
parallel-line phase-retrieval theory already supplies closely related
uniqueness mechanisms; see Wellershoff~\cite{p04:nxi:Wellershoff}.
The adaptive analytic-functional question of
\cite[Section~3]{p03:an:KNU} remains unresolved: our adaptive lower bound
uses linearity essentially, and does not extend by merely replacing
``linear'' with ``analytic''.

\section{An eighteen-target register and the remaining barriers}
\label{p04:nxo:register}

Here ``complete'' means that the precisely scoped statement has a proof
in this manuscript. It does not mean that every broader problem bearing
the same name has been resolved. ``Partial'' records a proved range or
one-sided bound. ``Open'' records a transfer for which the attempted
mechanism does not supply a proof. Supporting lemmas are not separate
targets.

\begingroup
\small
\setlength{\tabcolsep}{4pt}
\renewcommand{\arraystretch}{1.15}
\begin{longtable}{@{}p{0.035\textwidth}p{0.245\textwidth}p{0.105\textwidth}p{0.555\textwidth}@{}}
\toprule
No. & Target & Status & Exact outcome and provenance \\
\midrule
\endfirsthead
\toprule
No. & Target & Status & Exact outcome and provenance \\
\midrule
\endhead
\bottomrule
\endfoot
1 & Near-linear binomial-mixture MLE & Complete &
Theorem~\ref{p04:nxm:thm:main}; all mixing laws and all MLE solutions, for
every fixed exponent below one. Reaches the published conditional range
in \cite[Section~6.2]{p01:info:MMRS2025}. \\
2 & Simultaneous bounded-test fingerprint control & Complete, negative &
Proposition~\ref{p04:nxm:prop:TVobstruction}; the uniform-over-tests
strengthening fails. This is not a counterexample to the original
fixed deterministic test conjecture. \\
3 & Predetermined Bernoulli replication caps & Complete &
Theorem~\ref{p04:nxb:thm:fixed-caps}; matching upper and lower bounds for
all $1\le Q\le B$. Internal minimax problem, arbitrary latent laws. \\
4 & Fully adaptive Bernoulli caps & Partial &
Theorem~\ref{p04:nxb:thm:adaptive-caps}; exact for
$Q\le B^{2/3-\gamma}$, but not throughout the shrinking critical
window near $B^{2/3}$. \\
5 & Uniform-endpoint critical budget rate & Complete &
Theorem~\ref{p04:nxb:thm:uniform}; the logarithm is necessary even with
fully adaptive allocation. Closes an internal critical-rate gap. \\
6 & Maximal ID domain of the third McKay family & Complete &
Theorem~\ref{p04:nxs:full-id}; the entire admissible domain. Directly
addresses the question in \cite[Section~2.3]{p01:rare:Baricz2026}. \\
7 & Full real-order fourth McKay ID domain & Complete &
Theorem~\ref{p04:nxs:fourth-full}; every $\mu>-1/4$. An extension of the
published family and theorem, not separately labeled as a named conjecture. \\
8 & Universal strictness of the SD thresholds & Partial &
Exact spectral criteria and tail bounds are proved in
Corollaries~\ref{p04:nxs:sd-general}, \ref{p04:nxs:fourth-sd} and
Proposition~\ref{p04:nxs:tail}. Strictness above one in all negative-cut
intervals is not established. \\
9 & Logarithmic covariance-tail boundary & Complete &
Theorem~\ref{p04:nxg:thm:log-threshold}; sufficiency in all dimensions,
sharpness for the scalar dyadic experiment, exactly within the stated
log-power scale. Internal observation problem. \\
10 & Sharp converse with matrix covariance in fixed $d>1$ & Open &
The scalar construction does not preserve the extra anchored observations.
No matrix-dimensional converse is claimed. \\
11 & Sparse random-clock L\'evy snapshots & Complete &
Theorem~\ref{p04:nxg:thm:levy}; the sum-of-times criterion, and the
logarithmic tail boundary for the stated symbols and scalar clocks.
Internal identification problem. \\
12 & All modular periods of exponential sums & Complete &
Theorems~\ref{p04:nxz:thm:all-periods} and~\ref{p04:nxz:thm:infinite}; complete
finite-sum classification and a normally convergent infinite extension.
The earlier existence answer to \cite{p02:zp:Podesta2026} is not counted again. \\
13 & Spectrum from a modular-density transition & Complete &
Theorem~\ref{p04:nxz:thm:tomography}; exact continuum data recover the
finite positive power sum. A newly formulated inverse problem;
lower-frequency phases remain invisible. \\
14 & Transfer to Euler-product fluctuation and prime recurrence laws & Open &
Prime-power relations and half-plane convergence invalidate the
independent-frequency transfer. Neither Euler hyperuniformity nor an
exact-prime recurrence exponent follows. \\
15 & Exact heat information under one analytic envelope & Complete &
Theorem~\ref{p04:nxi:exact}; uniform exponential order, with a lower bound
also for deterministic adaptive linear measurements. The two
exponential constants need not match. \\
16 & A sampling-uniform analytic noise modulus & Complete, negative &
Theorem~\ref{p04:nxi:arbitrary}; one fixed analytic class admits no common
rate over all sampling sequences, while each fixed sequence remains
consistent. Internal minimax problem. \\
17 & Sharp noisy rate for geometric heat times & Partial &
Corollary~\ref{p04:nxi:geometric}; a logarithmic lower bound rules out
H\"older stability. A matching sharp upper bound is not proved. \\
18 & Nonlinear continuous-measurement transfer & Open &
The adaptive linear annihilator does not establish the distinct
continuous nonlinear measurement claims studied in \cite{p03:an:KNU}.
No theorem in that paper is strengthened here by assertion. \\
\end{longtable}
\endgroup

\subsection{Three concrete next obstructions}

\paragraph{The fully adaptive cap window.}
At $Q=B^{2/3}/\log B$, the predetermined rate is
$B^{-1/3}\sqrt{\log B/\log\log B}$, while the available adaptive
lower bound is only $cB^{-1/3}$. The fixed-design lower proof charges
whole groups using local moment matching. A fully adaptive policy may
revisit individuals selected by earlier outcomes. An exact extension
would require a conditional information charge stable under this
selection, or an explicit adaptive construction beating the
predetermined rate. The endpoint experiment supplies such a charge
for a different likelihood; it does not automatically supply it for
Bernoulli populations.

\paragraph{Matrix tail sharpness.}
The scalar logarithmic counterexample annihilates a family of
one-dimensional Laplace arguments. Anchored matrix observations probe
additional positive-semidefinite directions. The scalar embedding
therefore fails as a proof, even though the sufficiency argument
extends to arbitrary dimension. A sharp converse must annihilate the
whole anchored transform set while preserving its quantitative tail.

\paragraph{Arithmetic transfer.}
In the finite-sum theorem, a primitive character relation and a largest
frequency produce the complete list of missing arcs. In a genuine
Euler product, prime-power frequencies have exact multiplicative
relations and the relevant series need not converge normally on the
whole plane. Removing the largest frequency from an absolutely
convergent toy model does not solve this arithmetic obstacle. No
Riemann-hypothesis implication or recovery of all prime data from one
right-hand zero is established by this manuscript.

\subsection{Other distinctions that survive the new results}

The MLE proof does not establish the exact constant-one
Chebyshev--Bernstein coefficient estimate, and the negative
simultaneous-test result leaves the original fixed-test fingerprint
question intact. The positive mixture transport transfer is the route
around those assertions, not a proof of them.

An exact SD threshold written as $\sup K'/K$ is more informative once
the kernel and a certified tail are explicit, but it is not an
elementary formula for every maximizing point. Likewise, relative
stability of an entire recovered power-sum profile does not imply
stable extraction of colliding frequencies. Exponential recovery
from exact real numbers does not imply a comparable noisy rate.
These separations are part of the mathematical conclusions.

\paragraph{Verification and use of computation.}
The asserted theorems are analytic arguments. Numerical checks were
not substituted for positivity, uniqueness, lower bounds, or
quantifiers over adaptive policies. The drafting process included
separate reviews of the likelihood transport and approximation steps,
the adaptive endpoint information calculation, the signed
hypergeometric representations, the logarithmic inverse-Laplace
construction, and the infinite-torus limit. Such reviews can miss
errors and do not replace external refereeing. This document is a
standalone single-source \LaTeX{} manuscript; it requires no external
data files, figures, or bibliography database to compile.
\endgroup
% ===== END SOURCE 04 =====


% ===== BEGIN SOURCE 05: KEIO_Adaptive_Analytic_Information_and_Recovery.tex =====
\clearpage
\part{Adaptive Analytic Information and Sharp Recovery Boundaries}
\label{part:05}
\setcounter{section}{0}
\begin{center}\large\itshape Polynomial Measurements, Complete Adaptive Budgets, and Nonlinear Spectral Reconstruction\end{center}
\begin{center}\small Source manuscript: 14 October 2026.\end{center}
\begingroup
\renewcommand{\arraystretch}{1}
\begin{quote}\small\textbf{Source abstract.}\quad
We continue a program on observation and recovery by changing the admissible
measurement geometry. Every vector in $\mathbb R^m$ can be recovered to any prescribed accuracy
using at most $\min\{m,\lceil\log_2m\rceil+2\}$ adaptive measurements,
each a globally convex polynomial. The same count is possible with
globally convex, real-analytic, one-Lipschitz measurements. This gives a positive complexity answer to the analytic and
convex information directions raised by Krieg, Novak, and Ullrich;
it does not determine the optimal number of queries. For populations of
Bernoulli parameters and uniform endpoints, we determine complete minimax
replication-cap curves, including arbitrary adaptive policies. A global
background in the Bernoulli hard family controls information after every
possible observation history. Further results give the sharp deterministic
binomial inverse modulus, certified approximate likelihood at statistical
precision, the exact view threshold for random low-rank Gaussian covariance,
and all finite iterated-log refinements of a scalar tail boundary.
For McKay laws we classify which powers restore hyperbolic complete
monotonicity, with a first-Bessel-zero asymptotic at the critical power.
Any $2m$ prescribed transition thresholds recover an $m$-atomic positive
spectrum, with sharp collision-stable noise exponent $1/(2m-1)$.
Finally, lacunary heat observations obey an exact logarithmic-risk
information-depth law. A target register distinguishes complete internal
classifications, positive answers to external questions, partial advances,
and unresolved arithmetic transfers.
\end{quote}


\section{The new questions and the change of viewpoint}
\label{p05:fo:intro}

The 223-page manuscript \emph{Observation and Recovery in Arithmetic and
Probability}, its three intermediate sequels, and the preceding
\emph{Transport, Spectral Positivity, and Sharp Observation Boundaries}
form the starting point of this paper
\cite{p01:KEIOIntegrated,comp:part:01,comp:part:02,comp:part:03,comp:part:04}.
The present text is a new continuation. It does not count results already
proved in those manuscripts as new closures, and it does not treat
unreviewed earlier claims as established literature. Its principal
statements are supplied with proofs here. The signed spectral kernel in
the final SD extension is an explicitly identified dependency on the
preceding paper.

There are three useful changes of viewpoint. First, a measurement rule
may admit a robust discrete implementation even when its original exact
membership test cannot be analytic. This allows approximation of the
measurement functions themselves. Second, in an adaptive statistical
experiment the necessary lower-bound object is conditional information
at the next query, including rare histories. An unconditional comparison
of whole individuals is insufficient. Third, nonlinear threshold data
can sometimes be converted into separated samples by contact geometry,
then into moments by a compact differential-equation argument.
These mechanisms transfer between problems, but no single theorem is
claimed to encompass all the branches below.

\subsection{A direct answer in the analytic-information direction}

Krieg, Novak, and Ullrich ask for complexity results when continuous
information is restricted to analytic, convex, or homogeneous mappings
\cite[Section~3]{p03:an:KNU}. We prove an $O(\log m)$ upper bound with
polynomial queries, and also with globally convex polynomial queries,
for arbitrary prescribed reconstruction accuracy. The convex-polynomial
bound on all of $\mathbb R^m$ is
\[
                         \min\{m,\lceil\log_2m\rceil+2\}.
\]
For example, it is twelve queries in dimension one thousand.
The statement counts real-valued oracle queries. The degree, coefficients,
precision cost, and computational work of such a query are not bounded
by this count. A separate construction permits every query to be globally convex,
real analytic, and Euclidean one-Lipschitz, at the same query count.
The optimal minimum remains open; the result establishes
that analytic regularity and global convexity alone do not remove the
exponential separation from linear information.

The crucial proof step is to retain proximity to a union of admissible
cells rather than exact membership. A polynomial can then replace each
distance query with a controlled error. A final polynomial distinguishes
finitely many separated expanded cells. An initial norm measurement
provides a compact sphere, whose lower-dimensional colored cover improves
the count, and the radial value needed to subtract a convexifying polynomial. Thus the polynomial result is obtained by a
robustification of the known colored-cover construction, not by claiming
that exact membership indicators are analytic.

\subsection{Complete adaptive budget curves}

For a total budget $B$ and a per-individual cap $Q$, the Bernoulli
population risk is, uniformly over $1\le Q\le B$,
\[
 R^{\rm ad}_{B,Q}\asymp R^{\rm fix}_{B,Q}\asymp
 \max\left\{B^{-1/3},Q^{-1},
       [Q\log(e+B/Q^{3/2})]^{-1/2}\right\}.
\]
The preceding paper had left a critical adaptive window open.
The new lower bound uses localized moment cancellation together with a
positive global background. It bounds the next-toss Fisher information
for every possible success/failure history, closing that window.
At $Q\asymp B^{2/3}/\log B$, the remaining cost is
$B^{-1/3}\sqrt{\log B/\log\log B}$, even with adaptation.

For an arbitrary joint endpoint law on $[1,2]^d$, one observation is a
conditionally independent coordinate vector from the corresponding
uniform intervals. The exact capped risk is
\[
 U^{\rm ad}_{B,Q,d}\asymp U^{\rm fix}_{B,Q,d}\asymp
 \begin{cases}
 \max\{\log(eB)/\sqrt B,(BQ)^{-1/3}\},&d=1,\\
 \max\{B^{-1/(d+1)},(BQ^{2d-1})^{-1/(3d)}\},&d\ge2.
 \end{cases}
\]
Dependence among the latent endpoints is unrestricted. The lower bound
uses a stronger Poisson record experiment and the full vector of
posterior scores, including all face terms. Thus the scalar critical
logarithm and its disappearance in higher dimensions are consequences
of the same allocation and information calculations.

\subsection{Finite observations and sharp inverse boundaries}

The deterministic binomial inverse modulus is determined at every
replication size and error level, rather than just along a statistical
asymptotic regime. Its lower bound identifies a precise limitation of
the global count-Wasserstein approach to MLE. A separate posterior repair
argument shows that normalized likelihood deficit $O(N^{-1})$ is enough
to retain the preceding statistical guarantees; a rational finite-grid
algorithm certifies this deficit.

For Gaussian location-covariance mixtures whose covariance rank is at
most $r<d$, $r+1$ independent conditional blocks identify an arbitrary
mixing law without moment conditions. Every experiment with at most $r$
blocks, even all block lengths together, can fail on compact exact-rank
models. The obstruction is a determinant differential operator that
survives the orthogonal covariance-factor quotient. The positive theorem
uses the covariance range revealed by repeated Gaussian contrasts.

For a positive $m$-atomic spectrum with exponents in $(0,1)$, an entire
transition curve can be replaced by any $2m$ distinct prescribed positive
settings. When the spectrum is a probability law in a fixed compact
subinterval, the sharp noisy recovery rate is $\delta^{1/(2m-1)}$ without
atom separation. This transfers the classical collision exponent from
moment data to nonlinear optimized thresholds. These thresholds are
population observables; the theorem does not assert recovery from a
finite-height list of Riemann-zeta zeros or recover prime phases.

\subsection{Scope of the claimed advances}

The target register at the end supplies individual statuses and theorem
references. Some entries are newly formulated inverse problems, some
close gaps left in this research program, and the analytic/convex query
results address directions explicitly asked in an external paper.
A logarithmic-risk equivalence is distinguished from a sharp minimax
exponent, and a partial parameter strip from a full classification.
Neither a count of targets nor the length of the manuscript establishes
novelty, correctness, or publication level. Internal independent proof
reviews and a targeted primary-source comparison were performed;
external specialist review and a broader priority check remain necessary.


% ===== knu.tex =====
\section{Adaptive recovery by convex polynomial measurements}
\label{p05:fi:knu-section}

Krieg--Novak--Ullrich~\cite[Section~3]{p03:an:KNU} ask for complexity
results when adaptive continuous information is restricted to
homogeneous, convex, or analytic functionals.  Their continuous
construction uses distance oracles and exact membership decisions.
We prove that a robust version of the color search permits polynomial
oracles instead.  An initial energy measurement then converts every
subsequent polynomial oracle into a globally convex polynomial without
increasing the number of subsequent measurements.

The information model is the exact-real-number model of that question.
At each step an algorithm selects a scalar functional on
$\mathbb R^m$ using only earlier answers and obtains its exact value.
The selected functional must belong to the indicated class, but its
degree, coefficients, and the selection rule have no complexity bound.
The final reconstruction may be discontinuous.  These conventions
matter: neither finite-precision stability nor efficient evaluation of
the polynomials is asserted.

\begin{theorem}[Logarithmic complexity with convex polynomials]
\label{p05:fi:knu-poly}
Let $m\ge1$, let $\|\cdot\|$ be any norm on $\mathbb R^m$, and
let $\varepsilon>0$.  There is an adaptive procedure recovering every
$x\in\mathbb R^m$ with
\[
                         \|x-\widehat x\|\le\varepsilon
\]
using at most
\begin{equation}\label{p05:fi:knu-count}
             \min\{m,\lceil\log_2m\rceil+2\}
\end{equation}
scalar measurements, each of which is a globally convex real
polynomial on $\mathbb R^m$.

In particular, an exponential reduction in the number of exact
measurements remains possible under the simultaneous restrictions
``convex'' and ``real analytic''.  The theorem gives an upper bound;
it does not determine the optimal number of such measurements.
\end{theorem}

The following compact-space statement is the central step.

\begin{lemma}[Polynomial measurements from a finite colored cover]
\label{p05:fi:robust-colors}
Let $K\subset\mathbb R^m$ be compact.  Suppose finitely many
nonempty compact sets $D_i\subset\mathbb R^m$ cover $K$, have
diameter at most $\varepsilon/2$ in the chosen norm, and have colors
in $\{1,\ldots,M\}$.  Suppose distinct sets of the same color are
at distance at least $\sigma>0$.  Then $K$ can be recovered to
error at most $\varepsilon$ by at most
\[
                         \lceil\log_2M\rceil+1
\]
adaptive real polynomial measurements on $\mathbb R^m$.
No convexity of these intermediate polynomials is required.
\end{lemma}
\begin{proof}
Discard colors not used by any set and decrease $\sigma$, if
necessary, so that $\sigma\le\varepsilon$.
Let $r=\lceil\log_2M\rceil$.  If $M=1$, only the final coding
measurement below is needed.  Otherwise put
$\eta=\sigma/(8r)$.
For a nonempty set $J$ of colors define
\[
 E_J=\bigcup_{i:\,\mathrm{color}(D_i)\in J}D_i,
 \qquad d_J(x)=\operatorname{dist}(x,E_J).
\]
For every such $J$, select a real polynomial $p_J$ with
\begin{equation}\label{p05:fi:distance-approx}
             \sup_{x\in K}|p_J(x)-d_J(x)|\le\eta.
\end{equation}
These approximations exist by multivariate polynomial approximation
on a box containing $K$, since $d_J$ is continuous on the whole box.
For clarity, no approximation of a discontinuous membership indicator
is used.

Start with all colors retained and error allowance $e_0=0$.
After $j$ queries maintain the invariant
\begin{equation}\label{p05:fi:search-invariant}
              \operatorname{dist}(x,E_J)\le e_j,
              \qquad e_j=2j\eta.
\end{equation}
If more than one color remains, divide $J$ into two nonempty sets
$J_1,J_2$ with sizes differing by at most one and query $p_{J_1}(x)$.
Retain $J_1$ if this answer is at most $e_j+\eta$, and otherwise
retain $J_2$.
In the first case \eqref{p05:fi:distance-approx} gives
$d_{J_1}(x)\le e_j+2\eta$.
In the second case it gives $d_{J_1}(x)>e_j$; since
$\min(d_{J_1}(x),d_{J_2}(x))\le e_j$, necessarily
$d_{J_2}(x)\le e_j$.
Thus the invariant holds with allowance $e_{j+1}$ in both cases.
After at most $r$ queries there is one retained color $c$ and
\[
           \operatorname{dist}(x,E_{\{c\}})\le\sigma/4.
\]

For the finitely many sets of color $c$, write
\[
       F_i=\{y:\operatorname{dist}(y,D_i)\le\sigma/3\}.
\]
They are compact and pairwise separated by at least $\sigma/3$.
There is a continuous real function $g_c$ on $\mathbb R^m$ which
equals the distinct integer $i$ on $F_i$.  For an explicit choice,
put
\[
 g_c(y)=\sum_{i:\,\mathrm{color}(D_i)=c}
 i\max\{0,1-6\operatorname{dist}(y,F_i)/\sigma\}.
\]
Its supports do not intersect at points where both summands are
positive.  Approximate $g_c$ on a box containing $K$ by a polynomial
$q_c$ with error less than $1/4$, and query $q_c(x)$.
The point $x$ belongs to one of the $F_i$, so rounding the answer
to the nearest integer identifies that $i$.
Output a fixed $x_i\in D_i$.  Then
\[
 \|x-x_i\|\le\operatorname{dist}(x,D_i)
                  +\operatorname{diam}(D_i)
             \le\sigma/4+\varepsilon/2<\varepsilon.
\]
If $M=1$, $x$ already belongs to some $D_i$ and the same coding
query applies directly.

The polynomial approximation invoked twice above can be obtained
explicitly from multivariate Bernstein polynomials after rescaling
the box to $[0,1]^m$.  Their value at a point is the expectation of
the target function at independent binomial grid coordinates.
Uniform continuity and the uniform bound on their coordinate
variances prove uniform convergence.  Thus no regularity beyond
continuity of the two distance-based functions is needed.
\end{proof}

\begin{lemma}[A small colored cover of a sphere]
\label{p05:fi:sphere-cover}
For $m\ge2$, $R>0$, and $d>0$, the Euclidean sphere
$S_R^{m-1}\subset\mathbb R^m$ has a finite cover by nonempty
compact sets of diameter at most $d$ in any prescribed norm,
colored with $m$ colors, such that distinct sets of the same color
are a positive distance apart.
\end{lemma}
\begin{proof}
Take a finite triangulation of the sphere, refine it so that each
closed vertex star after barycentric subdivision has diameter at
most $d$, and use that subdivision.  Such a triangulation can be
obtained from the boundary of a cross-polytope by radial projection;
successive subdivision and uniform continuity give arbitrarily
small stars.
Vertices in the barycentric subdivision correspond to faces in the
preceding triangulation.  Color each vertex by the dimension of
that face, one of $0,\ldots,m-1$.
Vertices of the same color never belong to one simplex, since a
simplex of the subdivision corresponds to a strictly nested chain
of faces.

Let $\varphi_v$ be the continuous piecewise affine barycentric
coordinate associated with a vertex $v$, transferred to the sphere.
The sets
\[
                    D_v=\{x:\varphi_v(x)\ge1/m\}
\]
cover the sphere: at any point at most $m$ barycentric coordinates
are nonzero and their sum is one.
Each $D_v$ is compact and contained in its small closed star.
For distinct vertices of one color,
$\varphi_v\varphi_w=0$, so $D_v$ and $D_w$ are disjoint.
There are finitely many such pairs; compactness makes their minimum
distance strictly positive.  Unused or empty sets may be discarded.
\end{proof}

\begin{lemma}[Convexification using known energy]
\label{p05:fi:convexification}
For every real polynomial $P$ on $\mathbb R^m$ there are an integer
$D\ge1$ and a constant $C\ge0$ such that
\[
                 Q(x)=P(x)+C(1+\|x\|_2^2)^D
\]
is a globally convex polynomial.
If $E=\|x\|_2^2$ is already known, one evaluation of $Q(x)$
therefore determines $P(x)$ exactly.
\end{lemma}
\begin{proof}
If $d$ is the degree of $P$, its Hessian satisfies
\[
 \|\nabla^2P(x)\|_{2\to2}
       \le C_P(1+\|x\|_2^2)^{\max\{d-2,0\}/2}.
\]
Choose an integer $D$ with
$D-1\ge\max\{d-2,0\}/2$.  Direct differentiation gives
\[
 \nabla^2(1+\|x\|_2^2)^D
 \succeq 2D(1+\|x\|_2^2)^{D-1}I_m.
\]
A sufficiently large $C$ makes $\nabla^2Q$ positive semidefinite
everywhere, which proves convexity.  Subtract
$C(1+E)^D$ from the measured value to recover $P(x)$.
\end{proof}

\begin{proof}[Proof of Theorem~\ref{p05:fi:knu-poly}]
The coordinate functionals provide exact recovery in $m$ queries;
they are convex polynomials.  For the other bound, first query the
convex polynomial $E(x)=\|x\|_2^2$.
If the answer is zero, output zero.  Otherwise the input is now
known to belong to the compact sphere of radius $R=\sqrt E$.
Apply Lemma~\ref{p05:fi:sphere-cover} with $d=\varepsilon/2$, and
then Lemma~\ref{p05:fi:robust-colors}.  This requires at most
$\lceil\log_2m\rceil+1$ further polynomial measurements.
Each selected polynomial is converted into a convex polynomial by
Lemma~\ref{p05:fi:convexification}; the initial energy allows its
original value to be recovered exactly from the same single answer.
All later branch decisions therefore agree with the robust
polynomial algorithm.  Including the energy query gives
\eqref{p05:fi:knu-count}.  The case $m=1$ is already covered by the
coordinate procedure.
\end{proof}

\begin{corollary}[Convex analytic measurements with unit Lipschitz bound]
\label{p05:fi:knu-analytic-lipschitz}
The bound \eqref{p05:fi:knu-count} also holds with every measurement
globally convex, globally real analytic, and Lipschitz with constant
at most one in the Euclidean norm.  These measurements need not be
polynomials.
\end{corollary}
\begin{proof}
First query $A(x)=\sqrt{1+\|x\|_2^2}$, which has all three
properties and determines $E=\|x\|_2^2$ exactly.
If $E=0$, output zero and stop.
Run the same polynomial algorithm on the resulting sphere.
To simulate a selected polynomial $P$, put
$F(y)=P(y)e^{-\|y\|_2^2}$.
The function $F$ and all its derivatives are polynomials times a
Gaussian.  In particular
\[
 M=\sup_y(1+\|y\|_2^2)^{3/2}
              \|\nabla^2 F(y)\|_{2\to2}<\infty,
 \qquad L=\sup_y\|\nabla F(y)\|_2+M<\infty.
\]
Since
$\nabla^2 A(y)\succeq(1+\|y\|_2^2)^{-3/2}I_m$,
the function $Q=F+MA$ is convex.  It is globally real analytic and
has Lipschitz constant at most $L$.
Query $Q/\max\{1,L\}$.  Multiply by its known scaling factor,
subtract the already known $MA(x)$, and multiply by $e^E$.
This recovers $P(x)$ exactly from one measurement, with no new query.
The coordinate alternative is obtained by scaling each coordinate
functional if needed.
\end{proof}

\subsection{Positively homogeneous convex information}

There is a related conclusion for homogeneity, with its convention
specified explicitly.  A function $F$ is positively one-homogeneous
when $F(tx)=tF(x)$ for $t\ge0$.  We do not require
$F(-x)=-F(x)$.

\begin{theorem}[Convex positively homogeneous measurements]
\label{p05:fi:knu-homogeneous}
The recovery and measurement-count bound of
Theorem~\ref{p05:fi:knu-poly} also hold when every measurement is
required instead to be globally convex and positively
one-homogeneous.  The measurements can be continuous everywhere
and real analytic on $\mathbb R^m\setminus\{0\}$, with
Euclidean Lipschitz constant at most one after a known rescaling.
They are not asserted to be real analytic at the origin.
\end{theorem}
\begin{proof}
First query $R=\|x\|_2$ and output zero if $R=0$.
On the known sphere, use the same finite polynomial algorithm.
For any selected polynomial $P$, put
\[
 H(y)=\frac{\|y\|_2}{R}
       P\!\left(\frac{R y}{\|y\|_2}\right)
       \quad(y\ne0),\qquad H(0)=0.
\]
This function is continuous, positively one-homogeneous and real
analytic away from zero, and $H(x)=P(x)$ for the unknown input.
It can be convexified by adding $C\|y\|_2$.
Indeed Euler's identity gives $\nabla^2H(y)y=0$, and the Hessian
is homogeneous of degree $-1$.  Smoothness on the unit sphere
therefore gives a constant $C_1$ with
\[
 \nabla^2H(y)\succeq
       -\frac{C_1}{\|y\|_2}
       \left(I_m-\frac{yy^{\mathsf T}}{\|y\|_2^2}\right).
\]
The matrix on the right, without its negative coefficient, is the
Hessian of $\|y\|_2$.  Thus
$Q(y)=H(y)+C\|y\|_2$ has positive semidefinite Hessian away
from zero when $C\ge C_1$.
Increase $C$ so that $C\ge\sup_{\|v\|_2=1}|H(v)|$.
On a line passing through zero the one-sided slopes of $Q(tv)$
then have the correct order, because $Q(v)+Q(-v)\ge0$.
On every other line convexity follows from the Hessian inequality.
Hence $Q$ is globally convex.  It is positively one-homogeneous
and querying it determines $P(x)=Q(x)-CR$, since $R$ is known.
The number of queries is unchanged.
Each constructed $Q$ is continuous, convex, and positively
one-homogeneous, hence subadditive.  Its bounded values on the
unit sphere imply a finite global Lipschitz constant.  Dividing
by a known upper bound and undoing that scaling in the scalar
answer makes the constant at most one.
\end{proof}

\begin{corollary}[Homogeneous polynomials of unrestricted degree]
\label{p05:fi:knu-homogeneous-polynomials}
If ``homogeneous polynomial'' permits a different nonnegative integer
degree for each measurement, then recovery on $\mathbb R^m$ to any
prescribed error requires at most
\[
                  \min\{m,2\lceil\log_2m\rceil+3\}
\]
adaptive homogeneous polynomial measurements.  The homogeneity
identity holds for every real scalar.  Convexity of these polynomials
is not asserted.
\end{corollary}
\begin{proof}
Query $E=\|x\|_2^2$ first and stop if it is zero.  Run the
polynomial sphere algorithm.  For each polynomial
$P(y)=\sum_\alpha c_\alpha y^\alpha$, split its monomials
into their even and odd total-degree parts $P_0,P_1$.
For $e\in\{0,1\}$ with nonzero part, choose an integer $D_e$
of parity $e$ at least as large as the degree of $P_e$ and define
\[
 H_e(y)=\sum_{|\alpha|\equiv e\,({\rm mod}\,2)}
 c_\alpha y^\alpha
       \left(\frac{\|y\|_2^2}{E}\right)^{(D_e-|\alpha|)/2}.
\]
This is a real homogeneous polynomial of degree $D_e$.
At the actual input, $H_e(x)=P_e(x)$.
At most two queries therefore simulate each of the
$\lceil\log_2m\rceil+1$ polynomial queries after the energy.
The coordinate alternative gives the minimum with $m$.
\end{proof}

\paragraph{The resolved direction and the remaining question.}
In the exact information model, restricting adaptive measurements
to globally convex polynomials does not restore the $m$-measurement
requirement of linear information.  Logarithmically many queries
still suffice.  The same is true for convex positively homogeneous
information.  Thus these theorems answer the exponential-speedup
direction of the cited complexity questions.  They do not prove
an optimal count, nor do they provide a bounded-degree or stable
finite-precision implementation.

The global domain in Theorem~\ref{p05:fi:knu-poly} is handled by an
exact preliminary energy query, rather than by assuming a uniform
polynomial approximation on all of $\mathbb R^m$.
Finally, a globally real-analytic, positively one-homogeneous
function is linear: differentiability at zero gives
$F(x)=\lim_{t\downarrow0}F(tx)/t=DF(0)x$.
Accordingly the two theorems impose different simultaneous
restrictions; they do not claim logarithmic recovery using
functionals that are both globally analytic and one-homogeneous.

% ===== mle.tex =====
\section{The sharp binomial inverse modulus and certified approximate likelihood}
\label{p05:fm:sec}

The preceding manuscript proved that the unrestricted binomial mixture
maximum likelihood estimator is rate optimal throughout every fixed
polynomially sublinear replication range. Two distinct questions remain:
whether that estimator is optimal at the statistical transition, and how
accurately its likelihood must be optimized. We resolve a deterministic
inverse problem underlying the first question, and give a quantitative,
computable answer to the second. The first result also proves that one
particular strategy cannot settle the remaining statistical transition.
We do not claim to settle that transition itself.

For a probability measure $G$ on $[0,1]$, let
\[
 B_j^t(x)=\binom tjx^j(1-x)^{t-j},\qquad
 p_G(j)=\int B_j^t(x)\,dG(x),\qquad 0\le j\le t.
\]
Write $W_{1,t}$ for Wasserstein distance on the count grid with cost
$|j-l|/t$. All logarithms are natural. Constants below are absolute unless
specified otherwise.

\subsection{A sharp modulus, uniformly in replication and measurement error}

Define the deterministic inverse modulus
\[
 \omega_t(\delta)=\sup\{W_1(G,H):W_{1,t}(p_G,p_H)\le\delta\},
 \qquad 0\le\delta\le1.
\]
\begin{theorem}[Sharp binomial inverse modulus]\label{p05:fm:modulus}
For every integer $t\ge1$ and $0\le\delta\le1$,
\begin{equation}\label{p05:fm:modulusformula}
 c\Phi_t(\delta)\le\omega_t(\delta)\le C\Phi_t(\delta),\qquad
 \Phi_t(\delta)=\min\left\{1,\frac1t+
 \frac1{\sqrt{t\log(1+1/(t\delta^2))}}\right\}.
\end{equation}
At $\delta=0$ the second term is defined to be zero.
In particular, the order of the entire interpolation expression, including
its transition, is necessary for deterministic inversion from count
Wasserstein error alone.
\end{theorem}

The spectral calculation used below is the classical beta--binomial Gibbs
diagonalization of Diaconis, Khare and Saloff-Coste
\cite[Section~4.1]{p05:fm:DKS2008}. Its use here is to construct sharp
Wasserstein inverse witnesses. We provide the calculation to make the
result independent of normalization conventions for Jacobi and Hahn
polynomials. Neither those eigenvalues nor their diagonalization are
claimed as new.

\begin{lemma}[Oscillatory witnesses]\label{p05:fm:witness}
For every integer $k\ge1$, put
\[
 dG_k^{\pm}(x)=\{1\pm T_k(2x-1)\}
                 \frac{dx}{\pi\sqrt{x(1-x)}}.
\]
These are probability measures, and
\begin{equation}\label{p05:fm:witnesslatent}
 \frac c{k}\le W_1(G_k^+,G_k^-)
 =\frac1{\pi k}\int_0^\pi|\sin(k\theta)|\sin\theta\,d\theta
 \le\frac C{k}.
\end{equation}
For $1\le k\le t$,
\begin{equation}\label{p05:fm:witnesscount}
 W_{1,t}(p_{G_k^+},p_{G_k^-})
 \le\frac Ck\left\{\prod_{i=0}^{k-1}
                  \frac{t-i}{t+1+i}\right\}^{1/2}
 \le\frac Ck e^{-k^2/(4t)}.
\end{equation}
For $k>t$ the two count laws are identical.
\end{lemma}
\begin{proof}
Let $\pi$ be the arcsine law on $[0,1]$, and let
$h_k(x)=T_k(2x-1)$. Under $x=(1+\cos\theta)/2$, $\pi$ is uniform
on $[0,\pi]$ and $h_k(x)=\cos(k\theta)$. Thus $\int h_k\,d\pi=0$,
$|h_k|\le1$, and the signed cumulative distribution function is
$-2\sin(k\theta)/(\pi k)$. Integration gives
\eqref{p05:fm:witnesslatent}. Its lower bound follows, for example, by
restricting to $[\pi/4,3\pi/4]$, where $\sin\theta\ge1/\sqrt2$;
the integral of $|\sin(k\theta)|$ on this interval is bounded below
uniformly in the integer $k\ge1$.

Draw $X\sim\pi$ and $J\mid X\sim\operatorname{Binomial}(t,X)$.
Let $Qf(j)=\mathbb E[f(X)\mid J=j]$ and
$Q^*v(x)=\mathbb E[v(J)\mid X=x]$.
The operator $A=Q^*Q$ is self-adjoint on $L^2(\pi)$ and preserves
polynomials of degree at most $k\le t$. Indeed, the posterior law is
$\operatorname{Beta}(j+1/2,t-j+1/2)$, so
\[
 \mathbb E[X^k\mid J=j]=
 \frac{(j+1/2)(j+3/2)\cdots(j+k-1/2)}
      {(t+1)(t+2)\cdots(t+k)}.
\]
The leading coefficient of $Ax^k$ is therefore
\[
 \lambda_{t,k}=\frac{t(t-1)\cdots(t-k+1)}
                         {(t+1)(t+2)\cdots(t+k)}.
\]
Since $h_k$ is orthogonal to every lower-degree polynomial,
self-adjointness and degree preservation give $Ah_k=\lambda_{t,k}h_k$.
Consequently
\[
 \mathbb E[Qh_k(J)^2]=\langle h_k,Ah_k\rangle_{L^2(\pi)}
                    =\lambda_{t,k}/2.
\]
If $p_\pi$ denotes the beta--binomial count law, the signed difference of
the two count laws is $\Delta(j)=2p_\pi(j)Qh_k(j)$. Hence
$\sum_j|\Delta(j)|\le\sqrt{2\lambda_{t,k}}$.
Moreover, $\sum_j P(j/t)\Delta(j)=0$ whenever $P$ has degree less than
$k$: condition on $X$ and use polynomial orthogonality.
Approximate an arbitrary $1$-Lipschitz function on $[0,1]$ by a polynomial
of degree at most $k-1$ with uniform error $C/k$. Such a Jackson
approximation is constructed in Lemma~\ref{p05:fm:polyinverse} below; the
case $k=1$ uses a constant. Thus its integral against $\Delta$ is at most
$C\sqrt{\lambda_{t,k}}/k$. Wasserstein duality proves the first inequality
in \eqref{p05:fm:witnesscount}. For the second, use
\[
 \log\lambda_{t,k}
 =\sum_{i=0}^{k-1}\log\left(1-\frac{1+2i}{t+1+i}\right)
 \le-\frac{k^2}{t+k}\le-\frac{k^2}{2t}.
\]
Finally, for $k>t$, each $B_j^t$ has degree below $k$ and is orthogonal
to $h_k$, proving exact equality of the count laws.
\end{proof}

\begin{lemma}[A derivative-preserving inverse inequality]\label{p05:fm:polyinverse}
For any two mixing laws and any integer $2\le k\le t/2$,
\begin{equation}\label{p05:fm:inverseineq}
 W_1(G,H)\le C/k+Ce^{4k^2/t}W_{1,t}(p_G,p_H).
\end{equation}
\end{lemma}
\begin{proof}
We recall the quantitative ingredients, with proofs. A $1$-Lipschitz
function $f$ on $[0,1]$, after subtraction of $f(0)$, extends by reflection
to a period-two $1$-Lipschitz function of supremum norm at most one.
Convolution with the positive, normalized Jackson kernel
\[
 J_s(u)=c_s\left\{\frac{\sin(s\pi u/2)}{\sin(\pi u/2)}\right\}^4
\]
gives a trigonometric polynomial $Q_s$, of frequencies $|l|\le2s$,
with $\|Q_s-f\|_\infty\le C/s$ and $\|Q_s'\|_\infty\le1$.
The approximation estimate follows from the kernel's first absolute
moment $O(1/s)$: the unnormalized integral is comparable to $s^3$ and
the first moment is $O(s^2)$, using the bound
$C\min(s,|u|^{-1})$ for the sine quotient. Every Fourier coefficient
of $Q_s$ has modulus at most one. Replacing each
$e^{i\pi l(x-1/2)}$ by its degree-$100s$ Taylor polynomial changes both
the function and its derivative by at most $e^{-100s}$, by the factorial
bound $m!\ge(m/e)^m$. Choosing $s=\lfloor k/100\rfloor$ for $k\ge100$,
and a constant for smaller $k$, gives a polynomial $P_k$ with
\[
 \deg P_k\le k,\qquad \|P_k-f\|_\infty\le C/k,
 \qquad \|P_k'\|_\infty\le2.
\]

We also need the following Bernstein coefficient estimate: if
$P=\sum_{j=0}^n b_jB_j^n$ has degree $m\le n/2$, then
\begin{equation}\label{p05:fm:berncoeff}
 \max_j|b_j|\le Ce^{2m^2/n}\|P\|_\infty.
\end{equation}
Here is a short contour proof. The maximum principle applied to
$v^mP((2+v+v^{-1})/4)$ gives
$|P((2+w+w^{-1})/4)|\le\|P\|_\infty|w|^m$ for $|w|\ge1$.
For $0<j<n$ set $r=j/(n-j)$, $q=r/(1+r)^2$, and
$u=2\sqrt q|\sin(\phi/2)|$. The larger inverse Joukowski modulus
corresponding to $re^{i\phi}/(1+re^{i\phi})$ is
$(1+u)/\sqrt{1-u^2}$. Thus, for
$F(z)=(1+z)^nP(z/(1+z))$,
\[
 \frac{|F(re^{i\phi})|}{\|P\|_\infty(1+r)^n}
 \le(1-u^2)^{(n-m)/2}(1+u)^m
 \le e^{2m^2/n}e^{-nu^2/8}.
\]
The average of the last Gaussian factor is at most
$C/\sqrt{1+nq}$. Cauchy's coefficient formula and the elementary
Stirling lower bound
$\binom nj(j/n)^j(1-j/n)^{n-j}\ge c/\sqrt{1+nq}$ prove
\eqref{p05:fm:berncoeff}. Endpoint coefficients are just $P(0),P(1)$.

Write $P_k=\sum_jb_jB_j^t$. Since
$P_k'=t\sum_{j=0}^{t-1}(b_{j+1}-b_j)B_j^{t-1}$,
\eqref{p05:fm:berncoeff} gives
$\max_j t|b_{j+1}-b_j|\le Ce^{4k^2/t}$.
Apply Wasserstein duality to this Lipschitz coefficient vector; the two
polynomial approximation errors contribute $C/k$.
\end{proof}

\begin{proof}[Proof of Theorem~\ref{p05:fm:modulus}]
For the upper bound let $x=1/(t\delta^2)$ and $L=\log(1+x)$.
Choose a small absolute $a>0$, and
$k=\lfloor\min(t/2,\sqrt{atL})\rfloor$ whenever that minimum is at
least four. The inverse inequality then yields
$C/t+C/\sqrt{tL}$. In fact, in the uncapped case its exponential term,
multiplied by $\sqrt{tL}$, is at most
\[
 C(1+x)^{4a}\sqrt{L/x},
\]
which is bounded when $4a<1/2$. In the capped case $L\ge t/(4a)$,
and $\delta e^t\le C/t$. If the minimum is below four, the trivial
bound one gives the result. At $\delta=0$, use $k$ of order $t$ and
handle bounded $t$ trivially.

For the lower bound, Lemma~\ref{p05:fm:witness} with $k=t+1$ always gives
$\omega_t(\delta)\ge c/t$. Also, endpoint mixing laws differing by
mass $\delta$ give $\omega_t(\delta)\ge\delta$. If
$x=1/(t\delta^2)\le1$, then $L\ge x/2$ and
$1/\sqrt{tL}\le\sqrt2\delta$, so these two witnesses suffice.
Suppose $x>1$. Let $A$ be a sufficiently large absolute constant and
$k=\lceil A\sqrt{tL}\rceil$. If $k\le t$, then
\[
 \frac{C k^{-1}e^{-k^2/(4t)}}{\delta}
 \le\frac{C\sqrt{x}}{A\sqrt L}(1+x)^{-A^2/4}\le1.
\]
The final inequality holds for all $x>1$ after increasing $A$.
The same witness has latent distance at least $c/k\ge c'/\sqrt{tL}$.
If $k>t$, then, apart from bounded $t$, the inequality
$A\sqrt{tL}>t-1$ implies $1/\sqrt{tL}\le2A/t$, so the exact
moment-matching witness already suffices. Combining the lower bounds,
and clipping at one, proves the result.
\end{proof}

\begin{corollary}[A limitation of global transport inversion]\label{p05:fm:barrier}
At $t=\lfloor N/\log N\rfloor$,
\[
 \omega_t(N^{-1/2})\asymp
       \sqrt{\frac{\log N}{N\log\log N}}.
\]
Consequently a deterministic argument using only a count-Wasserstein
error bound of order $N^{-1/2}$ cannot guarantee latent error of order
$N^{-1/2}$ at this transition.
\end{corollary}
This does not lower-bound the statistical risk of the MLE: the witness
pairs need not represent its fitted-versus-true pairs on likely samples.
Closing the statistical transition would require additional likelihood
structure or a more informative, localized observation norm. The open
coefficient and fixed-test questions in \cite[Section~6.2]{p01:info:MMRS2025}
are not answered by this deterministic obstruction.

\subsection{Entropy--transport with an optimization certificate}

Let $(\mathcal X,d)$ be a finite metric space of diameter $D_0$, and
let $K_\theta$ be probability kernels. For a probability vector $q$ write
$\ell(H)=\sum_xq(x)\log p_H(x)$ and $\ell^*=\sup_H\ell(H)$.
The following statement extends the exact likelihood transport identity;
it does not assume that mixtures inherit the concentration of their
components.

\begin{theorem}[Posterior repair of approximate mixture likelihood]
\label{p05:fm:approxtransport}
Suppose each kernel satisfies
$W_{1,d}(r,K_\theta)^2\le2cD(r\Vert K_\theta)$ for all $r$.
Let $H$ have $p_H>0$ on the support of $q$, and put
$e=\ell^*-\ell(H)$. Define
\[
 S_H(\theta)=\sum_{x:q(x)>0}\frac{q(x)K_\theta(x)}{p_H(x)},
 \qquad dH^+(\theta)=S_H(\theta)\,dH(\theta),
 \qquad d_H=D(H^+\Vert H).
\]
Then $H^+$ is a probability measure and
\begin{equation}\label{p05:fm:repair}
 W_{1,d}(q,p_H)
 \le\sqrt{2c\{D(q\Vert p_H)-d_H\}}
       +D_0\|H^+-H\|_{\mathrm{TV}},\qquad
 d_H\le\ell(H^+)-\ell(H)\le e.
\end{equation}
Here total variation is one half the $L^1$ norm. In particular,
\begin{equation}\label{p05:fm:approxT1}
 W_{1,d}(q,p_H)
 \le\sqrt{2cD(q\Vert p_H)}+D_0\sqrt{e/2}.
\end{equation}
Alternatively, if $\sigma=\sup_\theta S_H(\theta)\le1+\eta$, then
\begin{equation}\label{p05:fm:scorecertificate}
 e\le\log(1+\eta),\qquad
 W_{1,d}(q,p_H)
 \le\sqrt{2cD(q\Vert p_H)}+D_0\frac{\eta}{1+\eta}.
\end{equation}
\end{theorem}
\begin{proof}
The identity $\int S_H\,dH=1$ is immediate. For $S_H(\theta)>0$ let
\[
 R_\theta(x)=\frac{q(x)K_\theta(x)}{p_H(x)S_H(\theta)}.
\]
This is a probability vector, and $\int R_\theta\,dH^+=q$.
Finite-sum calculation gives the exact entropy decomposition
\begin{equation}\label{p05:fm:entropydecomposition}
 D(q\Vert p_H)=D(H^+\Vert H)+
                  \int D(R_\theta\Vert K_\theta)\,dH^+(\theta).
\end{equation}
Terms on $S_H=0$ are absent. Convexity of transportation cost, the
kernel inequality and Cauchy--Schwarz bound $W_{1,d}(q,p_{H^+})$ by
the square root in \eqref{p05:fm:repair}. Coupling the common part of
$H,H^+$ gives $W_{1,d}(p_{H^+},p_H)\le D_0\|H^+-H\|_{\mathrm{TV}}$.

The posterior of $\theta$ at an observed $x$ under $H$ is
$K_\theta(x)dH(\theta)/p_H(x)$. Jensen's inequality yields
\begin{align*}
 \ell(H^+)-\ell(H)
 &=\sum_xq(x)\log\mathbb E_H[S_H(\theta)\mid x]\\
 &\ge\sum_xq(x)\mathbb E_H[\log S_H(\theta)\mid x]
 =\int S_H\log S_H\,dH=d_H.
\end{align*}
This is the familiar EM monotonicity calculation, here retaining the
specific divergence needed for transportation. Pinsker's inequality
gives \eqref{p05:fm:approxT1}. For any comparison mixing law $F$, a second
application of Jensen gives
\[
 \ell(F)-\ell(H)
 \le\log\sum_xq(x)\frac{p_F(x)}{p_H(x)}
 =\log\int S_H\,dF\le\log\sigma.
\]
Finally, when $0\le S_H\le1+\eta$ and $\int S_H\,dH=1$,
\[
 \|H^+-H\|_{\mathrm{TV}}
 =\int(S_H-1)_+\,dH
 \le\frac{\eta}{1+\eta}\int S_H\,dH.
\]
These prove the score statements.
\end{proof}

\begin{theorem}[Statistically sufficient likelihood precision]
\label{p05:fm:statisticalprecision}
Let $q$ be the empirical count law from $N$ independent observations
from $p_G$. Suppose a data-dependent mixing law $\widetilde G$ satisfies
$\ell^*-\ell(\widetilde G)\le e_0$ for every sample, where the likelihood
is normalized by $N$. Then, for $2\le k\le t/2$,
\begin{equation}\label{p05:fm:approxmaster}
 \mathbb E_G W_1(\widetilde G,G)
 \le\frac Ck+Ce^{4k^2/t}\{N^{-1/2}+\sqrt{e_0}\}.
\end{equation}
In particular $e_0\le C_0/N$ preserves the preceding manuscript's
uniform interpolation bound and its near-linear minimax rate. More
explicitly, for every fixed $a,\varepsilon,C_0>0$ and
$a\log N\le t\le N^{1-\varepsilon}$,
\[
 \mathbb E_G W_1(\widetilde G,G)
 \le\frac{C_{a,\varepsilon,C_0}}{\sqrt{t\log N}}.
\]
For all $t\ge N$, the same optimization tolerance gives risk
$O_{C_0}(N^{-1/2})$.
\end{theorem}
\begin{proof}
For a $1$-Lipschitz function on the normalized binomial grid, changing
one Bernoulli trial changes its value by at most $1/t$. Conditional
Hoeffding bounds give
$\log\mathbb E\exp\{\lambda(f-\mathbb Ef)\}\le\lambda^2/(8t)$.
The entropy variational inequality, optimized over $\lambda>0$, then
gives $W_{1,t}(r,K_\theta)^2\le D(r\Vert K_\theta)/(2t)$.
Theorem~\ref{p05:fm:approxtransport} applies with $c=1/(4t)$ and $D_0=1$.
Also
\[
 D(q\Vert p_{\widetilde G})\le D(q\Vert p_G)+e_0,
 \qquad \mathbb E D(q\Vert p_G)\le t/N.
\]
The last assertion follows from $D\le\chi^2$ and the elementary
multinomial variance identity. Moreover
$\mathbb E W_{1,t}(q,p_G)\le1/(2\sqrt N)$ by integrating
$\mathbb E|F_N-F|\le1/(2\sqrt N)$. Thus
\[
 \mathbb E W_{1,t}(p_{\widetilde G},p_G)
 \le C\{N^{-1/2}+\sqrt{e_0}\}.
\]
Use Lemma~\ref{p05:fm:polyinverse}. For the stated near-linear range choose
$k$ a sufficiently small fixed multiple of
$\sqrt{\varepsilon t\log N}$; the same optimization as in the upper
bound of Theorem~\ref{p05:fm:modulus} gives the all-range interpolation
bound. When $t\ge N$, that bound is $O(N^{-1/2})$.
\end{proof}

\begin{proposition}[Sharpness of root-sample likelihood precision]
\label{p05:fm:precisionlower}
There are constants $c,C,e_*>0$ such that, for all sufficiently large
$N$, all $t\ge1$, and $C/N\le e_0\le e_*$, there exists an
$e_0$-approximate likelihood maximizer whose risk at
$G=(\delta_0+\delta_1)/2$ is at least $c\sqrt{e_0}$.
Consequently likelihood deficit alone cannot guarantee root-$N$ risk
with a worst-case optimization tolerance larger in order than $N^{-1}$.
\end{proposition}
\begin{proof}
At this $G$ all observed counts are endpoints. Let $r$ be the empirical
mass at $t$. The endpoint law $H_r=(1-r)\delta_0+r\delta_1$ attains the
global maximum likelihood: for any mixture $p(0)+p(t)\le1$, and the
endpoint masses $1-r,r$ maximize the resulting two-coordinate
likelihood. Put $h=\sqrt{e_0}/4$. On the event $r\in[1/4,3/4]$, output
$H_{r+h}$; otherwise output $H_r$. For $e_*$ small enough,
$r+h\in[1/4,7/8]$, and
\[
 \ell(H_r)-\ell(H_{r+h})
 =D(\operatorname{Bernoulli}(r)\Vert\operatorname{Bernoulli}(r+h))
 \le\frac{h^2}{(r+h)(1-r-h)}\le e_0.
\]
Outside the endpoint-data set define the estimator to be any exact MLE.
On the stated event its loss is $|r+h-1/2|\ge h-|r-1/2|$.
Hoeffding's bound and $\mathbb E|r-1/2|\le1/(2\sqrt N)$ give expected
loss at least $h(1-2e^{-N/8})-1/(2\sqrt N)$, which is at least
$c\sqrt{e_0}$ when $C$ is large enough.
\end{proof}

\subsection{A polynomial-time, fully certified construction}

Approximate mixture likelihood by convex optimization is an established
computational approach \cite{p05:fm:FP2018}. The next result specifies an
elementary global certificate at the precision needed above. Its purpose
is a polynomial complexity guarantee, not a competitive runtime claim.

\begin{theorem}[Certified finite-grid likelihood optimization]
\label{p05:fm:algorithm}
For rational $q$ and rational $0<e_0\le1/4$, an $e_0$-approximate
unrestricted binomial-mixture MLE can be computed with a number of
rational arithmetic operations polynomial in $t,e_0^{-1}$ and the
input bit length. Its support lies on a grid of
$M+1=O(t^2/e_0)$ points. The algorithm terminates with a computable
upper bound on its unrestricted likelihood deficit.
\end{theorem}
\begin{proof}
Let $M\ge\max(2t,4t^2/e_0)$ be an integer and set
$\alpha=(M)_t/M^t$, where $(M)_t=M(M-1)\cdots(M-t+1)$.
For any $x\in[0,1]$, draw $Z\sim\operatorname{Binomial}(M,x)$.
Sampling $t$ indices with replacement from these $M$ Bernoulli trials
shows, coordinatewise,
\begin{equation}\label{p05:fm:griddomination}
 \mathbb E B_j^t(Z/M)\ge\alpha B_j^t(x).
\end{equation}
Indeed, the indices are distinct with probability $\alpha$, and on that
event the sampled values are independent Bernoulli$(x)$ variables;
the remaining event contributes a nonnegative probability.
Integrating over any mixing law gives a grid-supported law with
likelihood at least $\ell(H)+\log\alpha$. Furthermore
$-\log\alpha\le t(t-1)/M\le e_0/4$, using
$-\log(1-u)\le2u$ for $u\le1/2$.

Let $u$ denote uniform weights on the grid $\{0,1/M,\ldots,1\}$ and
put $\gamma=e_0/4$. The grid law generated by applying the above
randomization to uniform $x$ has exactly these weights, because
$\int_0^1B_m^M(x)\,dx=1/(M+1)$. Its predictive probabilities are
therefore at least $\alpha/(t+1)\ge1/(2(t+1))$.
Restrict weights to $w=(1-\gamma)v+\gamma u$, with $v$ in the simplex.
Replacing any grid weights by this mixture loses at most
$-\log(1-\gamma)\le e_0/3$ in normalized likelihood. Every feasible
predictive coordinate is at least $\gamma/(2(t+1))$.

Write $L(v)=\ell((1-\gamma)v+\gamma u)$. At a current $v$ compute
the usual simplex directional gap
\[
 g(v)=\max_m\partial_m L(v)-\sum_m v_m\partial_mL(v).
\]
Concavity gives $\max L-L(v)\le g(v)$. Hence the explicitly computable
quantity
\begin{equation}\label{p05:fm:computablecertificate}
 -\log\alpha-\log(1-\gamma)+g(v)
\end{equation}
upper-bounds the unrestricted likelihood deficit. Stop when
$g(v)\le e_0/4$; the total bound is then below $e_0$.

For completeness a simple fixed-step conditional-gradient method finds
such a $v$ in polynomial time. Along a segment from any feasible $v$
to a simplex vertex, $|L''|\le C_*=4(t+1)^2/\gamma^2$, because every
predictive coordinate has the preceding lower bound and its directional
change is at most one. Start at $v=u$. If $g(v)>e_0/4$, move toward a
maximizing vertex by the fixed rational step $s=e_0/(4C_*)$.
Taylor's theorem gives an increase at least $e_0^2/(32C_*)$.
The initial likelihood is at least $-\log(2(t+1))$, and all likelihoods
are at most zero. Therefore the number of iterations is at most
\[
 1+32C_*\log(2(t+1))/e_0^2
 =O(t^2e_0^{-4}\log(2t)).
\]
All kernel probabilities are rational numbers of bit length
$O(t\log M)$, and every update uses the same rational step. After
polynomially many iterations all weights have polynomial bit length.
The derivatives and the stopping comparison are rational, so no
logarithm evaluation is needed by the algorithm. The displayed
certificate may equivalently be replaced by its rational upper bound
$t(t-1)/M+\gamma/(1-\gamma)+g(v)$.
\end{proof}

\paragraph{Status and scope.}
The sharp deterministic inverse modulus and the certified approximate-MLE
extension are completed targets in the present program. They do not
constitute two newly resolved published conjectures. The published
near-linear likelihood question was addressed in the preceding
manuscript; its extension to the full statistical transition remains
unproved here. The inverse theorem proves a sharp limitation of one
proof strategy, while the optimization theorem removes the need for an
exact optimizer at statistically relevant precision. Classical
beta--binomial spectral theory, Jackson approximation, EM monotonicity
and conditional-gradient optimization are explicitly credited as
ingredients.

% ===== budget.tex =====
\section{Complete adaptive replication thresholds for endpoint populations}
\label{p05:fb:sec:budget}

The scalar uniform-endpoint experiment in the preceding paper has the
critical uncapped rate $(\log B)/\sqrt B$. We now determine its entire
replication-cap curve, and obtain a different, noncritical phase diagram
in every higher dimension. The mixing law may have arbitrary dependence
between its coordinates. The lower bound permits sequential decisions
within a strictly stronger continuous-time observation experiment.

Fix an integer $d\geq1$. An individual has a latent vector
$\theta\in[1,2]^d$ drawn independently from an unknown probability law
$G$. A raw observation is a vector whose coordinates, conditionally on
$\theta$, are independent $\operatorname{Uniform}[0,\theta_j]$ variables.
Raw vectors from the same individual are conditionally independent.
A policy spends at most $B$ raw observations and observes any individual
at most $Q$ times, where $1\leq Q\leq B$. Let $U^{\rm ad}_{B,Q,d}$ be
the minimax expected Euclidean $W_1$ loss over arbitrary adaptive policies
and probability-valued estimators; $U^{\rm fix}_{B,Q,d}$ restricts the
replication counts to a schedule chosen independently of the data.
Constants below may depend on the fixed dimension.

\begin{theorem}[The complete endpoint cap curve]
\label{p05:fb:thm:cap}
Uniformly for integers $B\geq2$ and $1\leq Q\leq B$,
\[
 U^{\rm ad}_{B,Q,1}\asymp U^{\rm fix}_{B,Q,1}
 \asymp \max\left\{\frac{\log(eB)}{\sqrt B},(BQ)^{-1/3}\right\}.
\]
For every fixed $d\geq2$,
\[
 U^{\rm ad}_{B,Q,d}\asymp U^{\rm fix}_{B,Q,d}
 \asymp\max\left\{B^{-1/(d+1)},
                  (B Q^{2d-1})^{-1/(3d)}\right\}.
\]
Consequently, the least cap that attains the uncapped minimax order,
within constants, is
\[
 Q\asymp \frac{\sqrt B}{\log^3(eB)}\quad(d=1),
 \qquad Q\asymp B^{1/(d+1)}\quad(d\geq2).
\]
Adaptive sampling changes none of these orders.
\end{theorem}

This theorem concerns a new classification problem for the specified
experiment, rather than a named conjecture from the literature. Its
$Q=1$ exponent $B^{-1/(3d)}$ is consistent with ordinary-smooth
deconvolution: after taking logarithms, each coordinate is contaminated
by a one-sided exponential error. Classical Wasserstein deconvolution
rates and inversion methods are developed in
\cite{p05:fb:DFM2015,p05:fb:RS2023}. Those results do not optimize replication
counts under a total budget or permit selection of the next individual
from the observed history. The new ingredient here is a uniform
information charge for every adaptive history. Recent multivariate CDF
deconvolution work \cite{p05:fb:Scricciolo2026} distinguishes integrated CDF
loss from $W_1$ for dependent laws. Accordingly, our multivariate proof
uses Lipschitz dual tests, and never identifies these two losses.

\subsection{Frequency-dependent replication gives the upper bound}

Put $a=\log2$ and $U_j=\log\theta_j/a$. The maps between
$[0,1]^d$ and $[1,2]^d$ are bi-Lipschitz, so it suffices to estimate
the law of $U$. After $r$ observations of an individual, the coordinate
log-maxima form a vector
\[
 Y_r=U-(E_{r,1},\ldots,E_{r,d}),
 \qquad E_{r,j}\stackrel{\rm iid}{\sim}\operatorname{Exp}(ar),
\]
with errors independent of $U$.

We use a torus of side length $4$, containing $[0,1]^d$ without a
metric wrap-around, and write $\omega=\pi/2$. For a probability law $H$
on the cube define $m_k(H)=\int e^{i\omega k\cdot u}\,dH(u)$.
For every integer $K\geq2$,
\begin{equation}\label{p05:fb:eq:fourier-modulus}
 W_1(H,H')\leq \frac{C_d}{K}
 +C_d\left\{\sum_{0<\|k\|_\infty\leq K}
             \frac{|m_k(H)-m_k(H')|^2}{\|k\|_2^2}\right\}^{1/2}.
\end{equation}
Indeed, a $1$-Lipschitz function on the cube extends to a $1$-Lipschitz
function on this torus by the McShane formula. Convolve it with the
product of positive one-dimensional Jackson kernels of degree at most
$K$. The uniform approximation error is $C_d/K$, while Parseval bounds
the weighted square sum of its Fourier coefficients by its bounded
Dirichlet integral. Cauchy--Schwarz and Kantorovich duality give
\eqref{p05:fb:eq:fourier-modulus}. The use of the larger torus is essential:
identifying the endpoints of $[0,1]$ itself would destroy the claimed
Euclidean transportation bound.

For each $k$ an unbiased statistic of $Y_r$ is
\[
 A_{r,k}(Y_r)=e^{i\omega k\cdot Y_r}
           \prod_{j=1}^d\left(1+\frac{i\omega k_j}{ar}\right),
 \qquad |A_{r,k}|\leq C_d(1+\|k\|_\infty/r)^d.
\]
Partition frequencies into dyadic bands of upper radius $k$, and use
$r_k=\min(Q,k)$ observations per individual in that band. With $n_k$
fresh individuals, its contribution to the expected squared norm in
\eqref{p05:fb:eq:fourier-modulus} is at most $a_k/n_k$, where
\[
 a_k=C_d k^{d-2}(1+k/r_k)^{2d}.
\]
Let $S=\sum_{k\leq K}\sqrt{r_k a_k}$ and take
$n_k=\lceil B\sqrt{a_k/r_k}/(2S)\rceil$. The cost is at most
$B/2+\sum r_k\leq B/2+C K$, and the total squared error at most
$2S^2/B$. Rounding is harmless for the choices of $K$ below, since
$K\leq\sqrt B$; adjust constants for bounded $B$.
Project the estimated moment vector onto the compact set of moment
vectors of probability laws on the cube, in the displayed weighted
Euclidean norm. A measurable approximate minimizer over a countable
dense set gives a probability-valued estimator and loses only a
factor two and an arbitrarily small additive error.

The geometric sums give
\begin{align}
 S&\leq C\{\log(eK)+\sqrt{K/Q}\}, &&d=1,\label{p05:fb:eq:S1}\\
 S&\leq C_d\{K^{(d-1)/2}
                  +K^{(3d-2)/2}Q^{-(2d-1)/2}\},&&d\geq2.
 \label{p05:fb:eq:Sd}
\end{align}
For $d=1$ choose a dyadic $K$ within constants of
\[
 \min\left\{\frac{\sqrt B}{\log(eB)},(BQ)^{1/3}\right\}.
\]
Then $K^{-1}+S/\sqrt B$ is at most a constant times the claimed
maximum. For $d\geq2$ take $K$ within constants of
\[
 \min\{B^{1/(d+1)},(B Q^{2d-1})^{1/(3d)}\}.
\]
Both inequalities $K^{d+1}\leq B$ and $K^{3d}\leq BQ^{2d-1}$
hold, and \eqref{p05:fb:eq:Sd} gives error at most $C_d/K$.
This proves every upper bound in Theorem~\ref{p05:fb:thm:cap}.

\subsection{A Poisson experiment that dominates every capped policy}

For one coordinate of one individual, expose a Poisson point process
on $[0,\lambda]\times[0,U_j]$ with unit intensity. The second coordinate
of its highest point, with value zero when there is no point, has the
same law as $\max(0,U_j-E)$, $E\sim\operatorname{Exp}(\lambda)$.
An exposure block of length $a$ therefore simulates one scaled
log-uniform observation: when no point occurs, attach independently
a negative $\operatorname{Exp}(a)$ value, and otherwise use the block
maximum. This reconstructed observation is $U_j-\operatorname{Exp}(a)$.
The full point configuration supplies at least as much information.

Allow the statistician to allocate exposure continuously and separately
to coordinates and individuals, with total exposure at most $daB$
and exposure of every coordinate of any individual at most $aQ$.
This stronger experiment simulates the original one, even with its
adaptation. It is enough to prove a lower bound in the stronger experiment.
Conditionally on a fixed adaptive history, an individual's likelihood
for $u\in[0,1]^d$ is proportional to
\begin{equation}\label{p05:fb:eq:posterior-density}
 e^{-\lambda\cdot u}\mathbf1_{\{u\geq b\}},
 \qquad 0\leq\lambda_j\leq aQ,
\end{equation}
where $b_j$ is the largest observed mark in its $j$th coordinate.
The observed choice probabilities have no parameter factor once their
history is fixed. Consequently \eqref{p05:fb:eq:posterior-density} remains
valid under arbitrary adaptive allocations.

\begin{lemma}[Localized information charge]
\label{p05:fb:lem:charge}
Fix a nonzero smooth function $\psi$ supported in $(0,1)$ with
$\int\psi=0$ and $\|\psi\|_\infty\leq1$. Pack $M\asymp_d h^{-d}$
disjoint cubes of side $h$ inside $[1/4,3/4]^d$, and set
\[
 \psi_j(u)=\prod_{\ell=1}^d
                  \psi((u_\ell-c_{j,\ell})/h),\qquad
 w_\beta(u)=1+\varepsilon\sum_{j=1}^M\beta_j\psi_j(u),
 \quad |\beta_j|\leq1,
\]
where $\varepsilon\in(0,1/4)$ is fixed. For every policy in the
stronger Poisson experiment, the trace of its Fisher information in
$\beta$ is at most
\begin{equation}\label{p05:fb:eq:cap-trace}
 C_d B Q^{2d-1}h^{2d}.
\end{equation}
\end{lemma}
\begin{proof}
Write $\Psi(t)=\int_{-\infty}^t\psi(s)\,ds$. The full tensor
antiderivative of $\psi_j$ is
$\Phi_j(u)=h^d\prod_\ell\Psi((u_\ell-c_{j,\ell})/h)$.
It is supported in its cube, and disjointness gives
$\|(\Phi_j(u))_{j=1}^M\|_2\leq C_d h^d$ at every $u$.
For a posterior state $(\lambda,b)$ set
\[
 Z(b)=\int_{u\geq b}w_\beta(u)e^{-\lambda\cdot u}\,du,
 \qquad H_j(b)=\varepsilon\int_{u\geq b}
                        \psi_j(u)e^{-\lambda\cdot u}\,du,
 \qquad s(b)=H(b)/Z(b).
\]
Integration by parts in each coordinate, followed by the triangle
inequality in $\ell_2$, gives
\[
 \|H(b)\|_2\leq C_d h^d e^{-\lambda\cdot b}.
\]
There are at most $2^d$ terms: each consists of a boundary evaluation
or an integral with factor $\lambda_\ell$, and
$\lambda_\ell\int_{b_\ell}^1e^{-\lambda_\ell u}\,du
\leq e^{-\lambda_\ell b_\ell}$. If every $b_\ell\leq3/4$, then
\[
 Z(b)\geq c_d e^{-\lambda\cdot b}
                          \prod_\ell(1+\lambda_\ell)^{-1},
 \qquad
 \|s(b)\|_2\leq C_d h^d\prod_\ell(1+\lambda_\ell).
 \tag{\(*\)}
\]
If any $b_\ell>3/4$, all the $H_j(b)$ vanish, so $s(b)=0$.

While exposing coordinate $i$, the conditional intensity of a new mark
$y>b_i$ is
$T_i(y)=Z(b^{i,y})/Z(b)$, where $b^{i,y}$ replaces $b_i$ by $y$.
Its score vector is $s(b^{i,y})-s(b)$. Marks below $b_i$ have
intensity one and zero score. The conditional Fisher trace per unit
exposure is therefore
\begin{equation}\label{p05:fb:eq:instant-info}
 \int_{b_i}^1 T_i(y)\|s(b^{i,y})-s(b)\|_2^2\,dy.
\end{equation}
For completeness, on a time interval $dt$ the probability of a mark
in $dy$ is $T_i(y)dy\,dt+o(dt)$; the squared score of the no-mark
outcome contributes $o(dt)$. Summing these conditional score
variances, or taking the corresponding compensated counting-process
martingale limit, gives \eqref{p05:fb:eq:instant-info}. Scores and
intensities here are bounded on each fixed finite-dimensional family,
so truncation and dominated convergence justify this calculation.

The posterior density is between fixed constant multiples of the
product of truncated exponential densities. Hence
\[
 \int_{b_i}^1T_i(y)\,dy
 =\mathbb E_\beta(U_i-b_i\mid\text{history})
 \leq\frac{C}{1+\lambda_i}.
\]
Using $(*)$, with zero score beyond the support when necessary, bounds
\eqref{p05:fb:eq:instant-info} by
\[
 C_d h^{2d}(1+\lambda_i)
                     \prod_{\ell\ne i}(1+\lambda_\ell)^2
 \leq C_d h^{2d}Q^{2d-1}.
\]
This is a conditional bound for every history, rather than an average
over predetermined sample sizes. Integrating over total exposure
at most $daB$ and using orthogonality of score increments proves
\eqref{p05:fb:eq:cap-trace}. Independent policy randomization adds no
information and may be included in the observed transcript.
\end{proof}

\subsection{Turning trace information into transportation loss}

We record the elementary absolute-loss information bound used below.
If a joint density $p(\beta,Y)$ has finite positive definite Fisher
matrix $J$ and integrable weak parameter derivatives, then every
estimator $T(Y)$ and vector $v$ satisfy
\begin{equation}\label{p05:fb:eq:abs-fisher}
 \mathbb E|T(Y)-v^T\beta|\geq\frac14\sqrt{v^TJ^{-1}v}.
\end{equation}
Indeed translation by $z$ changes total variation by at most
$\sqrt{z^TJz}/2$. Integrating the triangle inequality against the
minimum of the original and translated densities gives
$2\mathbb E|T-v^T\beta|\geq|v^Tz|(1-\operatorname{TV})$.
Take $z=J^{-1}v/\sqrt{v^TJ^{-1}v}$.

Give each $\beta_j$ independently the density
$\cos^2(\pi\beta_j/2)\mathbf1_{[-1,1]}(\beta_j)$.
Its Fisher information is $\pi^2$, so the joint prior-and-data matrix
satisfies
\[
 \operatorname{tr}J\leq C_d\{M+BQ^{2d-1}h^{2d}\}.
\]
Let $f_j$ be a smooth function supported in the $j$th cube, with
Lipschitz constant bounded independently of $h$, such that
$\int f_j\psi_j=c_\psi h^{d+1}$ for a nonzero fixed $c_\psi$.
For example, take $f_j=h\psi_j$, so
$c_\psi=(\int_0^1\psi(t)^2dt)^d>0$.
Disjoint support and zero boundary values ensure
that every sum $\sum_j\sigma_j f_j$, $|\sigma_j|\leq1$, is
$C_d$-Lipschitz. Thus every probability-valued estimator obeys
\[
 W_1(\widehat G,G_\beta)
 \geq c_d h^{d+1}\sum_{j=1}^M|T_j(\widehat G)-\beta_j|,
 \qquad
 T_j(H)=\frac{\int f_j\,dH-\int f_j(u)\,du}
                  {\varepsilon c_\psi h^{d+1}}.
\]
Apply \eqref{p05:fb:eq:abs-fisher} to each coordinate. Since
$(J^{-1})_{jj}\geq1/J_{jj}$ and
$\sum_jJ_{jj}^{-1/2}\geq M^{3/2}/\sqrt{\operatorname{tr}J}$,
the Bayes, and hence minimax, risk is at least
\[
 \frac{c_d h^{d+1}M^{3/2}}
              {\sqrt{M+BQ^{2d-1}h^{2d}}}.
\]
Choose $h$ within fixed small constants of
$(BQ^{2d-1})^{-1/(3d)}$, and use $M\asymp_d h^{-d}$.
The denominator is $O_d(\sqrt M)$, proving
\begin{equation}\label{p05:fb:eq:cap-lower}
 U^{\rm ad}_{B,Q,d}\geq c_d(BQ^{2d-1})^{-1/(3d)}.
\end{equation}
Every density in this lower-bound family is bounded above and below
by fixed positive constants. No density regularity is assumed in the
minimax problem; the derivatives of these localized bumps increase
as $h$ decreases.

\subsection{The uncapped obstruction and the scalar critical dimension}

For $d\geq2$ the uncapped lower bound $B^{-1/(d+1)}$ has a short
fully adaptive proof. Pack $M\asymp_d h^{-d}$ separated cells in
the interior of the latent cube. In each cell choose one of two point
parameters separated by $h$ in the first coordinate, and give every
cell mass $1/M$. Strengthen the experiment by allowing the statistician
to select the cell for every raw observation. A new individual in the
original experiment is simulated by an independent uniform cell label;
revisits retain the label. The oracle is therefore more informative.

The total variation between the two component observation laws at a
cell is at most $C_dh$. Couple two adaptive transcripts, differing
only at cell $j$, until their first different answer. Conditional on
agreement, each query to cell $j$ disagrees with probability at most
$C_dh$. Consequently their transcript total variation is bounded by
$C_dh\,\mathbb E_\omega N_j$, with expectation under either specified
starting vertex $\omega$. Averaging this bound over all directed
hypercube edges and using $\sum_jN_j\leq B$ gives mean edge variation
at most $C_dBh/M$. On the other hand localized Lipschitz functions
show that the $W_1$ separation between two vertices is at least
$c_dh/M$ times their Hamming distance. Binary testing and nearest-vertex
decoding give risk at least $c_dh$ when $Bh/M$ is a sufficiently
small constant. Taking $h=c_dB^{-1/(d+1)}$ proves the claimed term.
This argument also gives $B^{-1/2}$ in dimension one but misses its
critical logarithm.

The necessary scalar logarithm was proved in the preceding paper.
We include its information argument to specify the exact input.
For $U\in[0,1]$ take
\[
 w_\beta(u)=1+\sum_{k=1}^K\sqrt2\{
       \beta_{k,c}\cos(2\pi ku)+\beta_{k,s}\sin(2\pi ku)\},
 \quad |\beta_{k,*}|\leq cK^{-2}.
\]
These densities and their Lipschitz constants are uniformly bounded,
and the densities are bounded away from zero. Under baseline density
one, a posterior after $r$ raw observations is a truncated exponential
$q(u)\propto e^{-ar u}\mathbf1_{[b,1]}(u)$. If $q_Y$ is the posterior
after the next observation, direct integration gives
\[
 \mathbb E_0\|q_Y-q\|_2^2=a/2,
 \qquad \|q\|_2^2\operatorname{Var}_q(U)\leq C.
\]
For clarity, writing $\lambda=ar$, $D=1-b$, the current squared norm
is $\lambda\coth(\lambda D/2)/2$; the expectation of the next one is
this quantity plus $a/2$. The identity then follows from
$\mathbb E_0q_Y=q$, including the limit $\lambda=0$.
If $z=\int w_\beta q$ and $z_Y=\int w_\beta q_Y$, the one-step score
is the vector of Fourier coefficients of $q_Y/z_Y-q/z$.
Bessel's inequality, $z,z_Y\asymp1$, and
\[
 \mathbb E_0|z_Y-z|^2
 \leq\operatorname{Lip}(w_\beta)^2
            \mathbb E_0 W_1(q_Y,q)^2
 \leq2\operatorname{Lip}(w_\beta)^2\operatorname{Var}_q(U)
\]
bound its information trace by an absolute constant for every history.
Thus any adaptive $B$-query experiment has trace at most $CB$.

A product cosine-squared prior of width $cK^{-2}$ contributes $CK^5$.
For the centered CDF, the coefficient vector is
\[
 v_{k,c}(x)=\frac{\sqrt2\sin(2\pi kx)}{2\pi k},\qquad
 v_{k,s}(x)=-\frac{\sqrt2\cos(2\pi kx)}{2\pi k}.
\]
The integrated absolute error of the centered CDF is at most twice
$W_1$. Applying \eqref{p05:fb:eq:abs-fisher} at every $x$ gives a lower
bound by a constant times $\int_0^1\sqrt{v(x)^TJ^{-1}v(x)}dx$.
This functional of $J$ is convex: for fixed $v$ it is
$\sup_{z\ne0}|v^Tz|/\sqrt{z^TJz}$. Average $J$ over the block
rotations induced by translating $x$. The averaged matrix has blocks
$j_kI_2$, and Cauchy--Schwarz yields
\[
 \int_0^1\sqrt{v(x)^TJ^{-1}v(x)}dx
 \geq\frac{c\sum_{k=1}^K k^{-1}}{\sqrt{\operatorname{tr}J}}
 \geq\frac{c\log(eK)}{\sqrt{B+K^5}}.
\]
Taking $K\asymp B^{1/5}$ proves $c\log(eB)/\sqrt B$ even without
a cap, and hence for every cap. Together with
\eqref{p05:fb:eq:cap-lower} and the preceding upper bounds this completes
Theorem~\ref{p05:fb:thm:cap}.

\section{The complete adaptive Bernoulli replication-cap curve}
\label{p05:fb:sec:bernoulli}

The preceding paper determined the Bernoulli cap curve for predetermined
designs and for adaptive designs at polynomially subcritical caps. Its
remaining transition window is closed by a different least-favourable
family: localized moment-matching perturbations are immersed in a global
uniform background. The background controls posterior normalizers even
after unusually favourable sequences of tosses. This supplies the
conditional information estimate that a fixed-prefix divergence bound
alone cannot imply.

Let $R^{\rm ad}_{B,Q}$ and $R^{\rm fix}_{B,Q}$ denote the minimax expected
$W_1$ risk for estimating an arbitrary law of $\theta\in[0,1]$ from
conditionally independent Bernoulli$(\theta)$ observations, with total
budget $B$ and per-individual cap $Q$.

\begin{theorem}[No adaptive improvement at any Bernoulli cap]
\label{p05:fb:thm:bernoulli-cap}
For absolute constants, uniformly over integers $B\geq2$, $1\leq Q\leq B$,
\[
 R^{\rm ad}_{B,Q}\asymp R^{\rm fix}_{B,Q}\asymp
 \Psi(B,Q):=\max\left\{B^{-1/3},Q^{-1},
          [Q\log(e+B/Q^{3/2})]^{-1/2}\right\}.
\]
In particular, at $Q\asymp B^{2/3}/\log B$ the adaptive minimax
risk has order
$B^{-1/3}\sqrt{\log B/\log\log B}$.
The uncapped order $B^{-1/3}$ is attainable within constants if and
only if $Q\gtrsim B^{2/3}$; the implicit cap constant may depend on
the allowed risk constant.
\end{theorem}

We give the new adaptive lower bound first. All constants in its proof
are absolute and can be chosen in the stated order.

\begin{lemma}[A background-stabilized conditional information bound]
\label{p05:fb:lem:bern-charge}
There are constants $c_0,C>0$ and $\rho\in(0,1)$ such that the following
holds. Let $1\leq K<Q$, put $h^2=c_0K/Q$, and suppose $h\leq1/100$.
Choose $M\asymp h^{-1}$ intervals $I_j=[u_j-h,u_j+h]$ with centers
in $[3/8,5/8]$ separated by at least $4h$.
There exist probability laws $\nu_{j,+},\nu_{j,-}$ on $I_j$, matching
moments through degree $K$, with
$W_1(\nu_{j,+},\nu_{j,-})=2h/(K+1)$.
Set $\delta_j=\nu_{j,+}-\nu_{j,-}$,
$\mu_j=(\nu_{j,+}+\nu_{j,-})/2$, and
\begin{equation}\label{p05:fb:eq:bern-family}
 G_\beta=\tfrac12\operatorname{Uniform}[1/4,3/4]
       +\frac1{2M}\sum_{j=1}^M(\mu_j+\tfrac12\beta_j\delta_j),
 \qquad |\beta_j|\leq1.
\end{equation}
For any adaptive policy with total budget $B$ and cap $Q$, the trace
of its Fisher information matrix in $\beta$ is at most
\begin{equation}\label{p05:fb:eq:bern-trace}
 C B h^2\rho^K.
\end{equation}
\end{lemma}
\begin{proof}
Split
$2^{-K}\sum_{i=0}^{K+1}(-1)^i\binom{K+1}{i}\delta_{i/(K+1)}$
into positive and negative probability parts and rescale to $I_j$.
Finite differences prove moment matching. Integrating the absolute
cumulative signed mass, and using
$\sum_{i=0}^l(-1)^i\binom{K+1}{i}=(-1)^l\binom Kl$,
gives the stated $W_1$ distance.
Every law in \eqref{p05:fb:eq:bern-family} obeys, for every interval $I$,
\begin{equation}\label{p05:fb:eq:interval-mass}
 G_\beta(I)\leq C(|I|+h),
 \qquad G_\beta\geq\tfrac12\operatorname{Uniform}[1/4,3/4].
\end{equation}
The first bound follows by counting the $4h$-separated cells meeting
$I$; each cell has total mass $1/(2M)\asymp h$.

Consider an individual with $s$ successes in $r$ previous tosses, where
$0\leq r\leq Q-1$. Write
\[
 L(x)=x^s(1-x)^{r-s},\quad m=\int L\,dG_\beta,\quad
 p=\frac{\int xL(x)\,dG_\beta(x)}m.
\]
Its next outcome is Bernoulli$(p)$, with $p\in[1/4,3/4]$.
Conditioning on the entire adaptive history adds no further factor
to this individual's posterior once $\beta$ is fixed. Therefore its
conditional Fisher trace is at most
\begin{equation}\label{p05:fb:eq:bernoulli-gradient}
 C\sum_{j=1}^M|\partial_{\beta_j}p|^2,
 \qquad
 \partial_{\beta_j}p=\frac1{4Mm}\int(x-p)L(x)\,d\delta_j(x).
\end{equation}
For $r<K$ these derivatives vanish by moment matching.

Suppose $r\geq K$, set $v=r+1$, and fix one center $u=u_j$.
Let $A=(\log L)'(u)=(s-ru)/(u(1-u))$. Uniformly in all the
possible histories,
\begin{align}
 m&\geq c\,v^{-1/2}L(u)\exp(cA^2/v),
       \label{p05:fb:eq:normalizer}\\
 |p-u|&\leq C\{|A|/v+h+v^{-1/2}\}.
       \label{p05:fb:eq:post-mean}
\end{align}
Here are details that in particular cover extreme success counts.
On $[1/4,3/4]$, $(\log L)''\geq-Cr$ and $|A|\leq Cr$.
For a sufficiently large absolute $C_1$, integrate the lower Taylor
bound
$\log L(u+t)\geq\log L(u)+At-C_1vt^2$
over an interval of length $c/\sqrt v$ centered at
$t=A/(2C_1v)$. Taking $C_1$ large keeps this interval inside
$[1/4-u,3/4-u]$. Completing the square and using the uniform
background proves \eqref{p05:fb:eq:normalizer}.

If $z=s/r\in[1/4,3/4]$, then
$L(x)\leq L(z)e^{-cr(x-z)^2}$ on that interval, while the background
gives $m\geq cL(z)/\sqrt v$ by one-sided integration when needed.
The interval mass bound \eqref{p05:fb:eq:interval-mass}, summed over
annuli of width $v^{-1/2}$ around $z$, gives
\[
 \int |x-z|L(x)\,dG_\beta(x)
 \leq CL(z)(v^{-1}+hv^{-1/2}).
\]
Thus $|p-z|\leq C(v^{-1/2}+h)$, and
$|u-z|\leq C|A|/v$. If $z$ lies outside $[1/4,3/4]$, its distance
from $u\in[3/8,5/8]$ is at least $1/8$, so $|A|/v$ is bounded
below by a positive absolute constant. Then
\eqref{p05:fb:eq:post-mean} follows from $p,u\in[1/4,3/4]$ directly.

Fix a small absolute $\kappa>0$ and let
$R=\kappa\sqrt{(K+1)/v}\leq\kappa$.
For complex $z$ with $|z|\leq R$, Taylor expansion of the analytic
logarithms around $u$ gives
\[
 |L(u+z)|\leq L(u)\exp\{|A|R+C rR^2\}.
\]
Choose $c_0$ sufficiently small, after $\kappa$, so that $h/R<1/4$.
Cauchy's estimate for the Taylor remainder of the polynomial
$F(x)=(x-p)L(x)$, and annihilation of its terms through degree $K$,
give
\[
 \left|\int F\,d\delta_j\right|
 \leq C L(u)(|u-p|+R)e^{|A|R+C rR^2}(h/R)^{K+1}.
\]
Combining this with \eqref{p05:fb:eq:normalizer}--\eqref{p05:fb:eq:post-mean},
$M^{-1}\asymp h$, and $h\sqrt v\leq\sqrt{c_0K}$ gives
\[
 |\partial_{\beta_j}p|
 \leq Ch\sqrt{K+1}\,e^{-cA^2/v}
                \{\eta\sqrt{v/Q}\}^{K+1},
 \tag{\(**\)}
\]
where $\eta$ can be made any prescribed positive small constant by
reducing $c_0$. To check this last absorption, set $z_0=|A|/\sqrt v$.
Before absorption the extra factor is bounded by
$(z_0+\sqrt{K+1})
 e^{-c z_0^2+\kappa\sqrt{K+1}z_0+C\kappa^2(K+1)}$.
Young's inequality bounds it by
$C\sqrt{K+1}e^{-c' z_0^2}e^{C'(K+1)}$;
the last exponential is absorbed into
$(h/R)^{K+1}\leq\{\sqrt{c_0}\sqrt{v/Q}/\kappa\}^{K+1}$.
All constants precede the final choice of $c_0$.

The centers are $4h$-separated, and
$A_j=(s-ru_j)/(u_j(1-u_j))$. A Gaussian sum therefore gives
\[
 \sum_j e^{-cA_j^2/v}\leq C(1+(h\sqrt v)^{-1}).
\]
Indeed $|A_j|/\sqrt v\geq c\sqrt v\,|u_j-s/r|$ since $r\geq1$.
Let $x=v/Q\leq1$. Squaring $(**)$ and summing shows that
\eqref{p05:fb:eq:bernoulli-gradient} is at most
\[
 Ch^2(K+1)\eta^{2(K+1)}x^{K+1}
       \{1+(c_0Kx)^{-1/2}\}
 \leq Ch^2\rho^K,
\]
for some fixed $\rho<1$, after taking $\eta<1/4$ and absorbing the
polynomial factors in $K$. This bound holds for every adaptive history
and every $\beta$ in the parameter cube. Orthogonality of conditional
score increments through at most $B$ tosses proves
\eqref{p05:fb:eq:bern-trace}. Policy randomization has zero parameter score.
\end{proof}

\begin{proof}[Adaptive lower bound in Theorem~\ref{p05:fb:thm:bernoulli-cap}]
Use the family in Lemma~\ref{p05:fb:lem:bern-charge}, and independent
cosine-squared priors on the coordinates $\beta_j\in[-1,1]$.
The joint Fisher matrix has
\[
 \operatorname{tr}J\leq C\{M+Bh^2\rho^K\}.
\]
For each cell choose a $1$-Lipschitz Kantorovich potential for
$\delta_j$. Subtract its linear interpolant between the two cell
endpoints. Its integral against $\delta_j$ is unchanged because the
zeroth and first moments vanish; its Lipschitz constant is at most two
and its endpoint values are zero. Extend by zero outside the cell and
call the resulting function $f_j$. Then
$\int f_j\,d\delta_j=2h/(K+1)$ and all signed sums of the $f_j$
are $2$-Lipschitz. Defining the affine estimator coordinate $T_j$ from
$\int f_j\,d\widehat G$ as in the endpoint argument gives
\[
 W_1(\widehat G,G_\beta)
 \geq\frac{ch}{M(K+1)}\sum_{j=1}^M|T_j-\beta_j|.
\]
The absolute-loss inequality \eqref{p05:fb:eq:abs-fisher} and the trace
bound thus give risk at least
\begin{equation}\label{p05:fb:eq:bern-bayes}
 \frac{ch}{K+1}\left(1+Bh^3\rho^K\right)^{-1/2}.
\end{equation}

Set $L=\log(e+B/Q^{3/2})$ and $K=\lceil A L\rceil$, where $A$ is
a sufficiently large absolute constant. When $Q\geq C_2 L$, choose
$C_2$ large enough for $K<Q$ and $h\leq1/100$.
Since $h^2=c_0K/Q$,
\[
 Bh^3\rho^K=Cc_0^{3/2}(B/Q^{3/2})K^{3/2}\rho^K\leq C,
\]
uniformly, by increasing $A$. Equation \eqref{p05:fb:eq:bern-bayes}
therefore gives $c[QL]^{-1/2}$.
When $Q<C_2L$, this term is dominated, within a constant, by $Q^{-1}$.

The remaining two lower bounds also allow adaptation. Two probability
laws matching moments through degree $Q$ and separated in $W_1$ by
$1/(Q+1)$ have identical full $Q$-toss laws. An oracle of independent
full $Q$-toss vectors simulates every capped policy, proving the
$Q^{-1}$ floor. For $B^{-1/3}$, pack $M\asymp h^{-1}$ cells with
two point biases at distance $h$ in each, inside $[1/4,3/4]$, and
reveal or even allow selection of cell labels. Neighbouring hypercube
transcripts have KL at most $Ch^2\mathbb E N_j$. Averaging over edges,
using $\sum_jN_j\leq B$, and applying Pinsker gives mean edge variation
at most $C\sqrt{Bh^2/M}$. With $h=cB^{-1/3}$ it is below $1/2$,
and the separated-cell transportation loss yields risk at least $ch$.
Combining the three bounds proves the adaptive lower bound.
\end{proof}

\begin{proof}[Predetermined upper bound in Theorem~\ref{p05:fb:thm:bernoulli-cap}]
We recall the frequency allocation argument to make the relation to
the lower bound explicit. Let $\widetilde T_m(x)=T_m(2x-1)$.
The Chebyshev-moment inequality
\[
 W_1(G,H)\leq C/K+
 C\left\{\sum_{m=1}^K m^{-2}
          \left|\int\widetilde T_m\,d(G-H)\right|^2\right\}^{1/2}
\]
follows by applying the positive Jackson kernel and Parseval to
$f((1+\cos u)/2)$; see also \cite{p01:info:MMRS2025}.
For $m\leq r/2$, the degree-$r$ Bernstein coefficients of
$\widetilde T_m$ are bounded by $Ce^{2m^2/r}$, and hence define an
unbiased statistic of a binomial$(r,\theta)$ observation with that bound.
For clarity, on the coefficient contour $z=te^{iu}$,
$t=j/(r-j)$, put $v=2\sqrt t\,|\sin(u/2)|/(1+t)$. The generating
polynomial is
\[
 (1+z)^r\widetilde T_m\!\left(\frac z{1+z}\right)
 =\tfrac12(1+z)^{r-m}\{(\sqrt z+i)^{2m}+(\sqrt z-i)^{2m}\}.
\]
Its normalized modulus is bounded by
$(1-v^2)^{(r-m)/2}(1+v)^m
\leq e^{2m^2/r}e^{-rv^2/8}$.
The Gaussian contour integral contributes
$C(1+rt/(1+t)^2)^{-1/2}$, canceled by the reciprocal binomial
probability in Stirling's formula. The endpoint coefficients are one
in absolute value. This proves the required coefficient bound.

For the dyadic band ending at $k\leq K\leq Q/2$, use
$r_k=\min(Q,\max(2,k^2))$ trials per fresh individual. Its weighted
moment variance is at most $a_k/n_k$, where
$a_k=Ck^{-1}e^{4k^2/r_k}$.
The allocation used earlier in this section gives
$\sum\sqrt{r_ka_k}\leq C\sqrt K e^{2K^2/Q}$ and total cost at
most $B$ when $K^2\leq cB$. Moment projection therefore gives
\[
 \sup_G\mathbb EW_1(\widehat G,G)
 \leq C/K+C\sqrt{K/B}\,e^{2K^2/Q}.
\]
Choose a dyadic $K$ within a sufficiently small absolute factor of
$\min\{B^{1/3},Q,\sqrt{Q L}\}$, with $L$ as above.
If $Q\geq B^{2/3}$ the exponential is bounded. Otherwise put
$x=B/Q^{3/2}>1$; the stochastic term divided by $1/K$ is at most
$C\log(e+x)^{3/4}(e+x)^\eta/\sqrt x$ with any prescribed
small $\eta>0$, by making the factor defining $K$ small enough.
It is bounded for $\eta<1/4$. This proves the upper bound, with
bounded small parameters covered by the diameter bound.
\end{proof}

The new Bernoulli and endpoint classifications use different observation
kernels but share a principle: the lower bound must control the information
of the next adaptively selected query conditionally on its entire history.
In the endpoint model tensor antiderivatives provide this control. In
the Bernoulli model it is the uniform background, the complex likelihood
remainder, and the Gaussian sum over localized cells. A fixed-prefix
divergence estimate without this conditional step would not justify
either conclusion about adaptive designs.

% ===== bessel.tex =====
\section{Power restoration of hyperbolic complete monotonicity}
\label{p05:fs:section}

The preceding manuscripts classified the GGC and HCM properties of the
original McKay families and then proved their full ID domains.  Those
results are not counted again here.  The new question is whether a
power transformation can restore HCM after the untransformed density
has failed it.  This section proves an exact exponent boundary for
restoration under exponential tilting, gives a constructive sufficient
bound for a fixed law, and extends the strict SD threshold into part
of the previously unresolved negative-cut range.

A positive function $f$ is HCM if, for each $u>0$, the function
$f(uv)f(u/v)$ is completely monotone in $v+v^{-1}$.  Multiplication by
any real monomial preserves HCM.  We use the classical facts that an
HCM function with value one at zero is the Laplace transform of a GGC,
and hence has a Stieltjes logarithmic derivative, and that every HCM
density is GGC.  Equivalently, the normalized HCM functions considered
below have the representation
\begin{equation}\label{p05:fs:canonical}
 g(x)=\exp\left\{-d_0x-\int_0^\infty
             \log(1+x/t)\,U(dt)\right\},\qquad d_0\geq0,\quad U\geq0.
\end{equation}
The constructed functions with $p<1/2$ below have $d_0=0$.
See \cite[Chapter~5, especially Theorem~5.4.1]{p05:fs:Bondesson} and
\cite{p01:rare:BondessonSimon2016}.  Sufficiency in our result can also be checked
directly from finite products $(1+x/t)^{-c}$ and passage to their
pointwise limits.

For $\mu>-1$ write
\[
 \mathcal I_\mu(z)=2^\mu\Gamma(\mu+1)z^{-\mu}I_\mu(z)
       ={}_0F_1(;\mu+1;z^2/4),\qquad
 R_\mu(z)=I_{\mu+1}(z)/I_\mu(z).
\]
The first expression is understood by its entire even continuation.
Thus $\mathcal I_\mu(0)=1$.  Let $j_{\mu,n}$ be the positive zeros of
$J_\mu$.  The classical product and its logarithmic derivative give
\begin{equation}\label{p05:fs:product}
 \mathcal I_\mu(z)=\prod_{n\geq1}
           \left(1+z^2/j_{\mu,n}^2\right),\qquad
 R_\mu(z)=\sum_{n\geq1}\frac{2z}{z^2+j_{\mu,n}^2}.
\end{equation}
The series is locally uniform away from its poles.  For $\mu>-1$ all
nonzero zeros of $J_\mu$ are real and simple; see
\cite[Sections~10.21(i), 10.21(iii)]{p01:rare:DLMF}.

\subsection{A positive spectral measure after changing the power}

\begin{lemma}[A ray estimate from a tilted beta law]
\label{p05:fs:ray}
For $\mu\geq1/2$, $y>0$, and $0<\theta<\pi/2$,
\begin{equation}\label{p05:fs:ray-bound}
 0<\frac{\operatorname{Im}\{e^{i\theta}
                       R_\mu(ye^{i\theta})\}}{\sin\theta}
 \leq\sec\theta.
\end{equation}
For every fixed $\mu>-1$ and $\theta$ in the same range the middle
quantity is positive, continuous in $y$, bounded, tends to zero at
zero, and tends to one at infinity.
\end{lemma}
\begin{proof}
The product formula gives, term by term,
\begin{align*}
 \frac{\operatorname{Im}\{e^{i\theta}R_\mu(ye^{i\theta})\}}
      {\sin\theta}
 &=4y\cos\theta\sum_{n\geq1}
  \frac{j_{\mu,n}^2}
  {j_{\mu,n}^4+2y^2j_{\mu,n}^2\cos(2\theta)+y^4}\\
 &\leq\frac1{\cos\theta}
       \{R_\mu(y)+yR_\mu'(y)\}.
\end{align*}
Indeed the denominator in each summand is at least
$\cos^2\theta(j_{\mu,n}^2+y^2)^2$.

We give an elementary proof that $R_\mu(y)+yR_\mu'(y)\leq1$ for
$\mu\geq1/2$.  Put $c=\mu-1/2\geq0$.  The Bessel beta integral
represents $R_\mu(y)$ as the expectation of $T$ under the probability
density proportional to $e^{yt}(1-t^2)^c$ on $(-1,1)$.
Consequently $y\{1-R_\mu(y)\}$ is the mean of $Z=y(1-T)$, whose
density is proportional to
\[
 z^c(2-z/y)^c e^{-z}\,\mathbf1_{(0,2y)}(z).
\]
If $y_2>y_1$, the ratio of these densities on their common support,
apart from a positive constant, is
$\{(2-z/y_2)/(2-z/y_1)\}^c$.  It is nondecreasing in $z$, and the
second density has the larger right endpoint.  The elementary
single-crossing criterion for likelihood-ratio order therefore shows
that $Z$ is stochastically nondecreasing in $y$.  Its mean is
nondecreasing. Differentiating that mean gives the desired inequality.
The case $c=0$ is the same truncated-exponential ordering.

For arbitrary $\mu>-1$, positivity follows from the displayed sum.
There are no poles on the fixed ray.  The expansion at zero gives a
limit of zero, and the standard Bessel asymptotic in every closed
subsector of $|\arg z|<\pi/2$ gives $R_\mu(z)=1+O(1/z)$ at infinity.
Thus the stated limit is one, and continuity proves boundedness.
\end{proof}

\begin{theorem}[An exact exponent boundary for HCM restoration]
\label{p05:fs:restoration}
Let $\mu>-1$, $a>0$, $p>0$, and $\eta>0$.  Consider the normalized
density
\begin{equation}\label{p05:fs:powered-density}
 f_{p,b}(x)=C x^{\eta-1}e^{-b x^p}
                           \mathcal I_\mu(a x^p),\qquad x>0,
 \qquad b>a.
\end{equation}
There exists a $b>a$ for which this density is HCM if and only if
\begin{equation}\label{p05:fs:power-boundary}
                         0<p<\tfrac12.
\end{equation}
For $0<p<1/2$, put $\theta=\pi p$ and
\begin{equation}\label{p05:fs:c-threshold}
 c_\mu(p)=\sup_{y>0}
  \frac{\operatorname{Im}\{e^{i\theta}R_\mu(ye^{i\theta})\}}
                         {\sin\theta}.
\end{equation}
Then $1\leq c_\mu(p)<\infty$, and for every $b>a$ the exact test is
\[
 f_{p,b}\hbox{ is HCM}\quad\Longleftrightarrow\quad b/a\geq c_\mu(p).
\]
In particular, for $\mu\geq1/2$ the elementary sufficient bound is
\begin{equation}\label{p05:fs:explicit-bound}
                         b/a\geq\sec(\pi p).
\end{equation}
The upper bound in \eqref{p05:fs:explicit-bound} is not asserted to be the
optimal tilt threshold.
\end{theorem}
\begin{proof}
The density is integrable: at zero its non-power factor tends to one,
and at infinity its exponential decay is $e^{-(b-a)x^p}$ times a
power.  Define $g(x)=e^{-b x^p}\mathcal I_\mu(a x^p)$, so $g(0)=1$.
Since powers preserve HCM, the density is HCM exactly when $g$ is.

Assume first $0<p<1/2$.  On the slit plane
$D=\mathbb C\setminus(-\infty,0]$, with principal powers, $g$ is
holomorphic and nonzero by \eqref{p05:fs:product}.  Its negative
logarithmic derivative is
\begin{equation}\label{p05:fs:h}
 h(z)=-\frac{g'(z)}{g(z)}
      =p z^{p-1}\{b-aR_\mu(a z^p)\}.
\end{equation}
The upper-bank boundary value yields the real density
\begin{equation}\label{p05:fs:U-density}
 u_{p,b}(t)=-\frac1\pi\operatorname{Im}h(-t+i0)
   =\frac{p t^{p-1}}\pi\left[
    b\sin\theta-a\operatorname{Im}
       \{e^{i\theta}R_\mu(a t^p e^{i\theta})\}\right].
\end{equation}
A keyhole application of Cauchy's formula gives the signed identity
\begin{equation}\label{p05:fs:h-keyhole}
             h(z)=\int_0^\infty\frac{u_{p,b}(t)}{z+t}\,dt.
\end{equation}
Here no residue is omitted.  At zero, $h(z)=O(|z|^{p-1})$;
at infinity, uniformly up to both banks,
$h(z)=p(b-a)z^{p-1}+O(1/z)$.
The circle terms in the Cauchy formula vanish respectively as
$O(\varepsilon^p)$ and $O(R^{p-1})$.
The density is $O(t^{p-1})$ at both endpoints, so the signed integral
is absolutely convergent for $z\in D$.

The measure $U(dt)=u_{p,b}(t)\,dt$ is nonnegative exactly when
$b/a\geq c_\mu(p)$.  In that event integration of
\eqref{p05:fs:h-keyhole} from zero to $x$, justified by the just stated
endpoint bounds, proves \eqref{p05:fs:canonical}.  This gives HCM,
either by its classical representation or by positive finite-product
approximation.  Conversely, if $g$ is HCM, its normalization at zero
makes it a GGC Laplace transform.  Its logarithmic derivative is a
Stieltjes function.  Stieltjes inversion in
\eqref{p05:fs:h-keyhole} forces $u_{p,b}\geq0$; its continuity rules out
a negative value even at a single interior point.  This proves the
exact test.  Lemma~\ref{p05:fs:ray} proves finiteness, the lower bound
one, and \eqref{p05:fs:explicit-bound}.  Choosing
$b/a>c_\mu(p)$ supplies the claimed existence for every $\mu>-1$.

At the boundary $p=1/2$, the same product formula gives
\begin{equation}\label{p05:fs:half-obstruction}
 h(z)=\frac{b}{2\sqrt z}
       -\sum_{n\geq1}\frac{a^2}{a^2z+j_{\mu,n}^2}.
\end{equation}
At each $z=-j_{\mu,n}^2/a^2$ the residue of the second term is
$-1$, while $b/(2\sqrt z)$ has a finite one-sided boundary value
there.  A Stieltjes transform of a nonnegative measure can only have
nonnegative residues at negative-axis poles.  More explicitly,
$\lim_{\varepsilon\downarrow0}i\varepsilon
h(-j_{\mu,n}^2/a^2+i\varepsilon)=-1$, whereas that limit must be the
nonnegative mass of the corresponding atom.  Thus $g$ cannot be HCM.

For $p>1/2$, the point
\[
 z_n=(j_{\mu,n}/a)^{1/p}\exp\{i\pi/(2p)\}
\]
lies strictly inside $D$ and is a zero of the analytic continuation
of $g$.  A normalized HCM function is a GGC Laplace transform and
has a zero-free continuation to $D$, contradicting this zero.
This excludes every $b$ and completes the equivalence.
\end{proof}

\begin{corollary}[A constructive power for a fixed McKay law]
\label{p05:fs:mckay-power}
Suppose $X$ has density
\[C x^{\nu-1}e^{-bx}I_\mu(ax),\qquad x>0,\]
with $\mu>-1$, $\mu+\nu>0$,
and $b>a>0$.  No positive power $X^d$ with $d\leq2$ has an HCM
density.  If $\mu\geq1/2$, every power satisfying
\begin{equation}\label{p05:fs:power-sufficient}
                 d\geq\frac{\pi}{\arccos(a/b)}
\end{equation}
has an HCM density, and consequently is GGC, SD, and ID.
The assertion for a negative power $d$ is identical with
that for $|d|$.
\end{corollary}
\begin{proof}
The change of variables has the form \eqref{p05:fs:powered-density} with
$p=1/d$ and $\eta=p(\mu+\nu)>0$.
The explicit sufficient bound follows from
Theorem~\ref{p05:fs:restoration}.
For negative powers, the reciprocal-density transformation multiplies
$f(1/x)$ by $x^{-2}$; inversion and monomials preserve HCM.
\end{proof}

\subsection{Which fixed laws can be repaired by any power?}

\begin{theorem}[The fixed-tilt restoration criterion and its order boundary]
\label{p05:fs:fixed-tilt}
For $\mu>-1$ set
\[
 M_\mu=\sup_{y>0}\{R_\mu(y)+yR_\mu'(y)\}.
\]
Then $1\leq M_\mu<\infty$, and
\begin{equation}\label{p05:fs:M-boundary}
 M_\mu=1\quad\Longleftrightarrow\quad\mu\geq1/2.
\end{equation}
If $-1<\mu<1/2$, the strictly larger maximum is attained.
For a generalized McKay random variable $X$ with fixed admissible
parameters, the following are equivalent:
\begin{enumerate}
\item Some positive power $X^d$ has an HCM density.
\item Every sufficiently large positive power $X^d$ has an HCM density.
\item $b/a>M_\mu$.
\end{enumerate}
A constructive sufficient power, valid for every order, is
\begin{equation}\label{p05:fs:allorder-power}
       d\geq\frac{\pi}{\arccos(aM_\mu/b)}
              \qquad\hbox{when } b/a>M_\mu.
\end{equation}
Consequently $\mu\geq1/2$ is exactly the set of orders for which
\emph{every} admissible exponential tilt can be restored by a power.
\end{theorem}
\begin{proof}
Let $D_\mu(y)=R_\mu(y)+yR_\mu'(y)$.  Formula
\eqref{p05:fs:product} shows $D_\mu(y)>0$ and $D_\mu(y)\to0$ at zero.
The Bessel asymptotic, differentiated in a real subsector, gives
\begin{equation}\label{p05:fs:D-asymptotic}
 D_\mu(y)=1-\frac{4\mu^2-1}{8y^2}+O(y^{-3}).
\end{equation}
Thus $M_\mu$ is finite and at least one.
Lemma~\ref{p05:fs:ray}'s real estimate gives $M_\mu=1$ when
$\mu\geq1/2$.  If $-1/2<\mu<1/2$, \eqref{p05:fs:D-asymptotic}
implies $D_\mu(y)>1$ at large $y$.  If $-1<\mu\leq-1/2$,
the function $H(y)=y(1-R_\mu(y))$ is positive at sufficiently small
positive $y$ and tends to $\mu+1/2\leq0$ at infinity.  Its derivative
must be negative somewhere; since $H'=1-D_\mu$, again $M_\mu>1$.
The endpoint limits of $D_\mu$ imply attainment when $M_\mu>1$.

The ray expression in Lemma~\ref{p05:fs:ray} satisfies the two-sided
pointwise estimate
\[
 \cos\theta\,D_\mu(y)
 \leq\frac{\operatorname{Im}\{e^{i\theta}R_\mu(ye^{i\theta})\}}
              {\sin\theta}
 \leq\sec\theta\,D_\mu(y).
\]
The lower estimate uses the opposite elementary bound
$j^4+2j^2y^2\cos(2\theta)+y^4\leq(j^2+y^2)^2$.
In particular,
\begin{equation}\label{p05:fs:c-limit}
 \cos(\pi p)M_\mu\leq c_\mu(p)\leq\sec(\pi p)M_\mu,
 \qquad \lim_{p\downarrow0}c_\mu(p)=M_\mu.
\end{equation}
If $b/a>M_\mu$, choosing $p=1/d$ as in
\eqref{p05:fs:allorder-power} proves HCM by
Theorem~\ref{p05:fs:restoration}, for every sufficiently large $d$.

For necessity suppose a power is HCM.  Its normalized non-power
factor $g(x)=e^{-bx^p}\mathcal I_\mu(ax^p)$ is then a nonconstant
HCM function with $g(0)=1$.  The classical GGC representation,
allowing a nonnegative drift $d_0$, gives
\[
 \frac{d}{dx}\left(x\frac{g'(x)}{g(x)}\right)
   =-d_0-\int_0^\infty\frac{t}{(x+t)^2}\,U(dt)<0.
\]
Strictness holds because $g$ is nonconstant and its normalization
at zero excludes an atom of the Thorin measure at zero.
Direct differentiation, with $y=ax^p$, makes the same derivative
\[
       \frac{p^2y}{x}\{-b/a+D_\mu(y)\}.
\]
Therefore $b/a>D_\mu(y)$ for every $y>0$.
If $\mu<1/2$, the maximum is attained, forcing $b/a>M_\mu$.
If $\mu\geq1/2$, admissibility already gives $b/a>1=M_\mu$.
This proves all equivalences including their equality cases.
\end{proof}

For example $R_{-1/2}(y)=\tanh y$.
The derivative of $D_{-1/2}(y)=\tanh y+y\operatorname{sech}^2y$ is
$2\operatorname{sech}^2y(1-y\tanh y)$.  Since $y\tanh y$ increases
strictly from zero to infinity, there is a unique $y_*>0$ satisfying
$y_*\tanh y_*=1$, and at that point $D_{-1/2}(y_*)=y_*$.
Thus $M_{-1/2}=y_*$, approximately $1.19968$.
For this order, no power restores HCM when $1<b/a\leq y_*$,
whereas some power does when $b/a>y_*$.  This is an explicit
one-equation transition for a family whose original density never
has the HCM property.

\begin{corollary}[The cost of approaching the square-power barrier]
\label{p05:fs:critical-cost}
For fixed $\mu>-1$ and $0<p<1/2$,
\[
 \frac{1}{j_{\mu,1}\cos(\pi p)}
       \leq c_\mu(p)\leq\frac{M_\mu}{\cos(\pi p)}.
\]
More precisely,
\begin{equation}\label{p05:fs:critical-constant}
 \lim_{p\uparrow1/2}\cos(\pi p)c_\mu(p)=\frac1{j_{\mu,1}}.
\end{equation}
If $r=b/a>M_\mu$, let $d_*(r)$ be the infimum of the powers
$d>0$ for which the corresponding McKay law has an HCM density.
For all sufficiently large $r$,
\[
 \frac{\pi}{\arccos(1/(rj_{\mu,1}))}
   \leq d_*(r)\leq\frac{\pi}{\arccos(M_\mu/r)}.
\]
In particular the first Bessel zero determines the leading cost:
\begin{equation}\label{p05:fs:critical-power}
 d_*(r)=2+\frac{4}{\pi j_{\mu,1}r}+o(r^{-1}),\qquad r\to\infty.
\end{equation}
\end{corollary}
\begin{proof}
In the positive series of Lemma~\ref{p05:fs:ray}, set $y=j_{\mu,1}$
and retain the first term.  It equals
$1/(j_{\mu,1}\cos\theta)$ exactly.
The upper bound is \eqref{p05:fs:c-limit}.  These bounds give the
displayed arccos bounds for $d_*(r)$.

We supply the uniform estimate needed for the sharper constant.
Write $\varepsilon=\cos\theta$ and abbreviate $j_n=j_{\mu,n}$.
The exact ray formula becomes
\begin{equation}\label{p05:fs:resonance}
 \varepsilon\frac{\operatorname{Im}
                 \{e^{i\theta}R_\mu(ye^{i\theta})\}}{\sin\theta}
 =\frac1y\sum_{n\geq1}
 \frac{4y^2j_n^2\varepsilon^2}
      {(j_n^2-y^2)^2+4y^2j_n^2\varepsilon^2}.
\end{equation}
The Bessel zero asymptotic $j_n=\pi(n+\mu/2-1/4)+O(n^{-1})$
ensures a positive minimum gap $\delta_\mu$ and the counting bound
$\#\{n:j_n\leq t\}\leq C_\mu(1+t)$.
For $y\geq1$, split the sum at $y/2$ and $2y$.
The ranges $j_n<y/2$ and $j_n>2y$, after the factor $1/y$, contribute
at most $C_\mu\varepsilon^2$: use respectively
$C\varepsilon^2j_n^2/y^2$ and
$C\varepsilon^2y^2/j_n^2$ to bound the summand.
In the middle range, retain one closest zero at cost at most one.
The other terms have successive distances at least a constant times
$\delta_\mu k$, and their sum is bounded by
\[
 C\sum_{k\geq1}\frac1{1+(\delta_\mu k/(C y\varepsilon))^2}
       \leq C_\mu y\varepsilon.
\]
Consequently the right side of \eqref{p05:fs:resonance} is at most
$1/y+C_\mu\varepsilon$ for $y\geq1$.
On every bounded $y$ interval, the summands tend to zero uniformly
away from the finitely many zeros.  Near $y=j_n$ the corresponding
summand is at most one and the other summands vanish uniformly;
the prefactor tends to $1/j_n$.  The same uniform convergence to
zero near $y=0$ follows directly from
$\sum_nj_n^{-2}<\infty$.
Thus the supremum on a fixed bounded interval containing $j_1$
tends to $1/j_1$.  The just proved tail bound, followed by taking
its starting point to infinity, proves \eqref{p05:fs:critical-constant}.

Finally, as $d\downarrow2$,
$\cos(\pi/d)\sim\pi(d-2)/4$.
The exact HCM test $r\geq c_\mu(1/d)$ and
\eqref{p05:fs:critical-constant} imply
$c_\mu(1/d)\sim4/[\pi j_1(d-2)]$.
For every fixed $\epsilon>0$, the lower estimate
$c_\mu(1/d)\geq(4/(\pi j_1)-\epsilon)/(d-2)$ holds for every
sufficiently small $d-2>0$, excluding all smaller powers at once.
The corresponding upper estimate supplies a feasible power
$d=2+(4/(\pi j_1)+2\epsilon)/r$.  Together with the earlier upper
bound, which already gives $d_*(r)\downarrow2$, these estimates prove
\eqref{p05:fs:critical-power} without a monotonicity assumption.
\end{proof}

\begin{corollary}[Positive products and arbitrary real Bessel powers]
\label{p05:fs:products}
Let $m\geq1$, $\mu_j\geq1/2$, $a_j>0$, $\kappa_j>0$, $\eta>0$,
and $0<p<1/2$.  If
\[
                 b\geq\sec(\pi p)\sum_{j=1}^m\kappa_j a_j,
\]
then the normalized density
\[
 Cx^{\eta-1}e^{-b x^p}
             \prod_{j=1}^m\mathcal I_{\mu_j}(a_jx^p)^{\kappa_j}
\]
is HCM, hence GGC, SD, and ID.
\end{corollary}
\begin{proof}
The negative logarithmic derivative of the non-power factor is the
sum of the corresponding functions \eqref{p05:fs:h}, with boundary
measure
\[
 \frac{pt^{p-1}}\pi\left[b\sin(\pi p)-
 \sum_j\kappa_j a_j\operatorname{Im}
   \{e^{i\pi p}R_{\mu_j}(a_jt^p e^{i\pi p})\}\right]dt.
\]
Lemma~\ref{p05:fs:ray} makes this nonnegative.  Integrating its Stieltjes
representation proves \eqref{p05:fs:canonical} again.  Integrability
follows from $\eta>0$ and
$b>\sum_j\kappa_ja_j$, which the stated condition implies.
\end{proof}

\subsection{A strict SD boundary in the negative-cut intervals}

Write $\alpha=\mu+1/2$, $q=\nu-\mu-1>0$, $A=\alpha+q$, and
$M=2\alpha+q$.  Admissibility means $\alpha>-1/2$ and $M>0$;
in particular $A>0$.  The preceding manuscript's signed spectral
construction supplies a positive smooth even kernel $K$ for which
\begin{equation}\label{p05:fs:old-K}
 -J'(u)/J(u)=\int_0^\infty e^{-ux}K(x)\,dx,\quad
 K(0)=M,\quad e^{-x}K(x)\longrightarrow A,
 \qquad u>1,
\end{equation}
where
\[
 J(u)=C\frac{u^q V(u^{-2})}{(u^2-1)^A},\qquad
 V(z)={}_2F_1(-q/2,(1-q)/2;\alpha+1/2;z).
\]
Moreover $K'/K\to1$ and
$\beta_{\alpha,q}=\sup_{x\geq0}K'(x)/K(x)<\infty$ is the exact
SD boundary: the law is SD precisely when $b/a\geq\beta_{\alpha,q}$.
These background kernel identities, established in
\cite[Section~4]{comp:part:04}, are not new claims of this section.
The old strictness proof covered positive-cut intervals and integer
$q$.  The following new argument also works in the negative-cut
intervals, without having to locate a trigonometric oscillation.

\begin{proposition}[A regular-part obstruction to threshold one]
\label{p05:fs:sd-strict}
For every admissible $q>0$ and $-1/2<\alpha\leq1$,
\[
                         \beta_{\alpha,q}>1.
\]
The supremum is attained at a finite positive point.
Thus, for every $-1<\mu\leq1/2$ and every admissible excess power
$q>0$, there is a nonempty interval of admissible tilts $b>a$ for
which the law is ID but not SD.
\end{proposition}
\begin{proof}
Suppose instead $\beta_{\alpha,q}=1$.  Then
$h(x)=e^{-x}K(x)$ is nonincreasing and tends to $A$, so
$K(x)\geq Ae^x$ for every $x\geq0$.
If $\alpha<0$, the inequality at zero contradicts
$K(0)-A=\alpha$.  If $\alpha=0$, a nonincreasing function with
$h(0)=h(\infty)=A$ must be constant; but
$h'(0)=K'(0)-K(0)=-M<0$.

Let $0<\alpha\leq1$ and first suppose $A>1$.
Gauss's value and the derivative identity give
\[
 \frac{V'(1)}{V(1)}=\frac{q(q-1)}{4(A-1)}.
\]
Consequently, direct logarithmic differentiation yields
\begin{align}
 \lim_{u\downarrow1}\left\{-\frac{J'(u)}{J(u)}
                                      -\frac{A}{u-1}\right\}
 &=\frac A2-q+\frac{q(q-1)}{2(A-1)}\notag\\
 &=\frac{\alpha(\alpha-1)}{2(A-1)}\leq0.
 \label{p05:fs:regular-part}
\end{align}
On the other hand the expression in braces is
$\int_0^\infty e^{-ux}\{K(x)-Ae^x\}\,dx$.
Its integrand would be nonnegative and strictly positive near zero,
because $K(0)-A=\alpha>0$.
Monotone convergence as $u\downarrow1$ therefore makes its limit
strictly positive, a contradiction, including at $\alpha=1$.

The only uncovered case has $0<\alpha<1$ and $A\leq1$, hence
$0<q<1$.  Here there are no negative spectral zeros and the signed
measure is positive.  For completeness, the background kernel is
\[
 K(x)=2A\cosh x-2\int_0^1\cosh(\sqrt t\,x)\,\rho(dt),
\]
where $\rho$ has positive mass on every compact subinterval of $(0,1)$.
In $K'-K$, the integral contributes
$2\int(\cosh(\sqrt t\,x)-\sqrt t\sinh(\sqrt t\,x))\rho(dt)$;
this grows exponentially along any such subinterval, whereas the
remaining contribution is $-2Ae^{-x}$.  Thus $K'-K>0$ for large $x$,
again a contradiction.  This last case reproduces the previously
known positive-cut argument only to cover the statement's endpoints.
Finally, $K'/K$ tends to one and is continuous; any value strictly
above one ensures that its supremum is attained on a compact interval.
\end{proof}

\paragraph{Scope of the progress.}
The power-restoration problem is formulated here as a new transfer
problem.  Its existence boundary $p<1/2$ and fixed-tilt criterion
$b/a>M_\mu$ are complete, including the failure at equality.  The
explicit sufficient exponent need not be the smallest possible
exponent for a fixed law.  The SD proposition is a partial advance on
the external classification problem in
\cite[Section~4.3]{p01:rare:Baricz2026}; it does not close strictness for all
$\mu>1/2$ in the negative-cut intervals, nor give an elementary
formula for the full SD threshold.  The earlier full ID and GGC
classifications are prerequisites, not additional closures in this
round.

Power transformations as a tool for HCM are classical, including
Bosch's work on stable quotients~\cite{p05:fs:Bosch}; the contribution here
is the Bessel-specific phase boundary and fixed-tilt classification.
Sj\"odin's recent preprint~\cite{p05:fs:Sjodin} concerns preservation of
GGC under powers larger than one.  That conclusion is distinct from
HCM restoration, which is stronger and here also applies to original
McKay laws lying outside the GGC domain.

% ===== gaussian.tex =====
\section{Random covariance ranges and the exact number of Gaussian views}
\label{p05:fg:sec:rank}

Low covariance rank reduces the required number of repeated observations,
but uncertainty about the covariance range has an irreducible cost.
We work throughout with arbitrary Borel mixing laws: the number of mixture
components is not assumed finite. Conditional on $(A,\Sigma)$, let
$X_i$ be independent $N_d(A,\Sigma)$ vectors, allowing singular covariance
matrices. A conditional block of length $m>0$ means a vector with law
$N_d(mA,m\Sigma)$; disjoint blocks are conditionally independent.
Integer block lengths are the usual partial-sum experiments.

\begin{theorem}[Exact rank-adaptive view threshold]
\label{p05:fg:thm:rank}
Fix $1\le r<d$ and restrict the mixing law to
$\operatorname{rank}\Sigma\le r$.
\begin{enumerate}
\item One experiment consisting of $r+1$ conditionally independent
Gaussian blocks of any known positive lengths identifies the full joint
law of $(A,\Sigma)$. No moment assumptions are required.
In particular $(X_1,\ldots,X_r,S_n)$ identifies it for every $n\ge r+1$.
\item Even the population laws of \emph{all} experiments with at most
$r$ blocks, allowing arbitrary positive lengths, do not identify the
mixing law. Two distinct witnesses can have compact support, covariance
rank exactly $r$, equal mean marginals, equal covariance marginals, and
equal joint polynomial moments through any prescribed finite degree.
Their nonzero covariance eigenvalues can be bounded above and bounded
away from zero.
\end{enumerate}
Consequently $r$ anchors are necessary and sufficient in this class.
The failure with $r-1$ anchors persists even when every terminal length
is available and the covariance support is bounded.
\end{theorem}

Here anchors are the separately observed $X_i$ preceding the terminal
sum. Thus the theorem counts $r+1$ independent blocks, not $r+1$
independent draws of the mixing parameter. Its negative assertion is
stronger than failure at one terminal length.

\begin{lemma}[Covariance mixtures of bounded rank]
\label{p05:fg:lem:covariance}
For finite complex measures supported on
$\{\Sigma\succeq0:\operatorname{rank}\Sigma\le r\}$, the kernel
\[
 \Sigma\longmapsto
 \mathcal L(Y_1,\ldots,Y_r\mid\Sigma),
 \qquad Y_j\mid\Sigma\stackrel{\mathrm{iid}}{\sim}N_d(0,\Sigma),
\]
is injective.
\end{lemma}
\begin{proof}
Write $Y=[Y_1\ \cdots\ Y_r]$. Conditional on a covariance of rank
$k\le r$, one has $\operatorname{range}Y=\operatorname{range}\Sigma$
almost surely. Indeed in orthonormal coordinates on that range, $Y$
is an invertible square root of the covariance times a $k\times r$
standard Gaussian matrix, whose row rank is $k$ almost surely.

This observable range takes values in the finite disjoint union of
Grassmannians. Disintegrate the total variation of a candidate signed
or complex measure over that range, and disintegrate the measure itself
using its bounded Radon--Nikodym derivative. If its mixture kernel is
zero, the corresponding conditional mixture on almost every range
$E$ is zero. This follows by testing against products of an arbitrary
bounded measurable function of the observable range and a member of
a countable determining class of functions of $Y$.

Choose a measurable orthonormal frame on each Grassmannian chart.
On a fixed $k$-dimensional range the covariance becomes a positive
definite $k\times k$ matrix $B$. The characteristic function of its
$r$ Gaussian vectors is
\[
 \int \exp\{-\tfrac12\operatorname{tr}(WW^TB)\}\,d\nu_E(B),
 \qquad W\in\mathbb R^{k\times r}.
\]
As $r\ge k$, the matrices $WW^T/2$ fill the positive semidefinite
cone. Uniqueness of the matrix Laplace transform gives $\nu_E=0$.
For completeness, zero on the open positive cone implies zero in
its complex tube by successive one-variable identity theorems;
evaluation at $I+iV$ is then Fourier uniqueness for the finite
measure $e^{-\operatorname{tr}B}\nu_E$. The zero-dimensional range
is immediate. Integrating the zero conditional measures proves the
lemma.
\end{proof}

\begin{proof}[Proof of the positive assertion of
Theorem~\ref{p05:fg:thm:rank}]
Let the block lengths be $m_1,\ldots,m_{r+1}$ and set $n=\sum m_j$.
Divide block $j$ by $\sqrt{m_j}$, and apply an orthogonal change of
coordinates in the block index with first row
$(\sqrt{m_1},\ldots,\sqrt{m_{r+1}})/\sqrt n$.
The resulting vectors are conditionally independent, with
\[
 Z\mid(A,\Sigma)\sim N_d(\sqrt n A,\Sigma),\qquad
 Y_1,\ldots,Y_r\mid(A,\Sigma)
       \stackrel{\mathrm{iid}}{\sim}N_d(0,\Sigma).
\]
For the difference $\rho$ of two mixing laws giving the same
experiment, Fourier transformation in $Z/\sqrt n$ gives, for every
$u\in\mathbb R^d$, the zero covariance mixture associated with
\[
 \nu_u(d\Sigma)=
 \int e^{iu^Ta-u^T\Sigma u/(2n)}\,d\rho(a,\Sigma).
\]
Lemma~\ref{p05:fg:lem:covariance} implies $\nu_u=0$. Multiplying this
zero measure by $e^{u^T\Sigma u/(2n)}$, first on
$\operatorname{tr}\Sigma\le M$ and then passing to the limit, shows
that $\int e^{iu^Ta}\,d\rho(a,d\Sigma)=0$. Fourier uniqueness in
$a$, applied after restriction to each covariance Borel set, now gives
$\rho=0$. One may first use rational $u$ and a countable generating
class of covariance sets, then use continuity, to avoid any issue with
uncountably many null sets.
\end{proof}

\subsection{An exterior differential annihilator}

The obstruction is easiest to see before quotienting the covariance
factorization. The general implication from a parameter differential
annihilator to nonidentifiability is established by Le, Bariletto,
Rinaldo, and Ho~\cite{p01:mix:Le2026}; the preceding integrated manuscript
also used differential annihilators. The new points here are a single
exterior operator valid for every allowed block design, its nonzero
descent through the covariance quotient, and the matching positive
rank threshold. Write $\Sigma=VV^T$ with
$V=[v_1\ \cdots\ v_r]\in\mathbb R^{d\times r}$ of full column rank.
Choose $r+1$ distinct row indices $I=(i_0,\ldots,i_r)$, and form the
constant-coefficient differential operator
\begin{equation}\label{p05:fg:eq:operator}
 D_I=\det\bigl[\partial_{a_i}\ 
               \partial_{(v_1)_i}\ \cdots\
               \partial_{(v_r)_i}\bigr]_{i\in I}.
\end{equation}
The determinant is an ordinary polynomial in commuting derivatives.
For
\[
 K_{u,T}(a,V)=e^{iu^Ta-\operatorname{tr}(V^TTV)},
 \qquad T\succeq0,
\]
direct differentiation gives
\begin{equation}\label{p05:fg:eq:annihilation}
 D_I K_{u,T}
 =i(-2)^r\det[u_I\ (Tv_1)_I\ \cdots\ (Tv_r)_I]K_{u,T}.
\end{equation}
There are no mixed Gaussian derivative terms: each column derivative
acts on a different $v_j$.

\begin{proof}[Proof of the negative assertion of
Theorem~\ref{p05:fg:thm:rank}]
For $q\le r$ blocks with lengths $m_j>0$ and Fourier arguments $w_j$,
the parameter transform is $K_{u,T}$ with
\[
 u=\sum_{j=1}^q m_jw_j,
 \qquad T=\tfrac12\sum_{j=1}^q m_jw_jw_j^T.
\]
Thus $\operatorname{rank}T\le q\le r$ and $u\in\operatorname{range}T$.
All $r+1$ columns in the determinant in
\eqref{p05:fg:eq:annihilation} lie in a subspace of dimension at most
$r$, so $D_IK_{u,T}=0$ for every such observation.

Choose a nonzero real $C_c^\infty$ function $\chi(a,V)$ supported on
full-rank $V$ and invariant under $V\mapsto VO$, $O\in O(r)$.
For instance, average a nonnegative nonzero compact bump over that
compact group. Its support can be chosen inside
$cI\prec V^TV\prec CI$ for any fixed $0<c<C$.
Let
\[
 g=D_I^{2M}\chi,\qquad M\ge1.
\]
The polynomial symbol of $D_I^{2M}$ is not zero. A nonzero compactly
supported smooth function cannot be annihilated by a nonzero
constant-coefficient differential operator: its Fourier transform
would vanish on the open set where the polynomial symbol is nonzero,
and hence everywhere by real analyticity. Therefore $g\ne0$.

Under right multiplication by $O$, $D_I$ transforms by
$\det O$. Its even power consequently preserves $O(r)$ invariance,
so $g(a,V)$ is constant on the fibers of
$\pi(a,V)=(a,VV^T)$. Those fibers are exactly the right $O(r)$
orbits for full-rank $V$. It follows that the signed pushforward
$\rho=\pi_*(g\,da\,dV)$ is nonzero: its sign is a measurable
function of $(a,VV^T)$, so integrating that sign against $\rho$
gives $\int|g|>0$. This step prevents a nonzero factor-space
perturbation from disappearing on the covariance quotient.

Integration by parts and $D_I K_{u,T}=0$ show that every selected
observation transform of $\rho$ is zero. Its total mass is zero.
Write its positive and negative parts as $\rho_+$ and $\rho_-$;
their common positive mass is $m$, and $\rho_+/m,\rho_-/m$ are the
required distinct mixing laws. Alternatively push forward $g_+/m$
and $g_-/m$; invariance makes these the same Jordan decomposition.
They arise from smooth compact densities in the covariance factor
coordinates and are supported on the smooth rank-$r$ covariance
stratum.

Every term of $D_I^{2M}$ contains $2M$ derivatives in $a$ and
$2M$ derivatives in each column $v_j$. Hence integration over all
$a$ gives zero identically in $V$, and integration over all $V$
gives zero identically in $a$. The covariance and mean marginals
therefore agree. If $2M>L$, integration against a polynomial in
$(a,VV^T)$ of total degree at most $L$ is also zero, because its
degree in $a$ is smaller than $2M$. This gives the asserted finite
moment matching. Compactness and the chosen bounds on $V^TV$
give all remaining support assertions.
\end{proof}

\begin{remark}[Why the full-rank endpoint is different]
The determinant \eqref{p05:fg:eq:operator} requires $r+1\le d$.
For full rank $r=d$, the preceding manuscript showed that $d-1$
anchors can suffice when infinitely many terminal lengths are used;
without a covariance tail restriction the exact condition is
$\sum_{n\in\mathcal N}n^{-1}=\infty$. The current theorem proves
that substituting the intrinsic rank $r$ for $d$ in that result is
false when the covariance range is itself random. If the range is a
fixed known $r$-plane, reduction to that plane is valid instead.
Thus unknown orientation, not simply the number of nonzero eigenvalues,
accounts for the extra anchor.
\end{remark}

\subsection{A hierarchy beyond the first logarithmic boundary}

The previous logarithmic tail theorem has a sharp scalar converse.
Its endpoint is not the weakest possible sufficient tail envelope:
one can repeatedly insert another iterated logarithm. We give the
complete finite hierarchy, including joint observation counterexamples.
Fix an integer $m\ge1$. Choose $E$ so large that every logarithm below
is at least one, and put
\[
 L_1(x)=\log(E+x),\qquad L_{j+1}(x)=\log L_j(x).
\]
Define the envelope
\[
 W_{m,a}(v)=
 \frac{\sqrt v}{L_1(v)\cdots L_{m-1}(v)L_m(v)^a},\quad a\ge0.
\]
The empty product is one. Changing $E$ affects none of the assertions.

\begin{theorem}[Every finite iterated-log boundary]
\label{p05:fg:thm:hierarchy}
Consider the separate population laws of
$(X_1,\ldots,X_{d-1},S_n)$ for an infinite terminal set
$\mathcal N\subset\{d,d+1,\ldots\}$.
\begin{enumerate}
\item For any $d\ge1$, any $m\ge1$, and $0\le a\le1$, these laws
identify the joint mixing law within the class
\[
 \mathbb E\exp\{cW_{m,a}(\operatorname{tr}\Sigma)\}<\infty
 \quad\hbox{for some }c>0.
\]
There is no moment condition on $A$.
\item For $d=1$, $\mathcal N=\{2^j:j\ge0\}$, and each $a>1$,
identification fails. Distinct absolutely continuous witnesses can
have equal coordinate marginals and all equal joint polynomial
moments, while each satisfies the displayed exponential integrability
for \emph{every} $c>0$. A fixed positive lower bound on the variance
can be imposed.
\end{enumerate}
\end{theorem}

The moment-determinacy input is classical; general quasi-analytic
weight criteria already occur in de Jeu~\cite{p05:fg:deJeu}.
The assertion specific to the present experiment is the simultaneous
construction of equal-marginal joint mixing laws annihilating the
entire prescribed sequence of observation transforms.

\begin{lemma}[Iterated-log moment estimate]
\label{p05:fg:lem:moments}
If $V\ge0$ and $\mathbb E e^{cW_{m,a}(V)}<\infty$ for some $c>0$,
then
\[
 (\mathbb E V^k)^{1/(2k)}
 \le Ck L_1(k)\cdots L_{m-1}(k)L_m(k)^a,
 \qquad k\ge1.
\]
In particular the Stieltjes--Carleman series diverges for $a\le1$.
\end{lemma}
\begin{proof}
Set $D(x)=L_1(x)\cdots L_{m-1}(x)L_m(x)^a$ and
$b_k=kD(k)$. Iterated logarithms of $y^2$ and of $y$ differ by
bounded multiplicative factors, so the exponential moment implies
$\mathbb E e^{c_1\sqrt V/D(\sqrt V)}<\infty$.
For all $w\ge1$, elementary repeated logarithm inequalities give
\[
 D(b_kw)\le C D(k)(1+\log w)^{m-1+a}.
\]
For clarity, $L_j(b_k)\asymp L_j(k)$ at every fixed depth, since
$\log b_k=\log k+O(\log\log k)$; and
$L_j(b_kw)\le C L_j(b_k)(1+\log w)$ follows successively from
$\log(xw)\le\log x+\log w$. Therefore
\[
 \sup_{y\ge0}y^{2k}e^{-c_1y/D(y)}\le (Cb_k)^{2k},
\]
by substituting $y=b_kw$ and bounding
$w^2\exp\{-c_2w/(1+\log w)^{m-1+a}\}$.
The moment assertion follows. The integral test, with successive
logarithmic substitutions, gives divergence of
$\sum[k L_1(k)\cdots L_{m-1}(k)L_m(k)^a]^{-1}$ exactly for
$a\le1$.
\end{proof}

\begin{proof}[Proof of the positive part of
Theorem~\ref{p05:fg:thm:hierarchy}]
Let $\rho$ be the difference of two candidate laws. The Fourier
arguments of the independent blocks of lengths
$1,\ldots,1,n-d+1$ have the form
\[
 (u,T),\qquad
 T=C+\frac{uu^T}{2n},\quad C\succeq0,\quad\operatorname{rank}C\le d-1.
\]
This follows by projecting their $d$ frequency columns onto the
vector of square roots of the block lengths.
Fix $u\ne0$ and a rank-$d-1$ matrix $C$ with
$u\not\perp\ker C$. Push forward
$e^{iu^Ta-\operatorname{tr}(C\Sigma)}\rho$ under
$(a,\Sigma)\mapsto v=u^T\Sigma u$. Its total variation is dominated
by the pushforward of the sum of the candidate probability laws.
The preceding lemma, unchanged by fixed positive rescaling of $v$,
gives a Stieltjes--Carleman moment envelope.

Its Laplace transform vanishes at $1/(2n)$ for the selected $n$.
These zeros accumulate at zero. The transform is $C^\infty$ there;
successive Taylor expansions show that all its derivatives vanish.
Thus the complex measure has zero moments. If $\mu$ is its positive
dominating measure, $\mu\pm\operatorname{Re}\nu$ and
$\mu\pm\operatorname{Im}\nu$ have the same moments as $\mu$;
Stieltjes--Carleman uniqueness forces $\nu=0$.

It follows that the full mixing transform vanishes at
$T=C+zuu^T$ for all $z>0$. Every $T\succ0$ has this form, with
$z=(u^TT^{-1}u)^{-1}$ and $C=T-zuu^T$. Continuity includes
$u=0$, and matrix Laplace--Fourier uniqueness concludes the proof.
\end{proof}

\begin{lemma}[Lacunary zeros at any finite iterated-log depth]
\label{p05:fg:lem:iterated-construction}
Let $b_j=b_0 2^{-j}$, $b_0>0$. For every $m\ge1$ and $q>1$ there
is a nonzero real Schwartz function $f$, supported on $[0,\infty)$,
whose Laplace transform $F$ satisfies $F(b_j)=0$ for all $j$,
whose moments all vanish, and such that for each fixed $\ell\ge0$
\[
 |f^{(\ell)}(v)|\le C_\ell e^{-c_\ell W_{m,q}(v)},\qquad v\ge1.
\]
\end{lemma}
\begin{proof}
Put $p=2(q-1)>0$. For a sufficiently large fixed $K$, define nested
holomorphic logarithms in the cut plane by
\[
 J_1(w)=\log(K+w),\qquad
 J_{j+1}(w)=\log(K+J_j(w)),\qquad
 Q(w)=\frac{w}{J_m(w)^p}.
\]
All branches are principal. In the upper cut half-plane each $J_j$
has positive imaginary part, and conjugation handles the lower half;
there are no unintended zeros or cuts. On the positive axis all
$J_j$ are positive. Increasing $K$ if necessary gives
$0\le p\arg J_m(w)\le\arg w$ when
$0\le\arg w\le\pi/2$: each logarithm has real part at least
$\log K$ on that half-plane, and its argument is at most its input
argument divided by $\log K$. Hence $\operatorname{Re}Q(w)\ge0$
on the right half-plane.

Define
\[
 B(z)=\prod_{j\ge0}\frac{z-b_j}{z+b_j},\qquad
 F(z)=B(z)\exp\{-\sqrt z-Q(1/z)\}.
\]
The geometric product is bounded uniformly on every fixed sector
$|\arg z|\le\theta_0<\pi$, and by one on the right half-plane.
Indeed split the factors at $b_j\le|z|/2$, $b_j\ge2|z|$,
and the bounded number of intervening terms. The two tail sums of
the positive logarithms of the moduli are bounded by geometric sums
of $b_j/|z|$ and $|z|/b_j$.

We give the quantitative estimate needed for inversion. Write
$r=e^{-L}$ and, for large $L$, put
\[
 \lambda_1=L,\quad\lambda_{j+1}=\log\lambda_j,
 \quad P_m(L)=\lambda_1\cdots\lambda_m,
 \quad\delta(r)=\frac{p\pi}{8P_m(L)}.
\]
Uniformly for bounded real $\theta$, repeated expansion of the
nested logarithms gives
\[
 |J_m(e^{L+i\theta})|\sim\lambda_m,
 \qquad
 \arg J_m(e^{L+i\theta})
   =\frac{\theta}{P_m(L)}(1+o(1)).
\]
The second relation and its first $\theta$ derivative hold uniformly
on bounded intervals. This can also be checked inductively: the
imaginary part of the next logarithm is the previous imaginary part
divided by its real part, with a relative error tending to zero;
the fixed additive $K$ changes only lower-order terms.
Consequently for $|\arg z|\le\pi/2+\delta(r)$ and $|z|=r$
small enough,
\begin{equation}\label{p05:fg:eq:sector-estimate}
 \operatorname{Re}Q(1/z)
 \ge \frac{c}{r P_m(L)\lambda_m^p}.
\end{equation}
At the angular endpoints the argument of $Q$ is at most
$\pi/2-3p\pi/(8P_m(L))+o(P_m(L)^{-1})$ in modulus;
inside them the argument is monotone for large $L$. The modulus
of $Q$ is comparable to $1/(r\lambda_m^p)$, proving the bound.

The estimate and Cauchy bounds in fixed sectors show that $F(iy)$
extends to a Schwartz function flat at $y=0$; every fixed derivative
is controlled because the exponential in
\eqref{p05:fg:eq:sector-estimate} dominates all inverse powers of $r$
and all logarithmic factors. At infinity $e^{-\sqrt z}$ gives
rapid decay. Define $f$ by inverse Fourier transformation of
$F(iy)$. It is real and Schwartz. For $x>0$,
\[
 |F(x+iy)|\le e^{-\sqrt{|y|/2}},
\]
so vertical contour shifts give
$f_x(v)=e^{-xv}f(v)$ with $f_x$ uniformly bounded. Letting
$x\to\infty$ proves $f(v)=0$ for $v<0$ and proves the Laplace
identity. Flatness at zero gives all zero moments.

For the tail, bend the inverse contour near zero onto
$z=r\exp\{\pm i(\pi/2+\delta(r))\}$, then continue it to infinity
on fixed rays with arguments between $\pi/2$ and $\pi$ in modulus.
No negative pole or branch cut is crossed. Small arcs disappear by
\eqref{p05:fg:eq:sector-estimate}; large arcs disappear by the square
root exponential. Compact joining portions are chosen strictly in
the left half-plane and give $O(e^{-cv})$.
The small-radius integrand for the inverse of $z^\ell F(z)$ is
bounded by a constant times
\[
 r^\ell\exp\left\{-c\left(
    \frac{vr}{P_m(L)}+
    \frac{1}{rP_m(L)\lambda_m^p}\right)\right\}\,dr.
\]
On $\tfrac14\log v\le L\le\log v$, arithmetic--geometric
mean gives a lower bound of a constant times
\[
 \frac{\sqrt v}
 {\left(\prod_{j=1}^{m-1}\log_j v\right)
                    (\log_m v)^{1+p/2}},
\]
with the evident interpretation for $m=1$. Below this range the
first exponential term is at least $v^{3/4}$ divided by a fixed
product of logarithms; above it the second term is eventually
increasing and is at least $v$ divided by such a product. Both
dominate the same lower bound. Reserving half the exponent for an
integrable majorant proves the stated estimate, with $q=1+p/2$,
for every fixed $\ell$. Nonzeroness follows from that of $F$.
\end{proof}

\begin{proof}[Proof of the negative part of
Theorem~\ref{p05:fg:thm:hierarchy}]
Choose $1<q<a$, use the preceding lemma with $b_j=2^{-j-1}$,
and let $h\in C_c^\infty(\mathbb R\setminus\{0\})$ be real,
even, and nonzero. Define
\[
 g(a,v)=\frac1{2\pi}\int e^{-iua}h(u)u^2 f(u^2v)\,du.
\]
Its Fourier--Laplace transform is $h(u)F(z/u^2)$. At the
observation argument $z=nu_0^2/2$ and $u=nu_0$, this is zero
because $z/u^2=1/(2n)$. The case $u=0$ follows from the support
of $h$. Thus every selected sum has zero signed observation law.

Integration by parts in $u$ and the derivative bounds for $f$ give
\[
 |g(a,v)|\le C_N(1+|a|)^{-N}e^{-c_NW_{m,q}(v)}
 \quad(v\ge1)
\]
for every fixed $N$. The function is real, nonzero, and Schwartz.
Its $v$ marginal and all joint polynomial moments vanish because
its Fourier transform in $a$ vanishes near $u=0$; its $a$ marginal
vanishes because $\int f=0$. The normalized positive and negative
parts therefore give distinct probability densities with all the
required equalities. Since
$W_{m,q}(v)/W_{m,a}(v)=L_m(v)^{a-q}\to\infty$, their tails
satisfy every asserted exponential moment. A common translation
$v\mapsto v+v_0$ preserves all transform zeros and all equalities.
\end{proof}

\subsection{Scope, comparison, and an unresolved matrix converse}

Theorem~\ref{p05:fg:thm:rank} concerns arbitrary mixing measures with
bounded component covariance rank. It is distinct from finite-component
grouped-mixture identifiability, where the number of components controls
the required group size; see Vandermeulen and Scott~\cite{p02:frontmix:VS2015}.
The present sufficiency uses the observed random covariance range,
and the necessity uses an exterior differential operator that survives
the orthogonal factorization quotient. The literature search did not
locate this precise classification, but is not an exhaustive priority
certificate. Both targets here are internally formulated questions,
not newly located named published conjectures.

The converse question about absence of finite differential barriers
in~\cite{p01:mix:Le2026} was already answered negatively in the 223-page
integrated manuscript, including by a jointly entire positive kernel.
It is not counted again as a new closure in this section.

The fixed-dimensional matrix sharp-tail converse remains open here:
for $d>1$ and $d-1$ anchors we do not construct nonidentifiable laws
at every supercritical $W_{m,a}$ envelope. The old matrix-tube
annihilator composes a scalar zero product with $u^TT^{-1}u$.
Near indefinite imaginary matrix arguments, cancellation in that
quadratic form prevents the scalar bent-contour estimate from transferring
directly. Schwartz decay of the matrix inverse does not imply the
required nearly square-root exponential tail. The compact low-rank
witness above uses fewer than $d-1$ anchors and therefore does not
resolve or contradict the matrix sufficiency theorem.

% ===== zeros.tex =====
\section{Finite spectral recovery from modular-density thresholds}
\label{p05:fz:sec}

The preceding manuscript recovered a finite positive spectrum from an
entire transition curve.  Here we replace that curve by finitely many
scalar measurements, prove the exact sample count for unnormalized
spectra, and determine the optimal noise exponent without any separation
between the unknown frequencies.  These are two newly formulated inverse
problems, not two previously published conjectures.  The exponent in the
noise theorem is the familiar finite-moment exponent from mixture
estimation; see Wu and Yang~\cite{p05:fz:WuYang}.  The additional issue resolved
here is that the observations are nonlinear minima with unknown
minimizing points, rather than moments or values of an exponential sum.

Let
\begin{equation}\label{p05:fz:eq:profile}
 S_\mu(x)=\int_{(0,1)}x^r\,d\mu(r),\qquad
 T_\mu(C)=\min_{x>0}\frac{1+Cx}{S_\mu(x)},\quad C>0,
\end{equation}
where $\mu$ is a nonzero positive measure with finite support in $(0,1)$.
The normalization of $\mu$ will be specified separately for each theorem.
If $\mu=\sum_{j=1}^k B_j\delta_{r_j}$ and
$1,r_1,\ldots,r_k$ are rationally independent, the preceding modular
classification gives the following interpretation.  For any $b>0$ and
any phases, the function
\[
 -1+A\sum_{j=1}^kB_j e^{br_js+i\phi_j}+Ce^{bs+i\phi_*}
\]
has positive zero ordinates dense modulo every period if and only if
$A\ge T_\mu(C)$.  Thus $T_\mu(C)$ is a scalar transition threshold.
Our inverse statements concern measurements of this threshold.  They do
not assert that a threshold can be certified by observing zeros up to a
fixed finite height.

\subsection{The exact finite sample count}

\begin{theorem}[Finite determination and the necessary sample count]
\label{p05:fz:thm:finite}
Fix an integer $m\ge1$ and any $2m$ distinct positive numbers
$C_1,\ldots,C_{2m}$.  The vector
\[
             (T_\mu(C_1),\ldots,T_\mu(C_{2m}))
\]
determines every positive measure $\mu$ on $(0,1)$ with at most $m$
atoms.  It determines its actual number of atoms, locations, and masses.
More generally, two measures with at most $m_1$ and $m_2$ atoms are equal
if their thresholds agree at $m_1+m_2$ distinct positive settings.

On the open parameter space of exactly $m$ distinct atoms with positive,
unnormalized masses, $K<2m$ fixed distinct settings never give local
identifiability: through every parameter point there is a real-analytic
fiber of dimension $2m-K$ with the same data.  At $K=2m$ the parameter
map has a nonsingular Jacobian everywhere.  Consequently its inverse is
locally Lipschitz at every spectrum with distinct atoms and positive
masses.
\end{theorem}

\begin{lemma}[Contact points and common supporting lines]
\label{p05:fz:lem:contacts}
Write $S=S_\mu$ and $H(x)=S(x)-xS'(x)$.  Then $S'>0$, $S''<0$,
$H>0$, and
\[
 C(x)=\frac{S'(x)}{H(x)},\qquad
 C'(x)=\frac{S''(x)S(x)}{H(x)^2}<0
\]
is a bijection from $(0,\infty)$ onto $(0,\infty)$.
The minimum defining $T(C)$ has the unique contact point $x(C)$,
the inverse of this bijection, and
\begin{equation}\label{p05:fz:eq:contacts}
 T(C)=\frac1{H(x(C))},\quad
 T'(C)=\frac{x(C)}{S(x(C))},\quad
 T(C)-CT'(C)=\frac1{S(x(C))}.
\end{equation}
The line
\[
 L_C(x)=\frac{1+Cx}{T(C)}
\]
majorizes $S$ and touches it only at $x(C)$.

If two such profiles have the same threshold at $N$ increasing settings
$C_i$, their difference has at least $N$ distinct zeros on $(0,\infty)$.
\end{lemma}
\begin{proof}
All derivative signs follow by differentiating $x^r$, with $0<r<1$.
Moreover $C(x)x$ is the ratio of the $r$-weighted and
$(1-r)$-weighted positive sums, so it stays between two positive
constants.  Together with $C'<0$ this proves the bijection.  The tangent
inequality for a strictly concave function gives
$S(y)\le H(x)(1+C(x)y)$, with equality only at $y=x$.
It proves all minimum assertions; differentiation proves
\eqref{p05:fz:eq:contacts}.

For the last assertion, let $T_i$ be the common value and put
$L_i(x)=(1+C_ix)/T_i=a_i+b_ix$.
The slopes strictly increase with $C$, since
\[
 \frac{d}{dC}\frac C{T(C)}
       =\frac{T(C)-CT'(C)}{T(C)^2}>0.
\]
For $i<N$, let $z_i$ be the intersection of $L_i$ and $L_{i+1}$.
For either profile its contact with $L_i$ lies strictly to the right of
$z_i$, and its contact with $L_{i+1}$ lies strictly to the left.
Indeed the profile is below both lines, and equality with both at one
point would give two different tangent slopes to a differentiable
function.  Put $z_0=+\infty$ and $z_N=0$.  Both contacts of $L_i$
therefore belong to the interval $(z_i,z_{i-1})$; in particular these
intervals are nonempty and disjoint.

At the contact for the first profile its difference from the second is
nonnegative; at the contact for the second it is nonpositive.  There is
therefore a zero between the two contacts, including their common point
if they coincide.  These zeros belong to the $N$ disjoint intervals.
\end{proof}

\begin{lemma}[Real exponential interpolation, including collisions]
\label{p05:fz:lem:ode}
A nonzero solution of
\[
             \prod_{j=1}^n(D-r_j)h=0,
             \qquad r_j\in\mathbb R,
\]
has at most $n-1$ distinct real zeros.  Repeated roots are allowed.
Fix $n$, a compact interval of allowed roots, a compact interval of
sampling points, and a positive minimum separation between $n$ sampling
points.  On every fixed compact interval $J$ there is a uniform constant
$A$ such that all these solutions satisfy
\begin{equation}\label{p05:fz:eq:interpolation}
       \sup_{u\in J}|h(u)|\le A\max_{1\le i\le n}|h(u_i)|.
\end{equation}
\end{lemma}
\begin{proof}
For the zero count, conjugate by $e^{r_1u}$ and differentiate.  Rolle's
theorem reduces the asserted count to the equation of order $n-1$
with real roots $r_j-r_1$, $j\ge2$.  If the derivative is identically
zero, the original function is a nonzero multiple of $e^{r_1u}$ and has
no zeros.  Induction proves the result.

Write the order-$n$ equation as a first-order system with companion
matrix.  Its coefficients depend continuously on the roots, including
coincident roots.  The matrix taking the initial jet
$(h(0),\ldots,h^{(n-1)}(0))$ to $(h(u_i))_{i=1}^n$ is nonsingular by
the zero count.  The allowed roots and separated sampling tuples form a
compact set, so the inverses of these matrices are uniformly bounded.
The fundamental matrices on $J$ are uniformly bounded as well.  This
proves \eqref{p05:fz:eq:interpolation} without a frequency-separation
assumption.
\end{proof}

\begin{proof}[Proof of Theorem~\ref{p05:fz:thm:finite}]
Under $x=e^u$, the difference of two profiles is a sum of at most
$m_1+m_2$ real exponentials.  Unless it vanishes identically, it has at
most $m_1+m_2-1$ zeros by Lemma~\ref{p05:fz:lem:ode}.  The last assertion
of Lemma~\ref{p05:fz:lem:contacts} contradicts that bound.  Linear
independence of distinct real exponentials now identifies all atoms
and masses.

For local necessity, order the unknown locations and use coordinates
$(B_1,r_1,\ldots,B_m,r_m)$.  Write $x_i=x(C_i)$ and $u_i=\log x_i$.
The envelope derivative is
\[
 \frac{\partial T(C_i)}{\partial B_j}
    =-\frac{T(C_i)}{S(x_i)}e^{r_ju_i},\qquad
 \frac{\partial T(C_i)}{\partial r_j}
    =-\frac{T(C_i)}{S(x_i)}B_ju_i e^{r_ju_i}.
\]
The $u_i$ are distinct.  After nonzero row and column rescalings this
is the evaluation matrix of
$e^{r_ju},u e^{r_ju}$, $j=1,\ldots,m$.
Any initial set of $K$ functions in this order is the full solution
space of an order-$K$ differential equation with real roots, repeated
as needed.  Lemma~\ref{p05:fz:lem:ode} makes its square evaluation matrix
invertible.  Thus the full Jacobian has rank $K$ for $K\le2m$.
The analytic implicit-function theorem gives the fibers for $K<2m$;
the inverse-function theorem gives the remaining assertion.

The local nonidentifiability also persists if the initial frequencies
$1,r_1,\ldots,r_m$ are rationally independent and this restriction is
imposed on alternatives.  Choose a nonconstant analytic curve in a
fiber through the initial point.  Every rational affine relation in
the frequencies either has no zeros on the curve or has isolated zeros;
it cannot vanish identically since it does not vanish at the initial
point.  Removing the countable union of these discrete zero sets still
leaves distinct admissible alternatives arbitrarily nearby.
\end{proof}

The necessity statement uses unnormalized masses.  It does not assert
that $2m$ is the smallest global sample count under every additional
normalization.  For example, even a single probability atom can fail to
be identified from one threshold: for $\mu=\delta_r$,
\[
            T_\mu(C)=\frac{C^r}{r^r(1-r)^{1-r}},
\]
and this function of $r$ need not be injective.

\subsection{A sharp noise exponent from the same finite experiment}

Fix $0<a<b<1$, an integer $m\ge1$, and $2m$ distinct positive settings
$\mathcal C=(C_1,\ldots,C_{2m})$.  Let $\mathcal P_m[a,b]$ denote
probability measures on $[a,b]$ with at most $m$ atoms, and define
\[
 d_{\mathcal C}(\mu,\nu)
       =\max_i|T_\mu(C_i)-T_\nu(C_i)|.
\]
All constants below may depend on $m,a,b,\mathcal C$, but not on
frequency separations or individual atom masses.

\begin{theorem}[Optimal resolution without frequency separation]
\label{p05:fz:thm:noise}
For all $\mu,\nu\in\mathcal P_m[a,b]$,
\begin{equation}\label{p05:fz:eq:holder}
 W_1(\mu,\nu)\le K\,d_{\mathcal C}(\mu,\nu)^{1/(2m-1)}.
\end{equation}
The exponent cannot be increased.  More strongly, if $I$ is any fixed
compact interval in $(0,\infty)$, there are measures $\mu_h,\nu_h$,
each with $m$ atoms in $[a,b]$, such that, as $h\downarrow0$,
\begin{equation}\label{p05:fz:eq:lower}
 W_1(\mu_h,\nu_h)\ge c h,
 \qquad \sup_{C\in I}|T_{\mu_h}(C)-T_{\nu_h}(C)|\le K_I h^{2m-1}.
\end{equation}
Both spectra can be required to be rationally independent together
with $1$.

Consequently, for adversarial error at most $\delta$ in each of the
$2m$ thresholds, the minimax Wasserstein reconstruction error is
\begin{equation}\label{p05:fz:eq:minimax}
                       \asymp\delta^{1/(2m-1)}
                       \qquad(\delta\downarrow0).
\end{equation}
The same order remains a lower bound if the entire threshold curve on
a fixed compact interval is observed with uniform error $\delta$.
\end{theorem}

\begin{lemma}[Separated samples from noisy tangents]
\label{p05:fz:lem:noisycontacts}
There are constants $\delta_0,c_0,K_0>0$ and a compact interval
$X\subset(0,\infty)$ such that, whenever
$d_{\mathcal C}(\mu,\nu)\le\delta\le\delta_0$, there are $2m$
points $y_i\in X$ with pairwise distances at least $c_0$ and
\[
                |S_\mu(y_i)-S_\nu(y_i)|\le K_0\delta.
\]
These constants are uniform over $\mathcal P_m[a,b]$.
\end{lemma}
\begin{proof}
Arrange $C_1<\cdots<C_n$, $n=2m$.  Since
\[
 Cx=\frac{\int r x^r\,d\mu}{\int(1-r)x^r\,d\mu},
\]
every contact lies in the fixed interval
\[
 X=\left[\frac{a}{(1-a)C_n},
          \frac{b}{(1-b)C_1}\right].
\]
On $X$, the functions $S,H$ are uniformly bounded above and away from
zero, and $-S''$ is uniformly bounded above and away from zero.
The formula for $C'(x)$ in Lemma~\ref{p05:fz:lem:contacts} therefore
bounds $|dx/dC|$ above and away from zero.  Consecutive contacts for
any one profile are separated by a fixed $d_0>0$.

Write $L_i^\mu=(1+C_ix)/T_\mu(C_i)$ and similarly for $\nu$.
Their coefficients differ by $O(\delta)$, uniformly, since all
thresholds are uniformly bounded above and away from zero.
Let $z_i^\mu$ be the intersection of $L_i^\mu$ and $L_{i+1}^\mu$.
The consecutive slope differences are uniformly bounded above and away
from zero, by \eqref{p05:fz:eq:contacts}.  It follows that
$|z_i^\mu-z_i^\nu|\le K_1\delta$.

There is also a uniform positive margin between each intersection and
its two neighboring contacts.  To see this, if the contacts are $v>w$
and $-S''\ge\kappa$ on $X$, the tangent at $w$ exceeds $S(v)$ by
at least $\kappa(v-w)^2/2\ge\kappa d_0^2/2$.
Divide by the bounded difference of tangent slopes to bound $v-z_i$
away from zero.  The same argument gives a positive lower bound for
$z_i-w$.  Let a common lower bound be $\eta>0$.

For $\delta_0$ small enough, the contacts of \emph{both} profiles at
$C_i$ lie in the interval between $z_i^\mu$ and $z_{i-1}^\mu$,
with distance at least $\eta/2$ from each finite endpoint.  Here the
outer endpoints are interpreted as $0$ and $+\infty$.  These intervals
are ordered and separated by the specified margins.

Put $D=S_\mu-S_\nu$.  At the contact of $S_\mu$ with $L_i^\mu$,
the inequality $S_\nu\le L_i^\nu$ gives $D\ge-K_2\delta$.
At the contact of $S_\nu$, the reversed argument gives
$D\le K_2\delta$.  By continuity there is a point between the two
contacts with $|D|\le K_2\delta$: if either endpoint already has this
property use it, and otherwise the signs are opposite.  The resulting
points for successive $i$ are separated by at least $\eta$.
\end{proof}

\begin{lemma}[A finite-support moment bound]
\label{p05:fz:lem:moments}
Fix a compact real interval $J$ and $m\ge1$.  If probability measures
$\alpha,\beta$ on $J$ each have at most $m$ atoms, and
\[
 \left|\int z^k\,d(\alpha-\beta)(z)\right|\le\epsilon,
                   \qquad 1\le k\le2m-1,
\]
then $W_1(\alpha,\beta)\le A_{J,m}\epsilon^{1/(2m-1)}$.
\end{lemma}
\begin{proof}
For $0<\epsilon\le1$, let $F$ be the difference of the two cumulative
distribution functions and $E=\int_J|F|=W_1(\alpha,\beta)$.
There are at most $2m-1$ intervals on which $F$ is nonzero and constant,
and hence at most $d\le2m-2$ sign changes, after zero intervals are
ignored.  Choose a root $t_j\in J$ at each sign change, anywhere in
an intervening zero interval when necessary.  With a suitable sign,
$P(z)=\pm\prod_{j=1}^d(z-t_j)$ satisfies $P(z)F(z)\ge0$.
For $d=0$ take $P=1$ or $-1$ as appropriate.

If $d\ge1$, the union of the intervals of radius $E/(4d)$ about
the roots has length at most $E/2$.  Since $|F|\le1$, at least $E/2$
of the $|F|$ integral remains outside this union.  Therefore
\[
 \int_J|P(z)F(z)|\,dz\ge
          \frac E2\left(\frac E{4d}\right)^d.
\]
For $d=0$ the integral equals $E$.  If $Q'=P$, integration by parts
gives $\int PF=-\int Q\,d(\alpha-\beta)$.  The degree of $Q$ is
$d+1\le2m-1$, and its coefficients are uniformly bounded because
all roots belong to $J$; its constant term is irrelevant since the
signed measure has mass zero.  Thus $E^{d+1}\le A\epsilon$.
For $\epsilon\le1$ this implies the asserted exponent, and for larger
$\epsilon$ the diameter bound is sufficient.
\end{proof}

\begin{proof}[Proof of Theorem~\ref{p05:fz:thm:noise}]
First suppose $\delta=d_{\mathcal C}(\mu,\nu)$ is smaller than the
constant in Lemma~\ref{p05:fz:lem:noisycontacts}.
The function
\[
             h(u)=S_\mu(e^u)-S_\nu(e^u)
\]
is a sum of at most $2m$ real exponentials with exponents in $[a,b]$.
It solves an order-$2m$ equation as in Lemma~\ref{p05:fz:lem:ode}:
if necessary, pad its annihilating polynomial with further roots in
$[a,b]$.  The points $\log y_i$ lie in a fixed compact interval and
are uniformly separated.  Lemmas~\ref{p05:fz:lem:noisycontacts}
and~\ref{p05:fz:lem:ode}, with the target interval $[0,2m-1]$, give
\[
                 |h(k)|\le K_3\delta,
                       \qquad k=1,\ldots,2m-1.
\]
Push both measures forward by $r\mapsto z=e^r$.  These are exactly
their moment differences of orders $1,\ldots,2m-1$.
Lemma~\ref{p05:fz:lem:moments} proves the required Wasserstein bound for
the pushed-forward measures.  The map $\log z$ is Lipschitz on
$[e^a,e^b]$, so the same bound holds for $\mu,\nu$.
For $\delta\ge\delta_0$, increase the constant using
$W_1(\mu,\nu)\le b-a$.  This proves \eqref{p05:fz:eq:holder}.

For sharpness, let $q=2m-1$ and choose $r_0$ in the interior of
$[a,b]$.  For sufficiently small $h>0$, define probability measures
by placing masses $2^{-(q-1)}\binom qj$ at $r_0+jh$, assigning even
$j$ to $\mu_h$ and odd $j$ to $\nu_h$.  Each measure has $m$ atoms.
Their Laplace-transform difference is exactly
\begin{equation}\label{p05:fz:eq:binomialdifference}
 S_{\mu_h}(e^u)-S_{\nu_h}(e^u)
       =2^{-(q-1)}e^{r_0u}(1-e^{hu})^q.
\end{equation}
On any fixed compact $u$ interval this is $O(h^q)$.
For $C$ in a fixed compact interval, all minimizing $x$ belong to one
fixed compact subinterval of $(0,\infty)$, by the contact bound in
Lemma~\ref{p05:fz:lem:noisycontacts}.  Both denominators $S$ are bounded
away from zero there.  Taking minima in \eqref{p05:fz:eq:profile} thus
gives the second bound of \eqref{p05:fz:eq:lower} uniformly in $C$.

The elementary identity
$\sum_{j=0}^k(-1)^j\binom qj=(-1)^k\binom{q-1}k$
shows, by integrating the absolute difference of cumulative
distribution functions over the $q$ gaps, that
$W_1(\mu_h,\nu_h)=h$ exactly.
The arithmetic progression is not rationally independent.  To enforce
that requirement, perturb every location by less than $h^{q+1}$,
choosing locations outside the countable union of rational affine
hyperplanes for each spectrum.  This preserves the two masses and
the number of atoms.  It changes the Wasserstein distance by at most
$2h^{q+1}$ and each transform on a fixed compact window by
$O(h^{q+1})$.  The stated lower construction therefore remains valid
in the class interpreted by the modular-zero theorem.

Finally the space $\mathcal P_m[a,b]$ is compact and the finite
threshold map is continuous.  A minimum-distance fit to data within
$\delta$ of the truth has threshold discrepancy at most $2\delta$
from the truth, proving the upper minimax bound.  For the lower bound,
take $h$ a fixed small multiple of $\delta^{1/q}$.  The midpoint of
the two data vectors is within $\delta$ of both.  Every estimator must
make error at least half their Wasserstein distance on one of them.
The same midpoint argument works for the two entire curves on $I$.
\end{proof}

\subsection{What this transfer does and does not settle}

The preceding exact inverse used infinitely many amplitude settings.
Theorem~\ref{p05:fz:thm:finite} reduces them to $2m$ and identifies the
local obstruction to using fewer when masses are unnormalized.
Theorem~\ref{p05:fz:thm:noise} proves that this same finite experiment
already attains the best possible collision exponent, even compared
with the entire noisy curve on a bounded interval.  There is no
frequency-separation hypothesis in that theorem; its deterioration
with $m$ is unavoidable.  Neither theorem estimates the original
complex phases, which remain invisible to these threshold data.

The interpretation of a scalar threshold uses an infinite-height zero
closure.  A finite-height experiment requires an additional uniform
recurrence theorem.  In particular, these results supply no improved
small-divisor estimate for the exact prime vector $(\log p)_p$.
The fixed-strip Euler-zero variance and the entropy-sharp maximal-gap
exponent at the exact prime vector remain open in this round.
The almost-sure perturbed-frequency theorem and finite-spectrum
interpolation cannot remove that arithmetic gap.  Thus this section
closes two internal inverse problems and does not claim a new theorem
about the Riemann hypothesis, prime hyperuniformity, or actual prime
recurrence exponents.

% ===== inverse.tex =====
\section{Precision complexity of sparse heat observations}
\label{p05:fi:section}

The previous manuscript proved exponential recovery from exact heat
observations under a common analytic envelope, but only lower bounds for
geometrically sparse noisy observations.  We now obtain a two-sided
\emph{information-depth law} for every Hadamard-lacunary sampling sequence.
It determines the order of the logarithm of the minimax risk, although
not its sharp leading constant.  A second result shows that choosing
arithmetic inverse times changes logarithmic precision dependence into
power dependence, and that a power obstruction persists for every
adaptive heat experiment.  These are internally formulated problems.
The separate polynomial-information theorem
\ref{p05:fi:knu-poly} addresses the logarithmic upper-bound direction of
the analytic-information question of
Krieg--Novak--Ullrich~\cite{p03:an:KNU}; optimal measurement counts
remain open.

\subsection{A fixed analytic class and the two experiments}

Fix $I=[a,b]$ with $0<a<b$, $\ell=b-a$, and $\rho>1$.
The open Bernstein ellipse $E_\rho$ has foci $a,b$ and boundary
\[
 z=\frac{a+b}{2}+\frac{\ell}{4}(w+w^{-1}),\qquad |w|=\rho.
\]
Fix $\kappa>0$.  We use constants $m,M>0$ for which the class
\begin{equation}\label{p05:fi:class}
 \mathcal A=\left\{f\in H^\infty(E_\rho):
 f|_I\text{ real},\quad \int_I f=1,\quad f\ge m,
 \quad (\log f)''\le-\kappa,\quad \|f\|_{E_\rho}\le M\right\}
\end{equation}
contains a fixed analytic perturbation ball.  More precisely, normalize
$f_0(x)=c_0e^{-\kappa x^2}$ on $I$ and choose $\epsilon>0$ so small that
\begin{equation}\label{p05:fi:ball}
 \int_I h=0,\quad h|_I\text{ real},\quad
 \|h\|_{E_\rho}\le\epsilon
 \quad\Longrightarrow\quad f_0\pm h\in\mathcal A.
\end{equation}
Such $m,M,\epsilon$ exist: Cauchy estimates control the $C^2(I)$ norm
of $h$, whereas $f_0$ has a positive lower bound and
$(\log f_0)''=-2\kappa$.  All constants defining this class remain
fixed throughout the section.

The observation at inverse time $s>0$ is
\begin{equation}\label{p05:fi:heat}
 H_sf=\sqrt{s/\pi}\int_a^b e^{-sx^2}f(x)\,dx.
\end{equation}
For an increasing sequence $\mathbf s=(s_j)$, let
$R_{\mathbf s}^{\infty}(\delta)$ be the minimax $L^1(I)$ error
from the entire sequence $(H_{s_j}f+e_j)_{j\ge1}$ with
$\sup_j|e_j|\le\delta$.  Let $R_{\mathbf s}^{G}(\delta)$ be
the minimax expected $L^1$ error when $e_j$ are independent
$N(0,\delta^2)$ variables.  Reconstructions may be nonlinear; all
upper bounds below can be achieved by measurable approximate
minimum-residual estimators taking their values in $\mathcal A$.
Normal-family compactness makes $\mathcal A$ compact in $C^2(I)$,
so existence and measurability can also be obtained from a countable
dense subset.

\subsection{The information-depth law}

Suppose
\begin{equation}\label{p05:fi:lacunary}
 s_{j+1}\ge q s_j\quad(j\ge1),\qquad q>1,
 \qquad N(S)=\#\{j:s_j\le S\}.
\end{equation}
The count $N$ is taken in inverse time, not in physical heat time.

\begin{theorem}[Noise risk and usable observation depth]
\label{p05:fi:depth}
For every sequence satisfying \eqref{p05:fi:lacunary}, there are constants
$c,C>0$ and $\vartheta\in(0,1)$ such that, for all sufficiently small
$\delta$ and $L=\log(1/\delta)$,
\begin{equation}\label{p05:fi:depth-bounds}
 \frac{c\rho^{-N(L)}}{(N(L)+2)^2}
 \ \le\ \min\{R_{\mathbf s}^{\infty}(\delta),
                         R_{\mathbf s}^{G}(\delta)\}
 \ \le\ \max\{R_{\mathbf s}^{\infty}(\delta),
                         R_{\mathbf s}^{G}(\delta)\}
 \ \le\ C\vartheta^{N(L)}.
\end{equation}
Consequently, for either noise model,
\begin{equation}\label{p05:fi:depth-order}
 -\log R_{\mathbf s}(\delta)
       =\Theta\!\left(N(\log(1/\delta))\right).
\end{equation}
Constants may depend on the analytic class, $q$, and $s_1$, but not
on the later arrangement of the sampling sequence.  The theorem
matches the order of the logarithm of the risk.  It does not identify
a sharp exponent or assert that the two endpoints in
\eqref{p05:fi:depth-bounds} are equivalent up to multiplicative constants.
\end{theorem}

We give a self-contained proof.  The elementary matrix estimate below
is the Hilbert-space instance of the classical lacunary
M\"untz/Riesz-sequence mechanism; see
Fricain--Lef\`evre~\cite[Corollary~3.2]{p05:fi:FricainLefevre}
and Gaillard--Lef\`evre~\cite{p05:fi:GaillardLefevre}.

\begin{lemma}[Uniform conditioning on a finite interval]
\label{p05:fi:gram}
For $d>0$ and $q>1$ there are $S_0,c_q,C_q>0$ such that, whenever
$S_0\le t_1<\cdots<t_n$ and $t_{j+1}\ge qt_j$,
\[
 c_q\sum_{j=1}^n|v_j|^2
 \le\int_0^d\left|\sum_{j=1}^n
               v_j\sqrt{2t_j}\,e^{-t_jy}\right|^2dy
 \le C_q\sum_{j=1}^n|v_j|^2.
\]
The constants are independent of $n$ and the chosen $t_j$.
\end{lemma}
\begin{proof}
The Gram matrix on the half-line has entries
\[
 G_{ij}=\frac{2\sqrt{t_it_j}}{t_i+t_j}.
\]
For $i\ne j$, $G_{ij}\le2q^{-|i-j|/2}$, so its row sums are
bounded by $K_q=1+4\sum_{r\ge1}q^{-r/2}$.
The Cauchy inverse formula gives
\[
 (G^{-1})_{ij}=G_{ij}A_iA_j,
 \qquad A_i=\prod_{k\ne i}\frac{t_i+t_k}{t_i-t_k}.
\]
This formula also follows by partial fractions: the rational function
with poles at $-t_k$ and prescribed values at the $t_j$ gives the
columns of the inverse.  Separation yields
\[
 |A_i|\le A_q:=
 \prod_{r=1}^{\infty}\left(\frac{1+q^{-r}}{1-q^{-r}}\right)^2<\infty.
\]
Thus $\|G^{-1}\|_{2\to2}\le A_q^2K_q$ by the row-sum bound,
and $\lambda_{\min}(G)\ge(A_q^2K_q)^{-1}$.
The Gram matrix on $[0,d]$ is $G-DGD$, where
$D=\operatorname{diag}(e^{-dt_j})$.
Since $\|DGD\|\le K_qe^{-2dt_1}$, choosing $S_0$ sufficiently
large proves the result.
\end{proof}

\begin{lemma}[Correcting finitely many noisy observations]
\label{p05:fi:correction}
There are $S_0,K,C>0$, depending only on $I,E_\rho,q$, with the
following property.  If $S_0\le t_1<\cdots<t_n$ satisfy the same
ratio separation, then, for arbitrary real $d_1,\ldots,d_n$, there
is a real entire function $p$ with
\begin{equation}\label{p05:fi:correction-bound}
 H_{t_j}p=d_j\quad(1\le j\le n),\qquad
 \|p\|_{E_\rho}\le C e^{Kt_n}\max_j|d_j|.
\end{equation}
\end{lemma}
\begin{proof}
Put $u=x^2-a^2\in[0,b^2-a^2]$ and
$\phi_j(x)=\sqrt{2t_j}e^{-t_j(x^2-a^2)}$.
The matrix $Q_{ij}=\int_I\phi_i\phi_j$ satisfies $Q\ge c I_n$:
use Lemma~\ref{p05:fi:gram} and
$dx=du/(2\sqrt{a^2+u})\ge du/(2b)$.
For $p=\sum c_j\phi_j$, the interpolation equations become
\[
 Qc=b,\qquad b_i=\sqrt{2\pi}\,e^{a^2t_i}d_i.
\]
In particular $\|c\|_2\le C\sqrt n\,e^{a^2t_n}\max|d_i|$.
For
$D=\max\{0,\sup_{z\in E_\rho}\operatorname{Re}(a^2-z^2)\}$,
\[
 \|p\|_{E_\rho}
 \le C n\sqrt{t_n}\,e^{(a^2+D)t_n}\max|d_i|.
\]
Since $n\le1+\log_q(t_n/S_0)$, increasing $K,C$ absorbs the
polynomial factor and gives \eqref{p05:fi:correction-bound}.
\end{proof}

\begin{lemma}[An analytic fiber diameter estimate]
\label{p05:fi:diameter}
There is $\vartheta\in(0,1)$, depending only on $I,E_\rho$, such
that a real holomorphic function $h$ bounded by $B$ on $E_\rho$
and orthogonal on $I$ to $n$ exponentials
$e^{-t_1x^2},\ldots,e^{-t_nx^2}$ with distinct positive $t_j$
satisfies
\[
 \|h\|_{L^1(I)}\le B\ell\vartheta^n.
\]
Consequently, for $f,g\in\mathcal A$ and separated $t_j\ge S_0$,
\begin{equation}\label{p05:fi:modulus}
 \|f-g\|_1\le C\vartheta^n+
        Ce^{Kt_n}\max_{j\le n}|H_{t_j}(f-g)|.
\end{equation}
\end{lemma}
\begin{proof}
Exponentials with distinct real exponents form a Chebyshev system:
after $y=x^2$, dividing by one exponential, differentiating, and
using Rolle's theorem proves that a nonzero combination of $r+1$
of them has at most $r$ real zeros, counted with multiplicity.
If a nonzero $h$ had fewer than $n$ sign changes, a combination of
as many exponentials as needed could be chosen with simple zeros
exactly at those sign-change points and the same sign as $h$.
Their integral pairing would be strictly positive, a contradiction.
Thus $h$ has at least $n$ distinct zeros in $I$.

Let $\phi:E_\rho\to\mathbb D$ be conformal and set
\[
 \vartheta=\max_{x,y\in I}
 \left|\frac{\phi(x)-\phi(y)}{1-\overline{\phi(y)}\phi(x)}\right|<1.
\]
Repeated Schwarz factorization at the $n$ zeros bounds
$|h(x)|$ by $B\vartheta^n$ on $I$.
For the last assertion apply Lemma~\ref{p05:fi:correction} to the
data of $f-g$.  The corrected function $h=f-g-p$ has zero data
and analytic norm at most $2M+\|p\|_{E_\rho}$.
Adding back $p$ proves \eqref{p05:fi:modulus}.
\end{proof}

\begin{proof}[Proof of Theorem~\ref{p05:fi:depth}]
For the upper bound, discard the finitely many parameters below
$S_0$ and use all the remaining parameters at most $c_0L$, with
$c_0>0$ chosen so small that $Kc_0<1/2$.
Write $n$ for their number.  With deterministic noise, a class
member fitting those observations to error at most $\delta$
has discrepancy from the truth at most $2\delta$.  Formula
\eqref{p05:fi:modulus} gives error at most
$C\vartheta^n+C\delta^{1/2}$.

With Gaussian noise, minimize the maximum residual on the same
coordinates.  On the event
$\max_{j\le n}|Z_j|\le A\sqrt n$, its discrepancy is at most
$2A\delta\sqrt n$, up to an arbitrarily small optimization error.
Choose fixed $A$ large enough that
$2n e^{-A^2n/2}\le C\vartheta^n$.
Outside this event the loss is at most two.  Since
$n=O(\log L)$, the resulting expected risk is at most
$C\vartheta^n+C\sqrt n\,\delta^{1/2}$, hence at most
$C\vartheta^n$ for small $\delta$.
Every multiplicative interval $[cL,CL]$ contains at most
$1+\log_q(C/c)$ sample points.  Therefore
$n=N(L)+O(1)$, and the stated upper bound follows.

For completeness we reconstruct the lower bound in the same fixed
analytic class.  Put $n=N(C_0L)$, where $C_0>2/a^2$.
There is a nonzero real polynomial $P$ of degree at most $n+1$
annihilated by integration and the first $n$ observations.
Normalize $\|P\|_{I,\infty}=1$.
The Bernstein--Walsh inequality and the Legendre reproducing-kernel
bound give
\[
 \|P\|_{E_\rho}\le\rho^{n+1},\qquad
 \|P\|_1\ge\frac{\ell}{(n+2)^2}.
\]
The first inequality follows directly by the maximum principle
after the Joukowski parametrization of the ellipse.  For the second,
the degree-$(n+1)$ Legendre kernel has diagonal at most
$(n+2)^2/\ell$, so
$1\le(n+2)\|P\|_2/\sqrt\ell$ and
$\|P\|_2^2\le\|P\|_1$.
Set $h=\epsilon\rho^{-(n+1)}P$.  By \eqref{p05:fi:ball},
$f_\pm=f_0\pm h$ are in $\mathcal A$, agree in the first $n$
observations, and have $L^1$ separation at least
$c\rho^{-n}/(n+2)^2$.

For every $j>n$, their difference is bounded by
\[
 |H_{s_j}(f_+-f_-)|\le2\sqrt{s_j/\pi}\,e^{-a^2s_j}.
\]
After another finite initial segment, the sequence has spacing at
least one.  Eventual monotonicity of $s e^{-2a^2s}$ then gives,
with $S=s_{n+1}>C_0L$,
\[
 \sum_{j>n}|H_{s_j}(f_+-f_-)|^2
 \le C_a(S+1)e^{-2a^2S}\le\delta^2.
\]
For deterministic errors, the midpoint mean sequence is compatible
with both densities, giving half their separation as a lower bound.
For Gaussian errors, the two product laws have relative entropy at
most $1/2$.  Pinsker's inequality and the common part of the two laws
give one quarter of their separation.  Finally
$N(C_0L)=N(L)+O(1)$ proves \eqref{p05:fi:depth-bounds}.
Since $N(L)\to\infty$ and $\log(N(L)+2)=o(N(L))$,
\eqref{p05:fi:depth-order} follows.
\end{proof}

\begin{corollary}[Three different precision scales]\label{p05:fi:scales}
For either noise model, with the analytic class unchanged:
\begin{align*}
 s_j=q^j\ (q>1)\quad&\Longrightarrow\quad
 R_{\mathbf s}(\delta)=[\log(1/\delta)]^{-\Theta(1)},\\
 s_j=e^{j^p}\ (p>1)\quad&\Longrightarrow\quad
 R_{\mathbf s}(\delta)=
 \exp\{-\Theta((\log\log(1/\delta))^{1/p})\},\\
 s_j=e^{q^j}\ (q>1)\quad&\Longrightarrow\quad
 R_{\mathbf s}(\delta)=[\log\log(1/\delta)]^{-\Theta(1)}.
\end{align*}
Here $A^{-\Theta(1)}$ means two bounds with possibly different
positive exponents.  In particular it is not a claim about a single
sharp exponent.
\end{corollary}
\begin{proof}
Substitute the respective counting functions into
\eqref{p05:fi:depth-order}.
\end{proof}

\paragraph{What the depth law closes.}
The earlier construction established arbitrarily slow consistency by
making the sampling sequence increasingly sparse.  Theorem
\ref{p05:fi:depth} now computes that slowdown, at the level of
$-\log R$, for \emph{every} sequence with a fixed multiplicative gap.
For geometric sampling it supplies the missing logarithmic-type upper
bound.  The optimal exponent in that logarithmic rate is still open.

\subsection{Designing the times: power accuracy and a universal obstruction}

The next result separates a limitation of an excessively sparse design
from the limitation of the full heat-observation family.
For adversarial errors let $R_{\mathrm{all}}^\infty(\delta)$ be
the infimum of worst-case $L^1$ errors over all adaptive procedures
allowed to query $H_sf+e$ at any $s>0$, with $|e|\le\delta$ per
query and no upper bound on the number of queries.
Randomization does not invalidate the lower bound below if expected
loss is used.  Let $R_B^{\mathrm{ad},G}$ denote the corresponding
minimax expected loss with at most $B$ adaptive queries, each with
fresh independent unit-variance Gaussian noise.

Choose $\tau>1$ so close to one that
\begin{equation}\label{p05:fi:tau}
 \operatorname{Re}(z^2)\ge c_*>0\qquad(z\in E_\tau),
 \qquad \beta=\frac{\log\rho}{\log\rho+\log\tau}\in(0,1).
\end{equation}
The ellipse here has the same foci $a,b$ as before; such a choice
exists because $a>0$.

\begin{theorem}[Power-type accuracy under designed sampling]
\label{p05:fi:design}
For fixed $h>0$ there are constants $c,C,\alpha>0$ such that
\begin{equation}\label{p05:fi:design-deterministic}
 \frac{c\delta^\beta}{1+\log(1/\delta)}
 \le R_{\mathrm{all}}^\infty(\delta)
 \le C\delta^\alpha
\end{equation}
for small $\delta$.  The upper bound uses only
$O(\log(1/\delta))$ nonadaptive observations at
$s_j=jh$, $1\le j\le n$.
It also holds in expected loss for Gaussian errors of standard
deviation $\delta$ on those observations, after possibly reducing
$\alpha$.

For the unit-variance Gaussian budget experiment there are
$\alpha_0,c,C>0$ such that
\begin{equation}\label{p05:fi:budget-design}
 \frac{c B^{-\beta/2}}{1+\log B}
 \le R_B^{\mathrm{ad},G}\le C B^{-\alpha_0}
 \qquad(B\ge2).
\end{equation}
The upper bound is achieved by a fixed design with repeated
arithmetic inverse times.  Thus
$-\log R_B^{\mathrm{ad},G}=\Theta(\log B)$.
Neither statement matches the exact power exponent.
\end{theorem}

\begin{lemma}[Arithmetic data correction]\label{p05:fi:arithmetic}
For any fixed $h>0$, there are $A>1,C>0$ such that
\begin{equation}\label{p05:fi:arithmetic-modulus}
 \|f-g\|_1\le C\vartheta^{n+1}
       +CA^n\max_{1\le j\le n}|H_{jh}(f-g)|,
 \qquad f,g\in\mathcal A.
\end{equation}
\end{lemma}
\begin{proof}
Put $z=e^{-hx^2}\in J=[e^{-hb^2},e^{-ha^2}]$.
The measure $dx$ becomes $w(z)\,dz$, with $w$ continuous and
bounded above and below by positive constants on $J$.
For $v(z)=\sum_{j=0}^n c_jz^j$, the monomial Gram matrix
\[
 Q_{ij}=\int_J z^{i+j}w(z)\,dz,\qquad 0\le i,j\le n,
\]
satisfies $\|Q^{-1}\|_{2\to2}\le C_1^{n+1}$ for some
$C_1>1$.  Here is an elementary justification of the exponential
bound.  Expand $v$ in the shifted orthonormal Legendre polynomials
on $J$.  The Euclidean norm of the monomial coefficient vector of
the degree-$k$ polynomial is at most $C_2^{k+1}$: this follows
by induction from the three-term Legendre recurrence and the fixed
affine change of variable.  Cauchy--Schwarz then bounds
$\|c\|_2\le C_3^{n+1}\|v\|_{L^2(J)}$.
The positive lower bound on $w$ proves the asserted inverse bound,
after enlarging the constant.

Apply this Gram inverse to find $p(x)=\sum_{j=0}^n c_j e^{-jh x^2}$
such that $\int_Ip=0$ and
$H_{jh}p=H_{jh}(f-g)$ for $1\le j\le n$.
The right-hand side of the moment equations is zero at $j=0$ and
$\sqrt{\pi/(jh)}H_{jh}(f-g)$ at $j\ge1$.
The Gram estimate, followed by the bound
$|e^{-jh z^2}|\le e^{jh\sup_{E_\rho}(-\operatorname{Re}z^2)}$,
gives
\[
 \|p\|_{E_\rho}\le CA^n\max_{j\le n}|H_{jh}(f-g)|
\]
after increasing $A$ to absorb polynomial factors.
The corrected function $f-g-p$ is orthogonal to the $n+1$
functions $1,e^{-hx^2},\ldots,e^{-nhx^2}$.
The argument of Lemma~\ref{p05:fi:diameter}, which also permits the
exponent zero, proves \eqref{p05:fi:arithmetic-modulus}.
\end{proof}

\begin{proof}[Proof of Theorem~\ref{p05:fi:design}]
For the deterministic upper bound, take
$n=\lfloor c_1\log(1/\delta)\rfloor$ with $c_1>0$ small
enough that $c_1\log A<1/2$.  A minimum-residual estimate in the
class has data discrepancy at most $2\delta$.
Lemma~\ref{p05:fi:arithmetic} bounds its error by
$C\vartheta^n+C\delta A^n\le C\delta^\alpha$ for some
$\alpha>0$.  Under Gaussian noise, use the event
$\max_{j\le n}|Z_j|\le A_0\sqrt n$ exactly as in the depth
theorem.  Its complement has exponentially small probability in
$n$, and its extra factor $\sqrt n$ is absorbed by decreasing
$\alpha$.

For the lower bounds, let $P_n$ be the shifted Legendre polynomial
of degree $n\ge1$, normalized by $\|P_n\|_{I,\infty}=1$.
It is orthogonal to every polynomial of degree less than $n$,
including the constant.  Moreover,
\[
 \|P_n\|_{E_\rho}\le\rho^n,\qquad
 \|P_n\|_1\ge\|P_n\|_2^2=\frac{\ell}{2n+1}.
\]
Thus $f_\pm=f_0\pm\epsilon\rho^{-n}P_n$ belong to the fixed
class and have separation at least $c\rho^{-n}/(n+1)$.

The heat kernels
$k_s(z)=\sqrt{s/\pi}e^{-sz^2}$ are uniformly bounded on
$E_\tau$ for all $s>0$, by \eqref{p05:fi:tau} and
$\sup_{s>0}\sqrt{s}\,e^{-c_*s}<\infty$.
Their degree-$(n-1)$ Chebyshev truncations consequently approximate
them on $I$ with error at most $C\tau^{-n}$, uniformly in $s$.
To verify this standard estimate directly, substitute the
Joukowski map into $k_s$ and use the Laurent coefficients on
circles $|w|<\tau$; Cauchy's bound gives geometric coefficients
and summing the tail gives the estimate.
Orthogonality of $P_n$ therefore gives the simultaneous bound
\begin{equation}\label{p05:fi:invisible-all}
 \sup_{s>0}|H_s(f_+-f_-)|\le C(\rho\tau)^{-n}.
\end{equation}

Choose $n$ as the smallest integer with the right side at most
$2\delta$.  An adversarial oracle can always return
$H_sf_0$.  This response is consistent with both $f_+$ and $f_-$
at every query, including queries chosen adaptively from all earlier
responses.  Thus unlimited adaptive querying cannot distinguish
the pair.  Their separation proves
$R_{\mathrm{all}}^\infty(\delta)\ge
c\delta^\beta/(1+\log(1/\delta))$.
The same common-response argument works after conditioning on the
random seed of a randomized procedure.

For the Gaussian budget lower bound, the chain rule for relative
entropy applies also to an adaptive experiment: conditional on the
past, the algorithm chooses the same query rule under both models,
and the next conditional divergence is half the squared difference
of their two means.  Hence \eqref{p05:fi:invisible-all} gives
\[
 D(P_+^{(B)}\Vert P_-^{(B)})
 \le C B(\rho\tau)^{-2n}.
\]
Choose $n=\lceil(\log B+C_4)/(2\log(\rho\tau))\rceil$
with $C_4$ large enough that the divergence is at most $1/2$.
The two-point inequality proves the lower bound in
\eqref{p05:fi:budget-design}.

For its upper bound, take $n=\lfloor c_2\log B\rfloor$
arithmetic locations and repeat each at least $\lfloor B/n\rfloor$
times.  The averaged noise standard deviation is at most
$C\sqrt{n/B}$.  Except on an event of probability $Ce^{-c_3n}$,
the maximum averaged noise is at most $Cn/\sqrt B$.
Lemma~\ref{p05:fi:arithmetic} then bounds the expected loss by
\[
 C\vartheta^n+C A^n n/\sqrt B+Ce^{-c_3n}.
\]
Choose $c_2$ sufficiently small and absorb the factor $n$ into
a slightly smaller positive power of $B$.  This proves the upper
bound.  Bounded small values of $B$ are covered by adjusting $C$.
\end{proof}

\paragraph{Scope and unresolved boundaries.}
The new result is a sampling-dependent minimax classification under a
single analytic envelope: all multiplicatively separated sequences
obey the depth law, whereas a designed arithmetic sequence achieves
power dependence on precision.  Classical Riesz-sequence conditioning,
analytic approximation, and two-point lower bounds are credited as
ingredients, not new theorems.  General inverse-Laplace regularization
has a substantial literature; for example
Mar\'echal--Triki--Simo Tao Lee~\cite{p05:fi:Marechal} study a
different global logarithmic-stability problem.  Priority for the
specific fixed-envelope, infinite-noisy-sequence classification here
has not been established by the literature search.

The following remain open in this work: the exact exponent for
$s_j=q^j$; the optimal power exponent and constant for designed heat
sampling; and the optimal number of arbitrary adaptive analytic
measurements in the sense of \cite[Section~3]{p03:an:KNU}, beyond
the logarithmic upper bound proved in Theorem~\ref{p05:fi:knu-poly}.
In particular, the adaptive lower bound in Theorem~\ref{p05:fi:design}
is for the concrete family of heat functionals.  It does not extend
to all analytic nonlinear information merely because each heat
functional is analytic.

% ===== back.tex =====
\section{A target register and the unresolved transfers}
\label{p05:fo:register}

The entries below record distinct questions rather than count individual
lemmas as separate discoveries. ``Complete'' refers to the mathematical
scope stated in that row, subject to the internal-review status of this
manuscript. The externally posed analytic and convex directions overlap:
the same convex-polynomial theorem addresses both, so they are not
represented as independent proof breakthroughs. The homogeneous variant
is grouped with them. Similarly, the several Bessel power criteria are
grouped into one restoration problem.

\small
\renewcommand{\arraystretch}{1.15}
\begin{longtable}{@{}p{0.055\textwidth}p{0.255\textwidth}p{0.62\textwidth}@{}}
\toprule
No.&Target and origin&Result, status, and precise limit\\
\midrule
\endfirsthead
\toprule No.&Target and origin&Result, status, and precise limit\\\midrule
\endhead
\bottomrule\endfoot
1&Analytic and convex information; externally posed complexity directions&
Positive logarithmic-query answer, Theorem~\ref{p05:fi:knu-poly}.
At most $\min\{m,\lceil\log_2m\rceil+2\}$ globally convex polynomial
queries suffice on all of $\mathbb R^m$. A convex analytic one-Lipschitz
variant has the same count. Positively homogeneous variants
are also supplied. Optimal query count, bounded-degree complexity and
finite-precision costs remain separate questions.\\
2&Full adaptive Bernoulli replication cap; inherited gap&
Complete rate classification for every $1\le Q\le B$,
Theorem~\ref{p05:fb:thm:bernoulli-cap}. The previous critical adaptive window
is closed. No improvement in minimax order from adaptation in this model.\\
3&Scalar endpoint cap; new internal extension&
Complete curve $\max\{\log(eB)/\sqrt B,(BQ)^{-1/3}\}$,
Theorem~\ref{p05:fb:thm:cap}. The uncapped logarithm was already proved;
the capped classification is new here.\\
4&Dependent multivariate endpoint cap; new internal problem&
Complete curve for every fixed $d\ge2$,
Theorem~\ref{p05:fb:thm:cap}, without independence of the latent coordinates.
Constants may depend on $d$.\\
5&Deterministic binomial inverse modulus; new internal problem&
Complete two-parameter modulus, Theorem~\ref{p05:fm:modulus}.
It is a channel-inversion result and is not the statistical minimax
risk of the MLE.\\
6&Certified approximate MLE; new internal problem&
Likelihood deficit $O(N^{-1})$ preserves the stated statistical rates;
this order is necessary for a worst-case root-$N$ guarantee based on
deficit alone. Theorems~\ref{p05:fm:statisticalprecision} and
\ref{p05:fm:algorithm} give a polynomial-time certified construction.\\
7&Power restoration of HCM; new internal transfer&
Complete existence and fixed-tilt classification,
Theorems~\ref{p05:fs:restoration} and~\ref{p05:fs:fixed-tilt}.
The first Bessel zero controls the critical power cost,
Corollary~\ref{p05:fs:critical-cost}. The smallest exponent at an arbitrary
fixed tilt is specified by a supremum test, not a general elementary formula.\\
8&Random covariance rank versus repeated views; new internal problem&
Complete threshold for rank at most $r<d$,
Theorem~\ref{p05:fg:thm:rank}. $r+1$ blocks suffice; every collection
of designs with at most $r$ blocks can fail on compact exact-rank laws.\\
9&Finite iterated-log tail hierarchy; new internal extension&
Complete at every finite depth with an all-dimensional sufficiency
and a sharp scalar dyadic converse,
Theorem~\ref{p05:fg:thm:hierarchy}. A fixed-$d>1$ converse is not claimed.\\
10&Finite transition-threshold identification; new internal problem&
Any $2m$ distinct settings identify a positive at-most-$m$-atom
spectrum, Theorem~\ref{p05:fz:thm:finite}. Fewer are locally insufficient
for unnormalized exactly-$m$-atom spectra. No minimal count for
normalized masses is asserted.\\
11&Noisy finite-threshold recovery; new internal problem&
Complete minimax error exponent $1/(2m-1)$ without separation,
Theorem~\ref{p05:fz:thm:noise}. The classical moment exponent is credited;
the new transfer starts from nonlinear optimized thresholds.\\
12&Noisy lacunary heat information depth; inherited inverse gap&
Complete order of $-\log R$ for every Hadamard-lacunary sequence,
Theorem~\ref{p05:fi:depth}. Sharp risk exponents are not matched.\\
13&Designed versus sparse heat sampling; new internal transfer&
Complete distinction between power-type and logarithmic-type
precision dependence, Theorem~\ref{p05:fi:design}, including adaptive
lower bounds. The optimal power exponent remains open.\\
14&McKay SD strictness; external classification direction&
Partial extension to every admissible $q>0$ and
$-1<\mu\le1/2$, Proposition~\ref{p05:fs:sd-strict}.
The entire negative-cut, high-order classification is not closed.\\
15&Full statistical MLE transition; inherited boundary&
Open near $t\asymp N/\log N$. The new deterministic modulus proves
a limitation of the global count-Wasserstein argument, rather than a
statistical lower bound ruling out an optimal MLE.\\
16&Fixed-dimensional matrix tail converse; inherited transfer&
Open for each $d>1$ at the supercritical iterated-log envelopes.
Cancellation near indefinite Fourier matrices prevents the scalar
contour estimate from transferring directly.\\
17&Exact-prime recurrence exponent; inherited arithmetic target&
Open. Generic frequency perturbations and continuum inverse theorems
do not supply quantitative recurrence for the actual prime-log vector.\\
18&Euler-product zero variance or hyperuniformity; inherited target&
Open. The needed quantitative Fourier/small-divisor control is absent.
No Riemann-hypothesis or Speiser-type conclusion follows from the
spectral threshold results.\\
19&Exact optimal query count in analytic information; external refinement&
Open. The logarithmic upper bound answers the speedup direction, while
sharp lower bounds beyond the standard continuous-information obstruction
are not established.\\
20&Sharp heat inverse exponents; inherited precision refinement&
Partial. Both sparse and designed regimes now have matching
logarithmic-risk orders, but their leading exponents need not coincide.\\
\end{longtable}
\normalsize

\subsection{What is actually gained over the preceding paper}

The adaptive Bernoulli boundary is now complete, while the previous paper
only had a complete predetermined curve and polynomially subcritical
adaptive regimes. The endpoint result now includes every cap and arbitrary
fixed dimension. The continuum spectral inverse has become a finite-data
inverse with an optimal noise exponent. The Gaussian result now distinguishes
intrinsic covariance rank from the ambient dimension through unknown
orientation. The heat section computes how arbitrary lacunarity changes
the usable information depth, rather than giving a geometric lower bound
alone. The Bessel section changes the variable by a power and obtains
HCM restoration even where the original density lacks it.

Most significantly for the externally posed questions, the analytic
information restriction does not block the colored-cover method after
its decisions are made tolerant to approximation. Convexification using
a known radial measurement transfers this construction into still more
restricted query classes. This is an upper-complexity theorem in an exact
oracle model. It should be assessed on that scope, not advertised as a
practical scheme for recording an arbitrary high-dimensional object in
a few finite-bit numbers.

\subsection{Why the arithmetic targets do not follow}

An exact inverse theorem answers which latent object is determined by an
observable. A quantitative arithmetic theorem also needs to show that
the observable is reached, or estimated, at a controlled height for the
specific arithmetic frequencies. The latter step is not provided by
identifiability, continuity, generic rational independence, or a
worst-case statistical inverse modulus. In particular, the transition
thresholds used here belong to a finite exponential model. Replacing
them by a list of right-half-plane zeta zeros is a different observation
problem and is not justified by any result of this paper.

The prime recurrence and zero-variance entries are therefore retained as
open. A credible next advance there would require a new quantitative
arithmetic estimate, not another relabeling of an abstract uniqueness
statement. Similarly, the full statistical MLE transition requires more
than a global count-transport estimate, and the matrix tail converse
requires more than a scalar flat transform.

\subsection{Review record and provenance}

All main sections received an independent internal proof review in addition
to author self-review. Checks focused on conditional versus unconditional
information, zero-mass normalization, singular covariance strata, quotient
cancellation, atom collisions, boundary residues, and the distinction
between exact and noisy information. The combined manuscript was compiled
and its rendered pages inspected. These are internal checks, not a
formal proof certificate or external peer review.

The external comparisons use the cited primary papers. A targeted search
can find direct overlaps and verify the statement of an open question;
it cannot certify exhaustive priority. The colored-cover idea, classical
quasi-analytic criteria, beta--binomial eigenvalues, HCM/GGC representations,
finite-moment collision exponent, and lacunary exponential conditioning
are attributed to their respective sources. The contributions claimed
here are the stated transfers and complete boundaries built from those
ingredients. No journal-level or prize-level evaluation is implied.
\endgroup
% ===== END SOURCE 05 =====


% ===== BEGIN SOURCE 06: KEIO_Boundary_Laws_for_Observation_and_Recovery.tex =====
\clearpage
\part{Boundary Laws for Observation and Recovery}
\label{part:06}
\setcounter{section}{0}
\begin{center}\large\itshape Prime-Zeta Counting, Global Likelihood Optimality, and Sharp Inverse Thresholds\end{center}
\begin{center}\small Source manuscript: 15 October 2026.\end{center}
\begingroup
\renewcommand{\arraystretch}{1}
\begin{quote}\small\textbf{Source abstract.}\quad
We continue an observation-and-recovery program through several distinct
changes of analytic representation. For the actual prime zeta function,
we prove that the number of zeros in $\Re s>1$ up to height $T$ is
$\Theta(T)$, with a multiplicity upper bound and a distinct-zero lower
bound. This supplies the endpoint upper bound explicitly proposed in
recent literature and extends to every fixed derivative. For arbitrary
binomial mixing laws we prove that every unrestricted maximum-likelihood
estimator attains the complete minimax order
$\min\{1,t^{-1}+(t\log(eN))^{-1/2}+N^{-1/2}\}$ for all sample sizes
$N$ and replication counts $t$. A low-frequency transportation estimate
and a small-amplitude high-frequency Hellinger estimate close the
previous transition gap. Further results give exact and degree-dependent
adversarial-noise floors for restricted adaptive information, sharp
iterated-log covariance-tail boundaries in every fixed Gaussian
dimension, positivity certificates for nonlinear spectral thresholds,
and three critical regimes for Bessel power restoration. Fixed analytic
population classes exhibit double- and triple-logarithmic replication
penalties. Sixteen precisely scoped targets receive proofs; dependent
consequences, classical ingredients, one partial SD advance, and
remaining questions are identified explicitly. The statements are
research claims supported by proofs and independent automated readings,
not externally certified priority or solutions of sixteen named
classical conjectures.
\end{quote}


\section{The new step beyond the preceding manuscripts}
\label{p06:br:intro}

The starting point is the 223-page manuscript \emph{Observation and
Recovery in Arithmetic and Probability} and its sequels
\cite{p01:KEIOIntegrated,comp:part:01,comp:part:02,comp:part:03,comp:part:04,comp:part:05}.
The present paper contains new continuations rather than a consolidation
of their already claimed results. Earlier manuscripts are unpublished
research sources, not substitutes for external validation. New arguments
are supplied here; explicit older dependencies are stated where they
are used, particularly for the Bessel signed spectral kernel.

The principal new statement in arithmetic concerns an actual prime
Dirichlet series, without independent random phases. The earlier
fixed-$\gamma>1$ zero counts do not control accumulation toward the line
$\Re s=1$. We remove the prime-power tail from the Euler logarithm,
exponentiate, and count zeros of a single-valued function defined in a
larger half-plane. Classical second moments then suffice for an
averaged Jensen argument. The result addresses precisely the endpoint
question recorded in \cite[p.~16]{p06:ra:Tuzcu}; it does not concern the
Riemann hypothesis or a zero-counting formula on continuation branches
to the left of one.

The statistical advance is a different kind of boundary result. The
preceding paper obtained the optimal MLE order in polynomially
sublinear replication ranges, but its deterministic count-transport
modulus lost a logarithmic factor around $t=N/\log N$. That modulus
was a real obstruction for a proof using only that norm. A fitted
likelihood also supplies Hellinger information. Splitting the inverse
test polynomial at degree about $\sqrt t$ makes the high-frequency
remainder small in supremum norm, so this second observation norm
closes the transition. This extends the statistical range questions
in \cite{p01:info:VKVK2019,p01:info:MMRS2025}, without resolving their separate
constant-one coefficient conjecture.

\subsection{Four transfers with different mathematical content}

The common discipline is to retain the observation model while changing
the object on which an available theorem acts. The actual mechanisms
are different and no universal theorem subsuming them is asserted.

\begin{enumerate}
\item A logarithmic arithmetic series is lifted to a meromorphic
function; positive logarithmic means then control zeros at a boundary.
\item An inverse statistical test is split between two norms; the
amplitude of the difficult component determines which norm is useful.
\item A meromorphic zero product is replaced by an entire canonical
product. Its slower growth can be absorbed by a matrix regularizer,
allowing contour deformation through all matrix signatures.
\item Nonlinear transition values yield signed inequalities at contact
points. Positive exponential interpolation then turns those inequalities
into certificates against arbitrary positive competitors.
\end{enumerate}

The regularity and Bessel results additionally show why a boundary must
be analyzed before a supremum or minimization is taken. Analytic
population smoothness leaves an increasing number of active frequency
bands. At Bessel order $1/2$, exponentially small terms determine a
critical power that cannot be detected by a finite algebraic expansion.

\subsection{Sixteen-target register}

In the table, \textbf{E} means an externally stated target in its exact
scope, \textbf{F} a sharpened or restricted continuation of an external
question, \textbf{I} a problem formulated in this research program, and
\textbf{D} a dependent transfer or consequence. ``Proved'' means a
complete argument is presented in this manuscript, subject to external
mathematical review. The rows do not count parameter values, derivative
orders, or elementary lemmas as separate conjectures.

\begingroup\small
\setlength{\tabcolsep}{4pt}
\begin{longtable}{@{}p{.05\textwidth}p{.07\textwidth}>{\raggedright\arraybackslash}p{.37\textwidth}>{\raggedright\arraybackslash}p{.43\textwidth}@{}}
\toprule
No. & Type & Target & Result and location\\
\midrule\endfirsthead
\toprule No. & Type & Target & Result and location\\\midrule\endhead
1&E&Prime-zeta count at the boundary $\Re s>1$&$M_1(T)=\Theta(T)$; Theorem~\ref{p06:ra:thm:endpoint}.\\[4pt]
2&D&Derivative and fixed-level arithmetic transfer&Every fixed derivative has $\Theta(T)$ zeros; each level has an $O(T)$ upper bound; Theorem~\ref{p06:ra:thm:endpoint}.\\[4pt]
3&F&Complete unrestricted binomial-MLE curve&All $N,t$, including $t\asymp N/\log N$; Theorem~\ref{p06:bm:full}.\\[4pt]
4&D&Equal-replication Bernoulli budget&Exact order $(B\log(eB))^{-1/4}$; Theorem~\ref{p06:bm:equalbudget}.\\[4pt]
5&D&Analytic component-model budget transfer&Sufficient hypotheses and a verified factorial Legendre example; Theorem~\ref{p06:bm:transfer}.\\[4pt]
6&F&Dimension-three restricted exact information&Three queries are necessary and sufficient for polynomial or convex analytic information; Theorem~\ref{p06:ni:two-query}.\\[4pt]
7&I&Range-normalized convex noisy recovery&Exact infinite-budget floor $\delta/(2-\delta)$, including convex analytic and convex polynomial queries; Theorem~\ref{p06:ni:convex-floor}.\\[4pt]
8&I&Degree-limited polynomial noisy recovery&Sharp order $\delta/D$ with an adaptive construction; Theorem~\ref{p06:ni:degree-floor}.\\[4pt]
9&I&Matrix Gaussian tail converse at every finite logarithmic depth&Sharp threshold $a=1$ for polynomially sparse terminal sets in each fixed dimension; Theorem~\ref{p06:bg:thm:hierarchy}.\\[4pt]
10&I&Denser Gaussian terminal designs&The first-log converse holds for $n_j\ge c j\log^\beta(e+j)$, $\beta>1$; Theorem~\ref{p06:bg:thm:dense}.\\[4pt]
11&I&Normalized finite-spectrum threshold count&Exactly $2m-1$ suitably placed settings for at most $m$ atoms; Theorem~\ref{p06:rs:thm:counts}.\\[4pt]
12&I&Blind positivity certification&$2m$ normalized settings certify an $m$-atomic truth against every positive competitor; Theorem~\ref{p06:rs:thm:counts}.\\[4pt]
13&I&Blind spectral consistency under noise&Sharp uniform exponent $1/(2m)$; Theorem~\ref{p06:rs:thm:noise}.\\[4pt]
14&I&Bessel critical power and order endpoints&Inverse-square-root, logarithmic, and bounded regimes; Theorems~\ref{p06:bs:three-regimes} and~\ref{p06:bs:endpoints}.\\[4pt]
15&I&Stretched-exponential population regularity&Adaptive and fixed-design budget risk $\asymp\log\log B/\sqrt B$; Theorem~\ref{p06:sb:gevrey}.\\[4pt]
16&I&Double-exponential population regularity&Adaptive and fixed-design budget risk $\asymp\log\log\log B/\sqrt B$; Theorem~\ref{p06:sb:double}.\\
\bottomrule
\end{longtable}\endgroup

The approximate-MLE extension and the unknown-atom-count stopping rule
are useful further corollaries, not additional independent entries.
The SD moment hierarchy in Proposition~\ref{p06:bs:sd-hierarchy} is a
\emph{partial} advance, and is not included among the sixteen completed
targets. The linear positive-measure recovery principle and generic
moment exponents already occur in prior work
\cite{p06:rs:Eftekhari,p05:fz:WuYang}; the spectral contribution here is the
nonlinear threshold transfer with its precise hypotheses.

\subsection{Notation and verification status}

$W_1$ denotes Wasserstein distance with the specified ground metric.
The symbols $\asymp$, $O$, and $\Theta$ allow positive constants depending
only on the fixed parameters named in each statement. Dimension is
fixed in the matrix-tail results. Exact population thresholds, exact
real-valued queries, and finite noisy samples are distinct observation
models and are never silently interchanged.

All branches received independent automated proof readings, with
separate checks of adaptive conditional information, zero-cell
likelihood envelopes, singular matrix contours, confluent interpolation,
and parameter-uniform Bessel limits. These checks can detect errors
but are neither formal verification nor external peer review. Literature
comparisons were checked against primary sources available on
15 October 2026; absence of a matching result in that search does not
establish exhaustive priority. Journal-level significance must be judged
after independent expert evaluation of the proofs and their novelty.


% ===== arithmetic =====
\section{The prime-zeta counting problem at the boundary of absolute convergence}
\label{p06:ra:sec:endpoint}

Let
\[
 P(s)=\sum_{p\ {
m prime}}p^{-s}\qquad(\Re s>1).
\]
The question of Belovas, \v{C}epaityt\.{e}, and Sabaliauskas
\cite[Section~3]{p01:pf:BCS} includes the endpoint upper bound
\[
 M_1(T):=\#\{s:P(s)=0,\ \Re s>1,\ 0<\Im s<T\}=O(T).
\]
It is also recorded explicitly in \cite[p.~16, Conjecture~2.1.4 and the
sentence following it]{p06:ra:Tuzcu}. The preceding manuscripts proved
uniform counting statements in each fixed half-plane $\Re s>\gamma>1$,
but explicitly left this endpoint open. Those statements alone do not
justify sending $\gamma$ down to one. Here a different observation,
which retains the exact prime frequencies, gives the missing estimate.
The proof uses classical mean-square estimates for the Riemann zeta
function; it is an application of those estimates, not a new zeta
mean-value theorem or a claim about the Riemann hypothesis.

\begin{theorem}[Endpoint counting and the derivative tower]
\label{p06:ra:thm:endpoint}
For every fixed integer $m\ge0$, let $N_m(T)$ count the zeros of
$P^{(m)}$ in $\Re s>1$, $0<\Im s<T$, with multiplicity, and let
$N_m^{\rm dist}(T)$ count them without multiplicity. There are constants
$0<c_m<C_m<\infty$ such that, for all sufficiently large $T$,
\begin{equation}\label{p06:ra:eq:tower-count}
 c_m T\le N_m^{\rm dist}(T)\le N_m(T)\le C_mT.
\end{equation}
In particular the proposed endpoint estimate $M_1(T)=O(T)$ holds.
For every fixed $a\in\mathbb C$, the number of solutions of $P(s)=a$
in the same half-plane and height interval, counted with multiplicity,
is $O_a(T)$.
\end{theorem}

The case $m=0$ addresses the displayed external endpoint question.
The derivative tower and level-set upper bounds are transfers of the
same argument. They are not counted as infinitely many separately
solved published conjectures. No limiting coefficient $N_m(T)/T$ at
the endpoint, short-interval bound uniform in the initial height,
or zero-variance asymptotic is asserted.

\subsection{An averaged Jensen lemma}

We write $\log^+u=\max\{0,\log u\}$ for $u>0$, and set
$\log^+0=0$.

\begin{lemma}[Vertical logarithmic means control endpoint zero counts]
\label{p06:ra:lem:jensen}
Let $\alpha<\beta$, and let $G$ be a nonzero function holomorphic
on $\Re s>\alpha$. Suppose that there are $R>\beta$ and $c>0$ such
that $G$ has no zeros in $\Re s\ge R$ and
$|G(R+it)|\ge c$ for every real $t$. Suppose also that for every
compact interval $I\subset(\alpha,\infty)$,
\begin{equation}\label{p06:ra:eq:log-mean}
 \sup_{\sigma\in I}\int_{-T}^T
          \log^+|G(\sigma+it)|\,dt=O_I(T+1).
\end{equation}
Then the number of zeros of $G$ in $\Re s>\beta$, $0<\Im s<T$,
with multiplicity, is $O(T+1)$.
\end{lemma}

\begin{proof}
Choose fixed radii $R-\beta<r<\rho<R-\alpha$, and put
$h=(r^2-(R-\beta)^2)^{1/2}>0$. Let $n_r(t)$ count the zeros in
the closed disk of radius $r$ centered at $R+it$.
Jensen's formula gives
\[
 n_r(t)\log(\rho/r)
 \le \frac1{2\pi}\int_0^{2\pi}
       \log^+|G(R+it+\rho e^{i\theta})|\,d\theta-\log c.
\]
One may replace $-\log c$ by $|\log c|$ and use limiting radii at
the exceptional circles containing zeros. Integrating over
$-h<t<T+h$, applying Tonelli to the nonnegative integrand, and using
\eqref{p06:ra:eq:log-mean} for
$\sigma\in[R-\rho,R+\rho]$ shows
$\int_{-h}^{T+h}n_r(t)\,dt=O(T+1)$.
Every zero $x+iy$ with $\beta<x<R$, $0<y<T$, belongs to these disks
for a set of center ordinates of length
$2(r^2-(R-x)^2)^{1/2}\ge2h$. Thus its multiplicity contributes at
least $2h$ to that integral. There are no zeros farther right, which
proves the lemma.
\end{proof}

\subsection{Removal of the prime powers}

The function
\begin{equation}\label{p06:ra:eq:H}
 H(s)=\sum_p\sum_{k=2}^{\infty}\frac{p^{-ks}}k
\end{equation}
is holomorphic on $\Re s>1/2$. For each $j\ge0$ and each fixed
$\sigma_0>1/2$, its derivative $H^{(j)}$ is bounded uniformly on
$\Re s\ge\sigma_0$. Indeed its absolute derivative series is
majorized by
\[
 \sum_p(\log p)^j\sum_{k\ge2}k^{j-1}p^{-k\sigma_0}
 \le C_{j,\sigma_0}\sum_p(\log p)^jp^{-2\sigma_0}<\infty.
\]
The convention $k^{j-1}=k^{-1}$ applies when $j=0$.
On the original half-plane of absolute convergence,
\begin{equation}\label{p06:ra:eq:euler-lift}
 \log\zeta(s)=P(s)+H(s),\qquad
 \zeta(s)=\exp(P(s)+H(s)).
\end{equation}
Here the logarithm is the Euler-product branch on $\Re s>1$.
We never continue that branch across a zero of $\zeta$.

We use the classical mean-square estimate, locally uniform in
$\sigma>1/2$ away from the pole at height zero,
\begin{equation}\label{p06:ra:eq:zeta-ms}
 \sup_{\sigma\in I}\int_1^T|\zeta^{(j)}(\sigma+it)|^2\,dt
       =O_{I,j}(T+1),
 \qquad I\Subset(1/2,\infty),\quad j\ge0.
\end{equation}
The $j=0$ input is the standard off-critical-line mean-square theorem
\cite[Chapter~VII]{p06:ra:Titchmarsh}; see also
\cite[Lecture~II, Section~3]{p06:ra:Gonek} for its basic form.
For clarity, the derivative version needs no additional moment
conjecture. Choose a fixed Cauchy radius smaller than both $1/2$ and
the distance of $I$ from $1/2$, apply Cauchy's formula and Cauchy--Schwarz on the
circle, and then integrate in $t$. The circle's real parts remain in
a fixed larger compact interval, and its imaginary parts differ from
$t$ by a bounded amount. The $j=0$ bound therefore gives
\eqref{p06:ra:eq:zeta-ms}. A fixed initial interval is harmless, and the
corresponding negative-height estimate follows by conjugation.

\subsection{The endpoint upper bound}

For $a\in\mathbb C$, define the single-valued meromorphic function
\[
 F_a(s)=\zeta(s)-\exp(a+H(s)),\qquad \Re s>1/2.
\]
Its only possible pole in this half-plane is the simple pole at one.
On $\Re s>1$, identity \eqref{p06:ra:eq:euler-lift} gives
\begin{equation}\label{p06:ra:eq:overcount}
 F_a(s)=e^{H(s)}(e^{P(s)}-e^a).
\end{equation}
Consequently a solution of $P(s)=a$ is a zero of $F_a$ with exactly
the same multiplicity. The converse need not hold: $F_a=0$ also
allows $P(s)=a+2\pi i k$. Overcounting these additional solutions
is legitimate for an upper bound.

If $e^a\ne1$, then $F_a(s)\to1-e^a\ne0$ uniformly as
$\Re s\to\infty$. If $e^a=1$, the absolutely convergent expansions
give instead
\[
 2^sF_a(s)=1+o(1)
       \quad(\Re s\to\infty),
\]
uniformly in the ordinate: the leading term of $P$ is $2^{-s}$,
and $\zeta(s)(1-e^{-P(s)})=P(s)+O(2^{-2\Re s})$.
Set $b_a=0$ in the first case and $b_a=\log2$ in the second, and put
\[
 G_a(s)=e^{b_as}\frac{s-1}{s+1}F_a(s).
\]
The pole at one is removable, so $G_a$ is holomorphic on
$\Re s>1/2$. Its far-right limit is nonzero uniformly in the
ordinate: the extra rational factor tends to one uniformly there.
Thus for sufficiently large $R$ it is zero-free in $\Re s\ge R$
and bounded away from zero on $R+i\mathbb R$.
For real parts in a fixed compact interval and $|t|\ge1$,
\[
 \log^+|G_a(\sigma+it)|
       \le C+\log(1+|\zeta(\sigma+it)|)
       \le C+|\zeta(\sigma+it)|^2.
\]
The bounded factors include $e^{H}$ and $(s-1)/(s+1)$; the latter
has modulus at most one on $\Re s>0$. On the remaining compact
region, $G_a$ itself is bounded, including at the removed pole.
Equation~\eqref{p06:ra:eq:zeta-ms} verifies
\eqref{p06:ra:eq:log-mean}. Lemma~\ref{p06:ra:lem:jensen}, with
$\alpha=1/2$ and $\beta=1$, gives the asserted $O_a(T)$ upper bound.
For $a=0$ this proves $N_0(T)=O(T)$ without any interchange of limits.

\subsection{All fixed derivatives, using only second moments}

For $m\ge1$ define differential polynomials in $\zeta$ recursively by
\begin{equation}\label{p06:ra:eq:A-recursion}
 A_1=\zeta',\qquad A_{m+1}=\zeta A_m'-m\zeta'A_m.
\end{equation}
Where $\zeta\ne0$, induction gives
$A_m=\zeta^m(\log\zeta)^{(m)}$. In particular,
\begin{equation}\label{p06:ra:eq:derivative-lift}
 F_m:=A_m-\zeta^mH^{(m)}=\zeta^mP^{(m)}
          \qquad(\Re s>1).
\end{equation}
The definition on the left is single-valued and meromorphic on
$\Re s>1/2$, even at zeros of $\zeta$; no division by $\zeta$ is
used there. Its only possible pole is at one, of order at most $2m$.
This follows either from the recursion, or from the fact that every
monomial of $A_m$ has $m$ factors and total derivative order $m$.
Since $\zeta$ does not vanish for $\Re s>1$, the zeros of $F_m$ and
$P^{(m)}$ in that half-plane agree, including multiplicities.

Normalize by
\[
 G_m(s)=\frac{2^s}{(-\log2)^m}
              \left(\frac{s-1}{s+1}\right)^{2m}F_m(s).
\]
This is holomorphic on $\Re s>1/2$, and tends to one uniformly as
$\Re s\to\infty$. Indeed
$P^{(m)}(s)=(-\log2)^m2^{-s}(1+o(1))$ uniformly there and
$\zeta(s)\to1$.
The finite polynomial expression for $F_m$, boundedness of
$H^{(m)}$, and boundedness of the remaining factors on compact
real-part intervals give, for $|t|\ge1$,
\begin{equation}\label{p06:ra:eq:polynomial-log}
 \log^+|G_m(\sigma+it)|
       \le C_m+C_m\sum_{j=0}^m
                         \log(1+|\zeta^{(j)}(\sigma+it)|).
\end{equation}
Integration and the mean-square bound give the required logarithmic
mean estimate. The compact region around the removed pole contributes
a bounded amount. Lemma~\ref{p06:ra:lem:jensen} now proves $N_m(T)=O_m(T)$.
The logarithm in \eqref{p06:ra:eq:polynomial-log} is essential to the
economy of the argument: high powers in a differential polynomial do
not require high moments of $\zeta$ for this counting conclusion.

\subsection{The matching lower order from the exact prime flow}

For a fixed $m\ge0$, put $a_p=(\log p)^m$ and let $S_m>1$ be the
unique solution of
\[
 \sum_{p\ge3}\frac{a_p}{a_2}(2/p)^{S_m}=1.
\]
Existence follows because this sum decreases continuously from
infinity to zero as the real argument increases from one to infinity.
For $m=0$, divergence at one is Euler's divergence of the sum of
reciprocal primes; the other cases follow by comparison. Consider
\[
 Q_m(s)=a_22^{-s}-\sum_{p\ge3}a_pp^{-s}.
\]
It has a simple zero at $S_m$, because the derivative of its
normalization by $a_22^{-s}$ at that point is
\[
 \sum_{p\ge3}\frac{a_p}{a_2}(2/p)^{S_m}\log(p/2)>0.
\]
Choose a small closed disk $D$ about $S_m$, wholly inside
$\Re s>1$, containing exactly this one zero, and with no boundary
zeros. A sufficiently accurate approximation to the coefficient
phases $(1,-1,-1,\ldots)$ makes
$\sum_p a_p\omega_pp^{-s}$ uniformly close to $Q_m$ on $\partial D$.
This requires only a finite open cylinder of phase conditions: first
make the normally convergent tail uniformly small and then constrain
the finitely many remaining phases. Rouch\'e's theorem then gives
exactly one zero in $D$.

The exact prime flow $\omega_p(t)=e^{-it\log p}$ visits this cylinder
relatively densely. To see the precise input, finite nonzero integer
vectors satisfy $\sum_p k_p\log p\ne0$ by unique factorization.
Finite-dimensional Kronecker density, followed by compactness of the
finite torus, says that the visits to any nonempty open cylinder have
bounded gaps. More explicitly, translates of the cylinder by the
forward flow cover that compact torus; a finite subcover bounds the
waiting time from every starting point.
Thus some fixed $L_m$ ensures that every sufficiently long height
interval contains a zero of
$\sum_pa_pp^{-s}=(-1)^mP^{(m)}(s)$ in the fixed strip containing
$D$. Allowing for the disk's bounded imaginary radius only changes
$L_m$. Disjoint intervals of that length give at least $c_mT$
distinct zeros up to height $T$. This proves the lower bound in
\eqref{p06:ra:eq:tower-count} and completes the theorem.

\subsection{What this resolves, and what the arithmetic still does not give}

The external endpoint $M_1(T)=O(T)$ is obtained for the actual prime
zeta function, and the two-sided linear order extends to every fixed
derivative. This is stronger than a fixed $\gamma>1$ statement, but
it is not a claim that its proof supplies a new technique comparable
to a major zero-density breakthrough. The new step within this
program is to lift the logarithmic prime series to a single-valued
meromorphic function and then control its positive logarithmic mean.

The growing-dimensional recurrence problem uses a different input:
small values of $\sum_pk_p\log p$ with the support and coefficients
growing as the required phase error shrinks. The present argument uses
only qualitative finite-dimensional Kronecker density for the lower
bound, so it gives no entropy-sharp bound for those return times.
Likewise, a first-moment count of zeros provides no variance or
hyperuniformity theorem. Both the exact-prime recurrence exponent gap
and fixed-strip Euler-zero variance remain open here. No result is
claimed about zeros on different branches of $P$ to the left of one.


% ===== mle =====
\section{The complete binomial-mixture likelihood transition}
\label{p06:bm:sec}

The near-linear MLE theorem in the preceding manuscript left open the
transition $t\asymp N/\log N$. Its sharp deterministic inverse modulus
also showed that global count-Wasserstein error alone cannot close that
transition. We now use two different statistical norms on two different
parts of a test polynomial. The low-degree part is controlled by
likelihood transportation. The remaining polynomial has small amplitude,
and is controlled by a predictive Hellinger bound. This yields the full
replication curve, including its transition and its parametric plateau.

For a probability measure $G$ on $[0,1]$, put
\[
 K_x(j)=B_j^t(x)=\binom tj x^j(1-x)^{t-j},\qquad
 p_G(j)=\int B_j^t(x)\,dG(x),\quad 0\le j\le t.
\]
We observe $N$ independent counts from $p_G$. Their empirical law is $q$.
The unrestricted mixing-measure MLE is any maximizer $\widehat G$ of
$\ell(H)=\sum_jq_j\log p_H(j)$. Write $\widehat p=p_{\widehat G}$,
$L_N=\log(eN)$, and
\[
 \mathcal R_{N,t}=\inf_{\widetilde G}\sup_G
       \mathbb E_G W_1(\widetilde G,G),\qquad
 r_{N,t}=\min\left\{1,\frac1t+
           \frac1{\sqrt{tL_N}}+\frac1{\sqrt N}\right\}.
\]
The infimum allows randomized estimators. Wasserstein distance on the
count space uses the normalized cost $|j-l|/t$ and is denoted $W_{1,t}$.
Our Hellinger convention is
$h(p,v)^2=\sum_j(\sqrt{p_j}-\sqrt{v_j})^2$, also for nonnormalized
nonnegative vectors $v$.

\begin{theorem}[All-replication optimality of the unrestricted MLE]
\label{p06:bm:full}
There are absolute constants $0<c<C<\infty$ such that, for every pair of
integers $N,t\ge1$,
\begin{equation}\label{p06:bm:fullrate}
 c r_{N,t}\le \mathcal R_{N,t}
 \le \sup_G\mathbb E_G W_1(\widehat G,G)
 \le C r_{N,t}.
\end{equation}
The upper bound holds for every measurable choice of an unrestricted MLE.
In particular, at $t\asymp N/\log(eN)$ its risk is $O(N^{-1/2})$,
and the same order holds for every larger $t$.
\end{theorem}

The statistical question extends the replication-range conjectures of
Vinayak, Kong, Valiant and Kakade \cite{p01:info:VKVK2019} and the discussion
in Section~6.2 of Musco, Musco, Rosenblatt and Singh
\cite{p01:info:MMRS2025}. The preceding near-linear result already answered the
fixed-polynomial-range question in that discussion. The new assertion is
the complete transition, not a second claim of novelty for the old
near-linear range. We do not prove the separate constant-one
Chebyshev--Bernstein coefficient conjecture from those papers.

\subsection{An elementary finite-envelope bound for predictive likelihood}

The following entropy argument is stated with its details because it
must apply even when the true distribution has zero cells. Neither an
interior mixing density nor a regular parametric likelihood is assumed.

\begin{lemma}[Small covers of the binomial moment body]
\label{p06:bm:cover}
For $0<\beta<1$ the set $\{p_G:G\in\mathcal P([0,1])\}$ has a finite
$\ell_1$ cover of radius $\beta$, with centers in that same set, whose
cardinality $M$ satisfies
\begin{equation}\label{p06:bm:entropy}
 \log M\le m\log\frac{C(m+1)}\beta,
 \qquad
 m=\min\left\{t,
       \left\lceil\sqrt{2t\log(16t/\beta)}\right\rceil\right\}.
\end{equation}
\end{lemma}
\begin{proof}
View $K_x=(B_0^t(x),\ldots,B_t^t(x))$ as a vector-valued polynomial.
Its shifted Chebyshev expansion is
\[
 K_x=\sum_{k=0}^t v_kT_k(2x-1).
\]
We claim $\|v_k\|_1\le2\exp\{-k^2/(2t)\}$ for $k\ge1$.
Indeed, the even periodic function
$F(z)=K_{(1+\cos z)/2}$ has Fourier coefficient $v_k/2$.
For $z=\theta+iu$, the ellipse identity gives
\[
 \left|\frac{1+\cos z}{2}\right|
 +\left|\frac{1-\cos z}{2}\right|=\cosh u.
\]
The binomial formula therefore yields
$\|F(\theta+iu)\|_1=(\cosh u)^t\le e^{tu^2/2}$.
Shift the Fourier contour in the appropriate direction and take
$u=k/t$. This proves the claim. In particular, truncating after degree
$m<t$ has vector $\ell_1$ error at most
$2t e^{-(m+1)^2/(2t)}\le\beta/8$, uniformly in $x$.
For $m=t$ the error is zero.

Each mixing law contributes its first $m$ shifted Chebyshev moments in
$[-1,1]^m$. Partition that cube into cells of coordinate diameter at most
$\beta/[8(m+1)]$, and choose one mixing-law representative from every
nonempty cell. If $G,H$ lie in one cell, their truncated transforms differ
in $\ell_1$ norm by at most
$2m\beta/[8(m+1)]\le\beta/4$. The two tails contribute at most
$\beta/4$. Thus their predictive laws have distance at most $\beta$.
The number of cells is at most $[C(m+1)/\beta]^m$, as asserted.
\end{proof}

\begin{lemma}[Likelihood envelopes, including zero cells]
\label{p06:bm:envelopes}
Let $\varepsilon\ge0$ and $\eta>0$. Let $p$ be the true law on a finite alphabet with $d$ cells, and suppose a
model admits an $\ell_1$ cover of radius $\eta/d$, with $M$ probability
centers in the model. Let $\widetilde p$ be any fitted model law satisfying
\[
 \sum_{i=1}^N\log\widetilde p(X_i)
 \ge\sum_{i=1}^N\log p(X_i)-N\varepsilon.
\]
For every $x\ge0$, simultaneously for all such fitted laws, except on an
event of probability at most $e^{-x}$,
\begin{equation}\label{p06:bm:htail}
 h(p,\widetilde p)
 \le\sqrt\eta+
       \sqrt{\varepsilon+\eta+2(\log M+x)/N}.
\end{equation}
Consequently
\begin{equation}\label{p06:bm:hexpect}
 \mathbb E h(p,\widetilde p)^2
 \le C\left\{\varepsilon+\eta+\frac{\log M+1}{N}\right\}.
\end{equation}
\end{lemma}
\begin{proof}
For each center $p_a$, form the nonnegative vector
$u_a(j)=p_a(j)+\eta/d$. Every model law is bounded coordinatewise by at
least one $u_a$, and $\sum_j u_a(j)=1+\eta$. If $\widetilde p\le u_a$,
then
\[
 h(u_a,\widetilde p)^2\le\sum_j(u_a(j)-\widetilde p(j))=\eta.
\]
Thus $h(p,\widetilde p)>\sqrt\eta+r$ implies $h(p,u_a)>r$.
Moreover the likelihood comparison implies
$\prod_i u_a(X_i)/p(X_i)\ge e^{-N\varepsilon}$.
The ratios need only be defined on the support of $p$, where all sampled
observations lie almost surely. Markov's inequality with exponent $1/2$
gives
\begin{align*}
 \Pr\left\{\prod_i\frac{u_a(X_i)}{p(X_i)}
                      \ge e^{-N\varepsilon}\right\}
 &\le e^{N\varepsilon/2}
        \left(\sum_j\sqrt{p(j)u_a(j)}\right)^N\\
 &=e^{N\varepsilon/2}
        \left(1+\frac\eta2-\frac{h(p,u_a)^2}{2}\right)^N\\
 &\le\exp\{N(\varepsilon+\eta-r^2)/2\}
\end{align*}
when $h(p,u_a)>r$. Taking the union over the $M$ envelopes proves
\eqref{p06:bm:htail}. Integrating that tail, and using
$(a+b)^2\le2a^2+2b^2$, proves \eqref{p06:bm:hexpect}.
\end{proof}

\begin{corollary}[A sufficient predictive Hellinger rate]
\label{p06:bm:hellinger}
For $1\le t\le N$, every exact MLE satisfies
\begin{equation}\label{p06:bm:hbound}
 \mathbb E_G h(p_G,\widehat p)
       \le \frac{C t^{1/4}L_N^{3/4}}{\sqrt N}.
\end{equation}
More generally, if $\ell(\widetilde G)\ge\ell^*-\varepsilon$, the right
side can be increased by $C\sqrt\varepsilon$.
\end{corollary}
\begin{proof}
Set $\eta=(eN)^{-2}$ and $\beta=\eta/(t+1)$ in the preceding lemmas.
For $t\le N$, \eqref{p06:bm:entropy} gives
$m\le C\sqrt{tL_N}$ and
$\log M\le C\sqrt t\,L_N^{3/2}$. The MLE likelihood is at least that of
the true law. The approximate likelihood condition loses at most
$N\varepsilon$. Apply \eqref{p06:bm:hexpect} and Jensen's inequality.
\end{proof}

The logarithmic factors in \eqref{p06:bm:hbound} need not be sharp. The
argument below only needs the strict improvement from effective dimension
$t$ to a quantity of order $\sqrt t$ up to logarithms.

\subsection{The two-scale inversion inequality}

We record the two approximation ingredients used in the preceding
manuscripts. They are included with proofs to specify the uniformity
needed here.

\begin{lemma}[Stable approximation and bounded-polynomial inversion]
\label{p06:bm:approx}
There is an absolute $C$ with the following properties.
For every $1$-Lipschitz $f:[0,1]\to\mathbb R$ and integer $k\ge1$, there
is a polynomial $P_k$ of degree at most $k$ such that
\[
 \|P_k-f\|_\infty\le C/k,\qquad \|P_k'\|_\infty\le2.
\]
If $P=\sum_{j=0}^t b_jB_j^t$ has degree $k\le t/2$, then
\begin{equation}\label{p06:bm:coeff}
 \max_j|b_j|\le Ce^{2k^2/t}\|P\|_\infty.
\end{equation}
For the derivative-stable $P_k$, its coefficient vector also satisfies
\begin{equation}\label{p06:bm:coefflip}
 \max_{j<t}t|b_{j+1}-b_j|\le Ce^{4k^2/t}.
\end{equation}
\end{lemma}
\begin{proof}
Subtract $f(0)$, reflect $f$ on $[0,1]$ to a period-two $1$-Lipschitz
function, and convolve with the normalized positive Jackson kernel
$c_s[\sin(s\pi u/2)/\sin(\pi u/2)]^4$. Its first absolute moment is
$O(1/s)$, so the resulting trigonometric polynomial $Q_s$ has error
$O(1/s)$ and derivative norm at most one. Its frequencies have magnitude
at most $2s$, and its Fourier coefficients have modulus at most one.
Replacing each exponential by its Taylor polynomial at $1/2$ of degree
$100s$ changes the function and derivative by at most $e^{-100s}$.
Indeed $|\pi l(x-1/2)|\le\pi s$ and the factorial inequality
$M!\ge(M/e)^M$ bounds both remainders, including the $4s+1$ terms and the
extra derivative factor. Take $s=\lfloor k/100\rfloor$ for $k\ge100$;
use a constant polynomial and enlarge $C$ for smaller $k$.

For \eqref{p06:bm:coeff}, put $M=\|P\|_\infty$ and
$F(z)=(1+z)^tP(z/(1+z))$. The maximum principle in the inverse Joukowski
variable gives
$|P((2+w+w^{-1})/4)|\le M|w|^k$ for $|w|\ge1$.
For $0<j<t$, set $r=j/(t-j)$, $v=r/(1+r)^2$, and
$u=2\sqrt v|\sin(\phi/2)|$. The larger inverse Joukowski modulus for
$re^{i\phi}/(1+re^{i\phi})$ is $(1+u)/\sqrt{1-u^2}$. Hence
\[
 \frac{|F(re^{i\phi})|}{M(1+r)^t}
 \le(1-u^2)^{(t-k)/2}(1+u)^k
 \le e^{2k^2/t}e^{-tu^2/8}.
\]
The last inequality follows from $k\le t/2$ and completing the square.
The average of the final Gaussian factor is at most
$C/\sqrt{1+tv}$. Cauchy's coefficient formula, together with
$\binom tj(j/t)^j(1-j/t)^{t-j}\ge c/\sqrt{1+tv}$, proves
\eqref{p06:bm:coeff}. The last lower bound follows directly from Stirling's
factorial bounds; endpoint coefficients are $P(0)$ and $P(1)$.
Finally differentiate the Bernstein expansion:
$P_k'=t\sum_{j<t}(b_{j+1}-b_j)B_j^{t-1}$.
Apply \eqref{p06:bm:coeff} at degree $k-1\le(t-1)/2$ and use
$2(k-1)^2/(t-1)\le4k^2/t$.
\end{proof}

\begin{lemma}[Likelihood transport, recalled]
\label{p06:bm:transport}
For every binomial-mixture MLE,
\begin{equation}\label{p06:bm:trans}
 W_{1,t}(q,\widehat p)^2\le\frac{D(q\Vert\widehat p)}{2t}
                    \le\frac{D(q\Vert p_G)}{2t}.
\end{equation}
Consequently
\begin{equation}\label{p06:bm:predictive}
 \mathbb E_G W_{1,t}(\widehat p,p_G)\le C/\sqrt N.
\end{equation}
\end{lemma}
\begin{proof}
In score sums, terms with $q_j=0$ are set to zero; every observed
cell has positive fitted probability. The likelihood score condition gives
$S(x)=\sum_jq_jK_x(j)/\widehat p(j)\le1$, and its integral against
$\widehat G$ is one. Thus $S=1$ almost surely under $\widehat G$.
The posterior vectors $R_x(j)=q_jK_x(j)/\widehat p(j)$ are therefore
probability laws, with mixture $q$, and
\[
 \int D(R_x\Vert K_x)\,d\widehat G(x)=D(q\Vert\widehat p).
\]
Every normalized-count binomial kernel satisfies
$W_{1,t}(R,K_x)^2\le D(R\Vert K_x)/(2t)$: a $1$-Lipschitz test of the
normalized sum of $t$ Bernoulli variables has centered logarithmic moment
generating function at most $\lambda^2/(8t)$ by bounded differences;
entropy duality and optimization over $\lambda$ give the inequality.
Mix the couplings, use Cauchy--Schwarz, and then use likelihood optimality
to prove \eqref{p06:bm:trans}. Also
$\mathbb E D(q\Vert p_G)\le t/N$, by
$D(q\Vert p_G)\le\sum_{j:p_G(j)>0}(q_j-p_G(j))^2/p_G(j)$.
Finally $\mathbb E W_{1,t}(q,p_G)\le1/(2\sqrt N)$ by integrating the
standard deviation of the empirical distribution function. The triangle
inequality proves \eqref{p06:bm:predictive}.
\end{proof}

\begin{proposition}[Two-scale deterministic inversion]
\label{p06:bm:twoscale}
For integers $1\le s\le k\le t/2$ and any two mixing laws $G,H$,
\begin{equation}\label{p06:bm:twoscaleineq}
 W_1(G,H)\le \frac Ck
       +Ce^{4s^2/t}W_{1,t}(p_G,p_H)
       +\frac{Ce^{2k^2/t}}s h(p_G,p_H).
\end{equation}
\end{proposition}
\begin{proof}
For any $1$-Lipschitz test $f$, take the polynomials $P_k,P_s$ from
Lemma~\ref{p06:bm:approx}. The approximation error in integrating $P_k$
instead of $f$ against $G-H$ is at most $C/k$. The Bernstein coefficients
of $P_s$ have normalized-grid Lipschitz constant at most $Ce^{4s^2/t}$.
On the other hand,
\[
 \|P_k-P_s\|_\infty\le C/k+C/s\le C/s.
\]
Thus the Bernstein coefficients of $P_k-P_s$ have supremum norm at most
$Ce^{2k^2/t}/s$. Test $P_s$ by count transportation and $P_k-P_s$ by
$\|p_G-p_H\|_1\le2h(p_G,p_H)$. All constants are uniform over $f$,
so taking the Wasserstein dual supremum proves the result.
\end{proof}

The small amplitude $1/s$ is essential. Applying only the Hellinger
bound to the entire polynomial loses the low-frequency empirical
$N^{-1/2}$ scale; applying only transportation loses the critical
transition. Proposition~\ref{p06:bm:twoscale} retains both.

\begin{proof}[Upper bound in Theorem~\ref{p06:bm:full}]
It suffices first to treat $t\le N$. Bounded $N$ or bounded $t$ are handled
by the loss bound one. Fix a sufficiently small absolute $a>0$, for
example $a=1/64$, and take
\[
 k=\left\lfloor\min\{t/2,\sqrt{atL_N}\}\right\rfloor,
 \qquad s=\lfloor\sqrt t/4\rfloor.
\]
For sufficiently large $N$ and $t$, these satisfy
$1\le s\le k\le t/2$, $s\asymp\sqrt t$, and
\[
 1/k\le C\{1/t+(tL_N)^{-1/2}\},\qquad e^{4s^2/t}\le C.
\]
Use Proposition~\ref{p06:bm:twoscale}, Lemma~\ref{p06:bm:transport}, and
Corollary~\ref{p06:bm:hellinger}. This gives
\begin{equation}\label{p06:bm:master}
 \mathbb E_G W_1(\widehat G,G)
 \le \frac Ck+\frac C{\sqrt N}
       +\frac{Ce^{2k^2/t}L_N^{3/4}}{\sqrt N\,t^{1/4}}.
\end{equation}
The final term, multiplied by $\sqrt{tL_N}$, is at most
\[
 C(eN)^{2a}N^{-1/2}t^{1/4}L_N^{5/4}
 \le C(eN)^{2a}N^{-1/4}L_N^{5/4}\le C.
\]
The last bound is uniform because $2a<1/4$. This proves the desired
upper bound when $t\le N$.

For $t\ge N$, couple every mixing parameter $x$ to its normalized
binomial count. The expected absolute difference is at most
$\sqrt{x(1-x)/t}\le1/(2\sqrt t)$. Hence, after identifying count laws
with measures on $[0,1]$,
\[
 W_1(\widehat G,G)
 \le 1/\sqrt t+W_{1,t}(\widehat p,p_G).
\]
Take expectations and apply \eqref{p06:bm:predictive}. This is $O(N^{-1/2})$
for $t\ge N$, and completes the proof of the upper bound.
\end{proof}

\subsection{Matching lower bounds and the equal-replication budget}

\begin{proof}[Lower bound in Theorem~\ref{p06:bm:full}]
The $N^{-1/2}$ term follows from mixing laws supported on $\{0,1\}$ with
weights $1/2\pm c/\sqrt N$. Each count then reveals exactly one endpoint
label, regardless of $t$. Their $N$-sample total variation stays bounded
away from one for sufficiently small absolute $c$.

For the other terms let $\pi(dx)=dx/[\pi\sqrt{x(1-x)}]$ and
$h_k(x)=T_k(2x-1)$. Define
$G_k^\pm=(1\pm h_k/2)\pi$. Integration of their signed cumulative
function gives
\[
 W_1(G_k^+,G_k^-)
  =\frac1{2\pi k}\int_0^\pi|\sin(k\theta)|\sin\theta\,d\theta
  \asymp1/k.
\]
If $k=t+1$, the count laws are identical, proving the $1/t$ lower bound.
For $k\le t$, draw $X\sim\pi$, $J\mid X\sim\operatorname{Binomial}(t,X)$
and put $Qh(j)=\mathbb E[h(X)\mid J=j]$. The posterior is
$\operatorname{Beta}(j+1/2,t-j+1/2)$. The self-adjoint operator $Q^*Q$
preserves polynomial degrees and has leading degree-$k$ coefficient
\[
 \lambda_{t,k}=\prod_{i=0}^{k-1}\frac{t-i}{t+1+i}.
\]
Orthogonality of the Chebyshev polynomials therefore gives
$Q^*Qh_k=\lambda_{t,k}h_k$ and
$\mathbb E[Qh_k(J)^2]=\lambda_{t,k}/2$.
This is the classical beta--binomial diagonalization
\cite[Section~4.1]{p05:fm:DKS2008}.
The count laws are $p_\pi(1\pm Qh_k/2)$, so
\[
 \chi^2(p_{G_k^+}\Vert p_{G_k^-})
 \le2\mathbb E[Qh_k(J)^2]=\lambda_{t,k}
 \le e^{-k^2/(2t)}.
\]
Take $k=\lceil\sqrt{8t\log(8N)}\rceil$ if this is at most $t$.
The chi-squared divergence of the $N$-fold products is
$(1+\chi^2)^N-1\le e^{N(8N)^{-4}}-1$, so their total variation is
bounded by an absolute constant smaller than one. The two-point testing
bound yields risk at least $c/k\ge c'/\sqrt{tL_N}$.
If this choice has $k>t$, then $t\le C L_N$ apart from bounded $t$,
and $1/\sqrt{tL_N}\le C'/t$; the exact moment-matching lower bound
already suffices. The maximum of the three lower bounds is comparable
to their sum clipped at one, proving \eqref{p06:bm:fullrate}.
\end{proof}

\begin{theorem}[Exact equal-replication budget order]
\label{p06:bm:equalbudget}
Suppose a total of at most $B$ Bernoulli observations may be used, with
one common replication count $t$ for every sampled individual. Then
\begin{equation}\label{p06:bm:budgetrate}
 \inf_{N,t\ge1:\,Nt\le B}\mathcal R_{N,t}
       \asymp\{B\log(eB)\}^{-1/4},\qquad B\ge1.
\end{equation}
One order-optimal choice is
\[
 N\asymp\sqrt{B\log(eB)},\qquad
 t\asymp\sqrt{B/\log(eB)}.
\]
The unrestricted MLE attains this order. This closes the logarithmic gap
in the equal-group Bernoulli allocation problem left in the previous
manuscripts; it does not change the previously proved, smaller risk
available with unequal group sizes.
\end{theorem}
\begin{proof}
For any admissible $N,t$, the lower bound in Theorem~\ref{p06:bm:full}
contains both $cN^{-1/2}$ and $c(t\log(eN))^{-1/2}$. Their geometric
mean is at least
$c\{Nt\log(eN)\}^{-1/4}\ge c\{B\log(eB)\}^{-1/4}$.
The clipped case only alters the absolute constant. The displayed choices
of $N,t$, rounded down while preserving $Nt\le B$, give the matching
upper bound by Theorem~\ref{p06:bm:full}. Bounded budgets are immediate.
\end{proof}

\begin{theorem}[A matching budget transfer to analytic component models]
\label{p06:bm:transfer}
Let a component model $P_\theta$, $0\le\theta\le1$, admit a statistic
$Y$ based on $b_0$ independent component observations such that
$|Y|\le M$ and $\mathbb E_\theta Y=\theta$, where $b_0,M$ are fixed.
Then its equal-replication population problem has the upper bound
\[
 \mathcal R_B^{\mathrm{eq}}\le C_{M,b_0}
                  \{B\log(eB)\}^{-1/4}.
\]
Suppose in addition that, for some interior $\theta_0$ and $R>0$, the
likelihood ratios $f_z=dP_{\theta_0+z}/dP_{\theta_0}$ admit a holomorphic
extension to $|z|\le R$ in $L^2(P_{\theta_0};\mathbb C)$ satisfying
\begin{equation}\label{p06:bm:analyticcondition}
 \int f_z\,dP_{\theta_0}=1,\qquad
 \|f_z-1\|_2\le C_0|z|.
\end{equation}
Then $\mathcal R_B^{\mathrm{eq}}\asymp\{B\log(eB)\}^{-1/4}$.
Constants may depend on the displayed model parameters. The analytic
condition is a sufficient lower-bound condition, not a consequence of
existence of an unbiased statistic.
\end{theorem}
\begin{proof}
With an independent uniform $U$, the variable
$Z=\mathbf1_{\{U\le(Y+M)/(2M)\}}$ is Bernoulli with parameter
$(\theta+M)/(2M)$. Disjoint observation blocks give independent such
variables conditionally on $\theta$. Apply Theorem~\ref{p06:bm:equalbudget}
with budget $\lfloor B/b_0\rfloor$ and push the estimated mixing law
through the clipped inverse $z\mapsto\min(1,\max(0,2Mz-M))$, which is
$2M$-Lipschitz. This proves the upper bound.

We include the lower-bound argument, already used in the preceding
manuscript, to specify the scope of this transfer. Let $q$ be the common
number of component observations per individual. An oracle revealing
all $N$ latent parameters gives the lower bound $cN^{-1/2}$ by varying
the weights of a two-point mixing law. For the other lower bound, write
\[
 F_z=f_z^{\otimes q}=\sum_{j\ge0}T_jz^j
 \quad\hbox{in }L^2(P_{\theta_0}^{\otimes q};\mathbb C).
\]
The zero integral of $f_z-1$ cancels the cross term, so
$\|F_z\|_2\le e^{C_0^2q|z|^2/2}$. For $0<h\le R/2$, Cauchy's estimate
on $|z|=2h$ gives
$\|T_j\|_2\le e^{2C_0^2qh^2}(2h)^{-j}$.
Let $n=m+1$ and set $\theta_j=\theta_0-h+2hj/n$. The even and odd
binomial-weight laws
\[
 G=2^{1-n}\sum_{j\ {\rm even}}\binom nj\delta_{\theta_j},\qquad
 H=2^{1-n}\sum_{j\ {\rm odd}}\binom nj\delta_{\theta_j}
\]
match moments through degree $m$ and have $W_1(G,H)=2h/(m+1)$.
The latter identity follows from the signed cumulative values
$2^{1-n}(-1)^j\binom{n-1}{j}$ between successive grid points.
If $Q_q(G)=\int P_\theta^{\otimes q}\,dG(\theta)$, then
\[
 \left\|\frac{dQ_q(G)-dQ_q(H)}{dP_{\theta_0}^{\otimes q}}\right\|_2
 \le2e^{2C_0^2qh^2}\sum_{j>m}2^{-j}
 =2e^{2C_0^2qh^2}2^{-m}.
\]
Take
\[
 m=\left\lceil\frac4{3\log2}\log(4N)\right\rceil,\qquad
 h=\min\left\{R/2,
        \sqrt{\frac{m\log2}{8\max(1,C_0^2)q}}\right\},
\]
after decreasing $R$ to keep all parameters in $[0,1]$.
One-group total variation is at most $1/(4N)$; the $N$-group distance
is at most $1/4$. Therefore
\[
 \mathcal R_{N,q}\ge c\max\left\{N^{-1/2},
          \min\{L_N^{-1},(qL_N)^{-1/2}\}\right\}.
\]
If $q\ge L_N$, the geometric mean of the first and last terms is
$(NqL_N)^{-1/4}\ge\{B\log(eB)\}^{-1/4}$. If $q<L_N$, the lower bound
$c/L_N\ge c/\log(eB)$ is larger than the required budget rate for all
sufficiently large $B$. Bounded budgets are absorbed in the constant.
\end{proof}

\begin{corollary}[The factorial Legendre equal-group gap is closed]
\label{p06:bm:factorial}
Let $L_j$ denote shifted Legendre polynomials and consider
\[
 k_*(x\mid\theta)=1+\sum_{j\ge2}\frac{L_j(x)L_j(\theta)}{j!},
 \qquad x,\theta\in[0,1].
\]
For this model the optimal expected population $W_1$ risk under common
replication and total budget $B$ is
$\Theta(\{B\log(eB)\}^{-1/4})$.
\end{corollary}
\begin{proof}
The kernel is bounded below by $3-e>0$. The factorial series and its
parameter derivatives converge uniformly on compact complex parameter
sets, so \eqref{p06:bm:analyticcondition} holds at $\theta_0=1/2$.
For $j\ge2$ define $A_j(x)=j!(2j+1)L_j(x)$; orthogonality gives
$\mathbb E_\theta A_j(X)=L_j(\theta)$. The Legendre product identity
\[
 L_2L_3=\frac{10}{21}L_5+\frac4{15}L_3+\frac9{35}L_1
\]
then shows that the two-observation statistic
\[
 Y(x,y)=\frac12+\frac{35}{18}\left\{
 A_2(x)A_3(y)-\frac{10}{21}A_5(x)-\frac4{15}A_3(x)\right\}
\]
has expectation $\theta$ and absolute value less than $2100$.
Apply Theorem~\ref{p06:bm:transfer}. The unbiased statistic and the lower
bound were established previously; the new ingredient completing the
matching upper bound is Theorem~\ref{p06:bm:full}.
\end{proof}

\subsection{Approximate likelihood and the remaining coefficient question}

\begin{corollary}[The complete rate survives $N^{-1}$ optimization error]
\label{p06:bm:approxmle}
If a fitted mixing law satisfies
$\ell^*-\ell(\widetilde G)\le1/N$, then
\[
 \sup_G\mathbb E_GW_1(\widetilde G,G)\le C r_{N,t}
\]
for all $N,t\ge1$. The likelihood gap is measured per observation.
\end{corollary}
\begin{proof}
We give the needed approximate transport argument. As above, all score
sums are over observed cells, with zero terms elsewhere. Set $H=\widetilde G$,
$p_H=p$, and $S(x)=\sum_jq_jK_x(j)/p(j)$. The measure
$H^+=S H$ is a probability measure. For $S(x)>0$ set
$R_x(j)=q_jK_x(j)/(p(j)S(x))$. Then
\[
 \int D(R_x\Vert K_x)\,dH^+(x)
       =D(q\Vert p)-D(H^+\Vert H).
\]
The same component transport inequality as before implies
$W_{1,t}(q,p_{H^+})\le\sqrt{D(q\Vert p)/(2t)}$.
Furthermore Jensen's inequality applied to the posterior mixing law
at each observed cell gives
$D(H^+\Vert H)\le\ell(H^+)-\ell(H)\le1/N$.
Pinsker's inequality therefore gives
$W_{1,t}(p_{H^+},p_H)\le\|H^+-H\|_{\rm TV}\le1/\sqrt{2N}$.
Finally $D(q\Vert p_H)\le D(q\Vert p_G)+1/N$.
Thus $\mathbb E W_{1,t}(p_H,p_G)\le C/\sqrt N$ still holds.
Corollary~\ref{p06:bm:hellinger} adds at most $C/\sqrt N$ to its original
right side. Repeating \eqref{p06:bm:master} absorbs this addition because
$t\ge1$ and $L_N\ge1$. The high-replication kernel-coupling argument is
unchanged.
\end{proof}

The constant-one inequality
$|C(t,m,j)|\le e^{m^2/t}$ for the Bernstein coefficients of a shifted
Chebyshev polynomial remains open in this manuscript. The full statistical
curve does not require it. Nor do we obtain optimal leading constants,
sharp predictive Hellinger logarithms, or a finite-sample confidence band
for the entire mixing law. Those are separate questions from the
order-optimal unrestricted-MLE result proved above.

% ===== information =====
\section{Restricted adaptive information: exact lower counts and sharp noise floors}
\label{p06:ni:section}

Krieg--Novak--Ullrich ask how adaptive recovery changes when the
admissible scalar measurements are restricted to convex, homogeneous,
or analytic functions, and ask for lower bounds already in dimension
three~\cite[Section~3]{p03:an:KNU}.  Their question concerns exact real
answers.  The earlier manuscript established logarithmic upper bounds
using convex polynomial measurements.  We first give a lower bound
which is sharp in dimension three, and then show that range-normalized
convex information has a very different behavior under adversarial
noise.  For the latter model the relevant comparison is
Krieg--Novak--Plaskota--Ullrich~\cite{p06:ni:KNPU}: unrestricted continuous
information can correct bounded deterministic noise to arbitrarily
small reconstruction error.  Convex information cannot do so.

The results below concern deterministic adaptive algorithms.  A
measurement chosen after a transcript must have the stated properties
as a function of the unknown input, for that fixed transcript.  The
selection rule and decoder need not be continuous.  The exact-real
and noisy models are kept separate throughout.

\subsection{Two exact queries cannot suffice with a polynomial or quasiconvex first query}

A function is called quasiconvex if all its sublevel sets are convex.
No differentiability is part of this definition.

\begin{theorem}[A two-query obstruction]
\label{p06:ni:two-query}
Let $m\ge3$.  Consider an algorithm on $\mathbb R^m$ which first
observes $f(x)$ exactly, then observes $g_{f(x)}(x)$ exactly, where
every $g_t:\mathbb R^m\to\mathbb R$ is continuous.  Suppose that
$f$ is either a real polynomial or a continuous quasiconvex function.
For every decoder $\Phi:\mathbb R^2\to\mathbb R^m$,
\[
 \sup_{x\in\mathbb R^m}
 \bigl\|x-\Phi\bigl(f(x),g_{f(x)}(x)\bigr)\bigr\|_2=\infty.
\]
The same conclusion holds with any norm on $\mathbb R^m$.
In particular, the minimum number of adaptive exact measurements
needed for uniformly bounded-error recovery on $\mathbb R^3$ is
exactly three in each of the following classes: real polynomials;
globally convex polynomials; and globally convex real-analytic
functions.  In the last class the measurements may additionally be
required to have Lipschitz constant at most one.
\end{theorem}

\begin{proof}
We use an elementary consequence of connectedness.  Suppose that for
arbitrarily large $R$ some level of $f$ contains a continuous graph
\[
 \gamma(z)=z+\tau(z)v,\qquad
 z\in H,\quad \|z-z_0\|_2\le R,
 \qquad H=v^\perp,
\]
where $v$ is a unit vector.  On the sphere
$\{z\in H:\|z-z_0\|_2=R\}$ the function
\[
 g_t(\gamma(z_0+w))-g_t(\gamma(z_0-w))
\]
is continuous and odd in $w$.  This sphere has dimension $m-2\ge1$
and is connected, so the function vanishes somewhere.  The two
corresponding inputs have identical transcripts and are at distance
at least $2R$, because their orthogonal projections onto $H$ are
antipodal.  At least one reconstruction error is at least $R$.

Suppose first that $f$ is a nonconstant polynomial of degree $d$.
Choose a direction $v$ on which its highest homogeneous part is
nonzero, and negate $f$ if necessary.  In the coordinates $z+sv$,
with $z\in H$, we have
\[
 f(z+sv)=a s^d+\sum_{j<d}a_j(z)s^j,\qquad a>0.
\]
On any fixed transverse ball of radius $2R$ all coefficients and
their derivatives are bounded.  Thus there is a $T>0$ such that
$\partial_s f(z+sv)>0$ for all $s\ge T$ and all $z$ in that ball,
and $f(z+sv)\to\infty$ uniformly there as $s\to\infty$.
Choose $t>\max_{\|z\|_2\le2R}f(z+Tv)$.  The equation
$f(z+sv)=t$ has a unique solution $s=\tau(z)>T$ on the ball.
The implicit function theorem makes $\tau$ continuous, and the
preceding graph argument applies.

Now suppose that $f$ is continuous and quasiconvex.  Write
$s=\sup_{\mathbb R^m}f\in\mathbb R\cup\{\infty\}$ and
$U=\{x:f(x)<s\}$, interpreting $U=\mathbb R^m$ when $s=\infty$.
If $U$ is a proper subset, it is an open convex set.  Unless it is
empty, separation of a proper open convex set gives a halfspace
disjoint from $U$.  On that halfspace $f=s$, and it contains balls
of arbitrarily large radius.  A single continuous second query on
the sphere of such a ball identifies an antipodal pair.  The case
$U=\varnothing$, including constant $f$, is immediate in the same way.

It remains to consider $U=\mathbb R^m$.  For any $R>0$, choose
\[
 \max_{\|x\|_2\le2R}f(x)<t<s.
\]
Then $K=\{x:f(x)\le t\}$ is a proper closed convex set containing
the radius-$2R$ ball in its interior.  Separation gives a unit vector
$v$ and a finite $b$ such that $\langle v,x\rangle\le b$ on $K$.
For $z\in v^\perp$ with $\|z\|_2<2R$, set
\[
 \tau(z)=\sup\{a:z+av\in K\}.
\]
This is finite, positive and concave on that ball.  A finite concave
function is continuous on the interior of its domain.  Its graph
lies in the boundary of $K$, hence in $\{f=t\}$ by continuity of
$f$.  The graph argument again applies.

Equivalence of norms proves the assertion for every norm.  Three
coordinate measurements recover $\mathbb R^3$ exactly and belong
to all the stated classes, including the unit-Lipschitz class.
\end{proof}

\paragraph{Scope of the lower bound.}
The theorem resolves the three-dimensional count for the indicated
restricted classes.  It does not settle whether two arbitrary
continuous measurements can approximate every vector in
$\mathbb R^3$ with prescribed uniform accuracy, the unrestricted
question in~\cite{p03:an:KNU}.  Nor does it determine the optimal count
in higher dimensions.  The unbounded input domain is essential to
the argument; it is not a claim that every two-query method on a
fixed compact ball has infinite error.

\subsection{The exact noise floor for bounded convex measurements}

Let $C_m=[-1,1]^m$.  For a class $\Lambda$ of real functions on
$C_m$ with values in $[-1,1]$, let
$e_n(\Lambda,\delta)$ denote the infimum of the worst-case
$\ell_\infty$ recovery error over algorithms using at most $n$
adaptive observations
\begin{equation}\label{p06:ni:noise-model}
 y_j=\lambda_j(x;y_1,\ldots,y_{j-1})+\xi_j,
 \qquad |\xi_j|\le\delta.
\end{equation}
The supremum is over $x\in C_m$ and all admissible noise sequences.
The range bound applies to every selected query on the whole cube,
not merely on the currently feasible subset.  Without this
normalization, multiplying a query by an arbitrary constant would
make the noise restriction vacuous.

Write $\Lambda_{\rm cvx}$ for restrictions of globally convex
functions to $C_m$, subject to this range bound, and
$\Lambda_{\rm ca}\subset\Lambda_{\rm cvx}$ for those which are
also globally real analytic.  Let $\Lambda_{\rm cp}$ be the
subclass consisting of globally convex polynomials, with no common
degree restriction.

\begin{theorem}[Sharp convex noise floor]
\label{p06:ni:convex-floor}
For every $m\ge1$ and $0<\delta<1$, put
\[
 r_\delta=\frac{\delta}{2-\delta}.
\]
For every $n\ge0$,
\begin{equation}\label{p06:ni:convex-lower}
 e_n(\Lambda_{\rm cvx},\delta)\ge r_\delta.
\end{equation}
For every integer $k\ge0$,
\begin{equation}\label{p06:ni:convex-upper}
 e_{mk}(\Lambda_{\rm cp},\delta)
 \le r_\delta+
       \left(\frac\delta2\right)^k(1-r_\delta).
\end{equation}
Consequently,
\begin{equation}\label{p06:ni:convex-limit}
 \inf_{n\ge0}e_n(\Lambda_{\rm cvx},\delta)
 =\inf_{n\ge0}e_n(\Lambda_{\rm ca},\delta)
 =\inf_{n\ge0}e_n(\Lambda_{\rm cp},\delta)
 =\frac{\delta}{2-\delta}.
\end{equation}
The lower bound remains true if the admissible functions are only
required to be convex on the cube.
\end{theorem}

\begin{proof}[Lower bound]
For a convex $h:[-1,1]\to[-1,1]$ and $0\le t<1$, convexity on
the interval $[-t,1]$ gives
\[
 h(t)-h(-t)
 \le\frac{2t}{1+t}\{h(1)-h(-t)\}
 \le\frac{4t}{1+t}.
\]
Applying the same argument with the orientation reversed shows
\begin{equation}\label{p06:ni:convex-chord}
 |h(t)-h(-t)|\le\frac{4t}{1+t}.
\end{equation}
Set $x_\pm=\pm r_\delta e_1$.  Every admissible query, restricted
to the central first-coordinate line, therefore satisfies
\[
 |\lambda(x_+)-\lambda(x_-)|\le2\delta.
\]
At every step return the midpoint of these two query values.  The
answer is within $\delta$ of the correct answer under both inputs.
Inductively the algorithm has the same transcript under $x_+$ and
$x_-$.  Any output has error at least $r_\delta$ for one of them.
This works for every finite number of adaptive queries.
\end{proof}

\begin{proof}[Upper bound for convex queries]
It suffices to recover one coordinate $z\in[-1,1]$.  Maintain a
closed feasible interval $I=[l,u]$, with center $c$ and radius $r$.
If $I$ is a singleton, no further query is needed.  Otherwise choose
whichever of the following two convex functions has the smaller
denominator:
\begin{align*}
 g_+(z)&=-1+\frac{2(z-l)_+}{1-l},\\
 g_-(z)&=-1+\frac{2(u-z)_+}{1+u}.
\end{align*}
Both are globally convex and lie in $[-1,1]$ on $[-1,1]$.
On $I$ the chosen function is affine, with slope magnitude
$2/\min(1-l,1+u)$.  Intersect $I$ with the interval compatible
with its noisy answer.  The new feasible interval has radius at most
\[
 \frac\delta2\min(1-l,1+u)
 =\frac\delta2(1+r-|c|)
 \le\frac\delta2(1+r).
\]
Starting from $r_0=1$, iteration yields
\[
 r_k\le r_\delta+(\delta/2)^k(1-r_\delta).
\]
Run this construction for each coordinate and output the midpoints.
It uses $mk$ queries and proves the stated upper bound for convex
measurements.
\end{proof}

\begin{proof}[Real-analytic and polynomial convex queries]
The hinges can be approximated without abandoning global convexity
or the range normalization.  Set
\[
 H_\tau(s)=\tau\log(1+e^{s/\tau}),\qquad \tau>0.
\]
This function is globally real analytic and convex, with
$0\le H_\tau(s)-s_+\le\tau\log2$.  Given $\eta>0$, replace
$g_+$ by the raw function
\[
 \widetilde g(z)=-1+\frac{2H_\tau(z-l)}{1-l},
 \qquad \tau=\frac{\eta(1-l)}{2\log2},
\]
and then by
\[
 h(z)=\frac{\widetilde g(z)-\eta/2}{1+\eta/2}.
\]
On $[-1,1]$ this function lies in $[-1,1]$ and satisfies
$|h-g_+|\le\eta$.  The analogous construction applies to $g_-$.
Every resulting query is globally convex and real analytic.
An answer to $h$ with noise bound $\delta$ is an answer to the
hinge with error bound $\delta+\eta$.  Use the interval procedure
with that enlarged bound, taking $\eta<1-\delta$.
To obtain globally convex polynomials as well, fix one of these
smooth convex functions $h$ on $[-1,1]$.  Approximate the continuous
function $\sqrt{h''}$ uniformly there by a real polynomial $q$.
Define the polynomial $P$ by
\[
 P''=q^2,\qquad P(0)=h(0),\qquad P'(0)=h'(0).
\]
It is convex on all of $\mathbb R$, and
$\|P-h\|_{[-1,1]}\le\frac12\|q^2-h''\|_{[-1,1]}$.
This error can be made arbitrarily small.  If it is at most
$\zeta$, replace $P$ by $P/(1+\zeta)$ to restore its range to
$[-1,1]$; the additional uniform error is at most $\zeta$.
Thus either globally convex analytic functions or globally convex
polynomials approximate each hinge with arbitrarily small uniform
error while preserving normalization.  For fixed $k$ the enlarged
noise bound tends to $\delta$ as these approximation errors tend
to zero.  Since $e_{mk}$ is an infimum over algorithms,
\eqref{p06:ni:convex-upper} follows for the polynomial subclass.
Finally let $k\to\infty$ and combine with the lower bound.
\end{proof}

The limit in \eqref{p06:ni:convex-limit} is an infimum over finite-query
algorithms, not an assumption that an infinite transcript can be
read.  Formula~\eqref{p06:ni:convex-upper} supplies a finite algorithm
for every desired excess error above the floor.  The analytic
queries have no common imposed Lipschitz constant, and the polynomial
queries have no common degree bound; either would be an additional
restriction.

\paragraph{A geometric identity behind the obstruction.}
The sharp lower bound is a special case of a useful exact dual
formula.  Let $K$ be a compact convex body, let $x\ne y$ lie in
its interior, and let $z$ be the intersection of the ray
$x+\mathbb R_+(y-x)$ with $\partial K$.  Put
\[
 q_K(x,y)=\frac{\|y-x\|_2}{\|z-x\|_2}.
\]
Then
\begin{equation}\label{p06:ni:convex-dual}
 \sup_{\substack{f\text{ convex on }K\\-1\le f\le1}}
 \{f(y)-f(x)\}=2q_K(x,y).
\end{equation}
Indeed, $y=(1-q)x+qz$ proves the upper bound.  For equality choose
a supporting linear functional $a$ at $z$, put
$b=\langle a,z\rangle$, and use
\[
 f(w)=\max\left\{-1,
  -1+2\frac{\langle a,w-x\rangle}{b-\langle a,x\rangle}
              \right\}.
\]
The denominator is positive because $x$ is interior; this convex
function lies between $-1$ and $1$ on $K$ and attains the desired
difference.  Thus the symmetric discrepancy induced by all bounded
convex measurements is
$2\max\{q_K(x,y),q_K(y,x)\}$.  This identity is included as the
geometric explanation of the noise floor, not as a separately counted
resolution of an external open problem.

\subsection{Degree is a precision resource for nonconvex polynomial information}

Let $\Lambda_D$ consist of real polynomials on $\mathbb R^m$ of
total degree at most the integer $D\ge1$, with
$|P(x)|\le1$ on $C_m$.  The polynomial is allowed to depend on
earlier noisy answers, but its degree and range bounds are uniform.
Write $e_{n,D}(\delta)=e_n(\Lambda_D,\delta)$.

\begin{theorem}[Sharp order of the bounded-degree noise floor]
\label{p06:ni:degree-floor}
For every $m,D\ge1$, every $n\ge0$, and $0<\delta<1$,
\begin{equation}\label{p06:ni:degree-lower}
 e_{n,D}(\delta)\ge
 \frac{\delta}{\sqrt{D^2+\delta^2}}.
\end{equation}
If $0<\delta\le1/72$, set $a=18\delta\le1/4$ and
\[
 k=2+\left\lceil\frac{\log D}{\log(1/a)}\right\rceil.
\]
Then
\begin{equation}\label{p06:ni:degree-upper}
 e_{mk,D}(\delta)\le\frac{18\delta}{D}.
\end{equation}
In particular, uniformly in $m,D$ and $0<\delta\le1/72$,
\[
 \inf_{n\ge0}e_{n,D}(\delta)\asymp\frac\delta D.
\]
The constants in the two inequalities are not claimed optimal.
\end{theorem}

\begin{proof}[Lower bound]
The classical Bernstein inequality for algebraic polynomials states
that, if $p$ has degree at most $D$ and $|p|\le1$ on $[-1,1]$,
then
\[
 |p'(s)|\le\frac D{\sqrt{1-s^2}},\qquad -1<s<1;
\]
see, for example, the treatments in~\cite{p06:ni:Bernstein}.
For $0<t<1$ this implies
\[
 |p(t)-p(-t)|\le\frac{2Dt}{\sqrt{1-t^2}}.
\]
Take $t=\delta/\sqrt{D^2+\delta^2}$ and restrict each selected
multivariate query to the line $se_1$.  Its values at $\pm te_1$
differ by at most $2\delta$.  The common-midpoint transcript
argument from Theorem~\ref{p06:ni:convex-floor} gives error at least $t$.
\end{proof}

\begin{proof}[Upper bound]
Again work coordinate by coordinate.  One linear query gives a
feasible interval $I=[c-r,c+r]\subset[-1,1]$ with $r\le\delta$.
We show that whenever $0<r\le1/2$ one further degree-$D$ query
reduces its radius to at most
\begin{equation}\label{p06:ni:cheb-recurrence}
 a\max\{r,D^{-1}\}.
\end{equation}
Let
\[
 K=\min\{D,(2r)^{-1}\},
\]
and choose the largest positive odd integer $j\le K$.  Since
$K\ge1$, we have $j\ge K/3$.  The function
\[
 P_{j,c}(z)=\sin\left(j\arcsin\frac{z-c}{2}\right)
\]
is a real polynomial of degree $j$: for odd $j$ it is, up to sign,
a Chebyshev polynomial.  Since $c,z\in[-1,1]$, it is bounded by
one on the whole input interval.  For $z\in I$ we have
\[
 \left|j\arcsin\frac{z-c}{2}\right|
 \le j r\le\frac12,
\]
where $\arcsin(r/2)\le r$ for $r\le1/2$.  Therefore
\[
 P'_{j,c}(z)
 =\frac{j\cos(j\arcsin((z-c)/2))}
        {2\sqrt{1-(z-c)^2/4}}
 \ge\frac j3.
\]
The polynomial is strictly increasing on the current feasible
interval.  Intersecting with the preimage of the answer interval
of width $2\delta$ leaves a feasible interval of radius at most
\[
 \frac{3\delta}{j}
 \le\frac{9\delta}{K}
 \le18\delta\max\{r,D^{-1}\},
\]
which proves \eqref{p06:ni:cheb-recurrence}.

Write $r_q$ for the radius after $q$ queries in one coordinate.
From $r_1\le\delta$ and \eqref{p06:ni:cheb-recurrence}, induction gives
\[
 r_q\le\max\{a^{q-1}\delta,a/D\}\qquad(q\ge1).
\]
The induction uses $a\le1$; all radii stay below $1/2$.
With the stated $k$ we have $a^{k-2}D\le1$, and consequently
$a^{k-1}\delta\le a/D$.  Thus $r_k\le18\delta/D$.
Applying the procedure to all $m$ coordinates proves
\eqref{p06:ni:degree-upper}.
\end{proof}

\paragraph{What the degree theorem adds.}
Bernstein's derivative bound is classical, and the entropy bounds
for unrestricted noisy information in~\cite{p06:ni:KNPU} are not new
claims here.  The result is the matching recovery construction
within a fixed polynomial degree: degree $D$ supplies precisely
an order-$D$ improvement in the irreducible precision scale.  Its
upper bound uses nonconvex polynomials.  No corresponding
improvement is possible within the convex class, even if the degree
is unbounded and real analyticity is allowed.

\subsection{Audit and unresolved extensions}

The prior exact-real convexification argument passes a direct audit:
an exact energy answer permits subtracting a known radial convex
term; Gaussian damping permits globally analytic, globally
Lipschitz convex queries.  The new noise result explains why this
does not automatically yield a stable noise-correcting construction.
Range normalization and errors in the compensating answer are
substantive costs, and the sharp floor above prohibits eliminating
them by an arbitrary number of bounded convex queries.

Three limitations are material.  First, the two-query theorem does
not cover a general nonconvex real-analytic first measurement.
Second, the full finite-$n$ minimax curve for bounded convex queries
is not determined here; the exact result is its infinite-budget
infimum and the explicit convergent upper bound.  Third, the
degree theorem identifies the uniform order $\delta/D$, not its
optimal leading constant or the smallest query count achieving it.
These are research targets, not corollaries already proved.

% ===== gaussian =====
\section{Sharp covariance-tail boundaries in every fixed dimension}
\label{p06:bg:sec:matrix}

The preceding manuscripts established logarithmic tail sufficiency in
every dimension, but their sharp counterexamples were scalar. Embedding
a scalar witness in one coordinate does not preserve invisibility to the
additional Gaussian anchors. We resolve that matrix obstruction by
replacing a meromorphic zero product with an entire one. This replacement
permits a contour deformation across the imaginary matrix boundary even
when the matrix is indefinite and its inverse quadratic form vanishes.

Conditional on $(A,\Sigma)\in\mathbb R^d\times\mathbb S_d^+$, let
$X_i$ be independent $N_d(A,\Sigma)$ vectors, and write
\[
 \mathcal O_n=(X_1,\ldots,X_{d-1},S_n),\qquad
 S_n=\sum_{i=1}^nX_i,\quad n\ge d.
\]
Only the separate population laws of these experiments are supplied.
There is no cross-design coupling information. Fix a sufficiently large
$E$, put $L_1(v)=\log(E+v)$, $L_{j+1}(v)=\log L_j(v)$, and set
\[
 W_{m,a}(v)=\frac{\sqrt v}
 {L_1(v)\cdots L_{m-1}(v)L_m(v)^a},\qquad m\ge1,
 \quad a\ge0.
\]

\begin{theorem}[The fixed-dimension tail boundary]
\label{p06:bg:thm:hierarchy}
Fix any dimension $d\ge1$ and any finite logarithmic depth $m\ge1$.
\begin{enumerate}
\item If $0\le a\le1$, every infinite terminal set identifies the
joint mixing law in the class
\begin{equation}\label{p06:bg:eq:class}
 \mathbb E\exp\{cW_{m,a}(\operatorname{tr}\Sigma)\}<\infty
 \quad\text{for some }c>0.
\end{equation}
No restriction on the mean marginal is needed.
\item Suppose $n_1<n_2<\cdots$ and $n_j\ge c_0j^s$ for some
$s>1$. For every $a>1$, these terminal experiments fail to identify
the mixing law in \eqref{p06:bg:eq:class}. The two distinct witnesses
can be absolutely continuous, have the same complete mean and covariance
marginals and the same joint polynomial moments of every degree, and
satisfy \eqref{p06:bg:eq:class} for \emph{every} $c>0$. One can impose
any fixed lower bound $\Sigma\succeq b_0I$ with $b_0>0$.
\end{enumerate}
\end{theorem}

The sufficiency assertion is recalled to state a sharp boundary; the new
assertion is its converse in each fixed $d>1$, together with the enlarged
class of terminal designs. The proof also gives a substantially denser
design theorem at the first logarithmic depth.

\begin{theorem}[Nearly harmonic terminal designs]
\label{p06:bg:thm:dense}
Fix $d\ge1$ and $\beta>1$. If
\[
 n_j\ge c_0j\{\log(e+j)\}^{\beta},\qquad j\ge1,
\]
then for every $a>1$ there are two witnesses with all the properties
in the negative part of Theorem~\ref{p06:bg:thm:hierarchy}, with $m=1$.
Thus the logarithmic covariance-tail threshold remains exactly $a=1$
for each such fixed design sequence, even though it can be much denser
than a geometric sequence.
\end{theorem}

These are population nonidentifiability results for arbitrary mixing
laws. They do not contradict finite Gaussian-mixture identifiability.
The general differential-annihilator approach to infinite-mixture
nonidentifiability is developed by Le, Bariletto, Rinaldo, and
Ho~\cite{p01:mix:Le2026}. The construction here simultaneously cancels infinitely
many matrix observation surfaces and imposes a quantitative covariance
tail; it is not a consequence of a single finite-order annihilator.

\subsection{Entire products and quantitative matrix regularization}

For large positive $R$, abbreviate
\[
 P_m(R)=\prod_{j=1}^m\log_jR,
 \qquad
 D_{m,p}(R)=P_m(R)(\log_mR)^p,
 \qquad p>0.
\]
All uses below have $R$ large; on bounded intervals these functions
may be extended smoothly and positively. Choose $K$ sufficiently large
and define, on the cut plane, the principal holomorphic functions
\begin{equation}\label{p06:bg:eq:Q}
 J_1(w)=\log(K+w),\qquad
 J_{j+1}(w)=\log(K+J_j(w)),\qquad
 Q(w)=\frac{w}{J_m(w)^p}.
\end{equation}
On the upper cut half-plane all nested logarithms have positive
imaginary part, and on the positive axis they are positive; conjugation
handles the lower half-plane. Thus these formulas introduce no
additional cuts or zeros.

\begin{lemma}[Scalar sector bounds]
\label{p06:bg:lem:sector}
For suitable constants $c,C,R_0>0$ the following hold.
\begin{enumerate}
\item $\operatorname{Re}Q(w)\ge0$ on the right half-plane, and
\[
 \operatorname{Re}Q(w)\ge
 c\frac{|w|}{D_{m,p}(|w|)}
 \quad (\operatorname{Re}w\ge0,\ |w|\ge R_0).
\]
\item The latter estimate remains valid when $|w|\ge R_0$ and
\[
 |\arg w|\le\frac\pi2+\frac{c}{P_m(|w|)}.
\]
\item There is a smooth positive function $\eta(r)$, $r>0$,
with bounded derivative, bounded above by an arbitrarily small fixed
constant, and
\begin{equation}\label{p06:bg:eq:eta}
 \eta(r)\asymp\frac{r}{P_m(1/r)}\quad(r\downarrow0),
\end{equation}
such that, for every real $\lambda\ne0$ and $0<r\le|\lambda|$,
\[
 \operatorname{Re}Q\bigl(( -t\eta(r)+i\lambda)^{-1}\bigr)\ge0,
 \qquad 0\le t\le1.
\]
For $|\lambda|$ small the last expression is at least
$c/(|\lambda|D_{m,p}(1/|\lambda|))$.
\end{enumerate}
\end{lemma}
\begin{proof}
If $0\le\arg w\le\pi/2$, each nested logarithm has real part
at least $\log K$. Successive argument estimates give
$0\le p\arg J_m(w)\le\arg w$ once $K$ is large enough.
This proves the first nonnegativity assertion. Uniformly for bounded
real $\theta$, repeated logarithmic expansion gives
\[
 |J_m(Re^{i\theta})|\sim\log_mR,
 \qquad
 \arg J_m(Re^{i\theta})
  =\frac{\theta}{P_m(R)}(1+o(1)).
\]
The second relation and its first angular derivative are uniform.
They follow inductively because the imaginary part after each
logarithm is the previous imaginary part divided by its real part,
up to a relative error tending to zero. Fixed additive constants only
affect the lower-order terms. Near $\theta=\pi/2$, the argument
of $Q$ is consequently at most $\pi/2-c/P_m(R)$, even if
$\theta$ exceeds $\pi/2$ by a sufficiently small multiple of
$1/P_m(R)$. Away from those endpoints its real part is larger.
This proves the two large-modulus assertions.

For the last part, take a sufficiently small multiple of
$r/P_m(1/r)$ near zero and extend it smoothly to a small positive
constant for large $r$, with bounded derivative. For small $|\lambda|$,
the angular overshoot of $(-t\eta(r)+i\lambda)^{-1}$ is at most
$\eta(r)/|\lambda|\le c/P_m(1/|\lambda|)$; here the function
$r/P_m(1/r)$ is increasing near zero. The sector estimate applies.
On compact intervals of nonzero $\lambda$, strict positivity on the
imaginary axis and continuity allow a common small upper bound on
$\eta$. Finally, at zero the Taylor expansion is
$Q(w)=b_1w-b_2w^2+O(w^3)$ with $b_1,b_2>0$.
For large $|\lambda|$ it yields
\[
 \operatorname{Re}Q((-t\eta+i\lambda)^{-1})
 =\frac{b_2-b_1t\eta}{\lambda^2}+O(|\lambda|^{-3}),
\]
uniformly for bounded $\eta$. Reducing its upper bound once more
completes the proof. The expansion's positivity follows directly from
$J_m(0)>0$ and $J_m'(0)>0$.
\end{proof}

Let $a_j=2n_j$. Since $\sum_j a_j^{-1}<\infty$ in both
theorems, the canonical product
\begin{equation}\label{p06:bg:eq:entire}
 E_0(z)=\prod_{j\ge1}(1-z/a_j)
\end{equation}
is a nonzero entire function, real on the real axis, with the required
positive zeros. The absence of negative poles is essential.

\begin{lemma}[Growth of the design product]
\label{p06:bg:lem:product}
\begin{enumerate}
\item If $a_j\ge cj^s$, $s>1$, then
$\log|E_0(z)|\le C(1+|z|^{1/s})$ throughout the plane.
\item If $a_j\ge cj\log^\beta(e+j)$, $\beta>1$, put
\[
 J(R)=\frac12\sum_j\log(1+R^2/a_j^2),\qquad
 S(R)=\sum_j\frac{a_j}{a_j^2+R^2}.
\]
For large $R$,
\[
 J(R)\le\frac{CR}{\log^\beta R},\qquad
 S(R)\le\frac C{\log^{\beta-1}R}.
\]
Moreover, for every complex $z$,
\begin{equation}\label{p06:bg:eq:product-sector}
 \log|E_0(z)|\le J(|z|)+(-\operatorname{Re}z)_+S(|z|).
\end{equation}
Consequently, uniformly when $|z|\le C_0R$ and
$(-\operatorname{Re}z)_+\le C_0R/\log R$,
\begin{equation}\label{p06:bg:eq:product-left}
 \log|E_0(z)|\le\frac{C R}{\log^\beta R}.
\end{equation}
The same bounds, with altered constants, apply to each fixed derivative
of $E_0$ on the indicated sets.
\end{enumerate}
\end{lemma}
\begin{proof}
In the first case,
$\sum_j\log(1+|z|/(cj^s))\le C(1+|z|^{1/s})$, by splitting
at $j\asymp |z|^{1/s}$ or comparing with the corresponding
integral. In the second case, the counting function satisfies
$N(t)=\#\{j:a_j\le t\}\le Ct/\log^\beta(e+t)$.
Integration by parts gives
\[
 J(R)=\int_0^\infty\frac{R^2N(t)}{t(t^2+R^2)}\,dt.
\]
Split at $\sqrt R$ and $R$. The first portion is $O(\sqrt R)$;
on the middle portion use $\log t\asymp\log R$; on the last
portion use $R^2/(t^2+R^2)\le R^2/t^2$.
This proves the stated bound for $J$.
For $S$, the terms $a_j\le R$ sum to at most $N(R)/R$.
Partial summation bounds the remaining terms by
$C\int_R^\infty dt/(t\log^\beta t)$, proving the assertion.

For $x=(-\operatorname{Re}z)_+$ and $r=|z|$,
\[
 |1-z/a|^2\le(1+r^2/a^2)
       \exp\{2xa/(a^2+r^2)\}.
\]
Summing logarithms proves \eqref{p06:bg:eq:product-sector}.
If $r\ge\sqrt R$, its right side is bounded by
$CR/\log^\beta R$ under the extra hypotheses. If $r<\sqrt R$,
use instead $\log|E_0(z)|\le r\sum a_j^{-1}=O(\sqrt R)$.
Fixed derivatives follow from Cauchy's formula on circles of radius
one; enlarging a set by that radius preserves all the displayed bounds
after a harmless change in constants.
\end{proof}

On $\mathbb S_d$ use orthonormal coordinates for
$\langle V,W\rangle=\operatorname{tr}(VW)$, and set
$D=d(d+1)/2$. Choose a nonzero real even
$h\in C_c^\infty(\mathbb R^d\setminus\{0\})$ and define, for
complex symmetric $T$ with $\operatorname{Re}T\succ0$,
\begin{equation}\label{p06:bg:eq:H}
 \mathcal H(u,T)=h(u)E_0(u^TT^{-1}u)
       \exp\{-\operatorname{tr}T^{1/2}
                    -\operatorname{tr}Q(T^{-1})\}.
\end{equation}
Matrix functions here are principal holomorphic functions. In the
polynomially sparse case $m\ge1$ and $p>0$ are arbitrary. In the
nearly harmonic case take $m=1$ and $0<p<\beta-1$.

\begin{lemma}[A quantitatively decaying invisible matrix density]
\label{p06:bg:lem:realization}
Under either set of hypotheses just specified,
$\mathcal H$ is the Fourier--Laplace transform of a nonzero real
Schwartz function $g$ on $\mathbb R^d\times\mathbb S_d$ supported
where $\Sigma\succeq0$. Both coordinate marginals and all joint
polynomial moments of $g$ vanish. For every integer $M\ge0$ there
are $c_M,C_M>0$ such that
\begin{equation}\label{p06:bg:eq:decay}
 |g(a,\Sigma)|\le C_M(1+|a|)^{-M}
    \exp\{-c_M W_{m,1+p/2}(\operatorname{tr}\Sigma)\},
 \qquad \Sigma\succeq0.
\end{equation}
\end{lemma}
\begin{proof}
We prove the tube estimates, the support assertion, and the quantitative
contour estimate separately.

\emph{Tube estimates.}
If $T=X+iV$ with $X\succ0$, then
$R=T^{-1}$ is accretive because
$\operatorname{Re}R=T^{-*}XT^{-1}\succ0$.
For any accretive $d\times d$ matrix,
\begin{equation}\label{p06:bg:eq:nonnormal}
 \|R\|\le C_d\sum_{j=1}^d|\lambda_j(R)|.
\end{equation}
Indeed, put $R$ in unitary Schur form. Positivity of its Hermitian
part gives $|R_{ij}|\le2\sqrt{\operatorname{Re}\lambda_i
\operatorname{Re}\lambda_j}$ for each strictly upper triangular
entry. The Frobenius norm is therefore bounded by a constant times
$\sum|\lambda_i|$. This also handles repeated eigenvalues and
non-normal matrices.
Lemma~\ref{p06:bg:lem:sector}, the trace identity for holomorphic
matrix functions, and \eqref{p06:bg:eq:nonnormal} imply
\begin{equation}\label{p06:bg:eq:trace-Q}
 \operatorname{Re}\operatorname{tr}Q(R)
 \ge c\frac{\|R\|}{D_{m,p}(\|R\|)}-C
\end{equation}
when $\|R\|$ is large. To see this directly, at least one
eigenvalue has modulus at least $c_d\|R\|$, all other real
contributions are nonnegative, and fixed rescaling preserves the
displayed asymptotic order.

In the first design class, the product growth is $O(\|R\|^{1/s})$;
in the second, $u^TRu$ lies in the right half-plane and
Lemma~\ref{p06:bg:lem:product} bounds its logarithmic growth by
$C\|R\|/\log^\beta(e+\|R\|)$. These quantities are
$o(\|R\|/D_{m,p}(\|R\|))$ in their respective parameter
ranges. Thus the negative exponent in \eqref{p06:bg:eq:trace-Q}
absorbs the entire product, uniformly for $u$ in the support of $h$.

There is also a uniform bound at large matrix frequencies:
\begin{equation}\label{p06:bg:eq:uniform}
 |\mathcal H(u,X+iV)|\le C\exp\{-c\sqrt{\|V\|}\},
 \qquad X\succ0.
\end{equation}
For completeness, write $T^{1/2}=B+iC$ with real symmetric
$B,C$. The principal square root of an accretive matrix is accretive
(for example, use its convergent resolvent integral), so $B\succ0$.
The equations $B^2-C^2=X\succ0$ and $BC+CB=V$ give
$\|C\|\le\|B\|$ and
$\|V\|\le2(\operatorname{tr}B)^2$.
This proves \eqref{p06:bg:eq:uniform} after the preceding absorption.

On the boundary $T=iV$ with $V$ real symmetric and invertible,
functional calculus is diagonal in an orthonormal real eigenbasis.
Put
\begin{equation}\label{p06:bg:eq:r}
 r(V)=\{\operatorname{tr}(V^{-2})\}^{-1/2}.
\end{equation}
It is smooth on the invertible locus,
$\min|\lambda_j(V)|/\sqrt d\le r(V)\le\min|\lambda_j(V)|$,
and $|\nabla r|\le1$. The last assertion follows from
$\nabla r=r^3V^{-3}$ and
$\operatorname{tr}V^{-6}\le(\operatorname{tr}V^{-2})^3$.
The boundary value $H(u,V)=\mathcal H(u,iV)$ decays faster than
every power at both $\|V\|=\infty$ and $r(V)=0$.
The same is true of every derivative. Here are the details needed
when a Cauchy neighborhood leaves the positive tube. Around $iV$ take
complex symmetric matrix perturbations of operator norm at most
$cr/P_m(1/r)$, with $c$ sufficiently small. Since $iV$ is normal,
every perturbed eigenvalue is within this distance of its real-normal
reference spectrum $i\operatorname{spec}V$. Its inverse therefore
lies in the enlarged scalar sector of Lemma~\ref{p06:bg:lem:sector},
and the inverse matrix norm is $O(1/r)$ by the Neumann-series bound.
Grouping overlapping spectral disks preserves the eigenvalue count
of each cluster along the perturbation. The cluster containing an
eigenvalue of smallest modulus has diameter at most $2d$ times the
perturbation radius. At least one inverse eigenvalue therefore has
modulus comparable to $1/r$; all the real $Q$ contributions remain
nonnegative. Thus the trace
regularization still gives $c/(rD_{m,p}(1/r))$ damping.
For the nearly harmonic product one also needs its real argument
control. For any invertible $T$,
\[
 \operatorname{Re}(T^{-1})=T^{-*}(\operatorname{Re}T)T^{-1}.
\]
In these neighborhoods $\|\operatorname{Re}T\|=O(r/P_m(1/r))$,
so the negative part of the quadratic form $u^TT^{-1}u$ is at most
$C/(rP_m(1/r))$. When $m=1$ this is exactly the range of
\eqref{p06:bg:eq:product-left}; the polynomially sparse product needs
only the inverse norm bound. Cauchy's formula in these neighborhoods
therefore bounds each fixed matrix derivative by the same exponential
damping times powers of $P_m(1/r)/r$ and of $1+\|V\|$.
At infinity, neighborhoods can be reduced if needed; the principal
square-root trace retains its $\sqrt{\|V\|}$ damping, up to a
constant. For $u$ derivatives apply the scalar Cauchy formula for
$E_0$ on circles of radius one and the chain rule; their extra factors
are polynomial in $\|V^{-1}\|$. The exponential bounds absorb all
these factors. Consequently $H$ extends by zero across $\det V=0$
to a Schwartz function, compactly supported in $u$.

\emph{Fourier inversion and cone support.}
Set
\[
 g_u(\Sigma)=(2\pi)^{-D}\int_{\mathbb S_d}
 e^{i\operatorname{tr}(V\Sigma)}H(u,V)\,dV,
 \qquad
 g(a,\Sigma)=(2\pi)^{-d}\int e^{-iu^Ta}g_u(\Sigma)\,du.
\]
Both are Schwartz. For $X\succ0$ let $g_{u,X}$ be the analogous
inverse Fourier integral of $\mathcal H(u,X+iV)$. Cauchy--Riemann
equations and integration by parts along a compact path in the
positive tube give
\[
 g_{u,X}(\Sigma)=e^{-\operatorname{tr}(X\Sigma)}g_u(\Sigma).
\]
One first starts at $X=\epsilon I$ and then sends $\epsilon$
to zero; \eqref{p06:bg:eq:uniform} justifies dominated convergence.
The same bound shows that $g_{u,X}$ is uniformly bounded for all
$X\succ0$. If $\Sigma$ has a negative eigenvalue, choose
$X_0\succ0$ with $\operatorname{tr}(X_0\Sigma)<0$ and send
$X=tX_0$ to infinity. Thus $g_u(\Sigma)=0$ outside the positive
semidefinite cone. Fourier inversion now proves the asserted
Fourier--Laplace identity.

\emph{The variable matrix contour.}
For invertible real symmetric $V$ deform the contour from $iV$ to
\begin{equation}\label{p06:bg:eq:contour}
 T(V)=-\eta(r(V))I+iV.
\end{equation}
All matrices along the homotopy
$-t\eta(r(V))I+iV$, $0\le t\le1$, are normal and have
nonreal eigenvalues; their spectra and the spectra of their inverses
avoid the negative real axis. Thus \eqref{p06:bg:eq:H} is holomorphic
in a neighborhood of every point of the homotopy. The entire function
$E_0$ introduces no further pole hypersurfaces.

At least one eigenvalue of $V$ has modulus comparable to $r(V)$.
Lemma~\ref{p06:bg:lem:sector} therefore gives, uniformly along this
homotopy when $r=r(V)$ is small,
\[
 \operatorname{Re}\operatorname{tr}Q(T^{-1})
    \ge\frac{c}{rD_{m,p}(1/r)}.
\]
All other summands are nonnegative. Also
$\operatorname{Re}\operatorname{tr}T^{1/2}
\ge c\sqrt{\|V\|}$ outside a fixed bounded set: the real
shift is bounded, so each large imaginary eigenvalue has a principal
square root with real part comparable to its square-root modulus.
In the first design class the product growth is $O(r^{-1/s})$
and is absorbed by the first displayed damping term.

In the nearly harmonic case put $R=1/r$. Since $m=1$,
$\eta(r)\le Cr/\log(1/r)$. In a real eigenbasis of $V$,
for $z=u^TT^{-1}u$ one has
\[
 |z|\le C R,\qquad
 (-\operatorname{Re}z)_+
 =t\eta(r)\sum_j\frac{u_j^2}
 {t^2\eta(r)^2+\lambda_j(V)^2}
 \le\frac{CR}{\log R}.
\]
Equation~\eqref{p06:bg:eq:product-left} bounds its logarithmic product
growth by $CR/\log^\beta R$. This is absorbed because
$p+1<\beta$. Every fixed derivative in $u$ obeys the same
estimate, up to polynomial factors in $R$, by Cauchy's formula for
the scalar entire function and the chain rule.
We have therefore proved, for each fixed multi-index $\alpha$,
\begin{equation}\label{p06:bg:eq:shift-bound}
 |\partial_u^\alpha\mathcal H(u,T(V))|
 \le C_\alpha\exp\left\{-c_\alpha\sqrt{\|V\|}
              -\frac{c_\alpha}{rD_{m,p}(1/r)}\right\}
\end{equation}
for small $r$; on $r$ bounded away from zero only the first
exponential is needed.

Here the contour substitution has bounded Jacobian. In orthonormal
matrix coordinates it is a rank-one change of $iI_D$, since its
real displacement is always a scalar multiple of the fixed matrix
$I$. Its determinant has modulus at most
$1+\sqrt d\,\|\nabla\eta(r(V))\|$, which is bounded by
$1+\sqrt d\sup|\eta'|$.
The deformation can be justified using only one complex variable,
which avoids any assumption about the geometry of the singular
matrix hypersurface. Write
\[
 V=W+\frac{t}{\sqrt d}I,\qquad \operatorname{tr}W=0,
\]
using orthonormal coordinates, and first fix $W$. The singular real
$t$ values are the finitely many numbers
$t_j=-\sqrt d\,\lambda_j(W)$. On each interval between successive
distinct $t_j$, deform the real $t$ contour to
\[
 z=t+i\sqrt d\,\eta(r(W+tI/\sqrt d)).
\]
The corresponding matrix is $iW+izI/\sqrt d$, exactly
\eqref{p06:bg:eq:contour}. Every eigenvalue has imaginary part of a
fixed nonzero sign in the interior of this vertical strip. Its
spectrum and that of its inverse therefore avoid all branch cuts;
the integrand is holomorphic throughout the region swept out by the
contour. Near a finite endpoint the deformation height tends to
zero, and the short joining segment has a vanishing integral by
\eqref{p06:bg:eq:shift-bound}, because $r$ is comparable to the
endpoint distance up to constants depending only on $d$.
The unbounded intervals have vanishing joining integrals at infinity
by the square-root damping. Cauchy's theorem proves the identity
for each fixed $W$, with line element
$(1+i\sqrt d\,\partial_t\eta)\,dt$.
Its modulus is bounded uniformly, as already noted. The global
bounds above, and
$|e^{\operatorname{tr}(T\Sigma)}|\le1$ for
$\Sigma\succeq0$ along the homotopy, make both original and deformed
integrals absolutely integrable in $(t,W)$. Fubini's theorem therefore
proves the full matrix contour identity. The same reasoning applies
after any fixed number of $u$ derivatives. Repeated eigenvalues of
$W$ merely repeat an endpoint and cause no change in this argument.
No contour passes through a singular matrix.

Write $v=\operatorname{tr}\Sigma$. The modulus of the deformed
inverse integrand, for small $r$, is bounded by
\[
 C\exp\left\{-c\sqrt{\|V\|}
 -c\left(\frac{vr}{P_m(1/r)}
       +\frac1{rP_m(1/r)(\log_m(1/r))^p}\right)\right\}.
\]
For $v$ large, the expression in parentheses is at least
$cW_{m,1+p/2}(v)$ uniformly in small $r$.
Indeed, on $\frac14\log v\le\log(1/r)\le\log v$ all
the fixed-depth logarithms are comparable to the corresponding
logarithms of $v$, and the arithmetic--geometric mean inequality
gives the assertion. Below that interval the first term is at least
$v^{3/4}$ divided by a fixed product of logarithms; above it the
second is at least $v$ divided by such a product. Both are larger
than the required bound. On $r$ bounded away from zero,
$\eta(r)$ has a positive lower bound and the estimate is
$O(e^{-cv})$ instead. Reserve half of the exponents in
\eqref{p06:bg:eq:shift-bound} for integrability in $V$. We conclude
\[
 |\partial_u^\alpha g_u(\Sigma)|
 \le C_\alpha e^{-c_\alpha W_{m,1+p/2}(v)}.
\]
Integration by parts in the compactly supported $u$ integral
now proves \eqref{p06:bg:eq:decay}.

Finally $g$ is real because
$\mathcal H(-u,\overline T)=\overline{\mathcal H(u,T)}$.
It is nonzero since, at a point where $h(u)\ne0$, varying $T=tI$
avoids the discrete zeros of $E_0$.
The function $g_u$ vanishes near $u=0$, so
$\int a^\gamma g(a,\Sigma)\,da=0$ for every multi-index
$\gamma$. This gives the covariance marginal and all joint
polynomial moments. For the other marginal send $T=\epsilon I$
to zero in \eqref{p06:bg:eq:H}. Its $Q$ damping dominates the product
growth and forces a zero limit. Dominated convergence and Fourier
uniqueness in $a$ show $\int g(a,\Sigma)\,d\Sigma=0$.
\end{proof}

\subsection{Identification, counterexamples, and the exact scope}

\begin{proof}[Proof of the two main theorems]
We first recall the sufficient argument, including the classical
moment input. If $V\ge0$ satisfies
$\mathbb E e^{cW_{m,a}(V)}<\infty$, then
\[
 (\mathbb E V^k)^{1/(2k)}
 \le C k L_1(k)\cdots L_{m-1}(k)L_m(k)^a.
\]
To verify this estimate, put $Y=\sqrt V$,
$D(x)=L_1(x)\cdots L_{m-1}(x)L_m(x)^a$, and $b_k=kD(k)$.
For $w\ge1$, iterated logarithm inequalities give
$D(b_kw)\le C D(k)(1+\log w)^{m-1+a}$.
Substitution $Y=b_kw$ in
$\sup_y y^{2k}e^{-cy/D(y)}$ proves the displayed moment bound.
The Stieltjes--Carleman series consequently diverges when $a\le1$.
This moment-determinacy mechanism, and much more general
quasi-analytic weights, are classical; see de Jeu~\cite{p05:fg:deJeu}.

The Fourier arguments of the $d$ independent blocks with lengths
$1,\ldots,1,n-d+1$ are exactly
\begin{equation}\label{p06:bg:eq:surface}
 (u,T),\qquad
 T=C+\frac{uu^T}{2n},\quad C\succeq0,
 \quad\operatorname{rank}C\le d-1.
\end{equation}
This follows by projecting the matrix of weighted block frequencies
onto the vector of square roots of the lengths and its orthogonal
complement. Let $\rho$ be the difference of two mixing laws with
the same observations. Fix $u\ne0$ and a rank-$d-1$ matrix $C$
with $u$ not orthogonal to its kernel. The pushforward of
$e^{iu^Ta-\operatorname{tr}(C\Sigma)}\rho$ under
$(a,\Sigma)\mapsto u^T\Sigma u$ is a finite complex measure
whose Laplace transform vanishes at $1/(2n)$ for every selected $n$.
The sum of the candidate laws dominates it and supplies the common
Stieltjes--Carleman moment bound. Since the zeros accumulate at zero,
successive Taylor expansions there force all moments to vanish.
Moment determinacy applied to the positive dominating measure plus
and minus the real and imaginary parts makes the complex measure
zero. Every $T\succ0$ can be written as $C+zuu^T$ of the above
type: take $z=(u^TT^{-1}u)^{-1}$. Thus the full mixing transform
vanishes on $u\in\mathbb R^d$, $T\succ0$. Matrix
Laplace--Fourier uniqueness proves identification.

For the negative assertion of Theorem~\ref{p06:bg:thm:hierarchy},
choose $0<p<2(a-1)$ and use Lemma~\ref{p06:bg:lem:realization}.
For Theorem~\ref{p06:bg:thm:dense} choose
$0<p<\min\{\beta-1,2(a-1)\}$ and set $m=1$.
At each positive definite observation argument in
\eqref{p06:bg:eq:surface}, the matrix determinant lemma gives
$u^TT^{-1}u=2n$. Hence $\mathcal H(u,T)=0$ by
\eqref{p06:bg:eq:entire}. Singular observation arguments are limits
of positive definite ones on the same observation surface: perturb
the square $d\times d$ weighted frequency matrix to an invertible
one. Continuity of the finite-measure transform therefore includes
them as well. The signed observation law of $g$ is zero at every
chosen terminal length.

Write $g=g_+-g_-$. Its nonzero positive and negative parts have
the same positive total mass $b$; use $g_+/b$ and $g_-/b$ as the
two mixing densities. Their coordinate marginals and all joint
moments agree by the preceding lemma. Since $1+p/2<a$, the ratio
$W_{m,1+p/2}(v)/W_{m,a}(v)$ tends to infinity. Equation
\eqref{p06:bg:eq:decay}, integrated first in $a$ and then on the
finite-dimensional positive cone, proves the required exponential
integrability for every $c>0$. Translating both covariances by
$b_0I$ multiplies the transform by
$e^{-b_0\operatorname{tr}T}$; all zeros, equalities, and tail
properties persist. This imposes the positive lower bound and
completes both proofs.
\end{proof}

\begin{remark}[What has and has not been closed]
The fixed-$d>1$ logarithmic converse and its finite iterated-log
extensions were explicit gaps in the preceding manuscripts and are
closed here. The nearly harmonic design theorem is an additional
extension, not a restatement of the scalar lacunary example.
The exact joint classification for \emph{every} convergent harmonic
subseries and \emph{every} covariance-tail envelope is not asserted.
In particular, the present estimates do not settle highly irregular
design sequences or supply sharp finite-sample stability at the
quasi-analytic boundary. No named external conjecture is claimed as
resolved by these results, and priority requires a broader external
literature review.
\end{remark}

% ===== spectral =====
\section{Blind spectral certification from nonlinear transition thresholds}
\label{p06:rs:sec}

The earlier threshold theorem assumed an a priori bound on the number
of atoms in both competing spectra.  Here we remove that assumption
from the competitor.  One additional, appropriately designed scalar
measurement certifies an atomic spectrum against every positive
measure, including continuous spectra.  The corresponding collision
exponent changes from $1/(2m-1)$ to $1/(2m)$.  These are inverse
problems for the nonlinear threshold observable below.  They are not
asserted to be previously published open problems.  In particular,
positivity certificates and the $2m+1$ sample principle for \emph{linear}
Chebyshev measurements are established results of
Eftekhari--Tanner--Thompson--Toader--Tyagi~\cite{p06:rs:Eftekhari}; moment
comparison by interpolation is developed in Wu--Yang~\cite{p05:fz:WuYang}.
The new step here is the transfer from unobserved contacts of nonlinear
minima to those certificates, with uniform constants through collisions.

Fix $0<a<b<1$.  For a nonzero finite positive Borel measure on $[a,b]$,
put
\begin{equation}\label{p06:rs:eq:model}
 S_\mu(x)=\int x^r\,d\mu(r),\qquad
 T_\mu(C)=\min_{x>0}\frac{1+Cx}{S_\mu(x)}.
\end{equation}
Write $\mathcal P[a,b]$ for all probability measures and
$\mathcal P_m[a,b]$ for those with at most $m$ atoms.  A list of distinct
settings is called one-sided here if
\begin{equation}\label{p06:rs:eq:oneside}
                       0<C_i<\frac a{1-a}\quad\hbox{for every }i.
\end{equation}
This is a design condition, independent of the unknown measure.
It ensures that every contact in \eqref{p06:rs:eq:model} is strictly
larger than $1$.  Settings all larger than $b/(1-b)$ give the symmetric
version with contacts below $1$.

\begin{theorem}[Exact normalized count and blind certification]
\label{p06:rs:thm:counts}
The following assertions hold.
\begin{enumerate}
\item Any $2m-1$ one-sided settings identify every member of
$\mathcal P_m[a,b]$ among members of that same class.  Fewer fixed
settings cannot identify every exactly $m$-atomic probability measure
with distinct interior atoms and positive masses.
\item Any $2m$ one-sided settings identify each member of
$\mathcal P_m[a,b]$ among \emph{all} of $\mathcal P[a,b]$.
This count is necessary: for every exactly $m$-atomic probability
measure with distinct interior atoms and every list of at most $2m-1$
fixed positive settings, there is a different, positive
$(m+1)$-atomic probability measure giving exactly the same data.
\item Without a mass normalization, any $2m+1$ distinct positive
settings identify an at most $m$-atomic positive measure among all
finite positive measures on $[a,b]$.  At every spectrum with exactly
$m$ distinct interior atoms, $2m$ settings admit positive
$(m+1)$-atomic alternatives.
\end{enumerate}
The first count is optimal for a known sparsity bound; the second and
third are optimal for certifying that bound from the data themselves.
\end{theorem}

\begin{theorem}[Sharp blind stability and the cost of certification]
\label{p06:rs:thm:noise}
For a fixed one-sided list $\mathcal C$ of $2m$ settings there is a
constant $K$ such that
\begin{equation}\label{p06:rs:eq:blindholder}
 W_1(\mu,\nu)\le
 K\left(\max_{C\in\mathcal C}|T_\mu(C)-T_\nu(C)|\right)^{1/(2m)}
 \quad(\mu\in\mathcal P_m[a,b],\ \nu\in\mathcal P[a,b]).
\end{equation}
No atom separation or lower mass bound is needed.  The exponent is
optimal, even if the entire threshold curve on a fixed compact
interval of settings is observed.  In contrast, for two members of
$\mathcal P_m[a,b]$, already $2m-1$ one-sided settings give the sharp
exponent $1/(2m-1)$.

Equivalently, the largest Wasserstein diameter between an $m$-atomic
truth and an arbitrary positive probability measure consistent with
threshold errors of size $\delta$ is
$\asymp\delta^{1/(2m)}$.  This is a \emph{blind consistency radius};
it is not the minimax estimation rate when the estimator is told that
the truth has at most $m$ atoms.
\end{theorem}

\subsection{Contacts, separated near-zeros, and positivity}

We include the needed geometric facts so that the extension to
non-atomic competitors is explicit.  Put $H=S-xS'$.  Differentiation
under the integral gives
\begin{equation}\label{p06:rs:eq:contact}
 C(x)=\frac{S'(x)}{H(x)},\qquad
 C'(x)=\frac{S''(x)S(x)}{H(x)^2}<0,\qquad
 T(C(x))=\frac1{H(x)}.
\end{equation}
Moreover
\[
 \frac{a}{1-a}\le C(x)x\le\frac{b}{1-b}.
\]
Thus $C(x)$ is a decreasing bijection and
$L_C(x)=(1+Cx)/T(C)$ is the unique supporting line at its contact.
These assertions hold for every finite positive measure in the stated
class, not just atomic ones.

\begin{lemma}[Geometric conversion of threshold data]
\label{p06:rs:lem:contacts}
If two positive profiles have equal thresholds at $n$ distinct
settings, their difference has a zero in each of $n$ disjoint
intervals of contact points.  Under \eqref{p06:rs:eq:oneside}, these
intervals lie strictly above $1$.

For probability measures on $[a,b]$ and a fixed list of settings,
there are constants $\delta_0,c,K>0$ and a compact interval $J$ such
that threshold discrepancy at most $\delta\le\delta_0$ gives points
$u_1,\ldots,u_n\in J$, pairwise separated by at least $c$, with
\begin{equation}\label{p06:rs:eq:nearzero}
 \left|\int e^{u_i r}\,d(\mu-\nu)(r)\right|\le K\delta.
\end{equation}
For one-sided settings they are also uniformly separated from $0$.
\end{lemma}
\begin{proof}
Order the settings increasingly.  The slopes of their supporting
lines increase strictly, because
\[
 \frac{d}{dC}\frac C{T(C)}
 =\frac{T(C)-CT'(C)}{T(C)^2}>0,
 \qquad T(C)-CT'(C)=\frac1{S(x(C))}.
\]
The intersection of two consecutive supporting lines lies strictly
between their two contacts.  For common lines the contacts for both
profiles therefore belong to the same pairwise disjoint intervals.
At the first profile's contact its difference from the second is
nonnegative; at the second profile's contact it is nonpositive.
The intermediate value theorem proves the exact assertion.

All contacts for probability measures lie in the fixed interval
\[
 X=\left[\frac{a}{(1-a)\max C_i},
          \frac{b}{(1-b)\min C_i}\right].
\]
On $X$, $S,H,-S''$ are bounded above and away from zero uniformly.
Equation \eqref{p06:rs:eq:contact} therefore separates consecutive
contacts uniformly.  If two contacts are $v>w$, the tangent at $w$
exceeds $S(v)$ by at least $\kappa(v-w)^2/2$, where
$\kappa=\inf_X(-S'')>0$.  Division by the bounded difference of
the tangent slopes gives a uniform positive distance between their
intersection and either contact.

Threshold errors of size $\delta$ change the supporting-line
coefficients by $O(\delta)$; their consecutive slope differences are
bounded away from zero, so the intersections change by $O(\delta)$
as well.  For small $\delta$, both sets of contacts remain within
uniformly separated versions of the same intervals.  At the first
contact the profile difference is at least $-K\delta$, and at the
second it is at most $K\delta$.  Some intermediate point consequently
has difference of absolute value at most $K\delta$.  Taking its
logarithm proves \eqref{p06:rs:eq:nearzero}.  The bound on $X$ under
\eqref{p06:rs:eq:oneside} also separates these logarithms from zero.
\end{proof}

The following uniform certificate is the step that allows an
arbitrary positive competitor.  A finite-dimensional differential
equation for the competitor's transform would be unavailable.

\begin{lemma}[Uniform one-sided exponential interpolation]
\label{p06:rs:lem:certificate}
Fix compact intervals $I,J$, an integer $m\ge1$, and $c>0$.
Let $v_1<\cdots<v_{2m+1}$ belong to $J$ with successive gaps at least
$c$.  Suppose $\mu$ is positive with at most $m$ atoms in $I$, and
$\nu$ is any finite positive measure on $I$.  Then, for every
$\varphi\in C^{2m}(I)$,
\begin{equation}\label{p06:rs:eq:certificate}
 \left|\int\varphi\,d(\nu-\mu)\right|
 \le K\|\varphi\|_{C^{2m}(I)}
       \max_{1\le i\le2m+1}
       \left|\int e^{v_i r}\,d(\nu-\mu)(r)\right|.
\end{equation}
Here $K$ is uniform over collisions of atoms and over their masses.
\end{lemma}
\begin{proof}
Pad the atom list to $m$ nodes $t_1,\ldots,t_m\in I$ if necessary;
repetitions mean higher multiplicity in what follows.  Put $n=2m$
and $P(r)=\prod_{j=1}^m(r-t_j)^2$.  Let $R\varphi$ be the Hermite
interpolant of $\varphi$ at the node list
$t_1,t_1,\ldots,t_m,t_m$ in
$E=\operatorname{span}\{e^{v_1r},\ldots,e^{v_nr}\}$.

For completeness, this interpolation operator has uniformly bounded
coefficients even when the nodes coalesce.  Write the interpolation
equations in successive divided differences at the ordered node list.
Their matrix and right side extend continuously when nodes coincide.
The limiting matrix is always nonsingular: a nonzero member of $E$
cannot have $n$ zeros counted with multiplicity.  This zero count
follows by repeatedly dividing by one exponential, differentiating,
and applying Rolle's theorem with multiplicities.  Compactness of
the node and exponent tuples now bounds the inverse matrices.
Divided differences of $\varphi$ are bounded by its derivatives.
Thus the coefficients of $R\varphi$, and its $C^n$ norm, are at most
$K\|\varphi\|_{C^n}$.

Define $q=e^{v_{n+1}r}-R(e^{v_{n+1}r})$.  It vanishes at the
specified nodes with the prescribed even multiplicities.  The
$(n+1)$-dimensional exponential system has at most $n$ zeros, counted
with multiplicity; hence $q/P$ never vanishes on $I$, including at
coincident nodes where the quotient is interpreted by continuity.
Change its sign if necessary.  Uniformly in all parameters,
\begin{equation}\label{p06:rs:eq:qpositive}
                         q(r)\ge c_1P(r)\quad(r\in I),
\end{equation}
and the coefficients of $q$ in the exponential basis are bounded.
To justify uniformity precisely, the quotient is the ordinary
divided difference of $q$ at its $n$ zeros and the additional point
$r$.  It extends continuously to every collision.  At a limiting
tuple it is nonzero by the same multiplicity zero count.  Its absolute
value therefore has a positive minimum on the compact parameter set.

The ordinary Hermite remainder estimate applied to
$\varphi-R\varphi$ yields
\[
 |\varphi(r)-R\varphi(r)|\le
        K\|\varphi\|_{C^n}P(r)
 \le K'\|\varphi\|_{C^n}q(r).
\]
Consequently $R\varphi\pm K'\|\varphi\|_{C^n}q$ are exponential
majorants and minorants of $\varphi$, agreeing with it on the support
of $\mu$.  Integrate these inequalities against $\nu$, subtract the
equal integrals against $\mu$, and use the bounded exponential
coefficients.  This proves \eqref{p06:rs:eq:certificate}.
\end{proof}

\subsection{Proof of the exact counts}

\begin{proof}[Proof of Theorem~\ref{p06:rs:thm:counts}]
If $\mu,\nu\in\mathcal P_m[a,b]$, the function
$h(u)=\int e^{ur}\,d(\mu-\nu)(r)$ is a sum of at most $2m$ real
exponentials and satisfies $h(0)=0$.  Equal thresholds at $2m-1$
one-sided settings provide $2m-1$ further distinct zeros, all on one
side of $0$, by Lemma~\ref{p06:rs:lem:contacts}.  The exponential zero
count forces $h=0$, hence $\mu=\nu$.

For an $m$-atomic probability $\mu$ and arbitrary $\nu$, $2m$
one-sided settings provide $2m$ transform zeros, which, together with
$h(0)=0$, give $2m+1$ distinct samples.  Apply
Lemma~\ref{p06:rs:lem:certificate} with zero right side to obtain equality
of integrals against every smooth function and hence equality of
measures.  Without normalization, use $2m+1$ settings directly.

We prove the claimed lower bounds locally, including settings for
which a contact happens to equal $1$.  Write an interior atomic
measure as $\sum_{j=1}^mB_j\delta_{r_j}$ and use its $2m$ mass and
location coordinates.  If $u_i$ is the logarithm of the contact for
$C_i$, the envelope derivatives of $T(C_i)$ are, up to a nonzero row
factor,
\[
                  (e^{r_ju_i},\ B_ju_i e^{r_ju_i})_{j=1}^m.
\]
Thus $K\le2m$ distinct threshold rows have rank $K$.  This is again
the multiplicity zero count for the system
$e^{r_ju},u e^{r_ju}$, or its initial subsystems.

For normalized measures also impose the mass row, which is evaluation
at $u=0$.  If no $u_i=0$, these $K+1\le2m$ rows are independent.
If one $u_i=0$, replace that threshold constraint by the first moment
constraint.  This adds the derivative row at $0$; together with the
mass row and the other $K-1$ contact rows it again has rank $K+1$.
Preserving mass and first moment indeed preserves the omitted
threshold: for probability measures a contact at $1$ occurs exactly
when $\int r\,d\mu=C_i/(1+C_i)$, and its threshold is then $1+C_i$.

For $K<2m-1$ the analytic implicit-function theorem therefore gives
nonconstant normalized $m$-atomic curves with the same constraints,
proving necessity in part (i).  For the stronger lower bound, add an
atom of small positive mass $\epsilon$ at a fixed new interior point.
The same full row rank lets us adjust the original masses and
locations analytically in $\epsilon$ to preserve all $K$ thresholds
and the normalization, whenever $K\le2m-1$.  The original masses
remain positive and the atoms remain distinct, so the result has
exactly $m+1$ atoms.  In the case of a contact at $1$, use the replaced
first-moment constraint as above.  Without normalization, the
$K\le2m$ threshold rows alone allow the identical construction.
\end{proof}

\begin{corollary}[Exact discovery of an unknown finite spectrum]
\label{p06:rs:cor:sequential}
Choose any infinite list of distinct settings satisfying
\eqref{p06:rs:eq:oneside}.  An exact threshold oracle permits the following
information-theoretic procedure with no prior bound on atom count.
After the first $2k$ observations, test whether a member of
$\mathcal P_k[a,b]$ fits them exactly.  Stop at the first successful
test and return that measure.  If the truth has exactly $M$ atoms,
the procedure stops at $k=M$, after $2M$ scalar observations, and
returns the truth uniquely.
\end{corollary}
\begin{proof}
A successful fit with $k<M$ would contradict the blind uniqueness in
Theorem~\ref{p06:rs:thm:counts}.  At $k=M$ the truth is an admissible fit,
and that same uniqueness theorem identifies it.
\end{proof}
This is an exact-data existence procedure.  No polynomial-time claim
is made for its nonlinear feasibility tests.  Nor can exact atom
count be uniformly recovered with positive noise: adding an
arbitrarily small extra atom makes the threshold data arbitrarily
close.  The next result quantifies the meaningful distributional
guarantee that remains.

\subsection{Proof of the sharp blind noise exponent}

\begin{lemma}[One atomic side in a moment comparison]
\label{p06:rs:lem:moments}
On a fixed compact interval, if $\alpha$ is an at most $m$-atomic
probability and $\beta$ any probability, and their moments of orders
$1,\ldots,2m$ differ by at most $\epsilon$, then
\begin{equation}\label{p06:rs:eq:moments}
                   W_1(\alpha,\beta)\le K\epsilon^{1/(2m)}.
\end{equation}
If both have at most $m$ atoms, orders $1,\ldots,2m-1$ suffice with
exponent $1/(2m-1)$.
\end{lemma}
\begin{proof}
Let $F=F_\alpha-F_\beta$ and $E=\int|F|=W_1(\alpha,\beta)$.
Between consecutive atoms of $\alpha$, $F$ is nonincreasing.  Before
its first atom $F\le0$ and after its last atom $F\ge0$.  There are
therefore at most $2m-1$ changes of sign, after zero sets are ignored:
upward changes can occur only at its $m$ atoms, and the initial and
final signs allow at most one fewer downward change.  Choose a point $t_j$ at
each sign change, including a point of an intervening zero interval
when necessary.  For $d\le2m-1$ such points, a sign choice gives
$P(x)=\pm\prod_{j=1}^d(x-t_j)$ with $P(x)F(x)\ge0$ almost everywhere.

For $d\ge1$, remove intervals of radius $E/(4d)$ about these roots.
They remove at most $E/2$ from the integral of $|F|$, since $|F|\le1$.
On their complement $|P|\ge(E/(4d))^d$.  Thus
\[
 \int PF\ge\frac E2\left(\frac E{4d}\right)^d.
\]
For $d=0$ the integral is $E$.  An antiderivative of $P$ has degree
at most $2m$, with uniformly bounded coefficients.  Integration by
parts and the moment assumptions bound $\int PF$ by $K\epsilon$.
This proves \eqref{p06:rs:eq:moments} for $\epsilon\le1$; a diameter
bound handles larger errors.  If both measures are at most
$m$-atomic, their CDF difference has at most $2m-1$ nonzero constant
intervals and hence at most $2m-2$ sign changes.  The same proof
gives the final assertion.
\end{proof}

\begin{proof}[Proof of Theorem~\ref{p06:rs:thm:noise}]
Let $\delta$ be the threshold discrepancy.  For small $\delta$,
Lemma~\ref{p06:rs:lem:contacts} supplies $2m$ uniformly separated near-zeros
of the transform difference.  Adding its exact zero at $0$ gives the
$2m+1$ uniformly separated samples needed by
Lemma~\ref{p06:rs:lem:certificate}.  Apply that lemma to
$\varphi(r)=e^{kr}$, $k=1,\ldots,2m$.  Thus the first $2m$ moment
differences of the pushforwards under $r\mapsto e^r$ are $O(\delta)$.
Lemma~\ref{p06:rs:lem:moments} and the Lipschitz property of $\log$ on
$[e^a,e^b]$ prove \eqref{p06:rs:eq:blindholder}.  Larger $\delta$ are
handled by increasing the constant.

For two at most $m$-atomic measures, $2m-1$ near-zeros and the exact
zero at $0$ are enough.  Their transform difference solves an
order-$2m$ constant-coefficient differential equation whose real
roots are in $[a,b]$, padding the root list if needed.  Evaluation at
$2m$ separated points uniformly controls its initial jet: the
evaluation matrix in the initial-jet basis is continuous in the
roots, including collisions, and invertible by the real-exponential
zero count.  Compactness uniformly bounds its inverse.  Its values
at $1,\ldots,2m-1$ are consequently $O(\delta)$, and the two-atomic-side
part of Lemma~\ref{p06:rs:lem:moments} proves the stated exponent.

For sharpness in the blind case, put $q=2m$ and take sufficiently
small $h>0$ with $r_0,r_0+qh$ in the interior of $[a,b]$.  Split the
binomial masses
\[
                      2^{-(q-1)}\binom qj,\qquad 0\le j\le q,
\]
between their even and odd locations $r_0+jh$.  Let $\mu_h$ be the
odd part, with $m$ atoms, and $\nu_h$ the even part, with $m+1$ atoms.
Both have total mass one, and
\[
 \int e^{ur}\,d(\nu_h-\mu_h)(r)
       =2^{-(q-1)}e^{r_0u}(1-e^{hu})^q.
\]
This is $O(h^{2m})$ on each fixed compact $u$ interval.  The contact
bound in Lemma~\ref{p06:rs:lem:contacts} then gives, on every fixed
compact setting interval $I\subset(0,\infty)$,
\[
          \sup_{C\in I}|T_{\mu_h}(C)-T_{\nu_h}(C)|\le K_I h^{2m}.
\]
The identity
$\sum_{j=0}^k(-1)^j\binom qj=(-1)^k\binom{q-1}k$
gives $W_1(\mu_h,\nu_h)=h$ by integrating their CDF difference over
the $q$ gaps.  This proves the optimal blind exponent.  Taking
$q=2m-1$ instead gives $m$ atoms on both sides and proves the other
sharp exponent.  If rational independence with $1$ is needed for
the original modular-zero interpretation, perturb the locations by
$o(h^q)$ outside the countably many rational affine hyperplanes.
The bounds and atom counts persist.
\end{proof}

\subsection{The precise scope of this transfer}

For normalized compact spectra the distinction is now exact:
\[
\begin{array}{c|c|c}
\text{comparison class}&\text{designed threshold count}
                         &\text{sharp collision exponent}\\ \hline
\text{both at most }m\text{ atoms}&2m-1&1/(2m-1)\\
\text{one at most }m\text{ atoms, other arbitrary}&2m&1/(2m)
\end{array}
\]
The data are still the ideal scalar transitions $T(C)$, defined
through infinite-height modular-zero closure in their motivating
exponential-sum model.  These theorems do not turn finitely many
observed zeros into exact thresholds, do not recover the phases
invisible to those thresholds, and do not give a new recurrence
estimate for the exact prime frequencies.  A quantitative finite-height
threshold estimator is a separate remaining problem.

% ===== bessel =====
\section{Three critical regimes for Bessel power restoration}
\label{p06:bs:section}

This section develops the power-restoration criterion of the preceding
manuscript, rather than counting that criterion or the full McKay ID
domain a second time.  The new conclusions are a sharp three-regime
transition in the necessary power, uniqueness of the critical real
maximizer, and the endpoint behavior of the obstruction.  A separate
moment test advances, but does not finish, the remaining SD problem.
The published questions about the McKay families are in
\cite[Sections~4.3--4.4]{p01:rare:Baricz2026}.  Our power-transition question is a
newly formulated transfer problem; we do not identify it as a previously
named conjecture.

Write
\[
 R_\mu(y)=\frac{I_{\mu+1}(y)}{I_\mu(y)},\qquad
 D_\mu(y)=R_\mu(y)+yR_\mu'(y),\qquad
 M_\mu=\sup_{y>0}D_\mu(y),\qquad \mu>-1.
\]
For a generalized McKay variable with density
\[
 Cx^{\nu-1}e^{-bx}I_\mu(ax),\qquad
 \mu+\nu>0,\quad b>a>0,
\]
put $r=b/a$.  Its least admissible HCM power is
\[
 d_\mu^*(r)=\inf\{d>0:X^d\text{ has an HCM density}\},
 \qquad \inf\varnothing=+\infty.
\]
Monomial factors do not change HCM, so this definition does not depend
on $\nu$.  Negative powers give the same criterion with $|d|$.

We recall the exact criterion from
\cite[Theorems~6.2 and 6.4]{comp:part:05}, in a form that
also fixes all normalizations.  For $0<\theta<\pi/2$, define
\begin{equation}\label{p06:bs:ray-def}
 B_\mu(y,\theta)=
 \frac{\operatorname{Im}\{e^{i\theta}R_\mu(ye^{i\theta})\}}
      {\sin\theta},\qquad
 C_\mu(\theta)=\sup_{y>0}B_\mu(y,\theta).
\end{equation}
The HCM test is
\begin{equation}\label{p06:bs:background-test}
 X^d\text{ is HCM}\quad\Longleftrightarrow\quad
 d>2\ \text{ and }\ r\geq C_\mu(\pi/d).
\end{equation}
Moreover, $M_\mu=1$ for $\mu\geq1/2$, $M_\mu>1$ for $\mu<1/2$,
and some power exists exactly when $r>M_\mu$.
For reference, the product over the positive Bessel zeros gives
\begin{align}
 R_\mu(z)&=\sum_{n\geq1}\frac{2z}{z^2+j_{\mu,n}^2},\notag\\
 B_\mu(y,\theta)
 &=4y\cos\theta\sum_{n\geq1}
 \frac{j_{\mu,n}^2}
 {(j_{\mu,n}^2-y^2)^2+4y^2j_{\mu,n}^2\cos^2\theta}.
 \label{p06:bs:positive-series}
\end{align}
In particular $B_\mu>0$ and
\begin{equation}\label{p06:bs:comparison}
 \cos\theta D_\mu(y)\leq B_\mu(y,\theta)
 \leq\sec\theta D_\mu(y).
\end{equation}
These are classical product identities followed by the previous
Stieltjes boundary test, not new assertions of this section.  The
underlying Bessel formulas can be found in \cite{p01:rare:DLMF}; the HCM
representation and its relation to GGC are in \cite{p05:fs:Bondesson}.

Here is a concise verification of the criterion used below.  Put
$p=1/d$ and $\mathcal I_\mu(z)=2^\mu\Gamma(\mu+1)z^{-\mu}I_\mu(z)$.
Up to a monomial, the powered density is
$g(x)=e^{-b x^p}\mathcal I_\mu(a x^p)$, with $g(0)=1$.
For $0<p<1/2$, its negative logarithmic derivative on the principal
slit plane is
\[
 h(z)=p z^{p-1}\{b-aR_\mu(a z^p)\}.
\]
There are no poles in that plane.  A keyhole contour gives
\[
 h(z)=\int_0^\infty\frac{u(t)}{z+t}\,dt,\qquad
 u(t)=\frac{pa\sin(\pi p)}\pi t^{p-1}
            \{r-B_\mu(a t^p,\pi p)\}.
\]
The small and large contour terms vanish by the bounds
$h(z)=O(|z|^{p-1})$ at zero and
$h(z)=p(b-a)z^{p-1}+O(z^{-1})$ at infinity.
The signed integral converges absolutely.  The normalized HCM/GGC
representation and Stieltjes inversion say exactly that $g$ is HCM
if and only if $u\geq0$, proving the test for these powers.
At $p=1/2$, $h$ has residue $-1$ at every
$-j_{\mu,n}^2/a^2$, incompatible with a positive Stieltjes measure.
For $p>1/2$, $g$ has a zero at
$(j_{\mu,n}/a)^{1/p}e^{i\pi/(2p)}$ inside the slit plane, whereas
a normalized HCM function is zero-free there.
Finally the HCM representation implies
$(xg'/g)'<0$ for this nonconstant $g$; direct differentiation gives
the necessary strict inequalities $r>D_\mu(y)$ for every $y>0$.
When $M_\mu>1$ is attained they imply $r>M_\mu$; when $M_\mu=1$
the same strict inequality is already imposed by $b>a$.
Conversely \eqref{p06:bs:comparison} makes every $r>M_\mu$ feasible
at a sufficiently small angle.  This verifies the background facts
without counting them among the new conclusions.

\subsection{A unique maximizer and its exact curvature}

\begin{theorem}[The shape of the real obstruction]
\label{p06:bs:shape}
If $\mu\geq1/2$, $D_\mu$ is strictly increasing from zero to one.
If $-1<\mu<1/2$, it has a unique global maximum, at a point
$y_\mu>0$, and
\begin{equation}\label{p06:bs:curvature}
 D_\mu(y_\mu)=M_\mu>1,\qquad
 D_\mu''(y_\mu)=-\frac{2(M_\mu^2-1)}{y_\mu}<0.
\end{equation}
For $-1/2\leq\mu<1/2$ this is its only stationary point.
For $-1<\mu<-1/2$ it has exactly one further stationary point,
a minimum below one, after which it increases to one.
\end{theorem}
\begin{proof}
Set $\alpha=\mu+1/2$, $v(y)=R_\mu(y)+\alpha/y$, and
$D=D_\mu$.  The Bessel recurrence yields
\begin{equation}\label{p06:bs:phase-system}
 R'=1-2\alpha R/y-R^2,\qquad
 v'=1-v^2+\frac{\alpha(\alpha-1)}{y^2}
       =\frac{D-v}{y},\qquad
 D'=2(1-vD).
\end{equation}
At a stationary point of $v$,
\[
 v''=-2\alpha(\alpha-1)/y^3.
\]
At a stationary point of $D$, positivity of $D$, supplied by the
product formula, gives $v=1/D>0$ and hence
\begin{equation}\label{p06:bs:stationary-curvature}
 D''=-2D v'=-2(D^2-1)/y.
\end{equation}
No stationary point has $D=1$.  Such a point would have $v=1$ and
$v'=0$, forcing $\alpha(\alpha-1)=0$ in the Riccati equation;
the two exceptional solutions are $v=\tanh y$ for $\alpha=0$
and $v=\coth y$ for $\alpha=1$, neither of which equals one at a
finite positive point.

The endpoint expansions are
\begin{align*}
 D(y)&=\frac{y}{\mu+1}+O(y^3)\quad(y\downarrow0),\\
 v(y)&=1+\frac{\alpha(\alpha-1)}{2y^2}+O(y^{-3}),\\
 D(y)&=1-\frac{\alpha(\alpha-1)}{2y^2}+O(y^{-3})
                  \quad(y\to\infty),
\end{align*}
with differentiated expansions valid as well.
If $0<\alpha<1$, $v'$ is initially negative and eventually positive,
and every zero of $v'$ is a strict upward crossing.  Thus $v$ has
exactly one minimum.  Before it, $D<v$, so every possible stationary
point of $D$ would be a minimum below one.  There can be none because
$D'>0$ initially.  After it, $D>v$, so every stationary point of $D$
is a maximum above one.  Since $D'$ is eventually negative, there is
exactly one such maximum.

If $\alpha<0$, $v'$ is initially positive and eventually negative,
and every zero is a strict downward crossing.  Thus $v$ has one
maximum.  Before it every stationary point of $D$ is a maximum above
one; after it every stationary point is a minimum below one.
There is a value $D>1$: otherwise
$[y(1-R_\mu(y))]'=1-D\geq0$, whereas this expression starts at zero
and tends to $\alpha<0$.  Since $D$ is eventually below one and
increasing, the stated maximum and minimum both exist and are unique.
For $\alpha=0$, the explicit solution $R=\tanh y$ gives the unique
maximum from $y\tanh y=1$.

If $\alpha>1$, $v'$ is negative initially, and every possible zero
is a downward crossing; hence there is no zero and $v'<0$ throughout.
All possible stationary points of $D$ would be minima below one,
which cannot be the first zero of $D'$ because $D'>0$ initially.
Consequently $D'>0$ throughout.  The same conclusion follows from
$v=\coth y$ when $\alpha=1$.  The endpoint limits complete all cases.
Equation~\eqref{p06:bs:stationary-curvature} proves the asserted curvature.
\end{proof}

\subsection{The critical power has three different asymptotic regimes}

\begin{theorem}[The three-regime critical law]
\label{p06:bs:three-regimes}
For fixed $\mu>-1$, the following statements hold.
\begin{enumerate}
\item If $-1<\mu<1/2$, then, as $r\downarrow M_\mu$,
\begin{equation}\label{p06:bs:subcritical-order}
 d_\mu^*(r)\sim
 \pi\sqrt{\frac{y_\mu(M_\mu^2-1)}{3(r-M_\mu)}}.
\end{equation}
\item At $\mu=1/2$, as $r\downarrow1$,
\begin{equation}\label{p06:bs:critical-order}
 d_{1/2}^*(r)=\log\frac1{r-1}+\log2-2+o(1).
\end{equation}
\item If $\mu>1/2$, there is a finite $d_0(\mu)>2$ such that
every $d\geq d_0(\mu)$ works simultaneously for every $r>1$.
In particular
\begin{equation}\label{p06:bs:supercritical-order}
 \sup_{r>1}d_\mu^*(r)<\infty.
\end{equation}
\end{enumerate}
Thus the order $1/2$ separates an inverse-square-root divergence,
a logarithmic divergence, and a uniformly bounded power cost.
\end{theorem}
\begin{proof}
For the first case, analytic expansion at $\theta=0$ on each compact
positive $y$ interval gives
\begin{equation}\label{p06:bs:angular-expansion}
 B_\mu(y,\theta)=D_\mu(y)
 -\frac{\theta^2}{6}\{3yD_\mu'(y)+y^2D_\mu''(y)\}
 +O(\theta^4).
\end{equation}
One can obtain this by applying the operator $1+y\partial_y$ to
$R_\mu$ three times in the numerator of \eqref{p06:bs:ray-def} and
expanding $\sin\theta$.  The error and its first two $y$ derivatives
are uniform on such compact sets.

When $M_\mu>1$, all maximizers for sufficiently small $\theta$ lie
in a fixed compact interval.  Indeed, $D_\mu$ tends to zero at zero
and to one at infinity; \eqref{p06:bs:comparison} excludes both ends.
The unique maximum in Theorem~\ref{p06:bs:shape} and its negative second
derivative then give, by the implicit function theorem, a unique
maximizing point $y(\theta)=y_\mu+O(\theta^2)$.  Substitution in
\eqref{p06:bs:angular-expansion} gives
\begin{equation}\label{p06:bs:quadratic-tilt}
 C_\mu(\theta)=M_\mu+
       \frac{y_\mu(M_\mu^2-1)}3\theta^2+O(\theta^4).
\end{equation}
The coefficient is strictly positive.  To justify inversion at the
least feasible power, note also that $C_\mu(\theta)>M_\mu$ for each
$\theta>0$: equality would make the normalized non-power factor
HCM at tilt $r=M_\mu$, contrary to the strict fixed-tilt obstruction
at the attained maximum of $D_\mu$.
Continuity of $C_\mu$ on compact subintervals of $(0,\pi/2)$ follows
from its endpoint limits in $y$ and the locally uniform Bessel
asymptotics.  Finally \eqref{p06:bs:positive-series}, evaluated at
$y=j_{\mu,1}$, gives
$C_\mu(\theta)\geq1/(j_{\mu,1}\cos\theta)$.
Thus feasible powers must tend to infinity as $r\downarrow M_\mu$;
\eqref{p06:bs:background-test} and \eqref{p06:bs:quadratic-tilt} give
\eqref{p06:bs:subcritical-order}.

At $\mu=1/2$, $R_\mu(z)=\coth z-1/z$.  Put $k=\cot\theta$ and
$t=2y\sin\theta$.  Direct algebra gives
\begin{equation}\label{p06:bs:coth-ray}
 B_{1/2}(y,\theta)-1
 =\frac{\cos t-k\sin t-e^{-kt}}{\cosh(kt)-\cos t}.
\end{equation}
For $0<t\leq\pi$, the numerator is nonpositive: the derivative of
$e^{kt}(\cos t-k\sin t)$ is
$-(1+k^2)e^{kt}\sin t\leq0$, and its value at zero is one.
For $t\geq\pi$ the denominator is
$\tfrac12e^{kt}(1+O(e^{-\pi k}))$, uniformly as $k\to\infty$.
Consequently the supremum of \eqref{p06:bs:coth-ray} equals
\[
 2\sup_{t\geq\pi}e^{-kt}(\cos t-k\sin t)
       +O((1+k)e^{-2\pi k}).
\]
The first positive maximum of the last expression occurs at
$t=\pi+2\arctan(1/k)$ and has trigonometric factor exactly one.
Subsequent positive maxima are smaller by successive factors
$e^{-2\pi k}$.  It follows that
\begin{equation}\label{p06:bs:essential}
 C_{1/2}(\theta)-1
 \sim 2\exp\{-\pi\cot\theta-2\}
                         \qquad(\theta\downarrow0).
\end{equation}
For every fixed $\theta>0$ the right side of
\eqref{p06:bs:coth-ray} is positive somewhere: use sufficiently large
$t$ in a positive trigonometric interval.  Thus
$C_{1/2}(\theta)>1$.  Compactness and the same resonance lower bound
as above show that feasible powers tend to infinity when $r\downarrow1$.
Since $\pi\cot(\pi/d)=d+o(1)$, inversion of
\eqref{p06:bs:essential} proves \eqref{p06:bs:critical-order}.

For $\mu>1/2$, put $a_2=(4\mu^2-1)/8>0$.
The fixed-order Bessel expansion, uniform in a closed sector
$|\arg z|\leq\theta_1<\pi/2$, gives
\begin{equation}\label{p06:bs:uniform-tail}
 B_\mu(y,\theta)=1-a_2/y^2+O(y^{-3})
 \quad(y\to\infty),\qquad 0\leq\theta\leq\theta_1.
\end{equation}
The value at zero angle is the continuous extension $D_\mu$.
Uniformity after division by $\sin\theta$ follows by applying
Cauchy's estimates in a slightly wider sector to the remainder;
its imaginary part vanishes on the positive real axis and is
$O(\theta y^{-3})$ after multiplication by $e^{i\theta}$.
Choose $Y$ large enough that the expression is strictly below one
for $y\geq Y$ and all these angles.  At the other end the power
series at zero makes $B_\mu<1/2$ uniformly for small $y$ and small
angles.  On the intervening compact interval,
Theorem~\ref{p06:bs:shape} gives $D_\mu<1$; continuity therefore gives
$B_\mu<1$ there for all sufficiently small positive $\theta$.
We have proved $B_\mu(y,\theta)<1$ for every finite $y$ at all
sufficiently small angles.  Its limit as $y\to\infty$ is one, so
$C_\mu(\theta)=1$ at all those angles.  The exact test now proves
\eqref{p06:bs:supercritical-order}.
\end{proof}

The last part strengthens the previous statement that every fixed
admissible tilt can be repaired.  Above order one half, the power can
be chosen before the tilt is specified.  At and below one half such
a uniform choice is impossible near the critical tilt.

\subsection{The two order endpoints of the obstruction}

\begin{theorem}[Endpoint asymptotics of the critical tilt]
\label{p06:bs:endpoints}
As $\varepsilon\downarrow0$,
\begin{align}
 M_{-1+\varepsilon}&\sim
                  \frac{3\sqrt3}{8\sqrt\varepsilon},&
 y_{-1+\varepsilon}&\sim\frac{2\sqrt\varepsilon}{\sqrt3},
                    \label{p06:bs:minus-one}\\
 M_{1/2-\varepsilon}-1&\sim
                   \frac{2\varepsilon}{\log^2(1/\varepsilon)},&
 y_{1/2-\varepsilon}&\sim\frac12\log(1/\varepsilon).
                    \label{p06:bs:half-endpoint}
\end{align}
\end{theorem}
\begin{proof}
The positive product series gives
\[
 D_\mu(y)=\sum_{n\geq1}
           \frac{4y j_{\mu,n}^2}{(y^2+j_{\mu,n}^2)^2}.
\]
When $\mu=-1+\varepsilon$, the first zero satisfies
$j_{\mu,1}^2=4\varepsilon+O(\varepsilon^2)$.
For completeness, multiply the power series
${}_0F_1(;\varepsilon;-w/4)$ by $\varepsilon$; near $(w,\varepsilon)
=(0,0)$ it is $\varepsilon-w/4+O(w^2)$, so the analytic implicit
function theorem gives this small zero.  The remaining zeros stay
uniformly away from zero and have uniformly linear growth in their
index; they converge to the positive zeros of $J_1$, since
$J_{-1}=-J_1$.  These facts also follow directly from the standard
zero-continuity and large-index formulas in \cite[Section~10.21]{p01:rare:DLMF}.
Splitting the remaining sum at $j_n=y$ shows that its contribution
is $O(1)$ uniformly in $y>0$ and in small $\varepsilon$.
The first summand, with $j=j_{\mu,1}$ and $y=ju$, is
\[
 \frac1j\frac{4u}{(1+u^2)^2}.
\]
Its unique maximum is $3\sqrt3/(4j)$, at $u=1/\sqrt3$.
Uniform boundedness of the other terms and the separated strict
maximum of this function yield both assertions in
\eqref{p06:bs:minus-one}.

For the second endpoint, write $L=\log(1/\varepsilon)$.
We use the parameter-uniform estimate
\begin{equation}\label{p06:bs:order-perturbation}
 D_{1/2-\varepsilon}(y)-D_{1/2}(y)
 =\frac{\varepsilon(1-\varepsilon)}{2y^2}
                       +O(\varepsilon y^{-3}),\qquad y\geq Y,
\end{equation}
where $Y$ and the constant are independent of small $\varepsilon$.
This follows from the fixed-order expansion for $D_\mu$, locally
uniform and analytic in the order near $\mu=1/2$.  More explicitly,
after subtracting $1-(4\mu^2-1)/(8y^2)$ the remainder is $O(y^{-3})$
uniformly on a small complex disk in $\mu$; Cauchy's derivative
estimate and the mean value formula give an $O(\varepsilon y^{-3})$
difference.  Thus an error independent of $\varepsilon$ is not being
substituted into this singular limit.

The exact half-order formula gives
\[
 D_{1/2}(y)-1=\coth y-y\operatorname{csch}^2y-1
                 =-(4y-2)e^{-2y}+O(ye^{-4y}).
\]
At $y=L/2+2\log L$, these formulas give
\[
 D_{1/2-\varepsilon}(y)-1
                      \geq(2-o(1))\varepsilon/L^2.
\]
On every fixed compact positive $y$ interval, $D_{1/2}<1$, so the
maximizer tends to infinity.  For any fixed $0<\eta<1$, the negative
exponential term dominates the positive perturbation uniformly when
$Y\leq y\leq(1-\eta)L/2$, once $Y$ is sufficiently large and
$\varepsilon$ sufficiently small.  In particular the maximizing
point lies beyond $(1-\eta)L/2$.  For larger $y$,
\eqref{p06:bs:order-perturbation} and $D_{1/2}<1$ imply
\[
 D_{1/2-\varepsilon}(y)-1
 \leq\frac{\varepsilon}{2y^2}(1+O(1/y)).
\]
This bounds the maximum by
$(2+o(1))\varepsilon/[(1-\eta)^2L^2]$.
Letting $\eta\downarrow0$ proves the first assertion of
\eqref{p06:bs:half-endpoint}.  The same bound excludes
$y\geq(1+\eta)L/2$ from being a maximizing point, because it is
strictly smaller than the already established lower bound there.
The two exclusions prove the location asymptotic.
\end{proof}

\subsection{A finite moment obstruction for the remaining SD problem}

Write $\alpha=\mu+1/2$, $q=\nu-\mu-1>0$, $A=\alpha+q$, and
$M=2\alpha+q$.  Admissibility means $\alpha>-1/2$ and $M>0$.
The signed-spectral construction of
\cite[Theorem~4.3 and Corollary~4.4]{comp:part:04}
gives a positive even
kernel $K$ with
\[
 -J'(u)/J(u)=\int_0^\infty e^{-ux}K(x)\,dx,\quad
 J(u)=C\frac{u^q\,{}_2F_1(-q/2,(1-q)/2;\alpha+1/2;u^{-2})}
                     {(u^2-1)^A},
\]
and $e^{-x}K(x)\to A$, $K(0)=M$.
The exact SD boundary is
$\beta_{\alpha,q}=\sup_{x\geq0}K'(x)/K(x)\geq1$.
Positive integers and positive-cut intervals $q\in(2m,2m+1)$ were
already proved to have $\beta>1$.  The previous regular-part argument
also covered $\alpha\leq1$.  The following criterion is new progress
in the remaining negative-cut intervals.

\begin{proposition}[An explicit extension and a hierarchy of certificates]
\label{p06:bs:sd-hierarchy}
In a negative-cut interval $q\in(2m+1,2m+2)$,
\begin{equation}\label{p06:bs:sd-explicit}
 1<\alpha\leq\frac{3+\sqrt{1+8q}}2
             \quad\Longrightarrow\quad\beta_{\alpha,q}>1.
\end{equation}
More generally, for $\alpha>1$ in such an interval, define
\begin{align}
 h_0&=\frac{\alpha(\alpha-1)}{2(A-1)},\notag\\
 h_n&=\frac{(n+1)h_{n-1}
       -2\sum_{i+j=n-1}h_i h_j-\sum_{i+j=n-2}h_i h_j}
                      {2(A-n-1)},\qquad n<A-1,
                    \label{p06:bs:moment-recursion}
\end{align}
where an empty sum is zero.  If any defined value satisfies
$(-1)^n h_n\leq0$, then $\beta_{\alpha,q}>1$.
More generally, if a Hankel matrix built from
$m_n=(-1)^n n!h_n$ within this index range is not positive
semidefinite, the same conclusion follows.
\end{proposition}
\begin{proof}
Suppose $\beta=1$.  Then $k(x)=e^{-x}K(x)$ is nonincreasing to $A$,
and $k-A\geq0$, strictly positive near zero since $k(0)-A=\alpha>0$.
Let
\[
 H(s)=-J'(1+s)/J(1+s)-A/s
                 =\int_0^\infty e^{-sx}(k(x)-A)\,dx.
\]
In a negative-cut interval the signed spectral formula shows
$k(x)-A=O(x^{-A})$ at infinity; the finite zero terms and the
elementary terms decay exponentially.  Hence moments through every
integer $n<A-1$ are finite, and
$h_n=H^{(n)}(0)/n!$ satisfy $(-1)^nh_n>0$.
All moment Hankel matrices, both unshifted and shifted when their
indices are available, must be positive semidefinite.

The Bessel differential equation, integrated against the defining
Laplace density, gives
\[
 (u^2-1)J''+(2\nu+1)uJ'+(\nu^2-\mu^2)J=0.
\]
Putting $-J'/J=A/s+H(s)$ with $u=1+s$ yields the identity
\[
 s(s+2)(H^2-H')+(2A-2-2s)H+\alpha(1-\alpha)=0.
\]
Expansion through any available finite derivative order gives exactly
\eqref{p06:bs:moment-recursion}.  A failed sign or moment-matrix test
therefore contradicts $\beta=1$.

In particular $A>2$ in the stated range, and
\[
 h_1=\frac{h_0(1-h_0)}{A-2}.
\]
If $h_0\leq1$, this is nonnegative, contradicting the strict
negative derivative of the Laplace transform of $k-A$ at zero.
The inequality $h_0\leq1$ is precisely
$\alpha(\alpha-1)\leq2(\alpha+q-1)$, which gives
\eqref{p06:bs:sd-explicit}, including equality.
\end{proof}

The hierarchy can go beyond its first-derivative corollary.  For
example $(\alpha,q)=(5,3/2)$ is outside
\eqref{p06:bs:sd-explicit}, but the exact recurrence gives
$h_0=20/11$, $h_1=-40/121$, and $h_2=-360/1331$.
The negative second moment certifies $\beta_{5,3/2}>1$.
This illustrates the same test, not an additional problem closure.

\paragraph{Adversarial scope check.}
The hierarchy is a sufficient obstruction, not a proved terminating
decision procedure.  We have not proved that it succeeds for every
$\alpha>1$ and every negative-cut $q$, and we therefore do not claim
universal SD strictness.  The three-regime HCM theorem and its endpoint
asymptotics are complete within their stated model.  Their proof uses
uniform estimates before taking a supremum, accounts for the unattained
maximum $M_\mu=1$, and treats $\mu=1/2$ separately; replacing that case
by a power-series approximation would miss the exponential transition.

The new preprint \cite{p02:bes:BariczPrabhu2026} reports a proof of the
Ismail--Miller Bessel-ratio ID conjecture.  That is a different problem
and is not claimed as a closure here.  No published priority for our
three-regime statement was identified in the primary sources checked;
this is a limited literature check, not a guarantee of novelty.

% ===== smooth_budget =====
\section{Analytic populations and iterated critical budget penalties}
\label{p06:sb:section}

The preceding endpoint theorem showed that the unrestricted scalar risk is
$\log(B)/\sqrt B$, even under adaptive replication. We now ask whether
analytic regularity of the population removes this critical penalty.
It reduces it, but does not in general remove it. A spectral smoothness
class leaves a smaller, precisely calculable number of resolution levels.
The following results concern new internally formulated minimax problems.

An individual has latent $U\in[0,1]$, with density $w$, and observations
conditionally uniform on $[0,2^U]$. A policy may revisit individuals and
start new ones using at most $B$ scalar observations. There is no
additional replication cap. Loss is expected $W_1$ for the law of $U$;
by a bi-Lipschitz change of variable the same rates hold for $2^U$.
Let $\mathcal W_\phi(L)$ consist of real, one-periodic densities with
\[
 w\ge1/2,\qquad \int_0^1w=1,\qquad
 |\widehat w(k)|\le L e^{-\phi(|k|)}\quad(k\in\mathbb Z\setminus\{0\}),
\]
where $\widehat w(k)=\int_0^1e^{-2\pi iku}w(u)\,du$.
The periodic restriction is part of the class, not an identification of
the two endpoints in the transportation metric. All constants $L,\tau,\nu$
below are fixed positive numbers, independent of $B$.

\begin{theorem}[A double logarithm for every fixed stretched-exponential spectrum]
\label{p06:sb:gevrey}
For $\phi(k)=\tau k^\nu$, $\nu>0$, the adaptive and predetermined
minimax risks on $\mathcal W_\phi(L)$ satisfy
\[
 R_B^{\rm ad}\asymp R_B^{\rm fix}\asymp
                \frac{\log\log B}{\sqrt B},\qquad B\longrightarrow\infty.
\]
For $\nu=1$ this is an analytic class with a fixed strip-width envelope.
For $\nu>1$ all its densities extend to entire periodic functions.
Despite that stronger regularity, the same double logarithm remains.
\end{theorem}

\begin{theorem}[A triple logarithm for a double-exponential spectrum]
\label{p06:sb:double}
For $\phi(k)=\tau e^{\nu k}$ the corresponding risks satisfy
\[
 R_B^{\rm ad}\asymp R_B^{\rm fix}\asymp
             \frac{\log\log\log B}{\sqrt B},\qquad B\longrightarrow\infty.
\]
The equivalences are asymptotic; one may replace the iterated logarithms
by positive shifted versions to state bounds for all $B\ge2$.
\end{theorem}

These statements count the distinct regularity classes as two targets,
not every possible value of $\nu$ as a separate solved problem.
Classical scale mixtures of uniforms and fixed-noise deconvolution are
studied in \cite{p02:bud:PW2012,p05:fb:DFM2015}; the experiment here instead
optimizes repeated observations under a total budget and allows arbitrary
adaptive selection of the next individual.

\subsection{A mean correction avoids circular transportation}

\begin{lemma}[A spectral upper bound on an interval]
\label{p06:sb:interval}
For any $w,v\in\mathcal W_\phi(L)$ and $K\ge1$,
\[
 W_1(w(u)\,du,v(u)\,du)
 \le |\mu(w)-\mu(v)|
 +C\left(\sum_{1\le |k|\le K}
             \frac{|\widehat w(k)-\widehat v(k)|^2}{k^2}\right)^{1/2}
 +CL e^{-\phi(K)},
\]
where $\mu(w)=\int_0^1u w(u)\,du$ and $\phi$ is increasing.
\end{lemma}
\begin{proof}
Put $F(x)=\int_0^x(w-v)$ and $\bar F=\int_0^1F$.
The ordinary interval Wasserstein identity is $W_1=\int_0^1|F|$.
Integration by parts gives $\bar F=\mu(v)-\mu(w)$.
For $k\ne0$, the Fourier coefficient of $F$ is
$(\widehat w(k)-\widehat v(k))/(2\pi ik)$, because $F(0)=F(1)=0$.
Hence $W_1\le|\bar F|+\|F-\bar F\|_2$, and Parseval gives the
low-frequency term. The squared remaining tail is bounded by
$C L^2e^{-2\phi(K)}\sum_{k>K}k^{-2}$, which proves the assertion.
In particular a periodic Fourier metric alone is not substituted for
Euclidean transportation; the separately estimated mean is essential.
\end{proof}

Write $a=\log2$. With $r$ observations from an individual, the log maximum
has distribution $Y_r=U-\operatorname{Exp}(ar)$. Thus
\[
 A_{r,k}(Y_r)=e^{-2\pi ikY_r}\left(1-\frac{2\pi ik}{ar}\right)
\]
is unbiased for $\widehat w(k)$, and is uniformly bounded when $r\asymp k$.
The statistic $Y_1+1/a$ is unbiased for $\mu(w)$, with uniformly bounded
variance. Spend one quarter of the budget on this mean statistic.
Partition the frequencies up to a dyadic $K$ into $J=O(\log(eK))$ bands.
For the band ending at $k$, use $r=k$ observations on each of
$n_k=\lceil cB/(Jk)\rceil$ fresh individuals. Choose the fixed $c>0$
small enough. The total band cost is at most $B/2+CK$, and
\[
 \sum_{1\le |l|\le K}\frac{\mathbb E|\widehat m_l-\widehat w(l)|^2}{l^2}
 \le C\sum_{k\text{ dyadic}\le K}\frac1{k n_k}
 \le C J^2/B.
\]
The extra rounding cost is harmless for the polylogarithmic cutoffs used
below. Project the empirical mean and Fourier vector onto the set supplied
by $\mathcal W_\phi(L)$, using the weighted Euclidean norm appearing in
Lemma~\ref{p06:sb:interval}. Fourier summability makes this density class
compact in the uniform norm. A measurable approximate minimizer, or a
countable dense approximation, therefore supplies a probability estimator.
The projection increases its discrepancy from the truth by at most a
factor two. We have proved the upper bound
\begin{equation}\label{p06:sb:upper}
 R_B^{\rm fix}\le C\left\{e^{-\phi(K)}+
                           \frac{\log(eK)}{\sqrt B}\right\},
 \qquad K\le c'B.
\end{equation}

\subsection{A uniform information bound after every history}

For completeness we specify the adaptive lower-bound input, which also
underlies the unrestricted critical theorem in the preceding manuscript
\cite{comp:part:05}. For $K\ge1$ put
\[
 w_\beta(u)=1+\sqrt2\sum_{k=1}^K
       \{\beta_{k,c}\cos(2\pi ku)+\beta_{k,s}\sin(2\pi ku)\},
 \qquad |\beta_{k,*}|\le b_K:=c K^{-2}e^{-\phi(K)}.
\]
Choose $c$ depending only on $L$ and small enough. These densities belong
to $\mathcal W_\phi(L)$, are between $1/2$ and $3/2$, and have a uniformly
bounded Lipschitz constant. These bounds hold for every $K$.

\begin{lemma}[Adaptive trace bound for Fourier perturbations]
\label{p06:sb:trace}
For every adaptive $B$-observation policy, the Fisher information trace
in the $2K$ parameters $\beta$ is at most $CB$, uniformly in $K$ and
$\beta$ in the displayed cube.
\end{lemma}
\begin{proof}
Under the baseline density one, the posterior of an individual after
$r$ observations is
$q(u)\propto e^{-ar u}\mathbf1_{[b,1]}(u)$ for some $0\le b<1$.
This remains the posterior conditional on the full adaptive history,
since the known selection rule contributes no parameter-dependent
factor beyond the observed values. Denote by $q_Y$ the posterior after
the next observation. Two elementary truncated-exponential identities are
\[
 \mathbb E_0\|q_Y-q\|_2^2=a/2,
 \qquad \|q\|_2^2\operatorname{Var}_q(U)\le C.
\]
To check the first, put $\lambda=ar$, $D=1-b$, $\mu=\lambda+a$,
and $Z_\lambda=(1-e^{-\lambda D})/\lambda$, with
$Z_0=D$. The current squared norm is
$(1+e^{-\lambda D})/(2Z_\lambda)$.
The next observation has a case corresponding to no increase in $b$,
plus a density for a new record $b+v$, $0<v<D$. Their respective
contributions to the expected next squared norm give
\begin{align*}
 \mathbb E_0\|q_Y\|_2^2
 &=\frac{1+e^{-\mu D}}{2Z_\lambda}
   +\frac{a}{2Z_\lambda}\int_0^D
       e^{-\lambda v}\{1+e^{-\mu(D-v)}\}\,dv\\
 &=\frac a2+\frac{1+e^{-\lambda D}}{2Z_\lambda}.
\end{align*}
Now $\mathbb E_0q_Y=q$ proves the identity. The formulas extend by
continuity to $\lambda=0$. For the second identity, the squared norm
is $O(\lambda+D^{-1})$ and the variance is
$O(\min(D^2,\lambda^{-2}))$, with the evident convention at zero;
use $D\le1$.

Put $z=\int w_\beta q$ and $z_Y=\int w_\beta q_Y$.
Both are uniformly bounded above and away from zero. The parameter score
of the conditional next-observation density has coordinates
\[
 \int\varphi_j(u)\left\{q_Y(u)/z_Y-q(u)/z\right\}\,du,
\]
where $\varphi_j$ runs over the orthonormal sine and cosine functions.
The true conditional law has density $z_Y/z$ relative to the baseline
conditional law. Bessel's inequality and boundedness of these ratios
reduce the information trace to
\[
 C\mathbb E_0\|q_Y-q\|_2^2+
 C\|q\|_2^2\mathbb E_0|z_Y-z|^2.
\]
The second term is bounded because
\[
 \mathbb E_0|z_Y-z|^2\le C\mathbb E_0 W_1(q_Y,q)^2
          \le 2C\operatorname{Var}_q(U).
\]
For the last inequality couple independently a draw from $q_Y$ with a
draw from $q$, and use $\mathbb E_0q_Y=q$.
Every next-query information trace is therefore bounded by a constant.
Conditional score increments are orthogonal martingale differences,
so their traces add through at most $B$ queries. Independent policy
randomization has zero score. This proves the lemma.
\end{proof}

Use independent cosine-squared priors on $[-b_K,b_K]$. Their total
information is $C K b_K^{-2}$, so the joint prior-and-data Fisher matrix
satisfies
\[
                  \operatorname{tr}J\le C\{B+K^5e^{2\phi(K)}\}.
\]
We recall explicitly the absolute-loss information inequality. For any
fixed vector $v$ and any data estimator $T$,
\begin{equation}\label{p06:sb:absolute}
 \mathbb E|T-v^T\beta|\ge\tfrac14\sqrt{v^TJ^{-1}v}.
\end{equation}
Indeed, let $p(\beta,y)$ be the joint parameter-data density and
put $p_h(\beta,y)=p(\beta-hz,y)$. The total variation between
these two distributions is at most
$|h|\sqrt{z^TJz}/2$, by the fundamental theorem of calculus and
Cauchy--Schwarz. Translation is valid after extending the density by
zero; the prior vanishes quadratically at its boundary. The target
under $p$ is $v^T\beta$, and the translated target under $p_h$ is
$v^T(\beta-hz)$. The risks for these respective targets are equal
by change of variables. On the common part of the two densities,
the targets differ by $hv^Tz$, so the two losses sum to at least
that absolute difference. Choosing
$|h|=(z^TJz)^{-1/2}$ gives the bound with $|v^Tz|$ in place of the
quadratic expression. Optimizing $z$ proves \eqref{p06:sb:absolute}.

For the centered CDF the vector multiplying $\beta$ is
\[
 v_{k,c}(x)=\frac{\sqrt2\sin(2\pi kx)}{2\pi k},\qquad
 v_{k,s}(x)=-\frac{\sqrt2\cos(2\pi kx)}{2\pi k}.
\]
Centering a CDF error increases its integrated absolute value by at most
a factor two. Applying \eqref{p06:sb:absolute} at each $x$ therefore gives
a risk lower bound by $c\int_0^1\sqrt{v(x)^TJ^{-1}v(x)}\,dx$.
For fixed $v$, this is convex in $J$, since
\[
 \sqrt{v^TJ^{-1}v}=\sup_{z\ne0}
                    \frac{|v^Tz|}{\sqrt{z^TJz}}.
\]
Average $J$ over the simultaneous frequency rotations corresponding
to translation of $x$. This kills cross-frequency terms and gives blocks
$j_k I_2$; the integral is unchanged by each rotation and can only
decrease under averaging. The resulting integrand is independent of
$x$, and Cauchy--Schwarz gives
\[
 \sum_{k=1}^K\frac1{k^2j_k}
 \ge\frac{(\sum_{k=1}^K k^{-1})^2}{\sum_{k=1}^K j_k}.
\]
Consequently every adaptive procedure has risk at least
\begin{equation}\label{p06:sb:lower}
 R_B^{\rm ad}\ge
       \frac{c\log(eK)}{\sqrt{B+K^5e^{2\phi(K)}}}.
\end{equation}

\begin{proof}[Proof of Theorems~\ref{p06:sb:gevrey} and~\ref{p06:sb:double}]
For $\phi(k)=\tau k^\nu$, choose an upper-bound cutoff within a fixed
factor of $(\log B/\tau)^{1/\nu}$, large enough that
$\phi(K)\ge\log B$. Equation~\eqref{p06:sb:upper} is then
$O(\log\log B/\sqrt B)$. For the lower bound choose
$K=\lfloor(\log B/(4\tau))^{1/\nu}\rfloor$.
Then $K^5e^{2\phi(K)}\le B$ for all sufficiently large $B$, and
$\log K\asymp\log\log B$. Equation~\eqref{p06:sb:lower} matches the upper
bound. Taking a nearby dyadic integer for the upper bound changes only
constants depending on $\tau,\nu$.

For $\phi(k)=\tau e^{\nu k}$ choose an upper cutoff at least
$\nu^{-1}\log(\log B/\tau)$, rounding upward. The bias is at most
$B^{-1}$ and $\log K\asymp\log\log\log B$.
For the lower bound round down
$K=\nu^{-1}\log(\log B/(4\tau))$.
Again $K^5e^{2\phi(K)}\le B$ eventually, and the same third iterated
logarithm appears in \eqref{p06:sb:lower}. Dyadic rounding in the upper
bound preserves $K\asymp\log\log B$ and only improves its bias.
\end{proof}

\paragraph{Interpretation and scope.}
The penalty is the number of statistically relevant dyadic resolution
bands, divided by $\sqrt B$. Finitely many derivatives leave a polynomial
frequency cutoff and the old $\log B$ penalty. A stretched-exponential
spectral envelope leaves a polylogarithmic cutoff and a $\log\log B$
penalty. A double-exponential envelope leaves a logarithmic-logarithmic
cutoff and a $\log\log\log B$ penalty. Only the last two assertions are
new theorems of this section; the first is an explanation of the earlier
critical mechanism, not an additional claimed closure.

The results use a known, fixed periodic regularity envelope and do not
prove adaptation to unknown smoothness, a capped version, or a uniform
statement over all analytic densities without a common envelope. They
also do not imply that every individual analytic density incurs the
minimax penalty. The lower bound ranges over a fixed infinite-dimensional
class whose least-favourable elements depend on the budget.
\section{What the new results leave open}
\label{p06:br:open}

The preceding proofs close the sixteen scoped targets in the register,
not sixteen independent established conjectures. Several consequences
share one mechanism. This distinction is especially important for the
derivative tower and the equal-replication budget: a family of parameter
values is not an infinite collection of separate breakthroughs.

\begin{enumerate}
\item \textbf{Arithmetic precision.} The endpoint count does not give
an asymptotic density constant, a variance law, or an entropy-sharp
quantitative return time for growing collections of exact prime phases.
Finite-dimensional Kronecker density suffices for the present lower
bound and supplies none of those quantitative statements.
\item \textbf{The coefficient conjecture.} The full statistical MLE
curve bypasses the constant-one Chebyshev--Bernstein bound. It does not
prove that bound. Optimal leading risk constants and sharp predictive
Hellinger logarithms are also not determined here.
\item \textbf{Unrestricted exact information.} Three-query optimality
is proved for polynomial or convex analytic classes on all of
$\mathbb R^3$. Two arbitrary continuous queries, or two nonconvex
real-analytic queries, remain outside that lower-bound argument.
The global unbounded domain is part of the theorem.
\item \textbf{The full design-tail classification.} The matrix converse
includes polynomially sparse and nearly harmonic design sets, but
does not classify every irregular sequence with $\sum_jn_j^{-1}<\infty$
against every possible covariance-tail envelope. Population
nonidentifiability is also distinct from a sharp finite-sample rate.
\item \textbf{Finite-height spectral observations.} The spectral branch
uses ideal nonlinear transition values. It gives no theorem turning
finitely many arithmetic zeros below a prescribed height into those
values with a controlled error. Its blind consistency exponent is not
the restricted $m$-atomic minimax exponent.
\item \textbf{Complete SD strictness.} The finite moment hierarchy
provides new sufficient certificates in the remaining negative-cut
region. We have not proved that it must terminate successfully for
every admissible pair of parameters.
\end{enumerate}

These are different missing inputs. In particular the arithmetic and
spectral results cannot be composed into a finite-height prime recovery
theorem without a new quantitative link between the observation models.
The matrix-tail construction cannot establish a positive identification
theorem in a class where it already furnishes indistinguishable laws.
The MLE result is a statistical order theorem, not an assertion about
unlimited computational accuracy or every possible optimization routine.

\subsection{Research claims suitable for separate external scrutiny}

The arithmetic endpoint proof and the complete MLE proof can be read
independently of the longer research program. The first uses an analytic
lift and classical zeta second moments; the second includes its
approximation, entropy, transport, and matching lower-bound arguments.
These are the first two branches for which outside experts could assess
both correctness and novelty without accepting earlier unpublished
theorems. The matrix converse supplies a separate construction whose
crucial contour step is reduced to one complex trace coordinate.

The remaining branches supply exact model classifications and sharpened
boundaries. Their potential value depends on the observation model and
on how they compare with the full literature, not on the number of
theorem labels. No result in this manuscript proves the Riemann
hypothesis, a general Euler-phase recovery claim from a finite zero set,
or universal SD strictness. The appropriate next mathematical step is
independent expert review of the strongest complete arguments and their
specific prior-art claims.
\endgroup
% ===== END SOURCE 06 =====


% ===== BEGIN SOURCE 07: KEIO_Exact_Recovery_Laws_and_Coefficient_Conjecture.tex =====
\clearpage
\part{Exact Recovery Laws and a Coefficient Conjecture}
\label{part:07}
\setcounter{section}{0}
\begin{center}\large\itshape Endpoint Zero Densities, Robust Populations, and Structural Obstructions\end{center}
\begin{center}\small Source manuscript: 16 October 2026.\end{center}
\begingroup
\renewcommand{\arraystretch}{1}
\begin{quote}\small\textbf{Source abstract.}\quad
We continue a program on observation and recovery with sixteen precisely
scoped targets. A Hahn-polynomial inequality, an optimized complex
contour, and an exact rational finite certificate prove the constant-one
Chebyshev--Bernstein coefficient conjecture posed by Vinayak, Kong,
Valiant, and Kakade. For binomial populations we determine the full
total-variation inverse modulus and the minimax curve under an unknown
Huber contamination proportion. For the actual prime zeta function and
each fixed derivative, we strengthen linear-order counting in
$\Re s>1$ to $c_mT+o(T)$, control limiting boundary mass, and identify
every fixed-window counting moment. A definable-game argument proves
that fixed-complexity adaptive information cannot compress intrinsic
dimension; a separate geometric formula identifies convex-query noise
floors through the Jung constant and gives an exact Euclidean formula
at every noise level. Further classifications concern
arbitrary sparse scalar Gaussian designs, fixed-shape stable mixtures,
nonlinear spectral order detection, finite-height zero observations,
gamma deconvolution and McKay self-decomposability, a critical HCM power
cost, and complete replication-cap curves for analytic populations.
Dependencies on earlier unpublished manuscripts and classical results
are explicit. The sixteen targets are not sixteen named historical
conjectures: they include a published coefficient conjecture, internal
gap closures, and newly formulated transfer problems. Proofs and an
executable exact certificate support the research claims; external
refereeing and a complete priority assessment remain outstanding.
\end{quote}


\section{Scope, target ledger, and the change from the preceding work}
\label{p07:intro13:section}

The starting sources are the 223-page integrated manuscript
\emph{Observation and Recovery in Arithmetic and Probability} and its
subsequent papers, most recently the 48-page \emph{Boundary Laws for
Observation and Recovery} \cite{p01:KEIOIntegrated,comp:part:06}.
This paper supplies new continuations. Earlier linear-order prime-zeta
counting, clean binomial MLE optimality, uncapped analytic budget rates,
and signed Bessel kernel identities are not counted again as new
achievements. These earlier manuscripts are unpublished inputs, not
external confirmations of their claims.

The most directly traceable external target here is
\[
 T_m(2x-1)=\sum_{j=0}^tC(t,m,j)\binom tjx^j(1-x)^{t-j},\qquad
                      |C(t,m,j)|\le e^{m^2/t}.
\]
The exact inequality is explicitly conjectured in the 2019 source
\cite[Supplement, Remark 4.1, equation (23)]{p01:info:VKVK2019} and is still
discussed as a conjecture in the May 2026 version of
\cite[Section 6.2]{p01:info:MMRS2025}. We prove it by joining a classical Hahn
bound to a sharper contour exponent. An analytic overlap covers
$t\ge512$; a short verifier proves all remaining cases using rational
arithmetic. It is included verbatim in the single source file.
The historical Hahn bound is credited to Dette, not presented as our
new inequality. The searches performed for this draft do not constitute
an exhaustive priority determination.

The other main changes have distinct meanings. Exact prime-zeta
density replaces a pair of order bounds by a limiting law and local
moment convergence. Robust statistics asks what the inverse experiment
can retain when the reported counts are contaminated, even when the
contamination rate is unknown. The definable-information theorem shows
that the complexity of allowed queries is essential to dimension
compression. The McKay argument identifies an explicit signed gamma
remainder, whose forced negative interval closes a strictness gap left
by finite-moment tests. These are different conclusions linked by
changes of representation, not consequences of one universal theorem.

\subsection{Sixteen targets, with dependent conclusions grouped}

In the table, ``published'' refers to an explicitly located conjecture;
``gap'' means a question left by our earlier manuscripts; and
``transfer'' means a new problem formulated here. These labels are not
claims of publication priority. All entries are research theorems with
proofs below, subject to their stated hypotheses and inherited inputs.

\small
\begin{longtable}{@{}>{\raggedright\arraybackslash}p{.045\textwidth}>{\raggedright\arraybackslash}p{.18\textwidth}>{\raggedright\arraybackslash}p{.49\textwidth}>{\raggedright\arraybackslash}p{.19\textwidth}@{}}
\toprule
No. & Status of target & Result and exact scope & Location\\
\midrule\endfirsthead
\toprule No. & Status of target & Result and exact scope & Location\\
\midrule\endhead
1 & Published & Constant-one Chebyshev--Bernstein coefficient bound for every $0\le m\le t$ and every coefficient. & Thm.~\ref{p07:ml13:constantone}\\
2 & Transfer & Complete TV inverse modulus for arbitrary binomial mixing laws. & Thm.~\ref{p07:ml13:tvmodulus}\\
3 & Transfer & Full Huber minimax curve, attained without knowing the contamination proportion. & Thm.~\ref{p07:ml13:robust}\\
4 & Gap / strengthening & Exact endpoint density $N_m(T)=c_mT+o(T)$ for actual prime zeta and its fixed derivatives; linear limiting boundary-strip control. & Thms.~\ref{p07:ar13:thm:endpoint}, \ref{p07:ar13:thm:strip}\\
5 & Transfer & All positive moments of fixed-window arithmetic zero counts converge to random-prime moments. & Thm.~\ref{p07:ar13:thm:endpoint}\\
6 & Transfer & Fixed definable query families cannot recover below intrinsic dimension; fixed-degree polynomials need exactly $m$ queries on an $m$-ball. & Thm.~\ref{p07:in13:definable-obstruction}\\
7 & Transfer & Exact convex-body noise-radius formula, Jung coefficient, and complete finite-noise Euclidean formula. & Thms.~\ref{p07:in13:radius-formula}, \ref{p07:in13:Euclidean}\\
8 & Gap & Scalar Gaussian identification classified jointly by every terminal design and finite iterated-log tail class. & Thm.~\ref{p07:ga13:thm:scalar}\\
9 & Gap & Fixed non-Gaussian stable shape $\alpha\ne1$: necessity as well as sufficiency of $\sum n^{-|1-\alpha|}=\infty$. & Thm.~\ref{p07:ga13:thm:stable}\\
10 & Gap & Fixed skew index-one stable shape: exact criterion $\sum(1+\log^2n)^{-1}=\infty$. & Thm.~\ref{p07:ga13:thm:stable}\\
11 & Transfer & Sharp weak-mass/separation law for nonlinear spectral model-order detection. & Thm.~\ref{p07:sp13:thm:order}\\
12 & Gap / transfer & Quantitative finite-height zero-to-threshold bridge under an explicit Diophantine recurrence condition. & Thm.~\ref{p07:sp13:thm:finiteheight}\\
13 & Gap / transfer & Universal strict SD boundary for excess-power McKay laws, via an exact gamma-factor obstruction. & Thms.~\ref{p07:be13:gamma-factor}, \ref{p07:be13:strictness}\\
14 & Gap & Exact logarithmic critical cost of a tilt-uniform HCM power above Bessel order $1/2$. & Thm.~\ref{p07:be13:uniform-power}\\
15 & Transfer & Full adaptive replication-cap curve for stretched-exponential Fourier envelopes. & Thm.~\ref{p07:bu13:cap}\\
16 & Transfer & Full adaptive replication-cap curve for double-exponential Fourier envelopes. & Thm.~\ref{p07:bu13:cap}\\
\bottomrule
\end{longtable}
\normalsize

Items 4--5 use one common endpoint method; items 9--10 concern different
stable sampling geometries; items 15--16 specialize one cap theorem to
two regularity classes. We do not count every derivative, every parameter
value, the robust budget corollary, algebraicity of a fixed-degree error,
the Jung constant, or gamma extraction as extra independent closures.

\subsection{The six principal proof transfers}

\begin{enumerate}
\item A coefficient becomes a discrete orthogonal polynomial. Two
      estimates that are individually insufficient cover complementary
      degree ranges, leaving a finite exact certificate.
\item A single inverse norm becomes two predictive norms: inexpensive
      transportation controls low frequencies, while total variation
      controls a small-amplitude high-frequency remainder.
\item Arithmetic vertical translates become random analytic functions
      only after a joint weak-convergence argument; logarithmic Jensen
      bounds supply the uniform integrability needed for zero moments.
\item An unrestricted adaptive rule becomes a definable finite game.
      Definable choice exposes a dimension obstruction that is invisible
      if the decision rules are considered one at a time.
\item A putative positive convolution factor becomes an explicit signed
      inverse Laplace transform. A differential identity forces a
      negative interval, excluding every nonnegative remainder.
\item A replication cap becomes a conditional information constraint in
      a stronger Poisson experiment. Boundary-vanishing perturbations
      keep that constraint uniform even after exceptional histories.
\end{enumerate}

\paragraph{Conventions and verification status.}
Each section defines its own experiment and symbols; no unstated
identification of observation models is intended. Constants in
$\asymp$ may depend on explicitly fixed class parameters.
The finite-certificate code was executed using exact rational arithmetic.
Separate automated proof readings examined each section, including the
endpoint uniform-integrability argument, the definable game, the signed
gamma inverse, and the adaptive information bound. Such readings are
useful error checks, not formal proof verification or independent human
refereeing. No claim about acceptance by a journal, the Riemann hypothesis,
or a general solution of famous conjectures follows from this manuscript.


% ===== mle =====
\section{The constant-one coefficient conjecture and robust population recovery}
\label{p07:ml13:sec}

This section has two independent contributions. First, a classical Hahn
polynomial inequality, combined with a contour estimate and a small exact
certificate, proves the coefficient conjecture in Remark~4.1 of the
supplement to Vinayak--Kong--Valiant--Kakade \cite{p01:info:VKVK2019}.
Second, the two-scale inverse method from the preceding manuscript yields
the complete contamination curve for unrestricted binomial populations,
with an estimator that does not know the contamination proportion.
The all-replication optimality of the uncontaminated MLE is a prior result
\cite{comp:part:06}, not a new result of this section.

\subsection{A constant-one Chebyshev--Bernstein bound}

Write $B_j^t(x)=\binom tjx^j(1-x)^{t-j}$ and define $C(t,m,j)$ by
\[
 T_m(2x-1)=\sum_{j=0}^t C(t,m,j)B_j^t(x),\qquad 0\le m\le t.
\]

\begin{theorem}[The coefficient conjecture]
\label{p07:ml13:constantone}
For all integers $t\ge1$, $0\le m\le t$, and $0\le j\le t$,
\begin{equation}\label{p07:ml13:conjecture}
 |C(t,m,j)|\le \exp(m^2/t).
\end{equation}
The proof below is a finite-certificate proof: its only computational part
uses exact rational arithmetic for $t<512$. The complete verifier is
included in the text.
\end{theorem}

This is the precise constant-one assertion posed in
\cite[Supplement, Remark~4.1, equation~(23)]{p01:info:VKVK2019}, rather than a
bound with an unspecified multiplicative constant. The distinction was
retained in the subsequent discussion \cite[Section~6.2]{p01:info:MMRS2025} and
in our preceding manuscript. No optimality claim is made for the exponent
constant one.

\begin{lemma}[Three complementary coefficient bounds]
\label{p07:ml13:threebounds}
For $1\le m\le t$,
\begin{align}
 |C(t,m,j)|&\le 2^m,\label{p07:ml13:degreebound}\\
 |C(t,m,j)|&\le (t+1)\exp\{tH(m/t)\}
       \le(t+1)\exp\{(\log2)m^2/t\},\label{p07:ml13:contourbound}
\end{align}
where $H(a)=\frac12[(1+a)\log(1+a)+(1-a)\log(1-a)]$ and $0\log0=0$.
For $m<t$ there is also the bound
\begin{equation}\label{p07:ml13:dettebound}
 |C(t,m,j)|\le D(t,m):=
          \prod_{l=1}^{m-1}\frac{t+l}{t-l}.
\end{equation}
In particular, $D(t,m)\le e^{m^2/t}$ whenever
\begin{equation}\label{p07:ml13:integercriterion}
 (m+6)(m-1)^2\le6t^2.
\end{equation}
\end{lemma}
\begin{proof}
At degree $m$ the coefficients are
$(-1)^{m-i}\binom{2m}{2i}/\binom mi$. The successive absolute-value ratio
is $(2m-2i-1)/(2i+1)$, so their maximum occurs at a middle index.
Vandermonde's identity bounds the numerator at that index by
$2^m\max_i\binom mi$, proving \eqref{p07:ml13:degreebound} at degree $m$.
Degree elevation replaces coefficients by convex combinations and proves
the bound at every $t\ge m$.

For \eqref{p07:ml13:contourbound}, put
\[
 F(z)=(1+z)^tT_m\left(\frac{z-1}{z+1}\right)
     =\sum_{j=0}^t\binom tj C(t,m,j)z^j.
\]
Endpoint coefficients have absolute value one. For $0<j<t$, set
$r=j/(t-j)$ and, on $z=re^{i\theta}$, set
$u=2\sqrt r\,|\sin(\theta/2)|/(1+r)$. The larger inverse Joukowski
modulus of $(z-1)/(z+1)$ is
$\sqrt{(1+u)/(1-u)}$. Consequently
\[
 \frac{|F(re^{i\theta})|}{(1+r)^t}
 \le(1-u^2)^{t/2}
          \left(\frac{1+u}{1-u}\right)^{m/2}
 \le \exp\{tH(m/t)\}.
\]
The last maximum is attained at $u=m/t$; limiting values handle $u=1$.
Cauchy's formula and
$\binom tj r^j/(1+r)^t\ge1/(t+1)$ give the first inequality in
\eqref{p07:ml13:contourbound}. This probability is a mode of a
$\operatorname{Binomial}(t,j/t)$ law, which proves the stated lower bound.
Finally
\[
 H(a)=\sum_{k\ge1}\frac{a^{2k}}{2k(2k-1)}
       \le a^2\sum_{k\ge1}\frac1{2k(2k-1)}=a^2\log2.
\]

To apply the older bound \eqref{p07:ml13:dettebound}, combine the
hypergeometric formula with the identity
\[
 x^r=\sum_j(j)_{\underline r}B_j^t(x)/(t)_{\underline r}.
\]
This gives
\[
 C(t,m,j)=(-1)^m\,{}_3F_2
       \left(\begin{matrix}-m,m,-j\\1/2,-t\end{matrix};1\right)
       =(-1)^mQ_m(j;-1/2,-1/2,t).
\]
Dette's Hahn bound \cite[Corollary~3.6(b)]{p07:ml13:Dette} is exactly
\eqref{p07:ml13:dettebound}. Its stipulated range $m<t$ is respected here.
Thus this inequality itself is prior art. For the final assertion,
using the positive series for $\operatorname{arctanh}$ gives
\begin{align*}
 \log D(t,m)
 &=2\sum_{l=1}^{m-1}\operatorname{arctanh}(l/t)\\
 &\le \frac{m(m-1)}t+
 \frac{m^2(m-1)^2}
 {6t^3\{1-(m-1)^2/t^2\}}.
\end{align*}
The second term is at most $m/t$ precisely under
\eqref{p07:ml13:integercriterion}.
\end{proof}

\begin{proof}[Proof of Theorem~\ref{p07:ml13:constantone}]
The case $m=0$ is immediate. Suppose first $t\ge512$.
If $m^2/t\ge\frac{10}{3}\log(t+1)$, the contour bound proves the
assertion because $\log2<7/10$. Otherwise,
\[
 m^2<\frac{10}{3}t\log(t+1)\le\frac{25}{9}t^{4/3},
 \qquad m<\frac53t^{2/3}.
\]
Here $\log(t+1)\le(5/6)t^{1/3}$ for $t\ge512$: at $512$ it follows
from $513<729$, $\log3<11/10$, and $33/5<20/3$; the difference
$(5/6)t^{1/3}-\log(t+1)$ has positive derivative thereafter.
In particular $m<t$, and
\[
 (m+6)(m-1)^2\le m^3+6m^2
 \le\left(\frac{125}{27}+\frac{25}{96}\right)t^2<6t^2.
\]
The last estimate uses $t^{2/3}\ge64$. The Hahn bound therefore applies.
The elementary logarithmic inequalities just used can, for example, be
verified by positive Taylor lower bounds for $e^{7/10}$ and $e^{11/10}$.

For $1\le t<512$, the following verifier covers every pair $(t,m)$.
Its first branch uses \eqref{p07:ml13:integercriterion}; its second uses
$2^m\le e^{7m/10}\le e^{m^2/t}$; its third uses the contour bound with
$\log2<7/10$. Every other pair is checked with the exact identity
\begin{equation}\label{p07:ml13:exactcoeff}
 C(t,m,j)=\frac1{\binom tj}
 \sum_{i=\max(0,j-t+m)}^{\min(m,j)}
 (-1)^{m-i}\binom{2m}{2i}\binom{t-m}{j-i}.
\end{equation}
The function \texttt{exp\_lower} returns a rational number no larger
than $e^x$ for $x\ge0$, so each successful assertion is a rigorous
inequality in the required direction. There are no floating-point
comparisons. The complete output is
\[
 (34784,39518,56329,185),\qquad
 12734\ \hbox{direct coefficient comparisons},\quad
 \max t_{\rm direct}=145.
\]
The four counts sum to $511\cdot512/2$, as required.
\end{proof}

\noindent\begin{minipage}{\textwidth}\small
\begin{verbatim}
from fractions import Fraction
from math import comb

def exp_lower(x, degree=32):
    term = total = Fraction(1)
    for k in range(1, degree + 1):
        term *= x / k
        total += term
    return total

counts = [0, 0, 0, 0]
comparisons = largest_direct = 0
for t in range(1, 512):
    for m in range(1, t + 1):
        if m < t and (m+6)*(m-1)**2 <= 6*t*t:
            counts[0] += 1
            continue
        if 10*m >= 7*t:
            counts[1] += 1
            continue
        if exp_lower(Fraction(3*m*m, 10*t)) >= t+1:
            counts[2] += 1
            continue
        lower = exp_lower(Fraction(m*m, t))
        for j in range(t + 1):
            numerator = sum(
                (-1)**(m-i)*comb(2*m, 2*i)*comb(t-m, j-i)
                for i in range(max(0,j-t+m), min(m,j)+1)
            )
            assert Fraction(abs(numerator), comb(t,j)) <= lower
            comparisons += 1
        counts[3] += 1
        largest_direct = max(largest_direct, t)
print(counts, comparisons, largest_direct)
assert sum(counts) == 511*512//2
\end{verbatim}
\end{minipage}

\subsection{A sharp total-variation inverse modulus}

For $G\in\mathcal P([0,1])$, write $p_G(j)=\int B_j^t\,dG$.
Total variation is $\|p-q\|_{\rm TV}=\frac12\sum_j|p_j-q_j|$;
$W_{1,t}$ uses cost $|j-k|/t$ on the count alphabet. Define
\[
 \omega_t^{\rm TV}(\eta)=\sup\{W_1(G,H):
                 \|p_G-p_H\|_{\rm TV}\le\eta\}.
\]
Throughout, $\log(e/0)=+\infty$ and $1/\sqrt{+\infty}=0$.

\begin{theorem}[The complete TV inverse curve]
\label{p07:ml13:tvmodulus}
Uniformly for integers $t\ge1$ and $0\le\eta\le1/4$,
\begin{equation}\label{p07:ml13:tvcurve}
 \omega_t^{\rm TV}(\eta)\asymp
 \min\left\{1,\frac1t+
       \frac1{\sqrt{t\log(e/\eta)}}+\eta\right\}.
\end{equation}
All constants are absolute. In particular, infinite sample size does not
remove an arbitrary predictive contamination of size $\eta$: its inverse
cost is generally larger than $\eta$.
\end{theorem}

We use the following prior approximation ingredients, recording their
proofs to specify the uniform constants. They do not require the
constant-one theorem above.

\begin{lemma}[Two-scale inversion in two predictive metrics]
\label{p07:ml13:twoscale}
For $1\le s\le k\le t/2$,
\begin{equation}\label{p07:ml13:twoscaleineq}
 W_1(G,H)\le\frac Ck+
 Ce^{4s^2/t}W_{1,t}(p_G,p_H)+
 \frac{Ce^{2k^2/t}}s\|p_G-p_H\|_{\rm TV}.
\end{equation}
\end{lemma}
\begin{proof}
For every $1$-Lipschitz function $f$, there is a degree-$k$ polynomial
$P_k$ such that $\|P_k-f\|_\infty\le C/k$ and
$\|P_k'\|_\infty\le2$. One elementary construction reflects $f$ to a
period-two Lipschitz function, convolves it with a positive Jackson kernel
of order $a$, and replaces its $O(a)$ Fourier exponentials by their
Taylor polynomials of degree $100a$. The Jackson first-moment bound is
$O(1/a)$, convolution does not increase the Lipschitz constant, and the
Taylor error in both the function and its derivative is exponentially
small. Taking $a=\lfloor k/100\rfloor$ and treating bounded $k$ by a
constant polynomial gives the assertion.

We also need a bound for arbitrary bounded polynomials, not merely a
single Chebyshev polynomial. If $P=\sum b_jB_j^t$ has degree at most
$k\le t/2$, then
\begin{equation}\label{p07:ml13:generalcoeff}
 \max_j|b_j|\le Ce^{2k^2/t}\|P\|_\infty.
\end{equation}
For completeness, Bernstein's complex polynomial bound gives
$|P((2+w+w^{-1})/4)|\le\|P\|_\infty|w|^k$ for $|w|\ge1$.
Use the same coefficient contour as above, with
$v=r/(1+r)^2$ and $u=2\sqrt v|\sin(\theta/2)|$. Its normalized
integrand is at most
\[
 (1-u^2)^{(t-k)/2}(1+u)^k
 \le e^{2k^2/t}e^{-tu^2/8}.
\]
The average of the last factor is $O((1+tv)^{-1/2})$, whereas Stirling's
factorial inequalities give
$\binom tj(j/t)^j(1-j/t)^{t-j}\ge c(1+tv)^{-1/2}$.
This proves \eqref{p07:ml13:generalcoeff}, including endpoint coefficients
by direct evaluation. Differentiating the Bernstein representation shows
that the coefficient vector of $P_s$ has normalized-grid Lipschitz
constant at most $Ce^{4s^2/t}$. The polynomial $P_k-P_s$ has supremum
norm $O(1/s)$, so its coefficients are bounded by
$Ce^{2k^2/t}/s$. Test the first part by $W_{1,t}$, the second by TV,
and take the Kantorovich dual supremum over $f$.
\end{proof}

\begin{proof}[Proof of Theorem~\ref{p07:ml13:tvmodulus}]
For the upper bound, $W_{1,t}\le\|\cdot\|_{\rm TV}$ because the
normalized alphabet has diameter one. For sufficiently large $t$, take
\[
 k=\left\lfloor\min\{t/2,\sqrt{at\log(e/\eta)}\}\right\rfloor,
 \qquad s=\lfloor\sqrt t/16\rfloor,
\]
where $a=1/64$; at $\eta=0$ take $k=\lfloor t/2\rfloor$.
Bounded $t$ are absorbed by the loss bound one. These choices have
$1\le s\le k$, $s\asymp\sqrt t$, and $e^{4s^2/t}\le C$.
For $\eta>0$ the last term of \eqref{p07:ml13:twoscaleineq} is at most
\[
 C\eta e^{2a\log(e/\eta)}/\sqrt t
 \le C/\sqrt{t\log(e/\eta)}.
\]
The approximation term is bounded by the other two inverse terms in
\eqref{p07:ml13:tvcurve}. At $\eta=0$ both predictive terms vanish.

For the lower bounds, endpoint mixing laws differing by mass $\eta$
have both population $W_1$ and predictive TV equal to $\eta$.
For the inverse term let
\[
 \pi(dx)=\frac{dx}{\pi\sqrt{x(1-x)}},\qquad
 h_k(x)=T_k(2x-1),\qquad G_k^\pm=(1\pm h_k/2)\pi.
\]
Direct integration of their signed distribution function gives
$W_1(G_k^+,G_k^-)\asymp1/k$, uniformly in $k\ge1$.
For $k=t+1$ the predictive laws coincide, giving the $1/t$ term.
For $k\le t$, the beta--binomial conditional-expectation operator
$Qh(j)=\mathbb E_\pi[h(X)\mid J=j]$ satisfies
\[
 Q^*Qh_k=\lambda_{t,k}h_k,
 \qquad\lambda_{t,k}=\prod_{l=0}^{k-1}\frac{t-l}{t+1+l}
                  \le e^{-k^2/(2t)}.
\]
Indeed it is self-adjoint, preserves polynomial degree, and has the
displayed leading coefficient; orthogonality then proves the identity.
This classical diagonalization is discussed in \cite{p05:fm:DKS2008}.
Consequently
\[
 \|p_{G_k^+}-p_{G_k^-}\|_{\rm TV}
 =\tfrac12\mathbb E|Qh_k(J)|
 \le\tfrac1{2\sqrt2}\sqrt{\lambda_{t,k}}
 \le\tfrac1{2\sqrt2}e^{-k^2/(4t)}.
\]
Take $k=\lceil\sqrt{8t\log(e/\eta)}\rceil$ if it is at most $t$.
This gives the required inverse lower bound. If that degree exceeds
$t$, then $1/\sqrt{t\log(e/\eta)}\le C/t$, and the exactly
indistinguishable pair suffices. A maximum of the three lower bounds
is comparable to their clipped sum.
\end{proof}

\subsection{Unknown Huber contamination: the full minimax curve}

We observe $N$ independent counts with common law
\[
 R=(1-\varepsilon)p_G+\varepsilon Q,
 \qquad G\in\mathcal P([0,1]),\quad Q\in\mathcal P(\{0,\ldots,t\}),
 \qquad0\le\varepsilon\le1/4.
\]
Contamination occurs at the level of whole individuals' reported counts,
not independently within each individual's $t$ trials. Write
\[
 \mathcal R_{N,t}^{\rm Hub}(\varepsilon)
 =\inf_{\widehat G}\sup_{G,Q}
             \mathbb E_R W_1(\widehat G,G).
\]

\begin{theorem}[Adaptive robust population estimation]
\label{p07:ml13:robust}
There are absolute positive constants such that, for every $N,t\ge1$
and $0\le\varepsilon\le1/4$,
\begin{equation}\label{p07:ml13:robustrate}
 \mathcal R_{N,t}^{\rm Hub}(\varepsilon)\asymp
 \min\left\{1,\frac1t+
 \frac1{\sqrt{t\log(eN)}}+\frac1{\sqrt N}
 +\varepsilon+\frac1{\sqrt{t\log(e/\varepsilon)}}\right\}.
\end{equation}
One sequence of estimators, depending on $N,t$ but not on
$\varepsilon$, attains the upper bound simultaneously for all
$\varepsilon$ in this interval. The estimator is a finite-model
minimum-distance estimator; no robustness claim for the ordinary MLE
is made.
\end{theorem}

General contamination theory identifies TV moduli as the relevant lower
bound mechanism \cite{p07:ml13:CGR}. What requires a separate argument here
is attaining both the sampling transition and the contamination
transition without allowing the predictive TV entropy cost to spoil
the parametric low-frequency term.

\begin{lemma}[Simultaneous predictive control without knowing contamination]
\label{p07:ml13:jointfit}
Let a family of laws on an ordered alphabet of diameter one have a
finite TV cover $p_1,\ldots,p_M$ of radius $\beta$, with representatives
in the family. Let $q$ be the empirical law of $N$ independent draws
from $R$, and suppose $\|R-p\|_{\rm TV}\le\varepsilon$ for a member
$p$ of the family. There is an estimator $\widehat p$ in the cover,
not using $\varepsilon$, such that
\begin{align}
 \mathbb E W_1(\widehat p,p)&\le C(\varepsilon+\beta+N^{-1/2}),
       \label{p07:ml13:jointw}\\
 \mathbb E\|\widehat p-p\|_{\rm TV}
     &\le C\left(\varepsilon+\beta+
                       \sqrt{\log(2M)/N}\right).
       \label{p07:ml13:jointtv}
\end{align}
With probability at least $1-3e^{-x}$, the same two bounds hold with
$N^{-1/2}$ replaced by $\sqrt{(1+x)/N}$ and
$\sqrt{\log(2M)/N}$ by $\sqrt{(\log(2M)+x)/N}$, respectively.
\end{lemma}
\begin{proof}
Let $\mathcal A=\{\{j:p_b(j)>p_c(j)\}:1\le b,c\le M\}$,
$D_a=\max_{A\in\mathcal A}|p_a(A)-q(A)|$, and
\[
 b_N=\sqrt{\frac{\log(2M^2)}{2N}},\qquad
 F_a=W_1(p_a,q)+(D_a-b_N)_+.
\]
Choose the smallest index minimizing $F_a$. If
$S=\max_{A\in\mathcal A}|q(A)-R(A)|$, the union bound and Hoeffding's
inequality imply
\[
 \Pr\{S>b_N+u\}\le e^{-2Nu^2},\qquad
 \mathbb E(S-b_N)_+\le C/\sqrt N.
\]
Choose a true nearby center $p_{a_*}$. Then
$D_{a_*}\le\beta+\varepsilon+S$ and
$W_1(p_{a_*},q)\le\beta+\varepsilon+W_1(R,q)$.
Integration of empirical CDF standard deviations gives
$\mathbb EW_1(R,q)\le1/(2\sqrt N)$. Thus
$\mathbb EF_{a_*}\le2(\beta+\varepsilon)+C/\sqrt N$.
The minimizing index has $W_1(\widehat p,q)\le F_{a_*}$ and
$D_{\widehat a}\le b_N+F_{a_*}$.
The triangle inequality proves \eqref{p07:ml13:jointw}. Since the set
$\{\widehat p>p_{a_*}\}$ belongs to $\mathcal A$,
\[
 \|\widehat p-p_{a_*}\|_{\rm TV}\le
                  D_{\widehat a}+D_{a_*}.
\]
Together with $\mathbb ES\le b_N+C/\sqrt N$, this proves
\eqref{p07:ml13:jointtv}. The probability assertion follows from the same
inequalities, the displayed tail bound for $S$, and the
empirical-CDF bound
$\Pr\{W_1(R,q)>\sqrt{x/(2N)}\}\le2e^{-x}$.
\end{proof}

\begin{lemma}[Effective dimension of the binomial moment body]
\label{p07:ml13:entropy}
For $t\le N$ the set $\{p_G\}$ has a TV cover of radius
$\beta=(eN)^{-2}$, with centers in that set, satisfying
\[
 \log(2M)\le C\sqrt t\,\log(eN)^{3/2}.
\]
\end{lemma}
\begin{proof}
The vector polynomial $K_x=(B_0^t(x),\ldots,B_t^t(x))$ has shifted
Chebyshev coefficients $v_l$ satisfying
$\|v_l\|_1\le2e^{-l^2/(2t)}$ for $l\ge1$.
To see this, shift the Fourier contour of
$K_{(1+\cos z)/2}$ by $u=l/t$ and use
\[
 \|K_{(1+\cos(\theta+iu))/2}\|_1=(\cosh u)^t\le e^{tu^2/2}.
\]
Truncation at $m=\min(t,\lceil\sqrt{2t\log(16t/\beta)}\rceil)$
has uniform $\ell_1$ error $O(\beta)$. Partition the first $m$
Chebyshev moments, each in $[-1,1]$, into coordinate intervals of length
$c\beta/(m+1)$ and retain one mixing representative per occupied cell.
The resulting TV radius is at most $\beta$ after reducing the absolute
constant $c$, and
$\log M\le m\log(C(m+1)/\beta)$. The asserted estimate follows.
\end{proof}

\begin{proof}[Upper bound in Theorem~\ref{p07:ml13:robust}]
For $t\le N$ use the cover of Lemma~\ref{p07:ml13:entropy} and the estimator
of Lemma~\ref{p07:ml13:jointfit}; return its representative mixing law.
Writing $L_N=\log(eN)$ and $L_\varepsilon=\log(e/\varepsilon)$ gives
\[
 \mathbb E W_{1,t}(\widehat p,p_G)\le C(\varepsilon+N^{-1/2}),
 \quad
 \mathbb E\|\widehat p-p_G\|_{\rm TV}
 \le C\left(\varepsilon+
           \frac{t^{1/4}L_N^{3/4}}{\sqrt N}\right).
\]
In Lemma~\ref{p07:ml13:twoscale}, choose
\[
 k=\left\lfloor\min\{t/2,
       \sqrt{at\min(L_N,L_\varepsilon)}\}\right\rfloor,
 \qquad s=\lfloor\sqrt t/16\rfloor,\qquad a=1/64.
\]
As above, bounded $t,N$ are harmless. The approximation term has the
required three inverse contributions. The low-frequency term costs
$C(\varepsilon+N^{-1/2})$. The contamination part of the last term is
at most $C/\sqrt{tL_\varepsilon}$. Its sampling part is bounded by
\[
 \frac{Ce^{2aL_N}L_N^{3/4}}{\sqrt N\,t^{1/4}}
 \le\frac C{\sqrt{tL_N}},
\]
because $t\le N$ and
$(eN)^{2a}N^{-1/4}L_N^{5/4}$ is uniformly bounded.
The parameter $\varepsilon$ appears only in the proof's choice of
test degree, not in the estimator.

For $t>N$, return the empirical law of the normalized counts $J_i/t$
as the estimated mixing distribution. Coupling a binomial count to its
parameter gives $W_1(G,p_G)\le1/(2\sqrt t)$, and hence its expected
loss is at most $\varepsilon+1/(2\sqrt N)+1/(2\sqrt t)$.
This is $O(\varepsilon+N^{-1/2})$, as required in this range.
\end{proof}

\begin{proof}[Lower bound in Theorem~\ref{p07:ml13:robust}]
Taking $Q=p_G$ recovers the uncontaminated experiment. Its elementary
lower bound is
$c\min\{1,t^{-1}+(tL_N)^{-1/2}+N^{-1/2}\}$:
the exactly moment-matched pair above gives $t^{-1}$; the same
beta--binomial pair with $k=\lceil\sqrt{8t\log(8N)}\rceil\le t$
has bounded $N$-fold total variation and gives $(tL_N)^{-1/2}$;
endpoint mixtures with weights $1/2\pm c/\sqrt N$ give $N^{-1/2}$.
If the stated $k$ exceeds $t$, its contribution is absorbed by $t^{-1}$.

For the contamination lower bound, the neighborhoods
$\{(1-\varepsilon)p+\varepsilon Q\}$ and
$\{(1-\varepsilon)p'+\varepsilon Q'\}$ intersect whenever
$\|p-p'\|_{\rm TV}\le\varepsilon/(1-\varepsilon)$.
Indeed $(1-\varepsilon)\max(p,p')$ has total mass at most one;
add any nonnegative remaining mass to produce a common probability law.
Its differences from $(1-\varepsilon)p$ and
$(1-\varepsilon)p'$ are nonnegative and each has mass $\varepsilon$.
Every estimator then has maximum expected population loss at least
$W_1(G,G')/2$, regardless of $N$. Apply
Theorem~\ref{p07:ml13:tvmodulus} with predictive threshold
$\eta=\varepsilon$, which is no larger than
$\varepsilon/(1-\varepsilon)$. Combining the two lower bounds proves
the theorem. At $\varepsilon=0$ only the clean lower bound is needed.
\end{proof}

\begin{corollary}[Equal-replication budgets remain robust at their optimum]
\label{p07:ml13:robustbudget}
For common replication and total trial budget $Nt\le B$, the optimal
Huber-contaminated population risk is
\[
 \inf_{Nt\le B}\mathcal R_{N,t}^{\rm Hub}(\varepsilon)
       \asymp \min\{1,(B\log(eB))^{-1/4}+\varepsilon\}.
\]
The uncontaminated order-optimal allocation
$N\asymp\sqrt{B\log(eB)}$, $t\asymp\sqrt{B/\log(eB)}$
continues to attain this order with the adaptive robust estimator.
\end{corollary}
\begin{proof}
The clean budget lower bound and the endpoint contamination lower bound
give the assertion from below. Put $r=(B\log(eB))^{-1/4}$ and use the
displayed allocation. Its clean terms are $O(r)$. If
$\varepsilon\le r$, then
$\log(e/\varepsilon)\ge c\log(eB)$, so the remaining inverse term is
$O(r)$. If $r<\varepsilon\le1/4$, the function
$x^2\log(e/x)$ is increasing in this interval. Thus
$t\varepsilon^2\log(e/\varepsilon)\ge
tr^2\log(e/r)\ge c$, and the inverse term is $O(\varepsilon)$.
Integer rounding and bounded budgets only change absolute constants.
\end{proof}

The coefficient conjecture is a published target. The TV curve and the
adaptive contamination curve are newly formulated and resolved transfer
problems in this manuscript; they are not presented as previously named
conjectures. The budget statement is a corollary, not an additional
independent conjecture closure. The exact best exponent constant in
\eqref{p07:ml13:conjecture}, computationally efficient robust fitting, and
sharp confidence-radius constants remain separate questions.

% ===== arithmetic =====
\section{Exact endpoint densities and local moments for prime-zeta zeros}
\label{p07:ar13:sec:endpoint}

The endpoint theorem in the preceding manuscript established
$N_m(T)\asymp T$ for each fixed derivative of the prime zeta function.
It deliberately did not identify a limiting coefficient. We now prove
that the coefficient exists, that no positive density of zeros can
escape into a shrinking strip adjacent to $\Re s=1$, and that every
fixed-height local counting moment has a random-prime limit. These
are endpoint extensions of the earlier result, not repetitions of its
linear upper bound. The published questions motivating this direction
are those of Belovas, \v{C}epaityt\.{e}, and Sabaliauskas
\cite{p01:pf:BCS}, also discussed in \cite{p06:ra:Tuzcu}. We do not
attribute the stronger moment statement below to a published conjecture.

Write $P(s)=\sum_p p^{-s}$ for $\Re s>1$. Let
$(\omega_p)_p$ be independent random variables uniform on the unit
circle, and put
\begin{equation}\label{p07:ar13:eq:random-prime}
 P_\omega(s)=\sum_p\omega_p p^{-s},\qquad
 Q_{m,\omega}(s)=P_\omega^{(m)}(s).
\end{equation}
These random series converge locally uniformly almost surely on
$\Re s>1/2$; a proof is included below. Zeros are counted with
multiplicity unless a superscript $\mathrm{dist}$ is displayed.
For $h>0$ define
\begin{align*}
 X_{m,h}(t)&=\#\{s:P^{(m)}(s)=0,\ \Re s>1,\ t<\Im s<t+h\},\\
 X_{m,h}^{\omega}&=\#\{s:Q_{m,\omega}(s)=0,\ \Re s>1,
                                     \ 0<\Im s<h\}.
\end{align*}
Changing an open horizontal edge to a closed edge affects only a
null set of starting heights in any integral below.

\begin{theorem}[Endpoint density and every fixed-window moment]
\label{p07:ar13:thm:endpoint}
For every fixed integer $m\ge0$ there is a number
$0<c_m<\infty$ such that
\begin{equation}\label{p07:ar13:eq:asymptotic}
 N_m(T)=c_mT+o(T),\qquad
 N_m^{\mathrm{dist}}(T)=c_mT+o(T).
\end{equation}
Here both counts include all zeros in $\Re s>1$, $0<\Im s<T$.
For every fixed $h>0$, as $t$ is sampled uniformly from $[0,T]$,
\begin{equation}\label{p07:ar13:eq:local-law}
 X_{m,h}(t)\ \Longrightarrow\ X_{m,h}^{\omega}.
\end{equation}
Moreover, for every real $q>0$,
\begin{equation}\label{p07:ar13:eq:all-moments}
 \lim_{T\to\infty}\frac1T\int_0^T X_{m,h}(t)^q\,dt
       =\mathbb E[(X_{m,h}^{\omega})^q]<\infty.
\end{equation}
The same statements hold for distinct-zero counts. In particular,
$c_m=\mathbb E X_{m,1}^{\omega}$ and the limiting mean and variance
in a fixed window are $c_mh$ and
$\operatorname{Var}(X_{m,h}^{\omega})$, respectively.
\end{theorem}

\begin{theorem}[Quantitative absence of endpoint mass]
\label{p07:ar13:thm:strip}
Fix $m\ge0$. There are $\epsilon_0>0$ and $C_m<\infty$ such that
for $0<\epsilon<\epsilon_0$,
\begin{equation}\label{p07:ar13:eq:strip}
 \limsup_{T\to\infty}\frac1T
 \#\{s:P^{(m)}(s)=0,\ 1<\Re s\le1+\epsilon,
                                      \ 0<\Im s<T\}
       \le C_m\epsilon.
\end{equation}
If $c_m(\gamma)$ denotes the limiting density in $\Re s>\gamma>1$,
then
\[
 0\le c_m-c_m(1+\epsilon)\le C_m\epsilon.
\]
Thus the endpoint coefficient is a controlled one-sided limit of
interior coefficients. No rate in $T$ is asserted in
\eqref{p07:ar13:eq:strip}.
\end{theorem}

\subsection{The random analytic limit of the single-valued lifts}

Set
\[
 H(s)=\sum_p\sum_{k\ge2}\frac{p^{-ks}}k,\qquad
 H_\omega(s)=\sum_p\sum_{k\ge2}\frac{\omega_p^kp^{-ks}}k.
\]
Both are holomorphic and locally bounded, with every derivative,
on $\Re s>1/2$; the bounds for $H_\omega$ are uniform in $\omega$.
Let
\[
 Z_\omega(s)=\exp(P_\omega(s)+H_\omega(s)).
\]
It is almost surely holomorphic and nowhere zero on that half-plane.

\begin{lemma}[Joint weak convergence, with its approximation input]
\label{p07:ar13:lem:weak}
On every fixed bounded open rectangle whose closure lies in
$\Re s>1/2$, the joint laws of
\[
 (\zeta(s+it),H(s+it))
\]
for uniform $t\in[t_0,T]$ converge, in the topology of locally
uniform convergence, to $(Z_\omega(s),H_\omega(s))$.
Here $t_0$ is chosen so that the moving pole of $\zeta$ lies outside
the rectangle. The same assertion holds jointly with any fixed
number of derivatives.
\end{lemma}

\begin{proof}
We give the approximation argument, rather than assuming that a
random model automatically represents arithmetic averages.
The classical off-critical-line mean-square theorem and its
fixed-coefficient Fourier formula imply, for $\sigma>1/2$,
\begin{equation}\label{p07:ar13:eq:ms-approx}
 \lim_{T\to\infty}\frac1T\int_{t_0}^T
 \left|\zeta(\sigma+it)-\sum_{n\le M}n^{-\sigma-it}\right|^2dt
       =\sum_{n>M}n^{-2\sigma}.
\end{equation}
The estimates needed to integrate this identity over a compact
real-part interval are locally uniform. These are standard
consequences of the zeta mean-square formula
\cite[Chapter VII]{p06:ra:Titchmarsh}. To specify the Fourier input,
$T^{-1}\int_{t_0}^T\zeta(\sigma+it)n^{it}dt\to n^{-\sigma}$
for each fixed $n$; expanding the square together with
$T^{-1}\int|\zeta(\sigma+it)|^2dt\to\zeta(2\sigma)$ gives
\eqref{p07:ar13:eq:ms-approx}.
One can obtain that fixed Fourier formula by moving the vertical
integral to a line $\Re s>1$, where the Dirichlet series is absolute;
the horizontal connecting integrals are $o(T)$ by the classical
polynomial growth bound for $\zeta$ on a fixed strip in $\Re s>1/2$.

For compact rectangles $K\Subset U$ the submean inequality for
holomorphic functions gives
\[
 \sup_K|f|^2\le C_{K,U}\int_U|f(s)|^2\,dA(s).
\]
Apply this to the difference in \eqref{p07:ar13:eq:ms-approx}, then
average over $t$, integrate over $U$, and let $M\to\infty$.
It follows that the Dirichlet polynomials approximate the translated
zeta functions in mean square for the compact-uniform topology.
The finitely many initial heights have no effect.

For clarity, the random analytic objects require no conjectural
Euler product convergence. On a bounded rectangle in $\Re s>1/2$,
\[
 \sum_p\mathbb E\|\omega_pp^{-s}\|_{L^2(U)}^2<\infty.
\]
The independent centered Hilbert-space series therefore converges
almost surely and in mean square in $L^2(U)$. The submean inequality
on smaller rectangles promotes it to locally uniform convergence.
A countable exhaustion proves the assertion about $P_\omega$.
Also, writing $\omega(n)=\prod_{p^a\Vert n}\omega_p^a$, orthogonality
$\mathbb E\omega(n)\overline{\omega(k)}=\mathbf1_{n=k}$ gives
\[
 \mathbb E\left\|\sum_{M<n\le L}\omega(n)n^{-s}\right\|_{L^2(U)}^2
    =\int_U\sum_{M<n\le L}n^{-2\Re s}\,dA(s).
\]
Thus the random Dirichlet polynomials converge in mean square in
the compact-uniform topology. Their limit agrees with $Z_\omega$
on $\Re s>1$, by the absolutely convergent Euler product, and
hence on the whole connected half-plane by the identity theorem.

For each fixed truncation, the vertical phases
$(p^{-it})_{p\le M}$ are equidistributed on their finite torus:
no nonzero integer combination of distinct prime logarithms is
zero. Consequently the finite Dirichlet polynomials and finite-prime
truncations of $H$ converge jointly in law to their random versions.
The remaining tail of $H$ tends uniformly to zero on compact sets,
uniformly in the phases. The deterministic and random mean-square
approximations above now prove the joint weak convergence, first
for bounded Lipschitz functionals and then for the compact-uniform
metric. Derivatives are continuous maps on holomorphic function
spaces by Cauchy's formula. This proves the lemma. The argument is
the classical probabilistic limit-theorem method of
\cite{p07:ar13:Bagchi,p07:ar13:Laurincikas}, specialized here to the joint
prime-power correction that is needed for zero counts.
\end{proof}

Define $A_1=\zeta'$ and
$A_{m+1}=\zeta A_m'-m\zeta'A_m$. Set
\[
 F_0=\zeta-e^H,\qquad F_m=A_m-\zeta^mH^{(m)}\quad(m\ge1),
 \qquad r(s)=\frac{s-1}{s+1},
\]
and remove the sole possible pole by
\[
 L_0=rF_0,\qquad L_m=r^{2m}F_m\quad(m\ge1).
\]
Every $L_m$ is holomorphic on $\Re s>1/2$. On $\Re s>1$,
\begin{equation}\label{p07:ar13:eq:lift-identities}
 F_0=e^H(e^P-1),\qquad F_m=\zeta^mP^{(m)}\quad(m\ge1).
\end{equation}
Since $r(s+it)\to1$ uniformly on compact sets as $t\to\infty$,
Lemma~\ref{p07:ar13:lem:weak} and the differential-polynomial identities
give
\begin{equation}\label{p07:ar13:eq:lift-limit}
 L_m(s+it)\Longrightarrow V_{m,\omega}(s),\qquad
 V_{0,\omega}=e^{H_\omega}(e^{P_\omega}-1),\qquad
 V_{m,\omega}=Z_\omega^m Q_{m,\omega}\ (m\ge1).
\end{equation}
For $m=0$ this limit includes all levels $P_\omega=2\pi i k$,
not merely the level zero. We will use it only as an upper bound
near the boundary. This avoids any unproved interchange of
logarithm branches across $\Re s=1$.

\subsection{Uniform integrability of arithmetic local zero counts}

\begin{lemma}[Every logarithmic counting moment is bounded]
\label{p07:ar13:lem:UI}
Let $K$ be any fixed compact rectangle in $\Re s>1/2$, and let
$Y_m(t)$ be the number of zeros of $L_m(s+it)$ in $K$.
For every fixed real $q>0$,
\begin{equation}\label{p07:ar13:eq:UI}
 \sup_{T\ge T_0}\frac1T\int_{t_0}^T Y_m(t)^q\,dt<\infty.
\end{equation}
The corresponding random zero counts of $V_{m,\omega}$ have
all finite moments as well.
\end{lemma}

\begin{proof}
Choose $\alpha>1/2$ to the left of $K$. For sufficiently large
fixed $R$, the functions $L_m(R+it)$ are uniformly bounded away
from zero for all real $t$. Indeed, the leading term of $F_0$ is
$2^{-s}$ and that of $F_m$ is $(-\log2)^m2^{-s}$, uniformly as
$\Re s\to\infty$, while $r$ is bounded away from zero on
$\Re s=R>1$. Increasing $R$ if necessary, a disk centered at
$R+i y_0$, with fixed radii $r_1<r_2<R-\alpha$, contains $K$
in its smaller disk. This is possible for any fixed height of
$K$ because increasing $R$ makes the available vertical aperture
arbitrarily large.

Jensen's formula, applied after the vertical shift, bounds
\[
 Y_m(t)\le C+C\int_0^{2\pi}
       \log^+|L_m(R+i(t+y_0)+r_2e^{i\theta})|\,d\theta.
\]
For $|\Im s|\ge1$ and real parts in the fixed disk's range,
\[
 \log^+|L_m(s)|\le C_m+C_m\sum_{j=0}^m
                                  \log(1+|\zeta^{(j)}(s)|).
\]
All $H$ derivatives are uniformly bounded there. Jensen's inequality
for powers at least one, and the elementary inequality
$\log^q(1+x)\le C_q(1+x^2)$, reduce the average of $Y_m(t)^q$
to the off-critical-line mean-square bounds for finitely many
zeta derivatives. The derivative bounds follow from the ordinary
mean-square theorem by Cauchy's formula on fixed smaller circles.
The remaining bounded-height region is compact and contributes
only a finite constant after pole removal. Powers $0<q<1$
follow from the bound for $q=1$.

For the random limits the same argument applies. On the far-right
line the leading-term lower bounds hold uniformly in all phases,
and
$\mathbb E|Z_\omega^{(j)}(s)|^2
 =\sum_n(\log n)^{2j}n^{-2\Re s}<\infty$.
The differential-polynomial identities in the limit and the same
logarithmic inequality prove all the random moment bounds.
\end{proof}

\subsection{Random transversality and a linear strip bound}

\begin{lemma}[No exceptional boundary and bounded mean intensity]
\label{p07:ar13:lem:intensity}
On each compact substrip of $\Re s>1/2$, the random functions
$Q_{m,\omega}$ have almost surely no zero on any prescribed line
segment, and almost surely no multiple zero in any prescribed
compact rectangle. The same assertions hold for $V_{m,\omega}$.
For a rectangle $D$ in such a substrip,
\begin{equation}\label{p07:ar13:eq:intensity}
 \mathbb E\#\{Q_{m,\omega}=0\hbox{ in }D\}\le C_m\,\operatorname{area}(D),
 \qquad
 \mathbb E\#\{V_{m,\omega}=0\hbox{ in }D\}\le C_m\,\operatorname{area}(D),
\end{equation}
where the constants depend on the compact real-part interval but
not on the height or location of $D$ in that interval.
\end{lemma}

\begin{proof}
For fixed $s$, five prime summands in $Q_{m,\omega}(s)$ have a
bounded planar density, uniformly for the real part in a compact
interval. Their characteristic function is a product of five
factors $J_0(b_p|u|)$, with $b_p=(\log p)^mp^{-\Re s}$ bounded
away from zero and infinity on that interval. The bound
$|J_0(x)|\le C(1+|x|)^{-1/2}$ makes this product integrable in
two dimensions. Fourier inversion gives the asserted density,
and convolution with the independent remainder preserves the bound.
The five-prime sum also has uniformly bounded support.

The pair $(Q_{m,\omega},Q_{m,\omega}')$ has a bounded density in
four dimensions. Use ten primes, with $\lambda_p=\log p$.
For every $(u,v)\in\mathbb C^2$, at least nine of the ten numbers
$|u-\lambda_pv|$ are at least $c\|(u,v)\|$: two simultaneous
violations would contradict invertibility of the corresponding
pair of distinct linear forms when $c$ is sufficiently small.
The characteristic function is consequently
$O((1+\|(u,v)\|)^{-9/2})$, which is integrable in dimension four.

These density estimates imply transversality by a mesh argument.
First restrict to the event that the required derivatives on a
compact neighborhood have supremum at most $M$. A zero on a line
segment forces $|Q|\le CMh$ at a mesh point of spacing $h$.
The probability is $O_M(h^{-1}h^2)\to0$.
A multiple zero in a rectangle forces $|Q|\le CMh^2$ and
$|Q'|\le CMh$ at a square-mesh point. The four-dimensional
density bound gives probability $O_M(h^{-2}h^6)\to0$.
Then let $M\to\infty$, using almost sure analyticity on the
neighborhood. The same reasoning applies to any fixed level
$Q=z$. For $V_0$ use the countable levels $P=2\pi i k$;
for $m\ge1$ use the nowhere-vanishing factor $Z_\omega^m$.

We give a direct expected-count bound, including the infinitely
many levels in $V_0$. Split $Q=S+R$, with $S$ the five-prime sum.
On the fixed substrip, $|S|\le A$, $|S'|\le A'$, and the
planar density of $S$ is bounded by $C$. Conditional on the
remaining phases, for $0<\eta<1$ and any target $z$,
\[
 \mathbb E_S\big[|Q'|^2\mathbf1_{\{|Q-z|<\eta\}}\big]
 \le C(1+|R'|^2)\eta^2
                    \mathbf1_{\{|z-R|\le A+1\}}.
\]
The constants are uniform in $s$, and
$\mathbb E|R'(s)|^2$ is uniformly bounded by the convergent
prime square sum. For the single target $z=0$ this proves the
required estimate on the approximate zero-density integral.
For the targets $z=2\pi i k$, only a bounded number of integers
$k$ can satisfy $|2\pi i k-R|\le A+1$, independent of $R$.
Summing the same bound over $k$ therefore remains $O(\eta^2)$.

For a deterministic analytic function having no boundary zero,
its number of zeros in $D$ is the limit, as $\eta\downarrow0$,
of
\[
 \frac1{\pi\eta^2}\int_D
          \mathbf1_{\{|Q(s)-z|<\eta\}}|Q'(s)|^2\,dA(s).
\]
This follows from the local change of variables at simple zeros
(and holds with multiplicity more generally). For the lattice
of levels, sum the displayed integrals over $k$; only finitely
many levels meet the image of the compact closure of $D$.
Fatou's lemma and the conditional estimate give
\eqref{p07:ar13:eq:intensity}. Stationarity under imaginary translation
makes all constants independent of the ordinate.
\end{proof}

\subsection{Passage to the arithmetic endpoint}

Let $D$ be a fixed rectangle in $\Re s>1/2$. By
Lemma~\ref{p07:ar13:lem:intensity}, its boundary almost surely contains
no zero of $V_{m,\omega}$. Zero counting is therefore almost surely
a continuous functional of the limiting analytic function, by
Rouch\'e's theorem. Equations~\eqref{p07:ar13:eq:lift-limit} and
\eqref{p07:ar13:eq:UI} imply both weak convergence and convergence
of every fixed positive moment of the lift counts in $D$.
Uniform integrability is essential here; weak convergence alone
would not justify this assertion.

In particular, with $D_\epsilon=(1,1+\epsilon)\times(0,h)$,
\begin{equation}\label{p07:ar13:eq:thin-window}
 \limsup_{T\to\infty}\frac1T\int_0^T
 \#\{P^{(m)}=0:\ 1<\Re s<1+\epsilon,
                           \ t<\Im s<t+h\}\,dt
        \le C_mh\epsilon.
\end{equation}
The arithmetic count is bounded by the lift count because of
\eqref{p07:ar13:eq:lift-identities}; this remains true for $m=0$
although the lift includes additional levels. Enlarging the strip
slightly proves the same inequality when either vertical edge is
closed. A fixed initial height interval is harmless, since all
zeros there are bounded by those of a holomorphic lift on a compact
set.

For comparison, fix $\gamma>1$. Absolute convergence gives a compact
prime-phase hull for $P^{(m)}$ on $\Re s>1$. Its members are
\[
 \sum_p(-\log p)^m\omega_pp^{-s}.
\]
They have uniform local zero bounds on every compact rectangle:
a contrary sequence would have a compact-uniform convergent
subsequence with a nonzero analytic limit and unbounded zero
counts, contradicting Rouch\'e's theorem. The limit cannot be
identically zero because its far-right leading coefficient is
$(-\log2)^m\omega_2$. The finite-phase equidistribution argument,
compact approximation, and the boundary and multiple-zero null
sets in Lemma~\ref{p07:ar13:lem:intensity} therefore imply convergence
of counts in fixed rectangles to the random counts, with every
moment. Equivalently, bounded zero-count observables are Riemann
integrable on the compact phase torus and their time averages
converge to Haar averages. This is the classical interior
almost-periodic mechanism; it is not a new endpoint argument.

All arithmetic and random zeros in $\Re s>1$ lie to the left of
one fixed $R_m$. Indeed, for sufficiently large $R_m$ the first
prime term dominates the sum of all the others, uniformly in
phases and heights. Thus $X_{m,h}$ differs from its count restricted
to $\Re s>1+\epsilon$ only if a zero occurs in the thin strip.
Markov's inequality and \eqref{p07:ar13:eq:thin-window} show that the
upper asymptotic probability of this difference is $O_m(h\epsilon)$.
For the random functions the same probability bound follows from
\eqref{p07:ar13:eq:intensity}. Letting first $T\to\infty$ at fixed
$\epsilon$, and then $\epsilon\downarrow0$, proves
\eqref{p07:ar13:eq:local-law}. All powers of $X_{m,h}$ are uniformly
integrable by domination by the lift count in the fixed compact
rectangle $[1,R_m]\times[0,h]$ and Lemma~\ref{p07:ar13:lem:UI}.
This proves \eqref{p07:ar13:eq:all-moments}.

Almost surely all random zeros are simple. The same truncation
argument consequently gives the identical local laws and moments
for distinct arithmetic zeros. The random mean zero measure is
stationary under imaginary translations, so its mass in a height
$h$ rectangle equals $h$ times its mass in height one. Put
$c_m=\mathbb E X_{m,1}^{\omega}<\infty$.

To pass from sliding windows to counts up to height $T$, positivity
and Fubini suffice; no uniform short-interval estimate is needed.
For unit windows the integral from $0$ to $T$ counts every zero
with ordinate in $(1,T)$ with weight one, while zeros in the two
end windows receive a weight between zero and one. Consequently
\[
 \int_0^{T-1}X_{m,1}(t)\,dt\le N_m(T),\qquad
 N_m(T)-N_m(1)\le\int_0^T X_{m,1}(t)\,dt.
\]
The first-moment convergence now proves
\eqref{p07:ar13:eq:asymptotic}. The identical argument with a thin
vertical strip proves \eqref{p07:ar13:eq:strip}.

Finally $c_m>0$. Let $a_p=(\log p)^m$ and choose the unique $S_m>1$
satisfying
$a_22^{-S_m}=\sum_{p\ge3}a_pp^{-S_m}$.
The frozen-phase series
$a_22^{-s}-\sum_{p\ge3}a_pp^{-s}$ has a simple zero at $S_m$,
since the derivative of its normalization by $a_22^{-s}$ is
strictly positive there. Translate the zero to height $1/2$.
On a small disk around it, a sufficiently long finite phase
specification controls the absolutely convergent tail uniformly.
Rouch\'e's theorem then places a zero in the disk for a cylinder
set of phases of positive Haar probability. Thus
$\mathbb E X_{m,1}^{\omega}>0$, completing both theorems.

\begin{corollary}[Endpoint ordinates and macroscopic rectangles]
\label{p07:ar13:cor:ordinates}
For every fixed $x\in[0,1]$,
\[
 \frac{N_m(xT)}{N_m(T)}\longrightarrow x.
\]
The convergence is uniform in $x\in[0,1]$. The same conclusion
holds for distinct zeros. More generally, if $I\subset(1,R_m)$
is an interval whose left endpoint may equal one, and
$0\le a<b<\infty$ are fixed, the number of zeros with real part
in $I$ and ordinate in $(aT,bT)$, divided by $T$, converges to
$(b-a)$ times the random mean zero mass in $I\times(0,1)$.
\end{corollary}
\begin{proof}
For each fixed $x>0$ the first assertion follows from
$N_m(xT)=c_mxT+o(T)$. Monotonicity and a finite mesh in $[0,1]$
upgrade pointwise convergence to uniform convergence, with the
interval next to zero controlled by monotonicity. The rectangle
assertion follows from the interior count law and endpoint
truncation exactly as above, followed by subtraction of the
counts up to $aT$ and $bT$.
\end{proof}

\begin{corollary}[A boundary formula for the density constant]
\label{p07:ar13:cor:constant}
Let $Q=Q_{m,\omega}$ and define
$J_m(\sigma)=\mathbb E\log|Q(\sigma)|$ for $\sigma>1/2$.
Then $J_m$ is continuously differentiable, with locally Lipschitz
first derivative, and
\begin{equation}\label{p07:ar13:eq:constant-formula}
 c_m=\frac{-\log2-J_m'(1)}{2\pi}
     =\frac{-\log2-\mathbb E\operatorname{Re}
                       \{Q'(1)/Q(1)\}}{2\pi}.
\end{equation}
The expectation in the last expression is absolutely integrable.
It also has a finite-prime approximation with an explicit
convergence order. Put
\[
 Q_M=\sum_{p\le M}(-\log p)^m\omega_p/p,\qquad
 D_M=\sum_{p\le M}(-\log p)^{m+1}\omega_p/p,
 \qquad \delta_M=M^{-1/6}(\log M)^{m/3}.
\]
Then, as $M\to\infty$,
\begin{equation}\label{p07:ar13:eq:finite-constant}
 c_m=\frac1{2\pi}\left[-\log2-
   \mathbb E\operatorname{Re}
       \frac{D_M\overline{Q_M}}{|Q_M|^2+\delta_M^2}\right]
       +O_m\!\left(M^{-1/6}(\log M)^{m/3}\right).
\end{equation}
The finite expectation is a finite-dimensional torus integral;
no claim about the computational complexity of evaluating that
integral is made.
\end{corollary}
\begin{proof}
The bounded planar density and the finite second moment of $Q$
make $\mathbb E|\log|Q(s)||$ locally bounded. They also make
$J_m$ continuous: the negative logarithmic tail is uniformly
integrable by the planar small-ball bound, the positive logarithmic
tail is uniformly integrable by the local second-moment bound, and
$Q(\sigma)$ varies continuously in mean square. The distributional
identity $\Delta\log|Q|=2\pi\sum_{Q(z)=0}\delta_z$, followed by
Fubini and imaginary-translation stationarity, shows that
$J_m''$ is $2\pi$ times the mean zero measure per unit height.
Lemma~\ref{p07:ar13:lem:intensity} bounds this measure by a constant
multiple of Lebesgue measure on compact real-part intervals.
Consequently $J_m$ has a locally Lipschitz first derivative.

The five-prime decomposition in the proof of that lemma gives,
uniformly on a compact substrip,
\[
 \mathbb E\left[\frac{|Q'|}{|Q|}
                       \mathbf1_{\{|Q|<\eta\}}\right]\le C\eta,
 \qquad 0<\eta<1.
\]
Indeed, conditional on the remaining phases, $|Q'|\le A'+|R'|$,
and integration of $|z|^{-1}$ against the bounded planar density
on the disk $|z|<\eta$ is $O(\eta)$. Away from that disk use
$|Q|^{-1}\le\eta^{-1}$ and $\mathbb E|Q'|<\infty$.
This proves absolute integrability and justifies the derivative
formula $J_m'(\sigma)=\mathbb E\operatorname{Re}(Q'/Q)(\sigma)$,
first as a distributional identity, either by truncation or by
integrating along real segments and using Fubini. The expectation
on the right is continuous: truncate away from $Q=0$, use the
mean-square continuity of $Q,Q'$ and their uniform moment bounds,
and then remove the truncation using the displayed $O(\eta)$
estimate. Thus both continuous representatives agree at every
$\sigma>1/2$, in particular at $\sigma=1$.

For sufficiently large $\sigma$ the first prime term strictly
dominates all the others. Averaging first over $\omega_2$ gives
exactly
$J_m(\sigma)=m\log\log2-\sigma\log2$ there, by the mean-value
formula for the logarithm on a circle. Integrating $J_m''$ from
one to such a real part proves \eqref{p07:ar13:eq:constant-formula}.

For completeness the approximation error can be controlled
without a small-denominator assumption. For $0<\delta<1$,
the same conditional density argument gives
\[
 \mathbb E\left|\frac{Q'}{Q}
       -\frac{Q'\overline Q}{|Q|^2+\delta^2}\right|\le C_m\delta.
\]
On $|Q|\le1$ this follows by integrating
$\delta^2/(|z|(|z|^2+\delta^2))$ against a bounded planar density;
the polar integral is $O(\delta)$. On $|Q|>1$ the contribution
is $O(\delta^2\mathbb E|Q'|)$.
The map $z\mapsto\overline z/(|z|^2+\delta^2)$ has supremum
at most $1/(2\delta)$ and Lipschitz constant at most $C/\delta^2$.
Writing $Q=Q(1)$ and $D=Q'(1)$, Cauchy--Schwarz therefore yields
\[
 \mathbb E\left|
 \frac{D\overline Q}{|Q|^2+\delta^2}
 -\frac{D_M\overline{Q_M}}{|Q_M|^2+\delta^2}\right|
 \le C_m\left(
 \frac{M^{-1/2}(\log M)^{m+1}}{\delta}
 +\frac{M^{-1/2}(\log M)^m}{\delta^2}\right).
\]
Here orthogonality of prime phases and comparison with the
integer tails give
$\mathbb E|Q-Q_M|^2\le C_mM^{-1}(\log M)^{2m}$ and
$\mathbb E|D-D_M|^2\le C_mM^{-1}(\log M)^{2m+2}$;
also $\mathbb E|D_M|^2$ is bounded uniformly in $M$.
Choosing $\delta=\delta_M$ proves
\eqref{p07:ar13:eq:finite-constant}.
\end{proof}

\paragraph{What remains open in this section.}
The moment limits are for each \emph{fixed} window height $h$.
They do not give the asymptotic growth of
$\operatorname{Var}(X_{m,h}^{\omega})$ as $h\to\infty$, a central
limit theorem, GUE statistics, or a height-effective error term.
No conclusion is made about a chosen branch of prime zeta inside
$\Re s<1$, and none of the arguments addresses the Riemann hypothesis.
The present quantitative estimate is linear in the boundary-strip
width after taking the height limit, not a joint uniform estimate
for arbitrary $\epsilon=\epsilon(T)$.

% ===== information =====
\section{Fixed-complexity adaptive information and convex-noise geometry}
\label{p07:in13:section}

The exact-real recovery phenomenon of Krieg--Novak--Ullrich
\cite{p03:an:KNU} allows the individual measurements to become arbitrarily
complicated.  Earlier results in this series replaced their continuous
queries by convex polynomials, with no common degree bound.  We prove
that this lack of a degree bound is essential: with any fixed degree,
fewer measurements than the dimension cannot give arbitrarily accurate
recovery, even under completely unrestricted adaptive selection rules.
The underlying obstruction applies to every uniformly definable
finite-parameter family, rather than only to polynomials.

A separate transfer identifies the geometry controlling adversarial
noise for bounded convex queries.  On an arbitrary convex body the
limiting recovery error is an explicit extremal-radius problem.  On a
norm ball its first-order small-noise coefficient is exactly the Jung
constant of the normed space.  The sharp variation estimate for bounded
convex functions used in this transfer is due to Kwon
\cite{p07:in13:Kwon}; we do not claim that estimate as new.

\subsection{A dimension obstruction for a fixed definable family}

Fix an o-minimal expansion $\mathcal R$ of the ordered real field;
``definable'' permits real parameters in this one structure.  The
semialgebraic case suffices for all polynomial conclusions below.
We use three standard facts: definable choice; the dimension bound
$\dim f(A)\le\dim A$ for a definable map; and the small-frontier
inequality
\begin{equation}\label{p07:in13:frontier}
 \dim(\overline A\setminus A)<\dim A
 \qquad(A\ne\varnothing\text{ definable}).
\end{equation}
These are standard consequences of cell decomposition; see
\cite{p07:in13:Coste}.

Let $K\subset\mathbb R^m$ be nonempty, compact and definable, and
let $P\subset\mathbb R^p$ be a nonempty definable parameter set.
A uniformly definable scalar query family is one definable map
\[
                    F:P\times K\longrightarrow\mathbb R.
\]
An $n$-query algorithm selects $a_j\in P$ as an arbitrary function
of the preceding exact answers and observes $y_j=F(a_j,x)$.
The coefficient selection rules and decoder are not assumed to be
continuous, measurable or definable.  Write
\[
 e_n(K,F)=\inf_A\sup_{x\in K}\|x-A(x)\|_2,
\]
where the infimum runs over all such deterministic adaptive algorithms.
Queries may be padded with any fixed member of the nonempty family,
so ``at most $n$'' is equivalent to allowing $n$ queries.

\begin{theorem}[Uniformly definable information cannot compress dimension]
\label{p07:in13:definable-obstruction}
If $n<\dim K$, then
\begin{equation}\label{p07:in13:positive-floor}
                         e_n(K,F)>0.
\end{equation}
The positive constant can depend on the entire family $F$, its
parameter domain $P$, and $K$.  No continuity assumption on the
individual queries is needed.
\end{theorem}

\begin{proof}
Suppose instead that $e_n(K,F)=0$.  For every $0<\varepsilon<1$
there is an algorithm whose worst-case error is at most $\varepsilon$.
This statement can be written as a finite first-order game:
\begin{equation}\label{p07:in13:game}
\begin{split}
 &\exists a_1\in P\ \forall y_1\in\mathbb R\ \cdots\
  \exists a_n\in P\ \forall y_n\in\mathbb R\
  \exists z\in\mathbb R^m\ \forall x\in K:\\[-2pt]
 &\hspace{8mm}\left[\bigwedge_{j=1}^n F(a_j,x)=y_j\right]
       \ \Longrightarrow\ \|x-z\|_2^2\le\varepsilon^2 .
\end{split}
\end{equation}
Impossible answer histories cause no problem: the antecedent is then
false, and all later existential choices can be arbitrary.  Conversely,
a winning strategy for this formula is exactly an adaptive recovery
algorithm with the asserted error.  Thus the use of the quantified
formula does not restrict the algorithms being considered.

Starting with the first existential quantifier, apply definable choice
to the set of parameters for which the remaining quantified formula
holds.  Substitute the selected parameter into the remaining formula
and repeat.  This constructs a winning strategy simultaneously for all
$\varepsilon\in(0,1)$ in which every choice is definable in
$\varepsilon$ and the earlier answers.  In particular the resulting
decoder is a definable map
\[
          \Phi:(0,1)\times\mathbb R^n\longrightarrow\mathbb R^m.
\]
One may assign fixed values outside any intermediate winning domain;
these extensions are definable.  Form its output family
\[
 A=\{(\varepsilon,\Phi(\varepsilon,y)):
                  0<\varepsilon<1,\ y\in\mathbb R^n\}.
\]
It is definable, and $\dim A\le n+1$.  For every fixed $x\in K$
and every $\varepsilon$, feed $x$ to this strategy and let
$y(\varepsilon,x)$ be its actual answer vector.  Its winning property
gives
\[
       \|\Phi(\varepsilon,y(\varepsilon,x))-x\|_2
                        \le\varepsilon.
\]
Consequently
\[
             \{0\}\times K\subset\overline A\setminus A.
\]
Applying \eqref{p07:in13:frontier}, we obtain
\[
        \dim K\le\dim(\overline A\setminus A)
                      <\dim A\le n+1,
\]
so $\dim K\le n$, a contradiction.  The case $n=0$ follows
from the same proof with no query variables.
\end{proof}

\paragraph{Why arbitrary decision rules do not evade the proof.}
The theorem does not assume that the actual algorithm supplied by an
adversary is definable.  It uses the stronger statement that, if any
winning finite-horizon strategy exists for each accuracy, the definable
finite-horizon game has a uniformly definable winning strategy.  It is
this replacement strategy, not the originally supplied one, whose
output family is subjected to dimension theory.  The parameter set must
remain fixed; a union of polynomial families of all degrees is not a
single finite-parameter definable family.

\begin{corollary}[A common degree bound restores the full query count]
\label{p07:in13:degree-threshold}
Fix integers $m\ge1$, $D\ge1$ and $n<m$.  For adaptive real
polynomial queries of degree at most $D$, there is a constant
$c_{m,n,D}>0$ such that the worst-case Euclidean recovery error on
$B_2^m$ is at least $c_{m,n,D}$.  On $R B_2^m$ it is at least
$R c_{m,n,D}$.  Hence every such $n$-query algorithm on all of
$\mathbb R^m$ has infinite worst-case error.

The exact minimum number of degree-at-most-$D$ queries required for
arbitrary-precision recovery of a ball, or for uniformly bounded-error
recovery of $\mathbb R^m$, is therefore $m$.  The same count holds
if the queries must also be globally convex polynomials.
\end{corollary}

\begin{proof}
Write a polynomial of degree at most $D$ in the monomial basis,
using its coefficient vector in $\mathbb R^{\binom{m+D}{D}}$ as
parameter.  This is a single semialgebraic family, so
Theorem~\ref{p07:in13:definable-obstruction} gives the positive unit-ball
constant.  Replacing $x$ by $Rx$ preserves the degree bound, and
rescaling the decoder divides the error by $R$.  This proves the
scaling assertion and the infinite-error conclusion.  Restricting to
convex polynomials cannot weaken the lower bound.  The $m$ coordinate
queries are affine, globally convex, and recover the vector exactly.
\end{proof}

\begin{theorem}[Algebraicity and decidability of fixed-degree adaptive error]
\label{p07:in13:algebraic-error}
For fixed integers $m,n\ge0$ and $D\ge1$, with $m\ge1$, the
optimal worst-case Euclidean error on $B_2^m$ of $n$ adaptive
polynomial queries of degree at most $D$ is a real algebraic number.
It can in principle be computed exactly by quantifier elimination over
the real closed field of real algebraic numbers.  The same assertion
holds with the additional restriction of global convexity.
For $n<m$ the number is strictly positive; for $n\ge m$ it is zero.
No practical complexity bound for this decision procedure is asserted.
\end{theorem}

\begin{proof}
Use the coefficient parameterization from the preceding corollary.
The winning formula \eqref{p07:in13:game}, including the unit-ball
condition and the squared Euclidean error bound, is a first-order
formula with polynomial equalities and inequalities and rational
coefficients.  Real quantifier elimination therefore computes a
semialgebraic description over $\mathbb Q$ of the set $E$ of
accuracies $\varepsilon\ge0$ for which a winning strategy exists.
The set $E$ is upward closed, nonempty (the zero decoder gives
error one), and its infimum is the optimal error.  Endpoints of
one-dimensional semialgebraic sets over $\mathbb Q$ are real
algebraic numbers, and they can be isolated effectively.  The
infimum need not be attained for this argument.  Global convexity
of a polynomial is expressed by the additional first-order condition
\[
       \forall x,v\in\mathbb R^m:\quad
                  v^{\mathsf T}\nabla^2 p(x)v\ge0.
\]
Thus it can be included in the same quantifier-elimination procedure.
The zero/positive dichotomy follows from
Corollary~\ref{p07:in13:degree-threshold} and coordinate queries.
\end{proof}

\begin{corollary}[Fixed architectures and necessary complexity growth]
\label{p07:in13:architecture}
Let $K$ be a compact definable set of dimension $d$.  Scalar neural
networks of any fixed finite architecture with ReLU activations form a
semialgebraic query family.  Scalar networks of a fixed finite
architecture with logistic activations form a family definable in the
real exponential field.  For either family, fewer than $d$ adaptive
queries have a strictly positive worst-case recovery error on $K$,
even with unrestricted real weights and unrestricted adaptive decision
rules.

In particular, any sequence of fewer-than-$m$ polynomial-query
algorithms whose errors on $B_2^m$ tend to zero must have maximum
query degrees tending to infinity along every subsequence with
vanishing error.  Unbounded coefficients alone cannot substitute for
this growth.
\end{corollary}

\begin{proof}
A fixed ReLU network is a finite composition of affine maps and
$u\mapsto\max(0,u)$, hence is semialgebraic jointly in weights and
input.  Logistic activation $1/(1+e^{-u})$ is definable in the
real exponential field, which is o-minimal; see \cite{p07:in13:Wilkie}.
Apply Theorem~\ref{p07:in13:definable-obstruction}.  If the polynomial
conclusion failed, some degree bound $D$ would hold along a
vanishing-error subsequence, contradicting
Corollary~\ref{p07:in13:degree-threshold}.
\end{proof}

The fixed-architecture conclusion concerns recovery from scalar
observations of an unknown input, not the expressive power of a
network decoder supplied with the entire input.  It does not settle
the unrestricted continuous or real-analytic two-query question in
$\mathbb R^3$, and it does not contradict logarithmic query counts
when polynomial degrees or network architectures may grow.

\subsection{Convex noise as an extremal-radius problem}

Let $K\subset\mathbb R^m$ be a compact convex body with nonempty
interior, and fix any norm $\|\cdot\|$ for the reconstruction loss.
Let $\mathcal C(K)$ consist of the continuous convex functions on
$K$ with range in $[-1,1]$.  Observations have the form
\[
                y_j=f_j(x)+\xi_j,\qquad |\xi_j|\le\delta,
\]
where $0<\delta<1$ and every selected $f_j$ obeys the range
restriction on all of $K$.  Let $e_n^{\rm ad}$ and
$e_n^{\rm non}$ be the deterministic worst-case recovery errors for
adaptive and nonadaptive algorithms respectively.  For a nonempty
compact set $S$ put
\[
       \operatorname{rad}(S)=\inf_{z\in\mathbb R^m}
                                  \sup_{x\in S}\|x-z\|.
\]
The centers are allowed to lie anywhere in the ambient space.

\begin{lemma}[Convex indistinguishability criterion]
\label{p07:in13:homothety}
For $x,y\in K$ and $0<\delta<1$, the following conditions are
equivalent:
\begin{align}
 &|f(x)-f(y)|\le2\delta\quad\text{for every }f\in\mathcal C(K);
                                                   \label{p07:in13:indist}\\
 &y-(1-\delta)x\in\delta K,
        \qquad x-(1-\delta)y\in\delta K.             \label{p07:in13:homot}
\end{align}
In testing the first condition, restrictions of globally convex hinge
functions suffice.
\end{lemma}

\begin{proof}
If $y=(1-\delta)x+\delta z$ with $z\in K$, convexity and the
range restriction give
\[
             f(y)-f(x)\le\delta(f(z)-f(x))\le2\delta.
\]
The reversed containment gives the reverse inequality.
For the converse, suppose $y-(1-\delta)x\notin\delta K$.
Strict separation supplies a linear functional $u$ such that
\[
      u(y)-(1-\delta)u(x)>\delta h_K(u),
      \qquad h_K(u)=\max_{z\in K}u(z).
\]
Set $a=h_K(u)-u(x)$.  Here $a>0$: if $a=0$, the displayed
strict inequality would imply $u(y)>h_K(u)$, impossible for $y\in K$.
The globally convex function
\[
                f(z)=-1+\frac{2(u(z)-u(x))_+}{a}
\]
has range in $[-1,1]$ on $K$, satisfies $f(x)=-1$, and has
$f(y)-f(x)>2\delta$.  The other failed containment is treated by
interchanging $x$ and $y$.  This is also a direct proof of the
part of Kwon's sharp variation estimate needed here.
\end{proof}

\begin{theorem}[Exact limiting error on every convex body]
\label{p07:in13:radius-formula}
Define
\begin{equation}\label{p07:in13:radius-value}
 \mathcal R_K(\delta)=
 \max\left\{\operatorname{rad}(S):\varnothing\ne S\subset K
     \text{ compact},\quad S-(1-\delta)S\subset\delta K\right\}.
\end{equation}
Then
\begin{equation}\label{p07:in13:limit-equality}
 \inf_{n\ge0}e_n^{\rm ad}
       =\inf_{n\ge0}e_n^{\rm non}=\mathcal R_K(\delta).
\end{equation}
For every $\varepsilon>0$, a finite nonadaptive family of globally
convex hinge queries attains error at most
$\mathcal R_K(\delta)+\varepsilon$.  The same infimum holds if
queries are restricted to globally convex real-analytic functions,
or to globally convex polynomials of unrestricted degree.
\end{theorem}

\begin{proof}
The hyperspace of nonempty compact subsets of $K$, equipped with the
Hausdorff metric, is compact.  The radius is continuous (indeed
one-Lipschitz), and the containment in \eqref{p07:in13:radius-value}
is closed under Hausdorff limits.  Thus the maximum exists.

If $S$ satisfies that containment, the lemma shows that the range of
every admissible query on $S$ has width at most $2\delta$.  Against
an arbitrary adaptive algorithm, answer each selected query by the
midpoint of its maximum and minimum on $S$.  This answer is within
$\delta$ of the true answer for every $x\in S$, so the same
transcript remains possible for all of $S$.  The worst-case error is
at least $\operatorname{rad}(S)$, proving the lower bound.

For the upper bound, fix $\varepsilon>0$ and consider the compact
hyperspace subset
\[
 \mathcal B=\{S\subset K:\ S\ne\varnothing\text{ compact},\
          \operatorname{rad}(S)\ge\mathcal R_K(\delta)+\varepsilon\}.
\]
Every $S\in\mathcal B$ fails the containment, so by the lemma
there is a globally convex hinge $f_S$, with range in $[-1,1]$
on $K$, for which
\[
                     \max_S f_S-\min_S f_S>2\delta.
\]
For a fixed continuous function the oscillation is continuous on the
hyperspace.  The sets where these strict inequalities hold therefore
form an open cover of $\mathcal B$.  Choose a finite subcover,
corresponding to $f_1,\ldots,f_N$, and query all of them.
For any possible answer vector $y$, the feasible set
\[
             C_y=\{x\in K:|f_i(x)-y_i|\le\delta, 1\le i\le N\}
\]
is compact and nonempty.  Each $f_i$ has oscillation at most
$2\delta$ on $C_y$, so $C_y\notin\mathcal B$.  Outputting a
Chebyshev center of $C_y$ gives error less than
$\mathcal R_K(\delta)+\varepsilon$.  If $\mathcal B$ is empty,
no query is needed for this accuracy.

For the regularity variants it suffices to replace each hinge by
arbitrarily close uniformly approximating globally convex polynomials
on $K$, while retaining the range bound.  Here is a one-dimensional
construction, applied to the affine argument of the hinge.  Approximate
the hinge uniformly on a compact interval by a smooth globally convex
function $h$.  Approximate the continuous function $\sqrt{h''}$
on that interval by a polynomial $q$, and integrate $q^2$ twice,
matching the value and derivative of $h$ at one point.  The resulting
polynomial is globally convex and uniformly approximates $h$ on the
interval.  After composition with the linear functional it is globally
convex on $\mathbb R^m$.  If its approximation error to a
range-normalized hinge is at most $\eta$, division by $1+\eta$
puts its range back in $[-1,1]$ on $K$ and keeps the approximation
error tending to zero.  For each bad set choose an approximation close
enough to retain its strict oscillation inequality; the same compact
subcover argument then applies.  These polynomials are also globally
real analytic.
\end{proof}

\paragraph{Geometry rather than pairwise half-distance.}
The lower bound is an entire-set obstruction.  A common answer must be
compatible with every point of $S$, and the reconstruction error is
its circumradius, which can strictly exceed half its diameter.  This is
why the cube calculation alone does not determine Euclidean-ball noise
constants.  Equation~\eqref{p07:in13:limit-equality} also concerns the
infimum over all finite query counts: it does not say that adaptivity
never improves a fixed finite-query error.

\subsection{The Jung constant is the exact small-noise coefficient}

Let $X=(\mathbb R^m,\|\cdot\|)$ be any finite-dimensional normed
space and let $B_X$ be its closed unit ball.  Use the same norm for
the loss.  Define the Jung constant in the radius/diameter convention by
\[
 J(X)=\sup\left\{\operatorname{rad}(S):
           S\subset X\text{ nonempty compact},\
           \operatorname{diam}(S)\le1\right\}.
\]
This convention has $J(\ell_\infty^m)=1/2$ and
$J(\ell_2^m)=\sqrt{m/(2m+2)}$.  In the theorem below put
$e_\infty(\delta)=\inf_n e_n^{\rm ad}=\inf_n e_n^{\rm non}$
for bounded convex queries on $B_X$.

\begin{theorem}[Jung law for adversarial convex information]
\label{p07:in13:Jung-law}
For every $0<\delta<1$,
\begin{equation}\label{p07:in13:Jung-bounds}
 \frac{J(X)\delta}{1+J(X)\delta}
       \le e_\infty(\delta)
       \le\frac{2J(X)\delta}{2-\delta}.
\end{equation}
In particular,
\begin{equation}\label{p07:in13:Jung-limit}
       \lim_{\delta\downarrow0}\frac{e_\infty(\delta)}\delta
                            =J(X).
\end{equation}
The same result holds for globally convex polynomials of unrestricted
degree and for globally convex real-analytic queries.
\end{theorem}

\begin{proof}
Let $S\subset B_X$ satisfy the containment in
\eqref{p07:in13:radius-value}, and put $q=1-\delta$.  For $x,y\in S$,
\[
                  \|x-qy\|\le\delta,\qquad\|y-qx\|\le\delta.
\]
Subtracting these vectors and using the triangle inequality gives
\[
                (1+q)\|x-y\|\le2\delta.
\]
Thus $\operatorname{diam}(S)\le2\delta/(2-\delta)$, and the
upper bound follows from the definition of $J(X)$ and
Theorem~\ref{p07:in13:radius-formula}.

For the lower bound take any compact $S_0$ of diameter at most one
and radius $r$, and translate a Chebyshev center to zero.  Then
$\|x\|\le r$ for $x\in S_0$.  Set
\[
                     a=\frac{\delta}{1+\delta r},\qquad S=aS_0.
\]
Since $ar<1$, the set $S$ lies in $B_X$.  For $x,y\in S_0$,
\[
       \|ax-(1-\delta)ay\|
           \le a\bigl(\|x-y\|+\delta\|y\|\bigr)\le\delta.
\]
Hence $S$ is admissible in \eqref{p07:in13:radius-value}, giving
$e_\infty(\delta)\ge ar$.  Let $r\uparrow J(X)$.
The two bounds have the same first-order coefficient as
$\delta\downarrow0$.  The regularity variants follow from
Theorem~\ref{p07:in13:radius-formula}.
\end{proof}

\begin{theorem}[The exact Euclidean convex-noise floor]
\label{p07:in13:Euclidean}
For range-normalized convex queries on the Euclidean unit ball in
$\mathbb R^m$, with Euclidean reconstruction loss, the exact
unlimited-query minimax error is
\begin{equation}\label{p07:in13:Euclidean-exact}
 e_\infty(\delta)
   =\frac{\delta}
       {\sqrt{\delta^2+2(1-\delta)(m+1)/m}},
                   \qquad 0<\delta<1.
\end{equation}
A centered regular $m$-simplex realizes the extremal ambiguity.
The formula holds for the infimum over finite adaptive counts and for
the infimum over finite nonadaptive counts, including the globally
convex polynomial and real-analytic subclasses.
In particular,
\[
        e_\infty(\delta)=\sqrt{\frac{m}{2m+2}}\,\delta
                              +O_m(\delta^2).
\]
For $m>1$ and every $0<\delta<1$, this is strictly larger than
the two-antipodal-point obstruction $\delta/(2-\delta)$.
\end{theorem}

\begin{proof}
Put $q=1-\delta$.  For the lower bound let $u_0,\ldots,u_m$
be a centered regular simplex on the unit sphere, so
$\langle u_i,u_j\rangle=-1/m$ for $i\ne j$ and
$\sum_i u_i=0$.  Set
\[
 r=\frac{\delta}{\sqrt{1+q^2+2q/m}},\qquad
                         S=\{ru_0,\ldots,ru_m\}.
\]
For distinct vertices, $\|ru_i-qru_j\|_2=\delta$; for equal
vertices that norm is $\delta r\le\delta$.  Thus $S$ is admissible
in Theorem~\ref{p07:in13:radius-formula}, and its circumradius is $r$.

For the upper bound let $S\subset B_2^m$ be any nonempty compact
set satisfying $S-qS\subset\delta B_2^m$.  Write $c$ and $R$ for
its Euclidean Chebyshev center and radius.  If $R=0$ there is
nothing to prove.  The contact-point characterization of a smallest
enclosing ball, followed by Carath\'eodory's theorem, gives
$2\le N\le m+1$ contact points
\[
 x_i=c+R u_i\in S,\quad \|u_i\|_2=1,\qquad
 \alpha_i>0,\quad \sum_i\alpha_i=1,\quad\sum_i\alpha_i u_i=0.
\]
The contact characterization follows by separation: if zero were
outside the convex hull of the contact directions, moving the center
a small distance toward a separating direction would strictly decrease
the enclosing radius.

For each fixed $i$, multiply
$\|x_i-qx_j\|_2^2\le\delta^2$ by $\alpha_j$ and sum over
$j\ne i$.  Expansion, using $\sum_j\alpha_j u_j=0$, gives
\begin{equation}\label{p07:in13:weighted-contact}
\begin{split}
 &(1-\alpha_i)\delta^2\|c\|_2^2
 +R^2(1+q^2-\alpha_i\delta^2)\\
 &\quad+2\delta R(1-\delta\alpha_i)\langle c,u_i\rangle
                         \le(1-\alpha_i)\delta^2.
\end{split}
\end{equation}
Multiply this inequality by $\alpha_i/(1-\delta\alpha_i)$ and
sum over $i$.  These weights cancel the center cross term exactly.
Put
\[
 U=\sum_i\frac{\alpha_i}{1-\delta\alpha_i},\qquad
 V=\sum_i\frac{\alpha_i^2}{1-\delta\alpha_i}
       =\frac{U-1}{\delta},\qquad
 A=U-V,\quad B=(1+q^2)U-\delta^2 V.
\]
Here $A>0$ because $N\ge2$, and $B>0$.  Equation
\eqref{p07:in13:weighted-contact} becomes
\[
             R^2 B\le\delta^2 A(1-\|c\|_2^2)\le\delta^2 A.
\]
The expressions simplify to
\[
          A=\frac{1-qU}{\delta},\qquad
          B=q(1+q)U+\delta.
\]
The function $A/B$ is decreasing in $U$, since its derivative is
$-2q/(\delta B^2)<0$.  Convexity of $a\mapsto a/(1-\delta a)$
and $N\le m+1$ yield
\[
 U\ge\frac{N}{N-\delta}\ge\frac{m+1}{m+1-\delta}.
\]
At the last lower bound for $U$, direct substitution gives
\[
 \frac AB\le\frac{1}{1+q^2+2q/m},\qquad
 R^2\le\frac{\delta^2}{1+q^2+2q/m}.
\]
This matches the regular-simplex construction and proves the exact
formula via Theorem~\ref{p07:in13:radius-formula}.  Since
$1+q^2+2q/m=\delta^2+2(1-\delta)(m+1)/m$, it is the stated
expression.  The small-noise expansion follows immediately.  For
$m>1$, its squared denominator is strictly smaller than
$(2-\delta)^2=1+q^2+2q$, proving the strict comparison.
\end{proof}

\paragraph{Status and remaining problems.}
The fixed-definable-family theorem and the convex-noise radius formula
are new results claimed by this manuscript, with complete proofs above.
The bounded-convex variation estimate and the classical Jung constants
are existing ingredients.  Quantitative lower bounds for
$c_{m,n,D}$, finite-query attainment of the convex-noise infimum,
and the unrestricted two-real-analytic-query problem in
$\mathbb R^3$ remain open here.  None of these unresolved statements
is counted as a solved problem.

% ===== gaussian =====
\section{Irregular sampling and the scalar Gaussian tail boundary}
\label{p07:ga13:sec:gaussian}

The earlier matrix construction requires a quantitative upper bound on
the density of the selected terminal times.  In one dimension that
restriction can be removed completely.  The reason is geometric:
the poles of a scalar Blaschke product lie on the negative real axis,
so a contour with nonzero imaginary part can approach that axis without
crossing a pole.  A narrower contour absorbs the whole Blaschke product
using only the sum of its positive zeros.  No regular variation or
separation of the terminal times is needed.

Let $\mathcal N$ be an infinite subset of the positive integers.
Conditionally on $(A,V)\in\mathbb R\times[0,\infty)$ let the $X_i$
be independent $N(A,V)$ variables, and put $S_n=\sum_{i=1}^nX_i$.
We observe the separate population laws of the $S_n$, $n\in\mathcal N$.
Choose a sufficiently large constant $E$ and define
\[
 L_1(v)=\log(E+v),\qquad L_{j+1}(v)=\log L_j(v),\qquad
 W_{m,a}(v)=\frac{\sqrt v}
 {L_1(v)\cdots L_{m-1}(v)L_m(v)^a}.
\]
Let $\mathcal C_{m,a}$ consist of all joint mixing laws for which
$\mathbb E\exp\{cW_{m,a}(V)\}<\infty$ for some $c>0$.
There is no restriction on dependence between $A$ and $V$.

\begin{theorem}[Complete scalar design--tail classification]
\label{p07:ga13:thm:scalar}
For every fixed integer $m\ge1$, every $a\ge0$, and every infinite
terminal set $\mathcal N$, the selected partial-sum laws identify every
mixing law in $\mathcal C_{m,a}$ if and only if
\begin{equation}\label{p07:ga13:eq:dichotomy}
 a\le1\quad\hbox{or}\quad\sum_{n\in\mathcal N}\frac1n=\infty.
\end{equation}
When both conditions fail, two distinct absolutely continuous witnesses
can have the same mean marginal, the same variance marginal, and all
the same mixed polynomial moments.  All these moments are finite.
The witnesses satisfy $\mathbb E\exp\{cW_{m,a}(V)\}<\infty$ for
every $c>0$, and can obey any prescribed lower bound $V\ge b_0>0$.
\end{theorem}

The positive alternatives in \eqref{p07:ga13:eq:dichotomy} were proved in
the preceding manuscripts.  The new assertion is the negative
alternative for \emph{every} convergent harmonic subseries, including
highly irregular ones.  We give the positive argument as well so that
the classification is self-contained.  This is a closure of a gap in
the preceding manuscripts, not an assertion that a named external
conjecture has been solved.  The relevant classical moment theory
includes de Jeu's quasi-analytic weights~\cite{p05:fg:deJeu}.

\subsection{A flat Laplace transform with arbitrary summable zeros}

For large $R$, write
\[
 P_m(R)=\prod_{j=1}^m\log_jR,\qquad
 D_{m,p}(R)=P_m(R)(\log_mR)^p,\qquad p>0.
\]
These expressions are only used for sufficiently large arguments;
elsewhere they can be extended smoothly and positively.

\begin{lemma}[Scalar logarithmic regularizer]
\label{p07:ga13:lem:regularizer}
For every $m\ge1$ and $p>0$ there is a function $Q$, holomorphic
off the negative real axis, such that $Q(\overline w)=\overline{Q(w)}$,
$\operatorname{Re}Q(w)\ge0$ for $\operatorname{Re}w\ge0$, and
\begin{equation}\label{p07:ga13:eq:sector}
 \operatorname{Re}Q(w)\ge
 \frac{c|w|}{D_{m,p}(|w|)}
 \quad\left(|w|\ge R_0,\quad
 |\arg w|\le\frac\pi2+\frac{c}{P_m(|w|)}\right).
\end{equation}
Every fixed derivative of $Q$ has at most polynomial growth in
$|w|+|w|^{-1}$ on any slightly narrower such sector.
\end{lemma}
\begin{proof}
Choose $K$ sufficiently large, use principal logarithms, and put
\[
 J_1(w)=\log(K+w),\qquad
 J_{j+1}(w)=\log(K+J_j(w)),\qquad
 Q(w)=\frac{w}{J_m(w)^p}.
\]
All the logarithms have positive real part on the right half-plane.
Their arguments have the sign of $\arg w$.  More explicitly,
$|\arg J_1(w)|\le|\arg w|/\log K$, since
$\operatorname{Re}J_1\ge\log K$ and
$|\operatorname{Im}J_1|=|\arg(K+w)|\le|\arg w|$;
iteration gives $|\arg J_m(w)|\le|\arg w|/(\log K)^m$.
Enlarging $K$ therefore makes
$p|\arg J_m(w)|\le|\arg w|$ there.  Thus the argument of $Q(w)$
lies between zero and the argument of $w$, proving nonnegativity.
For bounded real $\theta$, successive expansions of the logarithms give,
uniformly together with their first angular derivatives,
\[
 |J_m(Re^{i\theta})|\sim\log_m R,\qquad
 \arg J_m(Re^{i\theta})
   =\frac{\theta}{P_m(R)}(1+o(1)).
\]
Near $|\theta|=\pi/2$ the subtraction of $p\arg J_m$ therefore
leaves an angular margin of a positive constant times $1/P_m(R)$.
Allowing a sufficiently small overshoot of $|\theta|$ preserves half
this margin.  Multiplying the resulting cosine bound by
$|Q(w)|\asymp R/(\log_mR)^p$ proves \eqref{p07:ga13:eq:sector}.
Away from the two sector endpoints the cosine has a larger lower
bound.  The derivative assertion follows either by differentiating
the displayed formulas or by Cauchy's formula on smaller sectors.
All branch cuts stay on the negative real axis because each nested
logarithm maps the upper half-plane into itself and is positive on
the positive real axis; reflection handles the lower half-plane.
\end{proof}

\begin{lemma}[Summable zeros with a near-critical tail]
\label{p07:ga13:lem:flat}
Let $0<\lambda_j\le M$ be distinct, with
$\Lambda=\sum_j\lambda_j<\infty$.  For every $m\ge1$ and $p>0$
there is a nonzero real Schwartz function $f$ supported in $[0,\infty)$
whose Laplace transform vanishes at every $\lambda_j$ and whose
integral is zero.  For each integer $k\ge0$ there are positive
constants $C_k,c_k$ such that
\begin{equation}\label{p07:ga13:eq:flat-tail}
 |f^{(k)}(v)|\le C_k\exp\{-c_kW_{m,1+p}(v)\},\qquad v\ge0.
\end{equation}
\end{lemma}
\begin{proof}
The product
\[
 B(z)=\prod_j\frac{z-\lambda_j}{z+\lambda_j}
\]
converges locally uniformly in the right half-plane, is nonzero
identically, and has modulus at most one there.  It continues
meromorphically off zero with poles only at $-\lambda_j$.  On
the imaginary axis, away from zero, its modulus is one.  Put
\[
 F(z)=B(z)\exp\{-\sqrt z-Q(1/z)\},\qquad\operatorname{Re}z>0.
\]
The principal square root and Lemma~\ref{p07:ga13:lem:regularizer} imply
$|F(z)|\le e^{-\operatorname{Re}\sqrt z}$.  In particular $F$ is
bounded and is not identically zero.

We first justify ordinary Fourier inversion.  Logarithmic
differentiation of $B$, followed by differentiation of each term,
gives
\[
 \left|\left(\frac{B'}B\right)^{(r)}(iy)\right|
 \le C_r\Lambda |y|^{-r-2},\qquad y\ne0.
\]
Indeed, each difference of inverse powers can be integrated over
$[-\lambda_j,\lambda_j]$, where the distance to $iy$ is at least
$|y|$.  The recursion for $B^{(k)}/B$ supplies polynomial bounds
for every derivative of $B(iy)$ in $|y|+|y|^{-1}$.
Equation~\eqref{p07:ga13:eq:sector} supplies damping
$\exp\{-c/(|y|D_{m,p}(1/|y|))\}$ as $y\to0$, while
$e^{-\sqrt z}$ supplies $e^{-c\sqrt{|y|}}$ at infinity.
Consequently $F(iy)$, extended by zero at $y=0$, is a Schwartz
function flat at zero.  Define
\[
 f(v)=\frac1{2\pi}\int_{\mathbb R}e^{iyv}F(iy)\,dy.
\]
It is real, Schwartz, and nonzero.  Shifting the line to
$x+i\mathbb R$, $x>0$, proves that $f(v)=0$ for $v<0$:
the bound $\operatorname{Re}\sqrt{x+iy}\ge\sqrt{|y|/2}$ makes
the integrals uniformly bounded before multiplication by $e^{xv}$,
and then $x\to\infty$.  The same contour identity, or Cauchy's
formula followed by Fubini, proves
$\int_0^\infty e^{-zv}f(v)\,dv=F(z)$ for $\operatorname{Re}z>0$.
It also gives $\int f=\lim_{z\downarrow0}F(z)=0$.

Here is the new estimate that removes all density assumptions on
the zero set.  If $z=-\eta+iy$, $r=|y|$, and $0\le\eta\le r/2$,
then
\begin{align}
 \log|B(z)|
 &=\frac12\sum_j\log\left(1+
       \frac{4\eta\lambda_j}{(\lambda_j-\eta)^2+r^2}\right)
       \notag\\
 &\le C\eta\sum_j\frac{\lambda_j}{\lambda_j^2+r^2}
 \le\frac{C\Lambda\eta}{r^2}.
 \label{p07:ga13:eq:product-shift}
\end{align}
The first inequality uses
$(\lambda-\eta)^2+r^2\ge c(\lambda^2+r^2)$ when $\eta\le r/2$.
Choose a smooth positive function $\eta(r)$ such that, near zero,
\begin{equation}\label{p07:ga13:eq:narrow}
 \eta(r)=\frac{\varepsilon r}{D_{m,p}(1/r)}.
\end{equation}
Extend it to a sufficiently small positive constant for large $r$,
with bounded derivative and $\eta(r)\le r/2$.  The constant
$\varepsilon>0$ may depend on $\Lambda,m,p$.
Since $D_{m,p}(R)\ge P_m(R)$ eventually, every inverse point
$(-t\eta(r)+iy)^{-1}$, $0\le t\le1$, lies in the sector of
\eqref{p07:ga13:eq:sector} for small $r$, and has modulus comparable
to $1/r$.  Equations~\eqref{p07:ga13:eq:product-shift} and
\eqref{p07:ga13:eq:narrow} show that choosing $\varepsilon$ small
enough absorbs the whole product into half the $Q$ damping.
Thus uniformly along this homotopy, for small $r$,
\[
 |F(-t\eta(r)+iy)|
 \le C\exp\left\{-\frac{c}{rD_{m,p}(1/r)}\right\}.
\]
For large $r$ the estimate is $Ce^{-c\sqrt r}$; on the remaining
compact set it is bounded.  These latter bounds follow directly
from \eqref{p07:ga13:eq:product-shift}, the boundedness of $Q(1/z)$
there, and the square-root factor.

Deform the upper and lower imaginary half-axes separately to
$z=-\eta(|y|)+iy$.  This deformation never crosses a negative
real pole or a branch cut.  The connecting arcs at zero vanish by
the displayed flatness estimate, and the arcs at infinity vanish
by square-root damping.  The line element has bounded modulus
because $\eta'$ is bounded.  For $v\ge0$ the deformed integral
therefore has, near zero, absolute integrand bounded by
\[
 C\exp\left\{-c\left(
 \frac{vr}{D_{m,p}(1/r)}+
 \frac1{rD_{m,p}(1/r)}\right)\right\}.
\]
Uniformly in small $r$ and large $v$, the expression in parentheses
is at least $cW_{m,1+p}(v)$.  On
$\frac14\log v\le\log(1/r)\le\log v$ all the fixed-depth
logarithms are comparable with the corresponding logarithms of $v$,
and $(vr+r^{-1})\ge2\sqrt v$ proves this assertion.  If
$\log(1/r)<\frac14\log v$, the first term is at least $v^{3/4}$
divided by logarithmic factors.  If $\log(1/r)>\log v$,
$r^{-1}/D_{m,p}(1/r)$ is increasing and is at least
$v/D_{m,p}(v)$.  Both are larger than the required bound.
For $r$ bounded away from zero, $\eta(r)$ has a positive lower
bound and gives exponential decay in $v$.  Reserving half the
damping for integrability proves \eqref{p07:ga13:eq:flat-tail} for
$k=0$.  Differentiating the inverse integral $k$ times introduces
only the factor $z^k$ and gives the same conclusion for every
fixed $k$.
\end{proof}

\begin{proof}[Proof of Theorem~\ref{p07:ga13:thm:scalar}]
For the negative assertion put $\lambda_n=(2n)^{-1}$ and choose
$0<p<a-1$.  Let $f$ be the preceding lemma's function and choose
a nonzero real even $h\in C_c^\infty(\{t:1<|t|<2\})$.  Define
\[
 g(a,v)=\frac1{2\pi}\int e^{-ita}h(t)t^2 f(vt^2)\,dt.
\]
This is a nonzero real Schwartz function supported on $v\ge0$.
Its Fourier--Laplace transform is
\[
 \int e^{ita-sv}g(a,v)\,da\,dv=h(t)F(s/t^2)
 \quad(t\ne0,s>0),
\]
and is zero at $t=0$.  At every observation argument
$s=t^2/(2n)$ it vanishes.  Its $v$ marginal and all joint
polynomial moments vanish because its partial Fourier transform
is zero near $t=0$.  Its $a$ marginal vanishes because $\int f=0$.
The normalized positive and negative parts of $g$ are therefore
the desired distinct mixing probabilities.

To check the quantitative tail, differentiate under the $t$ integral
and integrate by parts there.  Every resulting term is a bounded
coefficient times $v^j f^{(j)}(vt^2)$ on $1\le|t|\le2$.
Equation~\eqref{p07:ga13:eq:flat-tail} absorbs the polynomial factors
and gives, for every integer $M$, constants $C_M,c_M>0$ with
\[
 |g(a,v)|\le C_M(1+|a|)^{-M}
       e^{-c_MW_{m,1+p}(v)}.
\]
Since $W_{m,1+p}(v)/W_{m,a}(v)\to\infty$, the required
integrability holds for every $c>0$.  Translation of $v$ by $b_0$
multiplies the transform by $e^{-b_0s}$ and preserves all zeros,
equal marginals, equal moments, and tail properties.

For completeness, suppose two candidate laws agree and let $\rho$
be their difference.  For fixed $t\ne0$ the bounded holomorphic
function $H_t(s)=\int e^{ita-sv}\,d\rho(a,v)$ has zeros
$t^2/(2n)$, $n\in\mathcal N$.  If $\sum1/n=\infty$, the
half-plane Blaschke condition forces $H_t\equiv0$.  Indeed, map
the half-plane to the disk and apply Jensen's formula; the zeros
of a nonzero bounded holomorphic function must satisfy
$\sum_j\operatorname{Re}z_j/(1+|z_j|^2)<\infty$.

If $a\le1$, let $\tau$ be the sum of the two variance marginals.
The tail assumption gives
\[
 \left(\int v^k\,d\tau(v)\right)^{1/(2k)}
 \le C k\log k\cdots\log_{m-1}k\,(\log_m k)^a
\]
for large $k$.  One obtains this bound by maximizing
$y^{2k}\exp\{-cy/[L_1(y^2)\cdots
L_{m-1}(y^2)L_m(y^2)^a]\}$, with $y=\sqrt v$.
The Stieltjes--Carleman series diverges.  Since the zeros of $H_t$
accumulate at zero, Taylor expansion from the positive real axis
successively shows that every derivative of $H_t$ at zero vanishes;
all derivatives exist by the moment bound.  The resulting finite
complex measure in $v$, dominated by $\tau$, has every moment zero.
Apply moment determinacy to $\tau$ plus and minus its real and
imaginary parts to conclude it is zero.  Thus $H_t\equiv0$ again.
Continuity includes $t=0$, and Fourier--Laplace uniqueness gives
$\rho=0$ in both cases.
\end{proof}

\begin{remark}[Why the same proof does not settle the matrix gap]
\label{p07:ga13:rem:matrix}
The scalar inverse variable $1/z$ has nonzero imaginary part all
along the contour.  In dimension greater than one, an inverse
quadratic form $u^TT^{-1}u$ can be negative real on a shifted
matrix contour even though $T$ is nonsingular.  A scalar Blaschke
factor composed with this form then has genuine pole surfaces.
The proof above does not cross or remove those surfaces.  The exact
irregular-design classification for full matrix covariance mixtures
therefore remains open in this manuscript.
\end{remark}

\section{Exact sparse designs for every fixed stable shape}
\label{p07:ga13:sec:stable}

The 223-page manuscript classified full partial-sum observations for
fixed stable shape, but stated its sparse sampling criteria only as
sufficient.  The next theorem proves their necessity as well.  The
construction transfers the Gaussian Blaschke method to nonreal zero
sets and, at index one, to a frequency-dependent imaginary shift.

Fix $0<\alpha\le2$ and $-1\le\beta\le1$, and use the stable
parameterization
\[
 \widehat K_{a,b}(t)=e^{ita-b\psi_{\alpha,\beta}(t)},\qquad
 \psi_{\alpha,\beta}(t)=
 \begin{cases}
 |t|^\alpha[1-i\beta\tan(\pi\alpha/2)\operatorname{sgn}t],
       &\alpha\ne1,\\
 |t|[1+i(2\beta/\pi)\operatorname{sgn}t\log|t|],&\alpha=1,
 \end{cases}
\]
with $\psi(0)=0$, $a\in\mathbb R$, and $b\ge0$.
This is Nolan's $1$-parameterization~\cite{p07:ga13:Nolan}.
Draw $(A,B)$ from an arbitrary joint mixing law and conditionally draw
the $X_i$ independently from $K_{A,B}$.  Only the separate laws of
$S_n=\sum_{i=1}^nX_i$, $n\in\mathcal N$, are supplied.

\begin{theorem}[Sharp stable sampling classification]
\label{p07:ga13:thm:stable}
Let $\mathcal N\subset\mathbb N$.  Identification of every Borel
mixing law on $\mathbb R\times[0,\infty)$ holds under exactly the
following conditions:
\begin{enumerate}
\item If $\alpha\ne1$, it holds if and only if
\begin{equation}\label{p07:ga13:eq:stable-power}
 \sum_{n\in\mathcal N}n^{-|1-\alpha|}=\infty.
\end{equation}
\item If $\alpha=1$ and $\beta\ne0$, it holds if and only if
\begin{equation}\label{p07:ga13:eq:stable-log}
 \sum_{n\in\mathcal N}\frac1{1+(\log n)^2}=\infty.
\end{equation}
\item If $\alpha=1$ and $\beta=0$, no terminal set identifies every
mixing law, even $\mathcal N=\mathbb N$.
\end{enumerate}
In either of the first two cases, whenever the displayed series
converges the ambiguity has absolutely continuous witnesses with
the same coordinate marginals and all the same mixed polynomial
moments, all finite.  Any fixed constraint $B\ge b_0>0$ can be imposed.
\end{theorem}

The Gaussian case $\alpha=2$ and the sufficient directions are
inherited.  The new results are the converse for the other power-law
indices and the converse for the skew index-one family.  No mixing
over $\alpha$ or $\beta$ is allowed.  The observation problem belongs
to the exchangeable partial-sum framework of Lin~\cite{p01:mix:Lin2015};
we make no claim that this exact sparse-design statement was a named
conjecture there.

\subsection{A boundary-regular Blaschke construction}

\begin{lemma}[Three families of zero sets]
\label{p07:ga13:lem:blaschke}
Suppose the distinct points $a_j$ belong to one of the following
three classes:
\begin{enumerate}
\item $a_j=c\lambda_j$, where $\operatorname{Re}c>0$,
      $0<\lambda_j\le M$, and $\sum_j\lambda_j<\infty$;
\item $a_j=c\lambda_j$, where $\operatorname{Re}c>0$,
      $\lambda_j\ge c_0>0$, $\lambda_j\to\infty$, and
      $\sum_j\lambda_j^{-1}<\infty$;
\item $a_j=1+iy_j$, where $|y_j|\to\infty$ and
      $\sum_j(1+y_j^2)^{-1}<\infty$.
\end{enumerate}
Finite sets are also permitted.  There is a nonzero bounded
holomorphic function $B$ on the right half-plane with these zeros,
$|B|\le1$, and unimodular boundary values on the imaginary axis
except possibly at zero in the first case.  In the first case every
boundary derivative grows at most polynomially in $|y|+|y|^{-1}$;
in the other cases it grows at most polynomially in $1+|y|$.
\end{lemma}
\begin{proof}
All three assumptions imply the half-plane Blaschke condition
$\sum\operatorname{Re}a_j/(1+|a_j|^2)<\infty$.  The usual disk
Blaschke product after the map $(z-1)/(z+1)$ gives
\[
 B(z)=\prod_j\gamma_j\frac{z-a_j}{z+\overline{a_j}},
 \qquad |\gamma_j|=1,
\]
with the standard convergence normalization.  The same product
converges locally through each part of the imaginary axis away
from its possible accumulation point zero; its boundary modulus
is one.  These assertions can also be checked directly by expanding
each normalized factor for small $a_j$, large ray points $a_j$, or
large vertical points, respectively.

Here are explicit derivative estimates, which will justify all
subsequent inversions.  Write $Q_B=B'/B$.  On the imaginary axis
\[
 Q_B(z)=\sum_j\left(\frac1{z-a_j}
                        -\frac1{z+\overline{a_j}}\right).
\]
For each $r\ge0$, an individual derivative is at most
\[
 \frac{C_r\operatorname{Re}a_j}
 {[(\operatorname{Re}a_j)^2+(y-\operatorname{Im}a_j)^2]^{(r+2)/2}}
 \quad\text{at }z=iy.
\]
For $a_j=c\lambda_j$, the denominator inside brackets is
bounded below by a positive constant times $y^2+\lambda_j^2$.
In the first case summing gives a bound
$C_r(\sum\lambda_j)|y|^{-r-2}$.  In the second case it gives a
uniform constant, since each summand is at most
$C_r\lambda_j^{-r-1}$ and $\lambda_j\ge c_0$.
In the vertical case use
\[
 \frac1{1+(y-y_j)^2}\le
 \frac{2(1+y^2)}{1+y_j^2}
\]
to bound every derivative of $Q_B$ by $C_r(1+y^2)$.
Recursion from $B'=Q_BB$, together with $|B(iy)|=1$, proves the
asserted polynomial derivative bounds.  The convergence of the
displayed derivative series is uniform on the specified compact
boundary intervals, so these differentiations are legitimate.
\end{proof}

\begin{proof}[Proof of Theorem~\ref{p07:ga13:thm:stable}]
The data for $S_n/n$ are values of the mixing Fourier--Laplace
transform at
\[
 (t,z_n(t)),\qquad z_n(t)=n\psi_{\alpha,\beta}(t/n).
\]
For $\alpha\ne1$ these points are
\[
 z_n(t)=|t|^\alpha(1-i\kappa\operatorname{sgn}t)n^{1-\alpha},
 \qquad \kappa=\beta\tan(\pi\alpha/2).
\]
Their Blaschke sum diverges exactly when
\eqref{p07:ga13:eq:stable-power} holds.  For $\alpha=1$, put
$k=2\beta/\pi$.  The points are
\[
 z_n(t)=|t|+ikt\log|t|-ikt\log n,
\]
whose Blaschke sum diverges exactly under
\eqref{p07:ga13:eq:stable-log} when $k\ne0$.  The bounded-holomorphic
uniqueness argument from the preceding theorem, followed by
Fourier--Laplace uniqueness, proves both sufficient directions
without any moment assumption.

For necessity first suppose $\alpha\ne1$ and the series converges.
Let $B_+$ be the preceding lemma's Blaschke product with zeros
\[
 a_n=(1-i\kappa)n^{1-\alpha},\qquad n\in\mathcal N.
\]
These are the first family if $\alpha>1$ and the second family
if $\alpha<1$.  For $t>0$ and $\operatorname{Re}s>0$ define
\begin{equation}\label{p07:ga13:eq:stable-H}
 F_t(s)=B_+(s/t^\alpha)
       \exp\{-\sqrt{s/t^\alpha}-1/\sqrt{s/t^\alpha}\}.
\end{equation}
It is bounded and holomorphic in $s$, is not identically zero,
and vanishes at every $s=z_n(t)$.

If $\alpha=1$, $k\ne0$, and the logarithmic series converges,
use instead the product with zeros
$a_n=1-ik\log n$, the third family of the lemma, and define
\begin{equation}\label{p07:ga13:eq:stable-H-one}
 F_t(s)=B_+(s/t-ik\log t)
       \exp\{-\sqrt{s/t}-1/\sqrt{s/t}\},\qquad t>0.
\end{equation}
The shift inside $B_+$ is purely imaginary, so boundedness on
the right half-plane is preserved.  At an observation argument,
$z_n(t)/t-ik\log t=1-ik\log n$, proving the required zeros.
The regularizing factor is deliberately \emph{not} shifted: this
ensures flatness at $s=0$ and will preserve the mean marginal.

In both cases set
$F_{-t}(s)=\overline{F_t(\overline s)}$ for $t>0$ and choose a
nonzero real even $h\in C_c^\infty(\{t:1<|t|<2\})$.
For $t$ in this compact support, Lemma~\ref{p07:ga13:lem:blaschke}
and the chain rule show that every derivative in $t$ and $y$ of
$F_t(iy)$ has at most polynomial factors before multiplication
by
\[
 \exp\{-c\sqrt{|y|}-c/\sqrt{|y|}\}.
\]
Thus $h(t)F_t(iy)$ is Schwartz in $(t,y)$, compactly supported
in $t$, and flat at $y=0$.  In the ray case this uses the lemma's
possible polynomial singularity at zero; in the vertical case
the imaginary shift ranges over a fixed compact interval and the
same estimates remain uniform.  Define
\[
 G(t,b)=\frac{h(t)}{2\pi}\int_{\mathbb R}e^{iyb}F_t(iy)\,dy,
 \qquad
 g(a,b)=\frac1{2\pi}\int_{\mathbb R}e^{-ita}G(t,b)\,dt.
\]
The function $g$ is real, nonzero, and Schwartz.  To see the
support in $b\ge0$, for each fixed $t$ shift the $s$ contour to
$x+i\mathbb R$.  The Blaschke factor remains bounded by one,
including after the pure imaginary shift, while the square-root
factor gives an integrable bound independent of $x\ge0$.
Letting $x\to\infty$ forces the inverse transform to vanish
for $b<0$.  Fourier inversion and the same elementary half-plane
contour argument give
\[
 \int e^{ita-sb}g(a,b)\,da\,db=h(t)F_t(s),
 \qquad \operatorname{Re}s>0.
\]
Every selected observation transform is therefore zero.

The $b$ marginal and all mixed moments of $g$ vanish because
$G(t,b)$ is zero near $t=0$.  Its $a$ marginal vanishes as well:
boundedness of the Blaschke factor and the unshifted
$e^{-1/\sqrt{s/t^\alpha}}$ or $e^{-1/\sqrt{s/t}}$ force
$F_t(s)\to0$ as $s\downarrow0$.  Dominated convergence in the
Schwartz inverse transform justifies passing to this limit.
Normalizing the positive and negative parts of $g$ produces the
claimed witnesses.  Replacing $B$ by $B+b_0$ multiplies their
signed transform by $e^{-b_0s}$ and preserves all assertions.

Finally, for the symmetric Cauchy case take independent standard
symmetric Cauchy variables $Z_0,Z_1,\ldots$ and constants $b,c>0$.
The sequence $X_j=cZ_0+bZ_j$ and an iid Cauchy sequence of scale
$b+c$ have identical partial-sum laws for every $n$ but different
joint laws: their pairwise differences have scales $2b$ and
$2(b+c)$, respectively.  Thus no terminal subset suffices.
This last counterexample is recalled and is not counted as new.
\end{proof}

\begin{example}[Different sparse clocks distinguish the stable shapes]
For $n_j\asymp j^q$, $q\ge1$, and $\alpha\ne1$, the exact
condition is $q|1-\alpha|\le1$.  For skew index one, the same
polynomial clock always identifies, because
$\sum_j(\log j)^{-2}=\infty$.  If instead
$n_j\asymp\exp(j^q)$ with $q>0$, skew index one identifies
exactly when $q\le1/2$, whereas every $\alpha\ne1$ fails for
every such $q$.  These conclusions are direct consequences of
Theorem~\ref{p07:ga13:thm:stable}, not separate solved problems.
\end{example}

% ===== spectral =====
\section{Finite-height spectral observation and sharp noisy order detection}
\label{p07:sp13:sec}

Two gaps remained after exact threshold inversion: the transition was an
infinite-height observable, and positive noise prevented unconditional
identification of the atom count.  We give an explicit finite-height
experiment under a quantitative recurrence assumption, and determine the
noise scale at which a separated weak atom can be certified.  These are
newly formulated problems for the nonlinear observation below.  The
number-resolution exponent $2m-2$ is already familiar in linear
super-resolution, including Liu--Ammari~\cite{p07:sp13:LiuAmmari}; our claim
is its sharp transfer, including arbitrary small masses, to nonlinear
minimum observations.  Quantitative Kronecker approximation is classical
\cite{p07:sp13:GonekMontgomery}; no optimal recurrence exponent is claimed.

Fix $0<a<b<1$ and set
\begin{equation}\label{p07:sp13:eq:model}
 S_\mu(x)=\int_a^b x^r\,d\mu(r),\qquad
 T_\mu(C)=\min_{x>0}\frac{1+Cx}{S_\mu(x)},
 \qquad\mu\in\mathcal P[a,b].
\end{equation}
Here $\mathcal P_k[a,b]$ denotes the probabilities with at most $k$
atoms.  A list $\mathcal C$ is one-sided if its settings are distinct
and belong to $(0,a/(1-a))$.  Constants may depend on the fixed support
interval, setting list, and specified atom bound, but not on unknown
atom masses, separations, or phases.

\subsection{The exact scale for detecting one additional atom}

\begin{theorem}[Sharp nonlinear model-order boundary]
\label{p07:sp13:thm:order}
Fix $m\ge2$ and any $2m-2$ one-sided settings $\mathcal C$.  Let
$\mu=\sum_{j=1}^m w_j\delta_{r_j}$, where
$r_{j+1}-r_j\ge\Delta>0$ and $w_j\ge w>0$.  Then
\begin{equation}\label{p07:sp13:eq:orderlower}
 \inf_{\nu\in\mathcal P_{m-1}[a,b]}
 \max_{C\in\mathcal C}|T_\mu(C)-T_\nu(C)|
                      \ge c_m w\Delta^{2m-2}.
\end{equation}
The joint dependence on $w$ and $\Delta$ is sharp.  For every
$0<w\le1/m$ and every sufficiently small $\Delta>0$, there are such
an $m$-atomic $\mu$ and an $(m-1)$-atomic probability $\nu$ with
\begin{equation}\label{p07:sp13:eq:orderupper}
 \sup_{C\in I}|T_\mu(C)-T_\nu(C)|
                      \le K_{m,I}w\Delta^{2m-2}
\end{equation}
for each fixed compact setting interval $I\subset(0,\infty)$.
The constants are uniform as $w\downarrow0$.
\end{theorem}

Thus adversarial threshold noise $\delta$ permits uniformly reliable
rejection of every smaller positive spectrum when
$\delta<c w\Delta^{2m-2}$; noise $\delta\ge Cw\Delta^{2m-2}$ can
make an $m$-atom and an $(m-1)$-atom spectrum indistinguishable.  This
is a deterministic testing boundary, not a claim that a vanishingly
small extra atom can be counted at fixed noise.

We first record explicitly the nonlinear-to-linear step used in the
proof.  It also explains why the observation settings need not locate
the unobserved minimizers.

\begin{lemma}[Smooth tests against a sparse competitor]
\label{p07:sp13:lem:comparison}
Fix $n\ge1$ and $2n$ one-sided settings.  For
$\alpha\in\mathcal P_n[a,b]$, $\beta\in\mathcal P[a,b]$, and
$\varphi\in C^{2n}[a,b]$,
\begin{equation}\label{p07:sp13:eq:test}
 \left|\int\varphi\,d(\beta-\alpha)\right|
 \le K\|\varphi\|_{C^{2n}}
                \max_i|T_\beta(C_i)-T_\alpha(C_i)|.
\end{equation}
The constant does not depend on separation or masses in $\alpha$.
\end{lemma}
\begin{proof}
For any positive profile $S$, put $H=S-xS'$.  Its unique contact at
setting $C$ satisfies
\[
 C=\frac{S'(x)}{H(x)},\qquad
 \frac{dC}{dx}=\frac{S''(x)S(x)}{H(x)^2}<0,
 \qquad L_C(x)=\frac{1+Cx}{T(C)}.
\]
The line $L_C$ majorizes $S$ and is tangent at that contact.  Moreover
$a/(1-a)\le Cx\le b/(1-b)$.  Hence, for the fixed settings, all
contacts lie in a compact interval strictly above $1$.  On this
interval $S,H,-S''$ are uniformly bounded above and below by positive
constants.  Consecutive contacts, tangent slopes, and the distances
from their tangent-line intersections to either contact therefore
have uniform positive lower bounds.  For the last assertion, the
strict concavity gap between contacts $v>w$ is at least
$\inf(-S'')(v-w)^2/2$; divide it by the bounded tangent-slope difference.

Write $\delta$ for the threshold discrepancy.  For small $\delta$,
the two sets of tangent-line coefficients and their intersections
are within $O(\delta)$.  The contacts belonging to setting $i$ for
both profiles lie in the same interval between neighboring tangent
intersections, with uniform margins.  At the first profile's contact
$S_\alpha-S_\beta\ge-K\delta$, and at the second it is at most
$K\delta$.  An intermediate point has absolute difference at most
$K\delta$.  Taking logarithms gives $2n$ uniformly separated points
$u_i$ in a compact positive interval, where
\[
 \left|\int e^{u_i r}\,d(\beta-\alpha)(r)\right|\le K\delta.
\]
Add the exact sample at $u_0=0$, using equal total masses.

Here is the requisite uniform positive interpolation argument.  Pad
the support of $\alpha$ to $n$ nodes $t_j$, allowing repetitions.
Interpolate $\varphi$ with multiplicity two at each node by $R\varphi$
in the span of the first $2n$ exponentials among $e^{u_i r}$.
Divided-difference interpolation matrices extend continuously across
node collisions and are always invertible: a nonzero sum of $2n$
distinct real exponentials has at most $2n-1$ zeros counted with
multiplicity, by repeated conjugation, differentiation, and Rolle's
theorem.  Compactness therefore bounds the interpolant coefficients
and its $C^{2n}$ norm by $K\|\varphi\|_{C^{2n}}$.

The interpolation residual $q$ of the remaining exponential vanishes
to even multiplicity at every $t_j$.  With one choice of sign,
\[
 q(r)\ge c\prod_{j=1}^n(r-t_j)^2.
\]
Indeed the quotient is a continuous divided difference, including
collisions, and cannot vanish by the same zero count for $2n+1$
exponentials.  Compactness gives the uniform lower bound.  The
ordinary Hermite remainder bound now shows that
$R\varphi\pm K\|\varphi\|_{C^{2n}}q$ bracket $\varphi$ everywhere
and agree with it on the support of $\alpha$.  Integrating against
the positive measure $\beta$ proves \eqref{p07:sp13:eq:test} from the
$2n+1$ exponential discrepancies.  Large $\delta$ are absorbed by
increasing $K$, using the trivial bound $2\|\varphi\|_\infty$.
\end{proof}

\begin{proof}[Proof of Theorem~\ref{p07:sp13:thm:order}]
For a competitor $\nu$ pad its support to $m-1$ nodes $s_j$ and put
$Q(r)=\prod_{j=1}^{m-1}(r-s_j)^2$.  Its $C^{2m-2}$ norm is bounded
uniformly.  Each open ball of radius $\Delta/2$ around an $s_j$
contains at most one of the $m$ true atoms.  At least one true atom
therefore has distance at least $\Delta/2$ from every $s_j$.  Hence
\[
 \int Q\,d\mu\ge w(\Delta/2)^{2m-2},\qquad \int Q\,d\nu=0.
\]
Lemma~\ref{p07:sp13:lem:comparison}, with sparse side $\nu$, gives the
lower bound, uniformly over all competitors.

For the matching example, let $n=m-1$, choose an interior $r_0$, and
put the $m$ atoms at $r_0+j\Delta$, $0\le j\le n$.  Give the final
atom mass $w$ and each of the first $n$ atoms mass $(1-w)/n\ge w$.
Let $p$ be the monic orthogonal polynomial of degree $n$ for this
probability.  It exists because the moment matrix up to degree $n$
is positive definite.  Its $n$ roots are simple and in the interior
of the convex hull of the support: otherwise the product of its
sign-change factors has degree less than $n$, and its product with
$p$ has a constant sign on the support, contradicting orthogonality.

The Gaussian quadrature measure $\nu$ on those roots has positive
weights and agrees with $\mu$ through degree $2n-1$.  To verify this
without an existence citation, divide any polynomial of degree at
most $2n-1$ by $p$ and interpolate its remainder on the roots.
Orthogonality makes the integral equal to the corresponding weighted
sum.  The resulting weight at root $s_i$ equals the integral of the
square of its degree-$(n-1)$ Lagrange polynomial, so it is positive;
the weights sum to one.

Hermite interpolation of any $C^{2n}$ function at these roots gives
\begin{equation}\label{p07:sp13:eq:gausserror}
 \left|\int\varphi\,d(\mu-\nu)\right|
 \le\frac{\|\varphi^{(2n)}\|_\infty}{(2n)!}\int p(r)^2\,d\mu(r).
\end{equation}
The monic orthogonal polynomial minimizes this squared norm.
Comparing with $\prod_{j=0}^{n-1}(r-r_0-j\Delta)$ yields
\[
 \int p^2\,d\mu\le w(n!)^2\Delta^{2n}.
\]
For $x$ in a fixed compact subset of $(0,\infty)$ apply
\eqref{p07:sp13:eq:gausserror} to $\varphi(r)=x^r$.  Its derivatives are
uniformly bounded, so $|S_\mu(x)-S_\nu(x)|\le K w\Delta^{2n}$.
All contacts for $C\in I$ lie in one such compact interval, and the
denominators $S$ are uniformly positive there.  Comparing the two
minimum quotients gives \eqref{p07:sp13:eq:orderupper}.
\end{proof}

\begin{corollary}[Certified order from noisy data]
\label{p07:sp13:cor:orderprocedure}
Suppose the truth has $m\le M$ atoms and observe $2M-2$ fixed
one-sided thresholds with errors at most $\delta$.  Return the smallest
$k\ge1$ for which some member of $\mathcal P_k[a,b]$ is consistent
with those errors.  This information-theoretic procedure returns
$k=m$ whenever $m=1$, or whenever $m\ge2$ and
$2\delta<c_m w\Delta^{2m-2}$.  Conversely, at noise of order
$w\Delta^{2m-2}$ even the entire threshold curve on a compact setting
interval cannot uniformly distinguish $m$ from $m-1$ atoms.
\end{corollary}
\begin{proof}
The truth is feasible.  Every smaller feasible measure would have
threshold discrepancy at most $2\delta$, contradicting
\eqref{p07:sp13:eq:orderlower} on a sublist of $2m-2$ settings.
For the converse use the midpoint of the two curves constructed in
\eqref{p07:sp13:eq:orderupper}.  No complexity assertion is made for the
nonlinear feasibility search.
\end{proof}

\subsection{A finite-height replacement for the transition oracle}

Write the known largest frequency as $\ell>0$, to distinguish it
from the support endpoint $b$.  Suppose
$\mu=\sum_{j=1}^m w_j\delta_{r_j}$ and consider the controlled family
\begin{equation}\label{p07:sp13:eq:exponential}
 f_{A,C}(z)=-1+A\sum_{j=1}^m w_j e^{\ell r_jz+i\phi_j}
                            +Ce^{\ell z+i\phi_*}.
\end{equation}
The control amplitudes $A,C$, the frequency $\ell$, and the top
phase $\phi_*$ are known.  The smaller frequencies, weights and
phases need not be known.  An experiment at $(A,C)$ reports all zero
ordinates in $0<t<H$ exactly.  This is a finite-height observation
model, not an algorithmic claim about computing a certified zero list.

Assume the frequency vector
$\omega=(\ell r_1,\ldots,\ell r_m,\ell)\in\mathbb R^d$,
$d=m+1$, obeys the specified quantitative condition
\begin{equation}\label{p07:sp13:eq:dio}
 |k\cdot\omega|\ge\gamma\|k\|_\infty^{-\tau}
 \quad(0\ne k\in\mathbb Z^d),\qquad \gamma>0,\quad\tau\ge0.
\end{equation}
In particular, its entries are rationally independent.

\begin{theorem}[Finite amplitude grid and finite zero height]
\label{p07:sp13:thm:finiteheight}
Fix a compact setting interval $I\subset(0,\infty)$.  There are
known constants $A_-<A_+$, $h_0,R,K>0$, depending only on
$a,b,I,\ell,m,\tau$, with the following property.  For
$0<h<h_0$ take an equally spaced amplitude grid in $[A_-,A_+]$
with spacing in $[h/2,h]$, including its endpoints, and observe
\eqref{p07:sp13:eq:exponential} at its $O(h^{-1})$ settings.  Declare a grid amplitude
accepted if one of its observed zero ordinates satisfies
\begin{equation}\label{p07:sp13:eq:arc}
       \operatorname{dist}(\ell t+\phi_*,\pi+2\pi\mathbb Z)
                                                   \le R\sqrt h.
\end{equation}
If
\begin{equation}\label{p07:sp13:eq:height}
                  H\ge3+K\gamma^{-1}h^{-(d+\tau)},
\end{equation}
there is an accepted amplitude, and the least accepted one obeys
\begin{equation}\label{p07:sp13:eq:thresholdaccuracy}
                T_\mu(C)-Kh\le\widehat T_h(C)\le T_\mu(C)+h.
\end{equation}
The result is uniform in $C\in I$, all positive masses, and all
smaller phases.  It does not assume that a qualitative density
property can be recognized from a finite list.
\end{theorem}

\begin{lemma}[An elementary uniform recurrence bound]
\label{p07:sp13:lem:recurrence}
Under \eqref{p07:sp13:eq:dio}, every target in $(\mathbb R/2\pi\mathbb Z)^d$
is approached within coordinate distance $\eta$ by $t\omega$ in
every interval of length at least
$K_{d,\tau}\gamma^{-1}\eta^{-(d+\tau)}$, for $0<\eta<1$.
\end{lemma}
\begin{proof}
Choose a nonnegative smooth bump supported in the coordinate cube of
radius one, scale it to radius $\eta$, and periodize around the
target.  Its mean is $c\eta^d>0$ and its Fourier coefficients obey
$|\widehat q(k)|\le C_N\eta^d(1+\eta\|k\|)^{-N}$ for every fixed
$N$.  Averaging $q(t\omega)$ on an interval of length $L$ differs
from its mean by at most
\[
 \frac{2}{L}\sum_{k\ne0}\frac{|\widehat q(k)|}{|k\cdot\omega|}
 \le\frac{K}{\gamma L}\eta^{-\tau}.
\]
Here take $N>d+\tau$ and compare the lattice sum with its scaled
integral.  For the stated $L$, with a sufficiently large constant,
the average is positive.  Some point of the orbit therefore meets
the bump support.  The bound is independent of the interval's start.
\end{proof}

\begin{proof}[Proof of Theorem~\ref{p07:sp13:thm:finiteheight}]
All transition values lie in a fixed positive compact interval:
$T(C)\le1+C$ by $x=1$, while
$S(x)\le\max(x^a,x^b)$ gives a uniform positive lower bound.
Choose $A_-,A_+$ with positive margins around that interval.

Let $x_*$ be the contact at $T=T_\mu(C)$ and
$\sigma_*=(\log x_*)/\ell$.  The phase-aligned function
\[
                 g(z)=-1+T S_\mu(e^{\ell z})-Ce^{\ell z}
\]
has a double zero at $\sigma_*$.  Indeed $TS(x_*)=1+Cx_*$ and
$TS'(x_*)=C$.  Its second derivative is
\[
 g''(\sigma_*)=-\ell^2T\int r(1-r)x_*^r\,d\mu(r)<0.
\]
The contact bound $a/(1-a)\le Cx_*\le b/(1-b)$ makes its magnitude
uniformly bounded below, while the third derivative is uniformly
bounded in a fixed complex neighborhood.  Consequently, for a fixed
small $\rho_0$ and $|z-\sigma_*|=\rho\le\rho_0$,
\begin{equation}\label{p07:sp13:eq:doublezero}
                         |g(z)|\ge c\rho^2.
\end{equation}
The same Taylor estimate ensures that $g$ has exactly two zeros,
counted with multiplicity, inside that circle.

Choose a grid amplitude $A_0\in[T,T+h]$.  Apply
Lemma~\ref{p07:sp13:lem:recurrence}, with $\eta=h$, to target phases
$(-\phi_1,\ldots,-\phi_m,\pi-\phi_*)$.  For the height in
\eqref{p07:sp13:eq:height}, some $t_0\in[1,H-1]$ meets that target.
On the fixed neighborhood of $\sigma_*$,
\[
                         |f_{A_0,C}(z+it_0)-g(z)|\le K_0h.
\]
Choose a fixed $R_0$ large enough that $cR_0^2>K_0$ and use the
circle of radius $R_0\sqrt h$ in \eqref{p07:sp13:eq:doublezero}.
For small $h$, Rouche's theorem gives two zeros in that disk for
$f_{A_0,C}(z+it_0)$.  Their original ordinates lie in $(0,H)$, and
their top phases differ from $\pi$ by at most
$h+\ell R_0\sqrt h$.  Taking $R=1+\ell R_0$ proves acceptance
of $A_0$ and the upper bound in \eqref{p07:sp13:eq:thresholdaccuracy}.

Conversely, let any accepted amplitude have zero $\sigma+it$, put
$x=e^{\ell\sigma}$ and write its top angle as $\pi+v$ with
$|v|\le R\sqrt h$.  The zero equation and the triangle inequality
imply
\[
 AS_\mu(x)\ge |1-Cxe^{i(\pi+v)}|.
\]
For $y=Cx$,
\[
 |1-ye^{i(\pi+v)}|^2
   =(1+y)^2-2y(1-\cos v)
   \ge(1+y)^2-yv^2
   \ge(1+y)^2(1-v^2/4).
\]
Thus $A\ge(1-R^2h/4)(1+Cx)/S_\mu(x)
\ge(1-R^2h/4)T$.  The uniform upper bound on $T$ proves the lower
bound in \eqref{p07:sp13:eq:thresholdaccuracy}.
\end{proof}

The mechanism involves two quadratic effects: the critical aligned
function has a double zero, while the missing-arc deficit is quadratic
in the angular distance from $\pi$.  They balance, so phase recurrence
at accuracy $h$ yields transition accuracy $h$, although the nearby
zeros themselves move on the larger scale $\sqrt h$.

\begin{corollary}[Conditional finite-height spectral resolution]
\label{p07:sp13:cor:recovery}
Suppose the truth has at most $m$ atoms, and run the finite-height
experiment at $2m-1$ fixed one-sided $C$ settings.  A minimum-distance
fit in $\mathcal P_m[a,b]$ satisfies
\begin{equation}\label{p07:sp13:eq:recovery}
 W_1(\widehat\mu,\mu)\le K h^{1/(2m-1)}.
\end{equation}
For fixed Diophantine parameters this yields
$O(H^{-1/((d+\tau)(2m-1))})$ at $O(H^{1/(d+\tau)})$
amplitude settings per $C$, with constants depending on $\gamma$.
Using $2m$ settings and allowing any positive consistent competitor
instead gives the blind radius $O(h^{1/(2m)})$.
For order certification up to a known bound $M$, the analogous
$2M-2$ setting experiment succeeds if $h<c w\Delta^{2m-2}$.
\end{corollary}
\begin{proof}
The finite-threshold comparison of the preceding paper
\cite{comp:part:06} gives
$W_1(\alpha,\beta)\le K\max_i|T_\alpha(C_i)-T_\beta(C_i)|^{1/(2m-1)}$
for two at most $m$-atomic probabilities, with $2m-1$ one-sided
settings.  Its elementary proof uses the separated near-zeros from
Lemma~\ref{p07:sp13:lem:comparison}: the transform difference is a sum
of at most $2m$ real exponentials, so its values at $2m$ separated
points, including $0$, uniformly control its first $2m-1$ moments.
A polynomial with roots at the at most $2m-2$ sign changes of the
CDF difference then gives $W_1^{2m-1}\le K$ times the moment error.
For one sparse side and an arbitrary positive side, the same proof
uses \eqref{p07:sp13:eq:test} and at most $2m-1$ CDF sign changes,
giving exponent $1/(2m)$.  Apply these bounds to the threshold errors
from Theorem~\ref{p07:sp13:thm:finiteheight}.  Minimum-distance fitting
at most doubles them.  The final assertion follows from
Corollary~\ref{p07:sp13:cor:orderprocedure}.
\end{proof}

\begin{proposition}[Why rational independence alone supplies no height]
\label{p07:sp13:prop:noheight}
Even for one unknown frequency with support interval
$[a,b]=[1/3,2/3]$, no uniform finite observation height
can guarantee detection of the arc in \eqref{p07:sp13:eq:arc} over all
rationally independent frequency pairs.  More precisely, fix
$\ell=C=1$, zero phases, and any finite $H$.  There is an irrational
$r$ arbitrarily close to $1/2$ for which, simultaneously for every
$A\in[1,4]$, the function $-1+Ae^{rz}+e^z$ has no zero in
$0\le\operatorname{Im}z\le H$ whose ordinate is within $\pi/4$ of
$\pi+2\pi\mathbb Z$.  Nevertheless its transition $T_{\delta_r}(1)$
lies near $2$, in the interior of this amplitude interval.
\end{proposition}
\begin{proof}
At $r=1/2$, substitute $u=e^{z/2}$.  The quadratic
$u^2+Au-1$ has two nonzero real roots for every $A\in[1,4]$.
Consequently every zero has $\operatorname{Im}z\in2\pi\mathbb Z$.
Uniformly for $A\in[1,4]$ and $r$ near $1/2$, all zeros lie in one
fixed vertical strip: the constant dominates far left and $e^z$
dominates far right.  On the compact part of that strip satisfying
the stated height and arc restrictions, the modulus at $r=1/2$ has
a positive minimum, uniformly also in $A$.  Continuity preserves
that positive lower bound for all sufficiently close irrational $r$.
Finally $T_{\delta_r}(1)=r^{-r}(1-r)^{-(1-r)}\to2$.
\end{proof}

The finite-height theorem therefore closes a conditional observation
bridge, and the final proposition shows precisely why its quantitative
hypothesis cannot simply be omitted.  The height exponent above is
only an explicit sufficient bound.  Neither an optimal hitting-time
exponent nor a uniform quantitative assertion for an infinite prime
vector is claimed.  The model-order theorem closes a different,
sharp deterministic problem; its exponent is a transfer of known
super-resolution geometry, not a newly discovered universal exponent.

% ===== bessel =====
\section{Gamma deconvolution closes the remaining McKay SD gap}
\label{p07:be13:section}

The remaining strictness question in the generalized McKay
self-decomposability problem has a different solution from the finite
moment tests of the preceding manuscript.  Dividing out the gamma
factor suggested by the tail produces an explicit signed density.
A Sturm identity forces this density to change sign.  This gives
both a sharp gamma-factor obstruction and the missing strict SD
boundary for every admissible parameter.  We then determine the
singular cost of a tilt-uniform HCM power just above the critical
Bessel order.  The general McKay questions are posed in
\cite[Section~4.3]{p01:rare:Baricz2026}; our gamma-factor and critical-power
questions are newly formulated transfer problems.

\subsection{Notation and previously established inputs}

Scale $a=1$ initially.  Let
\[
 \alpha=\mu+\tfrac12,\qquad q=\nu-\mu-1,\qquad
 M=2\alpha+q=\mu+\nu,\qquad A=\alpha+q=\nu-\tfrac12,
\]
and write
\[
 J(u)=\int_0^\infty x^{\nu-1}e^{-ux}I_\mu(x)\,dx,
 \qquad u>1.
\]
The density is $x^{\nu-1}e^{-bx}I_\mu(x)/J(b)$, where
$\alpha>-1/2$, $M>0$, and $b>1$.
The preceding signed-spectral construction
\cite[Section~4]{comp:part:04} proves that every such law is
ID.  For $q>0$, it supplies a positive smooth even kernel $K$ with
\begin{equation}\label{p07:be13:kernel-input}
 -\frac{J'(u)}{J(u)}=\int_0^\infty e^{-ux}K(x)\,dx,
 \quad K(0)=M,\quad
 e^{-x}K(x)\longrightarrow A,\quad \frac{K'(x)}{K(x)}\longrightarrow1.
\end{equation}
Its SD boundary is
\begin{equation}\label{p07:be13:beta-input}
 \beta_{\alpha,q}=\sup_{x\geq0}\frac{K'(x)}{K(x)}<\infty:
 \qquad X\text{ is SD}\ \Longleftrightarrow\ b\geq\beta_{\alpha,q}.
\end{equation}
For $q\leq0$ the law is GGC, hence SD, for every admissible tilt.
These classifications and kernel identities are inputs, not new
results counted in this section.  In contrast, strictness of
\eqref{p07:be13:beta-input} for all $q>0$ was not established previously.

\subsection{An exact signed inverse for the critical gamma quotient}

We use the usual entire Kummer function ${}_1F_1(a;c;z)$, with
$c>0$.  The next sign lemma includes every positive noninteger $q$;
it does not rely on a numerical zero count.

\begin{lemma}[A forced negative interval]\label{p07:be13:kummer-sign}
For every $q>0$ and $c>0$, the real function
$y(z)={}_1F_1(-q;c;z)$ is strictly negative on a nonempty open
subinterval of $(0,\infty)$.
\end{lemma}
\begin{proof}
Kummer's differential equation \cite[Section~13.2]{p01:rare:DLMF} has the self-adjoint form
\[
 (z^c e^{-z}y'(z))'+qz^{c-1}e^{-z}y(z)=0,
 \qquad y(0)=1.
\]
The left boundary flux is zero, because $c>0$ and $y'$ is bounded
near zero.  If $y$ were positive on the whole positive axis, then
for $z\geq1$
\[
 y'(z)=-q e^z z^{-c}\int_0^z t^{c-1}e^{-t}y(t)\,dt
       \leq-Ce^z z^{-c}
\]
for some $C>0$.  Integrating forces $y$ to become negative, a
contradiction.  At its first zero the displayed integral is
strictly positive, so the derivative there is strictly negative.
Thus the first zero is followed by an open negative interval.
\end{proof}

\begin{theorem}[Critical gamma extraction is possible exactly at $q\leq0$]
\label{p07:be13:gamma-factor}
Assume $\alpha>0$, $M>0$, and $A>0$, and let $b>a>0$.
The generalized McKay law admits a decomposition
\[
 X\ \stackrel{\mathrm d}=\ G+Y,\qquad
 G\sim\operatorname{Gamma}(A,\text{rate }b-a),\qquad
 Y\geq0,\qquad G\perp Y,
\]
if and only if $q\leq0$.  No infinite-divisibility assumption is
made about the possible remainder $Y$.

For $a=1$, a signed inverse Laplace density of
$G_0(s):=s^A J(1+s)$ is
\begin{equation}\label{p07:be13:signed-inverse}
 g(x)=c_\alpha\,\frac{\Gamma(M)}{\Gamma(\alpha)}
 x^{\alpha-1}e^{-2x}{}_1F_1(-q;2\alpha;2x),\qquad x>0,
 \qquad
 c_\alpha=\frac{2^{1/2-\alpha}}{\Gamma(\alpha+1/2)}.
\end{equation}
In the positive case the remainder has density
$e^{-(b-1)x}g(x)/G_0(b-1)$.
\end{theorem}
\begin{proof}
The Euler integral for $I_{\alpha-1/2}$ \cite[Section~10.32]{p01:rare:DLMF} gives
\[
 I_{\alpha-1/2}(x)
 =c_\alpha x^{\alpha-1/2}e^x
 \frac1{B(\alpha,\alpha)}
 \int_0^1e^{-2xt}t^{\alpha-1}(1-t)^{\alpha-1}\,dt.
\]
Consequently
\begin{equation}\label{p07:be13:J-euler}
 J(1+s)=c_\alpha\Gamma(M)(s+2)^{-M}
 {}_2F_1\left(M,\alpha;2\alpha;\frac2{s+2}\right).
\end{equation}
For $s>0$, termwise integration of the Kummer series is justified
absolutely since $2/(s+2)<1$, and gives
\begin{align*}
 &\int_0^\infty e^{-sx}x^{\alpha-1}e^{-2x}
                  {}_1F_1(-q;2\alpha;2x)\,dx\\
 &\hspace{12mm}=\Gamma(\alpha)(s+2)^{-\alpha}
 {}_2F_1\left(-q,\alpha;2\alpha;\frac2{s+2}\right)\\
 &\hspace{12mm}=\Gamma(\alpha)s^A(s+2)^{-M}
 {}_2F_1\left(M,\alpha;2\alpha;\frac2{s+2}\right).
\end{align*}
The last equality is Euler's transformation; its exponent is
$2\alpha-(-q)-\alpha=A$.  Comparison with
\eqref{p07:be13:J-euler} proves \eqref{p07:be13:signed-inverse}.
The classical real-axis Kummer expansion, or its elementary
polynomial formula at nonnegative integral $q$, shows that $g$
is integrable at infinity: its algebraic part is
$O(x^{-A-1})$, and its other part has exponential decay times a
fixed power.  It is integrable at zero since $\alpha>0$.
All exponentially tilted versions are absolutely integrable directly
from the preceding series calculation as well.

The Laplace transform of the putative remainder is necessarily
\[
 \frac{J(b+s)}{J(b)}
       \left(\frac{b-1+s}{b-1}\right)^A
             =\frac{G_0(b-1+s)}{G_0(b-1)}.
\]
Uniqueness of Laplace transforms of finite signed measures makes
its only possible inverse the stated tilted density.  When $q\leq0$
all coefficients of ${}_1F_1(-q;2\alpha;2x)$ are nonnegative and its
constant coefficient is one.  The inverse is positive, has total
mass one, and gives the required independent convolution.
When $q>0$, Lemma~\ref{p07:be13:kummer-sign} makes this inverse negative
on an open interval, ruling out any probability remainder, including
one with atoms.  Rescaling $x$ proves the statement for general $a$.
\end{proof}

The positive half of Theorem~\ref{p07:be13:gamma-factor} also follows
from the endpoint atom in the previously known Thorin measure.
Its new content is the explicit inverse and the obstruction at
$q>0$ for an \emph{arbitrary} nonnegative remainder, a stronger
statement than failure of a GGC or ID remainder.  The obstruction
is independent of the exponential tilt $b/a$.

\subsection{Universal strictness of the SD boundary}

\begin{theorem}[The remaining SD strictness question is closed]
\label{p07:be13:strictness}
For every $q>0$, $\alpha>-1/2$, and $2\alpha+q>0$,
\[
                    1<\beta_{\alpha,q}<\infty.
\]
The supremum is attained at a finite positive point.  In particular,
for every admissible excess-power pair there is a nonempty interval
\[
                     1<b/a<\beta_{\alpha,q}
\]
of generalized McKay laws that are ID but are not SD.  Together with
the earlier GGC classification, this proves the exact dichotomy:
\[
 \text{the entire admissible tilt family is SD}
               \quad\Longleftrightarrow\quad q\leq0.
\]
\end{theorem}
\begin{proof}
The kernel limit already gives $\beta\geq1$.  Suppose $\beta=1$.
Then $k(x)=e^{-x}K(x)$ is nonincreasing, tends to $A$, and satisfies
$k(x)\geq A$.  If $\alpha<0$, this contradicts
$k(0)-A=\alpha$.  If $\alpha=0$, monotonicity and equal endpoint
values force $k$ to be constant, whereas
$k'(0)=K'(0)-K(0)=-M<0$ because $K$ is even.

It remains to treat $\alpha>0$.  For any fixed $b>1$, the L\'evy
spectral function of $X$ is $e^{-bx}K(x)$.  Since $k\geq A$,
subtracting the L\'evy density of
$\operatorname{Gamma}(A,\text{rate }b-1)$ leaves the nonnegative
L\'evy density
\[
             \frac{e^{-(b-1)x}}x\{k(x)-A\}.
\]
It is integrable in the L\'evy sense: near zero $k-A$ is bounded,
and at infinity the exponential tilt and boundedness of $k$ suffice.
The corresponding positive ID variable is an independent remainder
in the gamma decomposition excluded by
Theorem~\ref{p07:be13:gamma-factor}.  This contradiction proves
$\beta>1$ for every remaining parameter, with no restriction to
integer $q$ or to a cut-sign interval.
Finally $K'/K$ is continuous, is zero at the origin, and tends to
one.  Its value above one therefore attains its global supremum
at a finite positive point.  Restoring $a$ gives the stated tilt
boundary and the dichotomy.
\end{proof}

\paragraph{What this resolves and what it does not.}
Theorem~\ref{p07:be13:strictness} supplies the missing strictness in the
negative-cut intervals with $\alpha>1$; it replaces rather than
assumes termination of the finite moment hierarchy.  The proof does
not assert that that particular hierarchy always terminates.
The exact numerical boundary remains the explicit spectral
supremum \eqref{p07:be13:beta-input}; no elementary formula for its
maximizer is claimed.  The full ID classification was proved in the
preceding work and is not counted again.

\subsection{The singular cost of a tilt-uniform HCM power}

Write $R_\mu(z)=I_{\mu+1}(z)/I_\mu(z)$ and, for
$0<\theta<\pi/2$, put
\[
 B_\mu(y,\theta)=
 \frac{\operatorname{Im}\{e^{i\theta}R_\mu(ye^{i\theta})\}}
      {\sin\theta},\qquad
 C_\mu(\theta)=\sup_{y>0}B_\mu(y,\theta).
\]
The exact power test from the preceding papers, recalled in
\cite[Bessel section]{comp:part:06}, is
\begin{equation}\label{p07:be13:power-input}
 X^d\text{ is HCM}\quad\Longleftrightarrow\quad
                  d>2\quad\hbox{and}\quad b/a\geq C_\mu(\pi/d).
\end{equation}
It is independent of the monomial exponent in the density.
For $\mu>1/2$, define the least tilt-uniform power by
\[
 d_{\mathrm{unif}}(\mu)=
 \inf\{d>2:C_\mu(\pi/d)=1\}.
\]
The previous manuscript proved that this number is finite, but did
not determine its singular behavior as $\mu\downarrow1/2$.
Because the HCM class is preserved under positive powers at least
one, feasible values of $d$ form an upper interval.  Equivalently,
the standard HCM composition rule is that $f(x^c)$ is HCM when
$f$ is HCM and $0<c\leq1$; monomial density factors do not affect
this property \cite{p05:fs:Bondesson}.

\begin{theorem}[A logarithm--double-logarithm uniform-power boundary]
\label{p07:be13:uniform-power}
As $\varepsilon\downarrow0$, with $L=\log(1/\varepsilon)$,
\begin{equation}\label{p07:be13:uniform-law}
 d_{\mathrm{unif}}(1/2+\varepsilon)
                    =L+2\log L-2+o(1).
\end{equation}
More precisely, if $d\to\infty$, $\varepsilon\geq0$, and
$\varepsilon e^d/d^2\to\lambda\in[0,\infty)$, then
\begin{equation}\label{p07:be13:crossover}
 e^d\{C_{1/2+\varepsilon}(\pi/d)-1\}
                      \longrightarrow 2(e^{-2}-\lambda)_+.
\end{equation}
Thus the algebraic order perturbation suppresses the exponentially
small half-order resonance at the scale
$\varepsilon\asymp d^2e^{-d}$.
\end{theorem}
\begin{proof}
We give the uniform estimates needed to control the supremum;
pointwise expansions at fixed $y$ would not suffice.

\emph{Step 1: the exact half-order resonance.}
Put $k=\cot\theta$ and $t=2y\sin\theta$.  Since
$R_{1/2}(z)=\coth z-1/z$,
\begin{equation}\label{p07:be13:coth}
 B_{1/2}(y,\theta)-1
       =\frac{\cos t-k\sin t-e^{-kt}}{\cosh(kt)-\cos t}.
\end{equation}
For $0<t\leq\pi$ the numerator is nonpositive, since
$[e^{kt}(\cos t-k\sin t)]'=-(1+k^2)e^{kt}\sin t\leq0$.
For $t\geq\pi$, uniformly as $k\to\infty$,
\[
 B_{1/2}-1=2e^{-kt}(\cos t-k\sin t)
                 +O((1+k)e^{-2kt}).
\]
The first positive maximum of
$e^{-kt}(\cos t-k\sin t)$ occurs at
$t_k=\pi+2\arctan(1/k)$, has trigonometric factor exactly one,
and subsequent positive maxima are smaller by successive factors
$e^{-2\pi k}$.  This follows by differentiating the displayed
function; its derivative is
$e^{-kt}\{-2k\cos t+(k^2-1)\sin t\}$.
Therefore
\begin{align}
 C_{1/2}(\theta)-1
       &\sim2e^{-\pi k-2},\label{p07:be13:half-height}\\
 W(\theta):=\sup_{y>0}y^2(B_{1/2}(y,\theta)-1)
       &\sim\frac{\pi^2k^2}{2}e^{-\pi k-2}.
                                                    \label{p07:be13:weighted-height}
\end{align}
To justify the second supremum, observe that
$y^2=t^2(1+k^2)/4$.  For
$t\geq\pi+k^{-1/2}$ the absolute value of the main term after
multiplication by $t^2$ is
$O(k e^{-\pi k-\sqrt{k}})$, uniformly for large $k$.
On the remaining interval $t^2=\pi^2(1+o(1))$.
The weighted error is exponentially smaller.  This proves
\eqref{p07:be13:weighted-height}, and the point $t=t_k$ attains its
asymptotic value.  At $\theta=\pi/d$,
\[
 \pi k=d+O(d^{-1}),\qquad
 W(\pi/d)=\tfrac12d^2e^{-d-2}(1+o(1)),\qquad
 y_k\sim d/2.
\]

\emph{Step 2: an estimate uniform in order and angle.}
For some fixed $\theta_1<\pi/2$ and $Y_0<\infty$, uniformly for
$y\geq Y_0$, $0\leq\theta\leq\theta_1$, and sufficiently small
$\varepsilon\geq0$,
\begin{equation}\label{p07:be13:order-estimate}
 B_{1/2+\varepsilon}(y,\theta)-B_{1/2}(y,\theta)
      =-\frac{\varepsilon(1+\varepsilon)}{2y^2}
                                      +O(\varepsilon y^{-3}).
\end{equation}
Here zero angle means the continuous extension.  Indeed the standard
Bessel ratio expansion, uniform on a closed sector and on a small
complex order disk around $1/2$, is
\[
 R_\mu(z)=1-\frac{\mu+1/2}{z}
                +\frac{4\mu^2-1}{8z^2}+O(z^{-3}).
\]
Cauchy's estimate in the order parameter bounds the difference of
the two remainders by $O(\varepsilon|z|^{-3})$.  For real orders
this difference is real on the positive real axis.  A Cauchy estimate
in a slightly wider sector therefore bounds the imaginary part
after multiplication by $e^{i\theta}$ by
$O(\varepsilon\theta y^{-3})$.  Dividing by $\sin\theta$ is thus
uniformly legitimate.  The $z^{-1}$ difference contributes a real
number, while the $z^{-2}$ term contributes exactly the term in
\eqref{p07:be13:order-estimate}.  These uniform Bessel asymptotics are
standard; see \cite[Section~10.40]{p01:rare:DLMF}.
On each fixed bounded positive $y$ interval,
$B_{1/2+\varepsilon}(y,\theta)$ tends uniformly as
$(\varepsilon,\theta)\to(0,0)$ to
\[
 D_{1/2}(y)=\coth y-y\operatorname{csch}^2 y<1.
\]
The power series at zero supplies the same strict upper bound near
zero.  In particular no bounded $y$ can cause a failure of the
uniform-power criterion in this limit.

\emph{Step 3: inversion at the least uniform power.}
Fix $\delta>0$ and set
$d_\pm=L+2\log L-2\pm\delta$.  Then
\[
 \frac{2W(\pi/d_\pm)}{\varepsilon}\longrightarrow e^{\mp\delta}.
\]
For $d=d_+$ choose a fixed $Y\geq Y_0$ large enough that the
$O(1/Y)$ relative error in \eqref{p07:be13:order-estimate} is smaller
than the gap between $e^{-\delta}$ and one.  For all $y\geq Y$,
\[
 y^2(B_{1/2+\varepsilon}(y,\pi/d_+)-1)
       \leq W(\pi/d_+)-\frac\varepsilon2(1-O(1/Y))<0
\]
for small $\varepsilon$.  Step 2 excludes the remaining bounded
interval.  The limit as $y\to\infty$ is one, so
$C_{1/2+\varepsilon}(\pi/d_+)=1$.
For $d=d_-$ evaluate at the first positive peak $y_k\sim d_-/2$.
Equations~\eqref{p07:be13:weighted-height} and
\eqref{p07:be13:order-estimate} make the weighted difference strictly
positive.  Hence $d_-$ is infeasible.  The upper-interval property
of feasible powers gives $d_-\leq d_{\mathrm{unif}}\leq d_+$.
Letting $\delta\downarrow0$ proves \eqref{p07:be13:uniform-law}.

\emph{Step 4: the crossover.}
At the same first peak, \eqref{p07:be13:order-estimate} gives
\[
 e^d(B_{1/2+\varepsilon}(y_k,\pi/d)-1)
                         \longrightarrow 2e^{-2}-2\lambda.
\]
The limit at infinity gives the additional lower bound zero.
For the upper bound, choose a fixed large $Y$ for which the order
perturbation in \eqref{p07:be13:order-estimate} is nonpositive.
The bounded interval is excluded as before, and the region
$2y\sin(\pi/d)\leq\pi$ has $B_{1/2}\leq1$.
For any fixed $\eta>0$, the region
$2y\sin(\pi/d)\geq\pi(1+\eta)$ has positive half-order excess
$O(d e^{-(1+\eta)d+O(1/d)})=o(e^{-d})$, and the order perturbation
is nonpositive.  In the intervening region,
\[
 \frac d2(1+o(1))\leq y\leq\frac d2(1+\eta)(1+o(1)).
\]
Using \eqref{p07:be13:half-height} and
\eqref{p07:be13:order-estimate}, the scaled excess there has limsup
at most $2e^{-2}-2\lambda/(1+\eta)^2$.
Take the maximum with zero and then let $\eta\downarrow0$.
This proves \eqref{p07:be13:crossover}.
\end{proof}

\paragraph{Research status.}
The two structurally distinct advances are the gamma-deconvolution
obstruction, which closes universal SD strictness, and the sharp
uniform-HCM critical law.  Formula~\eqref{p07:be13:crossover} refines the
second advance rather than constituting another named problem.
The cited identities and HCM closure rules are classical.  We make
no priority claim for the separate Ismail--Miller conjecture discussed
in the October 2026 Bessel preprint; none of our proofs uses it.

% ===== budget =====
\section{Complete replication-cap curves for analytic populations}
\label{p07:bu13:section}

The preceding manuscript established uncapped analytic-population rates.
Here we determine how much replication is actually necessary to attain
them. The new lower bound works after every adaptive history, including
histories that nearly reveal an individual's endpoint.

An individual has latent $U\in[0,1]$ with density $w$. Conditional on $U$,
its observations are independent uniforms on $[0,2^U]$. A policy uses at
most $B$ scalar observations, with at most $Q$ observations of any one
individual, where $1\le Q\le B$. Write $R_{B,Q}^{\rm ad}$ and
$R_{B,Q}^{\rm fix}$ for minimax expected ordinary-interval Wasserstein
loss, allowing respectively adaptive policies and predetermined designs.
For fixed positive $L,\tau,\nu$, let
\[
 \mathcal W_\phi(L)=\left\{w:\ w\ge\tfrac12,\ \int_0^1w=1,
 \ w\text{ is one-periodic},\quad
 |\widehat w(k)|\le Le^{-\phi(|k|)}\ (k\ne0)\right\}.
\]
Periodicity is a density constraint; the transportation cost does not
identify the two endpoints.

\begin{theorem}[The full cap curve]\label{p07:bu13:cap}
For either $\phi(k)=\tau k^\nu$ or $\phi(k)=\tau e^{\nu k}$, put respectively
\[
 K_B=(\log B)^{1/\nu},\qquad K_B=\log\log B.
\]
As $B\to\infty$, uniformly over integers $1\le Q\le B$,
\begin{equation}\label{p07:bu13:rate}
 R_{B,Q}^{\rm ad}\asymp R_{B,Q}^{\rm fix}\asymp
 \frac1{\sqrt B}\max\left\{\log K_B,\sqrt{K_B/Q}\right\}.
\end{equation}
Constants may depend on the fixed class parameters, but not on $B,Q$.
Consequently the least order of replication cap permitting the uncapped
rate is
\[
 Q_{\rm crit}\asymp K_B/(\log K_B)^2.
\]
This statement concerns a rate threshold; it does not assert a sharp
leading constant or adaptation to an unknown smoothness envelope.
\end{theorem}

For the stretched-exponential class the two terms are
$\log\log B/\sqrt B$ and $(\log B)^{1/(2\nu)}/\sqrt{BQ}$.
For the double-exponential class they are
$\log\log\log B/\sqrt B$ and $\sqrt{\log\log B/(BQ)}$.
We count the two class-specific full curves as two internally formulated
targets. The cap threshold is their consequence, not a third independent
conjecture. The uncapped term was proved previously
\cite{comp:part:06}; it is retained explicitly to make the new comparison
checkable.

\subsection{A design attaining the curve}

Set $a=\log2$. The log maximum of $r$ observations is distributed as
$Y_r=U-\operatorname{Exp}(ar)$. Thus
\[
 A_{r,k}(Y_r)=e^{-2\pi ikY_r}\left(1-\frac{2\pi ik}{ar}\right)
\]
is unbiased for $\widehat w(k)$, with second moment at most
$C\{1+(k/r)^2\}$. Also $Y_1+1/a$ estimates the population mean with
bounded variance. For any two densities in the class, Fourier expansion
of their CDF difference gives
\begin{equation}\label{p07:bu13:modulus}
 W_1(w,v)\le |\mu(w)-\mu(v)|+
 C\left(\sum_{1\le |k|\le K}
 \frac{|\widehat w(k)-\widehat v(k)|^2}{k^2}\right)^{1/2}
 +CL e^{-\phi(K)}.
\end{equation}
Indeed the CDF difference has mean $\mu(v)-\mu(w)$ and nonzero Fourier
coefficients $(\widehat w(k)-\widehat v(k))/(2\pi ik)$; apply Parseval
and $\sum_{k>K}k^{-2}\le C$. This mean term is necessary for interval,
as opposed to circular, transportation.

Use a quarter of the budget to estimate the mean. Divide frequencies
into dyadic bands ending at $k\le K$, and choose $r_k=\min(Q,k)$.
For $n_k$ fresh individuals in this band its contribution to the squared
Fourier error in \eqref{p07:bu13:modulus} is at most $C a_k/n_k$, where
\[
 a_k=k^{-1}\{1+(k/r_k)^2\},\qquad
 S=\sum_{k\text{ dyadic}\le K}\sqrt{r_k a_k}.
\]
Choose $n_k=\lceil B\sqrt{a_k/r_k}/(4S)\rceil$. The cost is at most
$B/4+\sum r_k\le B/4+CK$ and the total variance at most $CS^2/B$.
Moreover
\[
 S\asymp\begin{cases}
 \log(eK),&Q\ge K,\\
 \log(eQ)+\sqrt{K/Q},&Q<K,
 \end{cases}
 \qquad
 S\le C\max\{\log(eK),\sqrt{K/Q}\}.
\]
Project the estimated mean and Fourier vector onto the vectors attained
by $\mathcal W_\phi(L)$, in the norm in \eqref{p07:bu13:modulus}. The class
is compact in the uniform norm by its summable Fourier envelope, so a
measurable approximate minimizer suffices. Projection costs at most a
factor two. Choose a dyadic $K\asymp K_B$ large enough that
$\phi(K)\ge\log B$. Since $K=o(B)$, rounding is harmless and proves
the upper bound in \eqref{p07:bu13:rate}.

\subsection{A boundary-vanishing perturbation family}

Put $\chi(u)=\sin^2(\pi u)$. For $K\le k\le2K$ define
\[
 \Phi_{k,c}(u)=\frac{\chi(u)\sin(2\pi ku)}{2\pi k},\qquad
 \Phi_{k,s}(u)=-\frac{\chi(u)\cos(2\pi ku)}{2\pi k},
 \qquad \varphi_{k,*}=\Phi_{k,*}'.
\]
Both endpoints of each $\Phi$ vanish, so $\int\varphi=0$. The
$d=2(K+1)$ functions $\varphi$ are uniformly bounded and have frequencies
between $K-1$ and $2K+1$. Each frequency receives contributions from at
most a fixed number of them. Consequently, for a sufficiently small
fixed $c>0$, the family
\begin{equation}\label{p07:bu13:family}
 w_\beta=1+\sum_j\beta_j\varphi_j,\qquad
 |\beta_j|\le b_K:=cK^{-2}e^{-\phi(2K+1)}
\end{equation}
lies in $\mathcal W_\phi(L)$ and satisfies $1/2\le w_\beta\le3/2$.
Furthermore
\begin{equation}\label{p07:bu13:phi-bound}
 \|\Phi(u)\|_2\le C\chi(u)/\sqrt K\le C(1-u)/\sqrt K.
\end{equation}
The last boundary factor is what prevents highly informative exceptional
histories from invalidating the cap lower bound.

\begin{lemma}[Information with a cap, after arbitrary histories]
\label{p07:bu13:trace}
For \eqref{p07:bu13:family}, every adaptive policy has Fisher-information
trace at most $CBQ/K$, uniformly in $\beta$.
\end{lemma}
\begin{proof}
We grant the statistician a stronger observation. For each individual
expose a rate-one Poisson process on $[0,\lambda]\times[0,U]$, where
exposure time $\lambda$ may be allocated predictably from the observed
history. Total exposure is
at most $aB$ and each individual receives at most $aQ$. A block of length
$a$ simulates an original observation: use its largest second coordinate
when the block is nonempty, and an independent $-\operatorname{Exp}(a)$
otherwise. The result is $U-\operatorname{Exp}(a)$, the logarithm of a
uniform observation. Thus a lower bound for this stronger experiment
applies to the original one.

After exposure $\lambda$ and largest observed mark $b$ (take $b=0$
if none), the posterior density is proportional to
$w_\beta(u)e^{-\lambda u}\mathbf1_{[b,1]}(u)$. Define
\[
 Z(b)=\int_b^1w_\beta(u)e^{-\lambda u}\,du,\quad
 H(b)=\int_b^1\Phi'(u)e^{-\lambda u}\,du,\quad s(b)=H(b)/Z(b).
\]
Write $D=1-b$. Integration by parts and \eqref{p07:bu13:phi-bound} give
\[
 H(b)=-e^{-\lambda b}\Phi(b)+\lambda\int_b^1e^{-\lambda u}\Phi(u)\,du,
 \qquad \|H(b)\|_2\le CD e^{-\lambda b}/\sqrt K.
\]
On the other hand
$Z(b)\ge c e^{-\lambda b}(1-e^{-\lambda D})/\lambda$,
with its continuous value at $\lambda=0$. Therefore
\begin{equation}\label{p07:bu13:score}
 \|s(b)\|_2\le C(1+\lambda D)/\sqrt K\le C(1+\lambda)/\sqrt K.
\end{equation}
The conditional intensity of a new record at $y>b$ per unit exposure
is $T(y)=Z(y)/Z(b)$, whose parameter score is $s(y)-s(b)$.
Marks below $b$ have parameter-independent intensity and zero score.
Thus the conditional information trace per unit exposure is
\[
 \int_b^1 T(y)\|s(y)-s(b)\|_2^2\,dy
 \le \frac{C(1+\lambda)^2}{K}\int_b^1T(y)\,dy
 \le \frac{C(1+\lambda)}{K}.
\]
The last integral is the posterior mean of $U-b$, at most
$C\min(1,\lambda^{-1})\le C'/(1+\lambda)$, because $w_\beta$ is
bounded above and below. Since $\lambda\le aQ$ and total exposure is
at most $aB$, integration proves the bound. The compensator formula
already includes the information in intervals with no points. Conditional
score increments are martingale differences; hence adaptive choices of
the next individual do not create additional terms. Randomized policies
can be conditioned on their independent random seed.
\end{proof}

\subsection{Turning information into interval transportation}

Use independent cosine-squared priors on $[-b_K,b_K]$. If $J$ is the
joint prior-and-data Fisher matrix, Lemma~\ref{p07:bu13:trace} gives
\begin{equation}\label{p07:bu13:J}
 \operatorname{tr}J\le C\{BQ/K+K^5e^{2\phi(2K+1)}\}.
\end{equation}
We use the following elementary absolute-loss information inequality:
for any vector $v$ and estimator $T$,
\begin{equation}\label{p07:bu13:abs}
 \mathbb E|T-v^T\beta|\ge\tfrac14\sqrt{v^TJ^{-1}v}.
\end{equation}
For clarity, extend the joint density by zero outside the prior cube.
Translating its parameter argument by $hz$ changes total variation by
at most $|h|\sqrt{z^TJz}/2$. The original target and the translated
target differ by $hv^Tz$, and their risks under the respective densities
are equal by change of variables. Integrating their losses over the
common part, then choosing $|h|=(z^TJz)^{-1/2}$, proves the lower bound
$|v^Tz|/(4\sqrt{z^TJz})$. Optimizing $z$ proves \eqref{p07:bu13:abs}.
The prior vanishes quadratically at its boundary, justifying the zero
extension and differentiation.

Here the CDF is exactly $u+\Phi(u)^T\beta$. The one-dimensional identity
$W_1=\int_0^1|F-\widehat F|$ therefore gives Bayes risk at least
$c\int\sqrt{\Phi(u)^TJ^{-1}\Phi(u)}\,du$.
Write $\Phi=\chi v$. The sine-cosine pairs imply
$\|v(u)\|_2^2\asymp K^{-1}$ independently of $u$. Hence
\[
 \sqrt{v^TJ^{-1}v}\ge\frac{\|v\|_2^2}{\sqrt{v^TJv}}.
\]
Weighted Jensen with weight $\chi$, followed by Fourier orthogonality,
shows
\[
 \int_0^1\sqrt{\Phi^TJ^{-1}\Phi}\,du
 \ge \frac{c/K}{\sqrt{\int_0^1\chi(u)v(u)^TJv(u)\,du}}
 \ge\frac{c}{\sqrt{\operatorname{tr}J}}.
\]
Indeed $\int\chi=1/2$, $0\le\chi\le1$, and
$\int vv^T$ is diagonal with entries at most $C/K^2$.
We have proved
\begin{equation}\label{p07:bu13:cap-lower}
 R_{B,Q}^{\rm ad}\ge
 c\{BQ/K+K^5e^{2\phi(2K+1)}\}^{-1/2}.
\end{equation}

Choose $K\asymp K_B$ sufficiently small that
$\phi(2K+1)\le(\log B)/8$. This is possible in both cases: for the
double-exponential envelope one may take
$K=\lfloor(2\nu)^{-1}\log(\log B/(8\tau e^\nu))\rfloor$.
Then $K^5e^{2\phi(2K+1)}\le K^5B^{1/4}=o(B/K)$, uniformly for
$Q\ge1$. Equation~\eqref{p07:bu13:cap-lower} becomes
$R_{B,Q}^{\rm ad}\ge c\sqrt{K_B/(BQ)}$.

\subsection{The uncapped term and completion}

We record the earlier uncapped information argument to specify the
second lower-bound input. Use perturbations
$w_\beta=1+\sqrt2\sum_{k=1}^K(\beta_{k,c}\cos(2\pi ku)+
\beta_{k,s}\sin(2\pi ku))$, with
$|\beta_{k,*}|\le cK^{-2}e^{-\phi(K)}$. They have uniformly bounded
Lipschitz constants. Under the constant baseline the posterior of an
individual is a truncated exponential $q\propto e^{-aru}\mathbf1_{[b,1]}$.
If $q_Y$ is its posterior after the next raw observation, direct integration
gives
\[
 \mathbb E_0\|q_Y-q\|_2^2=a/2,\qquad
 \|q\|_2^2\operatorname{Var}_q(U)\le C.
\]
For verification, putting $D=1-b$, $\lambda=ar$,
$Z_\lambda=(1-e^{-\lambda D})/\lambda$, and $\mu=\lambda+a$,
\[
 \mathbb E_0\|q_Y\|_2^2
 =\frac{1+e^{-\mu D}}{2Z_\lambda}
 +\frac{a}{2Z_\lambda}\int_0^D e^{-\lambda v}
       \{1+e^{-\mu(D-v)}\}\,dv
 =\frac a2+\frac{1+e^{-\lambda D}}{2Z_\lambda}.
\]
Also $\|q\|_2^2=O(\lambda+D^{-1})$ and
$\operatorname{Var}_q(U)=O(\min(D^2,\lambda^{-2}))$.
For $z=\int w_\beta q$, $z_Y=\int w_\beta q_Y$, Bessel's inequality
bounds the conditional score trace by
\[
 C\mathbb E_0\|q_Y-q\|_2^2+
 C\|q\|_2^2\mathbb E_0|z_Y-z|^2\le C.
\]
Here $z,z_Y$ are bounded away from zero and infinity, and the second
expectation is at most $C\operatorname{Var}_q(U)$: use the uniform
Lipschitz bound, an independent coupling of $q_Y$ with $q$, and
$\mathbb E_0q_Y=q$. Adding martingale score increments gives trace $CB$,
even without a cap. The joint matrix thus has trace
$C\{B+K^5e^{2\phi(K)}\}$.

For the centered CDF the parameter vector is
$(\sqrt2\sin(2\pi ku),-\sqrt2\cos(2\pi ku))/(2\pi k)$.
Centering multiplies the integrated error by at most two. Apply
\eqref{p07:bu13:abs}, then average $J$ over simultaneous frequency rotations.
The map $J\mapsto\sqrt{v^TJ^{-1}v}$ is convex, since it is
$\sup_z |v^Tz|/\sqrt{z^TJz}$. Rotation averaging over
$R(u)=\operatorname{diag}(R_{2\pi ku})_{k=1}^K$, where $R_\theta$ is
planar rotation through $\theta$, leaves blocks $j_kI_2$.
Cauchy--Schwarz gives
\[
 \sum_{k=1}^K(k^2j_k)^{-1}
 \ge\frac{(\sum_{k=1}^K k^{-1})^2}{\sum_{k=1}^Kj_k}.
\]
Thus every capped or uncapped procedure has risk at least
$c\log(eK)/\sqrt{B+K^5e^{2\phi(K)}}$. Taking $K\asymp K_B$ with
$\phi(K)\le(\log B)/8$ gives $c\log K_B/\sqrt B$. Combining this
with \eqref{p07:bu13:cap-lower} and the predetermined upper design completes
the proof of Theorem~\ref{p07:bu13:cap}.

\paragraph{What has and has not been transferred.}
The cap penalty comes from a boundary-vanishing high-frequency family
and a conditional Poisson-information bound. It is not obtained by
substituting a smoothness cutoff into an uncapped answer. The result
allows arbitrary adaptive allocations and arbitrarily small admissible
caps. It does not settle unknown-smoothness adaptation, nonperiodic
analytic classes, binomial-mixture budget optimization, or a universal
claim about every analytic inverse problem.
\section{Remaining obstructions and what the new theorems do not close}
\label{p07:end13:section}

Several ambitions that motivated the work remain open. Listing them
explicitly prevents a restricted theorem from being mistaken for a
solution of a broader problem.
\begin{enumerate}
\item \textbf{Prime-zeta fluctuations at growing scales.} We obtain all
      moments for each fixed window. A variance asymptotic as the window
      height tends to infinity, a central limit theorem, an effective
      error in $T$, and universality of pair correlations are not proved.
      The zeros counted are those of $P^{(m)}$ in $\Re s>1$, not those
      of the Riemann zeta function in its critical strip.
\item \textbf{Unrestricted continuous or analytic two-query recovery.}
      The definable obstruction fixes one finite-parameter query family.
      It does not apply to the union of all polynomial degrees or to all
      analytic functions. The unrestricted two-query question in three
      dimensions is not declared solved.
\item \textbf{General matrix designs at the sharp tail boundary.}
      The arbitrary convergent-design construction is scalar. The
      previous fixed-dimensional matrix results for regular sparse
      designs remain distinct; an arbitrary matrix-design extension is
      not inferred by substituting a quadratic form into a scalar
      Blaschke product.
\item \textbf{Robust computation and optimal constants.} The robust
      population estimator uses a finite predictive cover and attains
      minimax order. A practical polynomial-time implementation, sharp
      risk constants, and the best universal exponent constant in the
      coefficient inequality are separate questions. Ordinary MLE
      robustness is not proved.
\item \textbf{Finite-height arithmetic reconstruction.} The spectral
      bridge assumes finitely many frequencies, a known top frequency
      and phase, a quantitative Diophantine inequality, and exact finite
      zero lists. It does not give a finite-height reconstruction
      algorithm for all primes or a new equivalence to the Riemann
      hypothesis. Its height exponent is not asserted optimal.
\item \textbf{Further Bessel boundaries.} The SD strictness theorem
      uses the signed spectral kernel established in the preceding
      manuscripts. It does not prove termination of every earlier
      finite-moment obstruction algorithm, nor an elementary formula
      for the SD maximizer. Rechecking that inherited kernel remains
      part of any external review of the full program.
\item \textbf{Unknown analytic regularity.} The population cap curves
      assume one fixed periodic Fourier envelope. Adaptation to its
      unknown parameters and nonperiodic analytic classes is not
      supplied by the present proof.
\end{enumerate}

The strongest case for external verification is the coefficient theorem:
its exact target is published, its dependence on the classical Hahn
bound is short, and its remaining finite calculation is reproducible
from this TEX file. The gamma-factor obstruction and definable-game
theorem offer separate compact arguments to scrutinize. The larger
statistical and analytic classifications require their own expert
assessment. We regard these as research claims ready for such scrutiny,
not as a certification of historical priority or journal level.
\endgroup
% ===== END SOURCE 07 =====


% ===== BEGIN SOURCE 08: KEIO_Geometry_and_Rigidity_of_Inverse_Observation.tex =====
\clearpage
\part{Geometry and Rigidity of Inverse Observation}
\label{part:08}
\setcounter{section}{0}
\begin{center}\large\itshape Adaptive Width, Blind Zero Recovery, and Sharp Boundaries in Inverse Probability\end{center}
\begin{center}\small Source manuscript: 16 October 2026.\end{center}
\begingroup
\renewcommand{\arraystretch}{1}
\begin{quote}\small\textbf{Source abstract.}\quad
We investigate twelve targets arising from a program on observation and
recovery. The principal geometric result is a universal lower bound
$e_n(X)\ge \operatorname{UW}_{2^n-1}(X)/2$ for adaptive continuous
scalar information, with completely unrestricted selection and decoding
rules. Together with the upper bound of Krieg, Novak, and Ullrich, it
determines the exact number of measurements in every power-of-two
dimension and leaves at most one measurement undecided in other
dimensions. A second geometric transfer quantifies the topological and
polynomial-degree cost of two-query recovery in dimension three.
For binomial populations, an explicit rational linear program attains
the previously established optimal contamination curve without knowing
the contamination level; the same curve holds for adversaries that
inspect the whole clean sample. For prime-supported Dirichlet series,
one-sided exact zeros determine unknown positive coefficient magnitudes
and relative phases, and the zero hull is conjugate to a compact torus
with an explicitly identified dynamical spectrum. Further results close
the arbitrary-design matrix Gaussian tail criterion, retain stable
nonidentifiability inside every positive compact scale interval, resolve
individual spectral parameters from nonlinear minima, and improve
finite-zero experiments by adaptation and calibration. Finally, an exact
gamma extraction boundary exhibits an exponential separation between
arbitrary and infinitely divisible remainders. Classical ingredients,
previous unpublished dependencies, and unresolved stronger questions are
identified throughout. These are twelve scoped targets, including related
refinements, not twelve named historical conjectures.
\end{quote}

\section*{Reading guide and claim accounting}
\addcontentsline{toc}{section}{Reading guide and claim accounting}

This is a sequel to the 223-page integrated observation-and-recovery
manuscript and its subsequent papers, most immediately the 49-page
\emph{Exact Recovery Laws and a Coefficient Conjecture}
\cite{comp:part:07}. We do not re-count that paper's coefficient
conjecture or its statistical and zero-density results. The common
question here is how the geometry of an observation map determines
what can be recovered, at what precision, and with what computation.
The mechanisms differ across the six branches; no theorem asserting
that all six models are equivalent is intended.

The externally posed target is the dimension-dependent lower-bound
question in \cite[Section~3]{p03:an:KNU}. The new bound answers that
request and gives exact counts for $m=2^k$; it does not settle the
optimal count in every dimension. Most other targets are extensions,
gap closures, or new questions formulated within this research program.
In particular, the large-order gamma asymptotic refines the exact
capacity theorem and should not be advertised as a second independent
historical conjecture. A search of the cited primary literature is
not an exhaustive priority determination.

{\small\setlength{\tabcolsep}{4pt}
\begin{longtable}{@{}p{.055\textwidth}p{.31\textwidth}p{.48\textwidth}p{.09\textwidth}@{}}
\toprule
No. & Target & Result and scope & Proof\\
\midrule
\endfirsthead
\toprule No. & Target & Result and scope & Proof\\\midrule
\endhead
1 & Unrestricted adaptive query lower bounds & Logarithmic necessity; exact query count in power-of-two dimensions. External question, partially resolved in general dimension. & \ref{p08:in14:adaptive-width}, \ref{p08:in14:query-count}\\
2 & Cost of a tame first query & Explicit degree--error and fiber-topology obstructions for two queries; nonsharp degree exponent. & \ref{p08:in14:topological-price}, \ref{p08:in14:degree-barrier}\\
3 & Efficient robust binomial recovery & Polynomial-bit-time proper estimator with the previous sharp rate and unknown contamination. & \ref{p08:ml14:computational}\\
4 & Full-sample adversarial corruption & Matching upper and lower rates with a pathwise cap on replacements. & \ref{p08:ml14:strong}\\
5 & Blind prime-series inversion & One-sided zeros determine positive magnitudes and relative phases, up to the necessary scalar ambiguity. & \ref{p08:ar14:thm:blind}\\
6 & Complete zero-process structure & Torus conjugacy, complete outside rigidity, and exact balanced-prime eigenfrequencies. & \ref{p08:ar14:thm:conjugacy}\\
7 & Arbitrary matrix Gaussian designs & Exact design--tail classification for every convergent or divergent harmonic subseries. & \ref{p08:ga14:thm:matrix}\\
8 & Compact stable scales & Sparse failure persists in every positive compact scale interval for fixed index below one. & \ref{p08:ga14:thm:compact-stable}\\
9 & Individual nonlinear spectral parameters & Sharp mass, separation, and noise exponents, with matching examples. & \ref{p08:sp14:thm:parameters}\\
10 & Adaptive noisy-zero experiments & Logarithmically many amplitude settings, improved sufficient height, and top-frequency calibration. & \ref{p08:sp14:thm:adaptive}\\
11 & Maximal probability gamma factor & Exact one-point Kummer criterion on the first excess-power line at the tail rate. & \ref{p08:be14:capacity}\\
12 & Size of the gamma-capacity separation & Exponential unrestricted deficit versus algebraic ID deficit; refinement of target 11. & \ref{p08:be14:asymptotic}\\
\bottomrule
\end{longtable}
}

\paragraph{Evidence and dependencies.}
All twelve rows have proofs in this manuscript, including proofs of the
new reductions. The geometric branch invokes classical width, Betti,
and triangulation theorems; the statistical branch explicitly states the
inverse estimate and lower bounds imported from the previous unpublished
paper. The improved recurrence exponent in the spectral branch is an
application of an established filling theorem. The prime-series inverse
and gamma-capacity arguments do not rely on the earlier unpublished
spectral constructions. Numerical calculations, where mentioned, are
not substitutes for proofs. Automated cross-review found and prompted
editorial and hypothesis clarifications; it does not constitute human
refereeing or machine-checked formal verification.

\paragraph{What is not claimed.}
There is no proof here of the Riemann hypothesis, no recovery of unknown
prime frequencies from nontrivial zeta zeros, no unrestricted two-query
solution in dimension three, and no growing-window prime-zero central
limit theorem. Theorems about exact infinite zero data do not by
themselves provide stable finite-data algorithms. Publication-level
importance and priority require assessment beyond this manuscript.



% ===== information =====
\section{A logarithmic lower bound for unrestricted adaptive continuous information}
\label{p08:in14:section}

We connect adaptive information to the multiplicative dimension bound
for Urysohn widths.  The resulting lower bound applies to arbitrary
continuous measurements, with no definability, polynomial-degree,
Lipschitz, or continuity assumption on the selection rule or decoder.
It answers the request for a dimension-dependent lower bound explicitly
posed by Krieg--Novak--Ullrich \cite[Section~3]{p03:an:KNU}.  It determines
the exact query count when the dimension is a power of two.  The
three-dimensional two-query problem remains unresolved.

\subsection{Width of an adaptive decision tree}

For a nonempty compact metric space $(X,d)$ and integer $d_0\ge0$, let
\[
 \operatorname{UW}_{d_0}(X)=\inf_{\mathcal U}\max_{U\in\mathcal U}
                       \operatorname{diam}U,
\]
where the infimum is over finite relatively open covers of $X$ of
multiplicity at most $d_0+1$.  Multiplicity means the maximum number of
cover members containing the same point.  Distances on subsets are
always inherited from $X$.  In particular, the subscript is the
\emph{dimension} of an approximating complex, not the allowed cover
multiplicity.  This is the usual Urysohn-width convention
\cite[Section~1]{p08:in14:BB}.

We recall the following width inequality.  The special case needed
here also has a short covering proof, which we include to make clear
that singular and disconnected fibers cause no difficulty.

\begin{lemma}[Width of a scalar map]\label{p08:in14:scalar-width}
If $f:X\to\mathbb R$ is continuous and $X$ is compact, then for every
integer $d_0\ge0$,
\begin{equation}\label{p08:in14:scalar-width-eq}
 \operatorname{UW}_{2d_0+1}(X)
       \le \sup_{t\in f(X)}\operatorname{UW}_{d_0}(f^{-1}(t)).
\end{equation}
More generally, for a continuous map to a metrizable space of covering
dimension at most $q$, the subscript on the left can be replaced by
$(q+1)(d_0+1)-1$.
\end{lemma}

\begin{proof}
This is the multiplicative width inequality of Gromov, in the form
proved in \cite[Theorem~2.1]{p08:in14:BB}.  Here is its covering argument.
Choose $w$ strictly above the supremum on the right.  A finite cover of
a compact fiber of mesh less than $w$ and multiplicity at most $d_0+1$
extends to a cover with the same bounds on the inverse image of a
neighborhood of its target point.  Indeed, take a finite closed
shrinking of the fiber cover, then sufficiently small ambient open
neighborhoods of the shrinking sets.  Compactness preserves the mesh
bound and keeps every originally empty $(d_0+2)$-fold intersection empty.
Their union contains the fiber, so it contains the inverse image of a
sufficiently small target neighborhood, again by compactness.

Refine the resulting cover of $f(X)$ to a finite open cover of
multiplicity at most $q+1$.  Over each member use its corresponding
fiber-neighborhood cover.  The combined cover of $X$ has mesh less than
$w$ and multiplicity at most $(q+1)(d_0+1)$.  Let $w$ decrease to the
supremum.  For scalar targets, $q=1$.
\end{proof}

An adaptive $n$-query procedure on $X$ chooses a continuous real
function at each stage, based arbitrarily on all earlier answers.
It outputs a point of a metric space $Z$ containing $X$ isometrically.
The selection rules and the output map may be discontinuous and need
not be measurable.  Write $e_n(X;Z)$ for its optimal worst-case error.
This notation concerns exact measurements and deterministic algorithms.

\begin{theorem}[The adaptive width bound]\label{p08:in14:adaptive-width}
For every compact metric space $X\subset Z$ and every integer $n\ge0$,
\begin{equation}\label{p08:in14:main-width}
                 e_n(X;Z)\ge\frac12
                           \operatorname{UW}_{2^n-1}(X).
\end{equation}
The conclusion remains valid if a query is required to be continuous
only on the compact feasible set determined by the preceding answers.
Thus it also holds under the usual, stronger requirement that every
selected query be continuous on the entire original input space.
\end{theorem}

\begin{proof}
Fix one procedure with worst-case error at most $r$.  We prove by
induction on its maximum remaining depth $n$ that
$\operatorname{UW}_{2^n-1}(X)\le2r$.  With no remaining query the
output is one point $z\in Z$.  All inputs lie in its radius-$r$ ball;
therefore $\operatorname{diam}X\le2r$, and the one-set cover gives
$\operatorname{UW}_0(X)\le2r$.

For $n\ge1$, let $f:X\to\mathbb R$ be the first query.  Each nonempty
fiber $X_t=f^{-1}(t)$ is compact.  After its answer $t$ is fixed, the
remaining procedure on $X_t$ has at most $n-1$ queries and error at most
$r$.  Every later query restricts to a continuous function on the
corresponding compact feasible subset.  The induction hypothesis gives
\[
           \operatorname{UW}_{2^{n-1}-1}(X_t)\le2r
                             \quad(t\in f(X)).
\]
Apply Lemma~\ref{p08:in14:scalar-width} with
$d_0=2^{n-1}-1$.  Its left subscript is
$2(2^{n-1}-1)+1=2^n-1$, as required.  Padding a shorter branch with
constant queries is harmless.  No regularity in the parameter $t$ is
used: the width lemma combines covers whose existence has been proved
separately on each fiber, not the possibly discontinuous choices made
by the original procedure.  Finally take the infimum over procedures.
\end{proof}

\begin{corollary}[Vector-valued rounds]\label{p08:in14:batch-width}
Suppose round $j$ observes a continuous vector with $q_j$ real
coordinates, selected adaptively from the preceding rounds.  For $s$
rounds the worst-case error is at least
\[
 \frac12\operatorname{UW}_{\prod_{j=1}^s(q_j+1)-1}(X).
\]
\end{corollary}
\begin{proof}
Use the general target-dimension version of
Lemma~\ref{p08:in14:scalar-width} in the same induction, starting with
cover multiplicity one at the leaves.  The multiplicities are
multiplied by $q_j+1$ at the successive levels.
\end{proof}

\subsection{Consequences for the open query-count problem}

Let $N(m)$ be the minimum number of adaptive continuous scalar
measurements sufficient to approximate every vector of $\mathbb R^m$
with any prescribed positive uniform error.  One can fix any norm;
finite-dimensional norm equivalence leaves $N(m)$ unchanged.

\begin{theorem}[Logarithmic necessity and exact power-of-two counts]
\label{p08:in14:query-count}
For every integer $m\ge1$,
\begin{equation}\label{p08:in14:query-bounds}
 \left\lceil\log_2(m+1)\right\rceil
       \le N(m)\le\left\lceil\log_2 m\right\rceil+1.
\end{equation}
In particular,
\begin{equation}\label{p08:in14:power-two}
                         N(2^k)=k+1\qquad(k\ge0).
\end{equation}
More quantitatively, on the cube $[-R,R]^m$ with maximum-norm loss,
\begin{equation}\label{p08:in14:cube-plateau}
                    e_n([-R,R]^m;\mathbb R^m)=R
                      \quad\text{whenever }2^n\le m.
\end{equation}
Thus below this threshold adaptivity and unrestricted continuous
nonlinear queries do not improve on the zero-output procedure at all.
\end{theorem}

\begin{proof}
The Lebesgue covering theorem for the cube implies
\[
 \operatorname{UW}_{m-1}([-R,R]^m,\|\cdot\|_\infty)=2R.
\]
For the lower bound, a cover of mesh less than $2R$ cannot have a
member meeting opposite faces in any coordinate; the covering theorem
therefore forces multiplicity at least $m+1$.  For the upper bound,
coordinate projection to $m-1$ coordinates has fibers of diameter
$2R$ and gives the asserted width; equivalently use arbitrarily fine
covers of the projected cube.  Width is nonincreasing in its subscript.
When $2^n\le m$, Theorem~\ref{p08:in14:adaptive-width} consequently gives
error at least $R$, and the zero output gives error at most $R$.

If an $n$-query procedure had bounded uniform error on all of
$\mathbb R^m$ with $2^n\le m$, restricting it to cubes of arbitrarily
large $R$ would contradict \eqref{p08:in14:cube-plateau}.  Hence
$2^n>m$, which is precisely the lower integer bound in
\eqref{p08:in14:query-bounds}.  The upper bound is the construction of
Krieg--Novak--Ullrich \cite[Theorem~1, version~2]{p03:an:KNU}.
At $m=2^k$ the lower and upper bounds both equal $k+1$.
\end{proof}

\paragraph{What is and is not closed.}
The September 2026 version of \cite[Section~3]{p03:an:KNU} explicitly
asks for lower bounds beyond the Borsuk--Ulam bound $N(m)\ge2$.
Theorem~\ref{p08:in14:query-count} supplies a logarithmic lower bound and
closes the exact-count question in every power-of-two dimension.
For a dimension strictly between consecutive powers of two the bounds
in \eqref{p08:in14:query-bounds} differ by one.  In particular they give
$2\le N(3)\le3$: this argument does not solve the two-continuous-query
problem in dimension three.  The width inequality used as an ingredient
is classical; the claimed contribution is its application to an
unrestricted adaptive decision tree and the resulting query bounds.

\subsection{A quantitative cost for topologically tame first measurements}

The general multiplicative-width bound loses information in dimension
three.  A second theorem of Balitskiy--Berdnikov retains the first
homology of the first-query fibers and gives an explicit obstruction
there.  For a compact metric space $X$ and a continuous scalar function
$f$, put
\[
 \beta(f;X)=\sup_{t\in f(X)}\dim_{\mathbb Q}H_1(f^{-1}(t);\mathbb Q).
\]
Call $f$ \emph{triangulable} if there is a homeomorphism from a finite
polyhedron to $X$ under which $f$ is piecewise linear, allowing also a
homeomorphic change of its scalar target coordinate.  The metric on the
polyhedron is the pullback of the metric of $X$.  It need not be the
linear Euclidean metric of a chosen geometric realization.

\begin{theorem}[Two queries require fiber topology]
\label{p08:in14:topological-price}
Suppose $f:X\to\mathbb R$ is triangulable and
$\beta=\beta(f;X)<\infty$.  Every adaptive procedure which uses $f$ as
its first measurement and an arbitrary continuous scalar second
measurement has error at least
\begin{equation}\label{p08:in14:betti-error}
                \frac{\operatorname{UW}_2(X)}{2(2\beta+3)}.
\end{equation}
The decoder and the choice of second query are unrestricted.  For the
Euclidean ball $B_2^3(R)$ this implies the explicit lower bound
\begin{equation}\label{p08:in14:ball-betti}
                            \frac{R}{\sqrt3(2\beta+3)}.
\end{equation}
\end{theorem}

\begin{proof}
The bounded-topology waist theorem
\cite[Theorem~3.16, with one-dimensional target]{p08:in14:BB} gives a fiber
$F=f^{-1}(t)$ such that
\[
              \operatorname{UW}_1(F)
                       \ge\frac{\operatorname{UW}_2(X)}{2\beta+3}.
\]
Its hypothesis is a piecewise linear map from a metric polyhedron;
triangulating $f$ and pulling back the metric verifies that hypothesis
without changing any width, any fiber topology, or any reconstruction
distance.  A homeomorphic scalar-coordinate change does not change the
fibers.  On $F$ the second selected query is one continuous scalar
function.  Either the definition of width by maps to graphs, or the
$n=1$ case of Theorem~\ref{p08:in14:adaptive-width}, shows that its
reconstruction error is at least $\operatorname{UW}_1(F)/2$.
This proves \eqref{p08:in14:betti-error}.

For \eqref{p08:in14:ball-betti}, the ball contains the cube
$[-R/\sqrt3,R/\sqrt3]^3$.  Restriction of covers and the Lebesgue
covering theorem give
$\operatorname{UW}_2(B_2^3(R))\ge2R/\sqrt3$.  Here Euclidean mesh is
at least maximum-norm mesh, so the cube lower bound applies unchanged.
\end{proof}

\begin{theorem}[An explicit polynomial-degree barrier]
\label{p08:in14:degree-barrier}
Let $D\ge1$.  On $B_2^3(R)$, every two-query procedure whose first
query is a real polynomial of degree at most $D$ and whose second
selected query is continuous has worst-case Euclidean error at least
\begin{equation}\label{p08:in14:degree-explicit}
             \frac{R}{\sqrt3\{2D(2D-1)^2+3\}}.
\end{equation}
In particular, for $0<\varepsilon<R/(3\sqrt3)$, error at most
$\varepsilon$ requires
\begin{equation}\label{p08:in14:degree-growth}
 D\ge\left(\frac{R/(\sqrt3\varepsilon)-3}{8}\right)^{1/3}.
\end{equation}
No degree or definability restriction is imposed on the second query.
The exponent $1/3$ in this necessary degree growth is not claimed sharp.
\end{theorem}

\begin{proof}
A continuous semialgebraic real-valued function on a compact
semialgebraic set admits a finite triangulation as a function; this is
the semialgebraic case of Shiota's theorem \cite{p08:in14:Shiota}.
Hence the first query restricted to the ball is triangulable.
Milnor's Betti estimate, including the compact ball truncation proved
in the proof of his Theorem~2 \cite{p08:in14:Milnor}, gives
\[
 \sum_j\dim_{\mathbb Q}H_j\bigl(\{x:p(x)=t\}\cap B_2^3(R);
                              \mathbb Q\bigr)
                       \le D(2D-1)^2.
\]
The estimate also covers singular levels; compact semialgebraic
triangulation identifies the singular and \v{C}ech versions of Betti
numbers here.  An identically constant polynomial has fiber the whole
ball at its attained value, whose first Betti number is zero, and can
be handled separately.  Thus
$\beta(p;B_2^3(R))\le D(2D-1)^2$.  Substitution in
\eqref{p08:in14:ball-betti} proves \eqref{p08:in14:degree-explicit}.
Finally $2D(2D-1)^2\le8D^3$, giving
\eqref{p08:in14:degree-growth}.
\end{proof}

\begin{corollary}[Necessary topology growth for analytic two-query recovery]
\label{p08:in14:analytic-topology}
Let $f:\mathbb R^3\to\mathbb R$ be real analytic.  If a two-query
procedure beginning with $f$ has uniform error at most $\varepsilon>0$
on all of $\mathbb R^3$, then for every $R>0$,
\begin{equation}\label{p08:in14:analytic-growth}
 \sup_t\dim_{\mathbb Q}H_1\bigl(f^{-1}(t)\cap B_2^3(R);\mathbb Q\bigr)
                \ge \frac{R}{2\sqrt3\varepsilon}-\frac32.
\end{equation}
Consequently, a fixed analytic first query whose truncated fibers have
uniformly sublinear first-Betti growth cannot start a two-query
bounded-error recovery procedure.
\end{corollary}

\begin{proof}
The restriction of $f$ to any compact ball is a continuous subanalytic
function and is triangulable by \cite{p08:in14:Shiota}.  Its fiber Betti
numbers are uniformly finite on this ball, since the triangulation is
finite.  Apply \eqref{p08:in14:ball-betti} and rearrange.  If the left side
were $o(R)$, the resulting inequality would fail for large $R$.
\end{proof}

\paragraph{Accounting and limitations.}
There are two principal new transfers here: the universal adaptive
width bound, and the quantitative topological/degree obstruction for
two queries.  The power-of-two counts, vector-valued-round bound,
polynomial-degree estimate, and analytic topology-growth requirement
are consequences with distinct interpretations, not independent
solutions of four unrelated conjectures.  The original width and
Betti-number theorems are explicitly credited.  Neither the exact
asymptotic degree cost, nor arbitrary real-analytic two-query recovery
in dimension three, nor the optimal count in every dimension has been
settled by these arguments.

% ===== mle =====
\section{Polynomial-time recovery under unknown adversarial contamination}
\label{p08:ml14:sec}

The preceding manuscript established the full Huber-contamination curve
for unrestricted binomial populations, but its estimator searched a
potentially exponential finite cover. This section removes that
computational obstruction. A second result proves the same curve when an
adversary observes the entire clean sample before choosing which reports
to replace. The contamination proportion is not supplied to the algorithm.
The two advances concern computation and the corruption model,
respectively; the clean and Huber statistical rates themselves are prior
results \cite{comp:part:07}.

For a probability measure $G$ on $[0,1]$, put
\[
 p_G(j)=\int_0^1\binom tjx^j(1-x)^{t-j}\,dG(x),\qquad 0\le j\le t.
\]
On the count alphabet, $W_{1,t}$ denotes Wasserstein distance for cost
$|i-j|/t$. Set $\log(e/0)=+\infty$ and define
\begin{equation}\label{p08:ml14:rate}
 r_{N,t}(\varepsilon)=\min\left\{1,\frac1t+
 \frac1{\sqrt{t\log(eN)}}+\frac1{\sqrt N}
 +\varepsilon+\frac1{\sqrt{t\log(e/\varepsilon)}}\right\}.
\end{equation}
All unspecified constants in this section are absolute.

\begin{theorem}[Polynomial-time robust population recovery]
\label{p08:ml14:computational}
There is a single algorithm taking the positive integers $N,t$ and $N$ reported
counts, and returning a finitely supported probability measure
$\widehat G$ with rational support and weights, with the following
properties.
\begin{enumerate}
\item Its bit complexity is polynomial in the length of the explicit
sample input, and it does not use a contamination proportion.
\item If the reports are independent with common law
$R=(1-\varepsilon)p_G+\varepsilon Q$, where $Q$ is arbitrary and
$0\le\varepsilon\le1/4$, then
\[
 \sup_{G,Q}\mathbb E_R W_1(\widehat G,G)
       \le C r_{N,t}(\varepsilon).
\]
\end{enumerate}
The rate is minimax optimal. For $t\le N$ the algorithm is a rational
linear program with $O(Nt+t^2)$ variables and polynomially many
constraints. For $t>N$ it returns the empirical distribution of the
normalized reports.
\end{theorem}

\begin{theorem}[Full-sample adversarial replacements]
\label{p08:ml14:strong}
Let $J_1,\ldots,J_N$ be independent with law $p_G$. After seeing all
these counts, an adversary, possibly randomized, may replace at most
$\lfloor\varepsilon N\rfloor$ of them by arbitrary elements of
$\{0,\ldots,t\}$. Write $\mathcal A_{N}(\varepsilon)$ for this class
of replacement rules. Then
\begin{equation}\label{p08:ml14:strongrate}
 \inf_{\widetilde G}\sup_G\sup_{A\in\mathcal A_N(\varepsilon)}
 \mathbb E W_1\bigl(\widetilde G(A(J_1,\ldots,J_N)),G\bigr)
 \asymp r_{N,t}(\varepsilon)
\end{equation}
for $0\le\varepsilon\le1/4$. The algorithm in
Theorem~\ref{p08:ml14:computational} attains the upper bound, simultaneously
for every $\varepsilon$, without knowing the adversary's replacement
budget.
\end{theorem}

Efficient fitting by polynomial approximation, interval discrepancies,
and linear or dynamic programming has substantial precedent
\cite{p08:ml14:CDSS}. Binomial population recovery by optimization is also
established \cite{p08:ml14:TKV,p01:info:VKVK2019}. The contribution here is the
specific simultaneous control of a low-frequency transport metric and
a computationally accessible high-frequency discrepancy. This preserves
the $N^{-1/2}$ term while achieving the full inverse-problem rate.
Neither interval dynamic programming nor polynomial-time linear
programming is being claimed as new.

\subsection{Hahn truncation makes the separating tests computable}

Let $\pi(dx)=dx/(\pi\sqrt{x(1-x)})$ be the arcsine law and let
$\rho_j=p_\pi(j)>0$. Define $h_k(x)=T_k(2x-1)$ and the
conditional-expectation operator
\[
 Qh(j)=\mathbb E_\pi[h(X)\mid J=j].
\]
The classical beta--binomial orthogonal-polynomial calculation
\cite{p05:fm:DKS2008} is recorded in the form needed here.

\begin{lemma}[Uniform low-degree likelihood-ratio approximation]
\label{p08:ml14:truncation}
For every $G$ and $0\le m\le t$ there is a polynomial $P_{G,m}$ of
degree at most $m$ such that the signed vector
$\widetilde p_{G,m}(j)=\rho_jP_{G,m}(j)$ has total mass one and
\begin{equation}\label{p08:ml14:truncerror}
 \sum_{j=0}^t|p_G(j)-\widetilde p_{G,m}(j)|
 \le \sqrt2\,t\exp\{-m^2/(4t)\}\quad(m<t).
\end{equation}
For $m=t$ the approximation is exact. Positivity of the truncated
vector is not asserted or needed.
\end{lemma}
\begin{proof}
The arcsine orthogonality relations are
$\int h_k\,d\pi=0$ and $\int h_kh_l\,d\pi=\frac12\mathbf1_{k=l}$
for positive indices. The beta--binomial posterior shows that $Qh_k$
is a polynomial of degree $k$ in $j$. The self-adjoint operator $Q^*Q$
preserves polynomial degree, and on degree $k$ its leading coefficient
is multiplied by
\[
 \lambda_{t,k}=\prod_{l=0}^{k-1}\frac{t-l}{t+1+l}.
\]
Indeed $\mathbb E[X^k\mid J=j]=(j+1/2)_{\overline k}/
(t+1)_{\overline k}$, and the leading coefficient of
$\mathbb E[J^k\mid X=x]$ is $(t)_{\underline k}$.
Self-adjointness and orthogonality consequently give
$Q^*Qh_k=\lambda_{t,k}h_k$ and
\[
 \sum_j\rho_j|Qh_k(j)|^2=\lambda_{t,k}/2.
\]
Expanding the degree-$t$ polynomial
$\binom tjx^j(1-x)^{t-j}/\rho_j$ in the Chebyshev basis gives the
exact identity
\[
 \frac{p_G(j)}{\rho_j}
   =1+2\sum_{k=1}^t a_k(G)Qh_k(j),\qquad
 a_k(G)=\int h_k\,dG,\quad |a_k(G)|\le1.
\]
Truncate this expansion at $m$. Every retained nonconstant term has
zero $\rho$-mean, proving the mass assertion. Further,
\[
 \log\lambda_{t,k}
 =\sum_{l<k}\log\left(1-\frac{2l+1}{t+1+l}\right)
 \le-\frac{k^2}{2t}.
\]
Cauchy--Schwarz and the triangle inequality therefore bound the
$\ell_1$ tail by
$\sqrt2\sum_{k=m+1}^t e^{-k^2/(4t)}$, which implies the stated bound.
\end{proof}

For an integer $s\ge1$, let $\mathcal I_s$ consist of subsets of
$\{0,\ldots,t\}$ that are unions of at most $s$ intervals of consecutive
integers; the empty set is included. For probability vectors define
\[
 D_s(p,q)=\max_{A\in\mathcal I_s}|p(A)-q(A)|.
\]

\begin{lemma}[A computable surrogate for predictive total variation]
\label{p08:ml14:surrogate}
Suppose the approximation errors in
Lemma~\ref{p08:ml14:truncation} are at most $\beta$ for every mixing law,
and put $s=m+1$. Then
\begin{equation}\label{p08:ml14:tvfrominterval}
 \|p_G-p_H\|_{\rm TV}\le D_s(p_G,p_H)+4\beta.
\end{equation}
Moreover $D_s(p,q)$ is computable using $O(st)$ arithmetic operations,
and its epigraph has a linear-program formulation with $O(st)$
auxiliary variables and constraints.
\end{lemma}
\begin{proof}
The difference of the truncated likelihood ratios is a polynomial of
degree at most $m$. Unless it vanishes identically, its positive entries
on the integer grid form a union $A$ of at most $m+1$ intervals, since a
nonzero polynomial has at most $m$ real roots. This generous bound also
covers isolated zeros separating positive grid entries. The truncated
difference has total mass zero, so its half $\ell_1$ norm equals its
sum over $A$. The two approximation errors change the half $\ell_1$
norm by at most $\beta$ and the sum over $A$ by at most $2\beta$.
This proves even the slightly sharper bound with $3\beta$.

For completeness, interval maximization is a longest-path problem in
an acyclic graph. Process the grid entries from left to right. A state
records the current position, the number of intervals started (at most
$s$), and whether the preceding entry was selected. Selecting a new
entry adds its signed weight; changing from unselected to selected
starts one interval. There are $O(st)$ states and edges. Run this
recurrence for the weights $p_j-q_j$ and for their negatives, and take
the larger terminal value. To express a bound on that value linearly,
introduce one potential per reachable state, fix the initial potential
to zero, and impose
\[
 V_{\rm next}\ge V_{\rm previous}+\text{edge weight},\qquad
 V_{\rm terminal}\le d
\]
for every edge and terminal state. Such potentials exist precisely
when every path has weight at most $d$. Duplicating the construction
for the two signs gives the epigraph $D_s(p,q)\le d$.
\end{proof}

\subsection{The rational linear program and simultaneous fitting}

Assume first $2\le t\le N$; bounded exceptional cases can use the
empirical normalized reports, with a change of the universal constant.
Choose the following integers and rational threshold:
\begin{align*}
 M&=Nt, & L&=\lceil\log_2(8tN)\rceil,\\
 m&=\min\{t,\lceil\sqrt{4tL}\rceil\},
 &s&=m+1,\\
 L_0&=\lceil\log_2(2(t+2))\rceil,
 &b_N&=\frac{\lceil\sqrt{N(2s+1)L_0}\rceil}{N}.
\end{align*}
Integer logarithms and ceiling square roots are computed by integer
comparisons. Let $q$ denote the empirical distribution of the reports.
On the support grid $\{l/M:0\le l\le M\}$ minimize
\begin{equation}\label{p08:ml14:program}
 F(G)=W_{1,t}(p_G,q)+(D_s(p_G,q)-b_N)_+.
\end{equation}
The decision variables are the nonnegative mixing weights, summing to
one. On the ordered count alphabet,
\[
 W_{1,t}(p_G,q)=\frac1t\sum_{j=0}^{t-1}
         \left|\sum_{i=0}^j(p_G(i)-q_i)\right|.
\]
Use ordinary epigraph variables for these absolute values and the
longest-path formulation of Lemma~\ref{p08:ml14:surrogate}, with
$d=b_N+z$ and $z\ge0$, for the remaining term. This is a linear program
with rational coefficients. The binomial probability at grid point
$l/M$ has bit length $O(t\log(M+1))$. There are $M+1=Nt+1$ mixing
variables and $O(st+t)=O(t^2)$ further variables and constraints,
apart from the polynomially sized dense coefficient matrix. Standard
rational linear programming therefore finds an optimal rational vertex
in polynomial bit time. Since $t\le N$ on this branch, the bound is
polynomial in the explicit sample input length. No finite cover is
enumerated and no contamination parameter is used.

\begin{lemma}[Simultaneous transport and variation control]
\label{p08:ml14:joint}
For the estimator obtained from \eqref{p08:ml14:program}, in either the
Huber model or the replacement model of
Theorem~\ref{p08:ml14:strong},
\begin{align}
 \mathbb E W_{1,t}(p_{\widehat G},p_G)
       &\le C(\varepsilon+N^{-1/2}),\label{p08:ml14:predw}\\
 \mathbb E\|p_{\widehat G}-p_G\|_{\rm TV}
       &\le C\left(\varepsilon+
           \frac{t^{1/4}\log(eN)^{3/4}}{\sqrt N}\right).
            \label{p08:ml14:predtv}
\end{align}
The expectations and constants are uniform over all adversaries in the
replacement model.
\end{lemma}
\begin{proof}
First consider replacements, and let $q^0$ be the empirical law of the
clean sample. Deterministically
$\|q-q^0\|_{\rm TV}\le\varepsilon$ and
$W_{1,t}(q,q^0)\le\varepsilon$. Define
\[
 S=\max_{A\in\mathcal I_s}|q^0(A)-p_G(A)|,
 \qquad Z=W_{1,t}(q^0,p_G).
\]
There are at most $(t+2)^{2s}$ interval unions, by specifying ordered
pairs of endpoints and allowing empty padding. Hoeffding's inequality
and our choice of $b_N$ give
\begin{equation}\label{p08:ml14:empirical}
 \Pr\{S>b_N+u\}\le e^{-2Nu^2},\qquad
 \mathbb E(S-b_N)_+\le C/\sqrt N.
\end{equation}
Also $\mathbb E Z\le1/(2\sqrt N)$, by integrating the standard
deviation of each empirical CDF coordinate. Round the latent variable
under $G$ to the nearest grid point to obtain a comparator $G_*$. A
common-uniform coupling of $t$ Bernoulli variables gives
\[
 W_1(G_*,G)\le1/M,\qquad
 \|p_{G_*}-p_G\|_{\rm TV}\le t/M=1/N.
\]
Consequently
\[
 F(G_*)\le 2(\varepsilon+N^{-1})+Z+(S-b_N)_+.
\]
Optimality yields $W_{1,t}(p_{\widehat G},q)\le F(G_*)$ and
$D_s(p_{\widehat G},q)\le b_N+F(G_*)$. The triangle inequalities
now prove \eqref{p08:ml14:predw} and
\[
 \mathbb E D_s(p_{\widehat G},p_G)
     \le C(\varepsilon+b_N+N^{-1/2}).
\]
The choice of $m$ and Lemma~\ref{p08:ml14:truncation} make its error
$\beta\le1/N$, since $L\ge\log(8tN)$. Apply
Lemma~\ref{p08:ml14:surrogate}. Finally,
$m+1\le C\sqrt{t\log(eN)}$ and
$L_0\le C\log(eN)$ give
$b_N\le C t^{1/4}\log(eN)^{3/4}/\sqrt N$.

For independent Huber reports, define $S$ and $Z$ with the empirical
law $q$ and its mean $R$ instead. The same two bounds hold, and
$\|R-p_G\|_{\rm TV},W_{1,t}(R,p_G)\le\varepsilon$ replace the
deterministic comparison with $q^0$. Every other inequality is unchanged.
\end{proof}

\subsection{Transfer to the latent population and sharpness}

We state explicitly the quantitative inverse estimate from the preceding
manuscript \cite{comp:part:07} that supplies this transfer. It holds for
all probability measures $G,H$ and integers $1\le u\le k\le t/2$:
\begin{equation}\label{p08:ml14:inverse}
 W_1(G,H)\le\frac Ck+Ce^{4u^2/t}W_{1,t}(p_G,p_H)
       +\frac{Ce^{2k^2/t}}u\|p_G-p_H\|_{\rm TV}.
\end{equation}
Here is its mechanism, to make clear that no property of the estimator
is implicit in the transfer. Approximate each $1$-Lipschitz test function
by degree-$k$ and degree-$u$ polynomials $P_k,P_u$ with uniform errors
$O(1/k),O(1/u)$ and bounded derivatives. The degree-$t$ Bernstein
coefficients of a degree-$d$ polynomial have norm at most
$Ce^{2d^2/t}$ times its uniform norm for $d\le t/2$. Applied to the
derivative of $P_u$, this bounds the normalized-grid Lipschitz constant
of its coefficients by $Ce^{4u^2/t}$. The coefficients of $P_k-P_u$
have size at most $Ce^{2k^2/t}/u$. Testing these two pieces by
$W_{1,t}$ and TV proves \eqref{p08:ml14:inverse}. The coefficient estimate
follows from the Cauchy coefficient formula for
$(1+z)^tP(z/(1+z))$, the Bernstein complex-polynomial bound, and the
binomial saddle-point estimate; these details are proved in
\cite[Two-scale inversion lemma]{comp:part:07}. In particular, the
new algorithm does not require a stronger coefficient conjecture.

\begin{proof}[Upper bounds in Theorems~\ref{p08:ml14:computational}
and~\ref{p08:ml14:strong}]
For $t\le N$, write $L_N=\log(eN)$ and
$L_\varepsilon=\log(e/\varepsilon)$. In \eqref{p08:ml14:inverse}, take
\[
 u=\lfloor\sqrt t/16\rfloor,\qquad
 k=\left\lfloor\min\{t/2,
        \sqrt{a t\min(L_N,L_\varepsilon)}\}\right\rfloor,
        \qquad a=1/64.
\]
For sufficiently large $t$, these satisfy $1\le u\le k$;
bounded $t$ are harmless because the loss is at most one and the rate
contains $1/t$. The approximation term is bounded by the sum of the
three inverse terms in \eqref{p08:ml14:rate}. By
Lemma~\ref{p08:ml14:joint}, the low-frequency term is
$O(\varepsilon+N^{-1/2})$. The contamination contribution to the
last term is bounded by
\[
 \frac{C\varepsilon e^{2aL_\varepsilon}}{\sqrt t}
       \le \frac C{\sqrt{tL_\varepsilon}},
\]
with the zero-contamination convention. Its sampling contribution is
at most
\[
 \frac{C(eN)^{2a}L_N^{3/4}}{t^{1/4}\sqrt N}
       \le\frac C{\sqrt{tL_N}},
\]
because $t\le N$ and
$(eN)^{2a}N^{-1/4}L_N^{5/4}$ is uniformly bounded. The proof's test
degree depends on $\varepsilon$, but the estimator does not.

When $t>N$, return the empirical law of the normalized reports.
In the replacement experiment it is within $\varepsilon$ in $W_1$
of the clean normalized empirical law. Coupling a binomial count to
its latent parameter costs at most $1/(2\sqrt t)$, while the latent
empirical law has expected $W_1$ error at most $1/(2\sqrt N)$.
Thus the risk is at most $\varepsilon+1/(2\sqrt t)+1/(2\sqrt N)$.
For Huber reports, comparison with $R$ and then $p_G$ gives the same
bound. This is $O(r_{N,t}(\varepsilon))$ when $t>N$.
\end{proof}

\begin{proof}[Lower bound for full-sample replacements]
The identity replacement rule gives the clean lower bound
\[
 c\min\{1,t^{-1}+(t\log(eN))^{-1/2}+N^{-1/2}\}.
\]
For reference, the first term follows from matching the first $t$
moments of two population laws. The second follows by taking
$G_k^\pm=(1\pm h_k/2)\pi$, whose $W_1$ distance is $\asymp1/k$
and whose predictive TV is at most
$(2\sqrt2)^{-1}e^{-k^2/(4t)}$, and setting
$k\asymp\sqrt{t\log(eN)}$. If this degree exceeds $t$, the first
term absorbs it. Endpoint mixtures with weights differing by
$\asymp N^{-1/2}$ give the last term.

Let $q=\lfloor\varepsilon N\rfloor\ge1$ and set $\eta=q/(16N)$.
The TV inverse modulus proved in the preceding manuscript is
\begin{equation}\label{p08:ml14:prior-modulus}
 \sup_{\|p_G-p_H\|_{\rm TV}\le\eta}W_1(G,H)
 \asymp\min\{1,t^{-1}+(t\log(e/\eta))^{-1/2}+\eta\}.
\end{equation}
Its lower bounds are the same Chebyshev perturbations just described,
together with endpoint mass perturbations and exact moment matching;
thus this use does not assume an unproved statistical equivalence.
Choose a pair attaining at least a fixed positive fraction of the
lower bound in \eqref{p08:ml14:prior-modulus}, and independently maximally
couple $N$ observations
from $p_G$ and $p_H$. The number $K$ of unequal coordinates has
expectation at most $N\eta=q/16$, so
$\Pr\{K>q\}\le1/16$. In the $G$ experiment the adversary samples
the coupled $H$ vector conditional on the observed $G$ vector. If
$K\le q$ it makes exactly those replacements; otherwise it leaves the
sample unchanged. This is an admissible randomized adversary for every
realized sample. Its output law is within TV $1/16$ of $p_H^{\otimes N}$.
Compare it to the identity adversary in the $H$ experiment. Le Cam's
two-point inequality gives a lower bound of an absolute positive
constant times $W_1(G,H)$.

Since $q/N\ge\varepsilon/2$ whenever $q\ge1$, equation
\eqref{p08:ml14:prior-modulus} supplies the contamination terms in
\eqref{p08:ml14:rate}, with absolute constant changes inside logarithms.
If $q=0$, then $\varepsilon<1/N$, so both contamination terms are
already bounded by the clean terms. Combining the lower bounds proves
\eqref{p08:ml14:strongrate}. For the Huber theorem the sharp lower bound
was established in \cite{comp:part:07}, using intersection of Huber
neighborhoods at predictive TV at most
$\varepsilon/(1-\varepsilon)$.
\end{proof}

\subsection{What remains separate}

Polynomial time here means a bit-complexity guarantee for a rational
linear program; no claim of nearly linear runtime or demonstrated
large-data practicality is made. The ordinary unmodified nonparametric
MLE has not been proved to attain the robust curves by this argument.
The algorithm uses both terms in \eqref{p08:ml14:program}; replacing it by
pure predictive TV fitting would lose the needed low-frequency control.

The sharper universal coefficient inequality
\[
 |C(t,m,j)|\stackrel{?}{\le}2^{m^2/t}
\]
was also investigated but is not proved here. If $c_*$ denotes the
least exponent constant for $|C(t,m,j)|\le e^{c_*m^2/t}$ uniformly,
the previous constant-one theorem and the degree-$t$ middle coefficient
give only
\[
 \log2\le c_*\le1.
\]
Indeed at even $m=t$ and $j=t/2$, the coefficient magnitude is
$\binom{2t}{t}/\binom{t}{t/2}=2^{t-1/2}(1+o(1))$.
Numerical support for the stronger bound is not a proof and is not
counted as an additional closure.

% ===== arithmetic =====
\section{Blind recovery and complete rigidity of prime-series zero processes}
\label{p08:ar14:sec}

An exact density does not decide whether a zero process forgets its
past.  In the present non-Gaussian model it does not.  We prove two
stronger statements.  First, one half of a zero divisor determines both
the positive coefficient magnitudes and their relative phases, even
when the magnitudes were not supplied to the observer.  Second, the
entire zero process is topologically conjugate to an explicitly
specified compact-group rotation.  This identifies its complete
dynamical spectrum and yields persistent, recurrent correlations of
local zero counts.  No growing-window central limit theorem is inferred
from these statements.

The inverse theorem below concerns a prescribed prime support.  It
does not recover unknown primes from the nontrivial zeros of the
Riemann zeta function.  Its objects are zeros of a prime-supported
Dirichlet series in its half-plane of absolute convergence.

\subsection{The class and its frozen boundary}

Fix a real number $\sigma_0$ and positive numbers $a_p$, indexed by
the primes, such that
\begin{equation}\label{p08:ar14:eq:class}
 \sum_p a_p p^{-\sigma}<\infty\quad(\sigma>\sigma_0),\qquad
 \sum_{p\ge3}\frac{a_p}{a_2}(2/p)^S=1
 \quad\hbox{for some }S>\sigma_0.
\end{equation}
The second equation has at most one solution, because its left side
is strictly decreasing.  For
$x=(x_p)_{p\ge3}\in X:=\prod_{p\ge3}\mathbb T$, put
\[
 f_{a,x}(s)=a_2 2^{-s}+\sum_{p\ge3}a_p x_p p^{-s},\qquad
 D_{a,x}=\sum_{f_{a,x}(\rho)=0}\operatorname{ord}_{\rho}(f_{a,x})
                                      \,\delta_\rho
\]
on $\mathcal H=\{\Re s>\sigma_0\}$.  Thus a common phase has already
been removed.  The unnormalized series
$\sum_p a_p\omega_p p^{-s}$ gives the same divisor with
$x_p=\omega_p/\omega_2$.

The topology on divisors is vague convergence of locally finite
point measures.  The map $x\mapsto f_{a,x}$ is continuous for locally
uniform convergence: on $\Re s\ge\alpha>\sigma_0$, discard a uniformly
small absolutely summable tail and use continuity of the finitely
many remaining coordinates.  No member is identically zero, by its
leading term at $+\infty$.  Consequently
\begin{equation}\label{p08:ar14:eq:divisor-continuity}
 x\longmapsto D_{a,x}
 \quad\hbox{is continuous.}
\end{equation}
For completeness, local uniform convergence to a nonzero analytic
function preserves the numbers of zeros in finitely many small disks
around its zeros, by Rouch\'e's theorem, and gives no other zeros on
the remaining compact set.  This proves vague convergence, including
multiplicities.  The same argument applies after restriction to any
open subdomain of $\mathcal H$.

\begin{lemma}[Boundary recurrence and a coefficient-free decoder]
\label{p08:ar14:lem:decoder}
Every zero of $f_{a,x}$ lies in $\Re s\le S$.  In each of the two
directions $\Im s\to\pm\infty$, there are zeros whose real parts
tend to $S$.  If $\rho_n=\sigma_n+it_n$ is any sequence of zeros with
$\sigma_n\to S$, then, for each prime $p\ge3$,
\begin{equation}\label{p08:ar14:eq:decoder}
 x_p=-\lim_{n\to\infty}\exp\!\bigl(it_n\log(p/2)\bigr).
\end{equation}
This decoder uses the known frequencies and the observed right edge,
but not the unknown $a_p$.
\end{lemma}

\begin{proof}
Divide the series by $a_2 2^{-s}$ and write
$b_p(\sigma)=(a_p/a_2)(2/p)^\sigma$.  At a zero,
\[
 1+\sum_{p\ge3}b_p(\sigma)x_p e^{-it\log(p/2)}=0.
\]
The triangle inequality gives $\sum b_p(\sigma)\ge1$, and hence
$\sigma\le S$.  Taking real parts gives the exact nonnegative identity
\begin{equation}\label{p08:ar14:eq:defect}
 \sum_{p\ge3}b_p(\sigma)
       \left|x_p e^{-it\log(p/2)}+1\right|^2
       =2\left(\sum_{p\ge3}b_p(\sigma)-1\right).
\end{equation}
As $\sigma\to S$, the right side tends to zero and each individual
weight tends to a positive number.  Equation~\eqref{p08:ar14:eq:decoder}
follows.

For existence, use the frozen point $x_p^*=-1$.  The real function
$f_{a,x^*}$ vanishes at $S$, and
\begin{equation}\label{p08:ar14:eq:edge-derivative}
 f_{a,x^*}'(S)=
 D_a:=\sum_{p\ge3}a_p p^{-S}\log(p/2)>0.
\end{equation}
The derivative series converges since $S>\sigma_0$.  The zero is
simple.  For every small $\varepsilon>0$, a disk about $S$, contained
in $\mathcal H$ and of radius less than $\varepsilon$, has one zero
of every $f_{a,u}$ for $u$ in a suitable open neighborhood of $x^*$.
This follows from uniform convergence and Rouch\'e's theorem.

The flow
\begin{equation}\label{p08:ar14:eq:flow}
 (\Phi_t x)_p=x_p e^{-it\log(p/2)}
\end{equation}
is minimal on $X$.  Indeed, every finite set of the relative
frequencies is rationally independent: an integer relation among
$\log(p/2)$ would give an equality of a product of odd-prime powers
and a power of $2$, so all exponents vanish.  The finite-dimensional
Kronecker theorem therefore gives a dense orbit on every finite
coordinate projection, hence on $X$.  Returns to a nonempty open
set occur in both time directions with bounded gaps.  One elementary
proof of the last assertion covers the compact group by finitely
many translates of that set along a dense orbit; translating this
finite cover works from every starting point.

Finally,
$f_{a,x}(s+it)=2^{-it}f_{a,\Phi_t x}(s)$.
The preceding returns transfer the zero near $S$ to arbitrarily
large positive and negative heights of $f_{a,x}$.  Choosing radii
tending to zero proves the assertion.  At most one such zero can
lie exactly on $\Re s=S$: equality in
\eqref{p08:ar14:eq:defect} at two different heights would make both
$\log(3/2)$ and $\log(5/2)$ integer multiples of a common nonzero
frequency, contrary to their rational independence.  Thus, if desired,
all sufficiently chosen zeros approach the edge from its left.
\end{proof}

\subsection{A blind inverse theorem from one half of the divisor}

\begin{theorem}[Unknown magnitudes and phases from one-sided zeros]
\label{p08:ar14:thm:blind}
Let $a,b$ satisfy \eqref{p08:ar14:eq:class} on the same half-plane, and
let $x,y\in X$.  Suppose their zero divisors agree on
\[
 \{s:\Re s>\sigma_0,\ \Im s<A\}
\]
for some finite $A$.  Then
\begin{equation}\label{p08:ar14:eq:blind-result}
 b_p=c a_p\quad\hbox{for every prime }p,\qquad y=x,\qquad c>0.
\end{equation}
The same conclusion holds with $\Im s>A$ in place of $\Im s<A$.
More generally, agreement need only hold on a common open vertical
strip containing the two right edges and on one of these half-planes.
In unnormalized notation, two such series with the same one-sided
divisor differ by a nonzero constant.
\end{theorem}

\begin{proof}
The supremum of the real parts of the one-sided observed zeros is
the corresponding edge, by Lemma~\ref{p08:ar14:lem:decoder}.  Thus the
two edges coincide, say at $S$.  Applying the coefficient-free decoder
to their common zeros gives $x=y$.

For any $u\in X$, minimality supplies a sequence $t_n\to-\infty$
with $\Phi_{t_n}x\to u$.  On each compact subset of $\mathcal H$,
the translated original divisors agree for all sufficiently large
$n$, since its translates then lie below $A$.  By
\eqref{p08:ar14:eq:divisor-continuity}, their limits agree.  Hence
\begin{equation}\label{p08:ar14:eq:all-phase-equality}
 D_{a,u}=D_{b,u}\quad\hbox{for every }u\in X.
\end{equation}
In the strip version this equality is needed and obtained only on
that strip, which still contains a disk about $S$.

Fix an odd prime $p$ and change only its frozen phase, setting
$u_p(\theta)=-e^{i\theta}$ and $u_q(\theta)=-1$ for $q\ne p$.
The implicit-function theorem at the simple frozen zero gives a
unique analytic zero $\rho_{a,p}(\theta)$ near $S$.  Its derivative is
\begin{equation}\label{p08:ar14:eq:root-velocity}
 \rho_{a,p}'(0)=i\frac{a_p p^{-S}}{D_a}.
\end{equation}
There is an analogous zero for $b$.  Equation
\eqref{p08:ar14:eq:all-phase-equality}, applied for real $\theta$ near
zero, identifies the two zeros and thus their derivatives.  It follows
that $b_p/a_p=D_b/D_a=:c$ for every odd prime.  The common edge equation
then gives $b_2=c a_2$ as well.  This proves
\eqref{p08:ar14:eq:blind-result}.  Positive heights are identical with the
sign of the returning sequence reversed.
\end{proof}

\begin{theorem}[The support of the zero law identifies the magnitudes]
\label{p08:ar14:thm:law}
Let $\mu_a$ be the law of $D_{a,x}$ for Haar-distributed $x\in X$.
For two admissible positive coefficient sequences, the following are
equivalent:
\begin{enumerate}
\item $b_p=c a_p$ for all primes and some $c>0$;
\item $\mu_a=\mu_b$;
\item $\operatorname{supp}\mu_a=\operatorname{supp}\mu_b$.
\end{enumerate}
Thus the full zero law, and even its topological support, determines
every ratio $a_p/a_2$.  Knowledge of only the mean zero intensity is
not claimed to suffice.
\end{theorem}

\begin{proof}
The image $\mathcal Z_a=\{D_{a,x}:x\in X\}$ is compact by
\eqref{p08:ar14:eq:divisor-continuity}, and it is exactly the support of
$\mu_a$.  To check the latter equality, a neighborhood of any image
point has a nonempty open preimage in $X$, hence positive Haar measure.
If the two supports agree, choose any divisor in their common image.
Theorem~\ref{p08:ar14:thm:blind}, even with the full divisor available,
gives proportional coefficients.  The other implications are immediate.
\end{proof}

There is also an explicit reconstruction within the recovered graph.
The decoder labels every divisor in $\mathcal Z_a$ by its phase $x$.
In that graph let $\rho_p(\theta)$ be the frozen zero with only the
$p$-th odd phase changed, and let $\rho_{\rm all}(\theta)$ be the
frozen zero when every odd phase equals $-e^{i\theta}$.  Uniform
absolute convergence permits the latter simultaneous change.  Then
\begin{equation}\label{p08:ar14:eq:constructive}
 \boxed{\displaystyle
 \frac{a_p}{a_2}=(p/2)^S
          \frac{\rho_p'(0)}{\rho_{\rm all}'(0)}\quad(p\ge3).}
\end{equation}
Indeed, the denominator is
$i\sum_{q\ge3}a_q q^{-S}/D_a=i a_2 2^{-S}/D_a\ne0$.
This construction uses the exact support, including its limit points;
it is not an estimator from finitely many observed zeros.

For $a_p=(\log p)^m$, $m\ge0$, condition
\eqref{p08:ar14:eq:class} holds on $\Re s>1$, because the edge sum
diverges as $\sigma\downarrow1$ and tends to zero as
$\sigma\to\infty$.  Consequently the stationary zero laws of the
fixed derivatives $P^{(m)}$ are pairwise distinct.  Equality of two
such laws would force $(\log p/\log2)^m=(\log p/\log2)^n$ for
every odd prime, and hence $m=n$.  This assertion concerns the entire
law; it does not prove that their scalar density constants are distinct.

\subsection{Complete spatial rigidity and the exact dynamical spectrum}

\begin{theorem}[A compact torus is encoded in the zeros]
\label{p08:ar14:thm:conjugacy}
For a fixed admissible $a$, the map $x\mapsto D_{a,x}$ is a
homeomorphism from $X$ onto $\mathcal Z_a$.  Under vertical
translation of zero configurations it conjugates the flow
\eqref{p08:ar14:eq:flow}.  The same injectivity, and a continuous inverse
onto its image, hold for the restriction of the divisor to either
one-sided open half-plane, or to a one-sided open strip containing
$S$.

In particular, for every bounded Borel region $B\Subset\mathcal H$,
the complete configuration in $B$ is a measurable function of the
configuration outside $B$.  Its conditional law given the outside
configuration is a point mass.  In fact the entire configuration is
already determined by its sufficiently remote past.

With Haar phases, the dynamical system is ergodic and has simple pure
point spectrum.  Its additive eigenfrequency group is exactly
\begin{equation}\label{p08:ar14:eq:spectrum}
 \begin{aligned}
 \Lambda_0&=
 \left\{\sum_{p\ge3}k_p\log(p/2):
              k_p\in\mathbb Z\hbox{ and finitely many }k_p\ne0\right\}
 \\
 &=\{\log(r/s):r,s\in\mathbb N,\ \Omega(r)=\Omega(s)\}.
 \end{aligned}
\end{equation}
Here $\Omega(n)$ counts prime factors with multiplicity.  Every
frequency displayed occurs, not merely a subset of them.
\end{theorem}

\begin{proof}
Injectivity follows already from Lemma~\ref{p08:ar14:lem:decoder};
continuity was proved above.  A continuous injection from a compact
space to a Hausdorff space is a homeomorphism onto its image.  The
same proof uses the open restricted domains and their one-sided
decoders.  The covariance identity
$f_{a,x}(s+it)=2^{-it}f_{a,\Phi_t x}(s)$ proves conjugacy.

For the outside assertion choose $A$ below all heights in $B$.
Restriction to $\Im s<A$ is measurable from the outside data and
its continuous inverse recovers $x$.  The full divisor, and hence its
restriction to $B$, is then a measurable function of $x$.

The characters
$\chi_k(x)=\prod_{p\ge3}x_p^{k_p}$ form a complete orthonormal basis
of $L^2(X)$: cylinder functions are dense, and finite-dimensional
Fourier polynomials are dense on each cylinder.  Under the flow,
\[
 \chi_k(\Phi_t x)=e^{-it\lambda_k}\chi_k(x),\qquad
 \lambda_k=\sum_{p\ge3}k_p\log(p/2).
\]
Unique factorization makes $k\mapsto\lambda_k$ injective.  In
particular, the only invariant character is constant, proving
ergodicity, and every eigenspace is one-dimensional.  Completeness
proves that there is no continuous spectral component.  Conjugacy
transfers every character to a measurable observable of the zeros,
so none of the displayed eigenfrequencies is lost.

To obtain the last expression in \eqref{p08:ar14:eq:spectrum}, insert
$k_2=-\sum_{p\ge3}k_p$.  Splitting the positive and negative
exponents produces integers $r,s$ with equal total prime-factor
counts.  Conversely, the prime exponents of any such ratio give
these $k_p$.  Cancellation of a common divisor preserves the equality
of the two counts.  This proves the exact formula.
\end{proof}

The restriction in \eqref{p08:ar14:eq:spectrum} is substantive.  For
example, $\log3$ is absent, whereas $\log(3/2)$ occurs.  The zero
process forgets the common phase, and this removes all unbalanced
prime characters.  It does not forget any further phase information.
The spectrum in this theorem is the spectrum of all square-integrable
observables of the process; it is not a claim that each local counting
observable has a nonzero coefficient at every listed frequency.

\begin{proposition}[Exact recovery after sparse deletions]
\label{p08:ar14:prop:deletions}
Fix $a$ and retain only some of the zeros below height $0$.  If the
number deleted between heights $-T$ and $0$ is $o(T)$, then the retained
zeros still determine $x$ and the full divisor.  The same conclusion
holds if one directly assumes that the retained zeros contain a
sequence approaching $S$.
\end{proposition}

\begin{proof}
For each fixed $\varepsilon>0$, the Rouch\'e neighborhood and
bounded-gap returns in Lemma~\ref{p08:ar14:lem:decoder} give at least
$c_\varepsilon T-O(1)$ distinct zeros in
$S-\varepsilon<\Re s\le S$, $-T<\Im s<0$.  To see distinctness,
use separated returning windows, with spacing greater than the
fixed disk diameter.  Deleting $o(T)$ points cannot remove all these
zeros, or all but finitely many of them.  A diagonal choice gives a
retained sequence whose real parts tend to $S$.  The decoder then
recovers $x$.  Multiplicities do not affect the argument, since the
Rouch\'e disk has one simple zero.  No inserted false zeros are allowed
in this statement.
\end{proof}

\subsection{What the classification says about fluctuations}

For a real, centered $G\in L^2(\mu_a)$, identify it with a function
on $X$ and write
$G=\sum_{k\ne0}g_k\chi_k$ in $L^2$.  Orthogonality gives the
absolutely and uniformly convergent covariance expansion
\begin{equation}\label{p08:ar14:eq:covariance}
 C_G(t):=\mathbb E[G(D)G(D-it)]
       =\sum_{k\ne0}|g_k|^2e^{it\lambda_k}.
\end{equation}
Here $D-it$ is the vertically shifted configuration, with the sign
chosen consistently with \eqref{p08:ar14:eq:flow}; replacing $t$ by $-t$
would give the same real covariance.  In particular,
\begin{align}
 \lim_{H\to\infty}\frac1H\int_0^H C_G(t)\,dt&=0,
                                      \label{p08:ar14:eq:cov-mean}\\
 \lim_{H\to\infty}\frac1H\int_0^H|C_G(t)|^2\,dt
       &=\sum_{k\ne0}|g_k|^4>0\quad\hbox{if }G\ne0.
                                      \label{p08:ar14:eq:cov-energy}
\end{align}
Furthermore, there are $t_j\to\infty$ for which
$C_G(t_j)\to C_G(0)=\operatorname{Var}G$.
To prove these facts, first truncate the absolutely convergent sum.
Nonzero distinct frequencies have zero time average; this proves the
two limits, with the tails controlled uniformly.  Simultaneous
approximation of the finitely many remaining phases makes their
cosines tend to one along arbitrarily large times.  An increasing
sequence of truncations proves the recurrence statement.  Thus the
stationary zero process is not weakly mixing and does not have
decaying correlations for every local observable.

This last assertion can be witnessed by actual local zero counts,
not just by a decoder depending on the whole divisor.  Choose
$\sigma_0<\alpha<S<\beta$ so that the mean zero measure has no mass
on $\Re s=\alpha$.  Such $\alpha$ exist outside a countable set.
There are no zeros on $\Re s=\beta$.  On this fixed strip, counts
in bounded-height windows are uniformly bounded over $x$: otherwise
compactness, locally uniform convergence, and Rouch\'e's theorem
would give unboundedly many zeros of a nonzero limit on a compact
neighborhood.  Stationarity therefore gives a finite intensity
$0<c_{a,\alpha}<\infty$, with positivity from the frozen-zero
neighborhood.  For any $0<h<1/c_{a,\alpha}$, the count
\[
 X_h(x)=D_{a,x}\bigl((\alpha,\beta)\times(0,h)\bigr)
\]
has mean strictly between zero and one and is an integer-valued
bounded random variable, hence has positive variance.  It obeys
\eqref{p08:ar14:eq:covariance}--\eqref{p08:ar14:eq:cov-energy} and recurrent
maximal covariance.  Horizontal boundary lines have zero mean mass:
by stationarity such an atom at one height would occur at every
height, contradicting local finiteness of the mean measure.

For an ordinary, fixed-phase prime series the same covariance is
an arithmetic translation average.  The finite torus Fourier
argument proves unique ergodicity of \eqref{p08:ar14:eq:flow}.  A bounded
function with a Haar-null discontinuity set has the same time
average on every orbit: sandwich its upper and lower regularizations
by continuous functions whose integrals differ arbitrarily little.
The count $X_h$ is discontinuous only when a zero meets the boundary
of the rectangle, a Haar-null event by the preceding choice of its
edges.  Its product with a fixed translate has the same property.
Consequently, for every fixed real $\tau$ and every $x\in X$,
\begin{equation}\label{p08:ar14:eq:arithmetic-covariance}
 \lim_{T\to\infty}\frac1T\int_0^T
 [X_h(\Phi_t x)-c_{a,\alpha}h]
 [X_h(\Phi_{t+\tau}x)-c_{a,\alpha}h]\,dt=C_{X_h}(\tau).
\end{equation}
Thus persistent correlations apply to the actual translates of
$P^{(m)}$, as well as to a formal random model.  The order of limits
is explicit: $T\to\infty$ first for each fixed lag, followed by a
study of that limiting covariance as the lag grows.

Finally, for the integrated centered observable
$A_H=\int_0^H G(D-it)\,dt$, Parseval gives the exact identity
\begin{equation}\label{p08:ar14:eq:integrated-variance}
 \operatorname{Var} A_H
   =\sum_{k\ne0}|g_k|^2
        \frac{4\sin^2(H\lambda_k/2)}{\lambda_k^2}.
\end{equation}
The equality follows first for Fourier polynomials and then in
$L^2$, since integration over $[0,H]$ is a bounded operator of norm
at most $H$.  In particular $\operatorname{Var} A_H=o(H^2)$ by
dominated convergence.  Determining a sharper growth law requires
the actual small-frequency weights $|g_k|^2$ near zero.  A dense pure
point frequency group can accumulate at zero, so the absence of a
continuous spectral component alone neither rules out a central
limit theorem nor proves bounded number variance.

\subsection{Relation to earlier work and unresolved questions}

The qualitative freezing of Euler-product prime phases near the
rightmost critical-point edge goes back to Arias de Reyna and
van de Lune \cite{p02:zp:AriasLune}.  The preceding manuscripts developed
its quantitative prime-series versions.  We use the elementary
nonnegative defect behind that mechanism, but the blind theorem
also recovers \emph{unknown positive magnitudes} by differentiating
the phase-labeled zero graph.  That additional inverse step, the
one-sided formulation, and the exact identification of the zero
system are the claims made here.

There is a different and much more developed rigidity theory for
Gaussian analytic functions.  Sodin's form of Calabi rigidity
\cite{p08:ar14:Sodin} identifies Gaussian zero laws from their first
intensity in the appropriate setting.  Our coefficients have fixed
magnitudes and independent uniform phases, and the proof above uses
the full topological support of the zero law, not its first intensity.
Feldheim's stationary Gaussian zero results
\cite{p02:zp:FeldheimDensity} likewise concern a different probability
structure.  Favorov's pure-point Fourier results
\cite{p08:ar14:Favorov} concern real zero sets of Dirichlet series with
bounded spectrum.  The dynamical spectrum in
\eqref{p08:ar14:eq:spectrum} is not identified with the distributional
Fourier transform of an individual unweighted zero measure.

These are internally formulated inverse and fluctuation problems;
we do not identify them with a named external open conjecture.  The
following stronger requests remain open here: recovery from first
intensity alone; quantitative stable blind recovery from noisy finite
zero data; a growing-window variance asymptotic for the endpoint
prime-zeta count; and a central limit theorem or a rigorous contrary
normalization theorem.  In particular, recurrence of local
correlations is not a proof that every possible central limit theorem
fails.

% ===== gaussian =====
\section{Every sparse matrix design at the quasi-analytic tail boundary}
\label{p08:ga14:sec:matrix}

The previous two manuscripts left a specific obstruction: a scalar
Blaschke denominator, composed with an inverse quadratic form, acquires
pole hypersurfaces inside the matrix contour deformation. The following
construction removes that obstruction. The numerator still depends on
the quadratic form, but its compensating denominator depends on the
\emph{individual matrix eigenvalues}. This distinction permits arbitrary
convergent harmonic subseries, without regularity of their counting
function.

Conditionally on $(A,\Sigma)\in\mathbb R^d\times\mathbb S_d^+$,
let $X_j$ be independent $N_d(A,\Sigma)$ vectors and put
\[
 \mathcal O_n=(X_1,\ldots,X_{d-1},X_1+\cdots+X_n),\qquad n\ge d.
\]
Only the separate population laws of these experiments are given.
For sufficiently large $E$, define
\[
 L_1(v)=\log(E+v),\quad L_{j+1}(v)=\log L_j(v),\qquad
 W_{m,a}(v)=\frac{\sqrt v}
 {L_1(v)\cdots L_{m-1}(v)L_m(v)^a}.
\]
Let $\mathcal C_{m,a}^{(d)}$ be the class of mixing laws satisfying
$\mathbb E e^{cW_{m,a}(\operatorname{tr}\Sigma)}<\infty$ for some
$c>0$, with no assumption on the mean marginal or its dependence on
the covariance.

\begin{theorem}[Complete matrix design--tail classification]
\label{p08:ga14:thm:matrix}
Fix $d,m\ge1$, $a\ge0$, and an infinite set
$\mathcal N\subset\{d,d+1,\ldots\}$. The selected experiments
identify every mixing law in $\mathcal C_{m,a}^{(d)}$ if and only if
\begin{equation}\label{p08:ga14:eq:classification}
 a\le1\quad\hbox{or}\quad\sum_{n\in\mathcal N}\frac1n=\infty.
\end{equation}
If both conditions fail, there are distinct absolutely continuous
witnesses with identical complete mean and covariance marginals and
identical finite mixed moments of every order. Their covariance tails
satisfy the defining exponential integrability for \emph{every} $c>0$.
Any prescribed lower bound $\Sigma\succeq b_0I$, $b_0>0$, can be
imposed.
\end{theorem}

The positive directions and the negative direction for regular sparse
matrix designs are inherited. The new statement is the matrix converse
for \emph{every} convergent harmonic subseries. This closes an internal
gap, rather than a named external conjecture. Classical quasi-analytic
moment theory supplies the positive tail criterion
\cite{p05:fg:deJeu}; general differential-operator obstructions to
infinite-mixture identification are discussed in \cite{p01:mix:Le2026}.

\subsection{A denominator on the spectrum rather than the quadratic form}

Let $a_j>0$, $a_j\to\infty$, and
$\Lambda=\sum_j a_j^{-1}<\infty$. Put
\begin{equation}\label{p08:ga14:eq:products}
 E_0(z)=\prod_j(1-z/a_j),\quad
 L(z)=\sum_j\log(1+z/a_j),\quad
 J(R)=\frac12\sum_j\log(1+R^2/a_j^2).
\end{equation}
The product is entire and the principal logarithm sum is holomorphic
off the negative real axis. Both converge locally uniformly on their
respective domains. We use the elementary bounds
\begin{align}
 \log|E_0(z)|&\le J(|z|)+\Lambda(-\operatorname{Re}z)_+,
       \label{p08:ga14:eq:entirebound}\\
 \operatorname{Re}L(z)&\ge J(|z|)
       &&(\operatorname{Re}z\ge0),\label{p08:ga14:eq:logpositive}\\
 \operatorname{Re}L(z)&\ge J(|z|)-2\Lambda(-\operatorname{Re}z)_+
       &&(|\operatorname{Re}z|\le |z|/2).
       \label{p08:ga14:eq:logsector}
\end{align}
For the first bound, square each factor and use
$\log(1+x)\le x$ with
$x=2a_j(-\operatorname{Re}z)_+/(a_j^2+|z|^2)$.
The second follows directly from the squared factors of $1+z/a_j$.
For the third, the possible negative correction inside the logarithm
has magnitude at most $1/2$; use $\log(1-x)\ge-2x$ and sum.
In particular, no quantitative density bound on the $a_j$ is used.

For large $R$ let
\[
 P_m(R)=\prod_{j=1}^m\log_jR,\qquad
 D(R)=P_m(R)(\log_mR)^p,\qquad p>0.
\]
Choose $K$ sufficiently large and define
\[
 J_1^*(w)=\log(K+w),\quad
 J_{j+1}^*(w)=\log(K+J_j^*(w)),\qquad
 Q(w)=w/J_m^*(w)^p.
\]
Principal logarithms are used. These functions are holomorphic off
the negative real axis, with no additional cuts or denominator zeros.
The following sector facts will be used:
\begin{equation}\label{p08:ga14:eq:Qsector}
 \operatorname{Re}Q(w)\ge0\quad(\operatorname{Re}w\ge0),\qquad
 \operatorname{Re}Q(w)\ge \frac{c|w|}{D(|w|)}
 \quad\left(|w|\ge R_0,\quad
 |\arg w|\le\frac\pi2+\frac{c}{P_m(|w|)}\right).
\end{equation}
For clarity, these estimates follow from
\[
 |J_m^*(Re^{i\theta})|\sim\log_mR,\qquad
 \arg J_m^*(Re^{i\theta})=
 \frac{\theta}{P_m(R)}(1+o(1)),
\]
uniformly with first angular derivatives for bounded $\theta$.
The denominator subtracts an angular margin of order $1/P_m(R)$
near $\theta=\pm\pi/2$. On the right half-plane its argument has
the same sign as that of $w$ and, by increasing $K$, satisfies
$p|\arg J_m^*(w)|\le|\arg w|$. These observations prove
\eqref{p08:ga14:eq:Qsector}. Fixed derivatives satisfy polynomial bounds
on slightly narrower sectors by Cauchy's formula. At zero,
$Q(w)=b_1w-b_2w^2+O(w^3)$ with $b_1,b_2>0$.

\begin{lemma}[Matrix realization with arbitrary summable reciprocal zeros]
\label{p08:ga14:lem:realization}
Let $h\in C_c^\infty(\mathbb R^d\setminus\{0\})$ be nonzero,
real and even. For a sufficiently large constant $c_0$, the function
\begin{equation}\label{p08:ga14:eq:H}
 \mathcal H(u,T)=h(u)E_0(u^TT^{-1}u)
 \exp\{-2\operatorname{tr}L(c_0T^{-1})
             -\operatorname{tr}Q(T^{-1})
             -\operatorname{tr}T^{1/2}\},
 \quad\operatorname{Re}T\succ0,
\end{equation}
is the Fourier--Laplace transform of a nonzero real Schwartz function
$g(a,\Sigma)$ supported on $\Sigma\succeq0$. Its two coordinate
marginals and all joint polynomial moments vanish. For every $M\ge0$
there are positive $C_M,c_M$ such that
\begin{equation}\label{p08:ga14:eq:tail}
 |g(a,\Sigma)|\le C_M(1+|a|)^{-M}
  e^{-c_MW_{m,1+p}(\operatorname{tr}\Sigma)}.
\end{equation}
\end{lemma}
\begin{proof}
All matrix functions are principal holomorphic functions. We verify
the positive tube estimate, the boundary realization, and the contour
estimate, including the location of the possible poles.

\emph{Positive tube.}
Write $R=T^{-1}$. If $\operatorname{Re}T\succ0$ then
$\operatorname{Re}R=T^{-*}(\operatorname{Re}T)T^{-1}\succ0$.
For an accretive matrix,
\begin{equation}\label{p08:ga14:eq:schur}
 \|R\|\le C_d\sum_{i=1}^d|\lambda_i(R)|.
\end{equation}
To prove this, triangularize $R$ unitarily. Positivity of its Hermitian
part implies $|R_{ij}|\le2\sqrt{\operatorname{Re}\lambda_i
\operatorname{Re}\lambda_j}$ for the strictly upper triangular
entries. The Frobenius norm gives \eqref{p08:ga14:eq:schur}.
Let $U=\sup_{u\in\operatorname{supp}h}|u|^2$. Taking
$c_0\ge2UC_dd$ gives
$|u^TRu|\le(c_0/2)\max_i|\lambda_i(R)|$.
Since $\operatorname{Re}(u^TRu)\ge0$, equations
\eqref{p08:ga14:eq:entirebound}--\eqref{p08:ga14:eq:logpositive} show that
one copy of $\operatorname{Re}\operatorname{tr}L(c_0R)$ absorbs
the entire numerator. The other copy is nonnegative. Furthermore,
\eqref{p08:ga14:eq:schur} and \eqref{p08:ga14:eq:Qsector} imply
\[
 \operatorname{Re}\operatorname{tr}Q(R)
 \ge c\|R\|/D(\|R\|)-C
\]
when $\|R\|$ is large. For $T=X+iV$, $X\succ0$, write
$T^{1/2}=B+iC$ with real symmetric matrices $B,C$. The principal
square root is accretive, $B^2-C^2=X\succ0$, and $BC+CB=V$.
Thus $\|C\|\le\|B\|$ and
$\|V\|\le2(\operatorname{tr}B)^2$. Consequently
\begin{equation}\label{p08:ga14:eq:tubebound}
 |\mathcal H(u,X+iV)|\le C e^{-c\sqrt{\|V\|}},
 \qquad X\succ0,
\end{equation}
uniformly in $X$ and in the compact $u$ support.

\emph{A boundary neighborhood estimate.}
For invertible real symmetric $V$, put
$r=r(V)=\{\operatorname{tr}(V^{-2})\}^{-1/2}$. Then
$\min_i|\lambda_i(V)|/\sqrt d\le r\le\min_i|\lambda_i(V)|$
and $|\nabla r|\le1$. Choose a smooth positive $\eta$ with bounded
derivative, sufficiently small upper bound, equal to a fixed positive
constant for all sufficiently large $r$, and
\begin{equation}\label{p08:ga14:eq:eta}
 \eta(r)=\varepsilon r/D(1/r)\quad\hbox{for small }r.
\end{equation}
Consider $T=iV+Z$, where $\|Z\|\le C_1\eta(r)$ and $C_1$ is
any fixed constant. The inverse norm is $O(1/r)$ by a Neumann series.
Every eigenvalue of $T$ stays within $C_1\eta(r)$ of the spectrum
of the normal matrix $iV$. Grouping overlapping spectral disks
shows that some inverse eigenvalue has modulus at least $c_d/r$;
all inverse eigenvalues have arguments within
$O(\eta(r)/|\lambda_i(V)|)$ of $\pm\pi/2$. Their $Q$ terms
have nonnegative real parts and at least one supplies
$c/(rD(1/r))$ damping. To justify nonnegativity also for bounded
and small inverse eigenvalues, use compactness away from zero and
the expansion $Q(w)=b_1w-b_2w^2+O(w^3)$ at zero; a sufficiently
small upper bound on $\eta$ gives
$\operatorname{Re}Q((-t\eta+i\lambda)^{-1})\ge0$ for real
$\lambda\ne0$. For the complex neighborhood used here, an eigenvalue
of $T$ has the form $x+iy$, with $|x|\le C_1\eta(r)$ and
$|y|\ge |\lambda_i(V)|/2$. If $|(x+iy)^{-1}|$ is large, the
angular margin in \eqref{p08:ga14:eq:Qsector} applies, since
$\eta(r)/|y|\le c/P_m(1/|y|)$ after reducing $\varepsilon$.
In a compact annulus, strict positivity at the imaginary axis and
compactness apply. Finally, for large $|y|$, direct Taylor expansion
gives
\[
 \operatorname{Re}Q((x+iy)^{-1})
 =\frac{b_2+b_1x}{y^2}+O(|y|^{-3}),
\]
uniformly for small bounded $x$. Reducing the fixed upper bound of
$\eta$ makes this positive too. These three ranges justify the
nonnegativity assertion for every eigenvalue in the neighborhood.

In this neighborhood
\[
 (-\operatorname{Re}(u^TT^{-1}u))_+
       \le C\eta(r)/r^2,
 \qquad
 (-\operatorname{Re}\lambda_i(T^{-1}))_+
       \le C\eta(r)/r^2.
\]
Choose $c_0$ still larger if necessary. The largest inverse eigenvalue
then makes $J(c_0|\lambda_i|)$ dominate $J(|u^TT^{-1}u|)$.
Equations \eqref{p08:ga14:eq:entirebound} and
\eqref{p08:ga14:eq:logsector} therefore bound the logarithm of the
numerator and its $L$ regularizer together by
\begin{equation}\label{p08:ga14:eq:cancellation}
 C\Lambda\eta(r)/r^2
 =C\Lambda\varepsilon/(rD(1/r)).
\end{equation}
For sufficiently small $\varepsilon$, the $Q$ damping absorbs this
quantity. This is the step that eliminates every design-density
assumption. At large $\|V\|$ the square-root trace retains its
$c\sqrt{\|V\|}$ damping, up to a constant. Cauchy's formula on
these neighborhoods gives the same damping for each fixed matrix
derivative, times powers of $1/r$, $D(1/r)$, and $1+\|V\|$.
For $u$ derivatives, apply the scalar Cauchy formula to $E_0$ on
circles of radius one and use the chain rule. Enlarging $c_0$ absorbs
the change in $J$; the additional negative-real-part error is bounded
by $C\Lambda$. All polynomial factors are absorbed by the damping.
Thus $H(u,V)=\mathcal H(u,iV)$, extended by zero at $\det V=0$,
is Schwartz, compactly supported in $u$.

\emph{Inversion and cone support.}
In orthonormal coordinates on $\mathbb S_d$, with
$D_d=d(d+1)/2$, define
\[
 g_u(\Sigma)=(2\pi)^{-D_d}\int e^{i\operatorname{tr}(V\Sigma)}
             H(u,V)\,dV,
 \quad
 g(a,\Sigma)=(2\pi)^{-d}\int e^{-iu^Ta}g_u(\Sigma)\,du.
\]
These are Schwartz functions. If the matrix line is shifted to
$X+i\mathbb S_d$, $X\succ0$, the inverse integral equals
$e^{-\operatorname{tr}(X\Sigma)}g_u(\Sigma)$, by the
Cauchy--Riemann equations and integration by parts, first starting
at $X=\epsilon I$ and then using \eqref{p08:ga14:eq:tubebound} to let
$\epsilon\downarrow0$. The same estimate bounds these integrals
uniformly in $X$. If $\Sigma$ has a negative eigenvalue, choose
$X_0\succ0$ with $\operatorname{tr}(X_0\Sigma)<0$ and let
$X=tX_0$, $t\to\infty$. This proves cone support. Fourier inversion
now yields the required Laplace identity.

\emph{The contour and its poles.}
Deform $iV$ to
\[
 T(V)=-\eta(r(V))I+iV.
\]
Throughout the homotopy, every matrix is normal with nonreal
eigenvalues. Its inverse eigenvalues avoid the negative real axis.
Consequently the $L$ factors have no poles or cuts on the swept
region. Equivalently,
\[
 e^{-2\operatorname{tr}L(c_0T^{-1})}
     =\prod_j\det(I+c_0T^{-1}/a_j)^{-2};
\]
its poles require an individual eigenvalue of $T$ to equal
$-c_0/a_j$, which cannot occur. In contrast, a denominator depending
on $u^TT^{-1}u$ could vanish there. The numerator $E_0$ is entire.

For rigor, write $V=W+tI/\sqrt d$, $\operatorname{tr}W=0$, and
fix $W$. Break the real $t$ axis at the finitely many roots of
$\det(W+tI/\sqrt d)=0$. On each open interval, deform the scalar
contour to $t+i\sqrt d\,\eta(r(W+tI/\sqrt d))$. Every eigenvalue
has nonzero imaginary part throughout the vertical strip swept out
by this contour, so all the preceding functions are holomorphic.
The short joining segments vanish at finite endpoints by the
$\exp\{-c/(rD(1/r))\}$ estimate and at infinity by the
square-root estimate. The Jacobian has modulus at most
$1+\sqrt d\sup|\eta'|$. These uniform integrable bounds justify
Fubini in $W$ and prove the full contour identity.

Put $v=\operatorname{tr}\Sigma$. Along the deformed contour, the
absolute inverse integrand for small $r$ is bounded by
\[
 C\exp\left\{-c\sqrt{\|V\|}
 -c\left(\frac{vr}{D(1/r)}+\frac1{rD(1/r)}\right)\right\}.
\]
The expression in parentheses is at least
$cW_{m,1+p}(v)$ uniformly for large $v$. Indeed, on
$\frac14\log v\le\log(1/r)\le\log v$ the fixed-depth logarithms
are comparable with those of $v$, and $vr+r^{-1}\ge2\sqrt v$.
Below this range the first term is at least $v^{3/4}$ divided by
logarithmic factors; above it the second is at least $v/D(v)$.
Both dominate the required bound. For $r$ bounded away from zero,
$\eta(r)$ has a positive lower bound and gives exponential decay
in $v$. Reserve half the damping for integration in $V$ and integrate
by parts in $u$ to obtain \eqref{p08:ga14:eq:tail}.

Reality follows by conjugation and evenness. The transform is nonzero
because varying $T=tI$ avoids the discrete zeros of $E_0$. Its
vanishing near $u=0$ gives the covariance marginal and every joint
polynomial moment. Finally $\mathcal H(u,\epsilon I)\to0$ as
$\epsilon\downarrow0$, by the positive-tube $Q$ estimate; dominated
convergence gives the zero mean marginal.
\end{proof}

\begin{proof}[Proof of Theorem~\ref{p08:ga14:thm:matrix}]
The conditionally independent blocks of lengths
$1,\ldots,1,n-d+1$ yield exactly the mixing-transform arguments
\begin{equation}\label{p08:ga14:eq:surface}
 (u,T),\qquad T=C+uu^T/(2n),\quad C\succeq0,
                       \operatorname{rank}C\le d-1.
\end{equation}
This follows by projecting the weighted frequency matrix onto the
vector of square roots of these block lengths and its orthogonal
complement. If $T\succ0$ on this surface, the determinant lemma gives
$u^TT^{-1}u=2n$.

When $a>1$ and $\sum1/n<\infty$, apply
Lemma~\ref{p08:ga14:lem:realization} with $a_j=2n_j$ and
$0<p<a-1$. Its transform vanishes at all positive definite arguments
in \eqref{p08:ga14:eq:surface}. Singular arguments follow by perturbing
the square weighted frequency matrix to an invertible one and using
continuity. Normalize the positive and negative parts of $g$ to obtain
the two mixing laws. Their marginals and moments coincide; the tail
assertion follows from $W_{m,1+p}/W_{m,a}\to\infty$. Translation
by $b_0I$ imposes the covariance lower bound without changing any
transform zero.

For completeness, let $\rho$ be the difference of two candidate laws
with identical experiments. Fix $u\ne0$ and a rank-$d-1$ matrix
$C\succeq0$ with $u$ not orthogonal to its kernel. The bounded
holomorphic function
\[
 H(z)=\int e^{iu^TA-\operatorname{tr}(C\Sigma)
                       -z u^T\Sigma u}\,d\rho
\]
vanishes at $z=1/(2n)$, $n\in\mathcal N$. If the harmonic subseries
diverges, the half-plane Blaschke condition gives $H\equiv0$.
If $a\le1$, the covariance envelope instead yields
\[
 (\mathbb E(\operatorname{tr}\Sigma)^k)^{1/(2k)}
 \le Ck\log k\cdots\log_{m-1}k\,(\log_mk)^a,
\]
so the Stieltjes--Carleman series diverges. Taylor expansion at zero
along the accumulating zeros makes every moment of the associated
complex measure vanish; moment determinacy of a common positive
dominating measure implies that measure is zero. Thus $H\equiv0$
again. For arbitrary $T\succ0$, take
$z=(u^TT^{-1}u)^{-1}$ and $C=T-zuu^T$ to recover the full transform.
Continuity at $u=0$ and matrix Fourier--Laplace uniqueness give
$\rho=0$.
\end{proof}

\section{Compact scale does not repair sparse stable observation below index one}
\label{p08:ga14:sec:stable}

We next ask whether a bounded scale interval eliminates the sparse
sampling counterexamples. For stable index greater than one, it does.
For index less than one, it does not, even on an arbitrarily short
positive scale interval. The distinction comes from whether the
Laplace observation points accumulate at zero or escape to infinity.
The classical background is completeness of exponential systems and
M\"untz theory \cite{p08:ga14:Almira}; we prove the smooth signed-density
construction explicitly.

Fix $0<\alpha<1$, $-1\le\beta\le1$, and let
\[
 \widehat K_{a,b}(t)=
 \exp\{ita-b|t|^\alpha(1-i\kappa\operatorname{sgn}t)\},
 \qquad \kappa=\beta\tan(\pi\alpha/2).
\]
Conditionally on $(A,B)$, the $X_j$ are independent with this kernel.
We observe only the separate laws of $S_n=X_1+\cdots+X_n$.

\begin{theorem}[Exact compact-scale stable sampling criterion]
\label{p08:ga14:thm:compact-stable}
For any fixed $0<b_0<b_1<\infty$ and terminal set
$\mathcal N\subset\mathbb N$, these observations identify every
mixing law on $\mathbb R\times[b_0,b_1]$ if and only if
\begin{equation}\label{p08:ga14:eq:stable-criterion}
 \sum_{n\in\mathcal N}n^{-(1-\alpha)}=\infty.
\end{equation}
When the series converges there are distinct absolutely continuous
witnesses, with scale support in the interior of $[b_0,b_1]$, having
the same two coordinate marginals and the same finite mixed polynomial
moments of every order. Their signed difference is smooth, compactly
supported in scale, and Schwartz in location.
\end{theorem}

No mixing over the index or skewness is permitted. The unrestricted
scale criterion was already proved in the preceding manuscript; the
new claim is the persistence of its converse inside every prescribed
compact positive scale interval.

\begin{lemma}[Smooth compact support with summable ray zeros]
\label{p08:ga14:lem:compactzeros}
Suppose $z_j=c\lambda_j$, $\operatorname{Re}c>0$,
$\lambda_j\to\infty$, and $\sum_j\lambda_j^{-1}<\infty$.
For every nonempty bounded open real interval $I$ there is a nonzero
real $f\in C_c^\infty(I)$ whose bilateral Laplace transform $F$
satisfies $F(z_j)=F(\overline z_j)=F(0)=0$ for every $j$.
Finite zero sets are allowed as well.
\end{lemma}
\begin{proof}
Put $r_j=|z_j|$. For real $x$,
\begin{equation}\label{p08:ga14:eq:sinc}
 (1+x^2)\left|\frac{\sin4x}{4x}\right|^2\le1.
\end{equation}
If $|x|\ge1/\sqrt{15}$, use $|\sin4x|\le1$. Otherwise
$|4x|\le4/\sqrt{15}$ and the alternating Taylor bound gives
$\sin u/u\le1-u^2/7$ on this interval. Squaring shows
$|\sin4x/(4x)|^2\le1-x^2$, proving the assertion.
Moreover
\[
 |(1-iy/z_j)(1-iy/\overline z_j)|
 \le1+y^2/r_j^2.
\]
This follows by squaring: the difference from
$(1+y^2/r_j^2)^2$ is $-4y^2(\operatorname{Im}z_j)^2/r_j^4$.

Choose $J_0$ so large that
$A=8\sum_{j>J_0}r_j^{-1}$ is arbitrarily small, and form the entire
product
\[
 G(s)=\prod_{j>J_0}
 (1-s/z_j)(1-s/\overline z_j)
 \left\{\frac{\sinh(4s/r_j)}{4s/r_j}\right\}^{\!2}.
\]
It converges locally uniformly, since its deviations from one are
$O(r_j^{-1})$ on compact sets. It is real on the real axis and
$G(0)=1$. Equation \eqref{p08:ga14:eq:sinc} gives $|G(iy)|\le1$.
Its exponential type is at most $A$: the product of the hyperbolic
sine factors has type at most $A$, while
\[
 \log\prod_j(1+|s|/r_j)^2=o(|s|).
\]
For the last assertion, remove finitely many factors and bound the
remaining logarithms by $2|s|\sum_{j>J}r_j^{-1}$, then send
$J\to\infty$.

The elementary bounded-axis form of the Paley--Wiener theorem now
realizes $G$ as the transform of a compactly supported distribution
on $[-A,A]$. One may apply the usual theorem with $A+\epsilon$ and
then let $\epsilon\downarrow0$: the bound on the imaginary axis
and the Phragm\'en--Lindel\"of principle give
$|G(s)|\le e^{(A+\epsilon)|\operatorname{Re}s|}$ up to a constant.
Multiply $G(s)$ by the Laplace transform of a nonzero real smooth
bump supported in $[-\epsilon,\epsilon]$, by
$s\prod_{j\le J_0}(1-s/z_j)(1-s/\overline z_j)$, and by
$e^{-bs}$, where $b$ is the midpoint of $I$. In real space these
operations are smoothing, finitely many differentiations, and
translation. By first making $A$ and then $\epsilon$ sufficiently
small, the resulting nonzero smooth function is supported in $I$.
The resulting transform has all the prescribed zeros and the factor
$s$ makes its integral zero.
\end{proof}

\begin{proof}[Proof of Theorem~\ref{p08:ga14:thm:compact-stable}]
For a fixed nonzero $t$, the data for $S_n/n$ evaluate the mixing
Fourier--Laplace transform at
\[
 (t,s_n(t)),\qquad
 s_n(t)=|t|^\alpha(1-i\kappa\operatorname{sgn}t)n^{1-\alpha}.
\]
Their half-plane Blaschke sum diverges precisely under
\eqref{p08:ga14:eq:stable-criterion}, proving sufficiency without any
scale bound.

For necessity apply Lemma~\ref{p08:ga14:lem:compactzeros} to
$z_n=(1-i\kappa)n^{1-\alpha}$, choosing the support of $f$ in a
small interval strictly inside $(b_0,b_1)$. Choose a real nonzero
even function $h\in C_c^\infty(\mathbb R\setminus\{0\})$
supported sufficiently close to $\{-1,1\}$ that
\[
 G(t,b)=h(t)|t|^\alpha f(b|t|^\alpha)
\]
has its $b$ support inside $(b_0,b_1)$ for every $t$ in that support.
Set $g(a,b)=(2\pi)^{-1}\int e^{-ita}G(t,b)\,dt$.
It is real and nonzero, smooth, compactly supported in $b$, and
Schwartz in $a$. Its transform is
\[
 \int e^{ita-sb}g(a,b)\,da\,db=h(t)F(s/|t|^\alpha).
\]
For positive $t$ the observation argument is a prescribed zero
$z_n$, and for negative $t$ its conjugate is a prescribed zero.
The scale marginal and every mixed polynomial moment vanish because
$G$ is zero near $t=0$. The location marginal vanishes because
$F(0)=0$. Normalize the positive and negative parts of $g$ to finish.
\end{proof}

\begin{remark}[The contrasting index range and remaining targets]
For $1<\alpha\le2$ and fixed skewness, a compact upper bound on
$B$ makes the transform entire in $s$. For every infinite terminal
set the points $s_n(t)$ converge to zero, so the identity theorem
identifies the mixing law. This elementary positive comparison is
not counted as a third new target. At index one with nonzero
skewness, the points instead run along a vertical line with imaginary
parts proportional to $\log n$; compact-scale identification then
requires a different entire-function sampling analysis, not the
preceding ray-zero argument. Neither that classification nor joint
mixing over unknown stable shapes is claimed here. Quantitative
finite-data stability at the matrix quasi-analytic boundary also
remains open.
\end{remark}

% ===== spectral =====
\section{Adaptive zero observation and sharp recovery of individual spectral atoms}
\label{p08:sp14:sec}

The earlier spectral results determined an atom-count boundary and
constructed a finite-height transition experiment on a uniform amplitude
grid.  We address two different targets here.  First, a noisy zero test
can be turned into an adaptive comparison oracle, including calibration
of the previously known largest frequency and its phase.  Second, the
nonlinear threshold data determine individual nodes and masses with the
sharp separation exponents, not merely a Wasserstein error.  The latter
exponents are classical in linear super-resolution
\cite{p08:sp14:BGY,p08:sp14:KDB}; our contribution is their transfer, with an
explicit weak-mass dependence, to the nonlinear minimum observations.
The improvement of the recurrence exponent uses the classical optimal
filling-time theorem \cite{p08:sp14:DF}, and is not claimed as a new theorem
in Diophantine approximation.

Throughout this section fix $0<a<b<1$, and write
\begin{equation}\label{p08:sp14:eq:model}
 S_\mu(x)=\int_a^b x^r\,d\mu(r),\qquad
 T_\mu(C)=\min_{x>0}\frac{1+Cx}{S_\mu(x)},
 \qquad \mu\in\mathcal P[a,b].
\end{equation}
A one-sided setting list consists of distinct points of
$(0,a/(1-a))$.  Constants below may depend on the fixed support
interval, setting list, and atom bound.  They do not depend on the
minimum mass, node separation, or noise level.

\subsection{Recovery of individual masses and locations}

\begin{theorem}[Sharp local parameter recovery from nonlinear minima]
\label{p08:sp14:thm:parameters}
Fix $m\ge2$ and $2m-1$ one-sided settings $C_i$.  Suppose
$\mu=\sum_{j=1}^m w_j\delta_{r_j}$, with increasing nodes,
$r_{j+1}-r_j\ge\Delta>0$, and $w_j\ge w>0$.  Let $\nu$ be any
probability on $[a,b]$ with at most $m$ atoms, and put
\[
 \delta=\max_i|T_\nu(C_i)-T_\mu(C_i)|,
 \qquad D_j=\prod_{k\ne j}|r_j-r_k|,
 \qquad \kappa_j=\sum_{k\ne j}|r_j-r_k|^{-1}.
\]
There are constants $c,K>0$ such that, whenever
\begin{equation}\label{p08:sp14:eq:parameterregime}
                         \delta\le c w\Delta^{2m-1},
\end{equation}
$\nu$ has exactly $m$ atoms, which can be written in increasing order
as $\nu=\sum_j v_j\delta_{s_j}$, and
\begin{align}
 |s_j-r_j|&\le \frac{K\delta}{w_jD_j^2},
                                                   \label{p08:sp14:eq:noderate}\\
 |v_j-w_j|&\le \frac{K(1+\kappa_j)\delta}{D_j^2}.
                                                   \label{p08:sp14:eq:massrate}
\end{align}
In particular, the uniform bounds are
\[
 \max_j|s_j-r_j|\le\frac{K\delta}{w\Delta^{2m-2}},
 \qquad
 \max_j|v_j-w_j|\le\frac{K\delta}{\Delta^{2m-1}}.
\]
Their joint dependence on $\delta,w,\Delta$ is sharp: for every
$0<w\le1/(2m)$, every sufficiently small $\Delta$, and
$0<\delta\le c w\Delta^{2m-1}$, there are two $m$-atomic
probabilities, both with minimum mass at least $w$ and minimum
separation at least $\Delta$, whose complete threshold curves on
any fixed compact interval $I\subset(0,\infty)$ differ by at most
$\delta$, while one node moves at least
$c_I\delta/(w\Delta^{2m-2})$ and one mass changes by at least
$c_I\delta/\Delta^{2m-1}$.
\end{theorem}

The threshold required for a stable assignment of individual atoms is
of order $w\Delta^{2m-1}$.  This is distinct from the larger
$w\Delta^{2m-2}$ noise scale at which the preceding paper could
reject every smaller positive model.  Counting the atoms can therefore
be easier than resolving their individual locations.

\begin{lemma}[Uniform smooth-test conversion with two sparse sides]
\label{p08:sp14:lem:sparsetest}
For probabilities $\alpha,\beta$ having at most $m$ atoms each,
and any $p\in C^{2m-1}[a,b]$, one has
\begin{equation}\label{p08:sp14:eq:sparsetest}
 \left|\int p\,d(\alpha-\beta)\right|
 \le K\|p\|_{C^{2m-1}}
             \max_{1\le i\le2m-1}|T_\alpha(C_i)-T_\beta(C_i)|.
\end{equation}
No separation or mass bound is needed in this lemma.
\end{lemma}
\begin{proof}
Here are the geometric and interpolation details of the conversion.
For a profile $S$, put $H=S-xS'$.  Strict concavity gives a unique
contact $x(C)$, satisfying
\[
 C=\frac{S'(x)}{H(x)},\qquad
 \frac{dC}{dx}=\frac{S''(x)S(x)}{H(x)^2}<0,
 \qquad \frac a{1-a}\le Cx\le\frac b{1-b}.
\]
The tangent line is $L_C(x)=(1+Cx)/T(C)$.  At our fixed one-sided
settings, all contacts lie in a fixed compact interval above $1$;
$S,H,-S''$ are bounded above and below there.  Consequently distinct
contacts have a uniform positive separation.  The intersection of two
consecutive tangent lines lies strictly between their contacts, at a
uniform positive distance from both: this follows by integrating the
uniform lower bound for $-S''$ and dividing by the bounded difference
of their slopes.
For the two end contacts, close the corresponding brackets by fixed
external endpoints strictly outside the common compact contact
interval, choosing the lower external endpoint still larger than $1$.

Write $\delta$ for the maximum threshold difference.  For small
$\delta$, the two tangent-line systems, and their consecutive
intersections, differ by $O(\delta)$.  For each setting, the two
contacts lie between the corresponding neighboring intersections, with
uniform margins.  At the $\alpha$ contact,
$S_\alpha-S_\beta\ge-K\delta$, and at the $\beta$ contact the
reverse inequality holds.  An intermediate point therefore has
$|S_\alpha-S_\beta|\le K\delta$.  The margins imply that the
$2m-1$ points obtained in this way have a uniform separation.
Their logarithms, together with $u_0=0$, give $2m$ separated
points in a fixed compact set, with $u_i>0$ for $i>0$, such that
\[
 \left|\int e^{u_i r}\,d(\alpha-\beta)(r)\right|\le K\delta.
\]
The sample at $0$ is exactly zero.

Interpolate $p$ on the union of the supports by a linear combination
of the $2m$ functions $e^{u_i r}$; pad the node list, if necessary,
with repeated nodes and impose the corresponding derivatives.
In divided-difference coordinates the interpolation matrix extends
continuously across every collision.  Every limiting matrix is
invertible, since a nonzero sum of $2m$ real exponentials with
distinct exponents has at most $2m-1$ zeros counted with
multiplicity.  The latter assertion follows inductively from Rolle's
theorem after multiplication by the first exponential's reciprocal.
Compactness bounds the interpolation coefficients by
$K\|p\|_{C^{2m-1}}$.  Integrating the interpolant proves the
claim.  For large $\delta$ increase $K$ and use
$|\int p\,d(\alpha-\beta)|\le2\|p\|_\infty$.
\end{proof}

\begin{proof}[Upper bounds in Theorem~\ref{p08:sp14:thm:parameters}]
Write the nodes of $\nu$ as $s_1,\ldots,s_q$, $q\le m$,
padding to $m$ by repetitions for the following polynomial test.
For each $j$ use
\[
 Q_j(r)=\prod_{k\ne j}(r-r_k)\prod_{\ell=1}^{m}(r-s_\ell).
\]
Its degree is $2m-1$ and its $C^{2m-1}$ norm is uniformly
bounded.  Lemma~\ref{p08:sp14:lem:sparsetest} gives
\begin{equation}\label{p08:sp14:eq:productbound}
 w_jD_j\prod_{\ell=1}^m|r_j-s_\ell|\le K\delta.
\end{equation}
If no node of $\nu$ is within $\Delta/4$ of $r_j$, the left
side is at least $w\Delta^{2m-1}/4^m$.  Thus, for sufficiently
small $c$ in \eqref{p08:sp14:eq:parameterregime}, each of the $m$
disjoint $\Delta/4$ neighborhoods contains a node.  There are
exactly $m$ distinct nodes, matched in increasing order.  For
$k\ne j$,
\[
 |r_j-s_k|\ge |r_j-r_k|-\Delta/4
                       \ge\tfrac34|r_j-r_k|.
\]
Substitution into \eqref{p08:sp14:eq:productbound} gives
\eqref{p08:sp14:eq:noderate}.

For the masses put
\[
 P_j(r)=
 \frac{\prod_{k\ne j}(r-r_k)(r-s_k)}
      {\prod_{k\ne j}(r_j-r_k)(r_j-s_k)}.
\]
Its $C^{2m-1}$ norm is at most $K/D_j^2$, so
\begin{equation}\label{p08:sp14:eq:weighttest}
                  |w_j-v_jP_j(s_j)|\le K\delta/D_j^2.
\end{equation}
By taking $c$ smaller, the node bound ensures
$|s_j-r_j|\le c_0\Delta$ with an arbitrarily small fixed $c_0$.
On this short segment all the factors defining $P_j$ stay away
from zero, its logarithmic derivative is bounded by $K\kappa_j$,
and
\[
 \tfrac12\le P_j(s_j)\le2,\qquad
 |P_j(s_j)-1|\le K\kappa_j|s_j-r_j|.
\]
Moreover $\delta/D_j^2\le cw\Delta\le c w_j$.  Thus
\eqref{p08:sp14:eq:weighttest} first gives $v_j\le K w_j$, and then
\eqref{p08:sp14:eq:noderate} and \eqref{p08:sp14:eq:weighttest} give
\eqref{p08:sp14:eq:massrate}.
\end{proof}

\begin{proof}[Matching examples in Theorem~\ref{p08:sp14:thm:parameters}]
Choose nodes $r_j(0)=r_0+2(j-1)\Delta$ in the interior of
$[a,b]$.  Give the left endpoint mass $2w$ and divide the
remaining mass equally among the other $m-1$ nodes.  All masses
are at least $2w$.  For distinct moving nodes define
\[
 B_j=\prod_{k\ne j}(r_j-r_k)^{-2},\qquad
 A_j=-2B_j\sum_{k\ne j}(r_j-r_k)^{-1},
\]
and solve the smooth ordinary differential equations
\begin{equation}\label{p08:sp14:eq:momentflow}
                 \dot w_j=A_j,\qquad \dot r_j=B_j/w_j.
\end{equation}
For $0\le s\le c w\Delta^{2m-1}$, a direct bootstrap gives
node separations between fixed positive multiples of $\Delta$,
displacements at most $\Delta/4$, masses at least $w$, and
$w_1(s)\asymp w$.  Indeed, while these bounds hold,
$|\dot w_j|\le K\Delta^{-(2m-1)}$ and
$|\dot r_j|\le K/(w\Delta^{2m-2})$.  Integrating gives strict
versions of the assumed bounds for a sufficiently small $c$, so
standard continuation extends the solution throughout the interval.

The partial-fraction identity for
$\prod_j(z-r_j)^{-2}$, or the leading coefficient of a Hermite
interpolant, gives for every $f\in C^{2m-1}[a,b]$
\begin{equation}\label{p08:sp14:eq:hermitedifference}
 \sum_j\{A_j f(r_j)+B_j f'(r_j)\}
       =f[r_1,r_1,\ldots,r_m,r_m].
\end{equation}
Here the right side is the confluent divided difference of order
$2m-1$; its absolute value is at most
$\|f^{(2m-1)}\|_\infty/(2m-1)!$.  This identity proves that
the flow preserves every moment through degree $2m-2$, including
total mass, while its moment of degree $2m-1$ has derivative $1$.
For $f(r)=x^r$ on any fixed compact $x$ interval,
\eqref{p08:sp14:eq:hermitedifference} is uniformly bounded.  Therefore
\[
 \sup_x|S_{\mu_s}(x)-S_{\mu_0}(x)|\le K s,
 \qquad
 \sup_{C\in I}|T_{\mu_s}(C)-T_{\mu_0}(C)|\le K_I s.
\]
The second estimate follows by evaluating the two quotients at each
other's minimizing contacts, which all lie in a fixed compact
interval with positive denominators.

At the leftmost node all terms in the sum defining $A_1$ have the
same sign.  Along the bootstrap interval,
\[
 \dot w_1\ge c\Delta^{-(2m-1)},\qquad
 \dot r_1\ge c/(w\Delta^{2m-2}).
\]
Take $s=\delta/K_I$, decreasing the allowed constant in
\eqref{p08:sp14:eq:parameterregime} if necessary.  Integrating these
two inequalities proves both matching lower bounds, while the
minimum separation remains at least $\Delta$.
\end{proof}

\subsection{A self-calibrating noisy comparison oracle}

Consider the controlled exponential polynomial
\begin{equation}\label{p08:sp14:eq:exponential}
 f_{A,C}(z)=-1+A\sum_{j=1}^{m}w_j
          e^{\ell r_jz+i\phi_j}+Ce^{\ell z+i\phi_*}.
\end{equation}
The control amplitudes $A,C$ are known, but $\ell$, $\phi_*$,
and all the smaller phases may be unknown.  Assume only a known
compact bound $0<\ell_-\le\ell\le\ell_+$.  A zero observation
at height $H$ means a complete list, with multiplicity, of the zero
ordinates in $(0,H)$, each perturbed by at most $\epsilon$.
There are no unreported or spurious zeros.  This is an observation
model; it does not assert that a certified zero list can be computed
at negligible cost.

Put $d=m+1$ and
$\omega=(\ell r_1,\ldots,\ell r_m,\ell)$.  It is enough to
assume a finite collection of nonresonance inequalities at the
resolution of the intended experiment:
\begin{equation}\label{p08:sp14:eq:truncated}
 |k\cdot\omega|\ge\gamma\|k\|_\infty^{-\tau}
 \quad\left(0<\|k\|_\infty\le K_0/h\right),
 \qquad 0<\gamma\le1,\quad\tau>d-1.
\end{equation}
The fixed constant $K_0$ below depends on $d$ and the frequency
bounds.  In particular, the theorem does not require infinitely
many inequalities or even global rational independence.

\begin{theorem}[Adaptive calibrated recovery from noisy zero ordinates]
\label{p08:sp14:thm:adaptive}
Fix $I\Subset(0,\infty)$.  There are known constants
$A_-<A_+$, $h_0,\epsilon_0,K_0,K,H_0>0$, depending only on
$a,b,m,I,\ell_-,\ell_+,\tau$, with the following property.
Let $0<h<h_0$, $0\le\epsilon<\epsilon_0$, and assume
\eqref{p08:sp14:eq:truncated}.  At height
\begin{equation}\label{p08:sp14:eq:newheight}
                       H\ge H_0+K\gamma^{-1}h^{-\tau},
\end{equation}
one calibration experiment at $A=0$, and at most
$K\log(3/h)$ adaptively selected amplitudes in $[A_-,A_+]$
for each $C\in I$, suffice to construct
\begin{equation}\label{p08:sp14:eq:robustthreshold}
                |\widehat T(C)-T_\mu(C)|\le K(h+\epsilon^2).
\end{equation}
The result is uniform in all phases and positive masses.  With a
known top frequency and phase, omit the calibration experiment.
The smaller phases are not reconstructed by this procedure.
\end{theorem}

Compared with the preceding uniform-grid construction, the number
of amplitude experiments changes from $O(h^{-1})$ to
$O(\log(1/h))$, while the sufficient height changes from
$O(\gamma^{-1}h^{-(d+\tau)})$ to
$O(\gamma^{-1}h^{-\tau})$.  The second improvement comes from
using an established optimal recurrence estimate.  The statement
does not claim an optimal lower bound for the height needed to
observe zeros.

\begin{lemma}[Finite-order filling bound]
\label{p08:sp14:lem:filling}
Under \eqref{p08:sp14:eq:truncated}, every target on
$(\mathbb R/2\pi\mathbb Z)^d$ is approached within coordinate
distance $h$ by $t\omega$ in every interval of length
$K\gamma^{-1}h^{-\tau}$.
\end{lemma}
\begin{proof}
Apply the truncated filling-time theorem of Dumas and Fischler
\cite[Theorem~1]{p08:sp14:DF}, which recovers the optimal theorem of
Berti, Biasco, and Bolle.  On the unit torus, a unit direction
$\alpha$ satisfying
$|k\cdot\alpha|\ge\gamma_0\|k\|_2^{-\tau}$ up to order
$(1+d^2d!)/\eta$ fills to Euclidean accuracy $\eta$ in time
at most $C(d,\tau)\gamma_0^{-1}\eta^{-\tau}$.
Use $\alpha=\omega/\|\omega\|_2$ and
$\eta=h/(2\pi)$.  Norm equivalence and the fixed bounds on
$\|\omega\|_2$ are absorbed in $K_0,K$.  Decrease the
normalized Diophantine constant by a fixed factor, if necessary,
so that the nonempty-class condition in that theorem holds;
this is possible for every $\tau>d-1$ and changes only $K$.
Translating the
initial point proves the assertion for every interval.
\end{proof}

\begin{lemma}[Calibration with bounded ordinate errors]
\label{p08:sp14:lem:calibration}
For sufficiently large fixed $H_0$ and small fixed $\epsilon_0$,
the $A=0$ list yields $\widehat\ell,\widehat\phi_*$ such that
\[
 |\widehat\ell-\ell|\le K\epsilon/H,
 \qquad
 \operatorname{dist}\bigl(
 \widehat\ell\widetilde t+\widehat\phi_*
      -(\ell t+\phi_*),2\pi\mathbb Z\bigr)\le K\epsilon
\]
for every $t\in[0,H]$ and every reported $\widetilde t$ with
$|\widetilde t-t|\le\epsilon$.
\end{lemma}
\begin{proof}
At $A=0$ the ordinates form an arithmetic progression of spacing
$2\pi/\ell$.  If $\epsilon<\pi/(4\ell_+)$, their order is
preserved.  Let $\widetilde t_1,\widetilde t_N$ be the first
and last entries and set
\[
 \widehat\ell=\frac{2\pi(N-1)}
                        {\widetilde t_N-\widetilde t_1},
 \qquad \widehat\phi_*=-\widehat\ell\widetilde t_1
                                      \pmod{2\pi}.
\]
For $H\ge H_0$ the true first-to-last span is comparable to $H$.
Perturbing its two endpoints by at most $\epsilon$ proves the
frequency estimate.  Since $\ell t_1+\phi_*$ is an integral
multiple of $2\pi$, the phase difference for any $t$ is bounded
by $|\widehat\ell-\ell|H+K\epsilon$, as required.
\end{proof}

The step that permits adaptive amplitude selection is stronger than
the existence of a successful grid point near the transition.

\begin{lemma}[A uniform one-sided comparison test]
\label{p08:sp14:lem:comparisonoracle}
There is a known $R>0$ such that the following rule applies to every
$A\in[A_-,A_+]$.  Accept $A$ if its noisy zero list contains
$\widetilde t$ satisfying
\begin{equation}\label{p08:sp14:eq:accept}
 \operatorname{dist}(\widehat\ell\widetilde t+
       \widehat\phi_*,\pi+2\pi\mathbb Z)
                              \le R\sqrt h+K\epsilon.
\end{equation}
At the height \eqref{p08:sp14:eq:newheight}:
\begin{enumerate}
\item every $A\ge T_\mu(C)$ is accepted;
\item every accepted $A$ satisfies
      $A\ge T_\mu(C)-K(h+\epsilon^2)$.
\end{enumerate}
\end{lemma}
\begin{proof}
All thresholds lie in a fixed positive compact interval: evaluate
at $x=1$ for the upper bound, and use
$S(x)\le\max(x^a,x^b)$ for the lower bound.  Choose
$[A_-,A_+]$ containing this interval with positive margins.
Write $T=T_\mu(C)$, let $x_*$ be its minimizing contact, and set
\[
 g_A(z)=-1+A S_\mu(e^{\ell z})-Ce^{\ell z},
 \qquad \sigma_*=(\log x_*)/\ell.
\]
The function $g_T$ has a double zero at $\sigma_*$, and
\[
 g_T''(\sigma_*)
  =-\ell^2T\int r(1-r)x_*^r\,d\mu(r)\le-c<0.
\]
All third derivatives are uniformly bounded in a fixed complex
neighborhood.  Hence, on a sufficiently small circle of radius
$\rho$, $|g_T(z)|\ge c\rho^2$, and $g_T$ has exactly two zeros
there, counting multiplicity.

If $A>T$, the strictly concave function
$F_A(x)=AS_\mu(x)-1-Cx$ has two real positive zeros.  All these
zeros, for $A\in[T,A_+]$, lie in a fixed compact interval:
$F_A$ is uniformly negative near $0$ and near infinity.  On this
interval $F_A''\le-c$, and
$\max F_A\ge(A-T)S(x_*)\ge c(A-T)$.  The upper quadratic
bound for a concave function of curvature at most $-c$ now gives
\[
 |F_A'(x)|\ge c\sqrt{A-T}
\]
at either root.  The corresponding derivative of $g_A$ has the
same lower bound, and its second derivative is uniformly bounded.

Lemma~\ref{p08:sp14:lem:filling}, applied to the phase target
$(-\phi_1,\ldots,-\phi_m,\pi-\phi_*)$, gives
$t_0\in[1,H-1]$ such that, uniformly in the fixed complex
neighborhood of these real roots,
\begin{equation}\label{p08:sp14:eq:phaseperturb}
                   |f_{A,C}(z+it_0)-g_A(z)|\le K_1h.
\end{equation}
We show that the perturbed function has a zero within imaginary
distance $K\sqrt h$ of the real line, uniformly for every
$A\ge T$.

For $0\le A-T\le D h$, compare directly with $g_T$ on a circle
of radius $R_0\sqrt h$.  Choose $R_0$ so large that its quadratic
lower bound dominates both \eqref{p08:sp14:eq:phaseperturb} and the
$O(Dh)$ difference $g_A-g_T$.  Rouche's theorem gives two zeros
inside.  Here $D$ is the fixed constant chosen next.
For $A-T>D h$, use a circle about a real simple root of $g_A$
of radius $R_1h/\sqrt{A-T}$.  Taylor's formula gives the lower
bound
\[
 |g_A(z)|\ge
      \left(cR_1-\frac{K R_1^2}{D}\right)h.
\]
First choose $R_1$ so that $cR_1>2K_1$, then $D$ large enough
to make this bound exceed $K_1h$.  Linear-term dominance also
shows that $g_A$ has exactly one zero in this disk; Rouche's
theorem supplies the perturbed zero.  Its distance is at most
$(R_1/\sqrt D)\sqrt h$.  Finally choose $R_0$ for the first
case.  Both cases are valid for all sufficiently small $h$.

After translating back, the resulting zero lies in $(0,H)$ and
its top phase is within $K\sqrt h$ of $\pi$.  The calibration
lemma and the ordinate error prove acceptance with a suitable
$R$.  This proves the first assertion.

Conversely, any accepted entry comes from a true zero whose top
phase equals $\pi+v$ modulo $2\pi$, with
$|v|\le K(\sqrt h+\epsilon)$.  Put $x=e^{\ell\Re z}$ and
$y=Cx$.  The zero equation and the triangle inequality give
\[
 AS_\mu(x)\ge|1-ye^{i(\pi+v)}|,
 \qquad
 |1-ye^{i(\pi+v)}|^2
 \ge(1+y)^2-yv^2
 \ge(1+y)^2(1-v^2/4).
\]
For small $h,\epsilon$, it follows that
$A\ge(1-K(h+\epsilon^2))(1+Cx)/S_\mu(x)
\ge T-K(h+\epsilon^2)$.  This proves the second assertion.
\end{proof}

\begin{proof}[Proof of Theorem~\ref{p08:sp14:thm:adaptive}]
Calibrate once, and start with a known interval $[L,U]$ containing
all possible thresholds.  At each step test its midpoint $A$.
If rejected, replace $L$ by $A$; the first assertion of
Lemma~\ref{p08:sp14:lem:comparisonoracle} guarantees $A<T$.
If accepted, replace $U$ by $\min(U,A+E)$, where
$E=K(h+\epsilon^2)$ is the known bound in its second assertion.
The true threshold remains in the interval.  Its width after $J$
steps is at most
\[
                     2^{-J}(U_0-L_0)+2E.
\]
Taking $J=\lceil\log_2((U_0-L_0)/h)\rceil$ makes the first
term at most $h$, with $J\le K\log(3/h)$ for small $h$.
The midpoint of the final interval satisfies
\eqref{p08:sp14:eq:robustthreshold}.  No monotonicity of the finite
zero-list acceptance function is assumed or needed.
\end{proof}

\begin{corollary}[Individual spectral resolution from finite noisy zeros]
\label{p08:sp14:cor:combined}
Run the adaptive experiment at $2m-1$ fixed one-sided settings,
and fit any at most $m$-atomic probability to the resulting
threshold vector within its guaranteed error.  If the truth has
minimum mass $w$, separation $\Delta$, and
$h+\epsilon^2\le c w\Delta^{2m-1}$, then after ordering the
atoms its location and mass errors are at most
\[
 \frac{K(h+\epsilon^2)}{w\Delta^{2m-2}},
 \qquad
 \frac{K(h+\epsilon^2)}{\Delta^{2m-1}},
\]
respectively.  The sharper $D_j,\kappa_j$ bounds in
Theorem~\ref{p08:sp14:thm:parameters} also apply.
\end{corollary}
\begin{proof}
The truth is a feasible fit, and the threshold discrepancy between
the truth and any such fit is at most $K(h+\epsilon^2)$.
Apply Theorem~\ref{p08:sp14:thm:parameters}.
\end{proof}

The quadratic occurrence of ordinate noise in this construction
comes from the quadratic angular deficit at the antipodal phase.
It is a proved robustness bound for this experiment, not a general
minimax lower bound for all uses of a noisy zero divisor.  The
calibration experiment requires the ability to set $A=0$, and the
finite nonresonance hypothesis remains an external quantitative
assumption.  Neither arbitrary missing zeros nor a uniform
infinite-prime analogue is covered.

% ===== bessel =====
\section{Gamma extraction has two different sharp boundaries}
\label{p08:be14:section}

The preceding paper excluded extraction of the full tail gamma factor
from generalized McKay laws with positive excess power.  Exclusion of
one factor does not determine the largest factor that can actually be
extracted.  We solve that optimization problem on the first excess-power
line, allowing an arbitrary probability remainder.  The answer is
strictly different from the answer with an infinitely divisible (ID)
remainder.  Their deficits from the tail shape have respectively
exponential and algebraic size at large Bessel order.

These are new transfer questions, not named conjectures in
\cite{p01:rare:Baricz2026}.  All arguments in this section are independent of
the earlier unpublished signed-spectral constructions.  We use only
an elementary Laplace transform, uniqueness of Laplace transforms,
the L\'evy--Khintchine representation on the positive half-line, and
standard identities for Kummer's function \cite{p01:rare:DLMF,p05:fs:Bondesson}.

\subsection{The family and the two capacities}

Fix $\kappa>1/2$ and $b>1$, and let $X_{\kappa,b}$ have density
proportional to
\[
 x^{\kappa-1/2}e^{-bx}I_{\kappa-3/2}(x),\qquad x>0.
\]
This is the $q=1$ generalized McKay line, with
$\mu=\kappa-3/2$ and $\nu=\mu+2$.  Expanding the positive
$I_\mu$ series, integrating term by term, and applying the gamma
duplication formula gives
\[
 \int_0^\infty x^{\kappa-1/2}e^{-bx}I_{\kappa-3/2}(x)\,dx
 =\frac{2^{\kappa-1/2}\Gamma(\kappa)}{\sqrt\pi}
       \frac b{(b^2-1)^\kappa}.
\]
Consequently
\begin{equation}\label{p08:be14:laplace}
 L_{\kappa,b}(s)=\frac{b+s}{b}
       \left(\frac{b^2-1}{(b+s)^2-1}\right)^\kappa,
       \qquad s\geq0.
\end{equation}
The range $\kappa>1/2$ covers the full standard parameter range
$\mu>-1$ considered here, including Bessel orders
$-1<\mu\leq-1/2$.

For a rate $\rho>0$, let $h_{\rm P}(\rho)$ be the supremum of shapes
$h\geq0$ for which
\[
 X_{\kappa,b}\stackrel{\rm d}=G_{h,\rho}+Y,
 \qquad G_{h,\rho}\sim\operatorname{Gamma}(h,\text{rate }\rho),
 \qquad Y\geq0,\qquad G_{h,\rho}\perp Y.
\]
The convention at $h=0$ is a point mass at zero.  Define
$h_{\rm ID}(\rho)$ by additionally requiring $Y$ to be ID.
The admissible shapes form intervals starting at zero: a smaller
gamma factor is obtained by putting the remaining independent gamma
factor into $Y$.  We first calculate the ID capacity at every rate.

\begin{proposition}[All-rate ID extraction diagram]\label{p08:be14:id}
Put $\lambda=\rho-b$.  Then
\begin{equation}\label{p08:be14:id-all}
 h_{\rm ID}(\rho)=
 \begin{cases}
 0,&\lambda<-1,\\
 z_\lambda^{-1-\lambda}
       \{\kappa(1+z_\lambda^2)-z_\lambda\},&-1\leq\lambda<0,\\
 2\kappa-1,&\lambda\geq0,
 \end{cases}
\end{equation}
where
\[
 z_\lambda=
 \frac{-\lambda+\sqrt{\lambda^2+4\kappa^2(1-\lambda^2)}}
      {2\kappa(1-\lambda)}.
\]
Every displayed maximum is attained.  In particular
\begin{equation}\label{p08:be14:id-endpoint}
 h_{\rm ID}(b-1)=\kappa-\frac1{4\kappa}.
\end{equation}
For unrestricted probability remainders one also has
$h_{\rm P}(\rho)=0$ when $\rho<b-1$ and
$h_{\rm P}(\rho)=2\kappa-1$ when $\rho\geq b$.
\end{proposition}
\begin{proof}
Logarithmic differentiation of \eqref{p08:be14:laplace} gives the
positive L\'evy density
\[
 \frac{e^{-bx}}x\{2\kappa\cosh x-1\},\qquad x>0.
\]
Subtracting the gamma L\'evy density shows, by uniqueness, that an
ID remainder exists exactly when
\[
 h\leq\inf_{x\geq0}e^{\lambda x}(2\kappa\cosh x-1).
\]
The infimum is zero for $\lambda<-1$, while for $\lambda\geq0$
the function is nondecreasing and its minimum is $2\kappa-1$.
For $-1\leq\lambda<0$, its derivative vanishes at exactly one
positive point.  With $z=e^{-x}$, the equation is
\[
 \kappa(1-\lambda)z^2+\lambda z-\kappa(1+\lambda)=0.
\]
Its relevant root is $z_\lambda\in(0,1)$; at $\lambda=-1$ it
equals $1/(2\kappa)$.  The derivative changes from negative to
positive there, proving \eqref{p08:be14:id-all}.  The residual L\'evy
density is nonnegative at equality and is integrable in the L\'evy
sense, proving attainment.

If $\rho<b-1$, choose $t$ with $\rho<t<b-1$.  Formula
\eqref{p08:be14:laplace} gives $\mathbb E e^{tX}<\infty$, whereas
$\mathbb E e^{tG_{h,\rho}}=\infty$ for every $h>0$.  The pathwise
inequality $G_{h,\rho}\leq G_{h,\rho}+Y$ rules out extraction.
Finally $L_{\kappa,b}(s)\asymp s^{-(2\kappa-1)}$ at infinity.
The prospective remainder transform would grow without bound if
$h>2\kappa-1$, which is impossible for a probability transform.
The ID construction attains this upper bound when $\rho\geq b$.
\end{proof}

\subsection{A global sign problem reduces to one evaluation}

Put $\delta_{\min}=\max(0,1-\kappa)$.  For
$\delta_{\min}<\delta<1$ set $c=\kappa+\delta-1>0$ and define
\begin{align}
 g_{\kappa,\delta}(x)
   &=\frac{x^{c-1}e^{-2x}}{\Gamma(c)}
       B_{\kappa,\delta}(2x),\label{p08:be14:inverse}\\
 B_{\kappa,\delta}(z)
   &={}_1F_1(\delta-1;c;z)
       +\frac{z}{2c}{}_1F_1(\delta;c+1;z).
                                                        \label{p08:be14:B}
\end{align}
Termwise Laplace integration and Kummer's transformation yield
\begin{equation}\label{p08:be14:inverse-transform}
 \int_0^\infty e^{-sx}g_{\kappa,\delta}(x)\,dx
       =s^{-\delta}(s+1)(s+2)^{-\kappa},\qquad s>0.
\end{equation}
For example one can invert separately
$s^{1-\delta}(s+2)^{-\kappa}$ and
$s^{-\delta}(s+2)^{-\kappa}$; their inverses are respectively
$x^{c-1}{}_1F_1(\kappa;c;-2x)/\Gamma(c)$ and
$x^c{}_1F_1(\kappa;c+1;-2x)/\Gamma(c+1)$.
All tilted signed densities are absolutely integrable.  For
$0<\delta<1$, the real-axis Kummer expansion gives
\begin{equation}\label{p08:be14:tail}
 g_{\kappa,\delta}(x)
       \sim\frac{2^{-\kappa}}{\Gamma(\delta)}x^{\delta-1}
       \quad(x\to\infty).
\end{equation}

\begin{lemma}[A one-point positivity criterion]\label{p08:be14:onepoint}
If $\kappa>1/2$ and $\delta_{\min}<\delta<\min(1,\kappa)$, then
\begin{equation}\label{p08:be14:onepoint-equivalence}
 g_{\kappa,\delta}(x)\geq0\text{ for all }x>0
 \quad\Longleftrightarrow\quad
 B_{\kappa,\delta}(2(\kappa-\delta))\geq0.
\end{equation}
If equality holds on the right, the only zero of the density is
$x=\kappa-\delta$, and it is a zero of order exactly two.
\end{lemma}
\begin{proof}
Write $p=1-\delta\in(0,1)$ and
$y(z)={}_1F_1(-p;c;z)$.  The derivative identity for Kummer's
function gives $y'(z)=-(p/c){}_1F_1(1-p;c+1;z)<0$ for $z>0$.
Moreover
\[
 B=y-\frac{z}{2p}y',\qquad
 zy''+(c-z)y'+py=0.
\]
At a zero of $B$, one has $y=zy'/(2p)<0$.  Substitution into the
differential equation gives the exact crossing rule
\begin{equation}\label{p08:be14:crossing}
 B'(z)=\frac{y'(z)}{4p}\{2(\kappa-\delta)-z\}
                    \quad\text{whenever }B(z)=0.
\end{equation}
Consequently every zero below $2(\kappa-\delta)$ is a strict
downward crossing and every zero above it is a strict upward
crossing.  If a zero is exactly at that point, differentiating
Kummer's equation once gives
\begin{equation}\label{p08:be14:double}
 B'(z)=0,\qquad B''(z)=-\frac{y'(z)}{4p}>0.
\end{equation}
The function is positive near zero because $B(0)=1$, and is
positive for all sufficiently large $z$ by \eqref{p08:be14:tail}.
Any negative component must therefore start with a downward
crossing below $2(\kappa-\delta)$ and end with an upward crossing
above it.  The double-zero rule excludes a negative component
having that point as an endpoint.  Thus a negative component
exists exactly when the value at that point is negative.  If
the value is zero, no other zero is compatible with the crossing
directions and the positive endpoint signs.  Equation
\eqref{p08:be14:double} proves that the remaining zero has order two.
\end{proof}

\begin{theorem}[Exact capacity with an arbitrary remainder]
\label{p08:be14:capacity}
For each $\kappa>1/2$ there is a unique number
\[
             \max(0,1-\kappa)<\delta_*(\kappa)<\frac1{4\kappa}
\]
satisfying
\begin{equation}\label{p08:be14:scalar-root}
 {}_1F_1(\delta-1;\kappa+\delta-1;2(\kappa-\delta))
 +\frac{\kappa-\delta}{\kappa+\delta-1}
 {}_1F_1(\delta;\kappa+\delta;2(\kappa-\delta))=0.
\end{equation}
The endpoint-rate capacity is attained and equals
\begin{equation}\label{p08:be14:capacity-result}
 h_{\rm P}(b-1)=\kappa-\delta_*(\kappa)
              >\kappa-\frac1{4\kappa}=h_{\rm ID}(b-1).
\end{equation}
It is independent of $b$.  For $h=\kappa-\delta_*(\kappa)$ the
unique probability remainder has density
\[
 \frac{e^{-(b-1)x}g_{\kappa,\delta_*}(x)}
 {(b-1)^{-\delta_*}b(b+1)^{-\kappa}},\qquad x>0.
\]
This density is strictly positive except for its double zero at
$x=h$.  Every shape $h_{\rm ID}(b-1)<h\leq h_{\rm P}(b-1)$
therefore gives an independent probability remainder which is not ID.
\end{theorem}
\begin{proof}
The transform of a putative remainder after extracting shape
$h=\kappa-\delta$ and rate $b-1$ is the normalization of
$s^{-\delta}(s+1)(s+2)^{-\kappa}$ translated by $b-1$.
Uniqueness of finite signed Laplace transforms and
\eqref{p08:be14:inverse-transform} show that extraction is equivalent
to $g_{\kappa,\delta}\geq0$ and is independent of the tilt.
When $\kappa>1$, at $\delta=0$ the inverse simplifies to
\begin{equation}\label{p08:be14:gzero}
 g_{\kappa,0}(x)=
       \frac{e^{-2x}x^{\kappa-2}}{\Gamma(\kappa)}
                    (\kappa-1-x).
\end{equation}
It is negative at $x=\kappa$, so continuity excludes all
sufficiently small positive $\delta$.  Extraction of any $h>\kappa$
would imply extraction of $\kappa$ and is likewise impossible.
When $1/2<\kappa\leq1$, the general near-zero shape bound gives
$h\leq2\kappa-1$, or $\delta\geq1-\kappa$.  At equality the
unnormalized candidate transform has the expansion
\[
 s^{\kappa-1}(s+1)(s+2)^{-\kappa}
       =1-\frac{2\kappa-1}{s}+O(s^{-2}).
\]
Its derivative is positive for large $s$, impossible for a
translated Laplace transform.  Fixing such a large $s$ and using
continuity in $\delta$ excludes a neighborhood above
$\delta_{\min}=1-\kappa$ as well.

Set $\delta_0=1/(4\kappa)$.  Proposition~\ref{p08:be14:id} gives an
ID remainder at $\delta_0$.  Its L\'evy density is
\[
 \frac{\kappa e^{-(b-1)x}}x
          \left(e^{-x}-\frac1{2\kappa}\right)^2.
\]
Its continuous density is strictly positive at every $x>0$.
Here is a direct justification of this strictness.  Split its
L\'evy measure at a sufficiently small $\eta>0$.  The sum of
jumps at most $\eta$ has mean tending to zero with $\eta$ and
thus has positive probability below $x/2$.  The compound Poisson
part from jumps above $\eta$ has, on the event of exactly two
jumps, a density strictly positive on $(2\eta,\infty)$: its
one-jump density is positive except at one point.  Their
convolution gives a positive continuous lower bound in a
neighborhood of $x$, if $\eta<x/8$.  This proves the assertion
also at a specified point, not merely almost everywhere.
In particular the right side of \eqref{p08:be14:onepoint-equivalence}
is strictly positive at $\delta_0$.  By continuity and the lemma,
some $\delta<\delta_0$ is already feasible.

Feasibility is upward closed in $\delta$, by convolution with an
additional gamma factor.  A neighborhood above $\delta_{\min}$ has just been excluded.
On the remaining compact interval ending at $\delta_0$, feasibility
is closed by continuity of the one-point test, so it has a
least point $\delta_*\in(\delta_{\min},\delta_0)$.  At that point the test
equals zero; otherwise continuity would allow a smaller value.
The one-point lemma gives the stated unique double zero.
If $\delta>\delta_*$, then
\[
 g_{\kappa,\delta}(x)=
 \frac1{\Gamma(\delta-\delta_*)}
 \int_0^x(x-y)^{\delta-\delta_*-1}
                 g_{\kappa,\delta_*}(y)\,dy>0.
\]
This strict inequality follows from positivity of the integrand
on an interval near zero.  Thus \eqref{p08:be14:scalar-root} has no
second solution in $(\delta_{\min},\delta_0)$, and all larger $\delta$ are
feasible.  The explicit density and its normalization follow
from \eqref{p08:be14:inverse-transform}.  Finally the exact ID
capacity rules out an ID remainder above it.
\end{proof}

\subsection{Exponential versus algebraic distance from the tail shape}

\begin{theorem}[Sharp large-order separation]\label{p08:be14:asymptotic}
As $\kappa\to\infty$,
\begin{equation}\label{p08:be14:delta-asymptotic}
 \delta_*(\kappa)
 \sim\frac{2^{\kappa-1}\kappa^{\kappa-1}e^{-2\kappa}}
               {\Gamma(\kappa)}
 \sim\frac1{2\sqrt{2\pi\kappa}}\left(\frac2e\right)^\kappa.
\end{equation}
In contrast, $\kappa-h_{\rm ID}(b-1)=1/(4\kappa)$ exactly.
Thus
\[
 \frac{\kappa-h_{\rm P}(b-1)}
      {\kappa-h_{\rm ID}(b-1)}
 \sim\sqrt{\frac{2\kappa}{\pi}}\left(\frac2e\right)^\kappa
 \longrightarrow0.
\]
This is a quantitative refinement of the preceding capacity
theorem, not a second historical conjecture.
\end{theorem}
\begin{proof}
Let
\[
 f_\kappa(y)=\frac{e^{-2y}y^{\kappa-1}}{\Gamma(\kappa)},\quad
 g_0=f_\kappa'+f_\kappa,\quad
 A_\kappa=-g_0(\kappa)
             =\frac{e^{-2\kappa}\kappa^{\kappa-2}}{\Gamma(\kappa)}.
\]
For $\delta>0$, \eqref{p08:be14:inverse-transform} gives the
fractional-integral identity
\begin{equation}\label{p08:be14:fractional}
 g_{\kappa,\delta}(x)=
 \frac{x^\delta}{\Gamma(1+\delta)}g_0(x)
 +\frac1{\Gamma(\delta)}
 \int_0^x(x-y)^{\delta-1}\{g_0(y)-g_0(x)\}\,dy.
\end{equation}
The subtraction makes the integral regular at its upper limit.
We need the estimate
\begin{equation}\label{p08:be14:J-estimate}
 J_\kappa:=\int_0^\kappa
       \frac{g_0(y)-g_0(\kappa)}{\kappa-y}\,dy
       =\frac{2^{1-\kappa}}\kappa(1+O(\kappa^{-1})).
\end{equation}

To prove it, split at $3\kappa/4$.  On the first interval,
integration by parts in the $f_\kappa'$ term gives
\begin{align*}
 \int_0^{3\kappa/4}\frac{g_0(y)}{\kappa-y}\,dy
 &=\frac{f_\kappa(3\kappa/4)}{\kappa/4}\\
 &\quad+\int_0^{3\kappa/4}f_\kappa(y)
       \left\{\frac1{\kappa-y}-\frac1{(\kappa-y)^2}\right\}\,dy.
\end{align*}
The measure $2^\kappa f_\kappa(y)dy$ is
$\operatorname{Gamma}(\kappa,2)$, with mean $\kappa/2$ and
variance $\kappa/4$.  Taylor expansion of the denominators on
$y\leq3\kappa/4$, together with the gamma exponential tail
bound, evaluates the last integral as
$2^{-\kappa}\{2/\kappa+O(\kappa^{-2})\}$; the boundary term
is exponentially smaller.  Such a tail bound follows directly
from Markov's inequality and the gamma moment generating function.
On $[3\kappa/4,\kappa]$, the divided difference of $g_0$ is at
most the supremum of $|g_0'|$ there.  The explicit formula for
$f_\kappa$ and Stirling's formula bound this contribution by
$2^{-\kappa}e^{-c\kappa}$ times a fixed polynomial in $\kappa$,
for a fixed $c>0$.  Finally the subtracted constant on
$[0,3\kappa/4]$ contributes only $g_0(\kappa)\log4$, again
exponentially smaller than $2^{-\kappa}/\kappa$.  This proves
\eqref{p08:be14:J-estimate}.

Set $d_\kappa=A_\kappa/J_\kappa$.  For every fixed $C>0$, the
same split in \eqref{p08:be14:fractional} gives uniformly for
$0\leq\delta\leq C d_\kappa$, with $x=\kappa-\delta$,
\begin{equation}\label{p08:be14:uniform-expansion}
 g_{\kappa,\delta}(\kappa-\delta)
       =-A_\kappa(1+o(1))+\delta J_\kappa(1+o(1)).
\end{equation}
For completeness, on $y\leq3\kappa/4$ the kernel changes from
$(\kappa-y)^{-1}$ by the uniform relative amount
$O(\delta\log\kappa+\delta/\kappa)=o(1)$.  On the remaining
interval the subtracted divided difference is bounded by
$\sup|g_0'|$, and the integral of $v^{\delta}$, including its
logarithmic derivative in $\delta$, is bounded by a polynomial
in $\kappa$ and $\log\kappa$.  The exponential tail estimate
above therefore remains negligible uniformly.  Also
$g_0(\kappa-\delta)/g_0(\kappa)=1+o(1)$,
$x^\delta/\Gamma(1+\delta)=1+o(1)$, and
$1/\Gamma(\delta)=\delta(1+o(1))$.  These observations prove
\eqref{p08:be14:uniform-expansion} without interchanging an uncontrolled
limit with a supremum.

For each fixed $\varepsilon\in(0,1)$, the one-point test is
negative at $(1-\varepsilon)d_\kappa$ and positive at
$(1+\varepsilon)d_\kappa$ for all sufficiently large $\kappa$.
The exact capacity theorem brackets $\delta_*$ between these
two values.  Hence $\delta_*\sim d_\kappa$.  Inserting
\eqref{p08:be14:J-estimate} and then Stirling's formula gives
\eqref{p08:be14:delta-asymptotic}.
\end{proof}

\paragraph{Scope and remaining gaps.}
The new exact boundary concerns arbitrary probability remainders on
the first excess-power line and the slowest allowable gamma rate.
The all-rate ID formula is an elementary supporting classification;
the exponential law quantifies a large separation which the ID
criterion misses.  Neither the unrestricted capacity for every
noninteger excess power nor its full intermediate-rate diagram is
settled here.  An attempted transfer to universal SD strictness for
the squared-Bessel McKay family remains open.  In particular, no
termination claim is made for the earlier finite moment hierarchy.
The noncentral chi-square HCM classification is already proved in
\cite{p01:rare:Salazar2026} and is not counted as our new result.
\section{Remaining problems and the limits of the transfers}
\label{p08:r14:remaining}

The strongest external step is the universal adaptive-width bound.
For example it gives $N(4)=3$, $N(8)=4$, and $N(16)=5$;
for dimension three the bound remains $2\le N(3)\le3$.
The scalar-fiber recursion reaches $\operatorname{UW}_3$ after two
queries, which vanishes on a three-dimensional compact polyhedron.
Thus that recursion alone cannot decide the three-dimensional case.
The Betti-sensitive theorem adds a real restriction, but analytic
functions can have increasingly complicated fibers on expanding balls.
A necessary topology-growth condition is not a contradiction.

The statistical advance changes computation and the corruption model,
not the previously proved minimax order. Its linear program is
polynomial in the explicit input length but can be large. A nearly
linear algorithm, a practical implementation with comparable constants,
and a robustness theorem for the unmodified nonparametric MLE remain
separate problems. The best universal Chebyshev--Bernstein exponent
is also unresolved: the available interval remains $\log2\le c_*\le1$.

The prime-series inverse theorem identifies all positive weights on
a prescribed frequency support, using infinitely many exact zeros in
one height direction. Its proof combines phase freezing with the
derivative of a frozen simple zero. A stability theorem for finitely
many noisy zeros would need effective recurrence and conditioning
bounds that are absent here. For fluctuations, the exact pure-point
variance expansion makes the missing task explicit: control the total
Fourier weight at very small nonzero frequencies. Recurrent local
correlations do not alone prohibit a central limit theorem.

The matrix construction closes a qualitative arbitrary-design gap,
but it supplies no uniform inverse modulus near the quasi-analytic
boundary. In stable models, compact scale support sharply distinguishes
indices below and above one; skew index one leads to a different
entire-function sampling problem, still open here. Mixing over unknown
stable indices or skewness is not covered by the fixed-shape theorems.

The finite-zero spectral experiment assumes complete zero lists,
controlled amplitudes including $A=0$, and quantitative finite-order
nonresonance. Missing or spurious zeros, efficient certified fitting,
an optimal lower bound on the necessary height, and an infinite-prime
version require further work. The classical recurrence theorem improves
a sufficient height; it does not prove that this height is optimal for
the inverse zero problem.

The gamma-capacity calculation resolves the first excess-power line
at the slowest allowable gamma rate. Arbitrary noninteger excess powers,
the unrestricted intermediate-rate capacity diagram, universal strict
self-decomposability for the squared-Bessel family, and termination of
the earlier finite moment hierarchy remain unresolved. The scalar
crossing rule that reduces global positivity to one evaluation is the
specific structure behind the present exact solution.

\paragraph{Research status.}
The manuscript records twelve attained targets at the proof-draft
level and explicitly preserves the stronger questions above. Its claims
have undergone automated cross-review, with no blocking gap identified
in the final review round. This status is narrower than independent
external validation, and none of the target counts is a claim to an
Annals acceptance or a major-prize-level result.
\endgroup
% ===== END SOURCE 08 =====


% ===== BEGIN SOURCE 09: KEIO_Adaptive_Recovery_and_Mixture_Identification.tex =====
\clearpage
\part[Adaptive Recovery and Mixture Identification]{Sharp Boundaries for Adaptive Recovery and Mixture Identification}
\label{part:09}
\setcounter{section}{0}
\begin{center}\large\itshape Homogeneous Information, Transport Corruption, Stable Mixtures, and Spectral Edge Laws\end{center}
\begin{center}\small Source manuscript: 16 October 2026.\end{center}
\begingroup
\renewcommand{\arraystretch}{1}
\begin{quote}\small\textbf{Source abstract.}\quad
We develop thirteen scoped targets extending an observation-and-recovery
program. Positive homogeneity of the first continuous measurement forces
at least $\lceil\log_2 m\rceil+1$ queries for bounded-error recovery on
$\mathbb R^m$, while the fully positively homogeneous class still has logarithmic
complexity. Every finite adaptive budget has a strictly positive gap
above the exact Euclidean convex-noise floor in dimension at least two.
For a bounded Hilbert operator, the limiting adaptive continuous error
is exactly its essential norm; spectral bounds also prove optimal
exponential and double-exponential recovery orders. A new predictive
transport modulus yields the full binomial-population minimax curve
under a hard budget of individual trial flips, with a polynomial-time
estimator. Compact-scale skew stable mixtures at index one are classified
by the completeness radius of logarithmic terminal times. For symmetric
zero-location stable mixtures with continuously unknown index and scale,
a mild two-sided scale moment makes every infinite terminal design
identifying, while every finite design fails. Further results quantify
frequency-relation losses and zero intensities at spectral edges, make
nonlinear spectral fitting constructive, and determine the strict
intermediate-rate gamma factor region and its quartic separation from
infinitely divisible remainders. Classical theorems, prior unpublished
dependencies, related refinements, and unresolved stronger questions are
explicitly separated. The target count is not a count of historical
conjectures solved.
\end{quote}

\section*{Research scope and thirteen targets}
\addcontentsline{toc}{section}{Research scope and thirteen targets}

This paper continues the 223-page integrated manuscript and its sequels,
most immediately the 39-page \emph{Geometry and Rigidity of Inverse
Observation} \cite{comp:part:08}. It contains new arguments and
extensions rather than a reprint of that manuscript. The common question
is how restrictions on observations change identifiability, optimal
precision, and computational recovery. The sections use different
models; a general equivalence between them is not asserted.

The restricted-information question in
\cite[Section~3]{p03:an:KNU} is externally posed. Our homogeneous
result answers its order-of-growth aspect and several exact cases;
the exact count in every dimension remains open. The remaining rows
are quantitative transfers or open questions within this research
program. Closely related consequences are grouped: for example, the
two gamma asymptotic regimes occupy one row, and spectral decay
corollaries are grouped with the operator theorem. No claim is made
that these thirteen rows are thirteen independent major conjectures.

{\small\setlength{\tabcolsep}{4pt}
\begin{longtable}{@{}p{.045\textwidth}p{.29\textwidth}p{.48\textwidth}p{.12\textwidth}@{}}
\toprule
No. & Target & Result and limitation & Proof\\
\midrule
\endfirsthead
\toprule No. & Target & Result and limitation & Proof\\\midrule
\endhead
1 & Homogeneous information complexity & Logarithmic upper and lower counts; exact low-dimensional and hybrid cases. General exact count retains a gap. & \ref{p09:in15:homogeneous-first}, \ref{p09:in15:homogeneous-count}\\
2 & Attainment of the convex-noise floor & Every fixed finite adaptive budget has a uniform positive gap in dimension $m\ge2$. The gap's sharp rate is not obtained. & \ref{p09:in15:finite-gap}\\
3 & Infinite-dimensional adaptive recovery & Essential norm is the exact limiting error; two-sided spectral bounds prove sharp decay orders. & \ref{p09:op15:theorem}\\
4 & Predictive transport inversion & Sharp binomial-mixture inverse modulus, with logarithm $\log(e+1/(\eta\sqrt t))$. & \ref{p09:ml15:modulus}\\
5 & Individual-trial adversarial flips & Full sharp minimax curve and polynomial-time estimator under a known hard flip budget. & \ref{p09:ml15:bitrate}\\
6 & Frequency relations and edge loss & Explicit positive relation matrix and fourth-order remainder for the local boundary loss. & \ref{p09:ar15:thm:relation}\\
7 & Dimension visible in zero intensity & Sharp edge intensity and leading constant; compatible hulls reveal rational rank. & \ref{p09:ar15:thm:intensity}, \ref{p09:ar15:thm:rank}\\
8 & Compact-scale skew index one & Exact completeness-radius criterion, including smooth joint counterexamples. & \ref{p09:ga15:thm:skew}\\
9 & Jointly mixed stable index and scale & Every infinite design identifies under the stated moments; every finite design has counterexamples. & \ref{p09:ga15:thm:unknown}\\
10 & Constructive nonlinear spectral fitting & Threshold inversion and a finite moment reconstruction under separation and small noise. & \ref{p09:sp15:thm:algorithm}\\
11 & Controlled critical zero counts & A local near-threshold intensity law; no claim of a uniform finite-height discrepancy bound. & \ref{p09:sp15:thm:intensity}\\
12 & Intermediate-rate gamma factors & Attained continuous strictly monotone boundary, strictly above the ID boundary at every interior rate. & \ref{p09:be15:region}\\
13 & Singular gamma-capacity gaps & Exact quartic laws at the central rate and lowest Bessel order; refinement of row 12. & \ref{p09:be15:quartic}, \ref{p09:be15:loworder}\\
\bottomrule
\end{longtable}
}

\paragraph{How to interpret the results.}
The new arguments are supplied in full for their stated models. Some
inputs remain prior unpublished results: the adaptive-width recursion,
the exact convex-noise floor, the Bernstein approximation estimates and
clean binomial lower bounds, the previous spectral parameter stability
bound, and the endpoint gamma capacity. Those dependencies are identified
where used. The KNU construction, Euclidean widths, Beurling--Malliavin
theorem, moment-determinate projection uniqueness, and classical moment
reconstruction are credited as existing tools. Checked primary sources
support these comparisons; the literature search does not certify
historical priority of every transfer.

There is no proof here of the Riemann hypothesis, no determination of
all primes from a rightmost zeta zero, and no unrestricted solution of
the two-query problem in dimension three. Exact population laws and
exact real measurements are mathematical observation models, not claims
of finite-precision experimental feasibility. Proof review within this
research process is not external peer review or formal verification.



% ===== information =====
\section{Homogeneous first measurements and the necessity of infinitely many noisy convex queries}
\label{p09:in15:section}

The two results in this section address different losses of information.
In the exact-real model, positive homogeneity of the first query costs
a full dimension in a way that survives arbitrary subsequent adaptive
choices.  This gives a stronger lower count than the unrestricted
continuous bound in every dimension that is not a power of two.
In the bounded-noise model, the exact Euclidean convex-information floor
from the earlier manuscript is a genuine infinite-query limit: in
dimension at least two, every fixed finite query budget has a strictly
larger minimax error, even with unrestricted adaptivity.

\subsection{Large level patches and the homogeneous-query question}

A continuous function $f:\mathbb R^m\to\mathbb R$ is positively
homogeneous of degree $a>0$ if $f(rx)=r^a f(x)$ for all $r>0$.
This convention does not require oddness.  It includes norms and
even-degree homogeneous polynomials.  The distinction is relevant to
the homogeneous-information question in
\cite[Section~3]{p03:an:KNU}.

Krieg and Kritzer study homogeneous full information maps and recovery
maps, and show that on cones a finite-error procedure using homogeneous
measurements can be rescaled to arbitrarily small error at unchanged
cost \cite[Propositions~1 and~2]{p09:in15:KK}.  The positive-scaling proof
of their second statement applies to positive homogeneity as well.
Here we quantify the number of individually continuous, adaptively
selected measurements; their joint transcript need not be homogeneous.
We do not resolve their general optimality question for homogeneous
algorithms \cite[Open Problem~1]{p09:in15:KK}.

We use the adaptive-width bound from the preceding manuscript
\cite{comp:part:08}: on a compact metric space $X$, every $q$-query
procedure using arbitrary continuous scalar measurements has error at
least $\operatorname{UW}_{2^q-1}(X)/2$.  Here $\operatorname{UW}_d$
is the infimum of the mesh of finite open covers of multiplicity at
most $d+1$.  Its short proof iterates the classical inequality
\[
 \operatorname{UW}_{2d+1}(X)
       \le\sup_t\operatorname{UW}_d(f^{-1}(t))
\]
for a continuous scalar map on a compact metric space; see
\cite[Theorem~2.1]{p08:in14:BB}.  Compact fibers and arbitrary selection
rules are allowed.  The Lebesgue covering theorem gives
$\operatorname{UW}_{d-1}([-h,h]^d,\|\cdot\|_\infty)=2h$.

\begin{lemma}[A large-patch obstruction]\label{p09:in15:patch}
Let $f:\mathbb R^m\to\mathbb R$ be continuous.  Suppose there are
levels of $f$ containing compact sets $X_R$ and continuous bijections
$\gamma_R:[-1,1]^d\to X_R$ such that
\begin{equation}\label{p09:in15:patch-expansion}
 \|\gamma_R(u)-\gamma_R(v)\|_2\ge R\|u-v\|_\infty,
                         \qquad R\longrightarrow\infty.
\end{equation}
Any $n$-query procedure beginning with $f$ and using arbitrary
continuous subsequent queries has infinite uniform recovery error
whenever $2^{n-1}\le d$.
\end{lemma}
\begin{proof}
The inverse of $\gamma_R$ is continuous and $R^{-1}$-Lipschitz by
\eqref{p09:in15:patch-expansion}.  Pulling a cover of $X_R$ back to the
cube shows
$\operatorname{UW}_{d-1}(X_R)\ge2R$.  After the first answer is
fixed, the remaining procedure has $n-1$ queries on $X_R$.
Since $2^{n-1}-1\le d-1$, the adaptive-width bound gives error at
least $R$.  Let $R$ tend to infinity.
\end{proof}

\begin{theorem}[A homogeneous first query forces the stronger count]
\label{p09:in15:homogeneous-first}
Let $m\ge2$.  If the first query is continuous and positively
homogeneous of any degree $a>0$, uniformly bounded-error recovery on
$\mathbb R^m$ requires
\begin{equation}\label{p09:in15:homogeneous-lower}
                         n\ge\lceil\log_2 m\rceil+1.
\end{equation}
Every later query may be an arbitrary continuous function, with no
homogeneity, degree, or definability restriction.  In particular, two
queries cannot suffice in $\mathbb R^3$ under this first-query
assumption.
\end{theorem}
\begin{proof}
Suppose first that $f$ is not identically zero.  Change its sign and
rotate coordinates so that $f(e_m)>0$.  By continuity, on a sufficiently
small cube $C=[-h,h]^{m-1}$ the function $v\mapsto f(v,1)$ is bounded
above and below by strictly positive constants.  Put
\[
 b(v)=f(v,1)^{-1/a},\qquad
 \gamma_\lambda(v)=\lambda b(v)(v,1),\qquad \lambda>0.
\]
Positive homogeneity gives $f(\gamma_\lambda(v))=\lambda^a$.
The map $\gamma_\lambda$ is injective because its inverse is the ratio
map $x\mapsto x'/x_m$.  On its image, $x_m$ is between two positive
constant multiples of $\lambda$, and $\|x'\|_2\le C\lambda$.
Consequently, for $x=\gamma_\lambda(v)$ and
$y=\gamma_\lambda(w)$,
\[
 \|v-w\|_2
 \le \frac{\|x'-y'\|_2}{x_m}
     +\frac{\|y'\|_2|x_m-y_m|}{x_my_m}
 \le\frac C\lambda\|x-y\|_2.
\]
Rescaling $C$ to $[-1,1]^{m-1}$ verifies
Lemma~\ref{p09:in15:patch} with $d=m-1$ and an expansion factor tending
to infinity.  Thus $2^{n-1}\le m-1$ is impossible.  Since the left
side is an integer, its negation is exactly
$n\ge\lceil\log_2m\rceil+1$.
If $f\equiv0$, the first query contributes no information; the
unrestricted adaptive-width bound on expanding $m$-cubes gives at
least the same conclusion.
\end{proof}

\begin{corollary}[Logarithmic complexity within the homogeneous class]
\label{p09:in15:homogeneous-count}
Let $N_{\rm hom}(m)$ denote the minimum number of continuous,
positively degree-one homogeneous adaptive scalar queries sufficient
for arbitrarily accurate uniform recovery on $\mathbb R^m$.
For $m\ge2$,
\begin{equation}\label{p09:in15:homogeneous-sandwich}
 \lceil\log_2m\rceil+1\le N_{\rm hom}(m)
           \le\min\{m,\lceil\log_2m\rceil+2\}.
\end{equation}
Hence positive homogeneity preserves logarithmic query complexity.
The count is exactly two for $m=2$ and exactly three for $m=3$.
For larger dimensions these bounds generally retain a one-query gap.
\end{corollary}
\begin{proof}
The lower bound is Theorem~\ref{p09:in15:homogeneous-first}; coordinate
queries give the upper bound $m$.  For the other upper bound, first
query $r=\|x\|_2$.  If $r=0$, output zero.  Otherwise apply the
Krieg--Novak--Ullrich continuous algorithm \cite[Theorem~1]{p03:an:KNU}
to $u=x/r$, at accuracy $\varepsilon/r$, using at most
$\lceil\log_2m\rceil+1$ further queries.  Each particular continuous
query $g$ selected by that algorithm, restricted to the unit sphere,
can be implemented by
\[
 H_g(x)=\begin{cases}
      \|x\|_2 g(x/\|x\|_2),&x\ne0,\\
      0,&x=0.
      \end{cases}
\]
It is positively homogeneous of degree one.  It is continuous at zero
because $g$ is bounded on the compact sphere.  Dividing its answer by
the already known $r$ recovers the intended answer $g(u)$.  The selected
function $g$ may depend on $r$ and on previous answers, as allowed by
the adaptive model.  Multiplying the final approximation to $u$ by $r$
gives error at most $\varepsilon$.
\end{proof}

\paragraph{Why the sign convention matters.}
If every admissible query instead satisfies $f(rx)=rf(x)$ for every
real $r$, exactly $m$ queries are necessary and sufficient.  Indeed,
freeze the $n$ continuous odd functions selected along the all-zero
transcript.  For $n<m$, Borsuk--Ulam applied to their joint map on
$S^{m-1}$ gives a common zero $u$.  Every point $ru$, $r>0$, then
follows that same transcript by homogeneity, so its recovery error
is unbounded.  Coordinate measurements give the matching upper
bound.  Thus the logarithmic conclusion above uses the larger class
allowing positive homogeneity without oddness, which includes the
norm and intensity measurements motivating \cite[Section~3]{p03:an:KNU}.

\begin{corollary}[An exact hybrid count in shifted dyadic dimensions]
\label{p09:in15:shifted-dyadic}
If only the first query is required to be positively homogeneous and
all subsequent queries may be arbitrary continuous functions, the
minimum count in dimension $m=2^k+1$ is exactly $k+2$ for $k\ge0$.
\end{corollary}
\begin{proof}
The lower bound follows from \eqref{p09:in15:homogeneous-lower}.
For the upper bound first read the last coordinate, and then apply
\cite[Theorem~1]{p03:an:KNU} to the remaining $2^k$ coordinates with
$k+1$ continuous queries.  The first coordinate query is homogeneous.
No homogeneity is asserted for the later queries in this hybrid model.
\end{proof}

\paragraph{Other large-patch classes.}
The same lower count \eqref{p09:in15:homogeneous-lower} holds when the
first query is a real polynomial or a continuous quasiconvex function.
Here the required $(m-1)$-dimensional level graphs were established in
the previous two-query obstruction, and the new step is their use with
an arbitrary number of remaining queries.  Briefly, for a polynomial,
choose a direction where its leading homogeneous term is nonzero.  On
each transverse ball its sufficiently distant positive tail is
strictly monotone, so a sufficiently high level is a graph over that
ball.  Orthogonal projection is a one-Lipschitz inverse.  For a
quasiconvex function, either a maximal level contains a halfspace, or
a proper convex sublevel set contains an arbitrarily large ball.
An upper boundary graph in a supporting direction gives the same
one-Lipschitz inverse.  These are level sets of the first query, so the
later queries remain unrestricted.

\subsection{A finite-query gap above the Euclidean convex-noise floor}

Fix $m\ge2$ and $0<\delta<1$.  On the unit Euclidean ball $B$, an
admissible query is a continuous convex function with range in $[-1,1]$.
Its observed value may differ from its true value by at most $\delta$.
Write $e_n^{\rm ad}$ for the optimal worst-case Euclidean recovery
error after at most $n$ adaptive queries, allowing arbitrary selection
rules and decoders.  Put
\begin{equation}\label{p09:in15:rstar}
 q=1-\delta,\qquad
 r_* =\frac{\delta}{\sqrt{1+q^2+2q/m}}.
\end{equation}
The exact unlimited-query value proved in the earlier manuscript is
$\inf_n e_n^{\rm ad}=r_*$ \cite{comp:part:07}.  The issue here is
whether any fixed finite budget attains this infimum.

Let $u_0,\ldots,u_m$ be the vertices of a centered regular simplex on
the unit sphere, with $\langle u_i,u_j\rangle=-1/m$ for $i\ne j$.
For $U\in SO(m)$ put $S_U=\{r_*Uu_i:0\le i\le m\}$.
For distinct vertices $x,y\in S_U$,
\begin{equation}\label{p09:in15:simplex-relation}
              y=qx+\delta z,\qquad z\in\partial B.
\end{equation}
Convexity and the range bound imply
$\operatorname{osc}_{S_U}f\le2\delta$ for every admissible query.

\begin{lemma}[No query saturates an open set of simplex rotations]
\label{p09:in15:nowhere-dense}
Let $f$ be convex on the open unit ball with $-1\le f\le1$.
Then
\[
 Z_f=\{U\in SO(m):\operatorname{osc}_{S_U}f=2\delta\}
\]
is closed and has empty interior.
\end{lemma}
\begin{proof}
A finite convex function is continuous in the open ball, so $Z_f$ is
closed.  Suppose it had nonempty interior.  It is the union over
finitely many ordered pairs $i\ne j$ of the closed sets on which
$f(r_*Uu_j)-f(r_*Uu_i)=2\delta$.  The Baire property of the locally
compact rotation group implies that one ordered pair has this equality
on a nonempty open set $V$.

Fix $U\in V$ and write $x=r_*Uu_i$, $y=r_*Uu_j$, and
$z=(y-qx)/\delta$.  The range bound and convexity imply
$f(y)-f(x)\le\delta(1-f(x))\le2\delta$, so equality forces
$f(x)=-1$ and $f(y)=-1+2\delta$.  This argument uses only interior
values: replace $z$ by a point just short of it on the segment from
$x$, apply convexity, and pass to the endpoint limit.

Take a supporting affine function $\ell(w)=b+v\cdot w$ for $f$ at
$y$.  Since $\ell\le f\le1$ on the open ball, continuity of $\ell$
gives $\ell(z)\le1$.  Also $\ell(x)\le-1$.  The identity
$y=qx+\delta z$ and $\ell(y)=-1+2\delta$ force
$\ell(x)=-1$ and $\ell(z)=1$.  Thus the nonzero linear functional
$v$ attains its maximum over the unit ball at $z$, whence
$v=c z$ for some $c>0$.

For all rotations in $V$, the corresponding vertex $r_*Uu_i$ is a
minimum point of $f$.  The map $U\mapsto Uu_i$ is open onto the
sphere, so these points contain a relatively open sphere patch around
$x$.  Since $\ell\le f=-1$ on this patch and $\ell(x)=-1$, the
linear functional $v$ has a local maximum on that sphere patch at
$x$.  Consequently $v$ is parallel to $x$.  But
$z=(r_*/\delta)U(u_j-qu_i)$ is not parallel to $x$ when $m\ge2$:
the distinct regular-simplex directions $u_i,u_j$ are then not
collinear.  This contradiction proves the lemma.
\end{proof}

\begin{theorem}[Every finite adaptive budget has a strict positive gap]
\label{p09:in15:finite-gap}
For every $m\ge2$, $0<\delta<1$, and finite integer $n\ge0$, there is
$c_{m,\delta,n}>0$ such that
\begin{equation}\label{p09:in15:finite-gap-eq}
                  e_n^{\rm ad}\ge r_*+c_{m,\delta,n}>r_*.
\end{equation}
In particular, the exact unlimited-query convex-noise floor cannot be
attained at any finite query count, even by taking an infimum over all
adaptive algorithms with that count.  The statement also applies to
any subclass of these queries, including globally convex analytic
functions and globally convex polynomials of arbitrary degree.
\end{theorem}

\begin{proof}
The case $n=0$ is immediate from $e_0^{\rm ad}=1>r_*$.  Fix $n\ge1$
and choose $r_*<\rho<1$.  All convex functions bounded by one in
absolute value on the open unit ball are uniformly Lipschitz on
$\rho B$, with Lipschitz constant
\[
                            L=\frac{2}{1-\rho}.
\]
Their restrictions to $\rho B$ form a compact family $\mathcal F$
in the uniform norm.  To verify closedness, take an arbitrary sequence
and use the same bounds on every smaller concentric ball, followed by
a diagonal subsequence.  Its limit is convex and bounded by one on the
open ball, and restricts to the prescribed uniform limit on $\rho B$.
Thus Lemma~\ref{p09:in15:nowhere-dense} applies also to limits in this
compact family.

Equip $SO(m)$ with the distance $\|U-V\|_{\rm op}$.  For $a>0$,
write $E_a=\{V\in SO(m):\|V-I\|_{\rm op}\le a\}$ and set
\begin{equation}\label{p09:in15:kappa}
 \kappa(a)=\min_{f\in\mathcal F,\,U\in SO(m)}
       \max_{V\in E_a}
          \{2\delta-\operatorname{osc}_{S_{UV}}f\}.
\end{equation}
The minimized function is continuous in $(f,U)$, and the parameter
space is compact.  Lemma~\ref{p09:in15:nowhere-dense} says that its value
at every $(f,U)$ is positive, because $UE_a$ contains a nonempty open
set.  Therefore $\kappa(a)>0$ for every $a>0$.

We construct a common noisy transcript for an entire ball of simplex
rotations.  Start with any rotation center and radius $a_0=1$.
Define deterministic positive radii by
\[
 a_j=\min\left\{\frac{a_{j-1}}2,
                 \frac{\kappa(a_{j-1}/2)}{4Lr_*}\right\},
 \qquad
 \sigma_n=\min_{1\le j\le n}\frac{\kappa(a_{j-1}/2)}4>0.
\]
At stage $j$, the previous answers fix the selected query $f_j$.
Use \eqref{p09:in15:kappa} to find a rotation within distance
$a_{j-1}/2$ of the current center at which its oscillation is at most
$2\delta-\kappa(a_{j-1}/2)$.  Answer by the midpoint of its maximum
and minimum on that simplex.  Every vertex value is then within
$\delta-\kappa(a_{j-1}/2)/2$ of the answer.  If the rotation changes
by at most $a_j$, a vertex moves by at most $r_*a_j$ and its query
value changes by at most $Lr_*a_j$.  The same answer is consequently
compatible, with slack at least $\kappa(a_{j-1}/2)/4$, with every
simplex in the new rotation ball.  That ball lies in the previous one
by the triangle inequality, so every earlier answer and its slack are
preserved.  This induction uses no continuity of the query-selection
rule: at each stage the transcript has already fixed one query.

After $n$ stages, choose any rotation $U$ in the retained ball.  Every
query value on $S_U$ is within $\delta-\sigma_n$ of its corresponding
answer.  Increase the simplex radius by
\[
 h=\min\left\{\frac{\rho-r_*}{2},\frac{\sigma_n}{2L}\right\}>0.
\]
All enlarged vertices remain in $\rho B$, and each recorded query
value changes by at most $Lh\le\sigma_n/2$.  Thus the same transcript
is possible for every vertex of the regular simplex
$\{(r_*+h)Uu_i\}_{i=0}^m$.  Its circumradius is $r_*+h$, so every
decoder has worst-case error at least that value.  The number $h$
depends only on $m,\delta,n$ and the fixed choice of $\rho$, not on
the algorithm.  This proves \eqref{p09:in15:finite-gap-eq}.
\end{proof}

\paragraph{The one-dimensional exception.}
For $B=[-1,1]$ put $r=\delta/(2-\delta)$.  The two fixed convex
queries
\[
 f_+(x)=-1+\frac{2(x+r)_+}{1+r},\qquad
 f_-(x)=-1+\frac{2(r-x)_+}{1+r}
\]
have range in $[-1,1]$ and attain the exact floor $r$.  Indeed, if
$x<y$ are compatible with the same two noisy answers and $x\ge-r$,
the first query implies $y-x\le\delta(1+r)=2r$.  If $y\le r$,
the second gives the same bound.  In the only remaining case
$x<-r<r<y$, the first query difference already exceeds $2\delta$,
which is impossible.  Every feasible set therefore has diameter at
most $2r$, and its interval midpoint gives error at most $r$.  The
universal indistinguishable pair $\{-r,r\}$ gives the matching lower
bound.  Thus the finite-query obstruction is genuinely dimensional.

\paragraph{Status and remaining scope.}
The homogeneous-first lower bound and logarithmic homogeneous-class
construction address the externally posed restricted-information
question.  The exact hybrid count is a consequence, not a closure of
the unrestricted KNU problem.  The finite-query gap settles the
attainment question left open by the preceding exact convex-noise
formula.  Its positive constant is established by compactness; neither
its sharp dependence on $n$ nor a sharp finite-query complexity law is
claimed.  Unrestricted two-query recovery in $\mathbb R^3$, and the
exact positive-homogeneous count in higher dimensions, remain open.

% ===== operators =====
\section{The exact compactness barrier for adaptive Hilbert-space recovery}
\label{p09:op15:section}

The earlier adaptive-width bound supplies a missing converse to the
known exponential speed-up for Hilbert operators. Arbitrary continuous
adaptive information cannot remove the essential-norm obstruction.
For compact operators with regularly decaying singular values, the
same argument proves that the exponential speed-up has the optimal
order, including its exponent for algebraic decay.

Let $S:H\to G$ be a bounded linear operator between real Hilbert
spaces. Measurements are arbitrary continuous real functions on the
closed unit ball $B_H$, chosen adaptively; the decoder is unrestricted.
Write $e_n(S)$ for the infimum of their worst-case errors after $n$
exact measurements. Define the approximation numbers by
\[
 a_j(S)=\inf_{\operatorname{rank}R<j}\|S-R\|,\qquad j\ge1,
\]
and the essential norm by
$\|S\|_{\rm ess}=\inf_{K\text{ compact}}\|S-K\|$.
These conventions make $a_j$ the $j$-th singular value in the compact
case. In finite dimension extend the sequence by zeros.

\begin{theorem}[Two-sided spectral bounds and the exact error floor]
\label{p09:op15:theorem}
For every $n\ge1$,
\begin{equation}\label{p09:op15:bounds}
 \max\left\{
 \sqrt{\frac{1+2^{-n}}2}\,a_{2^n}(S),\quad
 a_{2^{n+1}}(S)\right\}
 \le e_n(S)\le a_{2^{n-1}+1}(S).
\end{equation}
Consequently
\begin{equation}\label{p09:op15:essential}
                  \lim_{n\to\infty}e_n(S)=\|S\|_{\rm ess}.
\end{equation}
Thus recovery to arbitrarily small uniform error by finitely many
adaptive continuous measurements is possible if and only if $S$ is
compact. This statement permits discontinuous selection rules and
decoders; imposing their continuity is not part of its proof.
\end{theorem}

The upper exponential index comes from the existing finite-dimensional
construction of Krieg--Novak--Ullrich \cite{p03:an:KNU}. Their
Corollary~11 already gives a Hilbert-operator upper bound with index
$2^{n-2}$. The improvement of that index here is a direct use of their
sharper finite-dimensional construction and is not the main novelty.
The contribution is the converse and the exact essential-norm floor.
The Euclidean width constants used below are classical
\cite{p09:op15:Gournay}, and the adaptive-width mechanism is from our
preceding manuscript \cite{comp:part:08}.

\begin{lemma}[The two Euclidean widths used in the proof]
\label{p09:op15:widths}
For the Euclidean unit ball $B_2^m$,
\[
 \operatorname{UW}_{m-1}(B_2^m)=\sqrt{2(1+1/m)}\quad(m\ge2),
 \qquad
 \operatorname{UW}_{d}(B_2^m)=2\quad(0\le d\le m/2-1).
\]
\end{lemma}
\begin{proof}
These are the $p=2$ case of Gournay's Theorem~1.4
\cite{p09:op15:Gournay}, with his width-dimension convention converted
to the diameter convention used here. Explicitly, his statement says
that below $\sqrt{2(1+1/m)}$ the width-dimension is $m$, and
at that threshold a map to an $(m-1)$-complex exists. It also says
that at every diameter threshold less than $2$ the width-dimension
is strictly larger than $m/2-1$. The upper bound $2$ follows from
the one-set cover. Open versus non-strict fiber thresholds give the
same infimum. No new Euclidean width computation is claimed.
\end{proof}

\begin{proof}[Proof of Theorem~\ref{p09:op15:theorem}]
We recall the width principle in the exact form required: for any
compact metric input $X$ embedded isometrically in the output space,
an $n$-query adaptive continuous procedure of error $r$ satisfies
\begin{equation}\label{p09:op15:adaptive}
                    \operatorname{UW}_{2^n-1}(X)\le2r.
\end{equation}
Indeed, a zero-query leaf has diameter at most $2r$. For a continuous
scalar query, the multiplicative fiber-width inequality bounds
$\operatorname{UW}_{2d+1}(X)$ by the supremum of the
$\operatorname{UW}_d$ of its compact fibers. Induction on remaining
depth gives \eqref{p09:op15:adaptive}. It does not require continuity of
the selected family as the preceding answers vary.

For an $m$-dimensional subspace $V\subset H$, set
$c_V=\inf_{x\in V,\|x\|=1}\|Sx\|$. If $c_V>0$, the image
$S(B_V)$ is a compact ellipsoid. Pulling back covers along $S|_V$
and using $\|Sx-Sy\|\ge c_V\|x-y\|$ shows
\[
 \operatorname{UW}_d(S(B_V))\ge c_V
                                  \operatorname{UW}_d(B_2^m).
\]
Every algorithm on $B_H$ restricts to one on $B_V$ and, through
the continuous inverse $S|_V^{-1}$, to one on this ellipsoid.
Apply \eqref{p09:op15:adaptive}. On Hilbert spaces the min--max identity
\begin{equation}\label{p09:op15:minmax}
 \sup_{\dim V=m}\inf_{x\in V,\|x\|=1}\|Sx\|=a_m(S)
\end{equation}
holds also for noncompact operators. It follows from the spectral
theorem for $S^*S$: a subspace on which its spectrum is bounded
below by $c^2$ supplies the lower estimate, while removing a
finite-dimensional spectral subspace above $c^2$ supplies a rank
approximation of error at most $c$. Taking limits at spectral
endpoints gives \eqref{p09:op15:minmax}. If $c_V=0$, the inequality
being asserted for that subspace is simply zero and needs no inverse.

If $m>\dim H$, the associated approximation number is zero and
the lower bound is trivial; restrict \eqref{p09:op15:minmax} to
$m\le\dim H$ otherwise. Choose $m=2^n$ for the first width in Lemma~\ref{p09:op15:widths},
and $m=2^{n+1}$ for the second. Taking the supremum over $V$
gives both lower bounds in \eqref{p09:op15:bounds}.

For the upper bound put $k=2^{n-1}$. Choose a linear operator $R$
of rank at most $k$ with
$\|S-R\|\le a_{k+1}(S)+\eta$. Its range is a Euclidean space
of dimension at most $k$. The construction in \cite[Theorem~1]{p03:an:KNU}
approximates every vector in that space to any prescribed positive
error using at most $\lceil\log_2 k\rceil+1=n$ continuous
adaptive measurements. Composing each selected measurement with $R$
gives admissible continuous measurements on $B_H$. Recovering $Rx$
within $\eta$ yields error at most $a_{k+1}(S)+2\eta$ for $Sx$.
Let $\eta\downarrow0$. A rank-zero case is immediate.

Finally $a_j(S)\downarrow\|S\|_{\rm ess}$. One inequality follows
because finite-rank operators are compact; the other follows by
approximating a compact operator in norm by finite-rank operators
on Hilbert spaces. The second lower bound and the upper bound in
\eqref{p09:op15:bounds} now converge to the same number, proving
\eqref{p09:op15:essential}. Compact operators form a norm-closed ideal,
so essential norm zero is equivalent to compactness.
\end{proof}

\begin{corollary}[Optimal exponential and double-exponential orders]
\label{p09:op15:rates}
If $a_j(S)\asymp j^{-s}$ for $s>0$, then
\[
                         e_n(S)\asymp2^{-sn}.
\]
If $a_j(S)\asymp\exp(-c j^\beta)$ with $c,\beta>0$, then
\[
                   \log\log\frac1{e_n(S)}
                                  =\beta n\log2+O(1).
\]
Writing $n(\varepsilon)$ for the least query count with error at
most $\varepsilon$, the corresponding complexities are
\[
 n(\varepsilon)=s^{-1}\log_2(1/\varepsilon)+O(1),\qquad
 n(\varepsilon)=\beta^{-1}\log_2\log(1/\varepsilon)+O(1),
\]
respectively, as $\varepsilon\downarrow0$.
\end{corollary}
\begin{proof}
Substitute the two spectral envelopes in \eqref{p09:op15:bounds} and
take logarithms. Replacing an infimal error by an attained procedure
is unnecessary: for the upper complexity choose an index whose
infimal error is strictly smaller than $\varepsilon$, which changes
the displayed query counts by at most a constant.
\end{proof}

For example the diagonal map $(x_j)\mapsto(j^{-s}x_j)$ on
$\ell_2$ has algebraic optimal error $a_{n+1}\asymp n^{-s}$
under nonadaptive continuous information, by the classical
Borsuk--Ulam argument and coordinate projection. Its unrestricted
adaptive error has exactly exponential order. Conversely the identity
on an infinite-dimensional Hilbert space has $e_n=1$ for every finite
$n$: the second lower bound is one, and output zero attains it.

This is one operator-theoretic transfer with several consequences,
not separate conjecture closures for every spectral envelope. It
does not give the exact finite-$n$ error for general ellipsoids or a
corresponding essential-norm theorem for arbitrary Banach operators.
It assumes exact real-valued measurements with unrestricted computational
precision; deterministic noisy information obeys different entropy
bounds \cite{p06:ni:KNPU}.

% ===== mle =====
\section{The transport inverse law and adversarial corruption of individual trials}
\label{p09:ml15:sec}

The geometry of corruption changes when a budget counts individual
Bernoulli flips rather than entirely replaced reports.  The relevant
predictive discrepancy is then transportation distance, and its sharp
inverse curve differs from the previously established total-variation
curve.  We prove that curve and obtain a matching, polynomial-time
minimax theorem for a strong adversary acting on individual trials.
Unlike the preceding whole-report theorem, the algorithm here is given
a valid upper bound on the flip budget.

For a mixing law $G\in\mathcal P([0,1])$, write
\[
 p_G(j)=\int\binom tjx^j(1-x)^{t-j}\,dG(x),\qquad j=0,\ldots,t.
\]
Let $W_{1,t}$ be Wasserstein distance on this count alphabet with
normalized cost $|i-j|/t$. Define
\[
 \omega_t^{\rm tr}(\eta)=
 \sup\{W_1(G,H):W_{1,t}(p_G,p_H)\le\eta\},\qquad
 L_t(\eta)=\log\left(e+\frac1{\eta\sqrt t}\right),
\]
with $L_t(0)=+\infty$.

\begin{theorem}[The sharp transportation inverse modulus]
\label{p09:ml15:modulus}
Uniformly for integers $t\ge1$ and $0\le\eta\le1/4$,
\begin{equation}\label{p09:ml15:modulus-eq}
 \omega_t^{\rm tr}(\eta)\asymp
 \min\left\{1,\frac1t+\frac1{\sqrt{tL_t(\eta)}}+\eta\right\}.
\end{equation}
The constants are absolute.  In particular the logarithm depends on
$\eta\sqrt t$, rather than just on $\eta$.
\end{theorem}

The observation model for the second theorem is as follows.  Draw
$X_1,\ldots,X_N$ independently from $G$, and, conditional on them,
observe $t$ independent Bernoulli trials for each individual.  A strong
adversary sees the entire clean trial array and may flip at most
$q=\lfloor\varepsilon Nt\rfloor$ bits in total.  Only the resulting
counts are passed to the estimator.  Equivalently, the adversary can
change clean counts $J_i$ into reports $J_i'$ subject to
\[
                   \sum_{i=1}^N|J_i'-J_i|\le q.
\]
Every such count change can be implemented by exactly that many bit
flips. Allowing the adversary to see the order of the clean bits adds
only randomization conditional on the counts, which is already allowed.
Denote this class of adversaries by $\mathcal A^{\rm bit}_{N,t}
(\varepsilon)$.

\begin{theorem}[The full bit-corruption minimax curve]
\label{p09:ml15:bitrate}
For $0\le\varepsilon\le1/4$, define
\begin{equation}\label{p09:ml15:riskcurve}
 r^{\rm bit}_{N,t}(\varepsilon)=\min\left\{1,
 \frac1t+\frac1{\sqrt{t\log(eN)}}+\frac1{\sqrt N}
 +\varepsilon+\frac1{\sqrt{tL_t(\varepsilon)}}\right\}.
\end{equation}
Then
\[
 \inf_{\widehat G}\sup_G\sup_{A\in\mathcal A^{\rm bit}_{N,t}
 (\varepsilon)}\mathbb E W_1(\widehat G(A(J)),G)
                 \asymp r^{\rm bit}_{N,t}(\varepsilon).
\]
The upper bound is attained by a rational linear program whose bit
complexity is polynomial in the explicit sample input length.  The
algorithm is supplied with $q$ or a rational upper bound
$\varepsilon$; no adaptation to an unknown bit-flip budget is claimed.
\end{theorem}

Wasserstein contamination, modulus-based lower bounds, and
computationally motivated projection estimators are established topics
\cite{p09:ml15:LiuLoh}. Coordinated Wasserstein perturbations and their
finite-sample minimax analysis were studied by Chao and Dobriban
\cite{p09:ml15:ChaoDobriban} for location, regression, and density estimation.
The new claims are the sharp binomial-mixture
transport modulus and the matching latent-population theorem under a
hard post-sample bit budget. The Hahn operator and interval-discrepancy
linear program are classical or prior ingredients, as specified below.

\subsection{A discrete energy estimate for the hidden oscillations}

Put $\pi(dx)=dx/(\pi\sqrt{x(1-x)})$, $\rho_j=p_\pi(j)$,
$h_k(x)=T_k(2x-1)$, and
\[
 Qh(j)=\mathbb E_\pi[h(X)\mid J=j],\qquad
 \lambda_{t,k}=\prod_{l=0}^{k-1}\frac{t-l}{t+1+l}.
\]
The beta--binomial orthogonal-polynomial diagonalization
\cite{p05:fm:DKS2008} gives
\begin{equation}\label{p09:ml15:qnorm}
 Qh_k\text{ is a degree-}k\text{ polynomial},\qquad
 \sum_j\rho_j(Qh_k(j))^2=\lambda_{t,k}/2,
 \qquad\lambda_{t,k}\le e^{-k^2/(2t)}.
\end{equation}
These identities follow directly from the posterior beta moments:
$Q^*Q$ is self-adjoint, preserves polynomial degree, and has leading
coefficient multiplier $\lambda_{t,k}$ on degree $k$. Arcsine
orthogonality proves the first two identities, and
$\log(1-u)\le-u$ proves the last one.

\begin{lemma}[Transportation attenuation of a Hahn mode]
\label{p09:ml15:attenuation}
For $1\le k\le t$, let $G_k^\pm=(1\pm h_k/2)\pi$. Then
\begin{equation}\label{p09:ml15:attenuation-eq}
 W_1(G_k^+,G_k^-)\asymp k^{-1},\qquad
 W_{1,t}(p_{G_k^+},p_{G_k^-})
                 \le\frac{e^{-k^2/(4t)}}{k}.
\end{equation}
The extra factor $1/k$ in the second inequality is the feature needed
for the transport inverse law.
\end{lemma}
\begin{proof}
Consider the reversible birth-and-death generator
\begin{align*}
 (\mathcal L f)(j)
 ={}&(j+1/2)(t-j)[f(j+1)-f(j)]\\
   &+j(t-j+1/2)[f(j-1)-f(j)],
\end{align*}
with the vanishing endpoint terms omitted. Detailed balance follows
from the beta--binomial ratio $\rho_{j+1}/\rho_j$.
The generator preserves polynomial degree and has leading multiplier
$-k^2$ on degree $k$. Since the polynomials $Qh_k$ are orthogonal under
$\rho$, self-adjointness gives
$\mathcal LQh_k=-k^2Qh_k$. Its Dirichlet form is
\[
 \mathcal E(f,g)=-\langle f,\mathcal Lg\rangle_\rho
   =\sum_{j=0}^{t-1}\rho_j(j+1/2)(t-j)
              [f(j+1)-f(j)][g(j+1)-g(j)].
\]
If $f$ is $1$-Lipschitz for the normalized count metric, then
$\mathcal E(f,f)\le2$, since its consecutive differences are at most
$1/t$ and $(j+1/2)(t-j)\le2t^2$. Thus
\[
 |\langle f,Qh_k\rangle_\rho|
 =k^{-2}|\mathcal E(f,Qh_k)|
 \le k^{-1}\sqrt{\mathcal E(f,f)}\sqrt{\lambda_{t,k}/2}
 \le\frac{\sqrt{\lambda_{t,k}}}{k}.
\]
The predictive signed difference is exactly $\rho_jQh_k(j)$.
Kantorovich duality and \eqref{p09:ml15:qnorm} prove the second assertion.
For the first, use $x=(1+\cos\theta)/2$. The signed population CDF
has magnitude $|\sin(k\theta)|/(\pi k)$, and hence
\[
 W_1(G_k^+,G_k^-)=\frac1{2\pi k}
          \int_0^\pi|\sin(k\theta)|\sin\theta\,d\theta
                         \asymp k^{-1}.
\]
Uniform lower and upper constants follow by integrating on
$[\pi/4,3\pi/4]$, with finitely many small $k$ treated directly.
\end{proof}

\subsection{Polynomial testing and proof of the modulus}

We use the following quantitative approximation ingredient, proved in
\cite{comp:part:07} and used in the preceding computational paper
\cite{comp:part:08}. For each $1$-Lipschitz $f:[0,1]\to\mathbb R$
and integer $d\ge1$ there is a degree-$d$ polynomial $P_d$ with
\[
 \|P_d-f\|_\infty\le C/d,\qquad\|P_d'\|_\infty\le C.
\]
If a polynomial $P$ of degree at most $d\le t/2$ has degree-$t$
Bernstein coefficients $b_j$, then
\begin{equation}\label{p09:ml15:coeffbounds}
 \max_j|b_j|\le Ce^{2d^2/t}\|P\|_\infty,
 \qquad
 t\max_j|b_{j+1}-b_j|\le Ce^{4d^2/t}\|P'\|_\infty.
\end{equation}
The second bound follows by differentiating the Bernstein expansion
and applying the first at degree $t-1$. The first follows from the
Cauchy coefficient contour for $(1+z)^tP(z/(1+z))$, with the
Bernstein complex-polynomial inequality and binomial saddle-point
normalization. We do not re-count these approximation bounds as new.

\begin{proof}[Proof of Theorem~\ref{p09:ml15:modulus}]
Testing the two populations with the polynomial $P_k$ and its
Bernstein coefficients gives
\begin{equation}\label{p09:ml15:transport-upper}
 W_1(G,H)\le C/k+Ce^{4k^2/t}W_{1,t}(p_G,p_H),
                         \qquad1\le k\le t/2.
\end{equation}
There is also the elementary estimate
\begin{equation}\label{p09:ml15:elementary}
 W_1(G,H)\le W_{1,t}(p_G,p_H)+t^{-1/2},
\end{equation}
since a normalized binomial count can be coupled to its parameter
with mean absolute error at most $1/(2\sqrt t)$.
If $\eta\sqrt t\ge1$, equation~\eqref{p09:ml15:elementary} proves the
required upper bound. Otherwise choose
\[
 k=\left\lfloor\min\{t/2,\sqrt{a t L_t(\eta)}\}\right\rfloor,
                                  \qquad a=1/128.
\]
At $\eta=0$ take $k=\lfloor t/2\rfloor$. For sufficiently large
$t$ these are admissible; bounded $t$ are absorbed by the $1/t$ term.
Writing $x=\eta\sqrt t\in(0,1)$, the quantity
$x(e+1/x)^{4a}\sqrt{\log(e+1/x)}$ is uniformly bounded.
Thus the second term of \eqref{p09:ml15:transport-upper} is
$O((tL_t(\eta))^{-1/2})$, while $1/k$ is bounded by that term plus
$C/t$. This proves the upper bound, with clipping at loss one.

For the lower bound, endpoint mixing laws differing by mass $\eta$
have both population and predictive transport distance $\eta$.
The pair $G_{t+1}^\pm$ has identical predictive laws by moment
orthogonality, and population distance $\asymp1/t$.
If $0<\eta\sqrt t<1$, put
$k=\lceil\sqrt{8tL_t(\eta)}\rceil$. When $k\le t$,
Lemma~\ref{p09:ml15:attenuation} gives
\[
 W_{1,t}(p_{G_k^+},p_{G_k^-})
 \le\frac{e^{-2L_t(\eta)}}{\sqrt{8tL_t(\eta)}}\le\eta,
\]
and the latent separation is $\asymp(tL_t(\eta))^{-1/2}$.
If $k>t$, this inverse term is $O(1/t)$ and exact moment matching
suffices. If $\eta\sqrt t\ge1$, it is instead $O(\eta)$.
Combining these lower bounds proves \eqref{p09:ml15:modulus-eq}.
\end{proof}

\subsection{A decontamination linear program}

For a signed vector $z$ on $\{0,\ldots,t\}$ of total mass zero,
let
\[
 D_s(z)=\max_{A\in\mathcal I_s}\left|\sum_{j\in A}z_j\right|,
\]
where $\mathcal I_s$ consists of unions of at most $s$ consecutive
integer intervals. We will need a property that holds even when $z$
is not a difference of model laws.

\begin{lemma}[Polynomial tests from interval discrepancies]
\label{p09:ml15:layercake}
If $b(j)$ is a polynomial of degree at most $k$ on the count grid,
and $s\ge k+1$, then
\begin{equation}\label{p09:ml15:layercake-eq}
             \left|\sum_jb(j)z_j\right|
                    \le2\max_j|b(j)|D_s(z).
\end{equation}
\end{lemma}
\begin{proof}
Subtract the minimum of $b$ and use the layer-cake formula. Every
superlevel set $\{j:b(j)>u\}$ is a union of at most $k+1$ intervals,
because a degree-$k$ nonconstant polynomial has at most $k$ real roots
at each level. The integral has length at most $2\max|b|$.
A constant polynomial gives zero because $z$ has total mass zero.
\end{proof}

Suppose $2\le t\le N$. Put
\[
 M=Nt,\qquad m=\min\{t,\lceil\sqrt{4t\lceil\log_2(8tN)\rceil}
                  \rceil\},\quad s=m+1,
\]
and choose the rational threshold
\[
 b_N=\frac{\lceil\sqrt{N(2s+1)
                    \lceil\log_2(2(t+2))\rceil}\rceil}{N}.
\]
Let $q'$ be the empirical law of the reported counts.  Jointly minimize
\begin{equation}\label{p09:ml15:lp}
 W_{1,t}(p_H,\widetilde q)
       +(D_s(p_H-\widetilde q)-b_N)_+
\end{equation}
over all probability laws $H$ supported on
$\{l/M:0\le l\le M\}$ and all probability vectors
$\widetilde q$ on the count alphabet, subject to
\begin{equation}\label{p09:ml15:feasible}
                      W_{1,t}(\widetilde q,q')\le\varepsilon.
\end{equation}
Return the minimizing $\widehat G=H$.  Here $\widetilde q$ is a fitted
clean empirical law, not a binomial-mixture law.

This is a polynomial-size rational linear program. Wasserstein
constraints and objectives on an ordered finite grid are sums of
absolute cumulative discrepancies. The epigraph of $D_s$ is a
longest-path dynamic program over states recording the position, the
number of selected intervals, and whether the last point was selected.
Introducing one potential for each state, and one inequality per edge,
expresses its epigraph with $O(st)$ variables and constraints, as proved
in \cite{comp:part:08}. The additional variable $\widetilde q$ and
constraint~\eqref{p09:ml15:feasible} add only $O(t)$ variables and
constraints. The support grid has $Nt+1$ points, with binomial
probabilities of bit length $O(t\log(Nt+1))$. Standard rational linear
programming therefore gives polynomial bit complexity when $t\le N$.
If $t>N$, or if $\varepsilon\sqrt t\ge1$, the estimator instead
returns the empirical law of the normalized reports; this branch is
also polynomial time. Bounded exceptional values are harmless.

\begin{proof}[Upper bound in Theorem~\ref{p09:ml15:bitrate}]
Let $q^0$ be the unobserved clean empirical law. The hard flip budget
implies $W_{1,t}(q^0,q')\le\varepsilon$ deterministically. Put
\[
 S=D_s(q^0-p_G),\qquad Z=W_{1,t}(q^0,p_G),
 \qquad A_N=2/N+Z+(S-b_N)_+.
\]
The number of interval unions is at most $(t+2)^{2s}$. Hoeffding's
inequality and the rational choice of $b_N$ give
\[
 \mathbb E(S-b_N)_+\le C/\sqrt N,\quad
 \mathbb ES\le b_N+C/\sqrt N,\quad
 \mathbb EZ\le1/(2\sqrt N),
 \quad b_N\le\frac{Ct^{1/4}\log(eN)^{3/4}}{\sqrt N}.
\]
Round $G$ to the support grid, obtaining $G_*$. Common-uniform
Bernoulli coupling gives
$\|p_{G_*}-p_G\|_{\rm TV}\le t/M=1/N$. The feasible comparator
$(G_*,q^0)$ has objective at most $A_N$. Thus the optimizer satisfies
\begin{align}
 W_{1,t}(p_{\widehat G},\widetilde q)&\le A_N,
                 &D_s(p_{\widehat G}-\widetilde q)&\le b_N+A_N,
                         \label{p09:ml15:fitbounds}\\
 W_{1,t}(\widetilde q,q^0)&\le2\varepsilon,
                 &\mathbb E A_N&\le C/\sqrt N.
                         \label{p09:ml15:transport-middle}
\end{align}
All these comparisons hold uniformly over the post-sample adversary.

To invert these bounds, split a Lipschitz test approximation into
$P_u+(P_k-P_u)$, with $u\asymp\sqrt t$ and $u\le k\le\min(m,t/2)$.
The low-degree Bernstein coefficient vector has normalized-grid
Lipschitz constant at most an absolute constant. The residual
coefficient vector has supremum norm at most
$Ce^{2k^2/t}/u$ and normalized-grid Lipschitz constant at most
$Ce^{4k^2/t}$, by \eqref{p09:ml15:coeffbounds}. Each coefficient vector
is a polynomial in the count index of degree at most $k$, as follows
from the falling-factorial representation of monomials.

Decompose the predictive signed difference as
\[
 p_{\widehat G}-p_G=(p_{\widehat G}-\widetilde q)
                   +(\widetilde q-q^0)+(q^0-p_G).
\]
Test the low-degree vector using transportation distance throughout.
For the residual vector use Lemma~\ref{p09:ml15:layercake} on the first
and third summands, and its Lipschitz bound on the middle summand.
After taking expectations and the Lipschitz dual supremum, this yields
\begin{equation}\label{p09:ml15:hybrid-inverse}
 \mathbb EW_1(\widehat G,G)\le\frac Ck+
 C(\varepsilon+N^{-1/2})
 +\frac{Ce^{2k^2/t}}{\sqrt t}(b_N+N^{-1/2})
 +Ce^{4k^2/t}\varepsilon.
\end{equation}
This hybrid estimate is why ordinary transport projection alone is
not substituted for \eqref{p09:ml15:lp}.

Take $u=\lfloor\sqrt t/16\rfloor$ and
\[
 k=\left\lfloor\min\{t/2,
    \sqrt{a t\min(\log(eN),L_t(\varepsilon))}\}\right\rfloor,
                                   \qquad a=1/128.
\]
For large enough $t$ these choices have $1\le u\le k\le m$;
bounded $t$ are absorbed by the $1/t$ term. The sampling contribution
in \eqref{p09:ml15:hybrid-inverse} is at most
$C(t\log(eN))^{-1/2}$, since $t\le N$ and
$(eN)^{2a}N^{-1/4}\log(eN)^{5/4}$ is bounded. The last term is
$O((tL_t(\varepsilon))^{-1/2})$ by the same elementary logarithmic
inequality as in the modulus proof. The approximation term has the
required two inverse scales plus $1/t$.

On either empirical-output branch, the risk is at most
$\varepsilon+(2\sqrt t)^{-1}+(2\sqrt N)^{-1}$, by the deterministic
transport budget and binomial coupling. If $t>N$, the second term is
absorbed by $N^{-1/2}$; if $\varepsilon\sqrt t\ge1$, it is absorbed
by $\varepsilon$. This proves the claimed upper bound everywhere.
\end{proof}

\begin{proof}[Lower bound in Theorem~\ref{p09:ml15:bitrate}]
No flips are permitted as an adversarial choice, so the established
clean lower bound supplies
$t^{-1}+(t\log(eN))^{-1/2}+N^{-1/2}$, clipped at one.
Its three elementary constructions are respectively exact moment
matching, the Chebyshev perturbations with
$k\asymp\sqrt{t\log(eN)}$, and endpoint mixing weights differing by
$\asymp N^{-1/2}$; see \cite{comp:part:07}.

Let $q=\lfloor\varepsilon Nt\rfloor\ge1$ and set
$\eta=q/(16Nt)$. Choose a pair from the lower-bound constructions in
Theorem~\ref{p09:ml15:modulus}, with predictive transportation distance
at most $\eta$ and latent separation at least a fixed positive
multiple of the modulus curve (using whichever term is largest). Independently couple each pair of counts optimally in
transport distance. The expected total integer displacement is at most
$Nt\eta=q/16$. By Markov's inequality the total displacement exceeds
$q$ with probability at most $1/16$. In the first experiment the
adversary samples the coupled second count vector conditional on its
observed vector. If its total displacement is at most $q$, it applies
these changes; otherwise it leaves the sample untouched. This obeys
the hard budget for every realization. Its output distribution is
within TV $1/16$ of the clean second experiment. Le Cam's inequality
therefore gives a fixed positive multiple of the latent separation as
a lower bound on maximum expected loss.

For $q\ge1$, $q/(Nt)\ge\varepsilon/2$, so scaling the modulus
argument by the absolute factor $1/32$ changes its curve only by
absolute constants. If $q=0$, then
$\varepsilon<1/(Nt)$, and both contamination terms are bounded by
the clean terms because $L_t(\varepsilon)\ge c\log(eN)$ for an absolute $c>0$.
This proves the complete lower bound.
\end{proof}

\subsection{An unbounded separation between two corruption geometries}

For comparison, corruption of an $\varepsilon$ fraction of whole
counts has the prior minimax curve
\[
 \min\{1,t^{-1}+(t\log(eN))^{-1/2}+N^{-1/2}
               +\varepsilon+(t\log(e/\varepsilon))^{-1/2}\}.
\]
Set $N=t^3$ and $\varepsilon=(t\log t)^{-1/2}$, along integer $t$
tending to infinity. The whole-report rate is
$\asymp(t\log t)^{-1/2}$, whereas the bit-budget rate is
$\asymp(t\log\log t)^{-1/2}$. Their ratio is therefore
\[
                       \asymp\sqrt{\frac{\log t}{\log\log t}}
                         \longrightarrow\infty.
\]
This comparison uses the same numerical fraction $\varepsilon$ but
different resource accounting: the second adversary may distribute its
flips across many individuals. It is not a claim that equal sets of
allowable corrupted data have different risks.

\paragraph{Scope and unresolved targets.}
The transport inverse theorem and the sharp bit-budget theorem are two
related new targets, not two previously named conjectures. Their proof
uses earlier clean-binomial lower bounds and Bernstein approximation
estimates, which remain explicitly cited dependencies. The stronger
Chebyshev--Bernstein exponent $\log2$, robustness of the unmodified
NPMLE, and simultaneous optimal adaptation to an unknown bit budget
remain unresolved here. The unknown-budget adaptation previously
proved for whole-report replacements cannot simply be imported: that
proof controls a total-variation contamination error, whereas cheap
individual flips may produce a much larger total-variation error.

% ===== arithmetic =====
\section{Relations, edge loss, and the dimension visible in a zero intensity}
\label{p09:ar15:sec}

The preceding blind inverse theorem requires independent relative
frequencies.  Here we determine what happens when that condition fails,
including the exact leading loss of the rightmost zero boundary.  A
second result computes a sharp edge-intensity constant in the independent
finite-frequency case.  It shows that a first intensity can reveal the
number of hidden terms even though identification of all their magnitudes
from that intensity remains open in this work.

We use the finite exponential sum
\begin{equation}\label{p09:ar15:eq:model}
 F_\theta(s)=1-\sum_{j=1}^{N}a_j e^{-\lambda_j s+i\theta_j},
 \qquad 0<\lambda_1<\cdots<\lambda_N,\quad a_j>0.
\end{equation}
There is a unique real $S$ such that
\begin{equation}\label{p09:ar15:eq:normalization}
 \sum_j a_j e^{-\lambda_jS}=1.
 \qquad b_j=a_j e^{-\lambda_jS},\quad
 B=\operatorname{diag}(b_1,\ldots,b_N),\quad
 M=\sum_j b_j\lambda_j>0.
\end{equation}
The point $S$ is the frozen right boundary, and $F_0'(S)=M$.
Write $s_+(\theta)=\sup\{\Re\rho:F_\theta(\rho)=0\}$.
All zeros lie in a fixed vertical strip, uniformly in $\theta$:
the constant dominates far to the right, and the term of largest
frequency dominates far to the left.

\subsection{An exact criterion and a sharp quadratic loss caused by relations}

Let
\[
 \mathcal L=\{k\in\mathbb Z^N:k\cdot\lambda=0\},\qquad
 H=\overline{\{t\lambda\bmod2\pi:t\in\mathbb R\}}.
\]
The standard character description of a compact torus gives
\begin{equation}\label{p09:ar15:eq:subtorus}
 H=\{\eta\bmod2\pi:k\cdot\eta\in2\pi\mathbb Z
                                    \text{ for all }k\in\mathcal L\}.
\end{equation}
It is a connected subtorus.  Its Lie algebra is
$V=\{v\in\mathbb R^N:k\cdot v=0\text{ for all }k\in\mathcal L\}$,
and $\lambda\in V$.  Choose linearly independent rows of an integer
matrix $K$ spanning the real relation space; put $q=\operatorname{rank}K$.
Locally near $0$, the invariant transverse phase coordinate is
$\xi=K\theta\in\mathbb R^q$, with no reduction modulo $2\pi$ needed.

\begin{theorem}[The relation matrix determines the leading edge loss]
\label{p09:ar15:thm:relation}
For every phase vector,
\begin{equation}\label{p09:ar15:eq:edge-criterion}
 s_+(\theta)=S\quad\Longleftrightarrow\quad\theta\in H.
\end{equation}
Assume $q\ge1$.  The matrix
\begin{equation}\label{p09:ar15:eq:relationmatrix}
 \mathcal A=K(B^{-1}-\mathbf1\mathbf1^{\mathsf T})K^{\mathsf T}
\end{equation}
is positive definite.  For $\theta$ sufficiently close to $H$, use
the local lift of its transverse coordinate near $0$.  Then
\begin{equation}\label{p09:ar15:eq:quadraticloss}
 \boxed{\displaystyle
 S-s_+(\theta)=\frac{1}{2M}\,
       \xi^{\mathsf T}\mathcal A^{-1}\xi+O(\|\xi\|^4).}
\end{equation}
The left side is a real-analytic, even function of $\xi$ in that
neighborhood.  In particular every sufficiently small nonzero
relation-phase defect creates a strictly positive zero-free gap.
If there are no relations, then $s_+(\theta)=S$ for all phases.
\end{theorem}

This is a statement about the characteristic roots themselves.  For
a delay-difference characteristic function of the form
\eqref{p09:ar15:eq:model}, with $\sum a_j=1$ and hence $S=0$, it gives
the exact leading leftward displacement of the spectral boundary.
It does not invoke an unstated equivalence between a root condition
and stability in a particular history space.

\begin{proof}
At a zero $s=S-\delta+it$, set
$\eta_j=\theta_j-t\lambda_j\bmod2\pi$.  The zero equation implies
\begin{equation}\label{p09:ar15:eq:phase-defect}
 \sum_j b_j e^{\lambda_j\delta}|e^{i\eta_j}-1|^2
       =2\left(\sum_jb_j e^{\lambda_j\delta}-1\right),
 \qquad\delta\ge0.
\end{equation}
Thus a sequence of zeros with $\delta\to0$ has
$\eta\to0$ in the torus.  Since $\eta\in\theta+H$, this implies
$\theta\in H$.  Conversely, if $\theta\in H$, the translation
hull contains $F_0$, whose zero at $S$ is simple.  Locally uniform
approximation and Rouch\'e's theorem put zeros of $F_\theta$
arbitrarily close to $S$.  This proves
\eqref{p09:ar15:eq:edge-criterion}.  It also proves strict loss away from
$H$: compactness of the coset in the left side of
\eqref{p09:ar15:eq:phase-defect} prevents a sequence with loss tending
to zero.

We now calculate the loss, rather than merely its sign.  Give
$\mathbb R^N$ the inner product
$\langle u,v\rangle_B=u^{\mathsf T}Bv$.  Let $P$ be the orthogonal
projection onto $V$, and choose the unique small representative
$z\in V^{\perp_B}$ of the phase coset.  Put
\[
 p=P\mathbf1,\quad c=\|p\|_B^2>0,\qquad
 \eta_0(z)=z-\frac{\langle\mathbf1,z\rangle_B}{c}\,p,
 \qquad W=V\cap\{w:\langle\mathbf1,w\rangle_B=0\}.
\]
Here $c>0$ because
$\langle\mathbf1,\lambda\rangle_B=M>0$ and $\lambda\in V$.
The vector $\eta_0(z)$ has weighted mean zero, lies in $z+V$,
and is $B$-orthogonal to $W$.  Therefore
\begin{equation}\label{p09:ar15:eq:Q}
 Q(z):=\|\eta_0(z)\|_B^2
 =\|z\|_B^2+
     \frac{\langle\mathbf1,z\rangle_B^2}{c}.
\end{equation}
Every small phase in the coset $z+H$ has a unique local expression
\[
                 \eta=\eta_0(z)+w+u\lambda,
                 \qquad w\in W,\quad u\in\mathbb R.
\]

Apply the real-analytic implicit-function theorem to
\begin{equation}\label{p09:ar15:eq:IFT}
 1-\sum_j b_j
     \exp\{\lambda_j\delta+i[\eta_0(z)+w+u\lambda]_j\}=0.
\end{equation}
At the origin, derivatives in $(\delta,u)$ are $-M$ and $-iM$,
so the two real equations have a unique analytic solution
$\delta=\delta(z,w)$, $u=u(z,w)$.  Complex conjugation shows that
$\delta(-z,-w)=\delta(z,w)$ and $u(-z,-w)=-u(z,w)$.
Weighted mean zero removes the first-order terms in $\delta$ and
$u$.  Expanding the real equation to second order gives
\begin{equation}\label{p09:ar15:eq:localmin}
 \delta(z,w)=\frac{Q(z)+\|w\|_B^2}{2M}
                     +O((\|z\|+\|w\|)^4).
\end{equation}
In particular its Hessian in $w$ is positive definite.  There is a
unique local minimizing point $w_*(z)$, analytic by another
implicit-function theorem, and $w_*(z)=O(\|z\|^3)$.
It follows that
\begin{equation}\label{p09:ar15:eq:minloss}
              \min_w\delta(z,w)=Q(z)/(2M)+O(\|z\|^4).
\end{equation}
This local minimum is the actual edge loss.  To check this point,
\eqref{p09:ar15:eq:phase-defect} forces every zero in a sufficiently
thin edge strip to have a small phase representative.  Those phases
are exactly parametrized by \eqref{p09:ar15:eq:IFT}.  Conversely every
zero produced there is simple, since its $s$ derivative is close
to $M$, and belongs to a member of the phase hull $z+H$.
Minimality of translation on this hull and Rouch\'e's theorem show
that its real part is a limit of real parts of actual zeros.  Since
\eqref{p09:ar15:eq:minloss} tends to zero with $z$, zeros outside the
thin strip cannot maximize the real part.  This proves the assertion.

It remains to write $Q$ in invariant coordinates.  Set
$G=KB^{-1}K^{\mathsf T}$ and $h=K\mathbf1$.  Then
\[
 z=B^{-1}K^{\mathsf T}G^{-1}\xi,
 \quad \|z\|_B^2=\xi^{\mathsf T}G^{-1}\xi,
 \quad\langle\mathbf1,z\rangle_B=h^{\mathsf T}G^{-1}\xi,
 \quad c=1-h^{\mathsf T}G^{-1}h.
\]
The Sherman--Morrison identity gives
$Q(z)=\xi^{\mathsf T}(G-hh^{\mathsf T})^{-1}\xi$.
For a direct check of positivity, Cauchy--Schwarz and $\sum b_j=1$
give $B^{-1}-\mathbf1\mathbf1^{\mathsf T}\succeq0$, with kernel
spanned by $b=(b_j)$.  If $K^{\mathsf T}v$ were a nonzero multiple
of $b$, taking its dot product with $\lambda$ would contradict
$K\lambda=0$ and $b\cdot\lambda=M>0$.  Thus the restriction
\eqref{p09:ar15:eq:relationmatrix} is positive definite.  Since the
coordinates $z$ and $\xi$ are linearly equivalent, substitution in
\eqref{p09:ar15:eq:minloss} proves \eqref{p09:ar15:eq:quadraticloss}.
\end{proof}

For example, take $N=2$, $(\lambda_1,\lambda_2)=(1,2)$,
$b_1=u$, $b_2=1-u$ and $K=(2,-1)$.  Then $M=2-u$, and
\begin{equation}\label{p09:ar15:eq:two-delay}
 S-s_+(\theta)=
 \frac{u(1-u)}{2(2-u)^3}(2\theta_1-\theta_2)^2
       +O((2\theta_1-\theta_2)^4).
\end{equation}
This agrees with expansion of the explicit quadratic equation in
$e^{-s}$, but the matrix theorem applies to arbitrary rank and
arbitrary finitely many integer relations.

\subsection{The exact first-intensity law at an independent frozen edge}

Now suppose $N\ge2$ and $\lambda_1,\ldots,\lambda_N$ are
linearly independent over $\mathbb Q$.  For fixed $\theta$, let
$n_\theta(\varepsilon,T)$ count zeros, with multiplicity, in
$S-\varepsilon<\Re s<S$, $0<\Im s<T$.
Let $V_d=\pi^{d/2}/\Gamma(1+d/2)$ denote the volume of the
Euclidean unit ball in dimension $d$.

\begin{theorem}[Edge intensity with its leading constant]
\label{p09:ar15:thm:intensity}
For all sufficiently small $\varepsilon>0$ and every $\theta$,
\begin{equation}\label{p09:ar15:eq:intensitylimit}
 \lim_{T\to\infty}\frac{n_\theta(\varepsilon,T)}T
       =c(\varepsilon),
\end{equation}
where
\begin{equation}\label{p09:ar15:eq:edgeconstant}
 \boxed{\displaystyle
 c(\varepsilon)=
 \frac{M V_{N-1}(2M)^{(N-1)/2}}
      {(2\pi)^N\sqrt{\prod_{j=1}^N b_j}}
       \,\varepsilon^{(N-1)/2}(1+O(\varepsilon)).}
\end{equation}
The same $c(\varepsilon)$ is the expected zero count per unit
height in that strip under independent uniform phases.
\end{theorem}

\begin{proof}
We compute the mean count by a local parametrization of every zero
in a thin edge strip.  Identity~\eqref{p09:ar15:eq:phase-defect} puts
all its zero phases in a small neighborhood of $0$.  Write
$v=(\eta_1,\ldots,\eta_{N-1})$.  Solve the real and imaginary
parts of
\[
             1-\sum_j b_j e^{\lambda_j\delta+i\eta_j}=0
\]
for $(\delta,\eta_N)$ as analytic functions of $v$.  The relevant
Jacobian at zero has absolute determinant $M b_N>0$.  Complex
conjugation implies that $\delta(v)$ is even and $\eta_N(v)$ is
odd.  Their leading terms are
\begin{align}
 \eta_N(v)&=-\frac1{b_N}\sum_{j<N}b_jv_j+O(\|v\|^3),
                                                     \label{p09:ar15:eq:etaN}\\
 \delta(v)&=\frac1{2M}v^{\mathsf T}Cv+O(\|v\|^4),
 \quad C=\operatorname{diag}(b_1,\ldots,b_{N-1})
             +\frac1{b_N}(b_i b_j)_{i,j<N}.
                                                     \label{p09:ar15:eq:ellipsoid}
\end{align}
The positive definite matrix satisfies
\begin{equation}\label{p09:ar15:eq:detC}
                  \det C=\frac{\prod_{j=1}^N b_j}{b_N^2}
\end{equation}
by the determinant lemma and $\sum b_j=1$.

Consider the map from local phase coordinates and height to original
phases,
\[
 (v,t)\longmapsto\Theta(v,t)=\eta(v)+t\lambda\pmod{2\pi},
                   \qquad 0<t<1.
\]
Each preimage gives exactly the zero
$s=S-\delta(v)+it$ of $F_{\Theta(v,t)}$.  Conversely every
zero in the thin strip and that height interval gives such a
preimage, since its phase has the unique small representative
provided above.  These zeros are simple: their $s$ derivatives tend
uniformly to $M$ as the strip shrinks.  The ordinary change-of-variables
formula with multiplicities, applied on finitely many coordinate
charts of the torus, therefore gives the expected count as
\begin{equation}\label{p09:ar15:eq:area}
 \frac1{(2\pi)^N}
    \int_{\{v:0<\delta(v)<\varepsilon\}}
                 J(v)\,dv,\qquad
 J(v)=\left|\det[\partial_{v_1}\eta,\ldots,
                    \partial_{v_{N-1}}\eta,\lambda]\right|.
\end{equation}
This formula also follows by partitioning the domain into locally
injective charts; its Jacobian is nonzero for small $v$, so the
usual finite-dimensional area formula counts all preimages correctly.
Equation~\eqref{p09:ar15:eq:etaN} yields
\[
 J(0)=\lambda_N+\sum_{j<N}(b_j/b_N)\lambda_j=M/b_N.
\]
Oddness of $\eta$ makes $J$ even, so
$J(v)=M/b_N+O(\|v\|^2)$.

Since $C$ is positive definite, the sublevel set in
\eqref{p09:ar15:eq:area} has $\|v\|=O(\sqrt\varepsilon)$ and volume
\[
 \frac{V_{N-1}(2M\varepsilon)^{(N-1)/2}}{\sqrt{\det C}}
                 (1+O(\varepsilon)).
\]
For a direct error bound, use polar coordinates after the linear
change $y=C^{1/2}v$ in \eqref{p09:ar15:eq:ellipsoid}; each radial boundary
is $\sqrt{2M\varepsilon}(1+O(\varepsilon))$, uniformly in direction.
Combining this with \eqref{p09:ar15:eq:detC} and the expansion of $J$
proves \eqref{p09:ar15:eq:edgeconstant} for the expected count.

It remains to justify the deterministic limit, rather than replacing
it by the random model.  The finite-dimensional phase flow is uniquely
ergodic by rational independence.  Counts in any fixed compact
rectangle are uniformly bounded over phases: a failure of boundedness,
followed by compactness of the phase torus and locally uniform
convergence, would contradict the finiteness of zeros of a nonzero
analytic limit in a slightly larger compact neighborhood.  No limit
function is identically zero because its constant term dominates at
$+\infty$.

The count in $(S-\varepsilon,S)\times(0,1)$ is continuous in phase
except when a zero meets its boundary.  The above parametrization
shows that the vertical boundary $S-\varepsilon$ has zero mean mass:
the level set $\delta(v)=\varepsilon$ is a smooth hypersurface for
small positive $\varepsilon$.  The boundary at $S$ corresponds to
$v=0$ and has measure zero as $N\ge2$.  Horizontal boundaries also
have zero mean mass by stationarity and local finiteness.  Thus this
bounded counting function is Riemann integrable on the phase torus.
Unique ergodicity, with continuous upper and lower approximations,
identifies its time average on every orbit with its Haar integral.
Integrating unit-height counts over starting heights in $[0,T]$
differs from $n_\theta(\varepsilon,T)$ by $O(1)$, using the same
uniform local bound at the two endpoints.  This proves
\eqref{p09:ar15:eq:intensitylimit}.
\end{proof}

\begin{corollary}[The first intensity identifies the hidden term count]
\label{p09:ar15:cor:dimension}
Within this independent-frequency finite model, the intensity of real
parts near its right support endpoint determines the number $N+1$
of terms:
\[
 N=1+2\lim_{\varepsilon\downarrow0}
                 \frac{\log c(\varepsilon)}{\log\varepsilon}.
\]
After $N$ and $S$ are known, it also determines the combination
$M^{(N+1)/2}/\sqrt{\prod b_j}$.  Neither fact by itself identifies
the individual frequencies or magnitudes.
\end{corollary}

The case $N=1$ is distinct.  Its zeros lie exactly on the line $S$
with intensity $\lambda_1/(2\pi)$, so the first intensity has an
atom at the edge rather than the continuous edge law above.  The
corollary counts the constant term as well as the $N$ exponentials.

\begin{theorem}[The intensity reads rational rank on a compatible hull]
\label{p09:ar15:thm:rank}
Allow integer frequency relations again, and suppose $\theta\in H$.
Put $r=\dim V=\dim_{\mathbb Q}\operatorname{span}_{\mathbb Q}
\{\lambda_1,\ldots,\lambda_N\}$.  Assume $r\ge2$.
Let $\mathfrak v_H$ be the Euclidean covolume in $V$ of the full
lattice $V\cap(2\pi\mathbb Z^N)$, let $b_V$ be the Euclidean
orthogonal projection of $b$ onto $V$, and put
\[
 W=\{w\in V:b\cdot w=0\},\qquad
 B_W=E^{\mathsf T}BE,
\]
where the columns of $E$ are any Euclidean orthonormal basis of $W$.
Then the vertical mean edge count exists for every $\theta\in H$,
and
\begin{equation}\label{p09:ar15:eq:rankconstant}
 \boxed{\displaystyle
 c_H(\varepsilon)=
 \frac{M V_{r-1}(2M)^{(r-1)/2}}
 {\mathfrak v_H\,\|b_V\|\sqrt{\det B_W}}
       \,\varepsilon^{(r-1)/2}(1+O(\varepsilon)).}
\end{equation}
Thus a non-atomic first intensity at the compatible edge determines
the rational rank $r$, rather than the number of terms in a possibly
redundant frequency representation.  For $r=1$, the edge intensity
is instead an atom of mass $1/T_0$, where $T_0$ is the least positive
period of $t\lambda\bmod2\pi$; a sufficiently thin strict strip to
its left has no zeros.
\end{theorem}

\begin{proof}
Since $V$ is a rational subspace, its displayed lattice has full rank
$r$, and Haar measure on $H$ is Euclidean volume divided by
$\mathfrak v_H$.  Also $b_V\ne0$ because $b\cdot\lambda=M>0$.
The flow is uniquely ergodic on $H$; rational independence in its
own character group is exactly the definition of its orbit closure.

For small $w\in W$, solve
\[
 1-\sum_jb_j\exp\{\lambda_j\delta+i(w_j+u\lambda_j)\}=0
\]
for $(\delta,u)$.  The same two-variable implicit-function argument
as in \eqref{p09:ar15:eq:IFT} gives
\[
 \delta(w)=\frac{w^{\mathsf T}Bw}{2M}+O(\|w\|^4),
 \qquad u(-w)=-u(w),\quad\delta(-w)=\delta(w).
\]
All edge-zero phases on $H$ are parametrized this way by the defect
identity.  The map
$(w,t)\mapsto w+u(w)\lambda+t\lambda\bmod2\pi$
has constant Euclidean Jacobian from $W\times\mathbb R$ to $V$,
namely
\[
 J_H=\left|\operatorname{proj}_{W^{\perp}\cap V}\lambda\right|
       =\frac{b_V\cdot\lambda}{\|b_V\|}
       =\frac{M}{\|b_V\|}.
\]
Indeed, adding multiples of the last column $\lambda$ to the other
columns does not change the determinant.  The area formula on $H$
therefore gives the expected edge count in a unit height interval as
$J_H/\mathfrak v_H$ times the Euclidean volume in $W$ of
$0<\delta(w)<\varepsilon$.  Its leading ellipsoid volume is
$V_{r-1}(2M\varepsilon)^{(r-1)/2}/\sqrt{\det B_W}$ with relative
error $O(\varepsilon)$.  Uniform local zero bounds, null boundary
sets, and unique ergodicity give the deterministic vertical mean
exactly as in Theorem~\ref{p09:ar15:thm:intensity}.

If $r=1$, write $\lambda_j=n_j\nu$ with positive integers $n_j$
having greatest common divisor $1$.  Then $T_0=2\pi/\nu$.
On a compatible hull the function is a translate of a polynomial
in $e^{-\nu s}$ with positive magnitudes and frozen phases.  Its
edge zero is simple and is the only zero on that edge modulo $iT_0$,
by equality in the triangle inequality.  Finitely many polynomial
roots give a positive gap to all other real parts.  This proves the
atomic assertion.

For consistency, when $r=N$ one has
$\mathfrak v_H=(2\pi)^N$, $b_V=b$, and
$\det B_W=(\prod b_j)/\|b\|^2$.  The last identity follows by
choosing the unit normal $b/\|b\|$ and using
$(b/\|b\|)^{\mathsf T}B^{-1}(b/\|b\|)=1/\|b\|^2$.
Equation~\eqref{p09:ar15:eq:rankconstant} then reduces to
\eqref{p09:ar15:eq:edgeconstant}.
\end{proof}

\subsection{An exact obstruction to recovery with unrestricted support}

\begin{proposition}[Unknown support can be invisible even in exact data]
\label{p09:ar15:prop:obstruction}
An exact zero divisor on $\Re s>0$, at every height, need not
determine an exponential polynomial when its frequency support is
allowed to change.  This failure occurs with strictly positive
coefficient magnitudes and distinct positive real frequencies.
\end{proposition}
\begin{proof}
For $0<c<1$ set
\[
 f(s)=1-2e^{-s},\qquad
 g_c(s)=(1-2e^{-s})(1-ce^{-s})
       =1-(2+c)e^{-s}+2ce^{-2s}.
\]
All extra zeros of $g_c$ lie on $\Re s=\log c<0$.
Thus the two divisors agree exactly on $\Re s>0$, although the
functions are not constant multiples and have different supports
and magnitudes.  Both have the form \eqref{p09:ar15:eq:model}, the
last coefficient of $g_c$ being represented by phase $\pi$.
The relation $(2,-1)\cdot(1,2)=0$ has invariant phase defect
$2\cdot0-\pi\ne0\bmod2\pi$, explaining why its triangle bound
is not the observable right edge.
\end{proof}

This obstruction has zero statistical distance and zero measurement
error.  Increasing the observation height or reducing ordinate noise
cannot remove it without a support or factor restriction.  It is an
elementary counterexample delimiting the target, not a claimed new
named impossibility theorem.

\subsection{Literature comparison and remaining targets}

The closure of real parts of zeros and its dependence on rational
relations have a long history.  Avellar and Hale
\cite{p09:ar15:AvellarHale} relate the spectral boundary to delay
equations.  Sepulcre \cite{p04:nxz:Sepulcre} gives a pointwise polygon
criterion in the independent case and discusses discontinuity with
respect to the delays.  We do not claim those qualitative principles
as new.  The specific quantitative target here is the explicit
positive matrix and fourth-order remainder in
\eqref{p09:ar15:eq:quadraticloss}.

The averaged logarithm for \eqref{p09:ar15:eq:model} is a restriction of
the Ronkin function of an affine linear polynomial.  Lundqvist
\cite{p09:ar15:Lundqvist} computes second derivatives in three variables
using elliptic and hypergeometric functions.  Our edge calculation is
a local, all-finite-dimensions formula with an explicit constant and
an interpretation as identification of term count.  The proof is
self-contained; historical priority beyond these checked sources is
not asserted.  It concerns a simple frozen edge, whereas the spectral
observation section concerns a controlled double-zero transition.

The original growing-window variance and central-limit questions
remain unresolved here.  So does full magnitude recovery from first
intensity alone when there are at least four terms.  The blind inverse
from the preceding manuscript still concerns exact infinite data on
prescribed independent support; no finite-noise blind algorithm is
claimed by this section.

% ===== gaussian =====
\section{The skew index-one boundary is a completeness-radius problem}
\label{p09:ga15:sec:skew}

The previous manuscript settled compact-scale sampling below stable
index one and left the skew index-one case open. Here the sampling
criterion changes its analytic nature: the observation points travel
along a vertical line, and the exact answer is the completeness radius
of the logarithms of the terminal times. The radius is a classical
object. The result below transfers it to joint location--scale mixing,
including the construction of a single joint signed density when the
radius is finite.

Fix $\beta\ne0$, $|\beta|\le1$, put $k=2\beta/\pi$, and use
the stable index-one kernels
\[
 \widehat K_{a,b}(t)
 =\exp\{ita-b|t|(1+ik\operatorname{sgn}(t)\log|t|)\},
 \qquad a\in\mathbb R,\ b>0.
\]
Conditionally on $(A,B)$, the $X_j$ are independent with this law;
$S_n=X_1+\cdots+X_n$. Only the separate population laws of the
selected sums are supplied. The location and scale may be dependent.

For a locally finite set $\Lambda\subset\mathbb R$, define
\begin{equation}\label{p09:ga15:eq:radius}
 R(\Lambda)=\sup\{a>0:\operatorname{span}
       \{e^{i\lambda x}:\lambda\in\Lambda\}
       \text{ is dense in }L^2(-a,a)\}.
\end{equation}
The supremum of the empty set is zero. Completeness on an interval
implies completeness on smaller intervals, and translation of the
interval does not affect completeness. Multiplication of $\Lambda$
by a nonzero constant preserves whether $R$ is infinite.

For reference, the Beurling--Malliavin theorem gives
\begin{equation}\label{p09:ga15:eq:BM}
 R(\Lambda)=\pi D_{\rm eff}(\Lambda),
\end{equation}
where $D_{\rm eff}$ is the supremum of $a\ge0$ for which there
are disjoint bounded intervals $I_j$ with
\[
 \sum_j\frac{|I_j|^2}{1+\operatorname{dist}(0,I_j)^2}=\infty,
 \qquad \#(\Lambda\cap I_j)\ge a|I_j|\quad\hbox{for every }j.
\]
This is the effective, or exterior, density convention used in
\cite[Section 1.1]{p09:ga15:MP}. Equation \eqref{p09:ga15:eq:BM} and its
density interpretation are classical and are not claimed as new.

\begin{theorem}[Exact compact-scale index-one criterion]
\label{p09:ga15:thm:skew}
Fix $0<b_0<b_1<\infty$ and a terminal set
$\mathcal N\subset\mathbb N$. The selected sums identify every
Borel mixing law on $\mathbb R\times[b_0,b_1]$ if and only if
\begin{equation}\label{p09:ga15:eq:skewcriterion}
 R(\{\log n:n\in\mathcal N\})=\infty,
\end{equation}
equivalently, the effective density of that logarithmic set is infinite.
If the radius is finite, there are distinct absolutely continuous
witnesses with scale support strictly inside $(b_0,b_1)$, equal
coordinate marginals, and equal finite mixed polynomial moments of
every order. Their signed difference is smooth, compactly supported
in scale, and Schwartz in location.
\end{theorem}

The interval width does not enter the criterion. This is because
location frequencies are unrestricted and can rescale a compact
annihilator into any positive scale interval. The theorem does not
assert that a finite-radius system is incomplete on every short
interval at one fixed frequency.

\subsection{From completeness to joint identification}

\begin{lemma}[Finite measures and smooth annihilators]
\label{p09:ga15:lem:radius}
For a locally finite real set $\Lambda$, the following hold.
\begin{enumerate}
\item If $R(\Lambda)=\infty$, every compactly supported finite
complex measure $\nu$ satisfying
$\int e^{i\lambda x}\,d\nu(x)=0$ for all $\lambda\in\Lambda$
is zero.
\item If $R(\Lambda)<\infty$, there is a nonzero complex
$f\in C_c^\infty(\mathbb R)$ whose bilateral Laplace transform
$F$ satisfies $F(1-i\lambda)=0$ for all $\lambda\in\Lambda$.
The support of $f$ may be translated to have any prescribed center.
\end{enumerate}
\end{lemma}
\begin{proof}
For the first assertion, convolve $\nu$ with a smooth compactly
supported approximate identity. Each convolution is a compactly
supported $L^2$ function with the same Fourier zeros. Completeness
on a large enough interval forces it to vanish; letting the smoothing
scale tend to zero gives $\nu=0$. The sign in the Fourier transform
does not matter: a function with $\int g(x)e^{i\lambda x}dx=0$
is orthogonal after conjugation to the given exponential system.

For the second, choose $a>R(\Lambda)$. Noncompleteness supplies
a nonzero $L^2(-a,a)$ function $q$ with the required Fourier zeros,
with conjugation if needed. Smoothing by a nonzero compact bump
preserves the zeros and does not annihilate $q$: the two entire
Fourier transforms have a nonzero product. The resulting smooth
compact function $q_0$ gives $f(x)=e^xq_0(x)$, since
\[
 F(1-i\lambda)=\int e^{i\lambda x}q_0(x)\,dx=0.
\]
Translation multiplies $F$ by a nonvanishing exponential, preserving
all zeros and allowing any support center.
\end{proof}

\begin{proof}[Proof of Theorem~\ref{p09:ga15:thm:skew}]
For $t>0$, the mixing transform of $S_n/n$ is evaluated at
\begin{equation}\label{p09:ga15:eq:skewpoints}
 (t,s_n(t)),\qquad
 s_n(t)=t+ikt\log t-ikt\log n.
\end{equation}
Suppose the radius in \eqref{p09:ga15:eq:skewcriterion} is infinite,
and let $\rho$ be the signed difference of two candidates. For
fixed $t>0$, define a complex measure on $x=tB$ by pushing forward
\[
 e^{itA-tB-iktB\log t}\,d\rho(A,B).
\]
It is supported in $[tb_0,tb_1]$ and its Fourier transform vanishes
at every $k\log n$ by \eqref{p09:ga15:eq:skewpoints}. The first part
of Lemma~\ref{p09:ga15:lem:radius} makes it zero. Undoing the bounded,
nonvanishing weight on this compact interval shows
\[
 \int e^{itA}\,d\rho(A,dB)=0.
\]
Conjugation gives negative $t$ and continuity gives $t=0$. Fourier
uniqueness in $A$, after restriction to arbitrary Borel sets in $B$,
then proves $\rho=0$.

For necessity let $\Lambda=\{k\log n:n\in\mathcal N\}$ and
choose $f,F$ from the second part of the lemma. Write the initial
support of $f$ as an interval of length $2a$. Choose $t_0$ large
enough that $2a/t_0<b_1-b_0$, then translate $f$ to be centered
at $t_0(b_0+b_1)/2$. Let $h\in C_c^\infty((0,\infty))$ be
real and nonzero with support sufficiently close to $t_0$ that
the support of $f(tb)$ lies strictly inside $b_0<b<b_1$ for all
$t\in\operatorname{supp}h$. For $t>0$ set
\begin{equation}\label{p09:ga15:eq:skewG}
 G(t,b)=h(t)\partial_b\{t f(tb)e^{iktb\log t}\},
 \qquad G(-t,b)=\overline{G(t,b)},
\end{equation}
and extend by zero near $t=0$. Define
$g(a,b)=(2\pi)^{-1}\int e^{-ita}G(t,b)\,dt$.
It is real, nonzero, smooth, compactly supported in $b$, and
Schwartz in $a$. Nonzeroness follows because a nonzero compactly
supported function of $b$ cannot have identically zero derivative.
Integration by parts, with no boundary terms, gives for $t>0$
\begin{equation}\label{p09:ga15:eq:skewH}
 \int e^{ita-sb}g(a,b)\,da\,db
       =h(t)sF(s/t-ik\log t).
\end{equation}
At \eqref{p09:ga15:eq:skewpoints}, its last argument is
$1-ik\log n$, a zero of $F$. Negative $t$ follows by conjugation.
The derivative in \eqref{p09:ga15:eq:skewG} makes the location marginal
zero; vanishing of $G$ near $t=0$ makes the scale marginal and all
mixed polynomial moments zero. Normalizing the positive and negative
parts of $g$ yields the asserted witnesses.
\end{proof}

\subsection{An explicit change of sampling threshold}

\begin{corollary}[Stretched-exponential terminal times]
\label{p09:ga15:cor:stretched}
Let $p>0$ and $n_j=\lceil\exp(j^p)\rceil$, discarding any repeated
initial terms. In the compact-scale index-one model above,
\[
 \{\mathcal L(S_{n_j}):j\ge1\}\text{ identifies every mixing law}
 \quad\Longleftrightarrow\quad p<1.
\]
The unrestricted-scale criterion in the preceding manuscript is
$p\le1/2$. Thus $1/2<p<1$ is a genuine interval of designs repaired
by compact scale support, whereas geometric and faster designs are
not repaired.
\end{corollary}
\begin{proof}
The successive differences of $e^{j^p}$ tend to infinity for every
$p>0$, so only finitely many repetitions can occur. Also
\[
 \log n_j=j^p+O(e^{-j^p}).
\]
If $p<1$, the logarithmic set has $\asymp R^{1/p}$ points in
$[0,R]$. A nonzero entire function of finite exponential type has
only $O(R)$ zeros in such a disk, by Jensen's formula centered at
a point where it is nonzero. Hence no nonzero compactly supported
$L^2$ function can have Fourier zeros at all these points, and
$R(\{\log n_j\})=\infty$.

If $p\ge1$, the logarithmic sequence is uniformly separated after
discarding finitely many terms. Any uniformly separated real
sequence has finite effective density. Indeed, if the separation is
$s>0$, an interval $I$ contains at most $1+|I|/s$ points. For a
fixed $a>2/s$, intervals with $\#(\Lambda\cap I)\ge a|I|$
have uniformly bounded length. Apart from finitely many intervals
near zero, assign such intervals to distinct points they contain;
separation implies that the sum of
$(1+\operatorname{dist}(0,I)^2)^{-1}$ is finite. They therefore
cannot form the long family in the definition of $D_{\rm eff}$.
Finite modifications do not change this conclusion. Equation
\eqref{p09:ga15:eq:BM} and Theorem~\ref{p09:ga15:thm:skew} prove failure.
Finally, the unrestricted series is
$\sum_j[1+(\log n_j)^2]^{-1}\asymp\sum_jj^{-2p}$,
which diverges exactly when $p\le1/2$.
\end{proof}

\section{Unknown stable index: every infinite design, no finite design}
\label{p09:ga15:sec:unknown}

The next result allows the stability index itself to vary continuously.
It concerns zero-location symmetric stable kernels, not arbitrary
location or skewness mixing. A logarithmic survival transform converts
the nonlinear stable mixture into a random-coefficient projection
problem, and a nonvanishing gamma factor can then be deconvolved.
The projection uniqueness argument is classical in spirit
\cite{p09:ga15:CFR}; the new assertions concern its exact realization
inside the stable-mixture model and the matching finite-design
counterexamples after inverse-gamma deconvolution.

Fix $0<\alpha_0<\alpha_1\le2$. Draw $(\alpha,B)$ with
$\alpha\in[\alpha_0,\alpha_1]$, $B>0$, and conditionally let
$X_j$ be independent with characteristic function
\[
 \mathbb E(e^{itX_j}\mid\alpha,B)=e^{-B|t|^\alpha}.
\]
Again write $S_n=\sum_{j=1}^nX_j$ and observe only its separate
population laws. ``Zero location'' does not assert existence of a
finite mean for the stable observations.

\begin{theorem}[Sharp finite-versus-infinite threshold for unknown index]
\label{p09:ga15:thm:unknown}
Consider the class of all mixing laws just described for which
\begin{equation}\label{p09:ga15:eq:twosided}
 \mathbb E(B^\eta+B^{-\eta})<\infty
 \quad\hbox{for some }\eta>0.
\end{equation}
The selected population laws $\{\mathcal L(S_n):n\in\mathcal N\}$
identify every joint law of $(\alpha,B)$ in this class if and only
if $\mathcal N$ is infinite.

For every finite terminal set and every integer $L\ge0$, failure
has two distinct absolutely continuous witnesses whose index support
lies in any prescribed nonempty open subinterval of
$(\alpha_0,\alpha_1)$. The witnesses have the same index marginal,
the same conditional moments of $\log B$ through degree $L$, and
satisfy \eqref{p09:ga15:eq:twosided} for every $\eta>0$.
\end{theorem}

The positive theorem includes all compact positive scale supports.
The finite-design witnesses below have unbounded scale support;
no finite-design nonidentifiability claim under a fixed compact
scale restriction is made.

\subsection{A nonlinear mixture becomes a family of linear projections}

Let $E$ be an independent exponential random variable of mean one,
and put
\begin{equation}\label{p09:ga15:eq:QR}
 R=1/\alpha,\qquad Q=-R\log B,
 \qquad C=Q+R\log E.
\end{equation}
Then for every real $x$,
\begin{equation}\label{p09:ga15:eq:survival}
 \mathbb P(C-R\log n>x)
 =\mathbb E e^{-nB e^{\alpha x}}
 =\mathbb E e^{i e^x S_n}.
\end{equation}
Thus the characteristic function of each observed stable sum supplies
the entire survival function of one linear projection of $(C,R)$.
No realization of $E$ is observed or required; it is a mathematical
device for describing this deterministic transform of the population
data.

\begin{proof}[Identification from every infinite design]
Condition \eqref{p09:ga15:eq:twosided} is equivalent, up to a constant,
to an exponential moment of $|\log B|$. Since $R$ is bounded above
and below and $\mathbb E e^{c|\log E|}<\infty$ for $0<c<1$,
there is $\epsilon>0$ with
\[
 \mathbb E e^{\epsilon(|C|+|R|)}<\infty.
\]
For each integer $j\ge0$ the projection moment is the polynomial
\begin{equation}\label{p09:ga15:eq:polynomial}
 \mathbb E(C-R\ell)^j
 =\sum_{q=0}^j\binom jq(-\ell)^q
                     \mathbb E(C^{j-q}R^q).
\end{equation}
The observations give this polynomial at every distinct value
$\ell=\log n$, $n\in\mathcal N$. For an infinite terminal set,
these values determine every coefficient. Hence two candidates with
equal observations have the same mixed moments of $(C,R)$.
The exponential moment makes their joint moment generating functions
analytic near zero, so moment equality identifies the joint law.
This polynomial projection argument is the familiar moment-determinate
form of sparse Cram\'er--Wold uniqueness; it is not a new general
projection theorem.

It remains to undo the noise in \eqref{p09:ga15:eq:QR}. For every real
$t$ and bounded measurable $\varphi$,
\begin{equation}\label{p09:ga15:eq:gamma-deconvolution}
 \mathbb E\{\varphi(R)e^{itC}\}
 =\mathbb E\{\varphi(R)\Gamma(1+itR)e^{itQ}\}.
\end{equation}
The gamma function has no zeros, and on the compact interval of
possible $R$ its reciprocal at $1+itR$ is bounded for fixed $t$.
Replacing $\varphi$ by $\varphi/\Gamma(1+itR)$ recovers
$\mathbb E\{\varphi(R)e^{itQ}\}$. Fourier uniqueness therefore
recovers the joint law of $(Q,R)$, and the invertible map
$\alpha=1/R$, $B=e^{-Q/R}$ recovers the desired mixing law.
\end{proof}

\subsection{Finite projection annihilators survive gamma deconvolution}

\begin{lemma}[An inverse-gamma realization]
\label{p09:ga15:lem:gamma}
Let $J\Subset(0,\infty)$ be an open interval, let
$q\in C_c^\infty(J)$ be real and nonzero, and let
$\ell_1,\ldots,\ell_M$ be finitely many real numbers. For any
integer $K\ge1$ define
\begin{equation}\label{p09:ga15:eq:ghat}
 \widehat g_Q(t,r)=
 \frac{(-it)^K e^{-t^2}}
      {\Gamma(1+itr)}
       \prod_{j=1}^M(\partial_r-it\ell_j)q(r),
 \qquad
 g_Q(v,r)=\frac1{2\pi}\int e^{-itv}\widehat g_Q(t,r)\,dt.
\end{equation}
Then $g_Q$ is a nonzero real smooth integrable function, compactly
supported in $r\in J$. For every fixed $a,b\ge0$ and $H>0$,
\begin{equation}\label{p09:ga15:eq:exptail}
 |\partial_v^a\partial_r^b g_Q(v,r)|\le C_{a,b,H}e^{-H|v|}.
\end{equation}
For every $r$, its $v$ moments of degrees below $K$ vanish.
After adding independent noise $r\log E$ to $v$, its signed joint
density becomes
\begin{equation}\label{p09:ga15:eq:gC}
 g_C(c,r)=\partial_c^K
       \prod_{j=1}^M(\partial_r+\ell_j\partial_c)
           \{\gamma(c)q(r)\},
 \qquad \widehat\gamma(t)=e^{-t^2}.
\end{equation}
Every signed projection $c-r\ell_j$ of $g_C$ is zero.
\end{lemma}
\begin{proof}
The reciprocal gamma function is entire. For a fixed real $y$,
vertical Stirling asymptotics, uniformly for $r$ in a compact positive
interval, give
\[
 |\Gamma(1+i(x+iy)r)^{-1}|
 \le C_y(1+|x|)^{C_y}e^{\pi r|x|/2}.
\]
The same bound with extra polynomial and logarithmic factors holds
after any fixed number of $r$ derivatives. These estimates follow
from the uniform gamma asymptotic in \cite[Equation 5.11.9]{p01:rare:DLMF}
and Cauchy's formula; bounded $x$ is harmless because $1/\Gamma$
is entire. The Gaussian factor has modulus
$e^{-x^2+y^2}$, so all resulting horizontal-line integrals converge
absolutely and uniformly in $r$, together with their derivatives.
Shifting the $t$ contour to $\operatorname{Im}t=-H$ for $v>0$
and to $\operatorname{Im}t=H$ for $v<0$ proves
\eqref{p09:ga15:eq:exptail}. The joining vertical segments vanish by
Gaussian decay. Compact support in $r$ follows directly from the
formula, as does smoothness across its support boundary.

Conjugation symmetry of \eqref{p09:ga15:eq:ghat} proves reality.
The function is nonzero: for fixed real $t\ne0$ the nonzero
constant-coefficient differential operator
$\prod_j(\partial_r-it\ell_j)$ cannot annihilate a nonzero
compactly supported smooth $q$, since its solutions are real analytic.
The reciprocal gamma factor has no zeros for real $t$ and $r>0$.
The factor $t^K$ makes all the stated conditional moments vanish.

Conditionally on $r$, the Fourier transform of $r\log E$ is
$\Gamma(1+itr)$. Multiplication cancels the denominator in
\eqref{p09:ga15:eq:ghat}, yielding exactly \eqref{p09:ga15:eq:gC}; all
convolutions are justified by the exponential decay. Finally the
kernel $e^{it(c-r\ell_j)}$ is annihilated by
$\partial_r+\ell_j\partial_c$. Integration by parts in
\eqref{p09:ga15:eq:gC} therefore makes each of its projection
characteristic functions zero, proving the last assertion.
\end{proof}

\begin{proof}[Finite-design failure in
Theorem~\ref{p09:ga15:thm:unknown}]
For a finite terminal set take $\ell_j=\log n_j$. Choose
$J\Subset(1/\alpha_1,1/\alpha_0)$ so that its reciprocal lies in
the prescribed index subinterval, and apply
Lemma~\ref{p09:ga15:lem:gamma} with $K>L$.
The signed density $g_Q$ has total mass zero, so its nonzero
positive and negative parts have the same positive mass. Normalize
them and push forward by
\[
 (v,r)\longmapsto(\alpha,B)=(1/r,e^{-v/r}).
\]
This is a smooth invertible map with nonzero Jacobian. It gives
distinct absolutely continuous mixing probabilities with the required
index support. Equation \eqref{p09:ga15:eq:exptail}, with arbitrary $H$,
implies $\mathbb E(B^\eta+B^{-\eta})<\infty$ for every $\eta>0$.
The conditional zero moments in $v$ show that the index marginals
and the conditional moments of $\log B=-v/r$ through degree $L$
agree, interpreted as equal finite signed moment measures in $r$.

After the independent auxiliary noise is added, the difference of
the laws of $(C,R)$ is the normalized signed density $g_C$ in
\eqref{p09:ga15:eq:gC}. Every selected projection $C-R\log n_j$ has
zero signed law. The survival identity \eqref{p09:ga15:eq:survival}
therefore gives equality of the characteristic functions of $S_{n_j}$
at every positive argument. Symmetry gives the negative arguments,
and normalization gives the value at zero. Thus all selected stable
sum laws agree, completing the proof.
\end{proof}

\paragraph{Classical inputs, new scope, and remaining gaps.}
The completeness-radius theorem, sparse moment-determinate projection
uniqueness, and the general differential-annihilator principle are
classical or established methods; see
\cite{p09:ga15:MP,p09:ga15:CFR,p01:mix:Le2026}. We do not count them as newly proved
conjectures. The two new targets here are the exact compact-scale
skew index-one classification, and the finite-versus-infinite
classification with a continuously mixed stable index. They refine
different models: the former allows arbitrary location dependence
but fixes the skewness and index, whereas the latter fixes location
zero and symmetry but mixes the index. Identification with index,
skewness, and location all unknown is not settled. Neither theorem
claims quantitative recovery from noisy finite samples. The finite
counterexamples in the unknown-index model do not have compact scale
support, and that stronger support-constrained question remains open.

% ===== spectral =====
\section{Constructive threshold inversion and the density of transition zeros}
\label{p09:sp15:sec}

Two issues remain after sharp stability of an inverse map: how to
compute the inverse without a nonlinear search, and how many usable
zeros survive near a transition.  We resolve these separately.  An
adaptive scalar minimization converts nonlinear threshold observations
to ordinary moments, after which a positive generalized-eigenvalue
construction recovers the atoms.  A second theorem computes an explicit
zero-density constant at a double-root transition.  The latter is
different from a simple rightmost frozen edge: it observes a shrinking
antipodal phase window at a perturbed double zero, and its exponent is
$(m+1)/2$, rather than the simple-edge exponent.

The eigenvalue construction is classical Gaussian quadrature; see,
for example, Blekherman et al.\ \cite{p09:sp15:quadrature}.  The claim
here is the noisy adaptive reduction from nonlinear minima and its
combination with sharp parameter bounds.  The zero-density proof uses
an area formula, within the general non-Gaussian Rice framework
\cite{p09:sp15:rice}; all local coordinates and constants needed here
are computed explicitly below.

Fix $0<a<b<1$, and write
\begin{equation}\label{p09:sp15:eq:model}
 S_\mu(x)=\int_a^b x^r\,d\mu(r),\qquad
 T_\mu(C)=\min_{x>0}\frac{1+Cx}{S_\mu(x)},
 \qquad \mu\in\mathcal P[a,b].
\end{equation}

\subsection{A constructive inverse from adaptive threshold queries}

\begin{theorem}[Adaptive reduction to a positive matrix problem]
\label{p09:sp15:thm:algorithm}
Fix $m\ge2$.  Suppose $\mu=\sum_{j=1}^m w_j\delta_{r_j}$,
the increasing nodes have separation at least $\Delta$, and
$w_j\ge w>0$.  An oracle returns, at a requested $C$ in the fixed
compact interval $J\subset(0,\infty)$ specified below, a real number
$Y(C)$ satisfying $|Y(C)-T_\mu(C)|\le\varepsilon$.
There are $c,K>0$, depending on $a,b,m$, such that if
\begin{equation}\label{p09:sp15:eq:regime}
                 0<\varepsilon\le c w\Delta^{2m-1},
\end{equation}
the following procedure uses at most
$Km\log(3/\varepsilon)$ oracle calls and returns positive
weights $\widehat w_j$ summing to one and ordered nodes
$\widehat r_j\in[a,b]$ with
\begin{equation}\label{p09:sp15:eq:algorithmrates}
 \max_j|\widehat r_j-r_j|
       \le\frac{K\varepsilon}{w\Delta^{2m-2}},\qquad
 \max_j|\widehat w_j-w_j|
       \le\frac{K\varepsilon}{\Delta^{2m-1}}.
\end{equation}
The procedure consists of scalar interval contractions, formation of
two $m\times m$ Hankel matrices, and a symmetric generalized
eigenvalue problem.  It requires no nonconvex search over atomic
measures and no semidefinite feasibility assumption.
\end{theorem}

This is an exact-arithmetic algorithmic statement with a quantitative
oracle-error guarantee.  It does not claim that unguarded computations
in a fixed machine precision remain accurate when atoms collide.
The separation and mass assumptions provide the explicit positive
matrix margins used in the proof.

\begin{lemma}[Noisy threshold-to-profile conversion]
\label{p09:sp15:lem:profile}
For every $x\in[1,e]$, $S_\mu(x)$ can be estimated to error
$K\varepsilon$ with $K\log(3/\varepsilon)$ threshold queries
in a fixed compact $J$.  The constants are uniform over all
probabilities $\mu\in\mathcal P[a,b]$ and $x\in[1,e]$.
\end{lemma}
\begin{proof}
Strict concavity of $S$ gives a unique contact $x_C$ at each
setting.  With $H=S-xS'$, the contact equation is
\[
 C=\frac{S'(x_C)}{H(x_C)},\qquad
 \frac{dC}{dx}=\frac{S''(x)S(x)}{H(x)^2}<0,
 \qquad \frac a{1-a}\le Cx_C\le\frac b{1-b}.
\]
Choose, for example,
\[
 J=\left[\frac{a}{2e(1-a)},\frac{2b}{1-b}\right].
\]
Every contact setting $C_x=S'(x)/(S(x)-xS'(x))$ for
$x\in[1,e]$ lies in the interior of $J$.  On the common compact
range of contacts for $C\in J$, the functions $S,H,-S''$ have
uniform positive upper and lower bounds.  Thus
$-K\le dx_C/dC\le-c<0$.

The tangent majorization identity gives
\begin{equation}\label{p09:sp15:eq:duality}
 S(x)=\min_{C\in J}g_x(C),\qquad
 g_x(C)=\frac{1+Cx}{T(C)},\qquad
 g_x'(C)=\frac{x-x_C}{T(C)(1+Cx_C)}.
\end{equation}
Consequently $g_x$ decreases up to $C_x$ and increases afterwards,
and
\begin{equation}\label{p09:sp15:eq:growth}
 c|C-C_x|\le |g_x'(C)|\le K|C-C_x|.
\end{equation}
The observed value $(1+Cx)/Y(C)$ has error at most
$\eta=K_0\varepsilon$, since all true thresholds have uniform
positive lower bounds on $J$.

Maintain an interval $[L,U]$ containing $C_x$, of width $q$,
and query its two trisection points $C_1,C_2$.  If the observed
difference $\widetilde g(C_2)-\widetilde g(C_1)$ exceeds
$2\eta$, retain $[L,C_2]$; if it is less than $-2\eta$, retain
$[C_1,U]$.  Unimodality makes both decisions valid.  In the
remaining case the true difference has modulus at most $4\eta$.
If $C_x\le C_1$, integration of \eqref{p09:sp15:eq:growth} gives
$g(C_2)-g(C_1)\ge c q^2/18$; the analogous bound holds when
$C_x\ge C_2$.  Thus, while $q^2>144\eta/c$, the ambiguous
case certifies $C_x\in(C_1,C_2)$, and we retain that middle
interval.  Each step contracts by at least a factor $2/3$.
Stop when $q^2\le144\eta/c$.  One final midpoint query, together
with the upper bound in \eqref{p09:sp15:eq:growth}, estimates
$g(C_x)=S(x)$ to $O(\eta)$.  The number of queries is
$O(\log(3/\varepsilon))$.
\end{proof}

\begin{proof}[Construction and proof of Theorem~\ref{p09:sp15:thm:algorithm}]
Set $\zeta=1/(2m)$ and $t_j=e^{\zeta r_j}$.  For
$0\le k\le2m-1$ the numbers
\[
 c_k=S_\mu(e^{k\zeta})=\sum_jw_jt_j^k
\]
are ordinary power moments.  Set $\widehat c_0=1$ exactly, and
apply Lemma~\ref{p09:sp15:lem:profile} at the other $2m-1$ points
to get $|\widehat c_k-c_k|\le K\varepsilon$.  Form
\begin{equation}\label{p09:sp15:eq:hankel}
 \widehat H_0=(\widehat c_{i+j})_{0\le i,j<m},\qquad
 \widehat H_1=(\widehat c_{i+j+1})_{0\le i,j<m}.
\end{equation}
Let $L=e^{\zeta a}/2$ and $U=2e^{\zeta b}$.  The true matrices
factor as $H_0=V\operatorname{diag}(w_j)V^\top$ and
$H_1=V\operatorname{diag}(w_jt_j)V^\top$, where
$V_{ij}=t_j^i$.  Coefficients of the Lagrange polynomials at the
$t_j$ bound $\|V^{-1}\|\le K\Delta^{-(m-1)}$.  Therefore
\begin{equation}\label{p09:sp15:eq:positivity}
 H_0\succeq c w\Delta^{2m-2}I,\quad
 H_1-LH_0\succeq c w\Delta^{2m-2}I,\quad
 UH_0-H_1\succeq c w\Delta^{2m-2}I.
\end{equation}
The enlarged interval supplies fixed positive endpoint margins.
The perturbation of each matrix has norm at most $K\varepsilon$.
Under \eqref{p09:sp15:eq:regime}, all three inequalities remain strict
for the estimated matrices.  These are directly checkable matrix
certificates for the construction.

On the polynomial space of degree at most $m-1$, put
$\langle p,q\rangle=\widehat L(pq)$, where
$\widehat L(t^k)=\widehat c_k$, and define the self-adjoint
operator $M$ by
\[
                  \langle Mp,q\rangle=\widehat L(tpq).
\]
Its coefficient matrix is $\widehat H_0^{-1}\widehat H_1$;
Cholesky reduction makes it a symmetric eigenvalue problem.  The
matrix certificates put every eigenvalue in $(L,U)$.  For a
polynomial of degree at most $m-2$, one has $Mp=tp$ exactly.
Hence $M^k1=t^k$ for $0\le k<m$, and $1$ is a cyclic vector.
Self-adjointness and cyclicity imply that the $m$ eigenvalues
$\widehat t_j$ are distinct and the spectral weights
$\widehat w_j=|\langle1,v_j\rangle|^2$ are strictly positive,
where $v_j$ are orthonormal eigenvectors in this inner product.
Their sum is $\langle1,1\rangle=1$.

The spectral measure matches all estimated moments through
degree $2m-1$.  For degrees at most $2m-2$, split $k=i+j$ with
$0\le i,j<m$ and use
$\langle M^i1,M^j1\rangle=\widehat c_{i+j}$.
At degree $2m-1$, use
\[
 \langle M M^{m-1}1,M^{m-1}1\rangle
      =\widehat L(t^{2m-1}).
\]
This proves the claim without appealing to a separate truncated
moment existence theorem.

We include the stability step to make the algorithmic claim
self-contained.  For two at most $m$-atomic probabilities whose
moments through degree $2m-1$ differ by $O(\varepsilon)$, any
polynomial of that degree with bounded coefficients has integral
discrepancy $O(\varepsilon)$.  At a true node $t_j$ apply this
to
\[
 Q_j(t)=\prod_{k\ne j}(t-t_k)
                   \prod_{\ell=1}^m(t-\widehat t_\ell).
\]
If no estimated node lies within a quarter of the minimum true
separation, its integral has modulus at least
$c w\Delta^{2m-1}$, contradicting \eqref{p09:sp15:eq:regime}.
Thus the nodes match, and division by the other factors gives
$|\widehat t_j-t_j|\le K\varepsilon/(w_j\Delta^{2m-2})$.
For the weights test
\[
 P_j(t)=
 \frac{\prod_{k\ne j}(t-t_k)(t-\widehat t_k)}
      {\prod_{k\ne j}(t_j-t_k)(t_j-\widehat t_k)}.
\]
Its coefficients are $O(\Delta^{-(2m-2)})$, and near $t_j$
its derivative is $O(\Delta^{-1})$.  The test first yields
$\widehat w_j\asymp w_j$ and then
$|\widehat w_j-w_j|\le K\varepsilon/\Delta^{2m-1}$.
Finally take $\log\widehat t_j/\zeta$ and project each result
onto $[a,b]$.  The logarithm is uniformly Lipschitz on $[L,U]$,
and projection cannot increase the location error.  The estimated
nodes remain distinct in the stated regime.  This proves
\eqref{p09:sp15:eq:algorithmrates} and the query bound.
\end{proof}

The separation exponents in \eqref{p09:sp15:eq:algorithmrates} are
sharp for this observation model by the positive moment-preserving
examples in the preceding paper \cite{comp:part:08}.  The present
advance is constructive attainment through an adaptive threshold
oracle, not a second claim to the already proved exponents.  For
rational moment estimates, the matrix certificates and the
characteristic polynomial $\det(\widehat H_1-t\widehat H_0)$
also permit interval-certified real-root calculations.  No theorem
about optimal bit complexity of the complete zero-observation
experiment is asserted.

\subsection{An explicit density law at the double-root transition}

Fix $m\ge1$ and distinct $r_1,\ldots,r_m\in(a,b)$, positive weights summing
to one, $C>0$, and $\ell>0$.  Consider
\begin{equation}\label{p09:sp15:eq:family}
 f_{A,\phi}(z)=-1+A\sum_{j=1}^m w_j
          e^{\ell r_jz+i\phi_j}+Ce^{\ell z+i\phi_*}.
\end{equation}
Write $T=T_\mu(C)$ and let $x_*$ be its contact.  Define
\begin{equation}\label{p09:sp15:eq:constants}
\begin{aligned}
 s_*&=S_\mu(x_*),\quad a_j=T w_j x_*^{r_j},\quad y=Cx_*,\\
 a_0&=1+y=\sum_j a_j,\quad B=\sum_j a_j r_j(1-r_j),\quad D=\frac{y}{1+y}.
\end{aligned}
\end{equation}
Both $B$ and $D$ are positive, and the contact equation gives
$\sum_j a_jr_j=y$.  Let $\lvert\mathbb S^{m-1}\rvert$ be
the surface area of the unit sphere in $\mathbb R^m$, with
$\lvert\mathbb S^0\rvert=2$.

Fix $\alpha>0$ and $R>0$.  At amplitude $A_h=T+\alpha h$,
call a zero $z=\sigma+it$ marked when
\begin{equation}\label{p09:sp15:eq:mark}
 \operatorname{dist}(\ell t+\phi_*,\pi+2\pi\mathbb Z)
                                                   <R\sqrt h.
\end{equation}
Let $N_h(H;\phi)$ count marked zeros in $0<t<H$, with
multiplicity and without restricting their real parts.

\begin{theorem}[Double-root caustic intensity]
\label{p09:sp15:thm:intensity}
For independent Haar phases $\phi$, there is an intensity
$\rho_h$ such that $\mathbb E N_h(H;\phi)=H\rho_h$ for
every $H>0$.  As $h\downarrow0$,
\begin{equation}\label{p09:sp15:eq:intensity}
 \rho_h=\mathcal K_{\mu,C,\ell}(\alpha,R)
                  h^{(m+1)/2}(1+O(\sqrt h)),
\end{equation}
where the explicitly positive constant is
\begin{align}
 \mathcal K_{\mu,C,\ell}(\alpha,R)
 &=\frac{\ell\lvert\mathbb S^{m-1}\rvert D}
 {(2\pi)^{m+1}m\sqrt{B a_0\prod_{j=1}^m a_j}}
                         \mathcal I_m(\alpha,R),
                                      \label{p09:sp15:eq:constant}\\
 \mathcal I_m(\alpha,R)
 &=\int_{-R}^{R}(2\alpha s_*+Dv^2)^{(m-2)/2}
                  [2\alpha s_*+(m+1)Dv^2],dv.
                                      \label{p09:sp15:eq:integral}
\end{align}
If $\ell r_1,\ldots,\ell r_m,\ell$ are rationally independent,
then for every fixed choice of phases and every sufficiently small
fixed $h$,
\begin{equation}\label{p09:sp15:eq:actualdensity}
                       \lim_{H\to\infty}\frac{N_h(H;\phi)}H
                                      =\rho_h.
\end{equation}
The constants in the error term may depend on the fixed spectrum,
$\alpha$, and $R$.  The theorem is not uniform as a weight tends
to zero, as $\alpha\downarrow0$, or as frequencies become resonant.
\end{theorem}

\begin{proof}
We give localization, the zero-count formula, and the constant
calculation separately.  All zeros of the family with $A$ in a
fixed compact interval lie in a uniform vertical strip: the constant
term dominates sufficiently far left, and the largest exponential
dominates sufficiently far right.  For a marked zero put
$x=e^{\ell\sigma}$ and write its top phase as $\pi+v$, with
$|v|<R\sqrt h$.  The zero equation gives
\[
 A_h S(x)\ge |1+Cx e^{iv}|
                     \ge(1-Kh)(1+Cx).
\]
The quotient $(1+Cx)/S(x)$ has a unique strict minimum at $x_*$,
with positive second derivative there.  Thus every marked zero
satisfies $x-x_*=O(\sqrt h)$.  The total magnitude of the smaller
terms exceeds the magnitude of their resultant by only $O(h)$.
The identity
\[
 \sum_j q_j|e^{i\theta_j}-e^{i\psi}|^2
                 =2\left(\sum_jq_j-\left|\sum_jq_j e^{i\theta_j}
                    \right|\right),
 \qquad e^{i\psi}=\frac{\sum_jq_j e^{i\theta_j}}
                         {|\sum_jq_j e^{i\theta_j}|},
\]
with $q_j=A_hw_jx^{r_j}$, shows that every smaller phase is
$O(\sqrt h)$ modulo $2\pi$.  The resultant is
$1+Cx e^{iv}$ and has argument $O(\sqrt h)$, so all marked zeros
are localized in this single phase chart.  Strict positivity of
the fixed weights is used here.

Use the coordinate $u=\ell\sigma-\log x_*$ and let
\[
 F_h(u,\theta,v)=-1+(1+\alpha h/T)
            \sum_j a_j e^{r_j u+i\theta_j}-y e^{u+iv}.
\]
With $q=(\sqrt{a_j})_{j=1}^m$, write the weighted phases uniquely as
\begin{equation}\label{p09:sp15:eq:phasecoords}
 (\sqrt{a_j}\theta_j)_{j=1}^m
       =D v q+t+\xi\frac q{\sqrt{a_0}},\qquad t\perp q.
\end{equation}
The linear imaginary part of $F_h$ is $\sqrt{a_0}\xi$.
The implicit-function theorem therefore eliminates $\xi$ on
$\Im F_h=0$, with
$\xi=O(h)$ on the localized region.  The measure factor in this
elimination is
\begin{equation}\label{p09:sp15:eq:measurefactor}
 d\theta_1\cdots d\theta_m\,
             \delta_0(\Im F_h)
   =\frac{1+O(\sqrt h)}{\sqrt{a_0\prod_j a_j}}\,dt.
\end{equation}
Here and below $\delta_0$ is shorthand for integration over the
regular level set with the coarea Jacobian, not a probabilistic
independence assertion.

Set $u=\sqrt h\,U$, $t=\sqrt h\,Z$, and
$v=\sqrt h\,V$.  Uniformly with two derivatives on the bounded
localized set, after the imaginary part has been eliminated,
\begin{equation}\label{p09:sp15:eq:quadratic}
 h^{-1}\Re F_h
       =\alpha s_*-\tfrac12 B U^2-\tfrac12|Z|^2
                            +\tfrac12D V^2+O(\sqrt h).
\end{equation}
Let $b\in q^\perp$ have coordinates
$b_j=\sqrt{a_j}(r_j-D)$.  Then
\begin{equation}\label{p09:sp15:eq:derivative}
 h^{-1/2}\partial_u F_h
              =-BU+i(b\cdot Z-DV)+O(\sqrt h),\qquad
                         |b|^2=D-B.
\end{equation}
Here $\partial_uF_h$ is the original partial derivative with the
phases held fixed, evaluated on the imaginary-zero surface; it is
not the total derivative after eliminating $\xi$. The last identity
follows from $\sum a_jr_j=y$ and
$\sum a_jr_j^2=y-B$.  In particular $D\ge B$.

Because $\alpha>0$, the limiting real-zero surface
\[
                  B U^2+|Z|^2=2\alpha s_*+DV^2
\]
has a uniformly nonzero gradient in $(U,Z)$ for $|V|\le R$.
The actual surfaces are uniformly regular perturbations of these
spheres after the linear change $Q=(\sqrt B U,Z)$.  Radial
parameterization and the implicit-function theorem give uniform
$C^1$ error $O(\sqrt h)$ in their surface measures and all the
bounded integrands below.  This verifies the regularity needed
for the coarea calculation, including the case $m=1$ of a
two-point sphere.

We justify the zero-count identity being used.  For a holomorphic
function, the real Jacobian of its value map is the squared
modulus of its complex derivative.  The area formula, followed
by averaging its level-counting function over a small disk about
zero, consequently counts its zeros with multiplicity.  On a
compact zero window the count is uniformly bounded for this
compact phase family, also after subtracting a sufficiently small
complex constant: locally uniform limits cannot be identically
zero, so Rouch\'e's theorem bounds the number of zeros.  This
uniform bound permits averaging over phases before taking the
disk radius to zero.  On the localized region just described the
combined map $(u,\theta,v)\mapsto(\Re F_h,\Im F_h)$ is a
submersion.  Its coarea formula thus gives the limit.  The exceptional
phases for which a localized zero is multiple have Haar measure
zero by Sard's theorem applied to projection of this regular zero
manifold onto phase space.  These observations also justify
indicators by continuous approximation when their boundary has
zero coarea measure.

Phase translation by $(\ell r_1t,\ldots,\ell r_mt,\ell t)$
preserves Haar measure.  Thus the expected marked count per unit
height is
\begin{equation}\label{p09:sp15:eq:coarea}
 \rho_h=\frac{\ell}{(2\pi)^{m+1}}
 \int_{|v|<R\sqrt h}\int
       |\partial_u F_h|^2\,
       \delta_0(\Re F_h)\delta_0(\Im F_h)\,du\,d\theta\,dv.
\end{equation}
The factor $\ell$ is $\ell^2$ from the complex derivative and
$\ell^{-1}$ from $d\sigma=du/\ell$.  Stationarity proves the
exact identity $\mathbb E N_h(H)=H\rho_h$.

Substitute \eqref{p09:sp15:eq:measurefactor}--\eqref{p09:sp15:eq:derivative}
into \eqref{p09:sp15:eq:coarea}.  The power of $h$ is
$h^{(m+1)/2}$: the variables $(u,t,v)$ contribute
$h^{(m+1)/2}$, the real level constraint contributes $h^{-1}$,
and the derivative squared contributes $h$.  The limiting integral
at a fixed $V$ is
\[
 \int_{\mathbb R\times\mathbb R^{m-1}}
 [B^2U^2+(b\cdot Z-DV)^2]
 \delta_0\left(\alpha s_*-\frac{BU^2+|Z|^2-DV^2}{2}\right)
                                                   \,dU\,dZ.
\]
Put $Q=(\sqrt B U,Z)$ and
$R_V^2=2\alpha s_*+DV^2$.  Radial integration and spherical
symmetry evaluate this integral as
\[
 \frac{\lvert\mathbb S^{m-1}\rvert}{\sqrt B}
 R_V^{m-2}
       \left[\frac{(B+|b|^2)R_V^2}{m}+D^2V^2\right]
 =\frac{\lvert\mathbb S^{m-1}\rvert D}{m\sqrt B}
 R_V^{m-2}[2\alpha s_*+(m+1)DV^2].
\]
This gives \eqref{p09:sp15:eq:constant}--\eqref{p09:sp15:eq:integral}.
The uniform surface perturbation already proved gives the stated
relative $O(\sqrt h)$ remainder.  Positivity is immediate.

Finally assume rational independence.  The phase flow is uniquely
ergodic on the full torus, by finite-dimensional Fourier
approximation.  The marked count in a unit-height window is a
bounded function of the initial phases.  Its discontinuities can
occur only at a zero on the horizontal boundary or at the two
angular boundaries in \eqref{p09:sp15:eq:mark}.  Horizontal boundaries
have zero mean mass by stationarity and local finiteness; the
coarea formula above gives zero mean mass to either angular
boundary.  Thus the function has Haar-null discontinuity set.
Continuous upper and lower approximations give its time average
on every orbit.  Integrating the sliding unit-window count differs
from $N_h(H)$ by a uniformly bounded endpoint contribution.
This proves \eqref{p09:sp15:eq:actualdensity}.
\end{proof}

\begin{corollary}[A deletion and insertion scale for marked zeros]
\label{p09:sp15:cor:contamination}
Under the rational-independence hypothesis, consider the reduced
record consisting only of the marked ordinates.  At $A_h=T+\alpha h$
it contains $\rho_hH+o(H)$ entries.  If an adversary deletes at
most $\eta H$ genuine entries and inserts at most $\eta H$
arbitrary entries, the test ``more than $\eta H$ reported marks''
distinguishes this record from an empty genuine marked record for
all sufficiently large $H$ whenever $2\eta<\rho_h$.
Conversely, when $2\eta>\rho_h$, these two reduced records can
be made identical within those budgets for all sufficiently large
$H$.  In particular, the critical density of corruptions is of
order $h^{(m+1)/2}$, with leading constant
$\mathcal K_{\mu,C,\ell}(\alpha,R)/2$ for this reduced-record
comparison.
\end{corollary}
\begin{proof}
The nonempty genuine record retains at least
$\rho_hH-\eta H+o(H)$ marks, whereas the empty genuine record
can acquire at most $\eta H$.  This proves the first assertion.
For the converse, keep half the genuine marks in the nonempty
record and insert those same marks into the empty record.  Each
budget is $\rho_hH/2+o(H)$.
\end{proof}

Empty marked records do occur below the transition: the triangle
inequality used in the proof gives
$A\ge(1-K_Rh)T$ for any marked zero, so amplitudes
$A<T-Kh$ have none for a sufficiently large fixed $K$.
The corollary deliberately concerns the reduced record after
discarding the other zeros.  It is not an impossibility theorem
for an observer allowed to exploit the entire zero divisor.
Nor does \eqref{p09:sp15:eq:actualdensity} provide a uniform finite
height for its asymptotic to become accurate; quantitative
recurrence information is still needed for that stronger request.

The two results close distinct internal targets: constructive
inversion of noisy nonlinear thresholds, and an explicit local
zero-density law with its marked-record corruption scale.
They do not resolve finite-data blind recovery of all phases,
optimal zero-observation height, or a general algorithm robust to
missing and spurious zeros without a density budget.

% ===== bessel =====
\section{The entire intermediate-rate gamma region and its quartic gap}
\label{p09:be15:section}

The preceding paper determined the largest endpoint-rate gamma factor
on the first excess-power McKay line.  It left open what happens when
the rate is moved inside the two original gamma rates.  We prove that
the unrestricted factor boundary is continuous and strictly monotone,
that it lies strictly above the ID-remainder boundary at every
intermediate rate, and that its boundary remainder has finitely many
interior zeros.  A local singularity analysis gives the exact quartic
separation when the central rate is approached.  The same analytic
construction resolves the singular limit at the lowest Bessel order.
These are new internal transfer questions, not additional named
conjectures from the external source \cite{p01:rare:Baricz2026}.

\subsection{An analytic inverse valid at every intermediate rate}

Fix $\kappa>1/2$, $b>1$, and write $M=2\kappa-1$.  We use the
same McKay line as in the preceding paper:
\begin{equation}\label{p09:be15:transform}
 L(s)=\frac{b+s}{b}
       \left(\frac{b^2-1}{(b+s)^2-1}\right)^\kappa.
\end{equation}
For an extracted gamma factor of shape $h$ and rate $\rho$, put
$c=b-\rho$.  The intermediate-rate interval is $0<c<1$.
The two capacities, with an arbitrary nonnegative remainder and an
ID remainder respectively, are denoted $H_{\rm P}(c)$ and
$H_{\rm ID}(c)$.  A positive normalizing constant and a translation
of the Laplace variable reduce the prospective remainder transform to
\begin{equation}\label{p09:be15:F}
 F_{c,h}(w)=w^h\frac{w+c}{((w+c)^2-1)^\kappa}.
\end{equation}
For $0\leq h<M$, write $\eta=M-h>0$ and define coefficients by
\begin{equation}\label{p09:be15:coefficients}
 (1+cz)\{1+2cz+(c^2-1)z^2\}^{-\kappa}
                  =\sum_{n=0}^\infty a_n(c)z^n.
\end{equation}
Its inverse on the positive axis is
\begin{equation}\label{p09:be15:inverse}
 g_{c,h}(x)=\sum_{n=0}^\infty
        \frac{a_n(c)x^{\eta+n-1}}{\Gamma(\eta+n)},\qquad x>0.
\end{equation}
This series and every fixed derivative converge locally uniformly.
Indeed Cauchy's estimate gives a geometric bound on $a_n$, and
the gamma denominator dominates it on every compact $x$ interval.
The actual probability remainder, if it exists, has density
$e^{-\rho x}g_{c,h}(x)/F_{c,h}(\rho)$.

\begin{lemma}[Uniform positivity at both ends]\label{p09:be15:ends}
For $0\leq c<1$ and $0\leq h<M$,
\begin{align}
 g_{c,h}(x)&=\frac{x^{M-h-1}}{\Gamma(M-h)}(1+O(x)),
                              &&x\downarrow0,\label{p09:be15:zero}\\
 g_{c,h}(x)&=\frac{2^{-\kappa}(1-c)^h}{\Gamma(\kappa)}
          e^{(1-c)x}x^{\kappa-1}(1+O(x^{-1})),
                              &&x\to\infty.\label{p09:be15:infinity}
\end{align}
The estimates are uniform when $(\kappa,c,h)$ ranges over a compact
subset of the displayed strict parameter domain.  The large-$x$
estimate also holds uniformly as $h\uparrow M$.  On compact
subintervals of $(0,\infty)$ the inverse extends continuously to
$h=M$, after removing the atom at zero that appears at that value.
Consequently positivity of $g_{c,h}$ on all of $(0,\infty)$ is an
open condition in the parameters whenever it is strict everywhere
and $h<M$.
\end{lemma}
\begin{proof}
The first assertion follows from \eqref{p09:be15:inverse}.  To justify
the second and the Laplace inversion, let $f_c$ be the positive
inverse of $(w+c)/((w+c)^2-1)^\kappa$.  The Bessel integral gives
\[
 f_c(x)=\frac{\sqrt\pi}{2^{\kappa-1/2}\Gamma(\kappa)}
       e^{-cx}x^{\kappa-1/2}I_{\kappa-3/2}(x).
\]
For an integer $N>h$, let $\gamma=N-h>0$.  Then
\[
 g_{c,h}(x)=\frac{d^N}{dx^N}
       \left\{\frac1{\Gamma(\gamma)}
       \int_0^x(x-y)^{\gamma-1}f_c(y)\,dy\right\}.
\]
Near zero the expression in braces has order $x^{M+N-h-1}$;
its derivatives through order $N-1$ vanish at zero when $h<M$.
Taking its Laplace transform therefore gives \eqref{p09:be15:F}.
The series inversion agrees first for large real $w$ and then
by uniqueness everywhere the tilted transform is defined.

The standard differentiated Bessel expansion gives
$f_c(x)=2^{-\kappa}e^{(1-c)x}x^{\kappa-1}/\Gamma(\kappa)$
times a full inverse-power expansion.  Write the fractional
integral with $v=x-y$, split it at $v=x/2$, and expand
$(x-v)^{\kappa-1}$ on the first part.  Integration against
$e^{-(1-c)v}v^{\gamma-1}/\Gamma(\gamma)$ supplies the leading
factor $(1-c)^{-\gamma}$.  The other part is exponentially
smaller, together with every fixed derivative.  Differentiating
$N$ times supplies $(1-c)^N$.  This proves
\eqref{p09:be15:infinity}, with uniform errors on the stated compact
sets, also if $h$ tends to $M$.  At that limiting value the
small-$x$ initial term represents an atom at zero; it does not
affect this calculation on $x>0$.  Positivity is uniform near
the two ends and persists on the intervening compact interval
by continuity, proving the last assertion.
\end{proof}

\subsection{Strict gain at every intermediate rate}

The ID capacity was computed previously by minimizing the L\'evy
kernel.  We recall its elementary expression:
\begin{equation}\label{p09:be15:ID}
 H_{\rm ID}(c)=\min_{x\geq0}e^{-cx}(2\kappa\cosh x-1),
 \qquad 0\leq c\leq1.
\end{equation}
For $0<c\leq1$ the minimum is attained at a unique positive
point and is strictly smaller than $M$.  These assertions follow
by putting $z=e^{-x}$ in the derivative equation:
\[
 \kappa(1+c)z^2-cz-\kappa(1-c)=0.
\]
At $c=1$ its value is $\kappa-1/(4\kappa)$.  At $c=0$ it equals
$M$.  These formulas are background, not newly counted results.

\begin{theorem}[The unrestricted factor region]\label{p09:be15:region}
For every $\kappa>1/2$ and every $0<c<1$,
\begin{equation}\label{p09:be15:strict-region}
       0<H_{\rm ID}(c)<H_{\rm P}(c)<M.
\end{equation}
The supremum defining $H_{\rm P}(c)$ is attained.  The function
$c\mapsto H_{\rm P}(c)$ is continuous and strictly decreasing
on $[0,1]$, with the endpoint values
\[
 H_{\rm P}(0)=M,\qquad
 H_{\rm P}(1)=\kappa-\delta_*(\kappa),
\]
where the endpoint Kummer root $\delta_*$ was determined in the
preceding paper \cite{comp:part:08}.  Every boundary remainder
for $0<c<1$ has a real analytic density on $(0,\infty)$, strictly
positive near both endpoints, with a nonempty finite set of
interior zeros, each of positive even order.
For $H_{\rm ID}(c)<h\leq H_{\rm P}(c)$ the unique remainder
is a probability law which is not ID.
\end{theorem}
\begin{proof}
The difference between the original and extracted gamma L\'evy
densities is
\[
 \frac{e^{-\rho x}}x
       \{e^{-cx}(2\kappa\cosh x-1)-h\}.
\]
At $h=H_{\rm ID}(c)$ it is nonnegative and vanishes at just
one positive point.  Its ID law has density strictly positive
at every $x>0$.  To see strict positivity directly, split the
jumps at $\epsilon>0$.  The small-jump sum has mean tending to
zero with $\epsilon$.  The event of exactly two larger jumps
has a density strictly positive on $(2\epsilon,\infty)$, since
the one-jump density is positive except at one point.  Convolution
with the small-jump law gives a positive continuous lower bound
near any prescribed $x>0$.  Lemma~\ref{p09:be15:ends} now permits
a small increase of $h$, proving the strict lower inequality.

As $w\to\infty$, the candidate transform has expansion
\[
 F_{c,h}(w)=w^{-(M-h)}\{1-Mc/w+O(w^{-2})\}.
\]
For $h>M$ it is unbounded; for $h=M$ it approaches its positive
limit from below and has positive derivative for all sufficiently
large $w$.  Neither behavior is possible for a translated
probability Laplace transform.  The admissible set is closed:
if $h_j$ are admissible and tend to $h$, their normalized
transforms converge pointwise to the normalized expression
\eqref{p09:be15:F}, which is continuous at the origin.  The
continuity theorem for Laplace transforms gives a probability
limit.  Thus a largest shape exists and is strictly less than $M$.

At this largest shape the density cannot be positive everywhere,
by the openness assertion in Lemma~\ref{p09:be15:ends}.  It is
analytic and positive near zero and infinity.  Its zeros therefore
lie in a compact interior interval, are isolated, and are finite
in number.  Nonnegativity forces each multiplicity to be even.

To prove strict monotonicity, take rates
$b-1\leq\rho_1<\rho_2<b$.  For every $h>0$,
\[
 G_{h,\rho_1}\stackrel{\rm d}=G_{h,\rho_2}+Z,
\]
where $Z$ is compound Poisson with L\'evy density
$h(e^{-\rho_1x}-e^{-\rho_2x})/x>0$.  Its one-jump density is
strictly positive on $(0,\infty)$.  If the maximal factor at
$\rho_1$ is extracted, convolution of its remainder with $Z$
makes the resulting remainder strictly positive everywhere at
$\rho_2$: the old remainder has positive mass arbitrarily close
to zero, and the one-jump term suffices.  The same endpoint
stability argument permits a strict increase in shape at $\rho_2$.
The known endpoint theorem supplies the required positive density
near zero also when $\rho_1=b-1$.  At $\rho_2=b$, the maximum
$M$ is strictly larger by the inequality already proved.

Closedness of admissible parameters gives upper semicontinuity.
For lower semicontinuity at an interior rate, fix a shape strictly
below the maximum.  Its remainder is the convolution of a maximal
remainder with a positive gamma density, hence is positive
everywhere.  Lemma~\ref{p09:be15:ends} preserves this shape at nearby
rates.  At $c=1$, monotonicity together with closedness gives
one-sided continuity.  At $c=0$, the squeeze
$H_{\rm ID}(c)\leq H_{\rm P}(c)\leq M$ and
$H_{\rm ID}(c)\to M$ gives continuity.  Finally uniqueness of
the L\'evy representation excludes an ID remainder for every
shape above \eqref{p09:be15:ID}.
\end{proof}

This theorem determines the topology and strict ordering of the
entire rate region.  It does not pretend to give an elementary
formula for the intermediate boundary.  In particular, multiple
interior contacts are not excluded away from the singular regimes
analyzed next.

\subsection{A double-root bifurcation at the central rate}

It is useful to remove the singular gamma prefactor.  Define
\begin{equation}\label{p09:be15:B}
 \mathcal B(\eta,c,x)=\eta+
       \sum_{n\geq1}\frac{a_n(c)x^n}{(1+\eta)_{n-1}}.
\end{equation}
Then $g_{c,h}(x)=x^{\eta-1}\mathcal B(\eta,c,x)/\Gamma(1+\eta)$.
The function $\mathcal B$ is analytic jointly near $\eta=c=x=0$.
At $c=\eta=0$ it is strictly positive for $x>0$, since
$a_{2j}(0)=(\kappa)_j/j!$ and $a_{2j+1}(0)=0$.

\begin{theorem}[The exact quartic cost of imposing ID]
\label{p09:be15:quartic}
For fixed $\kappa>1/2$, as $c\downarrow0$,
\begin{align}
 M-H_{\rm P}(c)
 &=\frac{M^2}{4\kappa}c^2
   -\frac{M^4(\kappa+1)}{192\kappa^3}c^4+O(c^6),
                                                    \label{p09:be15:P-exp}\\
 M-H_{\rm ID}(c)
 &=\frac{M^2}{4\kappa}c^2
   -\frac{M^4}{192\kappa^3}c^4+O(c^6).
                                                    \label{p09:be15:ID-exp}
\end{align}
Consequently
\begin{equation}\label{p09:be15:gap}
 H_{\rm P}(c)-H_{\rm ID}(c)
       =\frac{M^4}{192\kappa^2}c^4+O(c^6).
\end{equation}
For all sufficiently small positive $c$, the maximal probability
remainder has exactly one zero on $(0,\infty)$; it is a double
zero at
\[
 x_c=\frac{M}{2\kappa}c+O(c^3).
\]
Both $M-H_{\rm P}(c)$ and $x_c/c$ are analytic functions of $c^2$
near zero after the indicated leading factor is removed.
\end{theorem}
\begin{proof}
The first coefficients in \eqref{p09:be15:coefficients} are
\[
 a_1=-Mc,\quad a_2=\kappa(1+Mc^2),\quad
 a_3=-\kappa(M+2)c+O(c^3),\quad
 a_4=\kappa(\kappa+1)/2+O(c^2).
\]
Set $\eta=c^2 A$ and $x=cy$.  Parity of $a_n(c)$ implies that
$\mathcal B(c^2A,c,cy)/c^2$ is analytic in $t=c^2,A,y$ near
$t=0$ and equals
\begin{align*}
 Q(t,A,y)&=A-My+\kappa y^2+tR(A,y)+O(t^2),\\
 R(A,y)&=\kappa(M-A)y^2
       -\frac{\kappa(M+2)}2y^3
       +\frac{\kappa(\kappa+1)}{12}y^4.
\end{align*}
The double-root equations $Q=Q_y=0$ at $t=0$ have solution
\[
 A_0=\frac{M^2}{4\kappa},\qquad y_0=\frac{M}{2\kappa}.
\]
Their Jacobian in $(A,y)$ has determinant $2\kappa>0$.
The analytic implicit-function theorem gives unique analytic
functions $A(t),y(t)$ near this solution.  Moreover
\[
 A'(0)=-R(A_0,y_0)
       =-\frac{M^4(\kappa+1)}{192\kappa^3}.
\]

We verify that this local double root is the global factor boundary.
For $x$ bounded away from zero, positivity at $c=\eta=0$,
continuity on compact intervals, and the uniform large-$x$
estimate in Lemma~\ref{p09:be15:ends} exclude any other zero.
For small $x$ the convergent series gives, uniformly for
$\eta\asymp c^2$,
\[
 \mathcal B=\eta-Mcx+\kappa x^2
       +O(c^2x^2+cx^3+x^4+\eta x^2).
\]
For $x/c$ sufficiently small the positive $\eta$ term dominates;
for $x/c$ sufficiently large and $x$ still small the positive
$\kappa x^2$ term dominates.  Thus only a fixed compact $y$
interval remains.  There $Q(0,A_0,y)=\kappa(y-y_0)^2$, so
zeros are confined to a small neighborhood of $y_0$, where
$Q_{yy}>0$ and $Q_A>0$.  The double-root curve is consequently
nonnegative everywhere, and decreasing $A$ makes it negative
at its unique minimum.  It is exactly the capacity boundary.
This also proves uniqueness and order two of the contact.
Expansion of $c^2A(c^2)$ yields \eqref{p09:be15:P-exp}.

For the ID boundary minimize
$e^{-cx}(2\kappa\cosh x-1)$ at $x=cy$.  Its expansion is
\[
 M+c^2(\kappa y^2-My)
 +c^4\left(\frac{\kappa y^4}{12}-\kappa y^3+
                        \frac{My^2}{2}\right)+O(c^6).
\]
The leading minimizer is $y_0$, and its $O(c^2)$ correction
does not affect the fourth-order value.  Substitution gives
the coefficient $M^4/(192\kappa^3)$ for the minimum, proving
\eqref{p09:be15:ID-exp}.  Subtraction gives \eqref{p09:be15:gap}.
\end{proof}

\subsection{The singular lowest-order limit}

\begin{corollary}[The fourth-order gap at $\kappa=1/2$]
\label{p09:be15:loworder}
Fix $0<c\leq1$ and let $M=2\kappa-1\downarrow0$.  Then
\begin{align*}
 M-H_{\rm P}(c)
 &=\frac{c^2}{2}M^2-\frac{c^2}{2}M^3
    +\left(\frac{c^2}{2}-\frac{c^4}{16}+\frac{c^6}{24}\right)M^4
    +O(M^5),\\
 M-H_{\rm ID}(c)
 &=\frac{c^2}{2}M^2-\frac{c^2}{2}M^3
    +\left(\frac{c^2}{2}-\frac{c^4}{24}+\frac{c^6}{24}\right)M^4
    +O(M^5).
\end{align*}
In particular,
\[
 H_{\rm P}(c)-H_{\rm ID}(c)
       =\frac{c^4}{48}(2\kappa-1)^4+O((2\kappa-1)^5).
\]
At the endpoint rate $c=1$ this is the sharp singular gap left
undetermined by the preceding Kummer characterization.
\end{corollary}
\begin{proof}
Use the same analytic series with $x=My$ and $\eta=M^2 A$.
At $M=0$ the leading equation is
$A-cy+y^2/2=0$, with double root $y=c$, $A=c^2/2$.
The positive-axis inverse at $M=h=0$, away from its atom at zero,
is $e^{-cx}I_1(x)>0$, since the transform is
$(w+c)/\sqrt{(w+c)^2-1}$.  The global exclusion and the analytic
implicit-function argument from the preceding proof apply again.
For fixed $c<1$, the large-$x$ expansion remains uniform in this
limit.  At $c=1$, put $\delta=\kappa-h$ and define the positive
Euler convolution
\[
 k(x)=\frac1{\Gamma(\kappa)\Gamma(\delta)}
       \int_0^x e^{-2u}u^{\kappa-1}(x-u)^{\delta-1}\,du.
\]
Its Laplace transform is $w^{-\delta}(w+2)^{-\kappa}$.
Since $\kappa+\delta-1=\eta>0$, one has $k(0)=0$, and
\eqref{p09:be15:F} has inverse $g_{1,h}=k'+k$.
Both $\kappa$ and $\delta$ stay in a compact positive interval
about $1/2$.  Splitting the convolution with a smooth cutoff at
$u\asymp x/2$ and expanding on the side $u<x/2$ gives uniformly
\[
 k(x)=\frac{2^{-\kappa}}{\Gamma(\delta)}
               x^{\delta-1}(1+O(x^{-1})),
 \qquad k'(x)=O(x^{\delta-2}).
\]
Indeed the moments of $e^{-2u}u^{\kappa-1}$ are uniformly
finite.  On the other part set $v=x-u$ before differentiating;
the factor $v^{\delta-1}$ is uniformly integrable and the
exponential factor makes this part and its derivative
exponentially small.  Derivatives of the cutoff have the same
bound.  Thus $g_{1,h}$ has a uniformly positive algebraic tail,
without any singular lower hypergeometric parameter.  No root
can escape to infinity.

Here are the coefficients to make the expansion reproducible.
With $\kappa=(1+M)/2$, the normalized equation has the form
\[
 A-cy+y^2/2+M q_1(y)+M^2q_2(A,y)+O(M^3),
\]
where
\[
 q_1(y)=\frac{1+c^2}{2}y^2-\frac c2y^3,
\quad
 q_2(A,y)=\left(\frac{c^2}{2}-\frac A2\right)y^2
 -\left(\frac{3c}{4}+\frac{c^3}{6}\right)y^3
 +\left(\frac1{16}+\frac{c^2}{4}\right)y^4.
\]
At $y=c$, $A=c^2/2$, one has $q_1=c^2/2$ and
$q_1'=c-c^3/2$.  The next minimum coefficient is
$q_2-(q_1')^2/2=-c^2/2+c^4/16-c^6/24$.
This gives the first expansion.  Expanding the elementary ID
minimum in the same variables changes $c^4/16$ to $c^4/24$,
with the other coefficients unchanged.  Subtraction proves
the last assertion.
\end{proof}

\paragraph{What remains unresolved.}
The strict factor-region theorem and the singular quartic laws are
proved here.  They do not provide a one-point test at every intermediate
rate, nor a formula for arbitrary noninteger excess power.  Attempts
to transfer gamma deconvolution to universal SD strictness of the
squared-Bessel McKay family did not yield a complete sign proof.
Termination of the earlier finite moment hierarchy likewise remains
unproved.  These unresolved targets are recorded rather than counted
as closures.
\section{What the transfers change, and what remains open}
\label{p09:r15:remaining}

The main geometric transfer is from finite-dimensional fiber width to
an exact infinite-dimensional obstruction. The preceding paper supplied
a lower count on a compact input; the present operator theorem shows
that even arbitrary adaptive continuous information cannot recover the
noncompact part of a Hilbert operator. Homogeneity adds a different
restriction: one large level patch already forces the stronger count.
These results do not settle the unrestricted three-dimensional case,
the exact homogeneous count in every dimension, finite-query ellipsoid
constants, or an essential-norm theorem for general Banach operators.
The strict finite-budget convex-noise gap is nonquantitative; its sharp
decay in the query count remains an independent problem.

The statistical transfer changes the discrepancy induced by the
adversary. Whole-report replacement is governed by total variation,
whereas individual flips are governed by transport. The resulting
logarithmic inverse scales can differ by an unbounded factor, as the
explicit example in Section~\ref{p09:ml15:sec} shows. This is a change in
the corruption model, not an improvement over the optimal rate for
the same model. The present linear program uses the flip budget.
Simultaneous adaptation to an unknown budget, robustness of the ordinary
NPMLE, nearly linear computation, and the optimal Bernstein exponent
$\log 2$ remain unresolved.

The two stable-mixture classifications close distinct gaps. The skew
index-one case allows arbitrary location--scale dependence but fixes
the shape. The unknown-index case fixes symmetry and location zero but
allows the shape and scale to mix jointly. Its finite-design witnesses
have unbounded scale, even though every two-sided scale moment is
finite. Compact-scale finite-design failure in that latter model, fully
unknown location and skewness, and quantitative noisy-sample inversion
are not settled. The survival identity is a useful bridge to random
linear projections; it does not turn a classical projection theorem
into a newly discovered general principle.

In the spectral branch, rational relations create an observable
boundary defect and change the dimension read by an edge intensity.
Without a support restriction, even exact zero divisors on an entire
half-plane can fail to identify the function, as the explicit factor
counterexample demonstrates. The finite controlled experiment and its
critical intensity concern a different, double-zero geometry. Neither
local intensity law gives a quantitative ergodic discrepancy rate by
itself. Full recovery of all magnitudes from first intensity alone,
finite-noise blind prime-series inversion, optimal observation height,
and growing-window variance or central limit laws remain open here.

The gamma results extend the endpoint calculation across the whole
intermediate-rate interval at the structural level. They also identify
the exact local cost of insisting on an infinitely divisible remainder.
An elementary closed formula or a one-point positivity test at every
interior rate is not obtained. Arbitrary excess powers, universal strict
self-decomposability for the squared-Bessel family, and termination of
the earlier moment hierarchy remain unresolved.

\paragraph{Research status.}
The thirteen-row accounting describes the stated results and their
related refinements. It does not certify thirteen external open-problem
closures, or a particular journal or prize level. Independent automated
cross-checks within this research process examined every branch, and the
final text includes corrections and explicit scope restrictions from
those checks. External mathematical review, a broader priority search,
and further development are still needed before publication claims.
\endgroup
% ===== END SOURCE 09 =====


% ===== BEGIN SOURCE 10: KEIO_Observation_Budgets_and_Hidden_Parameters.tex =====
\clearpage
\part{Observation Budgets and Hidden Parameters}
\label{part:10}
\setcounter{section}{0}
\begin{center}\large\itshape Noncompact Recovery, Replication Thresholds, and Blind Spectral Identification\end{center}
\begin{center}\small Source manuscript: 16 October 2026.\end{center}
\begingroup
\renewcommand{\arraystretch}{1}
\begin{quote}\small\textbf{Source abstract.}\quad
We study thirteen scoped questions about the information required to
recover hidden geometric, probabilistic, and spectral structure.
Smooth strongly convex measurements with uniformly bounded Hessians
retain logarithmic adaptive query complexity in the dimension.
For the noncompact embedding $\ell_p\hookrightarrow\ell_q$, $p<q$,
the optimal continuous-information errors have orders
$n^{-(1/p-1/q)}$ nonadaptively and $2^{-n(1/p-1/q)}$ adaptively.
This gives a sharp boundary to the preceding Hilbert compactness theorem.
In latent Bernoulli populations, unequal replication yields the optimal
$B^{-1/3}$ total-budget rate, while honest uncertainty under an unknown corruption
budget has an unavoidable clean-data diameter. For stable mixtures,
removing all scale moments changes separate-terminal identification to
an exact exponential-completeness criterion; two joint replicates
nevertheless identify arbitrary continuously mixed index and scale.
Six replicates suffice when location and skewness are also mixed and
correlated. One-sided zero sets identify the unknown internally
rationally independent support of a finite exponential sum; finitely
many selected simple zeros also give a local analytic inverse at each
such support. Near-resonance creates an explicit finite-record
obstruction to uniform stability. Further
results give a sharp inverse for arbitrary positive spectra, a
Diophantine criterion for optimal marked-zero observation height, a
complete two-rate gamma factor region, and recovery of a boundary
remainder from its factor-exchange response. Full proofs, classical
inputs, prior unpublished dependencies, and remaining open questions
are distinguished. The thirteen questions are not thirteen named
historical conjectures.
\end{quote}

\section*{Scope, principal mechanisms, and claim accounting}
\addcontentsline{toc}{section}{Scope, principal mechanisms, and claim accounting}

This manuscript continues the 223-page \emph{Observation and Recovery
in Arithmetic and Probability: Prime Flows, Rare Events, and Infinite
Mixtures} and the subsequent papers, most immediately the 41-page
\emph{Sharp Boundaries for Adaptive Recovery and Mixture Identification}.
We do not re-count their results. In particular, the exact power-of-two
query count, Hilbert essential-norm floor, fixed-design binomial
corruption curve, prescribed-support zero inverse, and single-rate
gamma capacity are starting points rather than new closures here.

The organizing question is what changes when the observation resource
is changed. A norm measurement makes subsequent smooth codes strongly
convex without changing their information. A support-localizing query
makes an unstable sparse moment inverse usable in an infinite sequence
space. Joint replicates retain a common latent variable that separate
terminal laws lose. At an exponential-sum edge, shared zeros force two
initially unknown frequency tori to be the same integral graph. On a
gamma factor boundary, directional response becomes a Laplace transform
of a positive measure. These are distinct mathematical mechanisms;
no assertion of a universal equivalence between the models is intended.

\paragraph{External question versus internal extensions.}
Krieg--Novak--Ullrich explicitly ask for complexity bounds when
measurements are convex, homogeneous, or analytic
\cite[Section~3]{p03:an:KNU}. The smooth uniformly curved construction
here answers the order-of-growth part of the convex question and
leaves a one-query gap in general dimension. Our noncompact sequence
embedding gives a fully quantified family for the Banach-space
adaptive comparison discussed in their Section~4, not a general
solution of that comparison. The other questions in the table were
formulated within this observation-and-recovery program. Some are
sharp transfers of established approximation or number-theoretic
principles; they are not advertised as newly invented classical tools.

{\small\setlength{\tabcolsep}{4pt}
\begin{longtable}{@{}p{.045\textwidth}p{.29\textwidth}p{.48\textwidth}p{.12\textwidth}@{}}
\toprule
No. & Target & Result and remaining scope & Proof\\
\midrule
\endfirsthead
\toprule No. & Target & Result and remaining scope & Proof\\\midrule
\endhead
1 & Smooth strongly convex information & Logarithmic count with fixed two-sided Hessian bounds; exact general count remains open. & \ref{p10:in16:convex-upper}\\
2 & Quantitative approach to a noise floor & Explicit exponential adaptive and inverse-linear nonadaptive lower gaps; matching rates not claimed. & \ref{p10:in16:quantitative-gap}\\
3 & Noncompact Banach recovery & Sharp algebraic and exponential errors for all $\ell_p\hookrightarrow\ell_q$, $p<q$; finite-bit obstruction is a consequence. & \ref{p10:op16:main}\\
4 & Total Bernoulli observation budget & Optimal unequal-replication design and a lower bound allowing adaptive allocation. & \ref{p10:ml16:budget}\\
5 & Unknown corruption and honest uncertainty & Clean-law diameter lower bound for confidence sets honest over a specified corruption range. & \ref{p10:ml16:confidence}\\
6 & Unknown frequency support from zeros & One-sided zero-set rigidity and a finite local inverse for internally independent positive supports. & \ref{p10:ar16:thm:unknown}, \ref{p10:ar16:thm:local}\\
7 & Near-resonant observation delay & Exact first-arrival asymptotic and uniform noisy-record impossibility in a two-frequency family. & \ref{p10:ar16:thm:escape}, \ref{p10:ar16:prop:finite}\\
8 & Stable terminal designs without moments & Necessary and sufficient infinite completeness radius; matching smooth signed witnesses. & \ref{p10:ga16:thm:terminal}\\
9 & Replication of continuously mixed stable laws & Two observations are necessary and sufficient for index--scale mixing; five/six suffice for additional correlated parameters. & \ref{p10:ga16:thm:two}, \ref{p10:ga16:cor:six}\\
10 & Continuous positive spectra & Sharp $\log\log(1/\delta)/\log(1/\delta)$ inverse modulus and finite convex reconstruction. & \ref{p10:sp16:thm:nonparametric}\\
11 & Optimal marked-zero observation height & Order $h^{-1}$ holds exactly for bounded partial quotients in the one-frequency model. & \ref{p10:sp16:thm:height}\\
12 & Simultaneous gamma extraction & Complete two-rate region, curved and atomic faces, and root-free boundary remainders. & \ref{p10:be16:diagram}, \ref{p10:be16:rootless}\\
13 & Recovery from factor-exchange response & The full rate-response determines the boundary remainder and both residual exponents. & \ref{p10:be16:inverse-response}\\
\bottomrule
\end{longtable}
}

\paragraph{What the proof status means.}
Each table row is proved for its stated model below. Classical width
and completeness theorems, finite-dimensional continuous recovery,
multilevel estimation, moment matching, and stable subordination are
credited where used. Selected results also use the previous unpublished
convex-first obstruction, convex-noise floor, and binomial transport
modulus; these dependencies are identified explicitly. Separate
proof reviews within this research process are checks for gaps, not
external peer review or machine-checked formal certification. A bounded
search of primary literature does not establish universal priority.

The exact-real query results do not imply finite-bit algorithms.
Population-law identification does not give stable finite-sample
estimation automatically. The zero-set theorem concerns finite
exponential sums with explicit independence assumptions, not the
Riemann zeta function. The unrestricted value $N(3)$, the general
analytic-information problem, unrestricted infinite-support blind
inversion, and minimal replication for fully mixed stable parameters
remain open in this work.




% ===== information =====

\section{Smooth convex information and quantitative finite-query obstructions}
\label{p10:in16:section}

We address two questions that were not settled by the earlier query-count
or noise-floor results.  First, does convexity of every measurement
destroy the logarithmic advantage of adaptive continuous information?
It does not: logarithmically many measurements suffice even when all
measurements are smooth and their Hessians lie between two fixed
positive multiples of the identity.  Second, how rapidly can a finite
number of bounded convex measurements approach the exact convex-noise
floor?  We give uniform exponential lower bounds for adaptive budgets
and inverse-linear lower bounds for nonadaptive budgets.  These are
quantitative obstructions; matching convergence rates are not claimed.

\subsection{Logarithmic recovery with uniformly curved smooth queries}

In this subsection measurements are exact real numbers.  An algorithm
may choose its next function arbitrarily from the previous answers,
and its decoder may be discontinuous.  All errors are uniform over
the full space $\mathbb R^m$.

Write $N_{\rm cvx}(m)$ for the minimum query count allowing arbitrarily
accurate recovery when all selected functions are continuous and
convex.  Let $N_{\rm uc}(m)$ denote the corresponding count for the
smaller class
\begin{equation}\label{p10:in16:curvature-class}
 \mathcal C_m=\{F\in C^\infty(\mathbb R^m):
       \tfrac12 I\preceq D^2F(x)\preceq\tfrac32 I
                     \text{ for every }x\in\mathbb R^m\}.
\end{equation}
The range of a query is not bounded in this exact-information model.
The curvature bounds in \eqref{p10:in16:curvature-class} are independent
of the requested accuracy, the input, and the stage of the algorithm.

\begin{theorem}[Convexity and smoothness preserve logarithmic complexity]
\label{p10:in16:convex-upper}
For every $m\ge2$,
\begin{align}
 \lceil\log_2m\rceil+1
    &\le N_{\rm cvx}(m)
       \le\min\{m,\lceil\log_2m\rceil+2\},
                                      \label{p10:in16:convex-count}\\
 \lceil\log_2m\rceil+1
    &\le N_{\rm uc}(m)
       \le\lceil\log_2m\rceil+2.       \label{p10:in16:uc-count}
\end{align}
In particular, both counts divided by $\log_2m$ tend to one.
\end{theorem}

The convex class is one of the restricted-information questions posed
in \cite[Section~3]{p03:an:KNU}.  The new assertion is the upper bound,
including its uniform-curvature strengthening.  The first-query
quasiconvex obstruction in the preceding manuscript
\cite{comp:part:09} gives the displayed lower bounds.  Thus we do not
claim an exact count in every dimension.

We give the details of the smooth construction, since uniform smooth
approximation of a continuous information function would not suffice:
the membership tests below use exact zero values.

\begin{lemma}[Smooth membership and finite smooth codes]
\label{p10:in16:smooth-tools}
For every closed $C\subset\mathbb R^d$ there is a nonnegative
$g\in C^\infty(\mathbb R^d)$ with $g^{-1}(0)=C$.
If $K_1,\ldots,K_s$ are pairwise disjoint compact sets, there is a
compactly supported smooth function taking the constant value $i$ on
$K_i$, for $1\le i\le s$.
\end{lemma}
\begin{proof}
Cover the open complement of $C$ by a countable family of open balls
whose closures miss $C$, and choose nonnegative compactly supported
smooth functions $\phi_j$ in those balls such that their positivity
sets cover the complement.  Choose positive $a_j$ so small that
\[
 a_j\max_{|\alpha|\le j}\|D^\alpha\phi_j\|_\infty\le2^{-j}.
\]
For each derivative order the tail series converges uniformly.
Consequently $g=\sum_j a_j\phi_j$ is smooth, nonnegative, zero on
$C$, and positive off $C$.  This argument also handles an empty
complement by taking $g=0$.  For the second assertion choose disjoint
bounded open neighborhoods of the compact sets and smooth bump
functions $\psi_i$ supported in these neighborhoods and equal to one
near $K_i$.  Then $\sum_i i\psi_i$ is the required function.
\end{proof}

\begin{proof}[Proof of Theorem~\ref{p10:in16:convex-upper}]
We construct the upper bound for \eqref{p10:in16:curvature-class}.
Set $Q(x)=\|x\|_2^2/2$ and first read $Q(x)$.  This determines
$r=\|x\|_2$.  For $r=0$ no further information is needed.  Otherwise
all later construction choices may depend on this already recorded
radius.

The colored partition in \cite[Section~2]{p03:an:KNU}, applied in
$\mathbb R^{m-1}$ and scaled to accuracy $\varepsilon$, has the
following properties: its cells have Euclidean diameter at most
$\varepsilon$; there are $m$ colors; distinct cells of the same
color have a fixed positive separation; and the partition is locally
finite.  We use the closures of its cells.  Only finitely many such
closures meet $[-r,r]^{m-1}$.  For each color take the union of these
finitely many closed cells, intersected with $[-r,r]^{m-1}$, and call
the resulting compact set $E_a$.  The sets $E_a$, $1\le a\le m$,
cover that cube.  A point on a boundary may belong to several colors;
this causes no difficulty.

Lemma~\ref{p10:in16:smooth-tools} supplies an exact smooth zero-set test
for each union $\bigcup_{a\in J}E_a$.  Regard that test as a function
of $x'=(x_1,\ldots,x_{m-1})$ and multiply by a smooth cutoff in
$\mathbb R^m$ equal to one on the ball of radius $r+1$.
On the known sphere this is still an exact membership test and is now
compactly supported and smooth.  Binary search among the $m$ colors
uses at most $\lceil\log_2m\rceil$ such tests and selects a color
$a$ with $x'\in E_a$.

For this color, enumerate its compact cell pieces as
$K_1,\ldots,K_s$.  They are pairwise disjoint.  A smooth code
$c(x')$ takes the integer value $i$ on $K_i$ by
Lemma~\ref{p10:in16:smooth-tools}.  Fix a smooth bijection
$h:\mathbb R\to(0,1)$, for example
$h(t)=\tfrac12+\pi^{-1}\arctan t$.  The last intended measurement
is $c(x')+h(x_m)$, again multiplied by a cutoff equal to one on the
known sphere.  Its integer part identifies $K_i$, and its fractional
part determines $x_m$ exactly.  Any representative of $K_i$ then
approximates $x'$ within $\varepsilon$.

It remains to implement each intended compactly supported smooth
function $g$ using an admissible query.  Its Hessian has a finite
global bound $M_g$.  Choose $\alpha_g>0$ with
$\alpha_g M_g\le1/2$ and read
\[
                         F_g=Q+\alpha_g g.
\]
This belongs to \eqref{p10:in16:curvature-class}.  Since $Q(x)$ is already
known exactly, the answer determines $g(x)$ exactly by subtraction
and division by $\alpha_g$.  The argument uses no lower bound on
$\alpha_g$ and no robustness to rounding.  The total count is at
most $1+\lceil\log_2m\rceil+1$.  The convex upper bound $m$ is
given separately by coordinate queries, which need not belong to
the uniformly curved class.  The earlier convex-first lower bound
applies to both classes and completes the proof.
\end{proof}

\paragraph{What this does not establish.}
The construction contains smooth functions flat on open sets and is
not an analytic or polynomial construction.  The subtraction of two
exact real measurements is essential; the theorem does not imply the
same performance with a fixed amount of observation noise.  In
particular it neither closes the unrestricted two-query problem in
$\mathbb R^3$ nor the analytic-information question.

\subsection{A quantitative obstruction at the convex-noise floor}

We now return to continuous convex queries $f:B_2^m\to[-1,1]$,
each observed with adversarial additive error at most $\delta$, where
$m\ge2$ and $0<\delta<1$.  Let $e_n^{\rm ad}$ and
$e_n^{\rm non}$ be the optimal worst-case Euclidean errors for
adaptive and nonadaptive $n$-query procedures, respectively.
Selection rules and decoders remain arbitrary.  Put
\begin{equation}\label{p10:in16:noise-parameters}
 q=1-\delta,\qquad
 r=\frac{\delta}{\sqrt{1+q^2+2q/m}},\qquad
 \rho=\frac{1+r}{2},\qquad L=\frac{2}{1-\rho}.
\end{equation}
The earlier exact floor is $\inf_n e_n^{\rm ad}=r$
\cite{comp:part:07}.  The preceding manuscript established strict
inequality $e_n^{\rm ad}>r$ for each finite $n$ by compactness
\cite{comp:part:09}.  The next theorem replaces that existence
argument with quantitative estimates.

\begin{theorem}[Quantitative finite-budget gaps]
\label{p10:in16:quantitative-gap}
There are constants $C_{\rm ad},C_{\rm non}>0$ and $0<\theta<1$,
depending only on $m$ and $\delta$, such that for every $n\ge1$,
\begin{equation}\label{p10:in16:gaps}
 e_n^{\rm ad}\ge r+C_{\rm ad}\theta^{n-1},
 \qquad
 e_n^{\rm non}\ge r+\frac{C_{\rm non}}n.
\end{equation}
The constants are given explicitly below after choosing one generic
skew-symmetric matrix.  Consequently reaching error $r+\varepsilon$
requires at least a constant multiple of $\log(1/\varepsilon)$
adaptive queries, and at least a constant multiple of
$1/\varepsilon$ nonadaptive queries, for sufficiently small
$\varepsilon$.
\end{theorem}

The second inequality does not prove a separation between the two
models: no matching adaptive upper bound is asserted.  General noisy
continuous information and its entropy-number description are studied
in \cite{p06:ni:KNPU}; the additional convexity and bounded-range
constraints here create the positive floor around which we quantify
finite-budget behavior.

Let $u_0,\ldots,u_m$ be regular-simplex directions, with
$\|u_i\|_2=1$ and $\langle u_i,u_j\rangle=-1/m$ for $i\ne j$.
Choose a skew-symmetric matrix $A$ with $\|A\|_{\rm op}\le1$ and
\[
                    \omega:=\min_{i\ne j}|u_j\cdot Au_i|>0.
\]
Such an $A$ exists: for each noncollinear pair $u_i,u_j$ the
corresponding linear functional on the space of skew-symmetric
matrices is nonzero, and finitely many proper hyperplanes cannot
cover that space.  A rational matrix outside these hyperplanes can
be chosen and then normalized.  No optimization of this choice is
required.  Define the positive constants
\begin{align}
 d&=\sqrt{2+2/m},&
 v_0&=\frac{\delta}{rd},&
 b&=\frac{v_0r^2\omega}{2\delta},\nonumber\\
 \eta_0&=\min\left\{\delta,
                  \frac{v_0^3r^2\omega^2}{8\delta L^2}\right\},&
 T&=\min\left\{1,\frac{b}{2Lr}\right\},&
 B_0&=\frac{2}{b\delta q},                 \label{p10:in16:constants}\\
 P&=m(m+1),&
 c&=\min\left\{\frac{\eta_0}{T},\frac{1}{PB_0}\right\}.\nonumber
\end{align}

\begin{lemma}[Short near-saturation intervals]
\label{p10:in16:short-intervals}
Let $f$ be convex and bounded between $-1$ and $1$ on the open unit
ball.  Fix $U\in SO(m)$ and put
$x_i(t)=rU\exp(tA)u_i$ for $-T\le t\le T$.
For $0<\eta\le\eta_0$ and each ordered pair $i\ne j$, the set
\[
 E_{ij}(f,\eta)=
 \{t\in[-T,T]:f(x_j(t))-f(x_i(t))\ge2\delta-\eta\}
\]
has diameter at most $B_0\eta$.
\end{lemma}
\begin{proof}
At a time in this set write $x=x_i(t)$, $y=x_j(t)$ and
$z=(y-qx)/\delta$.  Equation~\eqref{p10:in16:noise-parameters} gives
$\|z\|_2=1$.  Let $\ell(w)=a+v\cdot w$ be a supporting affine
function for $f$ at $y$.  It satisfies $\ell\le f$ on the open
ball and $\ell\le1$ on its closure.  With
$E=f(x)-\ell(x)\ge0$ and $D=1-\ell(z)\ge0$, the affine identity
$\ell(y)=q\ell(x)+\delta\ell(z)$ gives
\begin{equation}\label{p10:in16:deficit-identity}
 \delta(f(x)+1)+qE+\delta D
               =2\delta-(f(y)-f(x))\le\eta.
\end{equation}
In particular,
\begin{equation}\label{p10:in16:near-minimum}
 f(x)\le-1+\eta/\delta,
 \qquad \ell(x)\ge-1-\eta/q,
 \qquad D\le\eta/\delta.
\end{equation}

Every supporting slope at a point of $\rho B_2^m$ has norm at most
$L$, by comparing its affine support with the range bound along
segments of length approaching $1-\rho$.  Also
\[
 v\cdot(y-x)=\ell(y)-\ell(x)
        \ge f(y)-f(x)\ge2\delta-\eta\ge\delta,
\]
so $\|v\|_2\ge v_0$.  Since $a+\|v\|_2\le1$,
\[
 \left\|\frac{v}{\|v\|_2}-z\right\|_2
       \le\sqrt{\frac{2D}{\|v\|_2}}
       \le\sqrt{\frac{2\eta}{\delta v_0}}.
\]
The simultaneous rotation preserves its skew coupling:
\[
 z\cdot x_i'(t)=\frac{r^2}{\delta}\,u_j\cdot Au_i.
\]
The choice of $\eta_0$ therefore implies
\begin{equation}\label{p10:in16:transverse-slope}
 |v\cdot x_i'(t)|\ge b,
 \qquad
 \operatorname{sgn}(v\cdot x_i'(t))
          =\operatorname{sgn}(u_j\cdot Au_i).
\end{equation}
Indeed, the directional error is at most
$Lr\sqrt{2\eta/(\delta v_0)}\le b$, whereas the absolute main
term is at least $2b$.

Take two times $t<s$ in $E_{ij}(f,\eta)$ and let $\Delta=s-t$.
If the sign in \eqref{p10:in16:transverse-slope} is positive, use the
supporting affine function at time $t$; if negative, use the one at
time $s$ and run the comparison backwards.  Since
$\|x_i''(t)\|_2\le r$, Taylor's formula for the curve and
\eqref{p10:in16:near-minimum} yield in either case
\[
 b\Delta-\frac{Lr}{2}\Delta^2
       \le\eta\left(\frac1\delta+\frac1q\right)
       =\frac{\eta}{\delta q}.
\]
Our choice $\Delta\le2T\le b/(Lr)$ makes the left side at least
$b\Delta/2$.  Thus $\Delta\le B_0\eta$, as asserted.
\end{proof}

\begin{corollary}[A uniform linear oscillation deficit]
\label{p10:in16:linear-deficit}
For every such $f$, every $U\in SO(m)$, and $0<a\le T$, there is
$V\in SO(m)$ with $\|V-I\|_{\rm op}\le a$ such that
\[
 \operatorname{osc}_{\{rUVu_i\}_{i=0}^m} f
                            <2\delta-ca.
\]
\end{corollary}
\begin{proof}
Apply Lemma~\ref{p10:in16:short-intervals} on $[-a,a]$ with
$\eta=ca\le\eta_0$.  Each of the $P$ near-saturation sets lies
in an interval of length at most $B_0ca$.  Their union has length
at most $PB_0ca\le a$, so it cannot cover $[-a,a]$.  At a point
outside the union every ordered difference is below
$2\delta-ca$.  Set $V=\exp(tA)$; then
$\|V-I\|_{\rm op}\le|t|\|A\|_{\rm op}\le a$.
\end{proof}

\begin{proof}[Proof of Theorem~\ref{p10:in16:quantitative-gap}]
For the adaptive bound define
\begin{equation}\label{p10:in16:final-constants}
 \theta=\min\left\{\frac12,\frac{c}{8Lr}\right\},
 \qquad
 C_{\rm ad}=\min\left\{\frac{\rho-r}{2},\frac{cT}{16L}\right\},
 \qquad
 C_{\rm non}=\min\left\{\frac{\rho-r}{2},\frac{cT}{4L}\right\}.
\end{equation}
Maintain a ball of simplex rotations all compatible with the same
recorded noisy transcript.  Start with any center and radius $T$,
using operator-norm distance on $SO(m)$.  If the current radius is
$a$, the transcript fixes the next selected function $f$.
Corollary~\ref{p10:in16:linear-deficit}, applied within distance $a/2$
of the current center, provides a simplex on which the oscillation
is below $2\delta-ca/2$.  Answer with the midpoint of the maximum
and minimum on that simplex.  Every vertex then has noise slack
at least $ca/4$.

Retain the ball of radius $\theta a$ about the new center.  It is
contained in the old ball, and a vertex moves by at most
$r\theta a$.  Since the query is $L$-Lipschitz on $\rho B_2^m$,
the answer remains feasible throughout the new ball with slack
at least $ca/8$.  Previously recorded answers retain their earlier
slack by nesting.  After $n$ stages the radius is $T\theta^n$ and
every answer has slack at least $cT\theta^{n-1}/8$.

Choose any retained simplex and increase its radius by
\[
 h=\min\left\{\frac{\rho-r}{2},
               \frac{cT\theta^{n-1}}{16L}\right\}
                   \ge C_{\rm ad}\theta^{n-1}.
\]
Every enlarged vertex remains in $\rho B_2^m$ and all query values
change by at most $Lh$, so the same transcript is feasible for all
these vertices.  Their circumradius is $r+h$.  More explicitly,
for every decoded point $w$, averaging squared distances over the
centered regular simplex gives
\[
 \frac1{m+1}\sum_i\|(r+h)Uu_i-w\|_2^2
                         =(r+h)^2+\|w\|_2^2.
\]
Thus some vertex has error at least $r+h$.  The construction is
valid for arbitrary, even discontinuous, query-selection rules
because each query is fixed only after its transcript is fixed.

For $n$ fixed nonadaptive functions, put $\eta=cT/n$ and use the
single rotation interval $[-T,T]$.  The union of all their $nP$
near-saturation sets has length at most
$nPB_0\eta\le T<2T$.  One rotation consequently gives
oscillation below $2\delta-\eta$ for every function at once.
Choose their midpoint answers and inflate by
\[
 h=\min\left\{\frac{\rho-r}{2},\frac{\eta}{4L}\right\}
                                      \ge C_{\rm non}/n.
\]
The same simplex argument proves the second bound.
Solving these inequalities for $n$ proves the stated necessary
query complexities.
\end{proof}

\paragraph{Open targets and failed extensions.}
The unrestricted exact value $N(3)$ remains unresolved here.  The
available width recursion reaches index $2^2-1=3$, which gives no
positive width obstruction on a three-dimensional cube; replacing
that recursion by a stronger assertion for scalar targets has not
been justified.  For analytic first queries, the large-level graphs
available for polynomials or convex functions need not follow from
analyticity alone.  Neither an unrestricted two-query construction
nor an analytic impossibility proof is claimed.  For the noisy problem
we have quantified the earlier strict gap, but have not proved that
the exponential adaptive bound or inverse-linear nonadaptive bound
is sharp.


% ===== operators =====

\section{Sharp adaptive recovery of a noncompact sequence embedding}
\label{p10:op16:section}

The Hilbert-space compactness criterion from the preceding paper
\cite{comp:part:09} does not extend to arbitrary Banach operators.
We prove this with a complete error-rate calculation for the canonical
embedding $I_{p,q}:\ell_p\to\ell_q$, where $1\le p<q\le\infty$.
The input is the closed unit ball $B_p$ with its $\ell_p$ topology.
For $q=\infty$ we may take the output space to be either $c_0$ or
$\ell_\infty$, with its supremum norm. Write
\[
 \alpha=\frac1p-\frac1q>0,
\]
with $1/\infty=0$. Let $e_n^{\rm ad}$ and $e_n^{\rm non}$ denote
infimal worst-case errors using at most $n$ continuous real-valued
measurements, adaptively and nonadaptively, respectively. The decoder
and the dependence of later queries on earlier answers are unrestricted.
All observations in these two definitions are exact real numbers.

\begin{theorem}[Noncompact recovery and its exact decay exponents]
\label{p10:op16:main}
For every $n\ge1$,
\[
 e_n^{\rm ad}(I_{p,q})\ge \tfrac12 2^{-\alpha n},
 \qquad e_n^{\rm non}(I_{p,q})\ge(n+1)^{-\alpha}.
\]
For every integer $s\ge1$,
\[
 e_{2s}^{\rm non}(I_{p,q})\le s^{-\alpha},\qquad
 e_{\lceil\log_2s\rceil+3}^{\rm ad}(I_{p,q})\le s^{-\alpha}.
\]
The nonadaptive upper bound can be attained with all queries
$1$-Lipschitz on $B_p$. Consequently
\begin{equation}\label{p10:op16:rates}
 e_n^{\rm non}(I_{p,q})\asymp n^{-\alpha},\qquad
 e_n^{\rm ad}(I_{p,q})\asymp2^{-\alpha n}.
\end{equation}
In particular, a bounded noncompact linear operator can have uniformly
vanishing error under continuous nonlinear information. Adaptivity
changes the optimal algebraic order into an exponential one, with
the same exponent $\alpha$.
\end{theorem}

The macroscopic-width behavior of finite-dimensional $\ell_p$ balls
in the $\ell_q$ metric is classical, and was studied by Tsukamoto
\cite{p10:op16:Tsukamoto}. The finite-dimensional adaptive compression
used below is the construction of Krieg--Novak--Ullrich
\cite{p03:an:KNU}. The matching adaptive converse uses the preceding
adaptive-width argument. Our contribution here is the explicit
infinite-sequence encoding, its adaptive localization, and the resulting
two-sided noncompact-operator rate theorem. The general Banach-space
comparison asked in \cite[Section~4]{p03:an:KNU} is not settled by this
one family of operators.

\subsection{A finite sparse code for an infinite sequence}

For $a>0$, define the soft threshold
\[
 (T_ax)_j=\operatorname{sgn}(x_j)(|x_j|-a)_+.
\]
If $x\in B_p$ and $a=s^{-1/p}$, the sequence $T_ax$ has at most
$s$ nonzero coordinates and $\|T_ax\|_p\le1$. Moreover,
\begin{equation}\label{p10:op16:threshold}
 \|x-T_ax\|_q\le s^{-\alpha}.
\end{equation}
For finite $q$, this follows from
$\min\{|x_j|,a\}^q\le a^{q-p}|x_j|^p$ and summation; for
$q=\infty$ it is the coordinatewise bound $a$.

Define the $2s$ moment queries
\begin{equation}\label{p10:op16:moments}
 L_l(x)=\sum_{j\ge1}j^{-l}(T_ax)_j,
                         \qquad l=0,\ldots,2s-1.
\end{equation}
These are finite sums at every $x\in\ell_p$ and are continuous in
the $\ell_p$ topology. Indeed, for any fixed $x$, all sufficiently
large coordinates have magnitude less than $a/2$. On a sufficiently
small $\ell_p$ neighborhood of $x$, all those coordinates remain
below $a$, so the displayed sum is locally a fixed finite sum of
continuous functions. On $B_p$ the stronger bound
\[
 |L_l(x)-L_l(z)|\le\|T_ax-T_az\|_1
 \le(2s)^{1-1/p}\|x-z\|_p
\]
holds: the difference is supported on at most $2s$ coordinates, and
soft thresholding contracts the $\ell_p$ distance. Dividing each
query by $(2s)^{1-1/p}$ makes it $1$-Lipschitz, without losing any
information because the responses are exact.

The map in \eqref{p10:op16:moments} determines every $s$-sparse sequence
uniquely. If two such sequences have the same moments, their difference
is supported on distinct integers $j_1,\ldots,j_r$, with $r\le2s$.
Its first $r$ zero moments at the distinct nodes $1/j_1,\ldots,1/j_r$
form a square Vandermonde system, whose
determinant is nonzero. All its coefficients therefore vanish.
An unrestricted decoder returns the unique sparse sequence when it
exists, and zero otherwise. Equation~\eqref{p10:op16:threshold} proves
the nonadaptive upper bound.

\subsection{Localizing the code before adaptive compression}

Direct approximation of the moment vector is insufficient: its inverse
need not have a uniform modulus over all integer supports. One additional
continuous query removes exactly that obstruction.

Choose a continuous function $\psi:[0,\infty)\to[0,1]$ that is zero
on $[0,1/2]$ and one on $[1,\infty)$. Query
\begin{equation}\label{p10:op16:localizer}
 H_a(x)=\sum_{j\ge1}j\,\psi(|x_j|/a).
\end{equation}
The same local-finiteness argument proves continuity on $\ell_p$.
Set $M=\max\{1,\lceil H_a(x)\rceil\}$. If $|x_j|\ge a$, then
$H_a(x)\ge j$; consequently $T_ax$ is supported in
$\{1,\ldots,M\}$.

For each recorded $M$, consider the compact set
\[
 K_{s,M}=\{y\in\R^M:\|y\|_p\le1,\ |\operatorname{supp}y|\le s\}
\]
and its moment map $V_M:K_{s,M}\to\R^{2s}$,
$(V_My)_l=\sum_{j=1}^M j^{-l} y_j$. The preceding Vandermonde argument
shows injectivity. More quantitatively, since there are finitely many
supports of size at most $2s$ in $\{1,\ldots,M\}$, there is a finite
$C_{s,M}>0$ such that
\begin{equation}\label{p10:op16:inverse}
 \|y-z\|_q\le C_{s,M}\|V_My-V_Mz\|_\infty
                   \qquad(y,z\in K_{s,M}).
\end{equation}
For each support this follows by a left inverse of its full-column-rank
Vandermonde matrix, and one takes the maximum of the finitely many
operator norms. The output vectors are padded by zeros when $M$ varies.

Fix any additional output tolerance $\eta>0$. After observing $M$,
apply the KNU procedure to the $2s$-dimensional vector
\[
 Z_M(x)=V_M\big((T_ax)_1,\ldots,(T_ax)_M\big)
\]
with accuracy $\eta/(2C_{s,M})$ in the maximum norm. For each fixed
history and hence fixed $M$, this vector and every query composed with
it are continuous functions of $x$ on the entire input space. The
number of additional measurements is at most
$\lceil\log_2(2s)\rceil+1=\lceil\log_2s\rceil+2$.
Given its approximation $\widehat Z$, choose a minimizer
$\widehat y\in K_{s,M}$ of $\|V_My-\widehat Z\|_\infty$.
Such a minimizer exists by compactness. Since the true $T_ax$ is
feasible, \eqref{p10:op16:inverse} yields
$\|\widehat y-T_ax\|_q\le\eta$. Combining with
\eqref{p10:op16:threshold} gives error at most $s^{-\alpha}+\eta$.
Taking the infimum over $\eta>0$ proves the adaptive upper bound.
The discontinuity of $M$ as a function of the recorded answer is
permitted by the model; it is not being used as a continuous query.

\subsection{Matching lower bounds}

For nonadaptive information, restrict to the first $m=n+1$
coordinates. Borsuk--Ulam on the $\ell_p$ unit sphere gives antipodal
$x,-x$ with identical data. Every common answer makes error at least
$\|x\|_q\ge m^{-\alpha}\|x\|_p=m^{-\alpha}$.

For the adaptive bound we recall, and justify, the width input.
Let $\operatorname{UW}_d(X)$ be the infimum of the mesh of finite
open covers of a compact metric space $X$ with multiplicity at most
$d+1$. The scalar fiber inequality
\cite[Theorem~2.1]{p08:in14:BB} is
\[
 \operatorname{UW}_{2d+1}(X)
 \le\sup_t\operatorname{UW}_d(f^{-1}(t))
                   \quad (f:X\to\R\text{ continuous}).
\]
If an $n$-query procedure has error $r$, a terminal feasible set has
diameter at most $2r$. Applying this inequality inductively to the
fixed-history continuous queries gives
\begin{equation}\label{p10:op16:width}
 \operatorname{UW}_{2^n-1}(X)\le2r.
\end{equation}
No continuity of the choices across histories is needed.

Take $m=2^n$ and $X=B_p^m$, with its $\ell_q$ metric. The classical
width-dimension inequality recalled in \cite[Eq.~(2)]{p10:op16:Tsukamoto}
implies $\operatorname{UW}_{m-1}(B_p^m,\ell_p)\ge1$.
Since $\|u-v\|_q\ge m^{-\alpha}\|u-v\|_p$ in $\R^m$, pulling
back the same covers gives
\[
 \operatorname{UW}_{m-1}(B_p^m,\ell_q)\ge m^{-\alpha}.
\]
Equation~\eqref{p10:op16:width} proves $r\ge\tfrac12m^{-\alpha}$.
This completes the proof of Theorem~\ref{p10:op16:main}.

\subsection{What finite precision removes}

The standard unit vectors have pairwise distance $2^{1/q}$ for
finite $q$, and one for $q=\infty$. Thus $I_{p,q}$ is noncompact.
The following consequence isolates the role of exact real information;
it is an elementary finite-codebook obstruction, not a separate claim
of historical novelty.

\begin{proposition}[Every finite transcript alphabet has error one]
\label{p10:op16:bits}
Take the output space to be $\ell_q$ for $q<\infty$, or $c_0$ for
$q=\infty$. If the complete transcript has a finite alphabet, then
the smallest worst-case error for $I_{p,q}$ is exactly one. This remains
true for randomized procedures with finitely many possible outputs
conditional on each realization of their random seed, when the risk
is worst-case expected norm error.
\end{proposition}
\begin{proof}
Every finite codebook $D$ in either output space satisfies
\[
 \liminf_{j\to\infty}\min_{z\in D}\|e_j-z\|_q\ge1,
\]
because its finitely many coordinate tails vanish uniformly, and
$\|e_j-z\|_q\ge|1-z_j|$. This proves the deterministic lower bound.
For a randomized procedure condition on its seed $\omega$, let
$D_\omega$ be its finite output codebook, and apply Fatou's lemma
to the nonnegative quantities
$\min_{z\in D_\omega}\|e_j-z\|_q$. Their limiting inferior is
at least one for every seed. Hence the supremum of expected errors
over the inputs $e_j$ is at least one. The algorithm returning zero
has error at most one and attains the bound. In the randomized case
all maps needed for the stated expectations are assumed measurable.
\end{proof}

The $c_0$ qualification matters: an arbitrary codeword in $\ell_\infty$
need not have vanishing coordinates. The proposition does not assert
the same exact constant for unrestricted $\ell_\infty$ codewords.
For instance $p=1,q=2$ gives
$e_n^{\rm ad}\asymp2^{-n/2}$ with exact continuous real observations,
while every fixed finite bit budget has error one. This is a statement
about observation models, not a finite-precision compression algorithm.


% ===== mle =====

\section{Optimal replication budgets and the price of honest robustness}
\label{p10:ml16:section}

We study two changes to the observation model. First, the experimenter
may distribute a total budget of Bernoulli observations among individuals,
including by an adaptive rule. We determine the minimax population
recovery rate for this experiment. Second, with a fixed rectangular
experiment, we ask whether confidence sets can adapt to the actual
number of corrupted bits. A sharp obstruction already occurs at one
fixed, uncorrupted population. These are distinct questions: the second
result does not rule out adaptation of point estimators.

\subsection{A total-budget theorem, including adaptive allocation}

Let $G$ be an arbitrary probability law on $[0,1]$. Each fresh individual
has an independent parameter $P\sim G$, and its observations are
conditionally independent Bernoulli$(P)$ trials. At every stage a
sampling policy may introduce a fresh individual or revisit any earlier
individual. The choice may depend arbitrarily on all previous choices
and observations and on private randomization. The number of trials is
at most an integer $B$ almost surely. Write
\[
 \mathfrak R_B=\inf_{\mathcal S,\widehat G}\sup_G
       \mathbb E_{G,\mathcal S}W_1(\widehat G,G),
\]
where the infimum includes all such policies and estimators. Introducing
an individual has no cost beyond the trials assigned to it.

\begin{theorem}[The unrestricted Bernoulli-budget law]
\label{p10:ml16:budget}
There are absolute positive constants $c,C$ such that, for every
integer $B\ge1$,
\begin{equation}\label{p10:ml16:budget-rate}
                  cB^{-1/3}\le\mathfrak R_B\le CB^{-1/3}.
\end{equation}
The upper bound is attained by a predetermined heterogeneous design and
a finite linear program of size polynomial in $B$. The lower bound
allows fully adaptive allocation, arbitrary revisits, and arbitrary
randomization. No smoothness, finite-support assumption, or separation
condition on $G$ is imposed.
\end{theorem}

The multilevel idea is classical. Bujok, Hambly and Reisinger
\cite{p10:ml16:BHR} use increasing Bernoulli replication levels to estimate
fixed functionals of a latent population. Giles, Nagapetyan and Ritter
\cite{p01:info:GNR2015} develop multilevel estimation of distribution functions
under regularity assumptions. Adaptive sampling of unknown coins is
studied by Lee and Valiant \cite{p01:info:LeeValiant2021} for estimating the proportion
of coins on either side of a prescribed gap. Here the target is the
entire arbitrary mixing law in $W_1$, and the lower bound covers all
adaptive designs. The theorem combines an integrated indicator-variance
calculation with an adaptive information lower bound. We do not claim
that either multilevel simulation or adaptive coin sampling originates
here. The primary sources checked do not state the exact unrestricted
population theorem \eqref{p10:ml16:budget-rate}; this is a bounded literature
comparison, not an assertion that every possible precedent is excluded.

\begin{proof}[Upper bound]
Put $t_l=2^l$, for $0\le l\le L$, and $T=2^L$. For every level use
an independent cohort of $N_l$ fresh individuals. At level zero observe
one trial from each individual. At level $l\ge1$ observe $t_l$ trials
and let $A_l$ be their full average and $C_l$ their first-half average.
For each individual define the random function
\[
 D_l(x)=\mathbf 1_{\{A_l\le x\}}-\mathbf 1_{\{C_l\le x\}},
                    \qquad 0\le x\le1.
\]
Let $G_t$ be the law of a normalized $t$-trial count under $G$,
write $F_t$ for its distribution function, and set
\[
 \widetilde F(x)=\widehat F_1(x)
       +\sum_{l=1}^L\frac1{N_l}\sum_{i=1}^{N_l}D_{l,i}(x).
\]
The expectation telescopes: $\mathbb E\widetilde F=F_T$.
Conditional on $P$, the full average minus the first-half average has
mean zero and variance $P(1-P)/t_l$. Consequently
\begin{equation}\label{p10:ml16:indicator-variance}
 \int_0^1\operatorname{Var}(D_l(x))\,dx
 \le\mathbb E|A_l-C_l|\le\frac1{2\sqrt{t_l}}.
\end{equation}
For the level-zero indicator the integrated variance is at most $1/4$.
Independence of the cohorts, followed by Cauchy--Schwarz on $[0,1]$,
gives
\begin{equation}\label{p10:ml16:stochastic}
 \mathbb E\|\widetilde F-F_T\|_{L^1}
       \le C\left(\sum_{l=0}^L\frac{2^{-l/2}}{N_l}\right)^{1/2}.
\end{equation}

Let $S_L=\sum_{l=0}^L2^{l/4}$ and choose
\[
 N_l=\left\lfloor\frac{B2^{-3l/4}}{2S_L}\right\rfloor.
\]
Choose $L$ so that $T$ is between two fixed positive multiples of
$B^{2/3}$. For all sufficiently large $B$, every expression inside the
floor is at least two. Indeed its minimum occurs at $l=L$ and is
bounded below by a constant times $B/T$. Thus rounding loses at most
a factor two and
\[
 \sum_lN_lt_l\le B/2,\qquad
 \sum_l\frac{2^{-l/2}}{N_l}\le\frac{4S_L^2}{B}
                                  \le\frac{CT^{1/2}}B.
\]
All quantities in $\widetilde F$ are constant between consecutive
points of the grid $\{0,1/T,\ldots,1\}$. Project $\widetilde F$ in
$L^1$ onto distribution functions supported on this grid. This is a
finite linear program: the grid values must be nondecreasing, between
zero and one, with terminal value one. Since $F_T$ is feasible, any
minimizer $\widehat F$ satisfies
\[
 \|\widehat F-F_T\|_{L^1}
                      \le2\|\widetilde F-F_T\|_{L^1}.
\]
Let $\widehat G$ be its probability law. Coupling $P$ to its normalized
$T$-trial count gives $W_1(G,G_T)\le1/(2\sqrt T)$.
The identity between one-dimensional $W_1$ and integrated distribution
function distance now yields
\[
 \mathbb EW_1(\widehat G,G)
       \le C\left(T^{-1/2}+\frac{T^{1/4}}{\sqrt B}\right)
       \le CB^{-1/3}.
\]
The finitely many smaller $B$ are covered by increasing $C$ and using a
fixed estimate. The grid has at most $B+1$ points for large $B$, and all
linear-program data are rational with polynomial bit length. This also
proves the computational assertion.
\end{proof}

\begin{proof}[Lower bound, including adaptive policies]
Fix an integer $m$ and set
\[
 c_j=\frac14+\frac{j-1/2}{2m},\quad h=\frac1{8m},\quad
 G_v=\frac1m\sum_{j=1}^m\delta_{c_j+hv_j},
                     \qquad v\in\{-1,1\}^m.
\]
The two possible locations in each cell are separated from those in
all other cells and belong to $[1/4,3/4]$. Monotone coupling gives
\begin{equation}\label{p10:ml16:hypercube-distance}
 W_1(G_v,G_w)=\frac{2h}{m}d_H(v,w).
\end{equation}
Enlarge the experiment: reveal the cell of every individual and allow
the learner to choose any cell for its next observation. In this
oracle experiment, querying cell $j$ returns an independent
Bernoulli$(c_j+hv_j)$ observation. It is more informative than the
original experiment restricted to these populations: a fresh person
in the original model is simulated by choosing a uniform cell and
remembering its label, and subsequent visits query that cell.

Let $\mathbb P_v$ be the complete transcript law of any adaptive policy
in this enlarged experiment, including its random seed. Let $T_j$
be the number of queries to cell $j$. If $v^{(j)}$ changes only sign
$j$, the chain rule for relative entropy gives
\begin{equation}\label{p10:ml16:adaptive-kl}
 D(\mathbb P_v\Vert\mathbb P_{v^{(j)}})
                     \le C h^2\mathbb E_vT_j.
\end{equation}
Indeed, conditional on an identical history the policy uses the same
choice kernel under both laws. Only a query to cell $j$ has a different
response distribution, and
$D(\operatorname{Ber}(p)\Vert\operatorname{Ber}(q))
\le(p-q)^2/[q(1-q)]\le Ch^2$ here. A stopping policy can be padded by
uninformative symbols, so the same chain rule applies to a hard budget.
Summing \eqref{p10:ml16:adaptive-kl} over $v,j$, using $\sum_jT_j\le B$,
and applying Pinsker's inequality and Jensen's inequality proves
\[
 \frac1{m2^m}\sum_{v,j}
 \|\mathbb P_v-\mathbb P_{v^{(j)}}\|_{\rm TV}
                                  \le\sqrt{CB/m^3}.
\]
Choose $m=\lceil K B^{1/3}\rceil$ for a sufficiently large absolute
$K$, making the right side at most $1/2$.

Decode an arbitrary estimate $\widehat G$ to a nearest member
$G_{\widehat v}$ of the finite family. The triangle inequality and
\eqref{p10:ml16:hypercube-distance} imply
$W_1(\widehat G,G_v)\ge(h/m)d_H(\widehat v,v)$.
For each coordinate, pairing the two neighboring hypotheses shows that
the average probability of a sign error is at least half of one minus
their total variation distance. Averaging over all coordinates therefore
gives
\[
 2^{-m}\sum_v\mathbb E_vd_H(\widehat v,v)\ge m/4.
\]
The average Wasserstein risk is at least $h/4\ge cB^{-1/3}$.
The oracle experiment was easier than the original one, so its lower
bound proves the required assertion.
\end{proof}

\paragraph{Comparison with equal replication.}
The earlier fixed-replication population theorem gives, after optimizing
$Nt\le B$ with the same $t$ for every individual, rate
$(B\log(eB))^{-1/4}$. The improvement in
Theorem~\ref{p10:ml16:budget} is a change in the power of the budget, made
possible by using different depths. It is not an improvement of that
earlier theorem within its own observation class. Adaptivity itself is
unnecessary for attaining the new minimax order.

\subsection{Honest confidence sets cannot learn the actual corruption budget}

We return to $N$ independent individuals, each observed $t$ times.
An adversary may, after observing the entire clean array, flip at most
$q=\lfloor\bar\varepsilon Nt\rfloor$ bits, where
$0\le\bar\varepsilon\le1/4$ is an upper bound. Equivalently, counts
$J_i$ may be changed to $J_i'$ with
$\sum_i|J_i-J_i'|\le q$. The upper bound is known; the actual number
of changed bits is not. Define
\[
 L_t(u)=\log\left(e+\frac1{u\sqrt t}\right),\quad L_t(0)=+\infty,
\]
\begin{equation}\label{p10:ml16:robust-rate}
 r_{N,t}(u)=\min\left\{1,\frac1t+
 \frac1{\sqrt{t\log(eN)}}+\frac1{\sqrt N}
              +u+\frac1{\sqrt{tL_t(u)}}\right\}.
\end{equation}
A random set $\mathcal C$ of mixing laws is honest at level $1-\alpha$
if, for every $G$ and every allowed adversary $A$,
\[
               \mathbb P_{G,A}\{G\in\mathcal C\}\ge1-\alpha.
\]
All sets and events below are assumed measurable; auxiliary
randomization in a confidence procedure is permitted.

\begin{theorem}[Sharp honest diameter at a fixed clean population]
\label{p10:ml16:confidence}
Let $0<\alpha\le1/8$ and let $\mathcal C$ be honest for the hard
bit budget above. For the fixed arcsine law
$\pi(dx)=dx/(\pi\sqrt{x(1-x)})$, with no corruption at all,
\begin{equation}\label{p10:ml16:confidence-lower}
 \mathbb P_{\pi,\mathrm{id}}
 \left\{\operatorname{diam}_{W_1}(\mathcal C)
                      \ge c\,r_{N,t}(\bar\varepsilon)\right\}
                              \ge1-2\alpha-\frac18.
\end{equation}
Here $c>0$ is absolute and does not depend on $N,t,\bar\varepsilon$.
Conversely, for each fixed $\alpha$ there are honest confidence sets
whose diameter is everywhere at most
$C_\alpha r_{N,t}(\bar\varepsilon)$. Thus the rate in
\eqref{p10:ml16:confidence-lower} is sharp even at this one clean law.
\end{theorem}

The distinction between adaptive point estimation and adaptive honest
confidence is established in robust statistics. Luo and Gao
\cite{p10:ml16:LG} obtain sharp adaptive confidence phenomena for Huber
contamination. Cho, Luo and Gao \cite{p10:ml16:CLG} study a single binomial
proportion under Huber contamination. The claim here concerns a
nonparametric binomial mixing law, a hard budget of individual bit
changes, the complete transport inverse curve, and a fixed clean
arcsine center. It is not a claim that the general obstruction to
robust confidence adaptation is new.

\begin{lemma}[Centered oscillatory alternatives]
\label{p10:ml16:centered}
Put $h_k(x)=T_k(2x-1)$ and $G_{k,a}=(1+ah_k)\pi$, where
$k\ge1$ and $0<a\le1/2$. There is an absolute $c>0$ such that
\[
 W_1(G_{k,a},\pi)\ge ca/k.
\]
Write $p_G$ for the $t$-trial predictive count law and $\rho=p_\pi$.
For $k\le t$,
\begin{align}
 \chi^2(p_{G_{k,a}},\rho)&=a^2\lambda_{t,k}/2,
 &\lambda_{t,k}&=\prod_{l=0}^{k-1}\frac{t-l}{t+1+l}
                                    \le e^{-k^2/(2t)},
                                      \label{p10:ml16:chi}\\
 W_{1,t}(p_{G_{k,a}},\rho)&\le
                   \frac{a}{k}e^{-k^2/(4t)}. &&
                                      \label{p10:ml16:attenuation}
\end{align}
Here $W_{1,t}$ uses count cost $|i-j|/t$. For $k>t$ the two predictive
laws are identical.
\end{lemma}
\begin{proof}
These are the centered forms of the oscillatory bounds in
\cite{comp:part:09}; we include their derivation. The change of
variable $x=(1+\cos\theta)/2$ gives
\[
 W_1(G_{k,a},\pi)=\frac{a}{2\pi k}
              \int_0^\pi|\sin(k\theta)|\sin\theta\,d\theta
                                      \asymp a/k.
\]
For a uniform lower bound, restrict to $[\pi/4,3\pi/4]$ and sum
complete periods of $|\sin(k\theta)|$; the remaining finitely many
small $k$ have positive integrals. The posterior operator
$Qh(j)=\mathbb E_\pi[h(P)\mid J=j]$ satisfies
$\|Qh_k\|_{L^2(\rho)}^2=\lambda_{t,k}/2$.
This classical beta--binomial diagonalization
\cite{p05:fm:DKS2008} also follows by observing that the self-adjoint
operator $Q^*Q$ preserves polynomial degree and has leading
multiplier $\lambda_{t,k}$ on degree $k$, and then using arcsine
orthogonality. Since $p_{G_{k,a}}/\rho=1+aQh_k$, this proves
\eqref{p10:ml16:chi}. Moreover
\[
 \log\lambda_{t,k}
 \le-\sum_{l=0}^{k-1}\frac{1+2l}{t+1+l}
 \le-\frac{k^2}{2t}.
\]
For completeness, the reversible birth-and-death generator with rates
$(j+1/2)(t-j)$ upward and $j(t-j+1/2)$ downward has eigenfunction
$Qh_k$ with eigenvalue $-k^2$. Its Dirichlet form is
\[
 \mathcal E(f,g)=\sum_{j=0}^{t-1}\rho_j(j+1/2)(t-j)
                    (f(j+1)-f(j))(g(j+1)-g(j)).
\]
For a function $f$ with normalized count Lipschitz constant one,
$\mathcal E(f,f)\le2$. Cauchy--Schwarz for this form implies
$|\langle f,Qh_k\rangle_\rho|\le\sqrt{\lambda_{t,k}}/k$.
Kantorovich duality proves \eqref{p10:ml16:attenuation}. If $k>t$,
each binomial probability is a polynomial of degree $t$ and is
orthogonal to $h_k$ under $\pi$, proving the last assertion.
\end{proof}

\begin{proof}[Proof of Theorem~\ref{p10:ml16:confidence}]
We use a two-experiment observation. Suppose an alternative $H$, after
an allowed attack, has data law within total variation $\tau$ of the
clean data law at $\pi$. Honesty at $H$ and at $\pi$ implies
\begin{equation}\label{p10:ml16:two-experiment}
 \mathbb P_{\pi,\mathrm{id}}\{\pi,H\in\mathcal C\}
                                    \ge1-2\alpha-\tau.
\end{equation}
Randomization can be appended to the data without increasing total
variation. On this event the diameter is at least $W_1(\pi,H)$.

First we produce clean alternatives for the three sampling terms.
Taking $k=t+1$ and $a=1/2$ gives identical predictive laws and
separation at least $c/t$. For the logarithmic sampling term, take
$k=\lceil\sqrt{16t\log(eN)}\rceil$ and $a=1/2$ if $k\le t$.
The product chi-squared identity gives
\[
 1+\chi^2(p_{G_{k,a}}^{\otimes N},\rho^{\otimes N})
          =(1+a^2\lambda_{t,k}/2)^N,
\]
so their total variation distance is less than $1/8$, with separation
at least $c/\sqrt{t\log(eN)}$. If $k>t$, this desired term is bounded
by $C/t$, and the preceding exact alternative suffices. Finally take
$k=1$ and $a=1/(8\sqrt N)$. The same product identity gives total
variation less than $1/8$ and separation at least $c/\sqrt N$.

For the corruption terms assume first that $q\ge1$ and put
$\eta=q/(16Nt)$. We claim that there is an alternative $H$ with
\begin{equation}\label{p10:ml16:centered-modulus}
 W_{1,t}(p_H,\rho)\le\eta,\qquad
 W_1(H,\pi)\ge c\max\left\{\frac1t,\eta,
                       \frac1{\sqrt{tL_t(\eta)}}\right\}.
\end{equation}
It is enough to construct an alternative separately for each term
and select the largest. The $1/t$ term again uses $k=t+1$.
For the linear term take $H=(1-\eta)\pi+\eta\delta_1$.
Its predictive distance is at most $\eta$, and its population distance
is exactly $\eta/2$. For the remaining term set $x=\eta\sqrt t$.
If $x\ge1$ the linear term already dominates it. Otherwise take
$k=\lceil\sqrt{16tL_t(\eta)}\rceil$ and $a=1/2$.
If $k\le t$, \eqref{p10:ml16:attenuation} implies
\[
 W_{1,t}(p_{G_{k,a}},\rho)
 \le\frac{(e+1/x)^{-4}}{8\sqrt{tL_t(\eta)}}\le\eta,
\]
and the latent separation has the desired order. If $k>t$, the
$1/t$ alternative dominates this term up to an absolute constant.
This proves \eqref{p10:ml16:centered-modulus}.

Couple independently each count from $p_H$ to a target count from
$\rho$ by an optimal transportation coupling. The expected sum of
absolute count changes is at most $Nt\eta=q/16$. Given the clean
counts from $H$, an adversary samples the conditional target counts.
It implements the changes when their total is at most $q$ and otherwise
leaves the array unchanged. Every count change is implemented by the
same number of bit flips. To cover procedures that observe the full
array, add the required ones uniformly among the zero positions, or
remove them uniformly among the one positions. Conditional on its
count, a clean row is uniform among strings with that count; this
symmetric operation produces a uniform row with the target count.
Thus the untruncated coupling produces exactly the clean full-array
law at $\pi$, as well as its count law. The product chi-squared
comparisons above likewise apply to full arrays, since conditional
row arrangements carry no additional information about $G$.
Markov's inequality shows that the resulting
output law differs from $\rho^{\otimes N}$ in total variation by at
most $1/16$. Hence \eqref{p10:ml16:two-experiment} applies.

When $q\ge1$, $q/(Nt)\ge\bar\varepsilon/2$, so $\eta$ and
$\bar\varepsilon$ differ by at most a fixed factor. In particular
$L_t(\eta)\le L_t(\bar\varepsilon)+\log32
\le C L_t(\bar\varepsilon)$.
When $q=0$, $\bar\varepsilon<1/(Nt)$, so the linear corruption term
is bounded by $N^{-1/2}$ and
$L_t(\bar\varepsilon)\ge c\log(eN)$; the other corruption term is
then absorbed by the clean logarithmic term. Thus in every case one
of the alternatives constructed above has separation at least
$c r_{N,t}(\bar\varepsilon)$ and data total variation at most $1/8$.
Equation~\eqref{p10:ml16:two-experiment} proves the lower bound.

For the upper bound, the earlier hard-bit-budget theorem
\cite{comp:part:09} supplies an estimator with uniform expected
loss at most $C r_{N,t}(\bar\varepsilon)$. The closed Wasserstein
ball of radius $C r_{N,t}(\bar\varepsilon)/\alpha$ about that
estimator is honest by Markov's inequality, and has diameter at most
twice this radius, or one if smaller. This proves the fixed-confidence
matching assertion. No optimal dependence on $\alpha$ is claimed.
\end{proof}

For example, let $N=t^3$ and $\bar\varepsilon=t^{-1/2}$, with
$t\to\infty$. The clean point-estimation rate is of order
$(t\log t)^{-1/2}$, whereas every confidence set honest for this
vanishing upper corruption budget has diameter at least $ct^{-1/2}$
with probability bounded away from zero at the clean law $\pi$.
The ratio diverges like $\sqrt{\log t}$. This is an obstruction to
honest adaptation to the actual budget zero, even at a fixed population;
it is not an impossibility theorem for adaptive point estimators.

\subsection{A remaining allocation problem under bit corruption}

For clarity, the clean adaptive-budget theorem does not close its
malicious-bit analogue. There is nevertheless a rigorous upper bound
for a predetermined design. If an adversary may change at most
$\lfloor\varepsilon B\rfloor$ observed bits after sampling, the
multilevel construction above, with its maximum level chosen using
$\varepsilon$, has worst-case expected loss
\begin{equation}\label{p10:ml16:corrupt-allocation-upper}
                 C\min\{1,B^{-1/3}+\varepsilon^{2/3}\}.
\end{equation}
To see this, a flip in a level-$l$ observation changes the integrated
absolute value of its full-average indicator by $1/t_l$, and that of
its half-average indicator by at most $2/t_l$. Thus it changes
$\widetilde F$ in $L^1$ by at most $3/(N_lt_l)$. With the allocation
above this is bounded by $CT^{1/4}/B$ uniformly over levels. Projection
and the triangle inequality therefore give
\[
 \mathbb EW_1(\widehat G,G)
 \le C\{T^{-1/2}+(B^{-1/2}+\varepsilon)T^{1/4}\}.
\]
Choosing a dyadic $T$ of order
$(B^{-1/2}+\varepsilon)^{-4/3}$ proves
\eqref{p10:ml16:corrupt-allocation-upper}; when this order is bounded,
a fixed estimate suffices. The rounding argument is valid for large
$B$, since this $T$ is at most a constant times $B^{2/3}$.

This bound is not claimed to be minimax, and we do not count it as a
third closed target. The exact corrupted adaptive-allocation law and
adaptation of point estimates to an unknown hard bit budget remain
open in this manuscript. Neither is resolved by
Theorem~\ref{p10:ml16:confidence}.


% ===== arithmetic =====

\section{Unknown frequency supports and the resonance barrier to observation}
\label{p10:ar16:sec}

The preceding inverse theorem prescribed the frequency support.  The
first result below removes that prescription for finite sums whose
frequencies are internally rationally independent.  A second result
shows why exact identification does not supply a uniform finite-height
recovery guarantee: a nearly rational relation can hide a fixed change
of coefficients and of the true right boundary.  We compute both the
finite-window ambiguity and the first height at which a shrinking
right-edge strip becomes visible.

These are two self-posed inverse-observation targets.  The qualitative
dependence of spectral boundaries on rational relations is classical
\cite{p09:ar15:AvellarHale,p04:nxz:Sepulcre}; it is not counted again here.
Nor is the first theorem a solution of the general common-zero problem
for exponential polynomials.  Its much stronger observation hypothesis
and its positive real, internally independent frequencies permit a
direct proof without a common-factor conjecture.

\subsection{One-sided zeros identify an unknown independent support}

Consider
\begin{equation}\label{p10:ar16:eq:two-models}
 F(s)=1-\sum_{j=1}^N a_j e^{-\lambda_js+i\theta_j},\qquad
 G(s)=1-\sum_{k=1}^{M} d_k e^{-\mu_ks+i\psi_k},
\end{equation}
where all magnitudes are strictly positive and
$0<\lambda_1<\cdots<\lambda_N$, $0<\mu_1<\cdots<\mu_M$.
Neither the two supports nor their sizes are assumed known or equal.
Let $S_F,S_G$ be the unique real solutions of
\[
 \sum_j a_j e^{-\lambda_jS_F}=1,
 \qquad \sum_k d_k e^{-\mu_kS_G}=1.
\]

\begin{theorem}[Unknown-support rigidity from a right-hand zero set]
\label{p10:ar16:thm:unknown}
Suppose each of the two frequency lists in
\eqref{p10:ar16:eq:two-models} is linearly independent over $\mathbb Q$.
Suppose there is a real $\sigma_0<\min(S_F,S_G)$ such that
\begin{equation}\label{p10:ar16:eq:samezeros}
 \{s:\Re s>\sigma_0,\ F(s)=0\}
   =\{s:\Re s>\sigma_0,\ G(s)=0\}.
\end{equation}
Equality of sets suffices; multiplicities need not be supplied.
Then $N=M$ and, after an ordering of the terms,
\[
 \lambda_j=\mu_j,\qquad a_j=d_j,\qquad
 \theta_j=\psi_j\pmod{2\pi}.
\]
Thus $F=G$ identically.  In particular the observation determines
the number of terms, the positive frequencies, and all complex
coefficients within this model class.
\end{theorem}

The normalization by the constant term and the positivity of all
frequencies are substantive.  Multiplication by a nonzero exponential
preserves every zero, so a normalization fixing the lowest frequency
is necessary.  The assumption that the observed half-plane contains
both right boundaries prevents an empty-record ambiguity.

\begin{lemma}[The joint hull is forced to be an integral graph]
\label{p10:ar16:lem:graph}
Under the hypotheses of Theorem~\ref{p10:ar16:thm:unknown}, one has
$S_F=S_G=:S$.  Moreover $N=M$, and there is a matrix
$K\in\operatorname{GL}_N(\mathbb Z)$ such that
\begin{equation}\label{p10:ar16:eq:graph}
 \mu=K\lambda,\qquad \psi=K\theta\pmod{2\pi}.
\end{equation}
For every $x\in\mathbb T^N$, the phase-frozen pair
$(F_x,G_{Kx})$ has the same zero set in $\Re s>\sigma_0$.
Here $F_x$ replaces $\theta$ by $x$, and similarly for $G$.
\end{lemma}
\begin{proof}
The triangle inequality excludes zeros to the right of $S_F$.
Rational independence makes $t\lambda$ dense modulo $2\pi$,
so vertical translates of $F$ approach $F_0$.  Its zero at $S_F$
is simple, with derivative
$\sum a_j\lambda_j e^{-\lambda_jS_F}>0$.  Rouch\'e's theorem
therefore places real parts of zeros of $F$ arbitrarily close to
$S_F$.  The same applies to $G$.  The observed zero sets thus
have right endpoints $S_F$ and $S_G$, proving their equality.
Put
\[
 b_j=a_j e^{-\lambda_jS},\qquad
 c_k=d_k e^{-\mu_kS},\qquad \sum b_j=\sum c_k=1.
\]

Let $\mathcal H$ be the closure in $\mathbb T^{N+M}$ of the
one-parameter subgroup $t(\lambda,\mu)$, and let
$\mathcal C=(\theta,\psi)+\mathcal H$.  Every pair of functions
with phases in $\mathcal C$ is a locally uniform limit of common
vertical translates of $(F,G)$.  Such a limit pair still has the
same zero set in the observed half-plane.  Indeed, a zero of one
limit function is approximated by zeros of its approximants inside
an arbitrarily small disk contained in the half-plane; those are
zeros of the other approximants by \eqref{p10:ar16:eq:samezeros}.
Passing to the limit makes the disk center a zero of the other
limit function.  Reverse the roles for the converse.  No limit
function is identically zero, since its constant term dominates as
$\Re s\to+\infty$.

The first projection of $\mathcal C$ is all of $\mathbb T^N$.
Consequently it contains some $(0,y)$.  The corresponding first
function vanishes at $S$, so the second one does also.  Hence
$\sum_k c_k e^{iy_k}=1$.  Equality in the triangle inequality,
with all $c_k>0$ and their sum one, forces $y=0$.  In particular
$(0,0)\in\mathcal C$, so $\mathcal C=\mathcal H$.
The same argument shows that the kernel of the first projection
$\mathcal H\to\mathbb T^N$ is trivial.  The second projection
is likewise surjective and has trivial kernel.

Both projections are therefore continuous isomorphisms of compact
groups.  The resulting isomorphism
$\mathbb T^N\to\mathbb T^M$ is given by an integer matrix:
each of its component characters is an integer character of
$\mathbb T^N$.  Applying the same argument to the inverse gives
integer matrices inverse to each other.  Thus $N=M$, the matrix
$K$ is unimodular, and $\mathcal H$ is its graph.  Differentiating
the identity $t\mu=K(t\lambda)\bmod2\pi$ at $t=0$ yields
$\mu=K\lambda$.  Membership of $(\theta,\psi)$ in the graph
gives the phase identity.  The limit-pair argument above proves
the last assertion.
\end{proof}

\begin{proof}[Proof of Theorem~\ref{p10:ar16:thm:unknown}]
Use the notation and matrix supplied by
Lemma~\ref{p10:ar16:lem:graph}.  In the Laurent polynomial ring
$\mathcal R=\mathbb C[z_1^{\pm1},\ldots,z_N^{\pm1}]$ set
\begin{equation}\label{p10:ar16:eq:Laurent}
 P(z)=1-\sum_j b_jz_j,\qquad
 Q(z)=1-\sum_k c_k z^{K_k},
\end{equation}
where $K_k$ is the $k$-th row of $K$.  We first show that $P$
divides $Q$ in $\mathcal R$.

For $z$ complex and near $\mathbf1$, the equation
\[
 \mathcal F(u,z)=1-\sum_j b_jz_j e^{-\lambda_ju}=0
\]
has a unique small holomorphic solution $u=\rho(z)$, since
$\partial_u\mathcal F(0,\mathbf1)=\sum b_j\lambda_j>0$.
For $z_j=e^{ix_j}$ and small real $x$, the corresponding root
$S+\rho(z)$ lies in the observed half-plane.  The graph lemma
gives
\[
 1-\sum_k c_k z^{K_k} e^{-\mu_k\rho(z)}=0.
\]
The expression on the left is holomorphic near $\mathbf1$ and
vanishes on an open part of the product of unit circles.  Successive
one-variable identity theorems, or the substitution $z=e^{iw}$,
show that it vanishes identically in a complex neighborhood.
If $P(z)=0$ in that neighborhood, then the uniqueness of the
implicit root gives $\rho(z)=0$.  Therefore $Q(z)=0$ on an open
part of the affine hyperplane $P(z)=0$.

Multiply $Q$ by a Laurent monomial to clear all negative exponents.
Substitute
$z_N=(1-\sum_{j<N}b_jz_j)/b_N$ in the resulting polynomial.
It vanishes on an open set and hence identically.  Polynomial
division by the linear polynomial $P$ proves divisibility.  The
clearing monomial is a unit in $\mathcal R$, so $P\mid Q$ in
that ring.  This argument includes $N=1$, when the hyperplane is
a single point and ordinary division by a linear factor applies.

The polynomial $P$ is irreducible in $\mathcal R$: it is linear
and irreducible in the ordinary polynomial ring, and it divides
no monomial, so localization by the coordinate variables preserves
its primality.  Likewise $1-\sum c_kz_k$ is irreducible.  The
unimodular substitution $z_k\mapsto z^{K_k}$ is an automorphism
of $\mathcal R$, so $Q$ is irreducible as well.  Consequently
\begin{equation}\label{p10:ar16:eq:associate}
                         Q(z)=\gamma z^vP(z)
\end{equation}
for some $\gamma\ne0$ and $v\in\mathbb Z^N$; the units of a
Laurent polynomial ring are precisely the nonzero monomials.

Evaluate exponent vectors by the real linear functional
$k\mapsto k\cdot\lambda$.  The smallest frequency in the
support of $Q$ is $0$, because the others are the strictly
positive numbers $\mu_k$.  The smallest frequency in the
right side of \eqref{p10:ar16:eq:associate} is $v\cdot\lambda$,
because those of $P$ are $0,\lambda_1,\ldots,\lambda_N$.
It follows that $v\cdot\lambda=0$.  Rational independence
gives $v=0$, and the constant coefficient gives $\gamma=1$.
Thus the rows of $K$ are the standard basis vectors in some
order, and the corresponding $c_k$ and $b_j$ coincide.
Equation~\eqref{p10:ar16:eq:graph} now identifies frequencies and
phases; normalization at the common $S$ identifies the original
magnitudes.  This proves the theorem.
\end{proof}

\begin{corollary}[Zero-density exact discrepancies cannot hide a support]
\label{p10:ar16:cor:sparse}
Keep the model and independence hypotheses of
Theorem~\ref{p10:ar16:thm:unknown}, with
$\sigma_0<\min(S_F,S_G)$.  Let $D(T)$ count the distinct points
in the symmetric difference of their zero sets satisfying
$\Re s>\sigma_0$ and $|\Im s|\le T$.  The conclusion of
the theorem still holds if exact equality is replaced by
$D(T)=o(T)$.  Finite changes or exact discrepancies of zero
vertical density therefore cannot conceal a different support.
\end{corollary}
\begin{proof}
We prove that $D(T)=o(T)$ already forces equality of the two
zero sets throughout the observed half-plane.  This step does
not require rational independence.  Suppose, to the contrary,
that $z_0$ is a zero of $F$ and not of $G$ there.  Choose a
closed disk of radius $r>0$ about $z_0$ contained in the
half-plane, with no zero of $G$ in the disk and no zero of $F$
on its boundary.  For every phase pair in a sufficiently small
relative open neighborhood $U$ of $(\theta,\psi)$ in its joint
translation hull, Rouch\'e's theorem preserves an $F$ zero in
the disk and the absence of any $G$ zero there.

Translation on the joint hull is uniquely ergodic for its Haar
measure: its characters have nonzero time frequencies except
for the characters trivial on that hull.  A continuous nonnegative
function supported in $U$ with positive Haar integral therefore
shows that the set $E_T$ of times $t\in[0,T]$ whose translated
phase lies in $U$ has $|E_T|\ge cT$ for all large $T$.
For each such time there is an unmatched zero of $F$ in the
translated disk centered at $z_0+it$.  A fixed unmatched zero
can account for a set of such times of length at most $2r$.
Consequently
\[
 |E_T|\le2rD(T+C)=o(T),
\]
where $C$ is fixed.  These counts are finite because the zeros
of each finite sum are locally finite and confined to a fixed
vertical strip.
This is a contradiction.  Interchanging $F$ and $G$ excludes
the opposite kind of discrepancy.  Exact equality now follows,
and Theorem~\ref{p10:ar16:thm:unknown} applies.
\end{proof}

\begin{theorem}[A finite local inverse at each independent support]
\label{p10:ar16:thm:local}
Fix one model $F$ from Theorem~\ref{p10:ar16:thm:unknown}, its
term count $N$, and a half-plane $\Re s>\sigma_0$ with
$\sigma_0<S_F$.  Parametrize nearby models by the real vector
\[
 p=(\lambda_1,\ldots,\lambda_N,
             \Re\alpha_1,\Im\alpha_1,\ldots,
             \Re\alpha_N,\Im\alpha_N),
 \qquad \alpha_j=a_je^{i\theta_j}.
\]
One can select at most $3N$ distinct simple zeros of $F$ in that
half-plane and pairwise disjoint disks about them such that their
locally matched zero coordinates determine $p$ real-analytically
on a neighborhood of the base model.  If $\Psi(p)$ is the vector
of these matched complex zero coordinates, this neighborhood can
be chosen so that
\begin{equation}\label{p10:ar16:eq:local-Lipschitz}
                    \|p-q\|\le C_F\|\Psi(p)-\Psi(q)\|.
\end{equation}
The nearby candidate models have the same term count and ordered
positive frequencies; they need not themselves be rationally
independent.  The selected disks, neighborhood, and constant
$C_F$ depend on the base model.  No uniform conditioning over
all independent supports is asserted.
\end{theorem}
\begin{proof}
We first show that the derivatives of all observed simple-zero
coordinates separate the tangent space.  Write a tangent vector
as $(\nu_j,\beta_j)$ with $\nu_j\in\mathbb R$ and
$\beta_j\in\mathbb C$.  The corresponding function derivative is
\[
 H(s)=-\sum_j(\beta_j-\alpha_j\nu_js)e^{-\lambda_js}.
\]
The derivative of a simple zero $\rho$ is
$-H(\rho)/F'(\rho)$.  Suppose it vanishes at every simple
zero in the observed half-plane.  Thus $H$ vanishes at all
those zeros.

For every real phase $x$ sufficiently close to zero, the frozen
function $F_x$ has a unique simple root $\rho(x)$ near $S_F$,
still strictly inside the observed half-plane.  Rational
independence supplies a sequence $t_n\to+\infty$ with
$\theta-t_n\lambda\to x$ modulo $2\pi$.  Arbitrarily large
times can be used in every phase neighborhood because the
translation flow is recurrent and dense.  Locally uniform
convergence and Rouch\'e's theorem give unique simple roots
$\rho_n\to\rho(x)$ of $F(s+it_n)$ in a small disk about
$\rho(x)$.  In particular $\rho_n+it_n$ is an observed
simple zero of the base function, so
$H(\rho_n+it_n)=0$.

On every fixed compact set of $s$, division by $t_n$ gives
the locally uniform limit
\[
 \frac{H(s+it_n)}{t_n}
       \longrightarrow i\sum_j a_j\nu_j e^{ix_j}e^{-\lambda_js}.
\]
Thus this last sum vanishes at $s=\rho(x)$ for every small
real $x$.  Normalize at $S_F$, putting
$b_j=a_je^{-\lambda_jS_F}$.  The same implicit-root
complexification used in the preceding theorem shows that the
linear polynomial $\sum b_j\nu_jz_j$ vanishes on an open
part of the hyperplane $1-\sum b_jz_j=0$.  It is therefore a
scalar multiple of that affine linear polynomial.  Its constant
term is zero, so that scalar and all its coefficients vanish:
$\nu_j=0$ for every $j$.

With the frequency part removed, repeat the same limiting
argument without dividing $H(s+it_n)$ by $t_n$.  Its limit is
$-\sum\beta_je^{-i\theta_j}e^{ix_j}e^{-\lambda_js}$.
The resulting linear polynomial on the same hyperplane has
coefficients $\beta_je^{-i\theta_j}e^{-\lambda_jS_F}$ and
zero constant term.  The identical argument gives $\beta_j=0$.
This proves separation of the full $3N$-dimensional real
tangent space.

Starting with that tangent space, select a simple-zero
derivative which is nonzero on the current common kernel, and
intersect with its kernel.  Each selection lowers the dimension
by at least one.  After at most $3N$ distinct zeros, the joint
derivative has trivial kernel.  Choose disjoint small disks
around the selected zeros.  Rouch\'e's theorem and the implicit
function theorem define their unique matched roots for all
parameters in a common neighborhood, and these roots depend
real-analytically on the parameters.

The Jacobian $A=D\Psi(p_0)$ at the base point is injective.
Some $3N$ real coordinate rows form an invertible square
minor, so the inverse function theorem gives a real-analytic
local inverse using those coordinates.  If $s_*>0$ is the
smallest singular value of $A$, shrink to a convex parameter
ball on which $\|D\Psi-A\|\le s_*/2$.  Integration along
a segment gives
\[
 \|\Psi(p)-\Psi(q)\|
 \ge\|A(p-q)\|-(s_*/2)\|p-q\|
 \ge(s_*/2)\|p-q\|.
\]
This is \eqref{p10:ar16:eq:local-Lipschitz} with $C_F=2/s_*$.
It also explains why neither a prescribed finite height nor
a uniform constant follows from this local theorem.
\end{proof}

The global inverse observes complex zero locations throughout a
half-plane, allowing the zero-density exception just proved.  The
finite local inverse instead fixes the term count and matches zeros
inside specified disks.  A right endpoint or the distribution of
real parts alone can fail to identify coefficients
\cite{p04:nxz:Sepulcre}.  The general factorization and common-zero
questions are discussed in \cite{p10:ar16:FornasieroTerzo}; in
particular its general irreducible common-zero statement carries
the Schanuel-conjecture assumption.  No such assumption or result
is used in the proof above.  Historical priority for this precise
restricted inverse statement remains to be assessed beyond the
checked primary sources.

\subsection{A sharp resonance escape law and finite-record ambiguity}

Fix $c\in I=[1/4,1/2]$ and a small positive detuning $d$.  Define
\begin{equation}\label{p10:ar16:eq:detuned}
 g_{c,d}(s)=1-(2+c)e^{-s}+2c e^{-(2+d)s}.
\end{equation}
For irrational $d$, the frequencies $(1,2+d)$ are rationally
independent.  Their separation is at least one, and all nonconstant
coefficient magnitudes lie in the fixed interval $[1/2,5/2]$.
Let $S_{c,d}$ solve
\begin{equation}\label{p10:ar16:eq:detuned-edge}
 (2+c)e^{-S_{c,d}}+2c e^{-(2+d)S_{c,d}}=1.
\end{equation}
For irrational $d$, this is the actual right zero boundary.
Put $u_c=(2+c)e^{-S_{c,0}}\in(0,1)$ and
\begin{equation}\label{p10:ar16:eq:kappa}
 \kappa_c=\left[\frac{2(2-u_c)^3}{u_c(1-u_c)}\right]^{1/2}.
\end{equation}

\begin{theorem}[Resonance escape at a precisely determined height]
\label{p10:ar16:thm:escape}
Let $T_{c,d}$ be the infimum of positive ordinates of zeros of
$g_{c,d}$ in the strict strip
$S_{c,d}-d<\Re s<S_{c,d}$.  Uniformly for $c\in I$, as
$d\downarrow0$ through irrational values,
\begin{equation}\label{p10:ar16:eq:escape}
              dT_{c,d}=\pi-\kappa_c\sqrt d+O(d).
\end{equation}
In particular $T_{c,d}\sim\pi/d$.  This is a first-arrival law,
not an average-density assertion or an effective ergodic theorem.
\end{theorem}

\begin{proof}
Write $S=S_{c,d}$ and
\[
 u=(2+c)e^{-S},\qquad 1-u=2c e^{-(2+d)S},\qquad
 M_d=u+(2+d)(1-u).
\]
All these quantities vary smoothly on a compact neighborhood of
$I\times\{0\}$; $u$ and $1-u$ are uniformly positive.  At a
zero $s=S-x+it$ in the strip, define its small phases, whenever
they exist, by
\[
 \eta_1=-t\pmod{2\pi},\qquad
 \eta_2=\pi-(2+d)t\pmod{2\pi}.
\]
The exact triangle-defect identity is
\[
 u e^x|e^{i\eta_1}-1|^2
 +(1-u)e^{(2+d)x}|e^{i\eta_2}-1|^2
       =2\{u e^x+(1-u)e^{(2+d)x}-1\}.
\]
For $0<x<d$ it forces the unique small representatives
$\eta_1,\eta_2=O(\sqrt d)$, uniformly.  Consequently their
relation coordinate
\begin{equation}\label{p10:ar16:eq:drifting-relation}
              \xi=2\eta_1-\eta_2
                    =dt-\pi\pmod{2\pi}
\end{equation}
is also $O(\sqrt d)$.

Parametrize all small zero phases by setting
$\eta_1=A$, $\eta_2=2A-\xi$ and solving
\begin{equation}\label{p10:ar16:eq:slow-chart}
 1-u e^{x+iA}-(1-u)e^{(2+d)x+i(2A-\xi)}=0
\end{equation}
for $(x,A)$.  The real and imaginary derivatives in those two
variables at zero are $-M_d$ and $-(2-u)$, respectively.  The
implicit-function theorem gives analytic functions $x_d(\xi)$
and $A_d(\xi)$ on a fixed neighborhood, uniformly in $c,d$.
Conjugation makes $x_d$ even and $A_d$ odd.  Expanding gives
\begin{align}
 A_d(\xi)&=\frac{1-u}{2-u}\xi+O(\xi^3),
                                              \label{p10:ar16:eq:Aexp}\\
 x_d(\xi)&=C_d\xi^2+O(\xi^4),\qquad
 C_d=\frac{u(1-u)}{2M_d(2-u)^2}>0.
                                              \label{p10:ar16:eq:xexp}
\end{align}
For example, the quadratic numerator is
$u[(1-u)/(2-u)]^2+(1-u)[u/(2-u)]^2$; division by $2M_d$
gives the displayed coefficient.  All remainders are uniform.

For sufficiently small positive $d$, the condition
$0<x_d(\xi)<d$ is therefore equivalent in this chart to
\[
                0<|\xi|<r_d,\qquad
 r_d=\sqrt{d/C_d}\,(1+O(d))
                    =\kappa_c\sqrt d+O(d^{3/2}).
\]
The second equality uses smooth dependence on $d$ and
$C_0=u_c(1-u_c)/(2(2-u_c)^3)$.  The defect identity proves
that this chart includes every zero in the strip, not just a
selected local branch.

The first possible positive-height group of such zeros has
$dt=\pi+\xi$, with $|\xi|<r_d$.  The phase equation
$\eta_1=A_d(\xi)=-t\bmod2\pi$ is equivalent to
\begin{equation}\label{p10:ar16:eq:lattice-hit}
            \xi+dA_d(\xi)=2\pi n d-\pi,\qquad n\in\mathbb Z.
\end{equation}
The function on the left is strictly increasing, with derivative
$1+O(d)$.  As $n$ increases, the right side advances in steps
of exactly $2\pi d$.  Since $r_d\asymp\sqrt d\gg d$, there
is an admissible first lattice point in the interval corresponding
to $-r_d<\xi<r_d$, and its solution satisfies
\[
                       \xi=-r_d+O(d).
\]
If a lattice point lands exactly on the excluded boundary, use
the next one, which changes this bound by at most $O(d)$.
For small $d$ this first point has negative $\xi$, so the excluded
value $\xi=0$ cannot interfere.  Equation
\eqref{p10:ar16:eq:drifting-relation} excludes all smaller positive
heights outside this first group; its previous possible group is
centered at $-\pi/d$.  Every solution of
\eqref{p10:ar16:eq:lattice-hit} gives an actual zero, since both
phase equations and \eqref{p10:ar16:eq:slow-chart} then hold.
Thus $dT_{c,d}=\pi-r_d+O(d)$, proving
\eqref{p10:ar16:eq:escape}.
\end{proof}

The same detuned family yields a quantitative obstruction to noisy
inversion even when its two frequencies are supplied to the observer.
Let
\[
 \rho_n=\log2+2\pi in,\qquad
 \mathcal C_n=\{s:\Re s>0, |\Im s-2\pi n|<\pi\}.
\]

\begin{proposition}[Uniform ambiguity of finite noisy zero records]
\label{p10:ar16:prop:finite}
There are absolute constants $k,K>0$ such that, for every
$c\in I$ and integer $L\ge0$, if $d(1+L)\le k$, then
$g_{c,d}$ has exactly one zero $\rho_{n,c,d}$ in each
$\mathcal C_n$, $|n|\le L$, and no other zeros in their union.
The zeros are simple and satisfy
\begin{equation}\label{p10:ar16:eq:finite-displacement}
 |\rho_{n,c,d}-\rho_n|\le Kd(1+|n|).
\end{equation}
Moreover, uniformly in that regime,
\begin{equation}\label{p10:ar16:eq:finite-expansion}
 \rho_{n,c,d}
 =\rho_n+\frac{c}{2-c}\,d\rho_n
                      +O(d^2(1+|n|)^2).
\end{equation}
For $c_-=1/4$ and $c_+=1/2$, the two models have a fixed positive
coefficient separation and a fixed positive separation of their
right boundaries as $d\downarrow0$, yet their records in these
cells admit a pairing at distance at most $2Kd(1+L)$.
\end{proposition}
\begin{proof}
At $d=0$ the exact factorization is
\begin{equation}\label{p10:ar16:eq:factor}
 g_{c,0}(s)=(1-2e^{-s})(1-ce^{-s}).
\end{equation}
The second factor has all its zeros in $\Re s=\log c<0$.
In $\Re s\ge0$ its modulus is at least $1-c\ge1/2$.
On $\Re s=0$, the first factor has modulus at least one;
on the horizontal sides of a cell it is $1+2e^{-\Re s}$.
Choose a fixed large right boundary $R$ on which the constant
term dominates every $g_{c,d}$ for all small $d\ge0$.
The modulus of $g_{c,0}$ is bounded below by a fixed positive
constant on all four sides of these rectangles.

For $s=\sigma+it$, $\sigma\ge0$, the elementary integral
bound $|e^{-ds}-1|\le d|s|$ gives
\begin{equation}\label{p10:ar16:eq:function-perturbation}
 |g_{c,d}(s)-g_{c,0}(s)|
 \le2c d e^{-2\sigma}(\sigma+|t|)
 \le Kd(1+|t|).
\end{equation}
If $k$ is sufficiently small, Rouch\'e's theorem on each
rectangle proves that it contains exactly one zero, counted with
multiplicity.  Constant-term domination excludes zeros to its
right.  This proves the exact count and simplicity, and also
excludes zeros on the horizontal or left boundaries.

On a fixed disk of radius at most $1/4$ around $\rho_n$,
the factorization \eqref{p10:ar16:eq:factor} gives
$|g_{c,0}(s)|\ge k_0|s-\rho_n|$, uniformly in $c,n$.
Applying Rouch\'e on the circle of radius
$K_0d(1+|n|)$, with $K_0$ large and then $k$ small, proves
\eqref{p10:ar16:eq:finite-displacement}.

For the expansion put $s=\rho_n+z$.  Periodicity of the
integer-frequency factors makes
\[
 g_{c,d}(\rho_n+z)
 =1-\frac{2+c}{2}e^{-z}
       +\frac c2 e^{-2z}e^{-d\rho_n-dz}.
\]
Use $z=O(d(1+|n|))$ in its Taylor expansion.  Its linear terms
are $(1-c/2)z-(c/2)d\rho_n$, and the remaining terms are
$O(d^2(1+|n|)^2)$.  Division by the uniformly positive
$1-c/2$ proves \eqref{p10:ar16:eq:finite-expansion}.

The magnitudes $(2+c,2c)$ differ by a fixed amount for $c_-$
and $c_+$.  At $d=0$, the left side of
\eqref{p10:ar16:eq:detuned-edge}, for fixed $S$, is strictly
increasing in $c$ and strictly decreasing in $S$.  Thus
$S_{c_+,0}>S_{c_-,0}$.  Smooth dependence on $d$ preserves
half that positive separation for all sufficiently small $d$.
Finally pair zeros with the same index and apply
\eqref{p10:ar16:eq:finite-displacement} twice.
\end{proof}

\begin{corollary}[No uniform consistency without a resonance margin]
\label{p10:ar16:cor:minimax}
Suppose an observer receives all zeros in the cells
$|n|\le L$, each with arbitrary complex location error at most
$\epsilon$, and is given the exact frequencies $(1,2+d)$,
the atom count two, and the fixed phases $(0,\pi)$.  There is a
constant $\gamma>0$ such that, whenever
$Kd(1+L)\le\epsilon$ and $d$ is sufficiently small and
irrational, some common reported record is consistent with both
$c_-$ and $c_+$.  Every estimator therefore makes coefficient
error at least $1/8$, or right-boundary error at least $\gamma$,
on one of these two models, when estimating the respective target.
For any finite observation schedule $L=L(\epsilon)$, uniform
consistency over this internally independent model class is
impossible as $\epsilon\downarrow0$.
\end{corollary}
\begin{proof}
Report the midpoint of each pair
$(\rho_{n,c_-,d},\rho_{n,c_+,d})$.  The preceding pairing
bound puts this record within the required error of either
truth, without any missing or inserted zero.  The two first
coefficient magnitudes differ by $1/4$, so any single estimate
has error at least $1/8$ for one truth.  Their right boundaries
have a uniform positive separation, giving the analogous
$\gamma$ lower bound.  For each positive $\epsilon$ and finite
$L(\epsilon)$ one can choose an irrational positive
$d<\epsilon/[K(1+L(\epsilon))]$ satisfying all the smallness
conditions.  This proves the uniform inconsistency assertion.
\end{proof}

There is no contradiction with Theorem~\ref{p10:ar16:thm:unknown}:
the finite records are close, not equal, while the infinite exact
records eventually reveal the distinct models.  Nor does the
corollary rule out exact finite-data identification in infinite
precision.  The obstruction concerns noise uniformly over supports
arbitrarily close to an integer relation.  Ordinary pairwise
frequency separation does not control this effect.  The detuning
$d=\lambda_2-2\lambda_1$ is the relevant small quantity.

\subsection{Dependencies, scope, and unresolved problems}

The unknown-support theorem extends the earlier prescribed-support
inverse result by introducing the joint-hull graph argument and the
Laurent-hyperplane rigidity step.  Its proof repeats the elementary
frozen-edge argument only as a dependency, not as a new target.
The zero-density corollary and the finite local inverse are dependent
extensions of that first target, not extra target counts.  The local
inverse is compatible with the nonuniform finite-noise obstruction:
its conditioning and selected observation heights may deteriorate
as the support approaches a rational relation.
The escape law and finite-noise obstruction form a single second
target: they quantify the observation cost of approaching a rational
relation.  The exact factorization at zero detuning was already an
elementary obstruction in the preceding manuscript; the first-arrival
asymptotic, its explicit square-root correction, and the uniform
noisy-record conclusion are the new assertions here.

Avellar--Hale and the discussion in Sepulcre
\cite{p09:ar15:AvellarHale,p04:nxz:Sepulcre} already establish that small
frequency changes can radically change the spectral boundary.
Our additional claims concern how long the change can remain
unobserved and what finite noisy zero data can distinguish.
These are proved for the stated two-frequency family, not for
all near-resonant exponential sums.  A general quantitative inverse
theorem with unknown independent support, a sharp dependence on
a suitable Diophantine margin, and extension to infinite Dirichlet
supports remain open in this work.  No assertion about the zeros
of the Riemann zeta function follows from these finite-sum results.


% ===== gaussian =====

\section{Without scale moments, terminal sampling has a different boundary}
\label{p10:ga16:sec:momentfree}

The preceding manuscript proved that every infinite collection of
separate sum laws identifies a continuously mixed symmetric stable
index and scale, provided the scale has a positive and a negative
moment. We now remove that assumption and obtain a different exact
criterion. We then change the observation experiment: two jointly
observed replicates recover the entire mixing law without any scale
moment. The latter result extends to correlated location, scale,
index, and skewness with six replicates.

These are two internal research targets: the necessity of the scale
moment assumption, and finite replication with continuously varying
stable parameters. They are not named conjectures in the external
literature. Classical exponential completeness, positive stable
subordination, and transform uniqueness are acknowledged inputs.
Historical priority of the specific mixture statements is not certified.

Fix $0<\alpha_0<\alpha_1\leq2$. In the first two parts, conditional
on $(\alpha,B)\in[\alpha_0,\alpha_1]\times(0,\infty)$, the variables
$X_j$ are independent with characteristic function
\begin{equation}\label{p10:ga16:eq:symmetric}
 \mathbb E(e^{itX_j}\mid\alpha,B)=e^{-B|t|^\alpha}.
\end{equation}
The law of $(\alpha,B)$ is otherwise unrestricted. Put
$S_n=X_1+\cdots+X_n$. Separate laws of selected $S_n$ do not include
the joint laws of observations from a common latent draw.

For a locally finite real set $\Lambda$, let
\[
 \mathcal R(\Lambda)=\sup\{a>0:\operatorname{span}
 \{e^{i\lambda x}:\lambda\in\Lambda\}
                  \text{ is dense in }L^2(-a,a)\}.
\]
The supremum of the empty set is zero. The classical
Beurling--Malliavin theorem identifies this radius with
$\pi D_{\rm eff}(\Lambda)$ in the effective-density convention
of \cite[Section 1.1]{p09:ga15:MP}.

\begin{theorem}[Exact moment-free terminal criterion]
\label{p10:ga16:thm:terminal}
For $\mathcal N\subset\mathbb N$, the separate population laws
$\{\mathcal L(S_n):n\in\mathcal N\}$ identify every mixing law in
\eqref{p10:ga16:eq:symmetric} if and only if
\begin{equation}\label{p10:ga16:eq:criterion}
             \mathcal R(\{\log n:n\in\mathcal N\})=\infty.
\end{equation}
If this radius is finite, there are distinct absolutely continuous
witnesses with index supported inside any prescribed nonempty open
subinterval of $(\alpha_0,\alpha_1)$. They have the same index
marginal and the same conditional moments of $\log B$ of every
nonnegative integer order. All absolute log-scale moments are finite,
but the witnesses can both be chosen to satisfy
\[
       \mathbb E B^\epsilon=\mathbb E B^{-\epsilon}=\infty
                         \quad\text{for every }\epsilon>0.
\]
Their signed difference is smooth after the change of coordinates
$(\alpha,B)\mapsto(1/\alpha,-\log B/\alpha)$, compact in the first
coordinate and Schwartz in the second.
\end{theorem}

The geometric design $n_j=2^j$ illustrates the change: its logarithms
are uniformly separated and have finite completeness radius, so it
fails in the unrestricted class. The preceding positive-and-negative
scale-moment theorem identifies the mixing law from this same infinite
design \cite{comp:part:09}. Having every log-scale moment, even with
those moments supplied as extra conditional data, does not repair the
failure proved here.

\subsection{A projection representation requiring no moments}

Let $E$ be an independent mean-one exponential variable and define
\begin{equation}\label{p10:ga16:eq:QR}
 R=1/\alpha,\qquad Q=-R\log B,\qquad C=Q+R\log E.
\end{equation}
For every real $x$ and positive integer $n$,
\begin{equation}\label{p10:ga16:eq:survival}
 \mathbb E e^{i e^x S_n}
 =\mathbb E e^{-nB e^{\alpha x}}
 =\mathbb P(C-R\log n>x).
\end{equation}
This is an identity between population transforms, not an assumption
that the auxiliary variable $E$ is observed. The selected sum laws
are equivalent to the indicated projection laws. In particular,
positive frequencies determine the symmetric characteristic functions
at negative frequencies as well.

The map from the joint law of $(Q,R)$ to that of $(C,R)$ is injective
on finite complex measures, not only on probabilities. For a fixed
real $t$, conditional Fourier transformation multiplies the complex
measure in $r$ by
\begin{equation}\label{p10:ga16:eq:gamma}
        \mathbb E E^{itr}=\Gamma(1+itr).
\end{equation}
This factor never vanishes and has a bounded reciprocal on the compact
interval $[1/\alpha_1,1/\alpha_0]$. Division for each fixed $t$,
followed by Fourier uniqueness, proves the assertion. No moment of
$Q$, $C$, or $B$ is used.

We record two immediate consequences of the definition of completeness
radius. If $\mathcal R(\Lambda)=\infty$, a compactly supported finite
complex measure whose Fourier transform vanishes on $\Lambda$ must be
zero: convolution with a compact smooth approximate identity gives
compactly supported $L^2$ annihilators, which vanish by completeness
on sufficiently large intervals. Conversely, if the radius is finite,
there is a nonzero complex $f\in C_c^\infty(\mathbb R)$ with
\begin{equation}\label{p10:ga16:eq:fzeros}
            \int e^{-i\lambda y}f(y)\,dy=0
                              \quad(\lambda\in\Lambda).
\end{equation}
Indeed choose an $L^2$ annihilator on an interval larger than the
radius and convolve it with a nonzero compact smooth bump. The
convolution is nonzero because its entire Fourier transform is a
product of two nonzero entire functions. Conjugation changes the
Fourier sign if needed, and translation of $f$ preserves all its zeros.

\begin{proof}[Proof of Theorem~\ref{p10:ga16:thm:terminal}]
Suppose first that \eqref{p10:ga16:eq:criterion} holds. Let $\sigma$ be
the difference of the induced laws of $(C,R)$ for two candidates.
For fixed $t\ne0$, the compact complex measure
\[
             \nu_t(dr)=\int e^{itc}\,\sigma(dc,dr)
\]
has Fourier zeros at $-t\log n$ by \eqref{p10:ga16:eq:survival}.
Scaling or reflecting frequencies preserves whether the completeness
radius is infinite. Thus $\nu_t=0$. Continuity gives the same
conclusion at $t=0$, and Fourier uniqueness after restriction to
arbitrary Borel sets in $r$ gives $\sigma=0$. Equation
\eqref{p10:ga16:eq:gamma} recovers the law of $(Q,R)$ and hence of
$(\alpha,B)$.

Now assume the radius is finite and choose $f$ satisfying
\eqref{p10:ga16:eq:fzeros} for $\Lambda=\{\log n:n\in\mathcal N\}$.
Let $J$ be a compact nondegenerate interval strictly inside the
prescribed reciprocal-index interval. Choose $t_0>0$ sufficiently
large, and translate $f$, so that $\operatorname{supp}f/t_0$
lies strictly inside $J$. Choose a nonzero real
$h\in C_c^\infty((0,\infty))$ with support sufficiently close to
$t_0$ that $\operatorname{supp}f/t\subset\operatorname{int}J$
throughout that support. For $t>0$ define
\begin{equation}\label{p10:ga16:eq:band}
 G_C(t,r)=h(t)t f(tr),\qquad
 G_Q(t,r)=\frac{G_C(t,r)}{\Gamma(1+itr)}.
\end{equation}
Extend both functions to negative $t$ by conjugation and by zero
near $t=0$. They are smooth with compact support in $(t,r)$.
Their inverse Fourier transforms in $t$, denoted $g_C(c,r)$ and
$g_Q(q,r)$, are real, nonzero, compactly supported in $r\in J$,
and Schwartz in their other coordinate, uniformly in $r$.
Equation \eqref{p10:ga16:eq:gamma} shows that adding $R\log E$ to
the signed measure $g_Q(q,r)\,dq\,dr$ gives exactly
$g_C(c,r)\,dc\,dr$.

For $t>0$ and every selected $n$,
\[
 \int e^{-it r\log n}G_C(t,r)\,dr
       =h(t)\int e^{-iy\log n}f(y)\,dy=0.
\]
Negative $t$ follows by conjugation. Consequently every selected
projection $C-R\log n$ has zero signed law, and
\eqref{p10:ga16:eq:survival} gives zero signed stable sum law.
Since $G_Q$ vanishes in a neighborhood of $t=0$,
\begin{equation}\label{p10:ga16:eq:conditionalmoments}
             \int q^j g_Q(q,r)\,dq=0
                      \quad(r\in J,\ j=0,1,2,\ldots).
\end{equation}
The positive and negative parts have the same nonzero total mass.
Normalize them to probabilities and push forward by
$\alpha=1/r$, $B=e^{-q/r}$. Their index marginals and all
conditional log-scale moments agree by
\eqref{p10:ga16:eq:conditionalmoments}; their absolute log-scale moments
are finite by the uniform Schwartz bound.

Finally mix each witness, with the same fixed positive weight, with
one common probability having a smooth index density in the allowed
interval and, independently in $(Q,R)$ coordinates, a $Q$ density
proportional to $\exp\{-(1+q^2)^{1/4}\}$. This common law has all
absolute log-scale moments and has infinite positive and negative
scale moments of every positive order. Both modified witnesses
therefore have these properties, while all observational equalities,
conditional moment equalities, and nonzeroness of their difference
are preserved.
\end{proof}

\section{Two joint replicates recover an arbitrary index--scale mixture}
\label{p10:ga16:sec:replication}

The change from separate sum laws to a joint replicated law supplies
information that the preceding sparse designs discard. The following
result permits arbitrary continuous mixing and does not bound the
number of mixture components. General finite-mixture grouped-sample
results, such as \cite{p02:frontmix:VS2015,p10:ga16:BJR}, address a different scope.

\begin{theorem}[Exact replication threshold at zero location]
\label{p10:ga16:thm:two}
For the class \eqref{p10:ga16:eq:symmetric}, the joint population law of
$(X_1,X_2)$ identifies every Borel mixing law of $(\alpha,B)$ without
any moment assumption. More strongly, the map
\begin{equation}\label{p10:ga16:eq:pairmap}
 \mu\longmapsto\left\{
 \int e^{-b(|u|^\alpha+|v|^\alpha)}\,d\mu(\alpha,b):
                                     u,v\in\mathbb R\right\}
\end{equation}
is injective on finite complex measures on
$[\alpha_0,\alpha_1]\times(0,\infty)$.
One replicate does not identify all probability mixtures on any
nondegenerate index interval. Thus two is the exact minimal number
of joint replicates for this class.
\end{theorem}

\subsection{A triangular dilation lemma}

\begin{lemma}[Recovery from integer dilations]
\label{p10:ga16:lem:dilation}
Let $0<a<b<\infty$, and let $c_m$ be bounded Borel functions on
$[a,b]$ satisfying $c_1=1$ and
$\sup_{\alpha\in[a,b]}|c_m(\alpha)|\leq C m^p$ for some fixed
$C,p$. If a finite complex measure $\eta$ on $[a,b]$ satisfies
\begin{equation}\label{p10:ga16:eq:dilationkernel}
 \sum_{m\geq1}\int c_m(\alpha)e^{-m\alpha d}\,d\eta(\alpha)=0
                                      \quad(d>0),
\end{equation}
then $\eta=0$.
\end{lemma}
\begin{proof}
Let $D_m(\alpha)=m\alpha$ and form the locally finite complex
measure on $[a,\infty)$
\[
              M_\eta=\sum_{m\geq1}(D_m)_*(c_m\eta).
\]
Only finitely many terms meet any bounded interval, and
$\int e^{-d x}\,d|M_\eta|(x)<\infty$ for every $d>0$ by the
polynomial coefficient bound. Equation
\eqref{p10:ga16:eq:dilationkernel} says that its Laplace transform is
zero. For completeness, fix $d_0>0$, multiply the measure by
$e^{-d_0x}$, and push forward by $y=e^{-x}$. Its moments of every
nonnegative integer order vanish. Polynomial density in the
continuous functions on $[0,e^{-a}]$ gives the zero measure, so
$M_\eta=0$.

On $[a,2a)$ only $m=1$ contributes, hence $\eta=0$ there.
If $\eta$ is already zero on $[a,L)$, then on $[a,2L)$ every
term with $m\geq2$ comes from $\alpha< L$ and is zero. Thus
$M_\eta=0$ forces $\eta=0$ on $[a,2L)$. Iteration reaches
past $b$ after finitely many steps. This proves the lemma and
also describes a recursive recovery of $\eta$ from $M_\eta$.
\end{proof}

\subsection{The two-replicate inversion}

\begin{proof}[Proof of Theorem~\ref{p10:ga16:thm:two}]
It suffices to prove that a finite complex measure in the kernel of
\eqref{p10:ga16:eq:pairmap} is zero. Push that measure to $(Q,R)$
coordinates by \eqref{p10:ga16:eq:QR}, writing it as $\rho$.
Let $E_1,E_2$ be independent mean-one exponential variables and set
\[
                  Y_j=Q+R\log E_j,\qquad j=1,2.
\]
The induced finite complex measure of $(Y_1,Y_2)$ has upper-orthant
values
\begin{equation}\label{p10:ga16:eq:jointsurvival}
 \int \exp\{-B(e^{\alpha x}+e^{\alpha y})\}\,d\mu(\alpha,B),
\end{equation}
which vanish by \eqref{p10:ga16:eq:pairmap}. Upper orthants determine
finite complex measures: rectangle values follow by inclusion and
exclusion, and their unbounded endpoints by finite-variation limits.
Thus the induced measure of $(Y_1,Y_2)$ is zero.

Fix a real $t$. Fourier-weight its first coordinate by $e^{itY_1}$
and project to $D=Y_2-Y_1$. This gives the zero complex measure
on the real line. Write
\[
                   \nu_t(dr)=\int e^{itq}\,\rho(dq,dr).
\]
The density of that weighted difference is
\begin{equation}\label{p10:ga16:eq:density}
 T_t(d)=\int\frac{\Gamma(2+itr)}r
       \frac{e^{d/r}}{(1+e^{d/r})^{2+itr}}\,d\nu_t(r).
\end{equation}
To verify the formula, set $E_1=z$ and $E_2=z e^{d/r}$ in
the two exponential densities. The Jacobian is
$z e^{d/r}/r$, while the Fourier weight contributes
$e^{itq}z^{itr}$. Integration in $z>0$ gives
\eqref{p10:ga16:eq:density}. Its integrand is continuous in $d$,
locally uniformly in the compact positive $r$ interval. Hence
$T_t$ is continuous and vanishes everywhere, not merely almost
everywhere.

For $d>0$, expand the denominator at $e^{-d/r}$ to obtain
\begin{equation}\label{p10:ga16:eq:expand}
 e^{itd}T_t(d)=\sum_{m\geq1}
   \int \frac{\Gamma(2+itr)}r
     \frac{(-1)^{m-1}(2+itr)_{m-1}}{(m-1)!}
                  e^{-md/r}\,d\nu_t(r).
\end{equation}
The cancellation of the factor $e^{-itd}$ is exact. For
$|tr|\leq T$, the coefficients obey the elementary uniform bound
\[
 \frac{|(2+itr)_{m-1}|}{(m-1)!}
  =m\prod_{j=2}^{m}\left(1+\frac{t^2r^2}{j^2}\right)^{1/2}
  \leq m\exp\left\{\frac{T^2}{2}\sum_{j=2}^\infty j^{-2}\right\}.
\]
Consequently the series can be integrated absolutely for every
$d>0$. Push $\nu_t$ forward by $\alpha=1/r$ and multiply it by
the bounded nonvanishing function
$\alpha\Gamma(2+it/\alpha)$, obtaining a finite complex measure
$\eta_t$ on $[\alpha_0,\alpha_1]$. Equation
\eqref{p10:ga16:eq:expand} has exactly the form
\eqref{p10:ga16:eq:dilationkernel}, with
\[
       c_m(\alpha)=
       \frac{(-1)^{m-1}(2+it/\alpha)_{m-1}}{(m-1)!}.
\]
Lemma~\ref{p10:ga16:lem:dilation} yields $\eta_t=0$, hence $\nu_t=0$.
This holds for every real $t$. Fourier uniqueness in $q$ after
restriction to arbitrary Borel sets in $r$ gives $\rho=0$ and
then $\mu=0$. No moment differentiation or analytic continuation
in the unbounded location coordinate $Q$ has been used.

To prove failure with one replicate, choose
$\alpha_0\leq\alpha_a<\alpha_b\leq\alpha_1$ and put
$\delta=\alpha_a/\alpha_b\in(0,1)$. The classical positive
stable variable $V$ has Laplace transform
$\mathbb E e^{-sV}=e^{-s^\delta}$; see
\cite[Equation (1.1)]{p10:ga16:Simon}. Its existence also follows
directly from the L\'evy measure
$\delta x^{-1-\delta}dx/\Gamma(1-\delta)$, since its Laplace
exponent is $s^\delta$. The variable is strictly positive because
its Laplace transform tends to zero at infinity. Therefore
\[
 e^{-|t|^{\alpha_a}}
          =\mathbb E e^{-V|t|^{\alpha_b}}.
\]
The mixing laws $(\alpha,B)=(\alpha_a,1)$ and
$(\alpha,B)=(\alpha_b,V)$ are distinct and give the same law
of one observation. This classical subordination obstruction
establishes minimality; it is not itself claimed as a new result.
\end{proof}

The inversion above also proves identification of arbitrary
location--scale mixtures of the reflected Gumbel kernel
$Q+R\log E$ from two replicates, with $R$ in a compact positive
interval. Equivalently it identifies continuously mixed Weibull
shape and rate from a joint pair of conditionally independent
lifetimes. These are transform-equivalent consequences of the same
proof, not separately counted discoveries. Transform transfer is a
classical methodological theme; compare \cite{p10:ga16:PX}.

\section{Contrasts and anchors recover correlated stable parameters}
\label{p10:ga16:sec:anchors}

The complex-measure formulation in Theorem~\ref{p10:ga16:thm:two} is
essential for recovering dependence on additional latent parameters.
The following extensions make no independence assumption between
those latent coordinates.

\begin{corollary}[Five replicates with arbitrary symmetric location]
\label{p10:ga16:cor:five}
Conditional on $(A,\alpha,B)$, let $X_j$ be independent with
characteristic function $e^{itA-B|t|^\alpha}$, where
$A\in\mathbb R$, $\alpha\in[\alpha_0,\alpha_1]$, and $B>0$.
The joint law of five observations identifies every joint mixing law
of $(A,\alpha,B)$, without moments of any coordinate. Five is a
sufficient bound; its minimality is not asserted.
\end{corollary}
\begin{proof}
Set $D_1=X_2-X_3$ and $D_2=X_4-X_5$. Given the latent draw,
these are independent symmetric stable observations with exponent
coefficient $2B$, independent also of $X_1$. For a difference
$\rho$ of two mixing laws giving the same joint observations and
a fixed real $t$, define
\[
 \sigma_t(d\alpha,dB)=
            \int e^{itA-B|t|^\alpha}\,\rho(dA,d\alpha,dB).
\]
The joint transform of $(X_1,D_1,D_2)$ shows that
\[
       \int e^{-2B(|u|^\alpha+|v|^\alpha)}\,d\sigma_t=0
                                      \quad(u,v\in\mathbb R).
\]
Theorem~\ref{p10:ga16:thm:two}, applied after the invertible change
$B\mapsto2B$, gives $\sigma_t=0$. The multiplier
$e^{-B|t|^\alpha}$ is positive everywhere, so it can be removed
by restricting first to bounded $B$ intervals and then taking their
union. Thus the Fourier transform in $A$ of $\rho$ restricted
to each Borel set in $(\alpha,B)$ is zero. Fourier uniqueness
proves $\rho=0$.
\end{proof}

For varying skewness we specify the parametrization, since a location
coordinate changes under different stable conventions. Put
\begin{equation}\label{p10:ga16:eq:psi}
 \psi_\alpha(t)=
 \begin{cases}
   \tan(\pi\alpha/2)\operatorname{sgn}(t)|t|^\alpha,
                                                   &\alpha\ne1,\ 2,\\
   -(2/\pi)t\log|t|,                              &\alpha=1,\\
   0,                                             &\alpha=2,
 \end{cases}
 \qquad \psi_\alpha(0)=0.
\end{equation}
Use the standard stable kernels
\begin{equation}\label{p10:ga16:eq:fullstable}
 \widehat K_{A,\alpha,B,\beta}(t)
          =\exp\{itA-B|t|^\alpha+iB\beta\psi_\alpha(t)\}.
\end{equation}
Here $B$ is the coefficient of $|t|^\alpha$ and
$\beta\in[-1,1]$. At $\alpha=2$ impose $\beta=0$, since
skewness is otherwise a redundant, unobservable Gaussian parameter.
The convention is allowed to be discontinuous in $\alpha$ at one;
the proof needs only continuity in the observation frequencies for
each fixed latent parameter, which holds.

\begin{corollary}[Six replicates with all four stable parameters mixed]
\label{p10:ga16:cor:six}
Let $0<\alpha_0<2$ and let the conditional kernels be
\eqref{p10:ga16:eq:fullstable}, with
$A\in\mathbb R$, $\alpha\in[\alpha_0,2]$, $B>0$, and
$\beta\in[-1,1]$, subject to $\beta=0$ at $\alpha=2$.
The joint population law of six conditionally independent observations
identifies every Borel joint mixing law of $(A,\alpha,B,\beta)$.
No coordinate moment and no latent-coordinate independence is required.
Six is a sufficient bound, not an asserted minimal bound.
\end{corollary}
\begin{proof}
Keep $X_1,X_2$ as anchors and set
$D_1=X_3-X_4$, $D_2=X_5-X_6$. The function $\psi_\alpha$ is
odd. Thus, conditional on all latent coordinates, $D_1,D_2$ are
independent symmetric stable observations of coefficient $2B$,
independent of the two anchors.

Let $\rho$ be the difference of two candidate mixing laws. For
fixed real $p,q$, integrate $\rho$ over $(A,\beta)$ with weight
\[
 e^{i(p+q)A-B(|p|^\alpha+|q|^\alpha)
                       +iB\beta(\psi_\alpha(p)+\psi_\alpha(q))}.
\]
This gives a finite complex measure on $(\alpha,B)$. The joint
transform with the two contrasts places it in the kernel of
\eqref{p10:ga16:eq:pairmap}, so it is zero by
Theorem~\ref{p10:ga16:thm:two}. Remove the strictly positive real
factor $e^{-B(|p|^\alpha+|q|^\alpha)}$ as in the preceding proof.

To make the ensuing conditional argument simultaneous in the
frequencies, dominate both $(\alpha,B)$ marginals by their sum
$\lambda$ and disintegrate $\rho$ with respect to $\lambda$,
obtaining finite signed conditional measures
$\rho_{\alpha,B}$ on $(A,\beta)$. First intersect the full-measure
sets furnished by the countably many rational pairs $(p,q)$.
On that common set the identities extend to all real $(p,q)$ by
bounded convergence and frequency continuity. Hence, for almost
every fixed $(\alpha,B)$,
\begin{equation}\label{p10:ga16:eq:conditionalanchor}
 \int e^{i(p+q)A+iB\beta(\psi_\alpha(p)+\psi_\alpha(q))}
                    \,d\rho_{\alpha,B}(A,\beta)=0
                                      \quad(p,q\in\mathbb R).
\end{equation}

Fix such a pair $(\alpha,B)$ with $\alpha<2$ and a nonzero
$u=p+q$. As $p$ varies strictly between $0$ and $u$ when $u>0$
(with reversed signs when $u<0$), the values
\[
               B\{\psi_\alpha(p)+\psi_\alpha(u-p)\}
\]
contain a nondegenerate interval. For $\alpha\ne1$, this follows
because $\tan(\pi\alpha/2)\ne0$ and the second derivative of
$p^\alpha+(u-p)^\alpha$ is
$\alpha(\alpha-1)\{p^{\alpha-2}+(u-p)^{\alpha-2}\}\ne0$.
For $\alpha=1$, the second derivative of
$p\log p+(u-p)\log(u-p)$ is $1/p+1/(u-p)>0$.

For each fixed $u$, the function
\[
             z\longmapsto
             \int e^{iuA+iz\beta}\,d\rho_{\alpha,B}(A,\beta)
\]
is entire because $|\beta|\leq1$. By
\eqref{p10:ga16:eq:conditionalanchor} it vanishes on that real
interval, hence everywhere. Continuity also supplies $u=0$.
Fourier uniqueness in $(A,\beta)$ gives
$\rho_{\alpha,B}=0$. If $\alpha=2$, the imposed condition
$\beta=0$ reduces \eqref{p10:ga16:eq:conditionalanchor} directly to
Fourier uniqueness in $A$. Therefore all conditional signed
measures vanish and $\rho=0$.
\end{proof}

\paragraph{Scope, previous work, and remaining barriers.}
The two main targets are the exact unrestricted-scale terminal
criterion and the exact two-replicate threshold for zero-location
symmetric mixtures. The five- and six-replicate results are extensions
of the second mechanism and are not counted as separate independent
closures. The classical BM theorem and the one-observation
subordination identity are not new results. The consulted primary
sources include the BM theorem \cite{p09:ga15:MP}, grouped finite-mixture
identification \cite{p02:frontmix:VS2015,p10:ga16:BJR}, and transform-based transfers
\cite{p10:ga16:PX}; they do not constitute an exhaustive priority search.

The lower bound $\alpha\geq\alpha_0>0$ is used essentially in
the displayed triangular inversion. If indices approach zero, integer
dilations have no positive initial interval from which to begin; this
proof does not settle that case. Minimal replication with unknown
location or skewness remains open, as does quantitative stability
under noisy estimated characteristic functions. The finite-stage
dilation recursion is an exact uniqueness mechanism after Laplace
inversion; it is not a claim that statistical inversion is well
conditioned. Nor do the results identify all four parameters from
separate sum marginals: the anchor argument explicitly uses joint
grouped observations.


% ===== spectral =====

\section{Continuous spectra and arithmetic observation costs}
\label{p10:sp16:sec}

The finite-atom theory leaves two different questions unanswered.
What can a threshold curve determine when no finite bound on the number
of frequencies is available?  When can a single experiment near a
transition be observed in a height of the best possible order?
We answer the first question for arbitrary probability spectra, and the
second for one hidden frequency.  The first answer is a sharp logarithmic
modulus, attained by a finite convex reconstruction.  The second is an
arithmetic classification of the observation height.

Neither underlying classical mechanism is claimed as a new principle.
Polynomial approximation and moment matching are standard tools for
Wasserstein recovery; see Han, Jiao, and Weissman
\cite{p10:sp16:momentmatching}.  The same factorial inverse scale also
appeared for a different kernel in the earlier integrated manuscript
\cite{p01:KEIOIntegrated}.  Likewise, the equivalence between optimal
covering of a circle rotation and bounded partial quotients belongs to
classical Diophantine approximation; Beck and Chen
\cite{p10:sp16:beckchen} describe its continuous version and extensions.
The results below establish the nonlinear reductions needed for these
mechanisms to answer the two present observation problems.

\subsection{An inverse without a bound on spectral complexity}

Fix $0<a<b<1$, and let $\mathcal P[a,b]$ be all Borel probability
measures on $[a,b]$.  Define
\begin{equation}\label{p10:sp16:eq:model}
 S_\mu(x)=\int_a^b x^r\,d\mu(r),\qquad
 T_\mu(C)=\min_{x>0}\frac{1+Cx}{S_\mu(x)},\qquad
 J=\left[\frac{a}{2e(1-a)},\frac{2b}{1-b}\right].
\end{equation}
Write $d_J(\mu,\nu)=\|T_\mu-T_\nu\|_{L^\infty(J)}$ and
\[
 \Omega_J(\delta)=\sup\{W_1(\mu,\nu):
          \mu,\nu\in\mathcal P[a,b],\ d_J(\mu,\nu)\le\delta\}.
\]
The Wasserstein distance uses the ordinary distance between exponents.

\begin{theorem}[Sharp nonparametric threshold inverse]
\label{p10:sp16:thm:nonparametric}
There are constants $c,K,\delta_0>0$, depending only on $a,b$, such
that, for $0<\delta<\delta_0$,
\begin{equation}\label{p10:sp16:eq:modulus}
 c\,\frac{\log\log(1/\delta)}{\log(1/\delta)}
 \le \Omega_J(\delta)
 \le K\,\frac{\log\log(1/\delta)}{\log(1/\delta)}.
\end{equation}
In particular the exact threshold curve on $J$ determines every
probability spectrum, including continuous spectra.  The lower bound
holds even though the entire curve, rather than finitely many settings,
is supplied.

If an oracle returns $Y(C)$ with
$|Y(C)-T_\mu(C)|\le\varepsilon$ at each requested $C\in J$, a finite
algorithm attains the upper bound in \eqref{p10:sp16:eq:modulus}, with
$\delta$ replaced by $\varepsilon$.  It uses at most
\[
 K\frac{\log^2(1/\varepsilon)}{\log\log(1/\varepsilon)}
\]
oracle calls and a linear feasibility problem with at most
$K/\varepsilon$ nonnegative variables.  There is no atom-count,
minimum-mass, or separation assumption.
\end{theorem}

The computational assertion uses exact arithmetic and bounded
adversarial oracle errors.  Its grid size is polynomial in inverse
accuracy, not in the number of bits of accuracy.  It therefore does not
assert optimal computational complexity.

\begin{lemma}[Uniform contact geometry]\label{p10:sp16:lem:duality}
For every $\mu\in\mathcal P[a,b]$,
\begin{equation}\label{p10:sp16:eq:duality}
 S_\mu(x)=\min_{C\in J}\frac{1+Cx}{T_\mu(C)}
                         \quad(1\le x\le e).
\end{equation}
Consequently $d_J(\mu,\nu)\le\delta$ implies
$|S_\mu(x)-S_\nu(x)|\le K\delta$ on $[1,e]$.
Moreover, for any prescribed $x\in[1,e]$, an oracle of error
$\varepsilon$ estimates $S_\mu(x)$ to error $K\varepsilon$ using
$K\log(3/\varepsilon)$ queries.
\end{lemma}
\begin{proof}
All derivatives of $S$ may be taken under the integral.  In particular
$S'>0$, $S''<0$, and $H=S-xS'>0$.  The unique minimizer $x_C$ solves
\[
 C=\frac{S'(x_C)}{H(x_C)},\qquad
 \frac{dC}{dx}=\frac{S''(x)S(x)}{H(x)^2}<0,\qquad
 \frac a{1-a}\le Cx_C\le\frac b{1-b}.
\]
The last inequality and the compactness of $J$ put every $x_C$ in
a common compact subset of $(0,\infty)$.  On this compact set,
$S,H,-S''$, and $T(C)$ have uniform positive lower and upper bounds.
The contact setting $C_x=S'(x)/H(x)$ for $x\in[1,e]$ lies in the
interior of $J$.  Every affine function $(1+Cx)/T(C)$ majorizes
$S(x)$, with equality at $x_C$, proving \eqref{p10:sp16:eq:duality}.
The assertion for two profiles follows by comparing the two minima;
the uniform lower bound on $T$ controls inversion of $T$.

For completeness the noisy scalar minimization is explicit.  For
$g_x(C)=(1+Cx)/T(C)$, differentiation of the contact formula gives
\[
 g_x'(C)=\frac{x-x_C}{T(C)(1+Cx_C)},\qquad
 c|C-C_x|\le |g_x'(C)|\le K|C-C_x|.
\]
The observed $g_x$ has error at most $\eta=K_0\varepsilon$.
Maintain an interval $[L,U]$ containing $C_x$, and compare its two
trisection points $C_1<C_2$.  Define the observed difference as
$\widetilde g_x(C_2)-\widetilde g_x(C_1)$.  If it is greater than
$2\eta$, discard the right third; if less than $-2\eta$, discard
the left third.  These decisions follow from unimodality.  Otherwise
the true difference has modulus at most $4\eta$.  When the interval
width $q$ satisfies $q^2>144\eta/c$, integration of the lower
derivative bound shows that a minimizer outside the middle third
would give a difference at least $cq^2/18>4\eta$ in modulus.
Thus in the ambiguous case retain the middle third.  Stop at
$q^2\le144\eta/c$ and query the midpoint.  The upper derivative
bound gives error $K\eta$ in the minimum.  Each step contracts by
at least $2/3$, proving the query bound.
\end{proof}

\begin{lemma}[Exponential samples control Wasserstein distance]
\label{p10:sp16:lem:samples}
For $k\ge2$ and $\mu,\nu\in\mathcal P[a,b]$, put
\[
 E_k=\max_{0\le l\le k}
       \left|\int e^{lr/k}\,d(\mu-\nu)(r)\right|.
\]
Then
\begin{equation}\label{p10:sp16:eq:samplebound}
                  W_1(\mu,\nu)\le\frac Kk+E_k(Kk)^k.
\end{equation}
\end{lemma}
\begin{proof}
Let $v=(e^{r/k}-e^{a/k})/d_k$, where
$d_k=e^{b/k}-e^{a/k}\asymp k^{-1}$.  This maps $[a,b]$ onto
$[0,1]$, and its inverse
$r_k(v)=k\log(e^{a/k}+d_kv)$ has a uniformly bounded derivative.
For $0\le j\le k$, expansion by the binomial theorem yields
\[
 \left|\int v(r)^j\,d(\mu-\nu)(r)\right|
 \le E_k d_k^{-j}(1+e^{a/k})^j\le E_k(Kk)^j.
\]
If $f$ is 1-Lipschitz and $f(a)=0$, the function
$g(v)=f(r_k(v))$ is uniformly bounded and Lipschitz.  There exists
a polynomial $p(v)=\sum_{j=0}^k p_jv^j$ with
\begin{equation}\label{p10:sp16:eq:jackson}
       \|g-p\|_\infty\le K/k,\qquad
       \sum_{j=0}^k|p_j|\le K^k.
\end{equation}
Here is a direct justification of the two bounds together.  The even
periodic function $g((1+\cos\theta)/2)$ has a uniform Lipschitz
constant.  Convolve it with the normalized nonnegative trigonometric
kernel
\[
 J_n(\theta)=c_n\left(\frac{\sin(n\theta/2)}
                              {\sin(\theta/2)}\right)^4,
 \qquad n=\lfloor k/2\rfloor+1,
\]
where normalization means integral one on $[-\pi,\pi]$.
Elementary bounds on the sine give $c_n\asymp n^{-3}$ and
$\int|\theta|J_n(\theta)\,d\theta\le K/n$.  The convolution
is an even trigonometric polynomial of degree $2(n-1)\le k$,
with approximation error $K/k$ and uniformly bounded cosine
coefficients.  Writing it as a polynomial in $v$ uses the
Chebyshev polynomials $T_j(2v-1)$.  Their recurrence bounds the sum
of absolute monomial coefficients by $K^j$, which proves
\eqref{p10:sp16:eq:jackson}.  Integrating the approximation and using
the moment bounds gives \eqref{p10:sp16:eq:samplebound}, after enlarging
$K$.  Taking the supremum over 1-Lipschitz $f$ is the
Kantorovich--Rubinstein formula on an interval.
\end{proof}

\begin{proof}[Proof of Theorem~\ref{p10:sp16:thm:nonparametric}]
Lemma~\ref{p10:sp16:lem:duality} gives $E_k\le K\delta$.  With
$L=\log(1/\delta)$ and
$k=\lfloor L/(4\log L)\rfloor$, the second term of
\eqref{p10:sp16:eq:samplebound} is smaller than $1/k$ for all
sufficiently small $\delta$.  This proves the upper bound.

For the lower bound let $d=b-a$, and define two probabilities by
the even and odd parts of
\begin{equation}\label{p10:sp16:eq:binomial}
 \mu_k-\nu_k=2^{1-k}\sum_{j=0}^k(-1)^j\binom kj
                                \delta_{a+jd/k}.
\end{equation}
The signed mass has total zero and both signs have mass one.
For any real $u$,
\[
 \int e^{ur}\,d(\mu_k-\nu_k)(r)
             =2^{1-k}e^{au}(1-e^{ud/k})^k.
\]
On any fixed bounded $u$ interval its modulus is at most
$(K/k)^k$.  The contact bounds in the preceding lemma put every
minimizer for $C\in J$ and every probability spectrum in a fixed
compact $x$ interval.  The profile there is uniformly positive.
Comparing the two objective functions at each other's minimizers
therefore gives
\[
                         d_J(\mu_k,\nu_k)\le(K_1/k)^k.
\]
The signed distribution function between the $j$th and $(j+1)$st
nodes has modulus $2^{1-k}\binom{k-1}{j}$, by the alternating
binomial identity.  The one-dimensional transport formula consequently
gives the exact value
\[
                         W_1(\mu_k,\nu_k)=d/k.
\]
Choose the least sufficiently large integer $k$ for which
$(K_1/k)^k\le\delta$.  Elementary comparison of
$k\log(k/K_1)$ with $L$ gives $k\le K L/\log L$ as
$\delta\downarrow0$.  This proves the lower bound at every
sufficiently small $\delta$, rather than only along a sequence.

For reconstruction choose the preceding upper-bound value of $k$,
now using $L=\log(1/\varepsilon)$.  Estimate
$S_\mu(e^{l/k})$, $1\le l\le k$, by
Lemma~\ref{p10:sp16:lem:duality}; set the zeroth sample equal to one.
Take a grid $\mathcal G_\varepsilon\subset[a,b]$ of mesh at most
$\varepsilon$ including the endpoints.  Solve the linear constraints
\[
 q_g\ge0,\quad\sum_{g\in\mathcal G_\varepsilon}q_g=1,
 \quad
 \left|\sum_{g\in\mathcal G_\varepsilon}q_g e^{lg/k}
                 -\widehat S_l\right|\le K_2\varepsilon
                         \quad(1\le l\le k).
\]
They are feasible: rounding every point of the true measure to a
nearest grid point changes each exponential sample by at most
$K\varepsilon$, uniformly in $l\le k$.  Choose $K_2$ larger than
this constant plus the profile-estimation constant.  Every feasible
probability $\widehat\mu=\sum q_g\delta_g$ has
$E_k(\widehat\mu,\mu)\le K\varepsilon$ and hence satisfies the
desired Wasserstein bound.  The grid has $O(\varepsilon^{-1})$
points and the query count is $O(k\log(1/\varepsilon))$.
\end{proof}

Thus replacing a finite atomic prior by arbitrary positive spectra
changes the stability class.  Positivity still permits recovery, but
there is no uniform power-law inverse on the enlarged class.  This is
an answer about a nonlinear threshold observation operator, not a
closure of a named conjecture about arbitrary Laplace inversion.

\subsection{The exact arithmetic obstruction to a short zero experiment}

Fix $0<r<1$, $C,\ell,\alpha,R>0$.  Set
\begin{equation}\label{p10:sp16:eq:oneatom}
 x_* =\frac{r}{C(1-r)},\quad
 T=\frac{1+Cx_*}{x_*^r},\quad A_h=T+\alpha h,
 \qquad
 f_h(z)=-1+A_h e^{\ell r z+i\phi}
                         +Ce^{\ell z+i\phi_*}.
\end{equation}
A zero $z=\sigma+it$ is \emph{marked} if
\begin{equation}\label{p10:sp16:eq:mark}
        \operatorname{dist}(\ell t+\phi_*,\pi+2\pi\mathbb Z)
                                                     \le R\sqrt h.
\end{equation}
Let $H_r(h)$ be the infimum of heights $H$ such that, for
\emph{every} pair of phases $(\phi,\phi_*)$, there is a marked zero
with $0\le t\le H$.  The real coordinate is unrestricted.
For an irrational $r$, write
\[
 \mathcal R_N(r)=\sup_{\theta\in\mathbb R/\mathbb Z}
          \min_{0\le n<N}\|nr-\theta\|_{\mathbb R/\mathbb Z},
 \qquad
 Q_r(\eta)=\min\{N\ge1:\mathcal R_N(r)\le\eta\}.
\]

\begin{theorem}[Marked-zero height and bounded partial quotients]
\label{p10:sp16:thm:height}
For fixed irrational $r$ and the other fixed parameters in
\eqref{p10:sp16:eq:oneatom}, there are $c,K,h_0>0$ such that
\begin{equation}\label{p10:sp16:eq:coverbridge}
 \frac{2\pi}{\ell}\bigl(Q_r(Kh)-3\bigr)
 \le H_r(h)
 \le\frac{2\pi}{\ell}\bigl(Q_r(ch)+3\bigr)
                    \qquad(0<h<h_0).
\end{equation}
Consequently the following are equivalent:
\begin{enumerate}
\item $r$ has bounded continued-fraction partial quotients;
\item there is $K_r<\infty$ such that $H_r(h)\le K_r/h$ for all
      sufficiently small $h$.
\end{enumerate}
In this case $c_r/h\le H_r(h)\le K_r/h$ for small $h$.  If the
partial quotients are unbounded, then
$\limsup_{h\downarrow0}hH_r(h)=\infty$.
\end{theorem}

This theorem concerns the presence of a marked zero at one fixed
critical-scale amplitude.  It is not a lower bound for every algorithm
that may use several amplitudes, unmarked zeros, or other statistics.
It is stronger than an asymptotic density statement because the
height bound is uniform over the initial phases.  The arithmetic
criterion itself is classical; the essential bridge is that the
effective target for the rotation has width $h$, despite the
$\sqrt h$ marking window.

\begin{lemma}[Quadratic phase cancellation]\label{p10:sp16:lem:phase}
Put $t_n=(\pi-\phi_*+2\pi n)/\ell$ and
$\kappa_n=\phi+r(\pi-\phi_*)+2\pi nr$ modulo $2\pi$.
For sufficiently small $h$, the following statements hold, uniformly
over the phases and the integer $n$:
\begin{enumerate}
\item Every marked zero satisfies
      $|t-t_n|\le R\sqrt h/\ell$ and
      $\operatorname{dist}(\kappa_n,2\pi\mathbb Z)\le Kh$
      for its nearest top-phase crossing $t_n$.
\item If $\operatorname{dist}(\kappa_n,2\pi\mathbb Z)\le ch$,
      then a marked zero exists with
      $|t-t_n|\le R\sqrt h/\ell$.
\end{enumerate}
\end{lemma}
\begin{proof}
At a zero put $x=e^{\ell\sigma}$ and
$v=\ell(t-t_n)$, so $|v|\le R\sqrt h$.  The zero equation reads
\begin{equation}\label{p10:sp16:eq:zerophase}
 A_h x^r e^{i(\kappa_n+rv)}=1+Cx e^{iv}.
\end{equation}
For all $x>0$,
\[
 |1+Cx e^{iv}|\ge(1-Kv^2)(1+Cx).
\]
Thus $(1+Cx)/x^r\le T+Kh$.  This objective tends to infinity
at both ends, has its unique minimum at $x_*$, and has strictly
positive second derivative in the variable $u=\log(x/x_*)$ at
$u=0$.  It follows that $|u|\le K\sqrt h$.  Let $y_*=Cx_*$.
Taking the argument in \eqref{p10:sp16:eq:zerophase} near $u=v=0$ gives
\[
 \kappa_n=\arg(1+y_*e^u e^{iv})-rv\pmod {2\pi}.
\]
The coefficient of $v$ at $u=0$ is
$y_*/(1+y_*)-r=0$.  More explicitly, smoothness and oddness in
$v$ give an absolute bound $K(|u||v|+|v|^3)$ for the right side.
This is at most $Kh$, proving the necessary condition.

For the sufficient condition use the complex coordinate
$z=(\log x_*+\zeta)/\ell+it_n$ and choose the representative
$|\kappa_n|\le ch$.  The function in this coordinate is
\[
 F_{h,\kappa}(\zeta)
      =-1+(T+\alpha h)x_*^r e^{r\zeta+i\kappa}
                                             -y_*e^\zeta.
\]
Write $s_*=x_*^r$ and $B=r>0$; the contact identities give,
uniformly for bounded complex $Z$,
\[
 h^{-1}F_{h,0}(\sqrt h Z)
                 =\alpha s_*-\tfrac12 BZ^2+O(\sqrt h).
\]
The limiting polynomial has two distinct real roots $\pm d$,
where $d=\sqrt{2\alpha s_*/B}$.  Choose a radius
$0<\rho<\min(R,d)/4$ around $d$.  Its boundary contains no
limiting root; on it $|F_{h,0}(\sqrt h Z)|\ge c_0h$ for small
$h$, and Rouch\'e's theorem gives one zero in the disk.
The perturbation obeys
$|F_{h,\kappa}-F_{h,0}|\le K_0|\kappa|$ there.
Choose $c<c_0/(2K_0)$.  A second application of Rouch\'e gives
a zero of $F_{h,\kappa}$ in the same disk.  Its imaginary
coordinate has modulus less than $\rho\sqrt h<R\sqrt h$,
so it is marked and has the claimed height proximity.
\end{proof}

\begin{proof}[Proof of \eqref{p10:sp16:eq:coverbridge}]
Absorb the conversion from angles to turns into $c,K$.
The numbers $\kappa_n/(2\pi)$ form a translate of the rotation
orbit $nr$ modulo one.  The spacing of the crossing times is
$L_*=2\pi/\ell$.  Take $N=Q_r(ch)$ consecutive crossings,
starting with the first one after $R\sqrt h/\ell$.  One of their
phase mismatches is at most $ch$.  By
Lemma~\ref{p10:sp16:lem:phase} it supplies a marked zero in
$[0,(N+3)L_*]$, for all initial phases.  This proves the upper
bound.

Conversely fix $\phi_*=\pi$, so crossings occur at $nL_*$.
For $h$ small, every marked zero in $[0,H]$ is associated to one
of the at most $\lfloor H/L_*\rfloor+3$ consecutive crossings
whose distance from this interval is at most $R\sqrt h/\ell$.
As $\phi$ varies, the phase offset ranges over the whole circle.
If every such offset supplies a marked zero, the necessary
condition in Lemma~\ref{p10:sp16:lem:phase} says those consecutive
rotation points form a $Kh$-net.  Translation of the starting
index does not change their covering radius.  Therefore
$Q_r(Kh)\le\lfloor H/L_*\rfloor+3$, proving the lower bound
by taking the infimum over admissible $H$.
\end{proof}

\begin{lemma}[Classical rotation covering criterion]
\label{p10:sp16:lem:rotation}
For irrational $r$, $Q_r(\eta)\ge1/(2\eta)$.
Moreover, $Q_r(\eta)=O_r(\eta^{-1})$ as $\eta\downarrow0$
if and only if the continued-fraction partial quotients of $r$
are bounded.
\end{lemma}
\begin{proof}
The union of $N$ arcs of radius $\eta$ has length at most
$2N\eta$, which proves the lower bound.  Let $p_j/q_j$ denote
the continued-fraction convergents.  We use their elementary
identities
\[
 q_{j+1}=a_{j+1}q_j+q_{j-1},\qquad
 \left|r-\frac{p_j}{q_j}\right|<\frac1{q_jq_{j+1}},
 \qquad (p_j,q_j)=1.
\]
For completeness the approximation inequality follows by writing
$r=(p_j\xi+p_{j-1})/(q_j\xi+q_{j-1})$, where the next complete
quotient satisfies $\xi>a_{j+1}$, and using the determinant
$|p_jq_{j-1}-p_{j-1}q_j|=1$.

If all $a_j\le A$, then $q_{j+1}\le(A+1)q_j$.
The first $q_j$ orbit points lie within $1/q_{j+1}$ of the
uniform $q_j$-point grid, since multiplication by $p_j$ permutes
the residue classes modulo $q_j$.  Hence
\[
            \mathcal R_{q_j}(r)\le\frac1{2q_j}
                                   +\frac1{q_{j+1}}
                              \le\frac3{2q_j}.
\]
Choose the first $q_j\ge3/(2\eta)$.  Its size is at most
$3(A+1)/(2\eta)$, up to the immaterial initial convergents,
proving the required upper bound.

If the partial quotients are unbounded, some ratios
$q_{j+1}/q_j$ tend to infinity.  Put
$N_j=\lfloor q_{j+1}/8\rfloor$.  For $0\le n<N_j$,
\[
 \left\|nr-\frac{np_j}{q_j}\right\|
                          <\frac1{8q_j}.
\]
Thus all these points lie within $1/(8q_j)$ of the uniform
$q_j$-grid.  The midpoint of any grid gap has distance at least
$3/(8q_j)$ from the orbit segment.  For
$\eta_j=1/(4q_j)$ it follows that $Q_r(\eta_j)>N_j$, and
$\eta_jQ_r(\eta_j)\to\infty$ along the chosen subsequence.
This excludes a bound of order $\eta^{-1}$.
\end{proof}

Combining the lemma with \eqref{p10:sp16:eq:coverbridge} proves the
remaining assertions of Theorem~\ref{p10:sp16:thm:height}.
In particular, if $hH_r(h)$ were bounded for all sufficiently
small $h$, the lower bridge would imply a bound for
$\eta Q_r(\eta)$ at every small $\eta$ by setting $h=\eta/K$.
The unbounded-partial-quotient case therefore gives the stated
limsup, not merely failure of one proposed constant.

These two investigations extend the earlier finite-dimensional
theory in different directions: a sharp inverse for unrestricted
positive spectra, and an exact criterion for optimal-height marked
zero detection in the one-frequency model.  Higher-dimensional
height optimality, recovery from incomplete unmarked divisors, and
joint phase-frequency estimation from finitely many noisy zeros
remain open here.


% ===== bessel =====

\section{Simultaneous gamma factors and reconstruction from boundary response}
\label{p10:be16:section}

The earlier endpoint extraction theorem involved one gamma factor.
We now classify simultaneous extraction at both original gamma rates,
including the atomic face where the total possible shape is attained.
The curved face has a qualitatively different property: none of its
remainders has any nontrivial convolution root.  A second investigation
turns factor geometry into an inverse problem.  The infinitesimal cost
of exchanging gamma factors at varying rates determines the entire
boundary remainder.  These are internally formulated transfer problems;
they are not claimed as additional named conjectures in
\cite{p01:rare:Baricz2026}, or as a priority determination for all related
factorization literature.

\subsection{The complete simultaneous extraction diagram}

Fix $\kappa>1/2$, $b>1$, put $\lambda=b-1$ and $M=2\kappa-1$,
and let $X$ have the previously studied McKay transform
\begin{equation}\label{p10:be16:original}
 L_X(s)=\frac{b+s}{b}
       \left\{\frac{b^2-1}{(b+s)^2-1}\right\}^{\kappa}.
\end{equation}
Gamma variables below are parametrized by shape and rate.  We ask for
all $h_-,h_+\geq0$ for which independent nonnegative variables satisfy
\begin{equation}\label{p10:be16:factorization}
 X\stackrel{\rm d}=G_{h_-,\lambda}+G_{h_+,\lambda+2}+R.
\end{equation}
A shape zero means the constant zero.  Define residual exponents
$a=\kappa-h_-$ and $d=\kappa-h_+$.  A prospective remainder has transform
\begin{equation}\label{p10:be16:remainder}
 L_R(s)=\frac{F_{a,d}(\lambda+s)}{F_{a,d}(\lambda)},\qquad
 F_{a,d}(u)=(u+1)u^{-a}(u+2)^{-d}.
\end{equation}
In particular, $d$ need not be positive: a fast-rate factor can have
shape larger than $\kappa$.

We first recall the scalar root from the endpoint theorem, with the
details needed for its extension and differentiation.  For
$\eta=a+d-1>0$ set
\begin{align}
 g_{a,d}(x)&=\frac{x^{\eta-1}e^{-2x}}{\Gamma(\eta)}
                 B_{a,d}(2x),\label{p10:be16:g}\\
 B_{a,d}(z)&={}_1F_1(a-1;\eta;z)
          +\frac{z}{2\eta}{}_1F_1(a;\eta+1;z).
                                                        \label{p10:be16:B}
\end{align}
Its Laplace transform is $F_{a,d}(u)$.  This follows by termwise
integration for large $u$ and Kummer's transformation, and then by
uniqueness.  The series has at most exponential growth and is real
analytic for $x>0$.  The Kummer identities used here and below are
recorded in \cite{p01:rare:DLMF}.

\begin{lemma}[The scalar ingredient and its exact slope]
\label{p10:be16:scalar}
For every $d>1/2$ there is a unique number
\[
 A(d)\in\bigl(\max(0,1-d),\,1/(4d)\bigr)
\]
such that $g_{a,d}\geq0$ exactly when $a\geq A(d)$, among
$a>0$ with $a+d>1$.  It is the unique root in the displayed interval of
\begin{equation}\label{p10:be16:root}
                 B_{a,d}\bigl(2(d-a)\bigr)=0.
\end{equation}
At $a=A(d)$ the density $g=g_{A(d),d}$ has precisely one zero,
at $x_*=d-A(d)$, and this zero has order two.  The function $A$ is
real analytic on $(1/2,\infty)$ and
\begin{equation}\label{p10:be16:slope}
 A'(d)=-\frac{\displaystyle\int_0^{x_*}
                  e^{-2y}g(x_*-y)\,dy/y}
                 {\displaystyle\int_0^{x_*}g(x_*-y)\,dy/y},
 \qquad -1<A'(d)<-e^{-2x_*}<0.
\end{equation}
The root and one-point criterion are previous inputs.  The analytic
tradeoff and formula \eqref{p10:be16:slope} are established here.
\end{lemma}
\begin{proof}
For completeness, the one-point argument is short.  When
$0<a<1$, put $p=1-a$ and $y(z)={}_1F_1(-p;\eta;z)$.
Then $y'(z)<0$ for $z>0$, since
$y'=-p\,{}_1F_1(a;\eta+1;z)/\eta$.
Kummer's differential equation gives the full identity
\[
 B'=\frac B2+\frac{2(d-a)-z}{4p}y'.
\]
Thus a zero before $2(d-a)$ crosses downward, whereas a zero after it
crosses upward.  A zero at that point satisfies
$B''=-y'/(4p)>0$.  If $d\leq a$, downward crossing is impossible
and $B$ is positive.  If $d>a$, positivity at both ends implies
that $B\geq0$ everywhere exactly when $B(2(d-a))\geq0$;
equality leaves just the stated double zero.  For $a\geq1$, both
terms in \eqref{p10:be16:B} are positive, so there is no obstruction.

Here positivity at infinity, uniformly on compact subsets of
$a,d>0$, $a+d>1$, follows from an elementary convolution.  Put
\[
 k_{a,d}(x)=\frac1{\Gamma(a)\Gamma(d)}
       \int_0^x e^{-2v}v^{d-1}(x-v)^{a-1}\,dv.
\]
Since $k_{a,d}(0)=0$, $g_{a,d}=k_{a,d}'+k_{a,d}$.  Splitting
the integral smoothly near $v=x/2$ and expanding on its first part
gives, also after differentiation,
\begin{equation}\label{p10:be16:ends}
 g_{a,d}(x)\sim\frac{x^{a+d-2}}{\Gamma(a+d-1)}\quad(x\downarrow0),
 \qquad
 g_{a,d}(x)\sim\frac{2^{-d}x^{a-1}}{\Gamma(a)}\quad(x\to\infty).
\end{equation}
On the other part, use $x-v$ as variable; its contribution and fixed
derivatives are exponentially small.  These estimates are uniform
on compact strict parameter sets.  Strict positivity everywhere is
therefore open in those parameters.

The logarithmic derivative of \eqref{p10:be16:remainder} has signed
L\'evy kernel
\begin{equation}\label{p10:be16:levy}
 \frac{e^{-(\lambda+1)x}}x\{ae^x+de^{-x}-1\}.
\end{equation}
At $a=1/(4d)$, $d>1/2$, this is nonnegative and has just one
positive zero.  The associated ID density is strictly positive
everywhere: split its jumps at a small $\epsilon$, use the positive
two-jump density on $(2\epsilon,\infty)$, and use the fact that the
small-jump sum has mean tending to zero.  This gives a positive lower
bound near any prescribed positive point.  Endpoint stability now
allows a strict decrease of $a$.

The lower endpoint is inadmissible with a nonempty inadmissible
neighborhood.  For $1/2<d<1$, at $a=1-d$ the transform expands as
$F_{a,d}(u)=1+(1-2d)/u+O(u^{-2})$ and increases for large $u$.
For $d=1$, $a=0$ gives $1-1/(u+2)$.  For $d>1$, at $a=0$ the
inverse on $x>0$ is
\[
 \frac{e^{-2x}x^{d-2}}{\Gamma(d-1)}
                       \left(1-\frac{x}{d-1}\right),
\]
which is negative on an interval.  Continuity preserves the
corresponding obstruction.  Increasing $a$ convolves the inverse with
$x^{\epsilon-1}/\Gamma(\epsilon)$; after tilting this is gamma
convolution.  Consequently the admissible set is an upper interval,
its lower endpoint is attained by the continuity theorem for Laplace
transforms, and above that endpoint its density is strictly positive.
The one-point test proves existence, uniqueness and the double-zero
description of $A(d)$.

At a contact point, differentiating the preceding convolution at
$\epsilon=0$ gives
\begin{align*}
 \partial_a g_{a,d}(x_*)&=\int_0^{x_*}g_{a,d}(x_*-y)\,dy/y,\\
 \partial_d g_{a,d}(x_*)&=\int_0^{x_*}e^{-2y}g_{a,d}(x_*-y)\,dy/y.
\end{align*}
The integrals converge: the double zero gives $O(y^2)$ at $y=0$,
and \eqref{p10:be16:ends} is integrable at the other endpoint.  To see
the derivative directly, divide the convolution by $\epsilon$ and
use $1/\Gamma(\epsilon)\sim\epsilon$; its constant term is zero
at the contact.  Analytic dependence in \eqref{p10:be16:g} makes this
one-sided calculation the ordinary derivative.
The first integral is positive.  Since $g_x(x_*)=0$, the implicit
function theorem applied to $g_{a,d}(d-a)=0$ proves analyticity and
\eqref{p10:be16:slope}.  The strict bounds follow because the quotient
is an average of $e^{-2y}$ over $0<y<x_*$.
\end{proof}

Extend $A$ to $d\leq1/2$ by $A(d)=1-d$, and define
\begin{equation}\label{p10:be16:ID-threshold}
 I(d)=\begin{cases}1-d,&d\leq1/2,\\[2pt]1/(4d),&d>1/2.\end{cases}
\end{equation}

\begin{theorem}[Two-rate extraction, including its atomic face]
\label{p10:be16:diagram}
For $h_-,h_+\geq0$, the factorization \eqref{p10:be16:factorization}
exists if and only if
\begin{equation}\label{p10:be16:P-region}
                  \kappa-h_-\geq A(\kappa-h_+).
\end{equation}
Its remainder is ID if and only if the same inequality holds with
$A$ replaced by $I$.  Equivalently, for $0\leq h_+\leq M$, the
largest admissible slow-rate shape is
\begin{equation}\label{p10:be16:frontier}
 h_-^{\max}(h_+)=
 \begin{cases}
  \kappa-A(\kappa-h_+),&0\leq h_+<M/2,\\
  M-h_+,&M/2\leq h_+\leq M.
 \end{cases}
\end{equation}
The curved part is analytic, with derivative strictly between $-1$
and $0$.  It joins the line segment continuously with slope $-1$.

The total-shape bound $h_-+h_+\leq M$ is attained precisely on the
line segment in \eqref{p10:be16:frontier}.  There the remainder is
compound Poisson and has atom
\begin{equation}\label{p10:be16:atom}
 \mathbb P(R=0)=\frac{\lambda^a(\lambda+2)^{1-a}}{\lambda+1},
 \qquad a=\kappa-h_-\geq1/2.
\end{equation}
On the curved part the remainder has no atom, is not ID, and its
density has exactly one interior double zero.
\end{theorem}
\begin{proof}
Since the original transform has order $s^{-M}$ at infinity,
every feasible pair satisfies $h_-+h_+\leq M$, or $a+d\geq1$.
Also $a\geq0$: otherwise the moment generating function of the
remainder, computed from \eqref{p10:be16:factorization}, tends to zero
as its argument increases to $\lambda$, although it must be at
least one.  The case $a=0$, $d\geq1$, is excluded by the explicit
negative inverse or increasing transform in the preceding proof.
Thus $a>0$ is necessary.

If $a+d>1$ and $d>1/2$, Lemma~\ref{p10:be16:scalar} gives exactly
$a\geq A(d)$.  If $d\leq1/2$, then $a+d\geq1$ forces
$a\geq1-d\geq d$.  The function $ae^x+de^{-x}-1$ is
nonnegative for $x\geq0$: its derivative is positive for $x>0$
when $d\leq a$, and its value at zero is $a+d-1$.
This argument permits negative $d$ and proves ID sufficiency.

On the boundary $a+d=1$, the expansion
$F_{a,d}(u)=1+(2a-1)/u+O(u^{-2})$ excludes $a<1/2$.
For $a\geq1/2$ the preceding nonnegative kernel applies.  Its
value at zero vanishes, so the L\'evy density in
\eqref{p10:be16:levy} is integrable at zero and at infinity.
The law is compound Poisson, and the limit of its Laplace transform
at infinity is \eqref{p10:be16:atom}.

For the full ID criterion minimize $ae^x+de^{-x}$ over $x\geq0$.
If $d\leq a$, its minimum is $a+d$; if $d>a>0$, its minimum is
$2\sqrt{ad}$.  Together with $a+d\geq1$ these are exactly
$a\geq I(d)$.  Necessity follows from uniqueness of the L\'evy
representation, since \eqref{p10:be16:levy} is integrable with the
usual L\'evy weight and is the inverse logarithmic derivative.

Translating the parameter inequalities gives
\eqref{p10:be16:frontier}.  Its curved derivative is $A'(\kappa-h_+)$.
As $d\downarrow1/2$, the two bounds on $A(d)$ squeeze it to
$1/2$, hence $x_*=d-A(d)\to0$.  Formula \eqref{p10:be16:slope}
then squeezes $A'(d)$ to $-1$.  The remaining assertions follow
from the strict scalar inequalities and the double-zero result.
\end{proof}

\subsection{Root obstruction on the curved factor boundary}

The distinction between the two boundary faces is stronger than
ID versus non-ID.  Here a convolution root means a probability law
$Q$ with $Q^{*n}=\mathcal L(R)$ for an integer $n\geq2$.

\begin{theorem}[No gamma factors and no convolution roots]
\label{p10:be16:rootless}
Every remainder on the curved boundary of
Theorem~\ref{p10:be16:diagram} has no nonzero gamma factor at any finite
positive rate, and has no convolution root of any integer order
$n\geq2$.  This assertion includes possible roots on the whole
real line.  In contrast, all remainders on the linear face have
convolution roots of every order, since they are compound Poisson.
\end{theorem}
\begin{proof}
Let $f(x)=e^{-\lambda x}g_{a,d}(x)/F_{a,d}(\lambda)$ be a curved
boundary density, let $x_*$ be its unique zero, and put
$\eta=a+d-1>0$.  If $R=G_{h,\rho}+Y$, with $h>0$, then $Y$
has support arbitrarily close to zero because $R$ does.  The part of
its convolution with the gamma density coming from
$Y\leq x_*/2$ has a strictly positive lower bound throughout a
neighborhood of $x_*$.  This contradicts continuity and $f(x_*)=0$.
This argument is valid for every $0<\rho<\infty$; it makes no
claim about a formal gamma factor at an infinite rate.

We next verify the absolute continuity needed for the root argument.
For any $t>0$, the only candidate Laplace transform for a positive
convolution power is $L_R(s)^t$.  At large $u=s+\lambda$ write
\[
 (1+z)^t(1+2z)^{-dt}=\sum_{j\geq0}c_j(t)z^j.
\]
Its candidate inverse is
\begin{equation}\label{p10:be16:power-inverse}
 q_t(x)=F_{a,d}(\lambda)^{-t}e^{-\lambda x}
       \sum_{j\geq0}\frac{c_j(t)x^{t\eta+j-1}}
                              {\Gamma(t\eta+j)}.
\end{equation}
Cauchy's bound on $c_j$ proves local convergence of this series and
of its derivatives on $x>0$, and gives at most exponential growth.
Termwise integration for sufficiently large $s$ gives $L_R(s)^t$.
If this transform belongs to a probability law, uniqueness after
a sufficiently large exponential tilt forces that law to have
density \eqref{p10:be16:power-inverse}.  Thus every putative positive
power has a real analytic density, with no atom, and a positive
leading term at zero.

If a root of order $n$ existed, its support would lie in
$[0,\infty)$: any positive mass on $(-\infty,0)$ would give
positive mass there for the sum of $n$ independent copies.
Its Laplace transform must be $L_R^{1/n}$, so the preceding density
is nonnegative and real analytic.  It is nonzero and hence positive
except at isolated points on $(0,\infty)$.  Its two-fold convolution
is strictly positive at every $x>0$, by integration over $(0,x)$;
the endpoint singularities are integrable since $\eta/n>0$.
The same is true of every higher convolution power.  The resulting
continuous density cannot equal $f$, which vanishes at $x_*$.
This is a contradiction for every $n\geq2$.
\end{proof}

\subsection{The inverse theorem for gamma-rate response}

Fix a curved boundary point $a=A(d)$ with $d>1/2$, and denote
its remainder transform by $L_R$ and density by $f$.  The original
slow-rate factor has positive shape.  It can therefore be reduced
slightly, adding a gamma component back to the remainder, while a
small new gamma factor at another rate is extracted.  For a fixed
$\rho\geq\lambda$ define the candidate transform
\begin{equation}\label{p10:be16:exchange}
 L_{t,\epsilon,\rho}(s)
  =L_R(s)\left(\frac{\lambda}{\lambda+s}\right)^t
               \left(\frac{\rho+s}{\rho}\right)^\epsilon.
\end{equation}
Here $\epsilon\geq0$ is the newly extracted shape, and $t\geq0$
is the amount returned at the slow rate.  This describes an actual
exchange of factors of $X$ for sufficiently small parameters.

\begin{theorem}[The boundary response determines the remainder]
\label{p10:be16:inverse-response}
For each finite $\rho\geq\lambda$ there is a real analytic local
boundary $t=T_\rho(\epsilon)$ with $T_\rho(0)=0$ such that
\eqref{p10:be16:exchange} is a probability Laplace transform locally
exactly when $t\geq T_\rho(\epsilon)$.  Its derivative is
\begin{equation}\label{p10:be16:response}
 T_\rho'(0)=
 \frac{\displaystyle\int_0^{x_*}e^{-\rho y}f(x_*-y)\,dy/y}
      {\displaystyle\int_0^{x_*}e^{-\lambda y}f(x_*-y)\,dy/y}.
\end{equation}
Write $C(r)=T_{\lambda+r}'(0)$ for $r\geq0$.  Then $C$ is the
Laplace transform of the probability density
\begin{equation}\label{p10:be16:response-density}
 w(y)=\frac{g_{a,d}(x_*-y)}{yJ},\qquad
 J=\int_0^{x_*}g_{a,d}(x_*-y)\,dy/y,
 \qquad 0<y<x_*.
\end{equation}
In particular $C$ is strictly completely monotone, $C(0)=1$, and
\begin{align}
 x_*&=\lim_{n\to\infty}
                \{(-1)^nC^{(n)}(0)\}^{1/n},\label{p10:be16:recover-zero}\\
 C(r)&\sim\frac{g_{a,d}''(x_*)}{2J}\,r^{-2}
                       \qquad(r\to\infty).\label{p10:be16:high-rate}
\end{align}
For the fixed known base rate $\lambda$, the entire response function
$C$ uniquely determines the remainder law.  Within the present family
it even recovers both residual
exponents from the endpoint behavior of its uniquely determined
density $w$:
\begin{equation}\label{p10:be16:recover-parameters}
 \eta=1+\lim_{z\downarrow0}
            \frac{\log w(x_*-z)}{\log z},\qquad
 a=\frac{\eta+1-x_*}{2},\qquad d=\frac{\eta+1+x_*}{2}.
\end{equation}
These are exact-data uniqueness assertions, not numerical stability
claims for analytic continuation or Laplace inversion.
\end{theorem}
\begin{proof}
We first establish that the local minimum controls positivity on
the whole half-line.  Put $r=\rho-\lambda$.  Apart from its positive
normalization, the inverse needed in \eqref{p10:be16:exchange} is that of
\[
 F_{a,d}(u)u^{-t}(u+r)^\epsilon.
\]
At large $u$, expand
$(1+z)(1+2z)^{-d}(1+rz)^\epsilon$ into a Taylor series and invert
termwise, with gamma denominator of initial order
$\eta+t-\epsilon>0$.  The same geometric-coefficient argument as
in \eqref{p10:be16:power-inverse} gives a jointly real analytic inverse
on $x>0$, depending analytically on $t,\epsilon$ near zero.
Its small-$x$ leading coefficient is positive and has power
$x^{\eta+t-\epsilon-1}$.

For $r>0$, its large-$x$ leading term is
\begin{equation}\label{p10:be16:exchange-tail}
 \frac{2^{-d}r^\epsilon}{\Gamma(a+t)}x^{a+t-1}.
\end{equation}
To justify uniformity in a neighborhood of $(t,\epsilon)=(0,0)$,
start with $g_{a+t,d}$ and convolve it with
$e^{-rx}x^{-\epsilon}/\Gamma(1-\epsilon)$, then apply $D+r$.
This has the required transform.  The initial value before
differentiation is zero since $\eta+t-\epsilon>0$.
Splitting the convolution at half its argument shows that its
leading factor is $r^{\epsilon-1}$; the part near the remote
endpoint and its derivative are exponentially small.  Applying
$D+r$ gives \eqref{p10:be16:exchange-tail}, with a uniform relative
$O(x^{-1})$ error.  For $r=0$ the inverse is exactly
$g_{a+t-\epsilon,d}$ and \eqref{p10:be16:ends} applies directly.
Thus both ends remain strictly positive, uniformly for nearby
parameters.  The normalized inverses are integrable after tilting
by $e^{-\lambda x}$ and have total mass one.

At the original density zero, adding a gamma shape of size $t$ at
rate $\lambda$ gives derivative
\[
 N_\lambda=\int_0^{x_*}e^{-\lambda y}f(x_*-y)\,dy/y>0.
\]
The same convolution calculation as in Lemma~\ref{p10:be16:scalar}
justifies the derivative; normalizing terms multiply $f(x_*)=0$.
Extracting shape $\epsilon$ at rate $\rho$ gives derivative
$-N_\rho$, where $N_\rho$ is the corresponding integral in
\eqref{p10:be16:response}.  This follows either by differentiating
the inverse multiplier or by analytically continuing the
gamma-convolution identity from negative $\epsilon$.

Since $f''(x_*)>0$, the implicit function theorem first supplies a
unique critical point near $x_*$ for the perturbed density.  Its
value $m(t,\epsilon)$ is analytic, with
$m_t(0,0)=N_\lambda$ and $m_\epsilon(0,0)=-N_\rho$.
A second application gives the unique local graph $m=0$ and
\eqref{p10:be16:response}.  Outside a fixed small neighborhood of
$x_*$, positivity persists by the endpoint estimates and positivity
on the remaining compact interval.  Inside that neighborhood the
critical point is the unique minimum.  Thus $m\geq0$ is exactly
global nonnegativity, proving the claimed local factor boundary.

Substitution of $f(x)=e^{-\lambda x}g_{a,d}(x)/F_{a,d}(\lambda)$
in \eqref{p10:be16:response} gives
$C(r)=\int_0^{x_*}e^{-ry}w(y)\,dy$.  This density is positive on
$(0,x_*)$, integrable at both ends, and has total mass one.
The strict derivative inequalities and
\eqref{p10:be16:recover-zero} now follow from differentiation and the
limit of the $L^n$ norms of a random variable whose essential
supremum is $x_*$.  At zero,
\[
 w(y)=\frac{g_{a,d}''(x_*)}{2J}y+O(y^2),
\]
which gives \eqref{p10:be16:high-rate} by a direct Laplace endpoint
expansion.  The singularity of $w$ at $x_*$ is integrable and
contributes only an exponentially small error to that expansion.

Laplace-transform uniqueness determines $w$ from $C$, and the
moment limit determines its right support endpoint.  Consequently
\eqref{p10:be16:response-density} recovers $g_{a,d}$ up to a constant
on $(0,x_*)$.  Real analyticity extends that determination to
$(0,\infty)$, and normalization after the known exponential tilt
fixes the constant.  More directly in the present family,
\eqref{p10:be16:ends} gives
$w(x_*-z)\sim z^{\eta-1}/\{x_*J\Gamma(\eta)\}$.
This proves the first identity in \eqref{p10:be16:recover-parameters};
the other two follow from $a+d=\eta+1$ and $d-a=x_*$.
\end{proof}

At the fast original rate, $\rho=\lambda+2$, the response theorem
reduces to $T_{\lambda+2}'(0)=-A'(d)$, so it includes the tradeoff
slope of the two-rate diagram.  Allowing all rates supplies strictly
more information: it reveals a whole compactly supported measure,
rather than a single directional derivative.  The linear atomic
face is excluded from this inverse theorem; it has no positive
interior contact to which this argument could apply.

The simultaneous classification and the response inverse theorem
are the two principal investigations of this section.  Root failure
is a structural consequence of the first, not an extra historical
conjecture.  Three-or-more-rate global regions, convexity of the
unrestricted extraction set, and quantitative stability of recovery
from finitely many noisy response values remain open here.

\section{Where the next obstructions lie}
\label{p10:r16:remaining}

The results identify several precise next steps. In continuous
information, the missing one-query gaps and the unrestricted
three-dimensional two-query case are not settled by the width recursion.
Smooth exact zero-set tests cannot simply be replaced by analytic ones.
For noisy convex information, a matching upper construction is needed
before the displayed budget gaps can be called sharp rates. For Banach
operators, the sequence-space example shows that compactness alone is
the wrong general criterion, but it does not identify the complete
operator class with vanishing adaptive error.

For mixtures, exact replication is much stronger than a collection of
separate marginals. The two-replicate stable inversion uses a positive
lower index bound and a triangular dilation recursion; it leaves the
case of indices accumulating at zero, sharp statistical stability,
and minimal replication for correlated location and skewness open.
Honest uncertainty under corruption and adaptive point estimation are
different questions. The confidence-set obstruction proved above should not
be read as an impossibility theorem for every adaptive estimator.

For spectra, infinitely many exact zeros can carry an unknown
frequency support while finite noisy records remain uniformly
ambiguous near a rational relation. A useful quantitative extension
must therefore state an arithmetic margin, not only pairwise
frequency separation. The marked-zero height theorem resolves one
frequency through circle rotations; multidimensional sharp covering
criteria require additional Diophantine analysis. Finally, the gamma
response inverse is exact but uses a full continuum of rates. A stable
finite-rate analogue and global regions for three or more rates are
not supplied here.

These remaining questions are recorded to prevent a proved restricted
statement from silently becoming a stronger conjecture. The intended
next research step is to attack those explicit obstructions, rather
than to count further special cases of the present theorems as
independent discoveries.
\endgroup
% ===== END SOURCE 10 =====

\clearpage\phantomsection
\addcontentsline{toc}{part}{Consolidated references}
\begin{thebibliography}{999}\small\raggedright
\bibitem{p01:KEIOIntegrated}
KEIO, \emph{Observation and Recovery in Arithmetic and Probability:
Prime Flows, Rare Events, and Infinite Mixtures}, integrated working
manuscript, 14 October 2026, 223 pp. AI-assisted, unpublished, and not
externally peer reviewed. This is the external predecessor of the present
compilation. Source: \path{KEIO_Integrated_Observation_and_Recovery.tex}.
Part~X cites its section on factorial smoothing and the sharp logarithmic modulus.

\bibitem{p01:Rainer2022}
A. Rainer, \emph{Ultradifferentiable extension theorems: a survey},
author manuscript dated 30 December 2021, Theorem 3.6 and Corollary 3.12.
\url{https://www.tuwien.at/mg/arainer/wp-content/uploads/2025/03/survey-uet.pdf}.

\bibitem{p01:rare:Baricz2026}
\'A. Baricz, Dhivya Prabhu K, S. Singh, and Antony Vijesh V,
\emph{Infinitely divisible modified Bessel distributions},
Pacific J. Math. \textbf{343} (2026), 261--313.
\href{https://doi.org/10.2140/pjm.2026.343.261}{doi:10.2140/pjm.2026.343.261}.
Preprint: \href{https://arxiv.org/abs/2406.17721}{arXiv:2406.17721}.
Published text: \url{https://msp.org/pjm/2026/343-2/pjm-v343-n2-p01-p.pdf}.
The maximal ID-domain question is in Section~2.3; the multipart McKay subclass problem is in Section~4.3, pp.~282--283 (locators retained from Part~III).

\bibitem{p01:rare:BondessonSimon2016}
L. Bondesson and T. Simon,
\emph{Stieltjes functions of finite order and hyperbolic monotonicity},
Trans. Amer. Math. Soc. \textbf{370} (2018), 4201--4222.
\href{https://doi.org/10.1090/tran/7123}{doi:10.1090/tran/7123}.
\href{https://arxiv.org/abs/1604.05267}{arXiv:1604.05267}.
Part~I specifically cites equation~(7) of the 2016 preprint.

\bibitem{p01:rare:Salazar2026}
D. S. P. Salazar,
\emph{A Bessel-zero obstruction to hyperbolic complete monotonicity of
noncentral chi-square densities},
\href{https://arxiv.org/abs/2606.22066}{arXiv:2606.22066} (2026).

\bibitem{p01:rare:DLMF}
F. W. J. Olver et al., eds., \emph{NIST Digital Library of Mathematical Functions},
\url{https://dlmf.nist.gov/}. The consolidated citations refer to Sections~5.11, 10.9, 10.21, 10.27, 10.29, 10.32, 10.39, 10.40, 10.42, 13.2--13.4, 13.7, 15.4, 15.8, and~15.10.
Specific equation locators retained from the source manuscripts:
\href{https://dlmf.nist.gov/5.11.E9}{5.11.9};
\href{https://dlmf.nist.gov/10.9.E30}{10.9.30 (Nicholson's integral)};
\href{https://dlmf.nist.gov/10.21.E15}{10.21.15 (the infinite product for $J_\nu$)};
and \href{https://dlmf.nist.gov/15.8.E1}{15.8.1 (Euler's hypergeometric transformation)}.
Parts~I--III record access on 14 October 2026.

\bibitem{p01:mix:Lin2015}
J. Lin,
\emph{A negative answer to a problem of Aldous on determination of
exchangeable sequences},
\href{https://arxiv.org/abs/1507.00678v1}{arXiv:1507.00678v1} (2015).
The open normal-class question audited here occurs in Section~2.6,
printed p.~17 (PDF page index~16); the exponential-moment result is
Proposition~2.6.4. Published in \emph{Electronic Journal of Probability}
\textbf{21} (2016), paper 30, 1--26,
\href{https://doi.org/10.1214/16-EJP4403}{doi:10.1214/16-EJP4403}.
The publisher PDF was not inspected in this audit.

\bibitem{p01:mix:Le2026}
D. Le, N. Bariletto, A. Rinaldo, and N. Ho,
\emph{Partial Differential Equation Barriers to Identifiability in
Infinite Mixture Models},
\href{https://arxiv.org/abs/2608.08597}{arXiv:2608.08597} (2026).
Question and definition locations refer to the
\href{https://nhatptnk8912.github.io/Identifiability_Infinite_Mixture_final.pdf}{author manuscript}
dated 8 August 2026: the remark following Proposition~6 and Section~6,
printed p.~13 (PDF page index~12); operator classes in Definitions~1
and~4.

\bibitem{p01:info:GNR2015}
M. B. Giles, T. Nagapetyan and K. Ritter,
Multilevel Monte Carlo approximation of distribution functions and densities,
SIAM/ASA Journal on Uncertainty Quantification 3 (2015), 267--295.
doi:10.1137/140960086.
\url{https://people.maths.ox.ac.uk/gilesm/files/juq_15.pdf}.

\bibitem{p01:info:Giles2018}
M. B. Giles, MLMC for nested expectations,
in Contemporary Computational Mathematics: A Celebration of the 80th
Birthday of Ian Sloan, Springer (2018), 425--442.
\url{https://people.maths.ox.ac.uk/gilesm/files/SLOAN80-056.pdf}.

\bibitem{p01:info:LeeValiant2021}
J. C. H. Lee and P. Valiant,
Uncertainty about uncertainty: Optimal adaptive algorithms for estimating
mixtures of unknown coins,
Proceedings of the 2021 ACM-SIAM Symposium on Discrete Algorithms,
414--433. doi:10.1137/1.9781611976465.25.
\url{https://www.ias.edu/sites/default/files/LeeValiant_Uncertainty.pdf}.
Preprint: \url{https://arxiv.org/abs/1904.09228}.

\bibitem{p01:info:Lee2021}
C. H. J. Lee, Classic and Modern Challenges in Statistical Estimation,
Ph.D. thesis, Brown University, May 2021, Chapter 4, especially pp.~65--67.
\url{https://cs.brown.edu/research/pubs/theses/phd/2021/chun.hin.jasper.lee.pdf}.

\bibitem{p01:info:VKVK2019}
R.~K.~Vinayak, W.~Kong, G.~Valiant and S.~Kakade,
\emph{Maximum likelihood estimation for learning populations of parameters},
Proceedings of the 36th International Conference on Machine Learning,
PMLR \textbf{97} (2019), 6448--6457.
\url{https://proceedings.mlr.press/v97/vinayak19a/vinayak19a.pdf}.
Theorem~3.2, Remark~3.2 and Theorem~3.3 appear on printed paper page~4;
Lemma~4.2 and its following discussion appear on printed page~5.
Remarks~3.2--3.3 distinguish the published replication-range MLE conjecture from the suggested localized moment-matching route; Theorem~3.3 gives the minimax lower bound (locators retained from Parts~V--VI).
Part~VII additionally cites the Supplement, Remark~4.1, equation~(23): \url{https://proceedings.mlr.press/v97/vinayak19a.html}.

\bibitem{p01:info:MMRS2025}
C.~Musco, C.~Musco, L.~Rosenblatt and A.~V.~Singh,
\emph{Sharper bounds for Chebyshev moment matching, with applications},
Proceedings of the 38th Conference on Learning Theory,
PMLR \textbf{291} (2025), 4309--4358;
author version arXiv:2408.12385v3, revised 18 May 2026.
\url{https://arxiv.org/html/2408.12385v3}.
Theorem~20 and Section~6.2 are the version-specific comparison points.
Proceedings record: \url{https://proceedings.mlr.press/v291/musco25a.html}.
The cited version's Section~6.2 records the coefficient and likelihood questions (Parts~V--VI).

\bibitem{p01:info:LauScarlett2026}
I. Lau and J. Scarlett,
Open problem: Is interaction necessary for order-optimal 1-bit mean estimation?,
Proceedings of Machine Learning Research 336 (2026), 7123--7128.
\url{https://proceedings.mlr.press/v336/lau26a.html}.

\bibitem{p01:info:LeeKennedy2026}
J. Lee and E. H. Kennedy,
Minimax rates for learning smooth populations of parameters,
arXiv:2610.04135v2, 9 October 2026.
\url{https://arxiv.org/abs/2610.04135}.
Part~IV uses the earlier version arXiv:2610.04135v1, 2 October 2026, Theorem~1 and Remark~1, for the distinct smooth-density target: \url{https://arxiv.org/html/2610.04135v1}.

\bibitem{p01:mom:BEW2022}
V. B\"urgin, J. Epperlein, and F. Wirth,
\emph{Remarks on the tail order on moment sequences},
Journal of Mathematical Analysis and Applications \textbf{512} (2022),
no.~1, 126135.
\href{https://doi.org/10.1016/j.jmaa.2022.126135}{doi:10.1016/j.jmaa.2022.126135}.
Locators refer to \href{https://arxiv.org/pdf/2104.10572v3}{arXiv:2104.10572v3}.

\bibitem{p01:mom:LK1971}
W. A. J. Luxemburg and J. Korevaar,
\emph{Entire functions and M\"untz--Sz\'asz type approximation},
Transactions of the American Mathematical Society \textbf{157} (1971),
23--37.
\href{https://doi.org/10.1090/S0002-9947-1971-0281929-0}{doi:10.1090/S0002-9947-1971-0281929-0}.
Theorem~5.2, p.~30, is used through the explicit restatement in
\cite[Theorem~21]{p01:mom:Wellershoff2021}.

\bibitem{p01:mom:Wellershoff2021}
M. Wellershoff, \emph{Injectivity of sampled Gabor phase retrieval in
spaces with general integrability conditions}, ETH Z\"urich SAM
Research Report 2021-43, 18 December 2021.
\href{https://www.sam.math.ethz.ch/sam_reports/reports_final/reports2021/2021-43.pdf}{Author-hosted report}.
Theorem~21 explicitly restates W. A. J. Luxemburg and J. Korevaar,
\emph{Entire functions and M\"untz--Sz\'asz type approximation},
Transactions of the American Mathematical Society \textbf{157}
(1971), 23--37, Theorem~5.2, p.~30.
The 2021 restatement, rather than the inaccessible original publisher
PDF, was checked for the product construction used here.

\bibitem{p01:mom:BE1995}
P. Borwein and T. Erd\'elyi,
\emph{M\"untz spaces and Remez inequalities},
Bulletin of the American Mathematical Society (N.S.) \textbf{32}
(1995), 38--42.
\href{https://arxiv.org/pdf/math/9501225}{arXiv:math/9501225}.

\bibitem{p01:pf:BCS}
I. Belovas, R. \v{C}epaityt\.{e}, and M. Sabaliauskas,
\emph{On the zero-free region and the distribution of zeros of the prime zeta function},
Analele \c{S}tiin\c{t}ifice ale Universit\u{a}\c{t}ii Ovidius Constan\c{t}a,
Seria Matematic\u{a} \textbf{33} (2025), no.~2, 27--44.
\href{https://doi.org/10.2478/auom-2025-0017}{doi:10.2478/auom-2025-0017}.
The conjectures cited here are in Section~3, p.~33;
Conjectures~2 and~3 impose both $0<\gamma<\sigma_0$ and $\gamma>1$.

\bibitem{p01:pf:SV}
J. M. Sepulcre and T. Vidal,
\emph{On the real projections of zeros of analytic almost periodic functions},
Carpathian Journal of Mathematics \textbf{38} (2022), no.~2, 489--501.
\href{https://semnul.com/carpathian/wp-content/uploads/2022/02/carpathian_2022_38_2_489_501.pdf}{Journal-hosted full text}.
Theorem~4.3 is on p.~494.

\bibitem{p01:pf:JT}
B. Jessen and H. Tornehave,
\emph{Mean motions and zeros of almost periodic functions},
Acta Mathematica \textbf{77} (1945), 137--279.
\href{https://doi.org/10.1007/BF02392225}{doi:10.1007/BF02392225}.
See the introductory Sections~25--26, pp.~165--166, and Theorems~5--7,
pp.~182--185.

\bibitem{p01:zet:AF}
L. Alp\"oge and R. Furman,
\emph{More than two thirds of the zeros of the Riemann zeta function
are simple and on the critical line}, arXiv:2608.13637v2 (2026).
\url{https://arxiv.org/html/2608.13637v2}.

\bibitem{p01:zet:FZ}
K. Ford and A. Zaharescu,
\emph{On the distribution of imaginary parts of zeros of the
Riemann zeta function}, J. reine angew. Math. 579 (2005), 145--158.
\url{https://ford126.web.illinois.edu/wwwpapers/fz.pdf}.

\bibitem{p01:zet:FSZ}
K. Ford, K. Soundararajan and A. Zaharescu,
\emph{On the distribution of imaginary parts of zeros of the
Riemann zeta function, II}, Math. Ann. 343 (2009), 487--505.
\url{https://arxiv.org/pdf/0805.2745}.

\bibitem{p01:zet:CG}
F. \c{C}i\c{c}ek and S. M. Gonek,
\emph{The uniform distribution modulo one of certain subsequences
of ordinates of zeros of the zeta function}, Math. Proc. Cambridge
Philos. Soc. 176 (2024), 593--608.
\url{https://arxiv.org/html/2310.10119v1}.

\bibitem{p01:zet:AR}
J. Arias de Reyna,
\emph{On the distribution (mod 1) of the normalized zeros of the
Riemann zeta-function}, J. Number Theory 153 (2015), 37--53.
\url{https://arxiv.org/pdf/1403.7126}.

\bibitem{comp:part:01}
KEIO, \emph{Exact Recovery, Optimal Replication, and Bessel Classification: Extensions of Observation and Recovery, with an Audit of Fifteen Published Questions}, research manuscript, 14 October 2026. Original edition: 42 pp. Reproduced as \hyperref[part:01]{Part~I} of this compilation. AI-assisted and not externally peer reviewed.

\bibitem{p02:frontmix:Haff2011}
L.~R.~Haff, P.~T.~Kim, J.-Y.~Koo, and D.~St.~P.~Richards,
Minimax estimation for mixtures of Wishart distributions,
\emph{Annals of Statistics} \textbf{39} (2011), no.~6, 3417--3440.
\href{https://doi.org/10.1214/11-AOS951}{doi:10.1214/11-AOS951};
inspected electronic reprint:
\href{https://arxiv.org/pdf/1203.3342}{arXiv:1203.3342}.

\bibitem{p02:frontmix:VS2015}
R.~A.~Vandermeulen and C.~Scott,
\emph{On the identifiability of mixture models from grouped samples},
\href{https://arxiv.org/abs/1502.06644}{arXiv:1502.06644} (2015).
Theorem~4 in the inspected 16-page preprint concerns mixtures with
finitely many components.
The revised preprint is arXiv:1502.06644v2 (2022). As noted in Part~V, the authors state that the work was subsumed and expanded in their Annals of Statistics article \emph{An operator theoretic approach to nonparametric mixture models}.

\bibitem{p02:frontmix:Yakowitz1968}
S.~J.~Yakowitz and J.~D.~Spragins,
On the identifiability of finite mixtures,
\emph{Annals of Mathematical Statistics} \textbf{39} (1968),
no.~1, 209--214.
\href{https://doi.org/10.1214/aoms/1177698520}{doi:10.1214/aoms/1177698520}.
The publisher full text was unavailable in this audit; the relevant
finite-Gaussian result is proved independently in the text.

\bibitem{p02:frontmix:BN1965}
O.~Barndorff-Nielsen,
Identifiability of mixtures of exponential families,
\emph{Journal of Mathematical Analysis and Applications}
\textbf{12} (1965), no.~1, 115--121.
\href{https://doi.org/10.1016/0022-247X(65)90059-4}{doi:10.1016/0022-247X(65)90059-4}.
Full text not inspected; included as an unresolved priority check,
not as support for an unverified theorem attribution.

\bibitem{p02:graph:BS2017}
G.~Blekherman and R.~Sinn,
\emph{Maximum likelihood threshold and generic completion rank of graphs},
Discrete \& Computational Geometry \textbf{61} (2019), 303--324.
\url{https://arxiv.org/abs/1703.07849}.
Inspected version: arXiv:1703.07849v2, 23 May 2017, \url{https://arxiv.org/html/1703.07849v2}.
Cited locators: Definition~1.1, Theorems~2.1 and~2.7, and Corollary~2.11 (Part~II). The question on bounding generic completion rank by maximum likelihood threshold is at the end of the Introduction (Part~III).

\bibitem{p02:graph:BDG2024}
D. I. Bernstein, S. Dewar, S. J. Gortler, A. Nixon, M. Sitharam,
and L. Theran, \emph{Maximum likelihood thresholds via graph rigidity},
Annals of Applied Probability \textbf{34} (2024), no.~3, 3288--3319.
\href{https://doi.org/10.1214/23-AAP2039}{doi:10.1214/23-AAP2039}.
\url{https://arxiv.org/abs/2108.02185}.
Part~III inspected author version \url{https://arxiv.org/html/2108.02185v3}.

\bibitem{p02:bud:WeedBerthet2019}
J. Weed and Q. Berthet,
Estimation of smooth densities in Wasserstein distance,
Proceedings of Machine Learning Research 99 (2019), 3118--3119
(extended abstract); full version arXiv:1902.01778.
\url{https://proceedings.mlr.press/v99/weed19a.html}.
\url{https://arxiv.org/abs/1902.01778}.

\bibitem{p02:bud:PW2012}
M. G. Pavlides and J. A. Wellner,
Nonparametric estimation of multivariate scale mixtures of uniform densities,
Journal of Multivariate Analysis 107 (2012), 71--89.
doi:10.1016/j.jmva.2012.01.001.
\url{https://sites.stat.washington.edu/jaw/JAW-papers/jaw-pavlides.jmva.12.pdf}.

\bibitem{p02:bud:KKG2025}
A. K. H. Kim, G. Kur and A. Guntuboyina,
Convergence rates for estimating multivariate scale mixtures of uniform densities,
Electronic Journal of Statistics 19 (2025), 3771--3834.
doi:10.1214/25-EJS2426.
Author version arXiv:2410.10251v2, revised 22 February 2026.
\url{https://arxiv.org/abs/2410.10251}.

\bibitem{p02:bes:BariczPrabhu2026}
\'A. Baricz and Dhivya Prabhu K,
\emph{On an infinitely divisible distribution involving modified Bessel functions},
\href{https://arxiv.org/abs/2610.09779}{arXiv:2610.09779},
7 October 2026, Theorem 1.1.

\bibitem{p02:mt:Prekopa1973}
A. Pr\'ekopa, \emph{On logarithmic concave measures and functions},
Acta Scientiarum Mathematicarum (Szeged) \textbf{34} (1973), 335--343.
\href{https://rutcor.rutgers.edu/~prekopa/SCIENT2.pdf}{Author-hosted text},
Theorem~6.

\bibitem{p02:mt:Phasebook2025}
M. Allain, S. Aslan, W. Coene, S. Dirksen, J. Dong, J. Flamant,
M. Iwen, F. Krahmer, T. van Leeuwen, O. Melnyk, A. Menzel,
A. P. Mosk, V. Nikitin, P. Salanevich, G. Plonka, and M. Wellershoff,
\emph{Phasebook: a survey of selected open problems in phase retrieval},
Sampling Theory, Signal Processing, and Data Analysis \textbf{23}
(2025), article~23.
\href{https://doi.org/10.1007/s43670-025-00117-0}{doi:10.1007/s43670-025-00117-0}.

\bibitem{p02:mt:Beinert2017}
R. Beinert, \emph{Non-negativity constraints in the one-dimensional
discrete-time phase retrieval problem}, Information and Inference
\textbf{6} (2017), 213--224.
\href{https://arxiv.org/abs/1605.05482}{arXiv:1605.05482}.

\bibitem{p02:mt:Choulli2016}
M. Choulli, \emph{Various stability estimates for the problem of
determining an initial heat distribution from a single measurement},
Rivista di Matematica della Universit\`a di Parma (N.S.)
\textbf{7} (2016), no.~2, 279--307.
\href{https://arxiv.org/abs/1512.07421}{arXiv:1512.07421}.

\bibitem{p02:pt:Bounemoura}
A. Bounemoura,
\emph{Ergodization time for linear flows on tori via geometry of numbers},
\href{https://arxiv.org/abs/1506.06469}{arXiv:1506.06469} (2015),
especially Theorem~1 and its account of preceding optimal bounds.

\bibitem{p02:pt:Aurzada}
F. Aurzada,
\emph{A short note on small deviations of sequences of i.i.d. random
variables with exponentially decreasing weights},
Statistics \& Probability Letters \textbf{78} (2008), no.~15, 2300--2307.
\href{https://arxiv.org/abs/0711.2576}{Author manuscript: arXiv:0711.2576}.
Theorem~4 contains the geometric-weight small-deviation constant.

\bibitem{p02:pt:Rozovsky}
L. V. Rozovsky,
\emph{Small deviation probabilities of weighted sums with fast decreasing
weights}, Probability and Mathematical Statistics \textbf{35} (2015),
no.~1, 161--178.
\href{https://www.math.uni.wroc.pl/~pms/files/35.1/Abstract/35.1.10.abs.pdf}{Publisher abstract}.
This is background on the established small-deviation problem;
no theorem from this paper is an unproved input to our argument.

\bibitem{p02:zp:AriasLune}
J.~Arias de Reyna and J.~van de Lune,
\emph{Some bounds and limits in the theory of Riemann's zeta function},
Journal of Mathematical Analysis and Applications \textbf{396} (2012),
199--214. \url{https://doi.org/10.1016/j.jmaa.2012.06.017}.
Preprint \url{https://arxiv.org/abs/1107.5134}.
Theorem~5.1 and Theorem~6.1 in the inspected preprint.

\bibitem{p02:zp:Podesta2026}
J.~J. Podesta, \emph{A Function Used to Locate the Zeros of an Exponential
Sum}, Monatshefte f\"ur Mathematik \textbf{210} (2026), 223--268,
\url{https://doi.org/10.1007/s00605-026-02174-6}.
The modular-ordinate density question and its deferred solution are
stated in Section~1.
Part~II refers to the same Section~1 formulation as a modular-ordinate existence question.

\bibitem{p02:zp:LenzSpindelerStrungaru}
D.~Lenz, T.~Spindeler, and N.~Strungaru,
\emph{Pure point diffraction and almost periodicity},
Israel Journal of Chemistry \textbf{64} (2024), e202300158.
\url{https://doi.org/10.1002/ijch.202300158};
\url{https://arxiv.org/abs/2312.12825}.

\bibitem{p02:zp:FeldheimDensity}
N.~Feldheim,
\emph{Zeroes of Gaussian analytic functions with translation-invariant
distribution}, Israel Journal of Mathematics \textbf{195} (2013),
317--345. \url{https://doi.org/10.1007/s11856-012-0130-0};
\url{https://arxiv.org/abs/1105.3929}.
Theorem~1 and Remark~2.3.

\bibitem{p02:zp:FeldheimVariance}
N.~D.~Feldheim,
\emph{Variance of the number of zeroes of shift-invariant Gaussian
analytic functions}, arXiv:1309.2111v5 (2016).
\url{https://arxiv.org/abs/1309.2111}.
Published as \emph{Variance of the number of zeroes of stationary
Gaussian analytic functions}, Israel Journal of Mathematics
\textbf{227} (2018), 753--792.
Theorem~1 and Section~2.

\bibitem{p02:zp:FeldheimStatement}
N.~D.~Feldheim, research statement, author-hosted document (2017),
p.~1, paragraph on the unresolved atomic-spectrum density problem.
\url{https://u.cs.biu.ac.il/~feldhen/CV\%20and\%20stat/Naomi_short_res_2017.pdf}.

\bibitem{comp:part:02}
KEIO, \emph{Observation Geometry and Exact Recovery: Replication Rank, Thorin Measures, and Zero Patterns}, research manuscript, 14 October 2026. Original edition: 56 pp. Reproduced as \hyperref[part:02]{Part~II} of this compilation. AI-assisted and not externally peer reviewed.
Part~III cites the theorem ``A rank-completing sequence of partial-sum experiments'' as the predecessor of its exact terminal criterion.

\bibitem{p03:bnew:RheeGlynn2015}
C.-H. Rhee and P. W. Glynn,
Unbiased estimation with square root convergence for SDE models,
Operations Research 63 (2015), 1026--1043.
\url{https://web.stanford.edu/~glynn/papers/2015/RheeG15.html}.

\bibitem{p03:mixnew:Garnett2007}
J.~B.~Garnett, \emph{Bounded Analytic Functions}, revised first edition,
Graduate Texts in Mathematics 236, Springer, New York, 2007.
\url{https://doi.org/10.1007/0-387-49763-3}.

\bibitem{p03:mixnew:Peled2026}
Y.~Peled, \emph{On the $d$-rigidity phase transition in random graphs},
arXiv:2605.25711v3, revised 24 August 2026.
\url{https://arxiv.org/pdf/2605.25711v3}.
Section~7, item~3, p.~30, explicitly records the remaining
global-rigidity route to the sparse-random-graph likelihood conjecture.

\bibitem{p03:zn:FeldheimWH}
N.~Feldheim, A.~Haimi, G.~Koliander and J.~L. Romero,
\emph{Hyperuniformity and non-hyperuniformity of zeros of Gaussian
Weyl--Heisenberg Functions}, arXiv:2406.20003v3 (8 June 2026),
\url{https://arxiv.org/abs/2406.20003v3}.

\bibitem{p03:an:Banaszczyk}
W. Banaszczyk,
\emph{New bounds in some transference theorems in the geometry of numbers},
Mathematische Annalen \textbf{296} (1993), 625--635.
\href{https://doi.org/10.1007/BF01445125}{doi:10.1007/BF01445125}.

\bibitem{p03:an:ASD}
D. Aggarwal and N. Stephens-Davidowitz,
\emph{An improved constant in Banaszczyk's transference theorem},
\href{https://arxiv.org/abs/1907.09020}{arXiv:1907.09020} (2019).
Theorems~1.1 and~1.2.

\bibitem{p03:an:KNU}
D.~Krieg, E.~Novak and M.~Ullrich,
\emph{How many continuous measurements are needed to learn a vector?},
Proc. Amer. Math. Soc., doi:10.1090/proc/17658;
arXiv:2412.06468v2, September 22, 2026.
\url{https://arxiv.org/html/2412.06468v2}.
Part~IX specifically cites Theorem~1 and Corollary~11 of arXiv:2412.06468v2.

\bibitem{comp:part:03}
KEIO, \emph{Binomial Mixture Likelihood and Exact Observation Thresholds}, research manuscript, 14 October 2026. Original edition: 39 pp. Reproduced as \hyperref[part:03]{Part~III} of this compilation. AI-assisted and not externally peer reviewed.

\bibitem{p04:nxm:Lindsay1983}
B.~G.~Lindsay,
\emph{The geometry of mixture likelihoods: a general theory},
Annals of Statistics \textbf{11} (1983), no.~1, 86--94.
\url{https://doi.org/10.1214/aos/1176346059}.

\bibitem{p04:nxm:BobkovGoetze1999}
S.~G.~Bobkov and F.~G\"otze,
\emph{Exponential integrability and transportation cost related to
logarithmic Sobolev inequalities},
Journal of Functional Analysis \textbf{163} (1999), no.~1, 1--28.
\url{https://doi.org/10.1006/jfan.1998.3326}.

\bibitem{p04:nxb:Clarkson2020}
E. Clarkson,
Probability of error for detecting a change in a parameter and Bayesian
Fisher information,
Journal of the Optical Society of America A 37 (2020), 174--181.
\url{https://arxiv.org/html/1902.00099v1}.

\bibitem{p04:nxs:DyachenkoKarp}
A. Dyachenko and D. Karp,
\emph{Ratios of the Gauss hypergeometric functions with parameters
shifted by integers: more on integral representations},
arXiv:2108.03595 (2021), especially Lemma~3 and Theorem~3,
\url{https://arxiv.org/abs/2108.03595}.

\bibitem{p04:nxg:StoyanovLin}
J.~Stoyanov and G.~D.~Lin,
\emph{Hardy's condition in the moment problem for probability distributions},
Theory of Probability and Its Applications \textbf{57} (2013), 699--708.
\url{https://doi.org/10.1137/S0040585X9798631X}.

\bibitem{p04:nxg:Lin}
G.~D.~Lin, \emph{Recent developments on the moment problem},
Journal of Statistical Distributions and Applications \textbf{4} (2017), article~5.
\url{https://arxiv.org/abs/1703.01027}.

\bibitem{p04:nxz:Sepulcre}
J.~M. Sepulcre,
\emph{On the Result of Invariance of the Closure Set of the Real
Projections of the Zeros of an Important Class of Exponential Polynomials},
Journal of Function Spaces \textbf{2016} (2016), Article 3605690, 9 pp.
\url{https://doi.org/10.1155/2016/3605690}.

\bibitem{p04:nxi:Wellershoff}
M.~Wellershoff,
\emph{Phase retrieval of entire functions and its implications for Gabor
phase retrieval},
Journal of Functional Analysis \textbf{286} (2024), no.~11, 110403.
\url{https://arxiv.org/abs/2202.03733}.

\bibitem{p05:fm:DKS2008}
P.~Diaconis, K.~Khare and L.~Saloff-Coste,
\emph{Gibbs sampling, exponential families and orthogonal polynomials},
Statistical Science \textbf{23} (2008), no.~2, 151--178.
\url{https://doi.org/10.1214/07-STS252}.
Author copy: \url{https://arxiv.org/pdf/0808.3852}.
Section~4.1, especially Propositions~4.1--4.2, contains the classical
beta--binomial spectral calculation used here.

\bibitem{p05:fm:FP2018}
L.~Feng and L.~H.~Dicker,
\emph{Approximate nonparametric maximum likelihood for mixture models:
A convex optimization approach to fitting arbitrary multivariate mixing
distributions}, Computational Statistics \& Data Analysis
\textbf{122} (2018), 80--91.
\url{https://doi.org/10.1016/j.csda.2018.01.006}.
Author version: \url{https://arxiv.org/abs/1606.02011}.

\bibitem{p05:fb:DFM2015}
J. Dedecker, A. Fischer and B. Michel,
Improved rates for Wasserstein deconvolution with ordinary smooth error
in dimension one, arXiv:1404.0646v2, revised 4 March 2015.
\url{https://arxiv.org/abs/1404.0646}.

\bibitem{p05:fb:RS2023}
J. Rousseau and C. Scricciolo,
Wasserstein convergence in Bayesian and frequentist deconvolution models,
arXiv:2309.15300v1, 26 September 2023.
\url{https://arxiv.org/abs/2309.15300}.

\bibitem{p05:fb:Scricciolo2026}
C. Scricciolo,
Adaptive minimax multivariate $L^1$-deconvolution: harmonic-mean rates
under noise filtering, and a $1$-Wasserstein equivalence,
arXiv:2610.11760v1, 8 October 2026.
\url{https://arxiv.org/abs/2610.11760}.

\bibitem{p05:fs:Bondesson}
L. Bondesson,
\emph{Generalized Gamma Convolutions and Related Classes of
Distributions and Densities}, Lecture Notes in Statistics 76,
Springer, New York, 1992, especially Chapter~5,
\url{https://doi.org/10.1007/978-1-4612-2948-3}.

\bibitem{comp:part:04}
KEIO, \emph{Transport, Spectral Positivity, and Sharp Observation Boundaries: Near-linear Mixture Likelihood, Full McKay Domains, and Critical Sampling Budgets}, research manuscript, 14 October 2026. Reproduced as \hyperref[part:04]{Part~IV} of this compilation. AI-assisted and not externally peer reviewed.
Cited locator: original Section~4.

\bibitem{p05:fs:Bosch}
P. Bosch,
\emph{HCM property and the half-Cauchy distribution},
Probab. Math. Statist. \textbf{35} (2015), 191--200,
\url{https://arxiv.org/abs/1402.1059}.

\bibitem{p05:fs:Sjodin}
T. Sj\"odin,
\emph{The Power Problem for Generalized Gamma Convolutions (GGC)
and Related Questions}, arXiv:2601.04038 (2026),
\url{https://arxiv.org/abs/2601.04038}.

\bibitem{p05:fg:deJeu}
M.~de Jeu, Determinate multidimensional measures, the extended
Carleman theorem and quasi-analytic weights,
\emph{Annals of Probability} \textbf{31} (2003), 1205--1227.
\href{https://doi.org/10.1214/aop/1055425776}{doi:10.1214/aop/1055425776}.
\href{https://arxiv.org/abs/math/0111019}{arXiv:math/0111019}.

\bibitem{p05:fz:WuYang}
Y. Wu and P. Yang, Optimal estimation of Gaussian mixtures via denoised
method of moments, \emph{Annals of Statistics} \textbf{48} (2020),
1981--2007. \url{https://doi.org/10.1214/19-AOS1873};
\url{https://arxiv.org/abs/1807.07237}.

\bibitem{p05:fi:FricainLefevre}
E.~Fricain and P.~Lef\`evre,
\emph{$L^2$-M\"untz spaces as model spaces},
author manuscript, 22 September 2017.
\url{https://pro.univ-lille.fr/fileadmin/user_upload/pages_pros/emmanuel_fricain/Muntz-model.pdf}.

\bibitem{p05:fi:GaillardLefevre}
L.~Gaillard and P.~Lef\`evre,
\emph{Lacunary M\"untz spaces: isomorphisms and Carleson embeddings},
Annales de l'Institut Fourier \textbf{68} (2018), no.~5, 2215--2251.
\url{https://doi.org/10.5802/aif.3207}.

\bibitem{p05:fi:Marechal}
P.~Mar\'echal, F.~Triki, and W.~C.~Simo Tao Lee,
\emph{Regularization of the inverse Laplace transform by Mollification},
arXiv:2304.08057, 2023.
\url{https://arxiv.org/abs/2304.08057}.

\bibitem{comp:part:05}
KEIO, \emph{Adaptive Analytic Information and Sharp Recovery Boundaries}, research manuscript, 14 October 2026. Original edition: 52 pp. Reproduced as \hyperref[part:05]{Part~V} of this compilation. AI-assisted and not externally peer reviewed.
Cited locators: original Theorems~6.2 and~6.4.

\bibitem{p06:ra:Tuzcu}
Y. E. Tuzcu,
\emph{Analysis of the prime zeta function zeros placement patterns},
Master's thesis, Vilnius University, 2026, p.~16.
\href{https://epublications.vu.lt/object/elaba:308153308/308153308.pdf}{University-hosted full text}.

\bibitem{p06:ra:Titchmarsh}
E. C. Titchmarsh, revised by D. R. Heath-Brown,
\emph{The Theory of the Riemann Zeta-Function}, second edition,
Clarendon Press, Oxford, 1986, Chapter~VII.

\bibitem{p06:ra:Gonek}
S. M. Gonek, \emph{Three lectures on the Riemann zeta-function},
arXiv:math/0401126, 2004, Lecture~II, Sections~3--4.
\href{https://arxiv.org/abs/math/0401126}{arXiv:math/0401126}.

\bibitem{p06:ni:KNPU}
D.~Krieg, E.~Novak, L.~Plaskota and M.~Ullrich,
\emph{Noisy nonlinear information and entropy numbers},
J. Fourier Anal. Appl. \textbf{32} (2026), article 36,
doi:10.1007/s00041-026-10249-z.
\url{https://arxiv.org/abs/2510.23213}.

\bibitem{p06:ni:Bernstein}
H.~Queff\'elec and R.~Zarouf,
\emph{On Bernstein's inequality for polynomials},
Analysis and Mathematical Physics (2019); arXiv:1903.10801.
\url{https://arxiv.org/abs/1903.10801}.

\bibitem{p06:rs:Eftekhari}
A. Eftekhari, J. Tanner, A. Thompson, B. Toader, and H. Tyagi,
\emph{Sparse non-negative super-resolution---simplified and stabilised},
arXiv:1804.01490, version 2 (2019).
\url{https://arxiv.org/abs/1804.01490}.

\bibitem{comp:part:06}
KEIO, \emph{Boundary Laws for Observation and Recovery: Prime-Zeta Counting, Global Likelihood Optimality, and Sharp Inverse Thresholds}, research manuscript, 15 October 2026. Original edition: 48 pp. Reproduced as \hyperref[part:06]{Part~VI} of this compilation. AI-assisted and not externally peer reviewed.
Part~VII cites the sections on the complete binomial-mixture likelihood transition, blind spectral certification from nonlinear transition thresholds, Bessel distributions, and analytic populations.

\bibitem{p07:ml13:Dette}
H.~Dette, New bounds for Hahn and Krawtchouk polynomials,
\emph{SIAM Journal on Mathematical Analysis} \textbf{26} (1995),
1647--1659. doi:10.1137/S0036141093255661.
Preprint arXiv:math/9406223, Corollary~3.6(b).
\url{https://arxiv.org/abs/math/9406223}.

\bibitem{p07:ml13:CGR}
M.~Chen, C.~Gao and Z.~Ren,
A general decision theory for Huber's $\varepsilon$-contamination model,
\emph{Electronic Journal of Statistics} \textbf{10} (2016), 3752--3774.
doi:10.1214/16-EJS1216.
\url{https://arxiv.org/abs/1511.04144}.

\bibitem{p07:ar13:Bagchi}
B. Bagchi,
\emph{The statistical behaviour and universality properties of the Riemann zeta-function and other allied Dirichlet series},
Ph.D. thesis, Indian Statistical Institute, Calcutta, 1981.
\href{https://digitalcommons.isical.ac.in/masters-dissertations/41/}{Institutional repository record}.

\bibitem{p07:ar13:Laurincikas}
A. Laurin\v{c}ikas,
\emph{Limit theorems for the Matsumoto zeta-function},
Journal de Th\'eorie des Nombres de Bordeaux \textbf{8} (1996), no.~1, 143--159.
\href{https://jtnb.centre-mersenne.org/article/JTNB_1996__8_1_143_0.pdf}{Journal-hosted full text}.

\bibitem{p07:in13:Coste}
M.~Coste, \emph{An Introduction to O-Minimal Geometry},
Dottorato di Ricerca in Matematica, Dipartimento di Matematica,
Universit\`a di Pisa, Istituti Editoriali e Poligrafici Internazionali,
Pisa, 2000 (lecture notes dated 1999).

\bibitem{p07:in13:Wilkie}
A.~J.~Wilkie, \emph{O-minimal structures},
Ast\'erisque \textbf{326} (2009), S\'eminaire Bourbaki,
Exp. No.~985, 131--142.
\url{https://www.numdam.org/article/AST_2009__326__131_0.pdf}.

\bibitem{p07:in13:Kwon}
J.~Kwon, \emph{A universal bound on the variations of bounded convex functions},
J. Convex Anal. \textbf{24} (2017), no.~1, 67--73;
arXiv:1401.2104.

\bibitem{p07:ga13:Nolan}
J.~P.~Nolan, \emph{Stable Distributions: Models for Heavy Tailed Data},
Chapter 1, author manuscript processed 30 January 2018,
Definition 1.8 and equation (1.6).
\href{https://edspace.american.edu/jpnolan/wp-content/uploads/sites/1720/2020/09/Chap1.pdf}{Author's chapter PDF}.

\bibitem{p07:sp13:LiuAmmari}
P.~Liu and H.~Ammari,
\emph{A mathematical theory of super-resolution and two-point resolution},
Forum of Mathematics, Sigma \textbf{12} (2024), e83.
\href{https://doi.org/10.1017/fms.2024.72}{doi:10.1017/fms.2024.72}.

\bibitem{p07:sp13:GonekMontgomery}
S.~M.~Gonek and H.~L.~Montgomery,
\emph{Kronecker's approximation theorem},
Indagationes Mathematicae \textbf{27} (2016), 506--523.
\href{https://doi.org/10.1016/j.indag.2016.02.002}{doi:10.1016/j.indag.2016.02.002}.

\bibitem{p08:in14:BB}
A.~Balitskiy and A.~Berdnikov,
\emph{Waist of maps measured via Urysohn width},
arXiv:2009.04558v2, July 25, 2021;
Theorems~2.1 and~3.16.
\url{https://arxiv.org/abs/2009.04558}.

\bibitem{p08:in14:Milnor}
J.~Milnor, \emph{On the Betti numbers of real varieties},
Proc. Amer. Math. Soc. \textbf{15} (1964), 275--280,
doi:10.1090/S0002-9939-1964-0161339-9.

\bibitem{p08:in14:Shiota}
M.~Shiota, \emph{Piecewise linearization of real-valued subanalytic
functions}, Trans. Amer. Math. Soc. \textbf{312} (1989), no.~2,
663--679, doi:10.1090/S0002-9947-1989-0958902-3.

\bibitem{comp:part:07}
KEIO, \emph{Exact Recovery Laws and a Coefficient Conjecture}, research manuscript, 16 October 2026. Reproduced as \hyperref[part:07]{Part~VII} of this compilation. AI-assisted and not externally peer reviewed.
Cited sections concern the constant-one Chebyshev--Bernstein bound, two-scale inversion and robust binomial population recovery, and the exact Euclidean convex-noise floor.

\bibitem{p08:ml14:CDSS}
S.-O.~Chan, I.~Diakonikolas, R.~A.~Servedio and X.~Sun,
Efficient density estimation via piecewise polynomial approximation,
\emph{Proceedings of the 46th Annual ACM Symposium on Theory of
Computing} (2014), 604--613. doi:10.1145/2591796.2591848.
\url{https://arxiv.org/abs/1305.3207}.

\bibitem{p08:ml14:TKV}
K.~Tian, W.~Kong and G.~Valiant, Learning populations of parameters,
\emph{Advances in Neural Information Processing Systems} \textbf{30}
(2017). \url{https://arxiv.org/abs/1709.02707}.

\bibitem{p08:ar14:Sodin}
M. Sodin,
\emph{Zeros of Gaussian analytic functions},
Mathematical Research Letters \textbf{7} (2000), 371--381.
\href{https://arxiv.org/abs/math/0007030}{arXiv:math/0007030}.

\bibitem{p08:ar14:Favorov}
S. Favorov,
\emph{Generalized Fourier quasicrystals, almost periodic sets, and zeros of Dirichlet series},
\href{https://arxiv.org/abs/2311.02728}{arXiv:2311.02728} (2023).

\bibitem{p08:ga14:Almira}
J.~M.~Almira, M\"untz type theorems I,
\emph{Surveys in Approximation Theory} \textbf{3} (2007), 152--194.
\href{https://arxiv.org/abs/0710.3570}{arXiv:0710.3570}.

\bibitem{p08:sp14:BGY}
D.~Batenkov, G.~Goldman, and Y.~Yomdin,
\emph{Super-resolution of near-colliding point sources},
Information and Inference \textbf{10} (2021), 515--572.
\href{https://arxiv.org/abs/1904.09186}{arXiv:1904.09186}.

\bibitem{p08:sp14:KDB}
R.~Katz, N.~Diab, and D.~Batenkov,
\emph{On the accuracy of Prony's method for recovery of exponential
sums with closely spaced exponents},
arXiv:2302.05883v2 (2024).
\href{https://arxiv.org/abs/2302.05883}{arXiv:2302.05883}.

\bibitem{p08:sp14:DF}
H.~S.~Dumas and S.~Fischler,
\emph{Filling times for linear flow on the torus with truncated
Diophantine conditions: a brief review and new proof},
arXiv:2112.11786 (2021).
\href{https://arxiv.org/abs/2112.11786}{arXiv:2112.11786}.

\bibitem{p09:in15:KK}
D.~Krieg and P.~Kritzer,
\emph{Homogeneous algorithms and solvable problems on cones},
arXiv:2311.15767v2, March 3, 2024,
Propositions~1--2 and Open Problem~1.
\url{https://arxiv.org/abs/2311.15767}.

\bibitem{comp:part:08}
KEIO, \emph{Geometry and Rigidity of Inverse Observation}, research manuscript, 16 October 2026. Original edition: 39 pp. Reproduced as \hyperref[part:08]{Part~VIII} of this compilation. AI-assisted and not externally peer reviewed.
Part~IX cites original Section~1 on logarithmic lower bounds for unrestricted adaptive continuous information, and the sections on polynomial-time population recovery under unknown adversarial contamination, spectral observation, and endpoint-rate gamma extraction.

\bibitem{p09:op15:Gournay}
A.~Gournay, \emph{Width of $l^p$ balls},
arXiv:0711.3081v2 (2008), Theorem~1.4.
\url{https://arxiv.org/abs/0711.3081}.

\bibitem{p09:ml15:LiuLoh}
Z.~Liu and P.-L.~Loh,
Robust W-GAN-based estimation under Wasserstein contamination,
\emph{Information and Inference} \textbf{12} (2023), 312--362.
doi:10.1093/imaiai/iaac020.
\url{https://arxiv.org/abs/2101.07969}.

\bibitem{p09:ml15:ChaoDobriban}
P.~Chao and E.~Dobriban,
\emph{Minimax Statistical Estimation under Wasserstein Contamination},
arXiv:2308.01853v3 (2025).
\url{https://arxiv.org/abs/2308.01853}.

\bibitem{p09:ar15:AvellarHale}
C.~E.~Avellar and J.~K.~Hale,
\emph{On the zeros of exponential polynomials},
Journal of Mathematical Analysis and Applications \textbf{73}
(1980), no.~2, 434--452.
\href{https://doi.org/10.1016/0022-247X(80)90289-9}
{doi:10.1016/0022-247X(80)90289-9}.

\bibitem{p09:ar15:Lundqvist}
J. Lundqvist,
\emph{An explicit calculation of the Ronkin function},
Annales de la Facult\'e des Sciences de Toulouse, S\'erie~6,
\textbf{24} (2015), no.~2, 227--250.
\href{https://www.numdam.org/item/AFST_2015_6_24_2_227_0/}{Journal-hosted article}.

\bibitem{p09:ga15:MP}
N. Makarov and A. Poltoratski,
\emph{Beurling--Malliavin theory for Toeplitz kernels},
Invent. Math. \textbf{180} (2010), 443--480;
\url{https://arxiv.org/abs/math/0702497}.
Part~IX specifically cites Section~1.1 of the 2007 preprint.

\bibitem{p09:ga15:CFR}
J.~A.~Cuesta-Albertos, R.~Fraiman, and T.~Ransford,
A sharp form of the Cram\'er--Wold theorem,
\emph{Journal of Theoretical Probability} \textbf{20} (2007), 201--209.
\href{https://doi.org/10.1007/s10959-007-0060-7}{doi:10.1007/s10959-007-0060-7}.

\bibitem{p09:sp15:quadrature}
G.~Blekherman, M.~Kummer, C.~Riener, M.~Schweighofer, and C.~Vinzant,
\emph{Generalized eigenvalue methods for Gaussian quadrature rules},
Annales Henri Lebesgue \textbf{3} (2020), 1327--1341.
\href{https://doi.org/10.5802/ahl.62}{doi:10.5802/ahl.62}.

\bibitem{p09:sp15:rice}
D.~Armentano, J.-M.~Aza\"is, and J.~R.~Le\'on,
\emph{On a general Kac--Rice formula for the measure of a level set},
arXiv:2304.07424v4 (2025).
\href{https://arxiv.org/abs/2304.07424}{arXiv:2304.07424}.

\bibitem{comp:part:09}
KEIO, \emph{Sharp Boundaries for Adaptive Recovery and Mixture Identification}, research manuscript, 16 October 2026. Original edition: 41 pp. Reproduced as \hyperref[part:09]{Part~IX} of this compilation. AI-assisted and not externally peer reviewed.

\bibitem{p10:op16:Tsukamoto}
M.~Tsukamoto, \emph{Macroscopic dimension of the $\ell^p$-ball with
respect to the $\ell^q$-norm}, J. Math. Kyoto Univ. 48 (2008),
445--454; arXiv:0801.4548.
\url{https://arxiv.org/abs/0801.4548}.

\bibitem{p10:ml16:BHR}
K.~Bujok, B.~M.~Hambly and C.~Reisinger,
Multilevel simulation of functionals of Bernoulli random variables
with application to basket credit derivatives,
\emph{Methodology and Computing in Applied Probability}
\textbf{17} (2015), 579--604. doi:10.1007/s11009-013-9380-5.
\url{https://arxiv.org/abs/1211.0707}.

\bibitem{p10:ml16:LG}
Y.~Luo and C.~Gao,
Adaptive robust confidence intervals,
\emph{Annals of Statistics} \textbf{54} (2026), 1128--1152.
doi:10.1214/25-AOS2574.
\url{https://arxiv.org/abs/2410.22647}.

\bibitem{p10:ml16:CLG}
M.~Cho, Y.~Luo and C.~Gao,
\emph{Robust confidence intervals for a binomial proportion:
local optimality and adaptivity}, arXiv:2509.05568v2 (2026).
\url{https://arxiv.org/abs/2509.05568}.

\bibitem{p10:ar16:FornasieroTerzo}
A.~Fornasiero and G.~Terzo,
\emph{A note on exponential polynomials},
Ricerche di Matematica \textbf{74} (2025), 301--308;
published online 4 July 2024.
\href{https://doi.org/10.1007/s11587-024-00871-8}
{doi:10.1007/s11587-024-00871-8}.

\bibitem{p10:ga16:BJR}
S. Bonhomme, K. Jochmans, and J.-M. Robin,
\emph{Non-parametric estimation of finite mixtures from repeated measurements},
J. R. Stat. Soc. Ser. B \textbf{78} (2016), 211--229,
\url{https://doi.org/10.1111/rssb.12110}.

\bibitem{p10:ga16:Simon}
T. Simon,
\emph{Comparing Fr\'echet and positive stable laws},
Electron. J. Probab. \textbf{19} (2014), paper 16, 1--25,
\url{https://doi.org/10.1214/EJP.v19-3058};
\url{https://arxiv.org/html/1310.1888v2}.

\bibitem{p10:ga16:PX}
M. R. Powers and J. Xu,
\emph{A criterion for extending continuous-mixture identifiability results},
J. Stat. Theory Appl. \textbf{24} (2025), 515--533,
\url{https://doi.org/10.1007/s44199-025-00124-z}.

\bibitem{p10:sp16:momentmatching}
Y.~Han, J.~Jiao, and T.~Weissman,
\emph{Local moment matching: A unified methodology for symmetric
functional estimation and distribution estimation under Wasserstein
distance}, Proceedings of the 31st Conference on Learning Theory,
PMLR \textbf{75} (2018), 3189--3221.
\href{https://proceedings.mlr.press/v75/han18b.html}{Publisher's article}.

\bibitem{p10:sp16:beckchen}
J.~Beck and W.~W.~L.~Chen,
\emph{Generalization of a density theorem of Khinchin and diophantine
approximation}, Journal de th\'eorie des nombres de Bordeaux
\textbf{35} (2023), no.~2, 511--542.
\href{https://doi.org/10.5802/jtnb.1255}{doi:10.5802/jtnb.1255}.

\end{thebibliography}
\end{document}
